\def\lakitdocumentclass{article}
\def\lakitclassoptions{11pt,a4paper}
\def\lakitgeometryoptions{margin=2.1cm,headheight=14pt}
\def\lakitlinespread{1.16}
% ============================================================================
% preamble.tex — tested LuaLaTeX preamble for Hebrew + math + figures PDFs (v4.32.0)
% Battle-tested: compiled an 80-page Hebrew quantum-computing summary AND a
% 197-page Hebrew computer-architecture anthology (TikZ diagrams, code listings,
% colored boxes) — both with 0 missing characters.
%
% Usage in your document:
%   % ============================================================================
% preamble.tex — tested LuaLaTeX preamble for Hebrew + math + figures PDFs (v4.32.0)
% Battle-tested: compiled an 80-page Hebrew quantum-computing summary AND a
% 197-page Hebrew computer-architecture anthology (TikZ diagrams, code listings,
% colored boxes) — both with 0 missing characters.
%
% Usage in your document:
%   % ============================================================================
% preamble.tex — tested LuaLaTeX preamble for Hebrew + math + figures PDFs (v4.32.0)
% Battle-tested: compiled an 80-page Hebrew quantum-computing summary AND a
% 197-page Hebrew computer-architecture anthology (TikZ diagrams, code listings,
% colored boxes) — both with 0 missing characters.
%
% Usage in your document:
%   \input{preamble.tex}     % (or paste these lines at the top)
%   \begin{document}  ...  \end{document}
%
% Compile with:  lualatex -interaction=nonstopmode doc.tex
%   Repeat until two consecutive PDFs are byte-identical (scripts/build.sh
%   does this) — a TOC document converges only on pass 3+; after two passes
%   the TOC page numbers are STALE (measured; see SKILL.md step 3).
% Font binaries are external dependencies, pinned by assets/fonts/SHA256SUMS.
% scripts/setup_fonts.sh imports a verified local installation or prior skill.
%
% PUBLIC DOCUMENT-SHELL / LAYOUT CONTRACT (introduced v4.3.1; architecture-owned v4.7.1). These hooks may be
% defined BEFORE \input{preamble.tex}; if omitted, the historical canonical
% defaults below are exact.  They deliberately control only the document shell
% and page layout.  RTL/math semantics stay in this preamble, while visual
% presentation themes (for example clean_theme.tex) load AFTER it.
%
%   \def\lakitdocumentclass{article}
%   \def\lakitclassoptions{11pt,a4paper}
%   \def\lakitgeometryoptions{margin=1.8cm,headheight=14pt}
%   \def\lakitlinespread{1.15}
%   \def\lakitparagraphstyle{block} % letters/briefs; indented for continuous prose
%   \def\lakittocdepth{1}
%   \def\lakitsecnumdepth{2}
%   \input{preamble.tex}
%
% Supported class contract: report or article.  Alternate paper sizes must also
% be declared to build.sh through EXPECTED_PAPER; the PDF gate never infers them.
% Do NOT use these hooks as a theme mechanism: theme overlays must not own page
% geometry, class selection, or direction semantics. A declared presentation-policy
% overlay may intentionally change numbering visibility after the preamble.
% ============================================================================
% Runtime architecture guard: generated from scripts/architecture_contract.py.
% The readiness sentinel remains delayed until the END of this file.
% <hllpdf:architecture-guard>
\ifdefined\hllpdfcoreloaded
  \errmessage{hllpdf: preamble.tex loaded more than once}
  \expandafter\endinput
\fi
% </hllpdf:architecture-guard>

\ifdefined\lakitdocumentclass\else\def\lakitdocumentclass{report}\fi
\ifdefined\lakitclassoptions\else\def\lakitclassoptions{12pt,a4paper}\fi
\ifdefined\lakitgeometryoptions\else\def\lakitgeometryoptions{margin=2.3cm,headheight=14pt}\fi
\ifdefined\lakitlinespread\else\def\lakitlinespread{1.3}\fi
\ifdefined\lakittocdepth\else\def\lakittocdepth{1}\fi
\ifdefined\lakitsecnumdepth\else\def\lakitsecnumdepth{2}\fi
% Paragraph organization belongs to the document layout, independently of its
% language and theme. Preserve historical class dimensions unless explicitly
% selecting block paragraphs; never suppress indentation with an everypar hack.
\ifdefined\lakitparagraphstyle\else\def\lakitparagraphstyle{indented}\fi
\edef\lakitparagraphstyle{\lakitparagraphstyle}
\def\hllpdfparagraphblock{block}
\def\hllpdfparagraphindented{indented}
\ifx\lakitparagraphstyle\hllpdfparagraphblock\else
  \ifx\lakitparagraphstyle\hllpdfparagraphindented\else
    \errmessage{hllpdf: paragraph style must be block or indented}
  \fi
\fi
% v4.3.1: the public class hook has a deliberately narrow support contract.
% Before this guard, \def\lakitdocumentclass{book} silently compiled even though
% every heading/TOC regression test targets report/article; an unsupported shell
% could therefore look "supported" until a subtle class-specific defect shipped.
% Fail before \documentclass rather than letting an undocumented combination enter
% the mechanism core. Plain-TeX comparison is used because LaTeX packages are not
% available yet at this point.
\edef\lakitvalidatedclass{\lakitdocumentclass}
\def\lakitreportclass{report}
\def\lakitarticleclass{article}
\ifx\lakitvalidatedclass\lakitreportclass
\else\ifx\lakitvalidatedclass\lakitarticleclass
\else
  \errmessage{hllpdf: unsupported class \lakitdocumentclass; use report or article}
\fi\fi
\let\lakitvalidatedclass\relax
\let\lakitreportclass\relax
\let\lakitarticleclass\relax
% Language intent is independent of document class, theme and content domain.
% The two dominant-language profiles differ in authoring intent, not direction.
% A balanced bilingual document must explicitly choose a navigation/base locale.
% <hllpdf:language-profile>
\providecommand{\lakitlanguageprofile}{hebrew}
\edef\lakitlanguageprofile{\lakitlanguageprofile}
\def\hllpdfhebrew{hebrew}\def\hllpdfenglish{english}
\def\hllpdfhebrewmixed{hebrew-with-english}
\def\hllpdfenglishmixed{english-with-hebrew}
\def\hllpdfbilingual{bilingual}
\let\hllpdfexpectedlanguage\relax
\ifx\lakitlanguageprofile\hllpdfhebrew
  \let\hllpdfexpectedlanguage\hllpdfhebrew
\else\ifx\lakitlanguageprofile\hllpdfhebrewmixed
  \let\hllpdfexpectedlanguage\hllpdfhebrew
\else\ifx\lakitlanguageprofile\hllpdfenglish
  \let\hllpdfexpectedlanguage\hllpdfenglish
\else\ifx\lakitlanguageprofile\hllpdfenglishmixed
  \let\hllpdfexpectedlanguage\hllpdfenglish
\else\ifx\lakitlanguageprofile\hllpdfbilingual
  \ifdefined\lakitmainlanguage\else
    \errmessage{hllpdf: bilingual profile requires explicit lakitmainlanguage}
  \fi
\else
  \errmessage{hllpdf: unsupported language profile \lakitlanguageprofile}
\fi\fi\fi\fi\fi
\ifdefined\lakitmainlanguage
  \edef\lakitmainlanguage{\lakitmainlanguage}
\else
  \let\lakitmainlanguage\hllpdfexpectedlanguage
\fi
\ifx\lakitmainlanguage\hllpdfhebrew\else
  \ifx\lakitmainlanguage\hllpdfenglish\else
    \errmessage{hllpdf: main language must be hebrew or english}
  \fi
\fi
\ifx\hllpdfexpectedlanguage\relax\else
  \ifx\lakitmainlanguage\hllpdfexpectedlanguage\else
    \errmessage{hllpdf: main language conflicts with language profile}
  \fi
\fi
\typeout{HLLLANG|\lakitlanguageprofile|\lakitmainlanguage}
% </hllpdf:language-profile>
\documentclass[\lakitclassoptions]{\lakitdocumentclass}
\usepackage[no-math]{fontspec}        % (1) ★ no-math — THE critical line (see recipe.md)
% layout=lists asks Babel to adapt list parshape when paragraph and page
% directions differ. Without it, minority-language list bodies lose their
% indentation and labels hang outside the text block, despite a clean log.
% The upstream mechanism preserves ordinary same-direction list geometry.
\usepackage[bidi=basic,layout=lists]{babel}
\ifx\lakitmainlanguage\hllpdfenglish
  \babelprovide[import]{hebrew}
  \babelprovide[import,main]{english}
\else
  \babelprovide[import,main]{hebrew}
  \babelprovide[import]{english}
\fi
% v4.22: bind canonical families and fallbacks to the selected installation.
% Lookup by basename or family alone is not proof of the file actually used.
% <hllpdf:font-binding>
\directlua{
  local dir = os.getenv("HEBREW_LATEX_FONT_DIR")
  if not dir or dir == "" then
    local data = os.getenv("XDG_DATA_HOME")
    if not data or data == "" or not (data:sub(1,1) == "/") then
      data = (os.getenv("HOME") or "") .. "/.local/share"
    end
    dir = data .. "/fonts/hebrew-lualatex-pdf"
  end
  if not (dir:sub(1,1) == "/") then dir = lfs.currentdir() .. "/" .. dir end
  if dir:find("[",1,true) or dir:find("]",1,true) or dir:find(string.char(0),1,true)
     or dir:find(string.char(10),1,true) or dir:find(string.char(13),1,true) then
    tex.error("hllpdf: unsupported character in font directory")
  end
  if not (lfs.attributes(dir,"mode") == "directory") then
    tex.error("hllpdf: font directory is missing")
  end
  hllpdf_font_dir = dir:gsub("/+$", "") .. "/"
  tex.sprint(string.char(92) .. "def" .. string.char(92) .. "hllpdffontdir{")
  tex.sprint(-2,hllpdf_font_dir)
  tex.sprint("}")
  luaotfload.add_fallback("hebfb",{"[" .. hllpdf_font_dir .. "DejaVuSerif.ttf]:mode=node;"})
  luaotfload.add_fallback("hebfbbf",{"[" .. hllpdf_font_dir .. "DejaVuSerif-Bold.ttf]:mode=node;"})
}
% </hllpdf:font-binding>
% v3.22: Frank Ruhl Libre ships NO italic face. Instead of silently-upright
% \emph (v<=3.20) or bold-as-emphasis (v3.21), synthesize a slanted face with
% FakeSlant — real, VISIBLE italic emphasis with correct \emph nesting
% semantics (inner \emph toggles back upright). Verified empirically at
% 300dpi: slant clearly visible, 0 errors, 0 missing glyphs.
\babelfont{rm}[Path={\hllpdffontdir},RawFeature={fallback=hebfb},
  BoldFont={FrankRuhlLibre-Bold.ttf},
  BoldFeatures={RawFeature={fallback=hebfbbf}},   % v3.56 — see hebfbbf above
  ItalicFont={FrankRuhlLibre-Regular.ttf}, ItalicFeatures={FakeSlant=0.22},
  % Bold-italic uses the weight-matched fallback and a synthetic slant.
  BoldItalicFont={FrankRuhlLibre-Bold.ttf}, BoldItalicFeatures={RawFeature={fallback=hebfbbf},FakeSlant=0.22}
  ]{FrankRuhlLibre-Regular.ttf}
% Sans fallbacks preserve family and weight for the few missing codepoints.
\directlua{luaotfload.add_fallback("hebfbsans",{"[" .. hllpdf_font_dir .. "DejaVuSans.ttf]:mode=node;"})}
\directlua{luaotfload.add_fallback("hebfbsansbf",{"[" .. hllpdf_font_dir .. "DejaVuSans-Bold.ttf]:mode=node;"})}
\babelfont{sf}[Path={\hllpdffontdir},RawFeature={fallback=hebfbsans},
  BoldFont={Heebo-Bold.ttf},
  BoldFeatures={RawFeature={fallback=hebfbsansbf}}]{Heebo-Regular.ttf}
% v3.24: Scale=MatchLowercase — DejaVu Sans Mono has an outsized x-height, so
% unscaled inline code visually SHOUTS next to the delicate Frank Ruhl body
% (measured on a 197-page build: scaled, a code atom sits in the line at the
% body's weight, and the document tightened by ~5%). Applies to listings too.
\babelfont{tt}[Path={\hllpdffontdir},RawFeature={fallback=hebfb},Scale=MatchLowercase,
  BoldFont={DejaVuSansMono.ttf},
  BoldFeatures={RawFeature={fallback=hebfb},FakeBold=1.5},
  ItalicFont={DejaVuSansMono.ttf},
  ItalicFeatures={RawFeature={fallback=hebfb},FakeSlant=0.2},
  BoldItalicFont={DejaVuSansMono.ttf},
  BoldItalicFeatures={RawFeature={fallback=hebfb},FakeBold=1.5,FakeSlant=0.2}
  ]{DejaVuSansMono.ttf}

% Explicit multilingual/English profiles use real Latin faces and italics.
% The legacy/default Hebrew profile retains its established measured typography.
% Latin Modern is provided by the TeX distribution, not the pinned font bundle;
% recorder/runtime evidence reports it as an additional toolchain resource.
\ifx\lakitlanguageprofile\hllpdfhebrew\else
  \babelfont[english]{rm}[Ligatures=TeX,RawFeature={fallback=hebfb},
    BoldFeatures={RawFeature={fallback=hebfbbf}},
    BoldItalicFeatures={RawFeature={fallback=hebfbbf}}]{Latin Modern Roman}
  \babelfont[english]{sf}[Ligatures=TeX,RawFeature={fallback=hebfbsans},
    BoldFeatures={RawFeature={fallback=hebfbsansbf}},
    BoldItalicFeatures={RawFeature={fallback=hebfbsansbf}}]{Latin Modern Sans}
\fi

\usepackage{amsmath,amssymb,mathtools}
% amsthm supplies theorem/proof machinery. Environment names remain
% document-specific; the generic proof label is localized below.
\usepackage{amsthm}
% Proof names follow the surrounding prose, including return from an island.
\addto\extrashebrew{\renewcommand{\proofname}{הוכחה}}
\addto\extrasenglish{\renewcommand{\proofname}{Proof}}
\ifx\lakitmainlanguage\hllpdfenglish
  \renewcommand{\proofname}{Proof}
\else
  \renewcommand{\proofname}{הוכחה}
\fi
% v3.24.1: equation tags under babel-bidi. A tag re-emitted by mathtools'
% showonlyrefs inside align inherits the RTL base and its parens MIRROR —
% ")2.1(" instead of "(2.1)" (plain align/equation tags are unaffected;
% verified on a minimal case + a 198-page build). LTR-isolate at the tag-BOX
% level: mathtools' show/hide + numbering logic is untouched. Do NOT override
% \tagform@ instead — that clobbers the showonlyrefs mechanism (measured:
% every tag shows, numbering breaks).
% v3.28.1: the v3.24.1 wrap must apply ONLY to tags of display environments.
% Inline \eqref in Hebrew PROSE also goes through \maketag@@@, and there the
% english-wrap MIRRORS the parens — ")1(" — while the unwrapped original is
% correct (proven on a 2x2 isolation matrix: patch x showonlyrefs; regressed
% every prose \eqref since v3.24.1, caught by the render gate's R4). The two
% call sites are indistinguishable by \ifmmode (tag boxes are text mode), so
% the display environments themselves raise a flag; \eqref in prose sees it
% false and gets the original untouched \maketag@@@.
\usepackage{etoolbox}

% babel bidi=basic (including upstream 26.9) assumes every automatic
% discretionary has a one-glyph replacement. LuaTeX can legally produce one
% with replace=nil after a hyphenated Latin value; Babel then indexes nil and the
% whole build dies. Do the narrow upstream-compatible correction: while Babel
% classifies the list, present only that malformed disc as an ordinary
% discretionary, then restore its subtype. Babel still reorders the ENTIRE box;
% there is no pcall and no successful fallback to an unreordered node list.
% If Babel throws for any other reason, the build remains fatal. The synthetic
% nil-disc probe and the real table gallery in test_all.sh pin both layers.
\AtBeginDocument{\directlua{
  if Babel and Babel.bidi and not Babel.heb_nil_disc_compat then
    Babel.heb_nil_disc_compat = true
    local original = Babel.bidi
    local disc_id = node.id('disc')
    Babel.bidi = function(head, ispar, hdir)
      local hidden = {}
      for item in node.traverse_id(disc_id, head) do
        if item.subtype == 2 and (not item.replace or not item.replace.char) then
          table.insert(hidden, item)
          item.subtype = 0
        end
      end
      local result = original(head, ispar, hdir)
      for _, item in ipairs(hidden) do item.subtype = 2 end
      return result
    end
  end
}}
\makeatletter
\newif\ifHEBdisptag
\let\ORIGmaketag@@@\maketag@@@
\def\maketag@@@#1{\ifHEBdisptag\ORIGmaketag@@@{\foreignlanguage{english}{#1}}\else\ORIGmaketag@@@{#1}\fi}
\@for\HEB@env:=equation,equation*,align,align*,gather,gather*,multline,multline*,flalign,flalign*,alignat,alignat*,eqnarray,eqnarray*\do{%
  \expandafter\AtBeginEnvironment\expandafter{\HEB@env}{\HEBdisptagtrue}}
\makeatother
% English page + Hebrew prose: Babel's equation-only tag flip uses a primitive
% leqno path that shifts the formula to the left edge and the tag into the body
% of an English page. Align's full-width alignment is unaffected. Give only this
% display a physical LTR container and native amsmath eqno/leqno. The engine then
% owns real tag width, collision handling and labels; no proxy tag box hides ink
% from its fit calculation. Babel's saved \bbl@eqnodir still restores Hebrew text
% inside math, and the prose locale/math direction are never switched.
% Close the prior RTL paragraph before normalization. After the environment's
% group restores its directions, close TeX's empty resumed display paragraph and
% use native LaTeX continuation handling: same-paragraph prose has no indent,
% while an authored blank line starts an ordinary indented RTL paragraph.
% Register after Babel's setup; remain local to equation(*) in a main-English
% document's actual RTL prose. Native Hebrew, ordinary English and explicit
% fleqn retain their original paths. No page/display dimensions are assigned.
% <hllpdf:minority-rtl-equation-geometry>
\makeatletter
\newcommand{\HEB@minorityprosecontinuation}{\par\@doendpe}
\newcommand{\HEB@minorityequationgeometry}{%
  \ifx\lakitmainlanguage\hllpdfenglish
    \ifnum\pardirection=\@ne
      \if@fleqn\else
        \par\aftergroup\HEB@minorityprosecontinuation
        \iftagsleft@\let\veqno\@@leqno
        \else\let\veqno\@@eqno\fi
        \bodydir TLT\pardir TLT\textdir TLT\relax
      \fi
    \fi
  \fi}
\AtBeginDocument{%
  \AddToHook{env/equation/begin}{\HEB@minorityequationgeometry}%
  \AddToHook{env/equation*/begin}{\HEB@minorityequationgeometry}}
\makeatother
% </hllpdf:minority-rtl-equation-geometry>
% v3.63: \shortintertext rendered Hebrew MIRROR-REVERSED and flush LEFT — an entire
% commentary line inside an align printed backwards. babel's bidi=basic patches
% \intertext@ only (luababel.def), and mathtools' \shortintertext uses the identical
% \vbox{\normalbaselines … \noindent\ignorespaces#1\par} skeleton but is never
% patched, so its paragraph keeps pardir/textdir LTR: no bidi reordering (logical
% order on the page) and flush-left. Measured against \intertext with the same
% Applying babel's own replacement to \shortintertext@ makes the two identical.
% \h{}/\mbox fix the letter order but leave the line flush LEFT, so they are not
% substitutes. Guarded + \AtBeginDocument because mathtools \let-copies at load time.
% NB \ExplSyntaxOn is REQUIRED: mathtools' internals carry expl3 names, and without
% the expl3 catcodes the `_` is read as a subscript and the patch detonates far away
% with "! Illegal unit of measure (pt inserted)" at the \end{align} (measured).
\makeatletter
\AtBeginDocument{%
  \@ifundefined{shortintertext@}{}{%
    \ExplSyntaxOn
    \bbl@sreplace\shortintertext@{\normalbaselines}%
      {\normalbaselines\ifx\bbl@eqnodir\relax\else\bbl@pardir\@ne\bbl@eqnodir\fi}%
    \ExplSyntaxOff}}
\makeatother
% geometry inherits the canonical A4 sheet from \documentclass; build.sh then
% verifies every rendered MediaBox and CropBox, so a silent Letter fallback is
% impossible to publish as a clean default build.
\usepackage[\lakitgeometryoptions]{geometry}
\linespread{\lakitlinespread}       % layout hook; canonical default remains 1.3
\ifx\lakitparagraphstyle\hllpdfparagraphblock
  \setlength{\parindent}{0pt}
  \setlength{\parskip}{6pt plus 1pt minus 1pt}
\fi
% These dimensions remain ordinary scoped TeX lengths: authors can customize
% them after core input. Language switches inherit them; lists keep their own
% indentation and narrow reference layouts establish their local paragraph gaps.
\typeout{HLLPAR|\lakitparagraphstyle}
\usepackage{xcolor}
\usepackage{mdframed}
\usepackage{listings}                 % code listings (assembly / C / pseudo-code)
\usepackage{tikz}
\usetikzlibrary{arrows.meta,positioning,calc,shapes.geometric,fit,backgrounds}
\usepackage{pgfplots}\pgfplotsset{compat=1.18}
\usepackage{quantikz}                 % quantum circuits (remove if unused)
\usepackage{booktabs}
\usepackage{array}
\usepackage{float}                    % [H] = place table exactly here (reference doc: no drifting)
\usepackage{enumitem}
\usepackage{needspace}

% ============================================================================
% ★★ v3.4 — THE DIRECTION-ISOLATION TRIO (the heart of mixing Hebrew + Latin)
% Mental model:
%   \h{...}    Hebrew (RTL) island      — Hebrew text inside an LTR context
%                                          (TikZ nodes, the english `tikzpic`).
%   \en{...}   English (LTR) isolate    — Latin text / parentheses inside a
%                                          Hebrew (RTL) context. KEEPS PARENS
%                                          AND ORDER CORRECT.
%   \code{...} \en + monospace          — code inside Hebrew / inside a frametitle.
% CRITICAL: never wrap Latin in \h{} (it is for HEBREW). Inside a TikZ node,
% Latin in \h{} renders REVERSED (Datapath->htapataD, CPU->UPC) and parens
% mirror. Isolate every Latin/paren chunk with \en{} or \code{}. See
% references/figures-boxes-listings.md.
% ---- TEXT INSIDE MATH (v3.8) ----------------------------------------------
% Put ANY text that sits inside math in a \text{}. The ONLY thing that is
% invalid in math is a bare \h/\en (= \foreignlanguage) used DIRECTLY, outside
% \text -> "\rmfamily invalid in math mode". Inside \text:
%   Hebrew label           -> \text{...}            (bare Hebrew renders fine;
%                              \h is optional — BUT NOT ALWAYS HARMLESS (v3.67):
%                              if the \h group OPENS with a nested island
%                              ($..$, \en{}, \code{}, \texttt{}, \mbox{}) the
%                              WHOLE group loses bidi and the Hebrew prints in
%                              LOGICAL order = gibberish. Measured:
%                                \text{\h{$f$ קבועה}}  -> reads "העובק"  BROKEN
%                                \text{\h{קבועה $f$}}  -> correct
%                                \text{$f$ קבועה}      -> correct (no \h)
%                              Hebrew first inside \h, or drop the \h. Flagged
%                              H-OPENS-WITH-LTR-GROUP.)
%   plain Latin word        -> \text{...}            (no marker needed)
%   Latin WITH ( ) or [ ]   -> \text{\en{...}}       (\en forces LTR; a plain
%                              \text mirrors the parens in RTL -> "Spec (TNR(" )
%   f(Hebrew/text arg)      -> \text{f}(\text{...})  -- keep the parens at MATH
%                              level, NOT \text{f(arg)} (parens inside the box mirror)
%   code in math            -> \mathtt{...}  (or \texttt{}, math-guarded below)
% NEVER write \en{...} or \h{...} straight into $...$ / \[..\] (outside \text);
% always via \text{}. (\mbox{...} also works but does not scale in scripts.)
% ============================================================================
\newcommand{\h}[1]{\foreignlanguage{hebrew}{#1}}
% v3.42.1 — \en islands do NOT hyphenate (auto OR at explicit hyphens). English hyphenation patterns were
% tearing identifiers and terms ACROSS THE RTL SEAM ("IO Regis-/ters",
% "In-/put", "Band-/width", "Per-/formance", "Multi-/ple" — all shipped on
% real pages; probe-proven: "ALUS-/rcB", "Mem-/ToReg"). An island is a label
% or a term, not English prose — a hyphen tear there is never right. A/B on
% a full 197pp lakit: identical overfull sets, same page count, ~11
% paragraphs healed, one mild underfull (badness 3525) that replaced a
% mutilated "Regis-/ters". Long English PROSE belongs in the englishprose
% environment below, which restores explicit-hyphen breaks while its paragraph
% is still active. A raw otherlanguage block inherits the global prohibition.
% (\l@nohyphenation needs \csname:
% '@' is not a letter here.)
% v3.43.2 — \exhyphenpenalty moved OUT of \en to a GLOBAL assignment below.
% Render-proven: a group-local \exhyphenpenalty is DEAD CODE — TeX reads the
% line-breaking parameters at \par time, after the \en group has closed, so
% \en{write-through} / plain not-taken still tore at the "-" ("write-/
% through" shipped on a real exam page). \language is different: it is
% recorded in the node list as text is added, so the in-group nohyphenation
% assignment above DOES work (and stays). The global penalty was A/B-measured
% on a 197pp lakit: identical overfull set (25), identical page count, the
% not-/taken tear healed. Hebrew is unaffected in practice — Hebrew compounds
% use maqaf (־), not ASCII "-".
\newcommand{\en}[1]{\foreignlanguage{english}{\language=\csname l@nohyphenation\endcsname #1}}
% Babel runs extras locally for inline islands and paragraph language changes.
% An English main document must not inherit the Hebrew no-hyphen paragraph rule.
\addto\extrashebrew{\exhyphenpenalty=10000\relax}
\addto\extrasenglish{\exhyphenpenalty=50\relax}
\ifx\lakitmainlanguage\hllpdfenglish
  \exhyphenpenalty=50\relax
\else
  \exhyphenpenalty=10000\relax
\fi
% Full English paragraphs need ordinary explicit-hyphen breaks. Use this
% environment instead of a raw otherlanguage block; the paragraph ends while
% the local penalty is still active.
\newenvironment{englishprose}
  {\par\begin{otherlanguage}{english}\exhyphenpenalty=50\relax}
  {\par\end{otherlanguage}}
\newenvironment{hebrewprose}
  {\par\begin{otherlanguage}{hebrew}\exhyphenpenalty=10000\relax}
  {\par\end{otherlanguage}}
% Use for shell text whose locale must not depend on the last body paragraph.
\DeclareRobustCommand{\doclang}[1]{\foreignlanguage{\lakitmainlanguage}{#1}}

% ============================================================================
% ★★ — fixes for defects that are INVISIBLE in the log and in the checks
% and only surface on a VISUAL RENDER of the PDF. (Each verified empirically.)
% ----------------------------------------------------------------------------
% v3.22: \emph is now REAL slanted emphasis (FakeSlant italic face defined on
% the \babelfont{rm} line above), replacing the v3.21 \emph->\textbf hack: it
% preserves \emph's toggle semantics and keeps bold free for its own role.
% Prefer bold-as-emphasis for your doc? one line: \renewcommand{\emph}[1]{\textbf{#1}}
% page number: force the digits LTR so they never reverse (18 -> 81) in the RTL
%   footer. That reversal is a NON-DETERMINISTIC bidi glitch (observed once in
%   ~80 pp); this LuaLaTeX build is not byte-reproducible, so the defensive
%   LTR wrap is the only reliable guard. Latin digits are always LTR anyway.
\renewcommand{\thepage}{\foreignlanguage{english}{\arabic{page}}}
% v3.70: \pagenumbering \gdef's \thepage and silently DESTROYS the wrap above —
% after \pagenumbering{arabic} the folio guard is gone for the whole rest of the
% document. Two collections discovered this independently and hand-restore the
% wrap in main.tex with the same comment; that knowledge now lives here instead.
% The wrapper re-applies the LTR wrap after the kernel call, for any style.
\makeatletter
\let\lakin@orig@pagenumbering\pagenumbering
\renewcommand{\pagenumbering}[1]{%
  \lakin@orig@pagenumbering{#1}%
  \renewcommand{\thepage}{\foreignlanguage{english}{\csname @#1\endcsname\c@page}}}
\makeatother
% emergencystretch: Hebrew does not hyphenate, so an unbreakable LTR island
%   (\en{...} / an inline math run) landing at the margin can leave a small
%   (~2-13pt) overfull with no legal breakpoint. This rescues ONLY those
%   inline-prose tie-breaks; it is a ceiling used only on offending lines and
%   does NOT touch display-math / listing / table overflow (those still surface
%   in the log for the content-level overflow pass). Verified: baseline overfull
%   -> 0 with this set, with no change to pagination or underfull.
\setlength{\emergencystretch}{3em}
% v3.40: inline math (almost) never breaks across lines. TeX's default
%   (binop 700 / rel 500) lets an inline run split at an operator —
%   $(1.f)\cdot 2^{e}$ shipped as "(1.f)·" ending one line and "2^e" opening
%   the next (reader-reported). In LTR that split is merely ugly; under an
%   RTL base the carried operator lands on the visually WRONG side of the
%   break and the formula reads backwards.
%   WHY 10000 (v3.45.0 — the display-first resolution). Inline math is
%   FORBIDDEN from breaking, so a torn island under an RTL base is
%   structurally impossible. The v3.40 cycle used 9999 ("break only as a last
%   resort") because a bare 10000 left ~30-39 WIDE inline equations
%   overflowing (worst 90.7pt) — but 9999's "last resort" fires OFTEN in
%   tight, non-hyphenating Hebrew: a reader found torn islands across the
%   fleet (~150 candidates; e.g. "P(y |" ending one line and "x_0)" opening
%   the next, read backwards). The step the 9999 cycle missed: every one of
%   those ~30 overflows is a wide equation that BELONGS in display. Convert
%   them to display FIRST, then 10000 gives 0 tears AND 0 overflows (measured,
%   12/12 fleet). Doctrine: inline math is ATOMIC; a wide equation goes in a
%   display, never inline. See content-style ("inline math is atomic").
\binoppenalty=10000 \relpenalty=10000
% v3.38: uniform sentence spacing. TeX's default (\nonfrenchspacing, sfcode
% 3000) is the OLD Anglo-American wide-space-after-period convention — Hebrew
% typography never had it, and modern Latin practice (Chicago et al.) dropped
% it too. Measured on real pages: post-period gaps up to 1.9-2.1x the line's
% mean word gap — the dominant "rivers" source in justified Hebrew prose.
% Survives babel language switches (probed: sfcode stays french inside \en{}
% islands and otherlanguage blocks).
\frenchspacing
% v3.38: widow/orphan protection. LaTeX's default penalties (150/150/50)
% allow a lone line at a page top/bottom almost freely. 10000 = forbid;
% this is SAFE here precisely because the document is \raggedbottom (the
% demoted line just moves on, no vertical stretch fight) — under flushbottom
% these values would cause underfull pages. raggedbottom is the report-class
% default; made explicit so a future class/option change cannot silently
% flip it and turn the 10000s harmful.
\raggedbottom
\widowpenalty=10000 \clubpenalty=10000 \displaywidowpenalty=10000
% em-dash: --- does NOT become — under Frank Ruhl (the font lacks tlig; verified
%   the width is identical with/without Ligatures=TeX, while the mechanism works
%   on Latin Modern). TYPE THE CHARACTER — (U+2014) DIRECTLY; likewise – (U+2013).
%   check_bidi_figures.py flags a --/--- that is adjacent to Hebrew.
% ============================================================================

% ---- STRUCTURAL DIRECTION-SAFETY for \texttt (v3.7) -------------------------
% A parenthesis, bracket, or $ INSIDE \texttt{} mirrors/migrates in RTL prose
% (MIPS 0($s0)->"(s0$)0", $rs->"rs$", arr[i]->"[arr[i") because \texttt sets the
% font but NOT the direction. Redefining \texttt to LTR-isolate its content fixes
% the WHOLE class at once (parens, brackets, $) so plain \texttt{lw $t0,0($s0)}
% is safe — no need to remember \code{}. Guarded by \ifmmode because
% \foreignlanguage is invalid in math (there, plain monospace is already LTR).
% NB (corrected v3.56 — the old note here said "stays tofu", which is WRONG):
% \texttt{Hebrew} is a FATAL build failure. The tt family is DejaVu Sans Mono
% (zero Hebrew) and the tt fallback hebfb is DejaVuSerif (also zero Hebrew), so
% the subsetter gets an empty glyph set and LuaTeX dies with
%   "! error: (file .../DejaVuSansMono.ttf) (type 2): there are no glyphs in the
%    subset"  -> no output PDF.
% The message names a FONT FILE, never Hebrew/\texttt/the source line. And it is
% non-local: add an unrelated Latin \code{} elsewhere and the fatal degrades into
% a SILENT DELETION of the Hebrew word. Put Hebrew in \h{}/normal text, never in
% \texttt{}/\code{}; flagged HEBREW-IN-MONO. To disable, delete this block.
\let\origtexttt\texttt
\protected\long\def\texttt#1{\ifmmode\origtexttt{#1}\else\foreignlanguage{english}{\origtexttt{#1}}\fi}
% v4.8.0: defer the direction test until TeX has inserted an alignment cell's
% math template. Unprotected expansion here tested \ifmmode too early at the
% FIRST token of aligned/array/cases/matrix cells. Engine-level \protected also
% preserves the original \texttt implementation without the self-recursion
% observed when retrofitting \DeclareRobustCommand to its saved alias.
\protected\long\def\code#1{\ifmmode\origtexttt{#1}\else\foreignlanguage{english}{\origtexttt{#1}}\fi}

% ---- \codeline — display-code line (v3.24): the code analogue of display math.
% Reader-friction rule (content-style.md §5b): inline \code{} is for a SINGLE
% ATOM only (one register/variable/mnemonic, a tight no-space expression).
% A multi-token expression — an instruction with operands, a spaced assignment,
% a register-binding list — zigzags the eye RTL<->LTR mid-sentence and, being
% an unbreakable island, line-breaks badly. Give it its own centered line with
% display-math vertical rhythm. \centerline = one box => the expression never
% wraps; \nopagebreak keeps it glued to its introducing sentence.
%   ...the Hebrew sentence ends with a colon:
%   \codeline{lw \$t1, -4(\$t0)}
%   and the prose continues (escape dollars as \$).
% v3.66 (first line of the body): flush a PENDING \item label before the \par.
% LaTeX defers the label box (\@labels) until the next \everypar, and an \hbox in
% vertical mode never fires \everypar — so `\item \codeline{...}` printed the code
% at y=69.1 and the item's own "1." BELOW it at y=103.6 (the number appears after
% its content, attached to nothing). \if@inlabel\leavevmode\fi is CONDITIONAL on
% purpose: an unconditional \leavevmode would open an empty paragraph line when
% \codeline follows ordinary prose. Verified both ways — in a list the label now
% sits on the code's line at its normal x, and after prose the geometry is
% byte-identical to before.
\makeatletter
\newcommand{\codeline}[1]{%
  \if@inlabel\leavevmode\fi
  \par\nopagebreak
  \vspace{\abovedisplayskip}\nopagebreak
  % v3.62: was \centerline{...}, i.e. \hbox to\hsize{\hss #1 \hss}. \hss has
  % INFINITE SHRINK, so TeX could always "fit" the box and NEVER reported an
  % Overfull — the same \hss mechanism as BLD-27's \item label. Measured: a wide
  % \codeline in a defbox painted its first token ON the box's left border bar and
  % sat 9.5pt outside the frame interior on both sides, with 0 Overfull, 0 Underfull
  % and all four gates PASS; in bare prose it reached 21.0pt into BOTH margins.
  % \hfil has no shrink, so the overflow is reported and build.sh's ovf_vis/ovf_big
  % counters see it; \linewidth (not \hsize) makes it box-aware, so a \codeline
  % inside a colored box centres on the box interior. Verified: a FITTING \codeline
  % keeps byte-identical glyph positions, and the box-crossing case now emits
  % "Overfull \hbox (19.03pt too wide)". A wide \codeline should be split in two.
  \hbox to\linewidth{\hfil\code{#1}\hfil}%
  \nopagebreak\vspace{\belowdisplayskip}}
\makeatother

% ---- a thin decorative rule helper ----
\newcommand{\sep}{\par\vspace{2pt}\noindent\textcolor{cRule}{\rule{\linewidth}{0.6pt}}\par\vspace{4pt}}

% ---- colors ----
\definecolor{cBlue}{HTML}{1F4E79}   \definecolor{cBlueBg}{HTML}{EAF1F8}
\definecolor{cGreen}{HTML}{2E6B4F}  \definecolor{cGreenBg}{HTML}{EAF4EE}
\definecolor{cAmber}{HTML}{9C6A00}  \definecolor{cAmberBg}{HTML}{FBF3E2}
\definecolor{cPurple}{HTML}{5B2C6F} \definecolor{cPurpleBg}{HTML}{F3ECF7}
\definecolor{cTeal}{HTML}{0E6E6E}   \definecolor{cTealBg}{HTML}{E6F4F4}
\definecolor{cRed}{HTML}{8E2B2B}    \definecolor{cRedBg}{HTML}{FBEAEA}
\definecolor{cGray}{HTML}{555555}   \definecolor{cRule}{HTML}{1F4E79}

% ---- math operators ----
\DeclareMathOperator{\Tr}{Tr}
\DeclareMathOperator{\rank}{rank}
\DeclareMathOperator{\diag}{diag}
\DeclareMathOperator{\sgn}{sgn}

% ---- bra-ket (quantikz defines \ket/\bra/\braket; add the rest safely) ----
\providecommand{\ket}[1]{\left|#1\right\rangle}
\providecommand{\bra}[1]{\left\langle#1\right|}
\providecommand{\braket}[2]{\left\langle#1\middle|#2\right\rangle}
\newcommand{\dyad}[1]{\left|#1\right\rangle\!\left\langle#1\right|}
% v3.21: general outer product |a><b| and expectation value <A> — agents reach
% for these by reflex (esp. \expval); quantikz defines neither, so their absence
% threw "Undefined control sequence" in a real build. \providecommand = safe.
\providecommand{\ketbra}[2]{\left|#1\right\rangle\!\left\langle#2\right|}
\providecommand{\expval}[1]{\left\langle#1\right\rangle}
\newcommand{\abs}[1]{\left|#1\right|}
\newcommand{\norm}[1]{\left\lVert#1\right\rVert}
\newcommand{\half}{\tfrac12}
\newcommand{\R}{\mathbb{R}}\newcommand{\C}{\mathbb{C}}\newcommand{\Z}{\mathbb{Z}}
% identity operator. NB: \mathbb{1} silently renders a WRONG glyph (msbm has no
% blackboard digits). Use \mathbb{I} = 𝕀 (works with amssymb, no extra package).
% Do NOT "fix" this back to \mathbb{1}.
\newcommand{\id}{\mathbb{I}}

% ---- physical units (v3.35): \punit — upright unit with thin space ---------
% Usage: $300\punit{K}$, $80\punit{mV/decade}$, $10^{7}\punit{cm/s}$ — the
% NUMBER AND UNIT LIVE IN ONE MATH ISLAND. Promoted verbatim from a 199-page
% production lakit (advanced-electronic-devices, dozens of uses).
% Two measured traps this design avoids (both render-proven):
%   1. siunitx SCRAMBLES under RTL — \SI{3.5}{eV} renders "eV 3.5", exponents
%      detach ("Hz 10 x 10^9"): its quantity assembly is not bidi-isolated.
%      Do NOT load siunitx in this pipeline.
%   2. A text-mode unit next to a bare-prose number is TWO adjacent LTR runs
%      and reorders ("5 mm" -> "mm5" — the BIDI-2/3 class). \punit is therefore
%      math-only ON PURPOSE: prose misuse fails LOUDLY at compile (\mathrm
%      outside math) instead of scrambling silently.
% Domain shorthands (\MeV, \fm...) stay per-document: \newcommand{\MeV}{\,\text{\en{MeV}}}
\providecommand{\punit}[1]{\,\mathrm{#1}}

% ============================================================================
% Colored boxes — mdframed (NOT tcolorbox, which breaks under bidi).
% ★★ v3.4: innerleftmargin/innerrightmargin give code listings room so their
% frame/line-numbers do not spill out the box border (see listing styles below).
% ============================================================================
% v3.61: everyline/split skips were applied to `tablebox` only (BLD-14, v3.37) —
% the six colored boxes below never got them, so they still split OPEN across a page
% break: the page-N segment gets top+sides and NO bottom rule, the page-N+1 segment
% sides+bottom and NO top rule. Measured on the shipped fleet this is a COVERAGE GAP,
% not a re-report of BLD-14: 0 split-open `tablebox` segments (that fix holds) vs 62
% split-open COLORED boxes across 7 documents. Verified after the fix: both segments
% close, and the frametitle is NOT duplicated on the continuation. BLD-14's caveat
% still applies — a box taller than a full page grinds mdframed's split-trial loop, so
% split such content manually.
% ============================================================================
% v3.72: frametitle={} leads EVERY \newmdenv option list below. mdframed options
% are GROUP-scoped: a box nested inside \begin{exbox}[frametitle={...}] used to
% INHERIT the outer title — measured, the identical title string rendered twice,
% the second riding a title-bar fill inside the INNER box's frame, for all seven
% environments in both directions. Both natural forms were broken (= preamble
% bug, not an authoring rule). The explicit reset makes every box titleless
% unless ITS OWN [frametitle={...}] says otherwise; verified both directions:
% no duplicate on the nested box, an explicit title still renders.
% Geometry is shared with the optional semantic-width table container. Themes
% select their own existing semantic style below; neither path fixes columns.
\mdfdefinestyle{hllpdfsemanticgeometry}{linewidth=1.1pt,
  innertopmargin=6pt,innerbottommargin=7pt,skipabove=9pt,skipbelow=9pt,
  innerleftmargin=9pt,innerrightmargin=9pt,everyline=true,
  splitbottomskip=6pt,splittopskip=12pt,roundcorner=3pt}
\newcommand{\hllpdfsemanticframestyle}{hllpdfsemanticgeometry}
\newmdenv[frametitle={},style=hllpdfsemanticgeometry,linecolor=cBlue,backgroundcolor=cBlueBg,
  frametitlebackgroundcolor=cBlue,frametitlefontcolor=white,
  ]{defbox}   % הגדרה
\newmdenv[frametitle={},style=hllpdfsemanticgeometry,linecolor=cGreen,backgroundcolor=cGreenBg,
  frametitlebackgroundcolor=cGreen,frametitlefontcolor=white,
  ]{thmbox}   % משפט/טענה
\newmdenv[frametitle={},style=hllpdfsemanticgeometry,linecolor=cAmber,backgroundcolor=cAmberBg,
  frametitlebackgroundcolor=cAmber,frametitlefontcolor=white,
  ]{notebox}  % הערה
\newmdenv[frametitle={},style=hllpdfsemanticgeometry,linecolor=cPurple,backgroundcolor=cPurpleBg,
  frametitlebackgroundcolor=cPurple,frametitlefontcolor=white,
  ]{exbox}    % דוגמה
\newmdenv[frametitle={},style=hllpdfsemanticgeometry,linecolor=cRed,backgroundcolor=cRedBg,
  frametitlebackgroundcolor=cRed,frametitlefontcolor=white,
  ]{warnbox}  % טעות נפוצה
\newmdenv[frametitle={},style=hllpdfsemanticgeometry,linecolor=cTeal,backgroundcolor=cTealBg,
  frametitlebackgroundcolor=cTeal,frametitlefontcolor=white,
  ]{keybox}   % רעיון מרכזי

% ---- codebox (v4.5.0): the dedicated container for code -------------------
% An OPTION, not a policy: asmblock unboxed stays canonical for a full listing,
% and a snippet inside any of the six colored boxes still works as before.
% What makes codebox worth its own environment is that it is the ONE framed box
% that keeps its margin line numbers.
%
% v4.4.0 shipped it WITHOUT them, on the strength of v3.37's "no margin knob can
% save them" plus four fresh probes that agreed. That was wrong: every probe had
% moved the code, the box, or both — never the key that carries the NUMBER
% COLUMN. Measured on the page:
%   resetmargins=false   the listing follows THIS box instead of the page,
%                        moving the code from x=75.26 to x=97.54
%   xrightmargin=26pt    under RTL this is the PHYSICAL LEFT indent, and the
%                        only key that carries the \llap'ed digits with it:
%                        x=64.28 -> 90.18, clearing the rule at 65.64 by 23.6pt
%   xleftmargin=26pt     moves the code and LEAVES THE DIGITS on the rule —
%                        the asymmetry that hid the fix for eleven versions
% resetmargins is hoisted out of the four styles (see the \lstset below them) so
% this override is possible at all; every other listing still resolves to true.
% Pinned by tests/codebox_numbers.tex + probe_lstnum.py --inside-rule, whose
% negative control strips xrightmargin and watches the digits fall back.
\definecolor{cSlate}{HTML}{3F4A56}  \definecolor{cSlateBg}{HTML}{F2F4F7}
\newmdenv[frametitle={},leftline=true,rightline=true,topline=true,bottomline=true,
  linecolor=cSlate,backgroundcolor=cSlateBg,linewidth=0.9pt,
  frametitlebackgroundcolor=cSlate,frametitlefontcolor=white,
  innertopmargin=6pt,innerbottommargin=7pt,skipabove=9pt,skipbelow=9pt,
  innerleftmargin=9pt,innerrightmargin=9pt,everyline=true,splitbottomskip=6pt,splittopskip=12pt,
  roundcorner=3pt]{codebox}   % קוד
% The two keys that make in-box numbering work, both measured on the page:
%   resetmargins=false  the listing follows THIS box instead of the page
%   xrightmargin        under RTL this is the PHYSICAL LEFT indent, and it is
%                       the only one that carries the \llap'ed number column
%                       with it. xleftmargin moves the code and leaves the
%                       digits behind on the rule (measured: code 97.54,
%                       digits 64.28 — the v3.37 defect, reproduced); with
%                       xrightmargin=26pt the digits land at x=90.18, a
%                       clearance of 23.6pt from the rule at 65.64.
% This asymmetry is why v3.4's note had the semantics backwards and v3.37
% concluded the digits could not be saved.
% <hllpdf:listing-container-state>
% v4.29: all state is local to the current environment group. A semantic
% callout masks number text, while a nested codebox restores the formatter
% saved before that mask. Save it once across successive semantic wrappers;
% an explicit numbers=none (or an explicitly empty formatter) remains valid.
\makeatletter
\newif\ifhllpdf@incodebox
\newif\ifhllpdf@englishcodebox
\newif\ifhllpdf@lstnumbermasked
\let\hllpdf@savedlstnumber\thelstnumber
\makeatother
% </hllpdf:listing-container-state>
% <hllpdf:codebox-gutter>
\makeatletter
\AtBeginEnvironment{codebox}{%
  \hllpdf@incodeboxtrue
  % Snapshot the OUTER language before asmblock selects English.
  \iflanguage{english}{\hllpdf@englishcodeboxtrue}{\hllpdf@englishcodeboxfalse}%
  \ifhllpdf@lstnumbermasked
    \let\thelstnumber\hllpdf@savedlstnumber
    \hllpdf@lstnumbermaskedfalse
  \fi
}
% listings' AddToHook is GLOBAL. Register once, and gate on local box state;
% adding an unconditional hook at every codebox entry leaks into later listings.
\lst@AddToHook{PreInit}{%
  \ifhllpdf@incodebox
    % Apply every listing override here, inside the LISTING's own group.
    % Setting margins at box entry leaks them into nested semantic callouts.
    \lstset{resetmargins=false,xrightmargin=26pt}%
    % Styles resolve before PreInit. In LTR the physical number gutter follows
    % xleftmargin; in Hebrew retain the measured inherited xrightmargin path.
    \ifhllpdf@englishcodebox
      \lstset{xleftmargin=26pt,xrightmargin=0pt}%
    \fi
    % ONE box, not two: suppress only the frame/background/skips of a listing
    % inside this codebox. The next unboxed listing keeps its own style.
    \lstset{frame=none,backgroundcolor=\color{cSlateBg},aboveskip=0pt,belowskip=0pt}%
  \fi
}
\makeatother
% </hllpdf:codebox-gutter>

% v3.37: line numbers inside a COLORED BOX are sliced by the box border —
% measured: digit span x=63.5-67.3 vs border line at x=65.2 (the border runs
% THROUGH the digits), and the inter-border gap (3.5pt) is narrower than a
% digit (3.8pt), so no margin knob can save them (under resetmargins,
% xleftmargin is inert in-box — measured: code at x=75.26 with 2.2em AND
% 3.0em). A 3-6 line boxed snippet does not need line numbers; blank the
% number TEXT inside every colored box (\thelstnumber is what gets printed;
% the style's numbers=left survives, so nothing else shifts — code position
% byte-frozen at x=75.26, verified). Unboxed listings keep their numbers.
% v3.78: the blanking is coupled to the FRAME, not to the box name —
% clean_theme.tex registers LATER hooks (registration order wins) that
% restore the number text inside the boxes it renders frameless.
% v4.5.0 RETRACTS the "no margin knob can save them" clause above. It holds only
% under resetmargins=true (the shipped default): there xleftmargin is inert
% in-box, as measured. With resetmargins=false, xrightmargin DOES carry the
% \llap'ed digit column clear of the rule — measured; the codebox \lstset that
% does it is the one above, in the codebox block. The loop below therefore
% covers SIX boxes, not "every colored box": codebox opts out and keeps its
% numbers.
% <hllpdf:semantic-listing-mask>
\makeatletter
\@for\HEB@lstbox:=defbox,thmbox,notebox,exbox,warnbox,keybox\do{%
  \expandafter\AtBeginEnvironment\expandafter{\HEB@lstbox}{%
    \ifhllpdf@lstnumbermasked\else
      \let\hllpdf@savedlstnumber\thelstnumber
    \fi
    \hllpdf@lstnumbermaskedtrue
    \def\thelstnumber{}%
    % The nearest semantic container owns its listing presentation. On exit,
    % TeX's group restores the enclosing codebox state and formatter.
    \hllpdf@incodeboxfalse}}
\makeatother
% </hllpdf:semantic-listing-mask>

% ============================================================================
% Unified TABLE presentation: subtle uniform slate frame + auto-numbered
% caption rendered INSIDE the frame. Use:
%   \begin{tablebox}\centering  <tabular>  \tabcap{<description>}\end{tablebox}
% Replaces the buggy \caption* (which emitted a stray "טבלה N : *" line) and
% provides an explicit framed option + numbered caption "טבלה N.M: תיאור".
% Optional intro text may be placed at the TOP, before the tabular, inside the box.
% ============================================================================
\definecolor{cTblLine}{HTML}{B9C4D2}  \definecolor{cTblBg}{HTML}{FAFBFC}
% v3.23: class-aware — under `article` there is no chapter counter (the promised
% report->article switch used to die with "No counter 'chapter' defined").
\makeatletter
\@ifundefined{c@chapter}{%
  % v3.31: in article the counter resets per SECTION, so an unprefixed number
  % repeated ("טבלה 1." in every section — measured on a two-section render).
  % Same conditional idea as the report branch: prefix the section number once
  % a real \section exists; a pre-section table stays "טבלה 1".
  \newcounter{tbl}[section]
  \renewcommand{\thetbl}{\ifnum\value{section}>0 \thesection.\fi\arabic{tbl}}%
}{%
  \newcounter{tbl}[chapter]
  % v3.30: a report-class doc that uses SECTIONS ONLY (chapter stays 0 —
  % the common short-doc pattern) used to print "טבלה 0.1"; the chapter
  % prefix now appears only once a real \chapter exists (render-verified
  % both ways: "טבלה 1." chapterless, "טבלה 1.1" with chapters).
  \renewcommand{\thetbl}{\ifnum\value{chapter}>0 \thechapter.\fi\arabic{tbl}}%
  % v3.64: the SAME leak in \thesection itself, left behind by v3.30. That fix
  % guarded the CAPTION counter and stopped there, so a chapterless report still
  % printed the heading "0.1 סעיף ראשון" and \ref to it printed "0.2" — measured:
  %   p1 y=65.2  '…0.1'   p1 y=144.9 '…0.2'   p1 y=183.4 '…0.2.1'
  % report.cls hardcodes \thesection as \thechapter.\@arabic\c@section, and the
  % leading 0 is the chapter counter that no \chapter ever incremented. Fixing
  % \thesection fixes the whole chain — \thesubsection and \thesubsubsection are
  % defined in terms of it. \theequation and \thetbl/\thefig were already guarded,
  % which is exactly why this survived: every NUMBER around the heading looked
  % right, so the heading's own 0 read as a style choice rather than a bug.
  % \appendix is safe: report.cls resets chapter to 0 and redefines \thechapter to
  % \@Alph\c@chapter, so an appendix \chapter makes the counter 1 and the guard
  % emits "A.1" unchanged. Render-verified in all four modes (chapterless,
  % chaptered, appendix+chapter, appendix+section-only).
  \renewcommand{\thesection}{\ifnum\value{chapter}>0 \thechapter.\fi\arabic{section}}%
}
\makeatother
% v3.37: everyline=true — if the box DOES split across pages, every segment
% gets a CLOSED frame (measured before: the split ran open-ended into the page
% edge with no bottom border, and the caption could strand alone at the next
% page top — "טבלה A.2" shipped that way). The split skips give the segments
% breathing room. In practice, with the caption GLUE below, every previously-
% split table in the corpus now moves WHOLE instead (the white-gap tradeoff,
% documented benign); everyline stays as the safety net for a genuinely
% page-tall table. NB: a table taller than ONE FULL PAGE can grind mdframed's
% split-trial loop (reproduced: minimal doc, minutes of 1pt-step Overfull
% probes) — break such a monster into two tables.
\mdfdefinestyle{hllpdftablegeometry}{leftline=true,rightline=true,topline=true,bottomline=true,
  linewidth=0.9pt,
  innertopmargin=9pt,innerbottommargin=8pt,skipabove=11pt,skipbelow=11pt,
  innerleftmargin=11pt,innerrightmargin=11pt,roundcorner=4pt}
\newcommand{\hllpdftableframestyle}{hllpdftablegeometry}
\newmdenv[frametitle={},style=hllpdftablegeometry,
  linecolor=cTblLine,backgroundcolor=cTblBg,
  % v4.2.0 (adopted): tablebox is ATOMIC. Short tableboxes near page bottoms
  % repeatedly entered mdframed's split trial, emitting benign Overfull-vbox
  % bursts and occasionally separating the intro prose from the tabular.
  % nobreak=true moves the complete short table to the next page instead. The
  % output-active Overfull-vbox gate still catches a box taller than a page;
  % everyline/split skips remain as defensive styling should that path ever
  % re-arm, but the canonical path never publishes a split tablebox.
  nobreak=true,
  everyline=true,splitbottomskip=6pt,splittopskip=12pt,
  % v3.79 (BLD-37 table side, correction): startinnercode/endinnercode run
  % INSIDE mdframed's body vbox (mdframed.sty:848 / :856+863) -- the second one
  % after \par\unskip, i.e. after line breaking, which is the only moment the
  % \tabcap anchor's true offset from the top of the box EXISTS. See
  % \lakintabboxclose below. These two keys are the ONLY tablebox change.
  startinnercode={\lakintabboxopen},endinnercode={\lakintabboxclose}]{tablebox}
% v3.37: \nopagebreak GLUE — \vsplit honors penalties, so the caption can no
% longer strand alone at the top of the next page (it takes its rows along, or
% the whole box moves). Render-verified on the real shipped orphan.
% v3.59: \NR@gettitle is what nameref injects into the stdlib title-bearing
% commands (\@sect, \@chapter, \caption) to fill \@currentlabelname. \figcap and
% \tabcap are INDEPENDENT macros — they never called it, so \@currentlabelname kept
% whatever the last heading left and \nameref{fig:x} silently printed the enclosing
% SECTION's title instead of the caption (proven in the .aux: field 3 of \newlabel
% was the heading text while field 1 and the anchor `fig.1.1` were correct — so this
% is a THIRD, distinct head next to BLD-13's counter fix and BIDI-38's autoref-name
% fix, neither of which writes that field). \@ifundefined keeps it inert if
% hyperref/nameref is ever dropped. Guarded helper because \figcap is defined
% after \makeatother.
% NB both helpers below need \makeatletter — this part of the preamble is OUTSIDE
% any letter-catcode region, so an unwrapped \@currentlabel parses as LaTeX's \@
% and dies with "You can't use \spacefactor in vertical mode" (measured).
\makeatletter
\newcommand{\lakinrtitle}[1]{\@ifundefined{NR@gettitle}{}{\NR@gettitle{#1}}}
% v3.59: \hitem{<letter>}{<key>} — a REFERENCEABLE Hebrew-lettered list item.
%   \begin{enumerate} \hitem{א}{it:a} הסעיף הראשון. \end{enumerate}  ->  \ref{it:a} = (א)
% Why this exists: automatic Hebrew list labels fail silently, so SKILL.md mandates
% an EXPLICIT \item[(\textbf{א})] — but LaTeX puts \refstepcounter INSIDE
% \if@noitemarg (latex.ltx:12964), and \item[...] never sets that flag. So an
% explicit item steps no counter and never touches \@currentlabel: a \label there
% binds to the AMBIENT one and \ref prints the CHAPTER/SECTION number where the
% item letter belongs (proven in the .aux: {chapter.1}). Same doctrine as \figcap —
% own the label state before \label. \phantomsection matters: without it the
% hyperlink still lands on the chapter head (measured); with it the anchor is
% section*.N and the bookmark outline stays clean.
\newcommand{\hitem}[2]{\item[(\textbf{#1})]%
  \def\@currentlabel{(#1)}\def\@currentlabelname{#1}\phantomsection\label{#2}}
% NB (v3.66): \hitem is \ref-only. \autoref{key} calls it a SECTION ("סעיף (א)")
% because \phantomsection brands the anchor section*.N, and \nameref{key} prints
% the bare letter. For \autoref-able items prefer \hebenumlabels below.
%
% v3.66: \hebenumlabels — OPT-IN automatic Hebrew enumerate letters, done right.
%   {\hebenumlabels\begin{enumerate}\item\label{q:a} ...\end{enumerate}}
%   -> labels (א) (ב) (ג); \ref{q:a} = א; anchor Item.N (real, per-item).
% This retires the belief that made \item[(\textbf{א})] + \hitem necessary.
% "Automatic Hebrew labels fail silently" was TRUE but unexplained, and the
% mechanism is now bisected (plain article vs +enumitem, measured): ENUMITEM
% PRE-EXPANDS the label expression in an \edef BEFORE \refstepcounter, with the
% counter still at 0. Tokens it shields by name (\arabic,\alph,\roman) survive;
% a plain \newcommand \ifcase wrapper evaluates at 0 -> the EMPTY zero-branch ->
% the label prints as "()" with the letter GONE (and \the\c@enumi prints a
% literal 0, \number\value dies fatally pointing at \item). The cure is one
% property, not a doctrine: the wrapper must be ROBUST (survives the \edef
% unexpanded), while \theenumi stays EXPANDABLE (so \refstepcounter writes the
% real letter to the .aux). Verified: labels (א)(ב)(ג), \ref -> א/ב/ג each item
% its own value, .aux = {א}{...}{Item.1} — versus the explicit-\item doctrine
% where BIDI-50/\hitem exist purely to patch the counter that never stepped.
\def\@hebalph#1{\ifcase#1\or א\or ב\or ג\or ד\or ה\or ו\or ז\or ח\or ט\or י\or
  יא\or יב\or יג\or יד\or טו\or טז\or יז\or יח\or יט\or כ\else\@ctrerr\fi}
\DeclareRobustCommand*\hebalph[1]{\expandafter\@hebalph\csname c@#1\endcsname}
\newcommand{\hebenumlabels}{%
  \renewcommand{\theenumi}{\@hebalph\c@enumi}%
  \renewcommand{\labelenumi}{(\textbf{\hebalph{enumi}})}}
\makeatother
% v3.63: \mbox the LABEL. \thetbl ("7.1") and a caption description that STARTS
% WITH A DIGIT ("32 הרגיסטרים…") are two European-Number runs separated only by
% neutrals (the label's own "." and \enspace). Under the RTL paragraph base the UBA
% resolves those neutrals to the RTL level, so both numbers sit in ONE reversed span
% and EXCHANGE VISUAL POSITIONS. Measured on the shipped PDF: a table whose real
% number is 7.1 printed "טבלה 32" — a different, EXISTING table number — while \ref
% correctly printed 7.1, so the caption silently contradicted its own cross-reference,
% and the description lost its leading count. The source contains no adjacency at all;
% it is manufactured by this macro, which is why every source linter was blind.
% \mbox makes the label one atomic box so nothing can reorder across it. Verified
% zero-regression: an ordinary Hebrew caption keeps identical x-geometry.
% One paragraph style covers label, number, separator and description.
% Explicit project/template overrides remain possible; the label hook is kept
% for compatibility, with an identity default. Counters/anchors stay unchanged.
\newcommand{\lakitcaptionstyle}{\normalfont\small\color{black}\hypersetup{allcolors=black}}
\newcommand{\lakitcaptionlabelstyle}[1]{#1}
\newcommand{\lakitcaptionpunctuation}{:}
% Native floats, captionof and longtable use the same presentation policy.
% subcaption is optional and inherits the same font rather than a second size.
\usepackage{caption}
\DeclareCaptionFont{lakit}{\lakitcaptionstyle}
\DeclareCaptionLabelFormat{lakit}{\lakitcaptionlabelstyle{#1~\en{#2}}}
\DeclareCaptionLabelSeparator{lakit}{\lakitcaptionpunctuation\enspace}
\captionsetup{font=lakit,labelfont={},textfont={},labelformat=lakit,
  labelsep=lakit,justification=centering}
\makeatletter
\AtBeginDocument{\@ifpackageloaded{subcaption}{%
  \captionsetup[sub]{font=lakit,labelfont={},textfont={},labelsep=lakit}}{}}
\makeatother
% v3.64: OPT-IN list-of-figures/tables entry via an optional short caption.
% \listoffigures and \listoftables are inherited from report.cls and LOOK available,
% but \figcap/\tabcap step the custom \c@fig/\c@tbl counters and never call
% \addcontentsline — so both printed their heading and NOTHING ELSE. Measured: a doc
% with one \figcap and one \tabcap rendered page 1 as 'רשימת האיורים\n1\n', and the
% .lof held only \@chapter's \addvspace, zero \contentsline entries. An empty list
% under a correct Hebrew heading reads as "this document has no figures".
%
% WHY OPTIONAL, AND NOT AUTOMATIC. Writing #1 unconditionally was implemented first
% and REJECTED: a caption containing \codeline compiles CLEANLY today (verified on
% the pre-v3.64 preamble: 0 errors, PDF produced) but becomes FATAL once it is also
% written to the .lof — \codeline is \hbox to\linewidth{...}, and a full-width box in
% a list entry dies with "\pdfextension endlink cannot be used in vertical mode", no
% output PDF. Gaining a feature nobody uses (0 of 13 collections call \listoffigures)
% at the price of breaking a construct that works is a bad trade. Stock \caption has
% exactly this limitation and exactly this escape hatch, so the convention is familiar.
% \en{}, \code{}, \h{} and inline math in the short text are all safe — verified, they
% survive the write and typeset correctly in the list (RTL geometry confirmed: number
% rightmost at x=509.9, title left of it, dots and folio at the far left).
% With no optional argument the macro is byte-identical to v3.63 — zero risk to every
% existing document, and \listoffigures lists exactly the captions that opted in.
% BLD-37, TABLE side (the \figcap v3.74 idea, adapted): clicking a \ref to a
% table used to land MID-table — hyperref's automatic anchor inside
% \refstepcounter was emitted on the still-open horizontal line whose TAIL is
% the tabular (an hbox whose baseline is the table's vertical CENTER), so the
% jump scrolled the top half of the table above the viewport. The v3.74 entry
% deferred this side claiming "tablebox splits — no single box to measure";
% that deferral dissolves: the measurement target is the TABULAR itself at
% \tabcap ENTRY, while it is still the tail of the open list — the same
% situation \lakin@grabfig handles at an island's end (and the fleet has 0
% split tableboxes). So \lakin@grabtab runs as the FIRST token, before any
% typesetting can push the box away, and the anchor is re-emitted raised above
% the measured tabular. \lakintabraise is a SEPARATE register from
% \lakinfigraise (defined with it, below the LTR-island block): a captionless
% island intentionally leaves the fig register set, and a table must never
% inherit that stale figure height.
%
% WHY the anchor stays AT THE ENTRY SPOT and is NOT deferred into the caption
% paragraph the way \figcap defers it: at \par TeX DROPS a final glue item,
% and the automatic anchor's smash box sitting BEHIND the interword glue kept
% that glue alive — so every centered tabular has always sat half a space off
% true center. Deferring the anchor removed that box and was measured as 982
% differing glyph records: EVERY tabular line shifted 1.30pt (2.61pt for
% pattern 7's two minipages) while captions stayed put. Emitting the raised
% anchor at the old spot — through hyperref's own \Hy@raisedlink, exactly the
% macro the automatic path uses — keeps the line's node structure, and so the
% A/B glyph diff at 0, BY CONSTRUCTION. \Hy@raisedlink reads the raise from
% \HyperRaiseLinkDefault (default \baselineskip = the old mid-table dest); the
% group-local \edef freezes measured-height+42pt into it, and \smash inside
% \Hy@raisedlink hides any raise from the line's extent. \HyCnt@ProvideTheH-
% Counter keeps the dest NAME byte-identical to the automatic one (it defines
% \theHtbl, which \hyper@makecurrent prefers — without it the first measured
% caption was named tbl.1 where the automatic path said tbl.0.1, measured).
% vmode entry, or anything but a box at the tail (no tabular to measure),
% falls back to the plain \refstepcounter — byte-identical to the old path.
\makeatletter
\newcommand{\tabcap}[2][]{\lakin@grabtab
  % v3.79: arm the box-close correction. Set unconditionally and in BOTH
  % branches below's scope, so the plain-\refstepcounter fallback is corrected
  % too; \lakintabboxopen clears it at every \begin{tablebox}, so a \tabcap
  % outside a tablebox can never leave it armed for a later box.
  \global\lakintabcapseentrue
  \ifdim\lakintabraise>0pt
    \H@refstepcounter{tbl}\HyCnt@ProvideTheHCounter{tbl}\hyper@makecurrent{tbl}%
    \begingroup
      % A frameless table targets its measured top. The frame's historical
      % cushion belongs only to tablebox and is corrected at box close.
      \edef\HyperRaiseLinkDefault{\the\lakintabraise}%
      \iflakintabframed
        \edef\HyperRaiseLinkDefault{\the\dimexpr\lakintabraise+42pt\relax}%
      \fi
      \Hy@raisedlink{\hyper@anchorstart{\@currentHref}\hyper@anchorend}%
    \endgroup
    \global\lakintabraise 0pt\relax
  \else
    \refstepcounter{tbl}%
  \fi
  \lakinrtitle{#2}\par\nopagebreak\addvspace{6pt}%
  \def\lakin@short{#1}\ifx\lakin@short\@empty\else
    \addcontentsline{lot}{table}{\protect\numberline{\thetbl}{\protect\ignorespaces #1}}\fi
  {\centering\lakitcaptionstyle\mbox{\lakitcaptionlabelstyle{\tablename~\en{\thetbl}\lakitcaptionpunctuation}}\enspace #2\par}}
\makeatother

% ============================================================================
% Unified FIGURE caption — the \tabcap twin for TikZ/plot figures (v3.34.0).
% A figure (tikzpic) sits inline in the prose; put \figcap BELOW it for a
% numbered, \ref-able caption "איור N.M: תיאור" in the same style as \tabcap:
%   \begin{tikzpic} ... \end{tikzpic}
%   \figcap{\label{fig:foo}description}     % \label FIRST, inside the argument
% WHY a macro, not a hand-written gray line: a manual
%   {\centering\small\textcolor{cGray}{איור: ...}} with a \label inside has NO
% \refstepcounter, so the \label binds to the AMBIENT \@currentlabel (the last
% equation) — \ref then prints an equation number AS the figure number
% (measured on a real page: "איור 48" where 48 is equation 48). \figcap steps a
% real \c@fig counter first, so \label is correct, numbering is consistent, and
% refs are stable. Class-aware like \thetbl (chapter- or section-prefixed).
% ============================================================================
\makeatletter
\@ifundefined{c@chapter}{%
  \newcounter{fig}[section]
  \renewcommand{\thefig}{\ifnum\value{section}>0 \thesection.\fi\arabic{fig}}%
}{%
  \newcounter{fig}[chapter]
  \renewcommand{\thefig}{\ifnum\value{chapter}>0 \thechapter.\fi\arabic{fig}}%
}
\makeatother
\makeatletter
\newcommand{\figcap}[2][]{%
  \ifdim\lakinfigraise>0pt
    % v3.74: step WITHOUT hyperref's automatic anchor, then emit the anchor
    % raised above the measured figure. \H@refstepcounter is hyperref's saved
    % pre-hook \refstepcounter (sets the counter AND \@currentlabel — \label
    % in the caption arg binds exactly as before); \hyper@makecurrent supplies
    % \@currentHref, so the dest NAME is identical to the old automatic one —
    % \ref/\autoref/bookmarks are unchanged, only the dest'S POSITION moves.
    \H@refstepcounter{fig}\hyper@makecurrent{fig}%
    % The anchor is emitted INSIDE the caption paragraph below, as a zero-size
    % raisebox — NOT here in vertical mode, where a \raisebox starts an empty
    % paragraph and adds a blank line (measured: every glyph after the caption
    % shifted +18.8pt). Inside an existing line a [0pt][0pt] box is inert.
    \edef\lakin@anch{\@currentHref}%
    % \edef FREEZES the measured height into the pending macro: the deferred
    % \raisebox is typeset inside the caption paragraph AFTER the register is
    % reset to the 0pt sentinel below — a plain \def read 0 at use time
    % (measured: LAKINRAISE=0.0, dest moved 14pt instead of above the figure).
    \edef\lakin@pend{\noexpand\raisebox{\the\dimexpr\lakinfigraise+42pt\relax}[0pt][0pt]{%
      \noexpand\hyper@anchorstart{\lakin@anch}\noexpand\hyper@anchorend}}%
    \global\lakinfigraise 0pt\relax
  \else
    \refstepcounter{fig}\let\lakin@pend\relax
  \fi
  \lakinrtitle{#2}\par\nopagebreak\addvspace{4pt}%
  \def\lakin@short{#1}\ifx\lakin@short\@empty\else
    \addcontentsline{lof}{figure}{\protect\numberline{\thefig}{\protect\ignorespaces #1}}\fi
  {\centering\lakin@pend\lakitcaptionstyle\mbox{\lakitcaptionlabelstyle{\figurename~\en{\thefig}\lakitcaptionpunctuation}}\enspace #2\par}}
\makeatother

% ============================================================================
% LTR islands.  EVERY tikzpicture / quantikz MUST live inside one, or the whole
% figure's coordinates and Latin labels flip. (Hebrew node text still uses \h{}.)
% ============================================================================
% v3.74 (BLD-37): the \figcap hyperlink DESTINATION used to land at the caption
% line — BELOW the figure — so clicking a figure reference scrolled the figure
% off-screen (measured fleet-wide: 97/97 fig dests in one collection sat 19-54pt
% under their figure's ink; a 347pt figure was fully hidden after the jump).
% hypcap-style fix in two parts: (1) at the island's END, \lakin@grabfig grabs
% the just-typeset centered line with \lastbox, measures it, and puts it back
% with \nointerlineskip — the interline glue above the line was already
% contributed and is NOT re-added, so the page geometry is unchanged to the
% glyph (verified); (2) \figcap then steps the counter WITHOUT hyperref's
% automatic anchor (\H@refstepcounter = the pre-hyperref \refstepcounter) and
% emits the anchor itself, raised by the measured height — the dest lands ABOVE
% the figure top. A captionless island leaves numbering untouched (nothing is
% stepped at the island; the grab flag is cleared at every island begin, so a
% stale height from a captionless figure cannot leak into a later caption).
\newsavebox{\lakinfigbox}
\newdimen\lakinfigraise  % real \dimen: tex.setdimen cannot target a \newlength SKIP register (measured lua error)
% BLD-37 table side: a SEPARATE register for \tabcap. The two must not share —
% a captionless island INTENTIONALLY leaves \lakinfigraise set (the grab runs
% at every island END; only the NEXT island's begin clears it), so a shared
% register would hand a \tabcap after a captionless figure that stale FIGURE
% height. \lakintabraise is written only by \lakin@grabtab at \tabcap entry
% and consumed/reset inside \tabcap itself, so it can never go stale.
\newdimen\lakintabraise
\makeatletter
% v3.74b: measure WITHOUT touching the list. The first draft grabbed the box
% with \lastbox and put it back — and \lastbox DISCARDS the node's baseline
% shift, so a vertically-anchored quantikz circuit rose 2.57pt and everything
% after it sank 5.32pt (measured on the A/B glyph diff). LuaTeX lets us READ
% the tail node's dimensions without removing anything: geometry is identical
% by construction. \lakinfigraise=0pt is the "no figure measured" sentinel.
\def\lakin@grabfig{%
  \ifhmode
    \directlua{
      local n = tex.nest[tex.nest.ptr].tail
      local GLUE, PEN = node.id('glue'), node.id('penalty')
      while n and (n.id == GLUE or n.id == PEN) do n = n.prev end
      local HL, VL = node.id('hlist'), node.id('vlist')
      if n and (n.id == HL or n.id == VL) then
        tex.setdimen('global', 'lakinfigraise',
                     (n.height or 0) + (n.depth or 0))
      else
        tex.setdimen('global', 'lakinfigraise', 0)
      end}%
  \else
    \global\lakinfigraise 0pt\relax
  \fi}
% BLD-37 table side: \lakin@grabfig's twin, NOT a shared helper — the Lua body
% differs in exactly one place, the skip set. Between \end{tabular} and \tabcap
% an otherlanguage wrapper (patterns 3/5 of example_tables) leaves a \write
% whatsit and a dir node ON TOP of the interword glue (measured node dump:
% "whatsit/write <- dir <- glue(2.62) <- hlist"), so the table reader must
% also step over whatsit/dir markers — all zero-size — to reach the tabular.
% \lakin@grabfig deliberately does NOT get the wider skip set: its glue/penalty
% body is the one the v3.74 A/B proof froze, and widening it would newly
% measure (and so newly raise) anchors in island shapes that today take the
% plain path — an unproven behavior change on the FIG side. Anything that is
% not a box behind those markers still hits the 0pt sentinel and \tabcap falls
% back to the plain \refstepcounter.
% SECOND deliberate difference: this stores the box's HEIGHT, not height+depth.
% \figcap's anchor is emitted in the caption paragraph BELOW the figure, so it
% must climb the full extent; \tabcap's anchor sits ON the tabular's own line,
% whose baseline is the vcentered tabular's MIDPOINT — height alone reaches
% the top edge, and adding the depth was measured to overshoot by half the
% table (dests 80-135pt above their frame, INSIDE the preceding table's box).
\def\lakin@grabtab{%
  \ifhmode
    \directlua{
      local n = tex.nest[tex.nest.ptr].tail
      local GLUE, PEN = node.id('glue'), node.id('penalty')
      local WHAT, DIR = node.id('whatsit'), node.id('dir')
      while n and (n.id == GLUE or n.id == PEN
                   or n.id == WHAT or n.id == DIR) do n = n.prev end
      local HL, VL = node.id('hlist'), node.id('vlist')
      if n and (n.id == HL or n.id == VL) then
        tex.setdimen('global', 'lakintabraise', n.height or 0)
      else
        tex.setdimen('global', 'lakintabraise', 0)
      end}%
  \else
    \global\lakintabraise 0pt\relax
  \fi}
\makeatother
% ============================================================================
% v3.79 (BLD-37 table side, CORRECTION): put the anchor at the TOP OF THE BOX,
% instead of guessing the distance to it.
% ============================================================================
% v3.77 raised the \tabcap anchor by (height of the box at the tail of the open
% horizontal list) + 42pt. That is a MEASUREMENT taken BEFORE line breaking, and
% it is short by every term it cannot see. Fleet-measured on the 12 built
% collections (299 tbl.* dests, ownership established through each table's own
% CAPTION, never through "which frame encloses the dest"): 3 dests landed INSIDE
% their own tablebox, below the frame top —
%   computer-architecture       tbl.14.1  -46.4pt   three stacked tabulars
%   computer-architecture       tbl.30.5  -13.5pt   tabular 2.85pt over \linewidth
%   advanced-electronic-devices tbl.B.1   -59.2pt   tabular 2.56pt over \linewidth
% Two missing terms, both reproduced in a minimal document (probe P2/P3):
%   (1) VERTICAL LIST ALREADY CONTRIBUTED inside the box — a second and third
%       tabular, or a heading above the table. The tail read sees only the LAST
%       box, so 77.7pt of stacked tabulars is invisible to the raise.
%   (2) WHICH LINE THE ANCHOR LANDS ON. When the tabular alone overfills
%       \linewidth (459.6121pt in a tablebox), the interword glue after
%       \end{tabular} becomes the chosen breakpoint and \Hy@raisedlink is
%       emitted on the NEXT line — one (depth + lineskip) lower. That term does
%       not EXIST at \tabcap time: line breaking has not happened yet. No
%       measurement taken there can ever recover it. (Proof: probe P3, the same
%       table 2.84912pt too wide, fails at -13.5pt; probe P4, the identical
%       table with \tabcolsep 2pt so it fits, is correct at +32.0pt.)
% So the raise is not computed at \tabcap at all. It is CORRECTED at
% \end{tablebox}, from inside mdframed's own body vbox (endinnercode runs after
% \par\unskip, i.e. after line breaking, while tex.nest[tex.nest.ptr] is still
% that vbox's vertical list). There the anchor's true offset from the top of the
% box content is not estimated but WALKED, and the \Hy@raisedlink shift is set
% so the dest lands 42pt above it — which is 42pt - innertopmargin(9pt) -
% linewidth(0.9pt) = 32.0pt above the FRAME top, the exact position the 295
% already-correct dests have today. That is why this is a correction and not a
% re-placement: the good sites do not move.
%
% WHY THE GLYPH DIFF IS 0 BY CONSTRUCTION — the only thing written is the
% `shift` field of the hlist that \Hy@raisedlink already built, and the guard
% below refuses to write it unless that box is 0x0x0. A zero-size box carries no
% dimension into its parent's packing, the packing already happened, and shift
% is not an input to it. Nothing else in the list is touched: no node is added,
% removed, resized or reordered. Measured, not assumed: assets/example_tables.tex
% and the whole computer-architecture collection A/B at 0 differing pymupdf
% rawdict glyph records.
%
% WHY NOT the alternatives (all prototyped and measured before this one):
%   * accumulate the vertical list at \tabcap time — fixes term (1), CANNOT fix
%     term (2); 2 of the 3 real failures survive. A fix that leaves the majority
%     of the defect standing is the false all-clear this repo exists to avoid.
%   * \zref-savepos two-pass — correct on all shapes, but it makes the dest
%     POSITION aux-derived (DEFECTS BLD-1/2), a single-pass compile silently
%     ships an unraised anchor, and the bare y-difference has no page guard.
%   * emit the anchor at \begin{tablebox} and let \tabcap adopt it — needs the
%     tbl counter stepped at box OPEN to keep the dest NAME byte-identical; a
%     captionless tablebox then either shifts every later table number or leaves
%     an orphan dest that the next table duplicates (a HARD build gate). This
%     correction touches neither the counter nor the name: \tabcap's naming path
%     is v3.77's, unchanged, so resolve_names() is byte-identical (verified).
%   * suppressing the line break (a penalty before the glue) — that MOVES the
%     tabular: the break currently discards the interword glue, so forbidding it
%     re-centers the two overfull tables. Reflow, rejected.
%
% SCOPE. The correction is armed only by \tabcap ITSELF (\lakintabcapseen),
% cleared at every box open, so a captionless tablebox is a no-op and a \tabcap
% hosted OUTSIDE a tablebox (the fleet has exactly one, computer-architecture
% tbl.14.2 in an exbox) keeps v3.77's measured raise. Both fallbacks are
% therefore never worse than today; and because hyperref's AUTOMATIC anchor uses
% the same \Hy@raisedlink shape, the plain-\refstepcounter path (vmode entry, or
% prose between \end{tabular} and \tabcap — probe P6, -51.5pt) is corrected too.
% KNOWN LIMIT — a tablebox that SPLITS across pages. The offset is measured from
% the top of the UNSPLIT body, so a dest in segment 2 gets raised past its own
% segment. MEASURED on a forced split (probe: 40 filler lines + a table): the
% dest went from y=585.0 on page 2 (inside segment 2, at the caption — the old
% BLD-37 symptom) to y=-113.4, i.e. ABOVE the page top. Viewers clamp that to the
% page top, which happens to show segment 2 whole, but an out-of-page coordinate
% is not something to ship on purpose. Hence the \textheight guard in
% lakintab.fix: a body TALLER THAN THE TEXT BLOCK cannot fit on any page, so it
% must split — the correction is skipped and v3.77's behaviour stands unchanged
% (verified: the probe returns to y=585.0). The guard is deliberately EXACT, not
% a page-fitting heuristic: it is sound (a taller-than-page body always splits)
% but PARTIAL (a shorter body that starts low on a page could still split). That
% residue is unreachable in practice — v3.37's everyline + the caption's
% \nopagebreak glue move such boxes WHOLE instead (probe N7: a 22-row table
% pushed down by 430pt moved whole, 0 segments), and the fleet has 0 split
% tableboxes, re-counted after this change. Reimplementing mdframed's own
% page-fitting arithmetic to close the residue would be a heuristic with its own
% failure mode, which is the trade this repo has repeatedly refused.
\newif\iflakintabcapseen
\newif\iflakintabframed
% The Lua lives in ONE \directlua at load time, and is written with NO
% backslash, %, # or ~ : \directlua expands its argument like \edef, so an
% active ~ would expand to \nobreakspace and a % would eat the rest of the line
% (the second is how the first draft of this block died). "not (a == b)" stands
% in for Lua's inequality operator for the same reason.
\directlua{
  lakintab = { }
  lakintab.HL = node.id("hlist")
  lakintab.VL = node.id("vlist")
  lakintab.RU = node.id("rule")
  lakintab.GL = node.id("glue")
  lakintab.KE = node.id("kern")
  lakintab.WH = node.id("whatsit")
  lakintab.TM = node.id("temp")
  lakintab.WD = node.subtype("pdf_dest")
  function lakintab.d(v)
    if v == nil then return 0 end
    if v < -1000000000 then return 0 end
    return v
  end
  function lakintab.scan(list, y, horiz, parent, target)
    local x = list
    while x do
      local i = x.id
      if i == lakintab.HL or i == lakintab.VL then
        if horiz then
          local b = y + x.shift
          if i == lakintab.HL then
            lakintab.scan(x.head, b, true, x, target)
          else
            lakintab.scan(x.head, b - lakintab.d(x.height), false, x, target)
          end
        else
          if i == lakintab.HL then
            lakintab.scan(x.head, y + lakintab.d(x.height), true, x, target)
          else
            lakintab.scan(x.head, y, false, x, target)
          end
          y = y + lakintab.d(x.height) + lakintab.d(x.depth)
        end
      elseif i == lakintab.WH then
        if x.subtype == lakintab.WD then
          local id = x.dest_id
          if type(id) == "string" and string.sub(id, 1, 4) == "tbl." then
            if horiz and parent and parent.width == 0
               and parent.height == 0 and parent.depth == 0 then
              parent.shift = parent.shift + (target - y)
            end
          end
        end
      elseif not horiz then
        if i == lakintab.RU then
          y = y + lakintab.d(x.height) + lakintab.d(x.depth)
        elseif i == lakintab.GL or i == lakintab.KE then
          y = y + lakintab.d(x.width)
        end
      end
      x = x.next
    end
  end
  function lakintab.fix(target, maxh)
    local nst = tex.nest[tex.nest.ptr]
    if nst == nil then return end
    local h = nst.head
    if h == nil then return end
    if h.id == lakintab.TM then h = h.next end
    local tot = 0
    local x = h
    while x do
      local i = x.id
      if i == lakintab.HL or i == lakintab.VL or i == lakintab.RU then
        tot = tot + lakintab.d(x.height) + lakintab.d(x.depth)
      elseif i == lakintab.GL or i == lakintab.KE then
        tot = tot + lakintab.d(x.width)
      end
      x = x.next
    end
    if tot > maxh then return end
    lakintab.scan(h, 0, false, nil, target)
  end
}
% @-FREE NAMES ON PURPOSE: these two are named in the `tablebox` key list at
% \newmdenv above, which is tokenised OUTSIDE any \makeatletter region — a
% \lakin@... there would tokenise as \lakin followed by the letters @tabopen.
\newcommand{\lakintabboxopen}{\lakintabframedtrue\global\lakintabcapseenfalse}
\newcommand{\lakintabboxclose}{%
  \iflakintabcapseen
    \directlua{lakintab.fix(-\number\dimexpr\lakintabcushion\relax,
                            \number\dimexpr\textheight\relax)}%
    \global\lakintabcapseenfalse
  \fi\lakintabframedfalse}
% The cushion above the box's CONTENT top. 42pt puts the dest 32.0pt above the
% frame top (42pt - innertopmargin 9pt - linewidth 0.9pt), i.e. exactly where
% the 295 correct dests already are — the correction moves only the broken ones.
\newcommand{\lakintabcushion}{42pt}
\newenvironment{qcirc}{\par\medskip\global\lakinfigraise 0pt\relax\begin{otherlanguage}{english}\begin{center}}
  {\csname lakin@grabfig\endcsname\end{center}\end{otherlanguage}\par\medskip}
\newenvironment{tikzpic}{\par\medskip\global\lakinfigraise 0pt\relax\begin{otherlanguage}{english}\begin{center}}
  {\csname lakin@grabfig\endcsname\end{center}\end{otherlanguage}\par\medskip}

% ============================================================================
% Code listings.  ★★ v3.4 THE BOX-OVERFLOW FIX:
%   resetmargins=true   → the listing IGNORES the enclosing environment's
%                         indentation and measures from the PAGE. v3.4 stated
%                         this backwards ("measured from the ENCLOSING box"),
%                         and v3.37 then blanked in-box line numbers because
%                         the page-anchored listing put its hanging digits on
%                         top of the box rule. Measured v4.5.0: with the key
%                         FALSE the listing follows the box and the code moves
%                         from x=75.26 to x=97.74, clearing the rule at 72.07.
%                         The key now lives in one \lstset below instead of in
%                         each style, so a box can override it — codebox does.
%   framexleftmargin=0  → the frame (left bar) does not extend left of the code.
%   xleftmargin keeps a small indent so the line numbers have room INSIDE.
% asmblock wraps the listing in an LTR island (code is LTR).
% ============================================================================
\lstdefinestyle{asm}{
  basicstyle=\ttfamily\small,
  backgroundcolor=\color{cBlueBg!55},
  frame=leftline,framerule=2.5pt,rulecolor=\color{cBlue},
  numbers=left,numberstyle=\tiny\color{cGray},numbersep=8pt,
  xleftmargin=2.2em,framexleftmargin=0pt,framexrightmargin=0pt,
  breaklines=true,columns=fullflexible,keepspaces=true,
  showstringspaces=false,aboveskip=8pt,belowskip=8pt,
  % v3.36: without a comment DELIMITER the commentstyle below never fires —
  % listings has no `language` here, so it does not know '#' opens a comment,
  % and every asm listing rendered flat black (measured: a whole-page pixel
  % sweep found ZERO cGreen glyphs across 135 shipped listings). morecomment
  % gives it the delimiter; color-only change, glyph widths untouched.
  morecomment={[l]{\#}},
  commentstyle=\color{cGreen},
}
\lstdefinestyle{cstyle}{
  language=C,basicstyle=\ttfamily\small,
  backgroundcolor=\color{cGreenBg!55},
  frame=leftline,framerule=2.5pt,rulecolor=\color{cGreen},
  numbers=left,numberstyle=\tiny\color{cGray},numbersep=8pt,
  xleftmargin=2.2em,framexleftmargin=0pt,framexrightmargin=0pt,
  breaklines=true,columns=fullflexible,keepspaces=true,
  showstringspaces=false,keywordstyle=\color{cBlue}\bfseries,
  commentstyle=\color{cGray}\itshape,aboveskip=8pt,belowskip=8pt,
}
% v3.36: pystyle/htmlstyle — colored twins of cstyle for the Python/HTML
% listings that used to repeat a ~200-char options blob with NO color
% (keywords rendered bold-black, comments italic-black — measured; C alone
% was colored). Geometry is copied VERBATIM from the shipped blob (frame,
% margins, skips) so swapping a call site to style=pystyle moves NOTHING —
% only keywordstyle/commentstyle/stringstyle are added (bold stays bold,
% italic stays italic, mono metrics identical => zero reflow, proven by
% per-page word-bbox diff). breaklines/breakatwhitespace stay per-site.
% v3.78: numbers=left joins both styles, with the exact numberstyle/numbersep
% keys asm/cstyle carry — fleet-wide consistency (author decision: all four
% styles numbered; "asm/cstyle numbered, py/html not" was an accident of the
% v3.36 geometry freeze, not a design). The numbers hang LEFT of the code box
% (listings sets them with \llap, so no overfull hbox can arise and the code
% x is untouched — measured: code frozen, digits at code_x-11.8pt); inside a
% colored box they are blanked like everyone else's (v3.37 hook), and the
% clean variant restores them in its frameless boxes (v3.78, above).
\lstdefinestyle{pystyle}{
  language=Python,basicstyle=\ttfamily\small,
  columns=fullflexible,keepspaces=true,
  frame=single,rulecolor=\color{cTblLine},
  numbers=left,numberstyle=\tiny\color{cGray},numbersep=8pt,
  xleftmargin=0.5em,framexleftmargin=0pt,framexrightmargin=0pt,
  showstringspaces=false,aboveskip=8pt,belowskip=8pt,
  keywordstyle=\color{cBlue}\bfseries,
  commentstyle=\color{cGray}\itshape,
  stringstyle=\color{cGreen},
}
% v4.5.0: hoisted out of the four styles so an enclosing box can override it.
% Every listing still resolves to true unless something later says otherwise,
% so shipped output is unchanged (verified byte-for-byte on all eight PDFs).
\lstset{resetmargins=true}
\lstdefinestyle{htmlstyle}{
  language=HTML,basicstyle=\ttfamily\small,
  columns=fullflexible,keepspaces=true,
  frame=single,rulecolor=\color{cTblLine},
  numbers=left,numberstyle=\tiny\color{cGray},numbersep=8pt,
  xleftmargin=0.5em,framexleftmargin=0pt,framexrightmargin=0pt,
  showstringspaces=false,aboveskip=8pt,belowskip=8pt,
  keywordstyle=\color{cBlue}\bfseries,
  commentstyle=\color{cGray}\itshape,
  stringstyle=\color{cGreen},
}
\newenvironment{asmblock}{\par\begin{otherlanguage}{english}}{\end{otherlanguage}\par}

% <hllpdf:listing-smart-quotes>
% v4.21.0: four Unicode punctuation characters must enter listings' token
% buffer. Unregistered characters bypass it: a closing quote was printed
% BEFORE the pending word, and an opening quote BEFORE a pending brace.
% Extend the standard character table, then activate only these codepoints
% in the listing's local character-table scope. Do this BEFORE the package's
% delimiter/alsoother/literate hooks so explicit user mappings still win.
% This is not universal Unicode support, nor a conversion to ASCII quotes.
% See listings developer guide 15.3.1 and references/listing-smart-quotes.md.
\makeatletter
\lst@CCPut \lst@ProcessOther {"2018}{"2019}{"201C}{"201D}\z@
\lst@AddToHookAtTop{SelectCharTable}{%
  \lst@ifactivechars
    \catcode"2018=\active \catcode"2019=\active
    \catcode"201C=\active \catcode"201D=\active
  \fi}
\makeatother
% </hllpdf:listing-smart-quotes>

% <hllpdf:listing-technical-symbols>
% v4.30: bounded extension for common technical punctuation in LTR listings.
% Without character-table registration these glyphs bypass the token buffer
% and can move before the preceding word even inside quoted strings.
% Keep the original four-quote block separate for its standing mutation test.
% This is source-faithful printing, not syntax validation or general Unicode.
\makeatletter
\lst@CCPut \lst@ProcessOther
  {"2013}{"2014}{"2026}{"2190}{"2192}{"2194}{"21D2}
  {"2212}{"221E}{"2260}{"2264}{"2265}\z@
\lst@AddToHookAtTop{SelectCharTable}{%
  \lst@ifactivechars
    \catcode"2013=\active \catcode"2014=\active \catcode"2026=\active
    \catcode"2190=\active \catcode"2192=\active \catcode"2194=\active
    \catcode"21D2=\active \catcode"2212=\active \catcode"221E=\active
    \catcode"2260=\active \catcode"2264=\active \catcode"2265=\active
  \fi}
\makeatother
% </hllpdf:listing-technical-symbols>


\usepackage[hidelinks,unicode]{hyperref}   % LAST, after colors. unicode = Hebrew bookmarks.
% v3.62: THE ISOLATION MACROS CORRUPTED EVERY PDF BOOKMARK THEY APPEARED IN.
% \en is defined (line ~155) as
%     \foreignlanguage{english}{\language=\csname l@nohyphenation\endcsname #1}
% and in a heading that body is written verbatim to the .toc and handed to
% hyperref's \pdfstringdef. hyperref strips the two CONTROL SEQUENCES (\language,
% \l@nohyphenation) but the plain `=` CHARACTER token between them is not a control
% sequence, so it survives into the PDF string: \chapter{שער \en{CNOT} והצמדה}
% produced the bookmark "שער =CNOT והצמדה". The printed TOC and the heading on the
% page are CORRECT — only the navigation pane and the viewer window title are wrong,
% which is why it went unnoticed. Measured on the shipped fleet: 70 corrupted
% bookmarks across 5 PDFs, from 177 unprotected \en-in-heading sites. Same corruption
% reached /Title and /Subject via the \hypersetup{pdftitle=...} the docs recommend.
% This hook tells hyperref to use the plain argument for all three isolation macros
% when building a PDF string. It does NOT touch the .toc (the printed TOC is
% byte-unchanged) and it does NOT mask the math signal — a heading with $...$ still
% logs "removing math shift", so BIDI-15 detection is intact.
% BONUS, measured: it also makes \code{} in an UNSTARRED heading COMPILE. BLD-10
% recorded that as needing \texorpdfstring or a starred heading; the real root cause
% is hyperref's \pdfstringdef scan of \code's \ifmmode...\else...\fi body — which is
% exactly why the identically-bodied \texttt survived (hyperref pre-neutralizes
% \texttt in PDF strings) and why \DeclareRobustCommand could never help.
% v3.64: \def\\{ } added. A manual line break in a heading WELDS the two words
% together in the bookmark: \chapter{שורה ראשונה\\שורה שנייה} produced the outline
% entry 'שורה ראשונהשורה שנייה' — "ראשונהשורה" is not a word. hyperref drops the
% \\ and puts nothing in its place, so the break that looks like a space on the
% page is a zero-width join in the navigation pane. The spaced form
% 'מילה אחת \\ מילה שנייה' escapes with a harmless DOUBLE space ('אחת  מילה'),
% which is why this never got noticed. Both are measured on a real outline.
% 0 live instances in the corpus — but SKILL.md and content-style.md both RECOMMEND
% a manual heading break, so this is a primed trap rather than a dormant one, and a
% one-token preamble fix retires it for good instead of asking authors to remember.
% Mapping \\ to a space fixes the WELD, which is the actual defect (a nonexistent
% word in the navigation pane). It does not normalise whitespace: measured after the
% fix, the welded form becomes 'שורה ראשונה שורה שנייה' (correct) and the already-
% spaced form goes from two spaces to three. That is cosmetic — viewers render it as
% a wider gap, and no word is corrupted. A per-heading \texorpdfstring is still the
% way to get an exactly-tuned bookmark string.
\pdfstringdefDisableCommands{\def\en#1{#1}\def\h#1{#1}\def\doclang#1{#1}\def\code#1{#1}\def\codeline#1{#1}\def\\{ }}
% v3.70: RESTORE \url LINE-BREAKING. url.sty freezes its break penalties at LOAD
% time (\mathchardef\UrlBreakPenalty=\binoppenalty), and v3.45's global
% \binoppenalty=10000 (inline-math atomicity) therefore priced EVERY breakpoint
% inside \Url at 10000 = never break. Measured: a long \url in Hebrew prose is
% one unbroken line, ~100pt overfull, invading the margin band the ink gates
% deliberately tolerate. With the two lines below: 0 overfull, the URL breaks
% after '/', the RTL split is CORRECT (fragment order forward), and each line
% gets its own link annotation exactly covering its fragment (measured).
% Confined to \Url's group — inline-math atomicity (TYP-7) is untouched.
\mathchardef\UrlBreakPenalty=200
\mathchardef\UrlBigBreakPenalty=100
\ifx\lakitmainlanguage\hllpdfenglish
  \hypersetup{pdflang=en}
\else
  \hypersetup{pdflang=he}
\fi  % /Lang is the document base locale (not per-span accessibility tagging).
% Set a per-document \hypersetup{pdftitle={...}} too — see SKILL.md.

% v3.52 (round-2 audit): two RTL bibliography/cross-ref defects, both invisible to
% every gate and both affecting the FIRST author to use the feature.
% (a) \@biblabel's default [#1] brackets are bidi-neutral and mirror to ]1[ under
% the RTL bibliography base — every \thebibliography reference list shows ]1[.
% LTR-isolating the label restores [1]. (The render gate's mirror rule fingerprints
% round parens )1(, so square brackets slipped through.)
\makeatletter
\renewcommand*{\@biblabel}[1]{\foreignlanguage{english}{[#1]}}
% v3.54: the SAME neutral-bracket trap on the INLINE side — a multi-key \cite{a,b}
% under the RTL base lays its European-number list right-to-left, so [1, 2] renders
% "[2,1]" and [1, 2, 3] renders "[3,2,1]" (verified by glyph geometry — the reader
% sees the sources in REVERSE order; a note form \cite[p.5]{a} also reorders to
% "[p. 5,1]"). LTR-isolate the whole bracket group (number list AND optional note),
% exactly like \@biblabel above: [1, 2] / [1, 2, 3] / [1, p. 5] all render correct.
% Placed AFTER hyperref so it wins; hyperref hyperlinks each number inside #1 (via
% \@citex), which this wrapper leaves intact. Mirrors the LaTeX kernel \@cite shape.
% v4.9.1: directional isolation must NOT switch the language here. Babel's
% safe-active citation wrapper passes its key list through \bbl@tempa;
% \foreignlanguage reuses that scratch macro before the deferred citation
% loop expands it. The old line therefore cited the literal key "english".
% If such a key existed, even the undefined-reference gate stayed green.
% \babelsublr isolates the bracket group without selecting a locale or
% changing that key list. Preserve native links, key order and optional notes.
\def\@cite#1#2{\babelsublr{[{#1\if@tempswa , #2\fi}]}}
\makeatother

% v4.11.0: a reference inside align's \text inherited HEBdisptag=true,
% so its enclosing hbox mirrored the brackets despite a clean compile.
% Isolate the complete printed reference without selecting a Babel language.
% Also neutralize the display-only normalfont direction reset locally. A
% direction-only patch inverted 12 to 21 inside equation in the measured probe.
% Plain maketag is local to the reference: mathtools' display-side wrapper
% otherwise globally consumes manual_tag while formatting a nested reference.
% Preserve each original command (including the selective auxiliary write),
% current tag forms, native links, and showonlyrefs on/off transitions.
\makeatletter
\NewCommandCopy{\HEB@nativeeqref}{\eqref}
\newcommand{\HEB@eqrefisolate}[1]{%
  \babelsublr{\HEBdisptagfalse\let\bbl@eqnodir\relax\let\maketag@@@\ORIGmaketag@@@#1}}
\protected\def\eqref#1{\HEB@eqrefisolate{\HEB@nativeeqref{#1}}}
\ExplSyntaxOn
\NewCommandCopy{\HEB@nativeMTeqref}{\MT_extended_eqref:n}
\cs_set_protected:Npn \MT_extended_eqref:n #1
  {\HEB@eqrefisolate{\HEB@nativeMTeqref{#1}}}
\ExplSyntaxOff
\makeatother

% v4.11.0: custom caption counters are not part of hyperref's English
% locale, and Hebrew has no upstream hyperref locale in this toolchain.
% Restore the appropriate names at every Babel language switch, including
% return from a nested inline island. Do not localize the destination keys.
\makeatletter
\newcommand{\HEB@refnames}{%
\def\chapterautorefname{פרק}%
\def\sectionautorefname{סעיף}%
\def\subsectionautorefname{תת־סעיף}%
\def\subsubsectionautorefname{תת־תת־סעיף}%
\def\paragraphautorefname{פסקה}%
% v4.31: hyperref uses a distinct counter name at the deepest standard level.
\def\subparagraphautorefname{תת־פסקה}%
\def\equationautorefname{משוואה}%
\def\figureautorefname{איור}%
\def\tableautorefname{טבלה}%
\def\partautorefname{חלק}%
\def\appendixautorefname{נספח}%
\def\footnoteautorefname{הערת שוליים}%
\def\itemautorefname{סעיף}%
\def\Itemautorefname{סעיף}%
\def\figautorefname{איור}%
\def\tblautorefname{טבלה}%
\def\pageautorefname{עמוד}% 
}
\addto\extrashebrew{\HEB@refnames}
\newcommand{\ENG@refnames}{%
  \def\subparagraphautorefname{subparagraph}%
  \def\chapterautorefname{chapter}%
  \def\sectionautorefname{section}%
  \def\subsectionautorefname{subsection}%
  \def\subsubsectionautorefname{subsubsection}%
  \def\paragraphautorefname{paragraph}%
  \def\equationautorefname{Equation}%
  \def\figureautorefname{Figure}%
  \def\tableautorefname{Table}%
  \def\partautorefname{Part}%
  \def\appendixautorefname{Appendix}%
  \def\footnoteautorefname{footnote}%
  \def\itemautorefname{item}%
  \def\figautorefname{Figure}%
  \def\tblautorefname{Table}%
  \def\Itemautorefname{item}%
  \def\pageautorefname{page}%
}
\addto\extrasenglish{\ENG@refnames}
\ifx\lakitmainlanguage\hllpdfenglish\ENG@refnames\else\HEB@refnames\fi
\makeatother

% ============================================================================
% Optional running head (v3.39): \lakithead{<book title>} — OFF by default.
% A document opts in with ONE line after \input{preamble.tex}:
%   \lakithead{מבנה המחשב}
% Every content page then carries the series' orientation line: chapter name
% at the reading edge, book title opposite, thin rule. (fancyhdr's [L]/[R]
% are PHYSICAL; under the RTL page direction [L] lands at the RIGHT/reading
% edge — measured, consistent across the corpus.) Packaged like the topic
% rail: always loaded, inert without the call. This centralizes the 8-line
% block that 8 lakits used to copy by hand — the same duplication vector
% that once let a lakit preamble go ~100 lines stale. Class-aware: report
% fills the mark from \chapter, article from \section. NB: a titlepage env
% self-empties; front matter that sets \pagestyle{empty} must restore with
% \pagestyle{fancy} afterwards (NOT plain — that discards the head).
% ============================================================================
\usepackage{fancyhdr}
% v3.61: \chapter* left a STALE running head — pages of an UNNUMBERED chapter
% carried the PREVIOUS chapter's title. report.cls routes \chapter* through
% \@schapter, which never calls \chaptermark, so \leftmark keeps whatever the last
% numbered chapter left. Measured: with \chapter{פרק ממוספר} then
% \chapter*{מבוא לא ממוספר}, page 4 (inside the starred chapter) printed
% "פרק ממוספר". Invisible on the chapter's FIRST page (\thispagestyle{plain}), so it
% only shows from the second page on — the pages an author is least likely to check.
% An unnumbered front-matter chapter (\chapter*{מבוא}, \chapter*{הקדמה}) is everyday
% Hebrew-academic structure, and the corpus already carries HAND-WRITTEN \markboth
% lines right after \chapter* — the exact per-document duplication \lakithead exists
% to retire. Patched INSIDE \lakithead so the macro stays inert when unused; a
% document's own later \markboth still wins, so those hand-patches do not regress.
\makeatletter
% <hllpdf:starred-heading-scope>
\newif\iflakin@starmarkpatched
\newcommand{\lakistarmark}{%
  \iflakin@starmarkpatched\else
  % Dispatch on the actual class counter, not hyperref's shared internals.
  \ifdefined\c@chapter
    \let\HEBorigschapter\@schapter
    % Record only after the native heading has reached its page. A mark
    % before the native routine can advertise the next heading too early.
    \def\@schapter##1{\HEBorigschapter{##1}\chaptermark{##1}}%
  \else
    % @ssect receives layout arguments, NOT the section name or level.
    % Capture the name in @startsection without introducing a TeX group:
    % grouping sectioning would break the native after-heading state.
    \def\HEB@sectionkind{}%
    \def\HEB@mainsection{section}%
    \let\HEBorigstartsection\@startsection
    \def\@startsection##1##2##3##4##5##6{%
      \def\HEB@sectionkind{##1}%
      \HEBorigstartsection{##1}{##2}{##3}{##4}{##5}{##6}}%
    \let\HEBorigssect\@ssect
    \def\@ssect##1##2##3##4##5{%
      \HEBorigssect{##1}{##2}{##3}{##4}{##5}%
      % Lower starred headings must not replace the article's section mark.
      \ifx\HEB@sectionkind\HEB@mainsection
        \sectionmark{##5}%
      \fi
      \ignorespaces}%
  \fi
  \global\lakin@starmarkpatchedtrue
  \fi}
% </hllpdf:starred-heading-scope>
% <hllpdf:running-language>
% Choose the wrapper when the mark is written, not when a later page ships.
% English punctuation must remain LTR even after returning to Hebrew prose.
\newcommand{\HEB@runningmark}[1]{%
  \iflanguage{english}{\markboth{\en{#1}}{}}{\markboth{\h{#1}}{}}}
% </hllpdf:running-language>
\makeatother
% v4.2.1 (NAV): every collection already names itself in \lakithead — feed the
% PDF /Title from the same declaration, so the navigation gate's
% PDF_REQUIRE_TITLE has a universal source and no fleet member ships
% title-less by missing a second declaration site. Measured before this: 14
% of 16 collections shipped /Title-empty while every one carried a perfectly
% good \lakithead string (plain Hebrew in all 16 — no math, no formatting, so
% the PDF-string path is clean). An explicit title always wins: set it with
% \lakitpdftitle{...} anywhere before \lakithead, or with a plain
% \hypersetup{pdftitle=...} — the fallback detects that via hyperref's
% \@pdftitle, probe-verified \@empty-equal until the first explicit set.
\makeatletter
\newif\iflakit@pdftitleset
\newcommand{\lakitpdftitle}[1]{%
  \hypersetup{pdftitle={#1}}\global\lakit@pdftitlesettrue}
\newcommand{\lakithead}[1]{%
  \iflakit@pdftitleset\else
    \ifx\@pdftitle\@empty
      \hypersetup{pdftitle={#1}}\global\lakit@pdftitlesettrue
    \fi
  \fi
  \pagestyle{fancy}\fancyhf{}%
  \lakistarmark

  \fancyhead[L]{\small\textcolor{cGray}{\leftmark}}%
  % <hllpdf:fixed-head-language>
  % The fixed document title has its OWN language snapshot, independent of
  % the current section/chapter mark. Choose the branch NOW, not at shipout:
  % the page may be Hebrew by the time an English title is finally printed.
  \iflanguage{english}{%
    \fancyhead[R]{\small\textcolor{cGray}{\en{#1}}}%
  }{%
    \fancyhead[R]{\small\textcolor{cGray}{\h{#1}}}%
  }%
  % </hllpdf:fixed-head-language>
  \fancyfoot[C]{\thepage}%
  \renewcommand{\headrulewidth}{0.4pt}%
  \ifdefined\chaptermark
    \renewcommand{\chaptermark}[1]{\HEB@runningmark{##1}}%
  \else
    \renewcommand{\sectionmark}[1]{\HEB@runningmark{##1}}%
  \fi}
\makeatother

% v3.39.1: path-independent byte-reproducibility. LuaTeX derives the trailer
% /ID from the OUTPUT PATH, so two clones of the same source at different
% directories produced different sha256 (measured: the ONLY delta was the
% 62-byte /ID line). Bit 512 suppresses exactly that line — measured: Creator/
% Producer/CreationDate/ModDate and all content untouched. /ID is optional
% per ISO 32000 (required only for ENCRYPTED PDFs and PDF/A — neither is a
% target here). If a future workflow needs an ID (PDF/A, Acrobat revision
% matching), delete this line and it returns.
\pdfvariable suppressoptionalinfo \numexpr 512\relax

% v3.75: BOLD MATH IN BOLD HEADINGS. Heading text is \bfseries but math inside
% it stayed regular-weight CM — a hairline $V_T$ inside a heavy 24.8pt chapter
% title (measured: CMMI12 fl=6 next to FrankRuhlLibre-Bold fl=20; 155 live
% heading-with-math sites fleet-wide). \boldmath rides along wherever the class
% sets \bfseries for a heading; A/B measured: $V_T$ -> CMMIB10 fl=22, 0 missing
% characters (the bold CM faces ship with this TeX Live). Guarded per-macro so
% the article class (no \@makechapterhead) and future class changes stay safe.
% v4.0.1 (adoption fix): the SEARCH patterns above carry no parameter token.
% v4.0.0 made them specific by including the class's own #1/#2, and measured on
% TeX Live 2023 that never matches: a \patchcmd search pattern containing a
% parameter token cannot match while the call sits inside \@ifundefined{...}{}{...}
% (probed both #1 and ##1 -> FAIL; the same pattern at top level succeeds). Every
% one of these patches therefore failed, and because v4.0.0 also replaced the
% empty failure branch with \PackageError, the failure became FATAL: all four
% bundled example assets stopped building ("Cannot patch chapter-title math
% weight -> Fatal error, no output PDF"). The \PackageError branches are KEPT —
% a silent no-op here is what let v3.75's fix rot unnoticed — only the parameter
% tokens are dropped, so the patches match and a real class change still errors.
\makeatletter
\@ifundefined{@makechapterhead}{}{%
  \patchcmd{\@makechapterhead}{\huge\bfseries \@chapapp}
    {\huge\bfseries\boldmath \@chapapp}{}
    {\PackageError{hebrew-lualatex-pdf}{Cannot patch chapter-number math weight}{Unsupported document class}}%
  \patchcmd{\@makechapterhead}{\Huge \bfseries}
    {\Huge \bfseries\boldmath}{}
    {\PackageError{hebrew-lualatex-pdf}{Cannot patch chapter-title math weight}{Unsupported document class}}%
  \patchcmd{\@makeschapterhead}{\Huge \bfseries}
    {\Huge \bfseries\boldmath}{}
    {\PackageError{hebrew-lualatex-pdf}{Cannot patch starred chapter math weight}{Unsupported document class}}}
% v4.0.1: \part sits under a THREE-DEEP wrapper chain, so the class-based anchor
% is not in \@part at all. Measured with \show: hyperref rewrites \@part to
% `...\phantomsection\fi \H@old@part`, and nameref rewrites THAT to
% `\NR@gettitle{#1}\NR@part[{#1}]{#2}` — the class body ends up in \NR@part.
% v4.0.0 patched \@part and made failure fatal, which is why all seven bundled
% examples stopped building. \part is NOT a dead heading level here: 188 live uses
% across the fleet, so the bold-math fix is worth getting right rather than
% dropping. We walk the chain innermost-first and stop at the first hit. Failure is
% a WARNING, not an error: the chain is a property of hyperref/nameref versions, so
% a future package release must degrade to regular-weight part math — a cosmetic
% loss — instead of refusing to build. Chapter and section keep \PackageError,
% because those targets are in the CLASS and were verified to patch here.
\newif\iflakin@partnum \newif\iflakin@parttitle \newif\iflakin@sparttitle
\def\lakin@partnum#1#2{\iflakin@partnum\else
  \patchcmd{#1}{#2\bfseries \partname}{#2\bfseries\boldmath \partname}
    {\lakin@partnumtrue}{}\fi}
\def\lakin@parttitle#1#2{\iflakin@parttitle\else
  \patchcmd{#1}{#2 \bfseries}{#2 \bfseries\boldmath}{\lakin@parttitletrue}{}\fi}
\def\lakin@sparttitle#1#2{\iflakin@sparttitle\else
  \patchcmd{#1}{#2 \bfseries}{#2 \bfseries\boldmath}{\lakin@sparttitletrue}{}\fi}
\@ifundefined{c@chapter}{\def\lakin@pnum{\Large}\def\lakin@ptit{\huge}}%
                        {\def\lakin@pnum{\huge}\def\lakin@ptit{\Huge}}
\@ifundefined{NR@part}{}{\expandafter\lakin@partnum\expandafter\NR@part\expandafter{\lakin@pnum}%
                         \expandafter\lakin@parttitle\expandafter\NR@part\expandafter{\lakin@ptit}}
\@ifundefined{H@old@part}{}{\expandafter\lakin@partnum\expandafter\H@old@part\expandafter{\lakin@pnum}%
                            \expandafter\lakin@parttitle\expandafter\H@old@part\expandafter{\lakin@ptit}}
\@ifundefined{@part}{}{\expandafter\lakin@partnum\expandafter\@part\expandafter{\lakin@pnum}%
                       \expandafter\lakin@parttitle\expandafter\@part\expandafter{\lakin@ptit}}
\@ifundefined{NR@spart}{}{\expandafter\lakin@sparttitle\expandafter\NR@spart\expandafter{\lakin@ptit}}
\@ifundefined{H@old@spart}{}{\expandafter\lakin@sparttitle\expandafter\H@old@spart\expandafter{\lakin@ptit}}
\@ifundefined{@spart}{}{\expandafter\lakin@sparttitle\expandafter\@spart\expandafter{\lakin@ptit}}
\iflakin@partnum\else\PackageWarning{hebrew-lualatex-pdf}{Part-number math stays
  regular weight: no known \noexpand\part wrapper matched (hyperref/nameref
  version change). Cosmetic only}\fi
\iflakin@parttitle\else\PackageWarning{hebrew-lualatex-pdf}{Part-title math stays
  regular weight: no known \noexpand\part wrapper matched. Cosmetic only}\fi
% v4.12.0: THE PART NUMBER IN THE TOC WAS NOT DIRECTIONALLY SEPARATED FROM THE
% TITLE. report.cls/article.cls compose the part entry as ONE text stream --
% \addcontentsline{toc}{part}{\thepart\hspace{1em}#1} -- while EVERY other level
% goes through \numberline, which puts the number in an \hbox. A box is a bidi
% OBJECT and cannot join a neighbouring run; a bare "IV" can. So in the RTL TOC
% paragraph a Latin number and a Latin-INITIAL title merge into a SINGLE LTR run
% and are laid out left-to-right as a unit: the number lands on the wrong side of
% the row, and a mixed title is scrambled outright. Measured on \part{MCP ...}:
% the printed order was "<tail> V MCP" -- title fragment, number, title fragment.
% The chapter/section rows one line below, with the SAME Latin-initial titles,
% were correct on the same page. That contrast is the whole diagnosis.
% \en{} IS NOT A WORKAROUND, measured both ways on one fixture: it fixes a
% PURE-Latin title (\part{\en{Hooks}} renders correctly) and does NOT fix a MIXED
% one (\part{\en{MCP} <hebrew>} stays scrambled). babel's language switch changes
% font and hyphenation, not the bidi run structure -- so the defect is not
% reachable from the source side at all, which is why it belongs here.
% Boxing the number is the entire fix and is a STRICT NO-OP where the bug never
% bit: a Hebrew-titled document compiles BYTE-IDENTICAL before and after (md5,
% 3 clean runs), because there the neutral glue already sat between an L run and
% an R run. The PDF outline is byte-identical too -- hyperref's \pdfstringdef
% neutralises \mbox, so bookmark strings keep their part numbers.
% \numberline{\thepart} would also fix the direction, but it re-geometries every
% part row in the fleet and DROPS the number from the part bookmark (hyperref
% strips \numberline). \mbox changes nothing that was already right.
% The search pattern deliberately carries no # token: etoolbox cannot match a
% parameter token in a patch pattern (measured -- all three chain members
% reported "not patchable" while the pattern included #1). "\thepart \hspace
% {1em}" is brace-balanced and occurs exactly once in the wrapper body.
\newif\iflakin@parttocnum
\def\lakin@parttocnum#1{\iflakin@parttocnum\else
  \patchcmd{#1}{\thepart \hspace {1em}}{\mbox {\thepart }\hspace {1em}}
    {\lakin@parttocnumtrue}{}\fi}
\@ifundefined{NR@part}{}{\lakin@parttocnum\NR@part}
\@ifundefined{H@old@part}{}{\lakin@parttocnum\H@old@part}
\@ifundefined{@part}{}{\lakin@parttocnum\@part}
\iflakin@parttocnum\else\PackageWarning{hebrew-lualatex-pdf}{Part TOC number is
  not boxed: no known \noexpand\part wrapper matched. A part title that BEGINS
  with Latin will print its number on the wrong side of its TOC row}\fi
\patchcmd{\section}{\bfseries}{\bfseries\boldmath}{}
  {\PackageError{hebrew-lualatex-pdf}{Cannot patch section math weight}{Unsupported document class}}
\patchcmd{\subsection}{\bfseries}{\bfseries\boldmath}{}
  {\PackageError{hebrew-lualatex-pdf}{Cannot patch subsection math weight}{Unsupported document class}}
\patchcmd{\subsubsection}{\bfseries}{\bfseries\boldmath}{}
  {\PackageError{hebrew-lualatex-pdf}{Cannot patch subsubsection math weight}{Unsupported document class}}
\patchcmd{\paragraph}{\bfseries}{\bfseries\boldmath}{}{}
\patchcmd{\subparagraph}{\bfseries}{\bfseries\boldmath}{}{}
\@ifundefined{l@chapter}{}{%
  \patchcmd{\l@chapter}{\bfseries}{\bfseries\boldmath}{}{}}
\makeatother
\setcounter{tocdepth}{\lakittocdepth}
% <hllpdf:measured-toc>
% v4.16.0: Measure actual number labels, not a guessed digit count. Read only
% canonical contentsline records; never execute titles or addtocontents code
% during the measuring pass. Widths are local to this individual list and are
% available before its first row: no per-row drift and no extra aux pass.
\makeatletter
\newread\hllpdf@tocstream
\newdimen\hllpdf@toc@chapter
\newdimen\hllpdf@toc@section
\newdimen\hllpdf@toc@subsection
\newdimen\hllpdf@toc@subsubsection
\newdimen\hllpdf@toc@paragraph
\newdimen\hllpdf@toc@subparagraph
\newdimen\hllpdf@toc@figure
\newdimen\hllpdf@toc@table
\newdimen\hllpdf@tocpage
\newdimen\hllpdf@toc@sectionindent
\newdimen\hllpdf@toc@subsectionindent
\newdimen\hllpdf@toc@subsubsectionindent
\newdimen\hllpdf@toc@paragraphindent
\newdimen\hllpdf@toc@subparagraphindent
\def\hllpdf@tocdefaults{%
  \hllpdf@toc@chapter=2em
  \hllpdf@toc@section=3.2em
  \hllpdf@toc@subsection=3.9em
  \@ifundefined{c@chapter}{%
    \hllpdf@toc@subsubsection=3.2em
    \hllpdf@toc@paragraph=4.1em
    \hllpdf@toc@subparagraph=5em
  }{%
    \hllpdf@toc@subsubsection=4.1em
    \hllpdf@toc@paragraph=5em
    \hllpdf@toc@subparagraph=6em
  }%
  \hllpdf@toc@figure=2.3em
  \hllpdf@toc@table=2.3em
  \hllpdf@tocpage=\@pnumwidth\relax
}
\def\hllpdf@tocindents{%
  % Preserve the established shallow layout when its slots already fit.
  \hllpdf@toc@sectionindent=1.5em
  \@tempdima=\hllpdf@toc@chapter \advance\@tempdima by -2em
  \advance\hllpdf@toc@sectionindent by \@tempdima
  \hllpdf@toc@subsectionindent=\dimexpr
    \hllpdf@toc@sectionindent+\hllpdf@toc@section\relax
  \hllpdf@toc@subsubsectionindent=\dimexpr
    \hllpdf@toc@subsectionindent+\hllpdf@toc@subsection\relax
  \hllpdf@toc@paragraphindent=\dimexpr
    \hllpdf@toc@subsubsectionindent+\hllpdf@toc@subsubsection\relax
  \hllpdf@toc@subparagraphindent=\dimexpr
    \hllpdf@toc@paragraphindent+\hllpdf@toc@paragraph\relax
}
\def\hllpdf@tocfont@chapter{\bfseries\boldmath}
\def\hllpdf@tocfont@part{\large\bfseries\boldmath}
\def\hllpdf@tocentryfont{%
  \ifcsname hllpdf@tocfont@\hllpdf@toclevel\endcsname
    \csname hllpdf@tocfont@\hllpdf@toclevel\endcsname
  \fi}
\long\def\hllpdf@tocdiscard#1\hllpdf@tocstop{}
\def\hllpdf@tocscanrecord{%
  \@ifnextchar\contentsline{\hllpdf@tocscanline}{%
    \@ifnextchar\babel@toc{\hllpdf@tocscanlanguage}{\hllpdf@tocdiscard}}}
\long\def\hllpdf@tocscanlanguage\babel@toc#1#2#3\hllpdf@tocstop{%
  % Babel records are declarations, not arbitrary code. Track the language
  % without executing the original directive or changing the surrounding list.
  \def\hllpdf@toclanguage{#1}}
\long\def\hllpdf@tocscanline\contentsline#1#2#3#4#5\hllpdf@tocstop{%
  \def\hllpdf@toclevel{#1}%
  \ifcsname hllpdf@toc@#1\endcsname
    \hllpdf@tocscantitle#2\hllpdf@tocstop
  \fi
  % The page-number column needs its own measurement (including bold chapter
  % rows). Its extra clearance separates the page label from dot leaders.
  \setbox\@tempboxa=\hbox{\foreignlanguage{\hllpdf@toclanguage}{\hllpdf@tocentryfont #3\kern.3em}}%
  \ifdim\wd\@tempboxa>\hllpdf@tocpage
    \hllpdf@tocpage=\wd\@tempboxa
  \fi}
\def\hllpdf@tocscantitle{%
  \@ifnextchar\numberline{\hllpdf@tocscanlabel}{\hllpdf@tocdiscard}}
\long\def\hllpdf@tocscanlabel\numberline#1#2\hllpdf@tocstop{%
  \setbox\@tempboxa=\hbox{\foreignlanguage{\hllpdf@toclanguage}{\hllpdf@tocentryfont #1\kern.5em}}%
  \ifdim\wd\@tempboxa>\csname hllpdf@toc@\hllpdf@toclevel\endcsname
    \csname hllpdf@toc@\hllpdf@toclevel\endcsname=\wd\@tempboxa
  \fi}
\def\hllpdf@tocread{%
  \ifeof\hllpdf@tocstream\else
    \read\hllpdf@tocstream to\hllpdf@tocrecord
    \expandafter\hllpdf@tocscanrecord\hllpdf@tocrecord\hllpdf@tocstop
    \expandafter\hllpdf@tocread
  \fi}
\def\hllpdf@tocmeasure#1{%
  \makeatletter
  \edef\hllpdf@toclanguage{\languagename}%
  \hllpdf@tocdefaults
  \openin\hllpdf@tocstream=\jobname.#1\relax
  \hllpdf@tocread
  \closein\hllpdf@tocstream
  \hllpdf@tocindents
  \edef\@pnumwidth{\the\hllpdf@tocpage}%
  \@tempdima=\dimexpr\hllpdf@tocpage+1em\relax
  \ifdim\@tempdima>\@tocrmarg\relax
    \edef\@tocrmarg{\the\@tempdima}%
  \fi}
\let\hllpdf@originalstarttoc\@starttoc
\def\@starttoc#1{%
  \begingroup
    \def\hllpdf@tocext{#1}%
    \def\hllpdf@tocwanted{toc}%
    \ifx\hllpdf@tocext\hllpdf@tocwanted \hllpdf@tocmeasure{#1}\else
      \def\hllpdf@tocwanted{lof}%
      \ifx\hllpdf@tocext\hllpdf@tocwanted \hllpdf@tocmeasure{#1}\else
        \def\hllpdf@tocwanted{lot}%
        \ifx\hllpdf@tocext\hllpdf@tocwanted \hllpdf@tocmeasure{#1}\fi
      \fi
    \fi
    \hllpdf@originalstarttoc{#1}%
  \endgroup}
\hllpdf@tocdefaults
\hllpdf@tocindents
\@ifundefined{l@chapter}{}{%
  \patchcmd{\l@chapter}{\setlength\@tempdima{1.5em}}%
    {\setlength\@tempdima{\hllpdf@toc@chapter}}{}%
    {\PackageError{hebrew-lualatex-pdf}{Cannot size chapter TOC numbers}%
       {Unsupported contents layout: the chapter number slot was not found}}}
\renewcommand*\l@section{\@dottedtocline{1}%
  {\hllpdf@toc@sectionindent}{\hllpdf@toc@section}}
\renewcommand*\l@subsection{\@dottedtocline{2}%
  {\hllpdf@toc@subsectionindent}{\hllpdf@toc@subsection}}
\renewcommand*\l@subsubsection{\@dottedtocline{3}%
  {\hllpdf@toc@subsubsectionindent}{\hllpdf@toc@subsubsection}}
\renewcommand*\l@paragraph{\@dottedtocline{4}%
  {\hllpdf@toc@paragraphindent}{\hllpdf@toc@paragraph}}
\renewcommand*\l@subparagraph{\@dottedtocline{5}%
  {\hllpdf@toc@subparagraphindent}{\hllpdf@toc@subparagraph}}
\renewcommand*\l@figure{\@dottedtocline{1}{1.5em}{\hllpdf@toc@figure}}
\renewcommand*\l@table{\@dottedtocline{1}{1.5em}{\hllpdf@toc@table}}
\makeatother
% </hllpdf:measured-toc>
\setcounter{secnumdepth}{\lakitsecnumdepth}
% Public preamble-only caption-numbering policy. The historical within-major
% default is untouched unless this command is called. "continuous" removes the
% reset edge as well as the printed prefix, so Figure 1/2 and Table 1/2 stay
% unique across sections/chapters. Returning to "within-major" restores the
% original class-aware displays and native destination conventions. The policy
% never replaces \figcap/\tabcap or their measured link-placement mechanism.
\makeatletter
\let\lakin@majorfig\thefig
\let\lakin@majortbl\thetbl
\let\lakin@majorHfig\theHfig
\let\lakin@majorHtbl\theHtbl
\newcommand{\lakitcaptionnumbering}[1]{%
  \edef\lakin@numbering{#1}%
  \def\lakin@continuous{continuous}\def\lakin@withinmajor{within-major}%
  \ifx\lakin@numbering\lakin@continuous
    \@ifundefined{c@chapter}{%
      \@removefromreset{fig}{section}\@removefromreset{tbl}{section}%
    }{\@removefromreset{fig}{chapter}\@removefromreset{tbl}{chapter}}%
    \renewcommand{\thefig}{\arabic{fig}}\renewcommand{\thetbl}{\arabic{tbl}}%
    \def\theHfig{\arabic{fig}}\def\theHtbl{\arabic{tbl}}%
  \else\ifx\lakin@numbering\lakin@withinmajor
    \@ifundefined{c@chapter}{%
      \@removefromreset{fig}{section}\@removefromreset{tbl}{section}%
      \@addtoreset{fig}{section}\@addtoreset{tbl}{section}%
    }{%
      \@removefromreset{fig}{chapter}\@removefromreset{tbl}{chapter}%
      \@addtoreset{fig}{chapter}\@addtoreset{tbl}{chapter}}%
    \let\thefig\lakin@majorfig\let\thetbl\lakin@majortbl
    \let\theHfig\lakin@majorHfig\let\theHtbl\lakin@majorHtbl
  \else
    \PackageError{hebrew-english-latex-studio}{Unknown caption-numbering policy '#1'}%
      {Use continuous or within-major before begin{document}.}%
  \fi\fi}
\@onlypreamble\lakitcaptionnumbering
\makeatother
% Architecture readiness sentinel. Optional runtime layers require the complete
% core, so the marker is defined only after every mechanism/baseline definition.
% <hllpdf:architecture-ready>
\def\hllpdfcoreloaded{1}
% </hllpdf:architecture-ready>

% v4.22: registered file resources, not an assertion about embedded subset bytes.
% The log framing avoids long-path wrapping and preserves UTF-8 as hex.
% <hllpdf:font-inventory>
\directlua{
  local function inventory()
    local names = {}
    local unnamed = 0
    for id, item in font.each() do
      local path = item.filename
      if path and not (path == "") then names[path] = true else unnamed = unnamed + 1 end
    end
    local sorted = {}
    for path in pairs(names) do table.insert(sorted,path) end
    table.sort(sorted)
    texio.write_nl("log","HLLFONT|BEGIN|1")
    for index,path in ipairs(sorted) do
      local hex = path:gsub(".",function(c) return string.format(string.char(37) .. "02x",string.byte(c)) end)
      local count = math.ceil(string.len(hex)/48)
      texio.write_nl("log","HLLFONT|FILE|" .. index .. "|" .. count)
      for offset=1,string.len(hex),48 do texio.write_nl("log","HLLFONT|HEX|" .. hex:sub(offset,offset+47)) end
    end
    texio.write_nl("log","HLLFONT|END|" .. rawlen(sorted) .. "|" .. unnamed)
    texio.write_nl("log","")
  end
  luatexbase.add_to_callback("finish_pdffile",inventory,"hllpdf.runtime.fonts")
}
% </hllpdf:font-inventory>
     % (or paste these lines at the top)
%   \begin{document}  ...  \end{document}
%
% Compile with:  lualatex -interaction=nonstopmode doc.tex
%   Repeat until two consecutive PDFs are byte-identical (scripts/build.sh
%   does this) — a TOC document converges only on pass 3+; after two passes
%   the TOC page numbers are STALE (measured; see SKILL.md step 3).
% Font binaries are external dependencies, pinned by assets/fonts/SHA256SUMS.
% scripts/setup_fonts.sh imports a verified local installation or prior skill.
%
% PUBLIC DOCUMENT-SHELL / LAYOUT CONTRACT (introduced v4.3.1; architecture-owned v4.7.1). These hooks may be
% defined BEFORE % ============================================================================
% preamble.tex — tested LuaLaTeX preamble for Hebrew + math + figures PDFs (v4.32.0)
% Battle-tested: compiled an 80-page Hebrew quantum-computing summary AND a
% 197-page Hebrew computer-architecture anthology (TikZ diagrams, code listings,
% colored boxes) — both with 0 missing characters.
%
% Usage in your document:
%   \input{preamble.tex}     % (or paste these lines at the top)
%   \begin{document}  ...  \end{document}
%
% Compile with:  lualatex -interaction=nonstopmode doc.tex
%   Repeat until two consecutive PDFs are byte-identical (scripts/build.sh
%   does this) — a TOC document converges only on pass 3+; after two passes
%   the TOC page numbers are STALE (measured; see SKILL.md step 3).
% Font binaries are external dependencies, pinned by assets/fonts/SHA256SUMS.
% scripts/setup_fonts.sh imports a verified local installation or prior skill.
%
% PUBLIC DOCUMENT-SHELL / LAYOUT CONTRACT (introduced v4.3.1; architecture-owned v4.7.1). These hooks may be
% defined BEFORE \input{preamble.tex}; if omitted, the historical canonical
% defaults below are exact.  They deliberately control only the document shell
% and page layout.  RTL/math semantics stay in this preamble, while visual
% presentation themes (for example clean_theme.tex) load AFTER it.
%
%   \def\lakitdocumentclass{article}
%   \def\lakitclassoptions{11pt,a4paper}
%   \def\lakitgeometryoptions{margin=1.8cm,headheight=14pt}
%   \def\lakitlinespread{1.15}
%   \def\lakitparagraphstyle{block} % letters/briefs; indented for continuous prose
%   \def\lakittocdepth{1}
%   \def\lakitsecnumdepth{2}
%   \input{preamble.tex}
%
% Supported class contract: report or article.  Alternate paper sizes must also
% be declared to build.sh through EXPECTED_PAPER; the PDF gate never infers them.
% Do NOT use these hooks as a theme mechanism: theme overlays must not own page
% geometry, class selection, or direction semantics. A declared presentation-policy
% overlay may intentionally change numbering visibility after the preamble.
% ============================================================================
% Runtime architecture guard: generated from scripts/architecture_contract.py.
% The readiness sentinel remains delayed until the END of this file.
% <hllpdf:architecture-guard>
\ifdefined\hllpdfcoreloaded
  \errmessage{hllpdf: preamble.tex loaded more than once}
  \expandafter\endinput
\fi
% </hllpdf:architecture-guard>

\ifdefined\lakitdocumentclass\else\def\lakitdocumentclass{report}\fi
\ifdefined\lakitclassoptions\else\def\lakitclassoptions{12pt,a4paper}\fi
\ifdefined\lakitgeometryoptions\else\def\lakitgeometryoptions{margin=2.3cm,headheight=14pt}\fi
\ifdefined\lakitlinespread\else\def\lakitlinespread{1.3}\fi
\ifdefined\lakittocdepth\else\def\lakittocdepth{1}\fi
\ifdefined\lakitsecnumdepth\else\def\lakitsecnumdepth{2}\fi
% Paragraph organization belongs to the document layout, independently of its
% language and theme. Preserve historical class dimensions unless explicitly
% selecting block paragraphs; never suppress indentation with an everypar hack.
\ifdefined\lakitparagraphstyle\else\def\lakitparagraphstyle{indented}\fi
\edef\lakitparagraphstyle{\lakitparagraphstyle}
\def\hllpdfparagraphblock{block}
\def\hllpdfparagraphindented{indented}
\ifx\lakitparagraphstyle\hllpdfparagraphblock\else
  \ifx\lakitparagraphstyle\hllpdfparagraphindented\else
    \errmessage{hllpdf: paragraph style must be block or indented}
  \fi
\fi
% v4.3.1: the public class hook has a deliberately narrow support contract.
% Before this guard, \def\lakitdocumentclass{book} silently compiled even though
% every heading/TOC regression test targets report/article; an unsupported shell
% could therefore look "supported" until a subtle class-specific defect shipped.
% Fail before \documentclass rather than letting an undocumented combination enter
% the mechanism core. Plain-TeX comparison is used because LaTeX packages are not
% available yet at this point.
\edef\lakitvalidatedclass{\lakitdocumentclass}
\def\lakitreportclass{report}
\def\lakitarticleclass{article}
\ifx\lakitvalidatedclass\lakitreportclass
\else\ifx\lakitvalidatedclass\lakitarticleclass
\else
  \errmessage{hllpdf: unsupported class \lakitdocumentclass; use report or article}
\fi\fi
\let\lakitvalidatedclass\relax
\let\lakitreportclass\relax
\let\lakitarticleclass\relax
% Language intent is independent of document class, theme and content domain.
% The two dominant-language profiles differ in authoring intent, not direction.
% A balanced bilingual document must explicitly choose a navigation/base locale.
% <hllpdf:language-profile>
\providecommand{\lakitlanguageprofile}{hebrew}
\edef\lakitlanguageprofile{\lakitlanguageprofile}
\def\hllpdfhebrew{hebrew}\def\hllpdfenglish{english}
\def\hllpdfhebrewmixed{hebrew-with-english}
\def\hllpdfenglishmixed{english-with-hebrew}
\def\hllpdfbilingual{bilingual}
\let\hllpdfexpectedlanguage\relax
\ifx\lakitlanguageprofile\hllpdfhebrew
  \let\hllpdfexpectedlanguage\hllpdfhebrew
\else\ifx\lakitlanguageprofile\hllpdfhebrewmixed
  \let\hllpdfexpectedlanguage\hllpdfhebrew
\else\ifx\lakitlanguageprofile\hllpdfenglish
  \let\hllpdfexpectedlanguage\hllpdfenglish
\else\ifx\lakitlanguageprofile\hllpdfenglishmixed
  \let\hllpdfexpectedlanguage\hllpdfenglish
\else\ifx\lakitlanguageprofile\hllpdfbilingual
  \ifdefined\lakitmainlanguage\else
    \errmessage{hllpdf: bilingual profile requires explicit lakitmainlanguage}
  \fi
\else
  \errmessage{hllpdf: unsupported language profile \lakitlanguageprofile}
\fi\fi\fi\fi\fi
\ifdefined\lakitmainlanguage
  \edef\lakitmainlanguage{\lakitmainlanguage}
\else
  \let\lakitmainlanguage\hllpdfexpectedlanguage
\fi
\ifx\lakitmainlanguage\hllpdfhebrew\else
  \ifx\lakitmainlanguage\hllpdfenglish\else
    \errmessage{hllpdf: main language must be hebrew or english}
  \fi
\fi
\ifx\hllpdfexpectedlanguage\relax\else
  \ifx\lakitmainlanguage\hllpdfexpectedlanguage\else
    \errmessage{hllpdf: main language conflicts with language profile}
  \fi
\fi
\typeout{HLLLANG|\lakitlanguageprofile|\lakitmainlanguage}
% </hllpdf:language-profile>
\documentclass[\lakitclassoptions]{\lakitdocumentclass}
\usepackage[no-math]{fontspec}        % (1) ★ no-math — THE critical line (see recipe.md)
% layout=lists asks Babel to adapt list parshape when paragraph and page
% directions differ. Without it, minority-language list bodies lose their
% indentation and labels hang outside the text block, despite a clean log.
% The upstream mechanism preserves ordinary same-direction list geometry.
\usepackage[bidi=basic,layout=lists]{babel}
\ifx\lakitmainlanguage\hllpdfenglish
  \babelprovide[import]{hebrew}
  \babelprovide[import,main]{english}
\else
  \babelprovide[import,main]{hebrew}
  \babelprovide[import]{english}
\fi
% v4.22: bind canonical families and fallbacks to the selected installation.
% Lookup by basename or family alone is not proof of the file actually used.
% <hllpdf:font-binding>
\directlua{
  local dir = os.getenv("HEBREW_LATEX_FONT_DIR")
  if not dir or dir == "" then
    local data = os.getenv("XDG_DATA_HOME")
    if not data or data == "" or not (data:sub(1,1) == "/") then
      data = (os.getenv("HOME") or "") .. "/.local/share"
    end
    dir = data .. "/fonts/hebrew-lualatex-pdf"
  end
  if not (dir:sub(1,1) == "/") then dir = lfs.currentdir() .. "/" .. dir end
  if dir:find("[",1,true) or dir:find("]",1,true) or dir:find(string.char(0),1,true)
     or dir:find(string.char(10),1,true) or dir:find(string.char(13),1,true) then
    tex.error("hllpdf: unsupported character in font directory")
  end
  if not (lfs.attributes(dir,"mode") == "directory") then
    tex.error("hllpdf: font directory is missing")
  end
  hllpdf_font_dir = dir:gsub("/+$", "") .. "/"
  tex.sprint(string.char(92) .. "def" .. string.char(92) .. "hllpdffontdir{")
  tex.sprint(-2,hllpdf_font_dir)
  tex.sprint("}")
  luaotfload.add_fallback("hebfb",{"[" .. hllpdf_font_dir .. "DejaVuSerif.ttf]:mode=node;"})
  luaotfload.add_fallback("hebfbbf",{"[" .. hllpdf_font_dir .. "DejaVuSerif-Bold.ttf]:mode=node;"})
}
% </hllpdf:font-binding>
% v3.22: Frank Ruhl Libre ships NO italic face. Instead of silently-upright
% \emph (v<=3.20) or bold-as-emphasis (v3.21), synthesize a slanted face with
% FakeSlant — real, VISIBLE italic emphasis with correct \emph nesting
% semantics (inner \emph toggles back upright). Verified empirically at
% 300dpi: slant clearly visible, 0 errors, 0 missing glyphs.
\babelfont{rm}[Path={\hllpdffontdir},RawFeature={fallback=hebfb},
  BoldFont={FrankRuhlLibre-Bold.ttf},
  BoldFeatures={RawFeature={fallback=hebfbbf}},   % v3.56 — see hebfbbf above
  ItalicFont={FrankRuhlLibre-Regular.ttf}, ItalicFeatures={FakeSlant=0.22},
  % Bold-italic uses the weight-matched fallback and a synthetic slant.
  BoldItalicFont={FrankRuhlLibre-Bold.ttf}, BoldItalicFeatures={RawFeature={fallback=hebfbbf},FakeSlant=0.22}
  ]{FrankRuhlLibre-Regular.ttf}
% Sans fallbacks preserve family and weight for the few missing codepoints.
\directlua{luaotfload.add_fallback("hebfbsans",{"[" .. hllpdf_font_dir .. "DejaVuSans.ttf]:mode=node;"})}
\directlua{luaotfload.add_fallback("hebfbsansbf",{"[" .. hllpdf_font_dir .. "DejaVuSans-Bold.ttf]:mode=node;"})}
\babelfont{sf}[Path={\hllpdffontdir},RawFeature={fallback=hebfbsans},
  BoldFont={Heebo-Bold.ttf},
  BoldFeatures={RawFeature={fallback=hebfbsansbf}}]{Heebo-Regular.ttf}
% v3.24: Scale=MatchLowercase — DejaVu Sans Mono has an outsized x-height, so
% unscaled inline code visually SHOUTS next to the delicate Frank Ruhl body
% (measured on a 197-page build: scaled, a code atom sits in the line at the
% body's weight, and the document tightened by ~5%). Applies to listings too.
\babelfont{tt}[Path={\hllpdffontdir},RawFeature={fallback=hebfb},Scale=MatchLowercase,
  BoldFont={DejaVuSansMono.ttf},
  BoldFeatures={RawFeature={fallback=hebfb},FakeBold=1.5},
  ItalicFont={DejaVuSansMono.ttf},
  ItalicFeatures={RawFeature={fallback=hebfb},FakeSlant=0.2},
  BoldItalicFont={DejaVuSansMono.ttf},
  BoldItalicFeatures={RawFeature={fallback=hebfb},FakeBold=1.5,FakeSlant=0.2}
  ]{DejaVuSansMono.ttf}

% Explicit multilingual/English profiles use real Latin faces and italics.
% The legacy/default Hebrew profile retains its established measured typography.
% Latin Modern is provided by the TeX distribution, not the pinned font bundle;
% recorder/runtime evidence reports it as an additional toolchain resource.
\ifx\lakitlanguageprofile\hllpdfhebrew\else
  \babelfont[english]{rm}[Ligatures=TeX,RawFeature={fallback=hebfb},
    BoldFeatures={RawFeature={fallback=hebfbbf}},
    BoldItalicFeatures={RawFeature={fallback=hebfbbf}}]{Latin Modern Roman}
  \babelfont[english]{sf}[Ligatures=TeX,RawFeature={fallback=hebfbsans},
    BoldFeatures={RawFeature={fallback=hebfbsansbf}},
    BoldItalicFeatures={RawFeature={fallback=hebfbsansbf}}]{Latin Modern Sans}
\fi

\usepackage{amsmath,amssymb,mathtools}
% amsthm supplies theorem/proof machinery. Environment names remain
% document-specific; the generic proof label is localized below.
\usepackage{amsthm}
% Proof names follow the surrounding prose, including return from an island.
\addto\extrashebrew{\renewcommand{\proofname}{הוכחה}}
\addto\extrasenglish{\renewcommand{\proofname}{Proof}}
\ifx\lakitmainlanguage\hllpdfenglish
  \renewcommand{\proofname}{Proof}
\else
  \renewcommand{\proofname}{הוכחה}
\fi
% v3.24.1: equation tags under babel-bidi. A tag re-emitted by mathtools'
% showonlyrefs inside align inherits the RTL base and its parens MIRROR —
% ")2.1(" instead of "(2.1)" (plain align/equation tags are unaffected;
% verified on a minimal case + a 198-page build). LTR-isolate at the tag-BOX
% level: mathtools' show/hide + numbering logic is untouched. Do NOT override
% \tagform@ instead — that clobbers the showonlyrefs mechanism (measured:
% every tag shows, numbering breaks).
% v3.28.1: the v3.24.1 wrap must apply ONLY to tags of display environments.
% Inline \eqref in Hebrew PROSE also goes through \maketag@@@, and there the
% english-wrap MIRRORS the parens — ")1(" — while the unwrapped original is
% correct (proven on a 2x2 isolation matrix: patch x showonlyrefs; regressed
% every prose \eqref since v3.24.1, caught by the render gate's R4). The two
% call sites are indistinguishable by \ifmmode (tag boxes are text mode), so
% the display environments themselves raise a flag; \eqref in prose sees it
% false and gets the original untouched \maketag@@@.
\usepackage{etoolbox}

% babel bidi=basic (including upstream 26.9) assumes every automatic
% discretionary has a one-glyph replacement. LuaTeX can legally produce one
% with replace=nil after a hyphenated Latin value; Babel then indexes nil and the
% whole build dies. Do the narrow upstream-compatible correction: while Babel
% classifies the list, present only that malformed disc as an ordinary
% discretionary, then restore its subtype. Babel still reorders the ENTIRE box;
% there is no pcall and no successful fallback to an unreordered node list.
% If Babel throws for any other reason, the build remains fatal. The synthetic
% nil-disc probe and the real table gallery in test_all.sh pin both layers.
\AtBeginDocument{\directlua{
  if Babel and Babel.bidi and not Babel.heb_nil_disc_compat then
    Babel.heb_nil_disc_compat = true
    local original = Babel.bidi
    local disc_id = node.id('disc')
    Babel.bidi = function(head, ispar, hdir)
      local hidden = {}
      for item in node.traverse_id(disc_id, head) do
        if item.subtype == 2 and (not item.replace or not item.replace.char) then
          table.insert(hidden, item)
          item.subtype = 0
        end
      end
      local result = original(head, ispar, hdir)
      for _, item in ipairs(hidden) do item.subtype = 2 end
      return result
    end
  end
}}
\makeatletter
\newif\ifHEBdisptag
\let\ORIGmaketag@@@\maketag@@@
\def\maketag@@@#1{\ifHEBdisptag\ORIGmaketag@@@{\foreignlanguage{english}{#1}}\else\ORIGmaketag@@@{#1}\fi}
\@for\HEB@env:=equation,equation*,align,align*,gather,gather*,multline,multline*,flalign,flalign*,alignat,alignat*,eqnarray,eqnarray*\do{%
  \expandafter\AtBeginEnvironment\expandafter{\HEB@env}{\HEBdisptagtrue}}
\makeatother
% English page + Hebrew prose: Babel's equation-only tag flip uses a primitive
% leqno path that shifts the formula to the left edge and the tag into the body
% of an English page. Align's full-width alignment is unaffected. Give only this
% display a physical LTR container and native amsmath eqno/leqno. The engine then
% owns real tag width, collision handling and labels; no proxy tag box hides ink
% from its fit calculation. Babel's saved \bbl@eqnodir still restores Hebrew text
% inside math, and the prose locale/math direction are never switched.
% Close the prior RTL paragraph before normalization. After the environment's
% group restores its directions, close TeX's empty resumed display paragraph and
% use native LaTeX continuation handling: same-paragraph prose has no indent,
% while an authored blank line starts an ordinary indented RTL paragraph.
% Register after Babel's setup; remain local to equation(*) in a main-English
% document's actual RTL prose. Native Hebrew, ordinary English and explicit
% fleqn retain their original paths. No page/display dimensions are assigned.
% <hllpdf:minority-rtl-equation-geometry>
\makeatletter
\newcommand{\HEB@minorityprosecontinuation}{\par\@doendpe}
\newcommand{\HEB@minorityequationgeometry}{%
  \ifx\lakitmainlanguage\hllpdfenglish
    \ifnum\pardirection=\@ne
      \if@fleqn\else
        \par\aftergroup\HEB@minorityprosecontinuation
        \iftagsleft@\let\veqno\@@leqno
        \else\let\veqno\@@eqno\fi
        \bodydir TLT\pardir TLT\textdir TLT\relax
      \fi
    \fi
  \fi}
\AtBeginDocument{%
  \AddToHook{env/equation/begin}{\HEB@minorityequationgeometry}%
  \AddToHook{env/equation*/begin}{\HEB@minorityequationgeometry}}
\makeatother
% </hllpdf:minority-rtl-equation-geometry>
% v3.63: \shortintertext rendered Hebrew MIRROR-REVERSED and flush LEFT — an entire
% commentary line inside an align printed backwards. babel's bidi=basic patches
% \intertext@ only (luababel.def), and mathtools' \shortintertext uses the identical
% \vbox{\normalbaselines … \noindent\ignorespaces#1\par} skeleton but is never
% patched, so its paragraph keeps pardir/textdir LTR: no bidi reordering (logical
% order on the page) and flush-left. Measured against \intertext with the same
% Applying babel's own replacement to \shortintertext@ makes the two identical.
% \h{}/\mbox fix the letter order but leave the line flush LEFT, so they are not
% substitutes. Guarded + \AtBeginDocument because mathtools \let-copies at load time.
% NB \ExplSyntaxOn is REQUIRED: mathtools' internals carry expl3 names, and without
% the expl3 catcodes the `_` is read as a subscript and the patch detonates far away
% with "! Illegal unit of measure (pt inserted)" at the \end{align} (measured).
\makeatletter
\AtBeginDocument{%
  \@ifundefined{shortintertext@}{}{%
    \ExplSyntaxOn
    \bbl@sreplace\shortintertext@{\normalbaselines}%
      {\normalbaselines\ifx\bbl@eqnodir\relax\else\bbl@pardir\@ne\bbl@eqnodir\fi}%
    \ExplSyntaxOff}}
\makeatother
% geometry inherits the canonical A4 sheet from \documentclass; build.sh then
% verifies every rendered MediaBox and CropBox, so a silent Letter fallback is
% impossible to publish as a clean default build.
\usepackage[\lakitgeometryoptions]{geometry}
\linespread{\lakitlinespread}       % layout hook; canonical default remains 1.3
\ifx\lakitparagraphstyle\hllpdfparagraphblock
  \setlength{\parindent}{0pt}
  \setlength{\parskip}{6pt plus 1pt minus 1pt}
\fi
% These dimensions remain ordinary scoped TeX lengths: authors can customize
% them after core input. Language switches inherit them; lists keep their own
% indentation and narrow reference layouts establish their local paragraph gaps.
\typeout{HLLPAR|\lakitparagraphstyle}
\usepackage{xcolor}
\usepackage{mdframed}
\usepackage{listings}                 % code listings (assembly / C / pseudo-code)
\usepackage{tikz}
\usetikzlibrary{arrows.meta,positioning,calc,shapes.geometric,fit,backgrounds}
\usepackage{pgfplots}\pgfplotsset{compat=1.18}
\usepackage{quantikz}                 % quantum circuits (remove if unused)
\usepackage{booktabs}
\usepackage{array}
\usepackage{float}                    % [H] = place table exactly here (reference doc: no drifting)
\usepackage{enumitem}
\usepackage{needspace}

% ============================================================================
% ★★ v3.4 — THE DIRECTION-ISOLATION TRIO (the heart of mixing Hebrew + Latin)
% Mental model:
%   \h{...}    Hebrew (RTL) island      — Hebrew text inside an LTR context
%                                          (TikZ nodes, the english `tikzpic`).
%   \en{...}   English (LTR) isolate    — Latin text / parentheses inside a
%                                          Hebrew (RTL) context. KEEPS PARENS
%                                          AND ORDER CORRECT.
%   \code{...} \en + monospace          — code inside Hebrew / inside a frametitle.
% CRITICAL: never wrap Latin in \h{} (it is for HEBREW). Inside a TikZ node,
% Latin in \h{} renders REVERSED (Datapath->htapataD, CPU->UPC) and parens
% mirror. Isolate every Latin/paren chunk with \en{} or \code{}. See
% references/figures-boxes-listings.md.
% ---- TEXT INSIDE MATH (v3.8) ----------------------------------------------
% Put ANY text that sits inside math in a \text{}. The ONLY thing that is
% invalid in math is a bare \h/\en (= \foreignlanguage) used DIRECTLY, outside
% \text -> "\rmfamily invalid in math mode". Inside \text:
%   Hebrew label           -> \text{...}            (bare Hebrew renders fine;
%                              \h is optional — BUT NOT ALWAYS HARMLESS (v3.67):
%                              if the \h group OPENS with a nested island
%                              ($..$, \en{}, \code{}, \texttt{}, \mbox{}) the
%                              WHOLE group loses bidi and the Hebrew prints in
%                              LOGICAL order = gibberish. Measured:
%                                \text{\h{$f$ קבועה}}  -> reads "העובק"  BROKEN
%                                \text{\h{קבועה $f$}}  -> correct
%                                \text{$f$ קבועה}      -> correct (no \h)
%                              Hebrew first inside \h, or drop the \h. Flagged
%                              H-OPENS-WITH-LTR-GROUP.)
%   plain Latin word        -> \text{...}            (no marker needed)
%   Latin WITH ( ) or [ ]   -> \text{\en{...}}       (\en forces LTR; a plain
%                              \text mirrors the parens in RTL -> "Spec (TNR(" )
%   f(Hebrew/text arg)      -> \text{f}(\text{...})  -- keep the parens at MATH
%                              level, NOT \text{f(arg)} (parens inside the box mirror)
%   code in math            -> \mathtt{...}  (or \texttt{}, math-guarded below)
% NEVER write \en{...} or \h{...} straight into $...$ / \[..\] (outside \text);
% always via \text{}. (\mbox{...} also works but does not scale in scripts.)
% ============================================================================
\newcommand{\h}[1]{\foreignlanguage{hebrew}{#1}}
% v3.42.1 — \en islands do NOT hyphenate (auto OR at explicit hyphens). English hyphenation patterns were
% tearing identifiers and terms ACROSS THE RTL SEAM ("IO Regis-/ters",
% "In-/put", "Band-/width", "Per-/formance", "Multi-/ple" — all shipped on
% real pages; probe-proven: "ALUS-/rcB", "Mem-/ToReg"). An island is a label
% or a term, not English prose — a hyphen tear there is never right. A/B on
% a full 197pp lakit: identical overfull sets, same page count, ~11
% paragraphs healed, one mild underfull (badness 3525) that replaced a
% mutilated "Regis-/ters". Long English PROSE belongs in the englishprose
% environment below, which restores explicit-hyphen breaks while its paragraph
% is still active. A raw otherlanguage block inherits the global prohibition.
% (\l@nohyphenation needs \csname:
% '@' is not a letter here.)
% v3.43.2 — \exhyphenpenalty moved OUT of \en to a GLOBAL assignment below.
% Render-proven: a group-local \exhyphenpenalty is DEAD CODE — TeX reads the
% line-breaking parameters at \par time, after the \en group has closed, so
% \en{write-through} / plain not-taken still tore at the "-" ("write-/
% through" shipped on a real exam page). \language is different: it is
% recorded in the node list as text is added, so the in-group nohyphenation
% assignment above DOES work (and stays). The global penalty was A/B-measured
% on a 197pp lakit: identical overfull set (25), identical page count, the
% not-/taken tear healed. Hebrew is unaffected in practice — Hebrew compounds
% use maqaf (־), not ASCII "-".
\newcommand{\en}[1]{\foreignlanguage{english}{\language=\csname l@nohyphenation\endcsname #1}}
% Babel runs extras locally for inline islands and paragraph language changes.
% An English main document must not inherit the Hebrew no-hyphen paragraph rule.
\addto\extrashebrew{\exhyphenpenalty=10000\relax}
\addto\extrasenglish{\exhyphenpenalty=50\relax}
\ifx\lakitmainlanguage\hllpdfenglish
  \exhyphenpenalty=50\relax
\else
  \exhyphenpenalty=10000\relax
\fi
% Full English paragraphs need ordinary explicit-hyphen breaks. Use this
% environment instead of a raw otherlanguage block; the paragraph ends while
% the local penalty is still active.
\newenvironment{englishprose}
  {\par\begin{otherlanguage}{english}\exhyphenpenalty=50\relax}
  {\par\end{otherlanguage}}
\newenvironment{hebrewprose}
  {\par\begin{otherlanguage}{hebrew}\exhyphenpenalty=10000\relax}
  {\par\end{otherlanguage}}
% Use for shell text whose locale must not depend on the last body paragraph.
\DeclareRobustCommand{\doclang}[1]{\foreignlanguage{\lakitmainlanguage}{#1}}

% ============================================================================
% ★★ — fixes for defects that are INVISIBLE in the log and in the checks
% and only surface on a VISUAL RENDER of the PDF. (Each verified empirically.)
% ----------------------------------------------------------------------------
% v3.22: \emph is now REAL slanted emphasis (FakeSlant italic face defined on
% the \babelfont{rm} line above), replacing the v3.21 \emph->\textbf hack: it
% preserves \emph's toggle semantics and keeps bold free for its own role.
% Prefer bold-as-emphasis for your doc? one line: \renewcommand{\emph}[1]{\textbf{#1}}
% page number: force the digits LTR so they never reverse (18 -> 81) in the RTL
%   footer. That reversal is a NON-DETERMINISTIC bidi glitch (observed once in
%   ~80 pp); this LuaLaTeX build is not byte-reproducible, so the defensive
%   LTR wrap is the only reliable guard. Latin digits are always LTR anyway.
\renewcommand{\thepage}{\foreignlanguage{english}{\arabic{page}}}
% v3.70: \pagenumbering \gdef's \thepage and silently DESTROYS the wrap above —
% after \pagenumbering{arabic} the folio guard is gone for the whole rest of the
% document. Two collections discovered this independently and hand-restore the
% wrap in main.tex with the same comment; that knowledge now lives here instead.
% The wrapper re-applies the LTR wrap after the kernel call, for any style.
\makeatletter
\let\lakin@orig@pagenumbering\pagenumbering
\renewcommand{\pagenumbering}[1]{%
  \lakin@orig@pagenumbering{#1}%
  \renewcommand{\thepage}{\foreignlanguage{english}{\csname @#1\endcsname\c@page}}}
\makeatother
% emergencystretch: Hebrew does not hyphenate, so an unbreakable LTR island
%   (\en{...} / an inline math run) landing at the margin can leave a small
%   (~2-13pt) overfull with no legal breakpoint. This rescues ONLY those
%   inline-prose tie-breaks; it is a ceiling used only on offending lines and
%   does NOT touch display-math / listing / table overflow (those still surface
%   in the log for the content-level overflow pass). Verified: baseline overfull
%   -> 0 with this set, with no change to pagination or underfull.
\setlength{\emergencystretch}{3em}
% v3.40: inline math (almost) never breaks across lines. TeX's default
%   (binop 700 / rel 500) lets an inline run split at an operator —
%   $(1.f)\cdot 2^{e}$ shipped as "(1.f)·" ending one line and "2^e" opening
%   the next (reader-reported). In LTR that split is merely ugly; under an
%   RTL base the carried operator lands on the visually WRONG side of the
%   break and the formula reads backwards.
%   WHY 10000 (v3.45.0 — the display-first resolution). Inline math is
%   FORBIDDEN from breaking, so a torn island under an RTL base is
%   structurally impossible. The v3.40 cycle used 9999 ("break only as a last
%   resort") because a bare 10000 left ~30-39 WIDE inline equations
%   overflowing (worst 90.7pt) — but 9999's "last resort" fires OFTEN in
%   tight, non-hyphenating Hebrew: a reader found torn islands across the
%   fleet (~150 candidates; e.g. "P(y |" ending one line and "x_0)" opening
%   the next, read backwards). The step the 9999 cycle missed: every one of
%   those ~30 overflows is a wide equation that BELONGS in display. Convert
%   them to display FIRST, then 10000 gives 0 tears AND 0 overflows (measured,
%   12/12 fleet). Doctrine: inline math is ATOMIC; a wide equation goes in a
%   display, never inline. See content-style ("inline math is atomic").
\binoppenalty=10000 \relpenalty=10000
% v3.38: uniform sentence spacing. TeX's default (\nonfrenchspacing, sfcode
% 3000) is the OLD Anglo-American wide-space-after-period convention — Hebrew
% typography never had it, and modern Latin practice (Chicago et al.) dropped
% it too. Measured on real pages: post-period gaps up to 1.9-2.1x the line's
% mean word gap — the dominant "rivers" source in justified Hebrew prose.
% Survives babel language switches (probed: sfcode stays french inside \en{}
% islands and otherlanguage blocks).
\frenchspacing
% v3.38: widow/orphan protection. LaTeX's default penalties (150/150/50)
% allow a lone line at a page top/bottom almost freely. 10000 = forbid;
% this is SAFE here precisely because the document is \raggedbottom (the
% demoted line just moves on, no vertical stretch fight) — under flushbottom
% these values would cause underfull pages. raggedbottom is the report-class
% default; made explicit so a future class/option change cannot silently
% flip it and turn the 10000s harmful.
\raggedbottom
\widowpenalty=10000 \clubpenalty=10000 \displaywidowpenalty=10000
% em-dash: --- does NOT become — under Frank Ruhl (the font lacks tlig; verified
%   the width is identical with/without Ligatures=TeX, while the mechanism works
%   on Latin Modern). TYPE THE CHARACTER — (U+2014) DIRECTLY; likewise – (U+2013).
%   check_bidi_figures.py flags a --/--- that is adjacent to Hebrew.
% ============================================================================

% ---- STRUCTURAL DIRECTION-SAFETY for \texttt (v3.7) -------------------------
% A parenthesis, bracket, or $ INSIDE \texttt{} mirrors/migrates in RTL prose
% (MIPS 0($s0)->"(s0$)0", $rs->"rs$", arr[i]->"[arr[i") because \texttt sets the
% font but NOT the direction. Redefining \texttt to LTR-isolate its content fixes
% the WHOLE class at once (parens, brackets, $) so plain \texttt{lw $t0,0($s0)}
% is safe — no need to remember \code{}. Guarded by \ifmmode because
% \foreignlanguage is invalid in math (there, plain monospace is already LTR).
% NB (corrected v3.56 — the old note here said "stays tofu", which is WRONG):
% \texttt{Hebrew} is a FATAL build failure. The tt family is DejaVu Sans Mono
% (zero Hebrew) and the tt fallback hebfb is DejaVuSerif (also zero Hebrew), so
% the subsetter gets an empty glyph set and LuaTeX dies with
%   "! error: (file .../DejaVuSansMono.ttf) (type 2): there are no glyphs in the
%    subset"  -> no output PDF.
% The message names a FONT FILE, never Hebrew/\texttt/the source line. And it is
% non-local: add an unrelated Latin \code{} elsewhere and the fatal degrades into
% a SILENT DELETION of the Hebrew word. Put Hebrew in \h{}/normal text, never in
% \texttt{}/\code{}; flagged HEBREW-IN-MONO. To disable, delete this block.
\let\origtexttt\texttt
\protected\long\def\texttt#1{\ifmmode\origtexttt{#1}\else\foreignlanguage{english}{\origtexttt{#1}}\fi}
% v4.8.0: defer the direction test until TeX has inserted an alignment cell's
% math template. Unprotected expansion here tested \ifmmode too early at the
% FIRST token of aligned/array/cases/matrix cells. Engine-level \protected also
% preserves the original \texttt implementation without the self-recursion
% observed when retrofitting \DeclareRobustCommand to its saved alias.
\protected\long\def\code#1{\ifmmode\origtexttt{#1}\else\foreignlanguage{english}{\origtexttt{#1}}\fi}

% ---- \codeline — display-code line (v3.24): the code analogue of display math.
% Reader-friction rule (content-style.md §5b): inline \code{} is for a SINGLE
% ATOM only (one register/variable/mnemonic, a tight no-space expression).
% A multi-token expression — an instruction with operands, a spaced assignment,
% a register-binding list — zigzags the eye RTL<->LTR mid-sentence and, being
% an unbreakable island, line-breaks badly. Give it its own centered line with
% display-math vertical rhythm. \centerline = one box => the expression never
% wraps; \nopagebreak keeps it glued to its introducing sentence.
%   ...the Hebrew sentence ends with a colon:
%   \codeline{lw \$t1, -4(\$t0)}
%   and the prose continues (escape dollars as \$).
% v3.66 (first line of the body): flush a PENDING \item label before the \par.
% LaTeX defers the label box (\@labels) until the next \everypar, and an \hbox in
% vertical mode never fires \everypar — so `\item \codeline{...}` printed the code
% at y=69.1 and the item's own "1." BELOW it at y=103.6 (the number appears after
% its content, attached to nothing). \if@inlabel\leavevmode\fi is CONDITIONAL on
% purpose: an unconditional \leavevmode would open an empty paragraph line when
% \codeline follows ordinary prose. Verified both ways — in a list the label now
% sits on the code's line at its normal x, and after prose the geometry is
% byte-identical to before.
\makeatletter
\newcommand{\codeline}[1]{%
  \if@inlabel\leavevmode\fi
  \par\nopagebreak
  \vspace{\abovedisplayskip}\nopagebreak
  % v3.62: was \centerline{...}, i.e. \hbox to\hsize{\hss #1 \hss}. \hss has
  % INFINITE SHRINK, so TeX could always "fit" the box and NEVER reported an
  % Overfull — the same \hss mechanism as BLD-27's \item label. Measured: a wide
  % \codeline in a defbox painted its first token ON the box's left border bar and
  % sat 9.5pt outside the frame interior on both sides, with 0 Overfull, 0 Underfull
  % and all four gates PASS; in bare prose it reached 21.0pt into BOTH margins.
  % \hfil has no shrink, so the overflow is reported and build.sh's ovf_vis/ovf_big
  % counters see it; \linewidth (not \hsize) makes it box-aware, so a \codeline
  % inside a colored box centres on the box interior. Verified: a FITTING \codeline
  % keeps byte-identical glyph positions, and the box-crossing case now emits
  % "Overfull \hbox (19.03pt too wide)". A wide \codeline should be split in two.
  \hbox to\linewidth{\hfil\code{#1}\hfil}%
  \nopagebreak\vspace{\belowdisplayskip}}
\makeatother

% ---- a thin decorative rule helper ----
\newcommand{\sep}{\par\vspace{2pt}\noindent\textcolor{cRule}{\rule{\linewidth}{0.6pt}}\par\vspace{4pt}}

% ---- colors ----
\definecolor{cBlue}{HTML}{1F4E79}   \definecolor{cBlueBg}{HTML}{EAF1F8}
\definecolor{cGreen}{HTML}{2E6B4F}  \definecolor{cGreenBg}{HTML}{EAF4EE}
\definecolor{cAmber}{HTML}{9C6A00}  \definecolor{cAmberBg}{HTML}{FBF3E2}
\definecolor{cPurple}{HTML}{5B2C6F} \definecolor{cPurpleBg}{HTML}{F3ECF7}
\definecolor{cTeal}{HTML}{0E6E6E}   \definecolor{cTealBg}{HTML}{E6F4F4}
\definecolor{cRed}{HTML}{8E2B2B}    \definecolor{cRedBg}{HTML}{FBEAEA}
\definecolor{cGray}{HTML}{555555}   \definecolor{cRule}{HTML}{1F4E79}

% ---- math operators ----
\DeclareMathOperator{\Tr}{Tr}
\DeclareMathOperator{\rank}{rank}
\DeclareMathOperator{\diag}{diag}
\DeclareMathOperator{\sgn}{sgn}

% ---- bra-ket (quantikz defines \ket/\bra/\braket; add the rest safely) ----
\providecommand{\ket}[1]{\left|#1\right\rangle}
\providecommand{\bra}[1]{\left\langle#1\right|}
\providecommand{\braket}[2]{\left\langle#1\middle|#2\right\rangle}
\newcommand{\dyad}[1]{\left|#1\right\rangle\!\left\langle#1\right|}
% v3.21: general outer product |a><b| and expectation value <A> — agents reach
% for these by reflex (esp. \expval); quantikz defines neither, so their absence
% threw "Undefined control sequence" in a real build. \providecommand = safe.
\providecommand{\ketbra}[2]{\left|#1\right\rangle\!\left\langle#2\right|}
\providecommand{\expval}[1]{\left\langle#1\right\rangle}
\newcommand{\abs}[1]{\left|#1\right|}
\newcommand{\norm}[1]{\left\lVert#1\right\rVert}
\newcommand{\half}{\tfrac12}
\newcommand{\R}{\mathbb{R}}\newcommand{\C}{\mathbb{C}}\newcommand{\Z}{\mathbb{Z}}
% identity operator. NB: \mathbb{1} silently renders a WRONG glyph (msbm has no
% blackboard digits). Use \mathbb{I} = 𝕀 (works with amssymb, no extra package).
% Do NOT "fix" this back to \mathbb{1}.
\newcommand{\id}{\mathbb{I}}

% ---- physical units (v3.35): \punit — upright unit with thin space ---------
% Usage: $300\punit{K}$, $80\punit{mV/decade}$, $10^{7}\punit{cm/s}$ — the
% NUMBER AND UNIT LIVE IN ONE MATH ISLAND. Promoted verbatim from a 199-page
% production lakit (advanced-electronic-devices, dozens of uses).
% Two measured traps this design avoids (both render-proven):
%   1. siunitx SCRAMBLES under RTL — \SI{3.5}{eV} renders "eV 3.5", exponents
%      detach ("Hz 10 x 10^9"): its quantity assembly is not bidi-isolated.
%      Do NOT load siunitx in this pipeline.
%   2. A text-mode unit next to a bare-prose number is TWO adjacent LTR runs
%      and reorders ("5 mm" -> "mm5" — the BIDI-2/3 class). \punit is therefore
%      math-only ON PURPOSE: prose misuse fails LOUDLY at compile (\mathrm
%      outside math) instead of scrambling silently.
% Domain shorthands (\MeV, \fm...) stay per-document: \newcommand{\MeV}{\,\text{\en{MeV}}}
\providecommand{\punit}[1]{\,\mathrm{#1}}

% ============================================================================
% Colored boxes — mdframed (NOT tcolorbox, which breaks under bidi).
% ★★ v3.4: innerleftmargin/innerrightmargin give code listings room so their
% frame/line-numbers do not spill out the box border (see listing styles below).
% ============================================================================
% v3.61: everyline/split skips were applied to `tablebox` only (BLD-14, v3.37) —
% the six colored boxes below never got them, so they still split OPEN across a page
% break: the page-N segment gets top+sides and NO bottom rule, the page-N+1 segment
% sides+bottom and NO top rule. Measured on the shipped fleet this is a COVERAGE GAP,
% not a re-report of BLD-14: 0 split-open `tablebox` segments (that fix holds) vs 62
% split-open COLORED boxes across 7 documents. Verified after the fix: both segments
% close, and the frametitle is NOT duplicated on the continuation. BLD-14's caveat
% still applies — a box taller than a full page grinds mdframed's split-trial loop, so
% split such content manually.
% ============================================================================
% v3.72: frametitle={} leads EVERY \newmdenv option list below. mdframed options
% are GROUP-scoped: a box nested inside \begin{exbox}[frametitle={...}] used to
% INHERIT the outer title — measured, the identical title string rendered twice,
% the second riding a title-bar fill inside the INNER box's frame, for all seven
% environments in both directions. Both natural forms were broken (= preamble
% bug, not an authoring rule). The explicit reset makes every box titleless
% unless ITS OWN [frametitle={...}] says otherwise; verified both directions:
% no duplicate on the nested box, an explicit title still renders.
% Geometry is shared with the optional semantic-width table container. Themes
% select their own existing semantic style below; neither path fixes columns.
\mdfdefinestyle{hllpdfsemanticgeometry}{linewidth=1.1pt,
  innertopmargin=6pt,innerbottommargin=7pt,skipabove=9pt,skipbelow=9pt,
  innerleftmargin=9pt,innerrightmargin=9pt,everyline=true,
  splitbottomskip=6pt,splittopskip=12pt,roundcorner=3pt}
\newcommand{\hllpdfsemanticframestyle}{hllpdfsemanticgeometry}
\newmdenv[frametitle={},style=hllpdfsemanticgeometry,linecolor=cBlue,backgroundcolor=cBlueBg,
  frametitlebackgroundcolor=cBlue,frametitlefontcolor=white,
  ]{defbox}   % הגדרה
\newmdenv[frametitle={},style=hllpdfsemanticgeometry,linecolor=cGreen,backgroundcolor=cGreenBg,
  frametitlebackgroundcolor=cGreen,frametitlefontcolor=white,
  ]{thmbox}   % משפט/טענה
\newmdenv[frametitle={},style=hllpdfsemanticgeometry,linecolor=cAmber,backgroundcolor=cAmberBg,
  frametitlebackgroundcolor=cAmber,frametitlefontcolor=white,
  ]{notebox}  % הערה
\newmdenv[frametitle={},style=hllpdfsemanticgeometry,linecolor=cPurple,backgroundcolor=cPurpleBg,
  frametitlebackgroundcolor=cPurple,frametitlefontcolor=white,
  ]{exbox}    % דוגמה
\newmdenv[frametitle={},style=hllpdfsemanticgeometry,linecolor=cRed,backgroundcolor=cRedBg,
  frametitlebackgroundcolor=cRed,frametitlefontcolor=white,
  ]{warnbox}  % טעות נפוצה
\newmdenv[frametitle={},style=hllpdfsemanticgeometry,linecolor=cTeal,backgroundcolor=cTealBg,
  frametitlebackgroundcolor=cTeal,frametitlefontcolor=white,
  ]{keybox}   % רעיון מרכזי

% ---- codebox (v4.5.0): the dedicated container for code -------------------
% An OPTION, not a policy: asmblock unboxed stays canonical for a full listing,
% and a snippet inside any of the six colored boxes still works as before.
% What makes codebox worth its own environment is that it is the ONE framed box
% that keeps its margin line numbers.
%
% v4.4.0 shipped it WITHOUT them, on the strength of v3.37's "no margin knob can
% save them" plus four fresh probes that agreed. That was wrong: every probe had
% moved the code, the box, or both — never the key that carries the NUMBER
% COLUMN. Measured on the page:
%   resetmargins=false   the listing follows THIS box instead of the page,
%                        moving the code from x=75.26 to x=97.54
%   xrightmargin=26pt    under RTL this is the PHYSICAL LEFT indent, and the
%                        only key that carries the \llap'ed digits with it:
%                        x=64.28 -> 90.18, clearing the rule at 65.64 by 23.6pt
%   xleftmargin=26pt     moves the code and LEAVES THE DIGITS on the rule —
%                        the asymmetry that hid the fix for eleven versions
% resetmargins is hoisted out of the four styles (see the \lstset below them) so
% this override is possible at all; every other listing still resolves to true.
% Pinned by tests/codebox_numbers.tex + probe_lstnum.py --inside-rule, whose
% negative control strips xrightmargin and watches the digits fall back.
\definecolor{cSlate}{HTML}{3F4A56}  \definecolor{cSlateBg}{HTML}{F2F4F7}
\newmdenv[frametitle={},leftline=true,rightline=true,topline=true,bottomline=true,
  linecolor=cSlate,backgroundcolor=cSlateBg,linewidth=0.9pt,
  frametitlebackgroundcolor=cSlate,frametitlefontcolor=white,
  innertopmargin=6pt,innerbottommargin=7pt,skipabove=9pt,skipbelow=9pt,
  innerleftmargin=9pt,innerrightmargin=9pt,everyline=true,splitbottomskip=6pt,splittopskip=12pt,
  roundcorner=3pt]{codebox}   % קוד
% The two keys that make in-box numbering work, both measured on the page:
%   resetmargins=false  the listing follows THIS box instead of the page
%   xrightmargin        under RTL this is the PHYSICAL LEFT indent, and it is
%                       the only one that carries the \llap'ed number column
%                       with it. xleftmargin moves the code and leaves the
%                       digits behind on the rule (measured: code 97.54,
%                       digits 64.28 — the v3.37 defect, reproduced); with
%                       xrightmargin=26pt the digits land at x=90.18, a
%                       clearance of 23.6pt from the rule at 65.64.
% This asymmetry is why v3.4's note had the semantics backwards and v3.37
% concluded the digits could not be saved.
% <hllpdf:listing-container-state>
% v4.29: all state is local to the current environment group. A semantic
% callout masks number text, while a nested codebox restores the formatter
% saved before that mask. Save it once across successive semantic wrappers;
% an explicit numbers=none (or an explicitly empty formatter) remains valid.
\makeatletter
\newif\ifhllpdf@incodebox
\newif\ifhllpdf@englishcodebox
\newif\ifhllpdf@lstnumbermasked
\let\hllpdf@savedlstnumber\thelstnumber
\makeatother
% </hllpdf:listing-container-state>
% <hllpdf:codebox-gutter>
\makeatletter
\AtBeginEnvironment{codebox}{%
  \hllpdf@incodeboxtrue
  % Snapshot the OUTER language before asmblock selects English.
  \iflanguage{english}{\hllpdf@englishcodeboxtrue}{\hllpdf@englishcodeboxfalse}%
  \ifhllpdf@lstnumbermasked
    \let\thelstnumber\hllpdf@savedlstnumber
    \hllpdf@lstnumbermaskedfalse
  \fi
}
% listings' AddToHook is GLOBAL. Register once, and gate on local box state;
% adding an unconditional hook at every codebox entry leaks into later listings.
\lst@AddToHook{PreInit}{%
  \ifhllpdf@incodebox
    % Apply every listing override here, inside the LISTING's own group.
    % Setting margins at box entry leaks them into nested semantic callouts.
    \lstset{resetmargins=false,xrightmargin=26pt}%
    % Styles resolve before PreInit. In LTR the physical number gutter follows
    % xleftmargin; in Hebrew retain the measured inherited xrightmargin path.
    \ifhllpdf@englishcodebox
      \lstset{xleftmargin=26pt,xrightmargin=0pt}%
    \fi
    % ONE box, not two: suppress only the frame/background/skips of a listing
    % inside this codebox. The next unboxed listing keeps its own style.
    \lstset{frame=none,backgroundcolor=\color{cSlateBg},aboveskip=0pt,belowskip=0pt}%
  \fi
}
\makeatother
% </hllpdf:codebox-gutter>

% v3.37: line numbers inside a COLORED BOX are sliced by the box border —
% measured: digit span x=63.5-67.3 vs border line at x=65.2 (the border runs
% THROUGH the digits), and the inter-border gap (3.5pt) is narrower than a
% digit (3.8pt), so no margin knob can save them (under resetmargins,
% xleftmargin is inert in-box — measured: code at x=75.26 with 2.2em AND
% 3.0em). A 3-6 line boxed snippet does not need line numbers; blank the
% number TEXT inside every colored box (\thelstnumber is what gets printed;
% the style's numbers=left survives, so nothing else shifts — code position
% byte-frozen at x=75.26, verified). Unboxed listings keep their numbers.
% v3.78: the blanking is coupled to the FRAME, not to the box name —
% clean_theme.tex registers LATER hooks (registration order wins) that
% restore the number text inside the boxes it renders frameless.
% v4.5.0 RETRACTS the "no margin knob can save them" clause above. It holds only
% under resetmargins=true (the shipped default): there xleftmargin is inert
% in-box, as measured. With resetmargins=false, xrightmargin DOES carry the
% \llap'ed digit column clear of the rule — measured; the codebox \lstset that
% does it is the one above, in the codebox block. The loop below therefore
% covers SIX boxes, not "every colored box": codebox opts out and keeps its
% numbers.
% <hllpdf:semantic-listing-mask>
\makeatletter
\@for\HEB@lstbox:=defbox,thmbox,notebox,exbox,warnbox,keybox\do{%
  \expandafter\AtBeginEnvironment\expandafter{\HEB@lstbox}{%
    \ifhllpdf@lstnumbermasked\else
      \let\hllpdf@savedlstnumber\thelstnumber
    \fi
    \hllpdf@lstnumbermaskedtrue
    \def\thelstnumber{}%
    % The nearest semantic container owns its listing presentation. On exit,
    % TeX's group restores the enclosing codebox state and formatter.
    \hllpdf@incodeboxfalse}}
\makeatother
% </hllpdf:semantic-listing-mask>

% ============================================================================
% Unified TABLE presentation: subtle uniform slate frame + auto-numbered
% caption rendered INSIDE the frame. Use:
%   \begin{tablebox}\centering  <tabular>  \tabcap{<description>}\end{tablebox}
% Replaces the buggy \caption* (which emitted a stray "טבלה N : *" line) and
% provides an explicit framed option + numbered caption "טבלה N.M: תיאור".
% Optional intro text may be placed at the TOP, before the tabular, inside the box.
% ============================================================================
\definecolor{cTblLine}{HTML}{B9C4D2}  \definecolor{cTblBg}{HTML}{FAFBFC}
% v3.23: class-aware — under `article` there is no chapter counter (the promised
% report->article switch used to die with "No counter 'chapter' defined").
\makeatletter
\@ifundefined{c@chapter}{%
  % v3.31: in article the counter resets per SECTION, so an unprefixed number
  % repeated ("טבלה 1." in every section — measured on a two-section render).
  % Same conditional idea as the report branch: prefix the section number once
  % a real \section exists; a pre-section table stays "טבלה 1".
  \newcounter{tbl}[section]
  \renewcommand{\thetbl}{\ifnum\value{section}>0 \thesection.\fi\arabic{tbl}}%
}{%
  \newcounter{tbl}[chapter]
  % v3.30: a report-class doc that uses SECTIONS ONLY (chapter stays 0 —
  % the common short-doc pattern) used to print "טבלה 0.1"; the chapter
  % prefix now appears only once a real \chapter exists (render-verified
  % both ways: "טבלה 1." chapterless, "טבלה 1.1" with chapters).
  \renewcommand{\thetbl}{\ifnum\value{chapter}>0 \thechapter.\fi\arabic{tbl}}%
  % v3.64: the SAME leak in \thesection itself, left behind by v3.30. That fix
  % guarded the CAPTION counter and stopped there, so a chapterless report still
  % printed the heading "0.1 סעיף ראשון" and \ref to it printed "0.2" — measured:
  %   p1 y=65.2  '…0.1'   p1 y=144.9 '…0.2'   p1 y=183.4 '…0.2.1'
  % report.cls hardcodes \thesection as \thechapter.\@arabic\c@section, and the
  % leading 0 is the chapter counter that no \chapter ever incremented. Fixing
  % \thesection fixes the whole chain — \thesubsection and \thesubsubsection are
  % defined in terms of it. \theequation and \thetbl/\thefig were already guarded,
  % which is exactly why this survived: every NUMBER around the heading looked
  % right, so the heading's own 0 read as a style choice rather than a bug.
  % \appendix is safe: report.cls resets chapter to 0 and redefines \thechapter to
  % \@Alph\c@chapter, so an appendix \chapter makes the counter 1 and the guard
  % emits "A.1" unchanged. Render-verified in all four modes (chapterless,
  % chaptered, appendix+chapter, appendix+section-only).
  \renewcommand{\thesection}{\ifnum\value{chapter}>0 \thechapter.\fi\arabic{section}}%
}
\makeatother
% v3.37: everyline=true — if the box DOES split across pages, every segment
% gets a CLOSED frame (measured before: the split ran open-ended into the page
% edge with no bottom border, and the caption could strand alone at the next
% page top — "טבלה A.2" shipped that way). The split skips give the segments
% breathing room. In practice, with the caption GLUE below, every previously-
% split table in the corpus now moves WHOLE instead (the white-gap tradeoff,
% documented benign); everyline stays as the safety net for a genuinely
% page-tall table. NB: a table taller than ONE FULL PAGE can grind mdframed's
% split-trial loop (reproduced: minimal doc, minutes of 1pt-step Overfull
% probes) — break such a monster into two tables.
\mdfdefinestyle{hllpdftablegeometry}{leftline=true,rightline=true,topline=true,bottomline=true,
  linewidth=0.9pt,
  innertopmargin=9pt,innerbottommargin=8pt,skipabove=11pt,skipbelow=11pt,
  innerleftmargin=11pt,innerrightmargin=11pt,roundcorner=4pt}
\newcommand{\hllpdftableframestyle}{hllpdftablegeometry}
\newmdenv[frametitle={},style=hllpdftablegeometry,
  linecolor=cTblLine,backgroundcolor=cTblBg,
  % v4.2.0 (adopted): tablebox is ATOMIC. Short tableboxes near page bottoms
  % repeatedly entered mdframed's split trial, emitting benign Overfull-vbox
  % bursts and occasionally separating the intro prose from the tabular.
  % nobreak=true moves the complete short table to the next page instead. The
  % output-active Overfull-vbox gate still catches a box taller than a page;
  % everyline/split skips remain as defensive styling should that path ever
  % re-arm, but the canonical path never publishes a split tablebox.
  nobreak=true,
  everyline=true,splitbottomskip=6pt,splittopskip=12pt,
  % v3.79 (BLD-37 table side, correction): startinnercode/endinnercode run
  % INSIDE mdframed's body vbox (mdframed.sty:848 / :856+863) -- the second one
  % after \par\unskip, i.e. after line breaking, which is the only moment the
  % \tabcap anchor's true offset from the top of the box EXISTS. See
  % \lakintabboxclose below. These two keys are the ONLY tablebox change.
  startinnercode={\lakintabboxopen},endinnercode={\lakintabboxclose}]{tablebox}
% v3.37: \nopagebreak GLUE — \vsplit honors penalties, so the caption can no
% longer strand alone at the top of the next page (it takes its rows along, or
% the whole box moves). Render-verified on the real shipped orphan.
% v3.59: \NR@gettitle is what nameref injects into the stdlib title-bearing
% commands (\@sect, \@chapter, \caption) to fill \@currentlabelname. \figcap and
% \tabcap are INDEPENDENT macros — they never called it, so \@currentlabelname kept
% whatever the last heading left and \nameref{fig:x} silently printed the enclosing
% SECTION's title instead of the caption (proven in the .aux: field 3 of \newlabel
% was the heading text while field 1 and the anchor `fig.1.1` were correct — so this
% is a THIRD, distinct head next to BLD-13's counter fix and BIDI-38's autoref-name
% fix, neither of which writes that field). \@ifundefined keeps it inert if
% hyperref/nameref is ever dropped. Guarded helper because \figcap is defined
% after \makeatother.
% NB both helpers below need \makeatletter — this part of the preamble is OUTSIDE
% any letter-catcode region, so an unwrapped \@currentlabel parses as LaTeX's \@
% and dies with "You can't use \spacefactor in vertical mode" (measured).
\makeatletter
\newcommand{\lakinrtitle}[1]{\@ifundefined{NR@gettitle}{}{\NR@gettitle{#1}}}
% v3.59: \hitem{<letter>}{<key>} — a REFERENCEABLE Hebrew-lettered list item.
%   \begin{enumerate} \hitem{א}{it:a} הסעיף הראשון. \end{enumerate}  ->  \ref{it:a} = (א)
% Why this exists: automatic Hebrew list labels fail silently, so SKILL.md mandates
% an EXPLICIT \item[(\textbf{א})] — but LaTeX puts \refstepcounter INSIDE
% \if@noitemarg (latex.ltx:12964), and \item[...] never sets that flag. So an
% explicit item steps no counter and never touches \@currentlabel: a \label there
% binds to the AMBIENT one and \ref prints the CHAPTER/SECTION number where the
% item letter belongs (proven in the .aux: {chapter.1}). Same doctrine as \figcap —
% own the label state before \label. \phantomsection matters: without it the
% hyperlink still lands on the chapter head (measured); with it the anchor is
% section*.N and the bookmark outline stays clean.
\newcommand{\hitem}[2]{\item[(\textbf{#1})]%
  \def\@currentlabel{(#1)}\def\@currentlabelname{#1}\phantomsection\label{#2}}
% NB (v3.66): \hitem is \ref-only. \autoref{key} calls it a SECTION ("סעיף (א)")
% because \phantomsection brands the anchor section*.N, and \nameref{key} prints
% the bare letter. For \autoref-able items prefer \hebenumlabels below.
%
% v3.66: \hebenumlabels — OPT-IN automatic Hebrew enumerate letters, done right.
%   {\hebenumlabels\begin{enumerate}\item\label{q:a} ...\end{enumerate}}
%   -> labels (א) (ב) (ג); \ref{q:a} = א; anchor Item.N (real, per-item).
% This retires the belief that made \item[(\textbf{א})] + \hitem necessary.
% "Automatic Hebrew labels fail silently" was TRUE but unexplained, and the
% mechanism is now bisected (plain article vs +enumitem, measured): ENUMITEM
% PRE-EXPANDS the label expression in an \edef BEFORE \refstepcounter, with the
% counter still at 0. Tokens it shields by name (\arabic,\alph,\roman) survive;
% a plain \newcommand \ifcase wrapper evaluates at 0 -> the EMPTY zero-branch ->
% the label prints as "()" with the letter GONE (and \the\c@enumi prints a
% literal 0, \number\value dies fatally pointing at \item). The cure is one
% property, not a doctrine: the wrapper must be ROBUST (survives the \edef
% unexpanded), while \theenumi stays EXPANDABLE (so \refstepcounter writes the
% real letter to the .aux). Verified: labels (א)(ב)(ג), \ref -> א/ב/ג each item
% its own value, .aux = {א}{...}{Item.1} — versus the explicit-\item doctrine
% where BIDI-50/\hitem exist purely to patch the counter that never stepped.
\def\@hebalph#1{\ifcase#1\or א\or ב\or ג\or ד\or ה\or ו\or ז\or ח\or ט\or י\or
  יא\or יב\or יג\or יד\or טו\or טז\or יז\or יח\or יט\or כ\else\@ctrerr\fi}
\DeclareRobustCommand*\hebalph[1]{\expandafter\@hebalph\csname c@#1\endcsname}
\newcommand{\hebenumlabels}{%
  \renewcommand{\theenumi}{\@hebalph\c@enumi}%
  \renewcommand{\labelenumi}{(\textbf{\hebalph{enumi}})}}
\makeatother
% v3.63: \mbox the LABEL. \thetbl ("7.1") and a caption description that STARTS
% WITH A DIGIT ("32 הרגיסטרים…") are two European-Number runs separated only by
% neutrals (the label's own "." and \enspace). Under the RTL paragraph base the UBA
% resolves those neutrals to the RTL level, so both numbers sit in ONE reversed span
% and EXCHANGE VISUAL POSITIONS. Measured on the shipped PDF: a table whose real
% number is 7.1 printed "טבלה 32" — a different, EXISTING table number — while \ref
% correctly printed 7.1, so the caption silently contradicted its own cross-reference,
% and the description lost its leading count. The source contains no adjacency at all;
% it is manufactured by this macro, which is why every source linter was blind.
% \mbox makes the label one atomic box so nothing can reorder across it. Verified
% zero-regression: an ordinary Hebrew caption keeps identical x-geometry.
% One paragraph style covers label, number, separator and description.
% Explicit project/template overrides remain possible; the label hook is kept
% for compatibility, with an identity default. Counters/anchors stay unchanged.
\newcommand{\lakitcaptionstyle}{\normalfont\small\color{black}\hypersetup{allcolors=black}}
\newcommand{\lakitcaptionlabelstyle}[1]{#1}
\newcommand{\lakitcaptionpunctuation}{:}
% Native floats, captionof and longtable use the same presentation policy.
% subcaption is optional and inherits the same font rather than a second size.
\usepackage{caption}
\DeclareCaptionFont{lakit}{\lakitcaptionstyle}
\DeclareCaptionLabelFormat{lakit}{\lakitcaptionlabelstyle{#1~\en{#2}}}
\DeclareCaptionLabelSeparator{lakit}{\lakitcaptionpunctuation\enspace}
\captionsetup{font=lakit,labelfont={},textfont={},labelformat=lakit,
  labelsep=lakit,justification=centering}
\makeatletter
\AtBeginDocument{\@ifpackageloaded{subcaption}{%
  \captionsetup[sub]{font=lakit,labelfont={},textfont={},labelsep=lakit}}{}}
\makeatother
% v3.64: OPT-IN list-of-figures/tables entry via an optional short caption.
% \listoffigures and \listoftables are inherited from report.cls and LOOK available,
% but \figcap/\tabcap step the custom \c@fig/\c@tbl counters and never call
% \addcontentsline — so both printed their heading and NOTHING ELSE. Measured: a doc
% with one \figcap and one \tabcap rendered page 1 as 'רשימת האיורים\n1\n', and the
% .lof held only \@chapter's \addvspace, zero \contentsline entries. An empty list
% under a correct Hebrew heading reads as "this document has no figures".
%
% WHY OPTIONAL, AND NOT AUTOMATIC. Writing #1 unconditionally was implemented first
% and REJECTED: a caption containing \codeline compiles CLEANLY today (verified on
% the pre-v3.64 preamble: 0 errors, PDF produced) but becomes FATAL once it is also
% written to the .lof — \codeline is \hbox to\linewidth{...}, and a full-width box in
% a list entry dies with "\pdfextension endlink cannot be used in vertical mode", no
% output PDF. Gaining a feature nobody uses (0 of 13 collections call \listoffigures)
% at the price of breaking a construct that works is a bad trade. Stock \caption has
% exactly this limitation and exactly this escape hatch, so the convention is familiar.
% \en{}, \code{}, \h{} and inline math in the short text are all safe — verified, they
% survive the write and typeset correctly in the list (RTL geometry confirmed: number
% rightmost at x=509.9, title left of it, dots and folio at the far left).
% With no optional argument the macro is byte-identical to v3.63 — zero risk to every
% existing document, and \listoffigures lists exactly the captions that opted in.
% BLD-37, TABLE side (the \figcap v3.74 idea, adapted): clicking a \ref to a
% table used to land MID-table — hyperref's automatic anchor inside
% \refstepcounter was emitted on the still-open horizontal line whose TAIL is
% the tabular (an hbox whose baseline is the table's vertical CENTER), so the
% jump scrolled the top half of the table above the viewport. The v3.74 entry
% deferred this side claiming "tablebox splits — no single box to measure";
% that deferral dissolves: the measurement target is the TABULAR itself at
% \tabcap ENTRY, while it is still the tail of the open list — the same
% situation \lakin@grabfig handles at an island's end (and the fleet has 0
% split tableboxes). So \lakin@grabtab runs as the FIRST token, before any
% typesetting can push the box away, and the anchor is re-emitted raised above
% the measured tabular. \lakintabraise is a SEPARATE register from
% \lakinfigraise (defined with it, below the LTR-island block): a captionless
% island intentionally leaves the fig register set, and a table must never
% inherit that stale figure height.
%
% WHY the anchor stays AT THE ENTRY SPOT and is NOT deferred into the caption
% paragraph the way \figcap defers it: at \par TeX DROPS a final glue item,
% and the automatic anchor's smash box sitting BEHIND the interword glue kept
% that glue alive — so every centered tabular has always sat half a space off
% true center. Deferring the anchor removed that box and was measured as 982
% differing glyph records: EVERY tabular line shifted 1.30pt (2.61pt for
% pattern 7's two minipages) while captions stayed put. Emitting the raised
% anchor at the old spot — through hyperref's own \Hy@raisedlink, exactly the
% macro the automatic path uses — keeps the line's node structure, and so the
% A/B glyph diff at 0, BY CONSTRUCTION. \Hy@raisedlink reads the raise from
% \HyperRaiseLinkDefault (default \baselineskip = the old mid-table dest); the
% group-local \edef freezes measured-height+42pt into it, and \smash inside
% \Hy@raisedlink hides any raise from the line's extent. \HyCnt@ProvideTheH-
% Counter keeps the dest NAME byte-identical to the automatic one (it defines
% \theHtbl, which \hyper@makecurrent prefers — without it the first measured
% caption was named tbl.1 where the automatic path said tbl.0.1, measured).
% vmode entry, or anything but a box at the tail (no tabular to measure),
% falls back to the plain \refstepcounter — byte-identical to the old path.
\makeatletter
\newcommand{\tabcap}[2][]{\lakin@grabtab
  % v3.79: arm the box-close correction. Set unconditionally and in BOTH
  % branches below's scope, so the plain-\refstepcounter fallback is corrected
  % too; \lakintabboxopen clears it at every \begin{tablebox}, so a \tabcap
  % outside a tablebox can never leave it armed for a later box.
  \global\lakintabcapseentrue
  \ifdim\lakintabraise>0pt
    \H@refstepcounter{tbl}\HyCnt@ProvideTheHCounter{tbl}\hyper@makecurrent{tbl}%
    \begingroup
      % A frameless table targets its measured top. The frame's historical
      % cushion belongs only to tablebox and is corrected at box close.
      \edef\HyperRaiseLinkDefault{\the\lakintabraise}%
      \iflakintabframed
        \edef\HyperRaiseLinkDefault{\the\dimexpr\lakintabraise+42pt\relax}%
      \fi
      \Hy@raisedlink{\hyper@anchorstart{\@currentHref}\hyper@anchorend}%
    \endgroup
    \global\lakintabraise 0pt\relax
  \else
    \refstepcounter{tbl}%
  \fi
  \lakinrtitle{#2}\par\nopagebreak\addvspace{6pt}%
  \def\lakin@short{#1}\ifx\lakin@short\@empty\else
    \addcontentsline{lot}{table}{\protect\numberline{\thetbl}{\protect\ignorespaces #1}}\fi
  {\centering\lakitcaptionstyle\mbox{\lakitcaptionlabelstyle{\tablename~\en{\thetbl}\lakitcaptionpunctuation}}\enspace #2\par}}
\makeatother

% ============================================================================
% Unified FIGURE caption — the \tabcap twin for TikZ/plot figures (v3.34.0).
% A figure (tikzpic) sits inline in the prose; put \figcap BELOW it for a
% numbered, \ref-able caption "איור N.M: תיאור" in the same style as \tabcap:
%   \begin{tikzpic} ... \end{tikzpic}
%   \figcap{\label{fig:foo}description}     % \label FIRST, inside the argument
% WHY a macro, not a hand-written gray line: a manual
%   {\centering\small\textcolor{cGray}{איור: ...}} with a \label inside has NO
% \refstepcounter, so the \label binds to the AMBIENT \@currentlabel (the last
% equation) — \ref then prints an equation number AS the figure number
% (measured on a real page: "איור 48" where 48 is equation 48). \figcap steps a
% real \c@fig counter first, so \label is correct, numbering is consistent, and
% refs are stable. Class-aware like \thetbl (chapter- or section-prefixed).
% ============================================================================
\makeatletter
\@ifundefined{c@chapter}{%
  \newcounter{fig}[section]
  \renewcommand{\thefig}{\ifnum\value{section}>0 \thesection.\fi\arabic{fig}}%
}{%
  \newcounter{fig}[chapter]
  \renewcommand{\thefig}{\ifnum\value{chapter}>0 \thechapter.\fi\arabic{fig}}%
}
\makeatother
\makeatletter
\newcommand{\figcap}[2][]{%
  \ifdim\lakinfigraise>0pt
    % v3.74: step WITHOUT hyperref's automatic anchor, then emit the anchor
    % raised above the measured figure. \H@refstepcounter is hyperref's saved
    % pre-hook \refstepcounter (sets the counter AND \@currentlabel — \label
    % in the caption arg binds exactly as before); \hyper@makecurrent supplies
    % \@currentHref, so the dest NAME is identical to the old automatic one —
    % \ref/\autoref/bookmarks are unchanged, only the dest'S POSITION moves.
    \H@refstepcounter{fig}\hyper@makecurrent{fig}%
    % The anchor is emitted INSIDE the caption paragraph below, as a zero-size
    % raisebox — NOT here in vertical mode, where a \raisebox starts an empty
    % paragraph and adds a blank line (measured: every glyph after the caption
    % shifted +18.8pt). Inside an existing line a [0pt][0pt] box is inert.
    \edef\lakin@anch{\@currentHref}%
    % \edef FREEZES the measured height into the pending macro: the deferred
    % \raisebox is typeset inside the caption paragraph AFTER the register is
    % reset to the 0pt sentinel below — a plain \def read 0 at use time
    % (measured: LAKINRAISE=0.0, dest moved 14pt instead of above the figure).
    \edef\lakin@pend{\noexpand\raisebox{\the\dimexpr\lakinfigraise+42pt\relax}[0pt][0pt]{%
      \noexpand\hyper@anchorstart{\lakin@anch}\noexpand\hyper@anchorend}}%
    \global\lakinfigraise 0pt\relax
  \else
    \refstepcounter{fig}\let\lakin@pend\relax
  \fi
  \lakinrtitle{#2}\par\nopagebreak\addvspace{4pt}%
  \def\lakin@short{#1}\ifx\lakin@short\@empty\else
    \addcontentsline{lof}{figure}{\protect\numberline{\thefig}{\protect\ignorespaces #1}}\fi
  {\centering\lakin@pend\lakitcaptionstyle\mbox{\lakitcaptionlabelstyle{\figurename~\en{\thefig}\lakitcaptionpunctuation}}\enspace #2\par}}
\makeatother

% ============================================================================
% LTR islands.  EVERY tikzpicture / quantikz MUST live inside one, or the whole
% figure's coordinates and Latin labels flip. (Hebrew node text still uses \h{}.)
% ============================================================================
% v3.74 (BLD-37): the \figcap hyperlink DESTINATION used to land at the caption
% line — BELOW the figure — so clicking a figure reference scrolled the figure
% off-screen (measured fleet-wide: 97/97 fig dests in one collection sat 19-54pt
% under their figure's ink; a 347pt figure was fully hidden after the jump).
% hypcap-style fix in two parts: (1) at the island's END, \lakin@grabfig grabs
% the just-typeset centered line with \lastbox, measures it, and puts it back
% with \nointerlineskip — the interline glue above the line was already
% contributed and is NOT re-added, so the page geometry is unchanged to the
% glyph (verified); (2) \figcap then steps the counter WITHOUT hyperref's
% automatic anchor (\H@refstepcounter = the pre-hyperref \refstepcounter) and
% emits the anchor itself, raised by the measured height — the dest lands ABOVE
% the figure top. A captionless island leaves numbering untouched (nothing is
% stepped at the island; the grab flag is cleared at every island begin, so a
% stale height from a captionless figure cannot leak into a later caption).
\newsavebox{\lakinfigbox}
\newdimen\lakinfigraise  % real \dimen: tex.setdimen cannot target a \newlength SKIP register (measured lua error)
% BLD-37 table side: a SEPARATE register for \tabcap. The two must not share —
% a captionless island INTENTIONALLY leaves \lakinfigraise set (the grab runs
% at every island END; only the NEXT island's begin clears it), so a shared
% register would hand a \tabcap after a captionless figure that stale FIGURE
% height. \lakintabraise is written only by \lakin@grabtab at \tabcap entry
% and consumed/reset inside \tabcap itself, so it can never go stale.
\newdimen\lakintabraise
\makeatletter
% v3.74b: measure WITHOUT touching the list. The first draft grabbed the box
% with \lastbox and put it back — and \lastbox DISCARDS the node's baseline
% shift, so a vertically-anchored quantikz circuit rose 2.57pt and everything
% after it sank 5.32pt (measured on the A/B glyph diff). LuaTeX lets us READ
% the tail node's dimensions without removing anything: geometry is identical
% by construction. \lakinfigraise=0pt is the "no figure measured" sentinel.
\def\lakin@grabfig{%
  \ifhmode
    \directlua{
      local n = tex.nest[tex.nest.ptr].tail
      local GLUE, PEN = node.id('glue'), node.id('penalty')
      while n and (n.id == GLUE or n.id == PEN) do n = n.prev end
      local HL, VL = node.id('hlist'), node.id('vlist')
      if n and (n.id == HL or n.id == VL) then
        tex.setdimen('global', 'lakinfigraise',
                     (n.height or 0) + (n.depth or 0))
      else
        tex.setdimen('global', 'lakinfigraise', 0)
      end}%
  \else
    \global\lakinfigraise 0pt\relax
  \fi}
% BLD-37 table side: \lakin@grabfig's twin, NOT a shared helper — the Lua body
% differs in exactly one place, the skip set. Between \end{tabular} and \tabcap
% an otherlanguage wrapper (patterns 3/5 of example_tables) leaves a \write
% whatsit and a dir node ON TOP of the interword glue (measured node dump:
% "whatsit/write <- dir <- glue(2.62) <- hlist"), so the table reader must
% also step over whatsit/dir markers — all zero-size — to reach the tabular.
% \lakin@grabfig deliberately does NOT get the wider skip set: its glue/penalty
% body is the one the v3.74 A/B proof froze, and widening it would newly
% measure (and so newly raise) anchors in island shapes that today take the
% plain path — an unproven behavior change on the FIG side. Anything that is
% not a box behind those markers still hits the 0pt sentinel and \tabcap falls
% back to the plain \refstepcounter.
% SECOND deliberate difference: this stores the box's HEIGHT, not height+depth.
% \figcap's anchor is emitted in the caption paragraph BELOW the figure, so it
% must climb the full extent; \tabcap's anchor sits ON the tabular's own line,
% whose baseline is the vcentered tabular's MIDPOINT — height alone reaches
% the top edge, and adding the depth was measured to overshoot by half the
% table (dests 80-135pt above their frame, INSIDE the preceding table's box).
\def\lakin@grabtab{%
  \ifhmode
    \directlua{
      local n = tex.nest[tex.nest.ptr].tail
      local GLUE, PEN = node.id('glue'), node.id('penalty')
      local WHAT, DIR = node.id('whatsit'), node.id('dir')
      while n and (n.id == GLUE or n.id == PEN
                   or n.id == WHAT or n.id == DIR) do n = n.prev end
      local HL, VL = node.id('hlist'), node.id('vlist')
      if n and (n.id == HL or n.id == VL) then
        tex.setdimen('global', 'lakintabraise', n.height or 0)
      else
        tex.setdimen('global', 'lakintabraise', 0)
      end}%
  \else
    \global\lakintabraise 0pt\relax
  \fi}
\makeatother
% ============================================================================
% v3.79 (BLD-37 table side, CORRECTION): put the anchor at the TOP OF THE BOX,
% instead of guessing the distance to it.
% ============================================================================
% v3.77 raised the \tabcap anchor by (height of the box at the tail of the open
% horizontal list) + 42pt. That is a MEASUREMENT taken BEFORE line breaking, and
% it is short by every term it cannot see. Fleet-measured on the 12 built
% collections (299 tbl.* dests, ownership established through each table's own
% CAPTION, never through "which frame encloses the dest"): 3 dests landed INSIDE
% their own tablebox, below the frame top —
%   computer-architecture       tbl.14.1  -46.4pt   three stacked tabulars
%   computer-architecture       tbl.30.5  -13.5pt   tabular 2.85pt over \linewidth
%   advanced-electronic-devices tbl.B.1   -59.2pt   tabular 2.56pt over \linewidth
% Two missing terms, both reproduced in a minimal document (probe P2/P3):
%   (1) VERTICAL LIST ALREADY CONTRIBUTED inside the box — a second and third
%       tabular, or a heading above the table. The tail read sees only the LAST
%       box, so 77.7pt of stacked tabulars is invisible to the raise.
%   (2) WHICH LINE THE ANCHOR LANDS ON. When the tabular alone overfills
%       \linewidth (459.6121pt in a tablebox), the interword glue after
%       \end{tabular} becomes the chosen breakpoint and \Hy@raisedlink is
%       emitted on the NEXT line — one (depth + lineskip) lower. That term does
%       not EXIST at \tabcap time: line breaking has not happened yet. No
%       measurement taken there can ever recover it. (Proof: probe P3, the same
%       table 2.84912pt too wide, fails at -13.5pt; probe P4, the identical
%       table with \tabcolsep 2pt so it fits, is correct at +32.0pt.)
% So the raise is not computed at \tabcap at all. It is CORRECTED at
% \end{tablebox}, from inside mdframed's own body vbox (endinnercode runs after
% \par\unskip, i.e. after line breaking, while tex.nest[tex.nest.ptr] is still
% that vbox's vertical list). There the anchor's true offset from the top of the
% box content is not estimated but WALKED, and the \Hy@raisedlink shift is set
% so the dest lands 42pt above it — which is 42pt - innertopmargin(9pt) -
% linewidth(0.9pt) = 32.0pt above the FRAME top, the exact position the 295
% already-correct dests have today. That is why this is a correction and not a
% re-placement: the good sites do not move.
%
% WHY THE GLYPH DIFF IS 0 BY CONSTRUCTION — the only thing written is the
% `shift` field of the hlist that \Hy@raisedlink already built, and the guard
% below refuses to write it unless that box is 0x0x0. A zero-size box carries no
% dimension into its parent's packing, the packing already happened, and shift
% is not an input to it. Nothing else in the list is touched: no node is added,
% removed, resized or reordered. Measured, not assumed: assets/example_tables.tex
% and the whole computer-architecture collection A/B at 0 differing pymupdf
% rawdict glyph records.
%
% WHY NOT the alternatives (all prototyped and measured before this one):
%   * accumulate the vertical list at \tabcap time — fixes term (1), CANNOT fix
%     term (2); 2 of the 3 real failures survive. A fix that leaves the majority
%     of the defect standing is the false all-clear this repo exists to avoid.
%   * \zref-savepos two-pass — correct on all shapes, but it makes the dest
%     POSITION aux-derived (DEFECTS BLD-1/2), a single-pass compile silently
%     ships an unraised anchor, and the bare y-difference has no page guard.
%   * emit the anchor at \begin{tablebox} and let \tabcap adopt it — needs the
%     tbl counter stepped at box OPEN to keep the dest NAME byte-identical; a
%     captionless tablebox then either shifts every later table number or leaves
%     an orphan dest that the next table duplicates (a HARD build gate). This
%     correction touches neither the counter nor the name: \tabcap's naming path
%     is v3.77's, unchanged, so resolve_names() is byte-identical (verified).
%   * suppressing the line break (a penalty before the glue) — that MOVES the
%     tabular: the break currently discards the interword glue, so forbidding it
%     re-centers the two overfull tables. Reflow, rejected.
%
% SCOPE. The correction is armed only by \tabcap ITSELF (\lakintabcapseen),
% cleared at every box open, so a captionless tablebox is a no-op and a \tabcap
% hosted OUTSIDE a tablebox (the fleet has exactly one, computer-architecture
% tbl.14.2 in an exbox) keeps v3.77's measured raise. Both fallbacks are
% therefore never worse than today; and because hyperref's AUTOMATIC anchor uses
% the same \Hy@raisedlink shape, the plain-\refstepcounter path (vmode entry, or
% prose between \end{tabular} and \tabcap — probe P6, -51.5pt) is corrected too.
% KNOWN LIMIT — a tablebox that SPLITS across pages. The offset is measured from
% the top of the UNSPLIT body, so a dest in segment 2 gets raised past its own
% segment. MEASURED on a forced split (probe: 40 filler lines + a table): the
% dest went from y=585.0 on page 2 (inside segment 2, at the caption — the old
% BLD-37 symptom) to y=-113.4, i.e. ABOVE the page top. Viewers clamp that to the
% page top, which happens to show segment 2 whole, but an out-of-page coordinate
% is not something to ship on purpose. Hence the \textheight guard in
% lakintab.fix: a body TALLER THAN THE TEXT BLOCK cannot fit on any page, so it
% must split — the correction is skipped and v3.77's behaviour stands unchanged
% (verified: the probe returns to y=585.0). The guard is deliberately EXACT, not
% a page-fitting heuristic: it is sound (a taller-than-page body always splits)
% but PARTIAL (a shorter body that starts low on a page could still split). That
% residue is unreachable in practice — v3.37's everyline + the caption's
% \nopagebreak glue move such boxes WHOLE instead (probe N7: a 22-row table
% pushed down by 430pt moved whole, 0 segments), and the fleet has 0 split
% tableboxes, re-counted after this change. Reimplementing mdframed's own
% page-fitting arithmetic to close the residue would be a heuristic with its own
% failure mode, which is the trade this repo has repeatedly refused.
\newif\iflakintabcapseen
\newif\iflakintabframed
% The Lua lives in ONE \directlua at load time, and is written with NO
% backslash, %, # or ~ : \directlua expands its argument like \edef, so an
% active ~ would expand to \nobreakspace and a % would eat the rest of the line
% (the second is how the first draft of this block died). "not (a == b)" stands
% in for Lua's inequality operator for the same reason.
\directlua{
  lakintab = { }
  lakintab.HL = node.id("hlist")
  lakintab.VL = node.id("vlist")
  lakintab.RU = node.id("rule")
  lakintab.GL = node.id("glue")
  lakintab.KE = node.id("kern")
  lakintab.WH = node.id("whatsit")
  lakintab.TM = node.id("temp")
  lakintab.WD = node.subtype("pdf_dest")
  function lakintab.d(v)
    if v == nil then return 0 end
    if v < -1000000000 then return 0 end
    return v
  end
  function lakintab.scan(list, y, horiz, parent, target)
    local x = list
    while x do
      local i = x.id
      if i == lakintab.HL or i == lakintab.VL then
        if horiz then
          local b = y + x.shift
          if i == lakintab.HL then
            lakintab.scan(x.head, b, true, x, target)
          else
            lakintab.scan(x.head, b - lakintab.d(x.height), false, x, target)
          end
        else
          if i == lakintab.HL then
            lakintab.scan(x.head, y + lakintab.d(x.height), true, x, target)
          else
            lakintab.scan(x.head, y, false, x, target)
          end
          y = y + lakintab.d(x.height) + lakintab.d(x.depth)
        end
      elseif i == lakintab.WH then
        if x.subtype == lakintab.WD then
          local id = x.dest_id
          if type(id) == "string" and string.sub(id, 1, 4) == "tbl." then
            if horiz and parent and parent.width == 0
               and parent.height == 0 and parent.depth == 0 then
              parent.shift = parent.shift + (target - y)
            end
          end
        end
      elseif not horiz then
        if i == lakintab.RU then
          y = y + lakintab.d(x.height) + lakintab.d(x.depth)
        elseif i == lakintab.GL or i == lakintab.KE then
          y = y + lakintab.d(x.width)
        end
      end
      x = x.next
    end
  end
  function lakintab.fix(target, maxh)
    local nst = tex.nest[tex.nest.ptr]
    if nst == nil then return end
    local h = nst.head
    if h == nil then return end
    if h.id == lakintab.TM then h = h.next end
    local tot = 0
    local x = h
    while x do
      local i = x.id
      if i == lakintab.HL or i == lakintab.VL or i == lakintab.RU then
        tot = tot + lakintab.d(x.height) + lakintab.d(x.depth)
      elseif i == lakintab.GL or i == lakintab.KE then
        tot = tot + lakintab.d(x.width)
      end
      x = x.next
    end
    if tot > maxh then return end
    lakintab.scan(h, 0, false, nil, target)
  end
}
% @-FREE NAMES ON PURPOSE: these two are named in the `tablebox` key list at
% \newmdenv above, which is tokenised OUTSIDE any \makeatletter region — a
% \lakin@... there would tokenise as \lakin followed by the letters @tabopen.
\newcommand{\lakintabboxopen}{\lakintabframedtrue\global\lakintabcapseenfalse}
\newcommand{\lakintabboxclose}{%
  \iflakintabcapseen
    \directlua{lakintab.fix(-\number\dimexpr\lakintabcushion\relax,
                            \number\dimexpr\textheight\relax)}%
    \global\lakintabcapseenfalse
  \fi\lakintabframedfalse}
% The cushion above the box's CONTENT top. 42pt puts the dest 32.0pt above the
% frame top (42pt - innertopmargin 9pt - linewidth 0.9pt), i.e. exactly where
% the 295 correct dests already are — the correction moves only the broken ones.
\newcommand{\lakintabcushion}{42pt}
\newenvironment{qcirc}{\par\medskip\global\lakinfigraise 0pt\relax\begin{otherlanguage}{english}\begin{center}}
  {\csname lakin@grabfig\endcsname\end{center}\end{otherlanguage}\par\medskip}
\newenvironment{tikzpic}{\par\medskip\global\lakinfigraise 0pt\relax\begin{otherlanguage}{english}\begin{center}}
  {\csname lakin@grabfig\endcsname\end{center}\end{otherlanguage}\par\medskip}

% ============================================================================
% Code listings.  ★★ v3.4 THE BOX-OVERFLOW FIX:
%   resetmargins=true   → the listing IGNORES the enclosing environment's
%                         indentation and measures from the PAGE. v3.4 stated
%                         this backwards ("measured from the ENCLOSING box"),
%                         and v3.37 then blanked in-box line numbers because
%                         the page-anchored listing put its hanging digits on
%                         top of the box rule. Measured v4.5.0: with the key
%                         FALSE the listing follows the box and the code moves
%                         from x=75.26 to x=97.74, clearing the rule at 72.07.
%                         The key now lives in one \lstset below instead of in
%                         each style, so a box can override it — codebox does.
%   framexleftmargin=0  → the frame (left bar) does not extend left of the code.
%   xleftmargin keeps a small indent so the line numbers have room INSIDE.
% asmblock wraps the listing in an LTR island (code is LTR).
% ============================================================================
\lstdefinestyle{asm}{
  basicstyle=\ttfamily\small,
  backgroundcolor=\color{cBlueBg!55},
  frame=leftline,framerule=2.5pt,rulecolor=\color{cBlue},
  numbers=left,numberstyle=\tiny\color{cGray},numbersep=8pt,
  xleftmargin=2.2em,framexleftmargin=0pt,framexrightmargin=0pt,
  breaklines=true,columns=fullflexible,keepspaces=true,
  showstringspaces=false,aboveskip=8pt,belowskip=8pt,
  % v3.36: without a comment DELIMITER the commentstyle below never fires —
  % listings has no `language` here, so it does not know '#' opens a comment,
  % and every asm listing rendered flat black (measured: a whole-page pixel
  % sweep found ZERO cGreen glyphs across 135 shipped listings). morecomment
  % gives it the delimiter; color-only change, glyph widths untouched.
  morecomment={[l]{\#}},
  commentstyle=\color{cGreen},
}
\lstdefinestyle{cstyle}{
  language=C,basicstyle=\ttfamily\small,
  backgroundcolor=\color{cGreenBg!55},
  frame=leftline,framerule=2.5pt,rulecolor=\color{cGreen},
  numbers=left,numberstyle=\tiny\color{cGray},numbersep=8pt,
  xleftmargin=2.2em,framexleftmargin=0pt,framexrightmargin=0pt,
  breaklines=true,columns=fullflexible,keepspaces=true,
  showstringspaces=false,keywordstyle=\color{cBlue}\bfseries,
  commentstyle=\color{cGray}\itshape,aboveskip=8pt,belowskip=8pt,
}
% v3.36: pystyle/htmlstyle — colored twins of cstyle for the Python/HTML
% listings that used to repeat a ~200-char options blob with NO color
% (keywords rendered bold-black, comments italic-black — measured; C alone
% was colored). Geometry is copied VERBATIM from the shipped blob (frame,
% margins, skips) so swapping a call site to style=pystyle moves NOTHING —
% only keywordstyle/commentstyle/stringstyle are added (bold stays bold,
% italic stays italic, mono metrics identical => zero reflow, proven by
% per-page word-bbox diff). breaklines/breakatwhitespace stay per-site.
% v3.78: numbers=left joins both styles, with the exact numberstyle/numbersep
% keys asm/cstyle carry — fleet-wide consistency (author decision: all four
% styles numbered; "asm/cstyle numbered, py/html not" was an accident of the
% v3.36 geometry freeze, not a design). The numbers hang LEFT of the code box
% (listings sets them with \llap, so no overfull hbox can arise and the code
% x is untouched — measured: code frozen, digits at code_x-11.8pt); inside a
% colored box they are blanked like everyone else's (v3.37 hook), and the
% clean variant restores them in its frameless boxes (v3.78, above).
\lstdefinestyle{pystyle}{
  language=Python,basicstyle=\ttfamily\small,
  columns=fullflexible,keepspaces=true,
  frame=single,rulecolor=\color{cTblLine},
  numbers=left,numberstyle=\tiny\color{cGray},numbersep=8pt,
  xleftmargin=0.5em,framexleftmargin=0pt,framexrightmargin=0pt,
  showstringspaces=false,aboveskip=8pt,belowskip=8pt,
  keywordstyle=\color{cBlue}\bfseries,
  commentstyle=\color{cGray}\itshape,
  stringstyle=\color{cGreen},
}
% v4.5.0: hoisted out of the four styles so an enclosing box can override it.
% Every listing still resolves to true unless something later says otherwise,
% so shipped output is unchanged (verified byte-for-byte on all eight PDFs).
\lstset{resetmargins=true}
\lstdefinestyle{htmlstyle}{
  language=HTML,basicstyle=\ttfamily\small,
  columns=fullflexible,keepspaces=true,
  frame=single,rulecolor=\color{cTblLine},
  numbers=left,numberstyle=\tiny\color{cGray},numbersep=8pt,
  xleftmargin=0.5em,framexleftmargin=0pt,framexrightmargin=0pt,
  showstringspaces=false,aboveskip=8pt,belowskip=8pt,
  keywordstyle=\color{cBlue}\bfseries,
  commentstyle=\color{cGray}\itshape,
  stringstyle=\color{cGreen},
}
\newenvironment{asmblock}{\par\begin{otherlanguage}{english}}{\end{otherlanguage}\par}

% <hllpdf:listing-smart-quotes>
% v4.21.0: four Unicode punctuation characters must enter listings' token
% buffer. Unregistered characters bypass it: a closing quote was printed
% BEFORE the pending word, and an opening quote BEFORE a pending brace.
% Extend the standard character table, then activate only these codepoints
% in the listing's local character-table scope. Do this BEFORE the package's
% delimiter/alsoother/literate hooks so explicit user mappings still win.
% This is not universal Unicode support, nor a conversion to ASCII quotes.
% See listings developer guide 15.3.1 and references/listing-smart-quotes.md.
\makeatletter
\lst@CCPut \lst@ProcessOther {"2018}{"2019}{"201C}{"201D}\z@
\lst@AddToHookAtTop{SelectCharTable}{%
  \lst@ifactivechars
    \catcode"2018=\active \catcode"2019=\active
    \catcode"201C=\active \catcode"201D=\active
  \fi}
\makeatother
% </hllpdf:listing-smart-quotes>

% <hllpdf:listing-technical-symbols>
% v4.30: bounded extension for common technical punctuation in LTR listings.
% Without character-table registration these glyphs bypass the token buffer
% and can move before the preceding word even inside quoted strings.
% Keep the original four-quote block separate for its standing mutation test.
% This is source-faithful printing, not syntax validation or general Unicode.
\makeatletter
\lst@CCPut \lst@ProcessOther
  {"2013}{"2014}{"2026}{"2190}{"2192}{"2194}{"21D2}
  {"2212}{"221E}{"2260}{"2264}{"2265}\z@
\lst@AddToHookAtTop{SelectCharTable}{%
  \lst@ifactivechars
    \catcode"2013=\active \catcode"2014=\active \catcode"2026=\active
    \catcode"2190=\active \catcode"2192=\active \catcode"2194=\active
    \catcode"21D2=\active \catcode"2212=\active \catcode"221E=\active
    \catcode"2260=\active \catcode"2264=\active \catcode"2265=\active
  \fi}
\makeatother
% </hllpdf:listing-technical-symbols>


\usepackage[hidelinks,unicode]{hyperref}   % LAST, after colors. unicode = Hebrew bookmarks.
% v3.62: THE ISOLATION MACROS CORRUPTED EVERY PDF BOOKMARK THEY APPEARED IN.
% \en is defined (line ~155) as
%     \foreignlanguage{english}{\language=\csname l@nohyphenation\endcsname #1}
% and in a heading that body is written verbatim to the .toc and handed to
% hyperref's \pdfstringdef. hyperref strips the two CONTROL SEQUENCES (\language,
% \l@nohyphenation) but the plain `=` CHARACTER token between them is not a control
% sequence, so it survives into the PDF string: \chapter{שער \en{CNOT} והצמדה}
% produced the bookmark "שער =CNOT והצמדה". The printed TOC and the heading on the
% page are CORRECT — only the navigation pane and the viewer window title are wrong,
% which is why it went unnoticed. Measured on the shipped fleet: 70 corrupted
% bookmarks across 5 PDFs, from 177 unprotected \en-in-heading sites. Same corruption
% reached /Title and /Subject via the \hypersetup{pdftitle=...} the docs recommend.
% This hook tells hyperref to use the plain argument for all three isolation macros
% when building a PDF string. It does NOT touch the .toc (the printed TOC is
% byte-unchanged) and it does NOT mask the math signal — a heading with $...$ still
% logs "removing math shift", so BIDI-15 detection is intact.
% BONUS, measured: it also makes \code{} in an UNSTARRED heading COMPILE. BLD-10
% recorded that as needing \texorpdfstring or a starred heading; the real root cause
% is hyperref's \pdfstringdef scan of \code's \ifmmode...\else...\fi body — which is
% exactly why the identically-bodied \texttt survived (hyperref pre-neutralizes
% \texttt in PDF strings) and why \DeclareRobustCommand could never help.
% v3.64: \def\\{ } added. A manual line break in a heading WELDS the two words
% together in the bookmark: \chapter{שורה ראשונה\\שורה שנייה} produced the outline
% entry 'שורה ראשונהשורה שנייה' — "ראשונהשורה" is not a word. hyperref drops the
% \\ and puts nothing in its place, so the break that looks like a space on the
% page is a zero-width join in the navigation pane. The spaced form
% 'מילה אחת \\ מילה שנייה' escapes with a harmless DOUBLE space ('אחת  מילה'),
% which is why this never got noticed. Both are measured on a real outline.
% 0 live instances in the corpus — but SKILL.md and content-style.md both RECOMMEND
% a manual heading break, so this is a primed trap rather than a dormant one, and a
% one-token preamble fix retires it for good instead of asking authors to remember.
% Mapping \\ to a space fixes the WELD, which is the actual defect (a nonexistent
% word in the navigation pane). It does not normalise whitespace: measured after the
% fix, the welded form becomes 'שורה ראשונה שורה שנייה' (correct) and the already-
% spaced form goes from two spaces to three. That is cosmetic — viewers render it as
% a wider gap, and no word is corrupted. A per-heading \texorpdfstring is still the
% way to get an exactly-tuned bookmark string.
\pdfstringdefDisableCommands{\def\en#1{#1}\def\h#1{#1}\def\doclang#1{#1}\def\code#1{#1}\def\codeline#1{#1}\def\\{ }}
% v3.70: RESTORE \url LINE-BREAKING. url.sty freezes its break penalties at LOAD
% time (\mathchardef\UrlBreakPenalty=\binoppenalty), and v3.45's global
% \binoppenalty=10000 (inline-math atomicity) therefore priced EVERY breakpoint
% inside \Url at 10000 = never break. Measured: a long \url in Hebrew prose is
% one unbroken line, ~100pt overfull, invading the margin band the ink gates
% deliberately tolerate. With the two lines below: 0 overfull, the URL breaks
% after '/', the RTL split is CORRECT (fragment order forward), and each line
% gets its own link annotation exactly covering its fragment (measured).
% Confined to \Url's group — inline-math atomicity (TYP-7) is untouched.
\mathchardef\UrlBreakPenalty=200
\mathchardef\UrlBigBreakPenalty=100
\ifx\lakitmainlanguage\hllpdfenglish
  \hypersetup{pdflang=en}
\else
  \hypersetup{pdflang=he}
\fi  % /Lang is the document base locale (not per-span accessibility tagging).
% Set a per-document \hypersetup{pdftitle={...}} too — see SKILL.md.

% v3.52 (round-2 audit): two RTL bibliography/cross-ref defects, both invisible to
% every gate and both affecting the FIRST author to use the feature.
% (a) \@biblabel's default [#1] brackets are bidi-neutral and mirror to ]1[ under
% the RTL bibliography base — every \thebibliography reference list shows ]1[.
% LTR-isolating the label restores [1]. (The render gate's mirror rule fingerprints
% round parens )1(, so square brackets slipped through.)
\makeatletter
\renewcommand*{\@biblabel}[1]{\foreignlanguage{english}{[#1]}}
% v3.54: the SAME neutral-bracket trap on the INLINE side — a multi-key \cite{a,b}
% under the RTL base lays its European-number list right-to-left, so [1, 2] renders
% "[2,1]" and [1, 2, 3] renders "[3,2,1]" (verified by glyph geometry — the reader
% sees the sources in REVERSE order; a note form \cite[p.5]{a} also reorders to
% "[p. 5,1]"). LTR-isolate the whole bracket group (number list AND optional note),
% exactly like \@biblabel above: [1, 2] / [1, 2, 3] / [1, p. 5] all render correct.
% Placed AFTER hyperref so it wins; hyperref hyperlinks each number inside #1 (via
% \@citex), which this wrapper leaves intact. Mirrors the LaTeX kernel \@cite shape.
% v4.9.1: directional isolation must NOT switch the language here. Babel's
% safe-active citation wrapper passes its key list through \bbl@tempa;
% \foreignlanguage reuses that scratch macro before the deferred citation
% loop expands it. The old line therefore cited the literal key "english".
% If such a key existed, even the undefined-reference gate stayed green.
% \babelsublr isolates the bracket group without selecting a locale or
% changing that key list. Preserve native links, key order and optional notes.
\def\@cite#1#2{\babelsublr{[{#1\if@tempswa , #2\fi}]}}
\makeatother

% v4.11.0: a reference inside align's \text inherited HEBdisptag=true,
% so its enclosing hbox mirrored the brackets despite a clean compile.
% Isolate the complete printed reference without selecting a Babel language.
% Also neutralize the display-only normalfont direction reset locally. A
% direction-only patch inverted 12 to 21 inside equation in the measured probe.
% Plain maketag is local to the reference: mathtools' display-side wrapper
% otherwise globally consumes manual_tag while formatting a nested reference.
% Preserve each original command (including the selective auxiliary write),
% current tag forms, native links, and showonlyrefs on/off transitions.
\makeatletter
\NewCommandCopy{\HEB@nativeeqref}{\eqref}
\newcommand{\HEB@eqrefisolate}[1]{%
  \babelsublr{\HEBdisptagfalse\let\bbl@eqnodir\relax\let\maketag@@@\ORIGmaketag@@@#1}}
\protected\def\eqref#1{\HEB@eqrefisolate{\HEB@nativeeqref{#1}}}
\ExplSyntaxOn
\NewCommandCopy{\HEB@nativeMTeqref}{\MT_extended_eqref:n}
\cs_set_protected:Npn \MT_extended_eqref:n #1
  {\HEB@eqrefisolate{\HEB@nativeMTeqref{#1}}}
\ExplSyntaxOff
\makeatother

% v4.11.0: custom caption counters are not part of hyperref's English
% locale, and Hebrew has no upstream hyperref locale in this toolchain.
% Restore the appropriate names at every Babel language switch, including
% return from a nested inline island. Do not localize the destination keys.
\makeatletter
\newcommand{\HEB@refnames}{%
\def\chapterautorefname{פרק}%
\def\sectionautorefname{סעיף}%
\def\subsectionautorefname{תת־סעיף}%
\def\subsubsectionautorefname{תת־תת־סעיף}%
\def\paragraphautorefname{פסקה}%
% v4.31: hyperref uses a distinct counter name at the deepest standard level.
\def\subparagraphautorefname{תת־פסקה}%
\def\equationautorefname{משוואה}%
\def\figureautorefname{איור}%
\def\tableautorefname{טבלה}%
\def\partautorefname{חלק}%
\def\appendixautorefname{נספח}%
\def\footnoteautorefname{הערת שוליים}%
\def\itemautorefname{סעיף}%
\def\Itemautorefname{סעיף}%
\def\figautorefname{איור}%
\def\tblautorefname{טבלה}%
\def\pageautorefname{עמוד}% 
}
\addto\extrashebrew{\HEB@refnames}
\newcommand{\ENG@refnames}{%
  \def\subparagraphautorefname{subparagraph}%
  \def\chapterautorefname{chapter}%
  \def\sectionautorefname{section}%
  \def\subsectionautorefname{subsection}%
  \def\subsubsectionautorefname{subsubsection}%
  \def\paragraphautorefname{paragraph}%
  \def\equationautorefname{Equation}%
  \def\figureautorefname{Figure}%
  \def\tableautorefname{Table}%
  \def\partautorefname{Part}%
  \def\appendixautorefname{Appendix}%
  \def\footnoteautorefname{footnote}%
  \def\itemautorefname{item}%
  \def\figautorefname{Figure}%
  \def\tblautorefname{Table}%
  \def\Itemautorefname{item}%
  \def\pageautorefname{page}%
}
\addto\extrasenglish{\ENG@refnames}
\ifx\lakitmainlanguage\hllpdfenglish\ENG@refnames\else\HEB@refnames\fi
\makeatother

% ============================================================================
% Optional running head (v3.39): \lakithead{<book title>} — OFF by default.
% A document opts in with ONE line after \input{preamble.tex}:
%   \lakithead{מבנה המחשב}
% Every content page then carries the series' orientation line: chapter name
% at the reading edge, book title opposite, thin rule. (fancyhdr's [L]/[R]
% are PHYSICAL; under the RTL page direction [L] lands at the RIGHT/reading
% edge — measured, consistent across the corpus.) Packaged like the topic
% rail: always loaded, inert without the call. This centralizes the 8-line
% block that 8 lakits used to copy by hand — the same duplication vector
% that once let a lakit preamble go ~100 lines stale. Class-aware: report
% fills the mark from \chapter, article from \section. NB: a titlepage env
% self-empties; front matter that sets \pagestyle{empty} must restore with
% \pagestyle{fancy} afterwards (NOT plain — that discards the head).
% ============================================================================
\usepackage{fancyhdr}
% v3.61: \chapter* left a STALE running head — pages of an UNNUMBERED chapter
% carried the PREVIOUS chapter's title. report.cls routes \chapter* through
% \@schapter, which never calls \chaptermark, so \leftmark keeps whatever the last
% numbered chapter left. Measured: with \chapter{פרק ממוספר} then
% \chapter*{מבוא לא ממוספר}, page 4 (inside the starred chapter) printed
% "פרק ממוספר". Invisible on the chapter's FIRST page (\thispagestyle{plain}), so it
% only shows from the second page on — the pages an author is least likely to check.
% An unnumbered front-matter chapter (\chapter*{מבוא}, \chapter*{הקדמה}) is everyday
% Hebrew-academic structure, and the corpus already carries HAND-WRITTEN \markboth
% lines right after \chapter* — the exact per-document duplication \lakithead exists
% to retire. Patched INSIDE \lakithead so the macro stays inert when unused; a
% document's own later \markboth still wins, so those hand-patches do not regress.
\makeatletter
% <hllpdf:starred-heading-scope>
\newif\iflakin@starmarkpatched
\newcommand{\lakistarmark}{%
  \iflakin@starmarkpatched\else
  % Dispatch on the actual class counter, not hyperref's shared internals.
  \ifdefined\c@chapter
    \let\HEBorigschapter\@schapter
    % Record only after the native heading has reached its page. A mark
    % before the native routine can advertise the next heading too early.
    \def\@schapter##1{\HEBorigschapter{##1}\chaptermark{##1}}%
  \else
    % @ssect receives layout arguments, NOT the section name or level.
    % Capture the name in @startsection without introducing a TeX group:
    % grouping sectioning would break the native after-heading state.
    \def\HEB@sectionkind{}%
    \def\HEB@mainsection{section}%
    \let\HEBorigstartsection\@startsection
    \def\@startsection##1##2##3##4##5##6{%
      \def\HEB@sectionkind{##1}%
      \HEBorigstartsection{##1}{##2}{##3}{##4}{##5}{##6}}%
    \let\HEBorigssect\@ssect
    \def\@ssect##1##2##3##4##5{%
      \HEBorigssect{##1}{##2}{##3}{##4}{##5}%
      % Lower starred headings must not replace the article's section mark.
      \ifx\HEB@sectionkind\HEB@mainsection
        \sectionmark{##5}%
      \fi
      \ignorespaces}%
  \fi
  \global\lakin@starmarkpatchedtrue
  \fi}
% </hllpdf:starred-heading-scope>
% <hllpdf:running-language>
% Choose the wrapper when the mark is written, not when a later page ships.
% English punctuation must remain LTR even after returning to Hebrew prose.
\newcommand{\HEB@runningmark}[1]{%
  \iflanguage{english}{\markboth{\en{#1}}{}}{\markboth{\h{#1}}{}}}
% </hllpdf:running-language>
\makeatother
% v4.2.1 (NAV): every collection already names itself in \lakithead — feed the
% PDF /Title from the same declaration, so the navigation gate's
% PDF_REQUIRE_TITLE has a universal source and no fleet member ships
% title-less by missing a second declaration site. Measured before this: 14
% of 16 collections shipped /Title-empty while every one carried a perfectly
% good \lakithead string (plain Hebrew in all 16 — no math, no formatting, so
% the PDF-string path is clean). An explicit title always wins: set it with
% \lakitpdftitle{...} anywhere before \lakithead, or with a plain
% \hypersetup{pdftitle=...} — the fallback detects that via hyperref's
% \@pdftitle, probe-verified \@empty-equal until the first explicit set.
\makeatletter
\newif\iflakit@pdftitleset
\newcommand{\lakitpdftitle}[1]{%
  \hypersetup{pdftitle={#1}}\global\lakit@pdftitlesettrue}
\newcommand{\lakithead}[1]{%
  \iflakit@pdftitleset\else
    \ifx\@pdftitle\@empty
      \hypersetup{pdftitle={#1}}\global\lakit@pdftitlesettrue
    \fi
  \fi
  \pagestyle{fancy}\fancyhf{}%
  \lakistarmark

  \fancyhead[L]{\small\textcolor{cGray}{\leftmark}}%
  % <hllpdf:fixed-head-language>
  % The fixed document title has its OWN language snapshot, independent of
  % the current section/chapter mark. Choose the branch NOW, not at shipout:
  % the page may be Hebrew by the time an English title is finally printed.
  \iflanguage{english}{%
    \fancyhead[R]{\small\textcolor{cGray}{\en{#1}}}%
  }{%
    \fancyhead[R]{\small\textcolor{cGray}{\h{#1}}}%
  }%
  % </hllpdf:fixed-head-language>
  \fancyfoot[C]{\thepage}%
  \renewcommand{\headrulewidth}{0.4pt}%
  \ifdefined\chaptermark
    \renewcommand{\chaptermark}[1]{\HEB@runningmark{##1}}%
  \else
    \renewcommand{\sectionmark}[1]{\HEB@runningmark{##1}}%
  \fi}
\makeatother

% v3.39.1: path-independent byte-reproducibility. LuaTeX derives the trailer
% /ID from the OUTPUT PATH, so two clones of the same source at different
% directories produced different sha256 (measured: the ONLY delta was the
% 62-byte /ID line). Bit 512 suppresses exactly that line — measured: Creator/
% Producer/CreationDate/ModDate and all content untouched. /ID is optional
% per ISO 32000 (required only for ENCRYPTED PDFs and PDF/A — neither is a
% target here). If a future workflow needs an ID (PDF/A, Acrobat revision
% matching), delete this line and it returns.
\pdfvariable suppressoptionalinfo \numexpr 512\relax

% v3.75: BOLD MATH IN BOLD HEADINGS. Heading text is \bfseries but math inside
% it stayed regular-weight CM — a hairline $V_T$ inside a heavy 24.8pt chapter
% title (measured: CMMI12 fl=6 next to FrankRuhlLibre-Bold fl=20; 155 live
% heading-with-math sites fleet-wide). \boldmath rides along wherever the class
% sets \bfseries for a heading; A/B measured: $V_T$ -> CMMIB10 fl=22, 0 missing
% characters (the bold CM faces ship with this TeX Live). Guarded per-macro so
% the article class (no \@makechapterhead) and future class changes stay safe.
% v4.0.1 (adoption fix): the SEARCH patterns above carry no parameter token.
% v4.0.0 made them specific by including the class's own #1/#2, and measured on
% TeX Live 2023 that never matches: a \patchcmd search pattern containing a
% parameter token cannot match while the call sits inside \@ifundefined{...}{}{...}
% (probed both #1 and ##1 -> FAIL; the same pattern at top level succeeds). Every
% one of these patches therefore failed, and because v4.0.0 also replaced the
% empty failure branch with \PackageError, the failure became FATAL: all four
% bundled example assets stopped building ("Cannot patch chapter-title math
% weight -> Fatal error, no output PDF"). The \PackageError branches are KEPT —
% a silent no-op here is what let v3.75's fix rot unnoticed — only the parameter
% tokens are dropped, so the patches match and a real class change still errors.
\makeatletter
\@ifundefined{@makechapterhead}{}{%
  \patchcmd{\@makechapterhead}{\huge\bfseries \@chapapp}
    {\huge\bfseries\boldmath \@chapapp}{}
    {\PackageError{hebrew-lualatex-pdf}{Cannot patch chapter-number math weight}{Unsupported document class}}%
  \patchcmd{\@makechapterhead}{\Huge \bfseries}
    {\Huge \bfseries\boldmath}{}
    {\PackageError{hebrew-lualatex-pdf}{Cannot patch chapter-title math weight}{Unsupported document class}}%
  \patchcmd{\@makeschapterhead}{\Huge \bfseries}
    {\Huge \bfseries\boldmath}{}
    {\PackageError{hebrew-lualatex-pdf}{Cannot patch starred chapter math weight}{Unsupported document class}}}
% v4.0.1: \part sits under a THREE-DEEP wrapper chain, so the class-based anchor
% is not in \@part at all. Measured with \show: hyperref rewrites \@part to
% `...\phantomsection\fi \H@old@part`, and nameref rewrites THAT to
% `\NR@gettitle{#1}\NR@part[{#1}]{#2}` — the class body ends up in \NR@part.
% v4.0.0 patched \@part and made failure fatal, which is why all seven bundled
% examples stopped building. \part is NOT a dead heading level here: 188 live uses
% across the fleet, so the bold-math fix is worth getting right rather than
% dropping. We walk the chain innermost-first and stop at the first hit. Failure is
% a WARNING, not an error: the chain is a property of hyperref/nameref versions, so
% a future package release must degrade to regular-weight part math — a cosmetic
% loss — instead of refusing to build. Chapter and section keep \PackageError,
% because those targets are in the CLASS and were verified to patch here.
\newif\iflakin@partnum \newif\iflakin@parttitle \newif\iflakin@sparttitle
\def\lakin@partnum#1#2{\iflakin@partnum\else
  \patchcmd{#1}{#2\bfseries \partname}{#2\bfseries\boldmath \partname}
    {\lakin@partnumtrue}{}\fi}
\def\lakin@parttitle#1#2{\iflakin@parttitle\else
  \patchcmd{#1}{#2 \bfseries}{#2 \bfseries\boldmath}{\lakin@parttitletrue}{}\fi}
\def\lakin@sparttitle#1#2{\iflakin@sparttitle\else
  \patchcmd{#1}{#2 \bfseries}{#2 \bfseries\boldmath}{\lakin@sparttitletrue}{}\fi}
\@ifundefined{c@chapter}{\def\lakin@pnum{\Large}\def\lakin@ptit{\huge}}%
                        {\def\lakin@pnum{\huge}\def\lakin@ptit{\Huge}}
\@ifundefined{NR@part}{}{\expandafter\lakin@partnum\expandafter\NR@part\expandafter{\lakin@pnum}%
                         \expandafter\lakin@parttitle\expandafter\NR@part\expandafter{\lakin@ptit}}
\@ifundefined{H@old@part}{}{\expandafter\lakin@partnum\expandafter\H@old@part\expandafter{\lakin@pnum}%
                            \expandafter\lakin@parttitle\expandafter\H@old@part\expandafter{\lakin@ptit}}
\@ifundefined{@part}{}{\expandafter\lakin@partnum\expandafter\@part\expandafter{\lakin@pnum}%
                       \expandafter\lakin@parttitle\expandafter\@part\expandafter{\lakin@ptit}}
\@ifundefined{NR@spart}{}{\expandafter\lakin@sparttitle\expandafter\NR@spart\expandafter{\lakin@ptit}}
\@ifundefined{H@old@spart}{}{\expandafter\lakin@sparttitle\expandafter\H@old@spart\expandafter{\lakin@ptit}}
\@ifundefined{@spart}{}{\expandafter\lakin@sparttitle\expandafter\@spart\expandafter{\lakin@ptit}}
\iflakin@partnum\else\PackageWarning{hebrew-lualatex-pdf}{Part-number math stays
  regular weight: no known \noexpand\part wrapper matched (hyperref/nameref
  version change). Cosmetic only}\fi
\iflakin@parttitle\else\PackageWarning{hebrew-lualatex-pdf}{Part-title math stays
  regular weight: no known \noexpand\part wrapper matched. Cosmetic only}\fi
% v4.12.0: THE PART NUMBER IN THE TOC WAS NOT DIRECTIONALLY SEPARATED FROM THE
% TITLE. report.cls/article.cls compose the part entry as ONE text stream --
% \addcontentsline{toc}{part}{\thepart\hspace{1em}#1} -- while EVERY other level
% goes through \numberline, which puts the number in an \hbox. A box is a bidi
% OBJECT and cannot join a neighbouring run; a bare "IV" can. So in the RTL TOC
% paragraph a Latin number and a Latin-INITIAL title merge into a SINGLE LTR run
% and are laid out left-to-right as a unit: the number lands on the wrong side of
% the row, and a mixed title is scrambled outright. Measured on \part{MCP ...}:
% the printed order was "<tail> V MCP" -- title fragment, number, title fragment.
% The chapter/section rows one line below, with the SAME Latin-initial titles,
% were correct on the same page. That contrast is the whole diagnosis.
% \en{} IS NOT A WORKAROUND, measured both ways on one fixture: it fixes a
% PURE-Latin title (\part{\en{Hooks}} renders correctly) and does NOT fix a MIXED
% one (\part{\en{MCP} <hebrew>} stays scrambled). babel's language switch changes
% font and hyphenation, not the bidi run structure -- so the defect is not
% reachable from the source side at all, which is why it belongs here.
% Boxing the number is the entire fix and is a STRICT NO-OP where the bug never
% bit: a Hebrew-titled document compiles BYTE-IDENTICAL before and after (md5,
% 3 clean runs), because there the neutral glue already sat between an L run and
% an R run. The PDF outline is byte-identical too -- hyperref's \pdfstringdef
% neutralises \mbox, so bookmark strings keep their part numbers.
% \numberline{\thepart} would also fix the direction, but it re-geometries every
% part row in the fleet and DROPS the number from the part bookmark (hyperref
% strips \numberline). \mbox changes nothing that was already right.
% The search pattern deliberately carries no # token: etoolbox cannot match a
% parameter token in a patch pattern (measured -- all three chain members
% reported "not patchable" while the pattern included #1). "\thepart \hspace
% {1em}" is brace-balanced and occurs exactly once in the wrapper body.
\newif\iflakin@parttocnum
\def\lakin@parttocnum#1{\iflakin@parttocnum\else
  \patchcmd{#1}{\thepart \hspace {1em}}{\mbox {\thepart }\hspace {1em}}
    {\lakin@parttocnumtrue}{}\fi}
\@ifundefined{NR@part}{}{\lakin@parttocnum\NR@part}
\@ifundefined{H@old@part}{}{\lakin@parttocnum\H@old@part}
\@ifundefined{@part}{}{\lakin@parttocnum\@part}
\iflakin@parttocnum\else\PackageWarning{hebrew-lualatex-pdf}{Part TOC number is
  not boxed: no known \noexpand\part wrapper matched. A part title that BEGINS
  with Latin will print its number on the wrong side of its TOC row}\fi
\patchcmd{\section}{\bfseries}{\bfseries\boldmath}{}
  {\PackageError{hebrew-lualatex-pdf}{Cannot patch section math weight}{Unsupported document class}}
\patchcmd{\subsection}{\bfseries}{\bfseries\boldmath}{}
  {\PackageError{hebrew-lualatex-pdf}{Cannot patch subsection math weight}{Unsupported document class}}
\patchcmd{\subsubsection}{\bfseries}{\bfseries\boldmath}{}
  {\PackageError{hebrew-lualatex-pdf}{Cannot patch subsubsection math weight}{Unsupported document class}}
\patchcmd{\paragraph}{\bfseries}{\bfseries\boldmath}{}{}
\patchcmd{\subparagraph}{\bfseries}{\bfseries\boldmath}{}{}
\@ifundefined{l@chapter}{}{%
  \patchcmd{\l@chapter}{\bfseries}{\bfseries\boldmath}{}{}}
\makeatother
\setcounter{tocdepth}{\lakittocdepth}
% <hllpdf:measured-toc>
% v4.16.0: Measure actual number labels, not a guessed digit count. Read only
% canonical contentsline records; never execute titles or addtocontents code
% during the measuring pass. Widths are local to this individual list and are
% available before its first row: no per-row drift and no extra aux pass.
\makeatletter
\newread\hllpdf@tocstream
\newdimen\hllpdf@toc@chapter
\newdimen\hllpdf@toc@section
\newdimen\hllpdf@toc@subsection
\newdimen\hllpdf@toc@subsubsection
\newdimen\hllpdf@toc@paragraph
\newdimen\hllpdf@toc@subparagraph
\newdimen\hllpdf@toc@figure
\newdimen\hllpdf@toc@table
\newdimen\hllpdf@tocpage
\newdimen\hllpdf@toc@sectionindent
\newdimen\hllpdf@toc@subsectionindent
\newdimen\hllpdf@toc@subsubsectionindent
\newdimen\hllpdf@toc@paragraphindent
\newdimen\hllpdf@toc@subparagraphindent
\def\hllpdf@tocdefaults{%
  \hllpdf@toc@chapter=2em
  \hllpdf@toc@section=3.2em
  \hllpdf@toc@subsection=3.9em
  \@ifundefined{c@chapter}{%
    \hllpdf@toc@subsubsection=3.2em
    \hllpdf@toc@paragraph=4.1em
    \hllpdf@toc@subparagraph=5em
  }{%
    \hllpdf@toc@subsubsection=4.1em
    \hllpdf@toc@paragraph=5em
    \hllpdf@toc@subparagraph=6em
  }%
  \hllpdf@toc@figure=2.3em
  \hllpdf@toc@table=2.3em
  \hllpdf@tocpage=\@pnumwidth\relax
}
\def\hllpdf@tocindents{%
  % Preserve the established shallow layout when its slots already fit.
  \hllpdf@toc@sectionindent=1.5em
  \@tempdima=\hllpdf@toc@chapter \advance\@tempdima by -2em
  \advance\hllpdf@toc@sectionindent by \@tempdima
  \hllpdf@toc@subsectionindent=\dimexpr
    \hllpdf@toc@sectionindent+\hllpdf@toc@section\relax
  \hllpdf@toc@subsubsectionindent=\dimexpr
    \hllpdf@toc@subsectionindent+\hllpdf@toc@subsection\relax
  \hllpdf@toc@paragraphindent=\dimexpr
    \hllpdf@toc@subsubsectionindent+\hllpdf@toc@subsubsection\relax
  \hllpdf@toc@subparagraphindent=\dimexpr
    \hllpdf@toc@paragraphindent+\hllpdf@toc@paragraph\relax
}
\def\hllpdf@tocfont@chapter{\bfseries\boldmath}
\def\hllpdf@tocfont@part{\large\bfseries\boldmath}
\def\hllpdf@tocentryfont{%
  \ifcsname hllpdf@tocfont@\hllpdf@toclevel\endcsname
    \csname hllpdf@tocfont@\hllpdf@toclevel\endcsname
  \fi}
\long\def\hllpdf@tocdiscard#1\hllpdf@tocstop{}
\def\hllpdf@tocscanrecord{%
  \@ifnextchar\contentsline{\hllpdf@tocscanline}{%
    \@ifnextchar\babel@toc{\hllpdf@tocscanlanguage}{\hllpdf@tocdiscard}}}
\long\def\hllpdf@tocscanlanguage\babel@toc#1#2#3\hllpdf@tocstop{%
  % Babel records are declarations, not arbitrary code. Track the language
  % without executing the original directive or changing the surrounding list.
  \def\hllpdf@toclanguage{#1}}
\long\def\hllpdf@tocscanline\contentsline#1#2#3#4#5\hllpdf@tocstop{%
  \def\hllpdf@toclevel{#1}%
  \ifcsname hllpdf@toc@#1\endcsname
    \hllpdf@tocscantitle#2\hllpdf@tocstop
  \fi
  % The page-number column needs its own measurement (including bold chapter
  % rows). Its extra clearance separates the page label from dot leaders.
  \setbox\@tempboxa=\hbox{\foreignlanguage{\hllpdf@toclanguage}{\hllpdf@tocentryfont #3\kern.3em}}%
  \ifdim\wd\@tempboxa>\hllpdf@tocpage
    \hllpdf@tocpage=\wd\@tempboxa
  \fi}
\def\hllpdf@tocscantitle{%
  \@ifnextchar\numberline{\hllpdf@tocscanlabel}{\hllpdf@tocdiscard}}
\long\def\hllpdf@tocscanlabel\numberline#1#2\hllpdf@tocstop{%
  \setbox\@tempboxa=\hbox{\foreignlanguage{\hllpdf@toclanguage}{\hllpdf@tocentryfont #1\kern.5em}}%
  \ifdim\wd\@tempboxa>\csname hllpdf@toc@\hllpdf@toclevel\endcsname
    \csname hllpdf@toc@\hllpdf@toclevel\endcsname=\wd\@tempboxa
  \fi}
\def\hllpdf@tocread{%
  \ifeof\hllpdf@tocstream\else
    \read\hllpdf@tocstream to\hllpdf@tocrecord
    \expandafter\hllpdf@tocscanrecord\hllpdf@tocrecord\hllpdf@tocstop
    \expandafter\hllpdf@tocread
  \fi}
\def\hllpdf@tocmeasure#1{%
  \makeatletter
  \edef\hllpdf@toclanguage{\languagename}%
  \hllpdf@tocdefaults
  \openin\hllpdf@tocstream=\jobname.#1\relax
  \hllpdf@tocread
  \closein\hllpdf@tocstream
  \hllpdf@tocindents
  \edef\@pnumwidth{\the\hllpdf@tocpage}%
  \@tempdima=\dimexpr\hllpdf@tocpage+1em\relax
  \ifdim\@tempdima>\@tocrmarg\relax
    \edef\@tocrmarg{\the\@tempdima}%
  \fi}
\let\hllpdf@originalstarttoc\@starttoc
\def\@starttoc#1{%
  \begingroup
    \def\hllpdf@tocext{#1}%
    \def\hllpdf@tocwanted{toc}%
    \ifx\hllpdf@tocext\hllpdf@tocwanted \hllpdf@tocmeasure{#1}\else
      \def\hllpdf@tocwanted{lof}%
      \ifx\hllpdf@tocext\hllpdf@tocwanted \hllpdf@tocmeasure{#1}\else
        \def\hllpdf@tocwanted{lot}%
        \ifx\hllpdf@tocext\hllpdf@tocwanted \hllpdf@tocmeasure{#1}\fi
      \fi
    \fi
    \hllpdf@originalstarttoc{#1}%
  \endgroup}
\hllpdf@tocdefaults
\hllpdf@tocindents
\@ifundefined{l@chapter}{}{%
  \patchcmd{\l@chapter}{\setlength\@tempdima{1.5em}}%
    {\setlength\@tempdima{\hllpdf@toc@chapter}}{}%
    {\PackageError{hebrew-lualatex-pdf}{Cannot size chapter TOC numbers}%
       {Unsupported contents layout: the chapter number slot was not found}}}
\renewcommand*\l@section{\@dottedtocline{1}%
  {\hllpdf@toc@sectionindent}{\hllpdf@toc@section}}
\renewcommand*\l@subsection{\@dottedtocline{2}%
  {\hllpdf@toc@subsectionindent}{\hllpdf@toc@subsection}}
\renewcommand*\l@subsubsection{\@dottedtocline{3}%
  {\hllpdf@toc@subsubsectionindent}{\hllpdf@toc@subsubsection}}
\renewcommand*\l@paragraph{\@dottedtocline{4}%
  {\hllpdf@toc@paragraphindent}{\hllpdf@toc@paragraph}}
\renewcommand*\l@subparagraph{\@dottedtocline{5}%
  {\hllpdf@toc@subparagraphindent}{\hllpdf@toc@subparagraph}}
\renewcommand*\l@figure{\@dottedtocline{1}{1.5em}{\hllpdf@toc@figure}}
\renewcommand*\l@table{\@dottedtocline{1}{1.5em}{\hllpdf@toc@table}}
\makeatother
% </hllpdf:measured-toc>
\setcounter{secnumdepth}{\lakitsecnumdepth}
% Public preamble-only caption-numbering policy. The historical within-major
% default is untouched unless this command is called. "continuous" removes the
% reset edge as well as the printed prefix, so Figure 1/2 and Table 1/2 stay
% unique across sections/chapters. Returning to "within-major" restores the
% original class-aware displays and native destination conventions. The policy
% never replaces \figcap/\tabcap or their measured link-placement mechanism.
\makeatletter
\let\lakin@majorfig\thefig
\let\lakin@majortbl\thetbl
\let\lakin@majorHfig\theHfig
\let\lakin@majorHtbl\theHtbl
\newcommand{\lakitcaptionnumbering}[1]{%
  \edef\lakin@numbering{#1}%
  \def\lakin@continuous{continuous}\def\lakin@withinmajor{within-major}%
  \ifx\lakin@numbering\lakin@continuous
    \@ifundefined{c@chapter}{%
      \@removefromreset{fig}{section}\@removefromreset{tbl}{section}%
    }{\@removefromreset{fig}{chapter}\@removefromreset{tbl}{chapter}}%
    \renewcommand{\thefig}{\arabic{fig}}\renewcommand{\thetbl}{\arabic{tbl}}%
    \def\theHfig{\arabic{fig}}\def\theHtbl{\arabic{tbl}}%
  \else\ifx\lakin@numbering\lakin@withinmajor
    \@ifundefined{c@chapter}{%
      \@removefromreset{fig}{section}\@removefromreset{tbl}{section}%
      \@addtoreset{fig}{section}\@addtoreset{tbl}{section}%
    }{%
      \@removefromreset{fig}{chapter}\@removefromreset{tbl}{chapter}%
      \@addtoreset{fig}{chapter}\@addtoreset{tbl}{chapter}}%
    \let\thefig\lakin@majorfig\let\thetbl\lakin@majortbl
    \let\theHfig\lakin@majorHfig\let\theHtbl\lakin@majorHtbl
  \else
    \PackageError{hebrew-english-latex-studio}{Unknown caption-numbering policy '#1'}%
      {Use continuous or within-major before begin{document}.}%
  \fi\fi}
\@onlypreamble\lakitcaptionnumbering
\makeatother
% Architecture readiness sentinel. Optional runtime layers require the complete
% core, so the marker is defined only after every mechanism/baseline definition.
% <hllpdf:architecture-ready>
\def\hllpdfcoreloaded{1}
% </hllpdf:architecture-ready>

% v4.22: registered file resources, not an assertion about embedded subset bytes.
% The log framing avoids long-path wrapping and preserves UTF-8 as hex.
% <hllpdf:font-inventory>
\directlua{
  local function inventory()
    local names = {}
    local unnamed = 0
    for id, item in font.each() do
      local path = item.filename
      if path and not (path == "") then names[path] = true else unnamed = unnamed + 1 end
    end
    local sorted = {}
    for path in pairs(names) do table.insert(sorted,path) end
    table.sort(sorted)
    texio.write_nl("log","HLLFONT|BEGIN|1")
    for index,path in ipairs(sorted) do
      local hex = path:gsub(".",function(c) return string.format(string.char(37) .. "02x",string.byte(c)) end)
      local count = math.ceil(string.len(hex)/48)
      texio.write_nl("log","HLLFONT|FILE|" .. index .. "|" .. count)
      for offset=1,string.len(hex),48 do texio.write_nl("log","HLLFONT|HEX|" .. hex:sub(offset,offset+47)) end
    end
    texio.write_nl("log","HLLFONT|END|" .. rawlen(sorted) .. "|" .. unnamed)
    texio.write_nl("log","")
  end
  luatexbase.add_to_callback("finish_pdffile",inventory,"hllpdf.runtime.fonts")
}
% </hllpdf:font-inventory>
; if omitted, the historical canonical
% defaults below are exact.  They deliberately control only the document shell
% and page layout.  RTL/math semantics stay in this preamble, while visual
% presentation themes (for example clean_theme.tex) load AFTER it.
%
%   \def\lakitdocumentclass{article}
%   \def\lakitclassoptions{11pt,a4paper}
%   \def\lakitgeometryoptions{margin=1.8cm,headheight=14pt}
%   \def\lakitlinespread{1.15}
%   \def\lakitparagraphstyle{block} % letters/briefs; indented for continuous prose
%   \def\lakittocdepth{1}
%   \def\lakitsecnumdepth{2}
%   % ============================================================================
% preamble.tex — tested LuaLaTeX preamble for Hebrew + math + figures PDFs (v4.32.0)
% Battle-tested: compiled an 80-page Hebrew quantum-computing summary AND a
% 197-page Hebrew computer-architecture anthology (TikZ diagrams, code listings,
% colored boxes) — both with 0 missing characters.
%
% Usage in your document:
%   \input{preamble.tex}     % (or paste these lines at the top)
%   \begin{document}  ...  \end{document}
%
% Compile with:  lualatex -interaction=nonstopmode doc.tex
%   Repeat until two consecutive PDFs are byte-identical (scripts/build.sh
%   does this) — a TOC document converges only on pass 3+; after two passes
%   the TOC page numbers are STALE (measured; see SKILL.md step 3).
% Font binaries are external dependencies, pinned by assets/fonts/SHA256SUMS.
% scripts/setup_fonts.sh imports a verified local installation or prior skill.
%
% PUBLIC DOCUMENT-SHELL / LAYOUT CONTRACT (introduced v4.3.1; architecture-owned v4.7.1). These hooks may be
% defined BEFORE \input{preamble.tex}; if omitted, the historical canonical
% defaults below are exact.  They deliberately control only the document shell
% and page layout.  RTL/math semantics stay in this preamble, while visual
% presentation themes (for example clean_theme.tex) load AFTER it.
%
%   \def\lakitdocumentclass{article}
%   \def\lakitclassoptions{11pt,a4paper}
%   \def\lakitgeometryoptions{margin=1.8cm,headheight=14pt}
%   \def\lakitlinespread{1.15}
%   \def\lakitparagraphstyle{block} % letters/briefs; indented for continuous prose
%   \def\lakittocdepth{1}
%   \def\lakitsecnumdepth{2}
%   \input{preamble.tex}
%
% Supported class contract: report or article.  Alternate paper sizes must also
% be declared to build.sh through EXPECTED_PAPER; the PDF gate never infers them.
% Do NOT use these hooks as a theme mechanism: theme overlays must not own page
% geometry, class selection, or direction semantics. A declared presentation-policy
% overlay may intentionally change numbering visibility after the preamble.
% ============================================================================
% Runtime architecture guard: generated from scripts/architecture_contract.py.
% The readiness sentinel remains delayed until the END of this file.
% <hllpdf:architecture-guard>
\ifdefined\hllpdfcoreloaded
  \errmessage{hllpdf: preamble.tex loaded more than once}
  \expandafter\endinput
\fi
% </hllpdf:architecture-guard>

\ifdefined\lakitdocumentclass\else\def\lakitdocumentclass{report}\fi
\ifdefined\lakitclassoptions\else\def\lakitclassoptions{12pt,a4paper}\fi
\ifdefined\lakitgeometryoptions\else\def\lakitgeometryoptions{margin=2.3cm,headheight=14pt}\fi
\ifdefined\lakitlinespread\else\def\lakitlinespread{1.3}\fi
\ifdefined\lakittocdepth\else\def\lakittocdepth{1}\fi
\ifdefined\lakitsecnumdepth\else\def\lakitsecnumdepth{2}\fi
% Paragraph organization belongs to the document layout, independently of its
% language and theme. Preserve historical class dimensions unless explicitly
% selecting block paragraphs; never suppress indentation with an everypar hack.
\ifdefined\lakitparagraphstyle\else\def\lakitparagraphstyle{indented}\fi
\edef\lakitparagraphstyle{\lakitparagraphstyle}
\def\hllpdfparagraphblock{block}
\def\hllpdfparagraphindented{indented}
\ifx\lakitparagraphstyle\hllpdfparagraphblock\else
  \ifx\lakitparagraphstyle\hllpdfparagraphindented\else
    \errmessage{hllpdf: paragraph style must be block or indented}
  \fi
\fi
% v4.3.1: the public class hook has a deliberately narrow support contract.
% Before this guard, \def\lakitdocumentclass{book} silently compiled even though
% every heading/TOC regression test targets report/article; an unsupported shell
% could therefore look "supported" until a subtle class-specific defect shipped.
% Fail before \documentclass rather than letting an undocumented combination enter
% the mechanism core. Plain-TeX comparison is used because LaTeX packages are not
% available yet at this point.
\edef\lakitvalidatedclass{\lakitdocumentclass}
\def\lakitreportclass{report}
\def\lakitarticleclass{article}
\ifx\lakitvalidatedclass\lakitreportclass
\else\ifx\lakitvalidatedclass\lakitarticleclass
\else
  \errmessage{hllpdf: unsupported class \lakitdocumentclass; use report or article}
\fi\fi
\let\lakitvalidatedclass\relax
\let\lakitreportclass\relax
\let\lakitarticleclass\relax
% Language intent is independent of document class, theme and content domain.
% The two dominant-language profiles differ in authoring intent, not direction.
% A balanced bilingual document must explicitly choose a navigation/base locale.
% <hllpdf:language-profile>
\providecommand{\lakitlanguageprofile}{hebrew}
\edef\lakitlanguageprofile{\lakitlanguageprofile}
\def\hllpdfhebrew{hebrew}\def\hllpdfenglish{english}
\def\hllpdfhebrewmixed{hebrew-with-english}
\def\hllpdfenglishmixed{english-with-hebrew}
\def\hllpdfbilingual{bilingual}
\let\hllpdfexpectedlanguage\relax
\ifx\lakitlanguageprofile\hllpdfhebrew
  \let\hllpdfexpectedlanguage\hllpdfhebrew
\else\ifx\lakitlanguageprofile\hllpdfhebrewmixed
  \let\hllpdfexpectedlanguage\hllpdfhebrew
\else\ifx\lakitlanguageprofile\hllpdfenglish
  \let\hllpdfexpectedlanguage\hllpdfenglish
\else\ifx\lakitlanguageprofile\hllpdfenglishmixed
  \let\hllpdfexpectedlanguage\hllpdfenglish
\else\ifx\lakitlanguageprofile\hllpdfbilingual
  \ifdefined\lakitmainlanguage\else
    \errmessage{hllpdf: bilingual profile requires explicit lakitmainlanguage}
  \fi
\else
  \errmessage{hllpdf: unsupported language profile \lakitlanguageprofile}
\fi\fi\fi\fi\fi
\ifdefined\lakitmainlanguage
  \edef\lakitmainlanguage{\lakitmainlanguage}
\else
  \let\lakitmainlanguage\hllpdfexpectedlanguage
\fi
\ifx\lakitmainlanguage\hllpdfhebrew\else
  \ifx\lakitmainlanguage\hllpdfenglish\else
    \errmessage{hllpdf: main language must be hebrew or english}
  \fi
\fi
\ifx\hllpdfexpectedlanguage\relax\else
  \ifx\lakitmainlanguage\hllpdfexpectedlanguage\else
    \errmessage{hllpdf: main language conflicts with language profile}
  \fi
\fi
\typeout{HLLLANG|\lakitlanguageprofile|\lakitmainlanguage}
% </hllpdf:language-profile>
\documentclass[\lakitclassoptions]{\lakitdocumentclass}
\usepackage[no-math]{fontspec}        % (1) ★ no-math — THE critical line (see recipe.md)
% layout=lists asks Babel to adapt list parshape when paragraph and page
% directions differ. Without it, minority-language list bodies lose their
% indentation and labels hang outside the text block, despite a clean log.
% The upstream mechanism preserves ordinary same-direction list geometry.
\usepackage[bidi=basic,layout=lists]{babel}
\ifx\lakitmainlanguage\hllpdfenglish
  \babelprovide[import]{hebrew}
  \babelprovide[import,main]{english}
\else
  \babelprovide[import,main]{hebrew}
  \babelprovide[import]{english}
\fi
% v4.22: bind canonical families and fallbacks to the selected installation.
% Lookup by basename or family alone is not proof of the file actually used.
% <hllpdf:font-binding>
\directlua{
  local dir = os.getenv("HEBREW_LATEX_FONT_DIR")
  if not dir or dir == "" then
    local data = os.getenv("XDG_DATA_HOME")
    if not data or data == "" or not (data:sub(1,1) == "/") then
      data = (os.getenv("HOME") or "") .. "/.local/share"
    end
    dir = data .. "/fonts/hebrew-lualatex-pdf"
  end
  if not (dir:sub(1,1) == "/") then dir = lfs.currentdir() .. "/" .. dir end
  if dir:find("[",1,true) or dir:find("]",1,true) or dir:find(string.char(0),1,true)
     or dir:find(string.char(10),1,true) or dir:find(string.char(13),1,true) then
    tex.error("hllpdf: unsupported character in font directory")
  end
  if not (lfs.attributes(dir,"mode") == "directory") then
    tex.error("hllpdf: font directory is missing")
  end
  hllpdf_font_dir = dir:gsub("/+$", "") .. "/"
  tex.sprint(string.char(92) .. "def" .. string.char(92) .. "hllpdffontdir{")
  tex.sprint(-2,hllpdf_font_dir)
  tex.sprint("}")
  luaotfload.add_fallback("hebfb",{"[" .. hllpdf_font_dir .. "DejaVuSerif.ttf]:mode=node;"})
  luaotfload.add_fallback("hebfbbf",{"[" .. hllpdf_font_dir .. "DejaVuSerif-Bold.ttf]:mode=node;"})
}
% </hllpdf:font-binding>
% v3.22: Frank Ruhl Libre ships NO italic face. Instead of silently-upright
% \emph (v<=3.20) or bold-as-emphasis (v3.21), synthesize a slanted face with
% FakeSlant — real, VISIBLE italic emphasis with correct \emph nesting
% semantics (inner \emph toggles back upright). Verified empirically at
% 300dpi: slant clearly visible, 0 errors, 0 missing glyphs.
\babelfont{rm}[Path={\hllpdffontdir},RawFeature={fallback=hebfb},
  BoldFont={FrankRuhlLibre-Bold.ttf},
  BoldFeatures={RawFeature={fallback=hebfbbf}},   % v3.56 — see hebfbbf above
  ItalicFont={FrankRuhlLibre-Regular.ttf}, ItalicFeatures={FakeSlant=0.22},
  % Bold-italic uses the weight-matched fallback and a synthetic slant.
  BoldItalicFont={FrankRuhlLibre-Bold.ttf}, BoldItalicFeatures={RawFeature={fallback=hebfbbf},FakeSlant=0.22}
  ]{FrankRuhlLibre-Regular.ttf}
% Sans fallbacks preserve family and weight for the few missing codepoints.
\directlua{luaotfload.add_fallback("hebfbsans",{"[" .. hllpdf_font_dir .. "DejaVuSans.ttf]:mode=node;"})}
\directlua{luaotfload.add_fallback("hebfbsansbf",{"[" .. hllpdf_font_dir .. "DejaVuSans-Bold.ttf]:mode=node;"})}
\babelfont{sf}[Path={\hllpdffontdir},RawFeature={fallback=hebfbsans},
  BoldFont={Heebo-Bold.ttf},
  BoldFeatures={RawFeature={fallback=hebfbsansbf}}]{Heebo-Regular.ttf}
% v3.24: Scale=MatchLowercase — DejaVu Sans Mono has an outsized x-height, so
% unscaled inline code visually SHOUTS next to the delicate Frank Ruhl body
% (measured on a 197-page build: scaled, a code atom sits in the line at the
% body's weight, and the document tightened by ~5%). Applies to listings too.
\babelfont{tt}[Path={\hllpdffontdir},RawFeature={fallback=hebfb},Scale=MatchLowercase,
  BoldFont={DejaVuSansMono.ttf},
  BoldFeatures={RawFeature={fallback=hebfb},FakeBold=1.5},
  ItalicFont={DejaVuSansMono.ttf},
  ItalicFeatures={RawFeature={fallback=hebfb},FakeSlant=0.2},
  BoldItalicFont={DejaVuSansMono.ttf},
  BoldItalicFeatures={RawFeature={fallback=hebfb},FakeBold=1.5,FakeSlant=0.2}
  ]{DejaVuSansMono.ttf}

% Explicit multilingual/English profiles use real Latin faces and italics.
% The legacy/default Hebrew profile retains its established measured typography.
% Latin Modern is provided by the TeX distribution, not the pinned font bundle;
% recorder/runtime evidence reports it as an additional toolchain resource.
\ifx\lakitlanguageprofile\hllpdfhebrew\else
  \babelfont[english]{rm}[Ligatures=TeX,RawFeature={fallback=hebfb},
    BoldFeatures={RawFeature={fallback=hebfbbf}},
    BoldItalicFeatures={RawFeature={fallback=hebfbbf}}]{Latin Modern Roman}
  \babelfont[english]{sf}[Ligatures=TeX,RawFeature={fallback=hebfbsans},
    BoldFeatures={RawFeature={fallback=hebfbsansbf}},
    BoldItalicFeatures={RawFeature={fallback=hebfbsansbf}}]{Latin Modern Sans}
\fi

\usepackage{amsmath,amssymb,mathtools}
% amsthm supplies theorem/proof machinery. Environment names remain
% document-specific; the generic proof label is localized below.
\usepackage{amsthm}
% Proof names follow the surrounding prose, including return from an island.
\addto\extrashebrew{\renewcommand{\proofname}{הוכחה}}
\addto\extrasenglish{\renewcommand{\proofname}{Proof}}
\ifx\lakitmainlanguage\hllpdfenglish
  \renewcommand{\proofname}{Proof}
\else
  \renewcommand{\proofname}{הוכחה}
\fi
% v3.24.1: equation tags under babel-bidi. A tag re-emitted by mathtools'
% showonlyrefs inside align inherits the RTL base and its parens MIRROR —
% ")2.1(" instead of "(2.1)" (plain align/equation tags are unaffected;
% verified on a minimal case + a 198-page build). LTR-isolate at the tag-BOX
% level: mathtools' show/hide + numbering logic is untouched. Do NOT override
% \tagform@ instead — that clobbers the showonlyrefs mechanism (measured:
% every tag shows, numbering breaks).
% v3.28.1: the v3.24.1 wrap must apply ONLY to tags of display environments.
% Inline \eqref in Hebrew PROSE also goes through \maketag@@@, and there the
% english-wrap MIRRORS the parens — ")1(" — while the unwrapped original is
% correct (proven on a 2x2 isolation matrix: patch x showonlyrefs; regressed
% every prose \eqref since v3.24.1, caught by the render gate's R4). The two
% call sites are indistinguishable by \ifmmode (tag boxes are text mode), so
% the display environments themselves raise a flag; \eqref in prose sees it
% false and gets the original untouched \maketag@@@.
\usepackage{etoolbox}

% babel bidi=basic (including upstream 26.9) assumes every automatic
% discretionary has a one-glyph replacement. LuaTeX can legally produce one
% with replace=nil after a hyphenated Latin value; Babel then indexes nil and the
% whole build dies. Do the narrow upstream-compatible correction: while Babel
% classifies the list, present only that malformed disc as an ordinary
% discretionary, then restore its subtype. Babel still reorders the ENTIRE box;
% there is no pcall and no successful fallback to an unreordered node list.
% If Babel throws for any other reason, the build remains fatal. The synthetic
% nil-disc probe and the real table gallery in test_all.sh pin both layers.
\AtBeginDocument{\directlua{
  if Babel and Babel.bidi and not Babel.heb_nil_disc_compat then
    Babel.heb_nil_disc_compat = true
    local original = Babel.bidi
    local disc_id = node.id('disc')
    Babel.bidi = function(head, ispar, hdir)
      local hidden = {}
      for item in node.traverse_id(disc_id, head) do
        if item.subtype == 2 and (not item.replace or not item.replace.char) then
          table.insert(hidden, item)
          item.subtype = 0
        end
      end
      local result = original(head, ispar, hdir)
      for _, item in ipairs(hidden) do item.subtype = 2 end
      return result
    end
  end
}}
\makeatletter
\newif\ifHEBdisptag
\let\ORIGmaketag@@@\maketag@@@
\def\maketag@@@#1{\ifHEBdisptag\ORIGmaketag@@@{\foreignlanguage{english}{#1}}\else\ORIGmaketag@@@{#1}\fi}
\@for\HEB@env:=equation,equation*,align,align*,gather,gather*,multline,multline*,flalign,flalign*,alignat,alignat*,eqnarray,eqnarray*\do{%
  \expandafter\AtBeginEnvironment\expandafter{\HEB@env}{\HEBdisptagtrue}}
\makeatother
% English page + Hebrew prose: Babel's equation-only tag flip uses a primitive
% leqno path that shifts the formula to the left edge and the tag into the body
% of an English page. Align's full-width alignment is unaffected. Give only this
% display a physical LTR container and native amsmath eqno/leqno. The engine then
% owns real tag width, collision handling and labels; no proxy tag box hides ink
% from its fit calculation. Babel's saved \bbl@eqnodir still restores Hebrew text
% inside math, and the prose locale/math direction are never switched.
% Close the prior RTL paragraph before normalization. After the environment's
% group restores its directions, close TeX's empty resumed display paragraph and
% use native LaTeX continuation handling: same-paragraph prose has no indent,
% while an authored blank line starts an ordinary indented RTL paragraph.
% Register after Babel's setup; remain local to equation(*) in a main-English
% document's actual RTL prose. Native Hebrew, ordinary English and explicit
% fleqn retain their original paths. No page/display dimensions are assigned.
% <hllpdf:minority-rtl-equation-geometry>
\makeatletter
\newcommand{\HEB@minorityprosecontinuation}{\par\@doendpe}
\newcommand{\HEB@minorityequationgeometry}{%
  \ifx\lakitmainlanguage\hllpdfenglish
    \ifnum\pardirection=\@ne
      \if@fleqn\else
        \par\aftergroup\HEB@minorityprosecontinuation
        \iftagsleft@\let\veqno\@@leqno
        \else\let\veqno\@@eqno\fi
        \bodydir TLT\pardir TLT\textdir TLT\relax
      \fi
    \fi
  \fi}
\AtBeginDocument{%
  \AddToHook{env/equation/begin}{\HEB@minorityequationgeometry}%
  \AddToHook{env/equation*/begin}{\HEB@minorityequationgeometry}}
\makeatother
% </hllpdf:minority-rtl-equation-geometry>
% v3.63: \shortintertext rendered Hebrew MIRROR-REVERSED and flush LEFT — an entire
% commentary line inside an align printed backwards. babel's bidi=basic patches
% \intertext@ only (luababel.def), and mathtools' \shortintertext uses the identical
% \vbox{\normalbaselines … \noindent\ignorespaces#1\par} skeleton but is never
% patched, so its paragraph keeps pardir/textdir LTR: no bidi reordering (logical
% order on the page) and flush-left. Measured against \intertext with the same
% Applying babel's own replacement to \shortintertext@ makes the two identical.
% \h{}/\mbox fix the letter order but leave the line flush LEFT, so they are not
% substitutes. Guarded + \AtBeginDocument because mathtools \let-copies at load time.
% NB \ExplSyntaxOn is REQUIRED: mathtools' internals carry expl3 names, and without
% the expl3 catcodes the `_` is read as a subscript and the patch detonates far away
% with "! Illegal unit of measure (pt inserted)" at the \end{align} (measured).
\makeatletter
\AtBeginDocument{%
  \@ifundefined{shortintertext@}{}{%
    \ExplSyntaxOn
    \bbl@sreplace\shortintertext@{\normalbaselines}%
      {\normalbaselines\ifx\bbl@eqnodir\relax\else\bbl@pardir\@ne\bbl@eqnodir\fi}%
    \ExplSyntaxOff}}
\makeatother
% geometry inherits the canonical A4 sheet from \documentclass; build.sh then
% verifies every rendered MediaBox and CropBox, so a silent Letter fallback is
% impossible to publish as a clean default build.
\usepackage[\lakitgeometryoptions]{geometry}
\linespread{\lakitlinespread}       % layout hook; canonical default remains 1.3
\ifx\lakitparagraphstyle\hllpdfparagraphblock
  \setlength{\parindent}{0pt}
  \setlength{\parskip}{6pt plus 1pt minus 1pt}
\fi
% These dimensions remain ordinary scoped TeX lengths: authors can customize
% them after core input. Language switches inherit them; lists keep their own
% indentation and narrow reference layouts establish their local paragraph gaps.
\typeout{HLLPAR|\lakitparagraphstyle}
\usepackage{xcolor}
\usepackage{mdframed}
\usepackage{listings}                 % code listings (assembly / C / pseudo-code)
\usepackage{tikz}
\usetikzlibrary{arrows.meta,positioning,calc,shapes.geometric,fit,backgrounds}
\usepackage{pgfplots}\pgfplotsset{compat=1.18}
\usepackage{quantikz}                 % quantum circuits (remove if unused)
\usepackage{booktabs}
\usepackage{array}
\usepackage{float}                    % [H] = place table exactly here (reference doc: no drifting)
\usepackage{enumitem}
\usepackage{needspace}

% ============================================================================
% ★★ v3.4 — THE DIRECTION-ISOLATION TRIO (the heart of mixing Hebrew + Latin)
% Mental model:
%   \h{...}    Hebrew (RTL) island      — Hebrew text inside an LTR context
%                                          (TikZ nodes, the english `tikzpic`).
%   \en{...}   English (LTR) isolate    — Latin text / parentheses inside a
%                                          Hebrew (RTL) context. KEEPS PARENS
%                                          AND ORDER CORRECT.
%   \code{...} \en + monospace          — code inside Hebrew / inside a frametitle.
% CRITICAL: never wrap Latin in \h{} (it is for HEBREW). Inside a TikZ node,
% Latin in \h{} renders REVERSED (Datapath->htapataD, CPU->UPC) and parens
% mirror. Isolate every Latin/paren chunk with \en{} or \code{}. See
% references/figures-boxes-listings.md.
% ---- TEXT INSIDE MATH (v3.8) ----------------------------------------------
% Put ANY text that sits inside math in a \text{}. The ONLY thing that is
% invalid in math is a bare \h/\en (= \foreignlanguage) used DIRECTLY, outside
% \text -> "\rmfamily invalid in math mode". Inside \text:
%   Hebrew label           -> \text{...}            (bare Hebrew renders fine;
%                              \h is optional — BUT NOT ALWAYS HARMLESS (v3.67):
%                              if the \h group OPENS with a nested island
%                              ($..$, \en{}, \code{}, \texttt{}, \mbox{}) the
%                              WHOLE group loses bidi and the Hebrew prints in
%                              LOGICAL order = gibberish. Measured:
%                                \text{\h{$f$ קבועה}}  -> reads "העובק"  BROKEN
%                                \text{\h{קבועה $f$}}  -> correct
%                                \text{$f$ קבועה}      -> correct (no \h)
%                              Hebrew first inside \h, or drop the \h. Flagged
%                              H-OPENS-WITH-LTR-GROUP.)
%   plain Latin word        -> \text{...}            (no marker needed)
%   Latin WITH ( ) or [ ]   -> \text{\en{...}}       (\en forces LTR; a plain
%                              \text mirrors the parens in RTL -> "Spec (TNR(" )
%   f(Hebrew/text arg)      -> \text{f}(\text{...})  -- keep the parens at MATH
%                              level, NOT \text{f(arg)} (parens inside the box mirror)
%   code in math            -> \mathtt{...}  (or \texttt{}, math-guarded below)
% NEVER write \en{...} or \h{...} straight into $...$ / \[..\] (outside \text);
% always via \text{}. (\mbox{...} also works but does not scale in scripts.)
% ============================================================================
\newcommand{\h}[1]{\foreignlanguage{hebrew}{#1}}
% v3.42.1 — \en islands do NOT hyphenate (auto OR at explicit hyphens). English hyphenation patterns were
% tearing identifiers and terms ACROSS THE RTL SEAM ("IO Regis-/ters",
% "In-/put", "Band-/width", "Per-/formance", "Multi-/ple" — all shipped on
% real pages; probe-proven: "ALUS-/rcB", "Mem-/ToReg"). An island is a label
% or a term, not English prose — a hyphen tear there is never right. A/B on
% a full 197pp lakit: identical overfull sets, same page count, ~11
% paragraphs healed, one mild underfull (badness 3525) that replaced a
% mutilated "Regis-/ters". Long English PROSE belongs in the englishprose
% environment below, which restores explicit-hyphen breaks while its paragraph
% is still active. A raw otherlanguage block inherits the global prohibition.
% (\l@nohyphenation needs \csname:
% '@' is not a letter here.)
% v3.43.2 — \exhyphenpenalty moved OUT of \en to a GLOBAL assignment below.
% Render-proven: a group-local \exhyphenpenalty is DEAD CODE — TeX reads the
% line-breaking parameters at \par time, after the \en group has closed, so
% \en{write-through} / plain not-taken still tore at the "-" ("write-/
% through" shipped on a real exam page). \language is different: it is
% recorded in the node list as text is added, so the in-group nohyphenation
% assignment above DOES work (and stays). The global penalty was A/B-measured
% on a 197pp lakit: identical overfull set (25), identical page count, the
% not-/taken tear healed. Hebrew is unaffected in practice — Hebrew compounds
% use maqaf (־), not ASCII "-".
\newcommand{\en}[1]{\foreignlanguage{english}{\language=\csname l@nohyphenation\endcsname #1}}
% Babel runs extras locally for inline islands and paragraph language changes.
% An English main document must not inherit the Hebrew no-hyphen paragraph rule.
\addto\extrashebrew{\exhyphenpenalty=10000\relax}
\addto\extrasenglish{\exhyphenpenalty=50\relax}
\ifx\lakitmainlanguage\hllpdfenglish
  \exhyphenpenalty=50\relax
\else
  \exhyphenpenalty=10000\relax
\fi
% Full English paragraphs need ordinary explicit-hyphen breaks. Use this
% environment instead of a raw otherlanguage block; the paragraph ends while
% the local penalty is still active.
\newenvironment{englishprose}
  {\par\begin{otherlanguage}{english}\exhyphenpenalty=50\relax}
  {\par\end{otherlanguage}}
\newenvironment{hebrewprose}
  {\par\begin{otherlanguage}{hebrew}\exhyphenpenalty=10000\relax}
  {\par\end{otherlanguage}}
% Use for shell text whose locale must not depend on the last body paragraph.
\DeclareRobustCommand{\doclang}[1]{\foreignlanguage{\lakitmainlanguage}{#1}}

% ============================================================================
% ★★ — fixes for defects that are INVISIBLE in the log and in the checks
% and only surface on a VISUAL RENDER of the PDF. (Each verified empirically.)
% ----------------------------------------------------------------------------
% v3.22: \emph is now REAL slanted emphasis (FakeSlant italic face defined on
% the \babelfont{rm} line above), replacing the v3.21 \emph->\textbf hack: it
% preserves \emph's toggle semantics and keeps bold free for its own role.
% Prefer bold-as-emphasis for your doc? one line: \renewcommand{\emph}[1]{\textbf{#1}}
% page number: force the digits LTR so they never reverse (18 -> 81) in the RTL
%   footer. That reversal is a NON-DETERMINISTIC bidi glitch (observed once in
%   ~80 pp); this LuaLaTeX build is not byte-reproducible, so the defensive
%   LTR wrap is the only reliable guard. Latin digits are always LTR anyway.
\renewcommand{\thepage}{\foreignlanguage{english}{\arabic{page}}}
% v3.70: \pagenumbering \gdef's \thepage and silently DESTROYS the wrap above —
% after \pagenumbering{arabic} the folio guard is gone for the whole rest of the
% document. Two collections discovered this independently and hand-restore the
% wrap in main.tex with the same comment; that knowledge now lives here instead.
% The wrapper re-applies the LTR wrap after the kernel call, for any style.
\makeatletter
\let\lakin@orig@pagenumbering\pagenumbering
\renewcommand{\pagenumbering}[1]{%
  \lakin@orig@pagenumbering{#1}%
  \renewcommand{\thepage}{\foreignlanguage{english}{\csname @#1\endcsname\c@page}}}
\makeatother
% emergencystretch: Hebrew does not hyphenate, so an unbreakable LTR island
%   (\en{...} / an inline math run) landing at the margin can leave a small
%   (~2-13pt) overfull with no legal breakpoint. This rescues ONLY those
%   inline-prose tie-breaks; it is a ceiling used only on offending lines and
%   does NOT touch display-math / listing / table overflow (those still surface
%   in the log for the content-level overflow pass). Verified: baseline overfull
%   -> 0 with this set, with no change to pagination or underfull.
\setlength{\emergencystretch}{3em}
% v3.40: inline math (almost) never breaks across lines. TeX's default
%   (binop 700 / rel 500) lets an inline run split at an operator —
%   $(1.f)\cdot 2^{e}$ shipped as "(1.f)·" ending one line and "2^e" opening
%   the next (reader-reported). In LTR that split is merely ugly; under an
%   RTL base the carried operator lands on the visually WRONG side of the
%   break and the formula reads backwards.
%   WHY 10000 (v3.45.0 — the display-first resolution). Inline math is
%   FORBIDDEN from breaking, so a torn island under an RTL base is
%   structurally impossible. The v3.40 cycle used 9999 ("break only as a last
%   resort") because a bare 10000 left ~30-39 WIDE inline equations
%   overflowing (worst 90.7pt) — but 9999's "last resort" fires OFTEN in
%   tight, non-hyphenating Hebrew: a reader found torn islands across the
%   fleet (~150 candidates; e.g. "P(y |" ending one line and "x_0)" opening
%   the next, read backwards). The step the 9999 cycle missed: every one of
%   those ~30 overflows is a wide equation that BELONGS in display. Convert
%   them to display FIRST, then 10000 gives 0 tears AND 0 overflows (measured,
%   12/12 fleet). Doctrine: inline math is ATOMIC; a wide equation goes in a
%   display, never inline. See content-style ("inline math is atomic").
\binoppenalty=10000 \relpenalty=10000
% v3.38: uniform sentence spacing. TeX's default (\nonfrenchspacing, sfcode
% 3000) is the OLD Anglo-American wide-space-after-period convention — Hebrew
% typography never had it, and modern Latin practice (Chicago et al.) dropped
% it too. Measured on real pages: post-period gaps up to 1.9-2.1x the line's
% mean word gap — the dominant "rivers" source in justified Hebrew prose.
% Survives babel language switches (probed: sfcode stays french inside \en{}
% islands and otherlanguage blocks).
\frenchspacing
% v3.38: widow/orphan protection. LaTeX's default penalties (150/150/50)
% allow a lone line at a page top/bottom almost freely. 10000 = forbid;
% this is SAFE here precisely because the document is \raggedbottom (the
% demoted line just moves on, no vertical stretch fight) — under flushbottom
% these values would cause underfull pages. raggedbottom is the report-class
% default; made explicit so a future class/option change cannot silently
% flip it and turn the 10000s harmful.
\raggedbottom
\widowpenalty=10000 \clubpenalty=10000 \displaywidowpenalty=10000
% em-dash: --- does NOT become — under Frank Ruhl (the font lacks tlig; verified
%   the width is identical with/without Ligatures=TeX, while the mechanism works
%   on Latin Modern). TYPE THE CHARACTER — (U+2014) DIRECTLY; likewise – (U+2013).
%   check_bidi_figures.py flags a --/--- that is adjacent to Hebrew.
% ============================================================================

% ---- STRUCTURAL DIRECTION-SAFETY for \texttt (v3.7) -------------------------
% A parenthesis, bracket, or $ INSIDE \texttt{} mirrors/migrates in RTL prose
% (MIPS 0($s0)->"(s0$)0", $rs->"rs$", arr[i]->"[arr[i") because \texttt sets the
% font but NOT the direction. Redefining \texttt to LTR-isolate its content fixes
% the WHOLE class at once (parens, brackets, $) so plain \texttt{lw $t0,0($s0)}
% is safe — no need to remember \code{}. Guarded by \ifmmode because
% \foreignlanguage is invalid in math (there, plain monospace is already LTR).
% NB (corrected v3.56 — the old note here said "stays tofu", which is WRONG):
% \texttt{Hebrew} is a FATAL build failure. The tt family is DejaVu Sans Mono
% (zero Hebrew) and the tt fallback hebfb is DejaVuSerif (also zero Hebrew), so
% the subsetter gets an empty glyph set and LuaTeX dies with
%   "! error: (file .../DejaVuSansMono.ttf) (type 2): there are no glyphs in the
%    subset"  -> no output PDF.
% The message names a FONT FILE, never Hebrew/\texttt/the source line. And it is
% non-local: add an unrelated Latin \code{} elsewhere and the fatal degrades into
% a SILENT DELETION of the Hebrew word. Put Hebrew in \h{}/normal text, never in
% \texttt{}/\code{}; flagged HEBREW-IN-MONO. To disable, delete this block.
\let\origtexttt\texttt
\protected\long\def\texttt#1{\ifmmode\origtexttt{#1}\else\foreignlanguage{english}{\origtexttt{#1}}\fi}
% v4.8.0: defer the direction test until TeX has inserted an alignment cell's
% math template. Unprotected expansion here tested \ifmmode too early at the
% FIRST token of aligned/array/cases/matrix cells. Engine-level \protected also
% preserves the original \texttt implementation without the self-recursion
% observed when retrofitting \DeclareRobustCommand to its saved alias.
\protected\long\def\code#1{\ifmmode\origtexttt{#1}\else\foreignlanguage{english}{\origtexttt{#1}}\fi}

% ---- \codeline — display-code line (v3.24): the code analogue of display math.
% Reader-friction rule (content-style.md §5b): inline \code{} is for a SINGLE
% ATOM only (one register/variable/mnemonic, a tight no-space expression).
% A multi-token expression — an instruction with operands, a spaced assignment,
% a register-binding list — zigzags the eye RTL<->LTR mid-sentence and, being
% an unbreakable island, line-breaks badly. Give it its own centered line with
% display-math vertical rhythm. \centerline = one box => the expression never
% wraps; \nopagebreak keeps it glued to its introducing sentence.
%   ...the Hebrew sentence ends with a colon:
%   \codeline{lw \$t1, -4(\$t0)}
%   and the prose continues (escape dollars as \$).
% v3.66 (first line of the body): flush a PENDING \item label before the \par.
% LaTeX defers the label box (\@labels) until the next \everypar, and an \hbox in
% vertical mode never fires \everypar — so `\item \codeline{...}` printed the code
% at y=69.1 and the item's own "1." BELOW it at y=103.6 (the number appears after
% its content, attached to nothing). \if@inlabel\leavevmode\fi is CONDITIONAL on
% purpose: an unconditional \leavevmode would open an empty paragraph line when
% \codeline follows ordinary prose. Verified both ways — in a list the label now
% sits on the code's line at its normal x, and after prose the geometry is
% byte-identical to before.
\makeatletter
\newcommand{\codeline}[1]{%
  \if@inlabel\leavevmode\fi
  \par\nopagebreak
  \vspace{\abovedisplayskip}\nopagebreak
  % v3.62: was \centerline{...}, i.e. \hbox to\hsize{\hss #1 \hss}. \hss has
  % INFINITE SHRINK, so TeX could always "fit" the box and NEVER reported an
  % Overfull — the same \hss mechanism as BLD-27's \item label. Measured: a wide
  % \codeline in a defbox painted its first token ON the box's left border bar and
  % sat 9.5pt outside the frame interior on both sides, with 0 Overfull, 0 Underfull
  % and all four gates PASS; in bare prose it reached 21.0pt into BOTH margins.
  % \hfil has no shrink, so the overflow is reported and build.sh's ovf_vis/ovf_big
  % counters see it; \linewidth (not \hsize) makes it box-aware, so a \codeline
  % inside a colored box centres on the box interior. Verified: a FITTING \codeline
  % keeps byte-identical glyph positions, and the box-crossing case now emits
  % "Overfull \hbox (19.03pt too wide)". A wide \codeline should be split in two.
  \hbox to\linewidth{\hfil\code{#1}\hfil}%
  \nopagebreak\vspace{\belowdisplayskip}}
\makeatother

% ---- a thin decorative rule helper ----
\newcommand{\sep}{\par\vspace{2pt}\noindent\textcolor{cRule}{\rule{\linewidth}{0.6pt}}\par\vspace{4pt}}

% ---- colors ----
\definecolor{cBlue}{HTML}{1F4E79}   \definecolor{cBlueBg}{HTML}{EAF1F8}
\definecolor{cGreen}{HTML}{2E6B4F}  \definecolor{cGreenBg}{HTML}{EAF4EE}
\definecolor{cAmber}{HTML}{9C6A00}  \definecolor{cAmberBg}{HTML}{FBF3E2}
\definecolor{cPurple}{HTML}{5B2C6F} \definecolor{cPurpleBg}{HTML}{F3ECF7}
\definecolor{cTeal}{HTML}{0E6E6E}   \definecolor{cTealBg}{HTML}{E6F4F4}
\definecolor{cRed}{HTML}{8E2B2B}    \definecolor{cRedBg}{HTML}{FBEAEA}
\definecolor{cGray}{HTML}{555555}   \definecolor{cRule}{HTML}{1F4E79}

% ---- math operators ----
\DeclareMathOperator{\Tr}{Tr}
\DeclareMathOperator{\rank}{rank}
\DeclareMathOperator{\diag}{diag}
\DeclareMathOperator{\sgn}{sgn}

% ---- bra-ket (quantikz defines \ket/\bra/\braket; add the rest safely) ----
\providecommand{\ket}[1]{\left|#1\right\rangle}
\providecommand{\bra}[1]{\left\langle#1\right|}
\providecommand{\braket}[2]{\left\langle#1\middle|#2\right\rangle}
\newcommand{\dyad}[1]{\left|#1\right\rangle\!\left\langle#1\right|}
% v3.21: general outer product |a><b| and expectation value <A> — agents reach
% for these by reflex (esp. \expval); quantikz defines neither, so their absence
% threw "Undefined control sequence" in a real build. \providecommand = safe.
\providecommand{\ketbra}[2]{\left|#1\right\rangle\!\left\langle#2\right|}
\providecommand{\expval}[1]{\left\langle#1\right\rangle}
\newcommand{\abs}[1]{\left|#1\right|}
\newcommand{\norm}[1]{\left\lVert#1\right\rVert}
\newcommand{\half}{\tfrac12}
\newcommand{\R}{\mathbb{R}}\newcommand{\C}{\mathbb{C}}\newcommand{\Z}{\mathbb{Z}}
% identity operator. NB: \mathbb{1} silently renders a WRONG glyph (msbm has no
% blackboard digits). Use \mathbb{I} = 𝕀 (works with amssymb, no extra package).
% Do NOT "fix" this back to \mathbb{1}.
\newcommand{\id}{\mathbb{I}}

% ---- physical units (v3.35): \punit — upright unit with thin space ---------
% Usage: $300\punit{K}$, $80\punit{mV/decade}$, $10^{7}\punit{cm/s}$ — the
% NUMBER AND UNIT LIVE IN ONE MATH ISLAND. Promoted verbatim from a 199-page
% production lakit (advanced-electronic-devices, dozens of uses).
% Two measured traps this design avoids (both render-proven):
%   1. siunitx SCRAMBLES under RTL — \SI{3.5}{eV} renders "eV 3.5", exponents
%      detach ("Hz 10 x 10^9"): its quantity assembly is not bidi-isolated.
%      Do NOT load siunitx in this pipeline.
%   2. A text-mode unit next to a bare-prose number is TWO adjacent LTR runs
%      and reorders ("5 mm" -> "mm5" — the BIDI-2/3 class). \punit is therefore
%      math-only ON PURPOSE: prose misuse fails LOUDLY at compile (\mathrm
%      outside math) instead of scrambling silently.
% Domain shorthands (\MeV, \fm...) stay per-document: \newcommand{\MeV}{\,\text{\en{MeV}}}
\providecommand{\punit}[1]{\,\mathrm{#1}}

% ============================================================================
% Colored boxes — mdframed (NOT tcolorbox, which breaks under bidi).
% ★★ v3.4: innerleftmargin/innerrightmargin give code listings room so their
% frame/line-numbers do not spill out the box border (see listing styles below).
% ============================================================================
% v3.61: everyline/split skips were applied to `tablebox` only (BLD-14, v3.37) —
% the six colored boxes below never got them, so they still split OPEN across a page
% break: the page-N segment gets top+sides and NO bottom rule, the page-N+1 segment
% sides+bottom and NO top rule. Measured on the shipped fleet this is a COVERAGE GAP,
% not a re-report of BLD-14: 0 split-open `tablebox` segments (that fix holds) vs 62
% split-open COLORED boxes across 7 documents. Verified after the fix: both segments
% close, and the frametitle is NOT duplicated on the continuation. BLD-14's caveat
% still applies — a box taller than a full page grinds mdframed's split-trial loop, so
% split such content manually.
% ============================================================================
% v3.72: frametitle={} leads EVERY \newmdenv option list below. mdframed options
% are GROUP-scoped: a box nested inside \begin{exbox}[frametitle={...}] used to
% INHERIT the outer title — measured, the identical title string rendered twice,
% the second riding a title-bar fill inside the INNER box's frame, for all seven
% environments in both directions. Both natural forms were broken (= preamble
% bug, not an authoring rule). The explicit reset makes every box titleless
% unless ITS OWN [frametitle={...}] says otherwise; verified both directions:
% no duplicate on the nested box, an explicit title still renders.
% Geometry is shared with the optional semantic-width table container. Themes
% select their own existing semantic style below; neither path fixes columns.
\mdfdefinestyle{hllpdfsemanticgeometry}{linewidth=1.1pt,
  innertopmargin=6pt,innerbottommargin=7pt,skipabove=9pt,skipbelow=9pt,
  innerleftmargin=9pt,innerrightmargin=9pt,everyline=true,
  splitbottomskip=6pt,splittopskip=12pt,roundcorner=3pt}
\newcommand{\hllpdfsemanticframestyle}{hllpdfsemanticgeometry}
\newmdenv[frametitle={},style=hllpdfsemanticgeometry,linecolor=cBlue,backgroundcolor=cBlueBg,
  frametitlebackgroundcolor=cBlue,frametitlefontcolor=white,
  ]{defbox}   % הגדרה
\newmdenv[frametitle={},style=hllpdfsemanticgeometry,linecolor=cGreen,backgroundcolor=cGreenBg,
  frametitlebackgroundcolor=cGreen,frametitlefontcolor=white,
  ]{thmbox}   % משפט/טענה
\newmdenv[frametitle={},style=hllpdfsemanticgeometry,linecolor=cAmber,backgroundcolor=cAmberBg,
  frametitlebackgroundcolor=cAmber,frametitlefontcolor=white,
  ]{notebox}  % הערה
\newmdenv[frametitle={},style=hllpdfsemanticgeometry,linecolor=cPurple,backgroundcolor=cPurpleBg,
  frametitlebackgroundcolor=cPurple,frametitlefontcolor=white,
  ]{exbox}    % דוגמה
\newmdenv[frametitle={},style=hllpdfsemanticgeometry,linecolor=cRed,backgroundcolor=cRedBg,
  frametitlebackgroundcolor=cRed,frametitlefontcolor=white,
  ]{warnbox}  % טעות נפוצה
\newmdenv[frametitle={},style=hllpdfsemanticgeometry,linecolor=cTeal,backgroundcolor=cTealBg,
  frametitlebackgroundcolor=cTeal,frametitlefontcolor=white,
  ]{keybox}   % רעיון מרכזי

% ---- codebox (v4.5.0): the dedicated container for code -------------------
% An OPTION, not a policy: asmblock unboxed stays canonical for a full listing,
% and a snippet inside any of the six colored boxes still works as before.
% What makes codebox worth its own environment is that it is the ONE framed box
% that keeps its margin line numbers.
%
% v4.4.0 shipped it WITHOUT them, on the strength of v3.37's "no margin knob can
% save them" plus four fresh probes that agreed. That was wrong: every probe had
% moved the code, the box, or both — never the key that carries the NUMBER
% COLUMN. Measured on the page:
%   resetmargins=false   the listing follows THIS box instead of the page,
%                        moving the code from x=75.26 to x=97.54
%   xrightmargin=26pt    under RTL this is the PHYSICAL LEFT indent, and the
%                        only key that carries the \llap'ed digits with it:
%                        x=64.28 -> 90.18, clearing the rule at 65.64 by 23.6pt
%   xleftmargin=26pt     moves the code and LEAVES THE DIGITS on the rule —
%                        the asymmetry that hid the fix for eleven versions
% resetmargins is hoisted out of the four styles (see the \lstset below them) so
% this override is possible at all; every other listing still resolves to true.
% Pinned by tests/codebox_numbers.tex + probe_lstnum.py --inside-rule, whose
% negative control strips xrightmargin and watches the digits fall back.
\definecolor{cSlate}{HTML}{3F4A56}  \definecolor{cSlateBg}{HTML}{F2F4F7}
\newmdenv[frametitle={},leftline=true,rightline=true,topline=true,bottomline=true,
  linecolor=cSlate,backgroundcolor=cSlateBg,linewidth=0.9pt,
  frametitlebackgroundcolor=cSlate,frametitlefontcolor=white,
  innertopmargin=6pt,innerbottommargin=7pt,skipabove=9pt,skipbelow=9pt,
  innerleftmargin=9pt,innerrightmargin=9pt,everyline=true,splitbottomskip=6pt,splittopskip=12pt,
  roundcorner=3pt]{codebox}   % קוד
% The two keys that make in-box numbering work, both measured on the page:
%   resetmargins=false  the listing follows THIS box instead of the page
%   xrightmargin        under RTL this is the PHYSICAL LEFT indent, and it is
%                       the only one that carries the \llap'ed number column
%                       with it. xleftmargin moves the code and leaves the
%                       digits behind on the rule (measured: code 97.54,
%                       digits 64.28 — the v3.37 defect, reproduced); with
%                       xrightmargin=26pt the digits land at x=90.18, a
%                       clearance of 23.6pt from the rule at 65.64.
% This asymmetry is why v3.4's note had the semantics backwards and v3.37
% concluded the digits could not be saved.
% <hllpdf:listing-container-state>
% v4.29: all state is local to the current environment group. A semantic
% callout masks number text, while a nested codebox restores the formatter
% saved before that mask. Save it once across successive semantic wrappers;
% an explicit numbers=none (or an explicitly empty formatter) remains valid.
\makeatletter
\newif\ifhllpdf@incodebox
\newif\ifhllpdf@englishcodebox
\newif\ifhllpdf@lstnumbermasked
\let\hllpdf@savedlstnumber\thelstnumber
\makeatother
% </hllpdf:listing-container-state>
% <hllpdf:codebox-gutter>
\makeatletter
\AtBeginEnvironment{codebox}{%
  \hllpdf@incodeboxtrue
  % Snapshot the OUTER language before asmblock selects English.
  \iflanguage{english}{\hllpdf@englishcodeboxtrue}{\hllpdf@englishcodeboxfalse}%
  \ifhllpdf@lstnumbermasked
    \let\thelstnumber\hllpdf@savedlstnumber
    \hllpdf@lstnumbermaskedfalse
  \fi
}
% listings' AddToHook is GLOBAL. Register once, and gate on local box state;
% adding an unconditional hook at every codebox entry leaks into later listings.
\lst@AddToHook{PreInit}{%
  \ifhllpdf@incodebox
    % Apply every listing override here, inside the LISTING's own group.
    % Setting margins at box entry leaks them into nested semantic callouts.
    \lstset{resetmargins=false,xrightmargin=26pt}%
    % Styles resolve before PreInit. In LTR the physical number gutter follows
    % xleftmargin; in Hebrew retain the measured inherited xrightmargin path.
    \ifhllpdf@englishcodebox
      \lstset{xleftmargin=26pt,xrightmargin=0pt}%
    \fi
    % ONE box, not two: suppress only the frame/background/skips of a listing
    % inside this codebox. The next unboxed listing keeps its own style.
    \lstset{frame=none,backgroundcolor=\color{cSlateBg},aboveskip=0pt,belowskip=0pt}%
  \fi
}
\makeatother
% </hllpdf:codebox-gutter>

% v3.37: line numbers inside a COLORED BOX are sliced by the box border —
% measured: digit span x=63.5-67.3 vs border line at x=65.2 (the border runs
% THROUGH the digits), and the inter-border gap (3.5pt) is narrower than a
% digit (3.8pt), so no margin knob can save them (under resetmargins,
% xleftmargin is inert in-box — measured: code at x=75.26 with 2.2em AND
% 3.0em). A 3-6 line boxed snippet does not need line numbers; blank the
% number TEXT inside every colored box (\thelstnumber is what gets printed;
% the style's numbers=left survives, so nothing else shifts — code position
% byte-frozen at x=75.26, verified). Unboxed listings keep their numbers.
% v3.78: the blanking is coupled to the FRAME, not to the box name —
% clean_theme.tex registers LATER hooks (registration order wins) that
% restore the number text inside the boxes it renders frameless.
% v4.5.0 RETRACTS the "no margin knob can save them" clause above. It holds only
% under resetmargins=true (the shipped default): there xleftmargin is inert
% in-box, as measured. With resetmargins=false, xrightmargin DOES carry the
% \llap'ed digit column clear of the rule — measured; the codebox \lstset that
% does it is the one above, in the codebox block. The loop below therefore
% covers SIX boxes, not "every colored box": codebox opts out and keeps its
% numbers.
% <hllpdf:semantic-listing-mask>
\makeatletter
\@for\HEB@lstbox:=defbox,thmbox,notebox,exbox,warnbox,keybox\do{%
  \expandafter\AtBeginEnvironment\expandafter{\HEB@lstbox}{%
    \ifhllpdf@lstnumbermasked\else
      \let\hllpdf@savedlstnumber\thelstnumber
    \fi
    \hllpdf@lstnumbermaskedtrue
    \def\thelstnumber{}%
    % The nearest semantic container owns its listing presentation. On exit,
    % TeX's group restores the enclosing codebox state and formatter.
    \hllpdf@incodeboxfalse}}
\makeatother
% </hllpdf:semantic-listing-mask>

% ============================================================================
% Unified TABLE presentation: subtle uniform slate frame + auto-numbered
% caption rendered INSIDE the frame. Use:
%   \begin{tablebox}\centering  <tabular>  \tabcap{<description>}\end{tablebox}
% Replaces the buggy \caption* (which emitted a stray "טבלה N : *" line) and
% provides an explicit framed option + numbered caption "טבלה N.M: תיאור".
% Optional intro text may be placed at the TOP, before the tabular, inside the box.
% ============================================================================
\definecolor{cTblLine}{HTML}{B9C4D2}  \definecolor{cTblBg}{HTML}{FAFBFC}
% v3.23: class-aware — under `article` there is no chapter counter (the promised
% report->article switch used to die with "No counter 'chapter' defined").
\makeatletter
\@ifundefined{c@chapter}{%
  % v3.31: in article the counter resets per SECTION, so an unprefixed number
  % repeated ("טבלה 1." in every section — measured on a two-section render).
  % Same conditional idea as the report branch: prefix the section number once
  % a real \section exists; a pre-section table stays "טבלה 1".
  \newcounter{tbl}[section]
  \renewcommand{\thetbl}{\ifnum\value{section}>0 \thesection.\fi\arabic{tbl}}%
}{%
  \newcounter{tbl}[chapter]
  % v3.30: a report-class doc that uses SECTIONS ONLY (chapter stays 0 —
  % the common short-doc pattern) used to print "טבלה 0.1"; the chapter
  % prefix now appears only once a real \chapter exists (render-verified
  % both ways: "טבלה 1." chapterless, "טבלה 1.1" with chapters).
  \renewcommand{\thetbl}{\ifnum\value{chapter}>0 \thechapter.\fi\arabic{tbl}}%
  % v3.64: the SAME leak in \thesection itself, left behind by v3.30. That fix
  % guarded the CAPTION counter and stopped there, so a chapterless report still
  % printed the heading "0.1 סעיף ראשון" and \ref to it printed "0.2" — measured:
  %   p1 y=65.2  '…0.1'   p1 y=144.9 '…0.2'   p1 y=183.4 '…0.2.1'
  % report.cls hardcodes \thesection as \thechapter.\@arabic\c@section, and the
  % leading 0 is the chapter counter that no \chapter ever incremented. Fixing
  % \thesection fixes the whole chain — \thesubsection and \thesubsubsection are
  % defined in terms of it. \theequation and \thetbl/\thefig were already guarded,
  % which is exactly why this survived: every NUMBER around the heading looked
  % right, so the heading's own 0 read as a style choice rather than a bug.
  % \appendix is safe: report.cls resets chapter to 0 and redefines \thechapter to
  % \@Alph\c@chapter, so an appendix \chapter makes the counter 1 and the guard
  % emits "A.1" unchanged. Render-verified in all four modes (chapterless,
  % chaptered, appendix+chapter, appendix+section-only).
  \renewcommand{\thesection}{\ifnum\value{chapter}>0 \thechapter.\fi\arabic{section}}%
}
\makeatother
% v3.37: everyline=true — if the box DOES split across pages, every segment
% gets a CLOSED frame (measured before: the split ran open-ended into the page
% edge with no bottom border, and the caption could strand alone at the next
% page top — "טבלה A.2" shipped that way). The split skips give the segments
% breathing room. In practice, with the caption GLUE below, every previously-
% split table in the corpus now moves WHOLE instead (the white-gap tradeoff,
% documented benign); everyline stays as the safety net for a genuinely
% page-tall table. NB: a table taller than ONE FULL PAGE can grind mdframed's
% split-trial loop (reproduced: minimal doc, minutes of 1pt-step Overfull
% probes) — break such a monster into two tables.
\mdfdefinestyle{hllpdftablegeometry}{leftline=true,rightline=true,topline=true,bottomline=true,
  linewidth=0.9pt,
  innertopmargin=9pt,innerbottommargin=8pt,skipabove=11pt,skipbelow=11pt,
  innerleftmargin=11pt,innerrightmargin=11pt,roundcorner=4pt}
\newcommand{\hllpdftableframestyle}{hllpdftablegeometry}
\newmdenv[frametitle={},style=hllpdftablegeometry,
  linecolor=cTblLine,backgroundcolor=cTblBg,
  % v4.2.0 (adopted): tablebox is ATOMIC. Short tableboxes near page bottoms
  % repeatedly entered mdframed's split trial, emitting benign Overfull-vbox
  % bursts and occasionally separating the intro prose from the tabular.
  % nobreak=true moves the complete short table to the next page instead. The
  % output-active Overfull-vbox gate still catches a box taller than a page;
  % everyline/split skips remain as defensive styling should that path ever
  % re-arm, but the canonical path never publishes a split tablebox.
  nobreak=true,
  everyline=true,splitbottomskip=6pt,splittopskip=12pt,
  % v3.79 (BLD-37 table side, correction): startinnercode/endinnercode run
  % INSIDE mdframed's body vbox (mdframed.sty:848 / :856+863) -- the second one
  % after \par\unskip, i.e. after line breaking, which is the only moment the
  % \tabcap anchor's true offset from the top of the box EXISTS. See
  % \lakintabboxclose below. These two keys are the ONLY tablebox change.
  startinnercode={\lakintabboxopen},endinnercode={\lakintabboxclose}]{tablebox}
% v3.37: \nopagebreak GLUE — \vsplit honors penalties, so the caption can no
% longer strand alone at the top of the next page (it takes its rows along, or
% the whole box moves). Render-verified on the real shipped orphan.
% v3.59: \NR@gettitle is what nameref injects into the stdlib title-bearing
% commands (\@sect, \@chapter, \caption) to fill \@currentlabelname. \figcap and
% \tabcap are INDEPENDENT macros — they never called it, so \@currentlabelname kept
% whatever the last heading left and \nameref{fig:x} silently printed the enclosing
% SECTION's title instead of the caption (proven in the .aux: field 3 of \newlabel
% was the heading text while field 1 and the anchor `fig.1.1` were correct — so this
% is a THIRD, distinct head next to BLD-13's counter fix and BIDI-38's autoref-name
% fix, neither of which writes that field). \@ifundefined keeps it inert if
% hyperref/nameref is ever dropped. Guarded helper because \figcap is defined
% after \makeatother.
% NB both helpers below need \makeatletter — this part of the preamble is OUTSIDE
% any letter-catcode region, so an unwrapped \@currentlabel parses as LaTeX's \@
% and dies with "You can't use \spacefactor in vertical mode" (measured).
\makeatletter
\newcommand{\lakinrtitle}[1]{\@ifundefined{NR@gettitle}{}{\NR@gettitle{#1}}}
% v3.59: \hitem{<letter>}{<key>} — a REFERENCEABLE Hebrew-lettered list item.
%   \begin{enumerate} \hitem{א}{it:a} הסעיף הראשון. \end{enumerate}  ->  \ref{it:a} = (א)
% Why this exists: automatic Hebrew list labels fail silently, so SKILL.md mandates
% an EXPLICIT \item[(\textbf{א})] — but LaTeX puts \refstepcounter INSIDE
% \if@noitemarg (latex.ltx:12964), and \item[...] never sets that flag. So an
% explicit item steps no counter and never touches \@currentlabel: a \label there
% binds to the AMBIENT one and \ref prints the CHAPTER/SECTION number where the
% item letter belongs (proven in the .aux: {chapter.1}). Same doctrine as \figcap —
% own the label state before \label. \phantomsection matters: without it the
% hyperlink still lands on the chapter head (measured); with it the anchor is
% section*.N and the bookmark outline stays clean.
\newcommand{\hitem}[2]{\item[(\textbf{#1})]%
  \def\@currentlabel{(#1)}\def\@currentlabelname{#1}\phantomsection\label{#2}}
% NB (v3.66): \hitem is \ref-only. \autoref{key} calls it a SECTION ("סעיף (א)")
% because \phantomsection brands the anchor section*.N, and \nameref{key} prints
% the bare letter. For \autoref-able items prefer \hebenumlabels below.
%
% v3.66: \hebenumlabels — OPT-IN automatic Hebrew enumerate letters, done right.
%   {\hebenumlabels\begin{enumerate}\item\label{q:a} ...\end{enumerate}}
%   -> labels (א) (ב) (ג); \ref{q:a} = א; anchor Item.N (real, per-item).
% This retires the belief that made \item[(\textbf{א})] + \hitem necessary.
% "Automatic Hebrew labels fail silently" was TRUE but unexplained, and the
% mechanism is now bisected (plain article vs +enumitem, measured): ENUMITEM
% PRE-EXPANDS the label expression in an \edef BEFORE \refstepcounter, with the
% counter still at 0. Tokens it shields by name (\arabic,\alph,\roman) survive;
% a plain \newcommand \ifcase wrapper evaluates at 0 -> the EMPTY zero-branch ->
% the label prints as "()" with the letter GONE (and \the\c@enumi prints a
% literal 0, \number\value dies fatally pointing at \item). The cure is one
% property, not a doctrine: the wrapper must be ROBUST (survives the \edef
% unexpanded), while \theenumi stays EXPANDABLE (so \refstepcounter writes the
% real letter to the .aux). Verified: labels (א)(ב)(ג), \ref -> א/ב/ג each item
% its own value, .aux = {א}{...}{Item.1} — versus the explicit-\item doctrine
% where BIDI-50/\hitem exist purely to patch the counter that never stepped.
\def\@hebalph#1{\ifcase#1\or א\or ב\or ג\or ד\or ה\or ו\or ז\or ח\or ט\or י\or
  יא\or יב\or יג\or יד\or טו\or טז\or יז\or יח\or יט\or כ\else\@ctrerr\fi}
\DeclareRobustCommand*\hebalph[1]{\expandafter\@hebalph\csname c@#1\endcsname}
\newcommand{\hebenumlabels}{%
  \renewcommand{\theenumi}{\@hebalph\c@enumi}%
  \renewcommand{\labelenumi}{(\textbf{\hebalph{enumi}})}}
\makeatother
% v3.63: \mbox the LABEL. \thetbl ("7.1") and a caption description that STARTS
% WITH A DIGIT ("32 הרגיסטרים…") are two European-Number runs separated only by
% neutrals (the label's own "." and \enspace). Under the RTL paragraph base the UBA
% resolves those neutrals to the RTL level, so both numbers sit in ONE reversed span
% and EXCHANGE VISUAL POSITIONS. Measured on the shipped PDF: a table whose real
% number is 7.1 printed "טבלה 32" — a different, EXISTING table number — while \ref
% correctly printed 7.1, so the caption silently contradicted its own cross-reference,
% and the description lost its leading count. The source contains no adjacency at all;
% it is manufactured by this macro, which is why every source linter was blind.
% \mbox makes the label one atomic box so nothing can reorder across it. Verified
% zero-regression: an ordinary Hebrew caption keeps identical x-geometry.
% One paragraph style covers label, number, separator and description.
% Explicit project/template overrides remain possible; the label hook is kept
% for compatibility, with an identity default. Counters/anchors stay unchanged.
\newcommand{\lakitcaptionstyle}{\normalfont\small\color{black}\hypersetup{allcolors=black}}
\newcommand{\lakitcaptionlabelstyle}[1]{#1}
\newcommand{\lakitcaptionpunctuation}{:}
% Native floats, captionof and longtable use the same presentation policy.
% subcaption is optional and inherits the same font rather than a second size.
\usepackage{caption}
\DeclareCaptionFont{lakit}{\lakitcaptionstyle}
\DeclareCaptionLabelFormat{lakit}{\lakitcaptionlabelstyle{#1~\en{#2}}}
\DeclareCaptionLabelSeparator{lakit}{\lakitcaptionpunctuation\enspace}
\captionsetup{font=lakit,labelfont={},textfont={},labelformat=lakit,
  labelsep=lakit,justification=centering}
\makeatletter
\AtBeginDocument{\@ifpackageloaded{subcaption}{%
  \captionsetup[sub]{font=lakit,labelfont={},textfont={},labelsep=lakit}}{}}
\makeatother
% v3.64: OPT-IN list-of-figures/tables entry via an optional short caption.
% \listoffigures and \listoftables are inherited from report.cls and LOOK available,
% but \figcap/\tabcap step the custom \c@fig/\c@tbl counters and never call
% \addcontentsline — so both printed their heading and NOTHING ELSE. Measured: a doc
% with one \figcap and one \tabcap rendered page 1 as 'רשימת האיורים\n1\n', and the
% .lof held only \@chapter's \addvspace, zero \contentsline entries. An empty list
% under a correct Hebrew heading reads as "this document has no figures".
%
% WHY OPTIONAL, AND NOT AUTOMATIC. Writing #1 unconditionally was implemented first
% and REJECTED: a caption containing \codeline compiles CLEANLY today (verified on
% the pre-v3.64 preamble: 0 errors, PDF produced) but becomes FATAL once it is also
% written to the .lof — \codeline is \hbox to\linewidth{...}, and a full-width box in
% a list entry dies with "\pdfextension endlink cannot be used in vertical mode", no
% output PDF. Gaining a feature nobody uses (0 of 13 collections call \listoffigures)
% at the price of breaking a construct that works is a bad trade. Stock \caption has
% exactly this limitation and exactly this escape hatch, so the convention is familiar.
% \en{}, \code{}, \h{} and inline math in the short text are all safe — verified, they
% survive the write and typeset correctly in the list (RTL geometry confirmed: number
% rightmost at x=509.9, title left of it, dots and folio at the far left).
% With no optional argument the macro is byte-identical to v3.63 — zero risk to every
% existing document, and \listoffigures lists exactly the captions that opted in.
% BLD-37, TABLE side (the \figcap v3.74 idea, adapted): clicking a \ref to a
% table used to land MID-table — hyperref's automatic anchor inside
% \refstepcounter was emitted on the still-open horizontal line whose TAIL is
% the tabular (an hbox whose baseline is the table's vertical CENTER), so the
% jump scrolled the top half of the table above the viewport. The v3.74 entry
% deferred this side claiming "tablebox splits — no single box to measure";
% that deferral dissolves: the measurement target is the TABULAR itself at
% \tabcap ENTRY, while it is still the tail of the open list — the same
% situation \lakin@grabfig handles at an island's end (and the fleet has 0
% split tableboxes). So \lakin@grabtab runs as the FIRST token, before any
% typesetting can push the box away, and the anchor is re-emitted raised above
% the measured tabular. \lakintabraise is a SEPARATE register from
% \lakinfigraise (defined with it, below the LTR-island block): a captionless
% island intentionally leaves the fig register set, and a table must never
% inherit that stale figure height.
%
% WHY the anchor stays AT THE ENTRY SPOT and is NOT deferred into the caption
% paragraph the way \figcap defers it: at \par TeX DROPS a final glue item,
% and the automatic anchor's smash box sitting BEHIND the interword glue kept
% that glue alive — so every centered tabular has always sat half a space off
% true center. Deferring the anchor removed that box and was measured as 982
% differing glyph records: EVERY tabular line shifted 1.30pt (2.61pt for
% pattern 7's two minipages) while captions stayed put. Emitting the raised
% anchor at the old spot — through hyperref's own \Hy@raisedlink, exactly the
% macro the automatic path uses — keeps the line's node structure, and so the
% A/B glyph diff at 0, BY CONSTRUCTION. \Hy@raisedlink reads the raise from
% \HyperRaiseLinkDefault (default \baselineskip = the old mid-table dest); the
% group-local \edef freezes measured-height+42pt into it, and \smash inside
% \Hy@raisedlink hides any raise from the line's extent. \HyCnt@ProvideTheH-
% Counter keeps the dest NAME byte-identical to the automatic one (it defines
% \theHtbl, which \hyper@makecurrent prefers — without it the first measured
% caption was named tbl.1 where the automatic path said tbl.0.1, measured).
% vmode entry, or anything but a box at the tail (no tabular to measure),
% falls back to the plain \refstepcounter — byte-identical to the old path.
\makeatletter
\newcommand{\tabcap}[2][]{\lakin@grabtab
  % v3.79: arm the box-close correction. Set unconditionally and in BOTH
  % branches below's scope, so the plain-\refstepcounter fallback is corrected
  % too; \lakintabboxopen clears it at every \begin{tablebox}, so a \tabcap
  % outside a tablebox can never leave it armed for a later box.
  \global\lakintabcapseentrue
  \ifdim\lakintabraise>0pt
    \H@refstepcounter{tbl}\HyCnt@ProvideTheHCounter{tbl}\hyper@makecurrent{tbl}%
    \begingroup
      % A frameless table targets its measured top. The frame's historical
      % cushion belongs only to tablebox and is corrected at box close.
      \edef\HyperRaiseLinkDefault{\the\lakintabraise}%
      \iflakintabframed
        \edef\HyperRaiseLinkDefault{\the\dimexpr\lakintabraise+42pt\relax}%
      \fi
      \Hy@raisedlink{\hyper@anchorstart{\@currentHref}\hyper@anchorend}%
    \endgroup
    \global\lakintabraise 0pt\relax
  \else
    \refstepcounter{tbl}%
  \fi
  \lakinrtitle{#2}\par\nopagebreak\addvspace{6pt}%
  \def\lakin@short{#1}\ifx\lakin@short\@empty\else
    \addcontentsline{lot}{table}{\protect\numberline{\thetbl}{\protect\ignorespaces #1}}\fi
  {\centering\lakitcaptionstyle\mbox{\lakitcaptionlabelstyle{\tablename~\en{\thetbl}\lakitcaptionpunctuation}}\enspace #2\par}}
\makeatother

% ============================================================================
% Unified FIGURE caption — the \tabcap twin for TikZ/plot figures (v3.34.0).
% A figure (tikzpic) sits inline in the prose; put \figcap BELOW it for a
% numbered, \ref-able caption "איור N.M: תיאור" in the same style as \tabcap:
%   \begin{tikzpic} ... \end{tikzpic}
%   \figcap{\label{fig:foo}description}     % \label FIRST, inside the argument
% WHY a macro, not a hand-written gray line: a manual
%   {\centering\small\textcolor{cGray}{איור: ...}} with a \label inside has NO
% \refstepcounter, so the \label binds to the AMBIENT \@currentlabel (the last
% equation) — \ref then prints an equation number AS the figure number
% (measured on a real page: "איור 48" where 48 is equation 48). \figcap steps a
% real \c@fig counter first, so \label is correct, numbering is consistent, and
% refs are stable. Class-aware like \thetbl (chapter- or section-prefixed).
% ============================================================================
\makeatletter
\@ifundefined{c@chapter}{%
  \newcounter{fig}[section]
  \renewcommand{\thefig}{\ifnum\value{section}>0 \thesection.\fi\arabic{fig}}%
}{%
  \newcounter{fig}[chapter]
  \renewcommand{\thefig}{\ifnum\value{chapter}>0 \thechapter.\fi\arabic{fig}}%
}
\makeatother
\makeatletter
\newcommand{\figcap}[2][]{%
  \ifdim\lakinfigraise>0pt
    % v3.74: step WITHOUT hyperref's automatic anchor, then emit the anchor
    % raised above the measured figure. \H@refstepcounter is hyperref's saved
    % pre-hook \refstepcounter (sets the counter AND \@currentlabel — \label
    % in the caption arg binds exactly as before); \hyper@makecurrent supplies
    % \@currentHref, so the dest NAME is identical to the old automatic one —
    % \ref/\autoref/bookmarks are unchanged, only the dest'S POSITION moves.
    \H@refstepcounter{fig}\hyper@makecurrent{fig}%
    % The anchor is emitted INSIDE the caption paragraph below, as a zero-size
    % raisebox — NOT here in vertical mode, where a \raisebox starts an empty
    % paragraph and adds a blank line (measured: every glyph after the caption
    % shifted +18.8pt). Inside an existing line a [0pt][0pt] box is inert.
    \edef\lakin@anch{\@currentHref}%
    % \edef FREEZES the measured height into the pending macro: the deferred
    % \raisebox is typeset inside the caption paragraph AFTER the register is
    % reset to the 0pt sentinel below — a plain \def read 0 at use time
    % (measured: LAKINRAISE=0.0, dest moved 14pt instead of above the figure).
    \edef\lakin@pend{\noexpand\raisebox{\the\dimexpr\lakinfigraise+42pt\relax}[0pt][0pt]{%
      \noexpand\hyper@anchorstart{\lakin@anch}\noexpand\hyper@anchorend}}%
    \global\lakinfigraise 0pt\relax
  \else
    \refstepcounter{fig}\let\lakin@pend\relax
  \fi
  \lakinrtitle{#2}\par\nopagebreak\addvspace{4pt}%
  \def\lakin@short{#1}\ifx\lakin@short\@empty\else
    \addcontentsline{lof}{figure}{\protect\numberline{\thefig}{\protect\ignorespaces #1}}\fi
  {\centering\lakin@pend\lakitcaptionstyle\mbox{\lakitcaptionlabelstyle{\figurename~\en{\thefig}\lakitcaptionpunctuation}}\enspace #2\par}}
\makeatother

% ============================================================================
% LTR islands.  EVERY tikzpicture / quantikz MUST live inside one, or the whole
% figure's coordinates and Latin labels flip. (Hebrew node text still uses \h{}.)
% ============================================================================
% v3.74 (BLD-37): the \figcap hyperlink DESTINATION used to land at the caption
% line — BELOW the figure — so clicking a figure reference scrolled the figure
% off-screen (measured fleet-wide: 97/97 fig dests in one collection sat 19-54pt
% under their figure's ink; a 347pt figure was fully hidden after the jump).
% hypcap-style fix in two parts: (1) at the island's END, \lakin@grabfig grabs
% the just-typeset centered line with \lastbox, measures it, and puts it back
% with \nointerlineskip — the interline glue above the line was already
% contributed and is NOT re-added, so the page geometry is unchanged to the
% glyph (verified); (2) \figcap then steps the counter WITHOUT hyperref's
% automatic anchor (\H@refstepcounter = the pre-hyperref \refstepcounter) and
% emits the anchor itself, raised by the measured height — the dest lands ABOVE
% the figure top. A captionless island leaves numbering untouched (nothing is
% stepped at the island; the grab flag is cleared at every island begin, so a
% stale height from a captionless figure cannot leak into a later caption).
\newsavebox{\lakinfigbox}
\newdimen\lakinfigraise  % real \dimen: tex.setdimen cannot target a \newlength SKIP register (measured lua error)
% BLD-37 table side: a SEPARATE register for \tabcap. The two must not share —
% a captionless island INTENTIONALLY leaves \lakinfigraise set (the grab runs
% at every island END; only the NEXT island's begin clears it), so a shared
% register would hand a \tabcap after a captionless figure that stale FIGURE
% height. \lakintabraise is written only by \lakin@grabtab at \tabcap entry
% and consumed/reset inside \tabcap itself, so it can never go stale.
\newdimen\lakintabraise
\makeatletter
% v3.74b: measure WITHOUT touching the list. The first draft grabbed the box
% with \lastbox and put it back — and \lastbox DISCARDS the node's baseline
% shift, so a vertically-anchored quantikz circuit rose 2.57pt and everything
% after it sank 5.32pt (measured on the A/B glyph diff). LuaTeX lets us READ
% the tail node's dimensions without removing anything: geometry is identical
% by construction. \lakinfigraise=0pt is the "no figure measured" sentinel.
\def\lakin@grabfig{%
  \ifhmode
    \directlua{
      local n = tex.nest[tex.nest.ptr].tail
      local GLUE, PEN = node.id('glue'), node.id('penalty')
      while n and (n.id == GLUE or n.id == PEN) do n = n.prev end
      local HL, VL = node.id('hlist'), node.id('vlist')
      if n and (n.id == HL or n.id == VL) then
        tex.setdimen('global', 'lakinfigraise',
                     (n.height or 0) + (n.depth or 0))
      else
        tex.setdimen('global', 'lakinfigraise', 0)
      end}%
  \else
    \global\lakinfigraise 0pt\relax
  \fi}
% BLD-37 table side: \lakin@grabfig's twin, NOT a shared helper — the Lua body
% differs in exactly one place, the skip set. Between \end{tabular} and \tabcap
% an otherlanguage wrapper (patterns 3/5 of example_tables) leaves a \write
% whatsit and a dir node ON TOP of the interword glue (measured node dump:
% "whatsit/write <- dir <- glue(2.62) <- hlist"), so the table reader must
% also step over whatsit/dir markers — all zero-size — to reach the tabular.
% \lakin@grabfig deliberately does NOT get the wider skip set: its glue/penalty
% body is the one the v3.74 A/B proof froze, and widening it would newly
% measure (and so newly raise) anchors in island shapes that today take the
% plain path — an unproven behavior change on the FIG side. Anything that is
% not a box behind those markers still hits the 0pt sentinel and \tabcap falls
% back to the plain \refstepcounter.
% SECOND deliberate difference: this stores the box's HEIGHT, not height+depth.
% \figcap's anchor is emitted in the caption paragraph BELOW the figure, so it
% must climb the full extent; \tabcap's anchor sits ON the tabular's own line,
% whose baseline is the vcentered tabular's MIDPOINT — height alone reaches
% the top edge, and adding the depth was measured to overshoot by half the
% table (dests 80-135pt above their frame, INSIDE the preceding table's box).
\def\lakin@grabtab{%
  \ifhmode
    \directlua{
      local n = tex.nest[tex.nest.ptr].tail
      local GLUE, PEN = node.id('glue'), node.id('penalty')
      local WHAT, DIR = node.id('whatsit'), node.id('dir')
      while n and (n.id == GLUE or n.id == PEN
                   or n.id == WHAT or n.id == DIR) do n = n.prev end
      local HL, VL = node.id('hlist'), node.id('vlist')
      if n and (n.id == HL or n.id == VL) then
        tex.setdimen('global', 'lakintabraise', n.height or 0)
      else
        tex.setdimen('global', 'lakintabraise', 0)
      end}%
  \else
    \global\lakintabraise 0pt\relax
  \fi}
\makeatother
% ============================================================================
% v3.79 (BLD-37 table side, CORRECTION): put the anchor at the TOP OF THE BOX,
% instead of guessing the distance to it.
% ============================================================================
% v3.77 raised the \tabcap anchor by (height of the box at the tail of the open
% horizontal list) + 42pt. That is a MEASUREMENT taken BEFORE line breaking, and
% it is short by every term it cannot see. Fleet-measured on the 12 built
% collections (299 tbl.* dests, ownership established through each table's own
% CAPTION, never through "which frame encloses the dest"): 3 dests landed INSIDE
% their own tablebox, below the frame top —
%   computer-architecture       tbl.14.1  -46.4pt   three stacked tabulars
%   computer-architecture       tbl.30.5  -13.5pt   tabular 2.85pt over \linewidth
%   advanced-electronic-devices tbl.B.1   -59.2pt   tabular 2.56pt over \linewidth
% Two missing terms, both reproduced in a minimal document (probe P2/P3):
%   (1) VERTICAL LIST ALREADY CONTRIBUTED inside the box — a second and third
%       tabular, or a heading above the table. The tail read sees only the LAST
%       box, so 77.7pt of stacked tabulars is invisible to the raise.
%   (2) WHICH LINE THE ANCHOR LANDS ON. When the tabular alone overfills
%       \linewidth (459.6121pt in a tablebox), the interword glue after
%       \end{tabular} becomes the chosen breakpoint and \Hy@raisedlink is
%       emitted on the NEXT line — one (depth + lineskip) lower. That term does
%       not EXIST at \tabcap time: line breaking has not happened yet. No
%       measurement taken there can ever recover it. (Proof: probe P3, the same
%       table 2.84912pt too wide, fails at -13.5pt; probe P4, the identical
%       table with \tabcolsep 2pt so it fits, is correct at +32.0pt.)
% So the raise is not computed at \tabcap at all. It is CORRECTED at
% \end{tablebox}, from inside mdframed's own body vbox (endinnercode runs after
% \par\unskip, i.e. after line breaking, while tex.nest[tex.nest.ptr] is still
% that vbox's vertical list). There the anchor's true offset from the top of the
% box content is not estimated but WALKED, and the \Hy@raisedlink shift is set
% so the dest lands 42pt above it — which is 42pt - innertopmargin(9pt) -
% linewidth(0.9pt) = 32.0pt above the FRAME top, the exact position the 295
% already-correct dests have today. That is why this is a correction and not a
% re-placement: the good sites do not move.
%
% WHY THE GLYPH DIFF IS 0 BY CONSTRUCTION — the only thing written is the
% `shift` field of the hlist that \Hy@raisedlink already built, and the guard
% below refuses to write it unless that box is 0x0x0. A zero-size box carries no
% dimension into its parent's packing, the packing already happened, and shift
% is not an input to it. Nothing else in the list is touched: no node is added,
% removed, resized or reordered. Measured, not assumed: assets/example_tables.tex
% and the whole computer-architecture collection A/B at 0 differing pymupdf
% rawdict glyph records.
%
% WHY NOT the alternatives (all prototyped and measured before this one):
%   * accumulate the vertical list at \tabcap time — fixes term (1), CANNOT fix
%     term (2); 2 of the 3 real failures survive. A fix that leaves the majority
%     of the defect standing is the false all-clear this repo exists to avoid.
%   * \zref-savepos two-pass — correct on all shapes, but it makes the dest
%     POSITION aux-derived (DEFECTS BLD-1/2), a single-pass compile silently
%     ships an unraised anchor, and the bare y-difference has no page guard.
%   * emit the anchor at \begin{tablebox} and let \tabcap adopt it — needs the
%     tbl counter stepped at box OPEN to keep the dest NAME byte-identical; a
%     captionless tablebox then either shifts every later table number or leaves
%     an orphan dest that the next table duplicates (a HARD build gate). This
%     correction touches neither the counter nor the name: \tabcap's naming path
%     is v3.77's, unchanged, so resolve_names() is byte-identical (verified).
%   * suppressing the line break (a penalty before the glue) — that MOVES the
%     tabular: the break currently discards the interword glue, so forbidding it
%     re-centers the two overfull tables. Reflow, rejected.
%
% SCOPE. The correction is armed only by \tabcap ITSELF (\lakintabcapseen),
% cleared at every box open, so a captionless tablebox is a no-op and a \tabcap
% hosted OUTSIDE a tablebox (the fleet has exactly one, computer-architecture
% tbl.14.2 in an exbox) keeps v3.77's measured raise. Both fallbacks are
% therefore never worse than today; and because hyperref's AUTOMATIC anchor uses
% the same \Hy@raisedlink shape, the plain-\refstepcounter path (vmode entry, or
% prose between \end{tabular} and \tabcap — probe P6, -51.5pt) is corrected too.
% KNOWN LIMIT — a tablebox that SPLITS across pages. The offset is measured from
% the top of the UNSPLIT body, so a dest in segment 2 gets raised past its own
% segment. MEASURED on a forced split (probe: 40 filler lines + a table): the
% dest went from y=585.0 on page 2 (inside segment 2, at the caption — the old
% BLD-37 symptom) to y=-113.4, i.e. ABOVE the page top. Viewers clamp that to the
% page top, which happens to show segment 2 whole, but an out-of-page coordinate
% is not something to ship on purpose. Hence the \textheight guard in
% lakintab.fix: a body TALLER THAN THE TEXT BLOCK cannot fit on any page, so it
% must split — the correction is skipped and v3.77's behaviour stands unchanged
% (verified: the probe returns to y=585.0). The guard is deliberately EXACT, not
% a page-fitting heuristic: it is sound (a taller-than-page body always splits)
% but PARTIAL (a shorter body that starts low on a page could still split). That
% residue is unreachable in practice — v3.37's everyline + the caption's
% \nopagebreak glue move such boxes WHOLE instead (probe N7: a 22-row table
% pushed down by 430pt moved whole, 0 segments), and the fleet has 0 split
% tableboxes, re-counted after this change. Reimplementing mdframed's own
% page-fitting arithmetic to close the residue would be a heuristic with its own
% failure mode, which is the trade this repo has repeatedly refused.
\newif\iflakintabcapseen
\newif\iflakintabframed
% The Lua lives in ONE \directlua at load time, and is written with NO
% backslash, %, # or ~ : \directlua expands its argument like \edef, so an
% active ~ would expand to \nobreakspace and a % would eat the rest of the line
% (the second is how the first draft of this block died). "not (a == b)" stands
% in for Lua's inequality operator for the same reason.
\directlua{
  lakintab = { }
  lakintab.HL = node.id("hlist")
  lakintab.VL = node.id("vlist")
  lakintab.RU = node.id("rule")
  lakintab.GL = node.id("glue")
  lakintab.KE = node.id("kern")
  lakintab.WH = node.id("whatsit")
  lakintab.TM = node.id("temp")
  lakintab.WD = node.subtype("pdf_dest")
  function lakintab.d(v)
    if v == nil then return 0 end
    if v < -1000000000 then return 0 end
    return v
  end
  function lakintab.scan(list, y, horiz, parent, target)
    local x = list
    while x do
      local i = x.id
      if i == lakintab.HL or i == lakintab.VL then
        if horiz then
          local b = y + x.shift
          if i == lakintab.HL then
            lakintab.scan(x.head, b, true, x, target)
          else
            lakintab.scan(x.head, b - lakintab.d(x.height), false, x, target)
          end
        else
          if i == lakintab.HL then
            lakintab.scan(x.head, y + lakintab.d(x.height), true, x, target)
          else
            lakintab.scan(x.head, y, false, x, target)
          end
          y = y + lakintab.d(x.height) + lakintab.d(x.depth)
        end
      elseif i == lakintab.WH then
        if x.subtype == lakintab.WD then
          local id = x.dest_id
          if type(id) == "string" and string.sub(id, 1, 4) == "tbl." then
            if horiz and parent and parent.width == 0
               and parent.height == 0 and parent.depth == 0 then
              parent.shift = parent.shift + (target - y)
            end
          end
        end
      elseif not horiz then
        if i == lakintab.RU then
          y = y + lakintab.d(x.height) + lakintab.d(x.depth)
        elseif i == lakintab.GL or i == lakintab.KE then
          y = y + lakintab.d(x.width)
        end
      end
      x = x.next
    end
  end
  function lakintab.fix(target, maxh)
    local nst = tex.nest[tex.nest.ptr]
    if nst == nil then return end
    local h = nst.head
    if h == nil then return end
    if h.id == lakintab.TM then h = h.next end
    local tot = 0
    local x = h
    while x do
      local i = x.id
      if i == lakintab.HL or i == lakintab.VL or i == lakintab.RU then
        tot = tot + lakintab.d(x.height) + lakintab.d(x.depth)
      elseif i == lakintab.GL or i == lakintab.KE then
        tot = tot + lakintab.d(x.width)
      end
      x = x.next
    end
    if tot > maxh then return end
    lakintab.scan(h, 0, false, nil, target)
  end
}
% @-FREE NAMES ON PURPOSE: these two are named in the `tablebox` key list at
% \newmdenv above, which is tokenised OUTSIDE any \makeatletter region — a
% \lakin@... there would tokenise as \lakin followed by the letters @tabopen.
\newcommand{\lakintabboxopen}{\lakintabframedtrue\global\lakintabcapseenfalse}
\newcommand{\lakintabboxclose}{%
  \iflakintabcapseen
    \directlua{lakintab.fix(-\number\dimexpr\lakintabcushion\relax,
                            \number\dimexpr\textheight\relax)}%
    \global\lakintabcapseenfalse
  \fi\lakintabframedfalse}
% The cushion above the box's CONTENT top. 42pt puts the dest 32.0pt above the
% frame top (42pt - innertopmargin 9pt - linewidth 0.9pt), i.e. exactly where
% the 295 correct dests already are — the correction moves only the broken ones.
\newcommand{\lakintabcushion}{42pt}
\newenvironment{qcirc}{\par\medskip\global\lakinfigraise 0pt\relax\begin{otherlanguage}{english}\begin{center}}
  {\csname lakin@grabfig\endcsname\end{center}\end{otherlanguage}\par\medskip}
\newenvironment{tikzpic}{\par\medskip\global\lakinfigraise 0pt\relax\begin{otherlanguage}{english}\begin{center}}
  {\csname lakin@grabfig\endcsname\end{center}\end{otherlanguage}\par\medskip}

% ============================================================================
% Code listings.  ★★ v3.4 THE BOX-OVERFLOW FIX:
%   resetmargins=true   → the listing IGNORES the enclosing environment's
%                         indentation and measures from the PAGE. v3.4 stated
%                         this backwards ("measured from the ENCLOSING box"),
%                         and v3.37 then blanked in-box line numbers because
%                         the page-anchored listing put its hanging digits on
%                         top of the box rule. Measured v4.5.0: with the key
%                         FALSE the listing follows the box and the code moves
%                         from x=75.26 to x=97.74, clearing the rule at 72.07.
%                         The key now lives in one \lstset below instead of in
%                         each style, so a box can override it — codebox does.
%   framexleftmargin=0  → the frame (left bar) does not extend left of the code.
%   xleftmargin keeps a small indent so the line numbers have room INSIDE.
% asmblock wraps the listing in an LTR island (code is LTR).
% ============================================================================
\lstdefinestyle{asm}{
  basicstyle=\ttfamily\small,
  backgroundcolor=\color{cBlueBg!55},
  frame=leftline,framerule=2.5pt,rulecolor=\color{cBlue},
  numbers=left,numberstyle=\tiny\color{cGray},numbersep=8pt,
  xleftmargin=2.2em,framexleftmargin=0pt,framexrightmargin=0pt,
  breaklines=true,columns=fullflexible,keepspaces=true,
  showstringspaces=false,aboveskip=8pt,belowskip=8pt,
  % v3.36: without a comment DELIMITER the commentstyle below never fires —
  % listings has no `language` here, so it does not know '#' opens a comment,
  % and every asm listing rendered flat black (measured: a whole-page pixel
  % sweep found ZERO cGreen glyphs across 135 shipped listings). morecomment
  % gives it the delimiter; color-only change, glyph widths untouched.
  morecomment={[l]{\#}},
  commentstyle=\color{cGreen},
}
\lstdefinestyle{cstyle}{
  language=C,basicstyle=\ttfamily\small,
  backgroundcolor=\color{cGreenBg!55},
  frame=leftline,framerule=2.5pt,rulecolor=\color{cGreen},
  numbers=left,numberstyle=\tiny\color{cGray},numbersep=8pt,
  xleftmargin=2.2em,framexleftmargin=0pt,framexrightmargin=0pt,
  breaklines=true,columns=fullflexible,keepspaces=true,
  showstringspaces=false,keywordstyle=\color{cBlue}\bfseries,
  commentstyle=\color{cGray}\itshape,aboveskip=8pt,belowskip=8pt,
}
% v3.36: pystyle/htmlstyle — colored twins of cstyle for the Python/HTML
% listings that used to repeat a ~200-char options blob with NO color
% (keywords rendered bold-black, comments italic-black — measured; C alone
% was colored). Geometry is copied VERBATIM from the shipped blob (frame,
% margins, skips) so swapping a call site to style=pystyle moves NOTHING —
% only keywordstyle/commentstyle/stringstyle are added (bold stays bold,
% italic stays italic, mono metrics identical => zero reflow, proven by
% per-page word-bbox diff). breaklines/breakatwhitespace stay per-site.
% v3.78: numbers=left joins both styles, with the exact numberstyle/numbersep
% keys asm/cstyle carry — fleet-wide consistency (author decision: all four
% styles numbered; "asm/cstyle numbered, py/html not" was an accident of the
% v3.36 geometry freeze, not a design). The numbers hang LEFT of the code box
% (listings sets them with \llap, so no overfull hbox can arise and the code
% x is untouched — measured: code frozen, digits at code_x-11.8pt); inside a
% colored box they are blanked like everyone else's (v3.37 hook), and the
% clean variant restores them in its frameless boxes (v3.78, above).
\lstdefinestyle{pystyle}{
  language=Python,basicstyle=\ttfamily\small,
  columns=fullflexible,keepspaces=true,
  frame=single,rulecolor=\color{cTblLine},
  numbers=left,numberstyle=\tiny\color{cGray},numbersep=8pt,
  xleftmargin=0.5em,framexleftmargin=0pt,framexrightmargin=0pt,
  showstringspaces=false,aboveskip=8pt,belowskip=8pt,
  keywordstyle=\color{cBlue}\bfseries,
  commentstyle=\color{cGray}\itshape,
  stringstyle=\color{cGreen},
}
% v4.5.0: hoisted out of the four styles so an enclosing box can override it.
% Every listing still resolves to true unless something later says otherwise,
% so shipped output is unchanged (verified byte-for-byte on all eight PDFs).
\lstset{resetmargins=true}
\lstdefinestyle{htmlstyle}{
  language=HTML,basicstyle=\ttfamily\small,
  columns=fullflexible,keepspaces=true,
  frame=single,rulecolor=\color{cTblLine},
  numbers=left,numberstyle=\tiny\color{cGray},numbersep=8pt,
  xleftmargin=0.5em,framexleftmargin=0pt,framexrightmargin=0pt,
  showstringspaces=false,aboveskip=8pt,belowskip=8pt,
  keywordstyle=\color{cBlue}\bfseries,
  commentstyle=\color{cGray}\itshape,
  stringstyle=\color{cGreen},
}
\newenvironment{asmblock}{\par\begin{otherlanguage}{english}}{\end{otherlanguage}\par}

% <hllpdf:listing-smart-quotes>
% v4.21.0: four Unicode punctuation characters must enter listings' token
% buffer. Unregistered characters bypass it: a closing quote was printed
% BEFORE the pending word, and an opening quote BEFORE a pending brace.
% Extend the standard character table, then activate only these codepoints
% in the listing's local character-table scope. Do this BEFORE the package's
% delimiter/alsoother/literate hooks so explicit user mappings still win.
% This is not universal Unicode support, nor a conversion to ASCII quotes.
% See listings developer guide 15.3.1 and references/listing-smart-quotes.md.
\makeatletter
\lst@CCPut \lst@ProcessOther {"2018}{"2019}{"201C}{"201D}\z@
\lst@AddToHookAtTop{SelectCharTable}{%
  \lst@ifactivechars
    \catcode"2018=\active \catcode"2019=\active
    \catcode"201C=\active \catcode"201D=\active
  \fi}
\makeatother
% </hllpdf:listing-smart-quotes>

% <hllpdf:listing-technical-symbols>
% v4.30: bounded extension for common technical punctuation in LTR listings.
% Without character-table registration these glyphs bypass the token buffer
% and can move before the preceding word even inside quoted strings.
% Keep the original four-quote block separate for its standing mutation test.
% This is source-faithful printing, not syntax validation or general Unicode.
\makeatletter
\lst@CCPut \lst@ProcessOther
  {"2013}{"2014}{"2026}{"2190}{"2192}{"2194}{"21D2}
  {"2212}{"221E}{"2260}{"2264}{"2265}\z@
\lst@AddToHookAtTop{SelectCharTable}{%
  \lst@ifactivechars
    \catcode"2013=\active \catcode"2014=\active \catcode"2026=\active
    \catcode"2190=\active \catcode"2192=\active \catcode"2194=\active
    \catcode"21D2=\active \catcode"2212=\active \catcode"221E=\active
    \catcode"2260=\active \catcode"2264=\active \catcode"2265=\active
  \fi}
\makeatother
% </hllpdf:listing-technical-symbols>


\usepackage[hidelinks,unicode]{hyperref}   % LAST, after colors. unicode = Hebrew bookmarks.
% v3.62: THE ISOLATION MACROS CORRUPTED EVERY PDF BOOKMARK THEY APPEARED IN.
% \en is defined (line ~155) as
%     \foreignlanguage{english}{\language=\csname l@nohyphenation\endcsname #1}
% and in a heading that body is written verbatim to the .toc and handed to
% hyperref's \pdfstringdef. hyperref strips the two CONTROL SEQUENCES (\language,
% \l@nohyphenation) but the plain `=` CHARACTER token between them is not a control
% sequence, so it survives into the PDF string: \chapter{שער \en{CNOT} והצמדה}
% produced the bookmark "שער =CNOT והצמדה". The printed TOC and the heading on the
% page are CORRECT — only the navigation pane and the viewer window title are wrong,
% which is why it went unnoticed. Measured on the shipped fleet: 70 corrupted
% bookmarks across 5 PDFs, from 177 unprotected \en-in-heading sites. Same corruption
% reached /Title and /Subject via the \hypersetup{pdftitle=...} the docs recommend.
% This hook tells hyperref to use the plain argument for all three isolation macros
% when building a PDF string. It does NOT touch the .toc (the printed TOC is
% byte-unchanged) and it does NOT mask the math signal — a heading with $...$ still
% logs "removing math shift", so BIDI-15 detection is intact.
% BONUS, measured: it also makes \code{} in an UNSTARRED heading COMPILE. BLD-10
% recorded that as needing \texorpdfstring or a starred heading; the real root cause
% is hyperref's \pdfstringdef scan of \code's \ifmmode...\else...\fi body — which is
% exactly why the identically-bodied \texttt survived (hyperref pre-neutralizes
% \texttt in PDF strings) and why \DeclareRobustCommand could never help.
% v3.64: \def\\{ } added. A manual line break in a heading WELDS the two words
% together in the bookmark: \chapter{שורה ראשונה\\שורה שנייה} produced the outline
% entry 'שורה ראשונהשורה שנייה' — "ראשונהשורה" is not a word. hyperref drops the
% \\ and puts nothing in its place, so the break that looks like a space on the
% page is a zero-width join in the navigation pane. The spaced form
% 'מילה אחת \\ מילה שנייה' escapes with a harmless DOUBLE space ('אחת  מילה'),
% which is why this never got noticed. Both are measured on a real outline.
% 0 live instances in the corpus — but SKILL.md and content-style.md both RECOMMEND
% a manual heading break, so this is a primed trap rather than a dormant one, and a
% one-token preamble fix retires it for good instead of asking authors to remember.
% Mapping \\ to a space fixes the WELD, which is the actual defect (a nonexistent
% word in the navigation pane). It does not normalise whitespace: measured after the
% fix, the welded form becomes 'שורה ראשונה שורה שנייה' (correct) and the already-
% spaced form goes from two spaces to three. That is cosmetic — viewers render it as
% a wider gap, and no word is corrupted. A per-heading \texorpdfstring is still the
% way to get an exactly-tuned bookmark string.
\pdfstringdefDisableCommands{\def\en#1{#1}\def\h#1{#1}\def\doclang#1{#1}\def\code#1{#1}\def\codeline#1{#1}\def\\{ }}
% v3.70: RESTORE \url LINE-BREAKING. url.sty freezes its break penalties at LOAD
% time (\mathchardef\UrlBreakPenalty=\binoppenalty), and v3.45's global
% \binoppenalty=10000 (inline-math atomicity) therefore priced EVERY breakpoint
% inside \Url at 10000 = never break. Measured: a long \url in Hebrew prose is
% one unbroken line, ~100pt overfull, invading the margin band the ink gates
% deliberately tolerate. With the two lines below: 0 overfull, the URL breaks
% after '/', the RTL split is CORRECT (fragment order forward), and each line
% gets its own link annotation exactly covering its fragment (measured).
% Confined to \Url's group — inline-math atomicity (TYP-7) is untouched.
\mathchardef\UrlBreakPenalty=200
\mathchardef\UrlBigBreakPenalty=100
\ifx\lakitmainlanguage\hllpdfenglish
  \hypersetup{pdflang=en}
\else
  \hypersetup{pdflang=he}
\fi  % /Lang is the document base locale (not per-span accessibility tagging).
% Set a per-document \hypersetup{pdftitle={...}} too — see SKILL.md.

% v3.52 (round-2 audit): two RTL bibliography/cross-ref defects, both invisible to
% every gate and both affecting the FIRST author to use the feature.
% (a) \@biblabel's default [#1] brackets are bidi-neutral and mirror to ]1[ under
% the RTL bibliography base — every \thebibliography reference list shows ]1[.
% LTR-isolating the label restores [1]. (The render gate's mirror rule fingerprints
% round parens )1(, so square brackets slipped through.)
\makeatletter
\renewcommand*{\@biblabel}[1]{\foreignlanguage{english}{[#1]}}
% v3.54: the SAME neutral-bracket trap on the INLINE side — a multi-key \cite{a,b}
% under the RTL base lays its European-number list right-to-left, so [1, 2] renders
% "[2,1]" and [1, 2, 3] renders "[3,2,1]" (verified by glyph geometry — the reader
% sees the sources in REVERSE order; a note form \cite[p.5]{a} also reorders to
% "[p. 5,1]"). LTR-isolate the whole bracket group (number list AND optional note),
% exactly like \@biblabel above: [1, 2] / [1, 2, 3] / [1, p. 5] all render correct.
% Placed AFTER hyperref so it wins; hyperref hyperlinks each number inside #1 (via
% \@citex), which this wrapper leaves intact. Mirrors the LaTeX kernel \@cite shape.
% v4.9.1: directional isolation must NOT switch the language here. Babel's
% safe-active citation wrapper passes its key list through \bbl@tempa;
% \foreignlanguage reuses that scratch macro before the deferred citation
% loop expands it. The old line therefore cited the literal key "english".
% If such a key existed, even the undefined-reference gate stayed green.
% \babelsublr isolates the bracket group without selecting a locale or
% changing that key list. Preserve native links, key order and optional notes.
\def\@cite#1#2{\babelsublr{[{#1\if@tempswa , #2\fi}]}}
\makeatother

% v4.11.0: a reference inside align's \text inherited HEBdisptag=true,
% so its enclosing hbox mirrored the brackets despite a clean compile.
% Isolate the complete printed reference without selecting a Babel language.
% Also neutralize the display-only normalfont direction reset locally. A
% direction-only patch inverted 12 to 21 inside equation in the measured probe.
% Plain maketag is local to the reference: mathtools' display-side wrapper
% otherwise globally consumes manual_tag while formatting a nested reference.
% Preserve each original command (including the selective auxiliary write),
% current tag forms, native links, and showonlyrefs on/off transitions.
\makeatletter
\NewCommandCopy{\HEB@nativeeqref}{\eqref}
\newcommand{\HEB@eqrefisolate}[1]{%
  \babelsublr{\HEBdisptagfalse\let\bbl@eqnodir\relax\let\maketag@@@\ORIGmaketag@@@#1}}
\protected\def\eqref#1{\HEB@eqrefisolate{\HEB@nativeeqref{#1}}}
\ExplSyntaxOn
\NewCommandCopy{\HEB@nativeMTeqref}{\MT_extended_eqref:n}
\cs_set_protected:Npn \MT_extended_eqref:n #1
  {\HEB@eqrefisolate{\HEB@nativeMTeqref{#1}}}
\ExplSyntaxOff
\makeatother

% v4.11.0: custom caption counters are not part of hyperref's English
% locale, and Hebrew has no upstream hyperref locale in this toolchain.
% Restore the appropriate names at every Babel language switch, including
% return from a nested inline island. Do not localize the destination keys.
\makeatletter
\newcommand{\HEB@refnames}{%
\def\chapterautorefname{פרק}%
\def\sectionautorefname{סעיף}%
\def\subsectionautorefname{תת־סעיף}%
\def\subsubsectionautorefname{תת־תת־סעיף}%
\def\paragraphautorefname{פסקה}%
% v4.31: hyperref uses a distinct counter name at the deepest standard level.
\def\subparagraphautorefname{תת־פסקה}%
\def\equationautorefname{משוואה}%
\def\figureautorefname{איור}%
\def\tableautorefname{טבלה}%
\def\partautorefname{חלק}%
\def\appendixautorefname{נספח}%
\def\footnoteautorefname{הערת שוליים}%
\def\itemautorefname{סעיף}%
\def\Itemautorefname{סעיף}%
\def\figautorefname{איור}%
\def\tblautorefname{טבלה}%
\def\pageautorefname{עמוד}% 
}
\addto\extrashebrew{\HEB@refnames}
\newcommand{\ENG@refnames}{%
  \def\subparagraphautorefname{subparagraph}%
  \def\chapterautorefname{chapter}%
  \def\sectionautorefname{section}%
  \def\subsectionautorefname{subsection}%
  \def\subsubsectionautorefname{subsubsection}%
  \def\paragraphautorefname{paragraph}%
  \def\equationautorefname{Equation}%
  \def\figureautorefname{Figure}%
  \def\tableautorefname{Table}%
  \def\partautorefname{Part}%
  \def\appendixautorefname{Appendix}%
  \def\footnoteautorefname{footnote}%
  \def\itemautorefname{item}%
  \def\figautorefname{Figure}%
  \def\tblautorefname{Table}%
  \def\Itemautorefname{item}%
  \def\pageautorefname{page}%
}
\addto\extrasenglish{\ENG@refnames}
\ifx\lakitmainlanguage\hllpdfenglish\ENG@refnames\else\HEB@refnames\fi
\makeatother

% ============================================================================
% Optional running head (v3.39): \lakithead{<book title>} — OFF by default.
% A document opts in with ONE line after \input{preamble.tex}:
%   \lakithead{מבנה המחשב}
% Every content page then carries the series' orientation line: chapter name
% at the reading edge, book title opposite, thin rule. (fancyhdr's [L]/[R]
% are PHYSICAL; under the RTL page direction [L] lands at the RIGHT/reading
% edge — measured, consistent across the corpus.) Packaged like the topic
% rail: always loaded, inert without the call. This centralizes the 8-line
% block that 8 lakits used to copy by hand — the same duplication vector
% that once let a lakit preamble go ~100 lines stale. Class-aware: report
% fills the mark from \chapter, article from \section. NB: a titlepage env
% self-empties; front matter that sets \pagestyle{empty} must restore with
% \pagestyle{fancy} afterwards (NOT plain — that discards the head).
% ============================================================================
\usepackage{fancyhdr}
% v3.61: \chapter* left a STALE running head — pages of an UNNUMBERED chapter
% carried the PREVIOUS chapter's title. report.cls routes \chapter* through
% \@schapter, which never calls \chaptermark, so \leftmark keeps whatever the last
% numbered chapter left. Measured: with \chapter{פרק ממוספר} then
% \chapter*{מבוא לא ממוספר}, page 4 (inside the starred chapter) printed
% "פרק ממוספר". Invisible on the chapter's FIRST page (\thispagestyle{plain}), so it
% only shows from the second page on — the pages an author is least likely to check.
% An unnumbered front-matter chapter (\chapter*{מבוא}, \chapter*{הקדמה}) is everyday
% Hebrew-academic structure, and the corpus already carries HAND-WRITTEN \markboth
% lines right after \chapter* — the exact per-document duplication \lakithead exists
% to retire. Patched INSIDE \lakithead so the macro stays inert when unused; a
% document's own later \markboth still wins, so those hand-patches do not regress.
\makeatletter
% <hllpdf:starred-heading-scope>
\newif\iflakin@starmarkpatched
\newcommand{\lakistarmark}{%
  \iflakin@starmarkpatched\else
  % Dispatch on the actual class counter, not hyperref's shared internals.
  \ifdefined\c@chapter
    \let\HEBorigschapter\@schapter
    % Record only after the native heading has reached its page. A mark
    % before the native routine can advertise the next heading too early.
    \def\@schapter##1{\HEBorigschapter{##1}\chaptermark{##1}}%
  \else
    % @ssect receives layout arguments, NOT the section name or level.
    % Capture the name in @startsection without introducing a TeX group:
    % grouping sectioning would break the native after-heading state.
    \def\HEB@sectionkind{}%
    \def\HEB@mainsection{section}%
    \let\HEBorigstartsection\@startsection
    \def\@startsection##1##2##3##4##5##6{%
      \def\HEB@sectionkind{##1}%
      \HEBorigstartsection{##1}{##2}{##3}{##4}{##5}{##6}}%
    \let\HEBorigssect\@ssect
    \def\@ssect##1##2##3##4##5{%
      \HEBorigssect{##1}{##2}{##3}{##4}{##5}%
      % Lower starred headings must not replace the article's section mark.
      \ifx\HEB@sectionkind\HEB@mainsection
        \sectionmark{##5}%
      \fi
      \ignorespaces}%
  \fi
  \global\lakin@starmarkpatchedtrue
  \fi}
% </hllpdf:starred-heading-scope>
% <hllpdf:running-language>
% Choose the wrapper when the mark is written, not when a later page ships.
% English punctuation must remain LTR even after returning to Hebrew prose.
\newcommand{\HEB@runningmark}[1]{%
  \iflanguage{english}{\markboth{\en{#1}}{}}{\markboth{\h{#1}}{}}}
% </hllpdf:running-language>
\makeatother
% v4.2.1 (NAV): every collection already names itself in \lakithead — feed the
% PDF /Title from the same declaration, so the navigation gate's
% PDF_REQUIRE_TITLE has a universal source and no fleet member ships
% title-less by missing a second declaration site. Measured before this: 14
% of 16 collections shipped /Title-empty while every one carried a perfectly
% good \lakithead string (plain Hebrew in all 16 — no math, no formatting, so
% the PDF-string path is clean). An explicit title always wins: set it with
% \lakitpdftitle{...} anywhere before \lakithead, or with a plain
% \hypersetup{pdftitle=...} — the fallback detects that via hyperref's
% \@pdftitle, probe-verified \@empty-equal until the first explicit set.
\makeatletter
\newif\iflakit@pdftitleset
\newcommand{\lakitpdftitle}[1]{%
  \hypersetup{pdftitle={#1}}\global\lakit@pdftitlesettrue}
\newcommand{\lakithead}[1]{%
  \iflakit@pdftitleset\else
    \ifx\@pdftitle\@empty
      \hypersetup{pdftitle={#1}}\global\lakit@pdftitlesettrue
    \fi
  \fi
  \pagestyle{fancy}\fancyhf{}%
  \lakistarmark

  \fancyhead[L]{\small\textcolor{cGray}{\leftmark}}%
  % <hllpdf:fixed-head-language>
  % The fixed document title has its OWN language snapshot, independent of
  % the current section/chapter mark. Choose the branch NOW, not at shipout:
  % the page may be Hebrew by the time an English title is finally printed.
  \iflanguage{english}{%
    \fancyhead[R]{\small\textcolor{cGray}{\en{#1}}}%
  }{%
    \fancyhead[R]{\small\textcolor{cGray}{\h{#1}}}%
  }%
  % </hllpdf:fixed-head-language>
  \fancyfoot[C]{\thepage}%
  \renewcommand{\headrulewidth}{0.4pt}%
  \ifdefined\chaptermark
    \renewcommand{\chaptermark}[1]{\HEB@runningmark{##1}}%
  \else
    \renewcommand{\sectionmark}[1]{\HEB@runningmark{##1}}%
  \fi}
\makeatother

% v3.39.1: path-independent byte-reproducibility. LuaTeX derives the trailer
% /ID from the OUTPUT PATH, so two clones of the same source at different
% directories produced different sha256 (measured: the ONLY delta was the
% 62-byte /ID line). Bit 512 suppresses exactly that line — measured: Creator/
% Producer/CreationDate/ModDate and all content untouched. /ID is optional
% per ISO 32000 (required only for ENCRYPTED PDFs and PDF/A — neither is a
% target here). If a future workflow needs an ID (PDF/A, Acrobat revision
% matching), delete this line and it returns.
\pdfvariable suppressoptionalinfo \numexpr 512\relax

% v3.75: BOLD MATH IN BOLD HEADINGS. Heading text is \bfseries but math inside
% it stayed regular-weight CM — a hairline $V_T$ inside a heavy 24.8pt chapter
% title (measured: CMMI12 fl=6 next to FrankRuhlLibre-Bold fl=20; 155 live
% heading-with-math sites fleet-wide). \boldmath rides along wherever the class
% sets \bfseries for a heading; A/B measured: $V_T$ -> CMMIB10 fl=22, 0 missing
% characters (the bold CM faces ship with this TeX Live). Guarded per-macro so
% the article class (no \@makechapterhead) and future class changes stay safe.
% v4.0.1 (adoption fix): the SEARCH patterns above carry no parameter token.
% v4.0.0 made them specific by including the class's own #1/#2, and measured on
% TeX Live 2023 that never matches: a \patchcmd search pattern containing a
% parameter token cannot match while the call sits inside \@ifundefined{...}{}{...}
% (probed both #1 and ##1 -> FAIL; the same pattern at top level succeeds). Every
% one of these patches therefore failed, and because v4.0.0 also replaced the
% empty failure branch with \PackageError, the failure became FATAL: all four
% bundled example assets stopped building ("Cannot patch chapter-title math
% weight -> Fatal error, no output PDF"). The \PackageError branches are KEPT —
% a silent no-op here is what let v3.75's fix rot unnoticed — only the parameter
% tokens are dropped, so the patches match and a real class change still errors.
\makeatletter
\@ifundefined{@makechapterhead}{}{%
  \patchcmd{\@makechapterhead}{\huge\bfseries \@chapapp}
    {\huge\bfseries\boldmath \@chapapp}{}
    {\PackageError{hebrew-lualatex-pdf}{Cannot patch chapter-number math weight}{Unsupported document class}}%
  \patchcmd{\@makechapterhead}{\Huge \bfseries}
    {\Huge \bfseries\boldmath}{}
    {\PackageError{hebrew-lualatex-pdf}{Cannot patch chapter-title math weight}{Unsupported document class}}%
  \patchcmd{\@makeschapterhead}{\Huge \bfseries}
    {\Huge \bfseries\boldmath}{}
    {\PackageError{hebrew-lualatex-pdf}{Cannot patch starred chapter math weight}{Unsupported document class}}}
% v4.0.1: \part sits under a THREE-DEEP wrapper chain, so the class-based anchor
% is not in \@part at all. Measured with \show: hyperref rewrites \@part to
% `...\phantomsection\fi \H@old@part`, and nameref rewrites THAT to
% `\NR@gettitle{#1}\NR@part[{#1}]{#2}` — the class body ends up in \NR@part.
% v4.0.0 patched \@part and made failure fatal, which is why all seven bundled
% examples stopped building. \part is NOT a dead heading level here: 188 live uses
% across the fleet, so the bold-math fix is worth getting right rather than
% dropping. We walk the chain innermost-first and stop at the first hit. Failure is
% a WARNING, not an error: the chain is a property of hyperref/nameref versions, so
% a future package release must degrade to regular-weight part math — a cosmetic
% loss — instead of refusing to build. Chapter and section keep \PackageError,
% because those targets are in the CLASS and were verified to patch here.
\newif\iflakin@partnum \newif\iflakin@parttitle \newif\iflakin@sparttitle
\def\lakin@partnum#1#2{\iflakin@partnum\else
  \patchcmd{#1}{#2\bfseries \partname}{#2\bfseries\boldmath \partname}
    {\lakin@partnumtrue}{}\fi}
\def\lakin@parttitle#1#2{\iflakin@parttitle\else
  \patchcmd{#1}{#2 \bfseries}{#2 \bfseries\boldmath}{\lakin@parttitletrue}{}\fi}
\def\lakin@sparttitle#1#2{\iflakin@sparttitle\else
  \patchcmd{#1}{#2 \bfseries}{#2 \bfseries\boldmath}{\lakin@sparttitletrue}{}\fi}
\@ifundefined{c@chapter}{\def\lakin@pnum{\Large}\def\lakin@ptit{\huge}}%
                        {\def\lakin@pnum{\huge}\def\lakin@ptit{\Huge}}
\@ifundefined{NR@part}{}{\expandafter\lakin@partnum\expandafter\NR@part\expandafter{\lakin@pnum}%
                         \expandafter\lakin@parttitle\expandafter\NR@part\expandafter{\lakin@ptit}}
\@ifundefined{H@old@part}{}{\expandafter\lakin@partnum\expandafter\H@old@part\expandafter{\lakin@pnum}%
                            \expandafter\lakin@parttitle\expandafter\H@old@part\expandafter{\lakin@ptit}}
\@ifundefined{@part}{}{\expandafter\lakin@partnum\expandafter\@part\expandafter{\lakin@pnum}%
                       \expandafter\lakin@parttitle\expandafter\@part\expandafter{\lakin@ptit}}
\@ifundefined{NR@spart}{}{\expandafter\lakin@sparttitle\expandafter\NR@spart\expandafter{\lakin@ptit}}
\@ifundefined{H@old@spart}{}{\expandafter\lakin@sparttitle\expandafter\H@old@spart\expandafter{\lakin@ptit}}
\@ifundefined{@spart}{}{\expandafter\lakin@sparttitle\expandafter\@spart\expandafter{\lakin@ptit}}
\iflakin@partnum\else\PackageWarning{hebrew-lualatex-pdf}{Part-number math stays
  regular weight: no known \noexpand\part wrapper matched (hyperref/nameref
  version change). Cosmetic only}\fi
\iflakin@parttitle\else\PackageWarning{hebrew-lualatex-pdf}{Part-title math stays
  regular weight: no known \noexpand\part wrapper matched. Cosmetic only}\fi
% v4.12.0: THE PART NUMBER IN THE TOC WAS NOT DIRECTIONALLY SEPARATED FROM THE
% TITLE. report.cls/article.cls compose the part entry as ONE text stream --
% \addcontentsline{toc}{part}{\thepart\hspace{1em}#1} -- while EVERY other level
% goes through \numberline, which puts the number in an \hbox. A box is a bidi
% OBJECT and cannot join a neighbouring run; a bare "IV" can. So in the RTL TOC
% paragraph a Latin number and a Latin-INITIAL title merge into a SINGLE LTR run
% and are laid out left-to-right as a unit: the number lands on the wrong side of
% the row, and a mixed title is scrambled outright. Measured on \part{MCP ...}:
% the printed order was "<tail> V MCP" -- title fragment, number, title fragment.
% The chapter/section rows one line below, with the SAME Latin-initial titles,
% were correct on the same page. That contrast is the whole diagnosis.
% \en{} IS NOT A WORKAROUND, measured both ways on one fixture: it fixes a
% PURE-Latin title (\part{\en{Hooks}} renders correctly) and does NOT fix a MIXED
% one (\part{\en{MCP} <hebrew>} stays scrambled). babel's language switch changes
% font and hyphenation, not the bidi run structure -- so the defect is not
% reachable from the source side at all, which is why it belongs here.
% Boxing the number is the entire fix and is a STRICT NO-OP where the bug never
% bit: a Hebrew-titled document compiles BYTE-IDENTICAL before and after (md5,
% 3 clean runs), because there the neutral glue already sat between an L run and
% an R run. The PDF outline is byte-identical too -- hyperref's \pdfstringdef
% neutralises \mbox, so bookmark strings keep their part numbers.
% \numberline{\thepart} would also fix the direction, but it re-geometries every
% part row in the fleet and DROPS the number from the part bookmark (hyperref
% strips \numberline). \mbox changes nothing that was already right.
% The search pattern deliberately carries no # token: etoolbox cannot match a
% parameter token in a patch pattern (measured -- all three chain members
% reported "not patchable" while the pattern included #1). "\thepart \hspace
% {1em}" is brace-balanced and occurs exactly once in the wrapper body.
\newif\iflakin@parttocnum
\def\lakin@parttocnum#1{\iflakin@parttocnum\else
  \patchcmd{#1}{\thepart \hspace {1em}}{\mbox {\thepart }\hspace {1em}}
    {\lakin@parttocnumtrue}{}\fi}
\@ifundefined{NR@part}{}{\lakin@parttocnum\NR@part}
\@ifundefined{H@old@part}{}{\lakin@parttocnum\H@old@part}
\@ifundefined{@part}{}{\lakin@parttocnum\@part}
\iflakin@parttocnum\else\PackageWarning{hebrew-lualatex-pdf}{Part TOC number is
  not boxed: no known \noexpand\part wrapper matched. A part title that BEGINS
  with Latin will print its number on the wrong side of its TOC row}\fi
\patchcmd{\section}{\bfseries}{\bfseries\boldmath}{}
  {\PackageError{hebrew-lualatex-pdf}{Cannot patch section math weight}{Unsupported document class}}
\patchcmd{\subsection}{\bfseries}{\bfseries\boldmath}{}
  {\PackageError{hebrew-lualatex-pdf}{Cannot patch subsection math weight}{Unsupported document class}}
\patchcmd{\subsubsection}{\bfseries}{\bfseries\boldmath}{}
  {\PackageError{hebrew-lualatex-pdf}{Cannot patch subsubsection math weight}{Unsupported document class}}
\patchcmd{\paragraph}{\bfseries}{\bfseries\boldmath}{}{}
\patchcmd{\subparagraph}{\bfseries}{\bfseries\boldmath}{}{}
\@ifundefined{l@chapter}{}{%
  \patchcmd{\l@chapter}{\bfseries}{\bfseries\boldmath}{}{}}
\makeatother
\setcounter{tocdepth}{\lakittocdepth}
% <hllpdf:measured-toc>
% v4.16.0: Measure actual number labels, not a guessed digit count. Read only
% canonical contentsline records; never execute titles or addtocontents code
% during the measuring pass. Widths are local to this individual list and are
% available before its first row: no per-row drift and no extra aux pass.
\makeatletter
\newread\hllpdf@tocstream
\newdimen\hllpdf@toc@chapter
\newdimen\hllpdf@toc@section
\newdimen\hllpdf@toc@subsection
\newdimen\hllpdf@toc@subsubsection
\newdimen\hllpdf@toc@paragraph
\newdimen\hllpdf@toc@subparagraph
\newdimen\hllpdf@toc@figure
\newdimen\hllpdf@toc@table
\newdimen\hllpdf@tocpage
\newdimen\hllpdf@toc@sectionindent
\newdimen\hllpdf@toc@subsectionindent
\newdimen\hllpdf@toc@subsubsectionindent
\newdimen\hllpdf@toc@paragraphindent
\newdimen\hllpdf@toc@subparagraphindent
\def\hllpdf@tocdefaults{%
  \hllpdf@toc@chapter=2em
  \hllpdf@toc@section=3.2em
  \hllpdf@toc@subsection=3.9em
  \@ifundefined{c@chapter}{%
    \hllpdf@toc@subsubsection=3.2em
    \hllpdf@toc@paragraph=4.1em
    \hllpdf@toc@subparagraph=5em
  }{%
    \hllpdf@toc@subsubsection=4.1em
    \hllpdf@toc@paragraph=5em
    \hllpdf@toc@subparagraph=6em
  }%
  \hllpdf@toc@figure=2.3em
  \hllpdf@toc@table=2.3em
  \hllpdf@tocpage=\@pnumwidth\relax
}
\def\hllpdf@tocindents{%
  % Preserve the established shallow layout when its slots already fit.
  \hllpdf@toc@sectionindent=1.5em
  \@tempdima=\hllpdf@toc@chapter \advance\@tempdima by -2em
  \advance\hllpdf@toc@sectionindent by \@tempdima
  \hllpdf@toc@subsectionindent=\dimexpr
    \hllpdf@toc@sectionindent+\hllpdf@toc@section\relax
  \hllpdf@toc@subsubsectionindent=\dimexpr
    \hllpdf@toc@subsectionindent+\hllpdf@toc@subsection\relax
  \hllpdf@toc@paragraphindent=\dimexpr
    \hllpdf@toc@subsubsectionindent+\hllpdf@toc@subsubsection\relax
  \hllpdf@toc@subparagraphindent=\dimexpr
    \hllpdf@toc@paragraphindent+\hllpdf@toc@paragraph\relax
}
\def\hllpdf@tocfont@chapter{\bfseries\boldmath}
\def\hllpdf@tocfont@part{\large\bfseries\boldmath}
\def\hllpdf@tocentryfont{%
  \ifcsname hllpdf@tocfont@\hllpdf@toclevel\endcsname
    \csname hllpdf@tocfont@\hllpdf@toclevel\endcsname
  \fi}
\long\def\hllpdf@tocdiscard#1\hllpdf@tocstop{}
\def\hllpdf@tocscanrecord{%
  \@ifnextchar\contentsline{\hllpdf@tocscanline}{%
    \@ifnextchar\babel@toc{\hllpdf@tocscanlanguage}{\hllpdf@tocdiscard}}}
\long\def\hllpdf@tocscanlanguage\babel@toc#1#2#3\hllpdf@tocstop{%
  % Babel records are declarations, not arbitrary code. Track the language
  % without executing the original directive or changing the surrounding list.
  \def\hllpdf@toclanguage{#1}}
\long\def\hllpdf@tocscanline\contentsline#1#2#3#4#5\hllpdf@tocstop{%
  \def\hllpdf@toclevel{#1}%
  \ifcsname hllpdf@toc@#1\endcsname
    \hllpdf@tocscantitle#2\hllpdf@tocstop
  \fi
  % The page-number column needs its own measurement (including bold chapter
  % rows). Its extra clearance separates the page label from dot leaders.
  \setbox\@tempboxa=\hbox{\foreignlanguage{\hllpdf@toclanguage}{\hllpdf@tocentryfont #3\kern.3em}}%
  \ifdim\wd\@tempboxa>\hllpdf@tocpage
    \hllpdf@tocpage=\wd\@tempboxa
  \fi}
\def\hllpdf@tocscantitle{%
  \@ifnextchar\numberline{\hllpdf@tocscanlabel}{\hllpdf@tocdiscard}}
\long\def\hllpdf@tocscanlabel\numberline#1#2\hllpdf@tocstop{%
  \setbox\@tempboxa=\hbox{\foreignlanguage{\hllpdf@toclanguage}{\hllpdf@tocentryfont #1\kern.5em}}%
  \ifdim\wd\@tempboxa>\csname hllpdf@toc@\hllpdf@toclevel\endcsname
    \csname hllpdf@toc@\hllpdf@toclevel\endcsname=\wd\@tempboxa
  \fi}
\def\hllpdf@tocread{%
  \ifeof\hllpdf@tocstream\else
    \read\hllpdf@tocstream to\hllpdf@tocrecord
    \expandafter\hllpdf@tocscanrecord\hllpdf@tocrecord\hllpdf@tocstop
    \expandafter\hllpdf@tocread
  \fi}
\def\hllpdf@tocmeasure#1{%
  \makeatletter
  \edef\hllpdf@toclanguage{\languagename}%
  \hllpdf@tocdefaults
  \openin\hllpdf@tocstream=\jobname.#1\relax
  \hllpdf@tocread
  \closein\hllpdf@tocstream
  \hllpdf@tocindents
  \edef\@pnumwidth{\the\hllpdf@tocpage}%
  \@tempdima=\dimexpr\hllpdf@tocpage+1em\relax
  \ifdim\@tempdima>\@tocrmarg\relax
    \edef\@tocrmarg{\the\@tempdima}%
  \fi}
\let\hllpdf@originalstarttoc\@starttoc
\def\@starttoc#1{%
  \begingroup
    \def\hllpdf@tocext{#1}%
    \def\hllpdf@tocwanted{toc}%
    \ifx\hllpdf@tocext\hllpdf@tocwanted \hllpdf@tocmeasure{#1}\else
      \def\hllpdf@tocwanted{lof}%
      \ifx\hllpdf@tocext\hllpdf@tocwanted \hllpdf@tocmeasure{#1}\else
        \def\hllpdf@tocwanted{lot}%
        \ifx\hllpdf@tocext\hllpdf@tocwanted \hllpdf@tocmeasure{#1}\fi
      \fi
    \fi
    \hllpdf@originalstarttoc{#1}%
  \endgroup}
\hllpdf@tocdefaults
\hllpdf@tocindents
\@ifundefined{l@chapter}{}{%
  \patchcmd{\l@chapter}{\setlength\@tempdima{1.5em}}%
    {\setlength\@tempdima{\hllpdf@toc@chapter}}{}%
    {\PackageError{hebrew-lualatex-pdf}{Cannot size chapter TOC numbers}%
       {Unsupported contents layout: the chapter number slot was not found}}}
\renewcommand*\l@section{\@dottedtocline{1}%
  {\hllpdf@toc@sectionindent}{\hllpdf@toc@section}}
\renewcommand*\l@subsection{\@dottedtocline{2}%
  {\hllpdf@toc@subsectionindent}{\hllpdf@toc@subsection}}
\renewcommand*\l@subsubsection{\@dottedtocline{3}%
  {\hllpdf@toc@subsubsectionindent}{\hllpdf@toc@subsubsection}}
\renewcommand*\l@paragraph{\@dottedtocline{4}%
  {\hllpdf@toc@paragraphindent}{\hllpdf@toc@paragraph}}
\renewcommand*\l@subparagraph{\@dottedtocline{5}%
  {\hllpdf@toc@subparagraphindent}{\hllpdf@toc@subparagraph}}
\renewcommand*\l@figure{\@dottedtocline{1}{1.5em}{\hllpdf@toc@figure}}
\renewcommand*\l@table{\@dottedtocline{1}{1.5em}{\hllpdf@toc@table}}
\makeatother
% </hllpdf:measured-toc>
\setcounter{secnumdepth}{\lakitsecnumdepth}
% Public preamble-only caption-numbering policy. The historical within-major
% default is untouched unless this command is called. "continuous" removes the
% reset edge as well as the printed prefix, so Figure 1/2 and Table 1/2 stay
% unique across sections/chapters. Returning to "within-major" restores the
% original class-aware displays and native destination conventions. The policy
% never replaces \figcap/\tabcap or their measured link-placement mechanism.
\makeatletter
\let\lakin@majorfig\thefig
\let\lakin@majortbl\thetbl
\let\lakin@majorHfig\theHfig
\let\lakin@majorHtbl\theHtbl
\newcommand{\lakitcaptionnumbering}[1]{%
  \edef\lakin@numbering{#1}%
  \def\lakin@continuous{continuous}\def\lakin@withinmajor{within-major}%
  \ifx\lakin@numbering\lakin@continuous
    \@ifundefined{c@chapter}{%
      \@removefromreset{fig}{section}\@removefromreset{tbl}{section}%
    }{\@removefromreset{fig}{chapter}\@removefromreset{tbl}{chapter}}%
    \renewcommand{\thefig}{\arabic{fig}}\renewcommand{\thetbl}{\arabic{tbl}}%
    \def\theHfig{\arabic{fig}}\def\theHtbl{\arabic{tbl}}%
  \else\ifx\lakin@numbering\lakin@withinmajor
    \@ifundefined{c@chapter}{%
      \@removefromreset{fig}{section}\@removefromreset{tbl}{section}%
      \@addtoreset{fig}{section}\@addtoreset{tbl}{section}%
    }{%
      \@removefromreset{fig}{chapter}\@removefromreset{tbl}{chapter}%
      \@addtoreset{fig}{chapter}\@addtoreset{tbl}{chapter}}%
    \let\thefig\lakin@majorfig\let\thetbl\lakin@majortbl
    \let\theHfig\lakin@majorHfig\let\theHtbl\lakin@majorHtbl
  \else
    \PackageError{hebrew-english-latex-studio}{Unknown caption-numbering policy '#1'}%
      {Use continuous or within-major before begin{document}.}%
  \fi\fi}
\@onlypreamble\lakitcaptionnumbering
\makeatother
% Architecture readiness sentinel. Optional runtime layers require the complete
% core, so the marker is defined only after every mechanism/baseline definition.
% <hllpdf:architecture-ready>
\def\hllpdfcoreloaded{1}
% </hllpdf:architecture-ready>

% v4.22: registered file resources, not an assertion about embedded subset bytes.
% The log framing avoids long-path wrapping and preserves UTF-8 as hex.
% <hllpdf:font-inventory>
\directlua{
  local function inventory()
    local names = {}
    local unnamed = 0
    for id, item in font.each() do
      local path = item.filename
      if path and not (path == "") then names[path] = true else unnamed = unnamed + 1 end
    end
    local sorted = {}
    for path in pairs(names) do table.insert(sorted,path) end
    table.sort(sorted)
    texio.write_nl("log","HLLFONT|BEGIN|1")
    for index,path in ipairs(sorted) do
      local hex = path:gsub(".",function(c) return string.format(string.char(37) .. "02x",string.byte(c)) end)
      local count = math.ceil(string.len(hex)/48)
      texio.write_nl("log","HLLFONT|FILE|" .. index .. "|" .. count)
      for offset=1,string.len(hex),48 do texio.write_nl("log","HLLFONT|HEX|" .. hex:sub(offset,offset+47)) end
    end
    texio.write_nl("log","HLLFONT|END|" .. rawlen(sorted) .. "|" .. unnamed)
    texio.write_nl("log","")
  end
  luatexbase.add_to_callback("finish_pdffile",inventory,"hllpdf.runtime.fonts")
}
% </hllpdf:font-inventory>

%
% Supported class contract: report or article.  Alternate paper sizes must also
% be declared to build.sh through EXPECTED_PAPER; the PDF gate never infers them.
% Do NOT use these hooks as a theme mechanism: theme overlays must not own page
% geometry, class selection, or direction semantics. A declared presentation-policy
% overlay may intentionally change numbering visibility after the preamble.
% ============================================================================
% Runtime architecture guard: generated from scripts/architecture_contract.py.
% The readiness sentinel remains delayed until the END of this file.
% <hllpdf:architecture-guard>
\ifdefined\hllpdfcoreloaded
  \errmessage{hllpdf: preamble.tex loaded more than once}
  \expandafter\endinput
\fi
% </hllpdf:architecture-guard>

\ifdefined\lakitdocumentclass\else\def\lakitdocumentclass{report}\fi
\ifdefined\lakitclassoptions\else\def\lakitclassoptions{12pt,a4paper}\fi
\ifdefined\lakitgeometryoptions\else\def\lakitgeometryoptions{margin=2.3cm,headheight=14pt}\fi
\ifdefined\lakitlinespread\else\def\lakitlinespread{1.3}\fi
\ifdefined\lakittocdepth\else\def\lakittocdepth{1}\fi
\ifdefined\lakitsecnumdepth\else\def\lakitsecnumdepth{2}\fi
% Paragraph organization belongs to the document layout, independently of its
% language and theme. Preserve historical class dimensions unless explicitly
% selecting block paragraphs; never suppress indentation with an everypar hack.
\ifdefined\lakitparagraphstyle\else\def\lakitparagraphstyle{indented}\fi
\edef\lakitparagraphstyle{\lakitparagraphstyle}
\def\hllpdfparagraphblock{block}
\def\hllpdfparagraphindented{indented}
\ifx\lakitparagraphstyle\hllpdfparagraphblock\else
  \ifx\lakitparagraphstyle\hllpdfparagraphindented\else
    \errmessage{hllpdf: paragraph style must be block or indented}
  \fi
\fi
% v4.3.1: the public class hook has a deliberately narrow support contract.
% Before this guard, \def\lakitdocumentclass{book} silently compiled even though
% every heading/TOC regression test targets report/article; an unsupported shell
% could therefore look "supported" until a subtle class-specific defect shipped.
% Fail before \documentclass rather than letting an undocumented combination enter
% the mechanism core. Plain-TeX comparison is used because LaTeX packages are not
% available yet at this point.
\edef\lakitvalidatedclass{\lakitdocumentclass}
\def\lakitreportclass{report}
\def\lakitarticleclass{article}
\ifx\lakitvalidatedclass\lakitreportclass
\else\ifx\lakitvalidatedclass\lakitarticleclass
\else
  \errmessage{hllpdf: unsupported class \lakitdocumentclass; use report or article}
\fi\fi
\let\lakitvalidatedclass\relax
\let\lakitreportclass\relax
\let\lakitarticleclass\relax
% Language intent is independent of document class, theme and content domain.
% The two dominant-language profiles differ in authoring intent, not direction.
% A balanced bilingual document must explicitly choose a navigation/base locale.
% <hllpdf:language-profile>
\providecommand{\lakitlanguageprofile}{hebrew}
\edef\lakitlanguageprofile{\lakitlanguageprofile}
\def\hllpdfhebrew{hebrew}\def\hllpdfenglish{english}
\def\hllpdfhebrewmixed{hebrew-with-english}
\def\hllpdfenglishmixed{english-with-hebrew}
\def\hllpdfbilingual{bilingual}
\let\hllpdfexpectedlanguage\relax
\ifx\lakitlanguageprofile\hllpdfhebrew
  \let\hllpdfexpectedlanguage\hllpdfhebrew
\else\ifx\lakitlanguageprofile\hllpdfhebrewmixed
  \let\hllpdfexpectedlanguage\hllpdfhebrew
\else\ifx\lakitlanguageprofile\hllpdfenglish
  \let\hllpdfexpectedlanguage\hllpdfenglish
\else\ifx\lakitlanguageprofile\hllpdfenglishmixed
  \let\hllpdfexpectedlanguage\hllpdfenglish
\else\ifx\lakitlanguageprofile\hllpdfbilingual
  \ifdefined\lakitmainlanguage\else
    \errmessage{hllpdf: bilingual profile requires explicit lakitmainlanguage}
  \fi
\else
  \errmessage{hllpdf: unsupported language profile \lakitlanguageprofile}
\fi\fi\fi\fi\fi
\ifdefined\lakitmainlanguage
  \edef\lakitmainlanguage{\lakitmainlanguage}
\else
  \let\lakitmainlanguage\hllpdfexpectedlanguage
\fi
\ifx\lakitmainlanguage\hllpdfhebrew\else
  \ifx\lakitmainlanguage\hllpdfenglish\else
    \errmessage{hllpdf: main language must be hebrew or english}
  \fi
\fi
\ifx\hllpdfexpectedlanguage\relax\else
  \ifx\lakitmainlanguage\hllpdfexpectedlanguage\else
    \errmessage{hllpdf: main language conflicts with language profile}
  \fi
\fi
\typeout{HLLLANG|\lakitlanguageprofile|\lakitmainlanguage}
% </hllpdf:language-profile>
\documentclass[\lakitclassoptions]{\lakitdocumentclass}
\usepackage[no-math]{fontspec}        % (1) ★ no-math — THE critical line (see recipe.md)
% layout=lists asks Babel to adapt list parshape when paragraph and page
% directions differ. Without it, minority-language list bodies lose their
% indentation and labels hang outside the text block, despite a clean log.
% The upstream mechanism preserves ordinary same-direction list geometry.
\usepackage[bidi=basic,layout=lists]{babel}
\ifx\lakitmainlanguage\hllpdfenglish
  \babelprovide[import]{hebrew}
  \babelprovide[import,main]{english}
\else
  \babelprovide[import,main]{hebrew}
  \babelprovide[import]{english}
\fi
% v4.22: bind canonical families and fallbacks to the selected installation.
% Lookup by basename or family alone is not proof of the file actually used.
% <hllpdf:font-binding>
\directlua{
  local dir = os.getenv("HEBREW_LATEX_FONT_DIR")
  if not dir or dir == "" then
    local data = os.getenv("XDG_DATA_HOME")
    if not data or data == "" or not (data:sub(1,1) == "/") then
      data = (os.getenv("HOME") or "") .. "/.local/share"
    end
    dir = data .. "/fonts/hebrew-lualatex-pdf"
  end
  if not (dir:sub(1,1) == "/") then dir = lfs.currentdir() .. "/" .. dir end
  if dir:find("[",1,true) or dir:find("]",1,true) or dir:find(string.char(0),1,true)
     or dir:find(string.char(10),1,true) or dir:find(string.char(13),1,true) then
    tex.error("hllpdf: unsupported character in font directory")
  end
  if not (lfs.attributes(dir,"mode") == "directory") then
    tex.error("hllpdf: font directory is missing")
  end
  hllpdf_font_dir = dir:gsub("/+$", "") .. "/"
  tex.sprint(string.char(92) .. "def" .. string.char(92) .. "hllpdffontdir{")
  tex.sprint(-2,hllpdf_font_dir)
  tex.sprint("}")
  luaotfload.add_fallback("hebfb",{"[" .. hllpdf_font_dir .. "DejaVuSerif.ttf]:mode=node;"})
  luaotfload.add_fallback("hebfbbf",{"[" .. hllpdf_font_dir .. "DejaVuSerif-Bold.ttf]:mode=node;"})
}
% </hllpdf:font-binding>
% v3.22: Frank Ruhl Libre ships NO italic face. Instead of silently-upright
% \emph (v<=3.20) or bold-as-emphasis (v3.21), synthesize a slanted face with
% FakeSlant — real, VISIBLE italic emphasis with correct \emph nesting
% semantics (inner \emph toggles back upright). Verified empirically at
% 300dpi: slant clearly visible, 0 errors, 0 missing glyphs.
\babelfont{rm}[Path={\hllpdffontdir},RawFeature={fallback=hebfb},
  BoldFont={FrankRuhlLibre-Bold.ttf},
  BoldFeatures={RawFeature={fallback=hebfbbf}},   % v3.56 — see hebfbbf above
  ItalicFont={FrankRuhlLibre-Regular.ttf}, ItalicFeatures={FakeSlant=0.22},
  % Bold-italic uses the weight-matched fallback and a synthetic slant.
  BoldItalicFont={FrankRuhlLibre-Bold.ttf}, BoldItalicFeatures={RawFeature={fallback=hebfbbf},FakeSlant=0.22}
  ]{FrankRuhlLibre-Regular.ttf}
% Sans fallbacks preserve family and weight for the few missing codepoints.
\directlua{luaotfload.add_fallback("hebfbsans",{"[" .. hllpdf_font_dir .. "DejaVuSans.ttf]:mode=node;"})}
\directlua{luaotfload.add_fallback("hebfbsansbf",{"[" .. hllpdf_font_dir .. "DejaVuSans-Bold.ttf]:mode=node;"})}
\babelfont{sf}[Path={\hllpdffontdir},RawFeature={fallback=hebfbsans},
  BoldFont={Heebo-Bold.ttf},
  BoldFeatures={RawFeature={fallback=hebfbsansbf}}]{Heebo-Regular.ttf}
% v3.24: Scale=MatchLowercase — DejaVu Sans Mono has an outsized x-height, so
% unscaled inline code visually SHOUTS next to the delicate Frank Ruhl body
% (measured on a 197-page build: scaled, a code atom sits in the line at the
% body's weight, and the document tightened by ~5%). Applies to listings too.
\babelfont{tt}[Path={\hllpdffontdir},RawFeature={fallback=hebfb},Scale=MatchLowercase,
  BoldFont={DejaVuSansMono.ttf},
  BoldFeatures={RawFeature={fallback=hebfb},FakeBold=1.5},
  ItalicFont={DejaVuSansMono.ttf},
  ItalicFeatures={RawFeature={fallback=hebfb},FakeSlant=0.2},
  BoldItalicFont={DejaVuSansMono.ttf},
  BoldItalicFeatures={RawFeature={fallback=hebfb},FakeBold=1.5,FakeSlant=0.2}
  ]{DejaVuSansMono.ttf}

% Explicit multilingual/English profiles use real Latin faces and italics.
% The legacy/default Hebrew profile retains its established measured typography.
% Latin Modern is provided by the TeX distribution, not the pinned font bundle;
% recorder/runtime evidence reports it as an additional toolchain resource.
\ifx\lakitlanguageprofile\hllpdfhebrew\else
  \babelfont[english]{rm}[Ligatures=TeX,RawFeature={fallback=hebfb},
    BoldFeatures={RawFeature={fallback=hebfbbf}},
    BoldItalicFeatures={RawFeature={fallback=hebfbbf}}]{Latin Modern Roman}
  \babelfont[english]{sf}[Ligatures=TeX,RawFeature={fallback=hebfbsans},
    BoldFeatures={RawFeature={fallback=hebfbsansbf}},
    BoldItalicFeatures={RawFeature={fallback=hebfbsansbf}}]{Latin Modern Sans}
\fi

\usepackage{amsmath,amssymb,mathtools}
% amsthm supplies theorem/proof machinery. Environment names remain
% document-specific; the generic proof label is localized below.
\usepackage{amsthm}
% Proof names follow the surrounding prose, including return from an island.
\addto\extrashebrew{\renewcommand{\proofname}{הוכחה}}
\addto\extrasenglish{\renewcommand{\proofname}{Proof}}
\ifx\lakitmainlanguage\hllpdfenglish
  \renewcommand{\proofname}{Proof}
\else
  \renewcommand{\proofname}{הוכחה}
\fi
% v3.24.1: equation tags under babel-bidi. A tag re-emitted by mathtools'
% showonlyrefs inside align inherits the RTL base and its parens MIRROR —
% ")2.1(" instead of "(2.1)" (plain align/equation tags are unaffected;
% verified on a minimal case + a 198-page build). LTR-isolate at the tag-BOX
% level: mathtools' show/hide + numbering logic is untouched. Do NOT override
% \tagform@ instead — that clobbers the showonlyrefs mechanism (measured:
% every tag shows, numbering breaks).
% v3.28.1: the v3.24.1 wrap must apply ONLY to tags of display environments.
% Inline \eqref in Hebrew PROSE also goes through \maketag@@@, and there the
% english-wrap MIRRORS the parens — ")1(" — while the unwrapped original is
% correct (proven on a 2x2 isolation matrix: patch x showonlyrefs; regressed
% every prose \eqref since v3.24.1, caught by the render gate's R4). The two
% call sites are indistinguishable by \ifmmode (tag boxes are text mode), so
% the display environments themselves raise a flag; \eqref in prose sees it
% false and gets the original untouched \maketag@@@.
\usepackage{etoolbox}

% babel bidi=basic (including upstream 26.9) assumes every automatic
% discretionary has a one-glyph replacement. LuaTeX can legally produce one
% with replace=nil after a hyphenated Latin value; Babel then indexes nil and the
% whole build dies. Do the narrow upstream-compatible correction: while Babel
% classifies the list, present only that malformed disc as an ordinary
% discretionary, then restore its subtype. Babel still reorders the ENTIRE box;
% there is no pcall and no successful fallback to an unreordered node list.
% If Babel throws for any other reason, the build remains fatal. The synthetic
% nil-disc probe and the real table gallery in test_all.sh pin both layers.
\AtBeginDocument{\directlua{
  if Babel and Babel.bidi and not Babel.heb_nil_disc_compat then
    Babel.heb_nil_disc_compat = true
    local original = Babel.bidi
    local disc_id = node.id('disc')
    Babel.bidi = function(head, ispar, hdir)
      local hidden = {}
      for item in node.traverse_id(disc_id, head) do
        if item.subtype == 2 and (not item.replace or not item.replace.char) then
          table.insert(hidden, item)
          item.subtype = 0
        end
      end
      local result = original(head, ispar, hdir)
      for _, item in ipairs(hidden) do item.subtype = 2 end
      return result
    end
  end
}}
\makeatletter
\newif\ifHEBdisptag
\let\ORIGmaketag@@@\maketag@@@
\def\maketag@@@#1{\ifHEBdisptag\ORIGmaketag@@@{\foreignlanguage{english}{#1}}\else\ORIGmaketag@@@{#1}\fi}
\@for\HEB@env:=equation,equation*,align,align*,gather,gather*,multline,multline*,flalign,flalign*,alignat,alignat*,eqnarray,eqnarray*\do{%
  \expandafter\AtBeginEnvironment\expandafter{\HEB@env}{\HEBdisptagtrue}}
\makeatother
% English page + Hebrew prose: Babel's equation-only tag flip uses a primitive
% leqno path that shifts the formula to the left edge and the tag into the body
% of an English page. Align's full-width alignment is unaffected. Give only this
% display a physical LTR container and native amsmath eqno/leqno. The engine then
% owns real tag width, collision handling and labels; no proxy tag box hides ink
% from its fit calculation. Babel's saved \bbl@eqnodir still restores Hebrew text
% inside math, and the prose locale/math direction are never switched.
% Close the prior RTL paragraph before normalization. After the environment's
% group restores its directions, close TeX's empty resumed display paragraph and
% use native LaTeX continuation handling: same-paragraph prose has no indent,
% while an authored blank line starts an ordinary indented RTL paragraph.
% Register after Babel's setup; remain local to equation(*) in a main-English
% document's actual RTL prose. Native Hebrew, ordinary English and explicit
% fleqn retain their original paths. No page/display dimensions are assigned.
% <hllpdf:minority-rtl-equation-geometry>
\makeatletter
\newcommand{\HEB@minorityprosecontinuation}{\par\@doendpe}
\newcommand{\HEB@minorityequationgeometry}{%
  \ifx\lakitmainlanguage\hllpdfenglish
    \ifnum\pardirection=\@ne
      \if@fleqn\else
        \par\aftergroup\HEB@minorityprosecontinuation
        \iftagsleft@\let\veqno\@@leqno
        \else\let\veqno\@@eqno\fi
        \bodydir TLT\pardir TLT\textdir TLT\relax
      \fi
    \fi
  \fi}
\AtBeginDocument{%
  \AddToHook{env/equation/begin}{\HEB@minorityequationgeometry}%
  \AddToHook{env/equation*/begin}{\HEB@minorityequationgeometry}}
\makeatother
% </hllpdf:minority-rtl-equation-geometry>
% v3.63: \shortintertext rendered Hebrew MIRROR-REVERSED and flush LEFT — an entire
% commentary line inside an align printed backwards. babel's bidi=basic patches
% \intertext@ only (luababel.def), and mathtools' \shortintertext uses the identical
% \vbox{\normalbaselines … \noindent\ignorespaces#1\par} skeleton but is never
% patched, so its paragraph keeps pardir/textdir LTR: no bidi reordering (logical
% order on the page) and flush-left. Measured against \intertext with the same
% Applying babel's own replacement to \shortintertext@ makes the two identical.
% \h{}/\mbox fix the letter order but leave the line flush LEFT, so they are not
% substitutes. Guarded + \AtBeginDocument because mathtools \let-copies at load time.
% NB \ExplSyntaxOn is REQUIRED: mathtools' internals carry expl3 names, and without
% the expl3 catcodes the `_` is read as a subscript and the patch detonates far away
% with "! Illegal unit of measure (pt inserted)" at the \end{align} (measured).
\makeatletter
\AtBeginDocument{%
  \@ifundefined{shortintertext@}{}{%
    \ExplSyntaxOn
    \bbl@sreplace\shortintertext@{\normalbaselines}%
      {\normalbaselines\ifx\bbl@eqnodir\relax\else\bbl@pardir\@ne\bbl@eqnodir\fi}%
    \ExplSyntaxOff}}
\makeatother
% geometry inherits the canonical A4 sheet from \documentclass; build.sh then
% verifies every rendered MediaBox and CropBox, so a silent Letter fallback is
% impossible to publish as a clean default build.
\usepackage[\lakitgeometryoptions]{geometry}
\linespread{\lakitlinespread}       % layout hook; canonical default remains 1.3
\ifx\lakitparagraphstyle\hllpdfparagraphblock
  \setlength{\parindent}{0pt}
  \setlength{\parskip}{6pt plus 1pt minus 1pt}
\fi
% These dimensions remain ordinary scoped TeX lengths: authors can customize
% them after core input. Language switches inherit them; lists keep their own
% indentation and narrow reference layouts establish their local paragraph gaps.
\typeout{HLLPAR|\lakitparagraphstyle}
\usepackage{xcolor}
\usepackage{mdframed}
\usepackage{listings}                 % code listings (assembly / C / pseudo-code)
\usepackage{tikz}
\usetikzlibrary{arrows.meta,positioning,calc,shapes.geometric,fit,backgrounds}
\usepackage{pgfplots}\pgfplotsset{compat=1.18}
\usepackage{quantikz}                 % quantum circuits (remove if unused)
\usepackage{booktabs}
\usepackage{array}
\usepackage{float}                    % [H] = place table exactly here (reference doc: no drifting)
\usepackage{enumitem}
\usepackage{needspace}

% ============================================================================
% ★★ v3.4 — THE DIRECTION-ISOLATION TRIO (the heart of mixing Hebrew + Latin)
% Mental model:
%   \h{...}    Hebrew (RTL) island      — Hebrew text inside an LTR context
%                                          (TikZ nodes, the english `tikzpic`).
%   \en{...}   English (LTR) isolate    — Latin text / parentheses inside a
%                                          Hebrew (RTL) context. KEEPS PARENS
%                                          AND ORDER CORRECT.
%   \code{...} \en + monospace          — code inside Hebrew / inside a frametitle.
% CRITICAL: never wrap Latin in \h{} (it is for HEBREW). Inside a TikZ node,
% Latin in \h{} renders REVERSED (Datapath->htapataD, CPU->UPC) and parens
% mirror. Isolate every Latin/paren chunk with \en{} or \code{}. See
% references/figures-boxes-listings.md.
% ---- TEXT INSIDE MATH (v3.8) ----------------------------------------------
% Put ANY text that sits inside math in a \text{}. The ONLY thing that is
% invalid in math is a bare \h/\en (= \foreignlanguage) used DIRECTLY, outside
% \text -> "\rmfamily invalid in math mode". Inside \text:
%   Hebrew label           -> \text{...}            (bare Hebrew renders fine;
%                              \h is optional — BUT NOT ALWAYS HARMLESS (v3.67):
%                              if the \h group OPENS with a nested island
%                              ($..$, \en{}, \code{}, \texttt{}, \mbox{}) the
%                              WHOLE group loses bidi and the Hebrew prints in
%                              LOGICAL order = gibberish. Measured:
%                                \text{\h{$f$ קבועה}}  -> reads "העובק"  BROKEN
%                                \text{\h{קבועה $f$}}  -> correct
%                                \text{$f$ קבועה}      -> correct (no \h)
%                              Hebrew first inside \h, or drop the \h. Flagged
%                              H-OPENS-WITH-LTR-GROUP.)
%   plain Latin word        -> \text{...}            (no marker needed)
%   Latin WITH ( ) or [ ]   -> \text{\en{...}}       (\en forces LTR; a plain
%                              \text mirrors the parens in RTL -> "Spec (TNR(" )
%   f(Hebrew/text arg)      -> \text{f}(\text{...})  -- keep the parens at MATH
%                              level, NOT \text{f(arg)} (parens inside the box mirror)
%   code in math            -> \mathtt{...}  (or \texttt{}, math-guarded below)
% NEVER write \en{...} or \h{...} straight into $...$ / \[..\] (outside \text);
% always via \text{}. (\mbox{...} also works but does not scale in scripts.)
% ============================================================================
\newcommand{\h}[1]{\foreignlanguage{hebrew}{#1}}
% v3.42.1 — \en islands do NOT hyphenate (auto OR at explicit hyphens). English hyphenation patterns were
% tearing identifiers and terms ACROSS THE RTL SEAM ("IO Regis-/ters",
% "In-/put", "Band-/width", "Per-/formance", "Multi-/ple" — all shipped on
% real pages; probe-proven: "ALUS-/rcB", "Mem-/ToReg"). An island is a label
% or a term, not English prose — a hyphen tear there is never right. A/B on
% a full 197pp lakit: identical overfull sets, same page count, ~11
% paragraphs healed, one mild underfull (badness 3525) that replaced a
% mutilated "Regis-/ters". Long English PROSE belongs in the englishprose
% environment below, which restores explicit-hyphen breaks while its paragraph
% is still active. A raw otherlanguage block inherits the global prohibition.
% (\l@nohyphenation needs \csname:
% '@' is not a letter here.)
% v3.43.2 — \exhyphenpenalty moved OUT of \en to a GLOBAL assignment below.
% Render-proven: a group-local \exhyphenpenalty is DEAD CODE — TeX reads the
% line-breaking parameters at \par time, after the \en group has closed, so
% \en{write-through} / plain not-taken still tore at the "-" ("write-/
% through" shipped on a real exam page). \language is different: it is
% recorded in the node list as text is added, so the in-group nohyphenation
% assignment above DOES work (and stays). The global penalty was A/B-measured
% on a 197pp lakit: identical overfull set (25), identical page count, the
% not-/taken tear healed. Hebrew is unaffected in practice — Hebrew compounds
% use maqaf (־), not ASCII "-".
\newcommand{\en}[1]{\foreignlanguage{english}{\language=\csname l@nohyphenation\endcsname #1}}
% Babel runs extras locally for inline islands and paragraph language changes.
% An English main document must not inherit the Hebrew no-hyphen paragraph rule.
\addto\extrashebrew{\exhyphenpenalty=10000\relax}
\addto\extrasenglish{\exhyphenpenalty=50\relax}
\ifx\lakitmainlanguage\hllpdfenglish
  \exhyphenpenalty=50\relax
\else
  \exhyphenpenalty=10000\relax
\fi
% Full English paragraphs need ordinary explicit-hyphen breaks. Use this
% environment instead of a raw otherlanguage block; the paragraph ends while
% the local penalty is still active.
\newenvironment{englishprose}
  {\par\begin{otherlanguage}{english}\exhyphenpenalty=50\relax}
  {\par\end{otherlanguage}}
\newenvironment{hebrewprose}
  {\par\begin{otherlanguage}{hebrew}\exhyphenpenalty=10000\relax}
  {\par\end{otherlanguage}}
% Use for shell text whose locale must not depend on the last body paragraph.
\DeclareRobustCommand{\doclang}[1]{\foreignlanguage{\lakitmainlanguage}{#1}}

% ============================================================================
% ★★ — fixes for defects that are INVISIBLE in the log and in the checks
% and only surface on a VISUAL RENDER of the PDF. (Each verified empirically.)
% ----------------------------------------------------------------------------
% v3.22: \emph is now REAL slanted emphasis (FakeSlant italic face defined on
% the \babelfont{rm} line above), replacing the v3.21 \emph->\textbf hack: it
% preserves \emph's toggle semantics and keeps bold free for its own role.
% Prefer bold-as-emphasis for your doc? one line: \renewcommand{\emph}[1]{\textbf{#1}}
% page number: force the digits LTR so they never reverse (18 -> 81) in the RTL
%   footer. That reversal is a NON-DETERMINISTIC bidi glitch (observed once in
%   ~80 pp); this LuaLaTeX build is not byte-reproducible, so the defensive
%   LTR wrap is the only reliable guard. Latin digits are always LTR anyway.
\renewcommand{\thepage}{\foreignlanguage{english}{\arabic{page}}}
% v3.70: \pagenumbering \gdef's \thepage and silently DESTROYS the wrap above —
% after \pagenumbering{arabic} the folio guard is gone for the whole rest of the
% document. Two collections discovered this independently and hand-restore the
% wrap in main.tex with the same comment; that knowledge now lives here instead.
% The wrapper re-applies the LTR wrap after the kernel call, for any style.
\makeatletter
\let\lakin@orig@pagenumbering\pagenumbering
\renewcommand{\pagenumbering}[1]{%
  \lakin@orig@pagenumbering{#1}%
  \renewcommand{\thepage}{\foreignlanguage{english}{\csname @#1\endcsname\c@page}}}
\makeatother
% emergencystretch: Hebrew does not hyphenate, so an unbreakable LTR island
%   (\en{...} / an inline math run) landing at the margin can leave a small
%   (~2-13pt) overfull with no legal breakpoint. This rescues ONLY those
%   inline-prose tie-breaks; it is a ceiling used only on offending lines and
%   does NOT touch display-math / listing / table overflow (those still surface
%   in the log for the content-level overflow pass). Verified: baseline overfull
%   -> 0 with this set, with no change to pagination or underfull.
\setlength{\emergencystretch}{3em}
% v3.40: inline math (almost) never breaks across lines. TeX's default
%   (binop 700 / rel 500) lets an inline run split at an operator —
%   $(1.f)\cdot 2^{e}$ shipped as "(1.f)·" ending one line and "2^e" opening
%   the next (reader-reported). In LTR that split is merely ugly; under an
%   RTL base the carried operator lands on the visually WRONG side of the
%   break and the formula reads backwards.
%   WHY 10000 (v3.45.0 — the display-first resolution). Inline math is
%   FORBIDDEN from breaking, so a torn island under an RTL base is
%   structurally impossible. The v3.40 cycle used 9999 ("break only as a last
%   resort") because a bare 10000 left ~30-39 WIDE inline equations
%   overflowing (worst 90.7pt) — but 9999's "last resort" fires OFTEN in
%   tight, non-hyphenating Hebrew: a reader found torn islands across the
%   fleet (~150 candidates; e.g. "P(y |" ending one line and "x_0)" opening
%   the next, read backwards). The step the 9999 cycle missed: every one of
%   those ~30 overflows is a wide equation that BELONGS in display. Convert
%   them to display FIRST, then 10000 gives 0 tears AND 0 overflows (measured,
%   12/12 fleet). Doctrine: inline math is ATOMIC; a wide equation goes in a
%   display, never inline. See content-style ("inline math is atomic").
\binoppenalty=10000 \relpenalty=10000
% v3.38: uniform sentence spacing. TeX's default (\nonfrenchspacing, sfcode
% 3000) is the OLD Anglo-American wide-space-after-period convention — Hebrew
% typography never had it, and modern Latin practice (Chicago et al.) dropped
% it too. Measured on real pages: post-period gaps up to 1.9-2.1x the line's
% mean word gap — the dominant "rivers" source in justified Hebrew prose.
% Survives babel language switches (probed: sfcode stays french inside \en{}
% islands and otherlanguage blocks).
\frenchspacing
% v3.38: widow/orphan protection. LaTeX's default penalties (150/150/50)
% allow a lone line at a page top/bottom almost freely. 10000 = forbid;
% this is SAFE here precisely because the document is \raggedbottom (the
% demoted line just moves on, no vertical stretch fight) — under flushbottom
% these values would cause underfull pages. raggedbottom is the report-class
% default; made explicit so a future class/option change cannot silently
% flip it and turn the 10000s harmful.
\raggedbottom
\widowpenalty=10000 \clubpenalty=10000 \displaywidowpenalty=10000
% em-dash: --- does NOT become — under Frank Ruhl (the font lacks tlig; verified
%   the width is identical with/without Ligatures=TeX, while the mechanism works
%   on Latin Modern). TYPE THE CHARACTER — (U+2014) DIRECTLY; likewise – (U+2013).
%   check_bidi_figures.py flags a --/--- that is adjacent to Hebrew.
% ============================================================================

% ---- STRUCTURAL DIRECTION-SAFETY for \texttt (v3.7) -------------------------
% A parenthesis, bracket, or $ INSIDE \texttt{} mirrors/migrates in RTL prose
% (MIPS 0($s0)->"(s0$)0", $rs->"rs$", arr[i]->"[arr[i") because \texttt sets the
% font but NOT the direction. Redefining \texttt to LTR-isolate its content fixes
% the WHOLE class at once (parens, brackets, $) so plain \texttt{lw $t0,0($s0)}
% is safe — no need to remember \code{}. Guarded by \ifmmode because
% \foreignlanguage is invalid in math (there, plain monospace is already LTR).
% NB (corrected v3.56 — the old note here said "stays tofu", which is WRONG):
% \texttt{Hebrew} is a FATAL build failure. The tt family is DejaVu Sans Mono
% (zero Hebrew) and the tt fallback hebfb is DejaVuSerif (also zero Hebrew), so
% the subsetter gets an empty glyph set and LuaTeX dies with
%   "! error: (file .../DejaVuSansMono.ttf) (type 2): there are no glyphs in the
%    subset"  -> no output PDF.
% The message names a FONT FILE, never Hebrew/\texttt/the source line. And it is
% non-local: add an unrelated Latin \code{} elsewhere and the fatal degrades into
% a SILENT DELETION of the Hebrew word. Put Hebrew in \h{}/normal text, never in
% \texttt{}/\code{}; flagged HEBREW-IN-MONO. To disable, delete this block.
\let\origtexttt\texttt
\protected\long\def\texttt#1{\ifmmode\origtexttt{#1}\else\foreignlanguage{english}{\origtexttt{#1}}\fi}
% v4.8.0: defer the direction test until TeX has inserted an alignment cell's
% math template. Unprotected expansion here tested \ifmmode too early at the
% FIRST token of aligned/array/cases/matrix cells. Engine-level \protected also
% preserves the original \texttt implementation without the self-recursion
% observed when retrofitting \DeclareRobustCommand to its saved alias.
\protected\long\def\code#1{\ifmmode\origtexttt{#1}\else\foreignlanguage{english}{\origtexttt{#1}}\fi}

% ---- \codeline — display-code line (v3.24): the code analogue of display math.
% Reader-friction rule (content-style.md §5b): inline \code{} is for a SINGLE
% ATOM only (one register/variable/mnemonic, a tight no-space expression).
% A multi-token expression — an instruction with operands, a spaced assignment,
% a register-binding list — zigzags the eye RTL<->LTR mid-sentence and, being
% an unbreakable island, line-breaks badly. Give it its own centered line with
% display-math vertical rhythm. \centerline = one box => the expression never
% wraps; \nopagebreak keeps it glued to its introducing sentence.
%   ...the Hebrew sentence ends with a colon:
%   \codeline{lw \$t1, -4(\$t0)}
%   and the prose continues (escape dollars as \$).
% v3.66 (first line of the body): flush a PENDING \item label before the \par.
% LaTeX defers the label box (\@labels) until the next \everypar, and an \hbox in
% vertical mode never fires \everypar — so `\item \codeline{...}` printed the code
% at y=69.1 and the item's own "1." BELOW it at y=103.6 (the number appears after
% its content, attached to nothing). \if@inlabel\leavevmode\fi is CONDITIONAL on
% purpose: an unconditional \leavevmode would open an empty paragraph line when
% \codeline follows ordinary prose. Verified both ways — in a list the label now
% sits on the code's line at its normal x, and after prose the geometry is
% byte-identical to before.
\makeatletter
\newcommand{\codeline}[1]{%
  \if@inlabel\leavevmode\fi
  \par\nopagebreak
  \vspace{\abovedisplayskip}\nopagebreak
  % v3.62: was \centerline{...}, i.e. \hbox to\hsize{\hss #1 \hss}. \hss has
  % INFINITE SHRINK, so TeX could always "fit" the box and NEVER reported an
  % Overfull — the same \hss mechanism as BLD-27's \item label. Measured: a wide
  % \codeline in a defbox painted its first token ON the box's left border bar and
  % sat 9.5pt outside the frame interior on both sides, with 0 Overfull, 0 Underfull
  % and all four gates PASS; in bare prose it reached 21.0pt into BOTH margins.
  % \hfil has no shrink, so the overflow is reported and build.sh's ovf_vis/ovf_big
  % counters see it; \linewidth (not \hsize) makes it box-aware, so a \codeline
  % inside a colored box centres on the box interior. Verified: a FITTING \codeline
  % keeps byte-identical glyph positions, and the box-crossing case now emits
  % "Overfull \hbox (19.03pt too wide)". A wide \codeline should be split in two.
  \hbox to\linewidth{\hfil\code{#1}\hfil}%
  \nopagebreak\vspace{\belowdisplayskip}}
\makeatother

% ---- a thin decorative rule helper ----
\newcommand{\sep}{\par\vspace{2pt}\noindent\textcolor{cRule}{\rule{\linewidth}{0.6pt}}\par\vspace{4pt}}

% ---- colors ----
\definecolor{cBlue}{HTML}{1F4E79}   \definecolor{cBlueBg}{HTML}{EAF1F8}
\definecolor{cGreen}{HTML}{2E6B4F}  \definecolor{cGreenBg}{HTML}{EAF4EE}
\definecolor{cAmber}{HTML}{9C6A00}  \definecolor{cAmberBg}{HTML}{FBF3E2}
\definecolor{cPurple}{HTML}{5B2C6F} \definecolor{cPurpleBg}{HTML}{F3ECF7}
\definecolor{cTeal}{HTML}{0E6E6E}   \definecolor{cTealBg}{HTML}{E6F4F4}
\definecolor{cRed}{HTML}{8E2B2B}    \definecolor{cRedBg}{HTML}{FBEAEA}
\definecolor{cGray}{HTML}{555555}   \definecolor{cRule}{HTML}{1F4E79}

% ---- math operators ----
\DeclareMathOperator{\Tr}{Tr}
\DeclareMathOperator{\rank}{rank}
\DeclareMathOperator{\diag}{diag}
\DeclareMathOperator{\sgn}{sgn}

% ---- bra-ket (quantikz defines \ket/\bra/\braket; add the rest safely) ----
\providecommand{\ket}[1]{\left|#1\right\rangle}
\providecommand{\bra}[1]{\left\langle#1\right|}
\providecommand{\braket}[2]{\left\langle#1\middle|#2\right\rangle}
\newcommand{\dyad}[1]{\left|#1\right\rangle\!\left\langle#1\right|}
% v3.21: general outer product |a><b| and expectation value <A> — agents reach
% for these by reflex (esp. \expval); quantikz defines neither, so their absence
% threw "Undefined control sequence" in a real build. \providecommand = safe.
\providecommand{\ketbra}[2]{\left|#1\right\rangle\!\left\langle#2\right|}
\providecommand{\expval}[1]{\left\langle#1\right\rangle}
\newcommand{\abs}[1]{\left|#1\right|}
\newcommand{\norm}[1]{\left\lVert#1\right\rVert}
\newcommand{\half}{\tfrac12}
\newcommand{\R}{\mathbb{R}}\newcommand{\C}{\mathbb{C}}\newcommand{\Z}{\mathbb{Z}}
% identity operator. NB: \mathbb{1} silently renders a WRONG glyph (msbm has no
% blackboard digits). Use \mathbb{I} = 𝕀 (works with amssymb, no extra package).
% Do NOT "fix" this back to \mathbb{1}.
\newcommand{\id}{\mathbb{I}}

% ---- physical units (v3.35): \punit — upright unit with thin space ---------
% Usage: $300\punit{K}$, $80\punit{mV/decade}$, $10^{7}\punit{cm/s}$ — the
% NUMBER AND UNIT LIVE IN ONE MATH ISLAND. Promoted verbatim from a 199-page
% production lakit (advanced-electronic-devices, dozens of uses).
% Two measured traps this design avoids (both render-proven):
%   1. siunitx SCRAMBLES under RTL — \SI{3.5}{eV} renders "eV 3.5", exponents
%      detach ("Hz 10 x 10^9"): its quantity assembly is not bidi-isolated.
%      Do NOT load siunitx in this pipeline.
%   2. A text-mode unit next to a bare-prose number is TWO adjacent LTR runs
%      and reorders ("5 mm" -> "mm5" — the BIDI-2/3 class). \punit is therefore
%      math-only ON PURPOSE: prose misuse fails LOUDLY at compile (\mathrm
%      outside math) instead of scrambling silently.
% Domain shorthands (\MeV, \fm...) stay per-document: \newcommand{\MeV}{\,\text{\en{MeV}}}
\providecommand{\punit}[1]{\,\mathrm{#1}}

% ============================================================================
% Colored boxes — mdframed (NOT tcolorbox, which breaks under bidi).
% ★★ v3.4: innerleftmargin/innerrightmargin give code listings room so their
% frame/line-numbers do not spill out the box border (see listing styles below).
% ============================================================================
% v3.61: everyline/split skips were applied to `tablebox` only (BLD-14, v3.37) —
% the six colored boxes below never got them, so they still split OPEN across a page
% break: the page-N segment gets top+sides and NO bottom rule, the page-N+1 segment
% sides+bottom and NO top rule. Measured on the shipped fleet this is a COVERAGE GAP,
% not a re-report of BLD-14: 0 split-open `tablebox` segments (that fix holds) vs 62
% split-open COLORED boxes across 7 documents. Verified after the fix: both segments
% close, and the frametitle is NOT duplicated on the continuation. BLD-14's caveat
% still applies — a box taller than a full page grinds mdframed's split-trial loop, so
% split such content manually.
% ============================================================================
% v3.72: frametitle={} leads EVERY \newmdenv option list below. mdframed options
% are GROUP-scoped: a box nested inside \begin{exbox}[frametitle={...}] used to
% INHERIT the outer title — measured, the identical title string rendered twice,
% the second riding a title-bar fill inside the INNER box's frame, for all seven
% environments in both directions. Both natural forms were broken (= preamble
% bug, not an authoring rule). The explicit reset makes every box titleless
% unless ITS OWN [frametitle={...}] says otherwise; verified both directions:
% no duplicate on the nested box, an explicit title still renders.
% Geometry is shared with the optional semantic-width table container. Themes
% select their own existing semantic style below; neither path fixes columns.
\mdfdefinestyle{hllpdfsemanticgeometry}{linewidth=1.1pt,
  innertopmargin=6pt,innerbottommargin=7pt,skipabove=9pt,skipbelow=9pt,
  innerleftmargin=9pt,innerrightmargin=9pt,everyline=true,
  splitbottomskip=6pt,splittopskip=12pt,roundcorner=3pt}
\newcommand{\hllpdfsemanticframestyle}{hllpdfsemanticgeometry}
\newmdenv[frametitle={},style=hllpdfsemanticgeometry,linecolor=cBlue,backgroundcolor=cBlueBg,
  frametitlebackgroundcolor=cBlue,frametitlefontcolor=white,
  ]{defbox}   % הגדרה
\newmdenv[frametitle={},style=hllpdfsemanticgeometry,linecolor=cGreen,backgroundcolor=cGreenBg,
  frametitlebackgroundcolor=cGreen,frametitlefontcolor=white,
  ]{thmbox}   % משפט/טענה
\newmdenv[frametitle={},style=hllpdfsemanticgeometry,linecolor=cAmber,backgroundcolor=cAmberBg,
  frametitlebackgroundcolor=cAmber,frametitlefontcolor=white,
  ]{notebox}  % הערה
\newmdenv[frametitle={},style=hllpdfsemanticgeometry,linecolor=cPurple,backgroundcolor=cPurpleBg,
  frametitlebackgroundcolor=cPurple,frametitlefontcolor=white,
  ]{exbox}    % דוגמה
\newmdenv[frametitle={},style=hllpdfsemanticgeometry,linecolor=cRed,backgroundcolor=cRedBg,
  frametitlebackgroundcolor=cRed,frametitlefontcolor=white,
  ]{warnbox}  % טעות נפוצה
\newmdenv[frametitle={},style=hllpdfsemanticgeometry,linecolor=cTeal,backgroundcolor=cTealBg,
  frametitlebackgroundcolor=cTeal,frametitlefontcolor=white,
  ]{keybox}   % רעיון מרכזי

% ---- codebox (v4.5.0): the dedicated container for code -------------------
% An OPTION, not a policy: asmblock unboxed stays canonical for a full listing,
% and a snippet inside any of the six colored boxes still works as before.
% What makes codebox worth its own environment is that it is the ONE framed box
% that keeps its margin line numbers.
%
% v4.4.0 shipped it WITHOUT them, on the strength of v3.37's "no margin knob can
% save them" plus four fresh probes that agreed. That was wrong: every probe had
% moved the code, the box, or both — never the key that carries the NUMBER
% COLUMN. Measured on the page:
%   resetmargins=false   the listing follows THIS box instead of the page,
%                        moving the code from x=75.26 to x=97.54
%   xrightmargin=26pt    under RTL this is the PHYSICAL LEFT indent, and the
%                        only key that carries the \llap'ed digits with it:
%                        x=64.28 -> 90.18, clearing the rule at 65.64 by 23.6pt
%   xleftmargin=26pt     moves the code and LEAVES THE DIGITS on the rule —
%                        the asymmetry that hid the fix for eleven versions
% resetmargins is hoisted out of the four styles (see the \lstset below them) so
% this override is possible at all; every other listing still resolves to true.
% Pinned by tests/codebox_numbers.tex + probe_lstnum.py --inside-rule, whose
% negative control strips xrightmargin and watches the digits fall back.
\definecolor{cSlate}{HTML}{3F4A56}  \definecolor{cSlateBg}{HTML}{F2F4F7}
\newmdenv[frametitle={},leftline=true,rightline=true,topline=true,bottomline=true,
  linecolor=cSlate,backgroundcolor=cSlateBg,linewidth=0.9pt,
  frametitlebackgroundcolor=cSlate,frametitlefontcolor=white,
  innertopmargin=6pt,innerbottommargin=7pt,skipabove=9pt,skipbelow=9pt,
  innerleftmargin=9pt,innerrightmargin=9pt,everyline=true,splitbottomskip=6pt,splittopskip=12pt,
  roundcorner=3pt]{codebox}   % קוד
% The two keys that make in-box numbering work, both measured on the page:
%   resetmargins=false  the listing follows THIS box instead of the page
%   xrightmargin        under RTL this is the PHYSICAL LEFT indent, and it is
%                       the only one that carries the \llap'ed number column
%                       with it. xleftmargin moves the code and leaves the
%                       digits behind on the rule (measured: code 97.54,
%                       digits 64.28 — the v3.37 defect, reproduced); with
%                       xrightmargin=26pt the digits land at x=90.18, a
%                       clearance of 23.6pt from the rule at 65.64.
% This asymmetry is why v3.4's note had the semantics backwards and v3.37
% concluded the digits could not be saved.
% <hllpdf:listing-container-state>
% v4.29: all state is local to the current environment group. A semantic
% callout masks number text, while a nested codebox restores the formatter
% saved before that mask. Save it once across successive semantic wrappers;
% an explicit numbers=none (or an explicitly empty formatter) remains valid.
\makeatletter
\newif\ifhllpdf@incodebox
\newif\ifhllpdf@englishcodebox
\newif\ifhllpdf@lstnumbermasked
\let\hllpdf@savedlstnumber\thelstnumber
\makeatother
% </hllpdf:listing-container-state>
% <hllpdf:codebox-gutter>
\makeatletter
\AtBeginEnvironment{codebox}{%
  \hllpdf@incodeboxtrue
  % Snapshot the OUTER language before asmblock selects English.
  \iflanguage{english}{\hllpdf@englishcodeboxtrue}{\hllpdf@englishcodeboxfalse}%
  \ifhllpdf@lstnumbermasked
    \let\thelstnumber\hllpdf@savedlstnumber
    \hllpdf@lstnumbermaskedfalse
  \fi
}
% listings' AddToHook is GLOBAL. Register once, and gate on local box state;
% adding an unconditional hook at every codebox entry leaks into later listings.
\lst@AddToHook{PreInit}{%
  \ifhllpdf@incodebox
    % Apply every listing override here, inside the LISTING's own group.
    % Setting margins at box entry leaks them into nested semantic callouts.
    \lstset{resetmargins=false,xrightmargin=26pt}%
    % Styles resolve before PreInit. In LTR the physical number gutter follows
    % xleftmargin; in Hebrew retain the measured inherited xrightmargin path.
    \ifhllpdf@englishcodebox
      \lstset{xleftmargin=26pt,xrightmargin=0pt}%
    \fi
    % ONE box, not two: suppress only the frame/background/skips of a listing
    % inside this codebox. The next unboxed listing keeps its own style.
    \lstset{frame=none,backgroundcolor=\color{cSlateBg},aboveskip=0pt,belowskip=0pt}%
  \fi
}
\makeatother
% </hllpdf:codebox-gutter>

% v3.37: line numbers inside a COLORED BOX are sliced by the box border —
% measured: digit span x=63.5-67.3 vs border line at x=65.2 (the border runs
% THROUGH the digits), and the inter-border gap (3.5pt) is narrower than a
% digit (3.8pt), so no margin knob can save them (under resetmargins,
% xleftmargin is inert in-box — measured: code at x=75.26 with 2.2em AND
% 3.0em). A 3-6 line boxed snippet does not need line numbers; blank the
% number TEXT inside every colored box (\thelstnumber is what gets printed;
% the style's numbers=left survives, so nothing else shifts — code position
% byte-frozen at x=75.26, verified). Unboxed listings keep their numbers.
% v3.78: the blanking is coupled to the FRAME, not to the box name —
% clean_theme.tex registers LATER hooks (registration order wins) that
% restore the number text inside the boxes it renders frameless.
% v4.5.0 RETRACTS the "no margin knob can save them" clause above. It holds only
% under resetmargins=true (the shipped default): there xleftmargin is inert
% in-box, as measured. With resetmargins=false, xrightmargin DOES carry the
% \llap'ed digit column clear of the rule — measured; the codebox \lstset that
% does it is the one above, in the codebox block. The loop below therefore
% covers SIX boxes, not "every colored box": codebox opts out and keeps its
% numbers.
% <hllpdf:semantic-listing-mask>
\makeatletter
\@for\HEB@lstbox:=defbox,thmbox,notebox,exbox,warnbox,keybox\do{%
  \expandafter\AtBeginEnvironment\expandafter{\HEB@lstbox}{%
    \ifhllpdf@lstnumbermasked\else
      \let\hllpdf@savedlstnumber\thelstnumber
    \fi
    \hllpdf@lstnumbermaskedtrue
    \def\thelstnumber{}%
    % The nearest semantic container owns its listing presentation. On exit,
    % TeX's group restores the enclosing codebox state and formatter.
    \hllpdf@incodeboxfalse}}
\makeatother
% </hllpdf:semantic-listing-mask>

% ============================================================================
% Unified TABLE presentation: subtle uniform slate frame + auto-numbered
% caption rendered INSIDE the frame. Use:
%   \begin{tablebox}\centering  <tabular>  \tabcap{<description>}\end{tablebox}
% Replaces the buggy \caption* (which emitted a stray "טבלה N : *" line) and
% provides an explicit framed option + numbered caption "טבלה N.M: תיאור".
% Optional intro text may be placed at the TOP, before the tabular, inside the box.
% ============================================================================
\definecolor{cTblLine}{HTML}{B9C4D2}  \definecolor{cTblBg}{HTML}{FAFBFC}
% v3.23: class-aware — under `article` there is no chapter counter (the promised
% report->article switch used to die with "No counter 'chapter' defined").
\makeatletter
\@ifundefined{c@chapter}{%
  % v3.31: in article the counter resets per SECTION, so an unprefixed number
  % repeated ("טבלה 1." in every section — measured on a two-section render).
  % Same conditional idea as the report branch: prefix the section number once
  % a real \section exists; a pre-section table stays "טבלה 1".
  \newcounter{tbl}[section]
  \renewcommand{\thetbl}{\ifnum\value{section}>0 \thesection.\fi\arabic{tbl}}%
}{%
  \newcounter{tbl}[chapter]
  % v3.30: a report-class doc that uses SECTIONS ONLY (chapter stays 0 —
  % the common short-doc pattern) used to print "טבלה 0.1"; the chapter
  % prefix now appears only once a real \chapter exists (render-verified
  % both ways: "טבלה 1." chapterless, "טבלה 1.1" with chapters).
  \renewcommand{\thetbl}{\ifnum\value{chapter}>0 \thechapter.\fi\arabic{tbl}}%
  % v3.64: the SAME leak in \thesection itself, left behind by v3.30. That fix
  % guarded the CAPTION counter and stopped there, so a chapterless report still
  % printed the heading "0.1 סעיף ראשון" and \ref to it printed "0.2" — measured:
  %   p1 y=65.2  '…0.1'   p1 y=144.9 '…0.2'   p1 y=183.4 '…0.2.1'
  % report.cls hardcodes \thesection as \thechapter.\@arabic\c@section, and the
  % leading 0 is the chapter counter that no \chapter ever incremented. Fixing
  % \thesection fixes the whole chain — \thesubsection and \thesubsubsection are
  % defined in terms of it. \theequation and \thetbl/\thefig were already guarded,
  % which is exactly why this survived: every NUMBER around the heading looked
  % right, so the heading's own 0 read as a style choice rather than a bug.
  % \appendix is safe: report.cls resets chapter to 0 and redefines \thechapter to
  % \@Alph\c@chapter, so an appendix \chapter makes the counter 1 and the guard
  % emits "A.1" unchanged. Render-verified in all four modes (chapterless,
  % chaptered, appendix+chapter, appendix+section-only).
  \renewcommand{\thesection}{\ifnum\value{chapter}>0 \thechapter.\fi\arabic{section}}%
}
\makeatother
% v3.37: everyline=true — if the box DOES split across pages, every segment
% gets a CLOSED frame (measured before: the split ran open-ended into the page
% edge with no bottom border, and the caption could strand alone at the next
% page top — "טבלה A.2" shipped that way). The split skips give the segments
% breathing room. In practice, with the caption GLUE below, every previously-
% split table in the corpus now moves WHOLE instead (the white-gap tradeoff,
% documented benign); everyline stays as the safety net for a genuinely
% page-tall table. NB: a table taller than ONE FULL PAGE can grind mdframed's
% split-trial loop (reproduced: minimal doc, minutes of 1pt-step Overfull
% probes) — break such a monster into two tables.
\mdfdefinestyle{hllpdftablegeometry}{leftline=true,rightline=true,topline=true,bottomline=true,
  linewidth=0.9pt,
  innertopmargin=9pt,innerbottommargin=8pt,skipabove=11pt,skipbelow=11pt,
  innerleftmargin=11pt,innerrightmargin=11pt,roundcorner=4pt}
\newcommand{\hllpdftableframestyle}{hllpdftablegeometry}
\newmdenv[frametitle={},style=hllpdftablegeometry,
  linecolor=cTblLine,backgroundcolor=cTblBg,
  % v4.2.0 (adopted): tablebox is ATOMIC. Short tableboxes near page bottoms
  % repeatedly entered mdframed's split trial, emitting benign Overfull-vbox
  % bursts and occasionally separating the intro prose from the tabular.
  % nobreak=true moves the complete short table to the next page instead. The
  % output-active Overfull-vbox gate still catches a box taller than a page;
  % everyline/split skips remain as defensive styling should that path ever
  % re-arm, but the canonical path never publishes a split tablebox.
  nobreak=true,
  everyline=true,splitbottomskip=6pt,splittopskip=12pt,
  % v3.79 (BLD-37 table side, correction): startinnercode/endinnercode run
  % INSIDE mdframed's body vbox (mdframed.sty:848 / :856+863) -- the second one
  % after \par\unskip, i.e. after line breaking, which is the only moment the
  % \tabcap anchor's true offset from the top of the box EXISTS. See
  % \lakintabboxclose below. These two keys are the ONLY tablebox change.
  startinnercode={\lakintabboxopen},endinnercode={\lakintabboxclose}]{tablebox}
% v3.37: \nopagebreak GLUE — \vsplit honors penalties, so the caption can no
% longer strand alone at the top of the next page (it takes its rows along, or
% the whole box moves). Render-verified on the real shipped orphan.
% v3.59: \NR@gettitle is what nameref injects into the stdlib title-bearing
% commands (\@sect, \@chapter, \caption) to fill \@currentlabelname. \figcap and
% \tabcap are INDEPENDENT macros — they never called it, so \@currentlabelname kept
% whatever the last heading left and \nameref{fig:x} silently printed the enclosing
% SECTION's title instead of the caption (proven in the .aux: field 3 of \newlabel
% was the heading text while field 1 and the anchor `fig.1.1` were correct — so this
% is a THIRD, distinct head next to BLD-13's counter fix and BIDI-38's autoref-name
% fix, neither of which writes that field). \@ifundefined keeps it inert if
% hyperref/nameref is ever dropped. Guarded helper because \figcap is defined
% after \makeatother.
% NB both helpers below need \makeatletter — this part of the preamble is OUTSIDE
% any letter-catcode region, so an unwrapped \@currentlabel parses as LaTeX's \@
% and dies with "You can't use \spacefactor in vertical mode" (measured).
\makeatletter
\newcommand{\lakinrtitle}[1]{\@ifundefined{NR@gettitle}{}{\NR@gettitle{#1}}}
% v3.59: \hitem{<letter>}{<key>} — a REFERENCEABLE Hebrew-lettered list item.
%   \begin{enumerate} \hitem{א}{it:a} הסעיף הראשון. \end{enumerate}  ->  \ref{it:a} = (א)
% Why this exists: automatic Hebrew list labels fail silently, so SKILL.md mandates
% an EXPLICIT \item[(\textbf{א})] — but LaTeX puts \refstepcounter INSIDE
% \if@noitemarg (latex.ltx:12964), and \item[...] never sets that flag. So an
% explicit item steps no counter and never touches \@currentlabel: a \label there
% binds to the AMBIENT one and \ref prints the CHAPTER/SECTION number where the
% item letter belongs (proven in the .aux: {chapter.1}). Same doctrine as \figcap —
% own the label state before \label. \phantomsection matters: without it the
% hyperlink still lands on the chapter head (measured); with it the anchor is
% section*.N and the bookmark outline stays clean.
\newcommand{\hitem}[2]{\item[(\textbf{#1})]%
  \def\@currentlabel{(#1)}\def\@currentlabelname{#1}\phantomsection\label{#2}}
% NB (v3.66): \hitem is \ref-only. \autoref{key} calls it a SECTION ("סעיף (א)")
% because \phantomsection brands the anchor section*.N, and \nameref{key} prints
% the bare letter. For \autoref-able items prefer \hebenumlabels below.
%
% v3.66: \hebenumlabels — OPT-IN automatic Hebrew enumerate letters, done right.
%   {\hebenumlabels\begin{enumerate}\item\label{q:a} ...\end{enumerate}}
%   -> labels (א) (ב) (ג); \ref{q:a} = א; anchor Item.N (real, per-item).
% This retires the belief that made \item[(\textbf{א})] + \hitem necessary.
% "Automatic Hebrew labels fail silently" was TRUE but unexplained, and the
% mechanism is now bisected (plain article vs +enumitem, measured): ENUMITEM
% PRE-EXPANDS the label expression in an \edef BEFORE \refstepcounter, with the
% counter still at 0. Tokens it shields by name (\arabic,\alph,\roman) survive;
% a plain \newcommand \ifcase wrapper evaluates at 0 -> the EMPTY zero-branch ->
% the label prints as "()" with the letter GONE (and \the\c@enumi prints a
% literal 0, \number\value dies fatally pointing at \item). The cure is one
% property, not a doctrine: the wrapper must be ROBUST (survives the \edef
% unexpanded), while \theenumi stays EXPANDABLE (so \refstepcounter writes the
% real letter to the .aux). Verified: labels (א)(ב)(ג), \ref -> א/ב/ג each item
% its own value, .aux = {א}{...}{Item.1} — versus the explicit-\item doctrine
% where BIDI-50/\hitem exist purely to patch the counter that never stepped.
\def\@hebalph#1{\ifcase#1\or א\or ב\or ג\or ד\or ה\or ו\or ז\or ח\or ט\or י\or
  יא\or יב\or יג\or יד\or טו\or טז\or יז\or יח\or יט\or כ\else\@ctrerr\fi}
\DeclareRobustCommand*\hebalph[1]{\expandafter\@hebalph\csname c@#1\endcsname}
\newcommand{\hebenumlabels}{%
  \renewcommand{\theenumi}{\@hebalph\c@enumi}%
  \renewcommand{\labelenumi}{(\textbf{\hebalph{enumi}})}}
\makeatother
% v3.63: \mbox the LABEL. \thetbl ("7.1") and a caption description that STARTS
% WITH A DIGIT ("32 הרגיסטרים…") are two European-Number runs separated only by
% neutrals (the label's own "." and \enspace). Under the RTL paragraph base the UBA
% resolves those neutrals to the RTL level, so both numbers sit in ONE reversed span
% and EXCHANGE VISUAL POSITIONS. Measured on the shipped PDF: a table whose real
% number is 7.1 printed "טבלה 32" — a different, EXISTING table number — while \ref
% correctly printed 7.1, so the caption silently contradicted its own cross-reference,
% and the description lost its leading count. The source contains no adjacency at all;
% it is manufactured by this macro, which is why every source linter was blind.
% \mbox makes the label one atomic box so nothing can reorder across it. Verified
% zero-regression: an ordinary Hebrew caption keeps identical x-geometry.
% One paragraph style covers label, number, separator and description.
% Explicit project/template overrides remain possible; the label hook is kept
% for compatibility, with an identity default. Counters/anchors stay unchanged.
\newcommand{\lakitcaptionstyle}{\normalfont\small\color{black}\hypersetup{allcolors=black}}
\newcommand{\lakitcaptionlabelstyle}[1]{#1}
\newcommand{\lakitcaptionpunctuation}{:}
% Native floats, captionof and longtable use the same presentation policy.
% subcaption is optional and inherits the same font rather than a second size.
\usepackage{caption}
\DeclareCaptionFont{lakit}{\lakitcaptionstyle}
\DeclareCaptionLabelFormat{lakit}{\lakitcaptionlabelstyle{#1~\en{#2}}}
\DeclareCaptionLabelSeparator{lakit}{\lakitcaptionpunctuation\enspace}
\captionsetup{font=lakit,labelfont={},textfont={},labelformat=lakit,
  labelsep=lakit,justification=centering}
\makeatletter
\AtBeginDocument{\@ifpackageloaded{subcaption}{%
  \captionsetup[sub]{font=lakit,labelfont={},textfont={},labelsep=lakit}}{}}
\makeatother
% v3.64: OPT-IN list-of-figures/tables entry via an optional short caption.
% \listoffigures and \listoftables are inherited from report.cls and LOOK available,
% but \figcap/\tabcap step the custom \c@fig/\c@tbl counters and never call
% \addcontentsline — so both printed their heading and NOTHING ELSE. Measured: a doc
% with one \figcap and one \tabcap rendered page 1 as 'רשימת האיורים\n1\n', and the
% .lof held only \@chapter's \addvspace, zero \contentsline entries. An empty list
% under a correct Hebrew heading reads as "this document has no figures".
%
% WHY OPTIONAL, AND NOT AUTOMATIC. Writing #1 unconditionally was implemented first
% and REJECTED: a caption containing \codeline compiles CLEANLY today (verified on
% the pre-v3.64 preamble: 0 errors, PDF produced) but becomes FATAL once it is also
% written to the .lof — \codeline is \hbox to\linewidth{...}, and a full-width box in
% a list entry dies with "\pdfextension endlink cannot be used in vertical mode", no
% output PDF. Gaining a feature nobody uses (0 of 13 collections call \listoffigures)
% at the price of breaking a construct that works is a bad trade. Stock \caption has
% exactly this limitation and exactly this escape hatch, so the convention is familiar.
% \en{}, \code{}, \h{} and inline math in the short text are all safe — verified, they
% survive the write and typeset correctly in the list (RTL geometry confirmed: number
% rightmost at x=509.9, title left of it, dots and folio at the far left).
% With no optional argument the macro is byte-identical to v3.63 — zero risk to every
% existing document, and \listoffigures lists exactly the captions that opted in.
% BLD-37, TABLE side (the \figcap v3.74 idea, adapted): clicking a \ref to a
% table used to land MID-table — hyperref's automatic anchor inside
% \refstepcounter was emitted on the still-open horizontal line whose TAIL is
% the tabular (an hbox whose baseline is the table's vertical CENTER), so the
% jump scrolled the top half of the table above the viewport. The v3.74 entry
% deferred this side claiming "tablebox splits — no single box to measure";
% that deferral dissolves: the measurement target is the TABULAR itself at
% \tabcap ENTRY, while it is still the tail of the open list — the same
% situation \lakin@grabfig handles at an island's end (and the fleet has 0
% split tableboxes). So \lakin@grabtab runs as the FIRST token, before any
% typesetting can push the box away, and the anchor is re-emitted raised above
% the measured tabular. \lakintabraise is a SEPARATE register from
% \lakinfigraise (defined with it, below the LTR-island block): a captionless
% island intentionally leaves the fig register set, and a table must never
% inherit that stale figure height.
%
% WHY the anchor stays AT THE ENTRY SPOT and is NOT deferred into the caption
% paragraph the way \figcap defers it: at \par TeX DROPS a final glue item,
% and the automatic anchor's smash box sitting BEHIND the interword glue kept
% that glue alive — so every centered tabular has always sat half a space off
% true center. Deferring the anchor removed that box and was measured as 982
% differing glyph records: EVERY tabular line shifted 1.30pt (2.61pt for
% pattern 7's two minipages) while captions stayed put. Emitting the raised
% anchor at the old spot — through hyperref's own \Hy@raisedlink, exactly the
% macro the automatic path uses — keeps the line's node structure, and so the
% A/B glyph diff at 0, BY CONSTRUCTION. \Hy@raisedlink reads the raise from
% \HyperRaiseLinkDefault (default \baselineskip = the old mid-table dest); the
% group-local \edef freezes measured-height+42pt into it, and \smash inside
% \Hy@raisedlink hides any raise from the line's extent. \HyCnt@ProvideTheH-
% Counter keeps the dest NAME byte-identical to the automatic one (it defines
% \theHtbl, which \hyper@makecurrent prefers — without it the first measured
% caption was named tbl.1 where the automatic path said tbl.0.1, measured).
% vmode entry, or anything but a box at the tail (no tabular to measure),
% falls back to the plain \refstepcounter — byte-identical to the old path.
\makeatletter
\newcommand{\tabcap}[2][]{\lakin@grabtab
  % v3.79: arm the box-close correction. Set unconditionally and in BOTH
  % branches below's scope, so the plain-\refstepcounter fallback is corrected
  % too; \lakintabboxopen clears it at every \begin{tablebox}, so a \tabcap
  % outside a tablebox can never leave it armed for a later box.
  \global\lakintabcapseentrue
  \ifdim\lakintabraise>0pt
    \H@refstepcounter{tbl}\HyCnt@ProvideTheHCounter{tbl}\hyper@makecurrent{tbl}%
    \begingroup
      % A frameless table targets its measured top. The frame's historical
      % cushion belongs only to tablebox and is corrected at box close.
      \edef\HyperRaiseLinkDefault{\the\lakintabraise}%
      \iflakintabframed
        \edef\HyperRaiseLinkDefault{\the\dimexpr\lakintabraise+42pt\relax}%
      \fi
      \Hy@raisedlink{\hyper@anchorstart{\@currentHref}\hyper@anchorend}%
    \endgroup
    \global\lakintabraise 0pt\relax
  \else
    \refstepcounter{tbl}%
  \fi
  \lakinrtitle{#2}\par\nopagebreak\addvspace{6pt}%
  \def\lakin@short{#1}\ifx\lakin@short\@empty\else
    \addcontentsline{lot}{table}{\protect\numberline{\thetbl}{\protect\ignorespaces #1}}\fi
  {\centering\lakitcaptionstyle\mbox{\lakitcaptionlabelstyle{\tablename~\en{\thetbl}\lakitcaptionpunctuation}}\enspace #2\par}}
\makeatother

% ============================================================================
% Unified FIGURE caption — the \tabcap twin for TikZ/plot figures (v3.34.0).
% A figure (tikzpic) sits inline in the prose; put \figcap BELOW it for a
% numbered, \ref-able caption "איור N.M: תיאור" in the same style as \tabcap:
%   \begin{tikzpic} ... \end{tikzpic}
%   \figcap{\label{fig:foo}description}     % \label FIRST, inside the argument
% WHY a macro, not a hand-written gray line: a manual
%   {\centering\small\textcolor{cGray}{איור: ...}} with a \label inside has NO
% \refstepcounter, so the \label binds to the AMBIENT \@currentlabel (the last
% equation) — \ref then prints an equation number AS the figure number
% (measured on a real page: "איור 48" where 48 is equation 48). \figcap steps a
% real \c@fig counter first, so \label is correct, numbering is consistent, and
% refs are stable. Class-aware like \thetbl (chapter- or section-prefixed).
% ============================================================================
\makeatletter
\@ifundefined{c@chapter}{%
  \newcounter{fig}[section]
  \renewcommand{\thefig}{\ifnum\value{section}>0 \thesection.\fi\arabic{fig}}%
}{%
  \newcounter{fig}[chapter]
  \renewcommand{\thefig}{\ifnum\value{chapter}>0 \thechapter.\fi\arabic{fig}}%
}
\makeatother
\makeatletter
\newcommand{\figcap}[2][]{%
  \ifdim\lakinfigraise>0pt
    % v3.74: step WITHOUT hyperref's automatic anchor, then emit the anchor
    % raised above the measured figure. \H@refstepcounter is hyperref's saved
    % pre-hook \refstepcounter (sets the counter AND \@currentlabel — \label
    % in the caption arg binds exactly as before); \hyper@makecurrent supplies
    % \@currentHref, so the dest NAME is identical to the old automatic one —
    % \ref/\autoref/bookmarks are unchanged, only the dest'S POSITION moves.
    \H@refstepcounter{fig}\hyper@makecurrent{fig}%
    % The anchor is emitted INSIDE the caption paragraph below, as a zero-size
    % raisebox — NOT here in vertical mode, where a \raisebox starts an empty
    % paragraph and adds a blank line (measured: every glyph after the caption
    % shifted +18.8pt). Inside an existing line a [0pt][0pt] box is inert.
    \edef\lakin@anch{\@currentHref}%
    % \edef FREEZES the measured height into the pending macro: the deferred
    % \raisebox is typeset inside the caption paragraph AFTER the register is
    % reset to the 0pt sentinel below — a plain \def read 0 at use time
    % (measured: LAKINRAISE=0.0, dest moved 14pt instead of above the figure).
    \edef\lakin@pend{\noexpand\raisebox{\the\dimexpr\lakinfigraise+42pt\relax}[0pt][0pt]{%
      \noexpand\hyper@anchorstart{\lakin@anch}\noexpand\hyper@anchorend}}%
    \global\lakinfigraise 0pt\relax
  \else
    \refstepcounter{fig}\let\lakin@pend\relax
  \fi
  \lakinrtitle{#2}\par\nopagebreak\addvspace{4pt}%
  \def\lakin@short{#1}\ifx\lakin@short\@empty\else
    \addcontentsline{lof}{figure}{\protect\numberline{\thefig}{\protect\ignorespaces #1}}\fi
  {\centering\lakin@pend\lakitcaptionstyle\mbox{\lakitcaptionlabelstyle{\figurename~\en{\thefig}\lakitcaptionpunctuation}}\enspace #2\par}}
\makeatother

% ============================================================================
% LTR islands.  EVERY tikzpicture / quantikz MUST live inside one, or the whole
% figure's coordinates and Latin labels flip. (Hebrew node text still uses \h{}.)
% ============================================================================
% v3.74 (BLD-37): the \figcap hyperlink DESTINATION used to land at the caption
% line — BELOW the figure — so clicking a figure reference scrolled the figure
% off-screen (measured fleet-wide: 97/97 fig dests in one collection sat 19-54pt
% under their figure's ink; a 347pt figure was fully hidden after the jump).
% hypcap-style fix in two parts: (1) at the island's END, \lakin@grabfig grabs
% the just-typeset centered line with \lastbox, measures it, and puts it back
% with \nointerlineskip — the interline glue above the line was already
% contributed and is NOT re-added, so the page geometry is unchanged to the
% glyph (verified); (2) \figcap then steps the counter WITHOUT hyperref's
% automatic anchor (\H@refstepcounter = the pre-hyperref \refstepcounter) and
% emits the anchor itself, raised by the measured height — the dest lands ABOVE
% the figure top. A captionless island leaves numbering untouched (nothing is
% stepped at the island; the grab flag is cleared at every island begin, so a
% stale height from a captionless figure cannot leak into a later caption).
\newsavebox{\lakinfigbox}
\newdimen\lakinfigraise  % real \dimen: tex.setdimen cannot target a \newlength SKIP register (measured lua error)
% BLD-37 table side: a SEPARATE register for \tabcap. The two must not share —
% a captionless island INTENTIONALLY leaves \lakinfigraise set (the grab runs
% at every island END; only the NEXT island's begin clears it), so a shared
% register would hand a \tabcap after a captionless figure that stale FIGURE
% height. \lakintabraise is written only by \lakin@grabtab at \tabcap entry
% and consumed/reset inside \tabcap itself, so it can never go stale.
\newdimen\lakintabraise
\makeatletter
% v3.74b: measure WITHOUT touching the list. The first draft grabbed the box
% with \lastbox and put it back — and \lastbox DISCARDS the node's baseline
% shift, so a vertically-anchored quantikz circuit rose 2.57pt and everything
% after it sank 5.32pt (measured on the A/B glyph diff). LuaTeX lets us READ
% the tail node's dimensions without removing anything: geometry is identical
% by construction. \lakinfigraise=0pt is the "no figure measured" sentinel.
\def\lakin@grabfig{%
  \ifhmode
    \directlua{
      local n = tex.nest[tex.nest.ptr].tail
      local GLUE, PEN = node.id('glue'), node.id('penalty')
      while n and (n.id == GLUE or n.id == PEN) do n = n.prev end
      local HL, VL = node.id('hlist'), node.id('vlist')
      if n and (n.id == HL or n.id == VL) then
        tex.setdimen('global', 'lakinfigraise',
                     (n.height or 0) + (n.depth or 0))
      else
        tex.setdimen('global', 'lakinfigraise', 0)
      end}%
  \else
    \global\lakinfigraise 0pt\relax
  \fi}
% BLD-37 table side: \lakin@grabfig's twin, NOT a shared helper — the Lua body
% differs in exactly one place, the skip set. Between \end{tabular} and \tabcap
% an otherlanguage wrapper (patterns 3/5 of example_tables) leaves a \write
% whatsit and a dir node ON TOP of the interword glue (measured node dump:
% "whatsit/write <- dir <- glue(2.62) <- hlist"), so the table reader must
% also step over whatsit/dir markers — all zero-size — to reach the tabular.
% \lakin@grabfig deliberately does NOT get the wider skip set: its glue/penalty
% body is the one the v3.74 A/B proof froze, and widening it would newly
% measure (and so newly raise) anchors in island shapes that today take the
% plain path — an unproven behavior change on the FIG side. Anything that is
% not a box behind those markers still hits the 0pt sentinel and \tabcap falls
% back to the plain \refstepcounter.
% SECOND deliberate difference: this stores the box's HEIGHT, not height+depth.
% \figcap's anchor is emitted in the caption paragraph BELOW the figure, so it
% must climb the full extent; \tabcap's anchor sits ON the tabular's own line,
% whose baseline is the vcentered tabular's MIDPOINT — height alone reaches
% the top edge, and adding the depth was measured to overshoot by half the
% table (dests 80-135pt above their frame, INSIDE the preceding table's box).
\def\lakin@grabtab{%
  \ifhmode
    \directlua{
      local n = tex.nest[tex.nest.ptr].tail
      local GLUE, PEN = node.id('glue'), node.id('penalty')
      local WHAT, DIR = node.id('whatsit'), node.id('dir')
      while n and (n.id == GLUE or n.id == PEN
                   or n.id == WHAT or n.id == DIR) do n = n.prev end
      local HL, VL = node.id('hlist'), node.id('vlist')
      if n and (n.id == HL or n.id == VL) then
        tex.setdimen('global', 'lakintabraise', n.height or 0)
      else
        tex.setdimen('global', 'lakintabraise', 0)
      end}%
  \else
    \global\lakintabraise 0pt\relax
  \fi}
\makeatother
% ============================================================================
% v3.79 (BLD-37 table side, CORRECTION): put the anchor at the TOP OF THE BOX,
% instead of guessing the distance to it.
% ============================================================================
% v3.77 raised the \tabcap anchor by (height of the box at the tail of the open
% horizontal list) + 42pt. That is a MEASUREMENT taken BEFORE line breaking, and
% it is short by every term it cannot see. Fleet-measured on the 12 built
% collections (299 tbl.* dests, ownership established through each table's own
% CAPTION, never through "which frame encloses the dest"): 3 dests landed INSIDE
% their own tablebox, below the frame top —
%   computer-architecture       tbl.14.1  -46.4pt   three stacked tabulars
%   computer-architecture       tbl.30.5  -13.5pt   tabular 2.85pt over \linewidth
%   advanced-electronic-devices tbl.B.1   -59.2pt   tabular 2.56pt over \linewidth
% Two missing terms, both reproduced in a minimal document (probe P2/P3):
%   (1) VERTICAL LIST ALREADY CONTRIBUTED inside the box — a second and third
%       tabular, or a heading above the table. The tail read sees only the LAST
%       box, so 77.7pt of stacked tabulars is invisible to the raise.
%   (2) WHICH LINE THE ANCHOR LANDS ON. When the tabular alone overfills
%       \linewidth (459.6121pt in a tablebox), the interword glue after
%       \end{tabular} becomes the chosen breakpoint and \Hy@raisedlink is
%       emitted on the NEXT line — one (depth + lineskip) lower. That term does
%       not EXIST at \tabcap time: line breaking has not happened yet. No
%       measurement taken there can ever recover it. (Proof: probe P3, the same
%       table 2.84912pt too wide, fails at -13.5pt; probe P4, the identical
%       table with \tabcolsep 2pt so it fits, is correct at +32.0pt.)
% So the raise is not computed at \tabcap at all. It is CORRECTED at
% \end{tablebox}, from inside mdframed's own body vbox (endinnercode runs after
% \par\unskip, i.e. after line breaking, while tex.nest[tex.nest.ptr] is still
% that vbox's vertical list). There the anchor's true offset from the top of the
% box content is not estimated but WALKED, and the \Hy@raisedlink shift is set
% so the dest lands 42pt above it — which is 42pt - innertopmargin(9pt) -
% linewidth(0.9pt) = 32.0pt above the FRAME top, the exact position the 295
% already-correct dests have today. That is why this is a correction and not a
% re-placement: the good sites do not move.
%
% WHY THE GLYPH DIFF IS 0 BY CONSTRUCTION — the only thing written is the
% `shift` field of the hlist that \Hy@raisedlink already built, and the guard
% below refuses to write it unless that box is 0x0x0. A zero-size box carries no
% dimension into its parent's packing, the packing already happened, and shift
% is not an input to it. Nothing else in the list is touched: no node is added,
% removed, resized or reordered. Measured, not assumed: assets/example_tables.tex
% and the whole computer-architecture collection A/B at 0 differing pymupdf
% rawdict glyph records.
%
% WHY NOT the alternatives (all prototyped and measured before this one):
%   * accumulate the vertical list at \tabcap time — fixes term (1), CANNOT fix
%     term (2); 2 of the 3 real failures survive. A fix that leaves the majority
%     of the defect standing is the false all-clear this repo exists to avoid.
%   * \zref-savepos two-pass — correct on all shapes, but it makes the dest
%     POSITION aux-derived (DEFECTS BLD-1/2), a single-pass compile silently
%     ships an unraised anchor, and the bare y-difference has no page guard.
%   * emit the anchor at \begin{tablebox} and let \tabcap adopt it — needs the
%     tbl counter stepped at box OPEN to keep the dest NAME byte-identical; a
%     captionless tablebox then either shifts every later table number or leaves
%     an orphan dest that the next table duplicates (a HARD build gate). This
%     correction touches neither the counter nor the name: \tabcap's naming path
%     is v3.77's, unchanged, so resolve_names() is byte-identical (verified).
%   * suppressing the line break (a penalty before the glue) — that MOVES the
%     tabular: the break currently discards the interword glue, so forbidding it
%     re-centers the two overfull tables. Reflow, rejected.
%
% SCOPE. The correction is armed only by \tabcap ITSELF (\lakintabcapseen),
% cleared at every box open, so a captionless tablebox is a no-op and a \tabcap
% hosted OUTSIDE a tablebox (the fleet has exactly one, computer-architecture
% tbl.14.2 in an exbox) keeps v3.77's measured raise. Both fallbacks are
% therefore never worse than today; and because hyperref's AUTOMATIC anchor uses
% the same \Hy@raisedlink shape, the plain-\refstepcounter path (vmode entry, or
% prose between \end{tabular} and \tabcap — probe P6, -51.5pt) is corrected too.
% KNOWN LIMIT — a tablebox that SPLITS across pages. The offset is measured from
% the top of the UNSPLIT body, so a dest in segment 2 gets raised past its own
% segment. MEASURED on a forced split (probe: 40 filler lines + a table): the
% dest went from y=585.0 on page 2 (inside segment 2, at the caption — the old
% BLD-37 symptom) to y=-113.4, i.e. ABOVE the page top. Viewers clamp that to the
% page top, which happens to show segment 2 whole, but an out-of-page coordinate
% is not something to ship on purpose. Hence the \textheight guard in
% lakintab.fix: a body TALLER THAN THE TEXT BLOCK cannot fit on any page, so it
% must split — the correction is skipped and v3.77's behaviour stands unchanged
% (verified: the probe returns to y=585.0). The guard is deliberately EXACT, not
% a page-fitting heuristic: it is sound (a taller-than-page body always splits)
% but PARTIAL (a shorter body that starts low on a page could still split). That
% residue is unreachable in practice — v3.37's everyline + the caption's
% \nopagebreak glue move such boxes WHOLE instead (probe N7: a 22-row table
% pushed down by 430pt moved whole, 0 segments), and the fleet has 0 split
% tableboxes, re-counted after this change. Reimplementing mdframed's own
% page-fitting arithmetic to close the residue would be a heuristic with its own
% failure mode, which is the trade this repo has repeatedly refused.
\newif\iflakintabcapseen
\newif\iflakintabframed
% The Lua lives in ONE \directlua at load time, and is written with NO
% backslash, %, # or ~ : \directlua expands its argument like \edef, so an
% active ~ would expand to \nobreakspace and a % would eat the rest of the line
% (the second is how the first draft of this block died). "not (a == b)" stands
% in for Lua's inequality operator for the same reason.
\directlua{
  lakintab = { }
  lakintab.HL = node.id("hlist")
  lakintab.VL = node.id("vlist")
  lakintab.RU = node.id("rule")
  lakintab.GL = node.id("glue")
  lakintab.KE = node.id("kern")
  lakintab.WH = node.id("whatsit")
  lakintab.TM = node.id("temp")
  lakintab.WD = node.subtype("pdf_dest")
  function lakintab.d(v)
    if v == nil then return 0 end
    if v < -1000000000 then return 0 end
    return v
  end
  function lakintab.scan(list, y, horiz, parent, target)
    local x = list
    while x do
      local i = x.id
      if i == lakintab.HL or i == lakintab.VL then
        if horiz then
          local b = y + x.shift
          if i == lakintab.HL then
            lakintab.scan(x.head, b, true, x, target)
          else
            lakintab.scan(x.head, b - lakintab.d(x.height), false, x, target)
          end
        else
          if i == lakintab.HL then
            lakintab.scan(x.head, y + lakintab.d(x.height), true, x, target)
          else
            lakintab.scan(x.head, y, false, x, target)
          end
          y = y + lakintab.d(x.height) + lakintab.d(x.depth)
        end
      elseif i == lakintab.WH then
        if x.subtype == lakintab.WD then
          local id = x.dest_id
          if type(id) == "string" and string.sub(id, 1, 4) == "tbl." then
            if horiz and parent and parent.width == 0
               and parent.height == 0 and parent.depth == 0 then
              parent.shift = parent.shift + (target - y)
            end
          end
        end
      elseif not horiz then
        if i == lakintab.RU then
          y = y + lakintab.d(x.height) + lakintab.d(x.depth)
        elseif i == lakintab.GL or i == lakintab.KE then
          y = y + lakintab.d(x.width)
        end
      end
      x = x.next
    end
  end
  function lakintab.fix(target, maxh)
    local nst = tex.nest[tex.nest.ptr]
    if nst == nil then return end
    local h = nst.head
    if h == nil then return end
    if h.id == lakintab.TM then h = h.next end
    local tot = 0
    local x = h
    while x do
      local i = x.id
      if i == lakintab.HL or i == lakintab.VL or i == lakintab.RU then
        tot = tot + lakintab.d(x.height) + lakintab.d(x.depth)
      elseif i == lakintab.GL or i == lakintab.KE then
        tot = tot + lakintab.d(x.width)
      end
      x = x.next
    end
    if tot > maxh then return end
    lakintab.scan(h, 0, false, nil, target)
  end
}
% @-FREE NAMES ON PURPOSE: these two are named in the `tablebox` key list at
% \newmdenv above, which is tokenised OUTSIDE any \makeatletter region — a
% \lakin@... there would tokenise as \lakin followed by the letters @tabopen.
\newcommand{\lakintabboxopen}{\lakintabframedtrue\global\lakintabcapseenfalse}
\newcommand{\lakintabboxclose}{%
  \iflakintabcapseen
    \directlua{lakintab.fix(-\number\dimexpr\lakintabcushion\relax,
                            \number\dimexpr\textheight\relax)}%
    \global\lakintabcapseenfalse
  \fi\lakintabframedfalse}
% The cushion above the box's CONTENT top. 42pt puts the dest 32.0pt above the
% frame top (42pt - innertopmargin 9pt - linewidth 0.9pt), i.e. exactly where
% the 295 correct dests already are — the correction moves only the broken ones.
\newcommand{\lakintabcushion}{42pt}
\newenvironment{qcirc}{\par\medskip\global\lakinfigraise 0pt\relax\begin{otherlanguage}{english}\begin{center}}
  {\csname lakin@grabfig\endcsname\end{center}\end{otherlanguage}\par\medskip}
\newenvironment{tikzpic}{\par\medskip\global\lakinfigraise 0pt\relax\begin{otherlanguage}{english}\begin{center}}
  {\csname lakin@grabfig\endcsname\end{center}\end{otherlanguage}\par\medskip}

% ============================================================================
% Code listings.  ★★ v3.4 THE BOX-OVERFLOW FIX:
%   resetmargins=true   → the listing IGNORES the enclosing environment's
%                         indentation and measures from the PAGE. v3.4 stated
%                         this backwards ("measured from the ENCLOSING box"),
%                         and v3.37 then blanked in-box line numbers because
%                         the page-anchored listing put its hanging digits on
%                         top of the box rule. Measured v4.5.0: with the key
%                         FALSE the listing follows the box and the code moves
%                         from x=75.26 to x=97.74, clearing the rule at 72.07.
%                         The key now lives in one \lstset below instead of in
%                         each style, so a box can override it — codebox does.
%   framexleftmargin=0  → the frame (left bar) does not extend left of the code.
%   xleftmargin keeps a small indent so the line numbers have room INSIDE.
% asmblock wraps the listing in an LTR island (code is LTR).
% ============================================================================
\lstdefinestyle{asm}{
  basicstyle=\ttfamily\small,
  backgroundcolor=\color{cBlueBg!55},
  frame=leftline,framerule=2.5pt,rulecolor=\color{cBlue},
  numbers=left,numberstyle=\tiny\color{cGray},numbersep=8pt,
  xleftmargin=2.2em,framexleftmargin=0pt,framexrightmargin=0pt,
  breaklines=true,columns=fullflexible,keepspaces=true,
  showstringspaces=false,aboveskip=8pt,belowskip=8pt,
  % v3.36: without a comment DELIMITER the commentstyle below never fires —
  % listings has no `language` here, so it does not know '#' opens a comment,
  % and every asm listing rendered flat black (measured: a whole-page pixel
  % sweep found ZERO cGreen glyphs across 135 shipped listings). morecomment
  % gives it the delimiter; color-only change, glyph widths untouched.
  morecomment={[l]{\#}},
  commentstyle=\color{cGreen},
}
\lstdefinestyle{cstyle}{
  language=C,basicstyle=\ttfamily\small,
  backgroundcolor=\color{cGreenBg!55},
  frame=leftline,framerule=2.5pt,rulecolor=\color{cGreen},
  numbers=left,numberstyle=\tiny\color{cGray},numbersep=8pt,
  xleftmargin=2.2em,framexleftmargin=0pt,framexrightmargin=0pt,
  breaklines=true,columns=fullflexible,keepspaces=true,
  showstringspaces=false,keywordstyle=\color{cBlue}\bfseries,
  commentstyle=\color{cGray}\itshape,aboveskip=8pt,belowskip=8pt,
}
% v3.36: pystyle/htmlstyle — colored twins of cstyle for the Python/HTML
% listings that used to repeat a ~200-char options blob with NO color
% (keywords rendered bold-black, comments italic-black — measured; C alone
% was colored). Geometry is copied VERBATIM from the shipped blob (frame,
% margins, skips) so swapping a call site to style=pystyle moves NOTHING —
% only keywordstyle/commentstyle/stringstyle are added (bold stays bold,
% italic stays italic, mono metrics identical => zero reflow, proven by
% per-page word-bbox diff). breaklines/breakatwhitespace stay per-site.
% v3.78: numbers=left joins both styles, with the exact numberstyle/numbersep
% keys asm/cstyle carry — fleet-wide consistency (author decision: all four
% styles numbered; "asm/cstyle numbered, py/html not" was an accident of the
% v3.36 geometry freeze, not a design). The numbers hang LEFT of the code box
% (listings sets them with \llap, so no overfull hbox can arise and the code
% x is untouched — measured: code frozen, digits at code_x-11.8pt); inside a
% colored box they are blanked like everyone else's (v3.37 hook), and the
% clean variant restores them in its frameless boxes (v3.78, above).
\lstdefinestyle{pystyle}{
  language=Python,basicstyle=\ttfamily\small,
  columns=fullflexible,keepspaces=true,
  frame=single,rulecolor=\color{cTblLine},
  numbers=left,numberstyle=\tiny\color{cGray},numbersep=8pt,
  xleftmargin=0.5em,framexleftmargin=0pt,framexrightmargin=0pt,
  showstringspaces=false,aboveskip=8pt,belowskip=8pt,
  keywordstyle=\color{cBlue}\bfseries,
  commentstyle=\color{cGray}\itshape,
  stringstyle=\color{cGreen},
}
% v4.5.0: hoisted out of the four styles so an enclosing box can override it.
% Every listing still resolves to true unless something later says otherwise,
% so shipped output is unchanged (verified byte-for-byte on all eight PDFs).
\lstset{resetmargins=true}
\lstdefinestyle{htmlstyle}{
  language=HTML,basicstyle=\ttfamily\small,
  columns=fullflexible,keepspaces=true,
  frame=single,rulecolor=\color{cTblLine},
  numbers=left,numberstyle=\tiny\color{cGray},numbersep=8pt,
  xleftmargin=0.5em,framexleftmargin=0pt,framexrightmargin=0pt,
  showstringspaces=false,aboveskip=8pt,belowskip=8pt,
  keywordstyle=\color{cBlue}\bfseries,
  commentstyle=\color{cGray}\itshape,
  stringstyle=\color{cGreen},
}
\newenvironment{asmblock}{\par\begin{otherlanguage}{english}}{\end{otherlanguage}\par}

% <hllpdf:listing-smart-quotes>
% v4.21.0: four Unicode punctuation characters must enter listings' token
% buffer. Unregistered characters bypass it: a closing quote was printed
% BEFORE the pending word, and an opening quote BEFORE a pending brace.
% Extend the standard character table, then activate only these codepoints
% in the listing's local character-table scope. Do this BEFORE the package's
% delimiter/alsoother/literate hooks so explicit user mappings still win.
% This is not universal Unicode support, nor a conversion to ASCII quotes.
% See listings developer guide 15.3.1 and references/listing-smart-quotes.md.
\makeatletter
\lst@CCPut \lst@ProcessOther {"2018}{"2019}{"201C}{"201D}\z@
\lst@AddToHookAtTop{SelectCharTable}{%
  \lst@ifactivechars
    \catcode"2018=\active \catcode"2019=\active
    \catcode"201C=\active \catcode"201D=\active
  \fi}
\makeatother
% </hllpdf:listing-smart-quotes>

% <hllpdf:listing-technical-symbols>
% v4.30: bounded extension for common technical punctuation in LTR listings.
% Without character-table registration these glyphs bypass the token buffer
% and can move before the preceding word even inside quoted strings.
% Keep the original four-quote block separate for its standing mutation test.
% This is source-faithful printing, not syntax validation or general Unicode.
\makeatletter
\lst@CCPut \lst@ProcessOther
  {"2013}{"2014}{"2026}{"2190}{"2192}{"2194}{"21D2}
  {"2212}{"221E}{"2260}{"2264}{"2265}\z@
\lst@AddToHookAtTop{SelectCharTable}{%
  \lst@ifactivechars
    \catcode"2013=\active \catcode"2014=\active \catcode"2026=\active
    \catcode"2190=\active \catcode"2192=\active \catcode"2194=\active
    \catcode"21D2=\active \catcode"2212=\active \catcode"221E=\active
    \catcode"2260=\active \catcode"2264=\active \catcode"2265=\active
  \fi}
\makeatother
% </hllpdf:listing-technical-symbols>


\usepackage[hidelinks,unicode]{hyperref}   % LAST, after colors. unicode = Hebrew bookmarks.
% v3.62: THE ISOLATION MACROS CORRUPTED EVERY PDF BOOKMARK THEY APPEARED IN.
% \en is defined (line ~155) as
%     \foreignlanguage{english}{\language=\csname l@nohyphenation\endcsname #1}
% and in a heading that body is written verbatim to the .toc and handed to
% hyperref's \pdfstringdef. hyperref strips the two CONTROL SEQUENCES (\language,
% \l@nohyphenation) but the plain `=` CHARACTER token between them is not a control
% sequence, so it survives into the PDF string: \chapter{שער \en{CNOT} והצמדה}
% produced the bookmark "שער =CNOT והצמדה". The printed TOC and the heading on the
% page are CORRECT — only the navigation pane and the viewer window title are wrong,
% which is why it went unnoticed. Measured on the shipped fleet: 70 corrupted
% bookmarks across 5 PDFs, from 177 unprotected \en-in-heading sites. Same corruption
% reached /Title and /Subject via the \hypersetup{pdftitle=...} the docs recommend.
% This hook tells hyperref to use the plain argument for all three isolation macros
% when building a PDF string. It does NOT touch the .toc (the printed TOC is
% byte-unchanged) and it does NOT mask the math signal — a heading with $...$ still
% logs "removing math shift", so BIDI-15 detection is intact.
% BONUS, measured: it also makes \code{} in an UNSTARRED heading COMPILE. BLD-10
% recorded that as needing \texorpdfstring or a starred heading; the real root cause
% is hyperref's \pdfstringdef scan of \code's \ifmmode...\else...\fi body — which is
% exactly why the identically-bodied \texttt survived (hyperref pre-neutralizes
% \texttt in PDF strings) and why \DeclareRobustCommand could never help.
% v3.64: \def\\{ } added. A manual line break in a heading WELDS the two words
% together in the bookmark: \chapter{שורה ראשונה\\שורה שנייה} produced the outline
% entry 'שורה ראשונהשורה שנייה' — "ראשונהשורה" is not a word. hyperref drops the
% \\ and puts nothing in its place, so the break that looks like a space on the
% page is a zero-width join in the navigation pane. The spaced form
% 'מילה אחת \\ מילה שנייה' escapes with a harmless DOUBLE space ('אחת  מילה'),
% which is why this never got noticed. Both are measured on a real outline.
% 0 live instances in the corpus — but SKILL.md and content-style.md both RECOMMEND
% a manual heading break, so this is a primed trap rather than a dormant one, and a
% one-token preamble fix retires it for good instead of asking authors to remember.
% Mapping \\ to a space fixes the WELD, which is the actual defect (a nonexistent
% word in the navigation pane). It does not normalise whitespace: measured after the
% fix, the welded form becomes 'שורה ראשונה שורה שנייה' (correct) and the already-
% spaced form goes from two spaces to three. That is cosmetic — viewers render it as
% a wider gap, and no word is corrupted. A per-heading \texorpdfstring is still the
% way to get an exactly-tuned bookmark string.
\pdfstringdefDisableCommands{\def\en#1{#1}\def\h#1{#1}\def\doclang#1{#1}\def\code#1{#1}\def\codeline#1{#1}\def\\{ }}
% v3.70: RESTORE \url LINE-BREAKING. url.sty freezes its break penalties at LOAD
% time (\mathchardef\UrlBreakPenalty=\binoppenalty), and v3.45's global
% \binoppenalty=10000 (inline-math atomicity) therefore priced EVERY breakpoint
% inside \Url at 10000 = never break. Measured: a long \url in Hebrew prose is
% one unbroken line, ~100pt overfull, invading the margin band the ink gates
% deliberately tolerate. With the two lines below: 0 overfull, the URL breaks
% after '/', the RTL split is CORRECT (fragment order forward), and each line
% gets its own link annotation exactly covering its fragment (measured).
% Confined to \Url's group — inline-math atomicity (TYP-7) is untouched.
\mathchardef\UrlBreakPenalty=200
\mathchardef\UrlBigBreakPenalty=100
\ifx\lakitmainlanguage\hllpdfenglish
  \hypersetup{pdflang=en}
\else
  \hypersetup{pdflang=he}
\fi  % /Lang is the document base locale (not per-span accessibility tagging).
% Set a per-document \hypersetup{pdftitle={...}} too — see SKILL.md.

% v3.52 (round-2 audit): two RTL bibliography/cross-ref defects, both invisible to
% every gate and both affecting the FIRST author to use the feature.
% (a) \@biblabel's default [#1] brackets are bidi-neutral and mirror to ]1[ under
% the RTL bibliography base — every \thebibliography reference list shows ]1[.
% LTR-isolating the label restores [1]. (The render gate's mirror rule fingerprints
% round parens )1(, so square brackets slipped through.)
\makeatletter
\renewcommand*{\@biblabel}[1]{\foreignlanguage{english}{[#1]}}
% v3.54: the SAME neutral-bracket trap on the INLINE side — a multi-key \cite{a,b}
% under the RTL base lays its European-number list right-to-left, so [1, 2] renders
% "[2,1]" and [1, 2, 3] renders "[3,2,1]" (verified by glyph geometry — the reader
% sees the sources in REVERSE order; a note form \cite[p.5]{a} also reorders to
% "[p. 5,1]"). LTR-isolate the whole bracket group (number list AND optional note),
% exactly like \@biblabel above: [1, 2] / [1, 2, 3] / [1, p. 5] all render correct.
% Placed AFTER hyperref so it wins; hyperref hyperlinks each number inside #1 (via
% \@citex), which this wrapper leaves intact. Mirrors the LaTeX kernel \@cite shape.
% v4.9.1: directional isolation must NOT switch the language here. Babel's
% safe-active citation wrapper passes its key list through \bbl@tempa;
% \foreignlanguage reuses that scratch macro before the deferred citation
% loop expands it. The old line therefore cited the literal key "english".
% If such a key existed, even the undefined-reference gate stayed green.
% \babelsublr isolates the bracket group without selecting a locale or
% changing that key list. Preserve native links, key order and optional notes.
\def\@cite#1#2{\babelsublr{[{#1\if@tempswa , #2\fi}]}}
\makeatother

% v4.11.0: a reference inside align's \text inherited HEBdisptag=true,
% so its enclosing hbox mirrored the brackets despite a clean compile.
% Isolate the complete printed reference without selecting a Babel language.
% Also neutralize the display-only normalfont direction reset locally. A
% direction-only patch inverted 12 to 21 inside equation in the measured probe.
% Plain maketag is local to the reference: mathtools' display-side wrapper
% otherwise globally consumes manual_tag while formatting a nested reference.
% Preserve each original command (including the selective auxiliary write),
% current tag forms, native links, and showonlyrefs on/off transitions.
\makeatletter
\NewCommandCopy{\HEB@nativeeqref}{\eqref}
\newcommand{\HEB@eqrefisolate}[1]{%
  \babelsublr{\HEBdisptagfalse\let\bbl@eqnodir\relax\let\maketag@@@\ORIGmaketag@@@#1}}
\protected\def\eqref#1{\HEB@eqrefisolate{\HEB@nativeeqref{#1}}}
\ExplSyntaxOn
\NewCommandCopy{\HEB@nativeMTeqref}{\MT_extended_eqref:n}
\cs_set_protected:Npn \MT_extended_eqref:n #1
  {\HEB@eqrefisolate{\HEB@nativeMTeqref{#1}}}
\ExplSyntaxOff
\makeatother

% v4.11.0: custom caption counters are not part of hyperref's English
% locale, and Hebrew has no upstream hyperref locale in this toolchain.
% Restore the appropriate names at every Babel language switch, including
% return from a nested inline island. Do not localize the destination keys.
\makeatletter
\newcommand{\HEB@refnames}{%
\def\chapterautorefname{פרק}%
\def\sectionautorefname{סעיף}%
\def\subsectionautorefname{תת־סעיף}%
\def\subsubsectionautorefname{תת־תת־סעיף}%
\def\paragraphautorefname{פסקה}%
% v4.31: hyperref uses a distinct counter name at the deepest standard level.
\def\subparagraphautorefname{תת־פסקה}%
\def\equationautorefname{משוואה}%
\def\figureautorefname{איור}%
\def\tableautorefname{טבלה}%
\def\partautorefname{חלק}%
\def\appendixautorefname{נספח}%
\def\footnoteautorefname{הערת שוליים}%
\def\itemautorefname{סעיף}%
\def\Itemautorefname{סעיף}%
\def\figautorefname{איור}%
\def\tblautorefname{טבלה}%
\def\pageautorefname{עמוד}% 
}
\addto\extrashebrew{\HEB@refnames}
\newcommand{\ENG@refnames}{%
  \def\subparagraphautorefname{subparagraph}%
  \def\chapterautorefname{chapter}%
  \def\sectionautorefname{section}%
  \def\subsectionautorefname{subsection}%
  \def\subsubsectionautorefname{subsubsection}%
  \def\paragraphautorefname{paragraph}%
  \def\equationautorefname{Equation}%
  \def\figureautorefname{Figure}%
  \def\tableautorefname{Table}%
  \def\partautorefname{Part}%
  \def\appendixautorefname{Appendix}%
  \def\footnoteautorefname{footnote}%
  \def\itemautorefname{item}%
  \def\figautorefname{Figure}%
  \def\tblautorefname{Table}%
  \def\Itemautorefname{item}%
  \def\pageautorefname{page}%
}
\addto\extrasenglish{\ENG@refnames}
\ifx\lakitmainlanguage\hllpdfenglish\ENG@refnames\else\HEB@refnames\fi
\makeatother

% ============================================================================
% Optional running head (v3.39): \lakithead{<book title>} — OFF by default.
% A document opts in with ONE line after % ============================================================================
% preamble.tex — tested LuaLaTeX preamble for Hebrew + math + figures PDFs (v4.32.0)
% Battle-tested: compiled an 80-page Hebrew quantum-computing summary AND a
% 197-page Hebrew computer-architecture anthology (TikZ diagrams, code listings,
% colored boxes) — both with 0 missing characters.
%
% Usage in your document:
%   \input{preamble.tex}     % (or paste these lines at the top)
%   \begin{document}  ...  \end{document}
%
% Compile with:  lualatex -interaction=nonstopmode doc.tex
%   Repeat until two consecutive PDFs are byte-identical (scripts/build.sh
%   does this) — a TOC document converges only on pass 3+; after two passes
%   the TOC page numbers are STALE (measured; see SKILL.md step 3).
% Font binaries are external dependencies, pinned by assets/fonts/SHA256SUMS.
% scripts/setup_fonts.sh imports a verified local installation or prior skill.
%
% PUBLIC DOCUMENT-SHELL / LAYOUT CONTRACT (introduced v4.3.1; architecture-owned v4.7.1). These hooks may be
% defined BEFORE \input{preamble.tex}; if omitted, the historical canonical
% defaults below are exact.  They deliberately control only the document shell
% and page layout.  RTL/math semantics stay in this preamble, while visual
% presentation themes (for example clean_theme.tex) load AFTER it.
%
%   \def\lakitdocumentclass{article}
%   \def\lakitclassoptions{11pt,a4paper}
%   \def\lakitgeometryoptions{margin=1.8cm,headheight=14pt}
%   \def\lakitlinespread{1.15}
%   \def\lakitparagraphstyle{block} % letters/briefs; indented for continuous prose
%   \def\lakittocdepth{1}
%   \def\lakitsecnumdepth{2}
%   \input{preamble.tex}
%
% Supported class contract: report or article.  Alternate paper sizes must also
% be declared to build.sh through EXPECTED_PAPER; the PDF gate never infers them.
% Do NOT use these hooks as a theme mechanism: theme overlays must not own page
% geometry, class selection, or direction semantics. A declared presentation-policy
% overlay may intentionally change numbering visibility after the preamble.
% ============================================================================
% Runtime architecture guard: generated from scripts/architecture_contract.py.
% The readiness sentinel remains delayed until the END of this file.
% <hllpdf:architecture-guard>
\ifdefined\hllpdfcoreloaded
  \errmessage{hllpdf: preamble.tex loaded more than once}
  \expandafter\endinput
\fi
% </hllpdf:architecture-guard>

\ifdefined\lakitdocumentclass\else\def\lakitdocumentclass{report}\fi
\ifdefined\lakitclassoptions\else\def\lakitclassoptions{12pt,a4paper}\fi
\ifdefined\lakitgeometryoptions\else\def\lakitgeometryoptions{margin=2.3cm,headheight=14pt}\fi
\ifdefined\lakitlinespread\else\def\lakitlinespread{1.3}\fi
\ifdefined\lakittocdepth\else\def\lakittocdepth{1}\fi
\ifdefined\lakitsecnumdepth\else\def\lakitsecnumdepth{2}\fi
% Paragraph organization belongs to the document layout, independently of its
% language and theme. Preserve historical class dimensions unless explicitly
% selecting block paragraphs; never suppress indentation with an everypar hack.
\ifdefined\lakitparagraphstyle\else\def\lakitparagraphstyle{indented}\fi
\edef\lakitparagraphstyle{\lakitparagraphstyle}
\def\hllpdfparagraphblock{block}
\def\hllpdfparagraphindented{indented}
\ifx\lakitparagraphstyle\hllpdfparagraphblock\else
  \ifx\lakitparagraphstyle\hllpdfparagraphindented\else
    \errmessage{hllpdf: paragraph style must be block or indented}
  \fi
\fi
% v4.3.1: the public class hook has a deliberately narrow support contract.
% Before this guard, \def\lakitdocumentclass{book} silently compiled even though
% every heading/TOC regression test targets report/article; an unsupported shell
% could therefore look "supported" until a subtle class-specific defect shipped.
% Fail before \documentclass rather than letting an undocumented combination enter
% the mechanism core. Plain-TeX comparison is used because LaTeX packages are not
% available yet at this point.
\edef\lakitvalidatedclass{\lakitdocumentclass}
\def\lakitreportclass{report}
\def\lakitarticleclass{article}
\ifx\lakitvalidatedclass\lakitreportclass
\else\ifx\lakitvalidatedclass\lakitarticleclass
\else
  \errmessage{hllpdf: unsupported class \lakitdocumentclass; use report or article}
\fi\fi
\let\lakitvalidatedclass\relax
\let\lakitreportclass\relax
\let\lakitarticleclass\relax
% Language intent is independent of document class, theme and content domain.
% The two dominant-language profiles differ in authoring intent, not direction.
% A balanced bilingual document must explicitly choose a navigation/base locale.
% <hllpdf:language-profile>
\providecommand{\lakitlanguageprofile}{hebrew}
\edef\lakitlanguageprofile{\lakitlanguageprofile}
\def\hllpdfhebrew{hebrew}\def\hllpdfenglish{english}
\def\hllpdfhebrewmixed{hebrew-with-english}
\def\hllpdfenglishmixed{english-with-hebrew}
\def\hllpdfbilingual{bilingual}
\let\hllpdfexpectedlanguage\relax
\ifx\lakitlanguageprofile\hllpdfhebrew
  \let\hllpdfexpectedlanguage\hllpdfhebrew
\else\ifx\lakitlanguageprofile\hllpdfhebrewmixed
  \let\hllpdfexpectedlanguage\hllpdfhebrew
\else\ifx\lakitlanguageprofile\hllpdfenglish
  \let\hllpdfexpectedlanguage\hllpdfenglish
\else\ifx\lakitlanguageprofile\hllpdfenglishmixed
  \let\hllpdfexpectedlanguage\hllpdfenglish
\else\ifx\lakitlanguageprofile\hllpdfbilingual
  \ifdefined\lakitmainlanguage\else
    \errmessage{hllpdf: bilingual profile requires explicit lakitmainlanguage}
  \fi
\else
  \errmessage{hllpdf: unsupported language profile \lakitlanguageprofile}
\fi\fi\fi\fi\fi
\ifdefined\lakitmainlanguage
  \edef\lakitmainlanguage{\lakitmainlanguage}
\else
  \let\lakitmainlanguage\hllpdfexpectedlanguage
\fi
\ifx\lakitmainlanguage\hllpdfhebrew\else
  \ifx\lakitmainlanguage\hllpdfenglish\else
    \errmessage{hllpdf: main language must be hebrew or english}
  \fi
\fi
\ifx\hllpdfexpectedlanguage\relax\else
  \ifx\lakitmainlanguage\hllpdfexpectedlanguage\else
    \errmessage{hllpdf: main language conflicts with language profile}
  \fi
\fi
\typeout{HLLLANG|\lakitlanguageprofile|\lakitmainlanguage}
% </hllpdf:language-profile>
\documentclass[\lakitclassoptions]{\lakitdocumentclass}
\usepackage[no-math]{fontspec}        % (1) ★ no-math — THE critical line (see recipe.md)
% layout=lists asks Babel to adapt list parshape when paragraph and page
% directions differ. Without it, minority-language list bodies lose their
% indentation and labels hang outside the text block, despite a clean log.
% The upstream mechanism preserves ordinary same-direction list geometry.
\usepackage[bidi=basic,layout=lists]{babel}
\ifx\lakitmainlanguage\hllpdfenglish
  \babelprovide[import]{hebrew}
  \babelprovide[import,main]{english}
\else
  \babelprovide[import,main]{hebrew}
  \babelprovide[import]{english}
\fi
% v4.22: bind canonical families and fallbacks to the selected installation.
% Lookup by basename or family alone is not proof of the file actually used.
% <hllpdf:font-binding>
\directlua{
  local dir = os.getenv("HEBREW_LATEX_FONT_DIR")
  if not dir or dir == "" then
    local data = os.getenv("XDG_DATA_HOME")
    if not data or data == "" or not (data:sub(1,1) == "/") then
      data = (os.getenv("HOME") or "") .. "/.local/share"
    end
    dir = data .. "/fonts/hebrew-lualatex-pdf"
  end
  if not (dir:sub(1,1) == "/") then dir = lfs.currentdir() .. "/" .. dir end
  if dir:find("[",1,true) or dir:find("]",1,true) or dir:find(string.char(0),1,true)
     or dir:find(string.char(10),1,true) or dir:find(string.char(13),1,true) then
    tex.error("hllpdf: unsupported character in font directory")
  end
  if not (lfs.attributes(dir,"mode") == "directory") then
    tex.error("hllpdf: font directory is missing")
  end
  hllpdf_font_dir = dir:gsub("/+$", "") .. "/"
  tex.sprint(string.char(92) .. "def" .. string.char(92) .. "hllpdffontdir{")
  tex.sprint(-2,hllpdf_font_dir)
  tex.sprint("}")
  luaotfload.add_fallback("hebfb",{"[" .. hllpdf_font_dir .. "DejaVuSerif.ttf]:mode=node;"})
  luaotfload.add_fallback("hebfbbf",{"[" .. hllpdf_font_dir .. "DejaVuSerif-Bold.ttf]:mode=node;"})
}
% </hllpdf:font-binding>
% v3.22: Frank Ruhl Libre ships NO italic face. Instead of silently-upright
% \emph (v<=3.20) or bold-as-emphasis (v3.21), synthesize a slanted face with
% FakeSlant — real, VISIBLE italic emphasis with correct \emph nesting
% semantics (inner \emph toggles back upright). Verified empirically at
% 300dpi: slant clearly visible, 0 errors, 0 missing glyphs.
\babelfont{rm}[Path={\hllpdffontdir},RawFeature={fallback=hebfb},
  BoldFont={FrankRuhlLibre-Bold.ttf},
  BoldFeatures={RawFeature={fallback=hebfbbf}},   % v3.56 — see hebfbbf above
  ItalicFont={FrankRuhlLibre-Regular.ttf}, ItalicFeatures={FakeSlant=0.22},
  % Bold-italic uses the weight-matched fallback and a synthetic slant.
  BoldItalicFont={FrankRuhlLibre-Bold.ttf}, BoldItalicFeatures={RawFeature={fallback=hebfbbf},FakeSlant=0.22}
  ]{FrankRuhlLibre-Regular.ttf}
% Sans fallbacks preserve family and weight for the few missing codepoints.
\directlua{luaotfload.add_fallback("hebfbsans",{"[" .. hllpdf_font_dir .. "DejaVuSans.ttf]:mode=node;"})}
\directlua{luaotfload.add_fallback("hebfbsansbf",{"[" .. hllpdf_font_dir .. "DejaVuSans-Bold.ttf]:mode=node;"})}
\babelfont{sf}[Path={\hllpdffontdir},RawFeature={fallback=hebfbsans},
  BoldFont={Heebo-Bold.ttf},
  BoldFeatures={RawFeature={fallback=hebfbsansbf}}]{Heebo-Regular.ttf}
% v3.24: Scale=MatchLowercase — DejaVu Sans Mono has an outsized x-height, so
% unscaled inline code visually SHOUTS next to the delicate Frank Ruhl body
% (measured on a 197-page build: scaled, a code atom sits in the line at the
% body's weight, and the document tightened by ~5%). Applies to listings too.
\babelfont{tt}[Path={\hllpdffontdir},RawFeature={fallback=hebfb},Scale=MatchLowercase,
  BoldFont={DejaVuSansMono.ttf},
  BoldFeatures={RawFeature={fallback=hebfb},FakeBold=1.5},
  ItalicFont={DejaVuSansMono.ttf},
  ItalicFeatures={RawFeature={fallback=hebfb},FakeSlant=0.2},
  BoldItalicFont={DejaVuSansMono.ttf},
  BoldItalicFeatures={RawFeature={fallback=hebfb},FakeBold=1.5,FakeSlant=0.2}
  ]{DejaVuSansMono.ttf}

% Explicit multilingual/English profiles use real Latin faces and italics.
% The legacy/default Hebrew profile retains its established measured typography.
% Latin Modern is provided by the TeX distribution, not the pinned font bundle;
% recorder/runtime evidence reports it as an additional toolchain resource.
\ifx\lakitlanguageprofile\hllpdfhebrew\else
  \babelfont[english]{rm}[Ligatures=TeX,RawFeature={fallback=hebfb},
    BoldFeatures={RawFeature={fallback=hebfbbf}},
    BoldItalicFeatures={RawFeature={fallback=hebfbbf}}]{Latin Modern Roman}
  \babelfont[english]{sf}[Ligatures=TeX,RawFeature={fallback=hebfbsans},
    BoldFeatures={RawFeature={fallback=hebfbsansbf}},
    BoldItalicFeatures={RawFeature={fallback=hebfbsansbf}}]{Latin Modern Sans}
\fi

\usepackage{amsmath,amssymb,mathtools}
% amsthm supplies theorem/proof machinery. Environment names remain
% document-specific; the generic proof label is localized below.
\usepackage{amsthm}
% Proof names follow the surrounding prose, including return from an island.
\addto\extrashebrew{\renewcommand{\proofname}{הוכחה}}
\addto\extrasenglish{\renewcommand{\proofname}{Proof}}
\ifx\lakitmainlanguage\hllpdfenglish
  \renewcommand{\proofname}{Proof}
\else
  \renewcommand{\proofname}{הוכחה}
\fi
% v3.24.1: equation tags under babel-bidi. A tag re-emitted by mathtools'
% showonlyrefs inside align inherits the RTL base and its parens MIRROR —
% ")2.1(" instead of "(2.1)" (plain align/equation tags are unaffected;
% verified on a minimal case + a 198-page build). LTR-isolate at the tag-BOX
% level: mathtools' show/hide + numbering logic is untouched. Do NOT override
% \tagform@ instead — that clobbers the showonlyrefs mechanism (measured:
% every tag shows, numbering breaks).
% v3.28.1: the v3.24.1 wrap must apply ONLY to tags of display environments.
% Inline \eqref in Hebrew PROSE also goes through \maketag@@@, and there the
% english-wrap MIRRORS the parens — ")1(" — while the unwrapped original is
% correct (proven on a 2x2 isolation matrix: patch x showonlyrefs; regressed
% every prose \eqref since v3.24.1, caught by the render gate's R4). The two
% call sites are indistinguishable by \ifmmode (tag boxes are text mode), so
% the display environments themselves raise a flag; \eqref in prose sees it
% false and gets the original untouched \maketag@@@.
\usepackage{etoolbox}

% babel bidi=basic (including upstream 26.9) assumes every automatic
% discretionary has a one-glyph replacement. LuaTeX can legally produce one
% with replace=nil after a hyphenated Latin value; Babel then indexes nil and the
% whole build dies. Do the narrow upstream-compatible correction: while Babel
% classifies the list, present only that malformed disc as an ordinary
% discretionary, then restore its subtype. Babel still reorders the ENTIRE box;
% there is no pcall and no successful fallback to an unreordered node list.
% If Babel throws for any other reason, the build remains fatal. The synthetic
% nil-disc probe and the real table gallery in test_all.sh pin both layers.
\AtBeginDocument{\directlua{
  if Babel and Babel.bidi and not Babel.heb_nil_disc_compat then
    Babel.heb_nil_disc_compat = true
    local original = Babel.bidi
    local disc_id = node.id('disc')
    Babel.bidi = function(head, ispar, hdir)
      local hidden = {}
      for item in node.traverse_id(disc_id, head) do
        if item.subtype == 2 and (not item.replace or not item.replace.char) then
          table.insert(hidden, item)
          item.subtype = 0
        end
      end
      local result = original(head, ispar, hdir)
      for _, item in ipairs(hidden) do item.subtype = 2 end
      return result
    end
  end
}}
\makeatletter
\newif\ifHEBdisptag
\let\ORIGmaketag@@@\maketag@@@
\def\maketag@@@#1{\ifHEBdisptag\ORIGmaketag@@@{\foreignlanguage{english}{#1}}\else\ORIGmaketag@@@{#1}\fi}
\@for\HEB@env:=equation,equation*,align,align*,gather,gather*,multline,multline*,flalign,flalign*,alignat,alignat*,eqnarray,eqnarray*\do{%
  \expandafter\AtBeginEnvironment\expandafter{\HEB@env}{\HEBdisptagtrue}}
\makeatother
% English page + Hebrew prose: Babel's equation-only tag flip uses a primitive
% leqno path that shifts the formula to the left edge and the tag into the body
% of an English page. Align's full-width alignment is unaffected. Give only this
% display a physical LTR container and native amsmath eqno/leqno. The engine then
% owns real tag width, collision handling and labels; no proxy tag box hides ink
% from its fit calculation. Babel's saved \bbl@eqnodir still restores Hebrew text
% inside math, and the prose locale/math direction are never switched.
% Close the prior RTL paragraph before normalization. After the environment's
% group restores its directions, close TeX's empty resumed display paragraph and
% use native LaTeX continuation handling: same-paragraph prose has no indent,
% while an authored blank line starts an ordinary indented RTL paragraph.
% Register after Babel's setup; remain local to equation(*) in a main-English
% document's actual RTL prose. Native Hebrew, ordinary English and explicit
% fleqn retain their original paths. No page/display dimensions are assigned.
% <hllpdf:minority-rtl-equation-geometry>
\makeatletter
\newcommand{\HEB@minorityprosecontinuation}{\par\@doendpe}
\newcommand{\HEB@minorityequationgeometry}{%
  \ifx\lakitmainlanguage\hllpdfenglish
    \ifnum\pardirection=\@ne
      \if@fleqn\else
        \par\aftergroup\HEB@minorityprosecontinuation
        \iftagsleft@\let\veqno\@@leqno
        \else\let\veqno\@@eqno\fi
        \bodydir TLT\pardir TLT\textdir TLT\relax
      \fi
    \fi
  \fi}
\AtBeginDocument{%
  \AddToHook{env/equation/begin}{\HEB@minorityequationgeometry}%
  \AddToHook{env/equation*/begin}{\HEB@minorityequationgeometry}}
\makeatother
% </hllpdf:minority-rtl-equation-geometry>
% v3.63: \shortintertext rendered Hebrew MIRROR-REVERSED and flush LEFT — an entire
% commentary line inside an align printed backwards. babel's bidi=basic patches
% \intertext@ only (luababel.def), and mathtools' \shortintertext uses the identical
% \vbox{\normalbaselines … \noindent\ignorespaces#1\par} skeleton but is never
% patched, so its paragraph keeps pardir/textdir LTR: no bidi reordering (logical
% order on the page) and flush-left. Measured against \intertext with the same
% Applying babel's own replacement to \shortintertext@ makes the two identical.
% \h{}/\mbox fix the letter order but leave the line flush LEFT, so they are not
% substitutes. Guarded + \AtBeginDocument because mathtools \let-copies at load time.
% NB \ExplSyntaxOn is REQUIRED: mathtools' internals carry expl3 names, and without
% the expl3 catcodes the `_` is read as a subscript and the patch detonates far away
% with "! Illegal unit of measure (pt inserted)" at the \end{align} (measured).
\makeatletter
\AtBeginDocument{%
  \@ifundefined{shortintertext@}{}{%
    \ExplSyntaxOn
    \bbl@sreplace\shortintertext@{\normalbaselines}%
      {\normalbaselines\ifx\bbl@eqnodir\relax\else\bbl@pardir\@ne\bbl@eqnodir\fi}%
    \ExplSyntaxOff}}
\makeatother
% geometry inherits the canonical A4 sheet from \documentclass; build.sh then
% verifies every rendered MediaBox and CropBox, so a silent Letter fallback is
% impossible to publish as a clean default build.
\usepackage[\lakitgeometryoptions]{geometry}
\linespread{\lakitlinespread}       % layout hook; canonical default remains 1.3
\ifx\lakitparagraphstyle\hllpdfparagraphblock
  \setlength{\parindent}{0pt}
  \setlength{\parskip}{6pt plus 1pt minus 1pt}
\fi
% These dimensions remain ordinary scoped TeX lengths: authors can customize
% them after core input. Language switches inherit them; lists keep their own
% indentation and narrow reference layouts establish their local paragraph gaps.
\typeout{HLLPAR|\lakitparagraphstyle}
\usepackage{xcolor}
\usepackage{mdframed}
\usepackage{listings}                 % code listings (assembly / C / pseudo-code)
\usepackage{tikz}
\usetikzlibrary{arrows.meta,positioning,calc,shapes.geometric,fit,backgrounds}
\usepackage{pgfplots}\pgfplotsset{compat=1.18}
\usepackage{quantikz}                 % quantum circuits (remove if unused)
\usepackage{booktabs}
\usepackage{array}
\usepackage{float}                    % [H] = place table exactly here (reference doc: no drifting)
\usepackage{enumitem}
\usepackage{needspace}

% ============================================================================
% ★★ v3.4 — THE DIRECTION-ISOLATION TRIO (the heart of mixing Hebrew + Latin)
% Mental model:
%   \h{...}    Hebrew (RTL) island      — Hebrew text inside an LTR context
%                                          (TikZ nodes, the english `tikzpic`).
%   \en{...}   English (LTR) isolate    — Latin text / parentheses inside a
%                                          Hebrew (RTL) context. KEEPS PARENS
%                                          AND ORDER CORRECT.
%   \code{...} \en + monospace          — code inside Hebrew / inside a frametitle.
% CRITICAL: never wrap Latin in \h{} (it is for HEBREW). Inside a TikZ node,
% Latin in \h{} renders REVERSED (Datapath->htapataD, CPU->UPC) and parens
% mirror. Isolate every Latin/paren chunk with \en{} or \code{}. See
% references/figures-boxes-listings.md.
% ---- TEXT INSIDE MATH (v3.8) ----------------------------------------------
% Put ANY text that sits inside math in a \text{}. The ONLY thing that is
% invalid in math is a bare \h/\en (= \foreignlanguage) used DIRECTLY, outside
% \text -> "\rmfamily invalid in math mode". Inside \text:
%   Hebrew label           -> \text{...}            (bare Hebrew renders fine;
%                              \h is optional — BUT NOT ALWAYS HARMLESS (v3.67):
%                              if the \h group OPENS with a nested island
%                              ($..$, \en{}, \code{}, \texttt{}, \mbox{}) the
%                              WHOLE group loses bidi and the Hebrew prints in
%                              LOGICAL order = gibberish. Measured:
%                                \text{\h{$f$ קבועה}}  -> reads "העובק"  BROKEN
%                                \text{\h{קבועה $f$}}  -> correct
%                                \text{$f$ קבועה}      -> correct (no \h)
%                              Hebrew first inside \h, or drop the \h. Flagged
%                              H-OPENS-WITH-LTR-GROUP.)
%   plain Latin word        -> \text{...}            (no marker needed)
%   Latin WITH ( ) or [ ]   -> \text{\en{...}}       (\en forces LTR; a plain
%                              \text mirrors the parens in RTL -> "Spec (TNR(" )
%   f(Hebrew/text arg)      -> \text{f}(\text{...})  -- keep the parens at MATH
%                              level, NOT \text{f(arg)} (parens inside the box mirror)
%   code in math            -> \mathtt{...}  (or \texttt{}, math-guarded below)
% NEVER write \en{...} or \h{...} straight into $...$ / \[..\] (outside \text);
% always via \text{}. (\mbox{...} also works but does not scale in scripts.)
% ============================================================================
\newcommand{\h}[1]{\foreignlanguage{hebrew}{#1}}
% v3.42.1 — \en islands do NOT hyphenate (auto OR at explicit hyphens). English hyphenation patterns were
% tearing identifiers and terms ACROSS THE RTL SEAM ("IO Regis-/ters",
% "In-/put", "Band-/width", "Per-/formance", "Multi-/ple" — all shipped on
% real pages; probe-proven: "ALUS-/rcB", "Mem-/ToReg"). An island is a label
% or a term, not English prose — a hyphen tear there is never right. A/B on
% a full 197pp lakit: identical overfull sets, same page count, ~11
% paragraphs healed, one mild underfull (badness 3525) that replaced a
% mutilated "Regis-/ters". Long English PROSE belongs in the englishprose
% environment below, which restores explicit-hyphen breaks while its paragraph
% is still active. A raw otherlanguage block inherits the global prohibition.
% (\l@nohyphenation needs \csname:
% '@' is not a letter here.)
% v3.43.2 — \exhyphenpenalty moved OUT of \en to a GLOBAL assignment below.
% Render-proven: a group-local \exhyphenpenalty is DEAD CODE — TeX reads the
% line-breaking parameters at \par time, after the \en group has closed, so
% \en{write-through} / plain not-taken still tore at the "-" ("write-/
% through" shipped on a real exam page). \language is different: it is
% recorded in the node list as text is added, so the in-group nohyphenation
% assignment above DOES work (and stays). The global penalty was A/B-measured
% on a 197pp lakit: identical overfull set (25), identical page count, the
% not-/taken tear healed. Hebrew is unaffected in practice — Hebrew compounds
% use maqaf (־), not ASCII "-".
\newcommand{\en}[1]{\foreignlanguage{english}{\language=\csname l@nohyphenation\endcsname #1}}
% Babel runs extras locally for inline islands and paragraph language changes.
% An English main document must not inherit the Hebrew no-hyphen paragraph rule.
\addto\extrashebrew{\exhyphenpenalty=10000\relax}
\addto\extrasenglish{\exhyphenpenalty=50\relax}
\ifx\lakitmainlanguage\hllpdfenglish
  \exhyphenpenalty=50\relax
\else
  \exhyphenpenalty=10000\relax
\fi
% Full English paragraphs need ordinary explicit-hyphen breaks. Use this
% environment instead of a raw otherlanguage block; the paragraph ends while
% the local penalty is still active.
\newenvironment{englishprose}
  {\par\begin{otherlanguage}{english}\exhyphenpenalty=50\relax}
  {\par\end{otherlanguage}}
\newenvironment{hebrewprose}
  {\par\begin{otherlanguage}{hebrew}\exhyphenpenalty=10000\relax}
  {\par\end{otherlanguage}}
% Use for shell text whose locale must not depend on the last body paragraph.
\DeclareRobustCommand{\doclang}[1]{\foreignlanguage{\lakitmainlanguage}{#1}}

% ============================================================================
% ★★ — fixes for defects that are INVISIBLE in the log and in the checks
% and only surface on a VISUAL RENDER of the PDF. (Each verified empirically.)
% ----------------------------------------------------------------------------
% v3.22: \emph is now REAL slanted emphasis (FakeSlant italic face defined on
% the \babelfont{rm} line above), replacing the v3.21 \emph->\textbf hack: it
% preserves \emph's toggle semantics and keeps bold free for its own role.
% Prefer bold-as-emphasis for your doc? one line: \renewcommand{\emph}[1]{\textbf{#1}}
% page number: force the digits LTR so they never reverse (18 -> 81) in the RTL
%   footer. That reversal is a NON-DETERMINISTIC bidi glitch (observed once in
%   ~80 pp); this LuaLaTeX build is not byte-reproducible, so the defensive
%   LTR wrap is the only reliable guard. Latin digits are always LTR anyway.
\renewcommand{\thepage}{\foreignlanguage{english}{\arabic{page}}}
% v3.70: \pagenumbering \gdef's \thepage and silently DESTROYS the wrap above —
% after \pagenumbering{arabic} the folio guard is gone for the whole rest of the
% document. Two collections discovered this independently and hand-restore the
% wrap in main.tex with the same comment; that knowledge now lives here instead.
% The wrapper re-applies the LTR wrap after the kernel call, for any style.
\makeatletter
\let\lakin@orig@pagenumbering\pagenumbering
\renewcommand{\pagenumbering}[1]{%
  \lakin@orig@pagenumbering{#1}%
  \renewcommand{\thepage}{\foreignlanguage{english}{\csname @#1\endcsname\c@page}}}
\makeatother
% emergencystretch: Hebrew does not hyphenate, so an unbreakable LTR island
%   (\en{...} / an inline math run) landing at the margin can leave a small
%   (~2-13pt) overfull with no legal breakpoint. This rescues ONLY those
%   inline-prose tie-breaks; it is a ceiling used only on offending lines and
%   does NOT touch display-math / listing / table overflow (those still surface
%   in the log for the content-level overflow pass). Verified: baseline overfull
%   -> 0 with this set, with no change to pagination or underfull.
\setlength{\emergencystretch}{3em}
% v3.40: inline math (almost) never breaks across lines. TeX's default
%   (binop 700 / rel 500) lets an inline run split at an operator —
%   $(1.f)\cdot 2^{e}$ shipped as "(1.f)·" ending one line and "2^e" opening
%   the next (reader-reported). In LTR that split is merely ugly; under an
%   RTL base the carried operator lands on the visually WRONG side of the
%   break and the formula reads backwards.
%   WHY 10000 (v3.45.0 — the display-first resolution). Inline math is
%   FORBIDDEN from breaking, so a torn island under an RTL base is
%   structurally impossible. The v3.40 cycle used 9999 ("break only as a last
%   resort") because a bare 10000 left ~30-39 WIDE inline equations
%   overflowing (worst 90.7pt) — but 9999's "last resort" fires OFTEN in
%   tight, non-hyphenating Hebrew: a reader found torn islands across the
%   fleet (~150 candidates; e.g. "P(y |" ending one line and "x_0)" opening
%   the next, read backwards). The step the 9999 cycle missed: every one of
%   those ~30 overflows is a wide equation that BELONGS in display. Convert
%   them to display FIRST, then 10000 gives 0 tears AND 0 overflows (measured,
%   12/12 fleet). Doctrine: inline math is ATOMIC; a wide equation goes in a
%   display, never inline. See content-style ("inline math is atomic").
\binoppenalty=10000 \relpenalty=10000
% v3.38: uniform sentence spacing. TeX's default (\nonfrenchspacing, sfcode
% 3000) is the OLD Anglo-American wide-space-after-period convention — Hebrew
% typography never had it, and modern Latin practice (Chicago et al.) dropped
% it too. Measured on real pages: post-period gaps up to 1.9-2.1x the line's
% mean word gap — the dominant "rivers" source in justified Hebrew prose.
% Survives babel language switches (probed: sfcode stays french inside \en{}
% islands and otherlanguage blocks).
\frenchspacing
% v3.38: widow/orphan protection. LaTeX's default penalties (150/150/50)
% allow a lone line at a page top/bottom almost freely. 10000 = forbid;
% this is SAFE here precisely because the document is \raggedbottom (the
% demoted line just moves on, no vertical stretch fight) — under flushbottom
% these values would cause underfull pages. raggedbottom is the report-class
% default; made explicit so a future class/option change cannot silently
% flip it and turn the 10000s harmful.
\raggedbottom
\widowpenalty=10000 \clubpenalty=10000 \displaywidowpenalty=10000
% em-dash: --- does NOT become — under Frank Ruhl (the font lacks tlig; verified
%   the width is identical with/without Ligatures=TeX, while the mechanism works
%   on Latin Modern). TYPE THE CHARACTER — (U+2014) DIRECTLY; likewise – (U+2013).
%   check_bidi_figures.py flags a --/--- that is adjacent to Hebrew.
% ============================================================================

% ---- STRUCTURAL DIRECTION-SAFETY for \texttt (v3.7) -------------------------
% A parenthesis, bracket, or $ INSIDE \texttt{} mirrors/migrates in RTL prose
% (MIPS 0($s0)->"(s0$)0", $rs->"rs$", arr[i]->"[arr[i") because \texttt sets the
% font but NOT the direction. Redefining \texttt to LTR-isolate its content fixes
% the WHOLE class at once (parens, brackets, $) so plain \texttt{lw $t0,0($s0)}
% is safe — no need to remember \code{}. Guarded by \ifmmode because
% \foreignlanguage is invalid in math (there, plain monospace is already LTR).
% NB (corrected v3.56 — the old note here said "stays tofu", which is WRONG):
% \texttt{Hebrew} is a FATAL build failure. The tt family is DejaVu Sans Mono
% (zero Hebrew) and the tt fallback hebfb is DejaVuSerif (also zero Hebrew), so
% the subsetter gets an empty glyph set and LuaTeX dies with
%   "! error: (file .../DejaVuSansMono.ttf) (type 2): there are no glyphs in the
%    subset"  -> no output PDF.
% The message names a FONT FILE, never Hebrew/\texttt/the source line. And it is
% non-local: add an unrelated Latin \code{} elsewhere and the fatal degrades into
% a SILENT DELETION of the Hebrew word. Put Hebrew in \h{}/normal text, never in
% \texttt{}/\code{}; flagged HEBREW-IN-MONO. To disable, delete this block.
\let\origtexttt\texttt
\protected\long\def\texttt#1{\ifmmode\origtexttt{#1}\else\foreignlanguage{english}{\origtexttt{#1}}\fi}
% v4.8.0: defer the direction test until TeX has inserted an alignment cell's
% math template. Unprotected expansion here tested \ifmmode too early at the
% FIRST token of aligned/array/cases/matrix cells. Engine-level \protected also
% preserves the original \texttt implementation without the self-recursion
% observed when retrofitting \DeclareRobustCommand to its saved alias.
\protected\long\def\code#1{\ifmmode\origtexttt{#1}\else\foreignlanguage{english}{\origtexttt{#1}}\fi}

% ---- \codeline — display-code line (v3.24): the code analogue of display math.
% Reader-friction rule (content-style.md §5b): inline \code{} is for a SINGLE
% ATOM only (one register/variable/mnemonic, a tight no-space expression).
% A multi-token expression — an instruction with operands, a spaced assignment,
% a register-binding list — zigzags the eye RTL<->LTR mid-sentence and, being
% an unbreakable island, line-breaks badly. Give it its own centered line with
% display-math vertical rhythm. \centerline = one box => the expression never
% wraps; \nopagebreak keeps it glued to its introducing sentence.
%   ...the Hebrew sentence ends with a colon:
%   \codeline{lw \$t1, -4(\$t0)}
%   and the prose continues (escape dollars as \$).
% v3.66 (first line of the body): flush a PENDING \item label before the \par.
% LaTeX defers the label box (\@labels) until the next \everypar, and an \hbox in
% vertical mode never fires \everypar — so `\item \codeline{...}` printed the code
% at y=69.1 and the item's own "1." BELOW it at y=103.6 (the number appears after
% its content, attached to nothing). \if@inlabel\leavevmode\fi is CONDITIONAL on
% purpose: an unconditional \leavevmode would open an empty paragraph line when
% \codeline follows ordinary prose. Verified both ways — in a list the label now
% sits on the code's line at its normal x, and after prose the geometry is
% byte-identical to before.
\makeatletter
\newcommand{\codeline}[1]{%
  \if@inlabel\leavevmode\fi
  \par\nopagebreak
  \vspace{\abovedisplayskip}\nopagebreak
  % v3.62: was \centerline{...}, i.e. \hbox to\hsize{\hss #1 \hss}. \hss has
  % INFINITE SHRINK, so TeX could always "fit" the box and NEVER reported an
  % Overfull — the same \hss mechanism as BLD-27's \item label. Measured: a wide
  % \codeline in a defbox painted its first token ON the box's left border bar and
  % sat 9.5pt outside the frame interior on both sides, with 0 Overfull, 0 Underfull
  % and all four gates PASS; in bare prose it reached 21.0pt into BOTH margins.
  % \hfil has no shrink, so the overflow is reported and build.sh's ovf_vis/ovf_big
  % counters see it; \linewidth (not \hsize) makes it box-aware, so a \codeline
  % inside a colored box centres on the box interior. Verified: a FITTING \codeline
  % keeps byte-identical glyph positions, and the box-crossing case now emits
  % "Overfull \hbox (19.03pt too wide)". A wide \codeline should be split in two.
  \hbox to\linewidth{\hfil\code{#1}\hfil}%
  \nopagebreak\vspace{\belowdisplayskip}}
\makeatother

% ---- a thin decorative rule helper ----
\newcommand{\sep}{\par\vspace{2pt}\noindent\textcolor{cRule}{\rule{\linewidth}{0.6pt}}\par\vspace{4pt}}

% ---- colors ----
\definecolor{cBlue}{HTML}{1F4E79}   \definecolor{cBlueBg}{HTML}{EAF1F8}
\definecolor{cGreen}{HTML}{2E6B4F}  \definecolor{cGreenBg}{HTML}{EAF4EE}
\definecolor{cAmber}{HTML}{9C6A00}  \definecolor{cAmberBg}{HTML}{FBF3E2}
\definecolor{cPurple}{HTML}{5B2C6F} \definecolor{cPurpleBg}{HTML}{F3ECF7}
\definecolor{cTeal}{HTML}{0E6E6E}   \definecolor{cTealBg}{HTML}{E6F4F4}
\definecolor{cRed}{HTML}{8E2B2B}    \definecolor{cRedBg}{HTML}{FBEAEA}
\definecolor{cGray}{HTML}{555555}   \definecolor{cRule}{HTML}{1F4E79}

% ---- math operators ----
\DeclareMathOperator{\Tr}{Tr}
\DeclareMathOperator{\rank}{rank}
\DeclareMathOperator{\diag}{diag}
\DeclareMathOperator{\sgn}{sgn}

% ---- bra-ket (quantikz defines \ket/\bra/\braket; add the rest safely) ----
\providecommand{\ket}[1]{\left|#1\right\rangle}
\providecommand{\bra}[1]{\left\langle#1\right|}
\providecommand{\braket}[2]{\left\langle#1\middle|#2\right\rangle}
\newcommand{\dyad}[1]{\left|#1\right\rangle\!\left\langle#1\right|}
% v3.21: general outer product |a><b| and expectation value <A> — agents reach
% for these by reflex (esp. \expval); quantikz defines neither, so their absence
% threw "Undefined control sequence" in a real build. \providecommand = safe.
\providecommand{\ketbra}[2]{\left|#1\right\rangle\!\left\langle#2\right|}
\providecommand{\expval}[1]{\left\langle#1\right\rangle}
\newcommand{\abs}[1]{\left|#1\right|}
\newcommand{\norm}[1]{\left\lVert#1\right\rVert}
\newcommand{\half}{\tfrac12}
\newcommand{\R}{\mathbb{R}}\newcommand{\C}{\mathbb{C}}\newcommand{\Z}{\mathbb{Z}}
% identity operator. NB: \mathbb{1} silently renders a WRONG glyph (msbm has no
% blackboard digits). Use \mathbb{I} = 𝕀 (works with amssymb, no extra package).
% Do NOT "fix" this back to \mathbb{1}.
\newcommand{\id}{\mathbb{I}}

% ---- physical units (v3.35): \punit — upright unit with thin space ---------
% Usage: $300\punit{K}$, $80\punit{mV/decade}$, $10^{7}\punit{cm/s}$ — the
% NUMBER AND UNIT LIVE IN ONE MATH ISLAND. Promoted verbatim from a 199-page
% production lakit (advanced-electronic-devices, dozens of uses).
% Two measured traps this design avoids (both render-proven):
%   1. siunitx SCRAMBLES under RTL — \SI{3.5}{eV} renders "eV 3.5", exponents
%      detach ("Hz 10 x 10^9"): its quantity assembly is not bidi-isolated.
%      Do NOT load siunitx in this pipeline.
%   2. A text-mode unit next to a bare-prose number is TWO adjacent LTR runs
%      and reorders ("5 mm" -> "mm5" — the BIDI-2/3 class). \punit is therefore
%      math-only ON PURPOSE: prose misuse fails LOUDLY at compile (\mathrm
%      outside math) instead of scrambling silently.
% Domain shorthands (\MeV, \fm...) stay per-document: \newcommand{\MeV}{\,\text{\en{MeV}}}
\providecommand{\punit}[1]{\,\mathrm{#1}}

% ============================================================================
% Colored boxes — mdframed (NOT tcolorbox, which breaks under bidi).
% ★★ v3.4: innerleftmargin/innerrightmargin give code listings room so their
% frame/line-numbers do not spill out the box border (see listing styles below).
% ============================================================================
% v3.61: everyline/split skips were applied to `tablebox` only (BLD-14, v3.37) —
% the six colored boxes below never got them, so they still split OPEN across a page
% break: the page-N segment gets top+sides and NO bottom rule, the page-N+1 segment
% sides+bottom and NO top rule. Measured on the shipped fleet this is a COVERAGE GAP,
% not a re-report of BLD-14: 0 split-open `tablebox` segments (that fix holds) vs 62
% split-open COLORED boxes across 7 documents. Verified after the fix: both segments
% close, and the frametitle is NOT duplicated on the continuation. BLD-14's caveat
% still applies — a box taller than a full page grinds mdframed's split-trial loop, so
% split such content manually.
% ============================================================================
% v3.72: frametitle={} leads EVERY \newmdenv option list below. mdframed options
% are GROUP-scoped: a box nested inside \begin{exbox}[frametitle={...}] used to
% INHERIT the outer title — measured, the identical title string rendered twice,
% the second riding a title-bar fill inside the INNER box's frame, for all seven
% environments in both directions. Both natural forms were broken (= preamble
% bug, not an authoring rule). The explicit reset makes every box titleless
% unless ITS OWN [frametitle={...}] says otherwise; verified both directions:
% no duplicate on the nested box, an explicit title still renders.
% Geometry is shared with the optional semantic-width table container. Themes
% select their own existing semantic style below; neither path fixes columns.
\mdfdefinestyle{hllpdfsemanticgeometry}{linewidth=1.1pt,
  innertopmargin=6pt,innerbottommargin=7pt,skipabove=9pt,skipbelow=9pt,
  innerleftmargin=9pt,innerrightmargin=9pt,everyline=true,
  splitbottomskip=6pt,splittopskip=12pt,roundcorner=3pt}
\newcommand{\hllpdfsemanticframestyle}{hllpdfsemanticgeometry}
\newmdenv[frametitle={},style=hllpdfsemanticgeometry,linecolor=cBlue,backgroundcolor=cBlueBg,
  frametitlebackgroundcolor=cBlue,frametitlefontcolor=white,
  ]{defbox}   % הגדרה
\newmdenv[frametitle={},style=hllpdfsemanticgeometry,linecolor=cGreen,backgroundcolor=cGreenBg,
  frametitlebackgroundcolor=cGreen,frametitlefontcolor=white,
  ]{thmbox}   % משפט/טענה
\newmdenv[frametitle={},style=hllpdfsemanticgeometry,linecolor=cAmber,backgroundcolor=cAmberBg,
  frametitlebackgroundcolor=cAmber,frametitlefontcolor=white,
  ]{notebox}  % הערה
\newmdenv[frametitle={},style=hllpdfsemanticgeometry,linecolor=cPurple,backgroundcolor=cPurpleBg,
  frametitlebackgroundcolor=cPurple,frametitlefontcolor=white,
  ]{exbox}    % דוגמה
\newmdenv[frametitle={},style=hllpdfsemanticgeometry,linecolor=cRed,backgroundcolor=cRedBg,
  frametitlebackgroundcolor=cRed,frametitlefontcolor=white,
  ]{warnbox}  % טעות נפוצה
\newmdenv[frametitle={},style=hllpdfsemanticgeometry,linecolor=cTeal,backgroundcolor=cTealBg,
  frametitlebackgroundcolor=cTeal,frametitlefontcolor=white,
  ]{keybox}   % רעיון מרכזי

% ---- codebox (v4.5.0): the dedicated container for code -------------------
% An OPTION, not a policy: asmblock unboxed stays canonical for a full listing,
% and a snippet inside any of the six colored boxes still works as before.
% What makes codebox worth its own environment is that it is the ONE framed box
% that keeps its margin line numbers.
%
% v4.4.0 shipped it WITHOUT them, on the strength of v3.37's "no margin knob can
% save them" plus four fresh probes that agreed. That was wrong: every probe had
% moved the code, the box, or both — never the key that carries the NUMBER
% COLUMN. Measured on the page:
%   resetmargins=false   the listing follows THIS box instead of the page,
%                        moving the code from x=75.26 to x=97.54
%   xrightmargin=26pt    under RTL this is the PHYSICAL LEFT indent, and the
%                        only key that carries the \llap'ed digits with it:
%                        x=64.28 -> 90.18, clearing the rule at 65.64 by 23.6pt
%   xleftmargin=26pt     moves the code and LEAVES THE DIGITS on the rule —
%                        the asymmetry that hid the fix for eleven versions
% resetmargins is hoisted out of the four styles (see the \lstset below them) so
% this override is possible at all; every other listing still resolves to true.
% Pinned by tests/codebox_numbers.tex + probe_lstnum.py --inside-rule, whose
% negative control strips xrightmargin and watches the digits fall back.
\definecolor{cSlate}{HTML}{3F4A56}  \definecolor{cSlateBg}{HTML}{F2F4F7}
\newmdenv[frametitle={},leftline=true,rightline=true,topline=true,bottomline=true,
  linecolor=cSlate,backgroundcolor=cSlateBg,linewidth=0.9pt,
  frametitlebackgroundcolor=cSlate,frametitlefontcolor=white,
  innertopmargin=6pt,innerbottommargin=7pt,skipabove=9pt,skipbelow=9pt,
  innerleftmargin=9pt,innerrightmargin=9pt,everyline=true,splitbottomskip=6pt,splittopskip=12pt,
  roundcorner=3pt]{codebox}   % קוד
% The two keys that make in-box numbering work, both measured on the page:
%   resetmargins=false  the listing follows THIS box instead of the page
%   xrightmargin        under RTL this is the PHYSICAL LEFT indent, and it is
%                       the only one that carries the \llap'ed number column
%                       with it. xleftmargin moves the code and leaves the
%                       digits behind on the rule (measured: code 97.54,
%                       digits 64.28 — the v3.37 defect, reproduced); with
%                       xrightmargin=26pt the digits land at x=90.18, a
%                       clearance of 23.6pt from the rule at 65.64.
% This asymmetry is why v3.4's note had the semantics backwards and v3.37
% concluded the digits could not be saved.
% <hllpdf:listing-container-state>
% v4.29: all state is local to the current environment group. A semantic
% callout masks number text, while a nested codebox restores the formatter
% saved before that mask. Save it once across successive semantic wrappers;
% an explicit numbers=none (or an explicitly empty formatter) remains valid.
\makeatletter
\newif\ifhllpdf@incodebox
\newif\ifhllpdf@englishcodebox
\newif\ifhllpdf@lstnumbermasked
\let\hllpdf@savedlstnumber\thelstnumber
\makeatother
% </hllpdf:listing-container-state>
% <hllpdf:codebox-gutter>
\makeatletter
\AtBeginEnvironment{codebox}{%
  \hllpdf@incodeboxtrue
  % Snapshot the OUTER language before asmblock selects English.
  \iflanguage{english}{\hllpdf@englishcodeboxtrue}{\hllpdf@englishcodeboxfalse}%
  \ifhllpdf@lstnumbermasked
    \let\thelstnumber\hllpdf@savedlstnumber
    \hllpdf@lstnumbermaskedfalse
  \fi
}
% listings' AddToHook is GLOBAL. Register once, and gate on local box state;
% adding an unconditional hook at every codebox entry leaks into later listings.
\lst@AddToHook{PreInit}{%
  \ifhllpdf@incodebox
    % Apply every listing override here, inside the LISTING's own group.
    % Setting margins at box entry leaks them into nested semantic callouts.
    \lstset{resetmargins=false,xrightmargin=26pt}%
    % Styles resolve before PreInit. In LTR the physical number gutter follows
    % xleftmargin; in Hebrew retain the measured inherited xrightmargin path.
    \ifhllpdf@englishcodebox
      \lstset{xleftmargin=26pt,xrightmargin=0pt}%
    \fi
    % ONE box, not two: suppress only the frame/background/skips of a listing
    % inside this codebox. The next unboxed listing keeps its own style.
    \lstset{frame=none,backgroundcolor=\color{cSlateBg},aboveskip=0pt,belowskip=0pt}%
  \fi
}
\makeatother
% </hllpdf:codebox-gutter>

% v3.37: line numbers inside a COLORED BOX are sliced by the box border —
% measured: digit span x=63.5-67.3 vs border line at x=65.2 (the border runs
% THROUGH the digits), and the inter-border gap (3.5pt) is narrower than a
% digit (3.8pt), so no margin knob can save them (under resetmargins,
% xleftmargin is inert in-box — measured: code at x=75.26 with 2.2em AND
% 3.0em). A 3-6 line boxed snippet does not need line numbers; blank the
% number TEXT inside every colored box (\thelstnumber is what gets printed;
% the style's numbers=left survives, so nothing else shifts — code position
% byte-frozen at x=75.26, verified). Unboxed listings keep their numbers.
% v3.78: the blanking is coupled to the FRAME, not to the box name —
% clean_theme.tex registers LATER hooks (registration order wins) that
% restore the number text inside the boxes it renders frameless.
% v4.5.0 RETRACTS the "no margin knob can save them" clause above. It holds only
% under resetmargins=true (the shipped default): there xleftmargin is inert
% in-box, as measured. With resetmargins=false, xrightmargin DOES carry the
% \llap'ed digit column clear of the rule — measured; the codebox \lstset that
% does it is the one above, in the codebox block. The loop below therefore
% covers SIX boxes, not "every colored box": codebox opts out and keeps its
% numbers.
% <hllpdf:semantic-listing-mask>
\makeatletter
\@for\HEB@lstbox:=defbox,thmbox,notebox,exbox,warnbox,keybox\do{%
  \expandafter\AtBeginEnvironment\expandafter{\HEB@lstbox}{%
    \ifhllpdf@lstnumbermasked\else
      \let\hllpdf@savedlstnumber\thelstnumber
    \fi
    \hllpdf@lstnumbermaskedtrue
    \def\thelstnumber{}%
    % The nearest semantic container owns its listing presentation. On exit,
    % TeX's group restores the enclosing codebox state and formatter.
    \hllpdf@incodeboxfalse}}
\makeatother
% </hllpdf:semantic-listing-mask>

% ============================================================================
% Unified TABLE presentation: subtle uniform slate frame + auto-numbered
% caption rendered INSIDE the frame. Use:
%   \begin{tablebox}\centering  <tabular>  \tabcap{<description>}\end{tablebox}
% Replaces the buggy \caption* (which emitted a stray "טבלה N : *" line) and
% provides an explicit framed option + numbered caption "טבלה N.M: תיאור".
% Optional intro text may be placed at the TOP, before the tabular, inside the box.
% ============================================================================
\definecolor{cTblLine}{HTML}{B9C4D2}  \definecolor{cTblBg}{HTML}{FAFBFC}
% v3.23: class-aware — under `article` there is no chapter counter (the promised
% report->article switch used to die with "No counter 'chapter' defined").
\makeatletter
\@ifundefined{c@chapter}{%
  % v3.31: in article the counter resets per SECTION, so an unprefixed number
  % repeated ("טבלה 1." in every section — measured on a two-section render).
  % Same conditional idea as the report branch: prefix the section number once
  % a real \section exists; a pre-section table stays "טבלה 1".
  \newcounter{tbl}[section]
  \renewcommand{\thetbl}{\ifnum\value{section}>0 \thesection.\fi\arabic{tbl}}%
}{%
  \newcounter{tbl}[chapter]
  % v3.30: a report-class doc that uses SECTIONS ONLY (chapter stays 0 —
  % the common short-doc pattern) used to print "טבלה 0.1"; the chapter
  % prefix now appears only once a real \chapter exists (render-verified
  % both ways: "טבלה 1." chapterless, "טבלה 1.1" with chapters).
  \renewcommand{\thetbl}{\ifnum\value{chapter}>0 \thechapter.\fi\arabic{tbl}}%
  % v3.64: the SAME leak in \thesection itself, left behind by v3.30. That fix
  % guarded the CAPTION counter and stopped there, so a chapterless report still
  % printed the heading "0.1 סעיף ראשון" and \ref to it printed "0.2" — measured:
  %   p1 y=65.2  '…0.1'   p1 y=144.9 '…0.2'   p1 y=183.4 '…0.2.1'
  % report.cls hardcodes \thesection as \thechapter.\@arabic\c@section, and the
  % leading 0 is the chapter counter that no \chapter ever incremented. Fixing
  % \thesection fixes the whole chain — \thesubsection and \thesubsubsection are
  % defined in terms of it. \theequation and \thetbl/\thefig were already guarded,
  % which is exactly why this survived: every NUMBER around the heading looked
  % right, so the heading's own 0 read as a style choice rather than a bug.
  % \appendix is safe: report.cls resets chapter to 0 and redefines \thechapter to
  % \@Alph\c@chapter, so an appendix \chapter makes the counter 1 and the guard
  % emits "A.1" unchanged. Render-verified in all four modes (chapterless,
  % chaptered, appendix+chapter, appendix+section-only).
  \renewcommand{\thesection}{\ifnum\value{chapter}>0 \thechapter.\fi\arabic{section}}%
}
\makeatother
% v3.37: everyline=true — if the box DOES split across pages, every segment
% gets a CLOSED frame (measured before: the split ran open-ended into the page
% edge with no bottom border, and the caption could strand alone at the next
% page top — "טבלה A.2" shipped that way). The split skips give the segments
% breathing room. In practice, with the caption GLUE below, every previously-
% split table in the corpus now moves WHOLE instead (the white-gap tradeoff,
% documented benign); everyline stays as the safety net for a genuinely
% page-tall table. NB: a table taller than ONE FULL PAGE can grind mdframed's
% split-trial loop (reproduced: minimal doc, minutes of 1pt-step Overfull
% probes) — break such a monster into two tables.
\mdfdefinestyle{hllpdftablegeometry}{leftline=true,rightline=true,topline=true,bottomline=true,
  linewidth=0.9pt,
  innertopmargin=9pt,innerbottommargin=8pt,skipabove=11pt,skipbelow=11pt,
  innerleftmargin=11pt,innerrightmargin=11pt,roundcorner=4pt}
\newcommand{\hllpdftableframestyle}{hllpdftablegeometry}
\newmdenv[frametitle={},style=hllpdftablegeometry,
  linecolor=cTblLine,backgroundcolor=cTblBg,
  % v4.2.0 (adopted): tablebox is ATOMIC. Short tableboxes near page bottoms
  % repeatedly entered mdframed's split trial, emitting benign Overfull-vbox
  % bursts and occasionally separating the intro prose from the tabular.
  % nobreak=true moves the complete short table to the next page instead. The
  % output-active Overfull-vbox gate still catches a box taller than a page;
  % everyline/split skips remain as defensive styling should that path ever
  % re-arm, but the canonical path never publishes a split tablebox.
  nobreak=true,
  everyline=true,splitbottomskip=6pt,splittopskip=12pt,
  % v3.79 (BLD-37 table side, correction): startinnercode/endinnercode run
  % INSIDE mdframed's body vbox (mdframed.sty:848 / :856+863) -- the second one
  % after \par\unskip, i.e. after line breaking, which is the only moment the
  % \tabcap anchor's true offset from the top of the box EXISTS. See
  % \lakintabboxclose below. These two keys are the ONLY tablebox change.
  startinnercode={\lakintabboxopen},endinnercode={\lakintabboxclose}]{tablebox}
% v3.37: \nopagebreak GLUE — \vsplit honors penalties, so the caption can no
% longer strand alone at the top of the next page (it takes its rows along, or
% the whole box moves). Render-verified on the real shipped orphan.
% v3.59: \NR@gettitle is what nameref injects into the stdlib title-bearing
% commands (\@sect, \@chapter, \caption) to fill \@currentlabelname. \figcap and
% \tabcap are INDEPENDENT macros — they never called it, so \@currentlabelname kept
% whatever the last heading left and \nameref{fig:x} silently printed the enclosing
% SECTION's title instead of the caption (proven in the .aux: field 3 of \newlabel
% was the heading text while field 1 and the anchor `fig.1.1` were correct — so this
% is a THIRD, distinct head next to BLD-13's counter fix and BIDI-38's autoref-name
% fix, neither of which writes that field). \@ifundefined keeps it inert if
% hyperref/nameref is ever dropped. Guarded helper because \figcap is defined
% after \makeatother.
% NB both helpers below need \makeatletter — this part of the preamble is OUTSIDE
% any letter-catcode region, so an unwrapped \@currentlabel parses as LaTeX's \@
% and dies with "You can't use \spacefactor in vertical mode" (measured).
\makeatletter
\newcommand{\lakinrtitle}[1]{\@ifundefined{NR@gettitle}{}{\NR@gettitle{#1}}}
% v3.59: \hitem{<letter>}{<key>} — a REFERENCEABLE Hebrew-lettered list item.
%   \begin{enumerate} \hitem{א}{it:a} הסעיף הראשון. \end{enumerate}  ->  \ref{it:a} = (א)
% Why this exists: automatic Hebrew list labels fail silently, so SKILL.md mandates
% an EXPLICIT \item[(\textbf{א})] — but LaTeX puts \refstepcounter INSIDE
% \if@noitemarg (latex.ltx:12964), and \item[...] never sets that flag. So an
% explicit item steps no counter and never touches \@currentlabel: a \label there
% binds to the AMBIENT one and \ref prints the CHAPTER/SECTION number where the
% item letter belongs (proven in the .aux: {chapter.1}). Same doctrine as \figcap —
% own the label state before \label. \phantomsection matters: without it the
% hyperlink still lands on the chapter head (measured); with it the anchor is
% section*.N and the bookmark outline stays clean.
\newcommand{\hitem}[2]{\item[(\textbf{#1})]%
  \def\@currentlabel{(#1)}\def\@currentlabelname{#1}\phantomsection\label{#2}}
% NB (v3.66): \hitem is \ref-only. \autoref{key} calls it a SECTION ("סעיף (א)")
% because \phantomsection brands the anchor section*.N, and \nameref{key} prints
% the bare letter. For \autoref-able items prefer \hebenumlabels below.
%
% v3.66: \hebenumlabels — OPT-IN automatic Hebrew enumerate letters, done right.
%   {\hebenumlabels\begin{enumerate}\item\label{q:a} ...\end{enumerate}}
%   -> labels (א) (ב) (ג); \ref{q:a} = א; anchor Item.N (real, per-item).
% This retires the belief that made \item[(\textbf{א})] + \hitem necessary.
% "Automatic Hebrew labels fail silently" was TRUE but unexplained, and the
% mechanism is now bisected (plain article vs +enumitem, measured): ENUMITEM
% PRE-EXPANDS the label expression in an \edef BEFORE \refstepcounter, with the
% counter still at 0. Tokens it shields by name (\arabic,\alph,\roman) survive;
% a plain \newcommand \ifcase wrapper evaluates at 0 -> the EMPTY zero-branch ->
% the label prints as "()" with the letter GONE (and \the\c@enumi prints a
% literal 0, \number\value dies fatally pointing at \item). The cure is one
% property, not a doctrine: the wrapper must be ROBUST (survives the \edef
% unexpanded), while \theenumi stays EXPANDABLE (so \refstepcounter writes the
% real letter to the .aux). Verified: labels (א)(ב)(ג), \ref -> א/ב/ג each item
% its own value, .aux = {א}{...}{Item.1} — versus the explicit-\item doctrine
% where BIDI-50/\hitem exist purely to patch the counter that never stepped.
\def\@hebalph#1{\ifcase#1\or א\or ב\or ג\or ד\or ה\or ו\or ז\or ח\or ט\or י\or
  יא\or יב\or יג\or יד\or טו\or טז\or יז\or יח\or יט\or כ\else\@ctrerr\fi}
\DeclareRobustCommand*\hebalph[1]{\expandafter\@hebalph\csname c@#1\endcsname}
\newcommand{\hebenumlabels}{%
  \renewcommand{\theenumi}{\@hebalph\c@enumi}%
  \renewcommand{\labelenumi}{(\textbf{\hebalph{enumi}})}}
\makeatother
% v3.63: \mbox the LABEL. \thetbl ("7.1") and a caption description that STARTS
% WITH A DIGIT ("32 הרגיסטרים…") are two European-Number runs separated only by
% neutrals (the label's own "." and \enspace). Under the RTL paragraph base the UBA
% resolves those neutrals to the RTL level, so both numbers sit in ONE reversed span
% and EXCHANGE VISUAL POSITIONS. Measured on the shipped PDF: a table whose real
% number is 7.1 printed "טבלה 32" — a different, EXISTING table number — while \ref
% correctly printed 7.1, so the caption silently contradicted its own cross-reference,
% and the description lost its leading count. The source contains no adjacency at all;
% it is manufactured by this macro, which is why every source linter was blind.
% \mbox makes the label one atomic box so nothing can reorder across it. Verified
% zero-regression: an ordinary Hebrew caption keeps identical x-geometry.
% One paragraph style covers label, number, separator and description.
% Explicit project/template overrides remain possible; the label hook is kept
% for compatibility, with an identity default. Counters/anchors stay unchanged.
\newcommand{\lakitcaptionstyle}{\normalfont\small\color{black}\hypersetup{allcolors=black}}
\newcommand{\lakitcaptionlabelstyle}[1]{#1}
\newcommand{\lakitcaptionpunctuation}{:}
% Native floats, captionof and longtable use the same presentation policy.
% subcaption is optional and inherits the same font rather than a second size.
\usepackage{caption}
\DeclareCaptionFont{lakit}{\lakitcaptionstyle}
\DeclareCaptionLabelFormat{lakit}{\lakitcaptionlabelstyle{#1~\en{#2}}}
\DeclareCaptionLabelSeparator{lakit}{\lakitcaptionpunctuation\enspace}
\captionsetup{font=lakit,labelfont={},textfont={},labelformat=lakit,
  labelsep=lakit,justification=centering}
\makeatletter
\AtBeginDocument{\@ifpackageloaded{subcaption}{%
  \captionsetup[sub]{font=lakit,labelfont={},textfont={},labelsep=lakit}}{}}
\makeatother
% v3.64: OPT-IN list-of-figures/tables entry via an optional short caption.
% \listoffigures and \listoftables are inherited from report.cls and LOOK available,
% but \figcap/\tabcap step the custom \c@fig/\c@tbl counters and never call
% \addcontentsline — so both printed their heading and NOTHING ELSE. Measured: a doc
% with one \figcap and one \tabcap rendered page 1 as 'רשימת האיורים\n1\n', and the
% .lof held only \@chapter's \addvspace, zero \contentsline entries. An empty list
% under a correct Hebrew heading reads as "this document has no figures".
%
% WHY OPTIONAL, AND NOT AUTOMATIC. Writing #1 unconditionally was implemented first
% and REJECTED: a caption containing \codeline compiles CLEANLY today (verified on
% the pre-v3.64 preamble: 0 errors, PDF produced) but becomes FATAL once it is also
% written to the .lof — \codeline is \hbox to\linewidth{...}, and a full-width box in
% a list entry dies with "\pdfextension endlink cannot be used in vertical mode", no
% output PDF. Gaining a feature nobody uses (0 of 13 collections call \listoffigures)
% at the price of breaking a construct that works is a bad trade. Stock \caption has
% exactly this limitation and exactly this escape hatch, so the convention is familiar.
% \en{}, \code{}, \h{} and inline math in the short text are all safe — verified, they
% survive the write and typeset correctly in the list (RTL geometry confirmed: number
% rightmost at x=509.9, title left of it, dots and folio at the far left).
% With no optional argument the macro is byte-identical to v3.63 — zero risk to every
% existing document, and \listoffigures lists exactly the captions that opted in.
% BLD-37, TABLE side (the \figcap v3.74 idea, adapted): clicking a \ref to a
% table used to land MID-table — hyperref's automatic anchor inside
% \refstepcounter was emitted on the still-open horizontal line whose TAIL is
% the tabular (an hbox whose baseline is the table's vertical CENTER), so the
% jump scrolled the top half of the table above the viewport. The v3.74 entry
% deferred this side claiming "tablebox splits — no single box to measure";
% that deferral dissolves: the measurement target is the TABULAR itself at
% \tabcap ENTRY, while it is still the tail of the open list — the same
% situation \lakin@grabfig handles at an island's end (and the fleet has 0
% split tableboxes). So \lakin@grabtab runs as the FIRST token, before any
% typesetting can push the box away, and the anchor is re-emitted raised above
% the measured tabular. \lakintabraise is a SEPARATE register from
% \lakinfigraise (defined with it, below the LTR-island block): a captionless
% island intentionally leaves the fig register set, and a table must never
% inherit that stale figure height.
%
% WHY the anchor stays AT THE ENTRY SPOT and is NOT deferred into the caption
% paragraph the way \figcap defers it: at \par TeX DROPS a final glue item,
% and the automatic anchor's smash box sitting BEHIND the interword glue kept
% that glue alive — so every centered tabular has always sat half a space off
% true center. Deferring the anchor removed that box and was measured as 982
% differing glyph records: EVERY tabular line shifted 1.30pt (2.61pt for
% pattern 7's two minipages) while captions stayed put. Emitting the raised
% anchor at the old spot — through hyperref's own \Hy@raisedlink, exactly the
% macro the automatic path uses — keeps the line's node structure, and so the
% A/B glyph diff at 0, BY CONSTRUCTION. \Hy@raisedlink reads the raise from
% \HyperRaiseLinkDefault (default \baselineskip = the old mid-table dest); the
% group-local \edef freezes measured-height+42pt into it, and \smash inside
% \Hy@raisedlink hides any raise from the line's extent. \HyCnt@ProvideTheH-
% Counter keeps the dest NAME byte-identical to the automatic one (it defines
% \theHtbl, which \hyper@makecurrent prefers — without it the first measured
% caption was named tbl.1 where the automatic path said tbl.0.1, measured).
% vmode entry, or anything but a box at the tail (no tabular to measure),
% falls back to the plain \refstepcounter — byte-identical to the old path.
\makeatletter
\newcommand{\tabcap}[2][]{\lakin@grabtab
  % v3.79: arm the box-close correction. Set unconditionally and in BOTH
  % branches below's scope, so the plain-\refstepcounter fallback is corrected
  % too; \lakintabboxopen clears it at every \begin{tablebox}, so a \tabcap
  % outside a tablebox can never leave it armed for a later box.
  \global\lakintabcapseentrue
  \ifdim\lakintabraise>0pt
    \H@refstepcounter{tbl}\HyCnt@ProvideTheHCounter{tbl}\hyper@makecurrent{tbl}%
    \begingroup
      % A frameless table targets its measured top. The frame's historical
      % cushion belongs only to tablebox and is corrected at box close.
      \edef\HyperRaiseLinkDefault{\the\lakintabraise}%
      \iflakintabframed
        \edef\HyperRaiseLinkDefault{\the\dimexpr\lakintabraise+42pt\relax}%
      \fi
      \Hy@raisedlink{\hyper@anchorstart{\@currentHref}\hyper@anchorend}%
    \endgroup
    \global\lakintabraise 0pt\relax
  \else
    \refstepcounter{tbl}%
  \fi
  \lakinrtitle{#2}\par\nopagebreak\addvspace{6pt}%
  \def\lakin@short{#1}\ifx\lakin@short\@empty\else
    \addcontentsline{lot}{table}{\protect\numberline{\thetbl}{\protect\ignorespaces #1}}\fi
  {\centering\lakitcaptionstyle\mbox{\lakitcaptionlabelstyle{\tablename~\en{\thetbl}\lakitcaptionpunctuation}}\enspace #2\par}}
\makeatother

% ============================================================================
% Unified FIGURE caption — the \tabcap twin for TikZ/plot figures (v3.34.0).
% A figure (tikzpic) sits inline in the prose; put \figcap BELOW it for a
% numbered, \ref-able caption "איור N.M: תיאור" in the same style as \tabcap:
%   \begin{tikzpic} ... \end{tikzpic}
%   \figcap{\label{fig:foo}description}     % \label FIRST, inside the argument
% WHY a macro, not a hand-written gray line: a manual
%   {\centering\small\textcolor{cGray}{איור: ...}} with a \label inside has NO
% \refstepcounter, so the \label binds to the AMBIENT \@currentlabel (the last
% equation) — \ref then prints an equation number AS the figure number
% (measured on a real page: "איור 48" where 48 is equation 48). \figcap steps a
% real \c@fig counter first, so \label is correct, numbering is consistent, and
% refs are stable. Class-aware like \thetbl (chapter- or section-prefixed).
% ============================================================================
\makeatletter
\@ifundefined{c@chapter}{%
  \newcounter{fig}[section]
  \renewcommand{\thefig}{\ifnum\value{section}>0 \thesection.\fi\arabic{fig}}%
}{%
  \newcounter{fig}[chapter]
  \renewcommand{\thefig}{\ifnum\value{chapter}>0 \thechapter.\fi\arabic{fig}}%
}
\makeatother
\makeatletter
\newcommand{\figcap}[2][]{%
  \ifdim\lakinfigraise>0pt
    % v3.74: step WITHOUT hyperref's automatic anchor, then emit the anchor
    % raised above the measured figure. \H@refstepcounter is hyperref's saved
    % pre-hook \refstepcounter (sets the counter AND \@currentlabel — \label
    % in the caption arg binds exactly as before); \hyper@makecurrent supplies
    % \@currentHref, so the dest NAME is identical to the old automatic one —
    % \ref/\autoref/bookmarks are unchanged, only the dest'S POSITION moves.
    \H@refstepcounter{fig}\hyper@makecurrent{fig}%
    % The anchor is emitted INSIDE the caption paragraph below, as a zero-size
    % raisebox — NOT here in vertical mode, where a \raisebox starts an empty
    % paragraph and adds a blank line (measured: every glyph after the caption
    % shifted +18.8pt). Inside an existing line a [0pt][0pt] box is inert.
    \edef\lakin@anch{\@currentHref}%
    % \edef FREEZES the measured height into the pending macro: the deferred
    % \raisebox is typeset inside the caption paragraph AFTER the register is
    % reset to the 0pt sentinel below — a plain \def read 0 at use time
    % (measured: LAKINRAISE=0.0, dest moved 14pt instead of above the figure).
    \edef\lakin@pend{\noexpand\raisebox{\the\dimexpr\lakinfigraise+42pt\relax}[0pt][0pt]{%
      \noexpand\hyper@anchorstart{\lakin@anch}\noexpand\hyper@anchorend}}%
    \global\lakinfigraise 0pt\relax
  \else
    \refstepcounter{fig}\let\lakin@pend\relax
  \fi
  \lakinrtitle{#2}\par\nopagebreak\addvspace{4pt}%
  \def\lakin@short{#1}\ifx\lakin@short\@empty\else
    \addcontentsline{lof}{figure}{\protect\numberline{\thefig}{\protect\ignorespaces #1}}\fi
  {\centering\lakin@pend\lakitcaptionstyle\mbox{\lakitcaptionlabelstyle{\figurename~\en{\thefig}\lakitcaptionpunctuation}}\enspace #2\par}}
\makeatother

% ============================================================================
% LTR islands.  EVERY tikzpicture / quantikz MUST live inside one, or the whole
% figure's coordinates and Latin labels flip. (Hebrew node text still uses \h{}.)
% ============================================================================
% v3.74 (BLD-37): the \figcap hyperlink DESTINATION used to land at the caption
% line — BELOW the figure — so clicking a figure reference scrolled the figure
% off-screen (measured fleet-wide: 97/97 fig dests in one collection sat 19-54pt
% under their figure's ink; a 347pt figure was fully hidden after the jump).
% hypcap-style fix in two parts: (1) at the island's END, \lakin@grabfig grabs
% the just-typeset centered line with \lastbox, measures it, and puts it back
% with \nointerlineskip — the interline glue above the line was already
% contributed and is NOT re-added, so the page geometry is unchanged to the
% glyph (verified); (2) \figcap then steps the counter WITHOUT hyperref's
% automatic anchor (\H@refstepcounter = the pre-hyperref \refstepcounter) and
% emits the anchor itself, raised by the measured height — the dest lands ABOVE
% the figure top. A captionless island leaves numbering untouched (nothing is
% stepped at the island; the grab flag is cleared at every island begin, so a
% stale height from a captionless figure cannot leak into a later caption).
\newsavebox{\lakinfigbox}
\newdimen\lakinfigraise  % real \dimen: tex.setdimen cannot target a \newlength SKIP register (measured lua error)
% BLD-37 table side: a SEPARATE register for \tabcap. The two must not share —
% a captionless island INTENTIONALLY leaves \lakinfigraise set (the grab runs
% at every island END; only the NEXT island's begin clears it), so a shared
% register would hand a \tabcap after a captionless figure that stale FIGURE
% height. \lakintabraise is written only by \lakin@grabtab at \tabcap entry
% and consumed/reset inside \tabcap itself, so it can never go stale.
\newdimen\lakintabraise
\makeatletter
% v3.74b: measure WITHOUT touching the list. The first draft grabbed the box
% with \lastbox and put it back — and \lastbox DISCARDS the node's baseline
% shift, so a vertically-anchored quantikz circuit rose 2.57pt and everything
% after it sank 5.32pt (measured on the A/B glyph diff). LuaTeX lets us READ
% the tail node's dimensions without removing anything: geometry is identical
% by construction. \lakinfigraise=0pt is the "no figure measured" sentinel.
\def\lakin@grabfig{%
  \ifhmode
    \directlua{
      local n = tex.nest[tex.nest.ptr].tail
      local GLUE, PEN = node.id('glue'), node.id('penalty')
      while n and (n.id == GLUE or n.id == PEN) do n = n.prev end
      local HL, VL = node.id('hlist'), node.id('vlist')
      if n and (n.id == HL or n.id == VL) then
        tex.setdimen('global', 'lakinfigraise',
                     (n.height or 0) + (n.depth or 0))
      else
        tex.setdimen('global', 'lakinfigraise', 0)
      end}%
  \else
    \global\lakinfigraise 0pt\relax
  \fi}
% BLD-37 table side: \lakin@grabfig's twin, NOT a shared helper — the Lua body
% differs in exactly one place, the skip set. Between \end{tabular} and \tabcap
% an otherlanguage wrapper (patterns 3/5 of example_tables) leaves a \write
% whatsit and a dir node ON TOP of the interword glue (measured node dump:
% "whatsit/write <- dir <- glue(2.62) <- hlist"), so the table reader must
% also step over whatsit/dir markers — all zero-size — to reach the tabular.
% \lakin@grabfig deliberately does NOT get the wider skip set: its glue/penalty
% body is the one the v3.74 A/B proof froze, and widening it would newly
% measure (and so newly raise) anchors in island shapes that today take the
% plain path — an unproven behavior change on the FIG side. Anything that is
% not a box behind those markers still hits the 0pt sentinel and \tabcap falls
% back to the plain \refstepcounter.
% SECOND deliberate difference: this stores the box's HEIGHT, not height+depth.
% \figcap's anchor is emitted in the caption paragraph BELOW the figure, so it
% must climb the full extent; \tabcap's anchor sits ON the tabular's own line,
% whose baseline is the vcentered tabular's MIDPOINT — height alone reaches
% the top edge, and adding the depth was measured to overshoot by half the
% table (dests 80-135pt above their frame, INSIDE the preceding table's box).
\def\lakin@grabtab{%
  \ifhmode
    \directlua{
      local n = tex.nest[tex.nest.ptr].tail
      local GLUE, PEN = node.id('glue'), node.id('penalty')
      local WHAT, DIR = node.id('whatsit'), node.id('dir')
      while n and (n.id == GLUE or n.id == PEN
                   or n.id == WHAT or n.id == DIR) do n = n.prev end
      local HL, VL = node.id('hlist'), node.id('vlist')
      if n and (n.id == HL or n.id == VL) then
        tex.setdimen('global', 'lakintabraise', n.height or 0)
      else
        tex.setdimen('global', 'lakintabraise', 0)
      end}%
  \else
    \global\lakintabraise 0pt\relax
  \fi}
\makeatother
% ============================================================================
% v3.79 (BLD-37 table side, CORRECTION): put the anchor at the TOP OF THE BOX,
% instead of guessing the distance to it.
% ============================================================================
% v3.77 raised the \tabcap anchor by (height of the box at the tail of the open
% horizontal list) + 42pt. That is a MEASUREMENT taken BEFORE line breaking, and
% it is short by every term it cannot see. Fleet-measured on the 12 built
% collections (299 tbl.* dests, ownership established through each table's own
% CAPTION, never through "which frame encloses the dest"): 3 dests landed INSIDE
% their own tablebox, below the frame top —
%   computer-architecture       tbl.14.1  -46.4pt   three stacked tabulars
%   computer-architecture       tbl.30.5  -13.5pt   tabular 2.85pt over \linewidth
%   advanced-electronic-devices tbl.B.1   -59.2pt   tabular 2.56pt over \linewidth
% Two missing terms, both reproduced in a minimal document (probe P2/P3):
%   (1) VERTICAL LIST ALREADY CONTRIBUTED inside the box — a second and third
%       tabular, or a heading above the table. The tail read sees only the LAST
%       box, so 77.7pt of stacked tabulars is invisible to the raise.
%   (2) WHICH LINE THE ANCHOR LANDS ON. When the tabular alone overfills
%       \linewidth (459.6121pt in a tablebox), the interword glue after
%       \end{tabular} becomes the chosen breakpoint and \Hy@raisedlink is
%       emitted on the NEXT line — one (depth + lineskip) lower. That term does
%       not EXIST at \tabcap time: line breaking has not happened yet. No
%       measurement taken there can ever recover it. (Proof: probe P3, the same
%       table 2.84912pt too wide, fails at -13.5pt; probe P4, the identical
%       table with \tabcolsep 2pt so it fits, is correct at +32.0pt.)
% So the raise is not computed at \tabcap at all. It is CORRECTED at
% \end{tablebox}, from inside mdframed's own body vbox (endinnercode runs after
% \par\unskip, i.e. after line breaking, while tex.nest[tex.nest.ptr] is still
% that vbox's vertical list). There the anchor's true offset from the top of the
% box content is not estimated but WALKED, and the \Hy@raisedlink shift is set
% so the dest lands 42pt above it — which is 42pt - innertopmargin(9pt) -
% linewidth(0.9pt) = 32.0pt above the FRAME top, the exact position the 295
% already-correct dests have today. That is why this is a correction and not a
% re-placement: the good sites do not move.
%
% WHY THE GLYPH DIFF IS 0 BY CONSTRUCTION — the only thing written is the
% `shift` field of the hlist that \Hy@raisedlink already built, and the guard
% below refuses to write it unless that box is 0x0x0. A zero-size box carries no
% dimension into its parent's packing, the packing already happened, and shift
% is not an input to it. Nothing else in the list is touched: no node is added,
% removed, resized or reordered. Measured, not assumed: assets/example_tables.tex
% and the whole computer-architecture collection A/B at 0 differing pymupdf
% rawdict glyph records.
%
% WHY NOT the alternatives (all prototyped and measured before this one):
%   * accumulate the vertical list at \tabcap time — fixes term (1), CANNOT fix
%     term (2); 2 of the 3 real failures survive. A fix that leaves the majority
%     of the defect standing is the false all-clear this repo exists to avoid.
%   * \zref-savepos two-pass — correct on all shapes, but it makes the dest
%     POSITION aux-derived (DEFECTS BLD-1/2), a single-pass compile silently
%     ships an unraised anchor, and the bare y-difference has no page guard.
%   * emit the anchor at \begin{tablebox} and let \tabcap adopt it — needs the
%     tbl counter stepped at box OPEN to keep the dest NAME byte-identical; a
%     captionless tablebox then either shifts every later table number or leaves
%     an orphan dest that the next table duplicates (a HARD build gate). This
%     correction touches neither the counter nor the name: \tabcap's naming path
%     is v3.77's, unchanged, so resolve_names() is byte-identical (verified).
%   * suppressing the line break (a penalty before the glue) — that MOVES the
%     tabular: the break currently discards the interword glue, so forbidding it
%     re-centers the two overfull tables. Reflow, rejected.
%
% SCOPE. The correction is armed only by \tabcap ITSELF (\lakintabcapseen),
% cleared at every box open, so a captionless tablebox is a no-op and a \tabcap
% hosted OUTSIDE a tablebox (the fleet has exactly one, computer-architecture
% tbl.14.2 in an exbox) keeps v3.77's measured raise. Both fallbacks are
% therefore never worse than today; and because hyperref's AUTOMATIC anchor uses
% the same \Hy@raisedlink shape, the plain-\refstepcounter path (vmode entry, or
% prose between \end{tabular} and \tabcap — probe P6, -51.5pt) is corrected too.
% KNOWN LIMIT — a tablebox that SPLITS across pages. The offset is measured from
% the top of the UNSPLIT body, so a dest in segment 2 gets raised past its own
% segment. MEASURED on a forced split (probe: 40 filler lines + a table): the
% dest went from y=585.0 on page 2 (inside segment 2, at the caption — the old
% BLD-37 symptom) to y=-113.4, i.e. ABOVE the page top. Viewers clamp that to the
% page top, which happens to show segment 2 whole, but an out-of-page coordinate
% is not something to ship on purpose. Hence the \textheight guard in
% lakintab.fix: a body TALLER THAN THE TEXT BLOCK cannot fit on any page, so it
% must split — the correction is skipped and v3.77's behaviour stands unchanged
% (verified: the probe returns to y=585.0). The guard is deliberately EXACT, not
% a page-fitting heuristic: it is sound (a taller-than-page body always splits)
% but PARTIAL (a shorter body that starts low on a page could still split). That
% residue is unreachable in practice — v3.37's everyline + the caption's
% \nopagebreak glue move such boxes WHOLE instead (probe N7: a 22-row table
% pushed down by 430pt moved whole, 0 segments), and the fleet has 0 split
% tableboxes, re-counted after this change. Reimplementing mdframed's own
% page-fitting arithmetic to close the residue would be a heuristic with its own
% failure mode, which is the trade this repo has repeatedly refused.
\newif\iflakintabcapseen
\newif\iflakintabframed
% The Lua lives in ONE \directlua at load time, and is written with NO
% backslash, %, # or ~ : \directlua expands its argument like \edef, so an
% active ~ would expand to \nobreakspace and a % would eat the rest of the line
% (the second is how the first draft of this block died). "not (a == b)" stands
% in for Lua's inequality operator for the same reason.
\directlua{
  lakintab = { }
  lakintab.HL = node.id("hlist")
  lakintab.VL = node.id("vlist")
  lakintab.RU = node.id("rule")
  lakintab.GL = node.id("glue")
  lakintab.KE = node.id("kern")
  lakintab.WH = node.id("whatsit")
  lakintab.TM = node.id("temp")
  lakintab.WD = node.subtype("pdf_dest")
  function lakintab.d(v)
    if v == nil then return 0 end
    if v < -1000000000 then return 0 end
    return v
  end
  function lakintab.scan(list, y, horiz, parent, target)
    local x = list
    while x do
      local i = x.id
      if i == lakintab.HL or i == lakintab.VL then
        if horiz then
          local b = y + x.shift
          if i == lakintab.HL then
            lakintab.scan(x.head, b, true, x, target)
          else
            lakintab.scan(x.head, b - lakintab.d(x.height), false, x, target)
          end
        else
          if i == lakintab.HL then
            lakintab.scan(x.head, y + lakintab.d(x.height), true, x, target)
          else
            lakintab.scan(x.head, y, false, x, target)
          end
          y = y + lakintab.d(x.height) + lakintab.d(x.depth)
        end
      elseif i == lakintab.WH then
        if x.subtype == lakintab.WD then
          local id = x.dest_id
          if type(id) == "string" and string.sub(id, 1, 4) == "tbl." then
            if horiz and parent and parent.width == 0
               and parent.height == 0 and parent.depth == 0 then
              parent.shift = parent.shift + (target - y)
            end
          end
        end
      elseif not horiz then
        if i == lakintab.RU then
          y = y + lakintab.d(x.height) + lakintab.d(x.depth)
        elseif i == lakintab.GL or i == lakintab.KE then
          y = y + lakintab.d(x.width)
        end
      end
      x = x.next
    end
  end
  function lakintab.fix(target, maxh)
    local nst = tex.nest[tex.nest.ptr]
    if nst == nil then return end
    local h = nst.head
    if h == nil then return end
    if h.id == lakintab.TM then h = h.next end
    local tot = 0
    local x = h
    while x do
      local i = x.id
      if i == lakintab.HL or i == lakintab.VL or i == lakintab.RU then
        tot = tot + lakintab.d(x.height) + lakintab.d(x.depth)
      elseif i == lakintab.GL or i == lakintab.KE then
        tot = tot + lakintab.d(x.width)
      end
      x = x.next
    end
    if tot > maxh then return end
    lakintab.scan(h, 0, false, nil, target)
  end
}
% @-FREE NAMES ON PURPOSE: these two are named in the `tablebox` key list at
% \newmdenv above, which is tokenised OUTSIDE any \makeatletter region — a
% \lakin@... there would tokenise as \lakin followed by the letters @tabopen.
\newcommand{\lakintabboxopen}{\lakintabframedtrue\global\lakintabcapseenfalse}
\newcommand{\lakintabboxclose}{%
  \iflakintabcapseen
    \directlua{lakintab.fix(-\number\dimexpr\lakintabcushion\relax,
                            \number\dimexpr\textheight\relax)}%
    \global\lakintabcapseenfalse
  \fi\lakintabframedfalse}
% The cushion above the box's CONTENT top. 42pt puts the dest 32.0pt above the
% frame top (42pt - innertopmargin 9pt - linewidth 0.9pt), i.e. exactly where
% the 295 correct dests already are — the correction moves only the broken ones.
\newcommand{\lakintabcushion}{42pt}
\newenvironment{qcirc}{\par\medskip\global\lakinfigraise 0pt\relax\begin{otherlanguage}{english}\begin{center}}
  {\csname lakin@grabfig\endcsname\end{center}\end{otherlanguage}\par\medskip}
\newenvironment{tikzpic}{\par\medskip\global\lakinfigraise 0pt\relax\begin{otherlanguage}{english}\begin{center}}
  {\csname lakin@grabfig\endcsname\end{center}\end{otherlanguage}\par\medskip}

% ============================================================================
% Code listings.  ★★ v3.4 THE BOX-OVERFLOW FIX:
%   resetmargins=true   → the listing IGNORES the enclosing environment's
%                         indentation and measures from the PAGE. v3.4 stated
%                         this backwards ("measured from the ENCLOSING box"),
%                         and v3.37 then blanked in-box line numbers because
%                         the page-anchored listing put its hanging digits on
%                         top of the box rule. Measured v4.5.0: with the key
%                         FALSE the listing follows the box and the code moves
%                         from x=75.26 to x=97.74, clearing the rule at 72.07.
%                         The key now lives in one \lstset below instead of in
%                         each style, so a box can override it — codebox does.
%   framexleftmargin=0  → the frame (left bar) does not extend left of the code.
%   xleftmargin keeps a small indent so the line numbers have room INSIDE.
% asmblock wraps the listing in an LTR island (code is LTR).
% ============================================================================
\lstdefinestyle{asm}{
  basicstyle=\ttfamily\small,
  backgroundcolor=\color{cBlueBg!55},
  frame=leftline,framerule=2.5pt,rulecolor=\color{cBlue},
  numbers=left,numberstyle=\tiny\color{cGray},numbersep=8pt,
  xleftmargin=2.2em,framexleftmargin=0pt,framexrightmargin=0pt,
  breaklines=true,columns=fullflexible,keepspaces=true,
  showstringspaces=false,aboveskip=8pt,belowskip=8pt,
  % v3.36: without a comment DELIMITER the commentstyle below never fires —
  % listings has no `language` here, so it does not know '#' opens a comment,
  % and every asm listing rendered flat black (measured: a whole-page pixel
  % sweep found ZERO cGreen glyphs across 135 shipped listings). morecomment
  % gives it the delimiter; color-only change, glyph widths untouched.
  morecomment={[l]{\#}},
  commentstyle=\color{cGreen},
}
\lstdefinestyle{cstyle}{
  language=C,basicstyle=\ttfamily\small,
  backgroundcolor=\color{cGreenBg!55},
  frame=leftline,framerule=2.5pt,rulecolor=\color{cGreen},
  numbers=left,numberstyle=\tiny\color{cGray},numbersep=8pt,
  xleftmargin=2.2em,framexleftmargin=0pt,framexrightmargin=0pt,
  breaklines=true,columns=fullflexible,keepspaces=true,
  showstringspaces=false,keywordstyle=\color{cBlue}\bfseries,
  commentstyle=\color{cGray}\itshape,aboveskip=8pt,belowskip=8pt,
}
% v3.36: pystyle/htmlstyle — colored twins of cstyle for the Python/HTML
% listings that used to repeat a ~200-char options blob with NO color
% (keywords rendered bold-black, comments italic-black — measured; C alone
% was colored). Geometry is copied VERBATIM from the shipped blob (frame,
% margins, skips) so swapping a call site to style=pystyle moves NOTHING —
% only keywordstyle/commentstyle/stringstyle are added (bold stays bold,
% italic stays italic, mono metrics identical => zero reflow, proven by
% per-page word-bbox diff). breaklines/breakatwhitespace stay per-site.
% v3.78: numbers=left joins both styles, with the exact numberstyle/numbersep
% keys asm/cstyle carry — fleet-wide consistency (author decision: all four
% styles numbered; "asm/cstyle numbered, py/html not" was an accident of the
% v3.36 geometry freeze, not a design). The numbers hang LEFT of the code box
% (listings sets them with \llap, so no overfull hbox can arise and the code
% x is untouched — measured: code frozen, digits at code_x-11.8pt); inside a
% colored box they are blanked like everyone else's (v3.37 hook), and the
% clean variant restores them in its frameless boxes (v3.78, above).
\lstdefinestyle{pystyle}{
  language=Python,basicstyle=\ttfamily\small,
  columns=fullflexible,keepspaces=true,
  frame=single,rulecolor=\color{cTblLine},
  numbers=left,numberstyle=\tiny\color{cGray},numbersep=8pt,
  xleftmargin=0.5em,framexleftmargin=0pt,framexrightmargin=0pt,
  showstringspaces=false,aboveskip=8pt,belowskip=8pt,
  keywordstyle=\color{cBlue}\bfseries,
  commentstyle=\color{cGray}\itshape,
  stringstyle=\color{cGreen},
}
% v4.5.0: hoisted out of the four styles so an enclosing box can override it.
% Every listing still resolves to true unless something later says otherwise,
% so shipped output is unchanged (verified byte-for-byte on all eight PDFs).
\lstset{resetmargins=true}
\lstdefinestyle{htmlstyle}{
  language=HTML,basicstyle=\ttfamily\small,
  columns=fullflexible,keepspaces=true,
  frame=single,rulecolor=\color{cTblLine},
  numbers=left,numberstyle=\tiny\color{cGray},numbersep=8pt,
  xleftmargin=0.5em,framexleftmargin=0pt,framexrightmargin=0pt,
  showstringspaces=false,aboveskip=8pt,belowskip=8pt,
  keywordstyle=\color{cBlue}\bfseries,
  commentstyle=\color{cGray}\itshape,
  stringstyle=\color{cGreen},
}
\newenvironment{asmblock}{\par\begin{otherlanguage}{english}}{\end{otherlanguage}\par}

% <hllpdf:listing-smart-quotes>
% v4.21.0: four Unicode punctuation characters must enter listings' token
% buffer. Unregistered characters bypass it: a closing quote was printed
% BEFORE the pending word, and an opening quote BEFORE a pending brace.
% Extend the standard character table, then activate only these codepoints
% in the listing's local character-table scope. Do this BEFORE the package's
% delimiter/alsoother/literate hooks so explicit user mappings still win.
% This is not universal Unicode support, nor a conversion to ASCII quotes.
% See listings developer guide 15.3.1 and references/listing-smart-quotes.md.
\makeatletter
\lst@CCPut \lst@ProcessOther {"2018}{"2019}{"201C}{"201D}\z@
\lst@AddToHookAtTop{SelectCharTable}{%
  \lst@ifactivechars
    \catcode"2018=\active \catcode"2019=\active
    \catcode"201C=\active \catcode"201D=\active
  \fi}
\makeatother
% </hllpdf:listing-smart-quotes>

% <hllpdf:listing-technical-symbols>
% v4.30: bounded extension for common technical punctuation in LTR listings.
% Without character-table registration these glyphs bypass the token buffer
% and can move before the preceding word even inside quoted strings.
% Keep the original four-quote block separate for its standing mutation test.
% This is source-faithful printing, not syntax validation or general Unicode.
\makeatletter
\lst@CCPut \lst@ProcessOther
  {"2013}{"2014}{"2026}{"2190}{"2192}{"2194}{"21D2}
  {"2212}{"221E}{"2260}{"2264}{"2265}\z@
\lst@AddToHookAtTop{SelectCharTable}{%
  \lst@ifactivechars
    \catcode"2013=\active \catcode"2014=\active \catcode"2026=\active
    \catcode"2190=\active \catcode"2192=\active \catcode"2194=\active
    \catcode"21D2=\active \catcode"2212=\active \catcode"221E=\active
    \catcode"2260=\active \catcode"2264=\active \catcode"2265=\active
  \fi}
\makeatother
% </hllpdf:listing-technical-symbols>


\usepackage[hidelinks,unicode]{hyperref}   % LAST, after colors. unicode = Hebrew bookmarks.
% v3.62: THE ISOLATION MACROS CORRUPTED EVERY PDF BOOKMARK THEY APPEARED IN.
% \en is defined (line ~155) as
%     \foreignlanguage{english}{\language=\csname l@nohyphenation\endcsname #1}
% and in a heading that body is written verbatim to the .toc and handed to
% hyperref's \pdfstringdef. hyperref strips the two CONTROL SEQUENCES (\language,
% \l@nohyphenation) but the plain `=` CHARACTER token between them is not a control
% sequence, so it survives into the PDF string: \chapter{שער \en{CNOT} והצמדה}
% produced the bookmark "שער =CNOT והצמדה". The printed TOC and the heading on the
% page are CORRECT — only the navigation pane and the viewer window title are wrong,
% which is why it went unnoticed. Measured on the shipped fleet: 70 corrupted
% bookmarks across 5 PDFs, from 177 unprotected \en-in-heading sites. Same corruption
% reached /Title and /Subject via the \hypersetup{pdftitle=...} the docs recommend.
% This hook tells hyperref to use the plain argument for all three isolation macros
% when building a PDF string. It does NOT touch the .toc (the printed TOC is
% byte-unchanged) and it does NOT mask the math signal — a heading with $...$ still
% logs "removing math shift", so BIDI-15 detection is intact.
% BONUS, measured: it also makes \code{} in an UNSTARRED heading COMPILE. BLD-10
% recorded that as needing \texorpdfstring or a starred heading; the real root cause
% is hyperref's \pdfstringdef scan of \code's \ifmmode...\else...\fi body — which is
% exactly why the identically-bodied \texttt survived (hyperref pre-neutralizes
% \texttt in PDF strings) and why \DeclareRobustCommand could never help.
% v3.64: \def\\{ } added. A manual line break in a heading WELDS the two words
% together in the bookmark: \chapter{שורה ראשונה\\שורה שנייה} produced the outline
% entry 'שורה ראשונהשורה שנייה' — "ראשונהשורה" is not a word. hyperref drops the
% \\ and puts nothing in its place, so the break that looks like a space on the
% page is a zero-width join in the navigation pane. The spaced form
% 'מילה אחת \\ מילה שנייה' escapes with a harmless DOUBLE space ('אחת  מילה'),
% which is why this never got noticed. Both are measured on a real outline.
% 0 live instances in the corpus — but SKILL.md and content-style.md both RECOMMEND
% a manual heading break, so this is a primed trap rather than a dormant one, and a
% one-token preamble fix retires it for good instead of asking authors to remember.
% Mapping \\ to a space fixes the WELD, which is the actual defect (a nonexistent
% word in the navigation pane). It does not normalise whitespace: measured after the
% fix, the welded form becomes 'שורה ראשונה שורה שנייה' (correct) and the already-
% spaced form goes from two spaces to three. That is cosmetic — viewers render it as
% a wider gap, and no word is corrupted. A per-heading \texorpdfstring is still the
% way to get an exactly-tuned bookmark string.
\pdfstringdefDisableCommands{\def\en#1{#1}\def\h#1{#1}\def\doclang#1{#1}\def\code#1{#1}\def\codeline#1{#1}\def\\{ }}
% v3.70: RESTORE \url LINE-BREAKING. url.sty freezes its break penalties at LOAD
% time (\mathchardef\UrlBreakPenalty=\binoppenalty), and v3.45's global
% \binoppenalty=10000 (inline-math atomicity) therefore priced EVERY breakpoint
% inside \Url at 10000 = never break. Measured: a long \url in Hebrew prose is
% one unbroken line, ~100pt overfull, invading the margin band the ink gates
% deliberately tolerate. With the two lines below: 0 overfull, the URL breaks
% after '/', the RTL split is CORRECT (fragment order forward), and each line
% gets its own link annotation exactly covering its fragment (measured).
% Confined to \Url's group — inline-math atomicity (TYP-7) is untouched.
\mathchardef\UrlBreakPenalty=200
\mathchardef\UrlBigBreakPenalty=100
\ifx\lakitmainlanguage\hllpdfenglish
  \hypersetup{pdflang=en}
\else
  \hypersetup{pdflang=he}
\fi  % /Lang is the document base locale (not per-span accessibility tagging).
% Set a per-document \hypersetup{pdftitle={...}} too — see SKILL.md.

% v3.52 (round-2 audit): two RTL bibliography/cross-ref defects, both invisible to
% every gate and both affecting the FIRST author to use the feature.
% (a) \@biblabel's default [#1] brackets are bidi-neutral and mirror to ]1[ under
% the RTL bibliography base — every \thebibliography reference list shows ]1[.
% LTR-isolating the label restores [1]. (The render gate's mirror rule fingerprints
% round parens )1(, so square brackets slipped through.)
\makeatletter
\renewcommand*{\@biblabel}[1]{\foreignlanguage{english}{[#1]}}
% v3.54: the SAME neutral-bracket trap on the INLINE side — a multi-key \cite{a,b}
% under the RTL base lays its European-number list right-to-left, so [1, 2] renders
% "[2,1]" and [1, 2, 3] renders "[3,2,1]" (verified by glyph geometry — the reader
% sees the sources in REVERSE order; a note form \cite[p.5]{a} also reorders to
% "[p. 5,1]"). LTR-isolate the whole bracket group (number list AND optional note),
% exactly like \@biblabel above: [1, 2] / [1, 2, 3] / [1, p. 5] all render correct.
% Placed AFTER hyperref so it wins; hyperref hyperlinks each number inside #1 (via
% \@citex), which this wrapper leaves intact. Mirrors the LaTeX kernel \@cite shape.
% v4.9.1: directional isolation must NOT switch the language here. Babel's
% safe-active citation wrapper passes its key list through \bbl@tempa;
% \foreignlanguage reuses that scratch macro before the deferred citation
% loop expands it. The old line therefore cited the literal key "english".
% If such a key existed, even the undefined-reference gate stayed green.
% \babelsublr isolates the bracket group without selecting a locale or
% changing that key list. Preserve native links, key order and optional notes.
\def\@cite#1#2{\babelsublr{[{#1\if@tempswa , #2\fi}]}}
\makeatother

% v4.11.0: a reference inside align's \text inherited HEBdisptag=true,
% so its enclosing hbox mirrored the brackets despite a clean compile.
% Isolate the complete printed reference without selecting a Babel language.
% Also neutralize the display-only normalfont direction reset locally. A
% direction-only patch inverted 12 to 21 inside equation in the measured probe.
% Plain maketag is local to the reference: mathtools' display-side wrapper
% otherwise globally consumes manual_tag while formatting a nested reference.
% Preserve each original command (including the selective auxiliary write),
% current tag forms, native links, and showonlyrefs on/off transitions.
\makeatletter
\NewCommandCopy{\HEB@nativeeqref}{\eqref}
\newcommand{\HEB@eqrefisolate}[1]{%
  \babelsublr{\HEBdisptagfalse\let\bbl@eqnodir\relax\let\maketag@@@\ORIGmaketag@@@#1}}
\protected\def\eqref#1{\HEB@eqrefisolate{\HEB@nativeeqref{#1}}}
\ExplSyntaxOn
\NewCommandCopy{\HEB@nativeMTeqref}{\MT_extended_eqref:n}
\cs_set_protected:Npn \MT_extended_eqref:n #1
  {\HEB@eqrefisolate{\HEB@nativeMTeqref{#1}}}
\ExplSyntaxOff
\makeatother

% v4.11.0: custom caption counters are not part of hyperref's English
% locale, and Hebrew has no upstream hyperref locale in this toolchain.
% Restore the appropriate names at every Babel language switch, including
% return from a nested inline island. Do not localize the destination keys.
\makeatletter
\newcommand{\HEB@refnames}{%
\def\chapterautorefname{פרק}%
\def\sectionautorefname{סעיף}%
\def\subsectionautorefname{תת־סעיף}%
\def\subsubsectionautorefname{תת־תת־סעיף}%
\def\paragraphautorefname{פסקה}%
% v4.31: hyperref uses a distinct counter name at the deepest standard level.
\def\subparagraphautorefname{תת־פסקה}%
\def\equationautorefname{משוואה}%
\def\figureautorefname{איור}%
\def\tableautorefname{טבלה}%
\def\partautorefname{חלק}%
\def\appendixautorefname{נספח}%
\def\footnoteautorefname{הערת שוליים}%
\def\itemautorefname{סעיף}%
\def\Itemautorefname{סעיף}%
\def\figautorefname{איור}%
\def\tblautorefname{טבלה}%
\def\pageautorefname{עמוד}% 
}
\addto\extrashebrew{\HEB@refnames}
\newcommand{\ENG@refnames}{%
  \def\subparagraphautorefname{subparagraph}%
  \def\chapterautorefname{chapter}%
  \def\sectionautorefname{section}%
  \def\subsectionautorefname{subsection}%
  \def\subsubsectionautorefname{subsubsection}%
  \def\paragraphautorefname{paragraph}%
  \def\equationautorefname{Equation}%
  \def\figureautorefname{Figure}%
  \def\tableautorefname{Table}%
  \def\partautorefname{Part}%
  \def\appendixautorefname{Appendix}%
  \def\footnoteautorefname{footnote}%
  \def\itemautorefname{item}%
  \def\figautorefname{Figure}%
  \def\tblautorefname{Table}%
  \def\Itemautorefname{item}%
  \def\pageautorefname{page}%
}
\addto\extrasenglish{\ENG@refnames}
\ifx\lakitmainlanguage\hllpdfenglish\ENG@refnames\else\HEB@refnames\fi
\makeatother

% ============================================================================
% Optional running head (v3.39): \lakithead{<book title>} — OFF by default.
% A document opts in with ONE line after \input{preamble.tex}:
%   \lakithead{מבנה המחשב}
% Every content page then carries the series' orientation line: chapter name
% at the reading edge, book title opposite, thin rule. (fancyhdr's [L]/[R]
% are PHYSICAL; under the RTL page direction [L] lands at the RIGHT/reading
% edge — measured, consistent across the corpus.) Packaged like the topic
% rail: always loaded, inert without the call. This centralizes the 8-line
% block that 8 lakits used to copy by hand — the same duplication vector
% that once let a lakit preamble go ~100 lines stale. Class-aware: report
% fills the mark from \chapter, article from \section. NB: a titlepage env
% self-empties; front matter that sets \pagestyle{empty} must restore with
% \pagestyle{fancy} afterwards (NOT plain — that discards the head).
% ============================================================================
\usepackage{fancyhdr}
% v3.61: \chapter* left a STALE running head — pages of an UNNUMBERED chapter
% carried the PREVIOUS chapter's title. report.cls routes \chapter* through
% \@schapter, which never calls \chaptermark, so \leftmark keeps whatever the last
% numbered chapter left. Measured: with \chapter{פרק ממוספר} then
% \chapter*{מבוא לא ממוספר}, page 4 (inside the starred chapter) printed
% "פרק ממוספר". Invisible on the chapter's FIRST page (\thispagestyle{plain}), so it
% only shows from the second page on — the pages an author is least likely to check.
% An unnumbered front-matter chapter (\chapter*{מבוא}, \chapter*{הקדמה}) is everyday
% Hebrew-academic structure, and the corpus already carries HAND-WRITTEN \markboth
% lines right after \chapter* — the exact per-document duplication \lakithead exists
% to retire. Patched INSIDE \lakithead so the macro stays inert when unused; a
% document's own later \markboth still wins, so those hand-patches do not regress.
\makeatletter
% <hllpdf:starred-heading-scope>
\newif\iflakin@starmarkpatched
\newcommand{\lakistarmark}{%
  \iflakin@starmarkpatched\else
  % Dispatch on the actual class counter, not hyperref's shared internals.
  \ifdefined\c@chapter
    \let\HEBorigschapter\@schapter
    % Record only after the native heading has reached its page. A mark
    % before the native routine can advertise the next heading too early.
    \def\@schapter##1{\HEBorigschapter{##1}\chaptermark{##1}}%
  \else
    % @ssect receives layout arguments, NOT the section name or level.
    % Capture the name in @startsection without introducing a TeX group:
    % grouping sectioning would break the native after-heading state.
    \def\HEB@sectionkind{}%
    \def\HEB@mainsection{section}%
    \let\HEBorigstartsection\@startsection
    \def\@startsection##1##2##3##4##5##6{%
      \def\HEB@sectionkind{##1}%
      \HEBorigstartsection{##1}{##2}{##3}{##4}{##5}{##6}}%
    \let\HEBorigssect\@ssect
    \def\@ssect##1##2##3##4##5{%
      \HEBorigssect{##1}{##2}{##3}{##4}{##5}%
      % Lower starred headings must not replace the article's section mark.
      \ifx\HEB@sectionkind\HEB@mainsection
        \sectionmark{##5}%
      \fi
      \ignorespaces}%
  \fi
  \global\lakin@starmarkpatchedtrue
  \fi}
% </hllpdf:starred-heading-scope>
% <hllpdf:running-language>
% Choose the wrapper when the mark is written, not when a later page ships.
% English punctuation must remain LTR even after returning to Hebrew prose.
\newcommand{\HEB@runningmark}[1]{%
  \iflanguage{english}{\markboth{\en{#1}}{}}{\markboth{\h{#1}}{}}}
% </hllpdf:running-language>
\makeatother
% v4.2.1 (NAV): every collection already names itself in \lakithead — feed the
% PDF /Title from the same declaration, so the navigation gate's
% PDF_REQUIRE_TITLE has a universal source and no fleet member ships
% title-less by missing a second declaration site. Measured before this: 14
% of 16 collections shipped /Title-empty while every one carried a perfectly
% good \lakithead string (plain Hebrew in all 16 — no math, no formatting, so
% the PDF-string path is clean). An explicit title always wins: set it with
% \lakitpdftitle{...} anywhere before \lakithead, or with a plain
% \hypersetup{pdftitle=...} — the fallback detects that via hyperref's
% \@pdftitle, probe-verified \@empty-equal until the first explicit set.
\makeatletter
\newif\iflakit@pdftitleset
\newcommand{\lakitpdftitle}[1]{%
  \hypersetup{pdftitle={#1}}\global\lakit@pdftitlesettrue}
\newcommand{\lakithead}[1]{%
  \iflakit@pdftitleset\else
    \ifx\@pdftitle\@empty
      \hypersetup{pdftitle={#1}}\global\lakit@pdftitlesettrue
    \fi
  \fi
  \pagestyle{fancy}\fancyhf{}%
  \lakistarmark

  \fancyhead[L]{\small\textcolor{cGray}{\leftmark}}%
  % <hllpdf:fixed-head-language>
  % The fixed document title has its OWN language snapshot, independent of
  % the current section/chapter mark. Choose the branch NOW, not at shipout:
  % the page may be Hebrew by the time an English title is finally printed.
  \iflanguage{english}{%
    \fancyhead[R]{\small\textcolor{cGray}{\en{#1}}}%
  }{%
    \fancyhead[R]{\small\textcolor{cGray}{\h{#1}}}%
  }%
  % </hllpdf:fixed-head-language>
  \fancyfoot[C]{\thepage}%
  \renewcommand{\headrulewidth}{0.4pt}%
  \ifdefined\chaptermark
    \renewcommand{\chaptermark}[1]{\HEB@runningmark{##1}}%
  \else
    \renewcommand{\sectionmark}[1]{\HEB@runningmark{##1}}%
  \fi}
\makeatother

% v3.39.1: path-independent byte-reproducibility. LuaTeX derives the trailer
% /ID from the OUTPUT PATH, so two clones of the same source at different
% directories produced different sha256 (measured: the ONLY delta was the
% 62-byte /ID line). Bit 512 suppresses exactly that line — measured: Creator/
% Producer/CreationDate/ModDate and all content untouched. /ID is optional
% per ISO 32000 (required only for ENCRYPTED PDFs and PDF/A — neither is a
% target here). If a future workflow needs an ID (PDF/A, Acrobat revision
% matching), delete this line and it returns.
\pdfvariable suppressoptionalinfo \numexpr 512\relax

% v3.75: BOLD MATH IN BOLD HEADINGS. Heading text is \bfseries but math inside
% it stayed regular-weight CM — a hairline $V_T$ inside a heavy 24.8pt chapter
% title (measured: CMMI12 fl=6 next to FrankRuhlLibre-Bold fl=20; 155 live
% heading-with-math sites fleet-wide). \boldmath rides along wherever the class
% sets \bfseries for a heading; A/B measured: $V_T$ -> CMMIB10 fl=22, 0 missing
% characters (the bold CM faces ship with this TeX Live). Guarded per-macro so
% the article class (no \@makechapterhead) and future class changes stay safe.
% v4.0.1 (adoption fix): the SEARCH patterns above carry no parameter token.
% v4.0.0 made them specific by including the class's own #1/#2, and measured on
% TeX Live 2023 that never matches: a \patchcmd search pattern containing a
% parameter token cannot match while the call sits inside \@ifundefined{...}{}{...}
% (probed both #1 and ##1 -> FAIL; the same pattern at top level succeeds). Every
% one of these patches therefore failed, and because v4.0.0 also replaced the
% empty failure branch with \PackageError, the failure became FATAL: all four
% bundled example assets stopped building ("Cannot patch chapter-title math
% weight -> Fatal error, no output PDF"). The \PackageError branches are KEPT —
% a silent no-op here is what let v3.75's fix rot unnoticed — only the parameter
% tokens are dropped, so the patches match and a real class change still errors.
\makeatletter
\@ifundefined{@makechapterhead}{}{%
  \patchcmd{\@makechapterhead}{\huge\bfseries \@chapapp}
    {\huge\bfseries\boldmath \@chapapp}{}
    {\PackageError{hebrew-lualatex-pdf}{Cannot patch chapter-number math weight}{Unsupported document class}}%
  \patchcmd{\@makechapterhead}{\Huge \bfseries}
    {\Huge \bfseries\boldmath}{}
    {\PackageError{hebrew-lualatex-pdf}{Cannot patch chapter-title math weight}{Unsupported document class}}%
  \patchcmd{\@makeschapterhead}{\Huge \bfseries}
    {\Huge \bfseries\boldmath}{}
    {\PackageError{hebrew-lualatex-pdf}{Cannot patch starred chapter math weight}{Unsupported document class}}}
% v4.0.1: \part sits under a THREE-DEEP wrapper chain, so the class-based anchor
% is not in \@part at all. Measured with \show: hyperref rewrites \@part to
% `...\phantomsection\fi \H@old@part`, and nameref rewrites THAT to
% `\NR@gettitle{#1}\NR@part[{#1}]{#2}` — the class body ends up in \NR@part.
% v4.0.0 patched \@part and made failure fatal, which is why all seven bundled
% examples stopped building. \part is NOT a dead heading level here: 188 live uses
% across the fleet, so the bold-math fix is worth getting right rather than
% dropping. We walk the chain innermost-first and stop at the first hit. Failure is
% a WARNING, not an error: the chain is a property of hyperref/nameref versions, so
% a future package release must degrade to regular-weight part math — a cosmetic
% loss — instead of refusing to build. Chapter and section keep \PackageError,
% because those targets are in the CLASS and were verified to patch here.
\newif\iflakin@partnum \newif\iflakin@parttitle \newif\iflakin@sparttitle
\def\lakin@partnum#1#2{\iflakin@partnum\else
  \patchcmd{#1}{#2\bfseries \partname}{#2\bfseries\boldmath \partname}
    {\lakin@partnumtrue}{}\fi}
\def\lakin@parttitle#1#2{\iflakin@parttitle\else
  \patchcmd{#1}{#2 \bfseries}{#2 \bfseries\boldmath}{\lakin@parttitletrue}{}\fi}
\def\lakin@sparttitle#1#2{\iflakin@sparttitle\else
  \patchcmd{#1}{#2 \bfseries}{#2 \bfseries\boldmath}{\lakin@sparttitletrue}{}\fi}
\@ifundefined{c@chapter}{\def\lakin@pnum{\Large}\def\lakin@ptit{\huge}}%
                        {\def\lakin@pnum{\huge}\def\lakin@ptit{\Huge}}
\@ifundefined{NR@part}{}{\expandafter\lakin@partnum\expandafter\NR@part\expandafter{\lakin@pnum}%
                         \expandafter\lakin@parttitle\expandafter\NR@part\expandafter{\lakin@ptit}}
\@ifundefined{H@old@part}{}{\expandafter\lakin@partnum\expandafter\H@old@part\expandafter{\lakin@pnum}%
                            \expandafter\lakin@parttitle\expandafter\H@old@part\expandafter{\lakin@ptit}}
\@ifundefined{@part}{}{\expandafter\lakin@partnum\expandafter\@part\expandafter{\lakin@pnum}%
                       \expandafter\lakin@parttitle\expandafter\@part\expandafter{\lakin@ptit}}
\@ifundefined{NR@spart}{}{\expandafter\lakin@sparttitle\expandafter\NR@spart\expandafter{\lakin@ptit}}
\@ifundefined{H@old@spart}{}{\expandafter\lakin@sparttitle\expandafter\H@old@spart\expandafter{\lakin@ptit}}
\@ifundefined{@spart}{}{\expandafter\lakin@sparttitle\expandafter\@spart\expandafter{\lakin@ptit}}
\iflakin@partnum\else\PackageWarning{hebrew-lualatex-pdf}{Part-number math stays
  regular weight: no known \noexpand\part wrapper matched (hyperref/nameref
  version change). Cosmetic only}\fi
\iflakin@parttitle\else\PackageWarning{hebrew-lualatex-pdf}{Part-title math stays
  regular weight: no known \noexpand\part wrapper matched. Cosmetic only}\fi
% v4.12.0: THE PART NUMBER IN THE TOC WAS NOT DIRECTIONALLY SEPARATED FROM THE
% TITLE. report.cls/article.cls compose the part entry as ONE text stream --
% \addcontentsline{toc}{part}{\thepart\hspace{1em}#1} -- while EVERY other level
% goes through \numberline, which puts the number in an \hbox. A box is a bidi
% OBJECT and cannot join a neighbouring run; a bare "IV" can. So in the RTL TOC
% paragraph a Latin number and a Latin-INITIAL title merge into a SINGLE LTR run
% and are laid out left-to-right as a unit: the number lands on the wrong side of
% the row, and a mixed title is scrambled outright. Measured on \part{MCP ...}:
% the printed order was "<tail> V MCP" -- title fragment, number, title fragment.
% The chapter/section rows one line below, with the SAME Latin-initial titles,
% were correct on the same page. That contrast is the whole diagnosis.
% \en{} IS NOT A WORKAROUND, measured both ways on one fixture: it fixes a
% PURE-Latin title (\part{\en{Hooks}} renders correctly) and does NOT fix a MIXED
% one (\part{\en{MCP} <hebrew>} stays scrambled). babel's language switch changes
% font and hyphenation, not the bidi run structure -- so the defect is not
% reachable from the source side at all, which is why it belongs here.
% Boxing the number is the entire fix and is a STRICT NO-OP where the bug never
% bit: a Hebrew-titled document compiles BYTE-IDENTICAL before and after (md5,
% 3 clean runs), because there the neutral glue already sat between an L run and
% an R run. The PDF outline is byte-identical too -- hyperref's \pdfstringdef
% neutralises \mbox, so bookmark strings keep their part numbers.
% \numberline{\thepart} would also fix the direction, but it re-geometries every
% part row in the fleet and DROPS the number from the part bookmark (hyperref
% strips \numberline). \mbox changes nothing that was already right.
% The search pattern deliberately carries no # token: etoolbox cannot match a
% parameter token in a patch pattern (measured -- all three chain members
% reported "not patchable" while the pattern included #1). "\thepart \hspace
% {1em}" is brace-balanced and occurs exactly once in the wrapper body.
\newif\iflakin@parttocnum
\def\lakin@parttocnum#1{\iflakin@parttocnum\else
  \patchcmd{#1}{\thepart \hspace {1em}}{\mbox {\thepart }\hspace {1em}}
    {\lakin@parttocnumtrue}{}\fi}
\@ifundefined{NR@part}{}{\lakin@parttocnum\NR@part}
\@ifundefined{H@old@part}{}{\lakin@parttocnum\H@old@part}
\@ifundefined{@part}{}{\lakin@parttocnum\@part}
\iflakin@parttocnum\else\PackageWarning{hebrew-lualatex-pdf}{Part TOC number is
  not boxed: no known \noexpand\part wrapper matched. A part title that BEGINS
  with Latin will print its number on the wrong side of its TOC row}\fi
\patchcmd{\section}{\bfseries}{\bfseries\boldmath}{}
  {\PackageError{hebrew-lualatex-pdf}{Cannot patch section math weight}{Unsupported document class}}
\patchcmd{\subsection}{\bfseries}{\bfseries\boldmath}{}
  {\PackageError{hebrew-lualatex-pdf}{Cannot patch subsection math weight}{Unsupported document class}}
\patchcmd{\subsubsection}{\bfseries}{\bfseries\boldmath}{}
  {\PackageError{hebrew-lualatex-pdf}{Cannot patch subsubsection math weight}{Unsupported document class}}
\patchcmd{\paragraph}{\bfseries}{\bfseries\boldmath}{}{}
\patchcmd{\subparagraph}{\bfseries}{\bfseries\boldmath}{}{}
\@ifundefined{l@chapter}{}{%
  \patchcmd{\l@chapter}{\bfseries}{\bfseries\boldmath}{}{}}
\makeatother
\setcounter{tocdepth}{\lakittocdepth}
% <hllpdf:measured-toc>
% v4.16.0: Measure actual number labels, not a guessed digit count. Read only
% canonical contentsline records; never execute titles or addtocontents code
% during the measuring pass. Widths are local to this individual list and are
% available before its first row: no per-row drift and no extra aux pass.
\makeatletter
\newread\hllpdf@tocstream
\newdimen\hllpdf@toc@chapter
\newdimen\hllpdf@toc@section
\newdimen\hllpdf@toc@subsection
\newdimen\hllpdf@toc@subsubsection
\newdimen\hllpdf@toc@paragraph
\newdimen\hllpdf@toc@subparagraph
\newdimen\hllpdf@toc@figure
\newdimen\hllpdf@toc@table
\newdimen\hllpdf@tocpage
\newdimen\hllpdf@toc@sectionindent
\newdimen\hllpdf@toc@subsectionindent
\newdimen\hllpdf@toc@subsubsectionindent
\newdimen\hllpdf@toc@paragraphindent
\newdimen\hllpdf@toc@subparagraphindent
\def\hllpdf@tocdefaults{%
  \hllpdf@toc@chapter=2em
  \hllpdf@toc@section=3.2em
  \hllpdf@toc@subsection=3.9em
  \@ifundefined{c@chapter}{%
    \hllpdf@toc@subsubsection=3.2em
    \hllpdf@toc@paragraph=4.1em
    \hllpdf@toc@subparagraph=5em
  }{%
    \hllpdf@toc@subsubsection=4.1em
    \hllpdf@toc@paragraph=5em
    \hllpdf@toc@subparagraph=6em
  }%
  \hllpdf@toc@figure=2.3em
  \hllpdf@toc@table=2.3em
  \hllpdf@tocpage=\@pnumwidth\relax
}
\def\hllpdf@tocindents{%
  % Preserve the established shallow layout when its slots already fit.
  \hllpdf@toc@sectionindent=1.5em
  \@tempdima=\hllpdf@toc@chapter \advance\@tempdima by -2em
  \advance\hllpdf@toc@sectionindent by \@tempdima
  \hllpdf@toc@subsectionindent=\dimexpr
    \hllpdf@toc@sectionindent+\hllpdf@toc@section\relax
  \hllpdf@toc@subsubsectionindent=\dimexpr
    \hllpdf@toc@subsectionindent+\hllpdf@toc@subsection\relax
  \hllpdf@toc@paragraphindent=\dimexpr
    \hllpdf@toc@subsubsectionindent+\hllpdf@toc@subsubsection\relax
  \hllpdf@toc@subparagraphindent=\dimexpr
    \hllpdf@toc@paragraphindent+\hllpdf@toc@paragraph\relax
}
\def\hllpdf@tocfont@chapter{\bfseries\boldmath}
\def\hllpdf@tocfont@part{\large\bfseries\boldmath}
\def\hllpdf@tocentryfont{%
  \ifcsname hllpdf@tocfont@\hllpdf@toclevel\endcsname
    \csname hllpdf@tocfont@\hllpdf@toclevel\endcsname
  \fi}
\long\def\hllpdf@tocdiscard#1\hllpdf@tocstop{}
\def\hllpdf@tocscanrecord{%
  \@ifnextchar\contentsline{\hllpdf@tocscanline}{%
    \@ifnextchar\babel@toc{\hllpdf@tocscanlanguage}{\hllpdf@tocdiscard}}}
\long\def\hllpdf@tocscanlanguage\babel@toc#1#2#3\hllpdf@tocstop{%
  % Babel records are declarations, not arbitrary code. Track the language
  % without executing the original directive or changing the surrounding list.
  \def\hllpdf@toclanguage{#1}}
\long\def\hllpdf@tocscanline\contentsline#1#2#3#4#5\hllpdf@tocstop{%
  \def\hllpdf@toclevel{#1}%
  \ifcsname hllpdf@toc@#1\endcsname
    \hllpdf@tocscantitle#2\hllpdf@tocstop
  \fi
  % The page-number column needs its own measurement (including bold chapter
  % rows). Its extra clearance separates the page label from dot leaders.
  \setbox\@tempboxa=\hbox{\foreignlanguage{\hllpdf@toclanguage}{\hllpdf@tocentryfont #3\kern.3em}}%
  \ifdim\wd\@tempboxa>\hllpdf@tocpage
    \hllpdf@tocpage=\wd\@tempboxa
  \fi}
\def\hllpdf@tocscantitle{%
  \@ifnextchar\numberline{\hllpdf@tocscanlabel}{\hllpdf@tocdiscard}}
\long\def\hllpdf@tocscanlabel\numberline#1#2\hllpdf@tocstop{%
  \setbox\@tempboxa=\hbox{\foreignlanguage{\hllpdf@toclanguage}{\hllpdf@tocentryfont #1\kern.5em}}%
  \ifdim\wd\@tempboxa>\csname hllpdf@toc@\hllpdf@toclevel\endcsname
    \csname hllpdf@toc@\hllpdf@toclevel\endcsname=\wd\@tempboxa
  \fi}
\def\hllpdf@tocread{%
  \ifeof\hllpdf@tocstream\else
    \read\hllpdf@tocstream to\hllpdf@tocrecord
    \expandafter\hllpdf@tocscanrecord\hllpdf@tocrecord\hllpdf@tocstop
    \expandafter\hllpdf@tocread
  \fi}
\def\hllpdf@tocmeasure#1{%
  \makeatletter
  \edef\hllpdf@toclanguage{\languagename}%
  \hllpdf@tocdefaults
  \openin\hllpdf@tocstream=\jobname.#1\relax
  \hllpdf@tocread
  \closein\hllpdf@tocstream
  \hllpdf@tocindents
  \edef\@pnumwidth{\the\hllpdf@tocpage}%
  \@tempdima=\dimexpr\hllpdf@tocpage+1em\relax
  \ifdim\@tempdima>\@tocrmarg\relax
    \edef\@tocrmarg{\the\@tempdima}%
  \fi}
\let\hllpdf@originalstarttoc\@starttoc
\def\@starttoc#1{%
  \begingroup
    \def\hllpdf@tocext{#1}%
    \def\hllpdf@tocwanted{toc}%
    \ifx\hllpdf@tocext\hllpdf@tocwanted \hllpdf@tocmeasure{#1}\else
      \def\hllpdf@tocwanted{lof}%
      \ifx\hllpdf@tocext\hllpdf@tocwanted \hllpdf@tocmeasure{#1}\else
        \def\hllpdf@tocwanted{lot}%
        \ifx\hllpdf@tocext\hllpdf@tocwanted \hllpdf@tocmeasure{#1}\fi
      \fi
    \fi
    \hllpdf@originalstarttoc{#1}%
  \endgroup}
\hllpdf@tocdefaults
\hllpdf@tocindents
\@ifundefined{l@chapter}{}{%
  \patchcmd{\l@chapter}{\setlength\@tempdima{1.5em}}%
    {\setlength\@tempdima{\hllpdf@toc@chapter}}{}%
    {\PackageError{hebrew-lualatex-pdf}{Cannot size chapter TOC numbers}%
       {Unsupported contents layout: the chapter number slot was not found}}}
\renewcommand*\l@section{\@dottedtocline{1}%
  {\hllpdf@toc@sectionindent}{\hllpdf@toc@section}}
\renewcommand*\l@subsection{\@dottedtocline{2}%
  {\hllpdf@toc@subsectionindent}{\hllpdf@toc@subsection}}
\renewcommand*\l@subsubsection{\@dottedtocline{3}%
  {\hllpdf@toc@subsubsectionindent}{\hllpdf@toc@subsubsection}}
\renewcommand*\l@paragraph{\@dottedtocline{4}%
  {\hllpdf@toc@paragraphindent}{\hllpdf@toc@paragraph}}
\renewcommand*\l@subparagraph{\@dottedtocline{5}%
  {\hllpdf@toc@subparagraphindent}{\hllpdf@toc@subparagraph}}
\renewcommand*\l@figure{\@dottedtocline{1}{1.5em}{\hllpdf@toc@figure}}
\renewcommand*\l@table{\@dottedtocline{1}{1.5em}{\hllpdf@toc@table}}
\makeatother
% </hllpdf:measured-toc>
\setcounter{secnumdepth}{\lakitsecnumdepth}
% Public preamble-only caption-numbering policy. The historical within-major
% default is untouched unless this command is called. "continuous" removes the
% reset edge as well as the printed prefix, so Figure 1/2 and Table 1/2 stay
% unique across sections/chapters. Returning to "within-major" restores the
% original class-aware displays and native destination conventions. The policy
% never replaces \figcap/\tabcap or their measured link-placement mechanism.
\makeatletter
\let\lakin@majorfig\thefig
\let\lakin@majortbl\thetbl
\let\lakin@majorHfig\theHfig
\let\lakin@majorHtbl\theHtbl
\newcommand{\lakitcaptionnumbering}[1]{%
  \edef\lakin@numbering{#1}%
  \def\lakin@continuous{continuous}\def\lakin@withinmajor{within-major}%
  \ifx\lakin@numbering\lakin@continuous
    \@ifundefined{c@chapter}{%
      \@removefromreset{fig}{section}\@removefromreset{tbl}{section}%
    }{\@removefromreset{fig}{chapter}\@removefromreset{tbl}{chapter}}%
    \renewcommand{\thefig}{\arabic{fig}}\renewcommand{\thetbl}{\arabic{tbl}}%
    \def\theHfig{\arabic{fig}}\def\theHtbl{\arabic{tbl}}%
  \else\ifx\lakin@numbering\lakin@withinmajor
    \@ifundefined{c@chapter}{%
      \@removefromreset{fig}{section}\@removefromreset{tbl}{section}%
      \@addtoreset{fig}{section}\@addtoreset{tbl}{section}%
    }{%
      \@removefromreset{fig}{chapter}\@removefromreset{tbl}{chapter}%
      \@addtoreset{fig}{chapter}\@addtoreset{tbl}{chapter}}%
    \let\thefig\lakin@majorfig\let\thetbl\lakin@majortbl
    \let\theHfig\lakin@majorHfig\let\theHtbl\lakin@majorHtbl
  \else
    \PackageError{hebrew-english-latex-studio}{Unknown caption-numbering policy '#1'}%
      {Use continuous or within-major before begin{document}.}%
  \fi\fi}
\@onlypreamble\lakitcaptionnumbering
\makeatother
% Architecture readiness sentinel. Optional runtime layers require the complete
% core, so the marker is defined only after every mechanism/baseline definition.
% <hllpdf:architecture-ready>
\def\hllpdfcoreloaded{1}
% </hllpdf:architecture-ready>

% v4.22: registered file resources, not an assertion about embedded subset bytes.
% The log framing avoids long-path wrapping and preserves UTF-8 as hex.
% <hllpdf:font-inventory>
\directlua{
  local function inventory()
    local names = {}
    local unnamed = 0
    for id, item in font.each() do
      local path = item.filename
      if path and not (path == "") then names[path] = true else unnamed = unnamed + 1 end
    end
    local sorted = {}
    for path in pairs(names) do table.insert(sorted,path) end
    table.sort(sorted)
    texio.write_nl("log","HLLFONT|BEGIN|1")
    for index,path in ipairs(sorted) do
      local hex = path:gsub(".",function(c) return string.format(string.char(37) .. "02x",string.byte(c)) end)
      local count = math.ceil(string.len(hex)/48)
      texio.write_nl("log","HLLFONT|FILE|" .. index .. "|" .. count)
      for offset=1,string.len(hex),48 do texio.write_nl("log","HLLFONT|HEX|" .. hex:sub(offset,offset+47)) end
    end
    texio.write_nl("log","HLLFONT|END|" .. rawlen(sorted) .. "|" .. unnamed)
    texio.write_nl("log","")
  end
  luatexbase.add_to_callback("finish_pdffile",inventory,"hllpdf.runtime.fonts")
}
% </hllpdf:font-inventory>
:
%   \lakithead{מבנה המחשב}
% Every content page then carries the series' orientation line: chapter name
% at the reading edge, book title opposite, thin rule. (fancyhdr's [L]/[R]
% are PHYSICAL; under the RTL page direction [L] lands at the RIGHT/reading
% edge — measured, consistent across the corpus.) Packaged like the topic
% rail: always loaded, inert without the call. This centralizes the 8-line
% block that 8 lakits used to copy by hand — the same duplication vector
% that once let a lakit preamble go ~100 lines stale. Class-aware: report
% fills the mark from \chapter, article from \section. NB: a titlepage env
% self-empties; front matter that sets \pagestyle{empty} must restore with
% \pagestyle{fancy} afterwards (NOT plain — that discards the head).
% ============================================================================
\usepackage{fancyhdr}
% v3.61: \chapter* left a STALE running head — pages of an UNNUMBERED chapter
% carried the PREVIOUS chapter's title. report.cls routes \chapter* through
% \@schapter, which never calls \chaptermark, so \leftmark keeps whatever the last
% numbered chapter left. Measured: with \chapter{פרק ממוספר} then
% \chapter*{מבוא לא ממוספר}, page 4 (inside the starred chapter) printed
% "פרק ממוספר". Invisible on the chapter's FIRST page (\thispagestyle{plain}), so it
% only shows from the second page on — the pages an author is least likely to check.
% An unnumbered front-matter chapter (\chapter*{מבוא}, \chapter*{הקדמה}) is everyday
% Hebrew-academic structure, and the corpus already carries HAND-WRITTEN \markboth
% lines right after \chapter* — the exact per-document duplication \lakithead exists
% to retire. Patched INSIDE \lakithead so the macro stays inert when unused; a
% document's own later \markboth still wins, so those hand-patches do not regress.
\makeatletter
% <hllpdf:starred-heading-scope>
\newif\iflakin@starmarkpatched
\newcommand{\lakistarmark}{%
  \iflakin@starmarkpatched\else
  % Dispatch on the actual class counter, not hyperref's shared internals.
  \ifdefined\c@chapter
    \let\HEBorigschapter\@schapter
    % Record only after the native heading has reached its page. A mark
    % before the native routine can advertise the next heading too early.
    \def\@schapter##1{\HEBorigschapter{##1}\chaptermark{##1}}%
  \else
    % @ssect receives layout arguments, NOT the section name or level.
    % Capture the name in @startsection without introducing a TeX group:
    % grouping sectioning would break the native after-heading state.
    \def\HEB@sectionkind{}%
    \def\HEB@mainsection{section}%
    \let\HEBorigstartsection\@startsection
    \def\@startsection##1##2##3##4##5##6{%
      \def\HEB@sectionkind{##1}%
      \HEBorigstartsection{##1}{##2}{##3}{##4}{##5}{##6}}%
    \let\HEBorigssect\@ssect
    \def\@ssect##1##2##3##4##5{%
      \HEBorigssect{##1}{##2}{##3}{##4}{##5}%
      % Lower starred headings must not replace the article's section mark.
      \ifx\HEB@sectionkind\HEB@mainsection
        \sectionmark{##5}%
      \fi
      \ignorespaces}%
  \fi
  \global\lakin@starmarkpatchedtrue
  \fi}
% </hllpdf:starred-heading-scope>
% <hllpdf:running-language>
% Choose the wrapper when the mark is written, not when a later page ships.
% English punctuation must remain LTR even after returning to Hebrew prose.
\newcommand{\HEB@runningmark}[1]{%
  \iflanguage{english}{\markboth{\en{#1}}{}}{\markboth{\h{#1}}{}}}
% </hllpdf:running-language>
\makeatother
% v4.2.1 (NAV): every collection already names itself in \lakithead — feed the
% PDF /Title from the same declaration, so the navigation gate's
% PDF_REQUIRE_TITLE has a universal source and no fleet member ships
% title-less by missing a second declaration site. Measured before this: 14
% of 16 collections shipped /Title-empty while every one carried a perfectly
% good \lakithead string (plain Hebrew in all 16 — no math, no formatting, so
% the PDF-string path is clean). An explicit title always wins: set it with
% \lakitpdftitle{...} anywhere before \lakithead, or with a plain
% \hypersetup{pdftitle=...} — the fallback detects that via hyperref's
% \@pdftitle, probe-verified \@empty-equal until the first explicit set.
\makeatletter
\newif\iflakit@pdftitleset
\newcommand{\lakitpdftitle}[1]{%
  \hypersetup{pdftitle={#1}}\global\lakit@pdftitlesettrue}
\newcommand{\lakithead}[1]{%
  \iflakit@pdftitleset\else
    \ifx\@pdftitle\@empty
      \hypersetup{pdftitle={#1}}\global\lakit@pdftitlesettrue
    \fi
  \fi
  \pagestyle{fancy}\fancyhf{}%
  \lakistarmark

  \fancyhead[L]{\small\textcolor{cGray}{\leftmark}}%
  % <hllpdf:fixed-head-language>
  % The fixed document title has its OWN language snapshot, independent of
  % the current section/chapter mark. Choose the branch NOW, not at shipout:
  % the page may be Hebrew by the time an English title is finally printed.
  \iflanguage{english}{%
    \fancyhead[R]{\small\textcolor{cGray}{\en{#1}}}%
  }{%
    \fancyhead[R]{\small\textcolor{cGray}{\h{#1}}}%
  }%
  % </hllpdf:fixed-head-language>
  \fancyfoot[C]{\thepage}%
  \renewcommand{\headrulewidth}{0.4pt}%
  \ifdefined\chaptermark
    \renewcommand{\chaptermark}[1]{\HEB@runningmark{##1}}%
  \else
    \renewcommand{\sectionmark}[1]{\HEB@runningmark{##1}}%
  \fi}
\makeatother

% v3.39.1: path-independent byte-reproducibility. LuaTeX derives the trailer
% /ID from the OUTPUT PATH, so two clones of the same source at different
% directories produced different sha256 (measured: the ONLY delta was the
% 62-byte /ID line). Bit 512 suppresses exactly that line — measured: Creator/
% Producer/CreationDate/ModDate and all content untouched. /ID is optional
% per ISO 32000 (required only for ENCRYPTED PDFs and PDF/A — neither is a
% target here). If a future workflow needs an ID (PDF/A, Acrobat revision
% matching), delete this line and it returns.
\pdfvariable suppressoptionalinfo \numexpr 512\relax

% v3.75: BOLD MATH IN BOLD HEADINGS. Heading text is \bfseries but math inside
% it stayed regular-weight CM — a hairline $V_T$ inside a heavy 24.8pt chapter
% title (measured: CMMI12 fl=6 next to FrankRuhlLibre-Bold fl=20; 155 live
% heading-with-math sites fleet-wide). \boldmath rides along wherever the class
% sets \bfseries for a heading; A/B measured: $V_T$ -> CMMIB10 fl=22, 0 missing
% characters (the bold CM faces ship with this TeX Live). Guarded per-macro so
% the article class (no \@makechapterhead) and future class changes stay safe.
% v4.0.1 (adoption fix): the SEARCH patterns above carry no parameter token.
% v4.0.0 made them specific by including the class's own #1/#2, and measured on
% TeX Live 2023 that never matches: a \patchcmd search pattern containing a
% parameter token cannot match while the call sits inside \@ifundefined{...}{}{...}
% (probed both #1 and ##1 -> FAIL; the same pattern at top level succeeds). Every
% one of these patches therefore failed, and because v4.0.0 also replaced the
% empty failure branch with \PackageError, the failure became FATAL: all four
% bundled example assets stopped building ("Cannot patch chapter-title math
% weight -> Fatal error, no output PDF"). The \PackageError branches are KEPT —
% a silent no-op here is what let v3.75's fix rot unnoticed — only the parameter
% tokens are dropped, so the patches match and a real class change still errors.
\makeatletter
\@ifundefined{@makechapterhead}{}{%
  \patchcmd{\@makechapterhead}{\huge\bfseries \@chapapp}
    {\huge\bfseries\boldmath \@chapapp}{}
    {\PackageError{hebrew-lualatex-pdf}{Cannot patch chapter-number math weight}{Unsupported document class}}%
  \patchcmd{\@makechapterhead}{\Huge \bfseries}
    {\Huge \bfseries\boldmath}{}
    {\PackageError{hebrew-lualatex-pdf}{Cannot patch chapter-title math weight}{Unsupported document class}}%
  \patchcmd{\@makeschapterhead}{\Huge \bfseries}
    {\Huge \bfseries\boldmath}{}
    {\PackageError{hebrew-lualatex-pdf}{Cannot patch starred chapter math weight}{Unsupported document class}}}
% v4.0.1: \part sits under a THREE-DEEP wrapper chain, so the class-based anchor
% is not in \@part at all. Measured with \show: hyperref rewrites \@part to
% `...\phantomsection\fi \H@old@part`, and nameref rewrites THAT to
% `\NR@gettitle{#1}\NR@part[{#1}]{#2}` — the class body ends up in \NR@part.
% v4.0.0 patched \@part and made failure fatal, which is why all seven bundled
% examples stopped building. \part is NOT a dead heading level here: 188 live uses
% across the fleet, so the bold-math fix is worth getting right rather than
% dropping. We walk the chain innermost-first and stop at the first hit. Failure is
% a WARNING, not an error: the chain is a property of hyperref/nameref versions, so
% a future package release must degrade to regular-weight part math — a cosmetic
% loss — instead of refusing to build. Chapter and section keep \PackageError,
% because those targets are in the CLASS and were verified to patch here.
\newif\iflakin@partnum \newif\iflakin@parttitle \newif\iflakin@sparttitle
\def\lakin@partnum#1#2{\iflakin@partnum\else
  \patchcmd{#1}{#2\bfseries \partname}{#2\bfseries\boldmath \partname}
    {\lakin@partnumtrue}{}\fi}
\def\lakin@parttitle#1#2{\iflakin@parttitle\else
  \patchcmd{#1}{#2 \bfseries}{#2 \bfseries\boldmath}{\lakin@parttitletrue}{}\fi}
\def\lakin@sparttitle#1#2{\iflakin@sparttitle\else
  \patchcmd{#1}{#2 \bfseries}{#2 \bfseries\boldmath}{\lakin@sparttitletrue}{}\fi}
\@ifundefined{c@chapter}{\def\lakin@pnum{\Large}\def\lakin@ptit{\huge}}%
                        {\def\lakin@pnum{\huge}\def\lakin@ptit{\Huge}}
\@ifundefined{NR@part}{}{\expandafter\lakin@partnum\expandafter\NR@part\expandafter{\lakin@pnum}%
                         \expandafter\lakin@parttitle\expandafter\NR@part\expandafter{\lakin@ptit}}
\@ifundefined{H@old@part}{}{\expandafter\lakin@partnum\expandafter\H@old@part\expandafter{\lakin@pnum}%
                            \expandafter\lakin@parttitle\expandafter\H@old@part\expandafter{\lakin@ptit}}
\@ifundefined{@part}{}{\expandafter\lakin@partnum\expandafter\@part\expandafter{\lakin@pnum}%
                       \expandafter\lakin@parttitle\expandafter\@part\expandafter{\lakin@ptit}}
\@ifundefined{NR@spart}{}{\expandafter\lakin@sparttitle\expandafter\NR@spart\expandafter{\lakin@ptit}}
\@ifundefined{H@old@spart}{}{\expandafter\lakin@sparttitle\expandafter\H@old@spart\expandafter{\lakin@ptit}}
\@ifundefined{@spart}{}{\expandafter\lakin@sparttitle\expandafter\@spart\expandafter{\lakin@ptit}}
\iflakin@partnum\else\PackageWarning{hebrew-lualatex-pdf}{Part-number math stays
  regular weight: no known \noexpand\part wrapper matched (hyperref/nameref
  version change). Cosmetic only}\fi
\iflakin@parttitle\else\PackageWarning{hebrew-lualatex-pdf}{Part-title math stays
  regular weight: no known \noexpand\part wrapper matched. Cosmetic only}\fi
% v4.12.0: THE PART NUMBER IN THE TOC WAS NOT DIRECTIONALLY SEPARATED FROM THE
% TITLE. report.cls/article.cls compose the part entry as ONE text stream --
% \addcontentsline{toc}{part}{\thepart\hspace{1em}#1} -- while EVERY other level
% goes through \numberline, which puts the number in an \hbox. A box is a bidi
% OBJECT and cannot join a neighbouring run; a bare "IV" can. So in the RTL TOC
% paragraph a Latin number and a Latin-INITIAL title merge into a SINGLE LTR run
% and are laid out left-to-right as a unit: the number lands on the wrong side of
% the row, and a mixed title is scrambled outright. Measured on \part{MCP ...}:
% the printed order was "<tail> V MCP" -- title fragment, number, title fragment.
% The chapter/section rows one line below, with the SAME Latin-initial titles,
% were correct on the same page. That contrast is the whole diagnosis.
% \en{} IS NOT A WORKAROUND, measured both ways on one fixture: it fixes a
% PURE-Latin title (\part{\en{Hooks}} renders correctly) and does NOT fix a MIXED
% one (\part{\en{MCP} <hebrew>} stays scrambled). babel's language switch changes
% font and hyphenation, not the bidi run structure -- so the defect is not
% reachable from the source side at all, which is why it belongs here.
% Boxing the number is the entire fix and is a STRICT NO-OP where the bug never
% bit: a Hebrew-titled document compiles BYTE-IDENTICAL before and after (md5,
% 3 clean runs), because there the neutral glue already sat between an L run and
% an R run. The PDF outline is byte-identical too -- hyperref's \pdfstringdef
% neutralises \mbox, so bookmark strings keep their part numbers.
% \numberline{\thepart} would also fix the direction, but it re-geometries every
% part row in the fleet and DROPS the number from the part bookmark (hyperref
% strips \numberline). \mbox changes nothing that was already right.
% The search pattern deliberately carries no # token: etoolbox cannot match a
% parameter token in a patch pattern (measured -- all three chain members
% reported "not patchable" while the pattern included #1). "\thepart \hspace
% {1em}" is brace-balanced and occurs exactly once in the wrapper body.
\newif\iflakin@parttocnum
\def\lakin@parttocnum#1{\iflakin@parttocnum\else
  \patchcmd{#1}{\thepart \hspace {1em}}{\mbox {\thepart }\hspace {1em}}
    {\lakin@parttocnumtrue}{}\fi}
\@ifundefined{NR@part}{}{\lakin@parttocnum\NR@part}
\@ifundefined{H@old@part}{}{\lakin@parttocnum\H@old@part}
\@ifundefined{@part}{}{\lakin@parttocnum\@part}
\iflakin@parttocnum\else\PackageWarning{hebrew-lualatex-pdf}{Part TOC number is
  not boxed: no known \noexpand\part wrapper matched. A part title that BEGINS
  with Latin will print its number on the wrong side of its TOC row}\fi
\patchcmd{\section}{\bfseries}{\bfseries\boldmath}{}
  {\PackageError{hebrew-lualatex-pdf}{Cannot patch section math weight}{Unsupported document class}}
\patchcmd{\subsection}{\bfseries}{\bfseries\boldmath}{}
  {\PackageError{hebrew-lualatex-pdf}{Cannot patch subsection math weight}{Unsupported document class}}
\patchcmd{\subsubsection}{\bfseries}{\bfseries\boldmath}{}
  {\PackageError{hebrew-lualatex-pdf}{Cannot patch subsubsection math weight}{Unsupported document class}}
\patchcmd{\paragraph}{\bfseries}{\bfseries\boldmath}{}{}
\patchcmd{\subparagraph}{\bfseries}{\bfseries\boldmath}{}{}
\@ifundefined{l@chapter}{}{%
  \patchcmd{\l@chapter}{\bfseries}{\bfseries\boldmath}{}{}}
\makeatother
\setcounter{tocdepth}{\lakittocdepth}
% <hllpdf:measured-toc>
% v4.16.0: Measure actual number labels, not a guessed digit count. Read only
% canonical contentsline records; never execute titles or addtocontents code
% during the measuring pass. Widths are local to this individual list and are
% available before its first row: no per-row drift and no extra aux pass.
\makeatletter
\newread\hllpdf@tocstream
\newdimen\hllpdf@toc@chapter
\newdimen\hllpdf@toc@section
\newdimen\hllpdf@toc@subsection
\newdimen\hllpdf@toc@subsubsection
\newdimen\hllpdf@toc@paragraph
\newdimen\hllpdf@toc@subparagraph
\newdimen\hllpdf@toc@figure
\newdimen\hllpdf@toc@table
\newdimen\hllpdf@tocpage
\newdimen\hllpdf@toc@sectionindent
\newdimen\hllpdf@toc@subsectionindent
\newdimen\hllpdf@toc@subsubsectionindent
\newdimen\hllpdf@toc@paragraphindent
\newdimen\hllpdf@toc@subparagraphindent
\def\hllpdf@tocdefaults{%
  \hllpdf@toc@chapter=2em
  \hllpdf@toc@section=3.2em
  \hllpdf@toc@subsection=3.9em
  \@ifundefined{c@chapter}{%
    \hllpdf@toc@subsubsection=3.2em
    \hllpdf@toc@paragraph=4.1em
    \hllpdf@toc@subparagraph=5em
  }{%
    \hllpdf@toc@subsubsection=4.1em
    \hllpdf@toc@paragraph=5em
    \hllpdf@toc@subparagraph=6em
  }%
  \hllpdf@toc@figure=2.3em
  \hllpdf@toc@table=2.3em
  \hllpdf@tocpage=\@pnumwidth\relax
}
\def\hllpdf@tocindents{%
  % Preserve the established shallow layout when its slots already fit.
  \hllpdf@toc@sectionindent=1.5em
  \@tempdima=\hllpdf@toc@chapter \advance\@tempdima by -2em
  \advance\hllpdf@toc@sectionindent by \@tempdima
  \hllpdf@toc@subsectionindent=\dimexpr
    \hllpdf@toc@sectionindent+\hllpdf@toc@section\relax
  \hllpdf@toc@subsubsectionindent=\dimexpr
    \hllpdf@toc@subsectionindent+\hllpdf@toc@subsection\relax
  \hllpdf@toc@paragraphindent=\dimexpr
    \hllpdf@toc@subsubsectionindent+\hllpdf@toc@subsubsection\relax
  \hllpdf@toc@subparagraphindent=\dimexpr
    \hllpdf@toc@paragraphindent+\hllpdf@toc@paragraph\relax
}
\def\hllpdf@tocfont@chapter{\bfseries\boldmath}
\def\hllpdf@tocfont@part{\large\bfseries\boldmath}
\def\hllpdf@tocentryfont{%
  \ifcsname hllpdf@tocfont@\hllpdf@toclevel\endcsname
    \csname hllpdf@tocfont@\hllpdf@toclevel\endcsname
  \fi}
\long\def\hllpdf@tocdiscard#1\hllpdf@tocstop{}
\def\hllpdf@tocscanrecord{%
  \@ifnextchar\contentsline{\hllpdf@tocscanline}{%
    \@ifnextchar\babel@toc{\hllpdf@tocscanlanguage}{\hllpdf@tocdiscard}}}
\long\def\hllpdf@tocscanlanguage\babel@toc#1#2#3\hllpdf@tocstop{%
  % Babel records are declarations, not arbitrary code. Track the language
  % without executing the original directive or changing the surrounding list.
  \def\hllpdf@toclanguage{#1}}
\long\def\hllpdf@tocscanline\contentsline#1#2#3#4#5\hllpdf@tocstop{%
  \def\hllpdf@toclevel{#1}%
  \ifcsname hllpdf@toc@#1\endcsname
    \hllpdf@tocscantitle#2\hllpdf@tocstop
  \fi
  % The page-number column needs its own measurement (including bold chapter
  % rows). Its extra clearance separates the page label from dot leaders.
  \setbox\@tempboxa=\hbox{\foreignlanguage{\hllpdf@toclanguage}{\hllpdf@tocentryfont #3\kern.3em}}%
  \ifdim\wd\@tempboxa>\hllpdf@tocpage
    \hllpdf@tocpage=\wd\@tempboxa
  \fi}
\def\hllpdf@tocscantitle{%
  \@ifnextchar\numberline{\hllpdf@tocscanlabel}{\hllpdf@tocdiscard}}
\long\def\hllpdf@tocscanlabel\numberline#1#2\hllpdf@tocstop{%
  \setbox\@tempboxa=\hbox{\foreignlanguage{\hllpdf@toclanguage}{\hllpdf@tocentryfont #1\kern.5em}}%
  \ifdim\wd\@tempboxa>\csname hllpdf@toc@\hllpdf@toclevel\endcsname
    \csname hllpdf@toc@\hllpdf@toclevel\endcsname=\wd\@tempboxa
  \fi}
\def\hllpdf@tocread{%
  \ifeof\hllpdf@tocstream\else
    \read\hllpdf@tocstream to\hllpdf@tocrecord
    \expandafter\hllpdf@tocscanrecord\hllpdf@tocrecord\hllpdf@tocstop
    \expandafter\hllpdf@tocread
  \fi}
\def\hllpdf@tocmeasure#1{%
  \makeatletter
  \edef\hllpdf@toclanguage{\languagename}%
  \hllpdf@tocdefaults
  \openin\hllpdf@tocstream=\jobname.#1\relax
  \hllpdf@tocread
  \closein\hllpdf@tocstream
  \hllpdf@tocindents
  \edef\@pnumwidth{\the\hllpdf@tocpage}%
  \@tempdima=\dimexpr\hllpdf@tocpage+1em\relax
  \ifdim\@tempdima>\@tocrmarg\relax
    \edef\@tocrmarg{\the\@tempdima}%
  \fi}
\let\hllpdf@originalstarttoc\@starttoc
\def\@starttoc#1{%
  \begingroup
    \def\hllpdf@tocext{#1}%
    \def\hllpdf@tocwanted{toc}%
    \ifx\hllpdf@tocext\hllpdf@tocwanted \hllpdf@tocmeasure{#1}\else
      \def\hllpdf@tocwanted{lof}%
      \ifx\hllpdf@tocext\hllpdf@tocwanted \hllpdf@tocmeasure{#1}\else
        \def\hllpdf@tocwanted{lot}%
        \ifx\hllpdf@tocext\hllpdf@tocwanted \hllpdf@tocmeasure{#1}\fi
      \fi
    \fi
    \hllpdf@originalstarttoc{#1}%
  \endgroup}
\hllpdf@tocdefaults
\hllpdf@tocindents
\@ifundefined{l@chapter}{}{%
  \patchcmd{\l@chapter}{\setlength\@tempdima{1.5em}}%
    {\setlength\@tempdima{\hllpdf@toc@chapter}}{}%
    {\PackageError{hebrew-lualatex-pdf}{Cannot size chapter TOC numbers}%
       {Unsupported contents layout: the chapter number slot was not found}}}
\renewcommand*\l@section{\@dottedtocline{1}%
  {\hllpdf@toc@sectionindent}{\hllpdf@toc@section}}
\renewcommand*\l@subsection{\@dottedtocline{2}%
  {\hllpdf@toc@subsectionindent}{\hllpdf@toc@subsection}}
\renewcommand*\l@subsubsection{\@dottedtocline{3}%
  {\hllpdf@toc@subsubsectionindent}{\hllpdf@toc@subsubsection}}
\renewcommand*\l@paragraph{\@dottedtocline{4}%
  {\hllpdf@toc@paragraphindent}{\hllpdf@toc@paragraph}}
\renewcommand*\l@subparagraph{\@dottedtocline{5}%
  {\hllpdf@toc@subparagraphindent}{\hllpdf@toc@subparagraph}}
\renewcommand*\l@figure{\@dottedtocline{1}{1.5em}{\hllpdf@toc@figure}}
\renewcommand*\l@table{\@dottedtocline{1}{1.5em}{\hllpdf@toc@table}}
\makeatother
% </hllpdf:measured-toc>
\setcounter{secnumdepth}{\lakitsecnumdepth}
% Public preamble-only caption-numbering policy. The historical within-major
% default is untouched unless this command is called. "continuous" removes the
% reset edge as well as the printed prefix, so Figure 1/2 and Table 1/2 stay
% unique across sections/chapters. Returning to "within-major" restores the
% original class-aware displays and native destination conventions. The policy
% never replaces \figcap/\tabcap or their measured link-placement mechanism.
\makeatletter
\let\lakin@majorfig\thefig
\let\lakin@majortbl\thetbl
\let\lakin@majorHfig\theHfig
\let\lakin@majorHtbl\theHtbl
\newcommand{\lakitcaptionnumbering}[1]{%
  \edef\lakin@numbering{#1}%
  \def\lakin@continuous{continuous}\def\lakin@withinmajor{within-major}%
  \ifx\lakin@numbering\lakin@continuous
    \@ifundefined{c@chapter}{%
      \@removefromreset{fig}{section}\@removefromreset{tbl}{section}%
    }{\@removefromreset{fig}{chapter}\@removefromreset{tbl}{chapter}}%
    \renewcommand{\thefig}{\arabic{fig}}\renewcommand{\thetbl}{\arabic{tbl}}%
    \def\theHfig{\arabic{fig}}\def\theHtbl{\arabic{tbl}}%
  \else\ifx\lakin@numbering\lakin@withinmajor
    \@ifundefined{c@chapter}{%
      \@removefromreset{fig}{section}\@removefromreset{tbl}{section}%
      \@addtoreset{fig}{section}\@addtoreset{tbl}{section}%
    }{%
      \@removefromreset{fig}{chapter}\@removefromreset{tbl}{chapter}%
      \@addtoreset{fig}{chapter}\@addtoreset{tbl}{chapter}}%
    \let\thefig\lakin@majorfig\let\thetbl\lakin@majortbl
    \let\theHfig\lakin@majorHfig\let\theHtbl\lakin@majorHtbl
  \else
    \PackageError{hebrew-english-latex-studio}{Unknown caption-numbering policy '#1'}%
      {Use continuous or within-major before begin{document}.}%
  \fi\fi}
\@onlypreamble\lakitcaptionnumbering
\makeatother
% Architecture readiness sentinel. Optional runtime layers require the complete
% core, so the marker is defined only after every mechanism/baseline definition.
% <hllpdf:architecture-ready>
\def\hllpdfcoreloaded{1}
% </hllpdf:architecture-ready>

% v4.22: registered file resources, not an assertion about embedded subset bytes.
% The log framing avoids long-path wrapping and preserves UTF-8 as hex.
% <hllpdf:font-inventory>
\directlua{
  local function inventory()
    local names = {}
    local unnamed = 0
    for id, item in font.each() do
      local path = item.filename
      if path and not (path == "") then names[path] = true else unnamed = unnamed + 1 end
    end
    local sorted = {}
    for path in pairs(names) do table.insert(sorted,path) end
    table.sort(sorted)
    texio.write_nl("log","HLLFONT|BEGIN|1")
    for index,path in ipairs(sorted) do
      local hex = path:gsub(".",function(c) return string.format(string.char(37) .. "02x",string.byte(c)) end)
      local count = math.ceil(string.len(hex)/48)
      texio.write_nl("log","HLLFONT|FILE|" .. index .. "|" .. count)
      for offset=1,string.len(hex),48 do texio.write_nl("log","HLLFONT|HEX|" .. hex:sub(offset,offset+47)) end
    end
    texio.write_nl("log","HLLFONT|END|" .. rawlen(sorted) .. "|" .. unnamed)
    texio.write_nl("log","")
  end
  luatexbase.add_to_callback("finish_pdffile",inventory,"hllpdf.runtime.fonts")
}
% </hllpdf:font-inventory>
     % (or paste these lines at the top)
%   \begin{document}  ...  \end{document}
%
% Compile with:  lualatex -interaction=nonstopmode doc.tex
%   Repeat until two consecutive PDFs are byte-identical (scripts/build.sh
%   does this) — a TOC document converges only on pass 3+; after two passes
%   the TOC page numbers are STALE (measured; see SKILL.md step 3).
% Font binaries are external dependencies, pinned by assets/fonts/SHA256SUMS.
% scripts/setup_fonts.sh imports a verified local installation or prior skill.
%
% PUBLIC DOCUMENT-SHELL / LAYOUT CONTRACT (introduced v4.3.1; architecture-owned v4.7.1). These hooks may be
% defined BEFORE % ============================================================================
% preamble.tex — tested LuaLaTeX preamble for Hebrew + math + figures PDFs (v4.32.0)
% Battle-tested: compiled an 80-page Hebrew quantum-computing summary AND a
% 197-page Hebrew computer-architecture anthology (TikZ diagrams, code listings,
% colored boxes) — both with 0 missing characters.
%
% Usage in your document:
%   % ============================================================================
% preamble.tex — tested LuaLaTeX preamble for Hebrew + math + figures PDFs (v4.32.0)
% Battle-tested: compiled an 80-page Hebrew quantum-computing summary AND a
% 197-page Hebrew computer-architecture anthology (TikZ diagrams, code listings,
% colored boxes) — both with 0 missing characters.
%
% Usage in your document:
%   \input{preamble.tex}     % (or paste these lines at the top)
%   \begin{document}  ...  \end{document}
%
% Compile with:  lualatex -interaction=nonstopmode doc.tex
%   Repeat until two consecutive PDFs are byte-identical (scripts/build.sh
%   does this) — a TOC document converges only on pass 3+; after two passes
%   the TOC page numbers are STALE (measured; see SKILL.md step 3).
% Font binaries are external dependencies, pinned by assets/fonts/SHA256SUMS.
% scripts/setup_fonts.sh imports a verified local installation or prior skill.
%
% PUBLIC DOCUMENT-SHELL / LAYOUT CONTRACT (introduced v4.3.1; architecture-owned v4.7.1). These hooks may be
% defined BEFORE \input{preamble.tex}; if omitted, the historical canonical
% defaults below are exact.  They deliberately control only the document shell
% and page layout.  RTL/math semantics stay in this preamble, while visual
% presentation themes (for example clean_theme.tex) load AFTER it.
%
%   \def\lakitdocumentclass{article}
%   \def\lakitclassoptions{11pt,a4paper}
%   \def\lakitgeometryoptions{margin=1.8cm,headheight=14pt}
%   \def\lakitlinespread{1.15}
%   \def\lakitparagraphstyle{block} % letters/briefs; indented for continuous prose
%   \def\lakittocdepth{1}
%   \def\lakitsecnumdepth{2}
%   \input{preamble.tex}
%
% Supported class contract: report or article.  Alternate paper sizes must also
% be declared to build.sh through EXPECTED_PAPER; the PDF gate never infers them.
% Do NOT use these hooks as a theme mechanism: theme overlays must not own page
% geometry, class selection, or direction semantics. A declared presentation-policy
% overlay may intentionally change numbering visibility after the preamble.
% ============================================================================
% Runtime architecture guard: generated from scripts/architecture_contract.py.
% The readiness sentinel remains delayed until the END of this file.
% <hllpdf:architecture-guard>
\ifdefined\hllpdfcoreloaded
  \errmessage{hllpdf: preamble.tex loaded more than once}
  \expandafter\endinput
\fi
% </hllpdf:architecture-guard>

\ifdefined\lakitdocumentclass\else\def\lakitdocumentclass{report}\fi
\ifdefined\lakitclassoptions\else\def\lakitclassoptions{12pt,a4paper}\fi
\ifdefined\lakitgeometryoptions\else\def\lakitgeometryoptions{margin=2.3cm,headheight=14pt}\fi
\ifdefined\lakitlinespread\else\def\lakitlinespread{1.3}\fi
\ifdefined\lakittocdepth\else\def\lakittocdepth{1}\fi
\ifdefined\lakitsecnumdepth\else\def\lakitsecnumdepth{2}\fi
% Paragraph organization belongs to the document layout, independently of its
% language and theme. Preserve historical class dimensions unless explicitly
% selecting block paragraphs; never suppress indentation with an everypar hack.
\ifdefined\lakitparagraphstyle\else\def\lakitparagraphstyle{indented}\fi
\edef\lakitparagraphstyle{\lakitparagraphstyle}
\def\hllpdfparagraphblock{block}
\def\hllpdfparagraphindented{indented}
\ifx\lakitparagraphstyle\hllpdfparagraphblock\else
  \ifx\lakitparagraphstyle\hllpdfparagraphindented\else
    \errmessage{hllpdf: paragraph style must be block or indented}
  \fi
\fi
% v4.3.1: the public class hook has a deliberately narrow support contract.
% Before this guard, \def\lakitdocumentclass{book} silently compiled even though
% every heading/TOC regression test targets report/article; an unsupported shell
% could therefore look "supported" until a subtle class-specific defect shipped.
% Fail before \documentclass rather than letting an undocumented combination enter
% the mechanism core. Plain-TeX comparison is used because LaTeX packages are not
% available yet at this point.
\edef\lakitvalidatedclass{\lakitdocumentclass}
\def\lakitreportclass{report}
\def\lakitarticleclass{article}
\ifx\lakitvalidatedclass\lakitreportclass
\else\ifx\lakitvalidatedclass\lakitarticleclass
\else
  \errmessage{hllpdf: unsupported class \lakitdocumentclass; use report or article}
\fi\fi
\let\lakitvalidatedclass\relax
\let\lakitreportclass\relax
\let\lakitarticleclass\relax
% Language intent is independent of document class, theme and content domain.
% The two dominant-language profiles differ in authoring intent, not direction.
% A balanced bilingual document must explicitly choose a navigation/base locale.
% <hllpdf:language-profile>
\providecommand{\lakitlanguageprofile}{hebrew}
\edef\lakitlanguageprofile{\lakitlanguageprofile}
\def\hllpdfhebrew{hebrew}\def\hllpdfenglish{english}
\def\hllpdfhebrewmixed{hebrew-with-english}
\def\hllpdfenglishmixed{english-with-hebrew}
\def\hllpdfbilingual{bilingual}
\let\hllpdfexpectedlanguage\relax
\ifx\lakitlanguageprofile\hllpdfhebrew
  \let\hllpdfexpectedlanguage\hllpdfhebrew
\else\ifx\lakitlanguageprofile\hllpdfhebrewmixed
  \let\hllpdfexpectedlanguage\hllpdfhebrew
\else\ifx\lakitlanguageprofile\hllpdfenglish
  \let\hllpdfexpectedlanguage\hllpdfenglish
\else\ifx\lakitlanguageprofile\hllpdfenglishmixed
  \let\hllpdfexpectedlanguage\hllpdfenglish
\else\ifx\lakitlanguageprofile\hllpdfbilingual
  \ifdefined\lakitmainlanguage\else
    \errmessage{hllpdf: bilingual profile requires explicit lakitmainlanguage}
  \fi
\else
  \errmessage{hllpdf: unsupported language profile \lakitlanguageprofile}
\fi\fi\fi\fi\fi
\ifdefined\lakitmainlanguage
  \edef\lakitmainlanguage{\lakitmainlanguage}
\else
  \let\lakitmainlanguage\hllpdfexpectedlanguage
\fi
\ifx\lakitmainlanguage\hllpdfhebrew\else
  \ifx\lakitmainlanguage\hllpdfenglish\else
    \errmessage{hllpdf: main language must be hebrew or english}
  \fi
\fi
\ifx\hllpdfexpectedlanguage\relax\else
  \ifx\lakitmainlanguage\hllpdfexpectedlanguage\else
    \errmessage{hllpdf: main language conflicts with language profile}
  \fi
\fi
\typeout{HLLLANG|\lakitlanguageprofile|\lakitmainlanguage}
% </hllpdf:language-profile>
\documentclass[\lakitclassoptions]{\lakitdocumentclass}
\usepackage[no-math]{fontspec}        % (1) ★ no-math — THE critical line (see recipe.md)
% layout=lists asks Babel to adapt list parshape when paragraph and page
% directions differ. Without it, minority-language list bodies lose their
% indentation and labels hang outside the text block, despite a clean log.
% The upstream mechanism preserves ordinary same-direction list geometry.
\usepackage[bidi=basic,layout=lists]{babel}
\ifx\lakitmainlanguage\hllpdfenglish
  \babelprovide[import]{hebrew}
  \babelprovide[import,main]{english}
\else
  \babelprovide[import,main]{hebrew}
  \babelprovide[import]{english}
\fi
% v4.22: bind canonical families and fallbacks to the selected installation.
% Lookup by basename or family alone is not proof of the file actually used.
% <hllpdf:font-binding>
\directlua{
  local dir = os.getenv("HEBREW_LATEX_FONT_DIR")
  if not dir or dir == "" then
    local data = os.getenv("XDG_DATA_HOME")
    if not data or data == "" or not (data:sub(1,1) == "/") then
      data = (os.getenv("HOME") or "") .. "/.local/share"
    end
    dir = data .. "/fonts/hebrew-lualatex-pdf"
  end
  if not (dir:sub(1,1) == "/") then dir = lfs.currentdir() .. "/" .. dir end
  if dir:find("[",1,true) or dir:find("]",1,true) or dir:find(string.char(0),1,true)
     or dir:find(string.char(10),1,true) or dir:find(string.char(13),1,true) then
    tex.error("hllpdf: unsupported character in font directory")
  end
  if not (lfs.attributes(dir,"mode") == "directory") then
    tex.error("hllpdf: font directory is missing")
  end
  hllpdf_font_dir = dir:gsub("/+$", "") .. "/"
  tex.sprint(string.char(92) .. "def" .. string.char(92) .. "hllpdffontdir{")
  tex.sprint(-2,hllpdf_font_dir)
  tex.sprint("}")
  luaotfload.add_fallback("hebfb",{"[" .. hllpdf_font_dir .. "DejaVuSerif.ttf]:mode=node;"})
  luaotfload.add_fallback("hebfbbf",{"[" .. hllpdf_font_dir .. "DejaVuSerif-Bold.ttf]:mode=node;"})
}
% </hllpdf:font-binding>
% v3.22: Frank Ruhl Libre ships NO italic face. Instead of silently-upright
% \emph (v<=3.20) or bold-as-emphasis (v3.21), synthesize a slanted face with
% FakeSlant — real, VISIBLE italic emphasis with correct \emph nesting
% semantics (inner \emph toggles back upright). Verified empirically at
% 300dpi: slant clearly visible, 0 errors, 0 missing glyphs.
\babelfont{rm}[Path={\hllpdffontdir},RawFeature={fallback=hebfb},
  BoldFont={FrankRuhlLibre-Bold.ttf},
  BoldFeatures={RawFeature={fallback=hebfbbf}},   % v3.56 — see hebfbbf above
  ItalicFont={FrankRuhlLibre-Regular.ttf}, ItalicFeatures={FakeSlant=0.22},
  % Bold-italic uses the weight-matched fallback and a synthetic slant.
  BoldItalicFont={FrankRuhlLibre-Bold.ttf}, BoldItalicFeatures={RawFeature={fallback=hebfbbf},FakeSlant=0.22}
  ]{FrankRuhlLibre-Regular.ttf}
% Sans fallbacks preserve family and weight for the few missing codepoints.
\directlua{luaotfload.add_fallback("hebfbsans",{"[" .. hllpdf_font_dir .. "DejaVuSans.ttf]:mode=node;"})}
\directlua{luaotfload.add_fallback("hebfbsansbf",{"[" .. hllpdf_font_dir .. "DejaVuSans-Bold.ttf]:mode=node;"})}
\babelfont{sf}[Path={\hllpdffontdir},RawFeature={fallback=hebfbsans},
  BoldFont={Heebo-Bold.ttf},
  BoldFeatures={RawFeature={fallback=hebfbsansbf}}]{Heebo-Regular.ttf}
% v3.24: Scale=MatchLowercase — DejaVu Sans Mono has an outsized x-height, so
% unscaled inline code visually SHOUTS next to the delicate Frank Ruhl body
% (measured on a 197-page build: scaled, a code atom sits in the line at the
% body's weight, and the document tightened by ~5%). Applies to listings too.
\babelfont{tt}[Path={\hllpdffontdir},RawFeature={fallback=hebfb},Scale=MatchLowercase,
  BoldFont={DejaVuSansMono.ttf},
  BoldFeatures={RawFeature={fallback=hebfb},FakeBold=1.5},
  ItalicFont={DejaVuSansMono.ttf},
  ItalicFeatures={RawFeature={fallback=hebfb},FakeSlant=0.2},
  BoldItalicFont={DejaVuSansMono.ttf},
  BoldItalicFeatures={RawFeature={fallback=hebfb},FakeBold=1.5,FakeSlant=0.2}
  ]{DejaVuSansMono.ttf}

% Explicit multilingual/English profiles use real Latin faces and italics.
% The legacy/default Hebrew profile retains its established measured typography.
% Latin Modern is provided by the TeX distribution, not the pinned font bundle;
% recorder/runtime evidence reports it as an additional toolchain resource.
\ifx\lakitlanguageprofile\hllpdfhebrew\else
  \babelfont[english]{rm}[Ligatures=TeX,RawFeature={fallback=hebfb},
    BoldFeatures={RawFeature={fallback=hebfbbf}},
    BoldItalicFeatures={RawFeature={fallback=hebfbbf}}]{Latin Modern Roman}
  \babelfont[english]{sf}[Ligatures=TeX,RawFeature={fallback=hebfbsans},
    BoldFeatures={RawFeature={fallback=hebfbsansbf}},
    BoldItalicFeatures={RawFeature={fallback=hebfbsansbf}}]{Latin Modern Sans}
\fi

\usepackage{amsmath,amssymb,mathtools}
% amsthm supplies theorem/proof machinery. Environment names remain
% document-specific; the generic proof label is localized below.
\usepackage{amsthm}
% Proof names follow the surrounding prose, including return from an island.
\addto\extrashebrew{\renewcommand{\proofname}{הוכחה}}
\addto\extrasenglish{\renewcommand{\proofname}{Proof}}
\ifx\lakitmainlanguage\hllpdfenglish
  \renewcommand{\proofname}{Proof}
\else
  \renewcommand{\proofname}{הוכחה}
\fi
% v3.24.1: equation tags under babel-bidi. A tag re-emitted by mathtools'
% showonlyrefs inside align inherits the RTL base and its parens MIRROR —
% ")2.1(" instead of "(2.1)" (plain align/equation tags are unaffected;
% verified on a minimal case + a 198-page build). LTR-isolate at the tag-BOX
% level: mathtools' show/hide + numbering logic is untouched. Do NOT override
% \tagform@ instead — that clobbers the showonlyrefs mechanism (measured:
% every tag shows, numbering breaks).
% v3.28.1: the v3.24.1 wrap must apply ONLY to tags of display environments.
% Inline \eqref in Hebrew PROSE also goes through \maketag@@@, and there the
% english-wrap MIRRORS the parens — ")1(" — while the unwrapped original is
% correct (proven on a 2x2 isolation matrix: patch x showonlyrefs; regressed
% every prose \eqref since v3.24.1, caught by the render gate's R4). The two
% call sites are indistinguishable by \ifmmode (tag boxes are text mode), so
% the display environments themselves raise a flag; \eqref in prose sees it
% false and gets the original untouched \maketag@@@.
\usepackage{etoolbox}

% babel bidi=basic (including upstream 26.9) assumes every automatic
% discretionary has a one-glyph replacement. LuaTeX can legally produce one
% with replace=nil after a hyphenated Latin value; Babel then indexes nil and the
% whole build dies. Do the narrow upstream-compatible correction: while Babel
% classifies the list, present only that malformed disc as an ordinary
% discretionary, then restore its subtype. Babel still reorders the ENTIRE box;
% there is no pcall and no successful fallback to an unreordered node list.
% If Babel throws for any other reason, the build remains fatal. The synthetic
% nil-disc probe and the real table gallery in test_all.sh pin both layers.
\AtBeginDocument{\directlua{
  if Babel and Babel.bidi and not Babel.heb_nil_disc_compat then
    Babel.heb_nil_disc_compat = true
    local original = Babel.bidi
    local disc_id = node.id('disc')
    Babel.bidi = function(head, ispar, hdir)
      local hidden = {}
      for item in node.traverse_id(disc_id, head) do
        if item.subtype == 2 and (not item.replace or not item.replace.char) then
          table.insert(hidden, item)
          item.subtype = 0
        end
      end
      local result = original(head, ispar, hdir)
      for _, item in ipairs(hidden) do item.subtype = 2 end
      return result
    end
  end
}}
\makeatletter
\newif\ifHEBdisptag
\let\ORIGmaketag@@@\maketag@@@
\def\maketag@@@#1{\ifHEBdisptag\ORIGmaketag@@@{\foreignlanguage{english}{#1}}\else\ORIGmaketag@@@{#1}\fi}
\@for\HEB@env:=equation,equation*,align,align*,gather,gather*,multline,multline*,flalign,flalign*,alignat,alignat*,eqnarray,eqnarray*\do{%
  \expandafter\AtBeginEnvironment\expandafter{\HEB@env}{\HEBdisptagtrue}}
\makeatother
% English page + Hebrew prose: Babel's equation-only tag flip uses a primitive
% leqno path that shifts the formula to the left edge and the tag into the body
% of an English page. Align's full-width alignment is unaffected. Give only this
% display a physical LTR container and native amsmath eqno/leqno. The engine then
% owns real tag width, collision handling and labels; no proxy tag box hides ink
% from its fit calculation. Babel's saved \bbl@eqnodir still restores Hebrew text
% inside math, and the prose locale/math direction are never switched.
% Close the prior RTL paragraph before normalization. After the environment's
% group restores its directions, close TeX's empty resumed display paragraph and
% use native LaTeX continuation handling: same-paragraph prose has no indent,
% while an authored blank line starts an ordinary indented RTL paragraph.
% Register after Babel's setup; remain local to equation(*) in a main-English
% document's actual RTL prose. Native Hebrew, ordinary English and explicit
% fleqn retain their original paths. No page/display dimensions are assigned.
% <hllpdf:minority-rtl-equation-geometry>
\makeatletter
\newcommand{\HEB@minorityprosecontinuation}{\par\@doendpe}
\newcommand{\HEB@minorityequationgeometry}{%
  \ifx\lakitmainlanguage\hllpdfenglish
    \ifnum\pardirection=\@ne
      \if@fleqn\else
        \par\aftergroup\HEB@minorityprosecontinuation
        \iftagsleft@\let\veqno\@@leqno
        \else\let\veqno\@@eqno\fi
        \bodydir TLT\pardir TLT\textdir TLT\relax
      \fi
    \fi
  \fi}
\AtBeginDocument{%
  \AddToHook{env/equation/begin}{\HEB@minorityequationgeometry}%
  \AddToHook{env/equation*/begin}{\HEB@minorityequationgeometry}}
\makeatother
% </hllpdf:minority-rtl-equation-geometry>
% v3.63: \shortintertext rendered Hebrew MIRROR-REVERSED and flush LEFT — an entire
% commentary line inside an align printed backwards. babel's bidi=basic patches
% \intertext@ only (luababel.def), and mathtools' \shortintertext uses the identical
% \vbox{\normalbaselines … \noindent\ignorespaces#1\par} skeleton but is never
% patched, so its paragraph keeps pardir/textdir LTR: no bidi reordering (logical
% order on the page) and flush-left. Measured against \intertext with the same
% Applying babel's own replacement to \shortintertext@ makes the two identical.
% \h{}/\mbox fix the letter order but leave the line flush LEFT, so they are not
% substitutes. Guarded + \AtBeginDocument because mathtools \let-copies at load time.
% NB \ExplSyntaxOn is REQUIRED: mathtools' internals carry expl3 names, and without
% the expl3 catcodes the `_` is read as a subscript and the patch detonates far away
% with "! Illegal unit of measure (pt inserted)" at the \end{align} (measured).
\makeatletter
\AtBeginDocument{%
  \@ifundefined{shortintertext@}{}{%
    \ExplSyntaxOn
    \bbl@sreplace\shortintertext@{\normalbaselines}%
      {\normalbaselines\ifx\bbl@eqnodir\relax\else\bbl@pardir\@ne\bbl@eqnodir\fi}%
    \ExplSyntaxOff}}
\makeatother
% geometry inherits the canonical A4 sheet from \documentclass; build.sh then
% verifies every rendered MediaBox and CropBox, so a silent Letter fallback is
% impossible to publish as a clean default build.
\usepackage[\lakitgeometryoptions]{geometry}
\linespread{\lakitlinespread}       % layout hook; canonical default remains 1.3
\ifx\lakitparagraphstyle\hllpdfparagraphblock
  \setlength{\parindent}{0pt}
  \setlength{\parskip}{6pt plus 1pt minus 1pt}
\fi
% These dimensions remain ordinary scoped TeX lengths: authors can customize
% them after core input. Language switches inherit them; lists keep their own
% indentation and narrow reference layouts establish their local paragraph gaps.
\typeout{HLLPAR|\lakitparagraphstyle}
\usepackage{xcolor}
\usepackage{mdframed}
\usepackage{listings}                 % code listings (assembly / C / pseudo-code)
\usepackage{tikz}
\usetikzlibrary{arrows.meta,positioning,calc,shapes.geometric,fit,backgrounds}
\usepackage{pgfplots}\pgfplotsset{compat=1.18}
\usepackage{quantikz}                 % quantum circuits (remove if unused)
\usepackage{booktabs}
\usepackage{array}
\usepackage{float}                    % [H] = place table exactly here (reference doc: no drifting)
\usepackage{enumitem}
\usepackage{needspace}

% ============================================================================
% ★★ v3.4 — THE DIRECTION-ISOLATION TRIO (the heart of mixing Hebrew + Latin)
% Mental model:
%   \h{...}    Hebrew (RTL) island      — Hebrew text inside an LTR context
%                                          (TikZ nodes, the english `tikzpic`).
%   \en{...}   English (LTR) isolate    — Latin text / parentheses inside a
%                                          Hebrew (RTL) context. KEEPS PARENS
%                                          AND ORDER CORRECT.
%   \code{...} \en + monospace          — code inside Hebrew / inside a frametitle.
% CRITICAL: never wrap Latin in \h{} (it is for HEBREW). Inside a TikZ node,
% Latin in \h{} renders REVERSED (Datapath->htapataD, CPU->UPC) and parens
% mirror. Isolate every Latin/paren chunk with \en{} or \code{}. See
% references/figures-boxes-listings.md.
% ---- TEXT INSIDE MATH (v3.8) ----------------------------------------------
% Put ANY text that sits inside math in a \text{}. The ONLY thing that is
% invalid in math is a bare \h/\en (= \foreignlanguage) used DIRECTLY, outside
% \text -> "\rmfamily invalid in math mode". Inside \text:
%   Hebrew label           -> \text{...}            (bare Hebrew renders fine;
%                              \h is optional — BUT NOT ALWAYS HARMLESS (v3.67):
%                              if the \h group OPENS with a nested island
%                              ($..$, \en{}, \code{}, \texttt{}, \mbox{}) the
%                              WHOLE group loses bidi and the Hebrew prints in
%                              LOGICAL order = gibberish. Measured:
%                                \text{\h{$f$ קבועה}}  -> reads "העובק"  BROKEN
%                                \text{\h{קבועה $f$}}  -> correct
%                                \text{$f$ קבועה}      -> correct (no \h)
%                              Hebrew first inside \h, or drop the \h. Flagged
%                              H-OPENS-WITH-LTR-GROUP.)
%   plain Latin word        -> \text{...}            (no marker needed)
%   Latin WITH ( ) or [ ]   -> \text{\en{...}}       (\en forces LTR; a plain
%                              \text mirrors the parens in RTL -> "Spec (TNR(" )
%   f(Hebrew/text arg)      -> \text{f}(\text{...})  -- keep the parens at MATH
%                              level, NOT \text{f(arg)} (parens inside the box mirror)
%   code in math            -> \mathtt{...}  (or \texttt{}, math-guarded below)
% NEVER write \en{...} or \h{...} straight into $...$ / \[..\] (outside \text);
% always via \text{}. (\mbox{...} also works but does not scale in scripts.)
% ============================================================================
\newcommand{\h}[1]{\foreignlanguage{hebrew}{#1}}
% v3.42.1 — \en islands do NOT hyphenate (auto OR at explicit hyphens). English hyphenation patterns were
% tearing identifiers and terms ACROSS THE RTL SEAM ("IO Regis-/ters",
% "In-/put", "Band-/width", "Per-/formance", "Multi-/ple" — all shipped on
% real pages; probe-proven: "ALUS-/rcB", "Mem-/ToReg"). An island is a label
% or a term, not English prose — a hyphen tear there is never right. A/B on
% a full 197pp lakit: identical overfull sets, same page count, ~11
% paragraphs healed, one mild underfull (badness 3525) that replaced a
% mutilated "Regis-/ters". Long English PROSE belongs in the englishprose
% environment below, which restores explicit-hyphen breaks while its paragraph
% is still active. A raw otherlanguage block inherits the global prohibition.
% (\l@nohyphenation needs \csname:
% '@' is not a letter here.)
% v3.43.2 — \exhyphenpenalty moved OUT of \en to a GLOBAL assignment below.
% Render-proven: a group-local \exhyphenpenalty is DEAD CODE — TeX reads the
% line-breaking parameters at \par time, after the \en group has closed, so
% \en{write-through} / plain not-taken still tore at the "-" ("write-/
% through" shipped on a real exam page). \language is different: it is
% recorded in the node list as text is added, so the in-group nohyphenation
% assignment above DOES work (and stays). The global penalty was A/B-measured
% on a 197pp lakit: identical overfull set (25), identical page count, the
% not-/taken tear healed. Hebrew is unaffected in practice — Hebrew compounds
% use maqaf (־), not ASCII "-".
\newcommand{\en}[1]{\foreignlanguage{english}{\language=\csname l@nohyphenation\endcsname #1}}
% Babel runs extras locally for inline islands and paragraph language changes.
% An English main document must not inherit the Hebrew no-hyphen paragraph rule.
\addto\extrashebrew{\exhyphenpenalty=10000\relax}
\addto\extrasenglish{\exhyphenpenalty=50\relax}
\ifx\lakitmainlanguage\hllpdfenglish
  \exhyphenpenalty=50\relax
\else
  \exhyphenpenalty=10000\relax
\fi
% Full English paragraphs need ordinary explicit-hyphen breaks. Use this
% environment instead of a raw otherlanguage block; the paragraph ends while
% the local penalty is still active.
\newenvironment{englishprose}
  {\par\begin{otherlanguage}{english}\exhyphenpenalty=50\relax}
  {\par\end{otherlanguage}}
\newenvironment{hebrewprose}
  {\par\begin{otherlanguage}{hebrew}\exhyphenpenalty=10000\relax}
  {\par\end{otherlanguage}}
% Use for shell text whose locale must not depend on the last body paragraph.
\DeclareRobustCommand{\doclang}[1]{\foreignlanguage{\lakitmainlanguage}{#1}}

% ============================================================================
% ★★ — fixes for defects that are INVISIBLE in the log and in the checks
% and only surface on a VISUAL RENDER of the PDF. (Each verified empirically.)
% ----------------------------------------------------------------------------
% v3.22: \emph is now REAL slanted emphasis (FakeSlant italic face defined on
% the \babelfont{rm} line above), replacing the v3.21 \emph->\textbf hack: it
% preserves \emph's toggle semantics and keeps bold free for its own role.
% Prefer bold-as-emphasis for your doc? one line: \renewcommand{\emph}[1]{\textbf{#1}}
% page number: force the digits LTR so they never reverse (18 -> 81) in the RTL
%   footer. That reversal is a NON-DETERMINISTIC bidi glitch (observed once in
%   ~80 pp); this LuaLaTeX build is not byte-reproducible, so the defensive
%   LTR wrap is the only reliable guard. Latin digits are always LTR anyway.
\renewcommand{\thepage}{\foreignlanguage{english}{\arabic{page}}}
% v3.70: \pagenumbering \gdef's \thepage and silently DESTROYS the wrap above —
% after \pagenumbering{arabic} the folio guard is gone for the whole rest of the
% document. Two collections discovered this independently and hand-restore the
% wrap in main.tex with the same comment; that knowledge now lives here instead.
% The wrapper re-applies the LTR wrap after the kernel call, for any style.
\makeatletter
\let\lakin@orig@pagenumbering\pagenumbering
\renewcommand{\pagenumbering}[1]{%
  \lakin@orig@pagenumbering{#1}%
  \renewcommand{\thepage}{\foreignlanguage{english}{\csname @#1\endcsname\c@page}}}
\makeatother
% emergencystretch: Hebrew does not hyphenate, so an unbreakable LTR island
%   (\en{...} / an inline math run) landing at the margin can leave a small
%   (~2-13pt) overfull with no legal breakpoint. This rescues ONLY those
%   inline-prose tie-breaks; it is a ceiling used only on offending lines and
%   does NOT touch display-math / listing / table overflow (those still surface
%   in the log for the content-level overflow pass). Verified: baseline overfull
%   -> 0 with this set, with no change to pagination or underfull.
\setlength{\emergencystretch}{3em}
% v3.40: inline math (almost) never breaks across lines. TeX's default
%   (binop 700 / rel 500) lets an inline run split at an operator —
%   $(1.f)\cdot 2^{e}$ shipped as "(1.f)·" ending one line and "2^e" opening
%   the next (reader-reported). In LTR that split is merely ugly; under an
%   RTL base the carried operator lands on the visually WRONG side of the
%   break and the formula reads backwards.
%   WHY 10000 (v3.45.0 — the display-first resolution). Inline math is
%   FORBIDDEN from breaking, so a torn island under an RTL base is
%   structurally impossible. The v3.40 cycle used 9999 ("break only as a last
%   resort") because a bare 10000 left ~30-39 WIDE inline equations
%   overflowing (worst 90.7pt) — but 9999's "last resort" fires OFTEN in
%   tight, non-hyphenating Hebrew: a reader found torn islands across the
%   fleet (~150 candidates; e.g. "P(y |" ending one line and "x_0)" opening
%   the next, read backwards). The step the 9999 cycle missed: every one of
%   those ~30 overflows is a wide equation that BELONGS in display. Convert
%   them to display FIRST, then 10000 gives 0 tears AND 0 overflows (measured,
%   12/12 fleet). Doctrine: inline math is ATOMIC; a wide equation goes in a
%   display, never inline. See content-style ("inline math is atomic").
\binoppenalty=10000 \relpenalty=10000
% v3.38: uniform sentence spacing. TeX's default (\nonfrenchspacing, sfcode
% 3000) is the OLD Anglo-American wide-space-after-period convention — Hebrew
% typography never had it, and modern Latin practice (Chicago et al.) dropped
% it too. Measured on real pages: post-period gaps up to 1.9-2.1x the line's
% mean word gap — the dominant "rivers" source in justified Hebrew prose.
% Survives babel language switches (probed: sfcode stays french inside \en{}
% islands and otherlanguage blocks).
\frenchspacing
% v3.38: widow/orphan protection. LaTeX's default penalties (150/150/50)
% allow a lone line at a page top/bottom almost freely. 10000 = forbid;
% this is SAFE here precisely because the document is \raggedbottom (the
% demoted line just moves on, no vertical stretch fight) — under flushbottom
% these values would cause underfull pages. raggedbottom is the report-class
% default; made explicit so a future class/option change cannot silently
% flip it and turn the 10000s harmful.
\raggedbottom
\widowpenalty=10000 \clubpenalty=10000 \displaywidowpenalty=10000
% em-dash: --- does NOT become — under Frank Ruhl (the font lacks tlig; verified
%   the width is identical with/without Ligatures=TeX, while the mechanism works
%   on Latin Modern). TYPE THE CHARACTER — (U+2014) DIRECTLY; likewise – (U+2013).
%   check_bidi_figures.py flags a --/--- that is adjacent to Hebrew.
% ============================================================================

% ---- STRUCTURAL DIRECTION-SAFETY for \texttt (v3.7) -------------------------
% A parenthesis, bracket, or $ INSIDE \texttt{} mirrors/migrates in RTL prose
% (MIPS 0($s0)->"(s0$)0", $rs->"rs$", arr[i]->"[arr[i") because \texttt sets the
% font but NOT the direction. Redefining \texttt to LTR-isolate its content fixes
% the WHOLE class at once (parens, brackets, $) so plain \texttt{lw $t0,0($s0)}
% is safe — no need to remember \code{}. Guarded by \ifmmode because
% \foreignlanguage is invalid in math (there, plain monospace is already LTR).
% NB (corrected v3.56 — the old note here said "stays tofu", which is WRONG):
% \texttt{Hebrew} is a FATAL build failure. The tt family is DejaVu Sans Mono
% (zero Hebrew) and the tt fallback hebfb is DejaVuSerif (also zero Hebrew), so
% the subsetter gets an empty glyph set and LuaTeX dies with
%   "! error: (file .../DejaVuSansMono.ttf) (type 2): there are no glyphs in the
%    subset"  -> no output PDF.
% The message names a FONT FILE, never Hebrew/\texttt/the source line. And it is
% non-local: add an unrelated Latin \code{} elsewhere and the fatal degrades into
% a SILENT DELETION of the Hebrew word. Put Hebrew in \h{}/normal text, never in
% \texttt{}/\code{}; flagged HEBREW-IN-MONO. To disable, delete this block.
\let\origtexttt\texttt
\protected\long\def\texttt#1{\ifmmode\origtexttt{#1}\else\foreignlanguage{english}{\origtexttt{#1}}\fi}
% v4.8.0: defer the direction test until TeX has inserted an alignment cell's
% math template. Unprotected expansion here tested \ifmmode too early at the
% FIRST token of aligned/array/cases/matrix cells. Engine-level \protected also
% preserves the original \texttt implementation without the self-recursion
% observed when retrofitting \DeclareRobustCommand to its saved alias.
\protected\long\def\code#1{\ifmmode\origtexttt{#1}\else\foreignlanguage{english}{\origtexttt{#1}}\fi}

% ---- \codeline — display-code line (v3.24): the code analogue of display math.
% Reader-friction rule (content-style.md §5b): inline \code{} is for a SINGLE
% ATOM only (one register/variable/mnemonic, a tight no-space expression).
% A multi-token expression — an instruction with operands, a spaced assignment,
% a register-binding list — zigzags the eye RTL<->LTR mid-sentence and, being
% an unbreakable island, line-breaks badly. Give it its own centered line with
% display-math vertical rhythm. \centerline = one box => the expression never
% wraps; \nopagebreak keeps it glued to its introducing sentence.
%   ...the Hebrew sentence ends with a colon:
%   \codeline{lw \$t1, -4(\$t0)}
%   and the prose continues (escape dollars as \$).
% v3.66 (first line of the body): flush a PENDING \item label before the \par.
% LaTeX defers the label box (\@labels) until the next \everypar, and an \hbox in
% vertical mode never fires \everypar — so `\item \codeline{...}` printed the code
% at y=69.1 and the item's own "1." BELOW it at y=103.6 (the number appears after
% its content, attached to nothing). \if@inlabel\leavevmode\fi is CONDITIONAL on
% purpose: an unconditional \leavevmode would open an empty paragraph line when
% \codeline follows ordinary prose. Verified both ways — in a list the label now
% sits on the code's line at its normal x, and after prose the geometry is
% byte-identical to before.
\makeatletter
\newcommand{\codeline}[1]{%
  \if@inlabel\leavevmode\fi
  \par\nopagebreak
  \vspace{\abovedisplayskip}\nopagebreak
  % v3.62: was \centerline{...}, i.e. \hbox to\hsize{\hss #1 \hss}. \hss has
  % INFINITE SHRINK, so TeX could always "fit" the box and NEVER reported an
  % Overfull — the same \hss mechanism as BLD-27's \item label. Measured: a wide
  % \codeline in a defbox painted its first token ON the box's left border bar and
  % sat 9.5pt outside the frame interior on both sides, with 0 Overfull, 0 Underfull
  % and all four gates PASS; in bare prose it reached 21.0pt into BOTH margins.
  % \hfil has no shrink, so the overflow is reported and build.sh's ovf_vis/ovf_big
  % counters see it; \linewidth (not \hsize) makes it box-aware, so a \codeline
  % inside a colored box centres on the box interior. Verified: a FITTING \codeline
  % keeps byte-identical glyph positions, and the box-crossing case now emits
  % "Overfull \hbox (19.03pt too wide)". A wide \codeline should be split in two.
  \hbox to\linewidth{\hfil\code{#1}\hfil}%
  \nopagebreak\vspace{\belowdisplayskip}}
\makeatother

% ---- a thin decorative rule helper ----
\newcommand{\sep}{\par\vspace{2pt}\noindent\textcolor{cRule}{\rule{\linewidth}{0.6pt}}\par\vspace{4pt}}

% ---- colors ----
\definecolor{cBlue}{HTML}{1F4E79}   \definecolor{cBlueBg}{HTML}{EAF1F8}
\definecolor{cGreen}{HTML}{2E6B4F}  \definecolor{cGreenBg}{HTML}{EAF4EE}
\definecolor{cAmber}{HTML}{9C6A00}  \definecolor{cAmberBg}{HTML}{FBF3E2}
\definecolor{cPurple}{HTML}{5B2C6F} \definecolor{cPurpleBg}{HTML}{F3ECF7}
\definecolor{cTeal}{HTML}{0E6E6E}   \definecolor{cTealBg}{HTML}{E6F4F4}
\definecolor{cRed}{HTML}{8E2B2B}    \definecolor{cRedBg}{HTML}{FBEAEA}
\definecolor{cGray}{HTML}{555555}   \definecolor{cRule}{HTML}{1F4E79}

% ---- math operators ----
\DeclareMathOperator{\Tr}{Tr}
\DeclareMathOperator{\rank}{rank}
\DeclareMathOperator{\diag}{diag}
\DeclareMathOperator{\sgn}{sgn}

% ---- bra-ket (quantikz defines \ket/\bra/\braket; add the rest safely) ----
\providecommand{\ket}[1]{\left|#1\right\rangle}
\providecommand{\bra}[1]{\left\langle#1\right|}
\providecommand{\braket}[2]{\left\langle#1\middle|#2\right\rangle}
\newcommand{\dyad}[1]{\left|#1\right\rangle\!\left\langle#1\right|}
% v3.21: general outer product |a><b| and expectation value <A> — agents reach
% for these by reflex (esp. \expval); quantikz defines neither, so their absence
% threw "Undefined control sequence" in a real build. \providecommand = safe.
\providecommand{\ketbra}[2]{\left|#1\right\rangle\!\left\langle#2\right|}
\providecommand{\expval}[1]{\left\langle#1\right\rangle}
\newcommand{\abs}[1]{\left|#1\right|}
\newcommand{\norm}[1]{\left\lVert#1\right\rVert}
\newcommand{\half}{\tfrac12}
\newcommand{\R}{\mathbb{R}}\newcommand{\C}{\mathbb{C}}\newcommand{\Z}{\mathbb{Z}}
% identity operator. NB: \mathbb{1} silently renders a WRONG glyph (msbm has no
% blackboard digits). Use \mathbb{I} = 𝕀 (works with amssymb, no extra package).
% Do NOT "fix" this back to \mathbb{1}.
\newcommand{\id}{\mathbb{I}}

% ---- physical units (v3.35): \punit — upright unit with thin space ---------
% Usage: $300\punit{K}$, $80\punit{mV/decade}$, $10^{7}\punit{cm/s}$ — the
% NUMBER AND UNIT LIVE IN ONE MATH ISLAND. Promoted verbatim from a 199-page
% production lakit (advanced-electronic-devices, dozens of uses).
% Two measured traps this design avoids (both render-proven):
%   1. siunitx SCRAMBLES under RTL — \SI{3.5}{eV} renders "eV 3.5", exponents
%      detach ("Hz 10 x 10^9"): its quantity assembly is not bidi-isolated.
%      Do NOT load siunitx in this pipeline.
%   2. A text-mode unit next to a bare-prose number is TWO adjacent LTR runs
%      and reorders ("5 mm" -> "mm5" — the BIDI-2/3 class). \punit is therefore
%      math-only ON PURPOSE: prose misuse fails LOUDLY at compile (\mathrm
%      outside math) instead of scrambling silently.
% Domain shorthands (\MeV, \fm...) stay per-document: \newcommand{\MeV}{\,\text{\en{MeV}}}
\providecommand{\punit}[1]{\,\mathrm{#1}}

% ============================================================================
% Colored boxes — mdframed (NOT tcolorbox, which breaks under bidi).
% ★★ v3.4: innerleftmargin/innerrightmargin give code listings room so their
% frame/line-numbers do not spill out the box border (see listing styles below).
% ============================================================================
% v3.61: everyline/split skips were applied to `tablebox` only (BLD-14, v3.37) —
% the six colored boxes below never got them, so they still split OPEN across a page
% break: the page-N segment gets top+sides and NO bottom rule, the page-N+1 segment
% sides+bottom and NO top rule. Measured on the shipped fleet this is a COVERAGE GAP,
% not a re-report of BLD-14: 0 split-open `tablebox` segments (that fix holds) vs 62
% split-open COLORED boxes across 7 documents. Verified after the fix: both segments
% close, and the frametitle is NOT duplicated on the continuation. BLD-14's caveat
% still applies — a box taller than a full page grinds mdframed's split-trial loop, so
% split such content manually.
% ============================================================================
% v3.72: frametitle={} leads EVERY \newmdenv option list below. mdframed options
% are GROUP-scoped: a box nested inside \begin{exbox}[frametitle={...}] used to
% INHERIT the outer title — measured, the identical title string rendered twice,
% the second riding a title-bar fill inside the INNER box's frame, for all seven
% environments in both directions. Both natural forms were broken (= preamble
% bug, not an authoring rule). The explicit reset makes every box titleless
% unless ITS OWN [frametitle={...}] says otherwise; verified both directions:
% no duplicate on the nested box, an explicit title still renders.
% Geometry is shared with the optional semantic-width table container. Themes
% select their own existing semantic style below; neither path fixes columns.
\mdfdefinestyle{hllpdfsemanticgeometry}{linewidth=1.1pt,
  innertopmargin=6pt,innerbottommargin=7pt,skipabove=9pt,skipbelow=9pt,
  innerleftmargin=9pt,innerrightmargin=9pt,everyline=true,
  splitbottomskip=6pt,splittopskip=12pt,roundcorner=3pt}
\newcommand{\hllpdfsemanticframestyle}{hllpdfsemanticgeometry}
\newmdenv[frametitle={},style=hllpdfsemanticgeometry,linecolor=cBlue,backgroundcolor=cBlueBg,
  frametitlebackgroundcolor=cBlue,frametitlefontcolor=white,
  ]{defbox}   % הגדרה
\newmdenv[frametitle={},style=hllpdfsemanticgeometry,linecolor=cGreen,backgroundcolor=cGreenBg,
  frametitlebackgroundcolor=cGreen,frametitlefontcolor=white,
  ]{thmbox}   % משפט/טענה
\newmdenv[frametitle={},style=hllpdfsemanticgeometry,linecolor=cAmber,backgroundcolor=cAmberBg,
  frametitlebackgroundcolor=cAmber,frametitlefontcolor=white,
  ]{notebox}  % הערה
\newmdenv[frametitle={},style=hllpdfsemanticgeometry,linecolor=cPurple,backgroundcolor=cPurpleBg,
  frametitlebackgroundcolor=cPurple,frametitlefontcolor=white,
  ]{exbox}    % דוגמה
\newmdenv[frametitle={},style=hllpdfsemanticgeometry,linecolor=cRed,backgroundcolor=cRedBg,
  frametitlebackgroundcolor=cRed,frametitlefontcolor=white,
  ]{warnbox}  % טעות נפוצה
\newmdenv[frametitle={},style=hllpdfsemanticgeometry,linecolor=cTeal,backgroundcolor=cTealBg,
  frametitlebackgroundcolor=cTeal,frametitlefontcolor=white,
  ]{keybox}   % רעיון מרכזי

% ---- codebox (v4.5.0): the dedicated container for code -------------------
% An OPTION, not a policy: asmblock unboxed stays canonical for a full listing,
% and a snippet inside any of the six colored boxes still works as before.
% What makes codebox worth its own environment is that it is the ONE framed box
% that keeps its margin line numbers.
%
% v4.4.0 shipped it WITHOUT them, on the strength of v3.37's "no margin knob can
% save them" plus four fresh probes that agreed. That was wrong: every probe had
% moved the code, the box, or both — never the key that carries the NUMBER
% COLUMN. Measured on the page:
%   resetmargins=false   the listing follows THIS box instead of the page,
%                        moving the code from x=75.26 to x=97.54
%   xrightmargin=26pt    under RTL this is the PHYSICAL LEFT indent, and the
%                        only key that carries the \llap'ed digits with it:
%                        x=64.28 -> 90.18, clearing the rule at 65.64 by 23.6pt
%   xleftmargin=26pt     moves the code and LEAVES THE DIGITS on the rule —
%                        the asymmetry that hid the fix for eleven versions
% resetmargins is hoisted out of the four styles (see the \lstset below them) so
% this override is possible at all; every other listing still resolves to true.
% Pinned by tests/codebox_numbers.tex + probe_lstnum.py --inside-rule, whose
% negative control strips xrightmargin and watches the digits fall back.
\definecolor{cSlate}{HTML}{3F4A56}  \definecolor{cSlateBg}{HTML}{F2F4F7}
\newmdenv[frametitle={},leftline=true,rightline=true,topline=true,bottomline=true,
  linecolor=cSlate,backgroundcolor=cSlateBg,linewidth=0.9pt,
  frametitlebackgroundcolor=cSlate,frametitlefontcolor=white,
  innertopmargin=6pt,innerbottommargin=7pt,skipabove=9pt,skipbelow=9pt,
  innerleftmargin=9pt,innerrightmargin=9pt,everyline=true,splitbottomskip=6pt,splittopskip=12pt,
  roundcorner=3pt]{codebox}   % קוד
% The two keys that make in-box numbering work, both measured on the page:
%   resetmargins=false  the listing follows THIS box instead of the page
%   xrightmargin        under RTL this is the PHYSICAL LEFT indent, and it is
%                       the only one that carries the \llap'ed number column
%                       with it. xleftmargin moves the code and leaves the
%                       digits behind on the rule (measured: code 97.54,
%                       digits 64.28 — the v3.37 defect, reproduced); with
%                       xrightmargin=26pt the digits land at x=90.18, a
%                       clearance of 23.6pt from the rule at 65.64.
% This asymmetry is why v3.4's note had the semantics backwards and v3.37
% concluded the digits could not be saved.
% <hllpdf:listing-container-state>
% v4.29: all state is local to the current environment group. A semantic
% callout masks number text, while a nested codebox restores the formatter
% saved before that mask. Save it once across successive semantic wrappers;
% an explicit numbers=none (or an explicitly empty formatter) remains valid.
\makeatletter
\newif\ifhllpdf@incodebox
\newif\ifhllpdf@englishcodebox
\newif\ifhllpdf@lstnumbermasked
\let\hllpdf@savedlstnumber\thelstnumber
\makeatother
% </hllpdf:listing-container-state>
% <hllpdf:codebox-gutter>
\makeatletter
\AtBeginEnvironment{codebox}{%
  \hllpdf@incodeboxtrue
  % Snapshot the OUTER language before asmblock selects English.
  \iflanguage{english}{\hllpdf@englishcodeboxtrue}{\hllpdf@englishcodeboxfalse}%
  \ifhllpdf@lstnumbermasked
    \let\thelstnumber\hllpdf@savedlstnumber
    \hllpdf@lstnumbermaskedfalse
  \fi
}
% listings' AddToHook is GLOBAL. Register once, and gate on local box state;
% adding an unconditional hook at every codebox entry leaks into later listings.
\lst@AddToHook{PreInit}{%
  \ifhllpdf@incodebox
    % Apply every listing override here, inside the LISTING's own group.
    % Setting margins at box entry leaks them into nested semantic callouts.
    \lstset{resetmargins=false,xrightmargin=26pt}%
    % Styles resolve before PreInit. In LTR the physical number gutter follows
    % xleftmargin; in Hebrew retain the measured inherited xrightmargin path.
    \ifhllpdf@englishcodebox
      \lstset{xleftmargin=26pt,xrightmargin=0pt}%
    \fi
    % ONE box, not two: suppress only the frame/background/skips of a listing
    % inside this codebox. The next unboxed listing keeps its own style.
    \lstset{frame=none,backgroundcolor=\color{cSlateBg},aboveskip=0pt,belowskip=0pt}%
  \fi
}
\makeatother
% </hllpdf:codebox-gutter>

% v3.37: line numbers inside a COLORED BOX are sliced by the box border —
% measured: digit span x=63.5-67.3 vs border line at x=65.2 (the border runs
% THROUGH the digits), and the inter-border gap (3.5pt) is narrower than a
% digit (3.8pt), so no margin knob can save them (under resetmargins,
% xleftmargin is inert in-box — measured: code at x=75.26 with 2.2em AND
% 3.0em). A 3-6 line boxed snippet does not need line numbers; blank the
% number TEXT inside every colored box (\thelstnumber is what gets printed;
% the style's numbers=left survives, so nothing else shifts — code position
% byte-frozen at x=75.26, verified). Unboxed listings keep their numbers.
% v3.78: the blanking is coupled to the FRAME, not to the box name —
% clean_theme.tex registers LATER hooks (registration order wins) that
% restore the number text inside the boxes it renders frameless.
% v4.5.0 RETRACTS the "no margin knob can save them" clause above. It holds only
% under resetmargins=true (the shipped default): there xleftmargin is inert
% in-box, as measured. With resetmargins=false, xrightmargin DOES carry the
% \llap'ed digit column clear of the rule — measured; the codebox \lstset that
% does it is the one above, in the codebox block. The loop below therefore
% covers SIX boxes, not "every colored box": codebox opts out and keeps its
% numbers.
% <hllpdf:semantic-listing-mask>
\makeatletter
\@for\HEB@lstbox:=defbox,thmbox,notebox,exbox,warnbox,keybox\do{%
  \expandafter\AtBeginEnvironment\expandafter{\HEB@lstbox}{%
    \ifhllpdf@lstnumbermasked\else
      \let\hllpdf@savedlstnumber\thelstnumber
    \fi
    \hllpdf@lstnumbermaskedtrue
    \def\thelstnumber{}%
    % The nearest semantic container owns its listing presentation. On exit,
    % TeX's group restores the enclosing codebox state and formatter.
    \hllpdf@incodeboxfalse}}
\makeatother
% </hllpdf:semantic-listing-mask>

% ============================================================================
% Unified TABLE presentation: subtle uniform slate frame + auto-numbered
% caption rendered INSIDE the frame. Use:
%   \begin{tablebox}\centering  <tabular>  \tabcap{<description>}\end{tablebox}
% Replaces the buggy \caption* (which emitted a stray "טבלה N : *" line) and
% provides an explicit framed option + numbered caption "טבלה N.M: תיאור".
% Optional intro text may be placed at the TOP, before the tabular, inside the box.
% ============================================================================
\definecolor{cTblLine}{HTML}{B9C4D2}  \definecolor{cTblBg}{HTML}{FAFBFC}
% v3.23: class-aware — under `article` there is no chapter counter (the promised
% report->article switch used to die with "No counter 'chapter' defined").
\makeatletter
\@ifundefined{c@chapter}{%
  % v3.31: in article the counter resets per SECTION, so an unprefixed number
  % repeated ("טבלה 1." in every section — measured on a two-section render).
  % Same conditional idea as the report branch: prefix the section number once
  % a real \section exists; a pre-section table stays "טבלה 1".
  \newcounter{tbl}[section]
  \renewcommand{\thetbl}{\ifnum\value{section}>0 \thesection.\fi\arabic{tbl}}%
}{%
  \newcounter{tbl}[chapter]
  % v3.30: a report-class doc that uses SECTIONS ONLY (chapter stays 0 —
  % the common short-doc pattern) used to print "טבלה 0.1"; the chapter
  % prefix now appears only once a real \chapter exists (render-verified
  % both ways: "טבלה 1." chapterless, "טבלה 1.1" with chapters).
  \renewcommand{\thetbl}{\ifnum\value{chapter}>0 \thechapter.\fi\arabic{tbl}}%
  % v3.64: the SAME leak in \thesection itself, left behind by v3.30. That fix
  % guarded the CAPTION counter and stopped there, so a chapterless report still
  % printed the heading "0.1 סעיף ראשון" and \ref to it printed "0.2" — measured:
  %   p1 y=65.2  '…0.1'   p1 y=144.9 '…0.2'   p1 y=183.4 '…0.2.1'
  % report.cls hardcodes \thesection as \thechapter.\@arabic\c@section, and the
  % leading 0 is the chapter counter that no \chapter ever incremented. Fixing
  % \thesection fixes the whole chain — \thesubsection and \thesubsubsection are
  % defined in terms of it. \theequation and \thetbl/\thefig were already guarded,
  % which is exactly why this survived: every NUMBER around the heading looked
  % right, so the heading's own 0 read as a style choice rather than a bug.
  % \appendix is safe: report.cls resets chapter to 0 and redefines \thechapter to
  % \@Alph\c@chapter, so an appendix \chapter makes the counter 1 and the guard
  % emits "A.1" unchanged. Render-verified in all four modes (chapterless,
  % chaptered, appendix+chapter, appendix+section-only).
  \renewcommand{\thesection}{\ifnum\value{chapter}>0 \thechapter.\fi\arabic{section}}%
}
\makeatother
% v3.37: everyline=true — if the box DOES split across pages, every segment
% gets a CLOSED frame (measured before: the split ran open-ended into the page
% edge with no bottom border, and the caption could strand alone at the next
% page top — "טבלה A.2" shipped that way). The split skips give the segments
% breathing room. In practice, with the caption GLUE below, every previously-
% split table in the corpus now moves WHOLE instead (the white-gap tradeoff,
% documented benign); everyline stays as the safety net for a genuinely
% page-tall table. NB: a table taller than ONE FULL PAGE can grind mdframed's
% split-trial loop (reproduced: minimal doc, minutes of 1pt-step Overfull
% probes) — break such a monster into two tables.
\mdfdefinestyle{hllpdftablegeometry}{leftline=true,rightline=true,topline=true,bottomline=true,
  linewidth=0.9pt,
  innertopmargin=9pt,innerbottommargin=8pt,skipabove=11pt,skipbelow=11pt,
  innerleftmargin=11pt,innerrightmargin=11pt,roundcorner=4pt}
\newcommand{\hllpdftableframestyle}{hllpdftablegeometry}
\newmdenv[frametitle={},style=hllpdftablegeometry,
  linecolor=cTblLine,backgroundcolor=cTblBg,
  % v4.2.0 (adopted): tablebox is ATOMIC. Short tableboxes near page bottoms
  % repeatedly entered mdframed's split trial, emitting benign Overfull-vbox
  % bursts and occasionally separating the intro prose from the tabular.
  % nobreak=true moves the complete short table to the next page instead. The
  % output-active Overfull-vbox gate still catches a box taller than a page;
  % everyline/split skips remain as defensive styling should that path ever
  % re-arm, but the canonical path never publishes a split tablebox.
  nobreak=true,
  everyline=true,splitbottomskip=6pt,splittopskip=12pt,
  % v3.79 (BLD-37 table side, correction): startinnercode/endinnercode run
  % INSIDE mdframed's body vbox (mdframed.sty:848 / :856+863) -- the second one
  % after \par\unskip, i.e. after line breaking, which is the only moment the
  % \tabcap anchor's true offset from the top of the box EXISTS. See
  % \lakintabboxclose below. These two keys are the ONLY tablebox change.
  startinnercode={\lakintabboxopen},endinnercode={\lakintabboxclose}]{tablebox}
% v3.37: \nopagebreak GLUE — \vsplit honors penalties, so the caption can no
% longer strand alone at the top of the next page (it takes its rows along, or
% the whole box moves). Render-verified on the real shipped orphan.
% v3.59: \NR@gettitle is what nameref injects into the stdlib title-bearing
% commands (\@sect, \@chapter, \caption) to fill \@currentlabelname. \figcap and
% \tabcap are INDEPENDENT macros — they never called it, so \@currentlabelname kept
% whatever the last heading left and \nameref{fig:x} silently printed the enclosing
% SECTION's title instead of the caption (proven in the .aux: field 3 of \newlabel
% was the heading text while field 1 and the anchor `fig.1.1` were correct — so this
% is a THIRD, distinct head next to BLD-13's counter fix and BIDI-38's autoref-name
% fix, neither of which writes that field). \@ifundefined keeps it inert if
% hyperref/nameref is ever dropped. Guarded helper because \figcap is defined
% after \makeatother.
% NB both helpers below need \makeatletter — this part of the preamble is OUTSIDE
% any letter-catcode region, so an unwrapped \@currentlabel parses as LaTeX's \@
% and dies with "You can't use \spacefactor in vertical mode" (measured).
\makeatletter
\newcommand{\lakinrtitle}[1]{\@ifundefined{NR@gettitle}{}{\NR@gettitle{#1}}}
% v3.59: \hitem{<letter>}{<key>} — a REFERENCEABLE Hebrew-lettered list item.
%   \begin{enumerate} \hitem{א}{it:a} הסעיף הראשון. \end{enumerate}  ->  \ref{it:a} = (א)
% Why this exists: automatic Hebrew list labels fail silently, so SKILL.md mandates
% an EXPLICIT \item[(\textbf{א})] — but LaTeX puts \refstepcounter INSIDE
% \if@noitemarg (latex.ltx:12964), and \item[...] never sets that flag. So an
% explicit item steps no counter and never touches \@currentlabel: a \label there
% binds to the AMBIENT one and \ref prints the CHAPTER/SECTION number where the
% item letter belongs (proven in the .aux: {chapter.1}). Same doctrine as \figcap —
% own the label state before \label. \phantomsection matters: without it the
% hyperlink still lands on the chapter head (measured); with it the anchor is
% section*.N and the bookmark outline stays clean.
\newcommand{\hitem}[2]{\item[(\textbf{#1})]%
  \def\@currentlabel{(#1)}\def\@currentlabelname{#1}\phantomsection\label{#2}}
% NB (v3.66): \hitem is \ref-only. \autoref{key} calls it a SECTION ("סעיף (א)")
% because \phantomsection brands the anchor section*.N, and \nameref{key} prints
% the bare letter. For \autoref-able items prefer \hebenumlabels below.
%
% v3.66: \hebenumlabels — OPT-IN automatic Hebrew enumerate letters, done right.
%   {\hebenumlabels\begin{enumerate}\item\label{q:a} ...\end{enumerate}}
%   -> labels (א) (ב) (ג); \ref{q:a} = א; anchor Item.N (real, per-item).
% This retires the belief that made \item[(\textbf{א})] + \hitem necessary.
% "Automatic Hebrew labels fail silently" was TRUE but unexplained, and the
% mechanism is now bisected (plain article vs +enumitem, measured): ENUMITEM
% PRE-EXPANDS the label expression in an \edef BEFORE \refstepcounter, with the
% counter still at 0. Tokens it shields by name (\arabic,\alph,\roman) survive;
% a plain \newcommand \ifcase wrapper evaluates at 0 -> the EMPTY zero-branch ->
% the label prints as "()" with the letter GONE (and \the\c@enumi prints a
% literal 0, \number\value dies fatally pointing at \item). The cure is one
% property, not a doctrine: the wrapper must be ROBUST (survives the \edef
% unexpanded), while \theenumi stays EXPANDABLE (so \refstepcounter writes the
% real letter to the .aux). Verified: labels (א)(ב)(ג), \ref -> א/ב/ג each item
% its own value, .aux = {א}{...}{Item.1} — versus the explicit-\item doctrine
% where BIDI-50/\hitem exist purely to patch the counter that never stepped.
\def\@hebalph#1{\ifcase#1\or א\or ב\or ג\or ד\or ה\or ו\or ז\or ח\or ט\or י\or
  יא\or יב\or יג\or יד\or טו\or טז\or יז\or יח\or יט\or כ\else\@ctrerr\fi}
\DeclareRobustCommand*\hebalph[1]{\expandafter\@hebalph\csname c@#1\endcsname}
\newcommand{\hebenumlabels}{%
  \renewcommand{\theenumi}{\@hebalph\c@enumi}%
  \renewcommand{\labelenumi}{(\textbf{\hebalph{enumi}})}}
\makeatother
% v3.63: \mbox the LABEL. \thetbl ("7.1") and a caption description that STARTS
% WITH A DIGIT ("32 הרגיסטרים…") are two European-Number runs separated only by
% neutrals (the label's own "." and \enspace). Under the RTL paragraph base the UBA
% resolves those neutrals to the RTL level, so both numbers sit in ONE reversed span
% and EXCHANGE VISUAL POSITIONS. Measured on the shipped PDF: a table whose real
% number is 7.1 printed "טבלה 32" — a different, EXISTING table number — while \ref
% correctly printed 7.1, so the caption silently contradicted its own cross-reference,
% and the description lost its leading count. The source contains no adjacency at all;
% it is manufactured by this macro, which is why every source linter was blind.
% \mbox makes the label one atomic box so nothing can reorder across it. Verified
% zero-regression: an ordinary Hebrew caption keeps identical x-geometry.
% One paragraph style covers label, number, separator and description.
% Explicit project/template overrides remain possible; the label hook is kept
% for compatibility, with an identity default. Counters/anchors stay unchanged.
\newcommand{\lakitcaptionstyle}{\normalfont\small\color{black}\hypersetup{allcolors=black}}
\newcommand{\lakitcaptionlabelstyle}[1]{#1}
\newcommand{\lakitcaptionpunctuation}{:}
% Native floats, captionof and longtable use the same presentation policy.
% subcaption is optional and inherits the same font rather than a second size.
\usepackage{caption}
\DeclareCaptionFont{lakit}{\lakitcaptionstyle}
\DeclareCaptionLabelFormat{lakit}{\lakitcaptionlabelstyle{#1~\en{#2}}}
\DeclareCaptionLabelSeparator{lakit}{\lakitcaptionpunctuation\enspace}
\captionsetup{font=lakit,labelfont={},textfont={},labelformat=lakit,
  labelsep=lakit,justification=centering}
\makeatletter
\AtBeginDocument{\@ifpackageloaded{subcaption}{%
  \captionsetup[sub]{font=lakit,labelfont={},textfont={},labelsep=lakit}}{}}
\makeatother
% v3.64: OPT-IN list-of-figures/tables entry via an optional short caption.
% \listoffigures and \listoftables are inherited from report.cls and LOOK available,
% but \figcap/\tabcap step the custom \c@fig/\c@tbl counters and never call
% \addcontentsline — so both printed their heading and NOTHING ELSE. Measured: a doc
% with one \figcap and one \tabcap rendered page 1 as 'רשימת האיורים\n1\n', and the
% .lof held only \@chapter's \addvspace, zero \contentsline entries. An empty list
% under a correct Hebrew heading reads as "this document has no figures".
%
% WHY OPTIONAL, AND NOT AUTOMATIC. Writing #1 unconditionally was implemented first
% and REJECTED: a caption containing \codeline compiles CLEANLY today (verified on
% the pre-v3.64 preamble: 0 errors, PDF produced) but becomes FATAL once it is also
% written to the .lof — \codeline is \hbox to\linewidth{...}, and a full-width box in
% a list entry dies with "\pdfextension endlink cannot be used in vertical mode", no
% output PDF. Gaining a feature nobody uses (0 of 13 collections call \listoffigures)
% at the price of breaking a construct that works is a bad trade. Stock \caption has
% exactly this limitation and exactly this escape hatch, so the convention is familiar.
% \en{}, \code{}, \h{} and inline math in the short text are all safe — verified, they
% survive the write and typeset correctly in the list (RTL geometry confirmed: number
% rightmost at x=509.9, title left of it, dots and folio at the far left).
% With no optional argument the macro is byte-identical to v3.63 — zero risk to every
% existing document, and \listoffigures lists exactly the captions that opted in.
% BLD-37, TABLE side (the \figcap v3.74 idea, adapted): clicking a \ref to a
% table used to land MID-table — hyperref's automatic anchor inside
% \refstepcounter was emitted on the still-open horizontal line whose TAIL is
% the tabular (an hbox whose baseline is the table's vertical CENTER), so the
% jump scrolled the top half of the table above the viewport. The v3.74 entry
% deferred this side claiming "tablebox splits — no single box to measure";
% that deferral dissolves: the measurement target is the TABULAR itself at
% \tabcap ENTRY, while it is still the tail of the open list — the same
% situation \lakin@grabfig handles at an island's end (and the fleet has 0
% split tableboxes). So \lakin@grabtab runs as the FIRST token, before any
% typesetting can push the box away, and the anchor is re-emitted raised above
% the measured tabular. \lakintabraise is a SEPARATE register from
% \lakinfigraise (defined with it, below the LTR-island block): a captionless
% island intentionally leaves the fig register set, and a table must never
% inherit that stale figure height.
%
% WHY the anchor stays AT THE ENTRY SPOT and is NOT deferred into the caption
% paragraph the way \figcap defers it: at \par TeX DROPS a final glue item,
% and the automatic anchor's smash box sitting BEHIND the interword glue kept
% that glue alive — so every centered tabular has always sat half a space off
% true center. Deferring the anchor removed that box and was measured as 982
% differing glyph records: EVERY tabular line shifted 1.30pt (2.61pt for
% pattern 7's two minipages) while captions stayed put. Emitting the raised
% anchor at the old spot — through hyperref's own \Hy@raisedlink, exactly the
% macro the automatic path uses — keeps the line's node structure, and so the
% A/B glyph diff at 0, BY CONSTRUCTION. \Hy@raisedlink reads the raise from
% \HyperRaiseLinkDefault (default \baselineskip = the old mid-table dest); the
% group-local \edef freezes measured-height+42pt into it, and \smash inside
% \Hy@raisedlink hides any raise from the line's extent. \HyCnt@ProvideTheH-
% Counter keeps the dest NAME byte-identical to the automatic one (it defines
% \theHtbl, which \hyper@makecurrent prefers — without it the first measured
% caption was named tbl.1 where the automatic path said tbl.0.1, measured).
% vmode entry, or anything but a box at the tail (no tabular to measure),
% falls back to the plain \refstepcounter — byte-identical to the old path.
\makeatletter
\newcommand{\tabcap}[2][]{\lakin@grabtab
  % v3.79: arm the box-close correction. Set unconditionally and in BOTH
  % branches below's scope, so the plain-\refstepcounter fallback is corrected
  % too; \lakintabboxopen clears it at every \begin{tablebox}, so a \tabcap
  % outside a tablebox can never leave it armed for a later box.
  \global\lakintabcapseentrue
  \ifdim\lakintabraise>0pt
    \H@refstepcounter{tbl}\HyCnt@ProvideTheHCounter{tbl}\hyper@makecurrent{tbl}%
    \begingroup
      % A frameless table targets its measured top. The frame's historical
      % cushion belongs only to tablebox and is corrected at box close.
      \edef\HyperRaiseLinkDefault{\the\lakintabraise}%
      \iflakintabframed
        \edef\HyperRaiseLinkDefault{\the\dimexpr\lakintabraise+42pt\relax}%
      \fi
      \Hy@raisedlink{\hyper@anchorstart{\@currentHref}\hyper@anchorend}%
    \endgroup
    \global\lakintabraise 0pt\relax
  \else
    \refstepcounter{tbl}%
  \fi
  \lakinrtitle{#2}\par\nopagebreak\addvspace{6pt}%
  \def\lakin@short{#1}\ifx\lakin@short\@empty\else
    \addcontentsline{lot}{table}{\protect\numberline{\thetbl}{\protect\ignorespaces #1}}\fi
  {\centering\lakitcaptionstyle\mbox{\lakitcaptionlabelstyle{\tablename~\en{\thetbl}\lakitcaptionpunctuation}}\enspace #2\par}}
\makeatother

% ============================================================================
% Unified FIGURE caption — the \tabcap twin for TikZ/plot figures (v3.34.0).
% A figure (tikzpic) sits inline in the prose; put \figcap BELOW it for a
% numbered, \ref-able caption "איור N.M: תיאור" in the same style as \tabcap:
%   \begin{tikzpic} ... \end{tikzpic}
%   \figcap{\label{fig:foo}description}     % \label FIRST, inside the argument
% WHY a macro, not a hand-written gray line: a manual
%   {\centering\small\textcolor{cGray}{איור: ...}} with a \label inside has NO
% \refstepcounter, so the \label binds to the AMBIENT \@currentlabel (the last
% equation) — \ref then prints an equation number AS the figure number
% (measured on a real page: "איור 48" where 48 is equation 48). \figcap steps a
% real \c@fig counter first, so \label is correct, numbering is consistent, and
% refs are stable. Class-aware like \thetbl (chapter- or section-prefixed).
% ============================================================================
\makeatletter
\@ifundefined{c@chapter}{%
  \newcounter{fig}[section]
  \renewcommand{\thefig}{\ifnum\value{section}>0 \thesection.\fi\arabic{fig}}%
}{%
  \newcounter{fig}[chapter]
  \renewcommand{\thefig}{\ifnum\value{chapter}>0 \thechapter.\fi\arabic{fig}}%
}
\makeatother
\makeatletter
\newcommand{\figcap}[2][]{%
  \ifdim\lakinfigraise>0pt
    % v3.74: step WITHOUT hyperref's automatic anchor, then emit the anchor
    % raised above the measured figure. \H@refstepcounter is hyperref's saved
    % pre-hook \refstepcounter (sets the counter AND \@currentlabel — \label
    % in the caption arg binds exactly as before); \hyper@makecurrent supplies
    % \@currentHref, so the dest NAME is identical to the old automatic one —
    % \ref/\autoref/bookmarks are unchanged, only the dest'S POSITION moves.
    \H@refstepcounter{fig}\hyper@makecurrent{fig}%
    % The anchor is emitted INSIDE the caption paragraph below, as a zero-size
    % raisebox — NOT here in vertical mode, where a \raisebox starts an empty
    % paragraph and adds a blank line (measured: every glyph after the caption
    % shifted +18.8pt). Inside an existing line a [0pt][0pt] box is inert.
    \edef\lakin@anch{\@currentHref}%
    % \edef FREEZES the measured height into the pending macro: the deferred
    % \raisebox is typeset inside the caption paragraph AFTER the register is
    % reset to the 0pt sentinel below — a plain \def read 0 at use time
    % (measured: LAKINRAISE=0.0, dest moved 14pt instead of above the figure).
    \edef\lakin@pend{\noexpand\raisebox{\the\dimexpr\lakinfigraise+42pt\relax}[0pt][0pt]{%
      \noexpand\hyper@anchorstart{\lakin@anch}\noexpand\hyper@anchorend}}%
    \global\lakinfigraise 0pt\relax
  \else
    \refstepcounter{fig}\let\lakin@pend\relax
  \fi
  \lakinrtitle{#2}\par\nopagebreak\addvspace{4pt}%
  \def\lakin@short{#1}\ifx\lakin@short\@empty\else
    \addcontentsline{lof}{figure}{\protect\numberline{\thefig}{\protect\ignorespaces #1}}\fi
  {\centering\lakin@pend\lakitcaptionstyle\mbox{\lakitcaptionlabelstyle{\figurename~\en{\thefig}\lakitcaptionpunctuation}}\enspace #2\par}}
\makeatother

% ============================================================================
% LTR islands.  EVERY tikzpicture / quantikz MUST live inside one, or the whole
% figure's coordinates and Latin labels flip. (Hebrew node text still uses \h{}.)
% ============================================================================
% v3.74 (BLD-37): the \figcap hyperlink DESTINATION used to land at the caption
% line — BELOW the figure — so clicking a figure reference scrolled the figure
% off-screen (measured fleet-wide: 97/97 fig dests in one collection sat 19-54pt
% under their figure's ink; a 347pt figure was fully hidden after the jump).
% hypcap-style fix in two parts: (1) at the island's END, \lakin@grabfig grabs
% the just-typeset centered line with \lastbox, measures it, and puts it back
% with \nointerlineskip — the interline glue above the line was already
% contributed and is NOT re-added, so the page geometry is unchanged to the
% glyph (verified); (2) \figcap then steps the counter WITHOUT hyperref's
% automatic anchor (\H@refstepcounter = the pre-hyperref \refstepcounter) and
% emits the anchor itself, raised by the measured height — the dest lands ABOVE
% the figure top. A captionless island leaves numbering untouched (nothing is
% stepped at the island; the grab flag is cleared at every island begin, so a
% stale height from a captionless figure cannot leak into a later caption).
\newsavebox{\lakinfigbox}
\newdimen\lakinfigraise  % real \dimen: tex.setdimen cannot target a \newlength SKIP register (measured lua error)
% BLD-37 table side: a SEPARATE register for \tabcap. The two must not share —
% a captionless island INTENTIONALLY leaves \lakinfigraise set (the grab runs
% at every island END; only the NEXT island's begin clears it), so a shared
% register would hand a \tabcap after a captionless figure that stale FIGURE
% height. \lakintabraise is written only by \lakin@grabtab at \tabcap entry
% and consumed/reset inside \tabcap itself, so it can never go stale.
\newdimen\lakintabraise
\makeatletter
% v3.74b: measure WITHOUT touching the list. The first draft grabbed the box
% with \lastbox and put it back — and \lastbox DISCARDS the node's baseline
% shift, so a vertically-anchored quantikz circuit rose 2.57pt and everything
% after it sank 5.32pt (measured on the A/B glyph diff). LuaTeX lets us READ
% the tail node's dimensions without removing anything: geometry is identical
% by construction. \lakinfigraise=0pt is the "no figure measured" sentinel.
\def\lakin@grabfig{%
  \ifhmode
    \directlua{
      local n = tex.nest[tex.nest.ptr].tail
      local GLUE, PEN = node.id('glue'), node.id('penalty')
      while n and (n.id == GLUE or n.id == PEN) do n = n.prev end
      local HL, VL = node.id('hlist'), node.id('vlist')
      if n and (n.id == HL or n.id == VL) then
        tex.setdimen('global', 'lakinfigraise',
                     (n.height or 0) + (n.depth or 0))
      else
        tex.setdimen('global', 'lakinfigraise', 0)
      end}%
  \else
    \global\lakinfigraise 0pt\relax
  \fi}
% BLD-37 table side: \lakin@grabfig's twin, NOT a shared helper — the Lua body
% differs in exactly one place, the skip set. Between \end{tabular} and \tabcap
% an otherlanguage wrapper (patterns 3/5 of example_tables) leaves a \write
% whatsit and a dir node ON TOP of the interword glue (measured node dump:
% "whatsit/write <- dir <- glue(2.62) <- hlist"), so the table reader must
% also step over whatsit/dir markers — all zero-size — to reach the tabular.
% \lakin@grabfig deliberately does NOT get the wider skip set: its glue/penalty
% body is the one the v3.74 A/B proof froze, and widening it would newly
% measure (and so newly raise) anchors in island shapes that today take the
% plain path — an unproven behavior change on the FIG side. Anything that is
% not a box behind those markers still hits the 0pt sentinel and \tabcap falls
% back to the plain \refstepcounter.
% SECOND deliberate difference: this stores the box's HEIGHT, not height+depth.
% \figcap's anchor is emitted in the caption paragraph BELOW the figure, so it
% must climb the full extent; \tabcap's anchor sits ON the tabular's own line,
% whose baseline is the vcentered tabular's MIDPOINT — height alone reaches
% the top edge, and adding the depth was measured to overshoot by half the
% table (dests 80-135pt above their frame, INSIDE the preceding table's box).
\def\lakin@grabtab{%
  \ifhmode
    \directlua{
      local n = tex.nest[tex.nest.ptr].tail
      local GLUE, PEN = node.id('glue'), node.id('penalty')
      local WHAT, DIR = node.id('whatsit'), node.id('dir')
      while n and (n.id == GLUE or n.id == PEN
                   or n.id == WHAT or n.id == DIR) do n = n.prev end
      local HL, VL = node.id('hlist'), node.id('vlist')
      if n and (n.id == HL or n.id == VL) then
        tex.setdimen('global', 'lakintabraise', n.height or 0)
      else
        tex.setdimen('global', 'lakintabraise', 0)
      end}%
  \else
    \global\lakintabraise 0pt\relax
  \fi}
\makeatother
% ============================================================================
% v3.79 (BLD-37 table side, CORRECTION): put the anchor at the TOP OF THE BOX,
% instead of guessing the distance to it.
% ============================================================================
% v3.77 raised the \tabcap anchor by (height of the box at the tail of the open
% horizontal list) + 42pt. That is a MEASUREMENT taken BEFORE line breaking, and
% it is short by every term it cannot see. Fleet-measured on the 12 built
% collections (299 tbl.* dests, ownership established through each table's own
% CAPTION, never through "which frame encloses the dest"): 3 dests landed INSIDE
% their own tablebox, below the frame top —
%   computer-architecture       tbl.14.1  -46.4pt   three stacked tabulars
%   computer-architecture       tbl.30.5  -13.5pt   tabular 2.85pt over \linewidth
%   advanced-electronic-devices tbl.B.1   -59.2pt   tabular 2.56pt over \linewidth
% Two missing terms, both reproduced in a minimal document (probe P2/P3):
%   (1) VERTICAL LIST ALREADY CONTRIBUTED inside the box — a second and third
%       tabular, or a heading above the table. The tail read sees only the LAST
%       box, so 77.7pt of stacked tabulars is invisible to the raise.
%   (2) WHICH LINE THE ANCHOR LANDS ON. When the tabular alone overfills
%       \linewidth (459.6121pt in a tablebox), the interword glue after
%       \end{tabular} becomes the chosen breakpoint and \Hy@raisedlink is
%       emitted on the NEXT line — one (depth + lineskip) lower. That term does
%       not EXIST at \tabcap time: line breaking has not happened yet. No
%       measurement taken there can ever recover it. (Proof: probe P3, the same
%       table 2.84912pt too wide, fails at -13.5pt; probe P4, the identical
%       table with \tabcolsep 2pt so it fits, is correct at +32.0pt.)
% So the raise is not computed at \tabcap at all. It is CORRECTED at
% \end{tablebox}, from inside mdframed's own body vbox (endinnercode runs after
% \par\unskip, i.e. after line breaking, while tex.nest[tex.nest.ptr] is still
% that vbox's vertical list). There the anchor's true offset from the top of the
% box content is not estimated but WALKED, and the \Hy@raisedlink shift is set
% so the dest lands 42pt above it — which is 42pt - innertopmargin(9pt) -
% linewidth(0.9pt) = 32.0pt above the FRAME top, the exact position the 295
% already-correct dests have today. That is why this is a correction and not a
% re-placement: the good sites do not move.
%
% WHY THE GLYPH DIFF IS 0 BY CONSTRUCTION — the only thing written is the
% `shift` field of the hlist that \Hy@raisedlink already built, and the guard
% below refuses to write it unless that box is 0x0x0. A zero-size box carries no
% dimension into its parent's packing, the packing already happened, and shift
% is not an input to it. Nothing else in the list is touched: no node is added,
% removed, resized or reordered. Measured, not assumed: assets/example_tables.tex
% and the whole computer-architecture collection A/B at 0 differing pymupdf
% rawdict glyph records.
%
% WHY NOT the alternatives (all prototyped and measured before this one):
%   * accumulate the vertical list at \tabcap time — fixes term (1), CANNOT fix
%     term (2); 2 of the 3 real failures survive. A fix that leaves the majority
%     of the defect standing is the false all-clear this repo exists to avoid.
%   * \zref-savepos two-pass — correct on all shapes, but it makes the dest
%     POSITION aux-derived (DEFECTS BLD-1/2), a single-pass compile silently
%     ships an unraised anchor, and the bare y-difference has no page guard.
%   * emit the anchor at \begin{tablebox} and let \tabcap adopt it — needs the
%     tbl counter stepped at box OPEN to keep the dest NAME byte-identical; a
%     captionless tablebox then either shifts every later table number or leaves
%     an orphan dest that the next table duplicates (a HARD build gate). This
%     correction touches neither the counter nor the name: \tabcap's naming path
%     is v3.77's, unchanged, so resolve_names() is byte-identical (verified).
%   * suppressing the line break (a penalty before the glue) — that MOVES the
%     tabular: the break currently discards the interword glue, so forbidding it
%     re-centers the two overfull tables. Reflow, rejected.
%
% SCOPE. The correction is armed only by \tabcap ITSELF (\lakintabcapseen),
% cleared at every box open, so a captionless tablebox is a no-op and a \tabcap
% hosted OUTSIDE a tablebox (the fleet has exactly one, computer-architecture
% tbl.14.2 in an exbox) keeps v3.77's measured raise. Both fallbacks are
% therefore never worse than today; and because hyperref's AUTOMATIC anchor uses
% the same \Hy@raisedlink shape, the plain-\refstepcounter path (vmode entry, or
% prose between \end{tabular} and \tabcap — probe P6, -51.5pt) is corrected too.
% KNOWN LIMIT — a tablebox that SPLITS across pages. The offset is measured from
% the top of the UNSPLIT body, so a dest in segment 2 gets raised past its own
% segment. MEASURED on a forced split (probe: 40 filler lines + a table): the
% dest went from y=585.0 on page 2 (inside segment 2, at the caption — the old
% BLD-37 symptom) to y=-113.4, i.e. ABOVE the page top. Viewers clamp that to the
% page top, which happens to show segment 2 whole, but an out-of-page coordinate
% is not something to ship on purpose. Hence the \textheight guard in
% lakintab.fix: a body TALLER THAN THE TEXT BLOCK cannot fit on any page, so it
% must split — the correction is skipped and v3.77's behaviour stands unchanged
% (verified: the probe returns to y=585.0). The guard is deliberately EXACT, not
% a page-fitting heuristic: it is sound (a taller-than-page body always splits)
% but PARTIAL (a shorter body that starts low on a page could still split). That
% residue is unreachable in practice — v3.37's everyline + the caption's
% \nopagebreak glue move such boxes WHOLE instead (probe N7: a 22-row table
% pushed down by 430pt moved whole, 0 segments), and the fleet has 0 split
% tableboxes, re-counted after this change. Reimplementing mdframed's own
% page-fitting arithmetic to close the residue would be a heuristic with its own
% failure mode, which is the trade this repo has repeatedly refused.
\newif\iflakintabcapseen
\newif\iflakintabframed
% The Lua lives in ONE \directlua at load time, and is written with NO
% backslash, %, # or ~ : \directlua expands its argument like \edef, so an
% active ~ would expand to \nobreakspace and a % would eat the rest of the line
% (the second is how the first draft of this block died). "not (a == b)" stands
% in for Lua's inequality operator for the same reason.
\directlua{
  lakintab = { }
  lakintab.HL = node.id("hlist")
  lakintab.VL = node.id("vlist")
  lakintab.RU = node.id("rule")
  lakintab.GL = node.id("glue")
  lakintab.KE = node.id("kern")
  lakintab.WH = node.id("whatsit")
  lakintab.TM = node.id("temp")
  lakintab.WD = node.subtype("pdf_dest")
  function lakintab.d(v)
    if v == nil then return 0 end
    if v < -1000000000 then return 0 end
    return v
  end
  function lakintab.scan(list, y, horiz, parent, target)
    local x = list
    while x do
      local i = x.id
      if i == lakintab.HL or i == lakintab.VL then
        if horiz then
          local b = y + x.shift
          if i == lakintab.HL then
            lakintab.scan(x.head, b, true, x, target)
          else
            lakintab.scan(x.head, b - lakintab.d(x.height), false, x, target)
          end
        else
          if i == lakintab.HL then
            lakintab.scan(x.head, y + lakintab.d(x.height), true, x, target)
          else
            lakintab.scan(x.head, y, false, x, target)
          end
          y = y + lakintab.d(x.height) + lakintab.d(x.depth)
        end
      elseif i == lakintab.WH then
        if x.subtype == lakintab.WD then
          local id = x.dest_id
          if type(id) == "string" and string.sub(id, 1, 4) == "tbl." then
            if horiz and parent and parent.width == 0
               and parent.height == 0 and parent.depth == 0 then
              parent.shift = parent.shift + (target - y)
            end
          end
        end
      elseif not horiz then
        if i == lakintab.RU then
          y = y + lakintab.d(x.height) + lakintab.d(x.depth)
        elseif i == lakintab.GL or i == lakintab.KE then
          y = y + lakintab.d(x.width)
        end
      end
      x = x.next
    end
  end
  function lakintab.fix(target, maxh)
    local nst = tex.nest[tex.nest.ptr]
    if nst == nil then return end
    local h = nst.head
    if h == nil then return end
    if h.id == lakintab.TM then h = h.next end
    local tot = 0
    local x = h
    while x do
      local i = x.id
      if i == lakintab.HL or i == lakintab.VL or i == lakintab.RU then
        tot = tot + lakintab.d(x.height) + lakintab.d(x.depth)
      elseif i == lakintab.GL or i == lakintab.KE then
        tot = tot + lakintab.d(x.width)
      end
      x = x.next
    end
    if tot > maxh then return end
    lakintab.scan(h, 0, false, nil, target)
  end
}
% @-FREE NAMES ON PURPOSE: these two are named in the `tablebox` key list at
% \newmdenv above, which is tokenised OUTSIDE any \makeatletter region — a
% \lakin@... there would tokenise as \lakin followed by the letters @tabopen.
\newcommand{\lakintabboxopen}{\lakintabframedtrue\global\lakintabcapseenfalse}
\newcommand{\lakintabboxclose}{%
  \iflakintabcapseen
    \directlua{lakintab.fix(-\number\dimexpr\lakintabcushion\relax,
                            \number\dimexpr\textheight\relax)}%
    \global\lakintabcapseenfalse
  \fi\lakintabframedfalse}
% The cushion above the box's CONTENT top. 42pt puts the dest 32.0pt above the
% frame top (42pt - innertopmargin 9pt - linewidth 0.9pt), i.e. exactly where
% the 295 correct dests already are — the correction moves only the broken ones.
\newcommand{\lakintabcushion}{42pt}
\newenvironment{qcirc}{\par\medskip\global\lakinfigraise 0pt\relax\begin{otherlanguage}{english}\begin{center}}
  {\csname lakin@grabfig\endcsname\end{center}\end{otherlanguage}\par\medskip}
\newenvironment{tikzpic}{\par\medskip\global\lakinfigraise 0pt\relax\begin{otherlanguage}{english}\begin{center}}
  {\csname lakin@grabfig\endcsname\end{center}\end{otherlanguage}\par\medskip}

% ============================================================================
% Code listings.  ★★ v3.4 THE BOX-OVERFLOW FIX:
%   resetmargins=true   → the listing IGNORES the enclosing environment's
%                         indentation and measures from the PAGE. v3.4 stated
%                         this backwards ("measured from the ENCLOSING box"),
%                         and v3.37 then blanked in-box line numbers because
%                         the page-anchored listing put its hanging digits on
%                         top of the box rule. Measured v4.5.0: with the key
%                         FALSE the listing follows the box and the code moves
%                         from x=75.26 to x=97.74, clearing the rule at 72.07.
%                         The key now lives in one \lstset below instead of in
%                         each style, so a box can override it — codebox does.
%   framexleftmargin=0  → the frame (left bar) does not extend left of the code.
%   xleftmargin keeps a small indent so the line numbers have room INSIDE.
% asmblock wraps the listing in an LTR island (code is LTR).
% ============================================================================
\lstdefinestyle{asm}{
  basicstyle=\ttfamily\small,
  backgroundcolor=\color{cBlueBg!55},
  frame=leftline,framerule=2.5pt,rulecolor=\color{cBlue},
  numbers=left,numberstyle=\tiny\color{cGray},numbersep=8pt,
  xleftmargin=2.2em,framexleftmargin=0pt,framexrightmargin=0pt,
  breaklines=true,columns=fullflexible,keepspaces=true,
  showstringspaces=false,aboveskip=8pt,belowskip=8pt,
  % v3.36: without a comment DELIMITER the commentstyle below never fires —
  % listings has no `language` here, so it does not know '#' opens a comment,
  % and every asm listing rendered flat black (measured: a whole-page pixel
  % sweep found ZERO cGreen glyphs across 135 shipped listings). morecomment
  % gives it the delimiter; color-only change, glyph widths untouched.
  morecomment={[l]{\#}},
  commentstyle=\color{cGreen},
}
\lstdefinestyle{cstyle}{
  language=C,basicstyle=\ttfamily\small,
  backgroundcolor=\color{cGreenBg!55},
  frame=leftline,framerule=2.5pt,rulecolor=\color{cGreen},
  numbers=left,numberstyle=\tiny\color{cGray},numbersep=8pt,
  xleftmargin=2.2em,framexleftmargin=0pt,framexrightmargin=0pt,
  breaklines=true,columns=fullflexible,keepspaces=true,
  showstringspaces=false,keywordstyle=\color{cBlue}\bfseries,
  commentstyle=\color{cGray}\itshape,aboveskip=8pt,belowskip=8pt,
}
% v3.36: pystyle/htmlstyle — colored twins of cstyle for the Python/HTML
% listings that used to repeat a ~200-char options blob with NO color
% (keywords rendered bold-black, comments italic-black — measured; C alone
% was colored). Geometry is copied VERBATIM from the shipped blob (frame,
% margins, skips) so swapping a call site to style=pystyle moves NOTHING —
% only keywordstyle/commentstyle/stringstyle are added (bold stays bold,
% italic stays italic, mono metrics identical => zero reflow, proven by
% per-page word-bbox diff). breaklines/breakatwhitespace stay per-site.
% v3.78: numbers=left joins both styles, with the exact numberstyle/numbersep
% keys asm/cstyle carry — fleet-wide consistency (author decision: all four
% styles numbered; "asm/cstyle numbered, py/html not" was an accident of the
% v3.36 geometry freeze, not a design). The numbers hang LEFT of the code box
% (listings sets them with \llap, so no overfull hbox can arise and the code
% x is untouched — measured: code frozen, digits at code_x-11.8pt); inside a
% colored box they are blanked like everyone else's (v3.37 hook), and the
% clean variant restores them in its frameless boxes (v3.78, above).
\lstdefinestyle{pystyle}{
  language=Python,basicstyle=\ttfamily\small,
  columns=fullflexible,keepspaces=true,
  frame=single,rulecolor=\color{cTblLine},
  numbers=left,numberstyle=\tiny\color{cGray},numbersep=8pt,
  xleftmargin=0.5em,framexleftmargin=0pt,framexrightmargin=0pt,
  showstringspaces=false,aboveskip=8pt,belowskip=8pt,
  keywordstyle=\color{cBlue}\bfseries,
  commentstyle=\color{cGray}\itshape,
  stringstyle=\color{cGreen},
}
% v4.5.0: hoisted out of the four styles so an enclosing box can override it.
% Every listing still resolves to true unless something later says otherwise,
% so shipped output is unchanged (verified byte-for-byte on all eight PDFs).
\lstset{resetmargins=true}
\lstdefinestyle{htmlstyle}{
  language=HTML,basicstyle=\ttfamily\small,
  columns=fullflexible,keepspaces=true,
  frame=single,rulecolor=\color{cTblLine},
  numbers=left,numberstyle=\tiny\color{cGray},numbersep=8pt,
  xleftmargin=0.5em,framexleftmargin=0pt,framexrightmargin=0pt,
  showstringspaces=false,aboveskip=8pt,belowskip=8pt,
  keywordstyle=\color{cBlue}\bfseries,
  commentstyle=\color{cGray}\itshape,
  stringstyle=\color{cGreen},
}
\newenvironment{asmblock}{\par\begin{otherlanguage}{english}}{\end{otherlanguage}\par}

% <hllpdf:listing-smart-quotes>
% v4.21.0: four Unicode punctuation characters must enter listings' token
% buffer. Unregistered characters bypass it: a closing quote was printed
% BEFORE the pending word, and an opening quote BEFORE a pending brace.
% Extend the standard character table, then activate only these codepoints
% in the listing's local character-table scope. Do this BEFORE the package's
% delimiter/alsoother/literate hooks so explicit user mappings still win.
% This is not universal Unicode support, nor a conversion to ASCII quotes.
% See listings developer guide 15.3.1 and references/listing-smart-quotes.md.
\makeatletter
\lst@CCPut \lst@ProcessOther {"2018}{"2019}{"201C}{"201D}\z@
\lst@AddToHookAtTop{SelectCharTable}{%
  \lst@ifactivechars
    \catcode"2018=\active \catcode"2019=\active
    \catcode"201C=\active \catcode"201D=\active
  \fi}
\makeatother
% </hllpdf:listing-smart-quotes>

% <hllpdf:listing-technical-symbols>
% v4.30: bounded extension for common technical punctuation in LTR listings.
% Without character-table registration these glyphs bypass the token buffer
% and can move before the preceding word even inside quoted strings.
% Keep the original four-quote block separate for its standing mutation test.
% This is source-faithful printing, not syntax validation or general Unicode.
\makeatletter
\lst@CCPut \lst@ProcessOther
  {"2013}{"2014}{"2026}{"2190}{"2192}{"2194}{"21D2}
  {"2212}{"221E}{"2260}{"2264}{"2265}\z@
\lst@AddToHookAtTop{SelectCharTable}{%
  \lst@ifactivechars
    \catcode"2013=\active \catcode"2014=\active \catcode"2026=\active
    \catcode"2190=\active \catcode"2192=\active \catcode"2194=\active
    \catcode"21D2=\active \catcode"2212=\active \catcode"221E=\active
    \catcode"2260=\active \catcode"2264=\active \catcode"2265=\active
  \fi}
\makeatother
% </hllpdf:listing-technical-symbols>


\usepackage[hidelinks,unicode]{hyperref}   % LAST, after colors. unicode = Hebrew bookmarks.
% v3.62: THE ISOLATION MACROS CORRUPTED EVERY PDF BOOKMARK THEY APPEARED IN.
% \en is defined (line ~155) as
%     \foreignlanguage{english}{\language=\csname l@nohyphenation\endcsname #1}
% and in a heading that body is written verbatim to the .toc and handed to
% hyperref's \pdfstringdef. hyperref strips the two CONTROL SEQUENCES (\language,
% \l@nohyphenation) but the plain `=` CHARACTER token between them is not a control
% sequence, so it survives into the PDF string: \chapter{שער \en{CNOT} והצמדה}
% produced the bookmark "שער =CNOT והצמדה". The printed TOC and the heading on the
% page are CORRECT — only the navigation pane and the viewer window title are wrong,
% which is why it went unnoticed. Measured on the shipped fleet: 70 corrupted
% bookmarks across 5 PDFs, from 177 unprotected \en-in-heading sites. Same corruption
% reached /Title and /Subject via the \hypersetup{pdftitle=...} the docs recommend.
% This hook tells hyperref to use the plain argument for all three isolation macros
% when building a PDF string. It does NOT touch the .toc (the printed TOC is
% byte-unchanged) and it does NOT mask the math signal — a heading with $...$ still
% logs "removing math shift", so BIDI-15 detection is intact.
% BONUS, measured: it also makes \code{} in an UNSTARRED heading COMPILE. BLD-10
% recorded that as needing \texorpdfstring or a starred heading; the real root cause
% is hyperref's \pdfstringdef scan of \code's \ifmmode...\else...\fi body — which is
% exactly why the identically-bodied \texttt survived (hyperref pre-neutralizes
% \texttt in PDF strings) and why \DeclareRobustCommand could never help.
% v3.64: \def\\{ } added. A manual line break in a heading WELDS the two words
% together in the bookmark: \chapter{שורה ראשונה\\שורה שנייה} produced the outline
% entry 'שורה ראשונהשורה שנייה' — "ראשונהשורה" is not a word. hyperref drops the
% \\ and puts nothing in its place, so the break that looks like a space on the
% page is a zero-width join in the navigation pane. The spaced form
% 'מילה אחת \\ מילה שנייה' escapes with a harmless DOUBLE space ('אחת  מילה'),
% which is why this never got noticed. Both are measured on a real outline.
% 0 live instances in the corpus — but SKILL.md and content-style.md both RECOMMEND
% a manual heading break, so this is a primed trap rather than a dormant one, and a
% one-token preamble fix retires it for good instead of asking authors to remember.
% Mapping \\ to a space fixes the WELD, which is the actual defect (a nonexistent
% word in the navigation pane). It does not normalise whitespace: measured after the
% fix, the welded form becomes 'שורה ראשונה שורה שנייה' (correct) and the already-
% spaced form goes from two spaces to three. That is cosmetic — viewers render it as
% a wider gap, and no word is corrupted. A per-heading \texorpdfstring is still the
% way to get an exactly-tuned bookmark string.
\pdfstringdefDisableCommands{\def\en#1{#1}\def\h#1{#1}\def\doclang#1{#1}\def\code#1{#1}\def\codeline#1{#1}\def\\{ }}
% v3.70: RESTORE \url LINE-BREAKING. url.sty freezes its break penalties at LOAD
% time (\mathchardef\UrlBreakPenalty=\binoppenalty), and v3.45's global
% \binoppenalty=10000 (inline-math atomicity) therefore priced EVERY breakpoint
% inside \Url at 10000 = never break. Measured: a long \url in Hebrew prose is
% one unbroken line, ~100pt overfull, invading the margin band the ink gates
% deliberately tolerate. With the two lines below: 0 overfull, the URL breaks
% after '/', the RTL split is CORRECT (fragment order forward), and each line
% gets its own link annotation exactly covering its fragment (measured).
% Confined to \Url's group — inline-math atomicity (TYP-7) is untouched.
\mathchardef\UrlBreakPenalty=200
\mathchardef\UrlBigBreakPenalty=100
\ifx\lakitmainlanguage\hllpdfenglish
  \hypersetup{pdflang=en}
\else
  \hypersetup{pdflang=he}
\fi  % /Lang is the document base locale (not per-span accessibility tagging).
% Set a per-document \hypersetup{pdftitle={...}} too — see SKILL.md.

% v3.52 (round-2 audit): two RTL bibliography/cross-ref defects, both invisible to
% every gate and both affecting the FIRST author to use the feature.
% (a) \@biblabel's default [#1] brackets are bidi-neutral and mirror to ]1[ under
% the RTL bibliography base — every \thebibliography reference list shows ]1[.
% LTR-isolating the label restores [1]. (The render gate's mirror rule fingerprints
% round parens )1(, so square brackets slipped through.)
\makeatletter
\renewcommand*{\@biblabel}[1]{\foreignlanguage{english}{[#1]}}
% v3.54: the SAME neutral-bracket trap on the INLINE side — a multi-key \cite{a,b}
% under the RTL base lays its European-number list right-to-left, so [1, 2] renders
% "[2,1]" and [1, 2, 3] renders "[3,2,1]" (verified by glyph geometry — the reader
% sees the sources in REVERSE order; a note form \cite[p.5]{a} also reorders to
% "[p. 5,1]"). LTR-isolate the whole bracket group (number list AND optional note),
% exactly like \@biblabel above: [1, 2] / [1, 2, 3] / [1, p. 5] all render correct.
% Placed AFTER hyperref so it wins; hyperref hyperlinks each number inside #1 (via
% \@citex), which this wrapper leaves intact. Mirrors the LaTeX kernel \@cite shape.
% v4.9.1: directional isolation must NOT switch the language here. Babel's
% safe-active citation wrapper passes its key list through \bbl@tempa;
% \foreignlanguage reuses that scratch macro before the deferred citation
% loop expands it. The old line therefore cited the literal key "english".
% If such a key existed, even the undefined-reference gate stayed green.
% \babelsublr isolates the bracket group without selecting a locale or
% changing that key list. Preserve native links, key order and optional notes.
\def\@cite#1#2{\babelsublr{[{#1\if@tempswa , #2\fi}]}}
\makeatother

% v4.11.0: a reference inside align's \text inherited HEBdisptag=true,
% so its enclosing hbox mirrored the brackets despite a clean compile.
% Isolate the complete printed reference without selecting a Babel language.
% Also neutralize the display-only normalfont direction reset locally. A
% direction-only patch inverted 12 to 21 inside equation in the measured probe.
% Plain maketag is local to the reference: mathtools' display-side wrapper
% otherwise globally consumes manual_tag while formatting a nested reference.
% Preserve each original command (including the selective auxiliary write),
% current tag forms, native links, and showonlyrefs on/off transitions.
\makeatletter
\NewCommandCopy{\HEB@nativeeqref}{\eqref}
\newcommand{\HEB@eqrefisolate}[1]{%
  \babelsublr{\HEBdisptagfalse\let\bbl@eqnodir\relax\let\maketag@@@\ORIGmaketag@@@#1}}
\protected\def\eqref#1{\HEB@eqrefisolate{\HEB@nativeeqref{#1}}}
\ExplSyntaxOn
\NewCommandCopy{\HEB@nativeMTeqref}{\MT_extended_eqref:n}
\cs_set_protected:Npn \MT_extended_eqref:n #1
  {\HEB@eqrefisolate{\HEB@nativeMTeqref{#1}}}
\ExplSyntaxOff
\makeatother

% v4.11.0: custom caption counters are not part of hyperref's English
% locale, and Hebrew has no upstream hyperref locale in this toolchain.
% Restore the appropriate names at every Babel language switch, including
% return from a nested inline island. Do not localize the destination keys.
\makeatletter
\newcommand{\HEB@refnames}{%
\def\chapterautorefname{פרק}%
\def\sectionautorefname{סעיף}%
\def\subsectionautorefname{תת־סעיף}%
\def\subsubsectionautorefname{תת־תת־סעיף}%
\def\paragraphautorefname{פסקה}%
% v4.31: hyperref uses a distinct counter name at the deepest standard level.
\def\subparagraphautorefname{תת־פסקה}%
\def\equationautorefname{משוואה}%
\def\figureautorefname{איור}%
\def\tableautorefname{טבלה}%
\def\partautorefname{חלק}%
\def\appendixautorefname{נספח}%
\def\footnoteautorefname{הערת שוליים}%
\def\itemautorefname{סעיף}%
\def\Itemautorefname{סעיף}%
\def\figautorefname{איור}%
\def\tblautorefname{טבלה}%
\def\pageautorefname{עמוד}% 
}
\addto\extrashebrew{\HEB@refnames}
\newcommand{\ENG@refnames}{%
  \def\subparagraphautorefname{subparagraph}%
  \def\chapterautorefname{chapter}%
  \def\sectionautorefname{section}%
  \def\subsectionautorefname{subsection}%
  \def\subsubsectionautorefname{subsubsection}%
  \def\paragraphautorefname{paragraph}%
  \def\equationautorefname{Equation}%
  \def\figureautorefname{Figure}%
  \def\tableautorefname{Table}%
  \def\partautorefname{Part}%
  \def\appendixautorefname{Appendix}%
  \def\footnoteautorefname{footnote}%
  \def\itemautorefname{item}%
  \def\figautorefname{Figure}%
  \def\tblautorefname{Table}%
  \def\Itemautorefname{item}%
  \def\pageautorefname{page}%
}
\addto\extrasenglish{\ENG@refnames}
\ifx\lakitmainlanguage\hllpdfenglish\ENG@refnames\else\HEB@refnames\fi
\makeatother

% ============================================================================
% Optional running head (v3.39): \lakithead{<book title>} — OFF by default.
% A document opts in with ONE line after \input{preamble.tex}:
%   \lakithead{מבנה המחשב}
% Every content page then carries the series' orientation line: chapter name
% at the reading edge, book title opposite, thin rule. (fancyhdr's [L]/[R]
% are PHYSICAL; under the RTL page direction [L] lands at the RIGHT/reading
% edge — measured, consistent across the corpus.) Packaged like the topic
% rail: always loaded, inert without the call. This centralizes the 8-line
% block that 8 lakits used to copy by hand — the same duplication vector
% that once let a lakit preamble go ~100 lines stale. Class-aware: report
% fills the mark from \chapter, article from \section. NB: a titlepage env
% self-empties; front matter that sets \pagestyle{empty} must restore with
% \pagestyle{fancy} afterwards (NOT plain — that discards the head).
% ============================================================================
\usepackage{fancyhdr}
% v3.61: \chapter* left a STALE running head — pages of an UNNUMBERED chapter
% carried the PREVIOUS chapter's title. report.cls routes \chapter* through
% \@schapter, which never calls \chaptermark, so \leftmark keeps whatever the last
% numbered chapter left. Measured: with \chapter{פרק ממוספר} then
% \chapter*{מבוא לא ממוספר}, page 4 (inside the starred chapter) printed
% "פרק ממוספר". Invisible on the chapter's FIRST page (\thispagestyle{plain}), so it
% only shows from the second page on — the pages an author is least likely to check.
% An unnumbered front-matter chapter (\chapter*{מבוא}, \chapter*{הקדמה}) is everyday
% Hebrew-academic structure, and the corpus already carries HAND-WRITTEN \markboth
% lines right after \chapter* — the exact per-document duplication \lakithead exists
% to retire. Patched INSIDE \lakithead so the macro stays inert when unused; a
% document's own later \markboth still wins, so those hand-patches do not regress.
\makeatletter
% <hllpdf:starred-heading-scope>
\newif\iflakin@starmarkpatched
\newcommand{\lakistarmark}{%
  \iflakin@starmarkpatched\else
  % Dispatch on the actual class counter, not hyperref's shared internals.
  \ifdefined\c@chapter
    \let\HEBorigschapter\@schapter
    % Record only after the native heading has reached its page. A mark
    % before the native routine can advertise the next heading too early.
    \def\@schapter##1{\HEBorigschapter{##1}\chaptermark{##1}}%
  \else
    % @ssect receives layout arguments, NOT the section name or level.
    % Capture the name in @startsection without introducing a TeX group:
    % grouping sectioning would break the native after-heading state.
    \def\HEB@sectionkind{}%
    \def\HEB@mainsection{section}%
    \let\HEBorigstartsection\@startsection
    \def\@startsection##1##2##3##4##5##6{%
      \def\HEB@sectionkind{##1}%
      \HEBorigstartsection{##1}{##2}{##3}{##4}{##5}{##6}}%
    \let\HEBorigssect\@ssect
    \def\@ssect##1##2##3##4##5{%
      \HEBorigssect{##1}{##2}{##3}{##4}{##5}%
      % Lower starred headings must not replace the article's section mark.
      \ifx\HEB@sectionkind\HEB@mainsection
        \sectionmark{##5}%
      \fi
      \ignorespaces}%
  \fi
  \global\lakin@starmarkpatchedtrue
  \fi}
% </hllpdf:starred-heading-scope>
% <hllpdf:running-language>
% Choose the wrapper when the mark is written, not when a later page ships.
% English punctuation must remain LTR even after returning to Hebrew prose.
\newcommand{\HEB@runningmark}[1]{%
  \iflanguage{english}{\markboth{\en{#1}}{}}{\markboth{\h{#1}}{}}}
% </hllpdf:running-language>
\makeatother
% v4.2.1 (NAV): every collection already names itself in \lakithead — feed the
% PDF /Title from the same declaration, so the navigation gate's
% PDF_REQUIRE_TITLE has a universal source and no fleet member ships
% title-less by missing a second declaration site. Measured before this: 14
% of 16 collections shipped /Title-empty while every one carried a perfectly
% good \lakithead string (plain Hebrew in all 16 — no math, no formatting, so
% the PDF-string path is clean). An explicit title always wins: set it with
% \lakitpdftitle{...} anywhere before \lakithead, or with a plain
% \hypersetup{pdftitle=...} — the fallback detects that via hyperref's
% \@pdftitle, probe-verified \@empty-equal until the first explicit set.
\makeatletter
\newif\iflakit@pdftitleset
\newcommand{\lakitpdftitle}[1]{%
  \hypersetup{pdftitle={#1}}\global\lakit@pdftitlesettrue}
\newcommand{\lakithead}[1]{%
  \iflakit@pdftitleset\else
    \ifx\@pdftitle\@empty
      \hypersetup{pdftitle={#1}}\global\lakit@pdftitlesettrue
    \fi
  \fi
  \pagestyle{fancy}\fancyhf{}%
  \lakistarmark

  \fancyhead[L]{\small\textcolor{cGray}{\leftmark}}%
  % <hllpdf:fixed-head-language>
  % The fixed document title has its OWN language snapshot, independent of
  % the current section/chapter mark. Choose the branch NOW, not at shipout:
  % the page may be Hebrew by the time an English title is finally printed.
  \iflanguage{english}{%
    \fancyhead[R]{\small\textcolor{cGray}{\en{#1}}}%
  }{%
    \fancyhead[R]{\small\textcolor{cGray}{\h{#1}}}%
  }%
  % </hllpdf:fixed-head-language>
  \fancyfoot[C]{\thepage}%
  \renewcommand{\headrulewidth}{0.4pt}%
  \ifdefined\chaptermark
    \renewcommand{\chaptermark}[1]{\HEB@runningmark{##1}}%
  \else
    \renewcommand{\sectionmark}[1]{\HEB@runningmark{##1}}%
  \fi}
\makeatother

% v3.39.1: path-independent byte-reproducibility. LuaTeX derives the trailer
% /ID from the OUTPUT PATH, so two clones of the same source at different
% directories produced different sha256 (measured: the ONLY delta was the
% 62-byte /ID line). Bit 512 suppresses exactly that line — measured: Creator/
% Producer/CreationDate/ModDate and all content untouched. /ID is optional
% per ISO 32000 (required only for ENCRYPTED PDFs and PDF/A — neither is a
% target here). If a future workflow needs an ID (PDF/A, Acrobat revision
% matching), delete this line and it returns.
\pdfvariable suppressoptionalinfo \numexpr 512\relax

% v3.75: BOLD MATH IN BOLD HEADINGS. Heading text is \bfseries but math inside
% it stayed regular-weight CM — a hairline $V_T$ inside a heavy 24.8pt chapter
% title (measured: CMMI12 fl=6 next to FrankRuhlLibre-Bold fl=20; 155 live
% heading-with-math sites fleet-wide). \boldmath rides along wherever the class
% sets \bfseries for a heading; A/B measured: $V_T$ -> CMMIB10 fl=22, 0 missing
% characters (the bold CM faces ship with this TeX Live). Guarded per-macro so
% the article class (no \@makechapterhead) and future class changes stay safe.
% v4.0.1 (adoption fix): the SEARCH patterns above carry no parameter token.
% v4.0.0 made them specific by including the class's own #1/#2, and measured on
% TeX Live 2023 that never matches: a \patchcmd search pattern containing a
% parameter token cannot match while the call sits inside \@ifundefined{...}{}{...}
% (probed both #1 and ##1 -> FAIL; the same pattern at top level succeeds). Every
% one of these patches therefore failed, and because v4.0.0 also replaced the
% empty failure branch with \PackageError, the failure became FATAL: all four
% bundled example assets stopped building ("Cannot patch chapter-title math
% weight -> Fatal error, no output PDF"). The \PackageError branches are KEPT —
% a silent no-op here is what let v3.75's fix rot unnoticed — only the parameter
% tokens are dropped, so the patches match and a real class change still errors.
\makeatletter
\@ifundefined{@makechapterhead}{}{%
  \patchcmd{\@makechapterhead}{\huge\bfseries \@chapapp}
    {\huge\bfseries\boldmath \@chapapp}{}
    {\PackageError{hebrew-lualatex-pdf}{Cannot patch chapter-number math weight}{Unsupported document class}}%
  \patchcmd{\@makechapterhead}{\Huge \bfseries}
    {\Huge \bfseries\boldmath}{}
    {\PackageError{hebrew-lualatex-pdf}{Cannot patch chapter-title math weight}{Unsupported document class}}%
  \patchcmd{\@makeschapterhead}{\Huge \bfseries}
    {\Huge \bfseries\boldmath}{}
    {\PackageError{hebrew-lualatex-pdf}{Cannot patch starred chapter math weight}{Unsupported document class}}}
% v4.0.1: \part sits under a THREE-DEEP wrapper chain, so the class-based anchor
% is not in \@part at all. Measured with \show: hyperref rewrites \@part to
% `...\phantomsection\fi \H@old@part`, and nameref rewrites THAT to
% `\NR@gettitle{#1}\NR@part[{#1}]{#2}` — the class body ends up in \NR@part.
% v4.0.0 patched \@part and made failure fatal, which is why all seven bundled
% examples stopped building. \part is NOT a dead heading level here: 188 live uses
% across the fleet, so the bold-math fix is worth getting right rather than
% dropping. We walk the chain innermost-first and stop at the first hit. Failure is
% a WARNING, not an error: the chain is a property of hyperref/nameref versions, so
% a future package release must degrade to regular-weight part math — a cosmetic
% loss — instead of refusing to build. Chapter and section keep \PackageError,
% because those targets are in the CLASS and were verified to patch here.
\newif\iflakin@partnum \newif\iflakin@parttitle \newif\iflakin@sparttitle
\def\lakin@partnum#1#2{\iflakin@partnum\else
  \patchcmd{#1}{#2\bfseries \partname}{#2\bfseries\boldmath \partname}
    {\lakin@partnumtrue}{}\fi}
\def\lakin@parttitle#1#2{\iflakin@parttitle\else
  \patchcmd{#1}{#2 \bfseries}{#2 \bfseries\boldmath}{\lakin@parttitletrue}{}\fi}
\def\lakin@sparttitle#1#2{\iflakin@sparttitle\else
  \patchcmd{#1}{#2 \bfseries}{#2 \bfseries\boldmath}{\lakin@sparttitletrue}{}\fi}
\@ifundefined{c@chapter}{\def\lakin@pnum{\Large}\def\lakin@ptit{\huge}}%
                        {\def\lakin@pnum{\huge}\def\lakin@ptit{\Huge}}
\@ifundefined{NR@part}{}{\expandafter\lakin@partnum\expandafter\NR@part\expandafter{\lakin@pnum}%
                         \expandafter\lakin@parttitle\expandafter\NR@part\expandafter{\lakin@ptit}}
\@ifundefined{H@old@part}{}{\expandafter\lakin@partnum\expandafter\H@old@part\expandafter{\lakin@pnum}%
                            \expandafter\lakin@parttitle\expandafter\H@old@part\expandafter{\lakin@ptit}}
\@ifundefined{@part}{}{\expandafter\lakin@partnum\expandafter\@part\expandafter{\lakin@pnum}%
                       \expandafter\lakin@parttitle\expandafter\@part\expandafter{\lakin@ptit}}
\@ifundefined{NR@spart}{}{\expandafter\lakin@sparttitle\expandafter\NR@spart\expandafter{\lakin@ptit}}
\@ifundefined{H@old@spart}{}{\expandafter\lakin@sparttitle\expandafter\H@old@spart\expandafter{\lakin@ptit}}
\@ifundefined{@spart}{}{\expandafter\lakin@sparttitle\expandafter\@spart\expandafter{\lakin@ptit}}
\iflakin@partnum\else\PackageWarning{hebrew-lualatex-pdf}{Part-number math stays
  regular weight: no known \noexpand\part wrapper matched (hyperref/nameref
  version change). Cosmetic only}\fi
\iflakin@parttitle\else\PackageWarning{hebrew-lualatex-pdf}{Part-title math stays
  regular weight: no known \noexpand\part wrapper matched. Cosmetic only}\fi
% v4.12.0: THE PART NUMBER IN THE TOC WAS NOT DIRECTIONALLY SEPARATED FROM THE
% TITLE. report.cls/article.cls compose the part entry as ONE text stream --
% \addcontentsline{toc}{part}{\thepart\hspace{1em}#1} -- while EVERY other level
% goes through \numberline, which puts the number in an \hbox. A box is a bidi
% OBJECT and cannot join a neighbouring run; a bare "IV" can. So in the RTL TOC
% paragraph a Latin number and a Latin-INITIAL title merge into a SINGLE LTR run
% and are laid out left-to-right as a unit: the number lands on the wrong side of
% the row, and a mixed title is scrambled outright. Measured on \part{MCP ...}:
% the printed order was "<tail> V MCP" -- title fragment, number, title fragment.
% The chapter/section rows one line below, with the SAME Latin-initial titles,
% were correct on the same page. That contrast is the whole diagnosis.
% \en{} IS NOT A WORKAROUND, measured both ways on one fixture: it fixes a
% PURE-Latin title (\part{\en{Hooks}} renders correctly) and does NOT fix a MIXED
% one (\part{\en{MCP} <hebrew>} stays scrambled). babel's language switch changes
% font and hyphenation, not the bidi run structure -- so the defect is not
% reachable from the source side at all, which is why it belongs here.
% Boxing the number is the entire fix and is a STRICT NO-OP where the bug never
% bit: a Hebrew-titled document compiles BYTE-IDENTICAL before and after (md5,
% 3 clean runs), because there the neutral glue already sat between an L run and
% an R run. The PDF outline is byte-identical too -- hyperref's \pdfstringdef
% neutralises \mbox, so bookmark strings keep their part numbers.
% \numberline{\thepart} would also fix the direction, but it re-geometries every
% part row in the fleet and DROPS the number from the part bookmark (hyperref
% strips \numberline). \mbox changes nothing that was already right.
% The search pattern deliberately carries no # token: etoolbox cannot match a
% parameter token in a patch pattern (measured -- all three chain members
% reported "not patchable" while the pattern included #1). "\thepart \hspace
% {1em}" is brace-balanced and occurs exactly once in the wrapper body.
\newif\iflakin@parttocnum
\def\lakin@parttocnum#1{\iflakin@parttocnum\else
  \patchcmd{#1}{\thepart \hspace {1em}}{\mbox {\thepart }\hspace {1em}}
    {\lakin@parttocnumtrue}{}\fi}
\@ifundefined{NR@part}{}{\lakin@parttocnum\NR@part}
\@ifundefined{H@old@part}{}{\lakin@parttocnum\H@old@part}
\@ifundefined{@part}{}{\lakin@parttocnum\@part}
\iflakin@parttocnum\else\PackageWarning{hebrew-lualatex-pdf}{Part TOC number is
  not boxed: no known \noexpand\part wrapper matched. A part title that BEGINS
  with Latin will print its number on the wrong side of its TOC row}\fi
\patchcmd{\section}{\bfseries}{\bfseries\boldmath}{}
  {\PackageError{hebrew-lualatex-pdf}{Cannot patch section math weight}{Unsupported document class}}
\patchcmd{\subsection}{\bfseries}{\bfseries\boldmath}{}
  {\PackageError{hebrew-lualatex-pdf}{Cannot patch subsection math weight}{Unsupported document class}}
\patchcmd{\subsubsection}{\bfseries}{\bfseries\boldmath}{}
  {\PackageError{hebrew-lualatex-pdf}{Cannot patch subsubsection math weight}{Unsupported document class}}
\patchcmd{\paragraph}{\bfseries}{\bfseries\boldmath}{}{}
\patchcmd{\subparagraph}{\bfseries}{\bfseries\boldmath}{}{}
\@ifundefined{l@chapter}{}{%
  \patchcmd{\l@chapter}{\bfseries}{\bfseries\boldmath}{}{}}
\makeatother
\setcounter{tocdepth}{\lakittocdepth}
% <hllpdf:measured-toc>
% v4.16.0: Measure actual number labels, not a guessed digit count. Read only
% canonical contentsline records; never execute titles or addtocontents code
% during the measuring pass. Widths are local to this individual list and are
% available before its first row: no per-row drift and no extra aux pass.
\makeatletter
\newread\hllpdf@tocstream
\newdimen\hllpdf@toc@chapter
\newdimen\hllpdf@toc@section
\newdimen\hllpdf@toc@subsection
\newdimen\hllpdf@toc@subsubsection
\newdimen\hllpdf@toc@paragraph
\newdimen\hllpdf@toc@subparagraph
\newdimen\hllpdf@toc@figure
\newdimen\hllpdf@toc@table
\newdimen\hllpdf@tocpage
\newdimen\hllpdf@toc@sectionindent
\newdimen\hllpdf@toc@subsectionindent
\newdimen\hllpdf@toc@subsubsectionindent
\newdimen\hllpdf@toc@paragraphindent
\newdimen\hllpdf@toc@subparagraphindent
\def\hllpdf@tocdefaults{%
  \hllpdf@toc@chapter=2em
  \hllpdf@toc@section=3.2em
  \hllpdf@toc@subsection=3.9em
  \@ifundefined{c@chapter}{%
    \hllpdf@toc@subsubsection=3.2em
    \hllpdf@toc@paragraph=4.1em
    \hllpdf@toc@subparagraph=5em
  }{%
    \hllpdf@toc@subsubsection=4.1em
    \hllpdf@toc@paragraph=5em
    \hllpdf@toc@subparagraph=6em
  }%
  \hllpdf@toc@figure=2.3em
  \hllpdf@toc@table=2.3em
  \hllpdf@tocpage=\@pnumwidth\relax
}
\def\hllpdf@tocindents{%
  % Preserve the established shallow layout when its slots already fit.
  \hllpdf@toc@sectionindent=1.5em
  \@tempdima=\hllpdf@toc@chapter \advance\@tempdima by -2em
  \advance\hllpdf@toc@sectionindent by \@tempdima
  \hllpdf@toc@subsectionindent=\dimexpr
    \hllpdf@toc@sectionindent+\hllpdf@toc@section\relax
  \hllpdf@toc@subsubsectionindent=\dimexpr
    \hllpdf@toc@subsectionindent+\hllpdf@toc@subsection\relax
  \hllpdf@toc@paragraphindent=\dimexpr
    \hllpdf@toc@subsubsectionindent+\hllpdf@toc@subsubsection\relax
  \hllpdf@toc@subparagraphindent=\dimexpr
    \hllpdf@toc@paragraphindent+\hllpdf@toc@paragraph\relax
}
\def\hllpdf@tocfont@chapter{\bfseries\boldmath}
\def\hllpdf@tocfont@part{\large\bfseries\boldmath}
\def\hllpdf@tocentryfont{%
  \ifcsname hllpdf@tocfont@\hllpdf@toclevel\endcsname
    \csname hllpdf@tocfont@\hllpdf@toclevel\endcsname
  \fi}
\long\def\hllpdf@tocdiscard#1\hllpdf@tocstop{}
\def\hllpdf@tocscanrecord{%
  \@ifnextchar\contentsline{\hllpdf@tocscanline}{%
    \@ifnextchar\babel@toc{\hllpdf@tocscanlanguage}{\hllpdf@tocdiscard}}}
\long\def\hllpdf@tocscanlanguage\babel@toc#1#2#3\hllpdf@tocstop{%
  % Babel records are declarations, not arbitrary code. Track the language
  % without executing the original directive or changing the surrounding list.
  \def\hllpdf@toclanguage{#1}}
\long\def\hllpdf@tocscanline\contentsline#1#2#3#4#5\hllpdf@tocstop{%
  \def\hllpdf@toclevel{#1}%
  \ifcsname hllpdf@toc@#1\endcsname
    \hllpdf@tocscantitle#2\hllpdf@tocstop
  \fi
  % The page-number column needs its own measurement (including bold chapter
  % rows). Its extra clearance separates the page label from dot leaders.
  \setbox\@tempboxa=\hbox{\foreignlanguage{\hllpdf@toclanguage}{\hllpdf@tocentryfont #3\kern.3em}}%
  \ifdim\wd\@tempboxa>\hllpdf@tocpage
    \hllpdf@tocpage=\wd\@tempboxa
  \fi}
\def\hllpdf@tocscantitle{%
  \@ifnextchar\numberline{\hllpdf@tocscanlabel}{\hllpdf@tocdiscard}}
\long\def\hllpdf@tocscanlabel\numberline#1#2\hllpdf@tocstop{%
  \setbox\@tempboxa=\hbox{\foreignlanguage{\hllpdf@toclanguage}{\hllpdf@tocentryfont #1\kern.5em}}%
  \ifdim\wd\@tempboxa>\csname hllpdf@toc@\hllpdf@toclevel\endcsname
    \csname hllpdf@toc@\hllpdf@toclevel\endcsname=\wd\@tempboxa
  \fi}
\def\hllpdf@tocread{%
  \ifeof\hllpdf@tocstream\else
    \read\hllpdf@tocstream to\hllpdf@tocrecord
    \expandafter\hllpdf@tocscanrecord\hllpdf@tocrecord\hllpdf@tocstop
    \expandafter\hllpdf@tocread
  \fi}
\def\hllpdf@tocmeasure#1{%
  \makeatletter
  \edef\hllpdf@toclanguage{\languagename}%
  \hllpdf@tocdefaults
  \openin\hllpdf@tocstream=\jobname.#1\relax
  \hllpdf@tocread
  \closein\hllpdf@tocstream
  \hllpdf@tocindents
  \edef\@pnumwidth{\the\hllpdf@tocpage}%
  \@tempdima=\dimexpr\hllpdf@tocpage+1em\relax
  \ifdim\@tempdima>\@tocrmarg\relax
    \edef\@tocrmarg{\the\@tempdima}%
  \fi}
\let\hllpdf@originalstarttoc\@starttoc
\def\@starttoc#1{%
  \begingroup
    \def\hllpdf@tocext{#1}%
    \def\hllpdf@tocwanted{toc}%
    \ifx\hllpdf@tocext\hllpdf@tocwanted \hllpdf@tocmeasure{#1}\else
      \def\hllpdf@tocwanted{lof}%
      \ifx\hllpdf@tocext\hllpdf@tocwanted \hllpdf@tocmeasure{#1}\else
        \def\hllpdf@tocwanted{lot}%
        \ifx\hllpdf@tocext\hllpdf@tocwanted \hllpdf@tocmeasure{#1}\fi
      \fi
    \fi
    \hllpdf@originalstarttoc{#1}%
  \endgroup}
\hllpdf@tocdefaults
\hllpdf@tocindents
\@ifundefined{l@chapter}{}{%
  \patchcmd{\l@chapter}{\setlength\@tempdima{1.5em}}%
    {\setlength\@tempdima{\hllpdf@toc@chapter}}{}%
    {\PackageError{hebrew-lualatex-pdf}{Cannot size chapter TOC numbers}%
       {Unsupported contents layout: the chapter number slot was not found}}}
\renewcommand*\l@section{\@dottedtocline{1}%
  {\hllpdf@toc@sectionindent}{\hllpdf@toc@section}}
\renewcommand*\l@subsection{\@dottedtocline{2}%
  {\hllpdf@toc@subsectionindent}{\hllpdf@toc@subsection}}
\renewcommand*\l@subsubsection{\@dottedtocline{3}%
  {\hllpdf@toc@subsubsectionindent}{\hllpdf@toc@subsubsection}}
\renewcommand*\l@paragraph{\@dottedtocline{4}%
  {\hllpdf@toc@paragraphindent}{\hllpdf@toc@paragraph}}
\renewcommand*\l@subparagraph{\@dottedtocline{5}%
  {\hllpdf@toc@subparagraphindent}{\hllpdf@toc@subparagraph}}
\renewcommand*\l@figure{\@dottedtocline{1}{1.5em}{\hllpdf@toc@figure}}
\renewcommand*\l@table{\@dottedtocline{1}{1.5em}{\hllpdf@toc@table}}
\makeatother
% </hllpdf:measured-toc>
\setcounter{secnumdepth}{\lakitsecnumdepth}
% Public preamble-only caption-numbering policy. The historical within-major
% default is untouched unless this command is called. "continuous" removes the
% reset edge as well as the printed prefix, so Figure 1/2 and Table 1/2 stay
% unique across sections/chapters. Returning to "within-major" restores the
% original class-aware displays and native destination conventions. The policy
% never replaces \figcap/\tabcap or their measured link-placement mechanism.
\makeatletter
\let\lakin@majorfig\thefig
\let\lakin@majortbl\thetbl
\let\lakin@majorHfig\theHfig
\let\lakin@majorHtbl\theHtbl
\newcommand{\lakitcaptionnumbering}[1]{%
  \edef\lakin@numbering{#1}%
  \def\lakin@continuous{continuous}\def\lakin@withinmajor{within-major}%
  \ifx\lakin@numbering\lakin@continuous
    \@ifundefined{c@chapter}{%
      \@removefromreset{fig}{section}\@removefromreset{tbl}{section}%
    }{\@removefromreset{fig}{chapter}\@removefromreset{tbl}{chapter}}%
    \renewcommand{\thefig}{\arabic{fig}}\renewcommand{\thetbl}{\arabic{tbl}}%
    \def\theHfig{\arabic{fig}}\def\theHtbl{\arabic{tbl}}%
  \else\ifx\lakin@numbering\lakin@withinmajor
    \@ifundefined{c@chapter}{%
      \@removefromreset{fig}{section}\@removefromreset{tbl}{section}%
      \@addtoreset{fig}{section}\@addtoreset{tbl}{section}%
    }{%
      \@removefromreset{fig}{chapter}\@removefromreset{tbl}{chapter}%
      \@addtoreset{fig}{chapter}\@addtoreset{tbl}{chapter}}%
    \let\thefig\lakin@majorfig\let\thetbl\lakin@majortbl
    \let\theHfig\lakin@majorHfig\let\theHtbl\lakin@majorHtbl
  \else
    \PackageError{hebrew-english-latex-studio}{Unknown caption-numbering policy '#1'}%
      {Use continuous or within-major before begin{document}.}%
  \fi\fi}
\@onlypreamble\lakitcaptionnumbering
\makeatother
% Architecture readiness sentinel. Optional runtime layers require the complete
% core, so the marker is defined only after every mechanism/baseline definition.
% <hllpdf:architecture-ready>
\def\hllpdfcoreloaded{1}
% </hllpdf:architecture-ready>

% v4.22: registered file resources, not an assertion about embedded subset bytes.
% The log framing avoids long-path wrapping and preserves UTF-8 as hex.
% <hllpdf:font-inventory>
\directlua{
  local function inventory()
    local names = {}
    local unnamed = 0
    for id, item in font.each() do
      local path = item.filename
      if path and not (path == "") then names[path] = true else unnamed = unnamed + 1 end
    end
    local sorted = {}
    for path in pairs(names) do table.insert(sorted,path) end
    table.sort(sorted)
    texio.write_nl("log","HLLFONT|BEGIN|1")
    for index,path in ipairs(sorted) do
      local hex = path:gsub(".",function(c) return string.format(string.char(37) .. "02x",string.byte(c)) end)
      local count = math.ceil(string.len(hex)/48)
      texio.write_nl("log","HLLFONT|FILE|" .. index .. "|" .. count)
      for offset=1,string.len(hex),48 do texio.write_nl("log","HLLFONT|HEX|" .. hex:sub(offset,offset+47)) end
    end
    texio.write_nl("log","HLLFONT|END|" .. rawlen(sorted) .. "|" .. unnamed)
    texio.write_nl("log","")
  end
  luatexbase.add_to_callback("finish_pdffile",inventory,"hllpdf.runtime.fonts")
}
% </hllpdf:font-inventory>
     % (or paste these lines at the top)
%   \begin{document}  ...  \end{document}
%
% Compile with:  lualatex -interaction=nonstopmode doc.tex
%   Repeat until two consecutive PDFs are byte-identical (scripts/build.sh
%   does this) — a TOC document converges only on pass 3+; after two passes
%   the TOC page numbers are STALE (measured; see SKILL.md step 3).
% Font binaries are external dependencies, pinned by assets/fonts/SHA256SUMS.
% scripts/setup_fonts.sh imports a verified local installation or prior skill.
%
% PUBLIC DOCUMENT-SHELL / LAYOUT CONTRACT (introduced v4.3.1; architecture-owned v4.7.1). These hooks may be
% defined BEFORE % ============================================================================
% preamble.tex — tested LuaLaTeX preamble for Hebrew + math + figures PDFs (v4.32.0)
% Battle-tested: compiled an 80-page Hebrew quantum-computing summary AND a
% 197-page Hebrew computer-architecture anthology (TikZ diagrams, code listings,
% colored boxes) — both with 0 missing characters.
%
% Usage in your document:
%   \input{preamble.tex}     % (or paste these lines at the top)
%   \begin{document}  ...  \end{document}
%
% Compile with:  lualatex -interaction=nonstopmode doc.tex
%   Repeat until two consecutive PDFs are byte-identical (scripts/build.sh
%   does this) — a TOC document converges only on pass 3+; after two passes
%   the TOC page numbers are STALE (measured; see SKILL.md step 3).
% Font binaries are external dependencies, pinned by assets/fonts/SHA256SUMS.
% scripts/setup_fonts.sh imports a verified local installation or prior skill.
%
% PUBLIC DOCUMENT-SHELL / LAYOUT CONTRACT (introduced v4.3.1; architecture-owned v4.7.1). These hooks may be
% defined BEFORE \input{preamble.tex}; if omitted, the historical canonical
% defaults below are exact.  They deliberately control only the document shell
% and page layout.  RTL/math semantics stay in this preamble, while visual
% presentation themes (for example clean_theme.tex) load AFTER it.
%
%   \def\lakitdocumentclass{article}
%   \def\lakitclassoptions{11pt,a4paper}
%   \def\lakitgeometryoptions{margin=1.8cm,headheight=14pt}
%   \def\lakitlinespread{1.15}
%   \def\lakitparagraphstyle{block} % letters/briefs; indented for continuous prose
%   \def\lakittocdepth{1}
%   \def\lakitsecnumdepth{2}
%   \input{preamble.tex}
%
% Supported class contract: report or article.  Alternate paper sizes must also
% be declared to build.sh through EXPECTED_PAPER; the PDF gate never infers them.
% Do NOT use these hooks as a theme mechanism: theme overlays must not own page
% geometry, class selection, or direction semantics. A declared presentation-policy
% overlay may intentionally change numbering visibility after the preamble.
% ============================================================================
% Runtime architecture guard: generated from scripts/architecture_contract.py.
% The readiness sentinel remains delayed until the END of this file.
% <hllpdf:architecture-guard>
\ifdefined\hllpdfcoreloaded
  \errmessage{hllpdf: preamble.tex loaded more than once}
  \expandafter\endinput
\fi
% </hllpdf:architecture-guard>

\ifdefined\lakitdocumentclass\else\def\lakitdocumentclass{report}\fi
\ifdefined\lakitclassoptions\else\def\lakitclassoptions{12pt,a4paper}\fi
\ifdefined\lakitgeometryoptions\else\def\lakitgeometryoptions{margin=2.3cm,headheight=14pt}\fi
\ifdefined\lakitlinespread\else\def\lakitlinespread{1.3}\fi
\ifdefined\lakittocdepth\else\def\lakittocdepth{1}\fi
\ifdefined\lakitsecnumdepth\else\def\lakitsecnumdepth{2}\fi
% Paragraph organization belongs to the document layout, independently of its
% language and theme. Preserve historical class dimensions unless explicitly
% selecting block paragraphs; never suppress indentation with an everypar hack.
\ifdefined\lakitparagraphstyle\else\def\lakitparagraphstyle{indented}\fi
\edef\lakitparagraphstyle{\lakitparagraphstyle}
\def\hllpdfparagraphblock{block}
\def\hllpdfparagraphindented{indented}
\ifx\lakitparagraphstyle\hllpdfparagraphblock\else
  \ifx\lakitparagraphstyle\hllpdfparagraphindented\else
    \errmessage{hllpdf: paragraph style must be block or indented}
  \fi
\fi
% v4.3.1: the public class hook has a deliberately narrow support contract.
% Before this guard, \def\lakitdocumentclass{book} silently compiled even though
% every heading/TOC regression test targets report/article; an unsupported shell
% could therefore look "supported" until a subtle class-specific defect shipped.
% Fail before \documentclass rather than letting an undocumented combination enter
% the mechanism core. Plain-TeX comparison is used because LaTeX packages are not
% available yet at this point.
\edef\lakitvalidatedclass{\lakitdocumentclass}
\def\lakitreportclass{report}
\def\lakitarticleclass{article}
\ifx\lakitvalidatedclass\lakitreportclass
\else\ifx\lakitvalidatedclass\lakitarticleclass
\else
  \errmessage{hllpdf: unsupported class \lakitdocumentclass; use report or article}
\fi\fi
\let\lakitvalidatedclass\relax
\let\lakitreportclass\relax
\let\lakitarticleclass\relax
% Language intent is independent of document class, theme and content domain.
% The two dominant-language profiles differ in authoring intent, not direction.
% A balanced bilingual document must explicitly choose a navigation/base locale.
% <hllpdf:language-profile>
\providecommand{\lakitlanguageprofile}{hebrew}
\edef\lakitlanguageprofile{\lakitlanguageprofile}
\def\hllpdfhebrew{hebrew}\def\hllpdfenglish{english}
\def\hllpdfhebrewmixed{hebrew-with-english}
\def\hllpdfenglishmixed{english-with-hebrew}
\def\hllpdfbilingual{bilingual}
\let\hllpdfexpectedlanguage\relax
\ifx\lakitlanguageprofile\hllpdfhebrew
  \let\hllpdfexpectedlanguage\hllpdfhebrew
\else\ifx\lakitlanguageprofile\hllpdfhebrewmixed
  \let\hllpdfexpectedlanguage\hllpdfhebrew
\else\ifx\lakitlanguageprofile\hllpdfenglish
  \let\hllpdfexpectedlanguage\hllpdfenglish
\else\ifx\lakitlanguageprofile\hllpdfenglishmixed
  \let\hllpdfexpectedlanguage\hllpdfenglish
\else\ifx\lakitlanguageprofile\hllpdfbilingual
  \ifdefined\lakitmainlanguage\else
    \errmessage{hllpdf: bilingual profile requires explicit lakitmainlanguage}
  \fi
\else
  \errmessage{hllpdf: unsupported language profile \lakitlanguageprofile}
\fi\fi\fi\fi\fi
\ifdefined\lakitmainlanguage
  \edef\lakitmainlanguage{\lakitmainlanguage}
\else
  \let\lakitmainlanguage\hllpdfexpectedlanguage
\fi
\ifx\lakitmainlanguage\hllpdfhebrew\else
  \ifx\lakitmainlanguage\hllpdfenglish\else
    \errmessage{hllpdf: main language must be hebrew or english}
  \fi
\fi
\ifx\hllpdfexpectedlanguage\relax\else
  \ifx\lakitmainlanguage\hllpdfexpectedlanguage\else
    \errmessage{hllpdf: main language conflicts with language profile}
  \fi
\fi
\typeout{HLLLANG|\lakitlanguageprofile|\lakitmainlanguage}
% </hllpdf:language-profile>
\documentclass[\lakitclassoptions]{\lakitdocumentclass}
\usepackage[no-math]{fontspec}        % (1) ★ no-math — THE critical line (see recipe.md)
% layout=lists asks Babel to adapt list parshape when paragraph and page
% directions differ. Without it, minority-language list bodies lose their
% indentation and labels hang outside the text block, despite a clean log.
% The upstream mechanism preserves ordinary same-direction list geometry.
\usepackage[bidi=basic,layout=lists]{babel}
\ifx\lakitmainlanguage\hllpdfenglish
  \babelprovide[import]{hebrew}
  \babelprovide[import,main]{english}
\else
  \babelprovide[import,main]{hebrew}
  \babelprovide[import]{english}
\fi
% v4.22: bind canonical families and fallbacks to the selected installation.
% Lookup by basename or family alone is not proof of the file actually used.
% <hllpdf:font-binding>
\directlua{
  local dir = os.getenv("HEBREW_LATEX_FONT_DIR")
  if not dir or dir == "" then
    local data = os.getenv("XDG_DATA_HOME")
    if not data or data == "" or not (data:sub(1,1) == "/") then
      data = (os.getenv("HOME") or "") .. "/.local/share"
    end
    dir = data .. "/fonts/hebrew-lualatex-pdf"
  end
  if not (dir:sub(1,1) == "/") then dir = lfs.currentdir() .. "/" .. dir end
  if dir:find("[",1,true) or dir:find("]",1,true) or dir:find(string.char(0),1,true)
     or dir:find(string.char(10),1,true) or dir:find(string.char(13),1,true) then
    tex.error("hllpdf: unsupported character in font directory")
  end
  if not (lfs.attributes(dir,"mode") == "directory") then
    tex.error("hllpdf: font directory is missing")
  end
  hllpdf_font_dir = dir:gsub("/+$", "") .. "/"
  tex.sprint(string.char(92) .. "def" .. string.char(92) .. "hllpdffontdir{")
  tex.sprint(-2,hllpdf_font_dir)
  tex.sprint("}")
  luaotfload.add_fallback("hebfb",{"[" .. hllpdf_font_dir .. "DejaVuSerif.ttf]:mode=node;"})
  luaotfload.add_fallback("hebfbbf",{"[" .. hllpdf_font_dir .. "DejaVuSerif-Bold.ttf]:mode=node;"})
}
% </hllpdf:font-binding>
% v3.22: Frank Ruhl Libre ships NO italic face. Instead of silently-upright
% \emph (v<=3.20) or bold-as-emphasis (v3.21), synthesize a slanted face with
% FakeSlant — real, VISIBLE italic emphasis with correct \emph nesting
% semantics (inner \emph toggles back upright). Verified empirically at
% 300dpi: slant clearly visible, 0 errors, 0 missing glyphs.
\babelfont{rm}[Path={\hllpdffontdir},RawFeature={fallback=hebfb},
  BoldFont={FrankRuhlLibre-Bold.ttf},
  BoldFeatures={RawFeature={fallback=hebfbbf}},   % v3.56 — see hebfbbf above
  ItalicFont={FrankRuhlLibre-Regular.ttf}, ItalicFeatures={FakeSlant=0.22},
  % Bold-italic uses the weight-matched fallback and a synthetic slant.
  BoldItalicFont={FrankRuhlLibre-Bold.ttf}, BoldItalicFeatures={RawFeature={fallback=hebfbbf},FakeSlant=0.22}
  ]{FrankRuhlLibre-Regular.ttf}
% Sans fallbacks preserve family and weight for the few missing codepoints.
\directlua{luaotfload.add_fallback("hebfbsans",{"[" .. hllpdf_font_dir .. "DejaVuSans.ttf]:mode=node;"})}
\directlua{luaotfload.add_fallback("hebfbsansbf",{"[" .. hllpdf_font_dir .. "DejaVuSans-Bold.ttf]:mode=node;"})}
\babelfont{sf}[Path={\hllpdffontdir},RawFeature={fallback=hebfbsans},
  BoldFont={Heebo-Bold.ttf},
  BoldFeatures={RawFeature={fallback=hebfbsansbf}}]{Heebo-Regular.ttf}
% v3.24: Scale=MatchLowercase — DejaVu Sans Mono has an outsized x-height, so
% unscaled inline code visually SHOUTS next to the delicate Frank Ruhl body
% (measured on a 197-page build: scaled, a code atom sits in the line at the
% body's weight, and the document tightened by ~5%). Applies to listings too.
\babelfont{tt}[Path={\hllpdffontdir},RawFeature={fallback=hebfb},Scale=MatchLowercase,
  BoldFont={DejaVuSansMono.ttf},
  BoldFeatures={RawFeature={fallback=hebfb},FakeBold=1.5},
  ItalicFont={DejaVuSansMono.ttf},
  ItalicFeatures={RawFeature={fallback=hebfb},FakeSlant=0.2},
  BoldItalicFont={DejaVuSansMono.ttf},
  BoldItalicFeatures={RawFeature={fallback=hebfb},FakeBold=1.5,FakeSlant=0.2}
  ]{DejaVuSansMono.ttf}

% Explicit multilingual/English profiles use real Latin faces and italics.
% The legacy/default Hebrew profile retains its established measured typography.
% Latin Modern is provided by the TeX distribution, not the pinned font bundle;
% recorder/runtime evidence reports it as an additional toolchain resource.
\ifx\lakitlanguageprofile\hllpdfhebrew\else
  \babelfont[english]{rm}[Ligatures=TeX,RawFeature={fallback=hebfb},
    BoldFeatures={RawFeature={fallback=hebfbbf}},
    BoldItalicFeatures={RawFeature={fallback=hebfbbf}}]{Latin Modern Roman}
  \babelfont[english]{sf}[Ligatures=TeX,RawFeature={fallback=hebfbsans},
    BoldFeatures={RawFeature={fallback=hebfbsansbf}},
    BoldItalicFeatures={RawFeature={fallback=hebfbsansbf}}]{Latin Modern Sans}
\fi

\usepackage{amsmath,amssymb,mathtools}
% amsthm supplies theorem/proof machinery. Environment names remain
% document-specific; the generic proof label is localized below.
\usepackage{amsthm}
% Proof names follow the surrounding prose, including return from an island.
\addto\extrashebrew{\renewcommand{\proofname}{הוכחה}}
\addto\extrasenglish{\renewcommand{\proofname}{Proof}}
\ifx\lakitmainlanguage\hllpdfenglish
  \renewcommand{\proofname}{Proof}
\else
  \renewcommand{\proofname}{הוכחה}
\fi
% v3.24.1: equation tags under babel-bidi. A tag re-emitted by mathtools'
% showonlyrefs inside align inherits the RTL base and its parens MIRROR —
% ")2.1(" instead of "(2.1)" (plain align/equation tags are unaffected;
% verified on a minimal case + a 198-page build). LTR-isolate at the tag-BOX
% level: mathtools' show/hide + numbering logic is untouched. Do NOT override
% \tagform@ instead — that clobbers the showonlyrefs mechanism (measured:
% every tag shows, numbering breaks).
% v3.28.1: the v3.24.1 wrap must apply ONLY to tags of display environments.
% Inline \eqref in Hebrew PROSE also goes through \maketag@@@, and there the
% english-wrap MIRRORS the parens — ")1(" — while the unwrapped original is
% correct (proven on a 2x2 isolation matrix: patch x showonlyrefs; regressed
% every prose \eqref since v3.24.1, caught by the render gate's R4). The two
% call sites are indistinguishable by \ifmmode (tag boxes are text mode), so
% the display environments themselves raise a flag; \eqref in prose sees it
% false and gets the original untouched \maketag@@@.
\usepackage{etoolbox}

% babel bidi=basic (including upstream 26.9) assumes every automatic
% discretionary has a one-glyph replacement. LuaTeX can legally produce one
% with replace=nil after a hyphenated Latin value; Babel then indexes nil and the
% whole build dies. Do the narrow upstream-compatible correction: while Babel
% classifies the list, present only that malformed disc as an ordinary
% discretionary, then restore its subtype. Babel still reorders the ENTIRE box;
% there is no pcall and no successful fallback to an unreordered node list.
% If Babel throws for any other reason, the build remains fatal. The synthetic
% nil-disc probe and the real table gallery in test_all.sh pin both layers.
\AtBeginDocument{\directlua{
  if Babel and Babel.bidi and not Babel.heb_nil_disc_compat then
    Babel.heb_nil_disc_compat = true
    local original = Babel.bidi
    local disc_id = node.id('disc')
    Babel.bidi = function(head, ispar, hdir)
      local hidden = {}
      for item in node.traverse_id(disc_id, head) do
        if item.subtype == 2 and (not item.replace or not item.replace.char) then
          table.insert(hidden, item)
          item.subtype = 0
        end
      end
      local result = original(head, ispar, hdir)
      for _, item in ipairs(hidden) do item.subtype = 2 end
      return result
    end
  end
}}
\makeatletter
\newif\ifHEBdisptag
\let\ORIGmaketag@@@\maketag@@@
\def\maketag@@@#1{\ifHEBdisptag\ORIGmaketag@@@{\foreignlanguage{english}{#1}}\else\ORIGmaketag@@@{#1}\fi}
\@for\HEB@env:=equation,equation*,align,align*,gather,gather*,multline,multline*,flalign,flalign*,alignat,alignat*,eqnarray,eqnarray*\do{%
  \expandafter\AtBeginEnvironment\expandafter{\HEB@env}{\HEBdisptagtrue}}
\makeatother
% English page + Hebrew prose: Babel's equation-only tag flip uses a primitive
% leqno path that shifts the formula to the left edge and the tag into the body
% of an English page. Align's full-width alignment is unaffected. Give only this
% display a physical LTR container and native amsmath eqno/leqno. The engine then
% owns real tag width, collision handling and labels; no proxy tag box hides ink
% from its fit calculation. Babel's saved \bbl@eqnodir still restores Hebrew text
% inside math, and the prose locale/math direction are never switched.
% Close the prior RTL paragraph before normalization. After the environment's
% group restores its directions, close TeX's empty resumed display paragraph and
% use native LaTeX continuation handling: same-paragraph prose has no indent,
% while an authored blank line starts an ordinary indented RTL paragraph.
% Register after Babel's setup; remain local to equation(*) in a main-English
% document's actual RTL prose. Native Hebrew, ordinary English and explicit
% fleqn retain their original paths. No page/display dimensions are assigned.
% <hllpdf:minority-rtl-equation-geometry>
\makeatletter
\newcommand{\HEB@minorityprosecontinuation}{\par\@doendpe}
\newcommand{\HEB@minorityequationgeometry}{%
  \ifx\lakitmainlanguage\hllpdfenglish
    \ifnum\pardirection=\@ne
      \if@fleqn\else
        \par\aftergroup\HEB@minorityprosecontinuation
        \iftagsleft@\let\veqno\@@leqno
        \else\let\veqno\@@eqno\fi
        \bodydir TLT\pardir TLT\textdir TLT\relax
      \fi
    \fi
  \fi}
\AtBeginDocument{%
  \AddToHook{env/equation/begin}{\HEB@minorityequationgeometry}%
  \AddToHook{env/equation*/begin}{\HEB@minorityequationgeometry}}
\makeatother
% </hllpdf:minority-rtl-equation-geometry>
% v3.63: \shortintertext rendered Hebrew MIRROR-REVERSED and flush LEFT — an entire
% commentary line inside an align printed backwards. babel's bidi=basic patches
% \intertext@ only (luababel.def), and mathtools' \shortintertext uses the identical
% \vbox{\normalbaselines … \noindent\ignorespaces#1\par} skeleton but is never
% patched, so its paragraph keeps pardir/textdir LTR: no bidi reordering (logical
% order on the page) and flush-left. Measured against \intertext with the same
% Applying babel's own replacement to \shortintertext@ makes the two identical.
% \h{}/\mbox fix the letter order but leave the line flush LEFT, so they are not
% substitutes. Guarded + \AtBeginDocument because mathtools \let-copies at load time.
% NB \ExplSyntaxOn is REQUIRED: mathtools' internals carry expl3 names, and without
% the expl3 catcodes the `_` is read as a subscript and the patch detonates far away
% with "! Illegal unit of measure (pt inserted)" at the \end{align} (measured).
\makeatletter
\AtBeginDocument{%
  \@ifundefined{shortintertext@}{}{%
    \ExplSyntaxOn
    \bbl@sreplace\shortintertext@{\normalbaselines}%
      {\normalbaselines\ifx\bbl@eqnodir\relax\else\bbl@pardir\@ne\bbl@eqnodir\fi}%
    \ExplSyntaxOff}}
\makeatother
% geometry inherits the canonical A4 sheet from \documentclass; build.sh then
% verifies every rendered MediaBox and CropBox, so a silent Letter fallback is
% impossible to publish as a clean default build.
\usepackage[\lakitgeometryoptions]{geometry}
\linespread{\lakitlinespread}       % layout hook; canonical default remains 1.3
\ifx\lakitparagraphstyle\hllpdfparagraphblock
  \setlength{\parindent}{0pt}
  \setlength{\parskip}{6pt plus 1pt minus 1pt}
\fi
% These dimensions remain ordinary scoped TeX lengths: authors can customize
% them after core input. Language switches inherit them; lists keep their own
% indentation and narrow reference layouts establish their local paragraph gaps.
\typeout{HLLPAR|\lakitparagraphstyle}
\usepackage{xcolor}
\usepackage{mdframed}
\usepackage{listings}                 % code listings (assembly / C / pseudo-code)
\usepackage{tikz}
\usetikzlibrary{arrows.meta,positioning,calc,shapes.geometric,fit,backgrounds}
\usepackage{pgfplots}\pgfplotsset{compat=1.18}
\usepackage{quantikz}                 % quantum circuits (remove if unused)
\usepackage{booktabs}
\usepackage{array}
\usepackage{float}                    % [H] = place table exactly here (reference doc: no drifting)
\usepackage{enumitem}
\usepackage{needspace}

% ============================================================================
% ★★ v3.4 — THE DIRECTION-ISOLATION TRIO (the heart of mixing Hebrew + Latin)
% Mental model:
%   \h{...}    Hebrew (RTL) island      — Hebrew text inside an LTR context
%                                          (TikZ nodes, the english `tikzpic`).
%   \en{...}   English (LTR) isolate    — Latin text / parentheses inside a
%                                          Hebrew (RTL) context. KEEPS PARENS
%                                          AND ORDER CORRECT.
%   \code{...} \en + monospace          — code inside Hebrew / inside a frametitle.
% CRITICAL: never wrap Latin in \h{} (it is for HEBREW). Inside a TikZ node,
% Latin in \h{} renders REVERSED (Datapath->htapataD, CPU->UPC) and parens
% mirror. Isolate every Latin/paren chunk with \en{} or \code{}. See
% references/figures-boxes-listings.md.
% ---- TEXT INSIDE MATH (v3.8) ----------------------------------------------
% Put ANY text that sits inside math in a \text{}. The ONLY thing that is
% invalid in math is a bare \h/\en (= \foreignlanguage) used DIRECTLY, outside
% \text -> "\rmfamily invalid in math mode". Inside \text:
%   Hebrew label           -> \text{...}            (bare Hebrew renders fine;
%                              \h is optional — BUT NOT ALWAYS HARMLESS (v3.67):
%                              if the \h group OPENS with a nested island
%                              ($..$, \en{}, \code{}, \texttt{}, \mbox{}) the
%                              WHOLE group loses bidi and the Hebrew prints in
%                              LOGICAL order = gibberish. Measured:
%                                \text{\h{$f$ קבועה}}  -> reads "העובק"  BROKEN
%                                \text{\h{קבועה $f$}}  -> correct
%                                \text{$f$ קבועה}      -> correct (no \h)
%                              Hebrew first inside \h, or drop the \h. Flagged
%                              H-OPENS-WITH-LTR-GROUP.)
%   plain Latin word        -> \text{...}            (no marker needed)
%   Latin WITH ( ) or [ ]   -> \text{\en{...}}       (\en forces LTR; a plain
%                              \text mirrors the parens in RTL -> "Spec (TNR(" )
%   f(Hebrew/text arg)      -> \text{f}(\text{...})  -- keep the parens at MATH
%                              level, NOT \text{f(arg)} (parens inside the box mirror)
%   code in math            -> \mathtt{...}  (or \texttt{}, math-guarded below)
% NEVER write \en{...} or \h{...} straight into $...$ / \[..\] (outside \text);
% always via \text{}. (\mbox{...} also works but does not scale in scripts.)
% ============================================================================
\newcommand{\h}[1]{\foreignlanguage{hebrew}{#1}}
% v3.42.1 — \en islands do NOT hyphenate (auto OR at explicit hyphens). English hyphenation patterns were
% tearing identifiers and terms ACROSS THE RTL SEAM ("IO Regis-/ters",
% "In-/put", "Band-/width", "Per-/formance", "Multi-/ple" — all shipped on
% real pages; probe-proven: "ALUS-/rcB", "Mem-/ToReg"). An island is a label
% or a term, not English prose — a hyphen tear there is never right. A/B on
% a full 197pp lakit: identical overfull sets, same page count, ~11
% paragraphs healed, one mild underfull (badness 3525) that replaced a
% mutilated "Regis-/ters". Long English PROSE belongs in the englishprose
% environment below, which restores explicit-hyphen breaks while its paragraph
% is still active. A raw otherlanguage block inherits the global prohibition.
% (\l@nohyphenation needs \csname:
% '@' is not a letter here.)
% v3.43.2 — \exhyphenpenalty moved OUT of \en to a GLOBAL assignment below.
% Render-proven: a group-local \exhyphenpenalty is DEAD CODE — TeX reads the
% line-breaking parameters at \par time, after the \en group has closed, so
% \en{write-through} / plain not-taken still tore at the "-" ("write-/
% through" shipped on a real exam page). \language is different: it is
% recorded in the node list as text is added, so the in-group nohyphenation
% assignment above DOES work (and stays). The global penalty was A/B-measured
% on a 197pp lakit: identical overfull set (25), identical page count, the
% not-/taken tear healed. Hebrew is unaffected in practice — Hebrew compounds
% use maqaf (־), not ASCII "-".
\newcommand{\en}[1]{\foreignlanguage{english}{\language=\csname l@nohyphenation\endcsname #1}}
% Babel runs extras locally for inline islands and paragraph language changes.
% An English main document must not inherit the Hebrew no-hyphen paragraph rule.
\addto\extrashebrew{\exhyphenpenalty=10000\relax}
\addto\extrasenglish{\exhyphenpenalty=50\relax}
\ifx\lakitmainlanguage\hllpdfenglish
  \exhyphenpenalty=50\relax
\else
  \exhyphenpenalty=10000\relax
\fi
% Full English paragraphs need ordinary explicit-hyphen breaks. Use this
% environment instead of a raw otherlanguage block; the paragraph ends while
% the local penalty is still active.
\newenvironment{englishprose}
  {\par\begin{otherlanguage}{english}\exhyphenpenalty=50\relax}
  {\par\end{otherlanguage}}
\newenvironment{hebrewprose}
  {\par\begin{otherlanguage}{hebrew}\exhyphenpenalty=10000\relax}
  {\par\end{otherlanguage}}
% Use for shell text whose locale must not depend on the last body paragraph.
\DeclareRobustCommand{\doclang}[1]{\foreignlanguage{\lakitmainlanguage}{#1}}

% ============================================================================
% ★★ — fixes for defects that are INVISIBLE in the log and in the checks
% and only surface on a VISUAL RENDER of the PDF. (Each verified empirically.)
% ----------------------------------------------------------------------------
% v3.22: \emph is now REAL slanted emphasis (FakeSlant italic face defined on
% the \babelfont{rm} line above), replacing the v3.21 \emph->\textbf hack: it
% preserves \emph's toggle semantics and keeps bold free for its own role.
% Prefer bold-as-emphasis for your doc? one line: \renewcommand{\emph}[1]{\textbf{#1}}
% page number: force the digits LTR so they never reverse (18 -> 81) in the RTL
%   footer. That reversal is a NON-DETERMINISTIC bidi glitch (observed once in
%   ~80 pp); this LuaLaTeX build is not byte-reproducible, so the defensive
%   LTR wrap is the only reliable guard. Latin digits are always LTR anyway.
\renewcommand{\thepage}{\foreignlanguage{english}{\arabic{page}}}
% v3.70: \pagenumbering \gdef's \thepage and silently DESTROYS the wrap above —
% after \pagenumbering{arabic} the folio guard is gone for the whole rest of the
% document. Two collections discovered this independently and hand-restore the
% wrap in main.tex with the same comment; that knowledge now lives here instead.
% The wrapper re-applies the LTR wrap after the kernel call, for any style.
\makeatletter
\let\lakin@orig@pagenumbering\pagenumbering
\renewcommand{\pagenumbering}[1]{%
  \lakin@orig@pagenumbering{#1}%
  \renewcommand{\thepage}{\foreignlanguage{english}{\csname @#1\endcsname\c@page}}}
\makeatother
% emergencystretch: Hebrew does not hyphenate, so an unbreakable LTR island
%   (\en{...} / an inline math run) landing at the margin can leave a small
%   (~2-13pt) overfull with no legal breakpoint. This rescues ONLY those
%   inline-prose tie-breaks; it is a ceiling used only on offending lines and
%   does NOT touch display-math / listing / table overflow (those still surface
%   in the log for the content-level overflow pass). Verified: baseline overfull
%   -> 0 with this set, with no change to pagination or underfull.
\setlength{\emergencystretch}{3em}
% v3.40: inline math (almost) never breaks across lines. TeX's default
%   (binop 700 / rel 500) lets an inline run split at an operator —
%   $(1.f)\cdot 2^{e}$ shipped as "(1.f)·" ending one line and "2^e" opening
%   the next (reader-reported). In LTR that split is merely ugly; under an
%   RTL base the carried operator lands on the visually WRONG side of the
%   break and the formula reads backwards.
%   WHY 10000 (v3.45.0 — the display-first resolution). Inline math is
%   FORBIDDEN from breaking, so a torn island under an RTL base is
%   structurally impossible. The v3.40 cycle used 9999 ("break only as a last
%   resort") because a bare 10000 left ~30-39 WIDE inline equations
%   overflowing (worst 90.7pt) — but 9999's "last resort" fires OFTEN in
%   tight, non-hyphenating Hebrew: a reader found torn islands across the
%   fleet (~150 candidates; e.g. "P(y |" ending one line and "x_0)" opening
%   the next, read backwards). The step the 9999 cycle missed: every one of
%   those ~30 overflows is a wide equation that BELONGS in display. Convert
%   them to display FIRST, then 10000 gives 0 tears AND 0 overflows (measured,
%   12/12 fleet). Doctrine: inline math is ATOMIC; a wide equation goes in a
%   display, never inline. See content-style ("inline math is atomic").
\binoppenalty=10000 \relpenalty=10000
% v3.38: uniform sentence spacing. TeX's default (\nonfrenchspacing, sfcode
% 3000) is the OLD Anglo-American wide-space-after-period convention — Hebrew
% typography never had it, and modern Latin practice (Chicago et al.) dropped
% it too. Measured on real pages: post-period gaps up to 1.9-2.1x the line's
% mean word gap — the dominant "rivers" source in justified Hebrew prose.
% Survives babel language switches (probed: sfcode stays french inside \en{}
% islands and otherlanguage blocks).
\frenchspacing
% v3.38: widow/orphan protection. LaTeX's default penalties (150/150/50)
% allow a lone line at a page top/bottom almost freely. 10000 = forbid;
% this is SAFE here precisely because the document is \raggedbottom (the
% demoted line just moves on, no vertical stretch fight) — under flushbottom
% these values would cause underfull pages. raggedbottom is the report-class
% default; made explicit so a future class/option change cannot silently
% flip it and turn the 10000s harmful.
\raggedbottom
\widowpenalty=10000 \clubpenalty=10000 \displaywidowpenalty=10000
% em-dash: --- does NOT become — under Frank Ruhl (the font lacks tlig; verified
%   the width is identical with/without Ligatures=TeX, while the mechanism works
%   on Latin Modern). TYPE THE CHARACTER — (U+2014) DIRECTLY; likewise – (U+2013).
%   check_bidi_figures.py flags a --/--- that is adjacent to Hebrew.
% ============================================================================

% ---- STRUCTURAL DIRECTION-SAFETY for \texttt (v3.7) -------------------------
% A parenthesis, bracket, or $ INSIDE \texttt{} mirrors/migrates in RTL prose
% (MIPS 0($s0)->"(s0$)0", $rs->"rs$", arr[i]->"[arr[i") because \texttt sets the
% font but NOT the direction. Redefining \texttt to LTR-isolate its content fixes
% the WHOLE class at once (parens, brackets, $) so plain \texttt{lw $t0,0($s0)}
% is safe — no need to remember \code{}. Guarded by \ifmmode because
% \foreignlanguage is invalid in math (there, plain monospace is already LTR).
% NB (corrected v3.56 — the old note here said "stays tofu", which is WRONG):
% \texttt{Hebrew} is a FATAL build failure. The tt family is DejaVu Sans Mono
% (zero Hebrew) and the tt fallback hebfb is DejaVuSerif (also zero Hebrew), so
% the subsetter gets an empty glyph set and LuaTeX dies with
%   "! error: (file .../DejaVuSansMono.ttf) (type 2): there are no glyphs in the
%    subset"  -> no output PDF.
% The message names a FONT FILE, never Hebrew/\texttt/the source line. And it is
% non-local: add an unrelated Latin \code{} elsewhere and the fatal degrades into
% a SILENT DELETION of the Hebrew word. Put Hebrew in \h{}/normal text, never in
% \texttt{}/\code{}; flagged HEBREW-IN-MONO. To disable, delete this block.
\let\origtexttt\texttt
\protected\long\def\texttt#1{\ifmmode\origtexttt{#1}\else\foreignlanguage{english}{\origtexttt{#1}}\fi}
% v4.8.0: defer the direction test until TeX has inserted an alignment cell's
% math template. Unprotected expansion here tested \ifmmode too early at the
% FIRST token of aligned/array/cases/matrix cells. Engine-level \protected also
% preserves the original \texttt implementation without the self-recursion
% observed when retrofitting \DeclareRobustCommand to its saved alias.
\protected\long\def\code#1{\ifmmode\origtexttt{#1}\else\foreignlanguage{english}{\origtexttt{#1}}\fi}

% ---- \codeline — display-code line (v3.24): the code analogue of display math.
% Reader-friction rule (content-style.md §5b): inline \code{} is for a SINGLE
% ATOM only (one register/variable/mnemonic, a tight no-space expression).
% A multi-token expression — an instruction with operands, a spaced assignment,
% a register-binding list — zigzags the eye RTL<->LTR mid-sentence and, being
% an unbreakable island, line-breaks badly. Give it its own centered line with
% display-math vertical rhythm. \centerline = one box => the expression never
% wraps; \nopagebreak keeps it glued to its introducing sentence.
%   ...the Hebrew sentence ends with a colon:
%   \codeline{lw \$t1, -4(\$t0)}
%   and the prose continues (escape dollars as \$).
% v3.66 (first line of the body): flush a PENDING \item label before the \par.
% LaTeX defers the label box (\@labels) until the next \everypar, and an \hbox in
% vertical mode never fires \everypar — so `\item \codeline{...}` printed the code
% at y=69.1 and the item's own "1." BELOW it at y=103.6 (the number appears after
% its content, attached to nothing). \if@inlabel\leavevmode\fi is CONDITIONAL on
% purpose: an unconditional \leavevmode would open an empty paragraph line when
% \codeline follows ordinary prose. Verified both ways — in a list the label now
% sits on the code's line at its normal x, and after prose the geometry is
% byte-identical to before.
\makeatletter
\newcommand{\codeline}[1]{%
  \if@inlabel\leavevmode\fi
  \par\nopagebreak
  \vspace{\abovedisplayskip}\nopagebreak
  % v3.62: was \centerline{...}, i.e. \hbox to\hsize{\hss #1 \hss}. \hss has
  % INFINITE SHRINK, so TeX could always "fit" the box and NEVER reported an
  % Overfull — the same \hss mechanism as BLD-27's \item label. Measured: a wide
  % \codeline in a defbox painted its first token ON the box's left border bar and
  % sat 9.5pt outside the frame interior on both sides, with 0 Overfull, 0 Underfull
  % and all four gates PASS; in bare prose it reached 21.0pt into BOTH margins.
  % \hfil has no shrink, so the overflow is reported and build.sh's ovf_vis/ovf_big
  % counters see it; \linewidth (not \hsize) makes it box-aware, so a \codeline
  % inside a colored box centres on the box interior. Verified: a FITTING \codeline
  % keeps byte-identical glyph positions, and the box-crossing case now emits
  % "Overfull \hbox (19.03pt too wide)". A wide \codeline should be split in two.
  \hbox to\linewidth{\hfil\code{#1}\hfil}%
  \nopagebreak\vspace{\belowdisplayskip}}
\makeatother

% ---- a thin decorative rule helper ----
\newcommand{\sep}{\par\vspace{2pt}\noindent\textcolor{cRule}{\rule{\linewidth}{0.6pt}}\par\vspace{4pt}}

% ---- colors ----
\definecolor{cBlue}{HTML}{1F4E79}   \definecolor{cBlueBg}{HTML}{EAF1F8}
\definecolor{cGreen}{HTML}{2E6B4F}  \definecolor{cGreenBg}{HTML}{EAF4EE}
\definecolor{cAmber}{HTML}{9C6A00}  \definecolor{cAmberBg}{HTML}{FBF3E2}
\definecolor{cPurple}{HTML}{5B2C6F} \definecolor{cPurpleBg}{HTML}{F3ECF7}
\definecolor{cTeal}{HTML}{0E6E6E}   \definecolor{cTealBg}{HTML}{E6F4F4}
\definecolor{cRed}{HTML}{8E2B2B}    \definecolor{cRedBg}{HTML}{FBEAEA}
\definecolor{cGray}{HTML}{555555}   \definecolor{cRule}{HTML}{1F4E79}

% ---- math operators ----
\DeclareMathOperator{\Tr}{Tr}
\DeclareMathOperator{\rank}{rank}
\DeclareMathOperator{\diag}{diag}
\DeclareMathOperator{\sgn}{sgn}

% ---- bra-ket (quantikz defines \ket/\bra/\braket; add the rest safely) ----
\providecommand{\ket}[1]{\left|#1\right\rangle}
\providecommand{\bra}[1]{\left\langle#1\right|}
\providecommand{\braket}[2]{\left\langle#1\middle|#2\right\rangle}
\newcommand{\dyad}[1]{\left|#1\right\rangle\!\left\langle#1\right|}
% v3.21: general outer product |a><b| and expectation value <A> — agents reach
% for these by reflex (esp. \expval); quantikz defines neither, so their absence
% threw "Undefined control sequence" in a real build. \providecommand = safe.
\providecommand{\ketbra}[2]{\left|#1\right\rangle\!\left\langle#2\right|}
\providecommand{\expval}[1]{\left\langle#1\right\rangle}
\newcommand{\abs}[1]{\left|#1\right|}
\newcommand{\norm}[1]{\left\lVert#1\right\rVert}
\newcommand{\half}{\tfrac12}
\newcommand{\R}{\mathbb{R}}\newcommand{\C}{\mathbb{C}}\newcommand{\Z}{\mathbb{Z}}
% identity operator. NB: \mathbb{1} silently renders a WRONG glyph (msbm has no
% blackboard digits). Use \mathbb{I} = 𝕀 (works with amssymb, no extra package).
% Do NOT "fix" this back to \mathbb{1}.
\newcommand{\id}{\mathbb{I}}

% ---- physical units (v3.35): \punit — upright unit with thin space ---------
% Usage: $300\punit{K}$, $80\punit{mV/decade}$, $10^{7}\punit{cm/s}$ — the
% NUMBER AND UNIT LIVE IN ONE MATH ISLAND. Promoted verbatim from a 199-page
% production lakit (advanced-electronic-devices, dozens of uses).
% Two measured traps this design avoids (both render-proven):
%   1. siunitx SCRAMBLES under RTL — \SI{3.5}{eV} renders "eV 3.5", exponents
%      detach ("Hz 10 x 10^9"): its quantity assembly is not bidi-isolated.
%      Do NOT load siunitx in this pipeline.
%   2. A text-mode unit next to a bare-prose number is TWO adjacent LTR runs
%      and reorders ("5 mm" -> "mm5" — the BIDI-2/3 class). \punit is therefore
%      math-only ON PURPOSE: prose misuse fails LOUDLY at compile (\mathrm
%      outside math) instead of scrambling silently.
% Domain shorthands (\MeV, \fm...) stay per-document: \newcommand{\MeV}{\,\text{\en{MeV}}}
\providecommand{\punit}[1]{\,\mathrm{#1}}

% ============================================================================
% Colored boxes — mdframed (NOT tcolorbox, which breaks under bidi).
% ★★ v3.4: innerleftmargin/innerrightmargin give code listings room so their
% frame/line-numbers do not spill out the box border (see listing styles below).
% ============================================================================
% v3.61: everyline/split skips were applied to `tablebox` only (BLD-14, v3.37) —
% the six colored boxes below never got them, so they still split OPEN across a page
% break: the page-N segment gets top+sides and NO bottom rule, the page-N+1 segment
% sides+bottom and NO top rule. Measured on the shipped fleet this is a COVERAGE GAP,
% not a re-report of BLD-14: 0 split-open `tablebox` segments (that fix holds) vs 62
% split-open COLORED boxes across 7 documents. Verified after the fix: both segments
% close, and the frametitle is NOT duplicated on the continuation. BLD-14's caveat
% still applies — a box taller than a full page grinds mdframed's split-trial loop, so
% split such content manually.
% ============================================================================
% v3.72: frametitle={} leads EVERY \newmdenv option list below. mdframed options
% are GROUP-scoped: a box nested inside \begin{exbox}[frametitle={...}] used to
% INHERIT the outer title — measured, the identical title string rendered twice,
% the second riding a title-bar fill inside the INNER box's frame, for all seven
% environments in both directions. Both natural forms were broken (= preamble
% bug, not an authoring rule). The explicit reset makes every box titleless
% unless ITS OWN [frametitle={...}] says otherwise; verified both directions:
% no duplicate on the nested box, an explicit title still renders.
% Geometry is shared with the optional semantic-width table container. Themes
% select their own existing semantic style below; neither path fixes columns.
\mdfdefinestyle{hllpdfsemanticgeometry}{linewidth=1.1pt,
  innertopmargin=6pt,innerbottommargin=7pt,skipabove=9pt,skipbelow=9pt,
  innerleftmargin=9pt,innerrightmargin=9pt,everyline=true,
  splitbottomskip=6pt,splittopskip=12pt,roundcorner=3pt}
\newcommand{\hllpdfsemanticframestyle}{hllpdfsemanticgeometry}
\newmdenv[frametitle={},style=hllpdfsemanticgeometry,linecolor=cBlue,backgroundcolor=cBlueBg,
  frametitlebackgroundcolor=cBlue,frametitlefontcolor=white,
  ]{defbox}   % הגדרה
\newmdenv[frametitle={},style=hllpdfsemanticgeometry,linecolor=cGreen,backgroundcolor=cGreenBg,
  frametitlebackgroundcolor=cGreen,frametitlefontcolor=white,
  ]{thmbox}   % משפט/טענה
\newmdenv[frametitle={},style=hllpdfsemanticgeometry,linecolor=cAmber,backgroundcolor=cAmberBg,
  frametitlebackgroundcolor=cAmber,frametitlefontcolor=white,
  ]{notebox}  % הערה
\newmdenv[frametitle={},style=hllpdfsemanticgeometry,linecolor=cPurple,backgroundcolor=cPurpleBg,
  frametitlebackgroundcolor=cPurple,frametitlefontcolor=white,
  ]{exbox}    % דוגמה
\newmdenv[frametitle={},style=hllpdfsemanticgeometry,linecolor=cRed,backgroundcolor=cRedBg,
  frametitlebackgroundcolor=cRed,frametitlefontcolor=white,
  ]{warnbox}  % טעות נפוצה
\newmdenv[frametitle={},style=hllpdfsemanticgeometry,linecolor=cTeal,backgroundcolor=cTealBg,
  frametitlebackgroundcolor=cTeal,frametitlefontcolor=white,
  ]{keybox}   % רעיון מרכזי

% ---- codebox (v4.5.0): the dedicated container for code -------------------
% An OPTION, not a policy: asmblock unboxed stays canonical for a full listing,
% and a snippet inside any of the six colored boxes still works as before.
% What makes codebox worth its own environment is that it is the ONE framed box
% that keeps its margin line numbers.
%
% v4.4.0 shipped it WITHOUT them, on the strength of v3.37's "no margin knob can
% save them" plus four fresh probes that agreed. That was wrong: every probe had
% moved the code, the box, or both — never the key that carries the NUMBER
% COLUMN. Measured on the page:
%   resetmargins=false   the listing follows THIS box instead of the page,
%                        moving the code from x=75.26 to x=97.54
%   xrightmargin=26pt    under RTL this is the PHYSICAL LEFT indent, and the
%                        only key that carries the \llap'ed digits with it:
%                        x=64.28 -> 90.18, clearing the rule at 65.64 by 23.6pt
%   xleftmargin=26pt     moves the code and LEAVES THE DIGITS on the rule —
%                        the asymmetry that hid the fix for eleven versions
% resetmargins is hoisted out of the four styles (see the \lstset below them) so
% this override is possible at all; every other listing still resolves to true.
% Pinned by tests/codebox_numbers.tex + probe_lstnum.py --inside-rule, whose
% negative control strips xrightmargin and watches the digits fall back.
\definecolor{cSlate}{HTML}{3F4A56}  \definecolor{cSlateBg}{HTML}{F2F4F7}
\newmdenv[frametitle={},leftline=true,rightline=true,topline=true,bottomline=true,
  linecolor=cSlate,backgroundcolor=cSlateBg,linewidth=0.9pt,
  frametitlebackgroundcolor=cSlate,frametitlefontcolor=white,
  innertopmargin=6pt,innerbottommargin=7pt,skipabove=9pt,skipbelow=9pt,
  innerleftmargin=9pt,innerrightmargin=9pt,everyline=true,splitbottomskip=6pt,splittopskip=12pt,
  roundcorner=3pt]{codebox}   % קוד
% The two keys that make in-box numbering work, both measured on the page:
%   resetmargins=false  the listing follows THIS box instead of the page
%   xrightmargin        under RTL this is the PHYSICAL LEFT indent, and it is
%                       the only one that carries the \llap'ed number column
%                       with it. xleftmargin moves the code and leaves the
%                       digits behind on the rule (measured: code 97.54,
%                       digits 64.28 — the v3.37 defect, reproduced); with
%                       xrightmargin=26pt the digits land at x=90.18, a
%                       clearance of 23.6pt from the rule at 65.64.
% This asymmetry is why v3.4's note had the semantics backwards and v3.37
% concluded the digits could not be saved.
% <hllpdf:listing-container-state>
% v4.29: all state is local to the current environment group. A semantic
% callout masks number text, while a nested codebox restores the formatter
% saved before that mask. Save it once across successive semantic wrappers;
% an explicit numbers=none (or an explicitly empty formatter) remains valid.
\makeatletter
\newif\ifhllpdf@incodebox
\newif\ifhllpdf@englishcodebox
\newif\ifhllpdf@lstnumbermasked
\let\hllpdf@savedlstnumber\thelstnumber
\makeatother
% </hllpdf:listing-container-state>
% <hllpdf:codebox-gutter>
\makeatletter
\AtBeginEnvironment{codebox}{%
  \hllpdf@incodeboxtrue
  % Snapshot the OUTER language before asmblock selects English.
  \iflanguage{english}{\hllpdf@englishcodeboxtrue}{\hllpdf@englishcodeboxfalse}%
  \ifhllpdf@lstnumbermasked
    \let\thelstnumber\hllpdf@savedlstnumber
    \hllpdf@lstnumbermaskedfalse
  \fi
}
% listings' AddToHook is GLOBAL. Register once, and gate on local box state;
% adding an unconditional hook at every codebox entry leaks into later listings.
\lst@AddToHook{PreInit}{%
  \ifhllpdf@incodebox
    % Apply every listing override here, inside the LISTING's own group.
    % Setting margins at box entry leaks them into nested semantic callouts.
    \lstset{resetmargins=false,xrightmargin=26pt}%
    % Styles resolve before PreInit. In LTR the physical number gutter follows
    % xleftmargin; in Hebrew retain the measured inherited xrightmargin path.
    \ifhllpdf@englishcodebox
      \lstset{xleftmargin=26pt,xrightmargin=0pt}%
    \fi
    % ONE box, not two: suppress only the frame/background/skips of a listing
    % inside this codebox. The next unboxed listing keeps its own style.
    \lstset{frame=none,backgroundcolor=\color{cSlateBg},aboveskip=0pt,belowskip=0pt}%
  \fi
}
\makeatother
% </hllpdf:codebox-gutter>

% v3.37: line numbers inside a COLORED BOX are sliced by the box border —
% measured: digit span x=63.5-67.3 vs border line at x=65.2 (the border runs
% THROUGH the digits), and the inter-border gap (3.5pt) is narrower than a
% digit (3.8pt), so no margin knob can save them (under resetmargins,
% xleftmargin is inert in-box — measured: code at x=75.26 with 2.2em AND
% 3.0em). A 3-6 line boxed snippet does not need line numbers; blank the
% number TEXT inside every colored box (\thelstnumber is what gets printed;
% the style's numbers=left survives, so nothing else shifts — code position
% byte-frozen at x=75.26, verified). Unboxed listings keep their numbers.
% v3.78: the blanking is coupled to the FRAME, not to the box name —
% clean_theme.tex registers LATER hooks (registration order wins) that
% restore the number text inside the boxes it renders frameless.
% v4.5.0 RETRACTS the "no margin knob can save them" clause above. It holds only
% under resetmargins=true (the shipped default): there xleftmargin is inert
% in-box, as measured. With resetmargins=false, xrightmargin DOES carry the
% \llap'ed digit column clear of the rule — measured; the codebox \lstset that
% does it is the one above, in the codebox block. The loop below therefore
% covers SIX boxes, not "every colored box": codebox opts out and keeps its
% numbers.
% <hllpdf:semantic-listing-mask>
\makeatletter
\@for\HEB@lstbox:=defbox,thmbox,notebox,exbox,warnbox,keybox\do{%
  \expandafter\AtBeginEnvironment\expandafter{\HEB@lstbox}{%
    \ifhllpdf@lstnumbermasked\else
      \let\hllpdf@savedlstnumber\thelstnumber
    \fi
    \hllpdf@lstnumbermaskedtrue
    \def\thelstnumber{}%
    % The nearest semantic container owns its listing presentation. On exit,
    % TeX's group restores the enclosing codebox state and formatter.
    \hllpdf@incodeboxfalse}}
\makeatother
% </hllpdf:semantic-listing-mask>

% ============================================================================
% Unified TABLE presentation: subtle uniform slate frame + auto-numbered
% caption rendered INSIDE the frame. Use:
%   \begin{tablebox}\centering  <tabular>  \tabcap{<description>}\end{tablebox}
% Replaces the buggy \caption* (which emitted a stray "טבלה N : *" line) and
% provides an explicit framed option + numbered caption "טבלה N.M: תיאור".
% Optional intro text may be placed at the TOP, before the tabular, inside the box.
% ============================================================================
\definecolor{cTblLine}{HTML}{B9C4D2}  \definecolor{cTblBg}{HTML}{FAFBFC}
% v3.23: class-aware — under `article` there is no chapter counter (the promised
% report->article switch used to die with "No counter 'chapter' defined").
\makeatletter
\@ifundefined{c@chapter}{%
  % v3.31: in article the counter resets per SECTION, so an unprefixed number
  % repeated ("טבלה 1." in every section — measured on a two-section render).
  % Same conditional idea as the report branch: prefix the section number once
  % a real \section exists; a pre-section table stays "טבלה 1".
  \newcounter{tbl}[section]
  \renewcommand{\thetbl}{\ifnum\value{section}>0 \thesection.\fi\arabic{tbl}}%
}{%
  \newcounter{tbl}[chapter]
  % v3.30: a report-class doc that uses SECTIONS ONLY (chapter stays 0 —
  % the common short-doc pattern) used to print "טבלה 0.1"; the chapter
  % prefix now appears only once a real \chapter exists (render-verified
  % both ways: "טבלה 1." chapterless, "טבלה 1.1" with chapters).
  \renewcommand{\thetbl}{\ifnum\value{chapter}>0 \thechapter.\fi\arabic{tbl}}%
  % v3.64: the SAME leak in \thesection itself, left behind by v3.30. That fix
  % guarded the CAPTION counter and stopped there, so a chapterless report still
  % printed the heading "0.1 סעיף ראשון" and \ref to it printed "0.2" — measured:
  %   p1 y=65.2  '…0.1'   p1 y=144.9 '…0.2'   p1 y=183.4 '…0.2.1'
  % report.cls hardcodes \thesection as \thechapter.\@arabic\c@section, and the
  % leading 0 is the chapter counter that no \chapter ever incremented. Fixing
  % \thesection fixes the whole chain — \thesubsection and \thesubsubsection are
  % defined in terms of it. \theequation and \thetbl/\thefig were already guarded,
  % which is exactly why this survived: every NUMBER around the heading looked
  % right, so the heading's own 0 read as a style choice rather than a bug.
  % \appendix is safe: report.cls resets chapter to 0 and redefines \thechapter to
  % \@Alph\c@chapter, so an appendix \chapter makes the counter 1 and the guard
  % emits "A.1" unchanged. Render-verified in all four modes (chapterless,
  % chaptered, appendix+chapter, appendix+section-only).
  \renewcommand{\thesection}{\ifnum\value{chapter}>0 \thechapter.\fi\arabic{section}}%
}
\makeatother
% v3.37: everyline=true — if the box DOES split across pages, every segment
% gets a CLOSED frame (measured before: the split ran open-ended into the page
% edge with no bottom border, and the caption could strand alone at the next
% page top — "טבלה A.2" shipped that way). The split skips give the segments
% breathing room. In practice, with the caption GLUE below, every previously-
% split table in the corpus now moves WHOLE instead (the white-gap tradeoff,
% documented benign); everyline stays as the safety net for a genuinely
% page-tall table. NB: a table taller than ONE FULL PAGE can grind mdframed's
% split-trial loop (reproduced: minimal doc, minutes of 1pt-step Overfull
% probes) — break such a monster into two tables.
\mdfdefinestyle{hllpdftablegeometry}{leftline=true,rightline=true,topline=true,bottomline=true,
  linewidth=0.9pt,
  innertopmargin=9pt,innerbottommargin=8pt,skipabove=11pt,skipbelow=11pt,
  innerleftmargin=11pt,innerrightmargin=11pt,roundcorner=4pt}
\newcommand{\hllpdftableframestyle}{hllpdftablegeometry}
\newmdenv[frametitle={},style=hllpdftablegeometry,
  linecolor=cTblLine,backgroundcolor=cTblBg,
  % v4.2.0 (adopted): tablebox is ATOMIC. Short tableboxes near page bottoms
  % repeatedly entered mdframed's split trial, emitting benign Overfull-vbox
  % bursts and occasionally separating the intro prose from the tabular.
  % nobreak=true moves the complete short table to the next page instead. The
  % output-active Overfull-vbox gate still catches a box taller than a page;
  % everyline/split skips remain as defensive styling should that path ever
  % re-arm, but the canonical path never publishes a split tablebox.
  nobreak=true,
  everyline=true,splitbottomskip=6pt,splittopskip=12pt,
  % v3.79 (BLD-37 table side, correction): startinnercode/endinnercode run
  % INSIDE mdframed's body vbox (mdframed.sty:848 / :856+863) -- the second one
  % after \par\unskip, i.e. after line breaking, which is the only moment the
  % \tabcap anchor's true offset from the top of the box EXISTS. See
  % \lakintabboxclose below. These two keys are the ONLY tablebox change.
  startinnercode={\lakintabboxopen},endinnercode={\lakintabboxclose}]{tablebox}
% v3.37: \nopagebreak GLUE — \vsplit honors penalties, so the caption can no
% longer strand alone at the top of the next page (it takes its rows along, or
% the whole box moves). Render-verified on the real shipped orphan.
% v3.59: \NR@gettitle is what nameref injects into the stdlib title-bearing
% commands (\@sect, \@chapter, \caption) to fill \@currentlabelname. \figcap and
% \tabcap are INDEPENDENT macros — they never called it, so \@currentlabelname kept
% whatever the last heading left and \nameref{fig:x} silently printed the enclosing
% SECTION's title instead of the caption (proven in the .aux: field 3 of \newlabel
% was the heading text while field 1 and the anchor `fig.1.1` were correct — so this
% is a THIRD, distinct head next to BLD-13's counter fix and BIDI-38's autoref-name
% fix, neither of which writes that field). \@ifundefined keeps it inert if
% hyperref/nameref is ever dropped. Guarded helper because \figcap is defined
% after \makeatother.
% NB both helpers below need \makeatletter — this part of the preamble is OUTSIDE
% any letter-catcode region, so an unwrapped \@currentlabel parses as LaTeX's \@
% and dies with "You can't use \spacefactor in vertical mode" (measured).
\makeatletter
\newcommand{\lakinrtitle}[1]{\@ifundefined{NR@gettitle}{}{\NR@gettitle{#1}}}
% v3.59: \hitem{<letter>}{<key>} — a REFERENCEABLE Hebrew-lettered list item.
%   \begin{enumerate} \hitem{א}{it:a} הסעיף הראשון. \end{enumerate}  ->  \ref{it:a} = (א)
% Why this exists: automatic Hebrew list labels fail silently, so SKILL.md mandates
% an EXPLICIT \item[(\textbf{א})] — but LaTeX puts \refstepcounter INSIDE
% \if@noitemarg (latex.ltx:12964), and \item[...] never sets that flag. So an
% explicit item steps no counter and never touches \@currentlabel: a \label there
% binds to the AMBIENT one and \ref prints the CHAPTER/SECTION number where the
% item letter belongs (proven in the .aux: {chapter.1}). Same doctrine as \figcap —
% own the label state before \label. \phantomsection matters: without it the
% hyperlink still lands on the chapter head (measured); with it the anchor is
% section*.N and the bookmark outline stays clean.
\newcommand{\hitem}[2]{\item[(\textbf{#1})]%
  \def\@currentlabel{(#1)}\def\@currentlabelname{#1}\phantomsection\label{#2}}
% NB (v3.66): \hitem is \ref-only. \autoref{key} calls it a SECTION ("סעיף (א)")
% because \phantomsection brands the anchor section*.N, and \nameref{key} prints
% the bare letter. For \autoref-able items prefer \hebenumlabels below.
%
% v3.66: \hebenumlabels — OPT-IN automatic Hebrew enumerate letters, done right.
%   {\hebenumlabels\begin{enumerate}\item\label{q:a} ...\end{enumerate}}
%   -> labels (א) (ב) (ג); \ref{q:a} = א; anchor Item.N (real, per-item).
% This retires the belief that made \item[(\textbf{א})] + \hitem necessary.
% "Automatic Hebrew labels fail silently" was TRUE but unexplained, and the
% mechanism is now bisected (plain article vs +enumitem, measured): ENUMITEM
% PRE-EXPANDS the label expression in an \edef BEFORE \refstepcounter, with the
% counter still at 0. Tokens it shields by name (\arabic,\alph,\roman) survive;
% a plain \newcommand \ifcase wrapper evaluates at 0 -> the EMPTY zero-branch ->
% the label prints as "()" with the letter GONE (and \the\c@enumi prints a
% literal 0, \number\value dies fatally pointing at \item). The cure is one
% property, not a doctrine: the wrapper must be ROBUST (survives the \edef
% unexpanded), while \theenumi stays EXPANDABLE (so \refstepcounter writes the
% real letter to the .aux). Verified: labels (א)(ב)(ג), \ref -> א/ב/ג each item
% its own value, .aux = {א}{...}{Item.1} — versus the explicit-\item doctrine
% where BIDI-50/\hitem exist purely to patch the counter that never stepped.
\def\@hebalph#1{\ifcase#1\or א\or ב\or ג\or ד\or ה\or ו\or ז\or ח\or ט\or י\or
  יא\or יב\or יג\or יד\or טו\or טז\or יז\or יח\or יט\or כ\else\@ctrerr\fi}
\DeclareRobustCommand*\hebalph[1]{\expandafter\@hebalph\csname c@#1\endcsname}
\newcommand{\hebenumlabels}{%
  \renewcommand{\theenumi}{\@hebalph\c@enumi}%
  \renewcommand{\labelenumi}{(\textbf{\hebalph{enumi}})}}
\makeatother
% v3.63: \mbox the LABEL. \thetbl ("7.1") and a caption description that STARTS
% WITH A DIGIT ("32 הרגיסטרים…") are two European-Number runs separated only by
% neutrals (the label's own "." and \enspace). Under the RTL paragraph base the UBA
% resolves those neutrals to the RTL level, so both numbers sit in ONE reversed span
% and EXCHANGE VISUAL POSITIONS. Measured on the shipped PDF: a table whose real
% number is 7.1 printed "טבלה 32" — a different, EXISTING table number — while \ref
% correctly printed 7.1, so the caption silently contradicted its own cross-reference,
% and the description lost its leading count. The source contains no adjacency at all;
% it is manufactured by this macro, which is why every source linter was blind.
% \mbox makes the label one atomic box so nothing can reorder across it. Verified
% zero-regression: an ordinary Hebrew caption keeps identical x-geometry.
% One paragraph style covers label, number, separator and description.
% Explicit project/template overrides remain possible; the label hook is kept
% for compatibility, with an identity default. Counters/anchors stay unchanged.
\newcommand{\lakitcaptionstyle}{\normalfont\small\color{black}\hypersetup{allcolors=black}}
\newcommand{\lakitcaptionlabelstyle}[1]{#1}
\newcommand{\lakitcaptionpunctuation}{:}
% Native floats, captionof and longtable use the same presentation policy.
% subcaption is optional and inherits the same font rather than a second size.
\usepackage{caption}
\DeclareCaptionFont{lakit}{\lakitcaptionstyle}
\DeclareCaptionLabelFormat{lakit}{\lakitcaptionlabelstyle{#1~\en{#2}}}
\DeclareCaptionLabelSeparator{lakit}{\lakitcaptionpunctuation\enspace}
\captionsetup{font=lakit,labelfont={},textfont={},labelformat=lakit,
  labelsep=lakit,justification=centering}
\makeatletter
\AtBeginDocument{\@ifpackageloaded{subcaption}{%
  \captionsetup[sub]{font=lakit,labelfont={},textfont={},labelsep=lakit}}{}}
\makeatother
% v3.64: OPT-IN list-of-figures/tables entry via an optional short caption.
% \listoffigures and \listoftables are inherited from report.cls and LOOK available,
% but \figcap/\tabcap step the custom \c@fig/\c@tbl counters and never call
% \addcontentsline — so both printed their heading and NOTHING ELSE. Measured: a doc
% with one \figcap and one \tabcap rendered page 1 as 'רשימת האיורים\n1\n', and the
% .lof held only \@chapter's \addvspace, zero \contentsline entries. An empty list
% under a correct Hebrew heading reads as "this document has no figures".
%
% WHY OPTIONAL, AND NOT AUTOMATIC. Writing #1 unconditionally was implemented first
% and REJECTED: a caption containing \codeline compiles CLEANLY today (verified on
% the pre-v3.64 preamble: 0 errors, PDF produced) but becomes FATAL once it is also
% written to the .lof — \codeline is \hbox to\linewidth{...}, and a full-width box in
% a list entry dies with "\pdfextension endlink cannot be used in vertical mode", no
% output PDF. Gaining a feature nobody uses (0 of 13 collections call \listoffigures)
% at the price of breaking a construct that works is a bad trade. Stock \caption has
% exactly this limitation and exactly this escape hatch, so the convention is familiar.
% \en{}, \code{}, \h{} and inline math in the short text are all safe — verified, they
% survive the write and typeset correctly in the list (RTL geometry confirmed: number
% rightmost at x=509.9, title left of it, dots and folio at the far left).
% With no optional argument the macro is byte-identical to v3.63 — zero risk to every
% existing document, and \listoffigures lists exactly the captions that opted in.
% BLD-37, TABLE side (the \figcap v3.74 idea, adapted): clicking a \ref to a
% table used to land MID-table — hyperref's automatic anchor inside
% \refstepcounter was emitted on the still-open horizontal line whose TAIL is
% the tabular (an hbox whose baseline is the table's vertical CENTER), so the
% jump scrolled the top half of the table above the viewport. The v3.74 entry
% deferred this side claiming "tablebox splits — no single box to measure";
% that deferral dissolves: the measurement target is the TABULAR itself at
% \tabcap ENTRY, while it is still the tail of the open list — the same
% situation \lakin@grabfig handles at an island's end (and the fleet has 0
% split tableboxes). So \lakin@grabtab runs as the FIRST token, before any
% typesetting can push the box away, and the anchor is re-emitted raised above
% the measured tabular. \lakintabraise is a SEPARATE register from
% \lakinfigraise (defined with it, below the LTR-island block): a captionless
% island intentionally leaves the fig register set, and a table must never
% inherit that stale figure height.
%
% WHY the anchor stays AT THE ENTRY SPOT and is NOT deferred into the caption
% paragraph the way \figcap defers it: at \par TeX DROPS a final glue item,
% and the automatic anchor's smash box sitting BEHIND the interword glue kept
% that glue alive — so every centered tabular has always sat half a space off
% true center. Deferring the anchor removed that box and was measured as 982
% differing glyph records: EVERY tabular line shifted 1.30pt (2.61pt for
% pattern 7's two minipages) while captions stayed put. Emitting the raised
% anchor at the old spot — through hyperref's own \Hy@raisedlink, exactly the
% macro the automatic path uses — keeps the line's node structure, and so the
% A/B glyph diff at 0, BY CONSTRUCTION. \Hy@raisedlink reads the raise from
% \HyperRaiseLinkDefault (default \baselineskip = the old mid-table dest); the
% group-local \edef freezes measured-height+42pt into it, and \smash inside
% \Hy@raisedlink hides any raise from the line's extent. \HyCnt@ProvideTheH-
% Counter keeps the dest NAME byte-identical to the automatic one (it defines
% \theHtbl, which \hyper@makecurrent prefers — without it the first measured
% caption was named tbl.1 where the automatic path said tbl.0.1, measured).
% vmode entry, or anything but a box at the tail (no tabular to measure),
% falls back to the plain \refstepcounter — byte-identical to the old path.
\makeatletter
\newcommand{\tabcap}[2][]{\lakin@grabtab
  % v3.79: arm the box-close correction. Set unconditionally and in BOTH
  % branches below's scope, so the plain-\refstepcounter fallback is corrected
  % too; \lakintabboxopen clears it at every \begin{tablebox}, so a \tabcap
  % outside a tablebox can never leave it armed for a later box.
  \global\lakintabcapseentrue
  \ifdim\lakintabraise>0pt
    \H@refstepcounter{tbl}\HyCnt@ProvideTheHCounter{tbl}\hyper@makecurrent{tbl}%
    \begingroup
      % A frameless table targets its measured top. The frame's historical
      % cushion belongs only to tablebox and is corrected at box close.
      \edef\HyperRaiseLinkDefault{\the\lakintabraise}%
      \iflakintabframed
        \edef\HyperRaiseLinkDefault{\the\dimexpr\lakintabraise+42pt\relax}%
      \fi
      \Hy@raisedlink{\hyper@anchorstart{\@currentHref}\hyper@anchorend}%
    \endgroup
    \global\lakintabraise 0pt\relax
  \else
    \refstepcounter{tbl}%
  \fi
  \lakinrtitle{#2}\par\nopagebreak\addvspace{6pt}%
  \def\lakin@short{#1}\ifx\lakin@short\@empty\else
    \addcontentsline{lot}{table}{\protect\numberline{\thetbl}{\protect\ignorespaces #1}}\fi
  {\centering\lakitcaptionstyle\mbox{\lakitcaptionlabelstyle{\tablename~\en{\thetbl}\lakitcaptionpunctuation}}\enspace #2\par}}
\makeatother

% ============================================================================
% Unified FIGURE caption — the \tabcap twin for TikZ/plot figures (v3.34.0).
% A figure (tikzpic) sits inline in the prose; put \figcap BELOW it for a
% numbered, \ref-able caption "איור N.M: תיאור" in the same style as \tabcap:
%   \begin{tikzpic} ... \end{tikzpic}
%   \figcap{\label{fig:foo}description}     % \label FIRST, inside the argument
% WHY a macro, not a hand-written gray line: a manual
%   {\centering\small\textcolor{cGray}{איור: ...}} with a \label inside has NO
% \refstepcounter, so the \label binds to the AMBIENT \@currentlabel (the last
% equation) — \ref then prints an equation number AS the figure number
% (measured on a real page: "איור 48" where 48 is equation 48). \figcap steps a
% real \c@fig counter first, so \label is correct, numbering is consistent, and
% refs are stable. Class-aware like \thetbl (chapter- or section-prefixed).
% ============================================================================
\makeatletter
\@ifundefined{c@chapter}{%
  \newcounter{fig}[section]
  \renewcommand{\thefig}{\ifnum\value{section}>0 \thesection.\fi\arabic{fig}}%
}{%
  \newcounter{fig}[chapter]
  \renewcommand{\thefig}{\ifnum\value{chapter}>0 \thechapter.\fi\arabic{fig}}%
}
\makeatother
\makeatletter
\newcommand{\figcap}[2][]{%
  \ifdim\lakinfigraise>0pt
    % v3.74: step WITHOUT hyperref's automatic anchor, then emit the anchor
    % raised above the measured figure. \H@refstepcounter is hyperref's saved
    % pre-hook \refstepcounter (sets the counter AND \@currentlabel — \label
    % in the caption arg binds exactly as before); \hyper@makecurrent supplies
    % \@currentHref, so the dest NAME is identical to the old automatic one —
    % \ref/\autoref/bookmarks are unchanged, only the dest'S POSITION moves.
    \H@refstepcounter{fig}\hyper@makecurrent{fig}%
    % The anchor is emitted INSIDE the caption paragraph below, as a zero-size
    % raisebox — NOT here in vertical mode, where a \raisebox starts an empty
    % paragraph and adds a blank line (measured: every glyph after the caption
    % shifted +18.8pt). Inside an existing line a [0pt][0pt] box is inert.
    \edef\lakin@anch{\@currentHref}%
    % \edef FREEZES the measured height into the pending macro: the deferred
    % \raisebox is typeset inside the caption paragraph AFTER the register is
    % reset to the 0pt sentinel below — a plain \def read 0 at use time
    % (measured: LAKINRAISE=0.0, dest moved 14pt instead of above the figure).
    \edef\lakin@pend{\noexpand\raisebox{\the\dimexpr\lakinfigraise+42pt\relax}[0pt][0pt]{%
      \noexpand\hyper@anchorstart{\lakin@anch}\noexpand\hyper@anchorend}}%
    \global\lakinfigraise 0pt\relax
  \else
    \refstepcounter{fig}\let\lakin@pend\relax
  \fi
  \lakinrtitle{#2}\par\nopagebreak\addvspace{4pt}%
  \def\lakin@short{#1}\ifx\lakin@short\@empty\else
    \addcontentsline{lof}{figure}{\protect\numberline{\thefig}{\protect\ignorespaces #1}}\fi
  {\centering\lakin@pend\lakitcaptionstyle\mbox{\lakitcaptionlabelstyle{\figurename~\en{\thefig}\lakitcaptionpunctuation}}\enspace #2\par}}
\makeatother

% ============================================================================
% LTR islands.  EVERY tikzpicture / quantikz MUST live inside one, or the whole
% figure's coordinates and Latin labels flip. (Hebrew node text still uses \h{}.)
% ============================================================================
% v3.74 (BLD-37): the \figcap hyperlink DESTINATION used to land at the caption
% line — BELOW the figure — so clicking a figure reference scrolled the figure
% off-screen (measured fleet-wide: 97/97 fig dests in one collection sat 19-54pt
% under their figure's ink; a 347pt figure was fully hidden after the jump).
% hypcap-style fix in two parts: (1) at the island's END, \lakin@grabfig grabs
% the just-typeset centered line with \lastbox, measures it, and puts it back
% with \nointerlineskip — the interline glue above the line was already
% contributed and is NOT re-added, so the page geometry is unchanged to the
% glyph (verified); (2) \figcap then steps the counter WITHOUT hyperref's
% automatic anchor (\H@refstepcounter = the pre-hyperref \refstepcounter) and
% emits the anchor itself, raised by the measured height — the dest lands ABOVE
% the figure top. A captionless island leaves numbering untouched (nothing is
% stepped at the island; the grab flag is cleared at every island begin, so a
% stale height from a captionless figure cannot leak into a later caption).
\newsavebox{\lakinfigbox}
\newdimen\lakinfigraise  % real \dimen: tex.setdimen cannot target a \newlength SKIP register (measured lua error)
% BLD-37 table side: a SEPARATE register for \tabcap. The two must not share —
% a captionless island INTENTIONALLY leaves \lakinfigraise set (the grab runs
% at every island END; only the NEXT island's begin clears it), so a shared
% register would hand a \tabcap after a captionless figure that stale FIGURE
% height. \lakintabraise is written only by \lakin@grabtab at \tabcap entry
% and consumed/reset inside \tabcap itself, so it can never go stale.
\newdimen\lakintabraise
\makeatletter
% v3.74b: measure WITHOUT touching the list. The first draft grabbed the box
% with \lastbox and put it back — and \lastbox DISCARDS the node's baseline
% shift, so a vertically-anchored quantikz circuit rose 2.57pt and everything
% after it sank 5.32pt (measured on the A/B glyph diff). LuaTeX lets us READ
% the tail node's dimensions without removing anything: geometry is identical
% by construction. \lakinfigraise=0pt is the "no figure measured" sentinel.
\def\lakin@grabfig{%
  \ifhmode
    \directlua{
      local n = tex.nest[tex.nest.ptr].tail
      local GLUE, PEN = node.id('glue'), node.id('penalty')
      while n and (n.id == GLUE or n.id == PEN) do n = n.prev end
      local HL, VL = node.id('hlist'), node.id('vlist')
      if n and (n.id == HL or n.id == VL) then
        tex.setdimen('global', 'lakinfigraise',
                     (n.height or 0) + (n.depth or 0))
      else
        tex.setdimen('global', 'lakinfigraise', 0)
      end}%
  \else
    \global\lakinfigraise 0pt\relax
  \fi}
% BLD-37 table side: \lakin@grabfig's twin, NOT a shared helper — the Lua body
% differs in exactly one place, the skip set. Between \end{tabular} and \tabcap
% an otherlanguage wrapper (patterns 3/5 of example_tables) leaves a \write
% whatsit and a dir node ON TOP of the interword glue (measured node dump:
% "whatsit/write <- dir <- glue(2.62) <- hlist"), so the table reader must
% also step over whatsit/dir markers — all zero-size — to reach the tabular.
% \lakin@grabfig deliberately does NOT get the wider skip set: its glue/penalty
% body is the one the v3.74 A/B proof froze, and widening it would newly
% measure (and so newly raise) anchors in island shapes that today take the
% plain path — an unproven behavior change on the FIG side. Anything that is
% not a box behind those markers still hits the 0pt sentinel and \tabcap falls
% back to the plain \refstepcounter.
% SECOND deliberate difference: this stores the box's HEIGHT, not height+depth.
% \figcap's anchor is emitted in the caption paragraph BELOW the figure, so it
% must climb the full extent; \tabcap's anchor sits ON the tabular's own line,
% whose baseline is the vcentered tabular's MIDPOINT — height alone reaches
% the top edge, and adding the depth was measured to overshoot by half the
% table (dests 80-135pt above their frame, INSIDE the preceding table's box).
\def\lakin@grabtab{%
  \ifhmode
    \directlua{
      local n = tex.nest[tex.nest.ptr].tail
      local GLUE, PEN = node.id('glue'), node.id('penalty')
      local WHAT, DIR = node.id('whatsit'), node.id('dir')
      while n and (n.id == GLUE or n.id == PEN
                   or n.id == WHAT or n.id == DIR) do n = n.prev end
      local HL, VL = node.id('hlist'), node.id('vlist')
      if n and (n.id == HL or n.id == VL) then
        tex.setdimen('global', 'lakintabraise', n.height or 0)
      else
        tex.setdimen('global', 'lakintabraise', 0)
      end}%
  \else
    \global\lakintabraise 0pt\relax
  \fi}
\makeatother
% ============================================================================
% v3.79 (BLD-37 table side, CORRECTION): put the anchor at the TOP OF THE BOX,
% instead of guessing the distance to it.
% ============================================================================
% v3.77 raised the \tabcap anchor by (height of the box at the tail of the open
% horizontal list) + 42pt. That is a MEASUREMENT taken BEFORE line breaking, and
% it is short by every term it cannot see. Fleet-measured on the 12 built
% collections (299 tbl.* dests, ownership established through each table's own
% CAPTION, never through "which frame encloses the dest"): 3 dests landed INSIDE
% their own tablebox, below the frame top —
%   computer-architecture       tbl.14.1  -46.4pt   three stacked tabulars
%   computer-architecture       tbl.30.5  -13.5pt   tabular 2.85pt over \linewidth
%   advanced-electronic-devices tbl.B.1   -59.2pt   tabular 2.56pt over \linewidth
% Two missing terms, both reproduced in a minimal document (probe P2/P3):
%   (1) VERTICAL LIST ALREADY CONTRIBUTED inside the box — a second and third
%       tabular, or a heading above the table. The tail read sees only the LAST
%       box, so 77.7pt of stacked tabulars is invisible to the raise.
%   (2) WHICH LINE THE ANCHOR LANDS ON. When the tabular alone overfills
%       \linewidth (459.6121pt in a tablebox), the interword glue after
%       \end{tabular} becomes the chosen breakpoint and \Hy@raisedlink is
%       emitted on the NEXT line — one (depth + lineskip) lower. That term does
%       not EXIST at \tabcap time: line breaking has not happened yet. No
%       measurement taken there can ever recover it. (Proof: probe P3, the same
%       table 2.84912pt too wide, fails at -13.5pt; probe P4, the identical
%       table with \tabcolsep 2pt so it fits, is correct at +32.0pt.)
% So the raise is not computed at \tabcap at all. It is CORRECTED at
% \end{tablebox}, from inside mdframed's own body vbox (endinnercode runs after
% \par\unskip, i.e. after line breaking, while tex.nest[tex.nest.ptr] is still
% that vbox's vertical list). There the anchor's true offset from the top of the
% box content is not estimated but WALKED, and the \Hy@raisedlink shift is set
% so the dest lands 42pt above it — which is 42pt - innertopmargin(9pt) -
% linewidth(0.9pt) = 32.0pt above the FRAME top, the exact position the 295
% already-correct dests have today. That is why this is a correction and not a
% re-placement: the good sites do not move.
%
% WHY THE GLYPH DIFF IS 0 BY CONSTRUCTION — the only thing written is the
% `shift` field of the hlist that \Hy@raisedlink already built, and the guard
% below refuses to write it unless that box is 0x0x0. A zero-size box carries no
% dimension into its parent's packing, the packing already happened, and shift
% is not an input to it. Nothing else in the list is touched: no node is added,
% removed, resized or reordered. Measured, not assumed: assets/example_tables.tex
% and the whole computer-architecture collection A/B at 0 differing pymupdf
% rawdict glyph records.
%
% WHY NOT the alternatives (all prototyped and measured before this one):
%   * accumulate the vertical list at \tabcap time — fixes term (1), CANNOT fix
%     term (2); 2 of the 3 real failures survive. A fix that leaves the majority
%     of the defect standing is the false all-clear this repo exists to avoid.
%   * \zref-savepos two-pass — correct on all shapes, but it makes the dest
%     POSITION aux-derived (DEFECTS BLD-1/2), a single-pass compile silently
%     ships an unraised anchor, and the bare y-difference has no page guard.
%   * emit the anchor at \begin{tablebox} and let \tabcap adopt it — needs the
%     tbl counter stepped at box OPEN to keep the dest NAME byte-identical; a
%     captionless tablebox then either shifts every later table number or leaves
%     an orphan dest that the next table duplicates (a HARD build gate). This
%     correction touches neither the counter nor the name: \tabcap's naming path
%     is v3.77's, unchanged, so resolve_names() is byte-identical (verified).
%   * suppressing the line break (a penalty before the glue) — that MOVES the
%     tabular: the break currently discards the interword glue, so forbidding it
%     re-centers the two overfull tables. Reflow, rejected.
%
% SCOPE. The correction is armed only by \tabcap ITSELF (\lakintabcapseen),
% cleared at every box open, so a captionless tablebox is a no-op and a \tabcap
% hosted OUTSIDE a tablebox (the fleet has exactly one, computer-architecture
% tbl.14.2 in an exbox) keeps v3.77's measured raise. Both fallbacks are
% therefore never worse than today; and because hyperref's AUTOMATIC anchor uses
% the same \Hy@raisedlink shape, the plain-\refstepcounter path (vmode entry, or
% prose between \end{tabular} and \tabcap — probe P6, -51.5pt) is corrected too.
% KNOWN LIMIT — a tablebox that SPLITS across pages. The offset is measured from
% the top of the UNSPLIT body, so a dest in segment 2 gets raised past its own
% segment. MEASURED on a forced split (probe: 40 filler lines + a table): the
% dest went from y=585.0 on page 2 (inside segment 2, at the caption — the old
% BLD-37 symptom) to y=-113.4, i.e. ABOVE the page top. Viewers clamp that to the
% page top, which happens to show segment 2 whole, but an out-of-page coordinate
% is not something to ship on purpose. Hence the \textheight guard in
% lakintab.fix: a body TALLER THAN THE TEXT BLOCK cannot fit on any page, so it
% must split — the correction is skipped and v3.77's behaviour stands unchanged
% (verified: the probe returns to y=585.0). The guard is deliberately EXACT, not
% a page-fitting heuristic: it is sound (a taller-than-page body always splits)
% but PARTIAL (a shorter body that starts low on a page could still split). That
% residue is unreachable in practice — v3.37's everyline + the caption's
% \nopagebreak glue move such boxes WHOLE instead (probe N7: a 22-row table
% pushed down by 430pt moved whole, 0 segments), and the fleet has 0 split
% tableboxes, re-counted after this change. Reimplementing mdframed's own
% page-fitting arithmetic to close the residue would be a heuristic with its own
% failure mode, which is the trade this repo has repeatedly refused.
\newif\iflakintabcapseen
\newif\iflakintabframed
% The Lua lives in ONE \directlua at load time, and is written with NO
% backslash, %, # or ~ : \directlua expands its argument like \edef, so an
% active ~ would expand to \nobreakspace and a % would eat the rest of the line
% (the second is how the first draft of this block died). "not (a == b)" stands
% in for Lua's inequality operator for the same reason.
\directlua{
  lakintab = { }
  lakintab.HL = node.id("hlist")
  lakintab.VL = node.id("vlist")
  lakintab.RU = node.id("rule")
  lakintab.GL = node.id("glue")
  lakintab.KE = node.id("kern")
  lakintab.WH = node.id("whatsit")
  lakintab.TM = node.id("temp")
  lakintab.WD = node.subtype("pdf_dest")
  function lakintab.d(v)
    if v == nil then return 0 end
    if v < -1000000000 then return 0 end
    return v
  end
  function lakintab.scan(list, y, horiz, parent, target)
    local x = list
    while x do
      local i = x.id
      if i == lakintab.HL or i == lakintab.VL then
        if horiz then
          local b = y + x.shift
          if i == lakintab.HL then
            lakintab.scan(x.head, b, true, x, target)
          else
            lakintab.scan(x.head, b - lakintab.d(x.height), false, x, target)
          end
        else
          if i == lakintab.HL then
            lakintab.scan(x.head, y + lakintab.d(x.height), true, x, target)
          else
            lakintab.scan(x.head, y, false, x, target)
          end
          y = y + lakintab.d(x.height) + lakintab.d(x.depth)
        end
      elseif i == lakintab.WH then
        if x.subtype == lakintab.WD then
          local id = x.dest_id
          if type(id) == "string" and string.sub(id, 1, 4) == "tbl." then
            if horiz and parent and parent.width == 0
               and parent.height == 0 and parent.depth == 0 then
              parent.shift = parent.shift + (target - y)
            end
          end
        end
      elseif not horiz then
        if i == lakintab.RU then
          y = y + lakintab.d(x.height) + lakintab.d(x.depth)
        elseif i == lakintab.GL or i == lakintab.KE then
          y = y + lakintab.d(x.width)
        end
      end
      x = x.next
    end
  end
  function lakintab.fix(target, maxh)
    local nst = tex.nest[tex.nest.ptr]
    if nst == nil then return end
    local h = nst.head
    if h == nil then return end
    if h.id == lakintab.TM then h = h.next end
    local tot = 0
    local x = h
    while x do
      local i = x.id
      if i == lakintab.HL or i == lakintab.VL or i == lakintab.RU then
        tot = tot + lakintab.d(x.height) + lakintab.d(x.depth)
      elseif i == lakintab.GL or i == lakintab.KE then
        tot = tot + lakintab.d(x.width)
      end
      x = x.next
    end
    if tot > maxh then return end
    lakintab.scan(h, 0, false, nil, target)
  end
}
% @-FREE NAMES ON PURPOSE: these two are named in the `tablebox` key list at
% \newmdenv above, which is tokenised OUTSIDE any \makeatletter region — a
% \lakin@... there would tokenise as \lakin followed by the letters @tabopen.
\newcommand{\lakintabboxopen}{\lakintabframedtrue\global\lakintabcapseenfalse}
\newcommand{\lakintabboxclose}{%
  \iflakintabcapseen
    \directlua{lakintab.fix(-\number\dimexpr\lakintabcushion\relax,
                            \number\dimexpr\textheight\relax)}%
    \global\lakintabcapseenfalse
  \fi\lakintabframedfalse}
% The cushion above the box's CONTENT top. 42pt puts the dest 32.0pt above the
% frame top (42pt - innertopmargin 9pt - linewidth 0.9pt), i.e. exactly where
% the 295 correct dests already are — the correction moves only the broken ones.
\newcommand{\lakintabcushion}{42pt}
\newenvironment{qcirc}{\par\medskip\global\lakinfigraise 0pt\relax\begin{otherlanguage}{english}\begin{center}}
  {\csname lakin@grabfig\endcsname\end{center}\end{otherlanguage}\par\medskip}
\newenvironment{tikzpic}{\par\medskip\global\lakinfigraise 0pt\relax\begin{otherlanguage}{english}\begin{center}}
  {\csname lakin@grabfig\endcsname\end{center}\end{otherlanguage}\par\medskip}

% ============================================================================
% Code listings.  ★★ v3.4 THE BOX-OVERFLOW FIX:
%   resetmargins=true   → the listing IGNORES the enclosing environment's
%                         indentation and measures from the PAGE. v3.4 stated
%                         this backwards ("measured from the ENCLOSING box"),
%                         and v3.37 then blanked in-box line numbers because
%                         the page-anchored listing put its hanging digits on
%                         top of the box rule. Measured v4.5.0: with the key
%                         FALSE the listing follows the box and the code moves
%                         from x=75.26 to x=97.74, clearing the rule at 72.07.
%                         The key now lives in one \lstset below instead of in
%                         each style, so a box can override it — codebox does.
%   framexleftmargin=0  → the frame (left bar) does not extend left of the code.
%   xleftmargin keeps a small indent so the line numbers have room INSIDE.
% asmblock wraps the listing in an LTR island (code is LTR).
% ============================================================================
\lstdefinestyle{asm}{
  basicstyle=\ttfamily\small,
  backgroundcolor=\color{cBlueBg!55},
  frame=leftline,framerule=2.5pt,rulecolor=\color{cBlue},
  numbers=left,numberstyle=\tiny\color{cGray},numbersep=8pt,
  xleftmargin=2.2em,framexleftmargin=0pt,framexrightmargin=0pt,
  breaklines=true,columns=fullflexible,keepspaces=true,
  showstringspaces=false,aboveskip=8pt,belowskip=8pt,
  % v3.36: without a comment DELIMITER the commentstyle below never fires —
  % listings has no `language` here, so it does not know '#' opens a comment,
  % and every asm listing rendered flat black (measured: a whole-page pixel
  % sweep found ZERO cGreen glyphs across 135 shipped listings). morecomment
  % gives it the delimiter; color-only change, glyph widths untouched.
  morecomment={[l]{\#}},
  commentstyle=\color{cGreen},
}
\lstdefinestyle{cstyle}{
  language=C,basicstyle=\ttfamily\small,
  backgroundcolor=\color{cGreenBg!55},
  frame=leftline,framerule=2.5pt,rulecolor=\color{cGreen},
  numbers=left,numberstyle=\tiny\color{cGray},numbersep=8pt,
  xleftmargin=2.2em,framexleftmargin=0pt,framexrightmargin=0pt,
  breaklines=true,columns=fullflexible,keepspaces=true,
  showstringspaces=false,keywordstyle=\color{cBlue}\bfseries,
  commentstyle=\color{cGray}\itshape,aboveskip=8pt,belowskip=8pt,
}
% v3.36: pystyle/htmlstyle — colored twins of cstyle for the Python/HTML
% listings that used to repeat a ~200-char options blob with NO color
% (keywords rendered bold-black, comments italic-black — measured; C alone
% was colored). Geometry is copied VERBATIM from the shipped blob (frame,
% margins, skips) so swapping a call site to style=pystyle moves NOTHING —
% only keywordstyle/commentstyle/stringstyle are added (bold stays bold,
% italic stays italic, mono metrics identical => zero reflow, proven by
% per-page word-bbox diff). breaklines/breakatwhitespace stay per-site.
% v3.78: numbers=left joins both styles, with the exact numberstyle/numbersep
% keys asm/cstyle carry — fleet-wide consistency (author decision: all four
% styles numbered; "asm/cstyle numbered, py/html not" was an accident of the
% v3.36 geometry freeze, not a design). The numbers hang LEFT of the code box
% (listings sets them with \llap, so no overfull hbox can arise and the code
% x is untouched — measured: code frozen, digits at code_x-11.8pt); inside a
% colored box they are blanked like everyone else's (v3.37 hook), and the
% clean variant restores them in its frameless boxes (v3.78, above).
\lstdefinestyle{pystyle}{
  language=Python,basicstyle=\ttfamily\small,
  columns=fullflexible,keepspaces=true,
  frame=single,rulecolor=\color{cTblLine},
  numbers=left,numberstyle=\tiny\color{cGray},numbersep=8pt,
  xleftmargin=0.5em,framexleftmargin=0pt,framexrightmargin=0pt,
  showstringspaces=false,aboveskip=8pt,belowskip=8pt,
  keywordstyle=\color{cBlue}\bfseries,
  commentstyle=\color{cGray}\itshape,
  stringstyle=\color{cGreen},
}
% v4.5.0: hoisted out of the four styles so an enclosing box can override it.
% Every listing still resolves to true unless something later says otherwise,
% so shipped output is unchanged (verified byte-for-byte on all eight PDFs).
\lstset{resetmargins=true}
\lstdefinestyle{htmlstyle}{
  language=HTML,basicstyle=\ttfamily\small,
  columns=fullflexible,keepspaces=true,
  frame=single,rulecolor=\color{cTblLine},
  numbers=left,numberstyle=\tiny\color{cGray},numbersep=8pt,
  xleftmargin=0.5em,framexleftmargin=0pt,framexrightmargin=0pt,
  showstringspaces=false,aboveskip=8pt,belowskip=8pt,
  keywordstyle=\color{cBlue}\bfseries,
  commentstyle=\color{cGray}\itshape,
  stringstyle=\color{cGreen},
}
\newenvironment{asmblock}{\par\begin{otherlanguage}{english}}{\end{otherlanguage}\par}

% <hllpdf:listing-smart-quotes>
% v4.21.0: four Unicode punctuation characters must enter listings' token
% buffer. Unregistered characters bypass it: a closing quote was printed
% BEFORE the pending word, and an opening quote BEFORE a pending brace.
% Extend the standard character table, then activate only these codepoints
% in the listing's local character-table scope. Do this BEFORE the package's
% delimiter/alsoother/literate hooks so explicit user mappings still win.
% This is not universal Unicode support, nor a conversion to ASCII quotes.
% See listings developer guide 15.3.1 and references/listing-smart-quotes.md.
\makeatletter
\lst@CCPut \lst@ProcessOther {"2018}{"2019}{"201C}{"201D}\z@
\lst@AddToHookAtTop{SelectCharTable}{%
  \lst@ifactivechars
    \catcode"2018=\active \catcode"2019=\active
    \catcode"201C=\active \catcode"201D=\active
  \fi}
\makeatother
% </hllpdf:listing-smart-quotes>

% <hllpdf:listing-technical-symbols>
% v4.30: bounded extension for common technical punctuation in LTR listings.
% Without character-table registration these glyphs bypass the token buffer
% and can move before the preceding word even inside quoted strings.
% Keep the original four-quote block separate for its standing mutation test.
% This is source-faithful printing, not syntax validation or general Unicode.
\makeatletter
\lst@CCPut \lst@ProcessOther
  {"2013}{"2014}{"2026}{"2190}{"2192}{"2194}{"21D2}
  {"2212}{"221E}{"2260}{"2264}{"2265}\z@
\lst@AddToHookAtTop{SelectCharTable}{%
  \lst@ifactivechars
    \catcode"2013=\active \catcode"2014=\active \catcode"2026=\active
    \catcode"2190=\active \catcode"2192=\active \catcode"2194=\active
    \catcode"21D2=\active \catcode"2212=\active \catcode"221E=\active
    \catcode"2260=\active \catcode"2264=\active \catcode"2265=\active
  \fi}
\makeatother
% </hllpdf:listing-technical-symbols>


\usepackage[hidelinks,unicode]{hyperref}   % LAST, after colors. unicode = Hebrew bookmarks.
% v3.62: THE ISOLATION MACROS CORRUPTED EVERY PDF BOOKMARK THEY APPEARED IN.
% \en is defined (line ~155) as
%     \foreignlanguage{english}{\language=\csname l@nohyphenation\endcsname #1}
% and in a heading that body is written verbatim to the .toc and handed to
% hyperref's \pdfstringdef. hyperref strips the two CONTROL SEQUENCES (\language,
% \l@nohyphenation) but the plain `=` CHARACTER token between them is not a control
% sequence, so it survives into the PDF string: \chapter{שער \en{CNOT} והצמדה}
% produced the bookmark "שער =CNOT והצמדה". The printed TOC and the heading on the
% page are CORRECT — only the navigation pane and the viewer window title are wrong,
% which is why it went unnoticed. Measured on the shipped fleet: 70 corrupted
% bookmarks across 5 PDFs, from 177 unprotected \en-in-heading sites. Same corruption
% reached /Title and /Subject via the \hypersetup{pdftitle=...} the docs recommend.
% This hook tells hyperref to use the plain argument for all three isolation macros
% when building a PDF string. It does NOT touch the .toc (the printed TOC is
% byte-unchanged) and it does NOT mask the math signal — a heading with $...$ still
% logs "removing math shift", so BIDI-15 detection is intact.
% BONUS, measured: it also makes \code{} in an UNSTARRED heading COMPILE. BLD-10
% recorded that as needing \texorpdfstring or a starred heading; the real root cause
% is hyperref's \pdfstringdef scan of \code's \ifmmode...\else...\fi body — which is
% exactly why the identically-bodied \texttt survived (hyperref pre-neutralizes
% \texttt in PDF strings) and why \DeclareRobustCommand could never help.
% v3.64: \def\\{ } added. A manual line break in a heading WELDS the two words
% together in the bookmark: \chapter{שורה ראשונה\\שורה שנייה} produced the outline
% entry 'שורה ראשונהשורה שנייה' — "ראשונהשורה" is not a word. hyperref drops the
% \\ and puts nothing in its place, so the break that looks like a space on the
% page is a zero-width join in the navigation pane. The spaced form
% 'מילה אחת \\ מילה שנייה' escapes with a harmless DOUBLE space ('אחת  מילה'),
% which is why this never got noticed. Both are measured on a real outline.
% 0 live instances in the corpus — but SKILL.md and content-style.md both RECOMMEND
% a manual heading break, so this is a primed trap rather than a dormant one, and a
% one-token preamble fix retires it for good instead of asking authors to remember.
% Mapping \\ to a space fixes the WELD, which is the actual defect (a nonexistent
% word in the navigation pane). It does not normalise whitespace: measured after the
% fix, the welded form becomes 'שורה ראשונה שורה שנייה' (correct) and the already-
% spaced form goes from two spaces to three. That is cosmetic — viewers render it as
% a wider gap, and no word is corrupted. A per-heading \texorpdfstring is still the
% way to get an exactly-tuned bookmark string.
\pdfstringdefDisableCommands{\def\en#1{#1}\def\h#1{#1}\def\doclang#1{#1}\def\code#1{#1}\def\codeline#1{#1}\def\\{ }}
% v3.70: RESTORE \url LINE-BREAKING. url.sty freezes its break penalties at LOAD
% time (\mathchardef\UrlBreakPenalty=\binoppenalty), and v3.45's global
% \binoppenalty=10000 (inline-math atomicity) therefore priced EVERY breakpoint
% inside \Url at 10000 = never break. Measured: a long \url in Hebrew prose is
% one unbroken line, ~100pt overfull, invading the margin band the ink gates
% deliberately tolerate. With the two lines below: 0 overfull, the URL breaks
% after '/', the RTL split is CORRECT (fragment order forward), and each line
% gets its own link annotation exactly covering its fragment (measured).
% Confined to \Url's group — inline-math atomicity (TYP-7) is untouched.
\mathchardef\UrlBreakPenalty=200
\mathchardef\UrlBigBreakPenalty=100
\ifx\lakitmainlanguage\hllpdfenglish
  \hypersetup{pdflang=en}
\else
  \hypersetup{pdflang=he}
\fi  % /Lang is the document base locale (not per-span accessibility tagging).
% Set a per-document \hypersetup{pdftitle={...}} too — see SKILL.md.

% v3.52 (round-2 audit): two RTL bibliography/cross-ref defects, both invisible to
% every gate and both affecting the FIRST author to use the feature.
% (a) \@biblabel's default [#1] brackets are bidi-neutral and mirror to ]1[ under
% the RTL bibliography base — every \thebibliography reference list shows ]1[.
% LTR-isolating the label restores [1]. (The render gate's mirror rule fingerprints
% round parens )1(, so square brackets slipped through.)
\makeatletter
\renewcommand*{\@biblabel}[1]{\foreignlanguage{english}{[#1]}}
% v3.54: the SAME neutral-bracket trap on the INLINE side — a multi-key \cite{a,b}
% under the RTL base lays its European-number list right-to-left, so [1, 2] renders
% "[2,1]" and [1, 2, 3] renders "[3,2,1]" (verified by glyph geometry — the reader
% sees the sources in REVERSE order; a note form \cite[p.5]{a} also reorders to
% "[p. 5,1]"). LTR-isolate the whole bracket group (number list AND optional note),
% exactly like \@biblabel above: [1, 2] / [1, 2, 3] / [1, p. 5] all render correct.
% Placed AFTER hyperref so it wins; hyperref hyperlinks each number inside #1 (via
% \@citex), which this wrapper leaves intact. Mirrors the LaTeX kernel \@cite shape.
% v4.9.1: directional isolation must NOT switch the language here. Babel's
% safe-active citation wrapper passes its key list through \bbl@tempa;
% \foreignlanguage reuses that scratch macro before the deferred citation
% loop expands it. The old line therefore cited the literal key "english".
% If such a key existed, even the undefined-reference gate stayed green.
% \babelsublr isolates the bracket group without selecting a locale or
% changing that key list. Preserve native links, key order and optional notes.
\def\@cite#1#2{\babelsublr{[{#1\if@tempswa , #2\fi}]}}
\makeatother

% v4.11.0: a reference inside align's \text inherited HEBdisptag=true,
% so its enclosing hbox mirrored the brackets despite a clean compile.
% Isolate the complete printed reference without selecting a Babel language.
% Also neutralize the display-only normalfont direction reset locally. A
% direction-only patch inverted 12 to 21 inside equation in the measured probe.
% Plain maketag is local to the reference: mathtools' display-side wrapper
% otherwise globally consumes manual_tag while formatting a nested reference.
% Preserve each original command (including the selective auxiliary write),
% current tag forms, native links, and showonlyrefs on/off transitions.
\makeatletter
\NewCommandCopy{\HEB@nativeeqref}{\eqref}
\newcommand{\HEB@eqrefisolate}[1]{%
  \babelsublr{\HEBdisptagfalse\let\bbl@eqnodir\relax\let\maketag@@@\ORIGmaketag@@@#1}}
\protected\def\eqref#1{\HEB@eqrefisolate{\HEB@nativeeqref{#1}}}
\ExplSyntaxOn
\NewCommandCopy{\HEB@nativeMTeqref}{\MT_extended_eqref:n}
\cs_set_protected:Npn \MT_extended_eqref:n #1
  {\HEB@eqrefisolate{\HEB@nativeMTeqref{#1}}}
\ExplSyntaxOff
\makeatother

% v4.11.0: custom caption counters are not part of hyperref's English
% locale, and Hebrew has no upstream hyperref locale in this toolchain.
% Restore the appropriate names at every Babel language switch, including
% return from a nested inline island. Do not localize the destination keys.
\makeatletter
\newcommand{\HEB@refnames}{%
\def\chapterautorefname{פרק}%
\def\sectionautorefname{סעיף}%
\def\subsectionautorefname{תת־סעיף}%
\def\subsubsectionautorefname{תת־תת־סעיף}%
\def\paragraphautorefname{פסקה}%
% v4.31: hyperref uses a distinct counter name at the deepest standard level.
\def\subparagraphautorefname{תת־פסקה}%
\def\equationautorefname{משוואה}%
\def\figureautorefname{איור}%
\def\tableautorefname{טבלה}%
\def\partautorefname{חלק}%
\def\appendixautorefname{נספח}%
\def\footnoteautorefname{הערת שוליים}%
\def\itemautorefname{סעיף}%
\def\Itemautorefname{סעיף}%
\def\figautorefname{איור}%
\def\tblautorefname{טבלה}%
\def\pageautorefname{עמוד}% 
}
\addto\extrashebrew{\HEB@refnames}
\newcommand{\ENG@refnames}{%
  \def\subparagraphautorefname{subparagraph}%
  \def\chapterautorefname{chapter}%
  \def\sectionautorefname{section}%
  \def\subsectionautorefname{subsection}%
  \def\subsubsectionautorefname{subsubsection}%
  \def\paragraphautorefname{paragraph}%
  \def\equationautorefname{Equation}%
  \def\figureautorefname{Figure}%
  \def\tableautorefname{Table}%
  \def\partautorefname{Part}%
  \def\appendixautorefname{Appendix}%
  \def\footnoteautorefname{footnote}%
  \def\itemautorefname{item}%
  \def\figautorefname{Figure}%
  \def\tblautorefname{Table}%
  \def\Itemautorefname{item}%
  \def\pageautorefname{page}%
}
\addto\extrasenglish{\ENG@refnames}
\ifx\lakitmainlanguage\hllpdfenglish\ENG@refnames\else\HEB@refnames\fi
\makeatother

% ============================================================================
% Optional running head (v3.39): \lakithead{<book title>} — OFF by default.
% A document opts in with ONE line after \input{preamble.tex}:
%   \lakithead{מבנה המחשב}
% Every content page then carries the series' orientation line: chapter name
% at the reading edge, book title opposite, thin rule. (fancyhdr's [L]/[R]
% are PHYSICAL; under the RTL page direction [L] lands at the RIGHT/reading
% edge — measured, consistent across the corpus.) Packaged like the topic
% rail: always loaded, inert without the call. This centralizes the 8-line
% block that 8 lakits used to copy by hand — the same duplication vector
% that once let a lakit preamble go ~100 lines stale. Class-aware: report
% fills the mark from \chapter, article from \section. NB: a titlepage env
% self-empties; front matter that sets \pagestyle{empty} must restore with
% \pagestyle{fancy} afterwards (NOT plain — that discards the head).
% ============================================================================
\usepackage{fancyhdr}
% v3.61: \chapter* left a STALE running head — pages of an UNNUMBERED chapter
% carried the PREVIOUS chapter's title. report.cls routes \chapter* through
% \@schapter, which never calls \chaptermark, so \leftmark keeps whatever the last
% numbered chapter left. Measured: with \chapter{פרק ממוספר} then
% \chapter*{מבוא לא ממוספר}, page 4 (inside the starred chapter) printed
% "פרק ממוספר". Invisible on the chapter's FIRST page (\thispagestyle{plain}), so it
% only shows from the second page on — the pages an author is least likely to check.
% An unnumbered front-matter chapter (\chapter*{מבוא}, \chapter*{הקדמה}) is everyday
% Hebrew-academic structure, and the corpus already carries HAND-WRITTEN \markboth
% lines right after \chapter* — the exact per-document duplication \lakithead exists
% to retire. Patched INSIDE \lakithead so the macro stays inert when unused; a
% document's own later \markboth still wins, so those hand-patches do not regress.
\makeatletter
% <hllpdf:starred-heading-scope>
\newif\iflakin@starmarkpatched
\newcommand{\lakistarmark}{%
  \iflakin@starmarkpatched\else
  % Dispatch on the actual class counter, not hyperref's shared internals.
  \ifdefined\c@chapter
    \let\HEBorigschapter\@schapter
    % Record only after the native heading has reached its page. A mark
    % before the native routine can advertise the next heading too early.
    \def\@schapter##1{\HEBorigschapter{##1}\chaptermark{##1}}%
  \else
    % @ssect receives layout arguments, NOT the section name or level.
    % Capture the name in @startsection without introducing a TeX group:
    % grouping sectioning would break the native after-heading state.
    \def\HEB@sectionkind{}%
    \def\HEB@mainsection{section}%
    \let\HEBorigstartsection\@startsection
    \def\@startsection##1##2##3##4##5##6{%
      \def\HEB@sectionkind{##1}%
      \HEBorigstartsection{##1}{##2}{##3}{##4}{##5}{##6}}%
    \let\HEBorigssect\@ssect
    \def\@ssect##1##2##3##4##5{%
      \HEBorigssect{##1}{##2}{##3}{##4}{##5}%
      % Lower starred headings must not replace the article's section mark.
      \ifx\HEB@sectionkind\HEB@mainsection
        \sectionmark{##5}%
      \fi
      \ignorespaces}%
  \fi
  \global\lakin@starmarkpatchedtrue
  \fi}
% </hllpdf:starred-heading-scope>
% <hllpdf:running-language>
% Choose the wrapper when the mark is written, not when a later page ships.
% English punctuation must remain LTR even after returning to Hebrew prose.
\newcommand{\HEB@runningmark}[1]{%
  \iflanguage{english}{\markboth{\en{#1}}{}}{\markboth{\h{#1}}{}}}
% </hllpdf:running-language>
\makeatother
% v4.2.1 (NAV): every collection already names itself in \lakithead — feed the
% PDF /Title from the same declaration, so the navigation gate's
% PDF_REQUIRE_TITLE has a universal source and no fleet member ships
% title-less by missing a second declaration site. Measured before this: 14
% of 16 collections shipped /Title-empty while every one carried a perfectly
% good \lakithead string (plain Hebrew in all 16 — no math, no formatting, so
% the PDF-string path is clean). An explicit title always wins: set it with
% \lakitpdftitle{...} anywhere before \lakithead, or with a plain
% \hypersetup{pdftitle=...} — the fallback detects that via hyperref's
% \@pdftitle, probe-verified \@empty-equal until the first explicit set.
\makeatletter
\newif\iflakit@pdftitleset
\newcommand{\lakitpdftitle}[1]{%
  \hypersetup{pdftitle={#1}}\global\lakit@pdftitlesettrue}
\newcommand{\lakithead}[1]{%
  \iflakit@pdftitleset\else
    \ifx\@pdftitle\@empty
      \hypersetup{pdftitle={#1}}\global\lakit@pdftitlesettrue
    \fi
  \fi
  \pagestyle{fancy}\fancyhf{}%
  \lakistarmark

  \fancyhead[L]{\small\textcolor{cGray}{\leftmark}}%
  % <hllpdf:fixed-head-language>
  % The fixed document title has its OWN language snapshot, independent of
  % the current section/chapter mark. Choose the branch NOW, not at shipout:
  % the page may be Hebrew by the time an English title is finally printed.
  \iflanguage{english}{%
    \fancyhead[R]{\small\textcolor{cGray}{\en{#1}}}%
  }{%
    \fancyhead[R]{\small\textcolor{cGray}{\h{#1}}}%
  }%
  % </hllpdf:fixed-head-language>
  \fancyfoot[C]{\thepage}%
  \renewcommand{\headrulewidth}{0.4pt}%
  \ifdefined\chaptermark
    \renewcommand{\chaptermark}[1]{\HEB@runningmark{##1}}%
  \else
    \renewcommand{\sectionmark}[1]{\HEB@runningmark{##1}}%
  \fi}
\makeatother

% v3.39.1: path-independent byte-reproducibility. LuaTeX derives the trailer
% /ID from the OUTPUT PATH, so two clones of the same source at different
% directories produced different sha256 (measured: the ONLY delta was the
% 62-byte /ID line). Bit 512 suppresses exactly that line — measured: Creator/
% Producer/CreationDate/ModDate and all content untouched. /ID is optional
% per ISO 32000 (required only for ENCRYPTED PDFs and PDF/A — neither is a
% target here). If a future workflow needs an ID (PDF/A, Acrobat revision
% matching), delete this line and it returns.
\pdfvariable suppressoptionalinfo \numexpr 512\relax

% v3.75: BOLD MATH IN BOLD HEADINGS. Heading text is \bfseries but math inside
% it stayed regular-weight CM — a hairline $V_T$ inside a heavy 24.8pt chapter
% title (measured: CMMI12 fl=6 next to FrankRuhlLibre-Bold fl=20; 155 live
% heading-with-math sites fleet-wide). \boldmath rides along wherever the class
% sets \bfseries for a heading; A/B measured: $V_T$ -> CMMIB10 fl=22, 0 missing
% characters (the bold CM faces ship with this TeX Live). Guarded per-macro so
% the article class (no \@makechapterhead) and future class changes stay safe.
% v4.0.1 (adoption fix): the SEARCH patterns above carry no parameter token.
% v4.0.0 made them specific by including the class's own #1/#2, and measured on
% TeX Live 2023 that never matches: a \patchcmd search pattern containing a
% parameter token cannot match while the call sits inside \@ifundefined{...}{}{...}
% (probed both #1 and ##1 -> FAIL; the same pattern at top level succeeds). Every
% one of these patches therefore failed, and because v4.0.0 also replaced the
% empty failure branch with \PackageError, the failure became FATAL: all four
% bundled example assets stopped building ("Cannot patch chapter-title math
% weight -> Fatal error, no output PDF"). The \PackageError branches are KEPT —
% a silent no-op here is what let v3.75's fix rot unnoticed — only the parameter
% tokens are dropped, so the patches match and a real class change still errors.
\makeatletter
\@ifundefined{@makechapterhead}{}{%
  \patchcmd{\@makechapterhead}{\huge\bfseries \@chapapp}
    {\huge\bfseries\boldmath \@chapapp}{}
    {\PackageError{hebrew-lualatex-pdf}{Cannot patch chapter-number math weight}{Unsupported document class}}%
  \patchcmd{\@makechapterhead}{\Huge \bfseries}
    {\Huge \bfseries\boldmath}{}
    {\PackageError{hebrew-lualatex-pdf}{Cannot patch chapter-title math weight}{Unsupported document class}}%
  \patchcmd{\@makeschapterhead}{\Huge \bfseries}
    {\Huge \bfseries\boldmath}{}
    {\PackageError{hebrew-lualatex-pdf}{Cannot patch starred chapter math weight}{Unsupported document class}}}
% v4.0.1: \part sits under a THREE-DEEP wrapper chain, so the class-based anchor
% is not in \@part at all. Measured with \show: hyperref rewrites \@part to
% `...\phantomsection\fi \H@old@part`, and nameref rewrites THAT to
% `\NR@gettitle{#1}\NR@part[{#1}]{#2}` — the class body ends up in \NR@part.
% v4.0.0 patched \@part and made failure fatal, which is why all seven bundled
% examples stopped building. \part is NOT a dead heading level here: 188 live uses
% across the fleet, so the bold-math fix is worth getting right rather than
% dropping. We walk the chain innermost-first and stop at the first hit. Failure is
% a WARNING, not an error: the chain is a property of hyperref/nameref versions, so
% a future package release must degrade to regular-weight part math — a cosmetic
% loss — instead of refusing to build. Chapter and section keep \PackageError,
% because those targets are in the CLASS and were verified to patch here.
\newif\iflakin@partnum \newif\iflakin@parttitle \newif\iflakin@sparttitle
\def\lakin@partnum#1#2{\iflakin@partnum\else
  \patchcmd{#1}{#2\bfseries \partname}{#2\bfseries\boldmath \partname}
    {\lakin@partnumtrue}{}\fi}
\def\lakin@parttitle#1#2{\iflakin@parttitle\else
  \patchcmd{#1}{#2 \bfseries}{#2 \bfseries\boldmath}{\lakin@parttitletrue}{}\fi}
\def\lakin@sparttitle#1#2{\iflakin@sparttitle\else
  \patchcmd{#1}{#2 \bfseries}{#2 \bfseries\boldmath}{\lakin@sparttitletrue}{}\fi}
\@ifundefined{c@chapter}{\def\lakin@pnum{\Large}\def\lakin@ptit{\huge}}%
                        {\def\lakin@pnum{\huge}\def\lakin@ptit{\Huge}}
\@ifundefined{NR@part}{}{\expandafter\lakin@partnum\expandafter\NR@part\expandafter{\lakin@pnum}%
                         \expandafter\lakin@parttitle\expandafter\NR@part\expandafter{\lakin@ptit}}
\@ifundefined{H@old@part}{}{\expandafter\lakin@partnum\expandafter\H@old@part\expandafter{\lakin@pnum}%
                            \expandafter\lakin@parttitle\expandafter\H@old@part\expandafter{\lakin@ptit}}
\@ifundefined{@part}{}{\expandafter\lakin@partnum\expandafter\@part\expandafter{\lakin@pnum}%
                       \expandafter\lakin@parttitle\expandafter\@part\expandafter{\lakin@ptit}}
\@ifundefined{NR@spart}{}{\expandafter\lakin@sparttitle\expandafter\NR@spart\expandafter{\lakin@ptit}}
\@ifundefined{H@old@spart}{}{\expandafter\lakin@sparttitle\expandafter\H@old@spart\expandafter{\lakin@ptit}}
\@ifundefined{@spart}{}{\expandafter\lakin@sparttitle\expandafter\@spart\expandafter{\lakin@ptit}}
\iflakin@partnum\else\PackageWarning{hebrew-lualatex-pdf}{Part-number math stays
  regular weight: no known \noexpand\part wrapper matched (hyperref/nameref
  version change). Cosmetic only}\fi
\iflakin@parttitle\else\PackageWarning{hebrew-lualatex-pdf}{Part-title math stays
  regular weight: no known \noexpand\part wrapper matched. Cosmetic only}\fi
% v4.12.0: THE PART NUMBER IN THE TOC WAS NOT DIRECTIONALLY SEPARATED FROM THE
% TITLE. report.cls/article.cls compose the part entry as ONE text stream --
% \addcontentsline{toc}{part}{\thepart\hspace{1em}#1} -- while EVERY other level
% goes through \numberline, which puts the number in an \hbox. A box is a bidi
% OBJECT and cannot join a neighbouring run; a bare "IV" can. So in the RTL TOC
% paragraph a Latin number and a Latin-INITIAL title merge into a SINGLE LTR run
% and are laid out left-to-right as a unit: the number lands on the wrong side of
% the row, and a mixed title is scrambled outright. Measured on \part{MCP ...}:
% the printed order was "<tail> V MCP" -- title fragment, number, title fragment.
% The chapter/section rows one line below, with the SAME Latin-initial titles,
% were correct on the same page. That contrast is the whole diagnosis.
% \en{} IS NOT A WORKAROUND, measured both ways on one fixture: it fixes a
% PURE-Latin title (\part{\en{Hooks}} renders correctly) and does NOT fix a MIXED
% one (\part{\en{MCP} <hebrew>} stays scrambled). babel's language switch changes
% font and hyphenation, not the bidi run structure -- so the defect is not
% reachable from the source side at all, which is why it belongs here.
% Boxing the number is the entire fix and is a STRICT NO-OP where the bug never
% bit: a Hebrew-titled document compiles BYTE-IDENTICAL before and after (md5,
% 3 clean runs), because there the neutral glue already sat between an L run and
% an R run. The PDF outline is byte-identical too -- hyperref's \pdfstringdef
% neutralises \mbox, so bookmark strings keep their part numbers.
% \numberline{\thepart} would also fix the direction, but it re-geometries every
% part row in the fleet and DROPS the number from the part bookmark (hyperref
% strips \numberline). \mbox changes nothing that was already right.
% The search pattern deliberately carries no # token: etoolbox cannot match a
% parameter token in a patch pattern (measured -- all three chain members
% reported "not patchable" while the pattern included #1). "\thepart \hspace
% {1em}" is brace-balanced and occurs exactly once in the wrapper body.
\newif\iflakin@parttocnum
\def\lakin@parttocnum#1{\iflakin@parttocnum\else
  \patchcmd{#1}{\thepart \hspace {1em}}{\mbox {\thepart }\hspace {1em}}
    {\lakin@parttocnumtrue}{}\fi}
\@ifundefined{NR@part}{}{\lakin@parttocnum\NR@part}
\@ifundefined{H@old@part}{}{\lakin@parttocnum\H@old@part}
\@ifundefined{@part}{}{\lakin@parttocnum\@part}
\iflakin@parttocnum\else\PackageWarning{hebrew-lualatex-pdf}{Part TOC number is
  not boxed: no known \noexpand\part wrapper matched. A part title that BEGINS
  with Latin will print its number on the wrong side of its TOC row}\fi
\patchcmd{\section}{\bfseries}{\bfseries\boldmath}{}
  {\PackageError{hebrew-lualatex-pdf}{Cannot patch section math weight}{Unsupported document class}}
\patchcmd{\subsection}{\bfseries}{\bfseries\boldmath}{}
  {\PackageError{hebrew-lualatex-pdf}{Cannot patch subsection math weight}{Unsupported document class}}
\patchcmd{\subsubsection}{\bfseries}{\bfseries\boldmath}{}
  {\PackageError{hebrew-lualatex-pdf}{Cannot patch subsubsection math weight}{Unsupported document class}}
\patchcmd{\paragraph}{\bfseries}{\bfseries\boldmath}{}{}
\patchcmd{\subparagraph}{\bfseries}{\bfseries\boldmath}{}{}
\@ifundefined{l@chapter}{}{%
  \patchcmd{\l@chapter}{\bfseries}{\bfseries\boldmath}{}{}}
\makeatother
\setcounter{tocdepth}{\lakittocdepth}
% <hllpdf:measured-toc>
% v4.16.0: Measure actual number labels, not a guessed digit count. Read only
% canonical contentsline records; never execute titles or addtocontents code
% during the measuring pass. Widths are local to this individual list and are
% available before its first row: no per-row drift and no extra aux pass.
\makeatletter
\newread\hllpdf@tocstream
\newdimen\hllpdf@toc@chapter
\newdimen\hllpdf@toc@section
\newdimen\hllpdf@toc@subsection
\newdimen\hllpdf@toc@subsubsection
\newdimen\hllpdf@toc@paragraph
\newdimen\hllpdf@toc@subparagraph
\newdimen\hllpdf@toc@figure
\newdimen\hllpdf@toc@table
\newdimen\hllpdf@tocpage
\newdimen\hllpdf@toc@sectionindent
\newdimen\hllpdf@toc@subsectionindent
\newdimen\hllpdf@toc@subsubsectionindent
\newdimen\hllpdf@toc@paragraphindent
\newdimen\hllpdf@toc@subparagraphindent
\def\hllpdf@tocdefaults{%
  \hllpdf@toc@chapter=2em
  \hllpdf@toc@section=3.2em
  \hllpdf@toc@subsection=3.9em
  \@ifundefined{c@chapter}{%
    \hllpdf@toc@subsubsection=3.2em
    \hllpdf@toc@paragraph=4.1em
    \hllpdf@toc@subparagraph=5em
  }{%
    \hllpdf@toc@subsubsection=4.1em
    \hllpdf@toc@paragraph=5em
    \hllpdf@toc@subparagraph=6em
  }%
  \hllpdf@toc@figure=2.3em
  \hllpdf@toc@table=2.3em
  \hllpdf@tocpage=\@pnumwidth\relax
}
\def\hllpdf@tocindents{%
  % Preserve the established shallow layout when its slots already fit.
  \hllpdf@toc@sectionindent=1.5em
  \@tempdima=\hllpdf@toc@chapter \advance\@tempdima by -2em
  \advance\hllpdf@toc@sectionindent by \@tempdima
  \hllpdf@toc@subsectionindent=\dimexpr
    \hllpdf@toc@sectionindent+\hllpdf@toc@section\relax
  \hllpdf@toc@subsubsectionindent=\dimexpr
    \hllpdf@toc@subsectionindent+\hllpdf@toc@subsection\relax
  \hllpdf@toc@paragraphindent=\dimexpr
    \hllpdf@toc@subsubsectionindent+\hllpdf@toc@subsubsection\relax
  \hllpdf@toc@subparagraphindent=\dimexpr
    \hllpdf@toc@paragraphindent+\hllpdf@toc@paragraph\relax
}
\def\hllpdf@tocfont@chapter{\bfseries\boldmath}
\def\hllpdf@tocfont@part{\large\bfseries\boldmath}
\def\hllpdf@tocentryfont{%
  \ifcsname hllpdf@tocfont@\hllpdf@toclevel\endcsname
    \csname hllpdf@tocfont@\hllpdf@toclevel\endcsname
  \fi}
\long\def\hllpdf@tocdiscard#1\hllpdf@tocstop{}
\def\hllpdf@tocscanrecord{%
  \@ifnextchar\contentsline{\hllpdf@tocscanline}{%
    \@ifnextchar\babel@toc{\hllpdf@tocscanlanguage}{\hllpdf@tocdiscard}}}
\long\def\hllpdf@tocscanlanguage\babel@toc#1#2#3\hllpdf@tocstop{%
  % Babel records are declarations, not arbitrary code. Track the language
  % without executing the original directive or changing the surrounding list.
  \def\hllpdf@toclanguage{#1}}
\long\def\hllpdf@tocscanline\contentsline#1#2#3#4#5\hllpdf@tocstop{%
  \def\hllpdf@toclevel{#1}%
  \ifcsname hllpdf@toc@#1\endcsname
    \hllpdf@tocscantitle#2\hllpdf@tocstop
  \fi
  % The page-number column needs its own measurement (including bold chapter
  % rows). Its extra clearance separates the page label from dot leaders.
  \setbox\@tempboxa=\hbox{\foreignlanguage{\hllpdf@toclanguage}{\hllpdf@tocentryfont #3\kern.3em}}%
  \ifdim\wd\@tempboxa>\hllpdf@tocpage
    \hllpdf@tocpage=\wd\@tempboxa
  \fi}
\def\hllpdf@tocscantitle{%
  \@ifnextchar\numberline{\hllpdf@tocscanlabel}{\hllpdf@tocdiscard}}
\long\def\hllpdf@tocscanlabel\numberline#1#2\hllpdf@tocstop{%
  \setbox\@tempboxa=\hbox{\foreignlanguage{\hllpdf@toclanguage}{\hllpdf@tocentryfont #1\kern.5em}}%
  \ifdim\wd\@tempboxa>\csname hllpdf@toc@\hllpdf@toclevel\endcsname
    \csname hllpdf@toc@\hllpdf@toclevel\endcsname=\wd\@tempboxa
  \fi}
\def\hllpdf@tocread{%
  \ifeof\hllpdf@tocstream\else
    \read\hllpdf@tocstream to\hllpdf@tocrecord
    \expandafter\hllpdf@tocscanrecord\hllpdf@tocrecord\hllpdf@tocstop
    \expandafter\hllpdf@tocread
  \fi}
\def\hllpdf@tocmeasure#1{%
  \makeatletter
  \edef\hllpdf@toclanguage{\languagename}%
  \hllpdf@tocdefaults
  \openin\hllpdf@tocstream=\jobname.#1\relax
  \hllpdf@tocread
  \closein\hllpdf@tocstream
  \hllpdf@tocindents
  \edef\@pnumwidth{\the\hllpdf@tocpage}%
  \@tempdima=\dimexpr\hllpdf@tocpage+1em\relax
  \ifdim\@tempdima>\@tocrmarg\relax
    \edef\@tocrmarg{\the\@tempdima}%
  \fi}
\let\hllpdf@originalstarttoc\@starttoc
\def\@starttoc#1{%
  \begingroup
    \def\hllpdf@tocext{#1}%
    \def\hllpdf@tocwanted{toc}%
    \ifx\hllpdf@tocext\hllpdf@tocwanted \hllpdf@tocmeasure{#1}\else
      \def\hllpdf@tocwanted{lof}%
      \ifx\hllpdf@tocext\hllpdf@tocwanted \hllpdf@tocmeasure{#1}\else
        \def\hllpdf@tocwanted{lot}%
        \ifx\hllpdf@tocext\hllpdf@tocwanted \hllpdf@tocmeasure{#1}\fi
      \fi
    \fi
    \hllpdf@originalstarttoc{#1}%
  \endgroup}
\hllpdf@tocdefaults
\hllpdf@tocindents
\@ifundefined{l@chapter}{}{%
  \patchcmd{\l@chapter}{\setlength\@tempdima{1.5em}}%
    {\setlength\@tempdima{\hllpdf@toc@chapter}}{}%
    {\PackageError{hebrew-lualatex-pdf}{Cannot size chapter TOC numbers}%
       {Unsupported contents layout: the chapter number slot was not found}}}
\renewcommand*\l@section{\@dottedtocline{1}%
  {\hllpdf@toc@sectionindent}{\hllpdf@toc@section}}
\renewcommand*\l@subsection{\@dottedtocline{2}%
  {\hllpdf@toc@subsectionindent}{\hllpdf@toc@subsection}}
\renewcommand*\l@subsubsection{\@dottedtocline{3}%
  {\hllpdf@toc@subsubsectionindent}{\hllpdf@toc@subsubsection}}
\renewcommand*\l@paragraph{\@dottedtocline{4}%
  {\hllpdf@toc@paragraphindent}{\hllpdf@toc@paragraph}}
\renewcommand*\l@subparagraph{\@dottedtocline{5}%
  {\hllpdf@toc@subparagraphindent}{\hllpdf@toc@subparagraph}}
\renewcommand*\l@figure{\@dottedtocline{1}{1.5em}{\hllpdf@toc@figure}}
\renewcommand*\l@table{\@dottedtocline{1}{1.5em}{\hllpdf@toc@table}}
\makeatother
% </hllpdf:measured-toc>
\setcounter{secnumdepth}{\lakitsecnumdepth}
% Public preamble-only caption-numbering policy. The historical within-major
% default is untouched unless this command is called. "continuous" removes the
% reset edge as well as the printed prefix, so Figure 1/2 and Table 1/2 stay
% unique across sections/chapters. Returning to "within-major" restores the
% original class-aware displays and native destination conventions. The policy
% never replaces \figcap/\tabcap or their measured link-placement mechanism.
\makeatletter
\let\lakin@majorfig\thefig
\let\lakin@majortbl\thetbl
\let\lakin@majorHfig\theHfig
\let\lakin@majorHtbl\theHtbl
\newcommand{\lakitcaptionnumbering}[1]{%
  \edef\lakin@numbering{#1}%
  \def\lakin@continuous{continuous}\def\lakin@withinmajor{within-major}%
  \ifx\lakin@numbering\lakin@continuous
    \@ifundefined{c@chapter}{%
      \@removefromreset{fig}{section}\@removefromreset{tbl}{section}%
    }{\@removefromreset{fig}{chapter}\@removefromreset{tbl}{chapter}}%
    \renewcommand{\thefig}{\arabic{fig}}\renewcommand{\thetbl}{\arabic{tbl}}%
    \def\theHfig{\arabic{fig}}\def\theHtbl{\arabic{tbl}}%
  \else\ifx\lakin@numbering\lakin@withinmajor
    \@ifundefined{c@chapter}{%
      \@removefromreset{fig}{section}\@removefromreset{tbl}{section}%
      \@addtoreset{fig}{section}\@addtoreset{tbl}{section}%
    }{%
      \@removefromreset{fig}{chapter}\@removefromreset{tbl}{chapter}%
      \@addtoreset{fig}{chapter}\@addtoreset{tbl}{chapter}}%
    \let\thefig\lakin@majorfig\let\thetbl\lakin@majortbl
    \let\theHfig\lakin@majorHfig\let\theHtbl\lakin@majorHtbl
  \else
    \PackageError{hebrew-english-latex-studio}{Unknown caption-numbering policy '#1'}%
      {Use continuous or within-major before begin{document}.}%
  \fi\fi}
\@onlypreamble\lakitcaptionnumbering
\makeatother
% Architecture readiness sentinel. Optional runtime layers require the complete
% core, so the marker is defined only after every mechanism/baseline definition.
% <hllpdf:architecture-ready>
\def\hllpdfcoreloaded{1}
% </hllpdf:architecture-ready>

% v4.22: registered file resources, not an assertion about embedded subset bytes.
% The log framing avoids long-path wrapping and preserves UTF-8 as hex.
% <hllpdf:font-inventory>
\directlua{
  local function inventory()
    local names = {}
    local unnamed = 0
    for id, item in font.each() do
      local path = item.filename
      if path and not (path == "") then names[path] = true else unnamed = unnamed + 1 end
    end
    local sorted = {}
    for path in pairs(names) do table.insert(sorted,path) end
    table.sort(sorted)
    texio.write_nl("log","HLLFONT|BEGIN|1")
    for index,path in ipairs(sorted) do
      local hex = path:gsub(".",function(c) return string.format(string.char(37) .. "02x",string.byte(c)) end)
      local count = math.ceil(string.len(hex)/48)
      texio.write_nl("log","HLLFONT|FILE|" .. index .. "|" .. count)
      for offset=1,string.len(hex),48 do texio.write_nl("log","HLLFONT|HEX|" .. hex:sub(offset,offset+47)) end
    end
    texio.write_nl("log","HLLFONT|END|" .. rawlen(sorted) .. "|" .. unnamed)
    texio.write_nl("log","")
  end
  luatexbase.add_to_callback("finish_pdffile",inventory,"hllpdf.runtime.fonts")
}
% </hllpdf:font-inventory>
; if omitted, the historical canonical
% defaults below are exact.  They deliberately control only the document shell
% and page layout.  RTL/math semantics stay in this preamble, while visual
% presentation themes (for example clean_theme.tex) load AFTER it.
%
%   \def\lakitdocumentclass{article}
%   \def\lakitclassoptions{11pt,a4paper}
%   \def\lakitgeometryoptions{margin=1.8cm,headheight=14pt}
%   \def\lakitlinespread{1.15}
%   \def\lakitparagraphstyle{block} % letters/briefs; indented for continuous prose
%   \def\lakittocdepth{1}
%   \def\lakitsecnumdepth{2}
%   % ============================================================================
% preamble.tex — tested LuaLaTeX preamble for Hebrew + math + figures PDFs (v4.32.0)
% Battle-tested: compiled an 80-page Hebrew quantum-computing summary AND a
% 197-page Hebrew computer-architecture anthology (TikZ diagrams, code listings,
% colored boxes) — both with 0 missing characters.
%
% Usage in your document:
%   \input{preamble.tex}     % (or paste these lines at the top)
%   \begin{document}  ...  \end{document}
%
% Compile with:  lualatex -interaction=nonstopmode doc.tex
%   Repeat until two consecutive PDFs are byte-identical (scripts/build.sh
%   does this) — a TOC document converges only on pass 3+; after two passes
%   the TOC page numbers are STALE (measured; see SKILL.md step 3).
% Font binaries are external dependencies, pinned by assets/fonts/SHA256SUMS.
% scripts/setup_fonts.sh imports a verified local installation or prior skill.
%
% PUBLIC DOCUMENT-SHELL / LAYOUT CONTRACT (introduced v4.3.1; architecture-owned v4.7.1). These hooks may be
% defined BEFORE \input{preamble.tex}; if omitted, the historical canonical
% defaults below are exact.  They deliberately control only the document shell
% and page layout.  RTL/math semantics stay in this preamble, while visual
% presentation themes (for example clean_theme.tex) load AFTER it.
%
%   \def\lakitdocumentclass{article}
%   \def\lakitclassoptions{11pt,a4paper}
%   \def\lakitgeometryoptions{margin=1.8cm,headheight=14pt}
%   \def\lakitlinespread{1.15}
%   \def\lakitparagraphstyle{block} % letters/briefs; indented for continuous prose
%   \def\lakittocdepth{1}
%   \def\lakitsecnumdepth{2}
%   \input{preamble.tex}
%
% Supported class contract: report or article.  Alternate paper sizes must also
% be declared to build.sh through EXPECTED_PAPER; the PDF gate never infers them.
% Do NOT use these hooks as a theme mechanism: theme overlays must not own page
% geometry, class selection, or direction semantics. A declared presentation-policy
% overlay may intentionally change numbering visibility after the preamble.
% ============================================================================
% Runtime architecture guard: generated from scripts/architecture_contract.py.
% The readiness sentinel remains delayed until the END of this file.
% <hllpdf:architecture-guard>
\ifdefined\hllpdfcoreloaded
  \errmessage{hllpdf: preamble.tex loaded more than once}
  \expandafter\endinput
\fi
% </hllpdf:architecture-guard>

\ifdefined\lakitdocumentclass\else\def\lakitdocumentclass{report}\fi
\ifdefined\lakitclassoptions\else\def\lakitclassoptions{12pt,a4paper}\fi
\ifdefined\lakitgeometryoptions\else\def\lakitgeometryoptions{margin=2.3cm,headheight=14pt}\fi
\ifdefined\lakitlinespread\else\def\lakitlinespread{1.3}\fi
\ifdefined\lakittocdepth\else\def\lakittocdepth{1}\fi
\ifdefined\lakitsecnumdepth\else\def\lakitsecnumdepth{2}\fi
% Paragraph organization belongs to the document layout, independently of its
% language and theme. Preserve historical class dimensions unless explicitly
% selecting block paragraphs; never suppress indentation with an everypar hack.
\ifdefined\lakitparagraphstyle\else\def\lakitparagraphstyle{indented}\fi
\edef\lakitparagraphstyle{\lakitparagraphstyle}
\def\hllpdfparagraphblock{block}
\def\hllpdfparagraphindented{indented}
\ifx\lakitparagraphstyle\hllpdfparagraphblock\else
  \ifx\lakitparagraphstyle\hllpdfparagraphindented\else
    \errmessage{hllpdf: paragraph style must be block or indented}
  \fi
\fi
% v4.3.1: the public class hook has a deliberately narrow support contract.
% Before this guard, \def\lakitdocumentclass{book} silently compiled even though
% every heading/TOC regression test targets report/article; an unsupported shell
% could therefore look "supported" until a subtle class-specific defect shipped.
% Fail before \documentclass rather than letting an undocumented combination enter
% the mechanism core. Plain-TeX comparison is used because LaTeX packages are not
% available yet at this point.
\edef\lakitvalidatedclass{\lakitdocumentclass}
\def\lakitreportclass{report}
\def\lakitarticleclass{article}
\ifx\lakitvalidatedclass\lakitreportclass
\else\ifx\lakitvalidatedclass\lakitarticleclass
\else
  \errmessage{hllpdf: unsupported class \lakitdocumentclass; use report or article}
\fi\fi
\let\lakitvalidatedclass\relax
\let\lakitreportclass\relax
\let\lakitarticleclass\relax
% Language intent is independent of document class, theme and content domain.
% The two dominant-language profiles differ in authoring intent, not direction.
% A balanced bilingual document must explicitly choose a navigation/base locale.
% <hllpdf:language-profile>
\providecommand{\lakitlanguageprofile}{hebrew}
\edef\lakitlanguageprofile{\lakitlanguageprofile}
\def\hllpdfhebrew{hebrew}\def\hllpdfenglish{english}
\def\hllpdfhebrewmixed{hebrew-with-english}
\def\hllpdfenglishmixed{english-with-hebrew}
\def\hllpdfbilingual{bilingual}
\let\hllpdfexpectedlanguage\relax
\ifx\lakitlanguageprofile\hllpdfhebrew
  \let\hllpdfexpectedlanguage\hllpdfhebrew
\else\ifx\lakitlanguageprofile\hllpdfhebrewmixed
  \let\hllpdfexpectedlanguage\hllpdfhebrew
\else\ifx\lakitlanguageprofile\hllpdfenglish
  \let\hllpdfexpectedlanguage\hllpdfenglish
\else\ifx\lakitlanguageprofile\hllpdfenglishmixed
  \let\hllpdfexpectedlanguage\hllpdfenglish
\else\ifx\lakitlanguageprofile\hllpdfbilingual
  \ifdefined\lakitmainlanguage\else
    \errmessage{hllpdf: bilingual profile requires explicit lakitmainlanguage}
  \fi
\else
  \errmessage{hllpdf: unsupported language profile \lakitlanguageprofile}
\fi\fi\fi\fi\fi
\ifdefined\lakitmainlanguage
  \edef\lakitmainlanguage{\lakitmainlanguage}
\else
  \let\lakitmainlanguage\hllpdfexpectedlanguage
\fi
\ifx\lakitmainlanguage\hllpdfhebrew\else
  \ifx\lakitmainlanguage\hllpdfenglish\else
    \errmessage{hllpdf: main language must be hebrew or english}
  \fi
\fi
\ifx\hllpdfexpectedlanguage\relax\else
  \ifx\lakitmainlanguage\hllpdfexpectedlanguage\else
    \errmessage{hllpdf: main language conflicts with language profile}
  \fi
\fi
\typeout{HLLLANG|\lakitlanguageprofile|\lakitmainlanguage}
% </hllpdf:language-profile>
\documentclass[\lakitclassoptions]{\lakitdocumentclass}
\usepackage[no-math]{fontspec}        % (1) ★ no-math — THE critical line (see recipe.md)
% layout=lists asks Babel to adapt list parshape when paragraph and page
% directions differ. Without it, minority-language list bodies lose their
% indentation and labels hang outside the text block, despite a clean log.
% The upstream mechanism preserves ordinary same-direction list geometry.
\usepackage[bidi=basic,layout=lists]{babel}
\ifx\lakitmainlanguage\hllpdfenglish
  \babelprovide[import]{hebrew}
  \babelprovide[import,main]{english}
\else
  \babelprovide[import,main]{hebrew}
  \babelprovide[import]{english}
\fi
% v4.22: bind canonical families and fallbacks to the selected installation.
% Lookup by basename or family alone is not proof of the file actually used.
% <hllpdf:font-binding>
\directlua{
  local dir = os.getenv("HEBREW_LATEX_FONT_DIR")
  if not dir or dir == "" then
    local data = os.getenv("XDG_DATA_HOME")
    if not data or data == "" or not (data:sub(1,1) == "/") then
      data = (os.getenv("HOME") or "") .. "/.local/share"
    end
    dir = data .. "/fonts/hebrew-lualatex-pdf"
  end
  if not (dir:sub(1,1) == "/") then dir = lfs.currentdir() .. "/" .. dir end
  if dir:find("[",1,true) or dir:find("]",1,true) or dir:find(string.char(0),1,true)
     or dir:find(string.char(10),1,true) or dir:find(string.char(13),1,true) then
    tex.error("hllpdf: unsupported character in font directory")
  end
  if not (lfs.attributes(dir,"mode") == "directory") then
    tex.error("hllpdf: font directory is missing")
  end
  hllpdf_font_dir = dir:gsub("/+$", "") .. "/"
  tex.sprint(string.char(92) .. "def" .. string.char(92) .. "hllpdffontdir{")
  tex.sprint(-2,hllpdf_font_dir)
  tex.sprint("}")
  luaotfload.add_fallback("hebfb",{"[" .. hllpdf_font_dir .. "DejaVuSerif.ttf]:mode=node;"})
  luaotfload.add_fallback("hebfbbf",{"[" .. hllpdf_font_dir .. "DejaVuSerif-Bold.ttf]:mode=node;"})
}
% </hllpdf:font-binding>
% v3.22: Frank Ruhl Libre ships NO italic face. Instead of silently-upright
% \emph (v<=3.20) or bold-as-emphasis (v3.21), synthesize a slanted face with
% FakeSlant — real, VISIBLE italic emphasis with correct \emph nesting
% semantics (inner \emph toggles back upright). Verified empirically at
% 300dpi: slant clearly visible, 0 errors, 0 missing glyphs.
\babelfont{rm}[Path={\hllpdffontdir},RawFeature={fallback=hebfb},
  BoldFont={FrankRuhlLibre-Bold.ttf},
  BoldFeatures={RawFeature={fallback=hebfbbf}},   % v3.56 — see hebfbbf above
  ItalicFont={FrankRuhlLibre-Regular.ttf}, ItalicFeatures={FakeSlant=0.22},
  % Bold-italic uses the weight-matched fallback and a synthetic slant.
  BoldItalicFont={FrankRuhlLibre-Bold.ttf}, BoldItalicFeatures={RawFeature={fallback=hebfbbf},FakeSlant=0.22}
  ]{FrankRuhlLibre-Regular.ttf}
% Sans fallbacks preserve family and weight for the few missing codepoints.
\directlua{luaotfload.add_fallback("hebfbsans",{"[" .. hllpdf_font_dir .. "DejaVuSans.ttf]:mode=node;"})}
\directlua{luaotfload.add_fallback("hebfbsansbf",{"[" .. hllpdf_font_dir .. "DejaVuSans-Bold.ttf]:mode=node;"})}
\babelfont{sf}[Path={\hllpdffontdir},RawFeature={fallback=hebfbsans},
  BoldFont={Heebo-Bold.ttf},
  BoldFeatures={RawFeature={fallback=hebfbsansbf}}]{Heebo-Regular.ttf}
% v3.24: Scale=MatchLowercase — DejaVu Sans Mono has an outsized x-height, so
% unscaled inline code visually SHOUTS next to the delicate Frank Ruhl body
% (measured on a 197-page build: scaled, a code atom sits in the line at the
% body's weight, and the document tightened by ~5%). Applies to listings too.
\babelfont{tt}[Path={\hllpdffontdir},RawFeature={fallback=hebfb},Scale=MatchLowercase,
  BoldFont={DejaVuSansMono.ttf},
  BoldFeatures={RawFeature={fallback=hebfb},FakeBold=1.5},
  ItalicFont={DejaVuSansMono.ttf},
  ItalicFeatures={RawFeature={fallback=hebfb},FakeSlant=0.2},
  BoldItalicFont={DejaVuSansMono.ttf},
  BoldItalicFeatures={RawFeature={fallback=hebfb},FakeBold=1.5,FakeSlant=0.2}
  ]{DejaVuSansMono.ttf}

% Explicit multilingual/English profiles use real Latin faces and italics.
% The legacy/default Hebrew profile retains its established measured typography.
% Latin Modern is provided by the TeX distribution, not the pinned font bundle;
% recorder/runtime evidence reports it as an additional toolchain resource.
\ifx\lakitlanguageprofile\hllpdfhebrew\else
  \babelfont[english]{rm}[Ligatures=TeX,RawFeature={fallback=hebfb},
    BoldFeatures={RawFeature={fallback=hebfbbf}},
    BoldItalicFeatures={RawFeature={fallback=hebfbbf}}]{Latin Modern Roman}
  \babelfont[english]{sf}[Ligatures=TeX,RawFeature={fallback=hebfbsans},
    BoldFeatures={RawFeature={fallback=hebfbsansbf}},
    BoldItalicFeatures={RawFeature={fallback=hebfbsansbf}}]{Latin Modern Sans}
\fi

\usepackage{amsmath,amssymb,mathtools}
% amsthm supplies theorem/proof machinery. Environment names remain
% document-specific; the generic proof label is localized below.
\usepackage{amsthm}
% Proof names follow the surrounding prose, including return from an island.
\addto\extrashebrew{\renewcommand{\proofname}{הוכחה}}
\addto\extrasenglish{\renewcommand{\proofname}{Proof}}
\ifx\lakitmainlanguage\hllpdfenglish
  \renewcommand{\proofname}{Proof}
\else
  \renewcommand{\proofname}{הוכחה}
\fi
% v3.24.1: equation tags under babel-bidi. A tag re-emitted by mathtools'
% showonlyrefs inside align inherits the RTL base and its parens MIRROR —
% ")2.1(" instead of "(2.1)" (plain align/equation tags are unaffected;
% verified on a minimal case + a 198-page build). LTR-isolate at the tag-BOX
% level: mathtools' show/hide + numbering logic is untouched. Do NOT override
% \tagform@ instead — that clobbers the showonlyrefs mechanism (measured:
% every tag shows, numbering breaks).
% v3.28.1: the v3.24.1 wrap must apply ONLY to tags of display environments.
% Inline \eqref in Hebrew PROSE also goes through \maketag@@@, and there the
% english-wrap MIRRORS the parens — ")1(" — while the unwrapped original is
% correct (proven on a 2x2 isolation matrix: patch x showonlyrefs; regressed
% every prose \eqref since v3.24.1, caught by the render gate's R4). The two
% call sites are indistinguishable by \ifmmode (tag boxes are text mode), so
% the display environments themselves raise a flag; \eqref in prose sees it
% false and gets the original untouched \maketag@@@.
\usepackage{etoolbox}

% babel bidi=basic (including upstream 26.9) assumes every automatic
% discretionary has a one-glyph replacement. LuaTeX can legally produce one
% with replace=nil after a hyphenated Latin value; Babel then indexes nil and the
% whole build dies. Do the narrow upstream-compatible correction: while Babel
% classifies the list, present only that malformed disc as an ordinary
% discretionary, then restore its subtype. Babel still reorders the ENTIRE box;
% there is no pcall and no successful fallback to an unreordered node list.
% If Babel throws for any other reason, the build remains fatal. The synthetic
% nil-disc probe and the real table gallery in test_all.sh pin both layers.
\AtBeginDocument{\directlua{
  if Babel and Babel.bidi and not Babel.heb_nil_disc_compat then
    Babel.heb_nil_disc_compat = true
    local original = Babel.bidi
    local disc_id = node.id('disc')
    Babel.bidi = function(head, ispar, hdir)
      local hidden = {}
      for item in node.traverse_id(disc_id, head) do
        if item.subtype == 2 and (not item.replace or not item.replace.char) then
          table.insert(hidden, item)
          item.subtype = 0
        end
      end
      local result = original(head, ispar, hdir)
      for _, item in ipairs(hidden) do item.subtype = 2 end
      return result
    end
  end
}}
\makeatletter
\newif\ifHEBdisptag
\let\ORIGmaketag@@@\maketag@@@
\def\maketag@@@#1{\ifHEBdisptag\ORIGmaketag@@@{\foreignlanguage{english}{#1}}\else\ORIGmaketag@@@{#1}\fi}
\@for\HEB@env:=equation,equation*,align,align*,gather,gather*,multline,multline*,flalign,flalign*,alignat,alignat*,eqnarray,eqnarray*\do{%
  \expandafter\AtBeginEnvironment\expandafter{\HEB@env}{\HEBdisptagtrue}}
\makeatother
% English page + Hebrew prose: Babel's equation-only tag flip uses a primitive
% leqno path that shifts the formula to the left edge and the tag into the body
% of an English page. Align's full-width alignment is unaffected. Give only this
% display a physical LTR container and native amsmath eqno/leqno. The engine then
% owns real tag width, collision handling and labels; no proxy tag box hides ink
% from its fit calculation. Babel's saved \bbl@eqnodir still restores Hebrew text
% inside math, and the prose locale/math direction are never switched.
% Close the prior RTL paragraph before normalization. After the environment's
% group restores its directions, close TeX's empty resumed display paragraph and
% use native LaTeX continuation handling: same-paragraph prose has no indent,
% while an authored blank line starts an ordinary indented RTL paragraph.
% Register after Babel's setup; remain local to equation(*) in a main-English
% document's actual RTL prose. Native Hebrew, ordinary English and explicit
% fleqn retain their original paths. No page/display dimensions are assigned.
% <hllpdf:minority-rtl-equation-geometry>
\makeatletter
\newcommand{\HEB@minorityprosecontinuation}{\par\@doendpe}
\newcommand{\HEB@minorityequationgeometry}{%
  \ifx\lakitmainlanguage\hllpdfenglish
    \ifnum\pardirection=\@ne
      \if@fleqn\else
        \par\aftergroup\HEB@minorityprosecontinuation
        \iftagsleft@\let\veqno\@@leqno
        \else\let\veqno\@@eqno\fi
        \bodydir TLT\pardir TLT\textdir TLT\relax
      \fi
    \fi
  \fi}
\AtBeginDocument{%
  \AddToHook{env/equation/begin}{\HEB@minorityequationgeometry}%
  \AddToHook{env/equation*/begin}{\HEB@minorityequationgeometry}}
\makeatother
% </hllpdf:minority-rtl-equation-geometry>
% v3.63: \shortintertext rendered Hebrew MIRROR-REVERSED and flush LEFT — an entire
% commentary line inside an align printed backwards. babel's bidi=basic patches
% \intertext@ only (luababel.def), and mathtools' \shortintertext uses the identical
% \vbox{\normalbaselines … \noindent\ignorespaces#1\par} skeleton but is never
% patched, so its paragraph keeps pardir/textdir LTR: no bidi reordering (logical
% order on the page) and flush-left. Measured against \intertext with the same
% Applying babel's own replacement to \shortintertext@ makes the two identical.
% \h{}/\mbox fix the letter order but leave the line flush LEFT, so they are not
% substitutes. Guarded + \AtBeginDocument because mathtools \let-copies at load time.
% NB \ExplSyntaxOn is REQUIRED: mathtools' internals carry expl3 names, and without
% the expl3 catcodes the `_` is read as a subscript and the patch detonates far away
% with "! Illegal unit of measure (pt inserted)" at the \end{align} (measured).
\makeatletter
\AtBeginDocument{%
  \@ifundefined{shortintertext@}{}{%
    \ExplSyntaxOn
    \bbl@sreplace\shortintertext@{\normalbaselines}%
      {\normalbaselines\ifx\bbl@eqnodir\relax\else\bbl@pardir\@ne\bbl@eqnodir\fi}%
    \ExplSyntaxOff}}
\makeatother
% geometry inherits the canonical A4 sheet from \documentclass; build.sh then
% verifies every rendered MediaBox and CropBox, so a silent Letter fallback is
% impossible to publish as a clean default build.
\usepackage[\lakitgeometryoptions]{geometry}
\linespread{\lakitlinespread}       % layout hook; canonical default remains 1.3
\ifx\lakitparagraphstyle\hllpdfparagraphblock
  \setlength{\parindent}{0pt}
  \setlength{\parskip}{6pt plus 1pt minus 1pt}
\fi
% These dimensions remain ordinary scoped TeX lengths: authors can customize
% them after core input. Language switches inherit them; lists keep their own
% indentation and narrow reference layouts establish their local paragraph gaps.
\typeout{HLLPAR|\lakitparagraphstyle}
\usepackage{xcolor}
\usepackage{mdframed}
\usepackage{listings}                 % code listings (assembly / C / pseudo-code)
\usepackage{tikz}
\usetikzlibrary{arrows.meta,positioning,calc,shapes.geometric,fit,backgrounds}
\usepackage{pgfplots}\pgfplotsset{compat=1.18}
\usepackage{quantikz}                 % quantum circuits (remove if unused)
\usepackage{booktabs}
\usepackage{array}
\usepackage{float}                    % [H] = place table exactly here (reference doc: no drifting)
\usepackage{enumitem}
\usepackage{needspace}

% ============================================================================
% ★★ v3.4 — THE DIRECTION-ISOLATION TRIO (the heart of mixing Hebrew + Latin)
% Mental model:
%   \h{...}    Hebrew (RTL) island      — Hebrew text inside an LTR context
%                                          (TikZ nodes, the english `tikzpic`).
%   \en{...}   English (LTR) isolate    — Latin text / parentheses inside a
%                                          Hebrew (RTL) context. KEEPS PARENS
%                                          AND ORDER CORRECT.
%   \code{...} \en + monospace          — code inside Hebrew / inside a frametitle.
% CRITICAL: never wrap Latin in \h{} (it is for HEBREW). Inside a TikZ node,
% Latin in \h{} renders REVERSED (Datapath->htapataD, CPU->UPC) and parens
% mirror. Isolate every Latin/paren chunk with \en{} or \code{}. See
% references/figures-boxes-listings.md.
% ---- TEXT INSIDE MATH (v3.8) ----------------------------------------------
% Put ANY text that sits inside math in a \text{}. The ONLY thing that is
% invalid in math is a bare \h/\en (= \foreignlanguage) used DIRECTLY, outside
% \text -> "\rmfamily invalid in math mode". Inside \text:
%   Hebrew label           -> \text{...}            (bare Hebrew renders fine;
%                              \h is optional — BUT NOT ALWAYS HARMLESS (v3.67):
%                              if the \h group OPENS with a nested island
%                              ($..$, \en{}, \code{}, \texttt{}, \mbox{}) the
%                              WHOLE group loses bidi and the Hebrew prints in
%                              LOGICAL order = gibberish. Measured:
%                                \text{\h{$f$ קבועה}}  -> reads "העובק"  BROKEN
%                                \text{\h{קבועה $f$}}  -> correct
%                                \text{$f$ קבועה}      -> correct (no \h)
%                              Hebrew first inside \h, or drop the \h. Flagged
%                              H-OPENS-WITH-LTR-GROUP.)
%   plain Latin word        -> \text{...}            (no marker needed)
%   Latin WITH ( ) or [ ]   -> \text{\en{...}}       (\en forces LTR; a plain
%                              \text mirrors the parens in RTL -> "Spec (TNR(" )
%   f(Hebrew/text arg)      -> \text{f}(\text{...})  -- keep the parens at MATH
%                              level, NOT \text{f(arg)} (parens inside the box mirror)
%   code in math            -> \mathtt{...}  (or \texttt{}, math-guarded below)
% NEVER write \en{...} or \h{...} straight into $...$ / \[..\] (outside \text);
% always via \text{}. (\mbox{...} also works but does not scale in scripts.)
% ============================================================================
\newcommand{\h}[1]{\foreignlanguage{hebrew}{#1}}
% v3.42.1 — \en islands do NOT hyphenate (auto OR at explicit hyphens). English hyphenation patterns were
% tearing identifiers and terms ACROSS THE RTL SEAM ("IO Regis-/ters",
% "In-/put", "Band-/width", "Per-/formance", "Multi-/ple" — all shipped on
% real pages; probe-proven: "ALUS-/rcB", "Mem-/ToReg"). An island is a label
% or a term, not English prose — a hyphen tear there is never right. A/B on
% a full 197pp lakit: identical overfull sets, same page count, ~11
% paragraphs healed, one mild underfull (badness 3525) that replaced a
% mutilated "Regis-/ters". Long English PROSE belongs in the englishprose
% environment below, which restores explicit-hyphen breaks while its paragraph
% is still active. A raw otherlanguage block inherits the global prohibition.
% (\l@nohyphenation needs \csname:
% '@' is not a letter here.)
% v3.43.2 — \exhyphenpenalty moved OUT of \en to a GLOBAL assignment below.
% Render-proven: a group-local \exhyphenpenalty is DEAD CODE — TeX reads the
% line-breaking parameters at \par time, after the \en group has closed, so
% \en{write-through} / plain not-taken still tore at the "-" ("write-/
% through" shipped on a real exam page). \language is different: it is
% recorded in the node list as text is added, so the in-group nohyphenation
% assignment above DOES work (and stays). The global penalty was A/B-measured
% on a 197pp lakit: identical overfull set (25), identical page count, the
% not-/taken tear healed. Hebrew is unaffected in practice — Hebrew compounds
% use maqaf (־), not ASCII "-".
\newcommand{\en}[1]{\foreignlanguage{english}{\language=\csname l@nohyphenation\endcsname #1}}
% Babel runs extras locally for inline islands and paragraph language changes.
% An English main document must not inherit the Hebrew no-hyphen paragraph rule.
\addto\extrashebrew{\exhyphenpenalty=10000\relax}
\addto\extrasenglish{\exhyphenpenalty=50\relax}
\ifx\lakitmainlanguage\hllpdfenglish
  \exhyphenpenalty=50\relax
\else
  \exhyphenpenalty=10000\relax
\fi
% Full English paragraphs need ordinary explicit-hyphen breaks. Use this
% environment instead of a raw otherlanguage block; the paragraph ends while
% the local penalty is still active.
\newenvironment{englishprose}
  {\par\begin{otherlanguage}{english}\exhyphenpenalty=50\relax}
  {\par\end{otherlanguage}}
\newenvironment{hebrewprose}
  {\par\begin{otherlanguage}{hebrew}\exhyphenpenalty=10000\relax}
  {\par\end{otherlanguage}}
% Use for shell text whose locale must not depend on the last body paragraph.
\DeclareRobustCommand{\doclang}[1]{\foreignlanguage{\lakitmainlanguage}{#1}}

% ============================================================================
% ★★ — fixes for defects that are INVISIBLE in the log and in the checks
% and only surface on a VISUAL RENDER of the PDF. (Each verified empirically.)
% ----------------------------------------------------------------------------
% v3.22: \emph is now REAL slanted emphasis (FakeSlant italic face defined on
% the \babelfont{rm} line above), replacing the v3.21 \emph->\textbf hack: it
% preserves \emph's toggle semantics and keeps bold free for its own role.
% Prefer bold-as-emphasis for your doc? one line: \renewcommand{\emph}[1]{\textbf{#1}}
% page number: force the digits LTR so they never reverse (18 -> 81) in the RTL
%   footer. That reversal is a NON-DETERMINISTIC bidi glitch (observed once in
%   ~80 pp); this LuaLaTeX build is not byte-reproducible, so the defensive
%   LTR wrap is the only reliable guard. Latin digits are always LTR anyway.
\renewcommand{\thepage}{\foreignlanguage{english}{\arabic{page}}}
% v3.70: \pagenumbering \gdef's \thepage and silently DESTROYS the wrap above —
% after \pagenumbering{arabic} the folio guard is gone for the whole rest of the
% document. Two collections discovered this independently and hand-restore the
% wrap in main.tex with the same comment; that knowledge now lives here instead.
% The wrapper re-applies the LTR wrap after the kernel call, for any style.
\makeatletter
\let\lakin@orig@pagenumbering\pagenumbering
\renewcommand{\pagenumbering}[1]{%
  \lakin@orig@pagenumbering{#1}%
  \renewcommand{\thepage}{\foreignlanguage{english}{\csname @#1\endcsname\c@page}}}
\makeatother
% emergencystretch: Hebrew does not hyphenate, so an unbreakable LTR island
%   (\en{...} / an inline math run) landing at the margin can leave a small
%   (~2-13pt) overfull with no legal breakpoint. This rescues ONLY those
%   inline-prose tie-breaks; it is a ceiling used only on offending lines and
%   does NOT touch display-math / listing / table overflow (those still surface
%   in the log for the content-level overflow pass). Verified: baseline overfull
%   -> 0 with this set, with no change to pagination or underfull.
\setlength{\emergencystretch}{3em}
% v3.40: inline math (almost) never breaks across lines. TeX's default
%   (binop 700 / rel 500) lets an inline run split at an operator —
%   $(1.f)\cdot 2^{e}$ shipped as "(1.f)·" ending one line and "2^e" opening
%   the next (reader-reported). In LTR that split is merely ugly; under an
%   RTL base the carried operator lands on the visually WRONG side of the
%   break and the formula reads backwards.
%   WHY 10000 (v3.45.0 — the display-first resolution). Inline math is
%   FORBIDDEN from breaking, so a torn island under an RTL base is
%   structurally impossible. The v3.40 cycle used 9999 ("break only as a last
%   resort") because a bare 10000 left ~30-39 WIDE inline equations
%   overflowing (worst 90.7pt) — but 9999's "last resort" fires OFTEN in
%   tight, non-hyphenating Hebrew: a reader found torn islands across the
%   fleet (~150 candidates; e.g. "P(y |" ending one line and "x_0)" opening
%   the next, read backwards). The step the 9999 cycle missed: every one of
%   those ~30 overflows is a wide equation that BELONGS in display. Convert
%   them to display FIRST, then 10000 gives 0 tears AND 0 overflows (measured,
%   12/12 fleet). Doctrine: inline math is ATOMIC; a wide equation goes in a
%   display, never inline. See content-style ("inline math is atomic").
\binoppenalty=10000 \relpenalty=10000
% v3.38: uniform sentence spacing. TeX's default (\nonfrenchspacing, sfcode
% 3000) is the OLD Anglo-American wide-space-after-period convention — Hebrew
% typography never had it, and modern Latin practice (Chicago et al.) dropped
% it too. Measured on real pages: post-period gaps up to 1.9-2.1x the line's
% mean word gap — the dominant "rivers" source in justified Hebrew prose.
% Survives babel language switches (probed: sfcode stays french inside \en{}
% islands and otherlanguage blocks).
\frenchspacing
% v3.38: widow/orphan protection. LaTeX's default penalties (150/150/50)
% allow a lone line at a page top/bottom almost freely. 10000 = forbid;
% this is SAFE here precisely because the document is \raggedbottom (the
% demoted line just moves on, no vertical stretch fight) — under flushbottom
% these values would cause underfull pages. raggedbottom is the report-class
% default; made explicit so a future class/option change cannot silently
% flip it and turn the 10000s harmful.
\raggedbottom
\widowpenalty=10000 \clubpenalty=10000 \displaywidowpenalty=10000
% em-dash: --- does NOT become — under Frank Ruhl (the font lacks tlig; verified
%   the width is identical with/without Ligatures=TeX, while the mechanism works
%   on Latin Modern). TYPE THE CHARACTER — (U+2014) DIRECTLY; likewise – (U+2013).
%   check_bidi_figures.py flags a --/--- that is adjacent to Hebrew.
% ============================================================================

% ---- STRUCTURAL DIRECTION-SAFETY for \texttt (v3.7) -------------------------
% A parenthesis, bracket, or $ INSIDE \texttt{} mirrors/migrates in RTL prose
% (MIPS 0($s0)->"(s0$)0", $rs->"rs$", arr[i]->"[arr[i") because \texttt sets the
% font but NOT the direction. Redefining \texttt to LTR-isolate its content fixes
% the WHOLE class at once (parens, brackets, $) so plain \texttt{lw $t0,0($s0)}
% is safe — no need to remember \code{}. Guarded by \ifmmode because
% \foreignlanguage is invalid in math (there, plain monospace is already LTR).
% NB (corrected v3.56 — the old note here said "stays tofu", which is WRONG):
% \texttt{Hebrew} is a FATAL build failure. The tt family is DejaVu Sans Mono
% (zero Hebrew) and the tt fallback hebfb is DejaVuSerif (also zero Hebrew), so
% the subsetter gets an empty glyph set and LuaTeX dies with
%   "! error: (file .../DejaVuSansMono.ttf) (type 2): there are no glyphs in the
%    subset"  -> no output PDF.
% The message names a FONT FILE, never Hebrew/\texttt/the source line. And it is
% non-local: add an unrelated Latin \code{} elsewhere and the fatal degrades into
% a SILENT DELETION of the Hebrew word. Put Hebrew in \h{}/normal text, never in
% \texttt{}/\code{}; flagged HEBREW-IN-MONO. To disable, delete this block.
\let\origtexttt\texttt
\protected\long\def\texttt#1{\ifmmode\origtexttt{#1}\else\foreignlanguage{english}{\origtexttt{#1}}\fi}
% v4.8.0: defer the direction test until TeX has inserted an alignment cell's
% math template. Unprotected expansion here tested \ifmmode too early at the
% FIRST token of aligned/array/cases/matrix cells. Engine-level \protected also
% preserves the original \texttt implementation without the self-recursion
% observed when retrofitting \DeclareRobustCommand to its saved alias.
\protected\long\def\code#1{\ifmmode\origtexttt{#1}\else\foreignlanguage{english}{\origtexttt{#1}}\fi}

% ---- \codeline — display-code line (v3.24): the code analogue of display math.
% Reader-friction rule (content-style.md §5b): inline \code{} is for a SINGLE
% ATOM only (one register/variable/mnemonic, a tight no-space expression).
% A multi-token expression — an instruction with operands, a spaced assignment,
% a register-binding list — zigzags the eye RTL<->LTR mid-sentence and, being
% an unbreakable island, line-breaks badly. Give it its own centered line with
% display-math vertical rhythm. \centerline = one box => the expression never
% wraps; \nopagebreak keeps it glued to its introducing sentence.
%   ...the Hebrew sentence ends with a colon:
%   \codeline{lw \$t1, -4(\$t0)}
%   and the prose continues (escape dollars as \$).
% v3.66 (first line of the body): flush a PENDING \item label before the \par.
% LaTeX defers the label box (\@labels) until the next \everypar, and an \hbox in
% vertical mode never fires \everypar — so `\item \codeline{...}` printed the code
% at y=69.1 and the item's own "1." BELOW it at y=103.6 (the number appears after
% its content, attached to nothing). \if@inlabel\leavevmode\fi is CONDITIONAL on
% purpose: an unconditional \leavevmode would open an empty paragraph line when
% \codeline follows ordinary prose. Verified both ways — in a list the label now
% sits on the code's line at its normal x, and after prose the geometry is
% byte-identical to before.
\makeatletter
\newcommand{\codeline}[1]{%
  \if@inlabel\leavevmode\fi
  \par\nopagebreak
  \vspace{\abovedisplayskip}\nopagebreak
  % v3.62: was \centerline{...}, i.e. \hbox to\hsize{\hss #1 \hss}. \hss has
  % INFINITE SHRINK, so TeX could always "fit" the box and NEVER reported an
  % Overfull — the same \hss mechanism as BLD-27's \item label. Measured: a wide
  % \codeline in a defbox painted its first token ON the box's left border bar and
  % sat 9.5pt outside the frame interior on both sides, with 0 Overfull, 0 Underfull
  % and all four gates PASS; in bare prose it reached 21.0pt into BOTH margins.
  % \hfil has no shrink, so the overflow is reported and build.sh's ovf_vis/ovf_big
  % counters see it; \linewidth (not \hsize) makes it box-aware, so a \codeline
  % inside a colored box centres on the box interior. Verified: a FITTING \codeline
  % keeps byte-identical glyph positions, and the box-crossing case now emits
  % "Overfull \hbox (19.03pt too wide)". A wide \codeline should be split in two.
  \hbox to\linewidth{\hfil\code{#1}\hfil}%
  \nopagebreak\vspace{\belowdisplayskip}}
\makeatother

% ---- a thin decorative rule helper ----
\newcommand{\sep}{\par\vspace{2pt}\noindent\textcolor{cRule}{\rule{\linewidth}{0.6pt}}\par\vspace{4pt}}

% ---- colors ----
\definecolor{cBlue}{HTML}{1F4E79}   \definecolor{cBlueBg}{HTML}{EAF1F8}
\definecolor{cGreen}{HTML}{2E6B4F}  \definecolor{cGreenBg}{HTML}{EAF4EE}
\definecolor{cAmber}{HTML}{9C6A00}  \definecolor{cAmberBg}{HTML}{FBF3E2}
\definecolor{cPurple}{HTML}{5B2C6F} \definecolor{cPurpleBg}{HTML}{F3ECF7}
\definecolor{cTeal}{HTML}{0E6E6E}   \definecolor{cTealBg}{HTML}{E6F4F4}
\definecolor{cRed}{HTML}{8E2B2B}    \definecolor{cRedBg}{HTML}{FBEAEA}
\definecolor{cGray}{HTML}{555555}   \definecolor{cRule}{HTML}{1F4E79}

% ---- math operators ----
\DeclareMathOperator{\Tr}{Tr}
\DeclareMathOperator{\rank}{rank}
\DeclareMathOperator{\diag}{diag}
\DeclareMathOperator{\sgn}{sgn}

% ---- bra-ket (quantikz defines \ket/\bra/\braket; add the rest safely) ----
\providecommand{\ket}[1]{\left|#1\right\rangle}
\providecommand{\bra}[1]{\left\langle#1\right|}
\providecommand{\braket}[2]{\left\langle#1\middle|#2\right\rangle}
\newcommand{\dyad}[1]{\left|#1\right\rangle\!\left\langle#1\right|}
% v3.21: general outer product |a><b| and expectation value <A> — agents reach
% for these by reflex (esp. \expval); quantikz defines neither, so their absence
% threw "Undefined control sequence" in a real build. \providecommand = safe.
\providecommand{\ketbra}[2]{\left|#1\right\rangle\!\left\langle#2\right|}
\providecommand{\expval}[1]{\left\langle#1\right\rangle}
\newcommand{\abs}[1]{\left|#1\right|}
\newcommand{\norm}[1]{\left\lVert#1\right\rVert}
\newcommand{\half}{\tfrac12}
\newcommand{\R}{\mathbb{R}}\newcommand{\C}{\mathbb{C}}\newcommand{\Z}{\mathbb{Z}}
% identity operator. NB: \mathbb{1} silently renders a WRONG glyph (msbm has no
% blackboard digits). Use \mathbb{I} = 𝕀 (works with amssymb, no extra package).
% Do NOT "fix" this back to \mathbb{1}.
\newcommand{\id}{\mathbb{I}}

% ---- physical units (v3.35): \punit — upright unit with thin space ---------
% Usage: $300\punit{K}$, $80\punit{mV/decade}$, $10^{7}\punit{cm/s}$ — the
% NUMBER AND UNIT LIVE IN ONE MATH ISLAND. Promoted verbatim from a 199-page
% production lakit (advanced-electronic-devices, dozens of uses).
% Two measured traps this design avoids (both render-proven):
%   1. siunitx SCRAMBLES under RTL — \SI{3.5}{eV} renders "eV 3.5", exponents
%      detach ("Hz 10 x 10^9"): its quantity assembly is not bidi-isolated.
%      Do NOT load siunitx in this pipeline.
%   2. A text-mode unit next to a bare-prose number is TWO adjacent LTR runs
%      and reorders ("5 mm" -> "mm5" — the BIDI-2/3 class). \punit is therefore
%      math-only ON PURPOSE: prose misuse fails LOUDLY at compile (\mathrm
%      outside math) instead of scrambling silently.
% Domain shorthands (\MeV, \fm...) stay per-document: \newcommand{\MeV}{\,\text{\en{MeV}}}
\providecommand{\punit}[1]{\,\mathrm{#1}}

% ============================================================================
% Colored boxes — mdframed (NOT tcolorbox, which breaks under bidi).
% ★★ v3.4: innerleftmargin/innerrightmargin give code listings room so their
% frame/line-numbers do not spill out the box border (see listing styles below).
% ============================================================================
% v3.61: everyline/split skips were applied to `tablebox` only (BLD-14, v3.37) —
% the six colored boxes below never got them, so they still split OPEN across a page
% break: the page-N segment gets top+sides and NO bottom rule, the page-N+1 segment
% sides+bottom and NO top rule. Measured on the shipped fleet this is a COVERAGE GAP,
% not a re-report of BLD-14: 0 split-open `tablebox` segments (that fix holds) vs 62
% split-open COLORED boxes across 7 documents. Verified after the fix: both segments
% close, and the frametitle is NOT duplicated on the continuation. BLD-14's caveat
% still applies — a box taller than a full page grinds mdframed's split-trial loop, so
% split such content manually.
% ============================================================================
% v3.72: frametitle={} leads EVERY \newmdenv option list below. mdframed options
% are GROUP-scoped: a box nested inside \begin{exbox}[frametitle={...}] used to
% INHERIT the outer title — measured, the identical title string rendered twice,
% the second riding a title-bar fill inside the INNER box's frame, for all seven
% environments in both directions. Both natural forms were broken (= preamble
% bug, not an authoring rule). The explicit reset makes every box titleless
% unless ITS OWN [frametitle={...}] says otherwise; verified both directions:
% no duplicate on the nested box, an explicit title still renders.
% Geometry is shared with the optional semantic-width table container. Themes
% select their own existing semantic style below; neither path fixes columns.
\mdfdefinestyle{hllpdfsemanticgeometry}{linewidth=1.1pt,
  innertopmargin=6pt,innerbottommargin=7pt,skipabove=9pt,skipbelow=9pt,
  innerleftmargin=9pt,innerrightmargin=9pt,everyline=true,
  splitbottomskip=6pt,splittopskip=12pt,roundcorner=3pt}
\newcommand{\hllpdfsemanticframestyle}{hllpdfsemanticgeometry}
\newmdenv[frametitle={},style=hllpdfsemanticgeometry,linecolor=cBlue,backgroundcolor=cBlueBg,
  frametitlebackgroundcolor=cBlue,frametitlefontcolor=white,
  ]{defbox}   % הגדרה
\newmdenv[frametitle={},style=hllpdfsemanticgeometry,linecolor=cGreen,backgroundcolor=cGreenBg,
  frametitlebackgroundcolor=cGreen,frametitlefontcolor=white,
  ]{thmbox}   % משפט/טענה
\newmdenv[frametitle={},style=hllpdfsemanticgeometry,linecolor=cAmber,backgroundcolor=cAmberBg,
  frametitlebackgroundcolor=cAmber,frametitlefontcolor=white,
  ]{notebox}  % הערה
\newmdenv[frametitle={},style=hllpdfsemanticgeometry,linecolor=cPurple,backgroundcolor=cPurpleBg,
  frametitlebackgroundcolor=cPurple,frametitlefontcolor=white,
  ]{exbox}    % דוגמה
\newmdenv[frametitle={},style=hllpdfsemanticgeometry,linecolor=cRed,backgroundcolor=cRedBg,
  frametitlebackgroundcolor=cRed,frametitlefontcolor=white,
  ]{warnbox}  % טעות נפוצה
\newmdenv[frametitle={},style=hllpdfsemanticgeometry,linecolor=cTeal,backgroundcolor=cTealBg,
  frametitlebackgroundcolor=cTeal,frametitlefontcolor=white,
  ]{keybox}   % רעיון מרכזי

% ---- codebox (v4.5.0): the dedicated container for code -------------------
% An OPTION, not a policy: asmblock unboxed stays canonical for a full listing,
% and a snippet inside any of the six colored boxes still works as before.
% What makes codebox worth its own environment is that it is the ONE framed box
% that keeps its margin line numbers.
%
% v4.4.0 shipped it WITHOUT them, on the strength of v3.37's "no margin knob can
% save them" plus four fresh probes that agreed. That was wrong: every probe had
% moved the code, the box, or both — never the key that carries the NUMBER
% COLUMN. Measured on the page:
%   resetmargins=false   the listing follows THIS box instead of the page,
%                        moving the code from x=75.26 to x=97.54
%   xrightmargin=26pt    under RTL this is the PHYSICAL LEFT indent, and the
%                        only key that carries the \llap'ed digits with it:
%                        x=64.28 -> 90.18, clearing the rule at 65.64 by 23.6pt
%   xleftmargin=26pt     moves the code and LEAVES THE DIGITS on the rule —
%                        the asymmetry that hid the fix for eleven versions
% resetmargins is hoisted out of the four styles (see the \lstset below them) so
% this override is possible at all; every other listing still resolves to true.
% Pinned by tests/codebox_numbers.tex + probe_lstnum.py --inside-rule, whose
% negative control strips xrightmargin and watches the digits fall back.
\definecolor{cSlate}{HTML}{3F4A56}  \definecolor{cSlateBg}{HTML}{F2F4F7}
\newmdenv[frametitle={},leftline=true,rightline=true,topline=true,bottomline=true,
  linecolor=cSlate,backgroundcolor=cSlateBg,linewidth=0.9pt,
  frametitlebackgroundcolor=cSlate,frametitlefontcolor=white,
  innertopmargin=6pt,innerbottommargin=7pt,skipabove=9pt,skipbelow=9pt,
  innerleftmargin=9pt,innerrightmargin=9pt,everyline=true,splitbottomskip=6pt,splittopskip=12pt,
  roundcorner=3pt]{codebox}   % קוד
% The two keys that make in-box numbering work, both measured on the page:
%   resetmargins=false  the listing follows THIS box instead of the page
%   xrightmargin        under RTL this is the PHYSICAL LEFT indent, and it is
%                       the only one that carries the \llap'ed number column
%                       with it. xleftmargin moves the code and leaves the
%                       digits behind on the rule (measured: code 97.54,
%                       digits 64.28 — the v3.37 defect, reproduced); with
%                       xrightmargin=26pt the digits land at x=90.18, a
%                       clearance of 23.6pt from the rule at 65.64.
% This asymmetry is why v3.4's note had the semantics backwards and v3.37
% concluded the digits could not be saved.
% <hllpdf:listing-container-state>
% v4.29: all state is local to the current environment group. A semantic
% callout masks number text, while a nested codebox restores the formatter
% saved before that mask. Save it once across successive semantic wrappers;
% an explicit numbers=none (or an explicitly empty formatter) remains valid.
\makeatletter
\newif\ifhllpdf@incodebox
\newif\ifhllpdf@englishcodebox
\newif\ifhllpdf@lstnumbermasked
\let\hllpdf@savedlstnumber\thelstnumber
\makeatother
% </hllpdf:listing-container-state>
% <hllpdf:codebox-gutter>
\makeatletter
\AtBeginEnvironment{codebox}{%
  \hllpdf@incodeboxtrue
  % Snapshot the OUTER language before asmblock selects English.
  \iflanguage{english}{\hllpdf@englishcodeboxtrue}{\hllpdf@englishcodeboxfalse}%
  \ifhllpdf@lstnumbermasked
    \let\thelstnumber\hllpdf@savedlstnumber
    \hllpdf@lstnumbermaskedfalse
  \fi
}
% listings' AddToHook is GLOBAL. Register once, and gate on local box state;
% adding an unconditional hook at every codebox entry leaks into later listings.
\lst@AddToHook{PreInit}{%
  \ifhllpdf@incodebox
    % Apply every listing override here, inside the LISTING's own group.
    % Setting margins at box entry leaks them into nested semantic callouts.
    \lstset{resetmargins=false,xrightmargin=26pt}%
    % Styles resolve before PreInit. In LTR the physical number gutter follows
    % xleftmargin; in Hebrew retain the measured inherited xrightmargin path.
    \ifhllpdf@englishcodebox
      \lstset{xleftmargin=26pt,xrightmargin=0pt}%
    \fi
    % ONE box, not two: suppress only the frame/background/skips of a listing
    % inside this codebox. The next unboxed listing keeps its own style.
    \lstset{frame=none,backgroundcolor=\color{cSlateBg},aboveskip=0pt,belowskip=0pt}%
  \fi
}
\makeatother
% </hllpdf:codebox-gutter>

% v3.37: line numbers inside a COLORED BOX are sliced by the box border —
% measured: digit span x=63.5-67.3 vs border line at x=65.2 (the border runs
% THROUGH the digits), and the inter-border gap (3.5pt) is narrower than a
% digit (3.8pt), so no margin knob can save them (under resetmargins,
% xleftmargin is inert in-box — measured: code at x=75.26 with 2.2em AND
% 3.0em). A 3-6 line boxed snippet does not need line numbers; blank the
% number TEXT inside every colored box (\thelstnumber is what gets printed;
% the style's numbers=left survives, so nothing else shifts — code position
% byte-frozen at x=75.26, verified). Unboxed listings keep their numbers.
% v3.78: the blanking is coupled to the FRAME, not to the box name —
% clean_theme.tex registers LATER hooks (registration order wins) that
% restore the number text inside the boxes it renders frameless.
% v4.5.0 RETRACTS the "no margin knob can save them" clause above. It holds only
% under resetmargins=true (the shipped default): there xleftmargin is inert
% in-box, as measured. With resetmargins=false, xrightmargin DOES carry the
% \llap'ed digit column clear of the rule — measured; the codebox \lstset that
% does it is the one above, in the codebox block. The loop below therefore
% covers SIX boxes, not "every colored box": codebox opts out and keeps its
% numbers.
% <hllpdf:semantic-listing-mask>
\makeatletter
\@for\HEB@lstbox:=defbox,thmbox,notebox,exbox,warnbox,keybox\do{%
  \expandafter\AtBeginEnvironment\expandafter{\HEB@lstbox}{%
    \ifhllpdf@lstnumbermasked\else
      \let\hllpdf@savedlstnumber\thelstnumber
    \fi
    \hllpdf@lstnumbermaskedtrue
    \def\thelstnumber{}%
    % The nearest semantic container owns its listing presentation. On exit,
    % TeX's group restores the enclosing codebox state and formatter.
    \hllpdf@incodeboxfalse}}
\makeatother
% </hllpdf:semantic-listing-mask>

% ============================================================================
% Unified TABLE presentation: subtle uniform slate frame + auto-numbered
% caption rendered INSIDE the frame. Use:
%   \begin{tablebox}\centering  <tabular>  \tabcap{<description>}\end{tablebox}
% Replaces the buggy \caption* (which emitted a stray "טבלה N : *" line) and
% provides an explicit framed option + numbered caption "טבלה N.M: תיאור".
% Optional intro text may be placed at the TOP, before the tabular, inside the box.
% ============================================================================
\definecolor{cTblLine}{HTML}{B9C4D2}  \definecolor{cTblBg}{HTML}{FAFBFC}
% v3.23: class-aware — under `article` there is no chapter counter (the promised
% report->article switch used to die with "No counter 'chapter' defined").
\makeatletter
\@ifundefined{c@chapter}{%
  % v3.31: in article the counter resets per SECTION, so an unprefixed number
  % repeated ("טבלה 1." in every section — measured on a two-section render).
  % Same conditional idea as the report branch: prefix the section number once
  % a real \section exists; a pre-section table stays "טבלה 1".
  \newcounter{tbl}[section]
  \renewcommand{\thetbl}{\ifnum\value{section}>0 \thesection.\fi\arabic{tbl}}%
}{%
  \newcounter{tbl}[chapter]
  % v3.30: a report-class doc that uses SECTIONS ONLY (chapter stays 0 —
  % the common short-doc pattern) used to print "טבלה 0.1"; the chapter
  % prefix now appears only once a real \chapter exists (render-verified
  % both ways: "טבלה 1." chapterless, "טבלה 1.1" with chapters).
  \renewcommand{\thetbl}{\ifnum\value{chapter}>0 \thechapter.\fi\arabic{tbl}}%
  % v3.64: the SAME leak in \thesection itself, left behind by v3.30. That fix
  % guarded the CAPTION counter and stopped there, so a chapterless report still
  % printed the heading "0.1 סעיף ראשון" and \ref to it printed "0.2" — measured:
  %   p1 y=65.2  '…0.1'   p1 y=144.9 '…0.2'   p1 y=183.4 '…0.2.1'
  % report.cls hardcodes \thesection as \thechapter.\@arabic\c@section, and the
  % leading 0 is the chapter counter that no \chapter ever incremented. Fixing
  % \thesection fixes the whole chain — \thesubsection and \thesubsubsection are
  % defined in terms of it. \theequation and \thetbl/\thefig were already guarded,
  % which is exactly why this survived: every NUMBER around the heading looked
  % right, so the heading's own 0 read as a style choice rather than a bug.
  % \appendix is safe: report.cls resets chapter to 0 and redefines \thechapter to
  % \@Alph\c@chapter, so an appendix \chapter makes the counter 1 and the guard
  % emits "A.1" unchanged. Render-verified in all four modes (chapterless,
  % chaptered, appendix+chapter, appendix+section-only).
  \renewcommand{\thesection}{\ifnum\value{chapter}>0 \thechapter.\fi\arabic{section}}%
}
\makeatother
% v3.37: everyline=true — if the box DOES split across pages, every segment
% gets a CLOSED frame (measured before: the split ran open-ended into the page
% edge with no bottom border, and the caption could strand alone at the next
% page top — "טבלה A.2" shipped that way). The split skips give the segments
% breathing room. In practice, with the caption GLUE below, every previously-
% split table in the corpus now moves WHOLE instead (the white-gap tradeoff,
% documented benign); everyline stays as the safety net for a genuinely
% page-tall table. NB: a table taller than ONE FULL PAGE can grind mdframed's
% split-trial loop (reproduced: minimal doc, minutes of 1pt-step Overfull
% probes) — break such a monster into two tables.
\mdfdefinestyle{hllpdftablegeometry}{leftline=true,rightline=true,topline=true,bottomline=true,
  linewidth=0.9pt,
  innertopmargin=9pt,innerbottommargin=8pt,skipabove=11pt,skipbelow=11pt,
  innerleftmargin=11pt,innerrightmargin=11pt,roundcorner=4pt}
\newcommand{\hllpdftableframestyle}{hllpdftablegeometry}
\newmdenv[frametitle={},style=hllpdftablegeometry,
  linecolor=cTblLine,backgroundcolor=cTblBg,
  % v4.2.0 (adopted): tablebox is ATOMIC. Short tableboxes near page bottoms
  % repeatedly entered mdframed's split trial, emitting benign Overfull-vbox
  % bursts and occasionally separating the intro prose from the tabular.
  % nobreak=true moves the complete short table to the next page instead. The
  % output-active Overfull-vbox gate still catches a box taller than a page;
  % everyline/split skips remain as defensive styling should that path ever
  % re-arm, but the canonical path never publishes a split tablebox.
  nobreak=true,
  everyline=true,splitbottomskip=6pt,splittopskip=12pt,
  % v3.79 (BLD-37 table side, correction): startinnercode/endinnercode run
  % INSIDE mdframed's body vbox (mdframed.sty:848 / :856+863) -- the second one
  % after \par\unskip, i.e. after line breaking, which is the only moment the
  % \tabcap anchor's true offset from the top of the box EXISTS. See
  % \lakintabboxclose below. These two keys are the ONLY tablebox change.
  startinnercode={\lakintabboxopen},endinnercode={\lakintabboxclose}]{tablebox}
% v3.37: \nopagebreak GLUE — \vsplit honors penalties, so the caption can no
% longer strand alone at the top of the next page (it takes its rows along, or
% the whole box moves). Render-verified on the real shipped orphan.
% v3.59: \NR@gettitle is what nameref injects into the stdlib title-bearing
% commands (\@sect, \@chapter, \caption) to fill \@currentlabelname. \figcap and
% \tabcap are INDEPENDENT macros — they never called it, so \@currentlabelname kept
% whatever the last heading left and \nameref{fig:x} silently printed the enclosing
% SECTION's title instead of the caption (proven in the .aux: field 3 of \newlabel
% was the heading text while field 1 and the anchor `fig.1.1` were correct — so this
% is a THIRD, distinct head next to BLD-13's counter fix and BIDI-38's autoref-name
% fix, neither of which writes that field). \@ifundefined keeps it inert if
% hyperref/nameref is ever dropped. Guarded helper because \figcap is defined
% after \makeatother.
% NB both helpers below need \makeatletter — this part of the preamble is OUTSIDE
% any letter-catcode region, so an unwrapped \@currentlabel parses as LaTeX's \@
% and dies with "You can't use \spacefactor in vertical mode" (measured).
\makeatletter
\newcommand{\lakinrtitle}[1]{\@ifundefined{NR@gettitle}{}{\NR@gettitle{#1}}}
% v3.59: \hitem{<letter>}{<key>} — a REFERENCEABLE Hebrew-lettered list item.
%   \begin{enumerate} \hitem{א}{it:a} הסעיף הראשון. \end{enumerate}  ->  \ref{it:a} = (א)
% Why this exists: automatic Hebrew list labels fail silently, so SKILL.md mandates
% an EXPLICIT \item[(\textbf{א})] — but LaTeX puts \refstepcounter INSIDE
% \if@noitemarg (latex.ltx:12964), and \item[...] never sets that flag. So an
% explicit item steps no counter and never touches \@currentlabel: a \label there
% binds to the AMBIENT one and \ref prints the CHAPTER/SECTION number where the
% item letter belongs (proven in the .aux: {chapter.1}). Same doctrine as \figcap —
% own the label state before \label. \phantomsection matters: without it the
% hyperlink still lands on the chapter head (measured); with it the anchor is
% section*.N and the bookmark outline stays clean.
\newcommand{\hitem}[2]{\item[(\textbf{#1})]%
  \def\@currentlabel{(#1)}\def\@currentlabelname{#1}\phantomsection\label{#2}}
% NB (v3.66): \hitem is \ref-only. \autoref{key} calls it a SECTION ("סעיף (א)")
% because \phantomsection brands the anchor section*.N, and \nameref{key} prints
% the bare letter. For \autoref-able items prefer \hebenumlabels below.
%
% v3.66: \hebenumlabels — OPT-IN automatic Hebrew enumerate letters, done right.
%   {\hebenumlabels\begin{enumerate}\item\label{q:a} ...\end{enumerate}}
%   -> labels (א) (ב) (ג); \ref{q:a} = א; anchor Item.N (real, per-item).
% This retires the belief that made \item[(\textbf{א})] + \hitem necessary.
% "Automatic Hebrew labels fail silently" was TRUE but unexplained, and the
% mechanism is now bisected (plain article vs +enumitem, measured): ENUMITEM
% PRE-EXPANDS the label expression in an \edef BEFORE \refstepcounter, with the
% counter still at 0. Tokens it shields by name (\arabic,\alph,\roman) survive;
% a plain \newcommand \ifcase wrapper evaluates at 0 -> the EMPTY zero-branch ->
% the label prints as "()" with the letter GONE (and \the\c@enumi prints a
% literal 0, \number\value dies fatally pointing at \item). The cure is one
% property, not a doctrine: the wrapper must be ROBUST (survives the \edef
% unexpanded), while \theenumi stays EXPANDABLE (so \refstepcounter writes the
% real letter to the .aux). Verified: labels (א)(ב)(ג), \ref -> א/ב/ג each item
% its own value, .aux = {א}{...}{Item.1} — versus the explicit-\item doctrine
% where BIDI-50/\hitem exist purely to patch the counter that never stepped.
\def\@hebalph#1{\ifcase#1\or א\or ב\or ג\or ד\or ה\or ו\or ז\or ח\or ט\or י\or
  יא\or יב\or יג\or יד\or טו\or טז\or יז\or יח\or יט\or כ\else\@ctrerr\fi}
\DeclareRobustCommand*\hebalph[1]{\expandafter\@hebalph\csname c@#1\endcsname}
\newcommand{\hebenumlabels}{%
  \renewcommand{\theenumi}{\@hebalph\c@enumi}%
  \renewcommand{\labelenumi}{(\textbf{\hebalph{enumi}})}}
\makeatother
% v3.63: \mbox the LABEL. \thetbl ("7.1") and a caption description that STARTS
% WITH A DIGIT ("32 הרגיסטרים…") are two European-Number runs separated only by
% neutrals (the label's own "." and \enspace). Under the RTL paragraph base the UBA
% resolves those neutrals to the RTL level, so both numbers sit in ONE reversed span
% and EXCHANGE VISUAL POSITIONS. Measured on the shipped PDF: a table whose real
% number is 7.1 printed "טבלה 32" — a different, EXISTING table number — while \ref
% correctly printed 7.1, so the caption silently contradicted its own cross-reference,
% and the description lost its leading count. The source contains no adjacency at all;
% it is manufactured by this macro, which is why every source linter was blind.
% \mbox makes the label one atomic box so nothing can reorder across it. Verified
% zero-regression: an ordinary Hebrew caption keeps identical x-geometry.
% One paragraph style covers label, number, separator and description.
% Explicit project/template overrides remain possible; the label hook is kept
% for compatibility, with an identity default. Counters/anchors stay unchanged.
\newcommand{\lakitcaptionstyle}{\normalfont\small\color{black}\hypersetup{allcolors=black}}
\newcommand{\lakitcaptionlabelstyle}[1]{#1}
\newcommand{\lakitcaptionpunctuation}{:}
% Native floats, captionof and longtable use the same presentation policy.
% subcaption is optional and inherits the same font rather than a second size.
\usepackage{caption}
\DeclareCaptionFont{lakit}{\lakitcaptionstyle}
\DeclareCaptionLabelFormat{lakit}{\lakitcaptionlabelstyle{#1~\en{#2}}}
\DeclareCaptionLabelSeparator{lakit}{\lakitcaptionpunctuation\enspace}
\captionsetup{font=lakit,labelfont={},textfont={},labelformat=lakit,
  labelsep=lakit,justification=centering}
\makeatletter
\AtBeginDocument{\@ifpackageloaded{subcaption}{%
  \captionsetup[sub]{font=lakit,labelfont={},textfont={},labelsep=lakit}}{}}
\makeatother
% v3.64: OPT-IN list-of-figures/tables entry via an optional short caption.
% \listoffigures and \listoftables are inherited from report.cls and LOOK available,
% but \figcap/\tabcap step the custom \c@fig/\c@tbl counters and never call
% \addcontentsline — so both printed their heading and NOTHING ELSE. Measured: a doc
% with one \figcap and one \tabcap rendered page 1 as 'רשימת האיורים\n1\n', and the
% .lof held only \@chapter's \addvspace, zero \contentsline entries. An empty list
% under a correct Hebrew heading reads as "this document has no figures".
%
% WHY OPTIONAL, AND NOT AUTOMATIC. Writing #1 unconditionally was implemented first
% and REJECTED: a caption containing \codeline compiles CLEANLY today (verified on
% the pre-v3.64 preamble: 0 errors, PDF produced) but becomes FATAL once it is also
% written to the .lof — \codeline is \hbox to\linewidth{...}, and a full-width box in
% a list entry dies with "\pdfextension endlink cannot be used in vertical mode", no
% output PDF. Gaining a feature nobody uses (0 of 13 collections call \listoffigures)
% at the price of breaking a construct that works is a bad trade. Stock \caption has
% exactly this limitation and exactly this escape hatch, so the convention is familiar.
% \en{}, \code{}, \h{} and inline math in the short text are all safe — verified, they
% survive the write and typeset correctly in the list (RTL geometry confirmed: number
% rightmost at x=509.9, title left of it, dots and folio at the far left).
% With no optional argument the macro is byte-identical to v3.63 — zero risk to every
% existing document, and \listoffigures lists exactly the captions that opted in.
% BLD-37, TABLE side (the \figcap v3.74 idea, adapted): clicking a \ref to a
% table used to land MID-table — hyperref's automatic anchor inside
% \refstepcounter was emitted on the still-open horizontal line whose TAIL is
% the tabular (an hbox whose baseline is the table's vertical CENTER), so the
% jump scrolled the top half of the table above the viewport. The v3.74 entry
% deferred this side claiming "tablebox splits — no single box to measure";
% that deferral dissolves: the measurement target is the TABULAR itself at
% \tabcap ENTRY, while it is still the tail of the open list — the same
% situation \lakin@grabfig handles at an island's end (and the fleet has 0
% split tableboxes). So \lakin@grabtab runs as the FIRST token, before any
% typesetting can push the box away, and the anchor is re-emitted raised above
% the measured tabular. \lakintabraise is a SEPARATE register from
% \lakinfigraise (defined with it, below the LTR-island block): a captionless
% island intentionally leaves the fig register set, and a table must never
% inherit that stale figure height.
%
% WHY the anchor stays AT THE ENTRY SPOT and is NOT deferred into the caption
% paragraph the way \figcap defers it: at \par TeX DROPS a final glue item,
% and the automatic anchor's smash box sitting BEHIND the interword glue kept
% that glue alive — so every centered tabular has always sat half a space off
% true center. Deferring the anchor removed that box and was measured as 982
% differing glyph records: EVERY tabular line shifted 1.30pt (2.61pt for
% pattern 7's two minipages) while captions stayed put. Emitting the raised
% anchor at the old spot — through hyperref's own \Hy@raisedlink, exactly the
% macro the automatic path uses — keeps the line's node structure, and so the
% A/B glyph diff at 0, BY CONSTRUCTION. \Hy@raisedlink reads the raise from
% \HyperRaiseLinkDefault (default \baselineskip = the old mid-table dest); the
% group-local \edef freezes measured-height+42pt into it, and \smash inside
% \Hy@raisedlink hides any raise from the line's extent. \HyCnt@ProvideTheH-
% Counter keeps the dest NAME byte-identical to the automatic one (it defines
% \theHtbl, which \hyper@makecurrent prefers — without it the first measured
% caption was named tbl.1 where the automatic path said tbl.0.1, measured).
% vmode entry, or anything but a box at the tail (no tabular to measure),
% falls back to the plain \refstepcounter — byte-identical to the old path.
\makeatletter
\newcommand{\tabcap}[2][]{\lakin@grabtab
  % v3.79: arm the box-close correction. Set unconditionally and in BOTH
  % branches below's scope, so the plain-\refstepcounter fallback is corrected
  % too; \lakintabboxopen clears it at every \begin{tablebox}, so a \tabcap
  % outside a tablebox can never leave it armed for a later box.
  \global\lakintabcapseentrue
  \ifdim\lakintabraise>0pt
    \H@refstepcounter{tbl}\HyCnt@ProvideTheHCounter{tbl}\hyper@makecurrent{tbl}%
    \begingroup
      % A frameless table targets its measured top. The frame's historical
      % cushion belongs only to tablebox and is corrected at box close.
      \edef\HyperRaiseLinkDefault{\the\lakintabraise}%
      \iflakintabframed
        \edef\HyperRaiseLinkDefault{\the\dimexpr\lakintabraise+42pt\relax}%
      \fi
      \Hy@raisedlink{\hyper@anchorstart{\@currentHref}\hyper@anchorend}%
    \endgroup
    \global\lakintabraise 0pt\relax
  \else
    \refstepcounter{tbl}%
  \fi
  \lakinrtitle{#2}\par\nopagebreak\addvspace{6pt}%
  \def\lakin@short{#1}\ifx\lakin@short\@empty\else
    \addcontentsline{lot}{table}{\protect\numberline{\thetbl}{\protect\ignorespaces #1}}\fi
  {\centering\lakitcaptionstyle\mbox{\lakitcaptionlabelstyle{\tablename~\en{\thetbl}\lakitcaptionpunctuation}}\enspace #2\par}}
\makeatother

% ============================================================================
% Unified FIGURE caption — the \tabcap twin for TikZ/plot figures (v3.34.0).
% A figure (tikzpic) sits inline in the prose; put \figcap BELOW it for a
% numbered, \ref-able caption "איור N.M: תיאור" in the same style as \tabcap:
%   \begin{tikzpic} ... \end{tikzpic}
%   \figcap{\label{fig:foo}description}     % \label FIRST, inside the argument
% WHY a macro, not a hand-written gray line: a manual
%   {\centering\small\textcolor{cGray}{איור: ...}} with a \label inside has NO
% \refstepcounter, so the \label binds to the AMBIENT \@currentlabel (the last
% equation) — \ref then prints an equation number AS the figure number
% (measured on a real page: "איור 48" where 48 is equation 48). \figcap steps a
% real \c@fig counter first, so \label is correct, numbering is consistent, and
% refs are stable. Class-aware like \thetbl (chapter- or section-prefixed).
% ============================================================================
\makeatletter
\@ifundefined{c@chapter}{%
  \newcounter{fig}[section]
  \renewcommand{\thefig}{\ifnum\value{section}>0 \thesection.\fi\arabic{fig}}%
}{%
  \newcounter{fig}[chapter]
  \renewcommand{\thefig}{\ifnum\value{chapter}>0 \thechapter.\fi\arabic{fig}}%
}
\makeatother
\makeatletter
\newcommand{\figcap}[2][]{%
  \ifdim\lakinfigraise>0pt
    % v3.74: step WITHOUT hyperref's automatic anchor, then emit the anchor
    % raised above the measured figure. \H@refstepcounter is hyperref's saved
    % pre-hook \refstepcounter (sets the counter AND \@currentlabel — \label
    % in the caption arg binds exactly as before); \hyper@makecurrent supplies
    % \@currentHref, so the dest NAME is identical to the old automatic one —
    % \ref/\autoref/bookmarks are unchanged, only the dest'S POSITION moves.
    \H@refstepcounter{fig}\hyper@makecurrent{fig}%
    % The anchor is emitted INSIDE the caption paragraph below, as a zero-size
    % raisebox — NOT here in vertical mode, where a \raisebox starts an empty
    % paragraph and adds a blank line (measured: every glyph after the caption
    % shifted +18.8pt). Inside an existing line a [0pt][0pt] box is inert.
    \edef\lakin@anch{\@currentHref}%
    % \edef FREEZES the measured height into the pending macro: the deferred
    % \raisebox is typeset inside the caption paragraph AFTER the register is
    % reset to the 0pt sentinel below — a plain \def read 0 at use time
    % (measured: LAKINRAISE=0.0, dest moved 14pt instead of above the figure).
    \edef\lakin@pend{\noexpand\raisebox{\the\dimexpr\lakinfigraise+42pt\relax}[0pt][0pt]{%
      \noexpand\hyper@anchorstart{\lakin@anch}\noexpand\hyper@anchorend}}%
    \global\lakinfigraise 0pt\relax
  \else
    \refstepcounter{fig}\let\lakin@pend\relax
  \fi
  \lakinrtitle{#2}\par\nopagebreak\addvspace{4pt}%
  \def\lakin@short{#1}\ifx\lakin@short\@empty\else
    \addcontentsline{lof}{figure}{\protect\numberline{\thefig}{\protect\ignorespaces #1}}\fi
  {\centering\lakin@pend\lakitcaptionstyle\mbox{\lakitcaptionlabelstyle{\figurename~\en{\thefig}\lakitcaptionpunctuation}}\enspace #2\par}}
\makeatother

% ============================================================================
% LTR islands.  EVERY tikzpicture / quantikz MUST live inside one, or the whole
% figure's coordinates and Latin labels flip. (Hebrew node text still uses \h{}.)
% ============================================================================
% v3.74 (BLD-37): the \figcap hyperlink DESTINATION used to land at the caption
% line — BELOW the figure — so clicking a figure reference scrolled the figure
% off-screen (measured fleet-wide: 97/97 fig dests in one collection sat 19-54pt
% under their figure's ink; a 347pt figure was fully hidden after the jump).
% hypcap-style fix in two parts: (1) at the island's END, \lakin@grabfig grabs
% the just-typeset centered line with \lastbox, measures it, and puts it back
% with \nointerlineskip — the interline glue above the line was already
% contributed and is NOT re-added, so the page geometry is unchanged to the
% glyph (verified); (2) \figcap then steps the counter WITHOUT hyperref's
% automatic anchor (\H@refstepcounter = the pre-hyperref \refstepcounter) and
% emits the anchor itself, raised by the measured height — the dest lands ABOVE
% the figure top. A captionless island leaves numbering untouched (nothing is
% stepped at the island; the grab flag is cleared at every island begin, so a
% stale height from a captionless figure cannot leak into a later caption).
\newsavebox{\lakinfigbox}
\newdimen\lakinfigraise  % real \dimen: tex.setdimen cannot target a \newlength SKIP register (measured lua error)
% BLD-37 table side: a SEPARATE register for \tabcap. The two must not share —
% a captionless island INTENTIONALLY leaves \lakinfigraise set (the grab runs
% at every island END; only the NEXT island's begin clears it), so a shared
% register would hand a \tabcap after a captionless figure that stale FIGURE
% height. \lakintabraise is written only by \lakin@grabtab at \tabcap entry
% and consumed/reset inside \tabcap itself, so it can never go stale.
\newdimen\lakintabraise
\makeatletter
% v3.74b: measure WITHOUT touching the list. The first draft grabbed the box
% with \lastbox and put it back — and \lastbox DISCARDS the node's baseline
% shift, so a vertically-anchored quantikz circuit rose 2.57pt and everything
% after it sank 5.32pt (measured on the A/B glyph diff). LuaTeX lets us READ
% the tail node's dimensions without removing anything: geometry is identical
% by construction. \lakinfigraise=0pt is the "no figure measured" sentinel.
\def\lakin@grabfig{%
  \ifhmode
    \directlua{
      local n = tex.nest[tex.nest.ptr].tail
      local GLUE, PEN = node.id('glue'), node.id('penalty')
      while n and (n.id == GLUE or n.id == PEN) do n = n.prev end
      local HL, VL = node.id('hlist'), node.id('vlist')
      if n and (n.id == HL or n.id == VL) then
        tex.setdimen('global', 'lakinfigraise',
                     (n.height or 0) + (n.depth or 0))
      else
        tex.setdimen('global', 'lakinfigraise', 0)
      end}%
  \else
    \global\lakinfigraise 0pt\relax
  \fi}
% BLD-37 table side: \lakin@grabfig's twin, NOT a shared helper — the Lua body
% differs in exactly one place, the skip set. Between \end{tabular} and \tabcap
% an otherlanguage wrapper (patterns 3/5 of example_tables) leaves a \write
% whatsit and a dir node ON TOP of the interword glue (measured node dump:
% "whatsit/write <- dir <- glue(2.62) <- hlist"), so the table reader must
% also step over whatsit/dir markers — all zero-size — to reach the tabular.
% \lakin@grabfig deliberately does NOT get the wider skip set: its glue/penalty
% body is the one the v3.74 A/B proof froze, and widening it would newly
% measure (and so newly raise) anchors in island shapes that today take the
% plain path — an unproven behavior change on the FIG side. Anything that is
% not a box behind those markers still hits the 0pt sentinel and \tabcap falls
% back to the plain \refstepcounter.
% SECOND deliberate difference: this stores the box's HEIGHT, not height+depth.
% \figcap's anchor is emitted in the caption paragraph BELOW the figure, so it
% must climb the full extent; \tabcap's anchor sits ON the tabular's own line,
% whose baseline is the vcentered tabular's MIDPOINT — height alone reaches
% the top edge, and adding the depth was measured to overshoot by half the
% table (dests 80-135pt above their frame, INSIDE the preceding table's box).
\def\lakin@grabtab{%
  \ifhmode
    \directlua{
      local n = tex.nest[tex.nest.ptr].tail
      local GLUE, PEN = node.id('glue'), node.id('penalty')
      local WHAT, DIR = node.id('whatsit'), node.id('dir')
      while n and (n.id == GLUE or n.id == PEN
                   or n.id == WHAT or n.id == DIR) do n = n.prev end
      local HL, VL = node.id('hlist'), node.id('vlist')
      if n and (n.id == HL or n.id == VL) then
        tex.setdimen('global', 'lakintabraise', n.height or 0)
      else
        tex.setdimen('global', 'lakintabraise', 0)
      end}%
  \else
    \global\lakintabraise 0pt\relax
  \fi}
\makeatother
% ============================================================================
% v3.79 (BLD-37 table side, CORRECTION): put the anchor at the TOP OF THE BOX,
% instead of guessing the distance to it.
% ============================================================================
% v3.77 raised the \tabcap anchor by (height of the box at the tail of the open
% horizontal list) + 42pt. That is a MEASUREMENT taken BEFORE line breaking, and
% it is short by every term it cannot see. Fleet-measured on the 12 built
% collections (299 tbl.* dests, ownership established through each table's own
% CAPTION, never through "which frame encloses the dest"): 3 dests landed INSIDE
% their own tablebox, below the frame top —
%   computer-architecture       tbl.14.1  -46.4pt   three stacked tabulars
%   computer-architecture       tbl.30.5  -13.5pt   tabular 2.85pt over \linewidth
%   advanced-electronic-devices tbl.B.1   -59.2pt   tabular 2.56pt over \linewidth
% Two missing terms, both reproduced in a minimal document (probe P2/P3):
%   (1) VERTICAL LIST ALREADY CONTRIBUTED inside the box — a second and third
%       tabular, or a heading above the table. The tail read sees only the LAST
%       box, so 77.7pt of stacked tabulars is invisible to the raise.
%   (2) WHICH LINE THE ANCHOR LANDS ON. When the tabular alone overfills
%       \linewidth (459.6121pt in a tablebox), the interword glue after
%       \end{tabular} becomes the chosen breakpoint and \Hy@raisedlink is
%       emitted on the NEXT line — one (depth + lineskip) lower. That term does
%       not EXIST at \tabcap time: line breaking has not happened yet. No
%       measurement taken there can ever recover it. (Proof: probe P3, the same
%       table 2.84912pt too wide, fails at -13.5pt; probe P4, the identical
%       table with \tabcolsep 2pt so it fits, is correct at +32.0pt.)
% So the raise is not computed at \tabcap at all. It is CORRECTED at
% \end{tablebox}, from inside mdframed's own body vbox (endinnercode runs after
% \par\unskip, i.e. after line breaking, while tex.nest[tex.nest.ptr] is still
% that vbox's vertical list). There the anchor's true offset from the top of the
% box content is not estimated but WALKED, and the \Hy@raisedlink shift is set
% so the dest lands 42pt above it — which is 42pt - innertopmargin(9pt) -
% linewidth(0.9pt) = 32.0pt above the FRAME top, the exact position the 295
% already-correct dests have today. That is why this is a correction and not a
% re-placement: the good sites do not move.
%
% WHY THE GLYPH DIFF IS 0 BY CONSTRUCTION — the only thing written is the
% `shift` field of the hlist that \Hy@raisedlink already built, and the guard
% below refuses to write it unless that box is 0x0x0. A zero-size box carries no
% dimension into its parent's packing, the packing already happened, and shift
% is not an input to it. Nothing else in the list is touched: no node is added,
% removed, resized or reordered. Measured, not assumed: assets/example_tables.tex
% and the whole computer-architecture collection A/B at 0 differing pymupdf
% rawdict glyph records.
%
% WHY NOT the alternatives (all prototyped and measured before this one):
%   * accumulate the vertical list at \tabcap time — fixes term (1), CANNOT fix
%     term (2); 2 of the 3 real failures survive. A fix that leaves the majority
%     of the defect standing is the false all-clear this repo exists to avoid.
%   * \zref-savepos two-pass — correct on all shapes, but it makes the dest
%     POSITION aux-derived (DEFECTS BLD-1/2), a single-pass compile silently
%     ships an unraised anchor, and the bare y-difference has no page guard.
%   * emit the anchor at \begin{tablebox} and let \tabcap adopt it — needs the
%     tbl counter stepped at box OPEN to keep the dest NAME byte-identical; a
%     captionless tablebox then either shifts every later table number or leaves
%     an orphan dest that the next table duplicates (a HARD build gate). This
%     correction touches neither the counter nor the name: \tabcap's naming path
%     is v3.77's, unchanged, so resolve_names() is byte-identical (verified).
%   * suppressing the line break (a penalty before the glue) — that MOVES the
%     tabular: the break currently discards the interword glue, so forbidding it
%     re-centers the two overfull tables. Reflow, rejected.
%
% SCOPE. The correction is armed only by \tabcap ITSELF (\lakintabcapseen),
% cleared at every box open, so a captionless tablebox is a no-op and a \tabcap
% hosted OUTSIDE a tablebox (the fleet has exactly one, computer-architecture
% tbl.14.2 in an exbox) keeps v3.77's measured raise. Both fallbacks are
% therefore never worse than today; and because hyperref's AUTOMATIC anchor uses
% the same \Hy@raisedlink shape, the plain-\refstepcounter path (vmode entry, or
% prose between \end{tabular} and \tabcap — probe P6, -51.5pt) is corrected too.
% KNOWN LIMIT — a tablebox that SPLITS across pages. The offset is measured from
% the top of the UNSPLIT body, so a dest in segment 2 gets raised past its own
% segment. MEASURED on a forced split (probe: 40 filler lines + a table): the
% dest went from y=585.0 on page 2 (inside segment 2, at the caption — the old
% BLD-37 symptom) to y=-113.4, i.e. ABOVE the page top. Viewers clamp that to the
% page top, which happens to show segment 2 whole, but an out-of-page coordinate
% is not something to ship on purpose. Hence the \textheight guard in
% lakintab.fix: a body TALLER THAN THE TEXT BLOCK cannot fit on any page, so it
% must split — the correction is skipped and v3.77's behaviour stands unchanged
% (verified: the probe returns to y=585.0). The guard is deliberately EXACT, not
% a page-fitting heuristic: it is sound (a taller-than-page body always splits)
% but PARTIAL (a shorter body that starts low on a page could still split). That
% residue is unreachable in practice — v3.37's everyline + the caption's
% \nopagebreak glue move such boxes WHOLE instead (probe N7: a 22-row table
% pushed down by 430pt moved whole, 0 segments), and the fleet has 0 split
% tableboxes, re-counted after this change. Reimplementing mdframed's own
% page-fitting arithmetic to close the residue would be a heuristic with its own
% failure mode, which is the trade this repo has repeatedly refused.
\newif\iflakintabcapseen
\newif\iflakintabframed
% The Lua lives in ONE \directlua at load time, and is written with NO
% backslash, %, # or ~ : \directlua expands its argument like \edef, so an
% active ~ would expand to \nobreakspace and a % would eat the rest of the line
% (the second is how the first draft of this block died). "not (a == b)" stands
% in for Lua's inequality operator for the same reason.
\directlua{
  lakintab = { }
  lakintab.HL = node.id("hlist")
  lakintab.VL = node.id("vlist")
  lakintab.RU = node.id("rule")
  lakintab.GL = node.id("glue")
  lakintab.KE = node.id("kern")
  lakintab.WH = node.id("whatsit")
  lakintab.TM = node.id("temp")
  lakintab.WD = node.subtype("pdf_dest")
  function lakintab.d(v)
    if v == nil then return 0 end
    if v < -1000000000 then return 0 end
    return v
  end
  function lakintab.scan(list, y, horiz, parent, target)
    local x = list
    while x do
      local i = x.id
      if i == lakintab.HL or i == lakintab.VL then
        if horiz then
          local b = y + x.shift
          if i == lakintab.HL then
            lakintab.scan(x.head, b, true, x, target)
          else
            lakintab.scan(x.head, b - lakintab.d(x.height), false, x, target)
          end
        else
          if i == lakintab.HL then
            lakintab.scan(x.head, y + lakintab.d(x.height), true, x, target)
          else
            lakintab.scan(x.head, y, false, x, target)
          end
          y = y + lakintab.d(x.height) + lakintab.d(x.depth)
        end
      elseif i == lakintab.WH then
        if x.subtype == lakintab.WD then
          local id = x.dest_id
          if type(id) == "string" and string.sub(id, 1, 4) == "tbl." then
            if horiz and parent and parent.width == 0
               and parent.height == 0 and parent.depth == 0 then
              parent.shift = parent.shift + (target - y)
            end
          end
        end
      elseif not horiz then
        if i == lakintab.RU then
          y = y + lakintab.d(x.height) + lakintab.d(x.depth)
        elseif i == lakintab.GL or i == lakintab.KE then
          y = y + lakintab.d(x.width)
        end
      end
      x = x.next
    end
  end
  function lakintab.fix(target, maxh)
    local nst = tex.nest[tex.nest.ptr]
    if nst == nil then return end
    local h = nst.head
    if h == nil then return end
    if h.id == lakintab.TM then h = h.next end
    local tot = 0
    local x = h
    while x do
      local i = x.id
      if i == lakintab.HL or i == lakintab.VL or i == lakintab.RU then
        tot = tot + lakintab.d(x.height) + lakintab.d(x.depth)
      elseif i == lakintab.GL or i == lakintab.KE then
        tot = tot + lakintab.d(x.width)
      end
      x = x.next
    end
    if tot > maxh then return end
    lakintab.scan(h, 0, false, nil, target)
  end
}
% @-FREE NAMES ON PURPOSE: these two are named in the `tablebox` key list at
% \newmdenv above, which is tokenised OUTSIDE any \makeatletter region — a
% \lakin@... there would tokenise as \lakin followed by the letters @tabopen.
\newcommand{\lakintabboxopen}{\lakintabframedtrue\global\lakintabcapseenfalse}
\newcommand{\lakintabboxclose}{%
  \iflakintabcapseen
    \directlua{lakintab.fix(-\number\dimexpr\lakintabcushion\relax,
                            \number\dimexpr\textheight\relax)}%
    \global\lakintabcapseenfalse
  \fi\lakintabframedfalse}
% The cushion above the box's CONTENT top. 42pt puts the dest 32.0pt above the
% frame top (42pt - innertopmargin 9pt - linewidth 0.9pt), i.e. exactly where
% the 295 correct dests already are — the correction moves only the broken ones.
\newcommand{\lakintabcushion}{42pt}
\newenvironment{qcirc}{\par\medskip\global\lakinfigraise 0pt\relax\begin{otherlanguage}{english}\begin{center}}
  {\csname lakin@grabfig\endcsname\end{center}\end{otherlanguage}\par\medskip}
\newenvironment{tikzpic}{\par\medskip\global\lakinfigraise 0pt\relax\begin{otherlanguage}{english}\begin{center}}
  {\csname lakin@grabfig\endcsname\end{center}\end{otherlanguage}\par\medskip}

% ============================================================================
% Code listings.  ★★ v3.4 THE BOX-OVERFLOW FIX:
%   resetmargins=true   → the listing IGNORES the enclosing environment's
%                         indentation and measures from the PAGE. v3.4 stated
%                         this backwards ("measured from the ENCLOSING box"),
%                         and v3.37 then blanked in-box line numbers because
%                         the page-anchored listing put its hanging digits on
%                         top of the box rule. Measured v4.5.0: with the key
%                         FALSE the listing follows the box and the code moves
%                         from x=75.26 to x=97.74, clearing the rule at 72.07.
%                         The key now lives in one \lstset below instead of in
%                         each style, so a box can override it — codebox does.
%   framexleftmargin=0  → the frame (left bar) does not extend left of the code.
%   xleftmargin keeps a small indent so the line numbers have room INSIDE.
% asmblock wraps the listing in an LTR island (code is LTR).
% ============================================================================
\lstdefinestyle{asm}{
  basicstyle=\ttfamily\small,
  backgroundcolor=\color{cBlueBg!55},
  frame=leftline,framerule=2.5pt,rulecolor=\color{cBlue},
  numbers=left,numberstyle=\tiny\color{cGray},numbersep=8pt,
  xleftmargin=2.2em,framexleftmargin=0pt,framexrightmargin=0pt,
  breaklines=true,columns=fullflexible,keepspaces=true,
  showstringspaces=false,aboveskip=8pt,belowskip=8pt,
  % v3.36: without a comment DELIMITER the commentstyle below never fires —
  % listings has no `language` here, so it does not know '#' opens a comment,
  % and every asm listing rendered flat black (measured: a whole-page pixel
  % sweep found ZERO cGreen glyphs across 135 shipped listings). morecomment
  % gives it the delimiter; color-only change, glyph widths untouched.
  morecomment={[l]{\#}},
  commentstyle=\color{cGreen},
}
\lstdefinestyle{cstyle}{
  language=C,basicstyle=\ttfamily\small,
  backgroundcolor=\color{cGreenBg!55},
  frame=leftline,framerule=2.5pt,rulecolor=\color{cGreen},
  numbers=left,numberstyle=\tiny\color{cGray},numbersep=8pt,
  xleftmargin=2.2em,framexleftmargin=0pt,framexrightmargin=0pt,
  breaklines=true,columns=fullflexible,keepspaces=true,
  showstringspaces=false,keywordstyle=\color{cBlue}\bfseries,
  commentstyle=\color{cGray}\itshape,aboveskip=8pt,belowskip=8pt,
}
% v3.36: pystyle/htmlstyle — colored twins of cstyle for the Python/HTML
% listings that used to repeat a ~200-char options blob with NO color
% (keywords rendered bold-black, comments italic-black — measured; C alone
% was colored). Geometry is copied VERBATIM from the shipped blob (frame,
% margins, skips) so swapping a call site to style=pystyle moves NOTHING —
% only keywordstyle/commentstyle/stringstyle are added (bold stays bold,
% italic stays italic, mono metrics identical => zero reflow, proven by
% per-page word-bbox diff). breaklines/breakatwhitespace stay per-site.
% v3.78: numbers=left joins both styles, with the exact numberstyle/numbersep
% keys asm/cstyle carry — fleet-wide consistency (author decision: all four
% styles numbered; "asm/cstyle numbered, py/html not" was an accident of the
% v3.36 geometry freeze, not a design). The numbers hang LEFT of the code box
% (listings sets them with \llap, so no overfull hbox can arise and the code
% x is untouched — measured: code frozen, digits at code_x-11.8pt); inside a
% colored box they are blanked like everyone else's (v3.37 hook), and the
% clean variant restores them in its frameless boxes (v3.78, above).
\lstdefinestyle{pystyle}{
  language=Python,basicstyle=\ttfamily\small,
  columns=fullflexible,keepspaces=true,
  frame=single,rulecolor=\color{cTblLine},
  numbers=left,numberstyle=\tiny\color{cGray},numbersep=8pt,
  xleftmargin=0.5em,framexleftmargin=0pt,framexrightmargin=0pt,
  showstringspaces=false,aboveskip=8pt,belowskip=8pt,
  keywordstyle=\color{cBlue}\bfseries,
  commentstyle=\color{cGray}\itshape,
  stringstyle=\color{cGreen},
}
% v4.5.0: hoisted out of the four styles so an enclosing box can override it.
% Every listing still resolves to true unless something later says otherwise,
% so shipped output is unchanged (verified byte-for-byte on all eight PDFs).
\lstset{resetmargins=true}
\lstdefinestyle{htmlstyle}{
  language=HTML,basicstyle=\ttfamily\small,
  columns=fullflexible,keepspaces=true,
  frame=single,rulecolor=\color{cTblLine},
  numbers=left,numberstyle=\tiny\color{cGray},numbersep=8pt,
  xleftmargin=0.5em,framexleftmargin=0pt,framexrightmargin=0pt,
  showstringspaces=false,aboveskip=8pt,belowskip=8pt,
  keywordstyle=\color{cBlue}\bfseries,
  commentstyle=\color{cGray}\itshape,
  stringstyle=\color{cGreen},
}
\newenvironment{asmblock}{\par\begin{otherlanguage}{english}}{\end{otherlanguage}\par}

% <hllpdf:listing-smart-quotes>
% v4.21.0: four Unicode punctuation characters must enter listings' token
% buffer. Unregistered characters bypass it: a closing quote was printed
% BEFORE the pending word, and an opening quote BEFORE a pending brace.
% Extend the standard character table, then activate only these codepoints
% in the listing's local character-table scope. Do this BEFORE the package's
% delimiter/alsoother/literate hooks so explicit user mappings still win.
% This is not universal Unicode support, nor a conversion to ASCII quotes.
% See listings developer guide 15.3.1 and references/listing-smart-quotes.md.
\makeatletter
\lst@CCPut \lst@ProcessOther {"2018}{"2019}{"201C}{"201D}\z@
\lst@AddToHookAtTop{SelectCharTable}{%
  \lst@ifactivechars
    \catcode"2018=\active \catcode"2019=\active
    \catcode"201C=\active \catcode"201D=\active
  \fi}
\makeatother
% </hllpdf:listing-smart-quotes>

% <hllpdf:listing-technical-symbols>
% v4.30: bounded extension for common technical punctuation in LTR listings.
% Without character-table registration these glyphs bypass the token buffer
% and can move before the preceding word even inside quoted strings.
% Keep the original four-quote block separate for its standing mutation test.
% This is source-faithful printing, not syntax validation or general Unicode.
\makeatletter
\lst@CCPut \lst@ProcessOther
  {"2013}{"2014}{"2026}{"2190}{"2192}{"2194}{"21D2}
  {"2212}{"221E}{"2260}{"2264}{"2265}\z@
\lst@AddToHookAtTop{SelectCharTable}{%
  \lst@ifactivechars
    \catcode"2013=\active \catcode"2014=\active \catcode"2026=\active
    \catcode"2190=\active \catcode"2192=\active \catcode"2194=\active
    \catcode"21D2=\active \catcode"2212=\active \catcode"221E=\active
    \catcode"2260=\active \catcode"2264=\active \catcode"2265=\active
  \fi}
\makeatother
% </hllpdf:listing-technical-symbols>


\usepackage[hidelinks,unicode]{hyperref}   % LAST, after colors. unicode = Hebrew bookmarks.
% v3.62: THE ISOLATION MACROS CORRUPTED EVERY PDF BOOKMARK THEY APPEARED IN.
% \en is defined (line ~155) as
%     \foreignlanguage{english}{\language=\csname l@nohyphenation\endcsname #1}
% and in a heading that body is written verbatim to the .toc and handed to
% hyperref's \pdfstringdef. hyperref strips the two CONTROL SEQUENCES (\language,
% \l@nohyphenation) but the plain `=` CHARACTER token between them is not a control
% sequence, so it survives into the PDF string: \chapter{שער \en{CNOT} והצמדה}
% produced the bookmark "שער =CNOT והצמדה". The printed TOC and the heading on the
% page are CORRECT — only the navigation pane and the viewer window title are wrong,
% which is why it went unnoticed. Measured on the shipped fleet: 70 corrupted
% bookmarks across 5 PDFs, from 177 unprotected \en-in-heading sites. Same corruption
% reached /Title and /Subject via the \hypersetup{pdftitle=...} the docs recommend.
% This hook tells hyperref to use the plain argument for all three isolation macros
% when building a PDF string. It does NOT touch the .toc (the printed TOC is
% byte-unchanged) and it does NOT mask the math signal — a heading with $...$ still
% logs "removing math shift", so BIDI-15 detection is intact.
% BONUS, measured: it also makes \code{} in an UNSTARRED heading COMPILE. BLD-10
% recorded that as needing \texorpdfstring or a starred heading; the real root cause
% is hyperref's \pdfstringdef scan of \code's \ifmmode...\else...\fi body — which is
% exactly why the identically-bodied \texttt survived (hyperref pre-neutralizes
% \texttt in PDF strings) and why \DeclareRobustCommand could never help.
% v3.64: \def\\{ } added. A manual line break in a heading WELDS the two words
% together in the bookmark: \chapter{שורה ראשונה\\שורה שנייה} produced the outline
% entry 'שורה ראשונהשורה שנייה' — "ראשונהשורה" is not a word. hyperref drops the
% \\ and puts nothing in its place, so the break that looks like a space on the
% page is a zero-width join in the navigation pane. The spaced form
% 'מילה אחת \\ מילה שנייה' escapes with a harmless DOUBLE space ('אחת  מילה'),
% which is why this never got noticed. Both are measured on a real outline.
% 0 live instances in the corpus — but SKILL.md and content-style.md both RECOMMEND
% a manual heading break, so this is a primed trap rather than a dormant one, and a
% one-token preamble fix retires it for good instead of asking authors to remember.
% Mapping \\ to a space fixes the WELD, which is the actual defect (a nonexistent
% word in the navigation pane). It does not normalise whitespace: measured after the
% fix, the welded form becomes 'שורה ראשונה שורה שנייה' (correct) and the already-
% spaced form goes from two spaces to three. That is cosmetic — viewers render it as
% a wider gap, and no word is corrupted. A per-heading \texorpdfstring is still the
% way to get an exactly-tuned bookmark string.
\pdfstringdefDisableCommands{\def\en#1{#1}\def\h#1{#1}\def\doclang#1{#1}\def\code#1{#1}\def\codeline#1{#1}\def\\{ }}
% v3.70: RESTORE \url LINE-BREAKING. url.sty freezes its break penalties at LOAD
% time (\mathchardef\UrlBreakPenalty=\binoppenalty), and v3.45's global
% \binoppenalty=10000 (inline-math atomicity) therefore priced EVERY breakpoint
% inside \Url at 10000 = never break. Measured: a long \url in Hebrew prose is
% one unbroken line, ~100pt overfull, invading the margin band the ink gates
% deliberately tolerate. With the two lines below: 0 overfull, the URL breaks
% after '/', the RTL split is CORRECT (fragment order forward), and each line
% gets its own link annotation exactly covering its fragment (measured).
% Confined to \Url's group — inline-math atomicity (TYP-7) is untouched.
\mathchardef\UrlBreakPenalty=200
\mathchardef\UrlBigBreakPenalty=100
\ifx\lakitmainlanguage\hllpdfenglish
  \hypersetup{pdflang=en}
\else
  \hypersetup{pdflang=he}
\fi  % /Lang is the document base locale (not per-span accessibility tagging).
% Set a per-document \hypersetup{pdftitle={...}} too — see SKILL.md.

% v3.52 (round-2 audit): two RTL bibliography/cross-ref defects, both invisible to
% every gate and both affecting the FIRST author to use the feature.
% (a) \@biblabel's default [#1] brackets are bidi-neutral and mirror to ]1[ under
% the RTL bibliography base — every \thebibliography reference list shows ]1[.
% LTR-isolating the label restores [1]. (The render gate's mirror rule fingerprints
% round parens )1(, so square brackets slipped through.)
\makeatletter
\renewcommand*{\@biblabel}[1]{\foreignlanguage{english}{[#1]}}
% v3.54: the SAME neutral-bracket trap on the INLINE side — a multi-key \cite{a,b}
% under the RTL base lays its European-number list right-to-left, so [1, 2] renders
% "[2,1]" and [1, 2, 3] renders "[3,2,1]" (verified by glyph geometry — the reader
% sees the sources in REVERSE order; a note form \cite[p.5]{a} also reorders to
% "[p. 5,1]"). LTR-isolate the whole bracket group (number list AND optional note),
% exactly like \@biblabel above: [1, 2] / [1, 2, 3] / [1, p. 5] all render correct.
% Placed AFTER hyperref so it wins; hyperref hyperlinks each number inside #1 (via
% \@citex), which this wrapper leaves intact. Mirrors the LaTeX kernel \@cite shape.
% v4.9.1: directional isolation must NOT switch the language here. Babel's
% safe-active citation wrapper passes its key list through \bbl@tempa;
% \foreignlanguage reuses that scratch macro before the deferred citation
% loop expands it. The old line therefore cited the literal key "english".
% If such a key existed, even the undefined-reference gate stayed green.
% \babelsublr isolates the bracket group without selecting a locale or
% changing that key list. Preserve native links, key order and optional notes.
\def\@cite#1#2{\babelsublr{[{#1\if@tempswa , #2\fi}]}}
\makeatother

% v4.11.0: a reference inside align's \text inherited HEBdisptag=true,
% so its enclosing hbox mirrored the brackets despite a clean compile.
% Isolate the complete printed reference without selecting a Babel language.
% Also neutralize the display-only normalfont direction reset locally. A
% direction-only patch inverted 12 to 21 inside equation in the measured probe.
% Plain maketag is local to the reference: mathtools' display-side wrapper
% otherwise globally consumes manual_tag while formatting a nested reference.
% Preserve each original command (including the selective auxiliary write),
% current tag forms, native links, and showonlyrefs on/off transitions.
\makeatletter
\NewCommandCopy{\HEB@nativeeqref}{\eqref}
\newcommand{\HEB@eqrefisolate}[1]{%
  \babelsublr{\HEBdisptagfalse\let\bbl@eqnodir\relax\let\maketag@@@\ORIGmaketag@@@#1}}
\protected\def\eqref#1{\HEB@eqrefisolate{\HEB@nativeeqref{#1}}}
\ExplSyntaxOn
\NewCommandCopy{\HEB@nativeMTeqref}{\MT_extended_eqref:n}
\cs_set_protected:Npn \MT_extended_eqref:n #1
  {\HEB@eqrefisolate{\HEB@nativeMTeqref{#1}}}
\ExplSyntaxOff
\makeatother

% v4.11.0: custom caption counters are not part of hyperref's English
% locale, and Hebrew has no upstream hyperref locale in this toolchain.
% Restore the appropriate names at every Babel language switch, including
% return from a nested inline island. Do not localize the destination keys.
\makeatletter
\newcommand{\HEB@refnames}{%
\def\chapterautorefname{פרק}%
\def\sectionautorefname{סעיף}%
\def\subsectionautorefname{תת־סעיף}%
\def\subsubsectionautorefname{תת־תת־סעיף}%
\def\paragraphautorefname{פסקה}%
% v4.31: hyperref uses a distinct counter name at the deepest standard level.
\def\subparagraphautorefname{תת־פסקה}%
\def\equationautorefname{משוואה}%
\def\figureautorefname{איור}%
\def\tableautorefname{טבלה}%
\def\partautorefname{חלק}%
\def\appendixautorefname{נספח}%
\def\footnoteautorefname{הערת שוליים}%
\def\itemautorefname{סעיף}%
\def\Itemautorefname{סעיף}%
\def\figautorefname{איור}%
\def\tblautorefname{טבלה}%
\def\pageautorefname{עמוד}% 
}
\addto\extrashebrew{\HEB@refnames}
\newcommand{\ENG@refnames}{%
  \def\subparagraphautorefname{subparagraph}%
  \def\chapterautorefname{chapter}%
  \def\sectionautorefname{section}%
  \def\subsectionautorefname{subsection}%
  \def\subsubsectionautorefname{subsubsection}%
  \def\paragraphautorefname{paragraph}%
  \def\equationautorefname{Equation}%
  \def\figureautorefname{Figure}%
  \def\tableautorefname{Table}%
  \def\partautorefname{Part}%
  \def\appendixautorefname{Appendix}%
  \def\footnoteautorefname{footnote}%
  \def\itemautorefname{item}%
  \def\figautorefname{Figure}%
  \def\tblautorefname{Table}%
  \def\Itemautorefname{item}%
  \def\pageautorefname{page}%
}
\addto\extrasenglish{\ENG@refnames}
\ifx\lakitmainlanguage\hllpdfenglish\ENG@refnames\else\HEB@refnames\fi
\makeatother

% ============================================================================
% Optional running head (v3.39): \lakithead{<book title>} — OFF by default.
% A document opts in with ONE line after \input{preamble.tex}:
%   \lakithead{מבנה המחשב}
% Every content page then carries the series' orientation line: chapter name
% at the reading edge, book title opposite, thin rule. (fancyhdr's [L]/[R]
% are PHYSICAL; under the RTL page direction [L] lands at the RIGHT/reading
% edge — measured, consistent across the corpus.) Packaged like the topic
% rail: always loaded, inert without the call. This centralizes the 8-line
% block that 8 lakits used to copy by hand — the same duplication vector
% that once let a lakit preamble go ~100 lines stale. Class-aware: report
% fills the mark from \chapter, article from \section. NB: a titlepage env
% self-empties; front matter that sets \pagestyle{empty} must restore with
% \pagestyle{fancy} afterwards (NOT plain — that discards the head).
% ============================================================================
\usepackage{fancyhdr}
% v3.61: \chapter* left a STALE running head — pages of an UNNUMBERED chapter
% carried the PREVIOUS chapter's title. report.cls routes \chapter* through
% \@schapter, which never calls \chaptermark, so \leftmark keeps whatever the last
% numbered chapter left. Measured: with \chapter{פרק ממוספר} then
% \chapter*{מבוא לא ממוספר}, page 4 (inside the starred chapter) printed
% "פרק ממוספר". Invisible on the chapter's FIRST page (\thispagestyle{plain}), so it
% only shows from the second page on — the pages an author is least likely to check.
% An unnumbered front-matter chapter (\chapter*{מבוא}, \chapter*{הקדמה}) is everyday
% Hebrew-academic structure, and the corpus already carries HAND-WRITTEN \markboth
% lines right after \chapter* — the exact per-document duplication \lakithead exists
% to retire. Patched INSIDE \lakithead so the macro stays inert when unused; a
% document's own later \markboth still wins, so those hand-patches do not regress.
\makeatletter
% <hllpdf:starred-heading-scope>
\newif\iflakin@starmarkpatched
\newcommand{\lakistarmark}{%
  \iflakin@starmarkpatched\else
  % Dispatch on the actual class counter, not hyperref's shared internals.
  \ifdefined\c@chapter
    \let\HEBorigschapter\@schapter
    % Record only after the native heading has reached its page. A mark
    % before the native routine can advertise the next heading too early.
    \def\@schapter##1{\HEBorigschapter{##1}\chaptermark{##1}}%
  \else
    % @ssect receives layout arguments, NOT the section name or level.
    % Capture the name in @startsection without introducing a TeX group:
    % grouping sectioning would break the native after-heading state.
    \def\HEB@sectionkind{}%
    \def\HEB@mainsection{section}%
    \let\HEBorigstartsection\@startsection
    \def\@startsection##1##2##3##4##5##6{%
      \def\HEB@sectionkind{##1}%
      \HEBorigstartsection{##1}{##2}{##3}{##4}{##5}{##6}}%
    \let\HEBorigssect\@ssect
    \def\@ssect##1##2##3##4##5{%
      \HEBorigssect{##1}{##2}{##3}{##4}{##5}%
      % Lower starred headings must not replace the article's section mark.
      \ifx\HEB@sectionkind\HEB@mainsection
        \sectionmark{##5}%
      \fi
      \ignorespaces}%
  \fi
  \global\lakin@starmarkpatchedtrue
  \fi}
% </hllpdf:starred-heading-scope>
% <hllpdf:running-language>
% Choose the wrapper when the mark is written, not when a later page ships.
% English punctuation must remain LTR even after returning to Hebrew prose.
\newcommand{\HEB@runningmark}[1]{%
  \iflanguage{english}{\markboth{\en{#1}}{}}{\markboth{\h{#1}}{}}}
% </hllpdf:running-language>
\makeatother
% v4.2.1 (NAV): every collection already names itself in \lakithead — feed the
% PDF /Title from the same declaration, so the navigation gate's
% PDF_REQUIRE_TITLE has a universal source and no fleet member ships
% title-less by missing a second declaration site. Measured before this: 14
% of 16 collections shipped /Title-empty while every one carried a perfectly
% good \lakithead string (plain Hebrew in all 16 — no math, no formatting, so
% the PDF-string path is clean). An explicit title always wins: set it with
% \lakitpdftitle{...} anywhere before \lakithead, or with a plain
% \hypersetup{pdftitle=...} — the fallback detects that via hyperref's
% \@pdftitle, probe-verified \@empty-equal until the first explicit set.
\makeatletter
\newif\iflakit@pdftitleset
\newcommand{\lakitpdftitle}[1]{%
  \hypersetup{pdftitle={#1}}\global\lakit@pdftitlesettrue}
\newcommand{\lakithead}[1]{%
  \iflakit@pdftitleset\else
    \ifx\@pdftitle\@empty
      \hypersetup{pdftitle={#1}}\global\lakit@pdftitlesettrue
    \fi
  \fi
  \pagestyle{fancy}\fancyhf{}%
  \lakistarmark

  \fancyhead[L]{\small\textcolor{cGray}{\leftmark}}%
  % <hllpdf:fixed-head-language>
  % The fixed document title has its OWN language snapshot, independent of
  % the current section/chapter mark. Choose the branch NOW, not at shipout:
  % the page may be Hebrew by the time an English title is finally printed.
  \iflanguage{english}{%
    \fancyhead[R]{\small\textcolor{cGray}{\en{#1}}}%
  }{%
    \fancyhead[R]{\small\textcolor{cGray}{\h{#1}}}%
  }%
  % </hllpdf:fixed-head-language>
  \fancyfoot[C]{\thepage}%
  \renewcommand{\headrulewidth}{0.4pt}%
  \ifdefined\chaptermark
    \renewcommand{\chaptermark}[1]{\HEB@runningmark{##1}}%
  \else
    \renewcommand{\sectionmark}[1]{\HEB@runningmark{##1}}%
  \fi}
\makeatother

% v3.39.1: path-independent byte-reproducibility. LuaTeX derives the trailer
% /ID from the OUTPUT PATH, so two clones of the same source at different
% directories produced different sha256 (measured: the ONLY delta was the
% 62-byte /ID line). Bit 512 suppresses exactly that line — measured: Creator/
% Producer/CreationDate/ModDate and all content untouched. /ID is optional
% per ISO 32000 (required only for ENCRYPTED PDFs and PDF/A — neither is a
% target here). If a future workflow needs an ID (PDF/A, Acrobat revision
% matching), delete this line and it returns.
\pdfvariable suppressoptionalinfo \numexpr 512\relax

% v3.75: BOLD MATH IN BOLD HEADINGS. Heading text is \bfseries but math inside
% it stayed regular-weight CM — a hairline $V_T$ inside a heavy 24.8pt chapter
% title (measured: CMMI12 fl=6 next to FrankRuhlLibre-Bold fl=20; 155 live
% heading-with-math sites fleet-wide). \boldmath rides along wherever the class
% sets \bfseries for a heading; A/B measured: $V_T$ -> CMMIB10 fl=22, 0 missing
% characters (the bold CM faces ship with this TeX Live). Guarded per-macro so
% the article class (no \@makechapterhead) and future class changes stay safe.
% v4.0.1 (adoption fix): the SEARCH patterns above carry no parameter token.
% v4.0.0 made them specific by including the class's own #1/#2, and measured on
% TeX Live 2023 that never matches: a \patchcmd search pattern containing a
% parameter token cannot match while the call sits inside \@ifundefined{...}{}{...}
% (probed both #1 and ##1 -> FAIL; the same pattern at top level succeeds). Every
% one of these patches therefore failed, and because v4.0.0 also replaced the
% empty failure branch with \PackageError, the failure became FATAL: all four
% bundled example assets stopped building ("Cannot patch chapter-title math
% weight -> Fatal error, no output PDF"). The \PackageError branches are KEPT —
% a silent no-op here is what let v3.75's fix rot unnoticed — only the parameter
% tokens are dropped, so the patches match and a real class change still errors.
\makeatletter
\@ifundefined{@makechapterhead}{}{%
  \patchcmd{\@makechapterhead}{\huge\bfseries \@chapapp}
    {\huge\bfseries\boldmath \@chapapp}{}
    {\PackageError{hebrew-lualatex-pdf}{Cannot patch chapter-number math weight}{Unsupported document class}}%
  \patchcmd{\@makechapterhead}{\Huge \bfseries}
    {\Huge \bfseries\boldmath}{}
    {\PackageError{hebrew-lualatex-pdf}{Cannot patch chapter-title math weight}{Unsupported document class}}%
  \patchcmd{\@makeschapterhead}{\Huge \bfseries}
    {\Huge \bfseries\boldmath}{}
    {\PackageError{hebrew-lualatex-pdf}{Cannot patch starred chapter math weight}{Unsupported document class}}}
% v4.0.1: \part sits under a THREE-DEEP wrapper chain, so the class-based anchor
% is not in \@part at all. Measured with \show: hyperref rewrites \@part to
% `...\phantomsection\fi \H@old@part`, and nameref rewrites THAT to
% `\NR@gettitle{#1}\NR@part[{#1}]{#2}` — the class body ends up in \NR@part.
% v4.0.0 patched \@part and made failure fatal, which is why all seven bundled
% examples stopped building. \part is NOT a dead heading level here: 188 live uses
% across the fleet, so the bold-math fix is worth getting right rather than
% dropping. We walk the chain innermost-first and stop at the first hit. Failure is
% a WARNING, not an error: the chain is a property of hyperref/nameref versions, so
% a future package release must degrade to regular-weight part math — a cosmetic
% loss — instead of refusing to build. Chapter and section keep \PackageError,
% because those targets are in the CLASS and were verified to patch here.
\newif\iflakin@partnum \newif\iflakin@parttitle \newif\iflakin@sparttitle
\def\lakin@partnum#1#2{\iflakin@partnum\else
  \patchcmd{#1}{#2\bfseries \partname}{#2\bfseries\boldmath \partname}
    {\lakin@partnumtrue}{}\fi}
\def\lakin@parttitle#1#2{\iflakin@parttitle\else
  \patchcmd{#1}{#2 \bfseries}{#2 \bfseries\boldmath}{\lakin@parttitletrue}{}\fi}
\def\lakin@sparttitle#1#2{\iflakin@sparttitle\else
  \patchcmd{#1}{#2 \bfseries}{#2 \bfseries\boldmath}{\lakin@sparttitletrue}{}\fi}
\@ifundefined{c@chapter}{\def\lakin@pnum{\Large}\def\lakin@ptit{\huge}}%
                        {\def\lakin@pnum{\huge}\def\lakin@ptit{\Huge}}
\@ifundefined{NR@part}{}{\expandafter\lakin@partnum\expandafter\NR@part\expandafter{\lakin@pnum}%
                         \expandafter\lakin@parttitle\expandafter\NR@part\expandafter{\lakin@ptit}}
\@ifundefined{H@old@part}{}{\expandafter\lakin@partnum\expandafter\H@old@part\expandafter{\lakin@pnum}%
                            \expandafter\lakin@parttitle\expandafter\H@old@part\expandafter{\lakin@ptit}}
\@ifundefined{@part}{}{\expandafter\lakin@partnum\expandafter\@part\expandafter{\lakin@pnum}%
                       \expandafter\lakin@parttitle\expandafter\@part\expandafter{\lakin@ptit}}
\@ifundefined{NR@spart}{}{\expandafter\lakin@sparttitle\expandafter\NR@spart\expandafter{\lakin@ptit}}
\@ifundefined{H@old@spart}{}{\expandafter\lakin@sparttitle\expandafter\H@old@spart\expandafter{\lakin@ptit}}
\@ifundefined{@spart}{}{\expandafter\lakin@sparttitle\expandafter\@spart\expandafter{\lakin@ptit}}
\iflakin@partnum\else\PackageWarning{hebrew-lualatex-pdf}{Part-number math stays
  regular weight: no known \noexpand\part wrapper matched (hyperref/nameref
  version change). Cosmetic only}\fi
\iflakin@parttitle\else\PackageWarning{hebrew-lualatex-pdf}{Part-title math stays
  regular weight: no known \noexpand\part wrapper matched. Cosmetic only}\fi
% v4.12.0: THE PART NUMBER IN THE TOC WAS NOT DIRECTIONALLY SEPARATED FROM THE
% TITLE. report.cls/article.cls compose the part entry as ONE text stream --
% \addcontentsline{toc}{part}{\thepart\hspace{1em}#1} -- while EVERY other level
% goes through \numberline, which puts the number in an \hbox. A box is a bidi
% OBJECT and cannot join a neighbouring run; a bare "IV" can. So in the RTL TOC
% paragraph a Latin number and a Latin-INITIAL title merge into a SINGLE LTR run
% and are laid out left-to-right as a unit: the number lands on the wrong side of
% the row, and a mixed title is scrambled outright. Measured on \part{MCP ...}:
% the printed order was "<tail> V MCP" -- title fragment, number, title fragment.
% The chapter/section rows one line below, with the SAME Latin-initial titles,
% were correct on the same page. That contrast is the whole diagnosis.
% \en{} IS NOT A WORKAROUND, measured both ways on one fixture: it fixes a
% PURE-Latin title (\part{\en{Hooks}} renders correctly) and does NOT fix a MIXED
% one (\part{\en{MCP} <hebrew>} stays scrambled). babel's language switch changes
% font and hyphenation, not the bidi run structure -- so the defect is not
% reachable from the source side at all, which is why it belongs here.
% Boxing the number is the entire fix and is a STRICT NO-OP where the bug never
% bit: a Hebrew-titled document compiles BYTE-IDENTICAL before and after (md5,
% 3 clean runs), because there the neutral glue already sat between an L run and
% an R run. The PDF outline is byte-identical too -- hyperref's \pdfstringdef
% neutralises \mbox, so bookmark strings keep their part numbers.
% \numberline{\thepart} would also fix the direction, but it re-geometries every
% part row in the fleet and DROPS the number from the part bookmark (hyperref
% strips \numberline). \mbox changes nothing that was already right.
% The search pattern deliberately carries no # token: etoolbox cannot match a
% parameter token in a patch pattern (measured -- all three chain members
% reported "not patchable" while the pattern included #1). "\thepart \hspace
% {1em}" is brace-balanced and occurs exactly once in the wrapper body.
\newif\iflakin@parttocnum
\def\lakin@parttocnum#1{\iflakin@parttocnum\else
  \patchcmd{#1}{\thepart \hspace {1em}}{\mbox {\thepart }\hspace {1em}}
    {\lakin@parttocnumtrue}{}\fi}
\@ifundefined{NR@part}{}{\lakin@parttocnum\NR@part}
\@ifundefined{H@old@part}{}{\lakin@parttocnum\H@old@part}
\@ifundefined{@part}{}{\lakin@parttocnum\@part}
\iflakin@parttocnum\else\PackageWarning{hebrew-lualatex-pdf}{Part TOC number is
  not boxed: no known \noexpand\part wrapper matched. A part title that BEGINS
  with Latin will print its number on the wrong side of its TOC row}\fi
\patchcmd{\section}{\bfseries}{\bfseries\boldmath}{}
  {\PackageError{hebrew-lualatex-pdf}{Cannot patch section math weight}{Unsupported document class}}
\patchcmd{\subsection}{\bfseries}{\bfseries\boldmath}{}
  {\PackageError{hebrew-lualatex-pdf}{Cannot patch subsection math weight}{Unsupported document class}}
\patchcmd{\subsubsection}{\bfseries}{\bfseries\boldmath}{}
  {\PackageError{hebrew-lualatex-pdf}{Cannot patch subsubsection math weight}{Unsupported document class}}
\patchcmd{\paragraph}{\bfseries}{\bfseries\boldmath}{}{}
\patchcmd{\subparagraph}{\bfseries}{\bfseries\boldmath}{}{}
\@ifundefined{l@chapter}{}{%
  \patchcmd{\l@chapter}{\bfseries}{\bfseries\boldmath}{}{}}
\makeatother
\setcounter{tocdepth}{\lakittocdepth}
% <hllpdf:measured-toc>
% v4.16.0: Measure actual number labels, not a guessed digit count. Read only
% canonical contentsline records; never execute titles or addtocontents code
% during the measuring pass. Widths are local to this individual list and are
% available before its first row: no per-row drift and no extra aux pass.
\makeatletter
\newread\hllpdf@tocstream
\newdimen\hllpdf@toc@chapter
\newdimen\hllpdf@toc@section
\newdimen\hllpdf@toc@subsection
\newdimen\hllpdf@toc@subsubsection
\newdimen\hllpdf@toc@paragraph
\newdimen\hllpdf@toc@subparagraph
\newdimen\hllpdf@toc@figure
\newdimen\hllpdf@toc@table
\newdimen\hllpdf@tocpage
\newdimen\hllpdf@toc@sectionindent
\newdimen\hllpdf@toc@subsectionindent
\newdimen\hllpdf@toc@subsubsectionindent
\newdimen\hllpdf@toc@paragraphindent
\newdimen\hllpdf@toc@subparagraphindent
\def\hllpdf@tocdefaults{%
  \hllpdf@toc@chapter=2em
  \hllpdf@toc@section=3.2em
  \hllpdf@toc@subsection=3.9em
  \@ifundefined{c@chapter}{%
    \hllpdf@toc@subsubsection=3.2em
    \hllpdf@toc@paragraph=4.1em
    \hllpdf@toc@subparagraph=5em
  }{%
    \hllpdf@toc@subsubsection=4.1em
    \hllpdf@toc@paragraph=5em
    \hllpdf@toc@subparagraph=6em
  }%
  \hllpdf@toc@figure=2.3em
  \hllpdf@toc@table=2.3em
  \hllpdf@tocpage=\@pnumwidth\relax
}
\def\hllpdf@tocindents{%
  % Preserve the established shallow layout when its slots already fit.
  \hllpdf@toc@sectionindent=1.5em
  \@tempdima=\hllpdf@toc@chapter \advance\@tempdima by -2em
  \advance\hllpdf@toc@sectionindent by \@tempdima
  \hllpdf@toc@subsectionindent=\dimexpr
    \hllpdf@toc@sectionindent+\hllpdf@toc@section\relax
  \hllpdf@toc@subsubsectionindent=\dimexpr
    \hllpdf@toc@subsectionindent+\hllpdf@toc@subsection\relax
  \hllpdf@toc@paragraphindent=\dimexpr
    \hllpdf@toc@subsubsectionindent+\hllpdf@toc@subsubsection\relax
  \hllpdf@toc@subparagraphindent=\dimexpr
    \hllpdf@toc@paragraphindent+\hllpdf@toc@paragraph\relax
}
\def\hllpdf@tocfont@chapter{\bfseries\boldmath}
\def\hllpdf@tocfont@part{\large\bfseries\boldmath}
\def\hllpdf@tocentryfont{%
  \ifcsname hllpdf@tocfont@\hllpdf@toclevel\endcsname
    \csname hllpdf@tocfont@\hllpdf@toclevel\endcsname
  \fi}
\long\def\hllpdf@tocdiscard#1\hllpdf@tocstop{}
\def\hllpdf@tocscanrecord{%
  \@ifnextchar\contentsline{\hllpdf@tocscanline}{%
    \@ifnextchar\babel@toc{\hllpdf@tocscanlanguage}{\hllpdf@tocdiscard}}}
\long\def\hllpdf@tocscanlanguage\babel@toc#1#2#3\hllpdf@tocstop{%
  % Babel records are declarations, not arbitrary code. Track the language
  % without executing the original directive or changing the surrounding list.
  \def\hllpdf@toclanguage{#1}}
\long\def\hllpdf@tocscanline\contentsline#1#2#3#4#5\hllpdf@tocstop{%
  \def\hllpdf@toclevel{#1}%
  \ifcsname hllpdf@toc@#1\endcsname
    \hllpdf@tocscantitle#2\hllpdf@tocstop
  \fi
  % The page-number column needs its own measurement (including bold chapter
  % rows). Its extra clearance separates the page label from dot leaders.
  \setbox\@tempboxa=\hbox{\foreignlanguage{\hllpdf@toclanguage}{\hllpdf@tocentryfont #3\kern.3em}}%
  \ifdim\wd\@tempboxa>\hllpdf@tocpage
    \hllpdf@tocpage=\wd\@tempboxa
  \fi}
\def\hllpdf@tocscantitle{%
  \@ifnextchar\numberline{\hllpdf@tocscanlabel}{\hllpdf@tocdiscard}}
\long\def\hllpdf@tocscanlabel\numberline#1#2\hllpdf@tocstop{%
  \setbox\@tempboxa=\hbox{\foreignlanguage{\hllpdf@toclanguage}{\hllpdf@tocentryfont #1\kern.5em}}%
  \ifdim\wd\@tempboxa>\csname hllpdf@toc@\hllpdf@toclevel\endcsname
    \csname hllpdf@toc@\hllpdf@toclevel\endcsname=\wd\@tempboxa
  \fi}
\def\hllpdf@tocread{%
  \ifeof\hllpdf@tocstream\else
    \read\hllpdf@tocstream to\hllpdf@tocrecord
    \expandafter\hllpdf@tocscanrecord\hllpdf@tocrecord\hllpdf@tocstop
    \expandafter\hllpdf@tocread
  \fi}
\def\hllpdf@tocmeasure#1{%
  \makeatletter
  \edef\hllpdf@toclanguage{\languagename}%
  \hllpdf@tocdefaults
  \openin\hllpdf@tocstream=\jobname.#1\relax
  \hllpdf@tocread
  \closein\hllpdf@tocstream
  \hllpdf@tocindents
  \edef\@pnumwidth{\the\hllpdf@tocpage}%
  \@tempdima=\dimexpr\hllpdf@tocpage+1em\relax
  \ifdim\@tempdima>\@tocrmarg\relax
    \edef\@tocrmarg{\the\@tempdima}%
  \fi}
\let\hllpdf@originalstarttoc\@starttoc
\def\@starttoc#1{%
  \begingroup
    \def\hllpdf@tocext{#1}%
    \def\hllpdf@tocwanted{toc}%
    \ifx\hllpdf@tocext\hllpdf@tocwanted \hllpdf@tocmeasure{#1}\else
      \def\hllpdf@tocwanted{lof}%
      \ifx\hllpdf@tocext\hllpdf@tocwanted \hllpdf@tocmeasure{#1}\else
        \def\hllpdf@tocwanted{lot}%
        \ifx\hllpdf@tocext\hllpdf@tocwanted \hllpdf@tocmeasure{#1}\fi
      \fi
    \fi
    \hllpdf@originalstarttoc{#1}%
  \endgroup}
\hllpdf@tocdefaults
\hllpdf@tocindents
\@ifundefined{l@chapter}{}{%
  \patchcmd{\l@chapter}{\setlength\@tempdima{1.5em}}%
    {\setlength\@tempdima{\hllpdf@toc@chapter}}{}%
    {\PackageError{hebrew-lualatex-pdf}{Cannot size chapter TOC numbers}%
       {Unsupported contents layout: the chapter number slot was not found}}}
\renewcommand*\l@section{\@dottedtocline{1}%
  {\hllpdf@toc@sectionindent}{\hllpdf@toc@section}}
\renewcommand*\l@subsection{\@dottedtocline{2}%
  {\hllpdf@toc@subsectionindent}{\hllpdf@toc@subsection}}
\renewcommand*\l@subsubsection{\@dottedtocline{3}%
  {\hllpdf@toc@subsubsectionindent}{\hllpdf@toc@subsubsection}}
\renewcommand*\l@paragraph{\@dottedtocline{4}%
  {\hllpdf@toc@paragraphindent}{\hllpdf@toc@paragraph}}
\renewcommand*\l@subparagraph{\@dottedtocline{5}%
  {\hllpdf@toc@subparagraphindent}{\hllpdf@toc@subparagraph}}
\renewcommand*\l@figure{\@dottedtocline{1}{1.5em}{\hllpdf@toc@figure}}
\renewcommand*\l@table{\@dottedtocline{1}{1.5em}{\hllpdf@toc@table}}
\makeatother
% </hllpdf:measured-toc>
\setcounter{secnumdepth}{\lakitsecnumdepth}
% Public preamble-only caption-numbering policy. The historical within-major
% default is untouched unless this command is called. "continuous" removes the
% reset edge as well as the printed prefix, so Figure 1/2 and Table 1/2 stay
% unique across sections/chapters. Returning to "within-major" restores the
% original class-aware displays and native destination conventions. The policy
% never replaces \figcap/\tabcap or their measured link-placement mechanism.
\makeatletter
\let\lakin@majorfig\thefig
\let\lakin@majortbl\thetbl
\let\lakin@majorHfig\theHfig
\let\lakin@majorHtbl\theHtbl
\newcommand{\lakitcaptionnumbering}[1]{%
  \edef\lakin@numbering{#1}%
  \def\lakin@continuous{continuous}\def\lakin@withinmajor{within-major}%
  \ifx\lakin@numbering\lakin@continuous
    \@ifundefined{c@chapter}{%
      \@removefromreset{fig}{section}\@removefromreset{tbl}{section}%
    }{\@removefromreset{fig}{chapter}\@removefromreset{tbl}{chapter}}%
    \renewcommand{\thefig}{\arabic{fig}}\renewcommand{\thetbl}{\arabic{tbl}}%
    \def\theHfig{\arabic{fig}}\def\theHtbl{\arabic{tbl}}%
  \else\ifx\lakin@numbering\lakin@withinmajor
    \@ifundefined{c@chapter}{%
      \@removefromreset{fig}{section}\@removefromreset{tbl}{section}%
      \@addtoreset{fig}{section}\@addtoreset{tbl}{section}%
    }{%
      \@removefromreset{fig}{chapter}\@removefromreset{tbl}{chapter}%
      \@addtoreset{fig}{chapter}\@addtoreset{tbl}{chapter}}%
    \let\thefig\lakin@majorfig\let\thetbl\lakin@majortbl
    \let\theHfig\lakin@majorHfig\let\theHtbl\lakin@majorHtbl
  \else
    \PackageError{hebrew-english-latex-studio}{Unknown caption-numbering policy '#1'}%
      {Use continuous or within-major before begin{document}.}%
  \fi\fi}
\@onlypreamble\lakitcaptionnumbering
\makeatother
% Architecture readiness sentinel. Optional runtime layers require the complete
% core, so the marker is defined only after every mechanism/baseline definition.
% <hllpdf:architecture-ready>
\def\hllpdfcoreloaded{1}
% </hllpdf:architecture-ready>

% v4.22: registered file resources, not an assertion about embedded subset bytes.
% The log framing avoids long-path wrapping and preserves UTF-8 as hex.
% <hllpdf:font-inventory>
\directlua{
  local function inventory()
    local names = {}
    local unnamed = 0
    for id, item in font.each() do
      local path = item.filename
      if path and not (path == "") then names[path] = true else unnamed = unnamed + 1 end
    end
    local sorted = {}
    for path in pairs(names) do table.insert(sorted,path) end
    table.sort(sorted)
    texio.write_nl("log","HLLFONT|BEGIN|1")
    for index,path in ipairs(sorted) do
      local hex = path:gsub(".",function(c) return string.format(string.char(37) .. "02x",string.byte(c)) end)
      local count = math.ceil(string.len(hex)/48)
      texio.write_nl("log","HLLFONT|FILE|" .. index .. "|" .. count)
      for offset=1,string.len(hex),48 do texio.write_nl("log","HLLFONT|HEX|" .. hex:sub(offset,offset+47)) end
    end
    texio.write_nl("log","HLLFONT|END|" .. rawlen(sorted) .. "|" .. unnamed)
    texio.write_nl("log","")
  end
  luatexbase.add_to_callback("finish_pdffile",inventory,"hllpdf.runtime.fonts")
}
% </hllpdf:font-inventory>

%
% Supported class contract: report or article.  Alternate paper sizes must also
% be declared to build.sh through EXPECTED_PAPER; the PDF gate never infers them.
% Do NOT use these hooks as a theme mechanism: theme overlays must not own page
% geometry, class selection, or direction semantics. A declared presentation-policy
% overlay may intentionally change numbering visibility after the preamble.
% ============================================================================
% Runtime architecture guard: generated from scripts/architecture_contract.py.
% The readiness sentinel remains delayed until the END of this file.
% <hllpdf:architecture-guard>
\ifdefined\hllpdfcoreloaded
  \errmessage{hllpdf: preamble.tex loaded more than once}
  \expandafter\endinput
\fi
% </hllpdf:architecture-guard>

\ifdefined\lakitdocumentclass\else\def\lakitdocumentclass{report}\fi
\ifdefined\lakitclassoptions\else\def\lakitclassoptions{12pt,a4paper}\fi
\ifdefined\lakitgeometryoptions\else\def\lakitgeometryoptions{margin=2.3cm,headheight=14pt}\fi
\ifdefined\lakitlinespread\else\def\lakitlinespread{1.3}\fi
\ifdefined\lakittocdepth\else\def\lakittocdepth{1}\fi
\ifdefined\lakitsecnumdepth\else\def\lakitsecnumdepth{2}\fi
% Paragraph organization belongs to the document layout, independently of its
% language and theme. Preserve historical class dimensions unless explicitly
% selecting block paragraphs; never suppress indentation with an everypar hack.
\ifdefined\lakitparagraphstyle\else\def\lakitparagraphstyle{indented}\fi
\edef\lakitparagraphstyle{\lakitparagraphstyle}
\def\hllpdfparagraphblock{block}
\def\hllpdfparagraphindented{indented}
\ifx\lakitparagraphstyle\hllpdfparagraphblock\else
  \ifx\lakitparagraphstyle\hllpdfparagraphindented\else
    \errmessage{hllpdf: paragraph style must be block or indented}
  \fi
\fi
% v4.3.1: the public class hook has a deliberately narrow support contract.
% Before this guard, \def\lakitdocumentclass{book} silently compiled even though
% every heading/TOC regression test targets report/article; an unsupported shell
% could therefore look "supported" until a subtle class-specific defect shipped.
% Fail before \documentclass rather than letting an undocumented combination enter
% the mechanism core. Plain-TeX comparison is used because LaTeX packages are not
% available yet at this point.
\edef\lakitvalidatedclass{\lakitdocumentclass}
\def\lakitreportclass{report}
\def\lakitarticleclass{article}
\ifx\lakitvalidatedclass\lakitreportclass
\else\ifx\lakitvalidatedclass\lakitarticleclass
\else
  \errmessage{hllpdf: unsupported class \lakitdocumentclass; use report or article}
\fi\fi
\let\lakitvalidatedclass\relax
\let\lakitreportclass\relax
\let\lakitarticleclass\relax
% Language intent is independent of document class, theme and content domain.
% The two dominant-language profiles differ in authoring intent, not direction.
% A balanced bilingual document must explicitly choose a navigation/base locale.
% <hllpdf:language-profile>
\providecommand{\lakitlanguageprofile}{hebrew}
\edef\lakitlanguageprofile{\lakitlanguageprofile}
\def\hllpdfhebrew{hebrew}\def\hllpdfenglish{english}
\def\hllpdfhebrewmixed{hebrew-with-english}
\def\hllpdfenglishmixed{english-with-hebrew}
\def\hllpdfbilingual{bilingual}
\let\hllpdfexpectedlanguage\relax
\ifx\lakitlanguageprofile\hllpdfhebrew
  \let\hllpdfexpectedlanguage\hllpdfhebrew
\else\ifx\lakitlanguageprofile\hllpdfhebrewmixed
  \let\hllpdfexpectedlanguage\hllpdfhebrew
\else\ifx\lakitlanguageprofile\hllpdfenglish
  \let\hllpdfexpectedlanguage\hllpdfenglish
\else\ifx\lakitlanguageprofile\hllpdfenglishmixed
  \let\hllpdfexpectedlanguage\hllpdfenglish
\else\ifx\lakitlanguageprofile\hllpdfbilingual
  \ifdefined\lakitmainlanguage\else
    \errmessage{hllpdf: bilingual profile requires explicit lakitmainlanguage}
  \fi
\else
  \errmessage{hllpdf: unsupported language profile \lakitlanguageprofile}
\fi\fi\fi\fi\fi
\ifdefined\lakitmainlanguage
  \edef\lakitmainlanguage{\lakitmainlanguage}
\else
  \let\lakitmainlanguage\hllpdfexpectedlanguage
\fi
\ifx\lakitmainlanguage\hllpdfhebrew\else
  \ifx\lakitmainlanguage\hllpdfenglish\else
    \errmessage{hllpdf: main language must be hebrew or english}
  \fi
\fi
\ifx\hllpdfexpectedlanguage\relax\else
  \ifx\lakitmainlanguage\hllpdfexpectedlanguage\else
    \errmessage{hllpdf: main language conflicts with language profile}
  \fi
\fi
\typeout{HLLLANG|\lakitlanguageprofile|\lakitmainlanguage}
% </hllpdf:language-profile>
\documentclass[\lakitclassoptions]{\lakitdocumentclass}
\usepackage[no-math]{fontspec}        % (1) ★ no-math — THE critical line (see recipe.md)
% layout=lists asks Babel to adapt list parshape when paragraph and page
% directions differ. Without it, minority-language list bodies lose their
% indentation and labels hang outside the text block, despite a clean log.
% The upstream mechanism preserves ordinary same-direction list geometry.
\usepackage[bidi=basic,layout=lists]{babel}
\ifx\lakitmainlanguage\hllpdfenglish
  \babelprovide[import]{hebrew}
  \babelprovide[import,main]{english}
\else
  \babelprovide[import,main]{hebrew}
  \babelprovide[import]{english}
\fi
% v4.22: bind canonical families and fallbacks to the selected installation.
% Lookup by basename or family alone is not proof of the file actually used.
% <hllpdf:font-binding>
\directlua{
  local dir = os.getenv("HEBREW_LATEX_FONT_DIR")
  if not dir or dir == "" then
    local data = os.getenv("XDG_DATA_HOME")
    if not data or data == "" or not (data:sub(1,1) == "/") then
      data = (os.getenv("HOME") or "") .. "/.local/share"
    end
    dir = data .. "/fonts/hebrew-lualatex-pdf"
  end
  if not (dir:sub(1,1) == "/") then dir = lfs.currentdir() .. "/" .. dir end
  if dir:find("[",1,true) or dir:find("]",1,true) or dir:find(string.char(0),1,true)
     or dir:find(string.char(10),1,true) or dir:find(string.char(13),1,true) then
    tex.error("hllpdf: unsupported character in font directory")
  end
  if not (lfs.attributes(dir,"mode") == "directory") then
    tex.error("hllpdf: font directory is missing")
  end
  hllpdf_font_dir = dir:gsub("/+$", "") .. "/"
  tex.sprint(string.char(92) .. "def" .. string.char(92) .. "hllpdffontdir{")
  tex.sprint(-2,hllpdf_font_dir)
  tex.sprint("}")
  luaotfload.add_fallback("hebfb",{"[" .. hllpdf_font_dir .. "DejaVuSerif.ttf]:mode=node;"})
  luaotfload.add_fallback("hebfbbf",{"[" .. hllpdf_font_dir .. "DejaVuSerif-Bold.ttf]:mode=node;"})
}
% </hllpdf:font-binding>
% v3.22: Frank Ruhl Libre ships NO italic face. Instead of silently-upright
% \emph (v<=3.20) or bold-as-emphasis (v3.21), synthesize a slanted face with
% FakeSlant — real, VISIBLE italic emphasis with correct \emph nesting
% semantics (inner \emph toggles back upright). Verified empirically at
% 300dpi: slant clearly visible, 0 errors, 0 missing glyphs.
\babelfont{rm}[Path={\hllpdffontdir},RawFeature={fallback=hebfb},
  BoldFont={FrankRuhlLibre-Bold.ttf},
  BoldFeatures={RawFeature={fallback=hebfbbf}},   % v3.56 — see hebfbbf above
  ItalicFont={FrankRuhlLibre-Regular.ttf}, ItalicFeatures={FakeSlant=0.22},
  % Bold-italic uses the weight-matched fallback and a synthetic slant.
  BoldItalicFont={FrankRuhlLibre-Bold.ttf}, BoldItalicFeatures={RawFeature={fallback=hebfbbf},FakeSlant=0.22}
  ]{FrankRuhlLibre-Regular.ttf}
% Sans fallbacks preserve family and weight for the few missing codepoints.
\directlua{luaotfload.add_fallback("hebfbsans",{"[" .. hllpdf_font_dir .. "DejaVuSans.ttf]:mode=node;"})}
\directlua{luaotfload.add_fallback("hebfbsansbf",{"[" .. hllpdf_font_dir .. "DejaVuSans-Bold.ttf]:mode=node;"})}
\babelfont{sf}[Path={\hllpdffontdir},RawFeature={fallback=hebfbsans},
  BoldFont={Heebo-Bold.ttf},
  BoldFeatures={RawFeature={fallback=hebfbsansbf}}]{Heebo-Regular.ttf}
% v3.24: Scale=MatchLowercase — DejaVu Sans Mono has an outsized x-height, so
% unscaled inline code visually SHOUTS next to the delicate Frank Ruhl body
% (measured on a 197-page build: scaled, a code atom sits in the line at the
% body's weight, and the document tightened by ~5%). Applies to listings too.
\babelfont{tt}[Path={\hllpdffontdir},RawFeature={fallback=hebfb},Scale=MatchLowercase,
  BoldFont={DejaVuSansMono.ttf},
  BoldFeatures={RawFeature={fallback=hebfb},FakeBold=1.5},
  ItalicFont={DejaVuSansMono.ttf},
  ItalicFeatures={RawFeature={fallback=hebfb},FakeSlant=0.2},
  BoldItalicFont={DejaVuSansMono.ttf},
  BoldItalicFeatures={RawFeature={fallback=hebfb},FakeBold=1.5,FakeSlant=0.2}
  ]{DejaVuSansMono.ttf}

% Explicit multilingual/English profiles use real Latin faces and italics.
% The legacy/default Hebrew profile retains its established measured typography.
% Latin Modern is provided by the TeX distribution, not the pinned font bundle;
% recorder/runtime evidence reports it as an additional toolchain resource.
\ifx\lakitlanguageprofile\hllpdfhebrew\else
  \babelfont[english]{rm}[Ligatures=TeX,RawFeature={fallback=hebfb},
    BoldFeatures={RawFeature={fallback=hebfbbf}},
    BoldItalicFeatures={RawFeature={fallback=hebfbbf}}]{Latin Modern Roman}
  \babelfont[english]{sf}[Ligatures=TeX,RawFeature={fallback=hebfbsans},
    BoldFeatures={RawFeature={fallback=hebfbsansbf}},
    BoldItalicFeatures={RawFeature={fallback=hebfbsansbf}}]{Latin Modern Sans}
\fi

\usepackage{amsmath,amssymb,mathtools}
% amsthm supplies theorem/proof machinery. Environment names remain
% document-specific; the generic proof label is localized below.
\usepackage{amsthm}
% Proof names follow the surrounding prose, including return from an island.
\addto\extrashebrew{\renewcommand{\proofname}{הוכחה}}
\addto\extrasenglish{\renewcommand{\proofname}{Proof}}
\ifx\lakitmainlanguage\hllpdfenglish
  \renewcommand{\proofname}{Proof}
\else
  \renewcommand{\proofname}{הוכחה}
\fi
% v3.24.1: equation tags under babel-bidi. A tag re-emitted by mathtools'
% showonlyrefs inside align inherits the RTL base and its parens MIRROR —
% ")2.1(" instead of "(2.1)" (plain align/equation tags are unaffected;
% verified on a minimal case + a 198-page build). LTR-isolate at the tag-BOX
% level: mathtools' show/hide + numbering logic is untouched. Do NOT override
% \tagform@ instead — that clobbers the showonlyrefs mechanism (measured:
% every tag shows, numbering breaks).
% v3.28.1: the v3.24.1 wrap must apply ONLY to tags of display environments.
% Inline \eqref in Hebrew PROSE also goes through \maketag@@@, and there the
% english-wrap MIRRORS the parens — ")1(" — while the unwrapped original is
% correct (proven on a 2x2 isolation matrix: patch x showonlyrefs; regressed
% every prose \eqref since v3.24.1, caught by the render gate's R4). The two
% call sites are indistinguishable by \ifmmode (tag boxes are text mode), so
% the display environments themselves raise a flag; \eqref in prose sees it
% false and gets the original untouched \maketag@@@.
\usepackage{etoolbox}

% babel bidi=basic (including upstream 26.9) assumes every automatic
% discretionary has a one-glyph replacement. LuaTeX can legally produce one
% with replace=nil after a hyphenated Latin value; Babel then indexes nil and the
% whole build dies. Do the narrow upstream-compatible correction: while Babel
% classifies the list, present only that malformed disc as an ordinary
% discretionary, then restore its subtype. Babel still reorders the ENTIRE box;
% there is no pcall and no successful fallback to an unreordered node list.
% If Babel throws for any other reason, the build remains fatal. The synthetic
% nil-disc probe and the real table gallery in test_all.sh pin both layers.
\AtBeginDocument{\directlua{
  if Babel and Babel.bidi and not Babel.heb_nil_disc_compat then
    Babel.heb_nil_disc_compat = true
    local original = Babel.bidi
    local disc_id = node.id('disc')
    Babel.bidi = function(head, ispar, hdir)
      local hidden = {}
      for item in node.traverse_id(disc_id, head) do
        if item.subtype == 2 and (not item.replace or not item.replace.char) then
          table.insert(hidden, item)
          item.subtype = 0
        end
      end
      local result = original(head, ispar, hdir)
      for _, item in ipairs(hidden) do item.subtype = 2 end
      return result
    end
  end
}}
\makeatletter
\newif\ifHEBdisptag
\let\ORIGmaketag@@@\maketag@@@
\def\maketag@@@#1{\ifHEBdisptag\ORIGmaketag@@@{\foreignlanguage{english}{#1}}\else\ORIGmaketag@@@{#1}\fi}
\@for\HEB@env:=equation,equation*,align,align*,gather,gather*,multline,multline*,flalign,flalign*,alignat,alignat*,eqnarray,eqnarray*\do{%
  \expandafter\AtBeginEnvironment\expandafter{\HEB@env}{\HEBdisptagtrue}}
\makeatother
% English page + Hebrew prose: Babel's equation-only tag flip uses a primitive
% leqno path that shifts the formula to the left edge and the tag into the body
% of an English page. Align's full-width alignment is unaffected. Give only this
% display a physical LTR container and native amsmath eqno/leqno. The engine then
% owns real tag width, collision handling and labels; no proxy tag box hides ink
% from its fit calculation. Babel's saved \bbl@eqnodir still restores Hebrew text
% inside math, and the prose locale/math direction are never switched.
% Close the prior RTL paragraph before normalization. After the environment's
% group restores its directions, close TeX's empty resumed display paragraph and
% use native LaTeX continuation handling: same-paragraph prose has no indent,
% while an authored blank line starts an ordinary indented RTL paragraph.
% Register after Babel's setup; remain local to equation(*) in a main-English
% document's actual RTL prose. Native Hebrew, ordinary English and explicit
% fleqn retain their original paths. No page/display dimensions are assigned.
% <hllpdf:minority-rtl-equation-geometry>
\makeatletter
\newcommand{\HEB@minorityprosecontinuation}{\par\@doendpe}
\newcommand{\HEB@minorityequationgeometry}{%
  \ifx\lakitmainlanguage\hllpdfenglish
    \ifnum\pardirection=\@ne
      \if@fleqn\else
        \par\aftergroup\HEB@minorityprosecontinuation
        \iftagsleft@\let\veqno\@@leqno
        \else\let\veqno\@@eqno\fi
        \bodydir TLT\pardir TLT\textdir TLT\relax
      \fi
    \fi
  \fi}
\AtBeginDocument{%
  \AddToHook{env/equation/begin}{\HEB@minorityequationgeometry}%
  \AddToHook{env/equation*/begin}{\HEB@minorityequationgeometry}}
\makeatother
% </hllpdf:minority-rtl-equation-geometry>
% v3.63: \shortintertext rendered Hebrew MIRROR-REVERSED and flush LEFT — an entire
% commentary line inside an align printed backwards. babel's bidi=basic patches
% \intertext@ only (luababel.def), and mathtools' \shortintertext uses the identical
% \vbox{\normalbaselines … \noindent\ignorespaces#1\par} skeleton but is never
% patched, so its paragraph keeps pardir/textdir LTR: no bidi reordering (logical
% order on the page) and flush-left. Measured against \intertext with the same
% Applying babel's own replacement to \shortintertext@ makes the two identical.
% \h{}/\mbox fix the letter order but leave the line flush LEFT, so they are not
% substitutes. Guarded + \AtBeginDocument because mathtools \let-copies at load time.
% NB \ExplSyntaxOn is REQUIRED: mathtools' internals carry expl3 names, and without
% the expl3 catcodes the `_` is read as a subscript and the patch detonates far away
% with "! Illegal unit of measure (pt inserted)" at the \end{align} (measured).
\makeatletter
\AtBeginDocument{%
  \@ifundefined{shortintertext@}{}{%
    \ExplSyntaxOn
    \bbl@sreplace\shortintertext@{\normalbaselines}%
      {\normalbaselines\ifx\bbl@eqnodir\relax\else\bbl@pardir\@ne\bbl@eqnodir\fi}%
    \ExplSyntaxOff}}
\makeatother
% geometry inherits the canonical A4 sheet from \documentclass; build.sh then
% verifies every rendered MediaBox and CropBox, so a silent Letter fallback is
% impossible to publish as a clean default build.
\usepackage[\lakitgeometryoptions]{geometry}
\linespread{\lakitlinespread}       % layout hook; canonical default remains 1.3
\ifx\lakitparagraphstyle\hllpdfparagraphblock
  \setlength{\parindent}{0pt}
  \setlength{\parskip}{6pt plus 1pt minus 1pt}
\fi
% These dimensions remain ordinary scoped TeX lengths: authors can customize
% them after core input. Language switches inherit them; lists keep their own
% indentation and narrow reference layouts establish their local paragraph gaps.
\typeout{HLLPAR|\lakitparagraphstyle}
\usepackage{xcolor}
\usepackage{mdframed}
\usepackage{listings}                 % code listings (assembly / C / pseudo-code)
\usepackage{tikz}
\usetikzlibrary{arrows.meta,positioning,calc,shapes.geometric,fit,backgrounds}
\usepackage{pgfplots}\pgfplotsset{compat=1.18}
\usepackage{quantikz}                 % quantum circuits (remove if unused)
\usepackage{booktabs}
\usepackage{array}
\usepackage{float}                    % [H] = place table exactly here (reference doc: no drifting)
\usepackage{enumitem}
\usepackage{needspace}

% ============================================================================
% ★★ v3.4 — THE DIRECTION-ISOLATION TRIO (the heart of mixing Hebrew + Latin)
% Mental model:
%   \h{...}    Hebrew (RTL) island      — Hebrew text inside an LTR context
%                                          (TikZ nodes, the english `tikzpic`).
%   \en{...}   English (LTR) isolate    — Latin text / parentheses inside a
%                                          Hebrew (RTL) context. KEEPS PARENS
%                                          AND ORDER CORRECT.
%   \code{...} \en + monospace          — code inside Hebrew / inside a frametitle.
% CRITICAL: never wrap Latin in \h{} (it is for HEBREW). Inside a TikZ node,
% Latin in \h{} renders REVERSED (Datapath->htapataD, CPU->UPC) and parens
% mirror. Isolate every Latin/paren chunk with \en{} or \code{}. See
% references/figures-boxes-listings.md.
% ---- TEXT INSIDE MATH (v3.8) ----------------------------------------------
% Put ANY text that sits inside math in a \text{}. The ONLY thing that is
% invalid in math is a bare \h/\en (= \foreignlanguage) used DIRECTLY, outside
% \text -> "\rmfamily invalid in math mode". Inside \text:
%   Hebrew label           -> \text{...}            (bare Hebrew renders fine;
%                              \h is optional — BUT NOT ALWAYS HARMLESS (v3.67):
%                              if the \h group OPENS with a nested island
%                              ($..$, \en{}, \code{}, \texttt{}, \mbox{}) the
%                              WHOLE group loses bidi and the Hebrew prints in
%                              LOGICAL order = gibberish. Measured:
%                                \text{\h{$f$ קבועה}}  -> reads "העובק"  BROKEN
%                                \text{\h{קבועה $f$}}  -> correct
%                                \text{$f$ קבועה}      -> correct (no \h)
%                              Hebrew first inside \h, or drop the \h. Flagged
%                              H-OPENS-WITH-LTR-GROUP.)
%   plain Latin word        -> \text{...}            (no marker needed)
%   Latin WITH ( ) or [ ]   -> \text{\en{...}}       (\en forces LTR; a plain
%                              \text mirrors the parens in RTL -> "Spec (TNR(" )
%   f(Hebrew/text arg)      -> \text{f}(\text{...})  -- keep the parens at MATH
%                              level, NOT \text{f(arg)} (parens inside the box mirror)
%   code in math            -> \mathtt{...}  (or \texttt{}, math-guarded below)
% NEVER write \en{...} or \h{...} straight into $...$ / \[..\] (outside \text);
% always via \text{}. (\mbox{...} also works but does not scale in scripts.)
% ============================================================================
\newcommand{\h}[1]{\foreignlanguage{hebrew}{#1}}
% v3.42.1 — \en islands do NOT hyphenate (auto OR at explicit hyphens). English hyphenation patterns were
% tearing identifiers and terms ACROSS THE RTL SEAM ("IO Regis-/ters",
% "In-/put", "Band-/width", "Per-/formance", "Multi-/ple" — all shipped on
% real pages; probe-proven: "ALUS-/rcB", "Mem-/ToReg"). An island is a label
% or a term, not English prose — a hyphen tear there is never right. A/B on
% a full 197pp lakit: identical overfull sets, same page count, ~11
% paragraphs healed, one mild underfull (badness 3525) that replaced a
% mutilated "Regis-/ters". Long English PROSE belongs in the englishprose
% environment below, which restores explicit-hyphen breaks while its paragraph
% is still active. A raw otherlanguage block inherits the global prohibition.
% (\l@nohyphenation needs \csname:
% '@' is not a letter here.)
% v3.43.2 — \exhyphenpenalty moved OUT of \en to a GLOBAL assignment below.
% Render-proven: a group-local \exhyphenpenalty is DEAD CODE — TeX reads the
% line-breaking parameters at \par time, after the \en group has closed, so
% \en{write-through} / plain not-taken still tore at the "-" ("write-/
% through" shipped on a real exam page). \language is different: it is
% recorded in the node list as text is added, so the in-group nohyphenation
% assignment above DOES work (and stays). The global penalty was A/B-measured
% on a 197pp lakit: identical overfull set (25), identical page count, the
% not-/taken tear healed. Hebrew is unaffected in practice — Hebrew compounds
% use maqaf (־), not ASCII "-".
\newcommand{\en}[1]{\foreignlanguage{english}{\language=\csname l@nohyphenation\endcsname #1}}
% Babel runs extras locally for inline islands and paragraph language changes.
% An English main document must not inherit the Hebrew no-hyphen paragraph rule.
\addto\extrashebrew{\exhyphenpenalty=10000\relax}
\addto\extrasenglish{\exhyphenpenalty=50\relax}
\ifx\lakitmainlanguage\hllpdfenglish
  \exhyphenpenalty=50\relax
\else
  \exhyphenpenalty=10000\relax
\fi
% Full English paragraphs need ordinary explicit-hyphen breaks. Use this
% environment instead of a raw otherlanguage block; the paragraph ends while
% the local penalty is still active.
\newenvironment{englishprose}
  {\par\begin{otherlanguage}{english}\exhyphenpenalty=50\relax}
  {\par\end{otherlanguage}}
\newenvironment{hebrewprose}
  {\par\begin{otherlanguage}{hebrew}\exhyphenpenalty=10000\relax}
  {\par\end{otherlanguage}}
% Use for shell text whose locale must not depend on the last body paragraph.
\DeclareRobustCommand{\doclang}[1]{\foreignlanguage{\lakitmainlanguage}{#1}}

% ============================================================================
% ★★ — fixes for defects that are INVISIBLE in the log and in the checks
% and only surface on a VISUAL RENDER of the PDF. (Each verified empirically.)
% ----------------------------------------------------------------------------
% v3.22: \emph is now REAL slanted emphasis (FakeSlant italic face defined on
% the \babelfont{rm} line above), replacing the v3.21 \emph->\textbf hack: it
% preserves \emph's toggle semantics and keeps bold free for its own role.
% Prefer bold-as-emphasis for your doc? one line: \renewcommand{\emph}[1]{\textbf{#1}}
% page number: force the digits LTR so they never reverse (18 -> 81) in the RTL
%   footer. That reversal is a NON-DETERMINISTIC bidi glitch (observed once in
%   ~80 pp); this LuaLaTeX build is not byte-reproducible, so the defensive
%   LTR wrap is the only reliable guard. Latin digits are always LTR anyway.
\renewcommand{\thepage}{\foreignlanguage{english}{\arabic{page}}}
% v3.70: \pagenumbering \gdef's \thepage and silently DESTROYS the wrap above —
% after \pagenumbering{arabic} the folio guard is gone for the whole rest of the
% document. Two collections discovered this independently and hand-restore the
% wrap in main.tex with the same comment; that knowledge now lives here instead.
% The wrapper re-applies the LTR wrap after the kernel call, for any style.
\makeatletter
\let\lakin@orig@pagenumbering\pagenumbering
\renewcommand{\pagenumbering}[1]{%
  \lakin@orig@pagenumbering{#1}%
  \renewcommand{\thepage}{\foreignlanguage{english}{\csname @#1\endcsname\c@page}}}
\makeatother
% emergencystretch: Hebrew does not hyphenate, so an unbreakable LTR island
%   (\en{...} / an inline math run) landing at the margin can leave a small
%   (~2-13pt) overfull with no legal breakpoint. This rescues ONLY those
%   inline-prose tie-breaks; it is a ceiling used only on offending lines and
%   does NOT touch display-math / listing / table overflow (those still surface
%   in the log for the content-level overflow pass). Verified: baseline overfull
%   -> 0 with this set, with no change to pagination or underfull.
\setlength{\emergencystretch}{3em}
% v3.40: inline math (almost) never breaks across lines. TeX's default
%   (binop 700 / rel 500) lets an inline run split at an operator —
%   $(1.f)\cdot 2^{e}$ shipped as "(1.f)·" ending one line and "2^e" opening
%   the next (reader-reported). In LTR that split is merely ugly; under an
%   RTL base the carried operator lands on the visually WRONG side of the
%   break and the formula reads backwards.
%   WHY 10000 (v3.45.0 — the display-first resolution). Inline math is
%   FORBIDDEN from breaking, so a torn island under an RTL base is
%   structurally impossible. The v3.40 cycle used 9999 ("break only as a last
%   resort") because a bare 10000 left ~30-39 WIDE inline equations
%   overflowing (worst 90.7pt) — but 9999's "last resort" fires OFTEN in
%   tight, non-hyphenating Hebrew: a reader found torn islands across the
%   fleet (~150 candidates; e.g. "P(y |" ending one line and "x_0)" opening
%   the next, read backwards). The step the 9999 cycle missed: every one of
%   those ~30 overflows is a wide equation that BELONGS in display. Convert
%   them to display FIRST, then 10000 gives 0 tears AND 0 overflows (measured,
%   12/12 fleet). Doctrine: inline math is ATOMIC; a wide equation goes in a
%   display, never inline. See content-style ("inline math is atomic").
\binoppenalty=10000 \relpenalty=10000
% v3.38: uniform sentence spacing. TeX's default (\nonfrenchspacing, sfcode
% 3000) is the OLD Anglo-American wide-space-after-period convention — Hebrew
% typography never had it, and modern Latin practice (Chicago et al.) dropped
% it too. Measured on real pages: post-period gaps up to 1.9-2.1x the line's
% mean word gap — the dominant "rivers" source in justified Hebrew prose.
% Survives babel language switches (probed: sfcode stays french inside \en{}
% islands and otherlanguage blocks).
\frenchspacing
% v3.38: widow/orphan protection. LaTeX's default penalties (150/150/50)
% allow a lone line at a page top/bottom almost freely. 10000 = forbid;
% this is SAFE here precisely because the document is \raggedbottom (the
% demoted line just moves on, no vertical stretch fight) — under flushbottom
% these values would cause underfull pages. raggedbottom is the report-class
% default; made explicit so a future class/option change cannot silently
% flip it and turn the 10000s harmful.
\raggedbottom
\widowpenalty=10000 \clubpenalty=10000 \displaywidowpenalty=10000
% em-dash: --- does NOT become — under Frank Ruhl (the font lacks tlig; verified
%   the width is identical with/without Ligatures=TeX, while the mechanism works
%   on Latin Modern). TYPE THE CHARACTER — (U+2014) DIRECTLY; likewise – (U+2013).
%   check_bidi_figures.py flags a --/--- that is adjacent to Hebrew.
% ============================================================================

% ---- STRUCTURAL DIRECTION-SAFETY for \texttt (v3.7) -------------------------
% A parenthesis, bracket, or $ INSIDE \texttt{} mirrors/migrates in RTL prose
% (MIPS 0($s0)->"(s0$)0", $rs->"rs$", arr[i]->"[arr[i") because \texttt sets the
% font but NOT the direction. Redefining \texttt to LTR-isolate its content fixes
% the WHOLE class at once (parens, brackets, $) so plain \texttt{lw $t0,0($s0)}
% is safe — no need to remember \code{}. Guarded by \ifmmode because
% \foreignlanguage is invalid in math (there, plain monospace is already LTR).
% NB (corrected v3.56 — the old note here said "stays tofu", which is WRONG):
% \texttt{Hebrew} is a FATAL build failure. The tt family is DejaVu Sans Mono
% (zero Hebrew) and the tt fallback hebfb is DejaVuSerif (also zero Hebrew), so
% the subsetter gets an empty glyph set and LuaTeX dies with
%   "! error: (file .../DejaVuSansMono.ttf) (type 2): there are no glyphs in the
%    subset"  -> no output PDF.
% The message names a FONT FILE, never Hebrew/\texttt/the source line. And it is
% non-local: add an unrelated Latin \code{} elsewhere and the fatal degrades into
% a SILENT DELETION of the Hebrew word. Put Hebrew in \h{}/normal text, never in
% \texttt{}/\code{}; flagged HEBREW-IN-MONO. To disable, delete this block.
\let\origtexttt\texttt
\protected\long\def\texttt#1{\ifmmode\origtexttt{#1}\else\foreignlanguage{english}{\origtexttt{#1}}\fi}
% v4.8.0: defer the direction test until TeX has inserted an alignment cell's
% math template. Unprotected expansion here tested \ifmmode too early at the
% FIRST token of aligned/array/cases/matrix cells. Engine-level \protected also
% preserves the original \texttt implementation without the self-recursion
% observed when retrofitting \DeclareRobustCommand to its saved alias.
\protected\long\def\code#1{\ifmmode\origtexttt{#1}\else\foreignlanguage{english}{\origtexttt{#1}}\fi}

% ---- \codeline — display-code line (v3.24): the code analogue of display math.
% Reader-friction rule (content-style.md §5b): inline \code{} is for a SINGLE
% ATOM only (one register/variable/mnemonic, a tight no-space expression).
% A multi-token expression — an instruction with operands, a spaced assignment,
% a register-binding list — zigzags the eye RTL<->LTR mid-sentence and, being
% an unbreakable island, line-breaks badly. Give it its own centered line with
% display-math vertical rhythm. \centerline = one box => the expression never
% wraps; \nopagebreak keeps it glued to its introducing sentence.
%   ...the Hebrew sentence ends with a colon:
%   \codeline{lw \$t1, -4(\$t0)}
%   and the prose continues (escape dollars as \$).
% v3.66 (first line of the body): flush a PENDING \item label before the \par.
% LaTeX defers the label box (\@labels) until the next \everypar, and an \hbox in
% vertical mode never fires \everypar — so `\item \codeline{...}` printed the code
% at y=69.1 and the item's own "1." BELOW it at y=103.6 (the number appears after
% its content, attached to nothing). \if@inlabel\leavevmode\fi is CONDITIONAL on
% purpose: an unconditional \leavevmode would open an empty paragraph line when
% \codeline follows ordinary prose. Verified both ways — in a list the label now
% sits on the code's line at its normal x, and after prose the geometry is
% byte-identical to before.
\makeatletter
\newcommand{\codeline}[1]{%
  \if@inlabel\leavevmode\fi
  \par\nopagebreak
  \vspace{\abovedisplayskip}\nopagebreak
  % v3.62: was \centerline{...}, i.e. \hbox to\hsize{\hss #1 \hss}. \hss has
  % INFINITE SHRINK, so TeX could always "fit" the box and NEVER reported an
  % Overfull — the same \hss mechanism as BLD-27's \item label. Measured: a wide
  % \codeline in a defbox painted its first token ON the box's left border bar and
  % sat 9.5pt outside the frame interior on both sides, with 0 Overfull, 0 Underfull
  % and all four gates PASS; in bare prose it reached 21.0pt into BOTH margins.
  % \hfil has no shrink, so the overflow is reported and build.sh's ovf_vis/ovf_big
  % counters see it; \linewidth (not \hsize) makes it box-aware, so a \codeline
  % inside a colored box centres on the box interior. Verified: a FITTING \codeline
  % keeps byte-identical glyph positions, and the box-crossing case now emits
  % "Overfull \hbox (19.03pt too wide)". A wide \codeline should be split in two.
  \hbox to\linewidth{\hfil\code{#1}\hfil}%
  \nopagebreak\vspace{\belowdisplayskip}}
\makeatother

% ---- a thin decorative rule helper ----
\newcommand{\sep}{\par\vspace{2pt}\noindent\textcolor{cRule}{\rule{\linewidth}{0.6pt}}\par\vspace{4pt}}

% ---- colors ----
\definecolor{cBlue}{HTML}{1F4E79}   \definecolor{cBlueBg}{HTML}{EAF1F8}
\definecolor{cGreen}{HTML}{2E6B4F}  \definecolor{cGreenBg}{HTML}{EAF4EE}
\definecolor{cAmber}{HTML}{9C6A00}  \definecolor{cAmberBg}{HTML}{FBF3E2}
\definecolor{cPurple}{HTML}{5B2C6F} \definecolor{cPurpleBg}{HTML}{F3ECF7}
\definecolor{cTeal}{HTML}{0E6E6E}   \definecolor{cTealBg}{HTML}{E6F4F4}
\definecolor{cRed}{HTML}{8E2B2B}    \definecolor{cRedBg}{HTML}{FBEAEA}
\definecolor{cGray}{HTML}{555555}   \definecolor{cRule}{HTML}{1F4E79}

% ---- math operators ----
\DeclareMathOperator{\Tr}{Tr}
\DeclareMathOperator{\rank}{rank}
\DeclareMathOperator{\diag}{diag}
\DeclareMathOperator{\sgn}{sgn}

% ---- bra-ket (quantikz defines \ket/\bra/\braket; add the rest safely) ----
\providecommand{\ket}[1]{\left|#1\right\rangle}
\providecommand{\bra}[1]{\left\langle#1\right|}
\providecommand{\braket}[2]{\left\langle#1\middle|#2\right\rangle}
\newcommand{\dyad}[1]{\left|#1\right\rangle\!\left\langle#1\right|}
% v3.21: general outer product |a><b| and expectation value <A> — agents reach
% for these by reflex (esp. \expval); quantikz defines neither, so their absence
% threw "Undefined control sequence" in a real build. \providecommand = safe.
\providecommand{\ketbra}[2]{\left|#1\right\rangle\!\left\langle#2\right|}
\providecommand{\expval}[1]{\left\langle#1\right\rangle}
\newcommand{\abs}[1]{\left|#1\right|}
\newcommand{\norm}[1]{\left\lVert#1\right\rVert}
\newcommand{\half}{\tfrac12}
\newcommand{\R}{\mathbb{R}}\newcommand{\C}{\mathbb{C}}\newcommand{\Z}{\mathbb{Z}}
% identity operator. NB: \mathbb{1} silently renders a WRONG glyph (msbm has no
% blackboard digits). Use \mathbb{I} = 𝕀 (works with amssymb, no extra package).
% Do NOT "fix" this back to \mathbb{1}.
\newcommand{\id}{\mathbb{I}}

% ---- physical units (v3.35): \punit — upright unit with thin space ---------
% Usage: $300\punit{K}$, $80\punit{mV/decade}$, $10^{7}\punit{cm/s}$ — the
% NUMBER AND UNIT LIVE IN ONE MATH ISLAND. Promoted verbatim from a 199-page
% production lakit (advanced-electronic-devices, dozens of uses).
% Two measured traps this design avoids (both render-proven):
%   1. siunitx SCRAMBLES under RTL — \SI{3.5}{eV} renders "eV 3.5", exponents
%      detach ("Hz 10 x 10^9"): its quantity assembly is not bidi-isolated.
%      Do NOT load siunitx in this pipeline.
%   2. A text-mode unit next to a bare-prose number is TWO adjacent LTR runs
%      and reorders ("5 mm" -> "mm5" — the BIDI-2/3 class). \punit is therefore
%      math-only ON PURPOSE: prose misuse fails LOUDLY at compile (\mathrm
%      outside math) instead of scrambling silently.
% Domain shorthands (\MeV, \fm...) stay per-document: \newcommand{\MeV}{\,\text{\en{MeV}}}
\providecommand{\punit}[1]{\,\mathrm{#1}}

% ============================================================================
% Colored boxes — mdframed (NOT tcolorbox, which breaks under bidi).
% ★★ v3.4: innerleftmargin/innerrightmargin give code listings room so their
% frame/line-numbers do not spill out the box border (see listing styles below).
% ============================================================================
% v3.61: everyline/split skips were applied to `tablebox` only (BLD-14, v3.37) —
% the six colored boxes below never got them, so they still split OPEN across a page
% break: the page-N segment gets top+sides and NO bottom rule, the page-N+1 segment
% sides+bottom and NO top rule. Measured on the shipped fleet this is a COVERAGE GAP,
% not a re-report of BLD-14: 0 split-open `tablebox` segments (that fix holds) vs 62
% split-open COLORED boxes across 7 documents. Verified after the fix: both segments
% close, and the frametitle is NOT duplicated on the continuation. BLD-14's caveat
% still applies — a box taller than a full page grinds mdframed's split-trial loop, so
% split such content manually.
% ============================================================================
% v3.72: frametitle={} leads EVERY \newmdenv option list below. mdframed options
% are GROUP-scoped: a box nested inside \begin{exbox}[frametitle={...}] used to
% INHERIT the outer title — measured, the identical title string rendered twice,
% the second riding a title-bar fill inside the INNER box's frame, for all seven
% environments in both directions. Both natural forms were broken (= preamble
% bug, not an authoring rule). The explicit reset makes every box titleless
% unless ITS OWN [frametitle={...}] says otherwise; verified both directions:
% no duplicate on the nested box, an explicit title still renders.
% Geometry is shared with the optional semantic-width table container. Themes
% select their own existing semantic style below; neither path fixes columns.
\mdfdefinestyle{hllpdfsemanticgeometry}{linewidth=1.1pt,
  innertopmargin=6pt,innerbottommargin=7pt,skipabove=9pt,skipbelow=9pt,
  innerleftmargin=9pt,innerrightmargin=9pt,everyline=true,
  splitbottomskip=6pt,splittopskip=12pt,roundcorner=3pt}
\newcommand{\hllpdfsemanticframestyle}{hllpdfsemanticgeometry}
\newmdenv[frametitle={},style=hllpdfsemanticgeometry,linecolor=cBlue,backgroundcolor=cBlueBg,
  frametitlebackgroundcolor=cBlue,frametitlefontcolor=white,
  ]{defbox}   % הגדרה
\newmdenv[frametitle={},style=hllpdfsemanticgeometry,linecolor=cGreen,backgroundcolor=cGreenBg,
  frametitlebackgroundcolor=cGreen,frametitlefontcolor=white,
  ]{thmbox}   % משפט/טענה
\newmdenv[frametitle={},style=hllpdfsemanticgeometry,linecolor=cAmber,backgroundcolor=cAmberBg,
  frametitlebackgroundcolor=cAmber,frametitlefontcolor=white,
  ]{notebox}  % הערה
\newmdenv[frametitle={},style=hllpdfsemanticgeometry,linecolor=cPurple,backgroundcolor=cPurpleBg,
  frametitlebackgroundcolor=cPurple,frametitlefontcolor=white,
  ]{exbox}    % דוגמה
\newmdenv[frametitle={},style=hllpdfsemanticgeometry,linecolor=cRed,backgroundcolor=cRedBg,
  frametitlebackgroundcolor=cRed,frametitlefontcolor=white,
  ]{warnbox}  % טעות נפוצה
\newmdenv[frametitle={},style=hllpdfsemanticgeometry,linecolor=cTeal,backgroundcolor=cTealBg,
  frametitlebackgroundcolor=cTeal,frametitlefontcolor=white,
  ]{keybox}   % רעיון מרכזי

% ---- codebox (v4.5.0): the dedicated container for code -------------------
% An OPTION, not a policy: asmblock unboxed stays canonical for a full listing,
% and a snippet inside any of the six colored boxes still works as before.
% What makes codebox worth its own environment is that it is the ONE framed box
% that keeps its margin line numbers.
%
% v4.4.0 shipped it WITHOUT them, on the strength of v3.37's "no margin knob can
% save them" plus four fresh probes that agreed. That was wrong: every probe had
% moved the code, the box, or both — never the key that carries the NUMBER
% COLUMN. Measured on the page:
%   resetmargins=false   the listing follows THIS box instead of the page,
%                        moving the code from x=75.26 to x=97.54
%   xrightmargin=26pt    under RTL this is the PHYSICAL LEFT indent, and the
%                        only key that carries the \llap'ed digits with it:
%                        x=64.28 -> 90.18, clearing the rule at 65.64 by 23.6pt
%   xleftmargin=26pt     moves the code and LEAVES THE DIGITS on the rule —
%                        the asymmetry that hid the fix for eleven versions
% resetmargins is hoisted out of the four styles (see the \lstset below them) so
% this override is possible at all; every other listing still resolves to true.
% Pinned by tests/codebox_numbers.tex + probe_lstnum.py --inside-rule, whose
% negative control strips xrightmargin and watches the digits fall back.
\definecolor{cSlate}{HTML}{3F4A56}  \definecolor{cSlateBg}{HTML}{F2F4F7}
\newmdenv[frametitle={},leftline=true,rightline=true,topline=true,bottomline=true,
  linecolor=cSlate,backgroundcolor=cSlateBg,linewidth=0.9pt,
  frametitlebackgroundcolor=cSlate,frametitlefontcolor=white,
  innertopmargin=6pt,innerbottommargin=7pt,skipabove=9pt,skipbelow=9pt,
  innerleftmargin=9pt,innerrightmargin=9pt,everyline=true,splitbottomskip=6pt,splittopskip=12pt,
  roundcorner=3pt]{codebox}   % קוד
% The two keys that make in-box numbering work, both measured on the page:
%   resetmargins=false  the listing follows THIS box instead of the page
%   xrightmargin        under RTL this is the PHYSICAL LEFT indent, and it is
%                       the only one that carries the \llap'ed number column
%                       with it. xleftmargin moves the code and leaves the
%                       digits behind on the rule (measured: code 97.54,
%                       digits 64.28 — the v3.37 defect, reproduced); with
%                       xrightmargin=26pt the digits land at x=90.18, a
%                       clearance of 23.6pt from the rule at 65.64.
% This asymmetry is why v3.4's note had the semantics backwards and v3.37
% concluded the digits could not be saved.
% <hllpdf:listing-container-state>
% v4.29: all state is local to the current environment group. A semantic
% callout masks number text, while a nested codebox restores the formatter
% saved before that mask. Save it once across successive semantic wrappers;
% an explicit numbers=none (or an explicitly empty formatter) remains valid.
\makeatletter
\newif\ifhllpdf@incodebox
\newif\ifhllpdf@englishcodebox
\newif\ifhllpdf@lstnumbermasked
\let\hllpdf@savedlstnumber\thelstnumber
\makeatother
% </hllpdf:listing-container-state>
% <hllpdf:codebox-gutter>
\makeatletter
\AtBeginEnvironment{codebox}{%
  \hllpdf@incodeboxtrue
  % Snapshot the OUTER language before asmblock selects English.
  \iflanguage{english}{\hllpdf@englishcodeboxtrue}{\hllpdf@englishcodeboxfalse}%
  \ifhllpdf@lstnumbermasked
    \let\thelstnumber\hllpdf@savedlstnumber
    \hllpdf@lstnumbermaskedfalse
  \fi
}
% listings' AddToHook is GLOBAL. Register once, and gate on local box state;
% adding an unconditional hook at every codebox entry leaks into later listings.
\lst@AddToHook{PreInit}{%
  \ifhllpdf@incodebox
    % Apply every listing override here, inside the LISTING's own group.
    % Setting margins at box entry leaks them into nested semantic callouts.
    \lstset{resetmargins=false,xrightmargin=26pt}%
    % Styles resolve before PreInit. In LTR the physical number gutter follows
    % xleftmargin; in Hebrew retain the measured inherited xrightmargin path.
    \ifhllpdf@englishcodebox
      \lstset{xleftmargin=26pt,xrightmargin=0pt}%
    \fi
    % ONE box, not two: suppress only the frame/background/skips of a listing
    % inside this codebox. The next unboxed listing keeps its own style.
    \lstset{frame=none,backgroundcolor=\color{cSlateBg},aboveskip=0pt,belowskip=0pt}%
  \fi
}
\makeatother
% </hllpdf:codebox-gutter>

% v3.37: line numbers inside a COLORED BOX are sliced by the box border —
% measured: digit span x=63.5-67.3 vs border line at x=65.2 (the border runs
% THROUGH the digits), and the inter-border gap (3.5pt) is narrower than a
% digit (3.8pt), so no margin knob can save them (under resetmargins,
% xleftmargin is inert in-box — measured: code at x=75.26 with 2.2em AND
% 3.0em). A 3-6 line boxed snippet does not need line numbers; blank the
% number TEXT inside every colored box (\thelstnumber is what gets printed;
% the style's numbers=left survives, so nothing else shifts — code position
% byte-frozen at x=75.26, verified). Unboxed listings keep their numbers.
% v3.78: the blanking is coupled to the FRAME, not to the box name —
% clean_theme.tex registers LATER hooks (registration order wins) that
% restore the number text inside the boxes it renders frameless.
% v4.5.0 RETRACTS the "no margin knob can save them" clause above. It holds only
% under resetmargins=true (the shipped default): there xleftmargin is inert
% in-box, as measured. With resetmargins=false, xrightmargin DOES carry the
% \llap'ed digit column clear of the rule — measured; the codebox \lstset that
% does it is the one above, in the codebox block. The loop below therefore
% covers SIX boxes, not "every colored box": codebox opts out and keeps its
% numbers.
% <hllpdf:semantic-listing-mask>
\makeatletter
\@for\HEB@lstbox:=defbox,thmbox,notebox,exbox,warnbox,keybox\do{%
  \expandafter\AtBeginEnvironment\expandafter{\HEB@lstbox}{%
    \ifhllpdf@lstnumbermasked\else
      \let\hllpdf@savedlstnumber\thelstnumber
    \fi
    \hllpdf@lstnumbermaskedtrue
    \def\thelstnumber{}%
    % The nearest semantic container owns its listing presentation. On exit,
    % TeX's group restores the enclosing codebox state and formatter.
    \hllpdf@incodeboxfalse}}
\makeatother
% </hllpdf:semantic-listing-mask>

% ============================================================================
% Unified TABLE presentation: subtle uniform slate frame + auto-numbered
% caption rendered INSIDE the frame. Use:
%   \begin{tablebox}\centering  <tabular>  \tabcap{<description>}\end{tablebox}
% Replaces the buggy \caption* (which emitted a stray "טבלה N : *" line) and
% provides an explicit framed option + numbered caption "טבלה N.M: תיאור".
% Optional intro text may be placed at the TOP, before the tabular, inside the box.
% ============================================================================
\definecolor{cTblLine}{HTML}{B9C4D2}  \definecolor{cTblBg}{HTML}{FAFBFC}
% v3.23: class-aware — under `article` there is no chapter counter (the promised
% report->article switch used to die with "No counter 'chapter' defined").
\makeatletter
\@ifundefined{c@chapter}{%
  % v3.31: in article the counter resets per SECTION, so an unprefixed number
  % repeated ("טבלה 1." in every section — measured on a two-section render).
  % Same conditional idea as the report branch: prefix the section number once
  % a real \section exists; a pre-section table stays "טבלה 1".
  \newcounter{tbl}[section]
  \renewcommand{\thetbl}{\ifnum\value{section}>0 \thesection.\fi\arabic{tbl}}%
}{%
  \newcounter{tbl}[chapter]
  % v3.30: a report-class doc that uses SECTIONS ONLY (chapter stays 0 —
  % the common short-doc pattern) used to print "טבלה 0.1"; the chapter
  % prefix now appears only once a real \chapter exists (render-verified
  % both ways: "טבלה 1." chapterless, "טבלה 1.1" with chapters).
  \renewcommand{\thetbl}{\ifnum\value{chapter}>0 \thechapter.\fi\arabic{tbl}}%
  % v3.64: the SAME leak in \thesection itself, left behind by v3.30. That fix
  % guarded the CAPTION counter and stopped there, so a chapterless report still
  % printed the heading "0.1 סעיף ראשון" and \ref to it printed "0.2" — measured:
  %   p1 y=65.2  '…0.1'   p1 y=144.9 '…0.2'   p1 y=183.4 '…0.2.1'
  % report.cls hardcodes \thesection as \thechapter.\@arabic\c@section, and the
  % leading 0 is the chapter counter that no \chapter ever incremented. Fixing
  % \thesection fixes the whole chain — \thesubsection and \thesubsubsection are
  % defined in terms of it. \theequation and \thetbl/\thefig were already guarded,
  % which is exactly why this survived: every NUMBER around the heading looked
  % right, so the heading's own 0 read as a style choice rather than a bug.
  % \appendix is safe: report.cls resets chapter to 0 and redefines \thechapter to
  % \@Alph\c@chapter, so an appendix \chapter makes the counter 1 and the guard
  % emits "A.1" unchanged. Render-verified in all four modes (chapterless,
  % chaptered, appendix+chapter, appendix+section-only).
  \renewcommand{\thesection}{\ifnum\value{chapter}>0 \thechapter.\fi\arabic{section}}%
}
\makeatother
% v3.37: everyline=true — if the box DOES split across pages, every segment
% gets a CLOSED frame (measured before: the split ran open-ended into the page
% edge with no bottom border, and the caption could strand alone at the next
% page top — "טבלה A.2" shipped that way). The split skips give the segments
% breathing room. In practice, with the caption GLUE below, every previously-
% split table in the corpus now moves WHOLE instead (the white-gap tradeoff,
% documented benign); everyline stays as the safety net for a genuinely
% page-tall table. NB: a table taller than ONE FULL PAGE can grind mdframed's
% split-trial loop (reproduced: minimal doc, minutes of 1pt-step Overfull
% probes) — break such a monster into two tables.
\mdfdefinestyle{hllpdftablegeometry}{leftline=true,rightline=true,topline=true,bottomline=true,
  linewidth=0.9pt,
  innertopmargin=9pt,innerbottommargin=8pt,skipabove=11pt,skipbelow=11pt,
  innerleftmargin=11pt,innerrightmargin=11pt,roundcorner=4pt}
\newcommand{\hllpdftableframestyle}{hllpdftablegeometry}
\newmdenv[frametitle={},style=hllpdftablegeometry,
  linecolor=cTblLine,backgroundcolor=cTblBg,
  % v4.2.0 (adopted): tablebox is ATOMIC. Short tableboxes near page bottoms
  % repeatedly entered mdframed's split trial, emitting benign Overfull-vbox
  % bursts and occasionally separating the intro prose from the tabular.
  % nobreak=true moves the complete short table to the next page instead. The
  % output-active Overfull-vbox gate still catches a box taller than a page;
  % everyline/split skips remain as defensive styling should that path ever
  % re-arm, but the canonical path never publishes a split tablebox.
  nobreak=true,
  everyline=true,splitbottomskip=6pt,splittopskip=12pt,
  % v3.79 (BLD-37 table side, correction): startinnercode/endinnercode run
  % INSIDE mdframed's body vbox (mdframed.sty:848 / :856+863) -- the second one
  % after \par\unskip, i.e. after line breaking, which is the only moment the
  % \tabcap anchor's true offset from the top of the box EXISTS. See
  % \lakintabboxclose below. These two keys are the ONLY tablebox change.
  startinnercode={\lakintabboxopen},endinnercode={\lakintabboxclose}]{tablebox}
% v3.37: \nopagebreak GLUE — \vsplit honors penalties, so the caption can no
% longer strand alone at the top of the next page (it takes its rows along, or
% the whole box moves). Render-verified on the real shipped orphan.
% v3.59: \NR@gettitle is what nameref injects into the stdlib title-bearing
% commands (\@sect, \@chapter, \caption) to fill \@currentlabelname. \figcap and
% \tabcap are INDEPENDENT macros — they never called it, so \@currentlabelname kept
% whatever the last heading left and \nameref{fig:x} silently printed the enclosing
% SECTION's title instead of the caption (proven in the .aux: field 3 of \newlabel
% was the heading text while field 1 and the anchor `fig.1.1` were correct — so this
% is a THIRD, distinct head next to BLD-13's counter fix and BIDI-38's autoref-name
% fix, neither of which writes that field). \@ifundefined keeps it inert if
% hyperref/nameref is ever dropped. Guarded helper because \figcap is defined
% after \makeatother.
% NB both helpers below need \makeatletter — this part of the preamble is OUTSIDE
% any letter-catcode region, so an unwrapped \@currentlabel parses as LaTeX's \@
% and dies with "You can't use \spacefactor in vertical mode" (measured).
\makeatletter
\newcommand{\lakinrtitle}[1]{\@ifundefined{NR@gettitle}{}{\NR@gettitle{#1}}}
% v3.59: \hitem{<letter>}{<key>} — a REFERENCEABLE Hebrew-lettered list item.
%   \begin{enumerate} \hitem{א}{it:a} הסעיף הראשון. \end{enumerate}  ->  \ref{it:a} = (א)
% Why this exists: automatic Hebrew list labels fail silently, so SKILL.md mandates
% an EXPLICIT \item[(\textbf{א})] — but LaTeX puts \refstepcounter INSIDE
% \if@noitemarg (latex.ltx:12964), and \item[...] never sets that flag. So an
% explicit item steps no counter and never touches \@currentlabel: a \label there
% binds to the AMBIENT one and \ref prints the CHAPTER/SECTION number where the
% item letter belongs (proven in the .aux: {chapter.1}). Same doctrine as \figcap —
% own the label state before \label. \phantomsection matters: without it the
% hyperlink still lands on the chapter head (measured); with it the anchor is
% section*.N and the bookmark outline stays clean.
\newcommand{\hitem}[2]{\item[(\textbf{#1})]%
  \def\@currentlabel{(#1)}\def\@currentlabelname{#1}\phantomsection\label{#2}}
% NB (v3.66): \hitem is \ref-only. \autoref{key} calls it a SECTION ("סעיף (א)")
% because \phantomsection brands the anchor section*.N, and \nameref{key} prints
% the bare letter. For \autoref-able items prefer \hebenumlabels below.
%
% v3.66: \hebenumlabels — OPT-IN automatic Hebrew enumerate letters, done right.
%   {\hebenumlabels\begin{enumerate}\item\label{q:a} ...\end{enumerate}}
%   -> labels (א) (ב) (ג); \ref{q:a} = א; anchor Item.N (real, per-item).
% This retires the belief that made \item[(\textbf{א})] + \hitem necessary.
% "Automatic Hebrew labels fail silently" was TRUE but unexplained, and the
% mechanism is now bisected (plain article vs +enumitem, measured): ENUMITEM
% PRE-EXPANDS the label expression in an \edef BEFORE \refstepcounter, with the
% counter still at 0. Tokens it shields by name (\arabic,\alph,\roman) survive;
% a plain \newcommand \ifcase wrapper evaluates at 0 -> the EMPTY zero-branch ->
% the label prints as "()" with the letter GONE (and \the\c@enumi prints a
% literal 0, \number\value dies fatally pointing at \item). The cure is one
% property, not a doctrine: the wrapper must be ROBUST (survives the \edef
% unexpanded), while \theenumi stays EXPANDABLE (so \refstepcounter writes the
% real letter to the .aux). Verified: labels (א)(ב)(ג), \ref -> א/ב/ג each item
% its own value, .aux = {א}{...}{Item.1} — versus the explicit-\item doctrine
% where BIDI-50/\hitem exist purely to patch the counter that never stepped.
\def\@hebalph#1{\ifcase#1\or א\or ב\or ג\or ד\or ה\or ו\or ז\or ח\or ט\or י\or
  יא\or יב\or יג\or יד\or טו\or טז\or יז\or יח\or יט\or כ\else\@ctrerr\fi}
\DeclareRobustCommand*\hebalph[1]{\expandafter\@hebalph\csname c@#1\endcsname}
\newcommand{\hebenumlabels}{%
  \renewcommand{\theenumi}{\@hebalph\c@enumi}%
  \renewcommand{\labelenumi}{(\textbf{\hebalph{enumi}})}}
\makeatother
% v3.63: \mbox the LABEL. \thetbl ("7.1") and a caption description that STARTS
% WITH A DIGIT ("32 הרגיסטרים…") are two European-Number runs separated only by
% neutrals (the label's own "." and \enspace). Under the RTL paragraph base the UBA
% resolves those neutrals to the RTL level, so both numbers sit in ONE reversed span
% and EXCHANGE VISUAL POSITIONS. Measured on the shipped PDF: a table whose real
% number is 7.1 printed "טבלה 32" — a different, EXISTING table number — while \ref
% correctly printed 7.1, so the caption silently contradicted its own cross-reference,
% and the description lost its leading count. The source contains no adjacency at all;
% it is manufactured by this macro, which is why every source linter was blind.
% \mbox makes the label one atomic box so nothing can reorder across it. Verified
% zero-regression: an ordinary Hebrew caption keeps identical x-geometry.
% One paragraph style covers label, number, separator and description.
% Explicit project/template overrides remain possible; the label hook is kept
% for compatibility, with an identity default. Counters/anchors stay unchanged.
\newcommand{\lakitcaptionstyle}{\normalfont\small\color{black}\hypersetup{allcolors=black}}
\newcommand{\lakitcaptionlabelstyle}[1]{#1}
\newcommand{\lakitcaptionpunctuation}{:}
% Native floats, captionof and longtable use the same presentation policy.
% subcaption is optional and inherits the same font rather than a second size.
\usepackage{caption}
\DeclareCaptionFont{lakit}{\lakitcaptionstyle}
\DeclareCaptionLabelFormat{lakit}{\lakitcaptionlabelstyle{#1~\en{#2}}}
\DeclareCaptionLabelSeparator{lakit}{\lakitcaptionpunctuation\enspace}
\captionsetup{font=lakit,labelfont={},textfont={},labelformat=lakit,
  labelsep=lakit,justification=centering}
\makeatletter
\AtBeginDocument{\@ifpackageloaded{subcaption}{%
  \captionsetup[sub]{font=lakit,labelfont={},textfont={},labelsep=lakit}}{}}
\makeatother
% v3.64: OPT-IN list-of-figures/tables entry via an optional short caption.
% \listoffigures and \listoftables are inherited from report.cls and LOOK available,
% but \figcap/\tabcap step the custom \c@fig/\c@tbl counters and never call
% \addcontentsline — so both printed their heading and NOTHING ELSE. Measured: a doc
% with one \figcap and one \tabcap rendered page 1 as 'רשימת האיורים\n1\n', and the
% .lof held only \@chapter's \addvspace, zero \contentsline entries. An empty list
% under a correct Hebrew heading reads as "this document has no figures".
%
% WHY OPTIONAL, AND NOT AUTOMATIC. Writing #1 unconditionally was implemented first
% and REJECTED: a caption containing \codeline compiles CLEANLY today (verified on
% the pre-v3.64 preamble: 0 errors, PDF produced) but becomes FATAL once it is also
% written to the .lof — \codeline is \hbox to\linewidth{...}, and a full-width box in
% a list entry dies with "\pdfextension endlink cannot be used in vertical mode", no
% output PDF. Gaining a feature nobody uses (0 of 13 collections call \listoffigures)
% at the price of breaking a construct that works is a bad trade. Stock \caption has
% exactly this limitation and exactly this escape hatch, so the convention is familiar.
% \en{}, \code{}, \h{} and inline math in the short text are all safe — verified, they
% survive the write and typeset correctly in the list (RTL geometry confirmed: number
% rightmost at x=509.9, title left of it, dots and folio at the far left).
% With no optional argument the macro is byte-identical to v3.63 — zero risk to every
% existing document, and \listoffigures lists exactly the captions that opted in.
% BLD-37, TABLE side (the \figcap v3.74 idea, adapted): clicking a \ref to a
% table used to land MID-table — hyperref's automatic anchor inside
% \refstepcounter was emitted on the still-open horizontal line whose TAIL is
% the tabular (an hbox whose baseline is the table's vertical CENTER), so the
% jump scrolled the top half of the table above the viewport. The v3.74 entry
% deferred this side claiming "tablebox splits — no single box to measure";
% that deferral dissolves: the measurement target is the TABULAR itself at
% \tabcap ENTRY, while it is still the tail of the open list — the same
% situation \lakin@grabfig handles at an island's end (and the fleet has 0
% split tableboxes). So \lakin@grabtab runs as the FIRST token, before any
% typesetting can push the box away, and the anchor is re-emitted raised above
% the measured tabular. \lakintabraise is a SEPARATE register from
% \lakinfigraise (defined with it, below the LTR-island block): a captionless
% island intentionally leaves the fig register set, and a table must never
% inherit that stale figure height.
%
% WHY the anchor stays AT THE ENTRY SPOT and is NOT deferred into the caption
% paragraph the way \figcap defers it: at \par TeX DROPS a final glue item,
% and the automatic anchor's smash box sitting BEHIND the interword glue kept
% that glue alive — so every centered tabular has always sat half a space off
% true center. Deferring the anchor removed that box and was measured as 982
% differing glyph records: EVERY tabular line shifted 1.30pt (2.61pt for
% pattern 7's two minipages) while captions stayed put. Emitting the raised
% anchor at the old spot — through hyperref's own \Hy@raisedlink, exactly the
% macro the automatic path uses — keeps the line's node structure, and so the
% A/B glyph diff at 0, BY CONSTRUCTION. \Hy@raisedlink reads the raise from
% \HyperRaiseLinkDefault (default \baselineskip = the old mid-table dest); the
% group-local \edef freezes measured-height+42pt into it, and \smash inside
% \Hy@raisedlink hides any raise from the line's extent. \HyCnt@ProvideTheH-
% Counter keeps the dest NAME byte-identical to the automatic one (it defines
% \theHtbl, which \hyper@makecurrent prefers — without it the first measured
% caption was named tbl.1 where the automatic path said tbl.0.1, measured).
% vmode entry, or anything but a box at the tail (no tabular to measure),
% falls back to the plain \refstepcounter — byte-identical to the old path.
\makeatletter
\newcommand{\tabcap}[2][]{\lakin@grabtab
  % v3.79: arm the box-close correction. Set unconditionally and in BOTH
  % branches below's scope, so the plain-\refstepcounter fallback is corrected
  % too; \lakintabboxopen clears it at every \begin{tablebox}, so a \tabcap
  % outside a tablebox can never leave it armed for a later box.
  \global\lakintabcapseentrue
  \ifdim\lakintabraise>0pt
    \H@refstepcounter{tbl}\HyCnt@ProvideTheHCounter{tbl}\hyper@makecurrent{tbl}%
    \begingroup
      % A frameless table targets its measured top. The frame's historical
      % cushion belongs only to tablebox and is corrected at box close.
      \edef\HyperRaiseLinkDefault{\the\lakintabraise}%
      \iflakintabframed
        \edef\HyperRaiseLinkDefault{\the\dimexpr\lakintabraise+42pt\relax}%
      \fi
      \Hy@raisedlink{\hyper@anchorstart{\@currentHref}\hyper@anchorend}%
    \endgroup
    \global\lakintabraise 0pt\relax
  \else
    \refstepcounter{tbl}%
  \fi
  \lakinrtitle{#2}\par\nopagebreak\addvspace{6pt}%
  \def\lakin@short{#1}\ifx\lakin@short\@empty\else
    \addcontentsline{lot}{table}{\protect\numberline{\thetbl}{\protect\ignorespaces #1}}\fi
  {\centering\lakitcaptionstyle\mbox{\lakitcaptionlabelstyle{\tablename~\en{\thetbl}\lakitcaptionpunctuation}}\enspace #2\par}}
\makeatother

% ============================================================================
% Unified FIGURE caption — the \tabcap twin for TikZ/plot figures (v3.34.0).
% A figure (tikzpic) sits inline in the prose; put \figcap BELOW it for a
% numbered, \ref-able caption "איור N.M: תיאור" in the same style as \tabcap:
%   \begin{tikzpic} ... \end{tikzpic}
%   \figcap{\label{fig:foo}description}     % \label FIRST, inside the argument
% WHY a macro, not a hand-written gray line: a manual
%   {\centering\small\textcolor{cGray}{איור: ...}} with a \label inside has NO
% \refstepcounter, so the \label binds to the AMBIENT \@currentlabel (the last
% equation) — \ref then prints an equation number AS the figure number
% (measured on a real page: "איור 48" where 48 is equation 48). \figcap steps a
% real \c@fig counter first, so \label is correct, numbering is consistent, and
% refs are stable. Class-aware like \thetbl (chapter- or section-prefixed).
% ============================================================================
\makeatletter
\@ifundefined{c@chapter}{%
  \newcounter{fig}[section]
  \renewcommand{\thefig}{\ifnum\value{section}>0 \thesection.\fi\arabic{fig}}%
}{%
  \newcounter{fig}[chapter]
  \renewcommand{\thefig}{\ifnum\value{chapter}>0 \thechapter.\fi\arabic{fig}}%
}
\makeatother
\makeatletter
\newcommand{\figcap}[2][]{%
  \ifdim\lakinfigraise>0pt
    % v3.74: step WITHOUT hyperref's automatic anchor, then emit the anchor
    % raised above the measured figure. \H@refstepcounter is hyperref's saved
    % pre-hook \refstepcounter (sets the counter AND \@currentlabel — \label
    % in the caption arg binds exactly as before); \hyper@makecurrent supplies
    % \@currentHref, so the dest NAME is identical to the old automatic one —
    % \ref/\autoref/bookmarks are unchanged, only the dest'S POSITION moves.
    \H@refstepcounter{fig}\hyper@makecurrent{fig}%
    % The anchor is emitted INSIDE the caption paragraph below, as a zero-size
    % raisebox — NOT here in vertical mode, where a \raisebox starts an empty
    % paragraph and adds a blank line (measured: every glyph after the caption
    % shifted +18.8pt). Inside an existing line a [0pt][0pt] box is inert.
    \edef\lakin@anch{\@currentHref}%
    % \edef FREEZES the measured height into the pending macro: the deferred
    % \raisebox is typeset inside the caption paragraph AFTER the register is
    % reset to the 0pt sentinel below — a plain \def read 0 at use time
    % (measured: LAKINRAISE=0.0, dest moved 14pt instead of above the figure).
    \edef\lakin@pend{\noexpand\raisebox{\the\dimexpr\lakinfigraise+42pt\relax}[0pt][0pt]{%
      \noexpand\hyper@anchorstart{\lakin@anch}\noexpand\hyper@anchorend}}%
    \global\lakinfigraise 0pt\relax
  \else
    \refstepcounter{fig}\let\lakin@pend\relax
  \fi
  \lakinrtitle{#2}\par\nopagebreak\addvspace{4pt}%
  \def\lakin@short{#1}\ifx\lakin@short\@empty\else
    \addcontentsline{lof}{figure}{\protect\numberline{\thefig}{\protect\ignorespaces #1}}\fi
  {\centering\lakin@pend\lakitcaptionstyle\mbox{\lakitcaptionlabelstyle{\figurename~\en{\thefig}\lakitcaptionpunctuation}}\enspace #2\par}}
\makeatother

% ============================================================================
% LTR islands.  EVERY tikzpicture / quantikz MUST live inside one, or the whole
% figure's coordinates and Latin labels flip. (Hebrew node text still uses \h{}.)
% ============================================================================
% v3.74 (BLD-37): the \figcap hyperlink DESTINATION used to land at the caption
% line — BELOW the figure — so clicking a figure reference scrolled the figure
% off-screen (measured fleet-wide: 97/97 fig dests in one collection sat 19-54pt
% under their figure's ink; a 347pt figure was fully hidden after the jump).
% hypcap-style fix in two parts: (1) at the island's END, \lakin@grabfig grabs
% the just-typeset centered line with \lastbox, measures it, and puts it back
% with \nointerlineskip — the interline glue above the line was already
% contributed and is NOT re-added, so the page geometry is unchanged to the
% glyph (verified); (2) \figcap then steps the counter WITHOUT hyperref's
% automatic anchor (\H@refstepcounter = the pre-hyperref \refstepcounter) and
% emits the anchor itself, raised by the measured height — the dest lands ABOVE
% the figure top. A captionless island leaves numbering untouched (nothing is
% stepped at the island; the grab flag is cleared at every island begin, so a
% stale height from a captionless figure cannot leak into a later caption).
\newsavebox{\lakinfigbox}
\newdimen\lakinfigraise  % real \dimen: tex.setdimen cannot target a \newlength SKIP register (measured lua error)
% BLD-37 table side: a SEPARATE register for \tabcap. The two must not share —
% a captionless island INTENTIONALLY leaves \lakinfigraise set (the grab runs
% at every island END; only the NEXT island's begin clears it), so a shared
% register would hand a \tabcap after a captionless figure that stale FIGURE
% height. \lakintabraise is written only by \lakin@grabtab at \tabcap entry
% and consumed/reset inside \tabcap itself, so it can never go stale.
\newdimen\lakintabraise
\makeatletter
% v3.74b: measure WITHOUT touching the list. The first draft grabbed the box
% with \lastbox and put it back — and \lastbox DISCARDS the node's baseline
% shift, so a vertically-anchored quantikz circuit rose 2.57pt and everything
% after it sank 5.32pt (measured on the A/B glyph diff). LuaTeX lets us READ
% the tail node's dimensions without removing anything: geometry is identical
% by construction. \lakinfigraise=0pt is the "no figure measured" sentinel.
\def\lakin@grabfig{%
  \ifhmode
    \directlua{
      local n = tex.nest[tex.nest.ptr].tail
      local GLUE, PEN = node.id('glue'), node.id('penalty')
      while n and (n.id == GLUE or n.id == PEN) do n = n.prev end
      local HL, VL = node.id('hlist'), node.id('vlist')
      if n and (n.id == HL or n.id == VL) then
        tex.setdimen('global', 'lakinfigraise',
                     (n.height or 0) + (n.depth or 0))
      else
        tex.setdimen('global', 'lakinfigraise', 0)
      end}%
  \else
    \global\lakinfigraise 0pt\relax
  \fi}
% BLD-37 table side: \lakin@grabfig's twin, NOT a shared helper — the Lua body
% differs in exactly one place, the skip set. Between \end{tabular} and \tabcap
% an otherlanguage wrapper (patterns 3/5 of example_tables) leaves a \write
% whatsit and a dir node ON TOP of the interword glue (measured node dump:
% "whatsit/write <- dir <- glue(2.62) <- hlist"), so the table reader must
% also step over whatsit/dir markers — all zero-size — to reach the tabular.
% \lakin@grabfig deliberately does NOT get the wider skip set: its glue/penalty
% body is the one the v3.74 A/B proof froze, and widening it would newly
% measure (and so newly raise) anchors in island shapes that today take the
% plain path — an unproven behavior change on the FIG side. Anything that is
% not a box behind those markers still hits the 0pt sentinel and \tabcap falls
% back to the plain \refstepcounter.
% SECOND deliberate difference: this stores the box's HEIGHT, not height+depth.
% \figcap's anchor is emitted in the caption paragraph BELOW the figure, so it
% must climb the full extent; \tabcap's anchor sits ON the tabular's own line,
% whose baseline is the vcentered tabular's MIDPOINT — height alone reaches
% the top edge, and adding the depth was measured to overshoot by half the
% table (dests 80-135pt above their frame, INSIDE the preceding table's box).
\def\lakin@grabtab{%
  \ifhmode
    \directlua{
      local n = tex.nest[tex.nest.ptr].tail
      local GLUE, PEN = node.id('glue'), node.id('penalty')
      local WHAT, DIR = node.id('whatsit'), node.id('dir')
      while n and (n.id == GLUE or n.id == PEN
                   or n.id == WHAT or n.id == DIR) do n = n.prev end
      local HL, VL = node.id('hlist'), node.id('vlist')
      if n and (n.id == HL or n.id == VL) then
        tex.setdimen('global', 'lakintabraise', n.height or 0)
      else
        tex.setdimen('global', 'lakintabraise', 0)
      end}%
  \else
    \global\lakintabraise 0pt\relax
  \fi}
\makeatother
% ============================================================================
% v3.79 (BLD-37 table side, CORRECTION): put the anchor at the TOP OF THE BOX,
% instead of guessing the distance to it.
% ============================================================================
% v3.77 raised the \tabcap anchor by (height of the box at the tail of the open
% horizontal list) + 42pt. That is a MEASUREMENT taken BEFORE line breaking, and
% it is short by every term it cannot see. Fleet-measured on the 12 built
% collections (299 tbl.* dests, ownership established through each table's own
% CAPTION, never through "which frame encloses the dest"): 3 dests landed INSIDE
% their own tablebox, below the frame top —
%   computer-architecture       tbl.14.1  -46.4pt   three stacked tabulars
%   computer-architecture       tbl.30.5  -13.5pt   tabular 2.85pt over \linewidth
%   advanced-electronic-devices tbl.B.1   -59.2pt   tabular 2.56pt over \linewidth
% Two missing terms, both reproduced in a minimal document (probe P2/P3):
%   (1) VERTICAL LIST ALREADY CONTRIBUTED inside the box — a second and third
%       tabular, or a heading above the table. The tail read sees only the LAST
%       box, so 77.7pt of stacked tabulars is invisible to the raise.
%   (2) WHICH LINE THE ANCHOR LANDS ON. When the tabular alone overfills
%       \linewidth (459.6121pt in a tablebox), the interword glue after
%       \end{tabular} becomes the chosen breakpoint and \Hy@raisedlink is
%       emitted on the NEXT line — one (depth + lineskip) lower. That term does
%       not EXIST at \tabcap time: line breaking has not happened yet. No
%       measurement taken there can ever recover it. (Proof: probe P3, the same
%       table 2.84912pt too wide, fails at -13.5pt; probe P4, the identical
%       table with \tabcolsep 2pt so it fits, is correct at +32.0pt.)
% So the raise is not computed at \tabcap at all. It is CORRECTED at
% \end{tablebox}, from inside mdframed's own body vbox (endinnercode runs after
% \par\unskip, i.e. after line breaking, while tex.nest[tex.nest.ptr] is still
% that vbox's vertical list). There the anchor's true offset from the top of the
% box content is not estimated but WALKED, and the \Hy@raisedlink shift is set
% so the dest lands 42pt above it — which is 42pt - innertopmargin(9pt) -
% linewidth(0.9pt) = 32.0pt above the FRAME top, the exact position the 295
% already-correct dests have today. That is why this is a correction and not a
% re-placement: the good sites do not move.
%
% WHY THE GLYPH DIFF IS 0 BY CONSTRUCTION — the only thing written is the
% `shift` field of the hlist that \Hy@raisedlink already built, and the guard
% below refuses to write it unless that box is 0x0x0. A zero-size box carries no
% dimension into its parent's packing, the packing already happened, and shift
% is not an input to it. Nothing else in the list is touched: no node is added,
% removed, resized or reordered. Measured, not assumed: assets/example_tables.tex
% and the whole computer-architecture collection A/B at 0 differing pymupdf
% rawdict glyph records.
%
% WHY NOT the alternatives (all prototyped and measured before this one):
%   * accumulate the vertical list at \tabcap time — fixes term (1), CANNOT fix
%     term (2); 2 of the 3 real failures survive. A fix that leaves the majority
%     of the defect standing is the false all-clear this repo exists to avoid.
%   * \zref-savepos two-pass — correct on all shapes, but it makes the dest
%     POSITION aux-derived (DEFECTS BLD-1/2), a single-pass compile silently
%     ships an unraised anchor, and the bare y-difference has no page guard.
%   * emit the anchor at \begin{tablebox} and let \tabcap adopt it — needs the
%     tbl counter stepped at box OPEN to keep the dest NAME byte-identical; a
%     captionless tablebox then either shifts every later table number or leaves
%     an orphan dest that the next table duplicates (a HARD build gate). This
%     correction touches neither the counter nor the name: \tabcap's naming path
%     is v3.77's, unchanged, so resolve_names() is byte-identical (verified).
%   * suppressing the line break (a penalty before the glue) — that MOVES the
%     tabular: the break currently discards the interword glue, so forbidding it
%     re-centers the two overfull tables. Reflow, rejected.
%
% SCOPE. The correction is armed only by \tabcap ITSELF (\lakintabcapseen),
% cleared at every box open, so a captionless tablebox is a no-op and a \tabcap
% hosted OUTSIDE a tablebox (the fleet has exactly one, computer-architecture
% tbl.14.2 in an exbox) keeps v3.77's measured raise. Both fallbacks are
% therefore never worse than today; and because hyperref's AUTOMATIC anchor uses
% the same \Hy@raisedlink shape, the plain-\refstepcounter path (vmode entry, or
% prose between \end{tabular} and \tabcap — probe P6, -51.5pt) is corrected too.
% KNOWN LIMIT — a tablebox that SPLITS across pages. The offset is measured from
% the top of the UNSPLIT body, so a dest in segment 2 gets raised past its own
% segment. MEASURED on a forced split (probe: 40 filler lines + a table): the
% dest went from y=585.0 on page 2 (inside segment 2, at the caption — the old
% BLD-37 symptom) to y=-113.4, i.e. ABOVE the page top. Viewers clamp that to the
% page top, which happens to show segment 2 whole, but an out-of-page coordinate
% is not something to ship on purpose. Hence the \textheight guard in
% lakintab.fix: a body TALLER THAN THE TEXT BLOCK cannot fit on any page, so it
% must split — the correction is skipped and v3.77's behaviour stands unchanged
% (verified: the probe returns to y=585.0). The guard is deliberately EXACT, not
% a page-fitting heuristic: it is sound (a taller-than-page body always splits)
% but PARTIAL (a shorter body that starts low on a page could still split). That
% residue is unreachable in practice — v3.37's everyline + the caption's
% \nopagebreak glue move such boxes WHOLE instead (probe N7: a 22-row table
% pushed down by 430pt moved whole, 0 segments), and the fleet has 0 split
% tableboxes, re-counted after this change. Reimplementing mdframed's own
% page-fitting arithmetic to close the residue would be a heuristic with its own
% failure mode, which is the trade this repo has repeatedly refused.
\newif\iflakintabcapseen
\newif\iflakintabframed
% The Lua lives in ONE \directlua at load time, and is written with NO
% backslash, %, # or ~ : \directlua expands its argument like \edef, so an
% active ~ would expand to \nobreakspace and a % would eat the rest of the line
% (the second is how the first draft of this block died). "not (a == b)" stands
% in for Lua's inequality operator for the same reason.
\directlua{
  lakintab = { }
  lakintab.HL = node.id("hlist")
  lakintab.VL = node.id("vlist")
  lakintab.RU = node.id("rule")
  lakintab.GL = node.id("glue")
  lakintab.KE = node.id("kern")
  lakintab.WH = node.id("whatsit")
  lakintab.TM = node.id("temp")
  lakintab.WD = node.subtype("pdf_dest")
  function lakintab.d(v)
    if v == nil then return 0 end
    if v < -1000000000 then return 0 end
    return v
  end
  function lakintab.scan(list, y, horiz, parent, target)
    local x = list
    while x do
      local i = x.id
      if i == lakintab.HL or i == lakintab.VL then
        if horiz then
          local b = y + x.shift
          if i == lakintab.HL then
            lakintab.scan(x.head, b, true, x, target)
          else
            lakintab.scan(x.head, b - lakintab.d(x.height), false, x, target)
          end
        else
          if i == lakintab.HL then
            lakintab.scan(x.head, y + lakintab.d(x.height), true, x, target)
          else
            lakintab.scan(x.head, y, false, x, target)
          end
          y = y + lakintab.d(x.height) + lakintab.d(x.depth)
        end
      elseif i == lakintab.WH then
        if x.subtype == lakintab.WD then
          local id = x.dest_id
          if type(id) == "string" and string.sub(id, 1, 4) == "tbl." then
            if horiz and parent and parent.width == 0
               and parent.height == 0 and parent.depth == 0 then
              parent.shift = parent.shift + (target - y)
            end
          end
        end
      elseif not horiz then
        if i == lakintab.RU then
          y = y + lakintab.d(x.height) + lakintab.d(x.depth)
        elseif i == lakintab.GL or i == lakintab.KE then
          y = y + lakintab.d(x.width)
        end
      end
      x = x.next
    end
  end
  function lakintab.fix(target, maxh)
    local nst = tex.nest[tex.nest.ptr]
    if nst == nil then return end
    local h = nst.head
    if h == nil then return end
    if h.id == lakintab.TM then h = h.next end
    local tot = 0
    local x = h
    while x do
      local i = x.id
      if i == lakintab.HL or i == lakintab.VL or i == lakintab.RU then
        tot = tot + lakintab.d(x.height) + lakintab.d(x.depth)
      elseif i == lakintab.GL or i == lakintab.KE then
        tot = tot + lakintab.d(x.width)
      end
      x = x.next
    end
    if tot > maxh then return end
    lakintab.scan(h, 0, false, nil, target)
  end
}
% @-FREE NAMES ON PURPOSE: these two are named in the `tablebox` key list at
% \newmdenv above, which is tokenised OUTSIDE any \makeatletter region — a
% \lakin@... there would tokenise as \lakin followed by the letters @tabopen.
\newcommand{\lakintabboxopen}{\lakintabframedtrue\global\lakintabcapseenfalse}
\newcommand{\lakintabboxclose}{%
  \iflakintabcapseen
    \directlua{lakintab.fix(-\number\dimexpr\lakintabcushion\relax,
                            \number\dimexpr\textheight\relax)}%
    \global\lakintabcapseenfalse
  \fi\lakintabframedfalse}
% The cushion above the box's CONTENT top. 42pt puts the dest 32.0pt above the
% frame top (42pt - innertopmargin 9pt - linewidth 0.9pt), i.e. exactly where
% the 295 correct dests already are — the correction moves only the broken ones.
\newcommand{\lakintabcushion}{42pt}
\newenvironment{qcirc}{\par\medskip\global\lakinfigraise 0pt\relax\begin{otherlanguage}{english}\begin{center}}
  {\csname lakin@grabfig\endcsname\end{center}\end{otherlanguage}\par\medskip}
\newenvironment{tikzpic}{\par\medskip\global\lakinfigraise 0pt\relax\begin{otherlanguage}{english}\begin{center}}
  {\csname lakin@grabfig\endcsname\end{center}\end{otherlanguage}\par\medskip}

% ============================================================================
% Code listings.  ★★ v3.4 THE BOX-OVERFLOW FIX:
%   resetmargins=true   → the listing IGNORES the enclosing environment's
%                         indentation and measures from the PAGE. v3.4 stated
%                         this backwards ("measured from the ENCLOSING box"),
%                         and v3.37 then blanked in-box line numbers because
%                         the page-anchored listing put its hanging digits on
%                         top of the box rule. Measured v4.5.0: with the key
%                         FALSE the listing follows the box and the code moves
%                         from x=75.26 to x=97.74, clearing the rule at 72.07.
%                         The key now lives in one \lstset below instead of in
%                         each style, so a box can override it — codebox does.
%   framexleftmargin=0  → the frame (left bar) does not extend left of the code.
%   xleftmargin keeps a small indent so the line numbers have room INSIDE.
% asmblock wraps the listing in an LTR island (code is LTR).
% ============================================================================
\lstdefinestyle{asm}{
  basicstyle=\ttfamily\small,
  backgroundcolor=\color{cBlueBg!55},
  frame=leftline,framerule=2.5pt,rulecolor=\color{cBlue},
  numbers=left,numberstyle=\tiny\color{cGray},numbersep=8pt,
  xleftmargin=2.2em,framexleftmargin=0pt,framexrightmargin=0pt,
  breaklines=true,columns=fullflexible,keepspaces=true,
  showstringspaces=false,aboveskip=8pt,belowskip=8pt,
  % v3.36: without a comment DELIMITER the commentstyle below never fires —
  % listings has no `language` here, so it does not know '#' opens a comment,
  % and every asm listing rendered flat black (measured: a whole-page pixel
  % sweep found ZERO cGreen glyphs across 135 shipped listings). morecomment
  % gives it the delimiter; color-only change, glyph widths untouched.
  morecomment={[l]{\#}},
  commentstyle=\color{cGreen},
}
\lstdefinestyle{cstyle}{
  language=C,basicstyle=\ttfamily\small,
  backgroundcolor=\color{cGreenBg!55},
  frame=leftline,framerule=2.5pt,rulecolor=\color{cGreen},
  numbers=left,numberstyle=\tiny\color{cGray},numbersep=8pt,
  xleftmargin=2.2em,framexleftmargin=0pt,framexrightmargin=0pt,
  breaklines=true,columns=fullflexible,keepspaces=true,
  showstringspaces=false,keywordstyle=\color{cBlue}\bfseries,
  commentstyle=\color{cGray}\itshape,aboveskip=8pt,belowskip=8pt,
}
% v3.36: pystyle/htmlstyle — colored twins of cstyle for the Python/HTML
% listings that used to repeat a ~200-char options blob with NO color
% (keywords rendered bold-black, comments italic-black — measured; C alone
% was colored). Geometry is copied VERBATIM from the shipped blob (frame,
% margins, skips) so swapping a call site to style=pystyle moves NOTHING —
% only keywordstyle/commentstyle/stringstyle are added (bold stays bold,
% italic stays italic, mono metrics identical => zero reflow, proven by
% per-page word-bbox diff). breaklines/breakatwhitespace stay per-site.
% v3.78: numbers=left joins both styles, with the exact numberstyle/numbersep
% keys asm/cstyle carry — fleet-wide consistency (author decision: all four
% styles numbered; "asm/cstyle numbered, py/html not" was an accident of the
% v3.36 geometry freeze, not a design). The numbers hang LEFT of the code box
% (listings sets them with \llap, so no overfull hbox can arise and the code
% x is untouched — measured: code frozen, digits at code_x-11.8pt); inside a
% colored box they are blanked like everyone else's (v3.37 hook), and the
% clean variant restores them in its frameless boxes (v3.78, above).
\lstdefinestyle{pystyle}{
  language=Python,basicstyle=\ttfamily\small,
  columns=fullflexible,keepspaces=true,
  frame=single,rulecolor=\color{cTblLine},
  numbers=left,numberstyle=\tiny\color{cGray},numbersep=8pt,
  xleftmargin=0.5em,framexleftmargin=0pt,framexrightmargin=0pt,
  showstringspaces=false,aboveskip=8pt,belowskip=8pt,
  keywordstyle=\color{cBlue}\bfseries,
  commentstyle=\color{cGray}\itshape,
  stringstyle=\color{cGreen},
}
% v4.5.0: hoisted out of the four styles so an enclosing box can override it.
% Every listing still resolves to true unless something later says otherwise,
% so shipped output is unchanged (verified byte-for-byte on all eight PDFs).
\lstset{resetmargins=true}
\lstdefinestyle{htmlstyle}{
  language=HTML,basicstyle=\ttfamily\small,
  columns=fullflexible,keepspaces=true,
  frame=single,rulecolor=\color{cTblLine},
  numbers=left,numberstyle=\tiny\color{cGray},numbersep=8pt,
  xleftmargin=0.5em,framexleftmargin=0pt,framexrightmargin=0pt,
  showstringspaces=false,aboveskip=8pt,belowskip=8pt,
  keywordstyle=\color{cBlue}\bfseries,
  commentstyle=\color{cGray}\itshape,
  stringstyle=\color{cGreen},
}
\newenvironment{asmblock}{\par\begin{otherlanguage}{english}}{\end{otherlanguage}\par}

% <hllpdf:listing-smart-quotes>
% v4.21.0: four Unicode punctuation characters must enter listings' token
% buffer. Unregistered characters bypass it: a closing quote was printed
% BEFORE the pending word, and an opening quote BEFORE a pending brace.
% Extend the standard character table, then activate only these codepoints
% in the listing's local character-table scope. Do this BEFORE the package's
% delimiter/alsoother/literate hooks so explicit user mappings still win.
% This is not universal Unicode support, nor a conversion to ASCII quotes.
% See listings developer guide 15.3.1 and references/listing-smart-quotes.md.
\makeatletter
\lst@CCPut \lst@ProcessOther {"2018}{"2019}{"201C}{"201D}\z@
\lst@AddToHookAtTop{SelectCharTable}{%
  \lst@ifactivechars
    \catcode"2018=\active \catcode"2019=\active
    \catcode"201C=\active \catcode"201D=\active
  \fi}
\makeatother
% </hllpdf:listing-smart-quotes>

% <hllpdf:listing-technical-symbols>
% v4.30: bounded extension for common technical punctuation in LTR listings.
% Without character-table registration these glyphs bypass the token buffer
% and can move before the preceding word even inside quoted strings.
% Keep the original four-quote block separate for its standing mutation test.
% This is source-faithful printing, not syntax validation or general Unicode.
\makeatletter
\lst@CCPut \lst@ProcessOther
  {"2013}{"2014}{"2026}{"2190}{"2192}{"2194}{"21D2}
  {"2212}{"221E}{"2260}{"2264}{"2265}\z@
\lst@AddToHookAtTop{SelectCharTable}{%
  \lst@ifactivechars
    \catcode"2013=\active \catcode"2014=\active \catcode"2026=\active
    \catcode"2190=\active \catcode"2192=\active \catcode"2194=\active
    \catcode"21D2=\active \catcode"2212=\active \catcode"221E=\active
    \catcode"2260=\active \catcode"2264=\active \catcode"2265=\active
  \fi}
\makeatother
% </hllpdf:listing-technical-symbols>


\usepackage[hidelinks,unicode]{hyperref}   % LAST, after colors. unicode = Hebrew bookmarks.
% v3.62: THE ISOLATION MACROS CORRUPTED EVERY PDF BOOKMARK THEY APPEARED IN.
% \en is defined (line ~155) as
%     \foreignlanguage{english}{\language=\csname l@nohyphenation\endcsname #1}
% and in a heading that body is written verbatim to the .toc and handed to
% hyperref's \pdfstringdef. hyperref strips the two CONTROL SEQUENCES (\language,
% \l@nohyphenation) but the plain `=` CHARACTER token between them is not a control
% sequence, so it survives into the PDF string: \chapter{שער \en{CNOT} והצמדה}
% produced the bookmark "שער =CNOT והצמדה". The printed TOC and the heading on the
% page are CORRECT — only the navigation pane and the viewer window title are wrong,
% which is why it went unnoticed. Measured on the shipped fleet: 70 corrupted
% bookmarks across 5 PDFs, from 177 unprotected \en-in-heading sites. Same corruption
% reached /Title and /Subject via the \hypersetup{pdftitle=...} the docs recommend.
% This hook tells hyperref to use the plain argument for all three isolation macros
% when building a PDF string. It does NOT touch the .toc (the printed TOC is
% byte-unchanged) and it does NOT mask the math signal — a heading with $...$ still
% logs "removing math shift", so BIDI-15 detection is intact.
% BONUS, measured: it also makes \code{} in an UNSTARRED heading COMPILE. BLD-10
% recorded that as needing \texorpdfstring or a starred heading; the real root cause
% is hyperref's \pdfstringdef scan of \code's \ifmmode...\else...\fi body — which is
% exactly why the identically-bodied \texttt survived (hyperref pre-neutralizes
% \texttt in PDF strings) and why \DeclareRobustCommand could never help.
% v3.64: \def\\{ } added. A manual line break in a heading WELDS the two words
% together in the bookmark: \chapter{שורה ראשונה\\שורה שנייה} produced the outline
% entry 'שורה ראשונהשורה שנייה' — "ראשונהשורה" is not a word. hyperref drops the
% \\ and puts nothing in its place, so the break that looks like a space on the
% page is a zero-width join in the navigation pane. The spaced form
% 'מילה אחת \\ מילה שנייה' escapes with a harmless DOUBLE space ('אחת  מילה'),
% which is why this never got noticed. Both are measured on a real outline.
% 0 live instances in the corpus — but SKILL.md and content-style.md both RECOMMEND
% a manual heading break, so this is a primed trap rather than a dormant one, and a
% one-token preamble fix retires it for good instead of asking authors to remember.
% Mapping \\ to a space fixes the WELD, which is the actual defect (a nonexistent
% word in the navigation pane). It does not normalise whitespace: measured after the
% fix, the welded form becomes 'שורה ראשונה שורה שנייה' (correct) and the already-
% spaced form goes from two spaces to three. That is cosmetic — viewers render it as
% a wider gap, and no word is corrupted. A per-heading \texorpdfstring is still the
% way to get an exactly-tuned bookmark string.
\pdfstringdefDisableCommands{\def\en#1{#1}\def\h#1{#1}\def\doclang#1{#1}\def\code#1{#1}\def\codeline#1{#1}\def\\{ }}
% v3.70: RESTORE \url LINE-BREAKING. url.sty freezes its break penalties at LOAD
% time (\mathchardef\UrlBreakPenalty=\binoppenalty), and v3.45's global
% \binoppenalty=10000 (inline-math atomicity) therefore priced EVERY breakpoint
% inside \Url at 10000 = never break. Measured: a long \url in Hebrew prose is
% one unbroken line, ~100pt overfull, invading the margin band the ink gates
% deliberately tolerate. With the two lines below: 0 overfull, the URL breaks
% after '/', the RTL split is CORRECT (fragment order forward), and each line
% gets its own link annotation exactly covering its fragment (measured).
% Confined to \Url's group — inline-math atomicity (TYP-7) is untouched.
\mathchardef\UrlBreakPenalty=200
\mathchardef\UrlBigBreakPenalty=100
\ifx\lakitmainlanguage\hllpdfenglish
  \hypersetup{pdflang=en}
\else
  \hypersetup{pdflang=he}
\fi  % /Lang is the document base locale (not per-span accessibility tagging).
% Set a per-document \hypersetup{pdftitle={...}} too — see SKILL.md.

% v3.52 (round-2 audit): two RTL bibliography/cross-ref defects, both invisible to
% every gate and both affecting the FIRST author to use the feature.
% (a) \@biblabel's default [#1] brackets are bidi-neutral and mirror to ]1[ under
% the RTL bibliography base — every \thebibliography reference list shows ]1[.
% LTR-isolating the label restores [1]. (The render gate's mirror rule fingerprints
% round parens )1(, so square brackets slipped through.)
\makeatletter
\renewcommand*{\@biblabel}[1]{\foreignlanguage{english}{[#1]}}
% v3.54: the SAME neutral-bracket trap on the INLINE side — a multi-key \cite{a,b}
% under the RTL base lays its European-number list right-to-left, so [1, 2] renders
% "[2,1]" and [1, 2, 3] renders "[3,2,1]" (verified by glyph geometry — the reader
% sees the sources in REVERSE order; a note form \cite[p.5]{a} also reorders to
% "[p. 5,1]"). LTR-isolate the whole bracket group (number list AND optional note),
% exactly like \@biblabel above: [1, 2] / [1, 2, 3] / [1, p. 5] all render correct.
% Placed AFTER hyperref so it wins; hyperref hyperlinks each number inside #1 (via
% \@citex), which this wrapper leaves intact. Mirrors the LaTeX kernel \@cite shape.
% v4.9.1: directional isolation must NOT switch the language here. Babel's
% safe-active citation wrapper passes its key list through \bbl@tempa;
% \foreignlanguage reuses that scratch macro before the deferred citation
% loop expands it. The old line therefore cited the literal key "english".
% If such a key existed, even the undefined-reference gate stayed green.
% \babelsublr isolates the bracket group without selecting a locale or
% changing that key list. Preserve native links, key order and optional notes.
\def\@cite#1#2{\babelsublr{[{#1\if@tempswa , #2\fi}]}}
\makeatother

% v4.11.0: a reference inside align's \text inherited HEBdisptag=true,
% so its enclosing hbox mirrored the brackets despite a clean compile.
% Isolate the complete printed reference without selecting a Babel language.
% Also neutralize the display-only normalfont direction reset locally. A
% direction-only patch inverted 12 to 21 inside equation in the measured probe.
% Plain maketag is local to the reference: mathtools' display-side wrapper
% otherwise globally consumes manual_tag while formatting a nested reference.
% Preserve each original command (including the selective auxiliary write),
% current tag forms, native links, and showonlyrefs on/off transitions.
\makeatletter
\NewCommandCopy{\HEB@nativeeqref}{\eqref}
\newcommand{\HEB@eqrefisolate}[1]{%
  \babelsublr{\HEBdisptagfalse\let\bbl@eqnodir\relax\let\maketag@@@\ORIGmaketag@@@#1}}
\protected\def\eqref#1{\HEB@eqrefisolate{\HEB@nativeeqref{#1}}}
\ExplSyntaxOn
\NewCommandCopy{\HEB@nativeMTeqref}{\MT_extended_eqref:n}
\cs_set_protected:Npn \MT_extended_eqref:n #1
  {\HEB@eqrefisolate{\HEB@nativeMTeqref{#1}}}
\ExplSyntaxOff
\makeatother

% v4.11.0: custom caption counters are not part of hyperref's English
% locale, and Hebrew has no upstream hyperref locale in this toolchain.
% Restore the appropriate names at every Babel language switch, including
% return from a nested inline island. Do not localize the destination keys.
\makeatletter
\newcommand{\HEB@refnames}{%
\def\chapterautorefname{פרק}%
\def\sectionautorefname{סעיף}%
\def\subsectionautorefname{תת־סעיף}%
\def\subsubsectionautorefname{תת־תת־סעיף}%
\def\paragraphautorefname{פסקה}%
% v4.31: hyperref uses a distinct counter name at the deepest standard level.
\def\subparagraphautorefname{תת־פסקה}%
\def\equationautorefname{משוואה}%
\def\figureautorefname{איור}%
\def\tableautorefname{טבלה}%
\def\partautorefname{חלק}%
\def\appendixautorefname{נספח}%
\def\footnoteautorefname{הערת שוליים}%
\def\itemautorefname{סעיף}%
\def\Itemautorefname{סעיף}%
\def\figautorefname{איור}%
\def\tblautorefname{טבלה}%
\def\pageautorefname{עמוד}% 
}
\addto\extrashebrew{\HEB@refnames}
\newcommand{\ENG@refnames}{%
  \def\subparagraphautorefname{subparagraph}%
  \def\chapterautorefname{chapter}%
  \def\sectionautorefname{section}%
  \def\subsectionautorefname{subsection}%
  \def\subsubsectionautorefname{subsubsection}%
  \def\paragraphautorefname{paragraph}%
  \def\equationautorefname{Equation}%
  \def\figureautorefname{Figure}%
  \def\tableautorefname{Table}%
  \def\partautorefname{Part}%
  \def\appendixautorefname{Appendix}%
  \def\footnoteautorefname{footnote}%
  \def\itemautorefname{item}%
  \def\figautorefname{Figure}%
  \def\tblautorefname{Table}%
  \def\Itemautorefname{item}%
  \def\pageautorefname{page}%
}
\addto\extrasenglish{\ENG@refnames}
\ifx\lakitmainlanguage\hllpdfenglish\ENG@refnames\else\HEB@refnames\fi
\makeatother

% ============================================================================
% Optional running head (v3.39): \lakithead{<book title>} — OFF by default.
% A document opts in with ONE line after % ============================================================================
% preamble.tex — tested LuaLaTeX preamble for Hebrew + math + figures PDFs (v4.32.0)
% Battle-tested: compiled an 80-page Hebrew quantum-computing summary AND a
% 197-page Hebrew computer-architecture anthology (TikZ diagrams, code listings,
% colored boxes) — both with 0 missing characters.
%
% Usage in your document:
%   \input{preamble.tex}     % (or paste these lines at the top)
%   \begin{document}  ...  \end{document}
%
% Compile with:  lualatex -interaction=nonstopmode doc.tex
%   Repeat until two consecutive PDFs are byte-identical (scripts/build.sh
%   does this) — a TOC document converges only on pass 3+; after two passes
%   the TOC page numbers are STALE (measured; see SKILL.md step 3).
% Font binaries are external dependencies, pinned by assets/fonts/SHA256SUMS.
% scripts/setup_fonts.sh imports a verified local installation or prior skill.
%
% PUBLIC DOCUMENT-SHELL / LAYOUT CONTRACT (introduced v4.3.1; architecture-owned v4.7.1). These hooks may be
% defined BEFORE \input{preamble.tex}; if omitted, the historical canonical
% defaults below are exact.  They deliberately control only the document shell
% and page layout.  RTL/math semantics stay in this preamble, while visual
% presentation themes (for example clean_theme.tex) load AFTER it.
%
%   \def\lakitdocumentclass{article}
%   \def\lakitclassoptions{11pt,a4paper}
%   \def\lakitgeometryoptions{margin=1.8cm,headheight=14pt}
%   \def\lakitlinespread{1.15}
%   \def\lakitparagraphstyle{block} % letters/briefs; indented for continuous prose
%   \def\lakittocdepth{1}
%   \def\lakitsecnumdepth{2}
%   \input{preamble.tex}
%
% Supported class contract: report or article.  Alternate paper sizes must also
% be declared to build.sh through EXPECTED_PAPER; the PDF gate never infers them.
% Do NOT use these hooks as a theme mechanism: theme overlays must not own page
% geometry, class selection, or direction semantics. A declared presentation-policy
% overlay may intentionally change numbering visibility after the preamble.
% ============================================================================
% Runtime architecture guard: generated from scripts/architecture_contract.py.
% The readiness sentinel remains delayed until the END of this file.
% <hllpdf:architecture-guard>
\ifdefined\hllpdfcoreloaded
  \errmessage{hllpdf: preamble.tex loaded more than once}
  \expandafter\endinput
\fi
% </hllpdf:architecture-guard>

\ifdefined\lakitdocumentclass\else\def\lakitdocumentclass{report}\fi
\ifdefined\lakitclassoptions\else\def\lakitclassoptions{12pt,a4paper}\fi
\ifdefined\lakitgeometryoptions\else\def\lakitgeometryoptions{margin=2.3cm,headheight=14pt}\fi
\ifdefined\lakitlinespread\else\def\lakitlinespread{1.3}\fi
\ifdefined\lakittocdepth\else\def\lakittocdepth{1}\fi
\ifdefined\lakitsecnumdepth\else\def\lakitsecnumdepth{2}\fi
% Paragraph organization belongs to the document layout, independently of its
% language and theme. Preserve historical class dimensions unless explicitly
% selecting block paragraphs; never suppress indentation with an everypar hack.
\ifdefined\lakitparagraphstyle\else\def\lakitparagraphstyle{indented}\fi
\edef\lakitparagraphstyle{\lakitparagraphstyle}
\def\hllpdfparagraphblock{block}
\def\hllpdfparagraphindented{indented}
\ifx\lakitparagraphstyle\hllpdfparagraphblock\else
  \ifx\lakitparagraphstyle\hllpdfparagraphindented\else
    \errmessage{hllpdf: paragraph style must be block or indented}
  \fi
\fi
% v4.3.1: the public class hook has a deliberately narrow support contract.
% Before this guard, \def\lakitdocumentclass{book} silently compiled even though
% every heading/TOC regression test targets report/article; an unsupported shell
% could therefore look "supported" until a subtle class-specific defect shipped.
% Fail before \documentclass rather than letting an undocumented combination enter
% the mechanism core. Plain-TeX comparison is used because LaTeX packages are not
% available yet at this point.
\edef\lakitvalidatedclass{\lakitdocumentclass}
\def\lakitreportclass{report}
\def\lakitarticleclass{article}
\ifx\lakitvalidatedclass\lakitreportclass
\else\ifx\lakitvalidatedclass\lakitarticleclass
\else
  \errmessage{hllpdf: unsupported class \lakitdocumentclass; use report or article}
\fi\fi
\let\lakitvalidatedclass\relax
\let\lakitreportclass\relax
\let\lakitarticleclass\relax
% Language intent is independent of document class, theme and content domain.
% The two dominant-language profiles differ in authoring intent, not direction.
% A balanced bilingual document must explicitly choose a navigation/base locale.
% <hllpdf:language-profile>
\providecommand{\lakitlanguageprofile}{hebrew}
\edef\lakitlanguageprofile{\lakitlanguageprofile}
\def\hllpdfhebrew{hebrew}\def\hllpdfenglish{english}
\def\hllpdfhebrewmixed{hebrew-with-english}
\def\hllpdfenglishmixed{english-with-hebrew}
\def\hllpdfbilingual{bilingual}
\let\hllpdfexpectedlanguage\relax
\ifx\lakitlanguageprofile\hllpdfhebrew
  \let\hllpdfexpectedlanguage\hllpdfhebrew
\else\ifx\lakitlanguageprofile\hllpdfhebrewmixed
  \let\hllpdfexpectedlanguage\hllpdfhebrew
\else\ifx\lakitlanguageprofile\hllpdfenglish
  \let\hllpdfexpectedlanguage\hllpdfenglish
\else\ifx\lakitlanguageprofile\hllpdfenglishmixed
  \let\hllpdfexpectedlanguage\hllpdfenglish
\else\ifx\lakitlanguageprofile\hllpdfbilingual
  \ifdefined\lakitmainlanguage\else
    \errmessage{hllpdf: bilingual profile requires explicit lakitmainlanguage}
  \fi
\else
  \errmessage{hllpdf: unsupported language profile \lakitlanguageprofile}
\fi\fi\fi\fi\fi
\ifdefined\lakitmainlanguage
  \edef\lakitmainlanguage{\lakitmainlanguage}
\else
  \let\lakitmainlanguage\hllpdfexpectedlanguage
\fi
\ifx\lakitmainlanguage\hllpdfhebrew\else
  \ifx\lakitmainlanguage\hllpdfenglish\else
    \errmessage{hllpdf: main language must be hebrew or english}
  \fi
\fi
\ifx\hllpdfexpectedlanguage\relax\else
  \ifx\lakitmainlanguage\hllpdfexpectedlanguage\else
    \errmessage{hllpdf: main language conflicts with language profile}
  \fi
\fi
\typeout{HLLLANG|\lakitlanguageprofile|\lakitmainlanguage}
% </hllpdf:language-profile>
\documentclass[\lakitclassoptions]{\lakitdocumentclass}
\usepackage[no-math]{fontspec}        % (1) ★ no-math — THE critical line (see recipe.md)
% layout=lists asks Babel to adapt list parshape when paragraph and page
% directions differ. Without it, minority-language list bodies lose their
% indentation and labels hang outside the text block, despite a clean log.
% The upstream mechanism preserves ordinary same-direction list geometry.
\usepackage[bidi=basic,layout=lists]{babel}
\ifx\lakitmainlanguage\hllpdfenglish
  \babelprovide[import]{hebrew}
  \babelprovide[import,main]{english}
\else
  \babelprovide[import,main]{hebrew}
  \babelprovide[import]{english}
\fi
% v4.22: bind canonical families and fallbacks to the selected installation.
% Lookup by basename or family alone is not proof of the file actually used.
% <hllpdf:font-binding>
\directlua{
  local dir = os.getenv("HEBREW_LATEX_FONT_DIR")
  if not dir or dir == "" then
    local data = os.getenv("XDG_DATA_HOME")
    if not data or data == "" or not (data:sub(1,1) == "/") then
      data = (os.getenv("HOME") or "") .. "/.local/share"
    end
    dir = data .. "/fonts/hebrew-lualatex-pdf"
  end
  if not (dir:sub(1,1) == "/") then dir = lfs.currentdir() .. "/" .. dir end
  if dir:find("[",1,true) or dir:find("]",1,true) or dir:find(string.char(0),1,true)
     or dir:find(string.char(10),1,true) or dir:find(string.char(13),1,true) then
    tex.error("hllpdf: unsupported character in font directory")
  end
  if not (lfs.attributes(dir,"mode") == "directory") then
    tex.error("hllpdf: font directory is missing")
  end
  hllpdf_font_dir = dir:gsub("/+$", "") .. "/"
  tex.sprint(string.char(92) .. "def" .. string.char(92) .. "hllpdffontdir{")
  tex.sprint(-2,hllpdf_font_dir)
  tex.sprint("}")
  luaotfload.add_fallback("hebfb",{"[" .. hllpdf_font_dir .. "DejaVuSerif.ttf]:mode=node;"})
  luaotfload.add_fallback("hebfbbf",{"[" .. hllpdf_font_dir .. "DejaVuSerif-Bold.ttf]:mode=node;"})
}
% </hllpdf:font-binding>
% v3.22: Frank Ruhl Libre ships NO italic face. Instead of silently-upright
% \emph (v<=3.20) or bold-as-emphasis (v3.21), synthesize a slanted face with
% FakeSlant — real, VISIBLE italic emphasis with correct \emph nesting
% semantics (inner \emph toggles back upright). Verified empirically at
% 300dpi: slant clearly visible, 0 errors, 0 missing glyphs.
\babelfont{rm}[Path={\hllpdffontdir},RawFeature={fallback=hebfb},
  BoldFont={FrankRuhlLibre-Bold.ttf},
  BoldFeatures={RawFeature={fallback=hebfbbf}},   % v3.56 — see hebfbbf above
  ItalicFont={FrankRuhlLibre-Regular.ttf}, ItalicFeatures={FakeSlant=0.22},
  % Bold-italic uses the weight-matched fallback and a synthetic slant.
  BoldItalicFont={FrankRuhlLibre-Bold.ttf}, BoldItalicFeatures={RawFeature={fallback=hebfbbf},FakeSlant=0.22}
  ]{FrankRuhlLibre-Regular.ttf}
% Sans fallbacks preserve family and weight for the few missing codepoints.
\directlua{luaotfload.add_fallback("hebfbsans",{"[" .. hllpdf_font_dir .. "DejaVuSans.ttf]:mode=node;"})}
\directlua{luaotfload.add_fallback("hebfbsansbf",{"[" .. hllpdf_font_dir .. "DejaVuSans-Bold.ttf]:mode=node;"})}
\babelfont{sf}[Path={\hllpdffontdir},RawFeature={fallback=hebfbsans},
  BoldFont={Heebo-Bold.ttf},
  BoldFeatures={RawFeature={fallback=hebfbsansbf}}]{Heebo-Regular.ttf}
% v3.24: Scale=MatchLowercase — DejaVu Sans Mono has an outsized x-height, so
% unscaled inline code visually SHOUTS next to the delicate Frank Ruhl body
% (measured on a 197-page build: scaled, a code atom sits in the line at the
% body's weight, and the document tightened by ~5%). Applies to listings too.
\babelfont{tt}[Path={\hllpdffontdir},RawFeature={fallback=hebfb},Scale=MatchLowercase,
  BoldFont={DejaVuSansMono.ttf},
  BoldFeatures={RawFeature={fallback=hebfb},FakeBold=1.5},
  ItalicFont={DejaVuSansMono.ttf},
  ItalicFeatures={RawFeature={fallback=hebfb},FakeSlant=0.2},
  BoldItalicFont={DejaVuSansMono.ttf},
  BoldItalicFeatures={RawFeature={fallback=hebfb},FakeBold=1.5,FakeSlant=0.2}
  ]{DejaVuSansMono.ttf}

% Explicit multilingual/English profiles use real Latin faces and italics.
% The legacy/default Hebrew profile retains its established measured typography.
% Latin Modern is provided by the TeX distribution, not the pinned font bundle;
% recorder/runtime evidence reports it as an additional toolchain resource.
\ifx\lakitlanguageprofile\hllpdfhebrew\else
  \babelfont[english]{rm}[Ligatures=TeX,RawFeature={fallback=hebfb},
    BoldFeatures={RawFeature={fallback=hebfbbf}},
    BoldItalicFeatures={RawFeature={fallback=hebfbbf}}]{Latin Modern Roman}
  \babelfont[english]{sf}[Ligatures=TeX,RawFeature={fallback=hebfbsans},
    BoldFeatures={RawFeature={fallback=hebfbsansbf}},
    BoldItalicFeatures={RawFeature={fallback=hebfbsansbf}}]{Latin Modern Sans}
\fi

\usepackage{amsmath,amssymb,mathtools}
% amsthm supplies theorem/proof machinery. Environment names remain
% document-specific; the generic proof label is localized below.
\usepackage{amsthm}
% Proof names follow the surrounding prose, including return from an island.
\addto\extrashebrew{\renewcommand{\proofname}{הוכחה}}
\addto\extrasenglish{\renewcommand{\proofname}{Proof}}
\ifx\lakitmainlanguage\hllpdfenglish
  \renewcommand{\proofname}{Proof}
\else
  \renewcommand{\proofname}{הוכחה}
\fi
% v3.24.1: equation tags under babel-bidi. A tag re-emitted by mathtools'
% showonlyrefs inside align inherits the RTL base and its parens MIRROR —
% ")2.1(" instead of "(2.1)" (plain align/equation tags are unaffected;
% verified on a minimal case + a 198-page build). LTR-isolate at the tag-BOX
% level: mathtools' show/hide + numbering logic is untouched. Do NOT override
% \tagform@ instead — that clobbers the showonlyrefs mechanism (measured:
% every tag shows, numbering breaks).
% v3.28.1: the v3.24.1 wrap must apply ONLY to tags of display environments.
% Inline \eqref in Hebrew PROSE also goes through \maketag@@@, and there the
% english-wrap MIRRORS the parens — ")1(" — while the unwrapped original is
% correct (proven on a 2x2 isolation matrix: patch x showonlyrefs; regressed
% every prose \eqref since v3.24.1, caught by the render gate's R4). The two
% call sites are indistinguishable by \ifmmode (tag boxes are text mode), so
% the display environments themselves raise a flag; \eqref in prose sees it
% false and gets the original untouched \maketag@@@.
\usepackage{etoolbox}

% babel bidi=basic (including upstream 26.9) assumes every automatic
% discretionary has a one-glyph replacement. LuaTeX can legally produce one
% with replace=nil after a hyphenated Latin value; Babel then indexes nil and the
% whole build dies. Do the narrow upstream-compatible correction: while Babel
% classifies the list, present only that malformed disc as an ordinary
% discretionary, then restore its subtype. Babel still reorders the ENTIRE box;
% there is no pcall and no successful fallback to an unreordered node list.
% If Babel throws for any other reason, the build remains fatal. The synthetic
% nil-disc probe and the real table gallery in test_all.sh pin both layers.
\AtBeginDocument{\directlua{
  if Babel and Babel.bidi and not Babel.heb_nil_disc_compat then
    Babel.heb_nil_disc_compat = true
    local original = Babel.bidi
    local disc_id = node.id('disc')
    Babel.bidi = function(head, ispar, hdir)
      local hidden = {}
      for item in node.traverse_id(disc_id, head) do
        if item.subtype == 2 and (not item.replace or not item.replace.char) then
          table.insert(hidden, item)
          item.subtype = 0
        end
      end
      local result = original(head, ispar, hdir)
      for _, item in ipairs(hidden) do item.subtype = 2 end
      return result
    end
  end
}}
\makeatletter
\newif\ifHEBdisptag
\let\ORIGmaketag@@@\maketag@@@
\def\maketag@@@#1{\ifHEBdisptag\ORIGmaketag@@@{\foreignlanguage{english}{#1}}\else\ORIGmaketag@@@{#1}\fi}
\@for\HEB@env:=equation,equation*,align,align*,gather,gather*,multline,multline*,flalign,flalign*,alignat,alignat*,eqnarray,eqnarray*\do{%
  \expandafter\AtBeginEnvironment\expandafter{\HEB@env}{\HEBdisptagtrue}}
\makeatother
% English page + Hebrew prose: Babel's equation-only tag flip uses a primitive
% leqno path that shifts the formula to the left edge and the tag into the body
% of an English page. Align's full-width alignment is unaffected. Give only this
% display a physical LTR container and native amsmath eqno/leqno. The engine then
% owns real tag width, collision handling and labels; no proxy tag box hides ink
% from its fit calculation. Babel's saved \bbl@eqnodir still restores Hebrew text
% inside math, and the prose locale/math direction are never switched.
% Close the prior RTL paragraph before normalization. After the environment's
% group restores its directions, close TeX's empty resumed display paragraph and
% use native LaTeX continuation handling: same-paragraph prose has no indent,
% while an authored blank line starts an ordinary indented RTL paragraph.
% Register after Babel's setup; remain local to equation(*) in a main-English
% document's actual RTL prose. Native Hebrew, ordinary English and explicit
% fleqn retain their original paths. No page/display dimensions are assigned.
% <hllpdf:minority-rtl-equation-geometry>
\makeatletter
\newcommand{\HEB@minorityprosecontinuation}{\par\@doendpe}
\newcommand{\HEB@minorityequationgeometry}{%
  \ifx\lakitmainlanguage\hllpdfenglish
    \ifnum\pardirection=\@ne
      \if@fleqn\else
        \par\aftergroup\HEB@minorityprosecontinuation
        \iftagsleft@\let\veqno\@@leqno
        \else\let\veqno\@@eqno\fi
        \bodydir TLT\pardir TLT\textdir TLT\relax
      \fi
    \fi
  \fi}
\AtBeginDocument{%
  \AddToHook{env/equation/begin}{\HEB@minorityequationgeometry}%
  \AddToHook{env/equation*/begin}{\HEB@minorityequationgeometry}}
\makeatother
% </hllpdf:minority-rtl-equation-geometry>
% v3.63: \shortintertext rendered Hebrew MIRROR-REVERSED and flush LEFT — an entire
% commentary line inside an align printed backwards. babel's bidi=basic patches
% \intertext@ only (luababel.def), and mathtools' \shortintertext uses the identical
% \vbox{\normalbaselines … \noindent\ignorespaces#1\par} skeleton but is never
% patched, so its paragraph keeps pardir/textdir LTR: no bidi reordering (logical
% order on the page) and flush-left. Measured against \intertext with the same
% Applying babel's own replacement to \shortintertext@ makes the two identical.
% \h{}/\mbox fix the letter order but leave the line flush LEFT, so they are not
% substitutes. Guarded + \AtBeginDocument because mathtools \let-copies at load time.
% NB \ExplSyntaxOn is REQUIRED: mathtools' internals carry expl3 names, and without
% the expl3 catcodes the `_` is read as a subscript and the patch detonates far away
% with "! Illegal unit of measure (pt inserted)" at the \end{align} (measured).
\makeatletter
\AtBeginDocument{%
  \@ifundefined{shortintertext@}{}{%
    \ExplSyntaxOn
    \bbl@sreplace\shortintertext@{\normalbaselines}%
      {\normalbaselines\ifx\bbl@eqnodir\relax\else\bbl@pardir\@ne\bbl@eqnodir\fi}%
    \ExplSyntaxOff}}
\makeatother
% geometry inherits the canonical A4 sheet from \documentclass; build.sh then
% verifies every rendered MediaBox and CropBox, so a silent Letter fallback is
% impossible to publish as a clean default build.
\usepackage[\lakitgeometryoptions]{geometry}
\linespread{\lakitlinespread}       % layout hook; canonical default remains 1.3
\ifx\lakitparagraphstyle\hllpdfparagraphblock
  \setlength{\parindent}{0pt}
  \setlength{\parskip}{6pt plus 1pt minus 1pt}
\fi
% These dimensions remain ordinary scoped TeX lengths: authors can customize
% them after core input. Language switches inherit them; lists keep their own
% indentation and narrow reference layouts establish their local paragraph gaps.
\typeout{HLLPAR|\lakitparagraphstyle}
\usepackage{xcolor}
\usepackage{mdframed}
\usepackage{listings}                 % code listings (assembly / C / pseudo-code)
\usepackage{tikz}
\usetikzlibrary{arrows.meta,positioning,calc,shapes.geometric,fit,backgrounds}
\usepackage{pgfplots}\pgfplotsset{compat=1.18}
\usepackage{quantikz}                 % quantum circuits (remove if unused)
\usepackage{booktabs}
\usepackage{array}
\usepackage{float}                    % [H] = place table exactly here (reference doc: no drifting)
\usepackage{enumitem}
\usepackage{needspace}

% ============================================================================
% ★★ v3.4 — THE DIRECTION-ISOLATION TRIO (the heart of mixing Hebrew + Latin)
% Mental model:
%   \h{...}    Hebrew (RTL) island      — Hebrew text inside an LTR context
%                                          (TikZ nodes, the english `tikzpic`).
%   \en{...}   English (LTR) isolate    — Latin text / parentheses inside a
%                                          Hebrew (RTL) context. KEEPS PARENS
%                                          AND ORDER CORRECT.
%   \code{...} \en + monospace          — code inside Hebrew / inside a frametitle.
% CRITICAL: never wrap Latin in \h{} (it is for HEBREW). Inside a TikZ node,
% Latin in \h{} renders REVERSED (Datapath->htapataD, CPU->UPC) and parens
% mirror. Isolate every Latin/paren chunk with \en{} or \code{}. See
% references/figures-boxes-listings.md.
% ---- TEXT INSIDE MATH (v3.8) ----------------------------------------------
% Put ANY text that sits inside math in a \text{}. The ONLY thing that is
% invalid in math is a bare \h/\en (= \foreignlanguage) used DIRECTLY, outside
% \text -> "\rmfamily invalid in math mode". Inside \text:
%   Hebrew label           -> \text{...}            (bare Hebrew renders fine;
%                              \h is optional — BUT NOT ALWAYS HARMLESS (v3.67):
%                              if the \h group OPENS with a nested island
%                              ($..$, \en{}, \code{}, \texttt{}, \mbox{}) the
%                              WHOLE group loses bidi and the Hebrew prints in
%                              LOGICAL order = gibberish. Measured:
%                                \text{\h{$f$ קבועה}}  -> reads "העובק"  BROKEN
%                                \text{\h{קבועה $f$}}  -> correct
%                                \text{$f$ קבועה}      -> correct (no \h)
%                              Hebrew first inside \h, or drop the \h. Flagged
%                              H-OPENS-WITH-LTR-GROUP.)
%   plain Latin word        -> \text{...}            (no marker needed)
%   Latin WITH ( ) or [ ]   -> \text{\en{...}}       (\en forces LTR; a plain
%                              \text mirrors the parens in RTL -> "Spec (TNR(" )
%   f(Hebrew/text arg)      -> \text{f}(\text{...})  -- keep the parens at MATH
%                              level, NOT \text{f(arg)} (parens inside the box mirror)
%   code in math            -> \mathtt{...}  (or \texttt{}, math-guarded below)
% NEVER write \en{...} or \h{...} straight into $...$ / \[..\] (outside \text);
% always via \text{}. (\mbox{...} also works but does not scale in scripts.)
% ============================================================================
\newcommand{\h}[1]{\foreignlanguage{hebrew}{#1}}
% v3.42.1 — \en islands do NOT hyphenate (auto OR at explicit hyphens). English hyphenation patterns were
% tearing identifiers and terms ACROSS THE RTL SEAM ("IO Regis-/ters",
% "In-/put", "Band-/width", "Per-/formance", "Multi-/ple" — all shipped on
% real pages; probe-proven: "ALUS-/rcB", "Mem-/ToReg"). An island is a label
% or a term, not English prose — a hyphen tear there is never right. A/B on
% a full 197pp lakit: identical overfull sets, same page count, ~11
% paragraphs healed, one mild underfull (badness 3525) that replaced a
% mutilated "Regis-/ters". Long English PROSE belongs in the englishprose
% environment below, which restores explicit-hyphen breaks while its paragraph
% is still active. A raw otherlanguage block inherits the global prohibition.
% (\l@nohyphenation needs \csname:
% '@' is not a letter here.)
% v3.43.2 — \exhyphenpenalty moved OUT of \en to a GLOBAL assignment below.
% Render-proven: a group-local \exhyphenpenalty is DEAD CODE — TeX reads the
% line-breaking parameters at \par time, after the \en group has closed, so
% \en{write-through} / plain not-taken still tore at the "-" ("write-/
% through" shipped on a real exam page). \language is different: it is
% recorded in the node list as text is added, so the in-group nohyphenation
% assignment above DOES work (and stays). The global penalty was A/B-measured
% on a 197pp lakit: identical overfull set (25), identical page count, the
% not-/taken tear healed. Hebrew is unaffected in practice — Hebrew compounds
% use maqaf (־), not ASCII "-".
\newcommand{\en}[1]{\foreignlanguage{english}{\language=\csname l@nohyphenation\endcsname #1}}
% Babel runs extras locally for inline islands and paragraph language changes.
% An English main document must not inherit the Hebrew no-hyphen paragraph rule.
\addto\extrashebrew{\exhyphenpenalty=10000\relax}
\addto\extrasenglish{\exhyphenpenalty=50\relax}
\ifx\lakitmainlanguage\hllpdfenglish
  \exhyphenpenalty=50\relax
\else
  \exhyphenpenalty=10000\relax
\fi
% Full English paragraphs need ordinary explicit-hyphen breaks. Use this
% environment instead of a raw otherlanguage block; the paragraph ends while
% the local penalty is still active.
\newenvironment{englishprose}
  {\par\begin{otherlanguage}{english}\exhyphenpenalty=50\relax}
  {\par\end{otherlanguage}}
\newenvironment{hebrewprose}
  {\par\begin{otherlanguage}{hebrew}\exhyphenpenalty=10000\relax}
  {\par\end{otherlanguage}}
% Use for shell text whose locale must not depend on the last body paragraph.
\DeclareRobustCommand{\doclang}[1]{\foreignlanguage{\lakitmainlanguage}{#1}}

% ============================================================================
% ★★ — fixes for defects that are INVISIBLE in the log and in the checks
% and only surface on a VISUAL RENDER of the PDF. (Each verified empirically.)
% ----------------------------------------------------------------------------
% v3.22: \emph is now REAL slanted emphasis (FakeSlant italic face defined on
% the \babelfont{rm} line above), replacing the v3.21 \emph->\textbf hack: it
% preserves \emph's toggle semantics and keeps bold free for its own role.
% Prefer bold-as-emphasis for your doc? one line: \renewcommand{\emph}[1]{\textbf{#1}}
% page number: force the digits LTR so they never reverse (18 -> 81) in the RTL
%   footer. That reversal is a NON-DETERMINISTIC bidi glitch (observed once in
%   ~80 pp); this LuaLaTeX build is not byte-reproducible, so the defensive
%   LTR wrap is the only reliable guard. Latin digits are always LTR anyway.
\renewcommand{\thepage}{\foreignlanguage{english}{\arabic{page}}}
% v3.70: \pagenumbering \gdef's \thepage and silently DESTROYS the wrap above —
% after \pagenumbering{arabic} the folio guard is gone for the whole rest of the
% document. Two collections discovered this independently and hand-restore the
% wrap in main.tex with the same comment; that knowledge now lives here instead.
% The wrapper re-applies the LTR wrap after the kernel call, for any style.
\makeatletter
\let\lakin@orig@pagenumbering\pagenumbering
\renewcommand{\pagenumbering}[1]{%
  \lakin@orig@pagenumbering{#1}%
  \renewcommand{\thepage}{\foreignlanguage{english}{\csname @#1\endcsname\c@page}}}
\makeatother
% emergencystretch: Hebrew does not hyphenate, so an unbreakable LTR island
%   (\en{...} / an inline math run) landing at the margin can leave a small
%   (~2-13pt) overfull with no legal breakpoint. This rescues ONLY those
%   inline-prose tie-breaks; it is a ceiling used only on offending lines and
%   does NOT touch display-math / listing / table overflow (those still surface
%   in the log for the content-level overflow pass). Verified: baseline overfull
%   -> 0 with this set, with no change to pagination or underfull.
\setlength{\emergencystretch}{3em}
% v3.40: inline math (almost) never breaks across lines. TeX's default
%   (binop 700 / rel 500) lets an inline run split at an operator —
%   $(1.f)\cdot 2^{e}$ shipped as "(1.f)·" ending one line and "2^e" opening
%   the next (reader-reported). In LTR that split is merely ugly; under an
%   RTL base the carried operator lands on the visually WRONG side of the
%   break and the formula reads backwards.
%   WHY 10000 (v3.45.0 — the display-first resolution). Inline math is
%   FORBIDDEN from breaking, so a torn island under an RTL base is
%   structurally impossible. The v3.40 cycle used 9999 ("break only as a last
%   resort") because a bare 10000 left ~30-39 WIDE inline equations
%   overflowing (worst 90.7pt) — but 9999's "last resort" fires OFTEN in
%   tight, non-hyphenating Hebrew: a reader found torn islands across the
%   fleet (~150 candidates; e.g. "P(y |" ending one line and "x_0)" opening
%   the next, read backwards). The step the 9999 cycle missed: every one of
%   those ~30 overflows is a wide equation that BELONGS in display. Convert
%   them to display FIRST, then 10000 gives 0 tears AND 0 overflows (measured,
%   12/12 fleet). Doctrine: inline math is ATOMIC; a wide equation goes in a
%   display, never inline. See content-style ("inline math is atomic").
\binoppenalty=10000 \relpenalty=10000
% v3.38: uniform sentence spacing. TeX's default (\nonfrenchspacing, sfcode
% 3000) is the OLD Anglo-American wide-space-after-period convention — Hebrew
% typography never had it, and modern Latin practice (Chicago et al.) dropped
% it too. Measured on real pages: post-period gaps up to 1.9-2.1x the line's
% mean word gap — the dominant "rivers" source in justified Hebrew prose.
% Survives babel language switches (probed: sfcode stays french inside \en{}
% islands and otherlanguage blocks).
\frenchspacing
% v3.38: widow/orphan protection. LaTeX's default penalties (150/150/50)
% allow a lone line at a page top/bottom almost freely. 10000 = forbid;
% this is SAFE here precisely because the document is \raggedbottom (the
% demoted line just moves on, no vertical stretch fight) — under flushbottom
% these values would cause underfull pages. raggedbottom is the report-class
% default; made explicit so a future class/option change cannot silently
% flip it and turn the 10000s harmful.
\raggedbottom
\widowpenalty=10000 \clubpenalty=10000 \displaywidowpenalty=10000
% em-dash: --- does NOT become — under Frank Ruhl (the font lacks tlig; verified
%   the width is identical with/without Ligatures=TeX, while the mechanism works
%   on Latin Modern). TYPE THE CHARACTER — (U+2014) DIRECTLY; likewise – (U+2013).
%   check_bidi_figures.py flags a --/--- that is adjacent to Hebrew.
% ============================================================================

% ---- STRUCTURAL DIRECTION-SAFETY for \texttt (v3.7) -------------------------
% A parenthesis, bracket, or $ INSIDE \texttt{} mirrors/migrates in RTL prose
% (MIPS 0($s0)->"(s0$)0", $rs->"rs$", arr[i]->"[arr[i") because \texttt sets the
% font but NOT the direction. Redefining \texttt to LTR-isolate its content fixes
% the WHOLE class at once (parens, brackets, $) so plain \texttt{lw $t0,0($s0)}
% is safe — no need to remember \code{}. Guarded by \ifmmode because
% \foreignlanguage is invalid in math (there, plain monospace is already LTR).
% NB (corrected v3.56 — the old note here said "stays tofu", which is WRONG):
% \texttt{Hebrew} is a FATAL build failure. The tt family is DejaVu Sans Mono
% (zero Hebrew) and the tt fallback hebfb is DejaVuSerif (also zero Hebrew), so
% the subsetter gets an empty glyph set and LuaTeX dies with
%   "! error: (file .../DejaVuSansMono.ttf) (type 2): there are no glyphs in the
%    subset"  -> no output PDF.
% The message names a FONT FILE, never Hebrew/\texttt/the source line. And it is
% non-local: add an unrelated Latin \code{} elsewhere and the fatal degrades into
% a SILENT DELETION of the Hebrew word. Put Hebrew in \h{}/normal text, never in
% \texttt{}/\code{}; flagged HEBREW-IN-MONO. To disable, delete this block.
\let\origtexttt\texttt
\protected\long\def\texttt#1{\ifmmode\origtexttt{#1}\else\foreignlanguage{english}{\origtexttt{#1}}\fi}
% v4.8.0: defer the direction test until TeX has inserted an alignment cell's
% math template. Unprotected expansion here tested \ifmmode too early at the
% FIRST token of aligned/array/cases/matrix cells. Engine-level \protected also
% preserves the original \texttt implementation without the self-recursion
% observed when retrofitting \DeclareRobustCommand to its saved alias.
\protected\long\def\code#1{\ifmmode\origtexttt{#1}\else\foreignlanguage{english}{\origtexttt{#1}}\fi}

% ---- \codeline — display-code line (v3.24): the code analogue of display math.
% Reader-friction rule (content-style.md §5b): inline \code{} is for a SINGLE
% ATOM only (one register/variable/mnemonic, a tight no-space expression).
% A multi-token expression — an instruction with operands, a spaced assignment,
% a register-binding list — zigzags the eye RTL<->LTR mid-sentence and, being
% an unbreakable island, line-breaks badly. Give it its own centered line with
% display-math vertical rhythm. \centerline = one box => the expression never
% wraps; \nopagebreak keeps it glued to its introducing sentence.
%   ...the Hebrew sentence ends with a colon:
%   \codeline{lw \$t1, -4(\$t0)}
%   and the prose continues (escape dollars as \$).
% v3.66 (first line of the body): flush a PENDING \item label before the \par.
% LaTeX defers the label box (\@labels) until the next \everypar, and an \hbox in
% vertical mode never fires \everypar — so `\item \codeline{...}` printed the code
% at y=69.1 and the item's own "1." BELOW it at y=103.6 (the number appears after
% its content, attached to nothing). \if@inlabel\leavevmode\fi is CONDITIONAL on
% purpose: an unconditional \leavevmode would open an empty paragraph line when
% \codeline follows ordinary prose. Verified both ways — in a list the label now
% sits on the code's line at its normal x, and after prose the geometry is
% byte-identical to before.
\makeatletter
\newcommand{\codeline}[1]{%
  \if@inlabel\leavevmode\fi
  \par\nopagebreak
  \vspace{\abovedisplayskip}\nopagebreak
  % v3.62: was \centerline{...}, i.e. \hbox to\hsize{\hss #1 \hss}. \hss has
  % INFINITE SHRINK, so TeX could always "fit" the box and NEVER reported an
  % Overfull — the same \hss mechanism as BLD-27's \item label. Measured: a wide
  % \codeline in a defbox painted its first token ON the box's left border bar and
  % sat 9.5pt outside the frame interior on both sides, with 0 Overfull, 0 Underfull
  % and all four gates PASS; in bare prose it reached 21.0pt into BOTH margins.
  % \hfil has no shrink, so the overflow is reported and build.sh's ovf_vis/ovf_big
  % counters see it; \linewidth (not \hsize) makes it box-aware, so a \codeline
  % inside a colored box centres on the box interior. Verified: a FITTING \codeline
  % keeps byte-identical glyph positions, and the box-crossing case now emits
  % "Overfull \hbox (19.03pt too wide)". A wide \codeline should be split in two.
  \hbox to\linewidth{\hfil\code{#1}\hfil}%
  \nopagebreak\vspace{\belowdisplayskip}}
\makeatother

% ---- a thin decorative rule helper ----
\newcommand{\sep}{\par\vspace{2pt}\noindent\textcolor{cRule}{\rule{\linewidth}{0.6pt}}\par\vspace{4pt}}

% ---- colors ----
\definecolor{cBlue}{HTML}{1F4E79}   \definecolor{cBlueBg}{HTML}{EAF1F8}
\definecolor{cGreen}{HTML}{2E6B4F}  \definecolor{cGreenBg}{HTML}{EAF4EE}
\definecolor{cAmber}{HTML}{9C6A00}  \definecolor{cAmberBg}{HTML}{FBF3E2}
\definecolor{cPurple}{HTML}{5B2C6F} \definecolor{cPurpleBg}{HTML}{F3ECF7}
\definecolor{cTeal}{HTML}{0E6E6E}   \definecolor{cTealBg}{HTML}{E6F4F4}
\definecolor{cRed}{HTML}{8E2B2B}    \definecolor{cRedBg}{HTML}{FBEAEA}
\definecolor{cGray}{HTML}{555555}   \definecolor{cRule}{HTML}{1F4E79}

% ---- math operators ----
\DeclareMathOperator{\Tr}{Tr}
\DeclareMathOperator{\rank}{rank}
\DeclareMathOperator{\diag}{diag}
\DeclareMathOperator{\sgn}{sgn}

% ---- bra-ket (quantikz defines \ket/\bra/\braket; add the rest safely) ----
\providecommand{\ket}[1]{\left|#1\right\rangle}
\providecommand{\bra}[1]{\left\langle#1\right|}
\providecommand{\braket}[2]{\left\langle#1\middle|#2\right\rangle}
\newcommand{\dyad}[1]{\left|#1\right\rangle\!\left\langle#1\right|}
% v3.21: general outer product |a><b| and expectation value <A> — agents reach
% for these by reflex (esp. \expval); quantikz defines neither, so their absence
% threw "Undefined control sequence" in a real build. \providecommand = safe.
\providecommand{\ketbra}[2]{\left|#1\right\rangle\!\left\langle#2\right|}
\providecommand{\expval}[1]{\left\langle#1\right\rangle}
\newcommand{\abs}[1]{\left|#1\right|}
\newcommand{\norm}[1]{\left\lVert#1\right\rVert}
\newcommand{\half}{\tfrac12}
\newcommand{\R}{\mathbb{R}}\newcommand{\C}{\mathbb{C}}\newcommand{\Z}{\mathbb{Z}}
% identity operator. NB: \mathbb{1} silently renders a WRONG glyph (msbm has no
% blackboard digits). Use \mathbb{I} = 𝕀 (works with amssymb, no extra package).
% Do NOT "fix" this back to \mathbb{1}.
\newcommand{\id}{\mathbb{I}}

% ---- physical units (v3.35): \punit — upright unit with thin space ---------
% Usage: $300\punit{K}$, $80\punit{mV/decade}$, $10^{7}\punit{cm/s}$ — the
% NUMBER AND UNIT LIVE IN ONE MATH ISLAND. Promoted verbatim from a 199-page
% production lakit (advanced-electronic-devices, dozens of uses).
% Two measured traps this design avoids (both render-proven):
%   1. siunitx SCRAMBLES under RTL — \SI{3.5}{eV} renders "eV 3.5", exponents
%      detach ("Hz 10 x 10^9"): its quantity assembly is not bidi-isolated.
%      Do NOT load siunitx in this pipeline.
%   2. A text-mode unit next to a bare-prose number is TWO adjacent LTR runs
%      and reorders ("5 mm" -> "mm5" — the BIDI-2/3 class). \punit is therefore
%      math-only ON PURPOSE: prose misuse fails LOUDLY at compile (\mathrm
%      outside math) instead of scrambling silently.
% Domain shorthands (\MeV, \fm...) stay per-document: \newcommand{\MeV}{\,\text{\en{MeV}}}
\providecommand{\punit}[1]{\,\mathrm{#1}}

% ============================================================================
% Colored boxes — mdframed (NOT tcolorbox, which breaks under bidi).
% ★★ v3.4: innerleftmargin/innerrightmargin give code listings room so their
% frame/line-numbers do not spill out the box border (see listing styles below).
% ============================================================================
% v3.61: everyline/split skips were applied to `tablebox` only (BLD-14, v3.37) —
% the six colored boxes below never got them, so they still split OPEN across a page
% break: the page-N segment gets top+sides and NO bottom rule, the page-N+1 segment
% sides+bottom and NO top rule. Measured on the shipped fleet this is a COVERAGE GAP,
% not a re-report of BLD-14: 0 split-open `tablebox` segments (that fix holds) vs 62
% split-open COLORED boxes across 7 documents. Verified after the fix: both segments
% close, and the frametitle is NOT duplicated on the continuation. BLD-14's caveat
% still applies — a box taller than a full page grinds mdframed's split-trial loop, so
% split such content manually.
% ============================================================================
% v3.72: frametitle={} leads EVERY \newmdenv option list below. mdframed options
% are GROUP-scoped: a box nested inside \begin{exbox}[frametitle={...}] used to
% INHERIT the outer title — measured, the identical title string rendered twice,
% the second riding a title-bar fill inside the INNER box's frame, for all seven
% environments in both directions. Both natural forms were broken (= preamble
% bug, not an authoring rule). The explicit reset makes every box titleless
% unless ITS OWN [frametitle={...}] says otherwise; verified both directions:
% no duplicate on the nested box, an explicit title still renders.
% Geometry is shared with the optional semantic-width table container. Themes
% select their own existing semantic style below; neither path fixes columns.
\mdfdefinestyle{hllpdfsemanticgeometry}{linewidth=1.1pt,
  innertopmargin=6pt,innerbottommargin=7pt,skipabove=9pt,skipbelow=9pt,
  innerleftmargin=9pt,innerrightmargin=9pt,everyline=true,
  splitbottomskip=6pt,splittopskip=12pt,roundcorner=3pt}
\newcommand{\hllpdfsemanticframestyle}{hllpdfsemanticgeometry}
\newmdenv[frametitle={},style=hllpdfsemanticgeometry,linecolor=cBlue,backgroundcolor=cBlueBg,
  frametitlebackgroundcolor=cBlue,frametitlefontcolor=white,
  ]{defbox}   % הגדרה
\newmdenv[frametitle={},style=hllpdfsemanticgeometry,linecolor=cGreen,backgroundcolor=cGreenBg,
  frametitlebackgroundcolor=cGreen,frametitlefontcolor=white,
  ]{thmbox}   % משפט/טענה
\newmdenv[frametitle={},style=hllpdfsemanticgeometry,linecolor=cAmber,backgroundcolor=cAmberBg,
  frametitlebackgroundcolor=cAmber,frametitlefontcolor=white,
  ]{notebox}  % הערה
\newmdenv[frametitle={},style=hllpdfsemanticgeometry,linecolor=cPurple,backgroundcolor=cPurpleBg,
  frametitlebackgroundcolor=cPurple,frametitlefontcolor=white,
  ]{exbox}    % דוגמה
\newmdenv[frametitle={},style=hllpdfsemanticgeometry,linecolor=cRed,backgroundcolor=cRedBg,
  frametitlebackgroundcolor=cRed,frametitlefontcolor=white,
  ]{warnbox}  % טעות נפוצה
\newmdenv[frametitle={},style=hllpdfsemanticgeometry,linecolor=cTeal,backgroundcolor=cTealBg,
  frametitlebackgroundcolor=cTeal,frametitlefontcolor=white,
  ]{keybox}   % רעיון מרכזי

% ---- codebox (v4.5.0): the dedicated container for code -------------------
% An OPTION, not a policy: asmblock unboxed stays canonical for a full listing,
% and a snippet inside any of the six colored boxes still works as before.
% What makes codebox worth its own environment is that it is the ONE framed box
% that keeps its margin line numbers.
%
% v4.4.0 shipped it WITHOUT them, on the strength of v3.37's "no margin knob can
% save them" plus four fresh probes that agreed. That was wrong: every probe had
% moved the code, the box, or both — never the key that carries the NUMBER
% COLUMN. Measured on the page:
%   resetmargins=false   the listing follows THIS box instead of the page,
%                        moving the code from x=75.26 to x=97.54
%   xrightmargin=26pt    under RTL this is the PHYSICAL LEFT indent, and the
%                        only key that carries the \llap'ed digits with it:
%                        x=64.28 -> 90.18, clearing the rule at 65.64 by 23.6pt
%   xleftmargin=26pt     moves the code and LEAVES THE DIGITS on the rule —
%                        the asymmetry that hid the fix for eleven versions
% resetmargins is hoisted out of the four styles (see the \lstset below them) so
% this override is possible at all; every other listing still resolves to true.
% Pinned by tests/codebox_numbers.tex + probe_lstnum.py --inside-rule, whose
% negative control strips xrightmargin and watches the digits fall back.
\definecolor{cSlate}{HTML}{3F4A56}  \definecolor{cSlateBg}{HTML}{F2F4F7}
\newmdenv[frametitle={},leftline=true,rightline=true,topline=true,bottomline=true,
  linecolor=cSlate,backgroundcolor=cSlateBg,linewidth=0.9pt,
  frametitlebackgroundcolor=cSlate,frametitlefontcolor=white,
  innertopmargin=6pt,innerbottommargin=7pt,skipabove=9pt,skipbelow=9pt,
  innerleftmargin=9pt,innerrightmargin=9pt,everyline=true,splitbottomskip=6pt,splittopskip=12pt,
  roundcorner=3pt]{codebox}   % קוד
% The two keys that make in-box numbering work, both measured on the page:
%   resetmargins=false  the listing follows THIS box instead of the page
%   xrightmargin        under RTL this is the PHYSICAL LEFT indent, and it is
%                       the only one that carries the \llap'ed number column
%                       with it. xleftmargin moves the code and leaves the
%                       digits behind on the rule (measured: code 97.54,
%                       digits 64.28 — the v3.37 defect, reproduced); with
%                       xrightmargin=26pt the digits land at x=90.18, a
%                       clearance of 23.6pt from the rule at 65.64.
% This asymmetry is why v3.4's note had the semantics backwards and v3.37
% concluded the digits could not be saved.
% <hllpdf:listing-container-state>
% v4.29: all state is local to the current environment group. A semantic
% callout masks number text, while a nested codebox restores the formatter
% saved before that mask. Save it once across successive semantic wrappers;
% an explicit numbers=none (or an explicitly empty formatter) remains valid.
\makeatletter
\newif\ifhllpdf@incodebox
\newif\ifhllpdf@englishcodebox
\newif\ifhllpdf@lstnumbermasked
\let\hllpdf@savedlstnumber\thelstnumber
\makeatother
% </hllpdf:listing-container-state>
% <hllpdf:codebox-gutter>
\makeatletter
\AtBeginEnvironment{codebox}{%
  \hllpdf@incodeboxtrue
  % Snapshot the OUTER language before asmblock selects English.
  \iflanguage{english}{\hllpdf@englishcodeboxtrue}{\hllpdf@englishcodeboxfalse}%
  \ifhllpdf@lstnumbermasked
    \let\thelstnumber\hllpdf@savedlstnumber
    \hllpdf@lstnumbermaskedfalse
  \fi
}
% listings' AddToHook is GLOBAL. Register once, and gate on local box state;
% adding an unconditional hook at every codebox entry leaks into later listings.
\lst@AddToHook{PreInit}{%
  \ifhllpdf@incodebox
    % Apply every listing override here, inside the LISTING's own group.
    % Setting margins at box entry leaks them into nested semantic callouts.
    \lstset{resetmargins=false,xrightmargin=26pt}%
    % Styles resolve before PreInit. In LTR the physical number gutter follows
    % xleftmargin; in Hebrew retain the measured inherited xrightmargin path.
    \ifhllpdf@englishcodebox
      \lstset{xleftmargin=26pt,xrightmargin=0pt}%
    \fi
    % ONE box, not two: suppress only the frame/background/skips of a listing
    % inside this codebox. The next unboxed listing keeps its own style.
    \lstset{frame=none,backgroundcolor=\color{cSlateBg},aboveskip=0pt,belowskip=0pt}%
  \fi
}
\makeatother
% </hllpdf:codebox-gutter>

% v3.37: line numbers inside a COLORED BOX are sliced by the box border —
% measured: digit span x=63.5-67.3 vs border line at x=65.2 (the border runs
% THROUGH the digits), and the inter-border gap (3.5pt) is narrower than a
% digit (3.8pt), so no margin knob can save them (under resetmargins,
% xleftmargin is inert in-box — measured: code at x=75.26 with 2.2em AND
% 3.0em). A 3-6 line boxed snippet does not need line numbers; blank the
% number TEXT inside every colored box (\thelstnumber is what gets printed;
% the style's numbers=left survives, so nothing else shifts — code position
% byte-frozen at x=75.26, verified). Unboxed listings keep their numbers.
% v3.78: the blanking is coupled to the FRAME, not to the box name —
% clean_theme.tex registers LATER hooks (registration order wins) that
% restore the number text inside the boxes it renders frameless.
% v4.5.0 RETRACTS the "no margin knob can save them" clause above. It holds only
% under resetmargins=true (the shipped default): there xleftmargin is inert
% in-box, as measured. With resetmargins=false, xrightmargin DOES carry the
% \llap'ed digit column clear of the rule — measured; the codebox \lstset that
% does it is the one above, in the codebox block. The loop below therefore
% covers SIX boxes, not "every colored box": codebox opts out and keeps its
% numbers.
% <hllpdf:semantic-listing-mask>
\makeatletter
\@for\HEB@lstbox:=defbox,thmbox,notebox,exbox,warnbox,keybox\do{%
  \expandafter\AtBeginEnvironment\expandafter{\HEB@lstbox}{%
    \ifhllpdf@lstnumbermasked\else
      \let\hllpdf@savedlstnumber\thelstnumber
    \fi
    \hllpdf@lstnumbermaskedtrue
    \def\thelstnumber{}%
    % The nearest semantic container owns its listing presentation. On exit,
    % TeX's group restores the enclosing codebox state and formatter.
    \hllpdf@incodeboxfalse}}
\makeatother
% </hllpdf:semantic-listing-mask>

% ============================================================================
% Unified TABLE presentation: subtle uniform slate frame + auto-numbered
% caption rendered INSIDE the frame. Use:
%   \begin{tablebox}\centering  <tabular>  \tabcap{<description>}\end{tablebox}
% Replaces the buggy \caption* (which emitted a stray "טבלה N : *" line) and
% provides an explicit framed option + numbered caption "טבלה N.M: תיאור".
% Optional intro text may be placed at the TOP, before the tabular, inside the box.
% ============================================================================
\definecolor{cTblLine}{HTML}{B9C4D2}  \definecolor{cTblBg}{HTML}{FAFBFC}
% v3.23: class-aware — under `article` there is no chapter counter (the promised
% report->article switch used to die with "No counter 'chapter' defined").
\makeatletter
\@ifundefined{c@chapter}{%
  % v3.31: in article the counter resets per SECTION, so an unprefixed number
  % repeated ("טבלה 1." in every section — measured on a two-section render).
  % Same conditional idea as the report branch: prefix the section number once
  % a real \section exists; a pre-section table stays "טבלה 1".
  \newcounter{tbl}[section]
  \renewcommand{\thetbl}{\ifnum\value{section}>0 \thesection.\fi\arabic{tbl}}%
}{%
  \newcounter{tbl}[chapter]
  % v3.30: a report-class doc that uses SECTIONS ONLY (chapter stays 0 —
  % the common short-doc pattern) used to print "טבלה 0.1"; the chapter
  % prefix now appears only once a real \chapter exists (render-verified
  % both ways: "טבלה 1." chapterless, "טבלה 1.1" with chapters).
  \renewcommand{\thetbl}{\ifnum\value{chapter}>0 \thechapter.\fi\arabic{tbl}}%
  % v3.64: the SAME leak in \thesection itself, left behind by v3.30. That fix
  % guarded the CAPTION counter and stopped there, so a chapterless report still
  % printed the heading "0.1 סעיף ראשון" and \ref to it printed "0.2" — measured:
  %   p1 y=65.2  '…0.1'   p1 y=144.9 '…0.2'   p1 y=183.4 '…0.2.1'
  % report.cls hardcodes \thesection as \thechapter.\@arabic\c@section, and the
  % leading 0 is the chapter counter that no \chapter ever incremented. Fixing
  % \thesection fixes the whole chain — \thesubsection and \thesubsubsection are
  % defined in terms of it. \theequation and \thetbl/\thefig were already guarded,
  % which is exactly why this survived: every NUMBER around the heading looked
  % right, so the heading's own 0 read as a style choice rather than a bug.
  % \appendix is safe: report.cls resets chapter to 0 and redefines \thechapter to
  % \@Alph\c@chapter, so an appendix \chapter makes the counter 1 and the guard
  % emits "A.1" unchanged. Render-verified in all four modes (chapterless,
  % chaptered, appendix+chapter, appendix+section-only).
  \renewcommand{\thesection}{\ifnum\value{chapter}>0 \thechapter.\fi\arabic{section}}%
}
\makeatother
% v3.37: everyline=true — if the box DOES split across pages, every segment
% gets a CLOSED frame (measured before: the split ran open-ended into the page
% edge with no bottom border, and the caption could strand alone at the next
% page top — "טבלה A.2" shipped that way). The split skips give the segments
% breathing room. In practice, with the caption GLUE below, every previously-
% split table in the corpus now moves WHOLE instead (the white-gap tradeoff,
% documented benign); everyline stays as the safety net for a genuinely
% page-tall table. NB: a table taller than ONE FULL PAGE can grind mdframed's
% split-trial loop (reproduced: minimal doc, minutes of 1pt-step Overfull
% probes) — break such a monster into two tables.
\mdfdefinestyle{hllpdftablegeometry}{leftline=true,rightline=true,topline=true,bottomline=true,
  linewidth=0.9pt,
  innertopmargin=9pt,innerbottommargin=8pt,skipabove=11pt,skipbelow=11pt,
  innerleftmargin=11pt,innerrightmargin=11pt,roundcorner=4pt}
\newcommand{\hllpdftableframestyle}{hllpdftablegeometry}
\newmdenv[frametitle={},style=hllpdftablegeometry,
  linecolor=cTblLine,backgroundcolor=cTblBg,
  % v4.2.0 (adopted): tablebox is ATOMIC. Short tableboxes near page bottoms
  % repeatedly entered mdframed's split trial, emitting benign Overfull-vbox
  % bursts and occasionally separating the intro prose from the tabular.
  % nobreak=true moves the complete short table to the next page instead. The
  % output-active Overfull-vbox gate still catches a box taller than a page;
  % everyline/split skips remain as defensive styling should that path ever
  % re-arm, but the canonical path never publishes a split tablebox.
  nobreak=true,
  everyline=true,splitbottomskip=6pt,splittopskip=12pt,
  % v3.79 (BLD-37 table side, correction): startinnercode/endinnercode run
  % INSIDE mdframed's body vbox (mdframed.sty:848 / :856+863) -- the second one
  % after \par\unskip, i.e. after line breaking, which is the only moment the
  % \tabcap anchor's true offset from the top of the box EXISTS. See
  % \lakintabboxclose below. These two keys are the ONLY tablebox change.
  startinnercode={\lakintabboxopen},endinnercode={\lakintabboxclose}]{tablebox}
% v3.37: \nopagebreak GLUE — \vsplit honors penalties, so the caption can no
% longer strand alone at the top of the next page (it takes its rows along, or
% the whole box moves). Render-verified on the real shipped orphan.
% v3.59: \NR@gettitle is what nameref injects into the stdlib title-bearing
% commands (\@sect, \@chapter, \caption) to fill \@currentlabelname. \figcap and
% \tabcap are INDEPENDENT macros — they never called it, so \@currentlabelname kept
% whatever the last heading left and \nameref{fig:x} silently printed the enclosing
% SECTION's title instead of the caption (proven in the .aux: field 3 of \newlabel
% was the heading text while field 1 and the anchor `fig.1.1` were correct — so this
% is a THIRD, distinct head next to BLD-13's counter fix and BIDI-38's autoref-name
% fix, neither of which writes that field). \@ifundefined keeps it inert if
% hyperref/nameref is ever dropped. Guarded helper because \figcap is defined
% after \makeatother.
% NB both helpers below need \makeatletter — this part of the preamble is OUTSIDE
% any letter-catcode region, so an unwrapped \@currentlabel parses as LaTeX's \@
% and dies with "You can't use \spacefactor in vertical mode" (measured).
\makeatletter
\newcommand{\lakinrtitle}[1]{\@ifundefined{NR@gettitle}{}{\NR@gettitle{#1}}}
% v3.59: \hitem{<letter>}{<key>} — a REFERENCEABLE Hebrew-lettered list item.
%   \begin{enumerate} \hitem{א}{it:a} הסעיף הראשון. \end{enumerate}  ->  \ref{it:a} = (א)
% Why this exists: automatic Hebrew list labels fail silently, so SKILL.md mandates
% an EXPLICIT \item[(\textbf{א})] — but LaTeX puts \refstepcounter INSIDE
% \if@noitemarg (latex.ltx:12964), and \item[...] never sets that flag. So an
% explicit item steps no counter and never touches \@currentlabel: a \label there
% binds to the AMBIENT one and \ref prints the CHAPTER/SECTION number where the
% item letter belongs (proven in the .aux: {chapter.1}). Same doctrine as \figcap —
% own the label state before \label. \phantomsection matters: without it the
% hyperlink still lands on the chapter head (measured); with it the anchor is
% section*.N and the bookmark outline stays clean.
\newcommand{\hitem}[2]{\item[(\textbf{#1})]%
  \def\@currentlabel{(#1)}\def\@currentlabelname{#1}\phantomsection\label{#2}}
% NB (v3.66): \hitem is \ref-only. \autoref{key} calls it a SECTION ("סעיף (א)")
% because \phantomsection brands the anchor section*.N, and \nameref{key} prints
% the bare letter. For \autoref-able items prefer \hebenumlabels below.
%
% v3.66: \hebenumlabels — OPT-IN automatic Hebrew enumerate letters, done right.
%   {\hebenumlabels\begin{enumerate}\item\label{q:a} ...\end{enumerate}}
%   -> labels (א) (ב) (ג); \ref{q:a} = א; anchor Item.N (real, per-item).
% This retires the belief that made \item[(\textbf{א})] + \hitem necessary.
% "Automatic Hebrew labels fail silently" was TRUE but unexplained, and the
% mechanism is now bisected (plain article vs +enumitem, measured): ENUMITEM
% PRE-EXPANDS the label expression in an \edef BEFORE \refstepcounter, with the
% counter still at 0. Tokens it shields by name (\arabic,\alph,\roman) survive;
% a plain \newcommand \ifcase wrapper evaluates at 0 -> the EMPTY zero-branch ->
% the label prints as "()" with the letter GONE (and \the\c@enumi prints a
% literal 0, \number\value dies fatally pointing at \item). The cure is one
% property, not a doctrine: the wrapper must be ROBUST (survives the \edef
% unexpanded), while \theenumi stays EXPANDABLE (so \refstepcounter writes the
% real letter to the .aux). Verified: labels (א)(ב)(ג), \ref -> א/ב/ג each item
% its own value, .aux = {א}{...}{Item.1} — versus the explicit-\item doctrine
% where BIDI-50/\hitem exist purely to patch the counter that never stepped.
\def\@hebalph#1{\ifcase#1\or א\or ב\or ג\or ד\or ה\or ו\or ז\or ח\or ט\or י\or
  יא\or יב\or יג\or יד\or טו\or טז\or יז\or יח\or יט\or כ\else\@ctrerr\fi}
\DeclareRobustCommand*\hebalph[1]{\expandafter\@hebalph\csname c@#1\endcsname}
\newcommand{\hebenumlabels}{%
  \renewcommand{\theenumi}{\@hebalph\c@enumi}%
  \renewcommand{\labelenumi}{(\textbf{\hebalph{enumi}})}}
\makeatother
% v3.63: \mbox the LABEL. \thetbl ("7.1") and a caption description that STARTS
% WITH A DIGIT ("32 הרגיסטרים…") are two European-Number runs separated only by
% neutrals (the label's own "." and \enspace). Under the RTL paragraph base the UBA
% resolves those neutrals to the RTL level, so both numbers sit in ONE reversed span
% and EXCHANGE VISUAL POSITIONS. Measured on the shipped PDF: a table whose real
% number is 7.1 printed "טבלה 32" — a different, EXISTING table number — while \ref
% correctly printed 7.1, so the caption silently contradicted its own cross-reference,
% and the description lost its leading count. The source contains no adjacency at all;
% it is manufactured by this macro, which is why every source linter was blind.
% \mbox makes the label one atomic box so nothing can reorder across it. Verified
% zero-regression: an ordinary Hebrew caption keeps identical x-geometry.
% One paragraph style covers label, number, separator and description.
% Explicit project/template overrides remain possible; the label hook is kept
% for compatibility, with an identity default. Counters/anchors stay unchanged.
\newcommand{\lakitcaptionstyle}{\normalfont\small\color{black}\hypersetup{allcolors=black}}
\newcommand{\lakitcaptionlabelstyle}[1]{#1}
\newcommand{\lakitcaptionpunctuation}{:}
% Native floats, captionof and longtable use the same presentation policy.
% subcaption is optional and inherits the same font rather than a second size.
\usepackage{caption}
\DeclareCaptionFont{lakit}{\lakitcaptionstyle}
\DeclareCaptionLabelFormat{lakit}{\lakitcaptionlabelstyle{#1~\en{#2}}}
\DeclareCaptionLabelSeparator{lakit}{\lakitcaptionpunctuation\enspace}
\captionsetup{font=lakit,labelfont={},textfont={},labelformat=lakit,
  labelsep=lakit,justification=centering}
\makeatletter
\AtBeginDocument{\@ifpackageloaded{subcaption}{%
  \captionsetup[sub]{font=lakit,labelfont={},textfont={},labelsep=lakit}}{}}
\makeatother
% v3.64: OPT-IN list-of-figures/tables entry via an optional short caption.
% \listoffigures and \listoftables are inherited from report.cls and LOOK available,
% but \figcap/\tabcap step the custom \c@fig/\c@tbl counters and never call
% \addcontentsline — so both printed their heading and NOTHING ELSE. Measured: a doc
% with one \figcap and one \tabcap rendered page 1 as 'רשימת האיורים\n1\n', and the
% .lof held only \@chapter's \addvspace, zero \contentsline entries. An empty list
% under a correct Hebrew heading reads as "this document has no figures".
%
% WHY OPTIONAL, AND NOT AUTOMATIC. Writing #1 unconditionally was implemented first
% and REJECTED: a caption containing \codeline compiles CLEANLY today (verified on
% the pre-v3.64 preamble: 0 errors, PDF produced) but becomes FATAL once it is also
% written to the .lof — \codeline is \hbox to\linewidth{...}, and a full-width box in
% a list entry dies with "\pdfextension endlink cannot be used in vertical mode", no
% output PDF. Gaining a feature nobody uses (0 of 13 collections call \listoffigures)
% at the price of breaking a construct that works is a bad trade. Stock \caption has
% exactly this limitation and exactly this escape hatch, so the convention is familiar.
% \en{}, \code{}, \h{} and inline math in the short text are all safe — verified, they
% survive the write and typeset correctly in the list (RTL geometry confirmed: number
% rightmost at x=509.9, title left of it, dots and folio at the far left).
% With no optional argument the macro is byte-identical to v3.63 — zero risk to every
% existing document, and \listoffigures lists exactly the captions that opted in.
% BLD-37, TABLE side (the \figcap v3.74 idea, adapted): clicking a \ref to a
% table used to land MID-table — hyperref's automatic anchor inside
% \refstepcounter was emitted on the still-open horizontal line whose TAIL is
% the tabular (an hbox whose baseline is the table's vertical CENTER), so the
% jump scrolled the top half of the table above the viewport. The v3.74 entry
% deferred this side claiming "tablebox splits — no single box to measure";
% that deferral dissolves: the measurement target is the TABULAR itself at
% \tabcap ENTRY, while it is still the tail of the open list — the same
% situation \lakin@grabfig handles at an island's end (and the fleet has 0
% split tableboxes). So \lakin@grabtab runs as the FIRST token, before any
% typesetting can push the box away, and the anchor is re-emitted raised above
% the measured tabular. \lakintabraise is a SEPARATE register from
% \lakinfigraise (defined with it, below the LTR-island block): a captionless
% island intentionally leaves the fig register set, and a table must never
% inherit that stale figure height.
%
% WHY the anchor stays AT THE ENTRY SPOT and is NOT deferred into the caption
% paragraph the way \figcap defers it: at \par TeX DROPS a final glue item,
% and the automatic anchor's smash box sitting BEHIND the interword glue kept
% that glue alive — so every centered tabular has always sat half a space off
% true center. Deferring the anchor removed that box and was measured as 982
% differing glyph records: EVERY tabular line shifted 1.30pt (2.61pt for
% pattern 7's two minipages) while captions stayed put. Emitting the raised
% anchor at the old spot — through hyperref's own \Hy@raisedlink, exactly the
% macro the automatic path uses — keeps the line's node structure, and so the
% A/B glyph diff at 0, BY CONSTRUCTION. \Hy@raisedlink reads the raise from
% \HyperRaiseLinkDefault (default \baselineskip = the old mid-table dest); the
% group-local \edef freezes measured-height+42pt into it, and \smash inside
% \Hy@raisedlink hides any raise from the line's extent. \HyCnt@ProvideTheH-
% Counter keeps the dest NAME byte-identical to the automatic one (it defines
% \theHtbl, which \hyper@makecurrent prefers — without it the first measured
% caption was named tbl.1 where the automatic path said tbl.0.1, measured).
% vmode entry, or anything but a box at the tail (no tabular to measure),
% falls back to the plain \refstepcounter — byte-identical to the old path.
\makeatletter
\newcommand{\tabcap}[2][]{\lakin@grabtab
  % v3.79: arm the box-close correction. Set unconditionally and in BOTH
  % branches below's scope, so the plain-\refstepcounter fallback is corrected
  % too; \lakintabboxopen clears it at every \begin{tablebox}, so a \tabcap
  % outside a tablebox can never leave it armed for a later box.
  \global\lakintabcapseentrue
  \ifdim\lakintabraise>0pt
    \H@refstepcounter{tbl}\HyCnt@ProvideTheHCounter{tbl}\hyper@makecurrent{tbl}%
    \begingroup
      % A frameless table targets its measured top. The frame's historical
      % cushion belongs only to tablebox and is corrected at box close.
      \edef\HyperRaiseLinkDefault{\the\lakintabraise}%
      \iflakintabframed
        \edef\HyperRaiseLinkDefault{\the\dimexpr\lakintabraise+42pt\relax}%
      \fi
      \Hy@raisedlink{\hyper@anchorstart{\@currentHref}\hyper@anchorend}%
    \endgroup
    \global\lakintabraise 0pt\relax
  \else
    \refstepcounter{tbl}%
  \fi
  \lakinrtitle{#2}\par\nopagebreak\addvspace{6pt}%
  \def\lakin@short{#1}\ifx\lakin@short\@empty\else
    \addcontentsline{lot}{table}{\protect\numberline{\thetbl}{\protect\ignorespaces #1}}\fi
  {\centering\lakitcaptionstyle\mbox{\lakitcaptionlabelstyle{\tablename~\en{\thetbl}\lakitcaptionpunctuation}}\enspace #2\par}}
\makeatother

% ============================================================================
% Unified FIGURE caption — the \tabcap twin for TikZ/plot figures (v3.34.0).
% A figure (tikzpic) sits inline in the prose; put \figcap BELOW it for a
% numbered, \ref-able caption "איור N.M: תיאור" in the same style as \tabcap:
%   \begin{tikzpic} ... \end{tikzpic}
%   \figcap{\label{fig:foo}description}     % \label FIRST, inside the argument
% WHY a macro, not a hand-written gray line: a manual
%   {\centering\small\textcolor{cGray}{איור: ...}} with a \label inside has NO
% \refstepcounter, so the \label binds to the AMBIENT \@currentlabel (the last
% equation) — \ref then prints an equation number AS the figure number
% (measured on a real page: "איור 48" where 48 is equation 48). \figcap steps a
% real \c@fig counter first, so \label is correct, numbering is consistent, and
% refs are stable. Class-aware like \thetbl (chapter- or section-prefixed).
% ============================================================================
\makeatletter
\@ifundefined{c@chapter}{%
  \newcounter{fig}[section]
  \renewcommand{\thefig}{\ifnum\value{section}>0 \thesection.\fi\arabic{fig}}%
}{%
  \newcounter{fig}[chapter]
  \renewcommand{\thefig}{\ifnum\value{chapter}>0 \thechapter.\fi\arabic{fig}}%
}
\makeatother
\makeatletter
\newcommand{\figcap}[2][]{%
  \ifdim\lakinfigraise>0pt
    % v3.74: step WITHOUT hyperref's automatic anchor, then emit the anchor
    % raised above the measured figure. \H@refstepcounter is hyperref's saved
    % pre-hook \refstepcounter (sets the counter AND \@currentlabel — \label
    % in the caption arg binds exactly as before); \hyper@makecurrent supplies
    % \@currentHref, so the dest NAME is identical to the old automatic one —
    % \ref/\autoref/bookmarks are unchanged, only the dest'S POSITION moves.
    \H@refstepcounter{fig}\hyper@makecurrent{fig}%
    % The anchor is emitted INSIDE the caption paragraph below, as a zero-size
    % raisebox — NOT here in vertical mode, where a \raisebox starts an empty
    % paragraph and adds a blank line (measured: every glyph after the caption
    % shifted +18.8pt). Inside an existing line a [0pt][0pt] box is inert.
    \edef\lakin@anch{\@currentHref}%
    % \edef FREEZES the measured height into the pending macro: the deferred
    % \raisebox is typeset inside the caption paragraph AFTER the register is
    % reset to the 0pt sentinel below — a plain \def read 0 at use time
    % (measured: LAKINRAISE=0.0, dest moved 14pt instead of above the figure).
    \edef\lakin@pend{\noexpand\raisebox{\the\dimexpr\lakinfigraise+42pt\relax}[0pt][0pt]{%
      \noexpand\hyper@anchorstart{\lakin@anch}\noexpand\hyper@anchorend}}%
    \global\lakinfigraise 0pt\relax
  \else
    \refstepcounter{fig}\let\lakin@pend\relax
  \fi
  \lakinrtitle{#2}\par\nopagebreak\addvspace{4pt}%
  \def\lakin@short{#1}\ifx\lakin@short\@empty\else
    \addcontentsline{lof}{figure}{\protect\numberline{\thefig}{\protect\ignorespaces #1}}\fi
  {\centering\lakin@pend\lakitcaptionstyle\mbox{\lakitcaptionlabelstyle{\figurename~\en{\thefig}\lakitcaptionpunctuation}}\enspace #2\par}}
\makeatother

% ============================================================================
% LTR islands.  EVERY tikzpicture / quantikz MUST live inside one, or the whole
% figure's coordinates and Latin labels flip. (Hebrew node text still uses \h{}.)
% ============================================================================
% v3.74 (BLD-37): the \figcap hyperlink DESTINATION used to land at the caption
% line — BELOW the figure — so clicking a figure reference scrolled the figure
% off-screen (measured fleet-wide: 97/97 fig dests in one collection sat 19-54pt
% under their figure's ink; a 347pt figure was fully hidden after the jump).
% hypcap-style fix in two parts: (1) at the island's END, \lakin@grabfig grabs
% the just-typeset centered line with \lastbox, measures it, and puts it back
% with \nointerlineskip — the interline glue above the line was already
% contributed and is NOT re-added, so the page geometry is unchanged to the
% glyph (verified); (2) \figcap then steps the counter WITHOUT hyperref's
% automatic anchor (\H@refstepcounter = the pre-hyperref \refstepcounter) and
% emits the anchor itself, raised by the measured height — the dest lands ABOVE
% the figure top. A captionless island leaves numbering untouched (nothing is
% stepped at the island; the grab flag is cleared at every island begin, so a
% stale height from a captionless figure cannot leak into a later caption).
\newsavebox{\lakinfigbox}
\newdimen\lakinfigraise  % real \dimen: tex.setdimen cannot target a \newlength SKIP register (measured lua error)
% BLD-37 table side: a SEPARATE register for \tabcap. The two must not share —
% a captionless island INTENTIONALLY leaves \lakinfigraise set (the grab runs
% at every island END; only the NEXT island's begin clears it), so a shared
% register would hand a \tabcap after a captionless figure that stale FIGURE
% height. \lakintabraise is written only by \lakin@grabtab at \tabcap entry
% and consumed/reset inside \tabcap itself, so it can never go stale.
\newdimen\lakintabraise
\makeatletter
% v3.74b: measure WITHOUT touching the list. The first draft grabbed the box
% with \lastbox and put it back — and \lastbox DISCARDS the node's baseline
% shift, so a vertically-anchored quantikz circuit rose 2.57pt and everything
% after it sank 5.32pt (measured on the A/B glyph diff). LuaTeX lets us READ
% the tail node's dimensions without removing anything: geometry is identical
% by construction. \lakinfigraise=0pt is the "no figure measured" sentinel.
\def\lakin@grabfig{%
  \ifhmode
    \directlua{
      local n = tex.nest[tex.nest.ptr].tail
      local GLUE, PEN = node.id('glue'), node.id('penalty')
      while n and (n.id == GLUE or n.id == PEN) do n = n.prev end
      local HL, VL = node.id('hlist'), node.id('vlist')
      if n and (n.id == HL or n.id == VL) then
        tex.setdimen('global', 'lakinfigraise',
                     (n.height or 0) + (n.depth or 0))
      else
        tex.setdimen('global', 'lakinfigraise', 0)
      end}%
  \else
    \global\lakinfigraise 0pt\relax
  \fi}
% BLD-37 table side: \lakin@grabfig's twin, NOT a shared helper — the Lua body
% differs in exactly one place, the skip set. Between \end{tabular} and \tabcap
% an otherlanguage wrapper (patterns 3/5 of example_tables) leaves a \write
% whatsit and a dir node ON TOP of the interword glue (measured node dump:
% "whatsit/write <- dir <- glue(2.62) <- hlist"), so the table reader must
% also step over whatsit/dir markers — all zero-size — to reach the tabular.
% \lakin@grabfig deliberately does NOT get the wider skip set: its glue/penalty
% body is the one the v3.74 A/B proof froze, and widening it would newly
% measure (and so newly raise) anchors in island shapes that today take the
% plain path — an unproven behavior change on the FIG side. Anything that is
% not a box behind those markers still hits the 0pt sentinel and \tabcap falls
% back to the plain \refstepcounter.
% SECOND deliberate difference: this stores the box's HEIGHT, not height+depth.
% \figcap's anchor is emitted in the caption paragraph BELOW the figure, so it
% must climb the full extent; \tabcap's anchor sits ON the tabular's own line,
% whose baseline is the vcentered tabular's MIDPOINT — height alone reaches
% the top edge, and adding the depth was measured to overshoot by half the
% table (dests 80-135pt above their frame, INSIDE the preceding table's box).
\def\lakin@grabtab{%
  \ifhmode
    \directlua{
      local n = tex.nest[tex.nest.ptr].tail
      local GLUE, PEN = node.id('glue'), node.id('penalty')
      local WHAT, DIR = node.id('whatsit'), node.id('dir')
      while n and (n.id == GLUE or n.id == PEN
                   or n.id == WHAT or n.id == DIR) do n = n.prev end
      local HL, VL = node.id('hlist'), node.id('vlist')
      if n and (n.id == HL or n.id == VL) then
        tex.setdimen('global', 'lakintabraise', n.height or 0)
      else
        tex.setdimen('global', 'lakintabraise', 0)
      end}%
  \else
    \global\lakintabraise 0pt\relax
  \fi}
\makeatother
% ============================================================================
% v3.79 (BLD-37 table side, CORRECTION): put the anchor at the TOP OF THE BOX,
% instead of guessing the distance to it.
% ============================================================================
% v3.77 raised the \tabcap anchor by (height of the box at the tail of the open
% horizontal list) + 42pt. That is a MEASUREMENT taken BEFORE line breaking, and
% it is short by every term it cannot see. Fleet-measured on the 12 built
% collections (299 tbl.* dests, ownership established through each table's own
% CAPTION, never through "which frame encloses the dest"): 3 dests landed INSIDE
% their own tablebox, below the frame top —
%   computer-architecture       tbl.14.1  -46.4pt   three stacked tabulars
%   computer-architecture       tbl.30.5  -13.5pt   tabular 2.85pt over \linewidth
%   advanced-electronic-devices tbl.B.1   -59.2pt   tabular 2.56pt over \linewidth
% Two missing terms, both reproduced in a minimal document (probe P2/P3):
%   (1) VERTICAL LIST ALREADY CONTRIBUTED inside the box — a second and third
%       tabular, or a heading above the table. The tail read sees only the LAST
%       box, so 77.7pt of stacked tabulars is invisible to the raise.
%   (2) WHICH LINE THE ANCHOR LANDS ON. When the tabular alone overfills
%       \linewidth (459.6121pt in a tablebox), the interword glue after
%       \end{tabular} becomes the chosen breakpoint and \Hy@raisedlink is
%       emitted on the NEXT line — one (depth + lineskip) lower. That term does
%       not EXIST at \tabcap time: line breaking has not happened yet. No
%       measurement taken there can ever recover it. (Proof: probe P3, the same
%       table 2.84912pt too wide, fails at -13.5pt; probe P4, the identical
%       table with \tabcolsep 2pt so it fits, is correct at +32.0pt.)
% So the raise is not computed at \tabcap at all. It is CORRECTED at
% \end{tablebox}, from inside mdframed's own body vbox (endinnercode runs after
% \par\unskip, i.e. after line breaking, while tex.nest[tex.nest.ptr] is still
% that vbox's vertical list). There the anchor's true offset from the top of the
% box content is not estimated but WALKED, and the \Hy@raisedlink shift is set
% so the dest lands 42pt above it — which is 42pt - innertopmargin(9pt) -
% linewidth(0.9pt) = 32.0pt above the FRAME top, the exact position the 295
% already-correct dests have today. That is why this is a correction and not a
% re-placement: the good sites do not move.
%
% WHY THE GLYPH DIFF IS 0 BY CONSTRUCTION — the only thing written is the
% `shift` field of the hlist that \Hy@raisedlink already built, and the guard
% below refuses to write it unless that box is 0x0x0. A zero-size box carries no
% dimension into its parent's packing, the packing already happened, and shift
% is not an input to it. Nothing else in the list is touched: no node is added,
% removed, resized or reordered. Measured, not assumed: assets/example_tables.tex
% and the whole computer-architecture collection A/B at 0 differing pymupdf
% rawdict glyph records.
%
% WHY NOT the alternatives (all prototyped and measured before this one):
%   * accumulate the vertical list at \tabcap time — fixes term (1), CANNOT fix
%     term (2); 2 of the 3 real failures survive. A fix that leaves the majority
%     of the defect standing is the false all-clear this repo exists to avoid.
%   * \zref-savepos two-pass — correct on all shapes, but it makes the dest
%     POSITION aux-derived (DEFECTS BLD-1/2), a single-pass compile silently
%     ships an unraised anchor, and the bare y-difference has no page guard.
%   * emit the anchor at \begin{tablebox} and let \tabcap adopt it — needs the
%     tbl counter stepped at box OPEN to keep the dest NAME byte-identical; a
%     captionless tablebox then either shifts every later table number or leaves
%     an orphan dest that the next table duplicates (a HARD build gate). This
%     correction touches neither the counter nor the name: \tabcap's naming path
%     is v3.77's, unchanged, so resolve_names() is byte-identical (verified).
%   * suppressing the line break (a penalty before the glue) — that MOVES the
%     tabular: the break currently discards the interword glue, so forbidding it
%     re-centers the two overfull tables. Reflow, rejected.
%
% SCOPE. The correction is armed only by \tabcap ITSELF (\lakintabcapseen),
% cleared at every box open, so a captionless tablebox is a no-op and a \tabcap
% hosted OUTSIDE a tablebox (the fleet has exactly one, computer-architecture
% tbl.14.2 in an exbox) keeps v3.77's measured raise. Both fallbacks are
% therefore never worse than today; and because hyperref's AUTOMATIC anchor uses
% the same \Hy@raisedlink shape, the plain-\refstepcounter path (vmode entry, or
% prose between \end{tabular} and \tabcap — probe P6, -51.5pt) is corrected too.
% KNOWN LIMIT — a tablebox that SPLITS across pages. The offset is measured from
% the top of the UNSPLIT body, so a dest in segment 2 gets raised past its own
% segment. MEASURED on a forced split (probe: 40 filler lines + a table): the
% dest went from y=585.0 on page 2 (inside segment 2, at the caption — the old
% BLD-37 symptom) to y=-113.4, i.e. ABOVE the page top. Viewers clamp that to the
% page top, which happens to show segment 2 whole, but an out-of-page coordinate
% is not something to ship on purpose. Hence the \textheight guard in
% lakintab.fix: a body TALLER THAN THE TEXT BLOCK cannot fit on any page, so it
% must split — the correction is skipped and v3.77's behaviour stands unchanged
% (verified: the probe returns to y=585.0). The guard is deliberately EXACT, not
% a page-fitting heuristic: it is sound (a taller-than-page body always splits)
% but PARTIAL (a shorter body that starts low on a page could still split). That
% residue is unreachable in practice — v3.37's everyline + the caption's
% \nopagebreak glue move such boxes WHOLE instead (probe N7: a 22-row table
% pushed down by 430pt moved whole, 0 segments), and the fleet has 0 split
% tableboxes, re-counted after this change. Reimplementing mdframed's own
% page-fitting arithmetic to close the residue would be a heuristic with its own
% failure mode, which is the trade this repo has repeatedly refused.
\newif\iflakintabcapseen
\newif\iflakintabframed
% The Lua lives in ONE \directlua at load time, and is written with NO
% backslash, %, # or ~ : \directlua expands its argument like \edef, so an
% active ~ would expand to \nobreakspace and a % would eat the rest of the line
% (the second is how the first draft of this block died). "not (a == b)" stands
% in for Lua's inequality operator for the same reason.
\directlua{
  lakintab = { }
  lakintab.HL = node.id("hlist")
  lakintab.VL = node.id("vlist")
  lakintab.RU = node.id("rule")
  lakintab.GL = node.id("glue")
  lakintab.KE = node.id("kern")
  lakintab.WH = node.id("whatsit")
  lakintab.TM = node.id("temp")
  lakintab.WD = node.subtype("pdf_dest")
  function lakintab.d(v)
    if v == nil then return 0 end
    if v < -1000000000 then return 0 end
    return v
  end
  function lakintab.scan(list, y, horiz, parent, target)
    local x = list
    while x do
      local i = x.id
      if i == lakintab.HL or i == lakintab.VL then
        if horiz then
          local b = y + x.shift
          if i == lakintab.HL then
            lakintab.scan(x.head, b, true, x, target)
          else
            lakintab.scan(x.head, b - lakintab.d(x.height), false, x, target)
          end
        else
          if i == lakintab.HL then
            lakintab.scan(x.head, y + lakintab.d(x.height), true, x, target)
          else
            lakintab.scan(x.head, y, false, x, target)
          end
          y = y + lakintab.d(x.height) + lakintab.d(x.depth)
        end
      elseif i == lakintab.WH then
        if x.subtype == lakintab.WD then
          local id = x.dest_id
          if type(id) == "string" and string.sub(id, 1, 4) == "tbl." then
            if horiz and parent and parent.width == 0
               and parent.height == 0 and parent.depth == 0 then
              parent.shift = parent.shift + (target - y)
            end
          end
        end
      elseif not horiz then
        if i == lakintab.RU then
          y = y + lakintab.d(x.height) + lakintab.d(x.depth)
        elseif i == lakintab.GL or i == lakintab.KE then
          y = y + lakintab.d(x.width)
        end
      end
      x = x.next
    end
  end
  function lakintab.fix(target, maxh)
    local nst = tex.nest[tex.nest.ptr]
    if nst == nil then return end
    local h = nst.head
    if h == nil then return end
    if h.id == lakintab.TM then h = h.next end
    local tot = 0
    local x = h
    while x do
      local i = x.id
      if i == lakintab.HL or i == lakintab.VL or i == lakintab.RU then
        tot = tot + lakintab.d(x.height) + lakintab.d(x.depth)
      elseif i == lakintab.GL or i == lakintab.KE then
        tot = tot + lakintab.d(x.width)
      end
      x = x.next
    end
    if tot > maxh then return end
    lakintab.scan(h, 0, false, nil, target)
  end
}
% @-FREE NAMES ON PURPOSE: these two are named in the `tablebox` key list at
% \newmdenv above, which is tokenised OUTSIDE any \makeatletter region — a
% \lakin@... there would tokenise as \lakin followed by the letters @tabopen.
\newcommand{\lakintabboxopen}{\lakintabframedtrue\global\lakintabcapseenfalse}
\newcommand{\lakintabboxclose}{%
  \iflakintabcapseen
    \directlua{lakintab.fix(-\number\dimexpr\lakintabcushion\relax,
                            \number\dimexpr\textheight\relax)}%
    \global\lakintabcapseenfalse
  \fi\lakintabframedfalse}
% The cushion above the box's CONTENT top. 42pt puts the dest 32.0pt above the
% frame top (42pt - innertopmargin 9pt - linewidth 0.9pt), i.e. exactly where
% the 295 correct dests already are — the correction moves only the broken ones.
\newcommand{\lakintabcushion}{42pt}
\newenvironment{qcirc}{\par\medskip\global\lakinfigraise 0pt\relax\begin{otherlanguage}{english}\begin{center}}
  {\csname lakin@grabfig\endcsname\end{center}\end{otherlanguage}\par\medskip}
\newenvironment{tikzpic}{\par\medskip\global\lakinfigraise 0pt\relax\begin{otherlanguage}{english}\begin{center}}
  {\csname lakin@grabfig\endcsname\end{center}\end{otherlanguage}\par\medskip}

% ============================================================================
% Code listings.  ★★ v3.4 THE BOX-OVERFLOW FIX:
%   resetmargins=true   → the listing IGNORES the enclosing environment's
%                         indentation and measures from the PAGE. v3.4 stated
%                         this backwards ("measured from the ENCLOSING box"),
%                         and v3.37 then blanked in-box line numbers because
%                         the page-anchored listing put its hanging digits on
%                         top of the box rule. Measured v4.5.0: with the key
%                         FALSE the listing follows the box and the code moves
%                         from x=75.26 to x=97.74, clearing the rule at 72.07.
%                         The key now lives in one \lstset below instead of in
%                         each style, so a box can override it — codebox does.
%   framexleftmargin=0  → the frame (left bar) does not extend left of the code.
%   xleftmargin keeps a small indent so the line numbers have room INSIDE.
% asmblock wraps the listing in an LTR island (code is LTR).
% ============================================================================
\lstdefinestyle{asm}{
  basicstyle=\ttfamily\small,
  backgroundcolor=\color{cBlueBg!55},
  frame=leftline,framerule=2.5pt,rulecolor=\color{cBlue},
  numbers=left,numberstyle=\tiny\color{cGray},numbersep=8pt,
  xleftmargin=2.2em,framexleftmargin=0pt,framexrightmargin=0pt,
  breaklines=true,columns=fullflexible,keepspaces=true,
  showstringspaces=false,aboveskip=8pt,belowskip=8pt,
  % v3.36: without a comment DELIMITER the commentstyle below never fires —
  % listings has no `language` here, so it does not know '#' opens a comment,
  % and every asm listing rendered flat black (measured: a whole-page pixel
  % sweep found ZERO cGreen glyphs across 135 shipped listings). morecomment
  % gives it the delimiter; color-only change, glyph widths untouched.
  morecomment={[l]{\#}},
  commentstyle=\color{cGreen},
}
\lstdefinestyle{cstyle}{
  language=C,basicstyle=\ttfamily\small,
  backgroundcolor=\color{cGreenBg!55},
  frame=leftline,framerule=2.5pt,rulecolor=\color{cGreen},
  numbers=left,numberstyle=\tiny\color{cGray},numbersep=8pt,
  xleftmargin=2.2em,framexleftmargin=0pt,framexrightmargin=0pt,
  breaklines=true,columns=fullflexible,keepspaces=true,
  showstringspaces=false,keywordstyle=\color{cBlue}\bfseries,
  commentstyle=\color{cGray}\itshape,aboveskip=8pt,belowskip=8pt,
}
% v3.36: pystyle/htmlstyle — colored twins of cstyle for the Python/HTML
% listings that used to repeat a ~200-char options blob with NO color
% (keywords rendered bold-black, comments italic-black — measured; C alone
% was colored). Geometry is copied VERBATIM from the shipped blob (frame,
% margins, skips) so swapping a call site to style=pystyle moves NOTHING —
% only keywordstyle/commentstyle/stringstyle are added (bold stays bold,
% italic stays italic, mono metrics identical => zero reflow, proven by
% per-page word-bbox diff). breaklines/breakatwhitespace stay per-site.
% v3.78: numbers=left joins both styles, with the exact numberstyle/numbersep
% keys asm/cstyle carry — fleet-wide consistency (author decision: all four
% styles numbered; "asm/cstyle numbered, py/html not" was an accident of the
% v3.36 geometry freeze, not a design). The numbers hang LEFT of the code box
% (listings sets them with \llap, so no overfull hbox can arise and the code
% x is untouched — measured: code frozen, digits at code_x-11.8pt); inside a
% colored box they are blanked like everyone else's (v3.37 hook), and the
% clean variant restores them in its frameless boxes (v3.78, above).
\lstdefinestyle{pystyle}{
  language=Python,basicstyle=\ttfamily\small,
  columns=fullflexible,keepspaces=true,
  frame=single,rulecolor=\color{cTblLine},
  numbers=left,numberstyle=\tiny\color{cGray},numbersep=8pt,
  xleftmargin=0.5em,framexleftmargin=0pt,framexrightmargin=0pt,
  showstringspaces=false,aboveskip=8pt,belowskip=8pt,
  keywordstyle=\color{cBlue}\bfseries,
  commentstyle=\color{cGray}\itshape,
  stringstyle=\color{cGreen},
}
% v4.5.0: hoisted out of the four styles so an enclosing box can override it.
% Every listing still resolves to true unless something later says otherwise,
% so shipped output is unchanged (verified byte-for-byte on all eight PDFs).
\lstset{resetmargins=true}
\lstdefinestyle{htmlstyle}{
  language=HTML,basicstyle=\ttfamily\small,
  columns=fullflexible,keepspaces=true,
  frame=single,rulecolor=\color{cTblLine},
  numbers=left,numberstyle=\tiny\color{cGray},numbersep=8pt,
  xleftmargin=0.5em,framexleftmargin=0pt,framexrightmargin=0pt,
  showstringspaces=false,aboveskip=8pt,belowskip=8pt,
  keywordstyle=\color{cBlue}\bfseries,
  commentstyle=\color{cGray}\itshape,
  stringstyle=\color{cGreen},
}
\newenvironment{asmblock}{\par\begin{otherlanguage}{english}}{\end{otherlanguage}\par}

% <hllpdf:listing-smart-quotes>
% v4.21.0: four Unicode punctuation characters must enter listings' token
% buffer. Unregistered characters bypass it: a closing quote was printed
% BEFORE the pending word, and an opening quote BEFORE a pending brace.
% Extend the standard character table, then activate only these codepoints
% in the listing's local character-table scope. Do this BEFORE the package's
% delimiter/alsoother/literate hooks so explicit user mappings still win.
% This is not universal Unicode support, nor a conversion to ASCII quotes.
% See listings developer guide 15.3.1 and references/listing-smart-quotes.md.
\makeatletter
\lst@CCPut \lst@ProcessOther {"2018}{"2019}{"201C}{"201D}\z@
\lst@AddToHookAtTop{SelectCharTable}{%
  \lst@ifactivechars
    \catcode"2018=\active \catcode"2019=\active
    \catcode"201C=\active \catcode"201D=\active
  \fi}
\makeatother
% </hllpdf:listing-smart-quotes>

% <hllpdf:listing-technical-symbols>
% v4.30: bounded extension for common technical punctuation in LTR listings.
% Without character-table registration these glyphs bypass the token buffer
% and can move before the preceding word even inside quoted strings.
% Keep the original four-quote block separate for its standing mutation test.
% This is source-faithful printing, not syntax validation or general Unicode.
\makeatletter
\lst@CCPut \lst@ProcessOther
  {"2013}{"2014}{"2026}{"2190}{"2192}{"2194}{"21D2}
  {"2212}{"221E}{"2260}{"2264}{"2265}\z@
\lst@AddToHookAtTop{SelectCharTable}{%
  \lst@ifactivechars
    \catcode"2013=\active \catcode"2014=\active \catcode"2026=\active
    \catcode"2190=\active \catcode"2192=\active \catcode"2194=\active
    \catcode"21D2=\active \catcode"2212=\active \catcode"221E=\active
    \catcode"2260=\active \catcode"2264=\active \catcode"2265=\active
  \fi}
\makeatother
% </hllpdf:listing-technical-symbols>


\usepackage[hidelinks,unicode]{hyperref}   % LAST, after colors. unicode = Hebrew bookmarks.
% v3.62: THE ISOLATION MACROS CORRUPTED EVERY PDF BOOKMARK THEY APPEARED IN.
% \en is defined (line ~155) as
%     \foreignlanguage{english}{\language=\csname l@nohyphenation\endcsname #1}
% and in a heading that body is written verbatim to the .toc and handed to
% hyperref's \pdfstringdef. hyperref strips the two CONTROL SEQUENCES (\language,
% \l@nohyphenation) but the plain `=` CHARACTER token between them is not a control
% sequence, so it survives into the PDF string: \chapter{שער \en{CNOT} והצמדה}
% produced the bookmark "שער =CNOT והצמדה". The printed TOC and the heading on the
% page are CORRECT — only the navigation pane and the viewer window title are wrong,
% which is why it went unnoticed. Measured on the shipped fleet: 70 corrupted
% bookmarks across 5 PDFs, from 177 unprotected \en-in-heading sites. Same corruption
% reached /Title and /Subject via the \hypersetup{pdftitle=...} the docs recommend.
% This hook tells hyperref to use the plain argument for all three isolation macros
% when building a PDF string. It does NOT touch the .toc (the printed TOC is
% byte-unchanged) and it does NOT mask the math signal — a heading with $...$ still
% logs "removing math shift", so BIDI-15 detection is intact.
% BONUS, measured: it also makes \code{} in an UNSTARRED heading COMPILE. BLD-10
% recorded that as needing \texorpdfstring or a starred heading; the real root cause
% is hyperref's \pdfstringdef scan of \code's \ifmmode...\else...\fi body — which is
% exactly why the identically-bodied \texttt survived (hyperref pre-neutralizes
% \texttt in PDF strings) and why \DeclareRobustCommand could never help.
% v3.64: \def\\{ } added. A manual line break in a heading WELDS the two words
% together in the bookmark: \chapter{שורה ראשונה\\שורה שנייה} produced the outline
% entry 'שורה ראשונהשורה שנייה' — "ראשונהשורה" is not a word. hyperref drops the
% \\ and puts nothing in its place, so the break that looks like a space on the
% page is a zero-width join in the navigation pane. The spaced form
% 'מילה אחת \\ מילה שנייה' escapes with a harmless DOUBLE space ('אחת  מילה'),
% which is why this never got noticed. Both are measured on a real outline.
% 0 live instances in the corpus — but SKILL.md and content-style.md both RECOMMEND
% a manual heading break, so this is a primed trap rather than a dormant one, and a
% one-token preamble fix retires it for good instead of asking authors to remember.
% Mapping \\ to a space fixes the WELD, which is the actual defect (a nonexistent
% word in the navigation pane). It does not normalise whitespace: measured after the
% fix, the welded form becomes 'שורה ראשונה שורה שנייה' (correct) and the already-
% spaced form goes from two spaces to three. That is cosmetic — viewers render it as
% a wider gap, and no word is corrupted. A per-heading \texorpdfstring is still the
% way to get an exactly-tuned bookmark string.
\pdfstringdefDisableCommands{\def\en#1{#1}\def\h#1{#1}\def\doclang#1{#1}\def\code#1{#1}\def\codeline#1{#1}\def\\{ }}
% v3.70: RESTORE \url LINE-BREAKING. url.sty freezes its break penalties at LOAD
% time (\mathchardef\UrlBreakPenalty=\binoppenalty), and v3.45's global
% \binoppenalty=10000 (inline-math atomicity) therefore priced EVERY breakpoint
% inside \Url at 10000 = never break. Measured: a long \url in Hebrew prose is
% one unbroken line, ~100pt overfull, invading the margin band the ink gates
% deliberately tolerate. With the two lines below: 0 overfull, the URL breaks
% after '/', the RTL split is CORRECT (fragment order forward), and each line
% gets its own link annotation exactly covering its fragment (measured).
% Confined to \Url's group — inline-math atomicity (TYP-7) is untouched.
\mathchardef\UrlBreakPenalty=200
\mathchardef\UrlBigBreakPenalty=100
\ifx\lakitmainlanguage\hllpdfenglish
  \hypersetup{pdflang=en}
\else
  \hypersetup{pdflang=he}
\fi  % /Lang is the document base locale (not per-span accessibility tagging).
% Set a per-document \hypersetup{pdftitle={...}} too — see SKILL.md.

% v3.52 (round-2 audit): two RTL bibliography/cross-ref defects, both invisible to
% every gate and both affecting the FIRST author to use the feature.
% (a) \@biblabel's default [#1] brackets are bidi-neutral and mirror to ]1[ under
% the RTL bibliography base — every \thebibliography reference list shows ]1[.
% LTR-isolating the label restores [1]. (The render gate's mirror rule fingerprints
% round parens )1(, so square brackets slipped through.)
\makeatletter
\renewcommand*{\@biblabel}[1]{\foreignlanguage{english}{[#1]}}
% v3.54: the SAME neutral-bracket trap on the INLINE side — a multi-key \cite{a,b}
% under the RTL base lays its European-number list right-to-left, so [1, 2] renders
% "[2,1]" and [1, 2, 3] renders "[3,2,1]" (verified by glyph geometry — the reader
% sees the sources in REVERSE order; a note form \cite[p.5]{a} also reorders to
% "[p. 5,1]"). LTR-isolate the whole bracket group (number list AND optional note),
% exactly like \@biblabel above: [1, 2] / [1, 2, 3] / [1, p. 5] all render correct.
% Placed AFTER hyperref so it wins; hyperref hyperlinks each number inside #1 (via
% \@citex), which this wrapper leaves intact. Mirrors the LaTeX kernel \@cite shape.
% v4.9.1: directional isolation must NOT switch the language here. Babel's
% safe-active citation wrapper passes its key list through \bbl@tempa;
% \foreignlanguage reuses that scratch macro before the deferred citation
% loop expands it. The old line therefore cited the literal key "english".
% If such a key existed, even the undefined-reference gate stayed green.
% \babelsublr isolates the bracket group without selecting a locale or
% changing that key list. Preserve native links, key order and optional notes.
\def\@cite#1#2{\babelsublr{[{#1\if@tempswa , #2\fi}]}}
\makeatother

% v4.11.0: a reference inside align's \text inherited HEBdisptag=true,
% so its enclosing hbox mirrored the brackets despite a clean compile.
% Isolate the complete printed reference without selecting a Babel language.
% Also neutralize the display-only normalfont direction reset locally. A
% direction-only patch inverted 12 to 21 inside equation in the measured probe.
% Plain maketag is local to the reference: mathtools' display-side wrapper
% otherwise globally consumes manual_tag while formatting a nested reference.
% Preserve each original command (including the selective auxiliary write),
% current tag forms, native links, and showonlyrefs on/off transitions.
\makeatletter
\NewCommandCopy{\HEB@nativeeqref}{\eqref}
\newcommand{\HEB@eqrefisolate}[1]{%
  \babelsublr{\HEBdisptagfalse\let\bbl@eqnodir\relax\let\maketag@@@\ORIGmaketag@@@#1}}
\protected\def\eqref#1{\HEB@eqrefisolate{\HEB@nativeeqref{#1}}}
\ExplSyntaxOn
\NewCommandCopy{\HEB@nativeMTeqref}{\MT_extended_eqref:n}
\cs_set_protected:Npn \MT_extended_eqref:n #1
  {\HEB@eqrefisolate{\HEB@nativeMTeqref{#1}}}
\ExplSyntaxOff
\makeatother

% v4.11.0: custom caption counters are not part of hyperref's English
% locale, and Hebrew has no upstream hyperref locale in this toolchain.
% Restore the appropriate names at every Babel language switch, including
% return from a nested inline island. Do not localize the destination keys.
\makeatletter
\newcommand{\HEB@refnames}{%
\def\chapterautorefname{פרק}%
\def\sectionautorefname{סעיף}%
\def\subsectionautorefname{תת־סעיף}%
\def\subsubsectionautorefname{תת־תת־סעיף}%
\def\paragraphautorefname{פסקה}%
% v4.31: hyperref uses a distinct counter name at the deepest standard level.
\def\subparagraphautorefname{תת־פסקה}%
\def\equationautorefname{משוואה}%
\def\figureautorefname{איור}%
\def\tableautorefname{טבלה}%
\def\partautorefname{חלק}%
\def\appendixautorefname{נספח}%
\def\footnoteautorefname{הערת שוליים}%
\def\itemautorefname{סעיף}%
\def\Itemautorefname{סעיף}%
\def\figautorefname{איור}%
\def\tblautorefname{טבלה}%
\def\pageautorefname{עמוד}% 
}
\addto\extrashebrew{\HEB@refnames}
\newcommand{\ENG@refnames}{%
  \def\subparagraphautorefname{subparagraph}%
  \def\chapterautorefname{chapter}%
  \def\sectionautorefname{section}%
  \def\subsectionautorefname{subsection}%
  \def\subsubsectionautorefname{subsubsection}%
  \def\paragraphautorefname{paragraph}%
  \def\equationautorefname{Equation}%
  \def\figureautorefname{Figure}%
  \def\tableautorefname{Table}%
  \def\partautorefname{Part}%
  \def\appendixautorefname{Appendix}%
  \def\footnoteautorefname{footnote}%
  \def\itemautorefname{item}%
  \def\figautorefname{Figure}%
  \def\tblautorefname{Table}%
  \def\Itemautorefname{item}%
  \def\pageautorefname{page}%
}
\addto\extrasenglish{\ENG@refnames}
\ifx\lakitmainlanguage\hllpdfenglish\ENG@refnames\else\HEB@refnames\fi
\makeatother

% ============================================================================
% Optional running head (v3.39): \lakithead{<book title>} — OFF by default.
% A document opts in with ONE line after \input{preamble.tex}:
%   \lakithead{מבנה המחשב}
% Every content page then carries the series' orientation line: chapter name
% at the reading edge, book title opposite, thin rule. (fancyhdr's [L]/[R]
% are PHYSICAL; under the RTL page direction [L] lands at the RIGHT/reading
% edge — measured, consistent across the corpus.) Packaged like the topic
% rail: always loaded, inert without the call. This centralizes the 8-line
% block that 8 lakits used to copy by hand — the same duplication vector
% that once let a lakit preamble go ~100 lines stale. Class-aware: report
% fills the mark from \chapter, article from \section. NB: a titlepage env
% self-empties; front matter that sets \pagestyle{empty} must restore with
% \pagestyle{fancy} afterwards (NOT plain — that discards the head).
% ============================================================================
\usepackage{fancyhdr}
% v3.61: \chapter* left a STALE running head — pages of an UNNUMBERED chapter
% carried the PREVIOUS chapter's title. report.cls routes \chapter* through
% \@schapter, which never calls \chaptermark, so \leftmark keeps whatever the last
% numbered chapter left. Measured: with \chapter{פרק ממוספר} then
% \chapter*{מבוא לא ממוספר}, page 4 (inside the starred chapter) printed
% "פרק ממוספר". Invisible on the chapter's FIRST page (\thispagestyle{plain}), so it
% only shows from the second page on — the pages an author is least likely to check.
% An unnumbered front-matter chapter (\chapter*{מבוא}, \chapter*{הקדמה}) is everyday
% Hebrew-academic structure, and the corpus already carries HAND-WRITTEN \markboth
% lines right after \chapter* — the exact per-document duplication \lakithead exists
% to retire. Patched INSIDE \lakithead so the macro stays inert when unused; a
% document's own later \markboth still wins, so those hand-patches do not regress.
\makeatletter
% <hllpdf:starred-heading-scope>
\newif\iflakin@starmarkpatched
\newcommand{\lakistarmark}{%
  \iflakin@starmarkpatched\else
  % Dispatch on the actual class counter, not hyperref's shared internals.
  \ifdefined\c@chapter
    \let\HEBorigschapter\@schapter
    % Record only after the native heading has reached its page. A mark
    % before the native routine can advertise the next heading too early.
    \def\@schapter##1{\HEBorigschapter{##1}\chaptermark{##1}}%
  \else
    % @ssect receives layout arguments, NOT the section name or level.
    % Capture the name in @startsection without introducing a TeX group:
    % grouping sectioning would break the native after-heading state.
    \def\HEB@sectionkind{}%
    \def\HEB@mainsection{section}%
    \let\HEBorigstartsection\@startsection
    \def\@startsection##1##2##3##4##5##6{%
      \def\HEB@sectionkind{##1}%
      \HEBorigstartsection{##1}{##2}{##3}{##4}{##5}{##6}}%
    \let\HEBorigssect\@ssect
    \def\@ssect##1##2##3##4##5{%
      \HEBorigssect{##1}{##2}{##3}{##4}{##5}%
      % Lower starred headings must not replace the article's section mark.
      \ifx\HEB@sectionkind\HEB@mainsection
        \sectionmark{##5}%
      \fi
      \ignorespaces}%
  \fi
  \global\lakin@starmarkpatchedtrue
  \fi}
% </hllpdf:starred-heading-scope>
% <hllpdf:running-language>
% Choose the wrapper when the mark is written, not when a later page ships.
% English punctuation must remain LTR even after returning to Hebrew prose.
\newcommand{\HEB@runningmark}[1]{%
  \iflanguage{english}{\markboth{\en{#1}}{}}{\markboth{\h{#1}}{}}}
% </hllpdf:running-language>
\makeatother
% v4.2.1 (NAV): every collection already names itself in \lakithead — feed the
% PDF /Title from the same declaration, so the navigation gate's
% PDF_REQUIRE_TITLE has a universal source and no fleet member ships
% title-less by missing a second declaration site. Measured before this: 14
% of 16 collections shipped /Title-empty while every one carried a perfectly
% good \lakithead string (plain Hebrew in all 16 — no math, no formatting, so
% the PDF-string path is clean). An explicit title always wins: set it with
% \lakitpdftitle{...} anywhere before \lakithead, or with a plain
% \hypersetup{pdftitle=...} — the fallback detects that via hyperref's
% \@pdftitle, probe-verified \@empty-equal until the first explicit set.
\makeatletter
\newif\iflakit@pdftitleset
\newcommand{\lakitpdftitle}[1]{%
  \hypersetup{pdftitle={#1}}\global\lakit@pdftitlesettrue}
\newcommand{\lakithead}[1]{%
  \iflakit@pdftitleset\else
    \ifx\@pdftitle\@empty
      \hypersetup{pdftitle={#1}}\global\lakit@pdftitlesettrue
    \fi
  \fi
  \pagestyle{fancy}\fancyhf{}%
  \lakistarmark

  \fancyhead[L]{\small\textcolor{cGray}{\leftmark}}%
  % <hllpdf:fixed-head-language>
  % The fixed document title has its OWN language snapshot, independent of
  % the current section/chapter mark. Choose the branch NOW, not at shipout:
  % the page may be Hebrew by the time an English title is finally printed.
  \iflanguage{english}{%
    \fancyhead[R]{\small\textcolor{cGray}{\en{#1}}}%
  }{%
    \fancyhead[R]{\small\textcolor{cGray}{\h{#1}}}%
  }%
  % </hllpdf:fixed-head-language>
  \fancyfoot[C]{\thepage}%
  \renewcommand{\headrulewidth}{0.4pt}%
  \ifdefined\chaptermark
    \renewcommand{\chaptermark}[1]{\HEB@runningmark{##1}}%
  \else
    \renewcommand{\sectionmark}[1]{\HEB@runningmark{##1}}%
  \fi}
\makeatother

% v3.39.1: path-independent byte-reproducibility. LuaTeX derives the trailer
% /ID from the OUTPUT PATH, so two clones of the same source at different
% directories produced different sha256 (measured: the ONLY delta was the
% 62-byte /ID line). Bit 512 suppresses exactly that line — measured: Creator/
% Producer/CreationDate/ModDate and all content untouched. /ID is optional
% per ISO 32000 (required only for ENCRYPTED PDFs and PDF/A — neither is a
% target here). If a future workflow needs an ID (PDF/A, Acrobat revision
% matching), delete this line and it returns.
\pdfvariable suppressoptionalinfo \numexpr 512\relax

% v3.75: BOLD MATH IN BOLD HEADINGS. Heading text is \bfseries but math inside
% it stayed regular-weight CM — a hairline $V_T$ inside a heavy 24.8pt chapter
% title (measured: CMMI12 fl=6 next to FrankRuhlLibre-Bold fl=20; 155 live
% heading-with-math sites fleet-wide). \boldmath rides along wherever the class
% sets \bfseries for a heading; A/B measured: $V_T$ -> CMMIB10 fl=22, 0 missing
% characters (the bold CM faces ship with this TeX Live). Guarded per-macro so
% the article class (no \@makechapterhead) and future class changes stay safe.
% v4.0.1 (adoption fix): the SEARCH patterns above carry no parameter token.
% v4.0.0 made them specific by including the class's own #1/#2, and measured on
% TeX Live 2023 that never matches: a \patchcmd search pattern containing a
% parameter token cannot match while the call sits inside \@ifundefined{...}{}{...}
% (probed both #1 and ##1 -> FAIL; the same pattern at top level succeeds). Every
% one of these patches therefore failed, and because v4.0.0 also replaced the
% empty failure branch with \PackageError, the failure became FATAL: all four
% bundled example assets stopped building ("Cannot patch chapter-title math
% weight -> Fatal error, no output PDF"). The \PackageError branches are KEPT —
% a silent no-op here is what let v3.75's fix rot unnoticed — only the parameter
% tokens are dropped, so the patches match and a real class change still errors.
\makeatletter
\@ifundefined{@makechapterhead}{}{%
  \patchcmd{\@makechapterhead}{\huge\bfseries \@chapapp}
    {\huge\bfseries\boldmath \@chapapp}{}
    {\PackageError{hebrew-lualatex-pdf}{Cannot patch chapter-number math weight}{Unsupported document class}}%
  \patchcmd{\@makechapterhead}{\Huge \bfseries}
    {\Huge \bfseries\boldmath}{}
    {\PackageError{hebrew-lualatex-pdf}{Cannot patch chapter-title math weight}{Unsupported document class}}%
  \patchcmd{\@makeschapterhead}{\Huge \bfseries}
    {\Huge \bfseries\boldmath}{}
    {\PackageError{hebrew-lualatex-pdf}{Cannot patch starred chapter math weight}{Unsupported document class}}}
% v4.0.1: \part sits under a THREE-DEEP wrapper chain, so the class-based anchor
% is not in \@part at all. Measured with \show: hyperref rewrites \@part to
% `...\phantomsection\fi \H@old@part`, and nameref rewrites THAT to
% `\NR@gettitle{#1}\NR@part[{#1}]{#2}` — the class body ends up in \NR@part.
% v4.0.0 patched \@part and made failure fatal, which is why all seven bundled
% examples stopped building. \part is NOT a dead heading level here: 188 live uses
% across the fleet, so the bold-math fix is worth getting right rather than
% dropping. We walk the chain innermost-first and stop at the first hit. Failure is
% a WARNING, not an error: the chain is a property of hyperref/nameref versions, so
% a future package release must degrade to regular-weight part math — a cosmetic
% loss — instead of refusing to build. Chapter and section keep \PackageError,
% because those targets are in the CLASS and were verified to patch here.
\newif\iflakin@partnum \newif\iflakin@parttitle \newif\iflakin@sparttitle
\def\lakin@partnum#1#2{\iflakin@partnum\else
  \patchcmd{#1}{#2\bfseries \partname}{#2\bfseries\boldmath \partname}
    {\lakin@partnumtrue}{}\fi}
\def\lakin@parttitle#1#2{\iflakin@parttitle\else
  \patchcmd{#1}{#2 \bfseries}{#2 \bfseries\boldmath}{\lakin@parttitletrue}{}\fi}
\def\lakin@sparttitle#1#2{\iflakin@sparttitle\else
  \patchcmd{#1}{#2 \bfseries}{#2 \bfseries\boldmath}{\lakin@sparttitletrue}{}\fi}
\@ifundefined{c@chapter}{\def\lakin@pnum{\Large}\def\lakin@ptit{\huge}}%
                        {\def\lakin@pnum{\huge}\def\lakin@ptit{\Huge}}
\@ifundefined{NR@part}{}{\expandafter\lakin@partnum\expandafter\NR@part\expandafter{\lakin@pnum}%
                         \expandafter\lakin@parttitle\expandafter\NR@part\expandafter{\lakin@ptit}}
\@ifundefined{H@old@part}{}{\expandafter\lakin@partnum\expandafter\H@old@part\expandafter{\lakin@pnum}%
                            \expandafter\lakin@parttitle\expandafter\H@old@part\expandafter{\lakin@ptit}}
\@ifundefined{@part}{}{\expandafter\lakin@partnum\expandafter\@part\expandafter{\lakin@pnum}%
                       \expandafter\lakin@parttitle\expandafter\@part\expandafter{\lakin@ptit}}
\@ifundefined{NR@spart}{}{\expandafter\lakin@sparttitle\expandafter\NR@spart\expandafter{\lakin@ptit}}
\@ifundefined{H@old@spart}{}{\expandafter\lakin@sparttitle\expandafter\H@old@spart\expandafter{\lakin@ptit}}
\@ifundefined{@spart}{}{\expandafter\lakin@sparttitle\expandafter\@spart\expandafter{\lakin@ptit}}
\iflakin@partnum\else\PackageWarning{hebrew-lualatex-pdf}{Part-number math stays
  regular weight: no known \noexpand\part wrapper matched (hyperref/nameref
  version change). Cosmetic only}\fi
\iflakin@parttitle\else\PackageWarning{hebrew-lualatex-pdf}{Part-title math stays
  regular weight: no known \noexpand\part wrapper matched. Cosmetic only}\fi
% v4.12.0: THE PART NUMBER IN THE TOC WAS NOT DIRECTIONALLY SEPARATED FROM THE
% TITLE. report.cls/article.cls compose the part entry as ONE text stream --
% \addcontentsline{toc}{part}{\thepart\hspace{1em}#1} -- while EVERY other level
% goes through \numberline, which puts the number in an \hbox. A box is a bidi
% OBJECT and cannot join a neighbouring run; a bare "IV" can. So in the RTL TOC
% paragraph a Latin number and a Latin-INITIAL title merge into a SINGLE LTR run
% and are laid out left-to-right as a unit: the number lands on the wrong side of
% the row, and a mixed title is scrambled outright. Measured on \part{MCP ...}:
% the printed order was "<tail> V MCP" -- title fragment, number, title fragment.
% The chapter/section rows one line below, with the SAME Latin-initial titles,
% were correct on the same page. That contrast is the whole diagnosis.
% \en{} IS NOT A WORKAROUND, measured both ways on one fixture: it fixes a
% PURE-Latin title (\part{\en{Hooks}} renders correctly) and does NOT fix a MIXED
% one (\part{\en{MCP} <hebrew>} stays scrambled). babel's language switch changes
% font and hyphenation, not the bidi run structure -- so the defect is not
% reachable from the source side at all, which is why it belongs here.
% Boxing the number is the entire fix and is a STRICT NO-OP where the bug never
% bit: a Hebrew-titled document compiles BYTE-IDENTICAL before and after (md5,
% 3 clean runs), because there the neutral glue already sat between an L run and
% an R run. The PDF outline is byte-identical too -- hyperref's \pdfstringdef
% neutralises \mbox, so bookmark strings keep their part numbers.
% \numberline{\thepart} would also fix the direction, but it re-geometries every
% part row in the fleet and DROPS the number from the part bookmark (hyperref
% strips \numberline). \mbox changes nothing that was already right.
% The search pattern deliberately carries no # token: etoolbox cannot match a
% parameter token in a patch pattern (measured -- all three chain members
% reported "not patchable" while the pattern included #1). "\thepart \hspace
% {1em}" is brace-balanced and occurs exactly once in the wrapper body.
\newif\iflakin@parttocnum
\def\lakin@parttocnum#1{\iflakin@parttocnum\else
  \patchcmd{#1}{\thepart \hspace {1em}}{\mbox {\thepart }\hspace {1em}}
    {\lakin@parttocnumtrue}{}\fi}
\@ifundefined{NR@part}{}{\lakin@parttocnum\NR@part}
\@ifundefined{H@old@part}{}{\lakin@parttocnum\H@old@part}
\@ifundefined{@part}{}{\lakin@parttocnum\@part}
\iflakin@parttocnum\else\PackageWarning{hebrew-lualatex-pdf}{Part TOC number is
  not boxed: no known \noexpand\part wrapper matched. A part title that BEGINS
  with Latin will print its number on the wrong side of its TOC row}\fi
\patchcmd{\section}{\bfseries}{\bfseries\boldmath}{}
  {\PackageError{hebrew-lualatex-pdf}{Cannot patch section math weight}{Unsupported document class}}
\patchcmd{\subsection}{\bfseries}{\bfseries\boldmath}{}
  {\PackageError{hebrew-lualatex-pdf}{Cannot patch subsection math weight}{Unsupported document class}}
\patchcmd{\subsubsection}{\bfseries}{\bfseries\boldmath}{}
  {\PackageError{hebrew-lualatex-pdf}{Cannot patch subsubsection math weight}{Unsupported document class}}
\patchcmd{\paragraph}{\bfseries}{\bfseries\boldmath}{}{}
\patchcmd{\subparagraph}{\bfseries}{\bfseries\boldmath}{}{}
\@ifundefined{l@chapter}{}{%
  \patchcmd{\l@chapter}{\bfseries}{\bfseries\boldmath}{}{}}
\makeatother
\setcounter{tocdepth}{\lakittocdepth}
% <hllpdf:measured-toc>
% v4.16.0: Measure actual number labels, not a guessed digit count. Read only
% canonical contentsline records; never execute titles or addtocontents code
% during the measuring pass. Widths are local to this individual list and are
% available before its first row: no per-row drift and no extra aux pass.
\makeatletter
\newread\hllpdf@tocstream
\newdimen\hllpdf@toc@chapter
\newdimen\hllpdf@toc@section
\newdimen\hllpdf@toc@subsection
\newdimen\hllpdf@toc@subsubsection
\newdimen\hllpdf@toc@paragraph
\newdimen\hllpdf@toc@subparagraph
\newdimen\hllpdf@toc@figure
\newdimen\hllpdf@toc@table
\newdimen\hllpdf@tocpage
\newdimen\hllpdf@toc@sectionindent
\newdimen\hllpdf@toc@subsectionindent
\newdimen\hllpdf@toc@subsubsectionindent
\newdimen\hllpdf@toc@paragraphindent
\newdimen\hllpdf@toc@subparagraphindent
\def\hllpdf@tocdefaults{%
  \hllpdf@toc@chapter=2em
  \hllpdf@toc@section=3.2em
  \hllpdf@toc@subsection=3.9em
  \@ifundefined{c@chapter}{%
    \hllpdf@toc@subsubsection=3.2em
    \hllpdf@toc@paragraph=4.1em
    \hllpdf@toc@subparagraph=5em
  }{%
    \hllpdf@toc@subsubsection=4.1em
    \hllpdf@toc@paragraph=5em
    \hllpdf@toc@subparagraph=6em
  }%
  \hllpdf@toc@figure=2.3em
  \hllpdf@toc@table=2.3em
  \hllpdf@tocpage=\@pnumwidth\relax
}
\def\hllpdf@tocindents{%
  % Preserve the established shallow layout when its slots already fit.
  \hllpdf@toc@sectionindent=1.5em
  \@tempdima=\hllpdf@toc@chapter \advance\@tempdima by -2em
  \advance\hllpdf@toc@sectionindent by \@tempdima
  \hllpdf@toc@subsectionindent=\dimexpr
    \hllpdf@toc@sectionindent+\hllpdf@toc@section\relax
  \hllpdf@toc@subsubsectionindent=\dimexpr
    \hllpdf@toc@subsectionindent+\hllpdf@toc@subsection\relax
  \hllpdf@toc@paragraphindent=\dimexpr
    \hllpdf@toc@subsubsectionindent+\hllpdf@toc@subsubsection\relax
  \hllpdf@toc@subparagraphindent=\dimexpr
    \hllpdf@toc@paragraphindent+\hllpdf@toc@paragraph\relax
}
\def\hllpdf@tocfont@chapter{\bfseries\boldmath}
\def\hllpdf@tocfont@part{\large\bfseries\boldmath}
\def\hllpdf@tocentryfont{%
  \ifcsname hllpdf@tocfont@\hllpdf@toclevel\endcsname
    \csname hllpdf@tocfont@\hllpdf@toclevel\endcsname
  \fi}
\long\def\hllpdf@tocdiscard#1\hllpdf@tocstop{}
\def\hllpdf@tocscanrecord{%
  \@ifnextchar\contentsline{\hllpdf@tocscanline}{%
    \@ifnextchar\babel@toc{\hllpdf@tocscanlanguage}{\hllpdf@tocdiscard}}}
\long\def\hllpdf@tocscanlanguage\babel@toc#1#2#3\hllpdf@tocstop{%
  % Babel records are declarations, not arbitrary code. Track the language
  % without executing the original directive or changing the surrounding list.
  \def\hllpdf@toclanguage{#1}}
\long\def\hllpdf@tocscanline\contentsline#1#2#3#4#5\hllpdf@tocstop{%
  \def\hllpdf@toclevel{#1}%
  \ifcsname hllpdf@toc@#1\endcsname
    \hllpdf@tocscantitle#2\hllpdf@tocstop
  \fi
  % The page-number column needs its own measurement (including bold chapter
  % rows). Its extra clearance separates the page label from dot leaders.
  \setbox\@tempboxa=\hbox{\foreignlanguage{\hllpdf@toclanguage}{\hllpdf@tocentryfont #3\kern.3em}}%
  \ifdim\wd\@tempboxa>\hllpdf@tocpage
    \hllpdf@tocpage=\wd\@tempboxa
  \fi}
\def\hllpdf@tocscantitle{%
  \@ifnextchar\numberline{\hllpdf@tocscanlabel}{\hllpdf@tocdiscard}}
\long\def\hllpdf@tocscanlabel\numberline#1#2\hllpdf@tocstop{%
  \setbox\@tempboxa=\hbox{\foreignlanguage{\hllpdf@toclanguage}{\hllpdf@tocentryfont #1\kern.5em}}%
  \ifdim\wd\@tempboxa>\csname hllpdf@toc@\hllpdf@toclevel\endcsname
    \csname hllpdf@toc@\hllpdf@toclevel\endcsname=\wd\@tempboxa
  \fi}
\def\hllpdf@tocread{%
  \ifeof\hllpdf@tocstream\else
    \read\hllpdf@tocstream to\hllpdf@tocrecord
    \expandafter\hllpdf@tocscanrecord\hllpdf@tocrecord\hllpdf@tocstop
    \expandafter\hllpdf@tocread
  \fi}
\def\hllpdf@tocmeasure#1{%
  \makeatletter
  \edef\hllpdf@toclanguage{\languagename}%
  \hllpdf@tocdefaults
  \openin\hllpdf@tocstream=\jobname.#1\relax
  \hllpdf@tocread
  \closein\hllpdf@tocstream
  \hllpdf@tocindents
  \edef\@pnumwidth{\the\hllpdf@tocpage}%
  \@tempdima=\dimexpr\hllpdf@tocpage+1em\relax
  \ifdim\@tempdima>\@tocrmarg\relax
    \edef\@tocrmarg{\the\@tempdima}%
  \fi}
\let\hllpdf@originalstarttoc\@starttoc
\def\@starttoc#1{%
  \begingroup
    \def\hllpdf@tocext{#1}%
    \def\hllpdf@tocwanted{toc}%
    \ifx\hllpdf@tocext\hllpdf@tocwanted \hllpdf@tocmeasure{#1}\else
      \def\hllpdf@tocwanted{lof}%
      \ifx\hllpdf@tocext\hllpdf@tocwanted \hllpdf@tocmeasure{#1}\else
        \def\hllpdf@tocwanted{lot}%
        \ifx\hllpdf@tocext\hllpdf@tocwanted \hllpdf@tocmeasure{#1}\fi
      \fi
    \fi
    \hllpdf@originalstarttoc{#1}%
  \endgroup}
\hllpdf@tocdefaults
\hllpdf@tocindents
\@ifundefined{l@chapter}{}{%
  \patchcmd{\l@chapter}{\setlength\@tempdima{1.5em}}%
    {\setlength\@tempdima{\hllpdf@toc@chapter}}{}%
    {\PackageError{hebrew-lualatex-pdf}{Cannot size chapter TOC numbers}%
       {Unsupported contents layout: the chapter number slot was not found}}}
\renewcommand*\l@section{\@dottedtocline{1}%
  {\hllpdf@toc@sectionindent}{\hllpdf@toc@section}}
\renewcommand*\l@subsection{\@dottedtocline{2}%
  {\hllpdf@toc@subsectionindent}{\hllpdf@toc@subsection}}
\renewcommand*\l@subsubsection{\@dottedtocline{3}%
  {\hllpdf@toc@subsubsectionindent}{\hllpdf@toc@subsubsection}}
\renewcommand*\l@paragraph{\@dottedtocline{4}%
  {\hllpdf@toc@paragraphindent}{\hllpdf@toc@paragraph}}
\renewcommand*\l@subparagraph{\@dottedtocline{5}%
  {\hllpdf@toc@subparagraphindent}{\hllpdf@toc@subparagraph}}
\renewcommand*\l@figure{\@dottedtocline{1}{1.5em}{\hllpdf@toc@figure}}
\renewcommand*\l@table{\@dottedtocline{1}{1.5em}{\hllpdf@toc@table}}
\makeatother
% </hllpdf:measured-toc>
\setcounter{secnumdepth}{\lakitsecnumdepth}
% Public preamble-only caption-numbering policy. The historical within-major
% default is untouched unless this command is called. "continuous" removes the
% reset edge as well as the printed prefix, so Figure 1/2 and Table 1/2 stay
% unique across sections/chapters. Returning to "within-major" restores the
% original class-aware displays and native destination conventions. The policy
% never replaces \figcap/\tabcap or their measured link-placement mechanism.
\makeatletter
\let\lakin@majorfig\thefig
\let\lakin@majortbl\thetbl
\let\lakin@majorHfig\theHfig
\let\lakin@majorHtbl\theHtbl
\newcommand{\lakitcaptionnumbering}[1]{%
  \edef\lakin@numbering{#1}%
  \def\lakin@continuous{continuous}\def\lakin@withinmajor{within-major}%
  \ifx\lakin@numbering\lakin@continuous
    \@ifundefined{c@chapter}{%
      \@removefromreset{fig}{section}\@removefromreset{tbl}{section}%
    }{\@removefromreset{fig}{chapter}\@removefromreset{tbl}{chapter}}%
    \renewcommand{\thefig}{\arabic{fig}}\renewcommand{\thetbl}{\arabic{tbl}}%
    \def\theHfig{\arabic{fig}}\def\theHtbl{\arabic{tbl}}%
  \else\ifx\lakin@numbering\lakin@withinmajor
    \@ifundefined{c@chapter}{%
      \@removefromreset{fig}{section}\@removefromreset{tbl}{section}%
      \@addtoreset{fig}{section}\@addtoreset{tbl}{section}%
    }{%
      \@removefromreset{fig}{chapter}\@removefromreset{tbl}{chapter}%
      \@addtoreset{fig}{chapter}\@addtoreset{tbl}{chapter}}%
    \let\thefig\lakin@majorfig\let\thetbl\lakin@majortbl
    \let\theHfig\lakin@majorHfig\let\theHtbl\lakin@majorHtbl
  \else
    \PackageError{hebrew-english-latex-studio}{Unknown caption-numbering policy '#1'}%
      {Use continuous or within-major before begin{document}.}%
  \fi\fi}
\@onlypreamble\lakitcaptionnumbering
\makeatother
% Architecture readiness sentinel. Optional runtime layers require the complete
% core, so the marker is defined only after every mechanism/baseline definition.
% <hllpdf:architecture-ready>
\def\hllpdfcoreloaded{1}
% </hllpdf:architecture-ready>

% v4.22: registered file resources, not an assertion about embedded subset bytes.
% The log framing avoids long-path wrapping and preserves UTF-8 as hex.
% <hllpdf:font-inventory>
\directlua{
  local function inventory()
    local names = {}
    local unnamed = 0
    for id, item in font.each() do
      local path = item.filename
      if path and not (path == "") then names[path] = true else unnamed = unnamed + 1 end
    end
    local sorted = {}
    for path in pairs(names) do table.insert(sorted,path) end
    table.sort(sorted)
    texio.write_nl("log","HLLFONT|BEGIN|1")
    for index,path in ipairs(sorted) do
      local hex = path:gsub(".",function(c) return string.format(string.char(37) .. "02x",string.byte(c)) end)
      local count = math.ceil(string.len(hex)/48)
      texio.write_nl("log","HLLFONT|FILE|" .. index .. "|" .. count)
      for offset=1,string.len(hex),48 do texio.write_nl("log","HLLFONT|HEX|" .. hex:sub(offset,offset+47)) end
    end
    texio.write_nl("log","HLLFONT|END|" .. rawlen(sorted) .. "|" .. unnamed)
    texio.write_nl("log","")
  end
  luatexbase.add_to_callback("finish_pdffile",inventory,"hllpdf.runtime.fonts")
}
% </hllpdf:font-inventory>
:
%   \lakithead{מבנה המחשב}
% Every content page then carries the series' orientation line: chapter name
% at the reading edge, book title opposite, thin rule. (fancyhdr's [L]/[R]
% are PHYSICAL; under the RTL page direction [L] lands at the RIGHT/reading
% edge — measured, consistent across the corpus.) Packaged like the topic
% rail: always loaded, inert without the call. This centralizes the 8-line
% block that 8 lakits used to copy by hand — the same duplication vector
% that once let a lakit preamble go ~100 lines stale. Class-aware: report
% fills the mark from \chapter, article from \section. NB: a titlepage env
% self-empties; front matter that sets \pagestyle{empty} must restore with
% \pagestyle{fancy} afterwards (NOT plain — that discards the head).
% ============================================================================
\usepackage{fancyhdr}
% v3.61: \chapter* left a STALE running head — pages of an UNNUMBERED chapter
% carried the PREVIOUS chapter's title. report.cls routes \chapter* through
% \@schapter, which never calls \chaptermark, so \leftmark keeps whatever the last
% numbered chapter left. Measured: with \chapter{פרק ממוספר} then
% \chapter*{מבוא לא ממוספר}, page 4 (inside the starred chapter) printed
% "פרק ממוספר". Invisible on the chapter's FIRST page (\thispagestyle{plain}), so it
% only shows from the second page on — the pages an author is least likely to check.
% An unnumbered front-matter chapter (\chapter*{מבוא}, \chapter*{הקדמה}) is everyday
% Hebrew-academic structure, and the corpus already carries HAND-WRITTEN \markboth
% lines right after \chapter* — the exact per-document duplication \lakithead exists
% to retire. Patched INSIDE \lakithead so the macro stays inert when unused; a
% document's own later \markboth still wins, so those hand-patches do not regress.
\makeatletter
% <hllpdf:starred-heading-scope>
\newif\iflakin@starmarkpatched
\newcommand{\lakistarmark}{%
  \iflakin@starmarkpatched\else
  % Dispatch on the actual class counter, not hyperref's shared internals.
  \ifdefined\c@chapter
    \let\HEBorigschapter\@schapter
    % Record only after the native heading has reached its page. A mark
    % before the native routine can advertise the next heading too early.
    \def\@schapter##1{\HEBorigschapter{##1}\chaptermark{##1}}%
  \else
    % @ssect receives layout arguments, NOT the section name or level.
    % Capture the name in @startsection without introducing a TeX group:
    % grouping sectioning would break the native after-heading state.
    \def\HEB@sectionkind{}%
    \def\HEB@mainsection{section}%
    \let\HEBorigstartsection\@startsection
    \def\@startsection##1##2##3##4##5##6{%
      \def\HEB@sectionkind{##1}%
      \HEBorigstartsection{##1}{##2}{##3}{##4}{##5}{##6}}%
    \let\HEBorigssect\@ssect
    \def\@ssect##1##2##3##4##5{%
      \HEBorigssect{##1}{##2}{##3}{##4}{##5}%
      % Lower starred headings must not replace the article's section mark.
      \ifx\HEB@sectionkind\HEB@mainsection
        \sectionmark{##5}%
      \fi
      \ignorespaces}%
  \fi
  \global\lakin@starmarkpatchedtrue
  \fi}
% </hllpdf:starred-heading-scope>
% <hllpdf:running-language>
% Choose the wrapper when the mark is written, not when a later page ships.
% English punctuation must remain LTR even after returning to Hebrew prose.
\newcommand{\HEB@runningmark}[1]{%
  \iflanguage{english}{\markboth{\en{#1}}{}}{\markboth{\h{#1}}{}}}
% </hllpdf:running-language>
\makeatother
% v4.2.1 (NAV): every collection already names itself in \lakithead — feed the
% PDF /Title from the same declaration, so the navigation gate's
% PDF_REQUIRE_TITLE has a universal source and no fleet member ships
% title-less by missing a second declaration site. Measured before this: 14
% of 16 collections shipped /Title-empty while every one carried a perfectly
% good \lakithead string (plain Hebrew in all 16 — no math, no formatting, so
% the PDF-string path is clean). An explicit title always wins: set it with
% \lakitpdftitle{...} anywhere before \lakithead, or with a plain
% \hypersetup{pdftitle=...} — the fallback detects that via hyperref's
% \@pdftitle, probe-verified \@empty-equal until the first explicit set.
\makeatletter
\newif\iflakit@pdftitleset
\newcommand{\lakitpdftitle}[1]{%
  \hypersetup{pdftitle={#1}}\global\lakit@pdftitlesettrue}
\newcommand{\lakithead}[1]{%
  \iflakit@pdftitleset\else
    \ifx\@pdftitle\@empty
      \hypersetup{pdftitle={#1}}\global\lakit@pdftitlesettrue
    \fi
  \fi
  \pagestyle{fancy}\fancyhf{}%
  \lakistarmark

  \fancyhead[L]{\small\textcolor{cGray}{\leftmark}}%
  % <hllpdf:fixed-head-language>
  % The fixed document title has its OWN language snapshot, independent of
  % the current section/chapter mark. Choose the branch NOW, not at shipout:
  % the page may be Hebrew by the time an English title is finally printed.
  \iflanguage{english}{%
    \fancyhead[R]{\small\textcolor{cGray}{\en{#1}}}%
  }{%
    \fancyhead[R]{\small\textcolor{cGray}{\h{#1}}}%
  }%
  % </hllpdf:fixed-head-language>
  \fancyfoot[C]{\thepage}%
  \renewcommand{\headrulewidth}{0.4pt}%
  \ifdefined\chaptermark
    \renewcommand{\chaptermark}[1]{\HEB@runningmark{##1}}%
  \else
    \renewcommand{\sectionmark}[1]{\HEB@runningmark{##1}}%
  \fi}
\makeatother

% v3.39.1: path-independent byte-reproducibility. LuaTeX derives the trailer
% /ID from the OUTPUT PATH, so two clones of the same source at different
% directories produced different sha256 (measured: the ONLY delta was the
% 62-byte /ID line). Bit 512 suppresses exactly that line — measured: Creator/
% Producer/CreationDate/ModDate and all content untouched. /ID is optional
% per ISO 32000 (required only for ENCRYPTED PDFs and PDF/A — neither is a
% target here). If a future workflow needs an ID (PDF/A, Acrobat revision
% matching), delete this line and it returns.
\pdfvariable suppressoptionalinfo \numexpr 512\relax

% v3.75: BOLD MATH IN BOLD HEADINGS. Heading text is \bfseries but math inside
% it stayed regular-weight CM — a hairline $V_T$ inside a heavy 24.8pt chapter
% title (measured: CMMI12 fl=6 next to FrankRuhlLibre-Bold fl=20; 155 live
% heading-with-math sites fleet-wide). \boldmath rides along wherever the class
% sets \bfseries for a heading; A/B measured: $V_T$ -> CMMIB10 fl=22, 0 missing
% characters (the bold CM faces ship with this TeX Live). Guarded per-macro so
% the article class (no \@makechapterhead) and future class changes stay safe.
% v4.0.1 (adoption fix): the SEARCH patterns above carry no parameter token.
% v4.0.0 made them specific by including the class's own #1/#2, and measured on
% TeX Live 2023 that never matches: a \patchcmd search pattern containing a
% parameter token cannot match while the call sits inside \@ifundefined{...}{}{...}
% (probed both #1 and ##1 -> FAIL; the same pattern at top level succeeds). Every
% one of these patches therefore failed, and because v4.0.0 also replaced the
% empty failure branch with \PackageError, the failure became FATAL: all four
% bundled example assets stopped building ("Cannot patch chapter-title math
% weight -> Fatal error, no output PDF"). The \PackageError branches are KEPT —
% a silent no-op here is what let v3.75's fix rot unnoticed — only the parameter
% tokens are dropped, so the patches match and a real class change still errors.
\makeatletter
\@ifundefined{@makechapterhead}{}{%
  \patchcmd{\@makechapterhead}{\huge\bfseries \@chapapp}
    {\huge\bfseries\boldmath \@chapapp}{}
    {\PackageError{hebrew-lualatex-pdf}{Cannot patch chapter-number math weight}{Unsupported document class}}%
  \patchcmd{\@makechapterhead}{\Huge \bfseries}
    {\Huge \bfseries\boldmath}{}
    {\PackageError{hebrew-lualatex-pdf}{Cannot patch chapter-title math weight}{Unsupported document class}}%
  \patchcmd{\@makeschapterhead}{\Huge \bfseries}
    {\Huge \bfseries\boldmath}{}
    {\PackageError{hebrew-lualatex-pdf}{Cannot patch starred chapter math weight}{Unsupported document class}}}
% v4.0.1: \part sits under a THREE-DEEP wrapper chain, so the class-based anchor
% is not in \@part at all. Measured with \show: hyperref rewrites \@part to
% `...\phantomsection\fi \H@old@part`, and nameref rewrites THAT to
% `\NR@gettitle{#1}\NR@part[{#1}]{#2}` — the class body ends up in \NR@part.
% v4.0.0 patched \@part and made failure fatal, which is why all seven bundled
% examples stopped building. \part is NOT a dead heading level here: 188 live uses
% across the fleet, so the bold-math fix is worth getting right rather than
% dropping. We walk the chain innermost-first and stop at the first hit. Failure is
% a WARNING, not an error: the chain is a property of hyperref/nameref versions, so
% a future package release must degrade to regular-weight part math — a cosmetic
% loss — instead of refusing to build. Chapter and section keep \PackageError,
% because those targets are in the CLASS and were verified to patch here.
\newif\iflakin@partnum \newif\iflakin@parttitle \newif\iflakin@sparttitle
\def\lakin@partnum#1#2{\iflakin@partnum\else
  \patchcmd{#1}{#2\bfseries \partname}{#2\bfseries\boldmath \partname}
    {\lakin@partnumtrue}{}\fi}
\def\lakin@parttitle#1#2{\iflakin@parttitle\else
  \patchcmd{#1}{#2 \bfseries}{#2 \bfseries\boldmath}{\lakin@parttitletrue}{}\fi}
\def\lakin@sparttitle#1#2{\iflakin@sparttitle\else
  \patchcmd{#1}{#2 \bfseries}{#2 \bfseries\boldmath}{\lakin@sparttitletrue}{}\fi}
\@ifundefined{c@chapter}{\def\lakin@pnum{\Large}\def\lakin@ptit{\huge}}%
                        {\def\lakin@pnum{\huge}\def\lakin@ptit{\Huge}}
\@ifundefined{NR@part}{}{\expandafter\lakin@partnum\expandafter\NR@part\expandafter{\lakin@pnum}%
                         \expandafter\lakin@parttitle\expandafter\NR@part\expandafter{\lakin@ptit}}
\@ifundefined{H@old@part}{}{\expandafter\lakin@partnum\expandafter\H@old@part\expandafter{\lakin@pnum}%
                            \expandafter\lakin@parttitle\expandafter\H@old@part\expandafter{\lakin@ptit}}
\@ifundefined{@part}{}{\expandafter\lakin@partnum\expandafter\@part\expandafter{\lakin@pnum}%
                       \expandafter\lakin@parttitle\expandafter\@part\expandafter{\lakin@ptit}}
\@ifundefined{NR@spart}{}{\expandafter\lakin@sparttitle\expandafter\NR@spart\expandafter{\lakin@ptit}}
\@ifundefined{H@old@spart}{}{\expandafter\lakin@sparttitle\expandafter\H@old@spart\expandafter{\lakin@ptit}}
\@ifundefined{@spart}{}{\expandafter\lakin@sparttitle\expandafter\@spart\expandafter{\lakin@ptit}}
\iflakin@partnum\else\PackageWarning{hebrew-lualatex-pdf}{Part-number math stays
  regular weight: no known \noexpand\part wrapper matched (hyperref/nameref
  version change). Cosmetic only}\fi
\iflakin@parttitle\else\PackageWarning{hebrew-lualatex-pdf}{Part-title math stays
  regular weight: no known \noexpand\part wrapper matched. Cosmetic only}\fi
% v4.12.0: THE PART NUMBER IN THE TOC WAS NOT DIRECTIONALLY SEPARATED FROM THE
% TITLE. report.cls/article.cls compose the part entry as ONE text stream --
% \addcontentsline{toc}{part}{\thepart\hspace{1em}#1} -- while EVERY other level
% goes through \numberline, which puts the number in an \hbox. A box is a bidi
% OBJECT and cannot join a neighbouring run; a bare "IV" can. So in the RTL TOC
% paragraph a Latin number and a Latin-INITIAL title merge into a SINGLE LTR run
% and are laid out left-to-right as a unit: the number lands on the wrong side of
% the row, and a mixed title is scrambled outright. Measured on \part{MCP ...}:
% the printed order was "<tail> V MCP" -- title fragment, number, title fragment.
% The chapter/section rows one line below, with the SAME Latin-initial titles,
% were correct on the same page. That contrast is the whole diagnosis.
% \en{} IS NOT A WORKAROUND, measured both ways on one fixture: it fixes a
% PURE-Latin title (\part{\en{Hooks}} renders correctly) and does NOT fix a MIXED
% one (\part{\en{MCP} <hebrew>} stays scrambled). babel's language switch changes
% font and hyphenation, not the bidi run structure -- so the defect is not
% reachable from the source side at all, which is why it belongs here.
% Boxing the number is the entire fix and is a STRICT NO-OP where the bug never
% bit: a Hebrew-titled document compiles BYTE-IDENTICAL before and after (md5,
% 3 clean runs), because there the neutral glue already sat between an L run and
% an R run. The PDF outline is byte-identical too -- hyperref's \pdfstringdef
% neutralises \mbox, so bookmark strings keep their part numbers.
% \numberline{\thepart} would also fix the direction, but it re-geometries every
% part row in the fleet and DROPS the number from the part bookmark (hyperref
% strips \numberline). \mbox changes nothing that was already right.
% The search pattern deliberately carries no # token: etoolbox cannot match a
% parameter token in a patch pattern (measured -- all three chain members
% reported "not patchable" while the pattern included #1). "\thepart \hspace
% {1em}" is brace-balanced and occurs exactly once in the wrapper body.
\newif\iflakin@parttocnum
\def\lakin@parttocnum#1{\iflakin@parttocnum\else
  \patchcmd{#1}{\thepart \hspace {1em}}{\mbox {\thepart }\hspace {1em}}
    {\lakin@parttocnumtrue}{}\fi}
\@ifundefined{NR@part}{}{\lakin@parttocnum\NR@part}
\@ifundefined{H@old@part}{}{\lakin@parttocnum\H@old@part}
\@ifundefined{@part}{}{\lakin@parttocnum\@part}
\iflakin@parttocnum\else\PackageWarning{hebrew-lualatex-pdf}{Part TOC number is
  not boxed: no known \noexpand\part wrapper matched. A part title that BEGINS
  with Latin will print its number on the wrong side of its TOC row}\fi
\patchcmd{\section}{\bfseries}{\bfseries\boldmath}{}
  {\PackageError{hebrew-lualatex-pdf}{Cannot patch section math weight}{Unsupported document class}}
\patchcmd{\subsection}{\bfseries}{\bfseries\boldmath}{}
  {\PackageError{hebrew-lualatex-pdf}{Cannot patch subsection math weight}{Unsupported document class}}
\patchcmd{\subsubsection}{\bfseries}{\bfseries\boldmath}{}
  {\PackageError{hebrew-lualatex-pdf}{Cannot patch subsubsection math weight}{Unsupported document class}}
\patchcmd{\paragraph}{\bfseries}{\bfseries\boldmath}{}{}
\patchcmd{\subparagraph}{\bfseries}{\bfseries\boldmath}{}{}
\@ifundefined{l@chapter}{}{%
  \patchcmd{\l@chapter}{\bfseries}{\bfseries\boldmath}{}{}}
\makeatother
\setcounter{tocdepth}{\lakittocdepth}
% <hllpdf:measured-toc>
% v4.16.0: Measure actual number labels, not a guessed digit count. Read only
% canonical contentsline records; never execute titles or addtocontents code
% during the measuring pass. Widths are local to this individual list and are
% available before its first row: no per-row drift and no extra aux pass.
\makeatletter
\newread\hllpdf@tocstream
\newdimen\hllpdf@toc@chapter
\newdimen\hllpdf@toc@section
\newdimen\hllpdf@toc@subsection
\newdimen\hllpdf@toc@subsubsection
\newdimen\hllpdf@toc@paragraph
\newdimen\hllpdf@toc@subparagraph
\newdimen\hllpdf@toc@figure
\newdimen\hllpdf@toc@table
\newdimen\hllpdf@tocpage
\newdimen\hllpdf@toc@sectionindent
\newdimen\hllpdf@toc@subsectionindent
\newdimen\hllpdf@toc@subsubsectionindent
\newdimen\hllpdf@toc@paragraphindent
\newdimen\hllpdf@toc@subparagraphindent
\def\hllpdf@tocdefaults{%
  \hllpdf@toc@chapter=2em
  \hllpdf@toc@section=3.2em
  \hllpdf@toc@subsection=3.9em
  \@ifundefined{c@chapter}{%
    \hllpdf@toc@subsubsection=3.2em
    \hllpdf@toc@paragraph=4.1em
    \hllpdf@toc@subparagraph=5em
  }{%
    \hllpdf@toc@subsubsection=4.1em
    \hllpdf@toc@paragraph=5em
    \hllpdf@toc@subparagraph=6em
  }%
  \hllpdf@toc@figure=2.3em
  \hllpdf@toc@table=2.3em
  \hllpdf@tocpage=\@pnumwidth\relax
}
\def\hllpdf@tocindents{%
  % Preserve the established shallow layout when its slots already fit.
  \hllpdf@toc@sectionindent=1.5em
  \@tempdima=\hllpdf@toc@chapter \advance\@tempdima by -2em
  \advance\hllpdf@toc@sectionindent by \@tempdima
  \hllpdf@toc@subsectionindent=\dimexpr
    \hllpdf@toc@sectionindent+\hllpdf@toc@section\relax
  \hllpdf@toc@subsubsectionindent=\dimexpr
    \hllpdf@toc@subsectionindent+\hllpdf@toc@subsection\relax
  \hllpdf@toc@paragraphindent=\dimexpr
    \hllpdf@toc@subsubsectionindent+\hllpdf@toc@subsubsection\relax
  \hllpdf@toc@subparagraphindent=\dimexpr
    \hllpdf@toc@paragraphindent+\hllpdf@toc@paragraph\relax
}
\def\hllpdf@tocfont@chapter{\bfseries\boldmath}
\def\hllpdf@tocfont@part{\large\bfseries\boldmath}
\def\hllpdf@tocentryfont{%
  \ifcsname hllpdf@tocfont@\hllpdf@toclevel\endcsname
    \csname hllpdf@tocfont@\hllpdf@toclevel\endcsname
  \fi}
\long\def\hllpdf@tocdiscard#1\hllpdf@tocstop{}
\def\hllpdf@tocscanrecord{%
  \@ifnextchar\contentsline{\hllpdf@tocscanline}{%
    \@ifnextchar\babel@toc{\hllpdf@tocscanlanguage}{\hllpdf@tocdiscard}}}
\long\def\hllpdf@tocscanlanguage\babel@toc#1#2#3\hllpdf@tocstop{%
  % Babel records are declarations, not arbitrary code. Track the language
  % without executing the original directive or changing the surrounding list.
  \def\hllpdf@toclanguage{#1}}
\long\def\hllpdf@tocscanline\contentsline#1#2#3#4#5\hllpdf@tocstop{%
  \def\hllpdf@toclevel{#1}%
  \ifcsname hllpdf@toc@#1\endcsname
    \hllpdf@tocscantitle#2\hllpdf@tocstop
  \fi
  % The page-number column needs its own measurement (including bold chapter
  % rows). Its extra clearance separates the page label from dot leaders.
  \setbox\@tempboxa=\hbox{\foreignlanguage{\hllpdf@toclanguage}{\hllpdf@tocentryfont #3\kern.3em}}%
  \ifdim\wd\@tempboxa>\hllpdf@tocpage
    \hllpdf@tocpage=\wd\@tempboxa
  \fi}
\def\hllpdf@tocscantitle{%
  \@ifnextchar\numberline{\hllpdf@tocscanlabel}{\hllpdf@tocdiscard}}
\long\def\hllpdf@tocscanlabel\numberline#1#2\hllpdf@tocstop{%
  \setbox\@tempboxa=\hbox{\foreignlanguage{\hllpdf@toclanguage}{\hllpdf@tocentryfont #1\kern.5em}}%
  \ifdim\wd\@tempboxa>\csname hllpdf@toc@\hllpdf@toclevel\endcsname
    \csname hllpdf@toc@\hllpdf@toclevel\endcsname=\wd\@tempboxa
  \fi}
\def\hllpdf@tocread{%
  \ifeof\hllpdf@tocstream\else
    \read\hllpdf@tocstream to\hllpdf@tocrecord
    \expandafter\hllpdf@tocscanrecord\hllpdf@tocrecord\hllpdf@tocstop
    \expandafter\hllpdf@tocread
  \fi}
\def\hllpdf@tocmeasure#1{%
  \makeatletter
  \edef\hllpdf@toclanguage{\languagename}%
  \hllpdf@tocdefaults
  \openin\hllpdf@tocstream=\jobname.#1\relax
  \hllpdf@tocread
  \closein\hllpdf@tocstream
  \hllpdf@tocindents
  \edef\@pnumwidth{\the\hllpdf@tocpage}%
  \@tempdima=\dimexpr\hllpdf@tocpage+1em\relax
  \ifdim\@tempdima>\@tocrmarg\relax
    \edef\@tocrmarg{\the\@tempdima}%
  \fi}
\let\hllpdf@originalstarttoc\@starttoc
\def\@starttoc#1{%
  \begingroup
    \def\hllpdf@tocext{#1}%
    \def\hllpdf@tocwanted{toc}%
    \ifx\hllpdf@tocext\hllpdf@tocwanted \hllpdf@tocmeasure{#1}\else
      \def\hllpdf@tocwanted{lof}%
      \ifx\hllpdf@tocext\hllpdf@tocwanted \hllpdf@tocmeasure{#1}\else
        \def\hllpdf@tocwanted{lot}%
        \ifx\hllpdf@tocext\hllpdf@tocwanted \hllpdf@tocmeasure{#1}\fi
      \fi
    \fi
    \hllpdf@originalstarttoc{#1}%
  \endgroup}
\hllpdf@tocdefaults
\hllpdf@tocindents
\@ifundefined{l@chapter}{}{%
  \patchcmd{\l@chapter}{\setlength\@tempdima{1.5em}}%
    {\setlength\@tempdima{\hllpdf@toc@chapter}}{}%
    {\PackageError{hebrew-lualatex-pdf}{Cannot size chapter TOC numbers}%
       {Unsupported contents layout: the chapter number slot was not found}}}
\renewcommand*\l@section{\@dottedtocline{1}%
  {\hllpdf@toc@sectionindent}{\hllpdf@toc@section}}
\renewcommand*\l@subsection{\@dottedtocline{2}%
  {\hllpdf@toc@subsectionindent}{\hllpdf@toc@subsection}}
\renewcommand*\l@subsubsection{\@dottedtocline{3}%
  {\hllpdf@toc@subsubsectionindent}{\hllpdf@toc@subsubsection}}
\renewcommand*\l@paragraph{\@dottedtocline{4}%
  {\hllpdf@toc@paragraphindent}{\hllpdf@toc@paragraph}}
\renewcommand*\l@subparagraph{\@dottedtocline{5}%
  {\hllpdf@toc@subparagraphindent}{\hllpdf@toc@subparagraph}}
\renewcommand*\l@figure{\@dottedtocline{1}{1.5em}{\hllpdf@toc@figure}}
\renewcommand*\l@table{\@dottedtocline{1}{1.5em}{\hllpdf@toc@table}}
\makeatother
% </hllpdf:measured-toc>
\setcounter{secnumdepth}{\lakitsecnumdepth}
% Public preamble-only caption-numbering policy. The historical within-major
% default is untouched unless this command is called. "continuous" removes the
% reset edge as well as the printed prefix, so Figure 1/2 and Table 1/2 stay
% unique across sections/chapters. Returning to "within-major" restores the
% original class-aware displays and native destination conventions. The policy
% never replaces \figcap/\tabcap or their measured link-placement mechanism.
\makeatletter
\let\lakin@majorfig\thefig
\let\lakin@majortbl\thetbl
\let\lakin@majorHfig\theHfig
\let\lakin@majorHtbl\theHtbl
\newcommand{\lakitcaptionnumbering}[1]{%
  \edef\lakin@numbering{#1}%
  \def\lakin@continuous{continuous}\def\lakin@withinmajor{within-major}%
  \ifx\lakin@numbering\lakin@continuous
    \@ifundefined{c@chapter}{%
      \@removefromreset{fig}{section}\@removefromreset{tbl}{section}%
    }{\@removefromreset{fig}{chapter}\@removefromreset{tbl}{chapter}}%
    \renewcommand{\thefig}{\arabic{fig}}\renewcommand{\thetbl}{\arabic{tbl}}%
    \def\theHfig{\arabic{fig}}\def\theHtbl{\arabic{tbl}}%
  \else\ifx\lakin@numbering\lakin@withinmajor
    \@ifundefined{c@chapter}{%
      \@removefromreset{fig}{section}\@removefromreset{tbl}{section}%
      \@addtoreset{fig}{section}\@addtoreset{tbl}{section}%
    }{%
      \@removefromreset{fig}{chapter}\@removefromreset{tbl}{chapter}%
      \@addtoreset{fig}{chapter}\@addtoreset{tbl}{chapter}}%
    \let\thefig\lakin@majorfig\let\thetbl\lakin@majortbl
    \let\theHfig\lakin@majorHfig\let\theHtbl\lakin@majorHtbl
  \else
    \PackageError{hebrew-english-latex-studio}{Unknown caption-numbering policy '#1'}%
      {Use continuous or within-major before begin{document}.}%
  \fi\fi}
\@onlypreamble\lakitcaptionnumbering
\makeatother
% Architecture readiness sentinel. Optional runtime layers require the complete
% core, so the marker is defined only after every mechanism/baseline definition.
% <hllpdf:architecture-ready>
\def\hllpdfcoreloaded{1}
% </hllpdf:architecture-ready>

% v4.22: registered file resources, not an assertion about embedded subset bytes.
% The log framing avoids long-path wrapping and preserves UTF-8 as hex.
% <hllpdf:font-inventory>
\directlua{
  local function inventory()
    local names = {}
    local unnamed = 0
    for id, item in font.each() do
      local path = item.filename
      if path and not (path == "") then names[path] = true else unnamed = unnamed + 1 end
    end
    local sorted = {}
    for path in pairs(names) do table.insert(sorted,path) end
    table.sort(sorted)
    texio.write_nl("log","HLLFONT|BEGIN|1")
    for index,path in ipairs(sorted) do
      local hex = path:gsub(".",function(c) return string.format(string.char(37) .. "02x",string.byte(c)) end)
      local count = math.ceil(string.len(hex)/48)
      texio.write_nl("log","HLLFONT|FILE|" .. index .. "|" .. count)
      for offset=1,string.len(hex),48 do texio.write_nl("log","HLLFONT|HEX|" .. hex:sub(offset,offset+47)) end
    end
    texio.write_nl("log","HLLFONT|END|" .. rawlen(sorted) .. "|" .. unnamed)
    texio.write_nl("log","")
  end
  luatexbase.add_to_callback("finish_pdffile",inventory,"hllpdf.runtime.fonts")
}
% </hllpdf:font-inventory>
; if omitted, the historical canonical
% defaults below are exact.  They deliberately control only the document shell
% and page layout.  RTL/math semantics stay in this preamble, while visual
% presentation themes (for example clean_theme.tex) load AFTER it.
%
%   \def\lakitdocumentclass{article}
%   \def\lakitclassoptions{11pt,a4paper}
%   \def\lakitgeometryoptions{margin=1.8cm,headheight=14pt}
%   \def\lakitlinespread{1.15}
%   \def\lakitparagraphstyle{block} % letters/briefs; indented for continuous prose
%   \def\lakittocdepth{1}
%   \def\lakitsecnumdepth{2}
%   % ============================================================================
% preamble.tex — tested LuaLaTeX preamble for Hebrew + math + figures PDFs (v4.32.0)
% Battle-tested: compiled an 80-page Hebrew quantum-computing summary AND a
% 197-page Hebrew computer-architecture anthology (TikZ diagrams, code listings,
% colored boxes) — both with 0 missing characters.
%
% Usage in your document:
%   % ============================================================================
% preamble.tex — tested LuaLaTeX preamble for Hebrew + math + figures PDFs (v4.32.0)
% Battle-tested: compiled an 80-page Hebrew quantum-computing summary AND a
% 197-page Hebrew computer-architecture anthology (TikZ diagrams, code listings,
% colored boxes) — both with 0 missing characters.
%
% Usage in your document:
%   \input{preamble.tex}     % (or paste these lines at the top)
%   \begin{document}  ...  \end{document}
%
% Compile with:  lualatex -interaction=nonstopmode doc.tex
%   Repeat until two consecutive PDFs are byte-identical (scripts/build.sh
%   does this) — a TOC document converges only on pass 3+; after two passes
%   the TOC page numbers are STALE (measured; see SKILL.md step 3).
% Font binaries are external dependencies, pinned by assets/fonts/SHA256SUMS.
% scripts/setup_fonts.sh imports a verified local installation or prior skill.
%
% PUBLIC DOCUMENT-SHELL / LAYOUT CONTRACT (introduced v4.3.1; architecture-owned v4.7.1). These hooks may be
% defined BEFORE \input{preamble.tex}; if omitted, the historical canonical
% defaults below are exact.  They deliberately control only the document shell
% and page layout.  RTL/math semantics stay in this preamble, while visual
% presentation themes (for example clean_theme.tex) load AFTER it.
%
%   \def\lakitdocumentclass{article}
%   \def\lakitclassoptions{11pt,a4paper}
%   \def\lakitgeometryoptions{margin=1.8cm,headheight=14pt}
%   \def\lakitlinespread{1.15}
%   \def\lakitparagraphstyle{block} % letters/briefs; indented for continuous prose
%   \def\lakittocdepth{1}
%   \def\lakitsecnumdepth{2}
%   \input{preamble.tex}
%
% Supported class contract: report or article.  Alternate paper sizes must also
% be declared to build.sh through EXPECTED_PAPER; the PDF gate never infers them.
% Do NOT use these hooks as a theme mechanism: theme overlays must not own page
% geometry, class selection, or direction semantics. A declared presentation-policy
% overlay may intentionally change numbering visibility after the preamble.
% ============================================================================
% Runtime architecture guard: generated from scripts/architecture_contract.py.
% The readiness sentinel remains delayed until the END of this file.
% <hllpdf:architecture-guard>
\ifdefined\hllpdfcoreloaded
  \errmessage{hllpdf: preamble.tex loaded more than once}
  \expandafter\endinput
\fi
% </hllpdf:architecture-guard>

\ifdefined\lakitdocumentclass\else\def\lakitdocumentclass{report}\fi
\ifdefined\lakitclassoptions\else\def\lakitclassoptions{12pt,a4paper}\fi
\ifdefined\lakitgeometryoptions\else\def\lakitgeometryoptions{margin=2.3cm,headheight=14pt}\fi
\ifdefined\lakitlinespread\else\def\lakitlinespread{1.3}\fi
\ifdefined\lakittocdepth\else\def\lakittocdepth{1}\fi
\ifdefined\lakitsecnumdepth\else\def\lakitsecnumdepth{2}\fi
% Paragraph organization belongs to the document layout, independently of its
% language and theme. Preserve historical class dimensions unless explicitly
% selecting block paragraphs; never suppress indentation with an everypar hack.
\ifdefined\lakitparagraphstyle\else\def\lakitparagraphstyle{indented}\fi
\edef\lakitparagraphstyle{\lakitparagraphstyle}
\def\hllpdfparagraphblock{block}
\def\hllpdfparagraphindented{indented}
\ifx\lakitparagraphstyle\hllpdfparagraphblock\else
  \ifx\lakitparagraphstyle\hllpdfparagraphindented\else
    \errmessage{hllpdf: paragraph style must be block or indented}
  \fi
\fi
% v4.3.1: the public class hook has a deliberately narrow support contract.
% Before this guard, \def\lakitdocumentclass{book} silently compiled even though
% every heading/TOC regression test targets report/article; an unsupported shell
% could therefore look "supported" until a subtle class-specific defect shipped.
% Fail before \documentclass rather than letting an undocumented combination enter
% the mechanism core. Plain-TeX comparison is used because LaTeX packages are not
% available yet at this point.
\edef\lakitvalidatedclass{\lakitdocumentclass}
\def\lakitreportclass{report}
\def\lakitarticleclass{article}
\ifx\lakitvalidatedclass\lakitreportclass
\else\ifx\lakitvalidatedclass\lakitarticleclass
\else
  \errmessage{hllpdf: unsupported class \lakitdocumentclass; use report or article}
\fi\fi
\let\lakitvalidatedclass\relax
\let\lakitreportclass\relax
\let\lakitarticleclass\relax
% Language intent is independent of document class, theme and content domain.
% The two dominant-language profiles differ in authoring intent, not direction.
% A balanced bilingual document must explicitly choose a navigation/base locale.
% <hllpdf:language-profile>
\providecommand{\lakitlanguageprofile}{hebrew}
\edef\lakitlanguageprofile{\lakitlanguageprofile}
\def\hllpdfhebrew{hebrew}\def\hllpdfenglish{english}
\def\hllpdfhebrewmixed{hebrew-with-english}
\def\hllpdfenglishmixed{english-with-hebrew}
\def\hllpdfbilingual{bilingual}
\let\hllpdfexpectedlanguage\relax
\ifx\lakitlanguageprofile\hllpdfhebrew
  \let\hllpdfexpectedlanguage\hllpdfhebrew
\else\ifx\lakitlanguageprofile\hllpdfhebrewmixed
  \let\hllpdfexpectedlanguage\hllpdfhebrew
\else\ifx\lakitlanguageprofile\hllpdfenglish
  \let\hllpdfexpectedlanguage\hllpdfenglish
\else\ifx\lakitlanguageprofile\hllpdfenglishmixed
  \let\hllpdfexpectedlanguage\hllpdfenglish
\else\ifx\lakitlanguageprofile\hllpdfbilingual
  \ifdefined\lakitmainlanguage\else
    \errmessage{hllpdf: bilingual profile requires explicit lakitmainlanguage}
  \fi
\else
  \errmessage{hllpdf: unsupported language profile \lakitlanguageprofile}
\fi\fi\fi\fi\fi
\ifdefined\lakitmainlanguage
  \edef\lakitmainlanguage{\lakitmainlanguage}
\else
  \let\lakitmainlanguage\hllpdfexpectedlanguage
\fi
\ifx\lakitmainlanguage\hllpdfhebrew\else
  \ifx\lakitmainlanguage\hllpdfenglish\else
    \errmessage{hllpdf: main language must be hebrew or english}
  \fi
\fi
\ifx\hllpdfexpectedlanguage\relax\else
  \ifx\lakitmainlanguage\hllpdfexpectedlanguage\else
    \errmessage{hllpdf: main language conflicts with language profile}
  \fi
\fi
\typeout{HLLLANG|\lakitlanguageprofile|\lakitmainlanguage}
% </hllpdf:language-profile>
\documentclass[\lakitclassoptions]{\lakitdocumentclass}
\usepackage[no-math]{fontspec}        % (1) ★ no-math — THE critical line (see recipe.md)
% layout=lists asks Babel to adapt list parshape when paragraph and page
% directions differ. Without it, minority-language list bodies lose their
% indentation and labels hang outside the text block, despite a clean log.
% The upstream mechanism preserves ordinary same-direction list geometry.
\usepackage[bidi=basic,layout=lists]{babel}
\ifx\lakitmainlanguage\hllpdfenglish
  \babelprovide[import]{hebrew}
  \babelprovide[import,main]{english}
\else
  \babelprovide[import,main]{hebrew}
  \babelprovide[import]{english}
\fi
% v4.22: bind canonical families and fallbacks to the selected installation.
% Lookup by basename or family alone is not proof of the file actually used.
% <hllpdf:font-binding>
\directlua{
  local dir = os.getenv("HEBREW_LATEX_FONT_DIR")
  if not dir or dir == "" then
    local data = os.getenv("XDG_DATA_HOME")
    if not data or data == "" or not (data:sub(1,1) == "/") then
      data = (os.getenv("HOME") or "") .. "/.local/share"
    end
    dir = data .. "/fonts/hebrew-lualatex-pdf"
  end
  if not (dir:sub(1,1) == "/") then dir = lfs.currentdir() .. "/" .. dir end
  if dir:find("[",1,true) or dir:find("]",1,true) or dir:find(string.char(0),1,true)
     or dir:find(string.char(10),1,true) or dir:find(string.char(13),1,true) then
    tex.error("hllpdf: unsupported character in font directory")
  end
  if not (lfs.attributes(dir,"mode") == "directory") then
    tex.error("hllpdf: font directory is missing")
  end
  hllpdf_font_dir = dir:gsub("/+$", "") .. "/"
  tex.sprint(string.char(92) .. "def" .. string.char(92) .. "hllpdffontdir{")
  tex.sprint(-2,hllpdf_font_dir)
  tex.sprint("}")
  luaotfload.add_fallback("hebfb",{"[" .. hllpdf_font_dir .. "DejaVuSerif.ttf]:mode=node;"})
  luaotfload.add_fallback("hebfbbf",{"[" .. hllpdf_font_dir .. "DejaVuSerif-Bold.ttf]:mode=node;"})
}
% </hllpdf:font-binding>
% v3.22: Frank Ruhl Libre ships NO italic face. Instead of silently-upright
% \emph (v<=3.20) or bold-as-emphasis (v3.21), synthesize a slanted face with
% FakeSlant — real, VISIBLE italic emphasis with correct \emph nesting
% semantics (inner \emph toggles back upright). Verified empirically at
% 300dpi: slant clearly visible, 0 errors, 0 missing glyphs.
\babelfont{rm}[Path={\hllpdffontdir},RawFeature={fallback=hebfb},
  BoldFont={FrankRuhlLibre-Bold.ttf},
  BoldFeatures={RawFeature={fallback=hebfbbf}},   % v3.56 — see hebfbbf above
  ItalicFont={FrankRuhlLibre-Regular.ttf}, ItalicFeatures={FakeSlant=0.22},
  % Bold-italic uses the weight-matched fallback and a synthetic slant.
  BoldItalicFont={FrankRuhlLibre-Bold.ttf}, BoldItalicFeatures={RawFeature={fallback=hebfbbf},FakeSlant=0.22}
  ]{FrankRuhlLibre-Regular.ttf}
% Sans fallbacks preserve family and weight for the few missing codepoints.
\directlua{luaotfload.add_fallback("hebfbsans",{"[" .. hllpdf_font_dir .. "DejaVuSans.ttf]:mode=node;"})}
\directlua{luaotfload.add_fallback("hebfbsansbf",{"[" .. hllpdf_font_dir .. "DejaVuSans-Bold.ttf]:mode=node;"})}
\babelfont{sf}[Path={\hllpdffontdir},RawFeature={fallback=hebfbsans},
  BoldFont={Heebo-Bold.ttf},
  BoldFeatures={RawFeature={fallback=hebfbsansbf}}]{Heebo-Regular.ttf}
% v3.24: Scale=MatchLowercase — DejaVu Sans Mono has an outsized x-height, so
% unscaled inline code visually SHOUTS next to the delicate Frank Ruhl body
% (measured on a 197-page build: scaled, a code atom sits in the line at the
% body's weight, and the document tightened by ~5%). Applies to listings too.
\babelfont{tt}[Path={\hllpdffontdir},RawFeature={fallback=hebfb},Scale=MatchLowercase,
  BoldFont={DejaVuSansMono.ttf},
  BoldFeatures={RawFeature={fallback=hebfb},FakeBold=1.5},
  ItalicFont={DejaVuSansMono.ttf},
  ItalicFeatures={RawFeature={fallback=hebfb},FakeSlant=0.2},
  BoldItalicFont={DejaVuSansMono.ttf},
  BoldItalicFeatures={RawFeature={fallback=hebfb},FakeBold=1.5,FakeSlant=0.2}
  ]{DejaVuSansMono.ttf}

% Explicit multilingual/English profiles use real Latin faces and italics.
% The legacy/default Hebrew profile retains its established measured typography.
% Latin Modern is provided by the TeX distribution, not the pinned font bundle;
% recorder/runtime evidence reports it as an additional toolchain resource.
\ifx\lakitlanguageprofile\hllpdfhebrew\else
  \babelfont[english]{rm}[Ligatures=TeX,RawFeature={fallback=hebfb},
    BoldFeatures={RawFeature={fallback=hebfbbf}},
    BoldItalicFeatures={RawFeature={fallback=hebfbbf}}]{Latin Modern Roman}
  \babelfont[english]{sf}[Ligatures=TeX,RawFeature={fallback=hebfbsans},
    BoldFeatures={RawFeature={fallback=hebfbsansbf}},
    BoldItalicFeatures={RawFeature={fallback=hebfbsansbf}}]{Latin Modern Sans}
\fi

\usepackage{amsmath,amssymb,mathtools}
% amsthm supplies theorem/proof machinery. Environment names remain
% document-specific; the generic proof label is localized below.
\usepackage{amsthm}
% Proof names follow the surrounding prose, including return from an island.
\addto\extrashebrew{\renewcommand{\proofname}{הוכחה}}
\addto\extrasenglish{\renewcommand{\proofname}{Proof}}
\ifx\lakitmainlanguage\hllpdfenglish
  \renewcommand{\proofname}{Proof}
\else
  \renewcommand{\proofname}{הוכחה}
\fi
% v3.24.1: equation tags under babel-bidi. A tag re-emitted by mathtools'
% showonlyrefs inside align inherits the RTL base and its parens MIRROR —
% ")2.1(" instead of "(2.1)" (plain align/equation tags are unaffected;
% verified on a minimal case + a 198-page build). LTR-isolate at the tag-BOX
% level: mathtools' show/hide + numbering logic is untouched. Do NOT override
% \tagform@ instead — that clobbers the showonlyrefs mechanism (measured:
% every tag shows, numbering breaks).
% v3.28.1: the v3.24.1 wrap must apply ONLY to tags of display environments.
% Inline \eqref in Hebrew PROSE also goes through \maketag@@@, and there the
% english-wrap MIRRORS the parens — ")1(" — while the unwrapped original is
% correct (proven on a 2x2 isolation matrix: patch x showonlyrefs; regressed
% every prose \eqref since v3.24.1, caught by the render gate's R4). The two
% call sites are indistinguishable by \ifmmode (tag boxes are text mode), so
% the display environments themselves raise a flag; \eqref in prose sees it
% false and gets the original untouched \maketag@@@.
\usepackage{etoolbox}

% babel bidi=basic (including upstream 26.9) assumes every automatic
% discretionary has a one-glyph replacement. LuaTeX can legally produce one
% with replace=nil after a hyphenated Latin value; Babel then indexes nil and the
% whole build dies. Do the narrow upstream-compatible correction: while Babel
% classifies the list, present only that malformed disc as an ordinary
% discretionary, then restore its subtype. Babel still reorders the ENTIRE box;
% there is no pcall and no successful fallback to an unreordered node list.
% If Babel throws for any other reason, the build remains fatal. The synthetic
% nil-disc probe and the real table gallery in test_all.sh pin both layers.
\AtBeginDocument{\directlua{
  if Babel and Babel.bidi and not Babel.heb_nil_disc_compat then
    Babel.heb_nil_disc_compat = true
    local original = Babel.bidi
    local disc_id = node.id('disc')
    Babel.bidi = function(head, ispar, hdir)
      local hidden = {}
      for item in node.traverse_id(disc_id, head) do
        if item.subtype == 2 and (not item.replace or not item.replace.char) then
          table.insert(hidden, item)
          item.subtype = 0
        end
      end
      local result = original(head, ispar, hdir)
      for _, item in ipairs(hidden) do item.subtype = 2 end
      return result
    end
  end
}}
\makeatletter
\newif\ifHEBdisptag
\let\ORIGmaketag@@@\maketag@@@
\def\maketag@@@#1{\ifHEBdisptag\ORIGmaketag@@@{\foreignlanguage{english}{#1}}\else\ORIGmaketag@@@{#1}\fi}
\@for\HEB@env:=equation,equation*,align,align*,gather,gather*,multline,multline*,flalign,flalign*,alignat,alignat*,eqnarray,eqnarray*\do{%
  \expandafter\AtBeginEnvironment\expandafter{\HEB@env}{\HEBdisptagtrue}}
\makeatother
% English page + Hebrew prose: Babel's equation-only tag flip uses a primitive
% leqno path that shifts the formula to the left edge and the tag into the body
% of an English page. Align's full-width alignment is unaffected. Give only this
% display a physical LTR container and native amsmath eqno/leqno. The engine then
% owns real tag width, collision handling and labels; no proxy tag box hides ink
% from its fit calculation. Babel's saved \bbl@eqnodir still restores Hebrew text
% inside math, and the prose locale/math direction are never switched.
% Close the prior RTL paragraph before normalization. After the environment's
% group restores its directions, close TeX's empty resumed display paragraph and
% use native LaTeX continuation handling: same-paragraph prose has no indent,
% while an authored blank line starts an ordinary indented RTL paragraph.
% Register after Babel's setup; remain local to equation(*) in a main-English
% document's actual RTL prose. Native Hebrew, ordinary English and explicit
% fleqn retain their original paths. No page/display dimensions are assigned.
% <hllpdf:minority-rtl-equation-geometry>
\makeatletter
\newcommand{\HEB@minorityprosecontinuation}{\par\@doendpe}
\newcommand{\HEB@minorityequationgeometry}{%
  \ifx\lakitmainlanguage\hllpdfenglish
    \ifnum\pardirection=\@ne
      \if@fleqn\else
        \par\aftergroup\HEB@minorityprosecontinuation
        \iftagsleft@\let\veqno\@@leqno
        \else\let\veqno\@@eqno\fi
        \bodydir TLT\pardir TLT\textdir TLT\relax
      \fi
    \fi
  \fi}
\AtBeginDocument{%
  \AddToHook{env/equation/begin}{\HEB@minorityequationgeometry}%
  \AddToHook{env/equation*/begin}{\HEB@minorityequationgeometry}}
\makeatother
% </hllpdf:minority-rtl-equation-geometry>
% v3.63: \shortintertext rendered Hebrew MIRROR-REVERSED and flush LEFT — an entire
% commentary line inside an align printed backwards. babel's bidi=basic patches
% \intertext@ only (luababel.def), and mathtools' \shortintertext uses the identical
% \vbox{\normalbaselines … \noindent\ignorespaces#1\par} skeleton but is never
% patched, so its paragraph keeps pardir/textdir LTR: no bidi reordering (logical
% order on the page) and flush-left. Measured against \intertext with the same
% Applying babel's own replacement to \shortintertext@ makes the two identical.
% \h{}/\mbox fix the letter order but leave the line flush LEFT, so they are not
% substitutes. Guarded + \AtBeginDocument because mathtools \let-copies at load time.
% NB \ExplSyntaxOn is REQUIRED: mathtools' internals carry expl3 names, and without
% the expl3 catcodes the `_` is read as a subscript and the patch detonates far away
% with "! Illegal unit of measure (pt inserted)" at the \end{align} (measured).
\makeatletter
\AtBeginDocument{%
  \@ifundefined{shortintertext@}{}{%
    \ExplSyntaxOn
    \bbl@sreplace\shortintertext@{\normalbaselines}%
      {\normalbaselines\ifx\bbl@eqnodir\relax\else\bbl@pardir\@ne\bbl@eqnodir\fi}%
    \ExplSyntaxOff}}
\makeatother
% geometry inherits the canonical A4 sheet from \documentclass; build.sh then
% verifies every rendered MediaBox and CropBox, so a silent Letter fallback is
% impossible to publish as a clean default build.
\usepackage[\lakitgeometryoptions]{geometry}
\linespread{\lakitlinespread}       % layout hook; canonical default remains 1.3
\ifx\lakitparagraphstyle\hllpdfparagraphblock
  \setlength{\parindent}{0pt}
  \setlength{\parskip}{6pt plus 1pt minus 1pt}
\fi
% These dimensions remain ordinary scoped TeX lengths: authors can customize
% them after core input. Language switches inherit them; lists keep their own
% indentation and narrow reference layouts establish their local paragraph gaps.
\typeout{HLLPAR|\lakitparagraphstyle}
\usepackage{xcolor}
\usepackage{mdframed}
\usepackage{listings}                 % code listings (assembly / C / pseudo-code)
\usepackage{tikz}
\usetikzlibrary{arrows.meta,positioning,calc,shapes.geometric,fit,backgrounds}
\usepackage{pgfplots}\pgfplotsset{compat=1.18}
\usepackage{quantikz}                 % quantum circuits (remove if unused)
\usepackage{booktabs}
\usepackage{array}
\usepackage{float}                    % [H] = place table exactly here (reference doc: no drifting)
\usepackage{enumitem}
\usepackage{needspace}

% ============================================================================
% ★★ v3.4 — THE DIRECTION-ISOLATION TRIO (the heart of mixing Hebrew + Latin)
% Mental model:
%   \h{...}    Hebrew (RTL) island      — Hebrew text inside an LTR context
%                                          (TikZ nodes, the english `tikzpic`).
%   \en{...}   English (LTR) isolate    — Latin text / parentheses inside a
%                                          Hebrew (RTL) context. KEEPS PARENS
%                                          AND ORDER CORRECT.
%   \code{...} \en + monospace          — code inside Hebrew / inside a frametitle.
% CRITICAL: never wrap Latin in \h{} (it is for HEBREW). Inside a TikZ node,
% Latin in \h{} renders REVERSED (Datapath->htapataD, CPU->UPC) and parens
% mirror. Isolate every Latin/paren chunk with \en{} or \code{}. See
% references/figures-boxes-listings.md.
% ---- TEXT INSIDE MATH (v3.8) ----------------------------------------------
% Put ANY text that sits inside math in a \text{}. The ONLY thing that is
% invalid in math is a bare \h/\en (= \foreignlanguage) used DIRECTLY, outside
% \text -> "\rmfamily invalid in math mode". Inside \text:
%   Hebrew label           -> \text{...}            (bare Hebrew renders fine;
%                              \h is optional — BUT NOT ALWAYS HARMLESS (v3.67):
%                              if the \h group OPENS with a nested island
%                              ($..$, \en{}, \code{}, \texttt{}, \mbox{}) the
%                              WHOLE group loses bidi and the Hebrew prints in
%                              LOGICAL order = gibberish. Measured:
%                                \text{\h{$f$ קבועה}}  -> reads "העובק"  BROKEN
%                                \text{\h{קבועה $f$}}  -> correct
%                                \text{$f$ קבועה}      -> correct (no \h)
%                              Hebrew first inside \h, or drop the \h. Flagged
%                              H-OPENS-WITH-LTR-GROUP.)
%   plain Latin word        -> \text{...}            (no marker needed)
%   Latin WITH ( ) or [ ]   -> \text{\en{...}}       (\en forces LTR; a plain
%                              \text mirrors the parens in RTL -> "Spec (TNR(" )
%   f(Hebrew/text arg)      -> \text{f}(\text{...})  -- keep the parens at MATH
%                              level, NOT \text{f(arg)} (parens inside the box mirror)
%   code in math            -> \mathtt{...}  (or \texttt{}, math-guarded below)
% NEVER write \en{...} or \h{...} straight into $...$ / \[..\] (outside \text);
% always via \text{}. (\mbox{...} also works but does not scale in scripts.)
% ============================================================================
\newcommand{\h}[1]{\foreignlanguage{hebrew}{#1}}
% v3.42.1 — \en islands do NOT hyphenate (auto OR at explicit hyphens). English hyphenation patterns were
% tearing identifiers and terms ACROSS THE RTL SEAM ("IO Regis-/ters",
% "In-/put", "Band-/width", "Per-/formance", "Multi-/ple" — all shipped on
% real pages; probe-proven: "ALUS-/rcB", "Mem-/ToReg"). An island is a label
% or a term, not English prose — a hyphen tear there is never right. A/B on
% a full 197pp lakit: identical overfull sets, same page count, ~11
% paragraphs healed, one mild underfull (badness 3525) that replaced a
% mutilated "Regis-/ters". Long English PROSE belongs in the englishprose
% environment below, which restores explicit-hyphen breaks while its paragraph
% is still active. A raw otherlanguage block inherits the global prohibition.
% (\l@nohyphenation needs \csname:
% '@' is not a letter here.)
% v3.43.2 — \exhyphenpenalty moved OUT of \en to a GLOBAL assignment below.
% Render-proven: a group-local \exhyphenpenalty is DEAD CODE — TeX reads the
% line-breaking parameters at \par time, after the \en group has closed, so
% \en{write-through} / plain not-taken still tore at the "-" ("write-/
% through" shipped on a real exam page). \language is different: it is
% recorded in the node list as text is added, so the in-group nohyphenation
% assignment above DOES work (and stays). The global penalty was A/B-measured
% on a 197pp lakit: identical overfull set (25), identical page count, the
% not-/taken tear healed. Hebrew is unaffected in practice — Hebrew compounds
% use maqaf (־), not ASCII "-".
\newcommand{\en}[1]{\foreignlanguage{english}{\language=\csname l@nohyphenation\endcsname #1}}
% Babel runs extras locally for inline islands and paragraph language changes.
% An English main document must not inherit the Hebrew no-hyphen paragraph rule.
\addto\extrashebrew{\exhyphenpenalty=10000\relax}
\addto\extrasenglish{\exhyphenpenalty=50\relax}
\ifx\lakitmainlanguage\hllpdfenglish
  \exhyphenpenalty=50\relax
\else
  \exhyphenpenalty=10000\relax
\fi
% Full English paragraphs need ordinary explicit-hyphen breaks. Use this
% environment instead of a raw otherlanguage block; the paragraph ends while
% the local penalty is still active.
\newenvironment{englishprose}
  {\par\begin{otherlanguage}{english}\exhyphenpenalty=50\relax}
  {\par\end{otherlanguage}}
\newenvironment{hebrewprose}
  {\par\begin{otherlanguage}{hebrew}\exhyphenpenalty=10000\relax}
  {\par\end{otherlanguage}}
% Use for shell text whose locale must not depend on the last body paragraph.
\DeclareRobustCommand{\doclang}[1]{\foreignlanguage{\lakitmainlanguage}{#1}}

% ============================================================================
% ★★ — fixes for defects that are INVISIBLE in the log and in the checks
% and only surface on a VISUAL RENDER of the PDF. (Each verified empirically.)
% ----------------------------------------------------------------------------
% v3.22: \emph is now REAL slanted emphasis (FakeSlant italic face defined on
% the \babelfont{rm} line above), replacing the v3.21 \emph->\textbf hack: it
% preserves \emph's toggle semantics and keeps bold free for its own role.
% Prefer bold-as-emphasis for your doc? one line: \renewcommand{\emph}[1]{\textbf{#1}}
% page number: force the digits LTR so they never reverse (18 -> 81) in the RTL
%   footer. That reversal is a NON-DETERMINISTIC bidi glitch (observed once in
%   ~80 pp); this LuaLaTeX build is not byte-reproducible, so the defensive
%   LTR wrap is the only reliable guard. Latin digits are always LTR anyway.
\renewcommand{\thepage}{\foreignlanguage{english}{\arabic{page}}}
% v3.70: \pagenumbering \gdef's \thepage and silently DESTROYS the wrap above —
% after \pagenumbering{arabic} the folio guard is gone for the whole rest of the
% document. Two collections discovered this independently and hand-restore the
% wrap in main.tex with the same comment; that knowledge now lives here instead.
% The wrapper re-applies the LTR wrap after the kernel call, for any style.
\makeatletter
\let\lakin@orig@pagenumbering\pagenumbering
\renewcommand{\pagenumbering}[1]{%
  \lakin@orig@pagenumbering{#1}%
  \renewcommand{\thepage}{\foreignlanguage{english}{\csname @#1\endcsname\c@page}}}
\makeatother
% emergencystretch: Hebrew does not hyphenate, so an unbreakable LTR island
%   (\en{...} / an inline math run) landing at the margin can leave a small
%   (~2-13pt) overfull with no legal breakpoint. This rescues ONLY those
%   inline-prose tie-breaks; it is a ceiling used only on offending lines and
%   does NOT touch display-math / listing / table overflow (those still surface
%   in the log for the content-level overflow pass). Verified: baseline overfull
%   -> 0 with this set, with no change to pagination or underfull.
\setlength{\emergencystretch}{3em}
% v3.40: inline math (almost) never breaks across lines. TeX's default
%   (binop 700 / rel 500) lets an inline run split at an operator —
%   $(1.f)\cdot 2^{e}$ shipped as "(1.f)·" ending one line and "2^e" opening
%   the next (reader-reported). In LTR that split is merely ugly; under an
%   RTL base the carried operator lands on the visually WRONG side of the
%   break and the formula reads backwards.
%   WHY 10000 (v3.45.0 — the display-first resolution). Inline math is
%   FORBIDDEN from breaking, so a torn island under an RTL base is
%   structurally impossible. The v3.40 cycle used 9999 ("break only as a last
%   resort") because a bare 10000 left ~30-39 WIDE inline equations
%   overflowing (worst 90.7pt) — but 9999's "last resort" fires OFTEN in
%   tight, non-hyphenating Hebrew: a reader found torn islands across the
%   fleet (~150 candidates; e.g. "P(y |" ending one line and "x_0)" opening
%   the next, read backwards). The step the 9999 cycle missed: every one of
%   those ~30 overflows is a wide equation that BELONGS in display. Convert
%   them to display FIRST, then 10000 gives 0 tears AND 0 overflows (measured,
%   12/12 fleet). Doctrine: inline math is ATOMIC; a wide equation goes in a
%   display, never inline. See content-style ("inline math is atomic").
\binoppenalty=10000 \relpenalty=10000
% v3.38: uniform sentence spacing. TeX's default (\nonfrenchspacing, sfcode
% 3000) is the OLD Anglo-American wide-space-after-period convention — Hebrew
% typography never had it, and modern Latin practice (Chicago et al.) dropped
% it too. Measured on real pages: post-period gaps up to 1.9-2.1x the line's
% mean word gap — the dominant "rivers" source in justified Hebrew prose.
% Survives babel language switches (probed: sfcode stays french inside \en{}
% islands and otherlanguage blocks).
\frenchspacing
% v3.38: widow/orphan protection. LaTeX's default penalties (150/150/50)
% allow a lone line at a page top/bottom almost freely. 10000 = forbid;
% this is SAFE here precisely because the document is \raggedbottom (the
% demoted line just moves on, no vertical stretch fight) — under flushbottom
% these values would cause underfull pages. raggedbottom is the report-class
% default; made explicit so a future class/option change cannot silently
% flip it and turn the 10000s harmful.
\raggedbottom
\widowpenalty=10000 \clubpenalty=10000 \displaywidowpenalty=10000
% em-dash: --- does NOT become — under Frank Ruhl (the font lacks tlig; verified
%   the width is identical with/without Ligatures=TeX, while the mechanism works
%   on Latin Modern). TYPE THE CHARACTER — (U+2014) DIRECTLY; likewise – (U+2013).
%   check_bidi_figures.py flags a --/--- that is adjacent to Hebrew.
% ============================================================================

% ---- STRUCTURAL DIRECTION-SAFETY for \texttt (v3.7) -------------------------
% A parenthesis, bracket, or $ INSIDE \texttt{} mirrors/migrates in RTL prose
% (MIPS 0($s0)->"(s0$)0", $rs->"rs$", arr[i]->"[arr[i") because \texttt sets the
% font but NOT the direction. Redefining \texttt to LTR-isolate its content fixes
% the WHOLE class at once (parens, brackets, $) so plain \texttt{lw $t0,0($s0)}
% is safe — no need to remember \code{}. Guarded by \ifmmode because
% \foreignlanguage is invalid in math (there, plain monospace is already LTR).
% NB (corrected v3.56 — the old note here said "stays tofu", which is WRONG):
% \texttt{Hebrew} is a FATAL build failure. The tt family is DejaVu Sans Mono
% (zero Hebrew) and the tt fallback hebfb is DejaVuSerif (also zero Hebrew), so
% the subsetter gets an empty glyph set and LuaTeX dies with
%   "! error: (file .../DejaVuSansMono.ttf) (type 2): there are no glyphs in the
%    subset"  -> no output PDF.
% The message names a FONT FILE, never Hebrew/\texttt/the source line. And it is
% non-local: add an unrelated Latin \code{} elsewhere and the fatal degrades into
% a SILENT DELETION of the Hebrew word. Put Hebrew in \h{}/normal text, never in
% \texttt{}/\code{}; flagged HEBREW-IN-MONO. To disable, delete this block.
\let\origtexttt\texttt
\protected\long\def\texttt#1{\ifmmode\origtexttt{#1}\else\foreignlanguage{english}{\origtexttt{#1}}\fi}
% v4.8.0: defer the direction test until TeX has inserted an alignment cell's
% math template. Unprotected expansion here tested \ifmmode too early at the
% FIRST token of aligned/array/cases/matrix cells. Engine-level \protected also
% preserves the original \texttt implementation without the self-recursion
% observed when retrofitting \DeclareRobustCommand to its saved alias.
\protected\long\def\code#1{\ifmmode\origtexttt{#1}\else\foreignlanguage{english}{\origtexttt{#1}}\fi}

% ---- \codeline — display-code line (v3.24): the code analogue of display math.
% Reader-friction rule (content-style.md §5b): inline \code{} is for a SINGLE
% ATOM only (one register/variable/mnemonic, a tight no-space expression).
% A multi-token expression — an instruction with operands, a spaced assignment,
% a register-binding list — zigzags the eye RTL<->LTR mid-sentence and, being
% an unbreakable island, line-breaks badly. Give it its own centered line with
% display-math vertical rhythm. \centerline = one box => the expression never
% wraps; \nopagebreak keeps it glued to its introducing sentence.
%   ...the Hebrew sentence ends with a colon:
%   \codeline{lw \$t1, -4(\$t0)}
%   and the prose continues (escape dollars as \$).
% v3.66 (first line of the body): flush a PENDING \item label before the \par.
% LaTeX defers the label box (\@labels) until the next \everypar, and an \hbox in
% vertical mode never fires \everypar — so `\item \codeline{...}` printed the code
% at y=69.1 and the item's own "1." BELOW it at y=103.6 (the number appears after
% its content, attached to nothing). \if@inlabel\leavevmode\fi is CONDITIONAL on
% purpose: an unconditional \leavevmode would open an empty paragraph line when
% \codeline follows ordinary prose. Verified both ways — in a list the label now
% sits on the code's line at its normal x, and after prose the geometry is
% byte-identical to before.
\makeatletter
\newcommand{\codeline}[1]{%
  \if@inlabel\leavevmode\fi
  \par\nopagebreak
  \vspace{\abovedisplayskip}\nopagebreak
  % v3.62: was \centerline{...}, i.e. \hbox to\hsize{\hss #1 \hss}. \hss has
  % INFINITE SHRINK, so TeX could always "fit" the box and NEVER reported an
  % Overfull — the same \hss mechanism as BLD-27's \item label. Measured: a wide
  % \codeline in a defbox painted its first token ON the box's left border bar and
  % sat 9.5pt outside the frame interior on both sides, with 0 Overfull, 0 Underfull
  % and all four gates PASS; in bare prose it reached 21.0pt into BOTH margins.
  % \hfil has no shrink, so the overflow is reported and build.sh's ovf_vis/ovf_big
  % counters see it; \linewidth (not \hsize) makes it box-aware, so a \codeline
  % inside a colored box centres on the box interior. Verified: a FITTING \codeline
  % keeps byte-identical glyph positions, and the box-crossing case now emits
  % "Overfull \hbox (19.03pt too wide)". A wide \codeline should be split in two.
  \hbox to\linewidth{\hfil\code{#1}\hfil}%
  \nopagebreak\vspace{\belowdisplayskip}}
\makeatother

% ---- a thin decorative rule helper ----
\newcommand{\sep}{\par\vspace{2pt}\noindent\textcolor{cRule}{\rule{\linewidth}{0.6pt}}\par\vspace{4pt}}

% ---- colors ----
\definecolor{cBlue}{HTML}{1F4E79}   \definecolor{cBlueBg}{HTML}{EAF1F8}
\definecolor{cGreen}{HTML}{2E6B4F}  \definecolor{cGreenBg}{HTML}{EAF4EE}
\definecolor{cAmber}{HTML}{9C6A00}  \definecolor{cAmberBg}{HTML}{FBF3E2}
\definecolor{cPurple}{HTML}{5B2C6F} \definecolor{cPurpleBg}{HTML}{F3ECF7}
\definecolor{cTeal}{HTML}{0E6E6E}   \definecolor{cTealBg}{HTML}{E6F4F4}
\definecolor{cRed}{HTML}{8E2B2B}    \definecolor{cRedBg}{HTML}{FBEAEA}
\definecolor{cGray}{HTML}{555555}   \definecolor{cRule}{HTML}{1F4E79}

% ---- math operators ----
\DeclareMathOperator{\Tr}{Tr}
\DeclareMathOperator{\rank}{rank}
\DeclareMathOperator{\diag}{diag}
\DeclareMathOperator{\sgn}{sgn}

% ---- bra-ket (quantikz defines \ket/\bra/\braket; add the rest safely) ----
\providecommand{\ket}[1]{\left|#1\right\rangle}
\providecommand{\bra}[1]{\left\langle#1\right|}
\providecommand{\braket}[2]{\left\langle#1\middle|#2\right\rangle}
\newcommand{\dyad}[1]{\left|#1\right\rangle\!\left\langle#1\right|}
% v3.21: general outer product |a><b| and expectation value <A> — agents reach
% for these by reflex (esp. \expval); quantikz defines neither, so their absence
% threw "Undefined control sequence" in a real build. \providecommand = safe.
\providecommand{\ketbra}[2]{\left|#1\right\rangle\!\left\langle#2\right|}
\providecommand{\expval}[1]{\left\langle#1\right\rangle}
\newcommand{\abs}[1]{\left|#1\right|}
\newcommand{\norm}[1]{\left\lVert#1\right\rVert}
\newcommand{\half}{\tfrac12}
\newcommand{\R}{\mathbb{R}}\newcommand{\C}{\mathbb{C}}\newcommand{\Z}{\mathbb{Z}}
% identity operator. NB: \mathbb{1} silently renders a WRONG glyph (msbm has no
% blackboard digits). Use \mathbb{I} = 𝕀 (works with amssymb, no extra package).
% Do NOT "fix" this back to \mathbb{1}.
\newcommand{\id}{\mathbb{I}}

% ---- physical units (v3.35): \punit — upright unit with thin space ---------
% Usage: $300\punit{K}$, $80\punit{mV/decade}$, $10^{7}\punit{cm/s}$ — the
% NUMBER AND UNIT LIVE IN ONE MATH ISLAND. Promoted verbatim from a 199-page
% production lakit (advanced-electronic-devices, dozens of uses).
% Two measured traps this design avoids (both render-proven):
%   1. siunitx SCRAMBLES under RTL — \SI{3.5}{eV} renders "eV 3.5", exponents
%      detach ("Hz 10 x 10^9"): its quantity assembly is not bidi-isolated.
%      Do NOT load siunitx in this pipeline.
%   2. A text-mode unit next to a bare-prose number is TWO adjacent LTR runs
%      and reorders ("5 mm" -> "mm5" — the BIDI-2/3 class). \punit is therefore
%      math-only ON PURPOSE: prose misuse fails LOUDLY at compile (\mathrm
%      outside math) instead of scrambling silently.
% Domain shorthands (\MeV, \fm...) stay per-document: \newcommand{\MeV}{\,\text{\en{MeV}}}
\providecommand{\punit}[1]{\,\mathrm{#1}}

% ============================================================================
% Colored boxes — mdframed (NOT tcolorbox, which breaks under bidi).
% ★★ v3.4: innerleftmargin/innerrightmargin give code listings room so their
% frame/line-numbers do not spill out the box border (see listing styles below).
% ============================================================================
% v3.61: everyline/split skips were applied to `tablebox` only (BLD-14, v3.37) —
% the six colored boxes below never got them, so they still split OPEN across a page
% break: the page-N segment gets top+sides and NO bottom rule, the page-N+1 segment
% sides+bottom and NO top rule. Measured on the shipped fleet this is a COVERAGE GAP,
% not a re-report of BLD-14: 0 split-open `tablebox` segments (that fix holds) vs 62
% split-open COLORED boxes across 7 documents. Verified after the fix: both segments
% close, and the frametitle is NOT duplicated on the continuation. BLD-14's caveat
% still applies — a box taller than a full page grinds mdframed's split-trial loop, so
% split such content manually.
% ============================================================================
% v3.72: frametitle={} leads EVERY \newmdenv option list below. mdframed options
% are GROUP-scoped: a box nested inside \begin{exbox}[frametitle={...}] used to
% INHERIT the outer title — measured, the identical title string rendered twice,
% the second riding a title-bar fill inside the INNER box's frame, for all seven
% environments in both directions. Both natural forms were broken (= preamble
% bug, not an authoring rule). The explicit reset makes every box titleless
% unless ITS OWN [frametitle={...}] says otherwise; verified both directions:
% no duplicate on the nested box, an explicit title still renders.
% Geometry is shared with the optional semantic-width table container. Themes
% select their own existing semantic style below; neither path fixes columns.
\mdfdefinestyle{hllpdfsemanticgeometry}{linewidth=1.1pt,
  innertopmargin=6pt,innerbottommargin=7pt,skipabove=9pt,skipbelow=9pt,
  innerleftmargin=9pt,innerrightmargin=9pt,everyline=true,
  splitbottomskip=6pt,splittopskip=12pt,roundcorner=3pt}
\newcommand{\hllpdfsemanticframestyle}{hllpdfsemanticgeometry}
\newmdenv[frametitle={},style=hllpdfsemanticgeometry,linecolor=cBlue,backgroundcolor=cBlueBg,
  frametitlebackgroundcolor=cBlue,frametitlefontcolor=white,
  ]{defbox}   % הגדרה
\newmdenv[frametitle={},style=hllpdfsemanticgeometry,linecolor=cGreen,backgroundcolor=cGreenBg,
  frametitlebackgroundcolor=cGreen,frametitlefontcolor=white,
  ]{thmbox}   % משפט/טענה
\newmdenv[frametitle={},style=hllpdfsemanticgeometry,linecolor=cAmber,backgroundcolor=cAmberBg,
  frametitlebackgroundcolor=cAmber,frametitlefontcolor=white,
  ]{notebox}  % הערה
\newmdenv[frametitle={},style=hllpdfsemanticgeometry,linecolor=cPurple,backgroundcolor=cPurpleBg,
  frametitlebackgroundcolor=cPurple,frametitlefontcolor=white,
  ]{exbox}    % דוגמה
\newmdenv[frametitle={},style=hllpdfsemanticgeometry,linecolor=cRed,backgroundcolor=cRedBg,
  frametitlebackgroundcolor=cRed,frametitlefontcolor=white,
  ]{warnbox}  % טעות נפוצה
\newmdenv[frametitle={},style=hllpdfsemanticgeometry,linecolor=cTeal,backgroundcolor=cTealBg,
  frametitlebackgroundcolor=cTeal,frametitlefontcolor=white,
  ]{keybox}   % רעיון מרכזי

% ---- codebox (v4.5.0): the dedicated container for code -------------------
% An OPTION, not a policy: asmblock unboxed stays canonical for a full listing,
% and a snippet inside any of the six colored boxes still works as before.
% What makes codebox worth its own environment is that it is the ONE framed box
% that keeps its margin line numbers.
%
% v4.4.0 shipped it WITHOUT them, on the strength of v3.37's "no margin knob can
% save them" plus four fresh probes that agreed. That was wrong: every probe had
% moved the code, the box, or both — never the key that carries the NUMBER
% COLUMN. Measured on the page:
%   resetmargins=false   the listing follows THIS box instead of the page,
%                        moving the code from x=75.26 to x=97.54
%   xrightmargin=26pt    under RTL this is the PHYSICAL LEFT indent, and the
%                        only key that carries the \llap'ed digits with it:
%                        x=64.28 -> 90.18, clearing the rule at 65.64 by 23.6pt
%   xleftmargin=26pt     moves the code and LEAVES THE DIGITS on the rule —
%                        the asymmetry that hid the fix for eleven versions
% resetmargins is hoisted out of the four styles (see the \lstset below them) so
% this override is possible at all; every other listing still resolves to true.
% Pinned by tests/codebox_numbers.tex + probe_lstnum.py --inside-rule, whose
% negative control strips xrightmargin and watches the digits fall back.
\definecolor{cSlate}{HTML}{3F4A56}  \definecolor{cSlateBg}{HTML}{F2F4F7}
\newmdenv[frametitle={},leftline=true,rightline=true,topline=true,bottomline=true,
  linecolor=cSlate,backgroundcolor=cSlateBg,linewidth=0.9pt,
  frametitlebackgroundcolor=cSlate,frametitlefontcolor=white,
  innertopmargin=6pt,innerbottommargin=7pt,skipabove=9pt,skipbelow=9pt,
  innerleftmargin=9pt,innerrightmargin=9pt,everyline=true,splitbottomskip=6pt,splittopskip=12pt,
  roundcorner=3pt]{codebox}   % קוד
% The two keys that make in-box numbering work, both measured on the page:
%   resetmargins=false  the listing follows THIS box instead of the page
%   xrightmargin        under RTL this is the PHYSICAL LEFT indent, and it is
%                       the only one that carries the \llap'ed number column
%                       with it. xleftmargin moves the code and leaves the
%                       digits behind on the rule (measured: code 97.54,
%                       digits 64.28 — the v3.37 defect, reproduced); with
%                       xrightmargin=26pt the digits land at x=90.18, a
%                       clearance of 23.6pt from the rule at 65.64.
% This asymmetry is why v3.4's note had the semantics backwards and v3.37
% concluded the digits could not be saved.
% <hllpdf:listing-container-state>
% v4.29: all state is local to the current environment group. A semantic
% callout masks number text, while a nested codebox restores the formatter
% saved before that mask. Save it once across successive semantic wrappers;
% an explicit numbers=none (or an explicitly empty formatter) remains valid.
\makeatletter
\newif\ifhllpdf@incodebox
\newif\ifhllpdf@englishcodebox
\newif\ifhllpdf@lstnumbermasked
\let\hllpdf@savedlstnumber\thelstnumber
\makeatother
% </hllpdf:listing-container-state>
% <hllpdf:codebox-gutter>
\makeatletter
\AtBeginEnvironment{codebox}{%
  \hllpdf@incodeboxtrue
  % Snapshot the OUTER language before asmblock selects English.
  \iflanguage{english}{\hllpdf@englishcodeboxtrue}{\hllpdf@englishcodeboxfalse}%
  \ifhllpdf@lstnumbermasked
    \let\thelstnumber\hllpdf@savedlstnumber
    \hllpdf@lstnumbermaskedfalse
  \fi
}
% listings' AddToHook is GLOBAL. Register once, and gate on local box state;
% adding an unconditional hook at every codebox entry leaks into later listings.
\lst@AddToHook{PreInit}{%
  \ifhllpdf@incodebox
    % Apply every listing override here, inside the LISTING's own group.
    % Setting margins at box entry leaks them into nested semantic callouts.
    \lstset{resetmargins=false,xrightmargin=26pt}%
    % Styles resolve before PreInit. In LTR the physical number gutter follows
    % xleftmargin; in Hebrew retain the measured inherited xrightmargin path.
    \ifhllpdf@englishcodebox
      \lstset{xleftmargin=26pt,xrightmargin=0pt}%
    \fi
    % ONE box, not two: suppress only the frame/background/skips of a listing
    % inside this codebox. The next unboxed listing keeps its own style.
    \lstset{frame=none,backgroundcolor=\color{cSlateBg},aboveskip=0pt,belowskip=0pt}%
  \fi
}
\makeatother
% </hllpdf:codebox-gutter>

% v3.37: line numbers inside a COLORED BOX are sliced by the box border —
% measured: digit span x=63.5-67.3 vs border line at x=65.2 (the border runs
% THROUGH the digits), and the inter-border gap (3.5pt) is narrower than a
% digit (3.8pt), so no margin knob can save them (under resetmargins,
% xleftmargin is inert in-box — measured: code at x=75.26 with 2.2em AND
% 3.0em). A 3-6 line boxed snippet does not need line numbers; blank the
% number TEXT inside every colored box (\thelstnumber is what gets printed;
% the style's numbers=left survives, so nothing else shifts — code position
% byte-frozen at x=75.26, verified). Unboxed listings keep their numbers.
% v3.78: the blanking is coupled to the FRAME, not to the box name —
% clean_theme.tex registers LATER hooks (registration order wins) that
% restore the number text inside the boxes it renders frameless.
% v4.5.0 RETRACTS the "no margin knob can save them" clause above. It holds only
% under resetmargins=true (the shipped default): there xleftmargin is inert
% in-box, as measured. With resetmargins=false, xrightmargin DOES carry the
% \llap'ed digit column clear of the rule — measured; the codebox \lstset that
% does it is the one above, in the codebox block. The loop below therefore
% covers SIX boxes, not "every colored box": codebox opts out and keeps its
% numbers.
% <hllpdf:semantic-listing-mask>
\makeatletter
\@for\HEB@lstbox:=defbox,thmbox,notebox,exbox,warnbox,keybox\do{%
  \expandafter\AtBeginEnvironment\expandafter{\HEB@lstbox}{%
    \ifhllpdf@lstnumbermasked\else
      \let\hllpdf@savedlstnumber\thelstnumber
    \fi
    \hllpdf@lstnumbermaskedtrue
    \def\thelstnumber{}%
    % The nearest semantic container owns its listing presentation. On exit,
    % TeX's group restores the enclosing codebox state and formatter.
    \hllpdf@incodeboxfalse}}
\makeatother
% </hllpdf:semantic-listing-mask>

% ============================================================================
% Unified TABLE presentation: subtle uniform slate frame + auto-numbered
% caption rendered INSIDE the frame. Use:
%   \begin{tablebox}\centering  <tabular>  \tabcap{<description>}\end{tablebox}
% Replaces the buggy \caption* (which emitted a stray "טבלה N : *" line) and
% provides an explicit framed option + numbered caption "טבלה N.M: תיאור".
% Optional intro text may be placed at the TOP, before the tabular, inside the box.
% ============================================================================
\definecolor{cTblLine}{HTML}{B9C4D2}  \definecolor{cTblBg}{HTML}{FAFBFC}
% v3.23: class-aware — under `article` there is no chapter counter (the promised
% report->article switch used to die with "No counter 'chapter' defined").
\makeatletter
\@ifundefined{c@chapter}{%
  % v3.31: in article the counter resets per SECTION, so an unprefixed number
  % repeated ("טבלה 1." in every section — measured on a two-section render).
  % Same conditional idea as the report branch: prefix the section number once
  % a real \section exists; a pre-section table stays "טבלה 1".
  \newcounter{tbl}[section]
  \renewcommand{\thetbl}{\ifnum\value{section}>0 \thesection.\fi\arabic{tbl}}%
}{%
  \newcounter{tbl}[chapter]
  % v3.30: a report-class doc that uses SECTIONS ONLY (chapter stays 0 —
  % the common short-doc pattern) used to print "טבלה 0.1"; the chapter
  % prefix now appears only once a real \chapter exists (render-verified
  % both ways: "טבלה 1." chapterless, "טבלה 1.1" with chapters).
  \renewcommand{\thetbl}{\ifnum\value{chapter}>0 \thechapter.\fi\arabic{tbl}}%
  % v3.64: the SAME leak in \thesection itself, left behind by v3.30. That fix
  % guarded the CAPTION counter and stopped there, so a chapterless report still
  % printed the heading "0.1 סעיף ראשון" and \ref to it printed "0.2" — measured:
  %   p1 y=65.2  '…0.1'   p1 y=144.9 '…0.2'   p1 y=183.4 '…0.2.1'
  % report.cls hardcodes \thesection as \thechapter.\@arabic\c@section, and the
  % leading 0 is the chapter counter that no \chapter ever incremented. Fixing
  % \thesection fixes the whole chain — \thesubsection and \thesubsubsection are
  % defined in terms of it. \theequation and \thetbl/\thefig were already guarded,
  % which is exactly why this survived: every NUMBER around the heading looked
  % right, so the heading's own 0 read as a style choice rather than a bug.
  % \appendix is safe: report.cls resets chapter to 0 and redefines \thechapter to
  % \@Alph\c@chapter, so an appendix \chapter makes the counter 1 and the guard
  % emits "A.1" unchanged. Render-verified in all four modes (chapterless,
  % chaptered, appendix+chapter, appendix+section-only).
  \renewcommand{\thesection}{\ifnum\value{chapter}>0 \thechapter.\fi\arabic{section}}%
}
\makeatother
% v3.37: everyline=true — if the box DOES split across pages, every segment
% gets a CLOSED frame (measured before: the split ran open-ended into the page
% edge with no bottom border, and the caption could strand alone at the next
% page top — "טבלה A.2" shipped that way). The split skips give the segments
% breathing room. In practice, with the caption GLUE below, every previously-
% split table in the corpus now moves WHOLE instead (the white-gap tradeoff,
% documented benign); everyline stays as the safety net for a genuinely
% page-tall table. NB: a table taller than ONE FULL PAGE can grind mdframed's
% split-trial loop (reproduced: minimal doc, minutes of 1pt-step Overfull
% probes) — break such a monster into two tables.
\mdfdefinestyle{hllpdftablegeometry}{leftline=true,rightline=true,topline=true,bottomline=true,
  linewidth=0.9pt,
  innertopmargin=9pt,innerbottommargin=8pt,skipabove=11pt,skipbelow=11pt,
  innerleftmargin=11pt,innerrightmargin=11pt,roundcorner=4pt}
\newcommand{\hllpdftableframestyle}{hllpdftablegeometry}
\newmdenv[frametitle={},style=hllpdftablegeometry,
  linecolor=cTblLine,backgroundcolor=cTblBg,
  % v4.2.0 (adopted): tablebox is ATOMIC. Short tableboxes near page bottoms
  % repeatedly entered mdframed's split trial, emitting benign Overfull-vbox
  % bursts and occasionally separating the intro prose from the tabular.
  % nobreak=true moves the complete short table to the next page instead. The
  % output-active Overfull-vbox gate still catches a box taller than a page;
  % everyline/split skips remain as defensive styling should that path ever
  % re-arm, but the canonical path never publishes a split tablebox.
  nobreak=true,
  everyline=true,splitbottomskip=6pt,splittopskip=12pt,
  % v3.79 (BLD-37 table side, correction): startinnercode/endinnercode run
  % INSIDE mdframed's body vbox (mdframed.sty:848 / :856+863) -- the second one
  % after \par\unskip, i.e. after line breaking, which is the only moment the
  % \tabcap anchor's true offset from the top of the box EXISTS. See
  % \lakintabboxclose below. These two keys are the ONLY tablebox change.
  startinnercode={\lakintabboxopen},endinnercode={\lakintabboxclose}]{tablebox}
% v3.37: \nopagebreak GLUE — \vsplit honors penalties, so the caption can no
% longer strand alone at the top of the next page (it takes its rows along, or
% the whole box moves). Render-verified on the real shipped orphan.
% v3.59: \NR@gettitle is what nameref injects into the stdlib title-bearing
% commands (\@sect, \@chapter, \caption) to fill \@currentlabelname. \figcap and
% \tabcap are INDEPENDENT macros — they never called it, so \@currentlabelname kept
% whatever the last heading left and \nameref{fig:x} silently printed the enclosing
% SECTION's title instead of the caption (proven in the .aux: field 3 of \newlabel
% was the heading text while field 1 and the anchor `fig.1.1` were correct — so this
% is a THIRD, distinct head next to BLD-13's counter fix and BIDI-38's autoref-name
% fix, neither of which writes that field). \@ifundefined keeps it inert if
% hyperref/nameref is ever dropped. Guarded helper because \figcap is defined
% after \makeatother.
% NB both helpers below need \makeatletter — this part of the preamble is OUTSIDE
% any letter-catcode region, so an unwrapped \@currentlabel parses as LaTeX's \@
% and dies with "You can't use \spacefactor in vertical mode" (measured).
\makeatletter
\newcommand{\lakinrtitle}[1]{\@ifundefined{NR@gettitle}{}{\NR@gettitle{#1}}}
% v3.59: \hitem{<letter>}{<key>} — a REFERENCEABLE Hebrew-lettered list item.
%   \begin{enumerate} \hitem{א}{it:a} הסעיף הראשון. \end{enumerate}  ->  \ref{it:a} = (א)
% Why this exists: automatic Hebrew list labels fail silently, so SKILL.md mandates
% an EXPLICIT \item[(\textbf{א})] — but LaTeX puts \refstepcounter INSIDE
% \if@noitemarg (latex.ltx:12964), and \item[...] never sets that flag. So an
% explicit item steps no counter and never touches \@currentlabel: a \label there
% binds to the AMBIENT one and \ref prints the CHAPTER/SECTION number where the
% item letter belongs (proven in the .aux: {chapter.1}). Same doctrine as \figcap —
% own the label state before \label. \phantomsection matters: without it the
% hyperlink still lands on the chapter head (measured); with it the anchor is
% section*.N and the bookmark outline stays clean.
\newcommand{\hitem}[2]{\item[(\textbf{#1})]%
  \def\@currentlabel{(#1)}\def\@currentlabelname{#1}\phantomsection\label{#2}}
% NB (v3.66): \hitem is \ref-only. \autoref{key} calls it a SECTION ("סעיף (א)")
% because \phantomsection brands the anchor section*.N, and \nameref{key} prints
% the bare letter. For \autoref-able items prefer \hebenumlabels below.
%
% v3.66: \hebenumlabels — OPT-IN automatic Hebrew enumerate letters, done right.
%   {\hebenumlabels\begin{enumerate}\item\label{q:a} ...\end{enumerate}}
%   -> labels (א) (ב) (ג); \ref{q:a} = א; anchor Item.N (real, per-item).
% This retires the belief that made \item[(\textbf{א})] + \hitem necessary.
% "Automatic Hebrew labels fail silently" was TRUE but unexplained, and the
% mechanism is now bisected (plain article vs +enumitem, measured): ENUMITEM
% PRE-EXPANDS the label expression in an \edef BEFORE \refstepcounter, with the
% counter still at 0. Tokens it shields by name (\arabic,\alph,\roman) survive;
% a plain \newcommand \ifcase wrapper evaluates at 0 -> the EMPTY zero-branch ->
% the label prints as "()" with the letter GONE (and \the\c@enumi prints a
% literal 0, \number\value dies fatally pointing at \item). The cure is one
% property, not a doctrine: the wrapper must be ROBUST (survives the \edef
% unexpanded), while \theenumi stays EXPANDABLE (so \refstepcounter writes the
% real letter to the .aux). Verified: labels (א)(ב)(ג), \ref -> א/ב/ג each item
% its own value, .aux = {א}{...}{Item.1} — versus the explicit-\item doctrine
% where BIDI-50/\hitem exist purely to patch the counter that never stepped.
\def\@hebalph#1{\ifcase#1\or א\or ב\or ג\or ד\or ה\or ו\or ז\or ח\or ט\or י\or
  יא\or יב\or יג\or יד\or טו\or טז\or יז\or יח\or יט\or כ\else\@ctrerr\fi}
\DeclareRobustCommand*\hebalph[1]{\expandafter\@hebalph\csname c@#1\endcsname}
\newcommand{\hebenumlabels}{%
  \renewcommand{\theenumi}{\@hebalph\c@enumi}%
  \renewcommand{\labelenumi}{(\textbf{\hebalph{enumi}})}}
\makeatother
% v3.63: \mbox the LABEL. \thetbl ("7.1") and a caption description that STARTS
% WITH A DIGIT ("32 הרגיסטרים…") are two European-Number runs separated only by
% neutrals (the label's own "." and \enspace). Under the RTL paragraph base the UBA
% resolves those neutrals to the RTL level, so both numbers sit in ONE reversed span
% and EXCHANGE VISUAL POSITIONS. Measured on the shipped PDF: a table whose real
% number is 7.1 printed "טבלה 32" — a different, EXISTING table number — while \ref
% correctly printed 7.1, so the caption silently contradicted its own cross-reference,
% and the description lost its leading count. The source contains no adjacency at all;
% it is manufactured by this macro, which is why every source linter was blind.
% \mbox makes the label one atomic box so nothing can reorder across it. Verified
% zero-regression: an ordinary Hebrew caption keeps identical x-geometry.
% One paragraph style covers label, number, separator and description.
% Explicit project/template overrides remain possible; the label hook is kept
% for compatibility, with an identity default. Counters/anchors stay unchanged.
\newcommand{\lakitcaptionstyle}{\normalfont\small\color{black}\hypersetup{allcolors=black}}
\newcommand{\lakitcaptionlabelstyle}[1]{#1}
\newcommand{\lakitcaptionpunctuation}{:}
% Native floats, captionof and longtable use the same presentation policy.
% subcaption is optional and inherits the same font rather than a second size.
\usepackage{caption}
\DeclareCaptionFont{lakit}{\lakitcaptionstyle}
\DeclareCaptionLabelFormat{lakit}{\lakitcaptionlabelstyle{#1~\en{#2}}}
\DeclareCaptionLabelSeparator{lakit}{\lakitcaptionpunctuation\enspace}
\captionsetup{font=lakit,labelfont={},textfont={},labelformat=lakit,
  labelsep=lakit,justification=centering}
\makeatletter
\AtBeginDocument{\@ifpackageloaded{subcaption}{%
  \captionsetup[sub]{font=lakit,labelfont={},textfont={},labelsep=lakit}}{}}
\makeatother
% v3.64: OPT-IN list-of-figures/tables entry via an optional short caption.
% \listoffigures and \listoftables are inherited from report.cls and LOOK available,
% but \figcap/\tabcap step the custom \c@fig/\c@tbl counters and never call
% \addcontentsline — so both printed their heading and NOTHING ELSE. Measured: a doc
% with one \figcap and one \tabcap rendered page 1 as 'רשימת האיורים\n1\n', and the
% .lof held only \@chapter's \addvspace, zero \contentsline entries. An empty list
% under a correct Hebrew heading reads as "this document has no figures".
%
% WHY OPTIONAL, AND NOT AUTOMATIC. Writing #1 unconditionally was implemented first
% and REJECTED: a caption containing \codeline compiles CLEANLY today (verified on
% the pre-v3.64 preamble: 0 errors, PDF produced) but becomes FATAL once it is also
% written to the .lof — \codeline is \hbox to\linewidth{...}, and a full-width box in
% a list entry dies with "\pdfextension endlink cannot be used in vertical mode", no
% output PDF. Gaining a feature nobody uses (0 of 13 collections call \listoffigures)
% at the price of breaking a construct that works is a bad trade. Stock \caption has
% exactly this limitation and exactly this escape hatch, so the convention is familiar.
% \en{}, \code{}, \h{} and inline math in the short text are all safe — verified, they
% survive the write and typeset correctly in the list (RTL geometry confirmed: number
% rightmost at x=509.9, title left of it, dots and folio at the far left).
% With no optional argument the macro is byte-identical to v3.63 — zero risk to every
% existing document, and \listoffigures lists exactly the captions that opted in.
% BLD-37, TABLE side (the \figcap v3.74 idea, adapted): clicking a \ref to a
% table used to land MID-table — hyperref's automatic anchor inside
% \refstepcounter was emitted on the still-open horizontal line whose TAIL is
% the tabular (an hbox whose baseline is the table's vertical CENTER), so the
% jump scrolled the top half of the table above the viewport. The v3.74 entry
% deferred this side claiming "tablebox splits — no single box to measure";
% that deferral dissolves: the measurement target is the TABULAR itself at
% \tabcap ENTRY, while it is still the tail of the open list — the same
% situation \lakin@grabfig handles at an island's end (and the fleet has 0
% split tableboxes). So \lakin@grabtab runs as the FIRST token, before any
% typesetting can push the box away, and the anchor is re-emitted raised above
% the measured tabular. \lakintabraise is a SEPARATE register from
% \lakinfigraise (defined with it, below the LTR-island block): a captionless
% island intentionally leaves the fig register set, and a table must never
% inherit that stale figure height.
%
% WHY the anchor stays AT THE ENTRY SPOT and is NOT deferred into the caption
% paragraph the way \figcap defers it: at \par TeX DROPS a final glue item,
% and the automatic anchor's smash box sitting BEHIND the interword glue kept
% that glue alive — so every centered tabular has always sat half a space off
% true center. Deferring the anchor removed that box and was measured as 982
% differing glyph records: EVERY tabular line shifted 1.30pt (2.61pt for
% pattern 7's two minipages) while captions stayed put. Emitting the raised
% anchor at the old spot — through hyperref's own \Hy@raisedlink, exactly the
% macro the automatic path uses — keeps the line's node structure, and so the
% A/B glyph diff at 0, BY CONSTRUCTION. \Hy@raisedlink reads the raise from
% \HyperRaiseLinkDefault (default \baselineskip = the old mid-table dest); the
% group-local \edef freezes measured-height+42pt into it, and \smash inside
% \Hy@raisedlink hides any raise from the line's extent. \HyCnt@ProvideTheH-
% Counter keeps the dest NAME byte-identical to the automatic one (it defines
% \theHtbl, which \hyper@makecurrent prefers — without it the first measured
% caption was named tbl.1 where the automatic path said tbl.0.1, measured).
% vmode entry, or anything but a box at the tail (no tabular to measure),
% falls back to the plain \refstepcounter — byte-identical to the old path.
\makeatletter
\newcommand{\tabcap}[2][]{\lakin@grabtab
  % v3.79: arm the box-close correction. Set unconditionally and in BOTH
  % branches below's scope, so the plain-\refstepcounter fallback is corrected
  % too; \lakintabboxopen clears it at every \begin{tablebox}, so a \tabcap
  % outside a tablebox can never leave it armed for a later box.
  \global\lakintabcapseentrue
  \ifdim\lakintabraise>0pt
    \H@refstepcounter{tbl}\HyCnt@ProvideTheHCounter{tbl}\hyper@makecurrent{tbl}%
    \begingroup
      % A frameless table targets its measured top. The frame's historical
      % cushion belongs only to tablebox and is corrected at box close.
      \edef\HyperRaiseLinkDefault{\the\lakintabraise}%
      \iflakintabframed
        \edef\HyperRaiseLinkDefault{\the\dimexpr\lakintabraise+42pt\relax}%
      \fi
      \Hy@raisedlink{\hyper@anchorstart{\@currentHref}\hyper@anchorend}%
    \endgroup
    \global\lakintabraise 0pt\relax
  \else
    \refstepcounter{tbl}%
  \fi
  \lakinrtitle{#2}\par\nopagebreak\addvspace{6pt}%
  \def\lakin@short{#1}\ifx\lakin@short\@empty\else
    \addcontentsline{lot}{table}{\protect\numberline{\thetbl}{\protect\ignorespaces #1}}\fi
  {\centering\lakitcaptionstyle\mbox{\lakitcaptionlabelstyle{\tablename~\en{\thetbl}\lakitcaptionpunctuation}}\enspace #2\par}}
\makeatother

% ============================================================================
% Unified FIGURE caption — the \tabcap twin for TikZ/plot figures (v3.34.0).
% A figure (tikzpic) sits inline in the prose; put \figcap BELOW it for a
% numbered, \ref-able caption "איור N.M: תיאור" in the same style as \tabcap:
%   \begin{tikzpic} ... \end{tikzpic}
%   \figcap{\label{fig:foo}description}     % \label FIRST, inside the argument
% WHY a macro, not a hand-written gray line: a manual
%   {\centering\small\textcolor{cGray}{איור: ...}} with a \label inside has NO
% \refstepcounter, so the \label binds to the AMBIENT \@currentlabel (the last
% equation) — \ref then prints an equation number AS the figure number
% (measured on a real page: "איור 48" where 48 is equation 48). \figcap steps a
% real \c@fig counter first, so \label is correct, numbering is consistent, and
% refs are stable. Class-aware like \thetbl (chapter- or section-prefixed).
% ============================================================================
\makeatletter
\@ifundefined{c@chapter}{%
  \newcounter{fig}[section]
  \renewcommand{\thefig}{\ifnum\value{section}>0 \thesection.\fi\arabic{fig}}%
}{%
  \newcounter{fig}[chapter]
  \renewcommand{\thefig}{\ifnum\value{chapter}>0 \thechapter.\fi\arabic{fig}}%
}
\makeatother
\makeatletter
\newcommand{\figcap}[2][]{%
  \ifdim\lakinfigraise>0pt
    % v3.74: step WITHOUT hyperref's automatic anchor, then emit the anchor
    % raised above the measured figure. \H@refstepcounter is hyperref's saved
    % pre-hook \refstepcounter (sets the counter AND \@currentlabel — \label
    % in the caption arg binds exactly as before); \hyper@makecurrent supplies
    % \@currentHref, so the dest NAME is identical to the old automatic one —
    % \ref/\autoref/bookmarks are unchanged, only the dest'S POSITION moves.
    \H@refstepcounter{fig}\hyper@makecurrent{fig}%
    % The anchor is emitted INSIDE the caption paragraph below, as a zero-size
    % raisebox — NOT here in vertical mode, where a \raisebox starts an empty
    % paragraph and adds a blank line (measured: every glyph after the caption
    % shifted +18.8pt). Inside an existing line a [0pt][0pt] box is inert.
    \edef\lakin@anch{\@currentHref}%
    % \edef FREEZES the measured height into the pending macro: the deferred
    % \raisebox is typeset inside the caption paragraph AFTER the register is
    % reset to the 0pt sentinel below — a plain \def read 0 at use time
    % (measured: LAKINRAISE=0.0, dest moved 14pt instead of above the figure).
    \edef\lakin@pend{\noexpand\raisebox{\the\dimexpr\lakinfigraise+42pt\relax}[0pt][0pt]{%
      \noexpand\hyper@anchorstart{\lakin@anch}\noexpand\hyper@anchorend}}%
    \global\lakinfigraise 0pt\relax
  \else
    \refstepcounter{fig}\let\lakin@pend\relax
  \fi
  \lakinrtitle{#2}\par\nopagebreak\addvspace{4pt}%
  \def\lakin@short{#1}\ifx\lakin@short\@empty\else
    \addcontentsline{lof}{figure}{\protect\numberline{\thefig}{\protect\ignorespaces #1}}\fi
  {\centering\lakin@pend\lakitcaptionstyle\mbox{\lakitcaptionlabelstyle{\figurename~\en{\thefig}\lakitcaptionpunctuation}}\enspace #2\par}}
\makeatother

% ============================================================================
% LTR islands.  EVERY tikzpicture / quantikz MUST live inside one, or the whole
% figure's coordinates and Latin labels flip. (Hebrew node text still uses \h{}.)
% ============================================================================
% v3.74 (BLD-37): the \figcap hyperlink DESTINATION used to land at the caption
% line — BELOW the figure — so clicking a figure reference scrolled the figure
% off-screen (measured fleet-wide: 97/97 fig dests in one collection sat 19-54pt
% under their figure's ink; a 347pt figure was fully hidden after the jump).
% hypcap-style fix in two parts: (1) at the island's END, \lakin@grabfig grabs
% the just-typeset centered line with \lastbox, measures it, and puts it back
% with \nointerlineskip — the interline glue above the line was already
% contributed and is NOT re-added, so the page geometry is unchanged to the
% glyph (verified); (2) \figcap then steps the counter WITHOUT hyperref's
% automatic anchor (\H@refstepcounter = the pre-hyperref \refstepcounter) and
% emits the anchor itself, raised by the measured height — the dest lands ABOVE
% the figure top. A captionless island leaves numbering untouched (nothing is
% stepped at the island; the grab flag is cleared at every island begin, so a
% stale height from a captionless figure cannot leak into a later caption).
\newsavebox{\lakinfigbox}
\newdimen\lakinfigraise  % real \dimen: tex.setdimen cannot target a \newlength SKIP register (measured lua error)
% BLD-37 table side: a SEPARATE register for \tabcap. The two must not share —
% a captionless island INTENTIONALLY leaves \lakinfigraise set (the grab runs
% at every island END; only the NEXT island's begin clears it), so a shared
% register would hand a \tabcap after a captionless figure that stale FIGURE
% height. \lakintabraise is written only by \lakin@grabtab at \tabcap entry
% and consumed/reset inside \tabcap itself, so it can never go stale.
\newdimen\lakintabraise
\makeatletter
% v3.74b: measure WITHOUT touching the list. The first draft grabbed the box
% with \lastbox and put it back — and \lastbox DISCARDS the node's baseline
% shift, so a vertically-anchored quantikz circuit rose 2.57pt and everything
% after it sank 5.32pt (measured on the A/B glyph diff). LuaTeX lets us READ
% the tail node's dimensions without removing anything: geometry is identical
% by construction. \lakinfigraise=0pt is the "no figure measured" sentinel.
\def\lakin@grabfig{%
  \ifhmode
    \directlua{
      local n = tex.nest[tex.nest.ptr].tail
      local GLUE, PEN = node.id('glue'), node.id('penalty')
      while n and (n.id == GLUE or n.id == PEN) do n = n.prev end
      local HL, VL = node.id('hlist'), node.id('vlist')
      if n and (n.id == HL or n.id == VL) then
        tex.setdimen('global', 'lakinfigraise',
                     (n.height or 0) + (n.depth or 0))
      else
        tex.setdimen('global', 'lakinfigraise', 0)
      end}%
  \else
    \global\lakinfigraise 0pt\relax
  \fi}
% BLD-37 table side: \lakin@grabfig's twin, NOT a shared helper — the Lua body
% differs in exactly one place, the skip set. Between \end{tabular} and \tabcap
% an otherlanguage wrapper (patterns 3/5 of example_tables) leaves a \write
% whatsit and a dir node ON TOP of the interword glue (measured node dump:
% "whatsit/write <- dir <- glue(2.62) <- hlist"), so the table reader must
% also step over whatsit/dir markers — all zero-size — to reach the tabular.
% \lakin@grabfig deliberately does NOT get the wider skip set: its glue/penalty
% body is the one the v3.74 A/B proof froze, and widening it would newly
% measure (and so newly raise) anchors in island shapes that today take the
% plain path — an unproven behavior change on the FIG side. Anything that is
% not a box behind those markers still hits the 0pt sentinel and \tabcap falls
% back to the plain \refstepcounter.
% SECOND deliberate difference: this stores the box's HEIGHT, not height+depth.
% \figcap's anchor is emitted in the caption paragraph BELOW the figure, so it
% must climb the full extent; \tabcap's anchor sits ON the tabular's own line,
% whose baseline is the vcentered tabular's MIDPOINT — height alone reaches
% the top edge, and adding the depth was measured to overshoot by half the
% table (dests 80-135pt above their frame, INSIDE the preceding table's box).
\def\lakin@grabtab{%
  \ifhmode
    \directlua{
      local n = tex.nest[tex.nest.ptr].tail
      local GLUE, PEN = node.id('glue'), node.id('penalty')
      local WHAT, DIR = node.id('whatsit'), node.id('dir')
      while n and (n.id == GLUE or n.id == PEN
                   or n.id == WHAT or n.id == DIR) do n = n.prev end
      local HL, VL = node.id('hlist'), node.id('vlist')
      if n and (n.id == HL or n.id == VL) then
        tex.setdimen('global', 'lakintabraise', n.height or 0)
      else
        tex.setdimen('global', 'lakintabraise', 0)
      end}%
  \else
    \global\lakintabraise 0pt\relax
  \fi}
\makeatother
% ============================================================================
% v3.79 (BLD-37 table side, CORRECTION): put the anchor at the TOP OF THE BOX,
% instead of guessing the distance to it.
% ============================================================================
% v3.77 raised the \tabcap anchor by (height of the box at the tail of the open
% horizontal list) + 42pt. That is a MEASUREMENT taken BEFORE line breaking, and
% it is short by every term it cannot see. Fleet-measured on the 12 built
% collections (299 tbl.* dests, ownership established through each table's own
% CAPTION, never through "which frame encloses the dest"): 3 dests landed INSIDE
% their own tablebox, below the frame top —
%   computer-architecture       tbl.14.1  -46.4pt   three stacked tabulars
%   computer-architecture       tbl.30.5  -13.5pt   tabular 2.85pt over \linewidth
%   advanced-electronic-devices tbl.B.1   -59.2pt   tabular 2.56pt over \linewidth
% Two missing terms, both reproduced in a minimal document (probe P2/P3):
%   (1) VERTICAL LIST ALREADY CONTRIBUTED inside the box — a second and third
%       tabular, or a heading above the table. The tail read sees only the LAST
%       box, so 77.7pt of stacked tabulars is invisible to the raise.
%   (2) WHICH LINE THE ANCHOR LANDS ON. When the tabular alone overfills
%       \linewidth (459.6121pt in a tablebox), the interword glue after
%       \end{tabular} becomes the chosen breakpoint and \Hy@raisedlink is
%       emitted on the NEXT line — one (depth + lineskip) lower. That term does
%       not EXIST at \tabcap time: line breaking has not happened yet. No
%       measurement taken there can ever recover it. (Proof: probe P3, the same
%       table 2.84912pt too wide, fails at -13.5pt; probe P4, the identical
%       table with \tabcolsep 2pt so it fits, is correct at +32.0pt.)
% So the raise is not computed at \tabcap at all. It is CORRECTED at
% \end{tablebox}, from inside mdframed's own body vbox (endinnercode runs after
% \par\unskip, i.e. after line breaking, while tex.nest[tex.nest.ptr] is still
% that vbox's vertical list). There the anchor's true offset from the top of the
% box content is not estimated but WALKED, and the \Hy@raisedlink shift is set
% so the dest lands 42pt above it — which is 42pt - innertopmargin(9pt) -
% linewidth(0.9pt) = 32.0pt above the FRAME top, the exact position the 295
% already-correct dests have today. That is why this is a correction and not a
% re-placement: the good sites do not move.
%
% WHY THE GLYPH DIFF IS 0 BY CONSTRUCTION — the only thing written is the
% `shift` field of the hlist that \Hy@raisedlink already built, and the guard
% below refuses to write it unless that box is 0x0x0. A zero-size box carries no
% dimension into its parent's packing, the packing already happened, and shift
% is not an input to it. Nothing else in the list is touched: no node is added,
% removed, resized or reordered. Measured, not assumed: assets/example_tables.tex
% and the whole computer-architecture collection A/B at 0 differing pymupdf
% rawdict glyph records.
%
% WHY NOT the alternatives (all prototyped and measured before this one):
%   * accumulate the vertical list at \tabcap time — fixes term (1), CANNOT fix
%     term (2); 2 of the 3 real failures survive. A fix that leaves the majority
%     of the defect standing is the false all-clear this repo exists to avoid.
%   * \zref-savepos two-pass — correct on all shapes, but it makes the dest
%     POSITION aux-derived (DEFECTS BLD-1/2), a single-pass compile silently
%     ships an unraised anchor, and the bare y-difference has no page guard.
%   * emit the anchor at \begin{tablebox} and let \tabcap adopt it — needs the
%     tbl counter stepped at box OPEN to keep the dest NAME byte-identical; a
%     captionless tablebox then either shifts every later table number or leaves
%     an orphan dest that the next table duplicates (a HARD build gate). This
%     correction touches neither the counter nor the name: \tabcap's naming path
%     is v3.77's, unchanged, so resolve_names() is byte-identical (verified).
%   * suppressing the line break (a penalty before the glue) — that MOVES the
%     tabular: the break currently discards the interword glue, so forbidding it
%     re-centers the two overfull tables. Reflow, rejected.
%
% SCOPE. The correction is armed only by \tabcap ITSELF (\lakintabcapseen),
% cleared at every box open, so a captionless tablebox is a no-op and a \tabcap
% hosted OUTSIDE a tablebox (the fleet has exactly one, computer-architecture
% tbl.14.2 in an exbox) keeps v3.77's measured raise. Both fallbacks are
% therefore never worse than today; and because hyperref's AUTOMATIC anchor uses
% the same \Hy@raisedlink shape, the plain-\refstepcounter path (vmode entry, or
% prose between \end{tabular} and \tabcap — probe P6, -51.5pt) is corrected too.
% KNOWN LIMIT — a tablebox that SPLITS across pages. The offset is measured from
% the top of the UNSPLIT body, so a dest in segment 2 gets raised past its own
% segment. MEASURED on a forced split (probe: 40 filler lines + a table): the
% dest went from y=585.0 on page 2 (inside segment 2, at the caption — the old
% BLD-37 symptom) to y=-113.4, i.e. ABOVE the page top. Viewers clamp that to the
% page top, which happens to show segment 2 whole, but an out-of-page coordinate
% is not something to ship on purpose. Hence the \textheight guard in
% lakintab.fix: a body TALLER THAN THE TEXT BLOCK cannot fit on any page, so it
% must split — the correction is skipped and v3.77's behaviour stands unchanged
% (verified: the probe returns to y=585.0). The guard is deliberately EXACT, not
% a page-fitting heuristic: it is sound (a taller-than-page body always splits)
% but PARTIAL (a shorter body that starts low on a page could still split). That
% residue is unreachable in practice — v3.37's everyline + the caption's
% \nopagebreak glue move such boxes WHOLE instead (probe N7: a 22-row table
% pushed down by 430pt moved whole, 0 segments), and the fleet has 0 split
% tableboxes, re-counted after this change. Reimplementing mdframed's own
% page-fitting arithmetic to close the residue would be a heuristic with its own
% failure mode, which is the trade this repo has repeatedly refused.
\newif\iflakintabcapseen
\newif\iflakintabframed
% The Lua lives in ONE \directlua at load time, and is written with NO
% backslash, %, # or ~ : \directlua expands its argument like \edef, so an
% active ~ would expand to \nobreakspace and a % would eat the rest of the line
% (the second is how the first draft of this block died). "not (a == b)" stands
% in for Lua's inequality operator for the same reason.
\directlua{
  lakintab = { }
  lakintab.HL = node.id("hlist")
  lakintab.VL = node.id("vlist")
  lakintab.RU = node.id("rule")
  lakintab.GL = node.id("glue")
  lakintab.KE = node.id("kern")
  lakintab.WH = node.id("whatsit")
  lakintab.TM = node.id("temp")
  lakintab.WD = node.subtype("pdf_dest")
  function lakintab.d(v)
    if v == nil then return 0 end
    if v < -1000000000 then return 0 end
    return v
  end
  function lakintab.scan(list, y, horiz, parent, target)
    local x = list
    while x do
      local i = x.id
      if i == lakintab.HL or i == lakintab.VL then
        if horiz then
          local b = y + x.shift
          if i == lakintab.HL then
            lakintab.scan(x.head, b, true, x, target)
          else
            lakintab.scan(x.head, b - lakintab.d(x.height), false, x, target)
          end
        else
          if i == lakintab.HL then
            lakintab.scan(x.head, y + lakintab.d(x.height), true, x, target)
          else
            lakintab.scan(x.head, y, false, x, target)
          end
          y = y + lakintab.d(x.height) + lakintab.d(x.depth)
        end
      elseif i == lakintab.WH then
        if x.subtype == lakintab.WD then
          local id = x.dest_id
          if type(id) == "string" and string.sub(id, 1, 4) == "tbl." then
            if horiz and parent and parent.width == 0
               and parent.height == 0 and parent.depth == 0 then
              parent.shift = parent.shift + (target - y)
            end
          end
        end
      elseif not horiz then
        if i == lakintab.RU then
          y = y + lakintab.d(x.height) + lakintab.d(x.depth)
        elseif i == lakintab.GL or i == lakintab.KE then
          y = y + lakintab.d(x.width)
        end
      end
      x = x.next
    end
  end
  function lakintab.fix(target, maxh)
    local nst = tex.nest[tex.nest.ptr]
    if nst == nil then return end
    local h = nst.head
    if h == nil then return end
    if h.id == lakintab.TM then h = h.next end
    local tot = 0
    local x = h
    while x do
      local i = x.id
      if i == lakintab.HL or i == lakintab.VL or i == lakintab.RU then
        tot = tot + lakintab.d(x.height) + lakintab.d(x.depth)
      elseif i == lakintab.GL or i == lakintab.KE then
        tot = tot + lakintab.d(x.width)
      end
      x = x.next
    end
    if tot > maxh then return end
    lakintab.scan(h, 0, false, nil, target)
  end
}
% @-FREE NAMES ON PURPOSE: these two are named in the `tablebox` key list at
% \newmdenv above, which is tokenised OUTSIDE any \makeatletter region — a
% \lakin@... there would tokenise as \lakin followed by the letters @tabopen.
\newcommand{\lakintabboxopen}{\lakintabframedtrue\global\lakintabcapseenfalse}
\newcommand{\lakintabboxclose}{%
  \iflakintabcapseen
    \directlua{lakintab.fix(-\number\dimexpr\lakintabcushion\relax,
                            \number\dimexpr\textheight\relax)}%
    \global\lakintabcapseenfalse
  \fi\lakintabframedfalse}
% The cushion above the box's CONTENT top. 42pt puts the dest 32.0pt above the
% frame top (42pt - innertopmargin 9pt - linewidth 0.9pt), i.e. exactly where
% the 295 correct dests already are — the correction moves only the broken ones.
\newcommand{\lakintabcushion}{42pt}
\newenvironment{qcirc}{\par\medskip\global\lakinfigraise 0pt\relax\begin{otherlanguage}{english}\begin{center}}
  {\csname lakin@grabfig\endcsname\end{center}\end{otherlanguage}\par\medskip}
\newenvironment{tikzpic}{\par\medskip\global\lakinfigraise 0pt\relax\begin{otherlanguage}{english}\begin{center}}
  {\csname lakin@grabfig\endcsname\end{center}\end{otherlanguage}\par\medskip}

% ============================================================================
% Code listings.  ★★ v3.4 THE BOX-OVERFLOW FIX:
%   resetmargins=true   → the listing IGNORES the enclosing environment's
%                         indentation and measures from the PAGE. v3.4 stated
%                         this backwards ("measured from the ENCLOSING box"),
%                         and v3.37 then blanked in-box line numbers because
%                         the page-anchored listing put its hanging digits on
%                         top of the box rule. Measured v4.5.0: with the key
%                         FALSE the listing follows the box and the code moves
%                         from x=75.26 to x=97.74, clearing the rule at 72.07.
%                         The key now lives in one \lstset below instead of in
%                         each style, so a box can override it — codebox does.
%   framexleftmargin=0  → the frame (left bar) does not extend left of the code.
%   xleftmargin keeps a small indent so the line numbers have room INSIDE.
% asmblock wraps the listing in an LTR island (code is LTR).
% ============================================================================
\lstdefinestyle{asm}{
  basicstyle=\ttfamily\small,
  backgroundcolor=\color{cBlueBg!55},
  frame=leftline,framerule=2.5pt,rulecolor=\color{cBlue},
  numbers=left,numberstyle=\tiny\color{cGray},numbersep=8pt,
  xleftmargin=2.2em,framexleftmargin=0pt,framexrightmargin=0pt,
  breaklines=true,columns=fullflexible,keepspaces=true,
  showstringspaces=false,aboveskip=8pt,belowskip=8pt,
  % v3.36: without a comment DELIMITER the commentstyle below never fires —
  % listings has no `language` here, so it does not know '#' opens a comment,
  % and every asm listing rendered flat black (measured: a whole-page pixel
  % sweep found ZERO cGreen glyphs across 135 shipped listings). morecomment
  % gives it the delimiter; color-only change, glyph widths untouched.
  morecomment={[l]{\#}},
  commentstyle=\color{cGreen},
}
\lstdefinestyle{cstyle}{
  language=C,basicstyle=\ttfamily\small,
  backgroundcolor=\color{cGreenBg!55},
  frame=leftline,framerule=2.5pt,rulecolor=\color{cGreen},
  numbers=left,numberstyle=\tiny\color{cGray},numbersep=8pt,
  xleftmargin=2.2em,framexleftmargin=0pt,framexrightmargin=0pt,
  breaklines=true,columns=fullflexible,keepspaces=true,
  showstringspaces=false,keywordstyle=\color{cBlue}\bfseries,
  commentstyle=\color{cGray}\itshape,aboveskip=8pt,belowskip=8pt,
}
% v3.36: pystyle/htmlstyle — colored twins of cstyle for the Python/HTML
% listings that used to repeat a ~200-char options blob with NO color
% (keywords rendered bold-black, comments italic-black — measured; C alone
% was colored). Geometry is copied VERBATIM from the shipped blob (frame,
% margins, skips) so swapping a call site to style=pystyle moves NOTHING —
% only keywordstyle/commentstyle/stringstyle are added (bold stays bold,
% italic stays italic, mono metrics identical => zero reflow, proven by
% per-page word-bbox diff). breaklines/breakatwhitespace stay per-site.
% v3.78: numbers=left joins both styles, with the exact numberstyle/numbersep
% keys asm/cstyle carry — fleet-wide consistency (author decision: all four
% styles numbered; "asm/cstyle numbered, py/html not" was an accident of the
% v3.36 geometry freeze, not a design). The numbers hang LEFT of the code box
% (listings sets them with \llap, so no overfull hbox can arise and the code
% x is untouched — measured: code frozen, digits at code_x-11.8pt); inside a
% colored box they are blanked like everyone else's (v3.37 hook), and the
% clean variant restores them in its frameless boxes (v3.78, above).
\lstdefinestyle{pystyle}{
  language=Python,basicstyle=\ttfamily\small,
  columns=fullflexible,keepspaces=true,
  frame=single,rulecolor=\color{cTblLine},
  numbers=left,numberstyle=\tiny\color{cGray},numbersep=8pt,
  xleftmargin=0.5em,framexleftmargin=0pt,framexrightmargin=0pt,
  showstringspaces=false,aboveskip=8pt,belowskip=8pt,
  keywordstyle=\color{cBlue}\bfseries,
  commentstyle=\color{cGray}\itshape,
  stringstyle=\color{cGreen},
}
% v4.5.0: hoisted out of the four styles so an enclosing box can override it.
% Every listing still resolves to true unless something later says otherwise,
% so shipped output is unchanged (verified byte-for-byte on all eight PDFs).
\lstset{resetmargins=true}
\lstdefinestyle{htmlstyle}{
  language=HTML,basicstyle=\ttfamily\small,
  columns=fullflexible,keepspaces=true,
  frame=single,rulecolor=\color{cTblLine},
  numbers=left,numberstyle=\tiny\color{cGray},numbersep=8pt,
  xleftmargin=0.5em,framexleftmargin=0pt,framexrightmargin=0pt,
  showstringspaces=false,aboveskip=8pt,belowskip=8pt,
  keywordstyle=\color{cBlue}\bfseries,
  commentstyle=\color{cGray}\itshape,
  stringstyle=\color{cGreen},
}
\newenvironment{asmblock}{\par\begin{otherlanguage}{english}}{\end{otherlanguage}\par}

% <hllpdf:listing-smart-quotes>
% v4.21.0: four Unicode punctuation characters must enter listings' token
% buffer. Unregistered characters bypass it: a closing quote was printed
% BEFORE the pending word, and an opening quote BEFORE a pending brace.
% Extend the standard character table, then activate only these codepoints
% in the listing's local character-table scope. Do this BEFORE the package's
% delimiter/alsoother/literate hooks so explicit user mappings still win.
% This is not universal Unicode support, nor a conversion to ASCII quotes.
% See listings developer guide 15.3.1 and references/listing-smart-quotes.md.
\makeatletter
\lst@CCPut \lst@ProcessOther {"2018}{"2019}{"201C}{"201D}\z@
\lst@AddToHookAtTop{SelectCharTable}{%
  \lst@ifactivechars
    \catcode"2018=\active \catcode"2019=\active
    \catcode"201C=\active \catcode"201D=\active
  \fi}
\makeatother
% </hllpdf:listing-smart-quotes>

% <hllpdf:listing-technical-symbols>
% v4.30: bounded extension for common technical punctuation in LTR listings.
% Without character-table registration these glyphs bypass the token buffer
% and can move before the preceding word even inside quoted strings.
% Keep the original four-quote block separate for its standing mutation test.
% This is source-faithful printing, not syntax validation or general Unicode.
\makeatletter
\lst@CCPut \lst@ProcessOther
  {"2013}{"2014}{"2026}{"2190}{"2192}{"2194}{"21D2}
  {"2212}{"221E}{"2260}{"2264}{"2265}\z@
\lst@AddToHookAtTop{SelectCharTable}{%
  \lst@ifactivechars
    \catcode"2013=\active \catcode"2014=\active \catcode"2026=\active
    \catcode"2190=\active \catcode"2192=\active \catcode"2194=\active
    \catcode"21D2=\active \catcode"2212=\active \catcode"221E=\active
    \catcode"2260=\active \catcode"2264=\active \catcode"2265=\active
  \fi}
\makeatother
% </hllpdf:listing-technical-symbols>


\usepackage[hidelinks,unicode]{hyperref}   % LAST, after colors. unicode = Hebrew bookmarks.
% v3.62: THE ISOLATION MACROS CORRUPTED EVERY PDF BOOKMARK THEY APPEARED IN.
% \en is defined (line ~155) as
%     \foreignlanguage{english}{\language=\csname l@nohyphenation\endcsname #1}
% and in a heading that body is written verbatim to the .toc and handed to
% hyperref's \pdfstringdef. hyperref strips the two CONTROL SEQUENCES (\language,
% \l@nohyphenation) but the plain `=` CHARACTER token between them is not a control
% sequence, so it survives into the PDF string: \chapter{שער \en{CNOT} והצמדה}
% produced the bookmark "שער =CNOT והצמדה". The printed TOC and the heading on the
% page are CORRECT — only the navigation pane and the viewer window title are wrong,
% which is why it went unnoticed. Measured on the shipped fleet: 70 corrupted
% bookmarks across 5 PDFs, from 177 unprotected \en-in-heading sites. Same corruption
% reached /Title and /Subject via the \hypersetup{pdftitle=...} the docs recommend.
% This hook tells hyperref to use the plain argument for all three isolation macros
% when building a PDF string. It does NOT touch the .toc (the printed TOC is
% byte-unchanged) and it does NOT mask the math signal — a heading with $...$ still
% logs "removing math shift", so BIDI-15 detection is intact.
% BONUS, measured: it also makes \code{} in an UNSTARRED heading COMPILE. BLD-10
% recorded that as needing \texorpdfstring or a starred heading; the real root cause
% is hyperref's \pdfstringdef scan of \code's \ifmmode...\else...\fi body — which is
% exactly why the identically-bodied \texttt survived (hyperref pre-neutralizes
% \texttt in PDF strings) and why \DeclareRobustCommand could never help.
% v3.64: \def\\{ } added. A manual line break in a heading WELDS the two words
% together in the bookmark: \chapter{שורה ראשונה\\שורה שנייה} produced the outline
% entry 'שורה ראשונהשורה שנייה' — "ראשונהשורה" is not a word. hyperref drops the
% \\ and puts nothing in its place, so the break that looks like a space on the
% page is a zero-width join in the navigation pane. The spaced form
% 'מילה אחת \\ מילה שנייה' escapes with a harmless DOUBLE space ('אחת  מילה'),
% which is why this never got noticed. Both are measured on a real outline.
% 0 live instances in the corpus — but SKILL.md and content-style.md both RECOMMEND
% a manual heading break, so this is a primed trap rather than a dormant one, and a
% one-token preamble fix retires it for good instead of asking authors to remember.
% Mapping \\ to a space fixes the WELD, which is the actual defect (a nonexistent
% word in the navigation pane). It does not normalise whitespace: measured after the
% fix, the welded form becomes 'שורה ראשונה שורה שנייה' (correct) and the already-
% spaced form goes from two spaces to three. That is cosmetic — viewers render it as
% a wider gap, and no word is corrupted. A per-heading \texorpdfstring is still the
% way to get an exactly-tuned bookmark string.
\pdfstringdefDisableCommands{\def\en#1{#1}\def\h#1{#1}\def\doclang#1{#1}\def\code#1{#1}\def\codeline#1{#1}\def\\{ }}
% v3.70: RESTORE \url LINE-BREAKING. url.sty freezes its break penalties at LOAD
% time (\mathchardef\UrlBreakPenalty=\binoppenalty), and v3.45's global
% \binoppenalty=10000 (inline-math atomicity) therefore priced EVERY breakpoint
% inside \Url at 10000 = never break. Measured: a long \url in Hebrew prose is
% one unbroken line, ~100pt overfull, invading the margin band the ink gates
% deliberately tolerate. With the two lines below: 0 overfull, the URL breaks
% after '/', the RTL split is CORRECT (fragment order forward), and each line
% gets its own link annotation exactly covering its fragment (measured).
% Confined to \Url's group — inline-math atomicity (TYP-7) is untouched.
\mathchardef\UrlBreakPenalty=200
\mathchardef\UrlBigBreakPenalty=100
\ifx\lakitmainlanguage\hllpdfenglish
  \hypersetup{pdflang=en}
\else
  \hypersetup{pdflang=he}
\fi  % /Lang is the document base locale (not per-span accessibility tagging).
% Set a per-document \hypersetup{pdftitle={...}} too — see SKILL.md.

% v3.52 (round-2 audit): two RTL bibliography/cross-ref defects, both invisible to
% every gate and both affecting the FIRST author to use the feature.
% (a) \@biblabel's default [#1] brackets are bidi-neutral and mirror to ]1[ under
% the RTL bibliography base — every \thebibliography reference list shows ]1[.
% LTR-isolating the label restores [1]. (The render gate's mirror rule fingerprints
% round parens )1(, so square brackets slipped through.)
\makeatletter
\renewcommand*{\@biblabel}[1]{\foreignlanguage{english}{[#1]}}
% v3.54: the SAME neutral-bracket trap on the INLINE side — a multi-key \cite{a,b}
% under the RTL base lays its European-number list right-to-left, so [1, 2] renders
% "[2,1]" and [1, 2, 3] renders "[3,2,1]" (verified by glyph geometry — the reader
% sees the sources in REVERSE order; a note form \cite[p.5]{a} also reorders to
% "[p. 5,1]"). LTR-isolate the whole bracket group (number list AND optional note),
% exactly like \@biblabel above: [1, 2] / [1, 2, 3] / [1, p. 5] all render correct.
% Placed AFTER hyperref so it wins; hyperref hyperlinks each number inside #1 (via
% \@citex), which this wrapper leaves intact. Mirrors the LaTeX kernel \@cite shape.
% v4.9.1: directional isolation must NOT switch the language here. Babel's
% safe-active citation wrapper passes its key list through \bbl@tempa;
% \foreignlanguage reuses that scratch macro before the deferred citation
% loop expands it. The old line therefore cited the literal key "english".
% If such a key existed, even the undefined-reference gate stayed green.
% \babelsublr isolates the bracket group without selecting a locale or
% changing that key list. Preserve native links, key order and optional notes.
\def\@cite#1#2{\babelsublr{[{#1\if@tempswa , #2\fi}]}}
\makeatother

% v4.11.0: a reference inside align's \text inherited HEBdisptag=true,
% so its enclosing hbox mirrored the brackets despite a clean compile.
% Isolate the complete printed reference without selecting a Babel language.
% Also neutralize the display-only normalfont direction reset locally. A
% direction-only patch inverted 12 to 21 inside equation in the measured probe.
% Plain maketag is local to the reference: mathtools' display-side wrapper
% otherwise globally consumes manual_tag while formatting a nested reference.
% Preserve each original command (including the selective auxiliary write),
% current tag forms, native links, and showonlyrefs on/off transitions.
\makeatletter
\NewCommandCopy{\HEB@nativeeqref}{\eqref}
\newcommand{\HEB@eqrefisolate}[1]{%
  \babelsublr{\HEBdisptagfalse\let\bbl@eqnodir\relax\let\maketag@@@\ORIGmaketag@@@#1}}
\protected\def\eqref#1{\HEB@eqrefisolate{\HEB@nativeeqref{#1}}}
\ExplSyntaxOn
\NewCommandCopy{\HEB@nativeMTeqref}{\MT_extended_eqref:n}
\cs_set_protected:Npn \MT_extended_eqref:n #1
  {\HEB@eqrefisolate{\HEB@nativeMTeqref{#1}}}
\ExplSyntaxOff
\makeatother

% v4.11.0: custom caption counters are not part of hyperref's English
% locale, and Hebrew has no upstream hyperref locale in this toolchain.
% Restore the appropriate names at every Babel language switch, including
% return from a nested inline island. Do not localize the destination keys.
\makeatletter
\newcommand{\HEB@refnames}{%
\def\chapterautorefname{פרק}%
\def\sectionautorefname{סעיף}%
\def\subsectionautorefname{תת־סעיף}%
\def\subsubsectionautorefname{תת־תת־סעיף}%
\def\paragraphautorefname{פסקה}%
% v4.31: hyperref uses a distinct counter name at the deepest standard level.
\def\subparagraphautorefname{תת־פסקה}%
\def\equationautorefname{משוואה}%
\def\figureautorefname{איור}%
\def\tableautorefname{טבלה}%
\def\partautorefname{חלק}%
\def\appendixautorefname{נספח}%
\def\footnoteautorefname{הערת שוליים}%
\def\itemautorefname{סעיף}%
\def\Itemautorefname{סעיף}%
\def\figautorefname{איור}%
\def\tblautorefname{טבלה}%
\def\pageautorefname{עמוד}% 
}
\addto\extrashebrew{\HEB@refnames}
\newcommand{\ENG@refnames}{%
  \def\subparagraphautorefname{subparagraph}%
  \def\chapterautorefname{chapter}%
  \def\sectionautorefname{section}%
  \def\subsectionautorefname{subsection}%
  \def\subsubsectionautorefname{subsubsection}%
  \def\paragraphautorefname{paragraph}%
  \def\equationautorefname{Equation}%
  \def\figureautorefname{Figure}%
  \def\tableautorefname{Table}%
  \def\partautorefname{Part}%
  \def\appendixautorefname{Appendix}%
  \def\footnoteautorefname{footnote}%
  \def\itemautorefname{item}%
  \def\figautorefname{Figure}%
  \def\tblautorefname{Table}%
  \def\Itemautorefname{item}%
  \def\pageautorefname{page}%
}
\addto\extrasenglish{\ENG@refnames}
\ifx\lakitmainlanguage\hllpdfenglish\ENG@refnames\else\HEB@refnames\fi
\makeatother

% ============================================================================
% Optional running head (v3.39): \lakithead{<book title>} — OFF by default.
% A document opts in with ONE line after \input{preamble.tex}:
%   \lakithead{מבנה המחשב}
% Every content page then carries the series' orientation line: chapter name
% at the reading edge, book title opposite, thin rule. (fancyhdr's [L]/[R]
% are PHYSICAL; under the RTL page direction [L] lands at the RIGHT/reading
% edge — measured, consistent across the corpus.) Packaged like the topic
% rail: always loaded, inert without the call. This centralizes the 8-line
% block that 8 lakits used to copy by hand — the same duplication vector
% that once let a lakit preamble go ~100 lines stale. Class-aware: report
% fills the mark from \chapter, article from \section. NB: a titlepage env
% self-empties; front matter that sets \pagestyle{empty} must restore with
% \pagestyle{fancy} afterwards (NOT plain — that discards the head).
% ============================================================================
\usepackage{fancyhdr}
% v3.61: \chapter* left a STALE running head — pages of an UNNUMBERED chapter
% carried the PREVIOUS chapter's title. report.cls routes \chapter* through
% \@schapter, which never calls \chaptermark, so \leftmark keeps whatever the last
% numbered chapter left. Measured: with \chapter{פרק ממוספר} then
% \chapter*{מבוא לא ממוספר}, page 4 (inside the starred chapter) printed
% "פרק ממוספר". Invisible on the chapter's FIRST page (\thispagestyle{plain}), so it
% only shows from the second page on — the pages an author is least likely to check.
% An unnumbered front-matter chapter (\chapter*{מבוא}, \chapter*{הקדמה}) is everyday
% Hebrew-academic structure, and the corpus already carries HAND-WRITTEN \markboth
% lines right after \chapter* — the exact per-document duplication \lakithead exists
% to retire. Patched INSIDE \lakithead so the macro stays inert when unused; a
% document's own later \markboth still wins, so those hand-patches do not regress.
\makeatletter
% <hllpdf:starred-heading-scope>
\newif\iflakin@starmarkpatched
\newcommand{\lakistarmark}{%
  \iflakin@starmarkpatched\else
  % Dispatch on the actual class counter, not hyperref's shared internals.
  \ifdefined\c@chapter
    \let\HEBorigschapter\@schapter
    % Record only after the native heading has reached its page. A mark
    % before the native routine can advertise the next heading too early.
    \def\@schapter##1{\HEBorigschapter{##1}\chaptermark{##1}}%
  \else
    % @ssect receives layout arguments, NOT the section name or level.
    % Capture the name in @startsection without introducing a TeX group:
    % grouping sectioning would break the native after-heading state.
    \def\HEB@sectionkind{}%
    \def\HEB@mainsection{section}%
    \let\HEBorigstartsection\@startsection
    \def\@startsection##1##2##3##4##5##6{%
      \def\HEB@sectionkind{##1}%
      \HEBorigstartsection{##1}{##2}{##3}{##4}{##5}{##6}}%
    \let\HEBorigssect\@ssect
    \def\@ssect##1##2##3##4##5{%
      \HEBorigssect{##1}{##2}{##3}{##4}{##5}%
      % Lower starred headings must not replace the article's section mark.
      \ifx\HEB@sectionkind\HEB@mainsection
        \sectionmark{##5}%
      \fi
      \ignorespaces}%
  \fi
  \global\lakin@starmarkpatchedtrue
  \fi}
% </hllpdf:starred-heading-scope>
% <hllpdf:running-language>
% Choose the wrapper when the mark is written, not when a later page ships.
% English punctuation must remain LTR even after returning to Hebrew prose.
\newcommand{\HEB@runningmark}[1]{%
  \iflanguage{english}{\markboth{\en{#1}}{}}{\markboth{\h{#1}}{}}}
% </hllpdf:running-language>
\makeatother
% v4.2.1 (NAV): every collection already names itself in \lakithead — feed the
% PDF /Title from the same declaration, so the navigation gate's
% PDF_REQUIRE_TITLE has a universal source and no fleet member ships
% title-less by missing a second declaration site. Measured before this: 14
% of 16 collections shipped /Title-empty while every one carried a perfectly
% good \lakithead string (plain Hebrew in all 16 — no math, no formatting, so
% the PDF-string path is clean). An explicit title always wins: set it with
% \lakitpdftitle{...} anywhere before \lakithead, or with a plain
% \hypersetup{pdftitle=...} — the fallback detects that via hyperref's
% \@pdftitle, probe-verified \@empty-equal until the first explicit set.
\makeatletter
\newif\iflakit@pdftitleset
\newcommand{\lakitpdftitle}[1]{%
  \hypersetup{pdftitle={#1}}\global\lakit@pdftitlesettrue}
\newcommand{\lakithead}[1]{%
  \iflakit@pdftitleset\else
    \ifx\@pdftitle\@empty
      \hypersetup{pdftitle={#1}}\global\lakit@pdftitlesettrue
    \fi
  \fi
  \pagestyle{fancy}\fancyhf{}%
  \lakistarmark

  \fancyhead[L]{\small\textcolor{cGray}{\leftmark}}%
  % <hllpdf:fixed-head-language>
  % The fixed document title has its OWN language snapshot, independent of
  % the current section/chapter mark. Choose the branch NOW, not at shipout:
  % the page may be Hebrew by the time an English title is finally printed.
  \iflanguage{english}{%
    \fancyhead[R]{\small\textcolor{cGray}{\en{#1}}}%
  }{%
    \fancyhead[R]{\small\textcolor{cGray}{\h{#1}}}%
  }%
  % </hllpdf:fixed-head-language>
  \fancyfoot[C]{\thepage}%
  \renewcommand{\headrulewidth}{0.4pt}%
  \ifdefined\chaptermark
    \renewcommand{\chaptermark}[1]{\HEB@runningmark{##1}}%
  \else
    \renewcommand{\sectionmark}[1]{\HEB@runningmark{##1}}%
  \fi}
\makeatother

% v3.39.1: path-independent byte-reproducibility. LuaTeX derives the trailer
% /ID from the OUTPUT PATH, so two clones of the same source at different
% directories produced different sha256 (measured: the ONLY delta was the
% 62-byte /ID line). Bit 512 suppresses exactly that line — measured: Creator/
% Producer/CreationDate/ModDate and all content untouched. /ID is optional
% per ISO 32000 (required only for ENCRYPTED PDFs and PDF/A — neither is a
% target here). If a future workflow needs an ID (PDF/A, Acrobat revision
% matching), delete this line and it returns.
\pdfvariable suppressoptionalinfo \numexpr 512\relax

% v3.75: BOLD MATH IN BOLD HEADINGS. Heading text is \bfseries but math inside
% it stayed regular-weight CM — a hairline $V_T$ inside a heavy 24.8pt chapter
% title (measured: CMMI12 fl=6 next to FrankRuhlLibre-Bold fl=20; 155 live
% heading-with-math sites fleet-wide). \boldmath rides along wherever the class
% sets \bfseries for a heading; A/B measured: $V_T$ -> CMMIB10 fl=22, 0 missing
% characters (the bold CM faces ship with this TeX Live). Guarded per-macro so
% the article class (no \@makechapterhead) and future class changes stay safe.
% v4.0.1 (adoption fix): the SEARCH patterns above carry no parameter token.
% v4.0.0 made them specific by including the class's own #1/#2, and measured on
% TeX Live 2023 that never matches: a \patchcmd search pattern containing a
% parameter token cannot match while the call sits inside \@ifundefined{...}{}{...}
% (probed both #1 and ##1 -> FAIL; the same pattern at top level succeeds). Every
% one of these patches therefore failed, and because v4.0.0 also replaced the
% empty failure branch with \PackageError, the failure became FATAL: all four
% bundled example assets stopped building ("Cannot patch chapter-title math
% weight -> Fatal error, no output PDF"). The \PackageError branches are KEPT —
% a silent no-op here is what let v3.75's fix rot unnoticed — only the parameter
% tokens are dropped, so the patches match and a real class change still errors.
\makeatletter
\@ifundefined{@makechapterhead}{}{%
  \patchcmd{\@makechapterhead}{\huge\bfseries \@chapapp}
    {\huge\bfseries\boldmath \@chapapp}{}
    {\PackageError{hebrew-lualatex-pdf}{Cannot patch chapter-number math weight}{Unsupported document class}}%
  \patchcmd{\@makechapterhead}{\Huge \bfseries}
    {\Huge \bfseries\boldmath}{}
    {\PackageError{hebrew-lualatex-pdf}{Cannot patch chapter-title math weight}{Unsupported document class}}%
  \patchcmd{\@makeschapterhead}{\Huge \bfseries}
    {\Huge \bfseries\boldmath}{}
    {\PackageError{hebrew-lualatex-pdf}{Cannot patch starred chapter math weight}{Unsupported document class}}}
% v4.0.1: \part sits under a THREE-DEEP wrapper chain, so the class-based anchor
% is not in \@part at all. Measured with \show: hyperref rewrites \@part to
% `...\phantomsection\fi \H@old@part`, and nameref rewrites THAT to
% `\NR@gettitle{#1}\NR@part[{#1}]{#2}` — the class body ends up in \NR@part.
% v4.0.0 patched \@part and made failure fatal, which is why all seven bundled
% examples stopped building. \part is NOT a dead heading level here: 188 live uses
% across the fleet, so the bold-math fix is worth getting right rather than
% dropping. We walk the chain innermost-first and stop at the first hit. Failure is
% a WARNING, not an error: the chain is a property of hyperref/nameref versions, so
% a future package release must degrade to regular-weight part math — a cosmetic
% loss — instead of refusing to build. Chapter and section keep \PackageError,
% because those targets are in the CLASS and were verified to patch here.
\newif\iflakin@partnum \newif\iflakin@parttitle \newif\iflakin@sparttitle
\def\lakin@partnum#1#2{\iflakin@partnum\else
  \patchcmd{#1}{#2\bfseries \partname}{#2\bfseries\boldmath \partname}
    {\lakin@partnumtrue}{}\fi}
\def\lakin@parttitle#1#2{\iflakin@parttitle\else
  \patchcmd{#1}{#2 \bfseries}{#2 \bfseries\boldmath}{\lakin@parttitletrue}{}\fi}
\def\lakin@sparttitle#1#2{\iflakin@sparttitle\else
  \patchcmd{#1}{#2 \bfseries}{#2 \bfseries\boldmath}{\lakin@sparttitletrue}{}\fi}
\@ifundefined{c@chapter}{\def\lakin@pnum{\Large}\def\lakin@ptit{\huge}}%
                        {\def\lakin@pnum{\huge}\def\lakin@ptit{\Huge}}
\@ifundefined{NR@part}{}{\expandafter\lakin@partnum\expandafter\NR@part\expandafter{\lakin@pnum}%
                         \expandafter\lakin@parttitle\expandafter\NR@part\expandafter{\lakin@ptit}}
\@ifundefined{H@old@part}{}{\expandafter\lakin@partnum\expandafter\H@old@part\expandafter{\lakin@pnum}%
                            \expandafter\lakin@parttitle\expandafter\H@old@part\expandafter{\lakin@ptit}}
\@ifundefined{@part}{}{\expandafter\lakin@partnum\expandafter\@part\expandafter{\lakin@pnum}%
                       \expandafter\lakin@parttitle\expandafter\@part\expandafter{\lakin@ptit}}
\@ifundefined{NR@spart}{}{\expandafter\lakin@sparttitle\expandafter\NR@spart\expandafter{\lakin@ptit}}
\@ifundefined{H@old@spart}{}{\expandafter\lakin@sparttitle\expandafter\H@old@spart\expandafter{\lakin@ptit}}
\@ifundefined{@spart}{}{\expandafter\lakin@sparttitle\expandafter\@spart\expandafter{\lakin@ptit}}
\iflakin@partnum\else\PackageWarning{hebrew-lualatex-pdf}{Part-number math stays
  regular weight: no known \noexpand\part wrapper matched (hyperref/nameref
  version change). Cosmetic only}\fi
\iflakin@parttitle\else\PackageWarning{hebrew-lualatex-pdf}{Part-title math stays
  regular weight: no known \noexpand\part wrapper matched. Cosmetic only}\fi
% v4.12.0: THE PART NUMBER IN THE TOC WAS NOT DIRECTIONALLY SEPARATED FROM THE
% TITLE. report.cls/article.cls compose the part entry as ONE text stream --
% \addcontentsline{toc}{part}{\thepart\hspace{1em}#1} -- while EVERY other level
% goes through \numberline, which puts the number in an \hbox. A box is a bidi
% OBJECT and cannot join a neighbouring run; a bare "IV" can. So in the RTL TOC
% paragraph a Latin number and a Latin-INITIAL title merge into a SINGLE LTR run
% and are laid out left-to-right as a unit: the number lands on the wrong side of
% the row, and a mixed title is scrambled outright. Measured on \part{MCP ...}:
% the printed order was "<tail> V MCP" -- title fragment, number, title fragment.
% The chapter/section rows one line below, with the SAME Latin-initial titles,
% were correct on the same page. That contrast is the whole diagnosis.
% \en{} IS NOT A WORKAROUND, measured both ways on one fixture: it fixes a
% PURE-Latin title (\part{\en{Hooks}} renders correctly) and does NOT fix a MIXED
% one (\part{\en{MCP} <hebrew>} stays scrambled). babel's language switch changes
% font and hyphenation, not the bidi run structure -- so the defect is not
% reachable from the source side at all, which is why it belongs here.
% Boxing the number is the entire fix and is a STRICT NO-OP where the bug never
% bit: a Hebrew-titled document compiles BYTE-IDENTICAL before and after (md5,
% 3 clean runs), because there the neutral glue already sat between an L run and
% an R run. The PDF outline is byte-identical too -- hyperref's \pdfstringdef
% neutralises \mbox, so bookmark strings keep their part numbers.
% \numberline{\thepart} would also fix the direction, but it re-geometries every
% part row in the fleet and DROPS the number from the part bookmark (hyperref
% strips \numberline). \mbox changes nothing that was already right.
% The search pattern deliberately carries no # token: etoolbox cannot match a
% parameter token in a patch pattern (measured -- all three chain members
% reported "not patchable" while the pattern included #1). "\thepart \hspace
% {1em}" is brace-balanced and occurs exactly once in the wrapper body.
\newif\iflakin@parttocnum
\def\lakin@parttocnum#1{\iflakin@parttocnum\else
  \patchcmd{#1}{\thepart \hspace {1em}}{\mbox {\thepart }\hspace {1em}}
    {\lakin@parttocnumtrue}{}\fi}
\@ifundefined{NR@part}{}{\lakin@parttocnum\NR@part}
\@ifundefined{H@old@part}{}{\lakin@parttocnum\H@old@part}
\@ifundefined{@part}{}{\lakin@parttocnum\@part}
\iflakin@parttocnum\else\PackageWarning{hebrew-lualatex-pdf}{Part TOC number is
  not boxed: no known \noexpand\part wrapper matched. A part title that BEGINS
  with Latin will print its number on the wrong side of its TOC row}\fi
\patchcmd{\section}{\bfseries}{\bfseries\boldmath}{}
  {\PackageError{hebrew-lualatex-pdf}{Cannot patch section math weight}{Unsupported document class}}
\patchcmd{\subsection}{\bfseries}{\bfseries\boldmath}{}
  {\PackageError{hebrew-lualatex-pdf}{Cannot patch subsection math weight}{Unsupported document class}}
\patchcmd{\subsubsection}{\bfseries}{\bfseries\boldmath}{}
  {\PackageError{hebrew-lualatex-pdf}{Cannot patch subsubsection math weight}{Unsupported document class}}
\patchcmd{\paragraph}{\bfseries}{\bfseries\boldmath}{}{}
\patchcmd{\subparagraph}{\bfseries}{\bfseries\boldmath}{}{}
\@ifundefined{l@chapter}{}{%
  \patchcmd{\l@chapter}{\bfseries}{\bfseries\boldmath}{}{}}
\makeatother
\setcounter{tocdepth}{\lakittocdepth}
% <hllpdf:measured-toc>
% v4.16.0: Measure actual number labels, not a guessed digit count. Read only
% canonical contentsline records; never execute titles or addtocontents code
% during the measuring pass. Widths are local to this individual list and are
% available before its first row: no per-row drift and no extra aux pass.
\makeatletter
\newread\hllpdf@tocstream
\newdimen\hllpdf@toc@chapter
\newdimen\hllpdf@toc@section
\newdimen\hllpdf@toc@subsection
\newdimen\hllpdf@toc@subsubsection
\newdimen\hllpdf@toc@paragraph
\newdimen\hllpdf@toc@subparagraph
\newdimen\hllpdf@toc@figure
\newdimen\hllpdf@toc@table
\newdimen\hllpdf@tocpage
\newdimen\hllpdf@toc@sectionindent
\newdimen\hllpdf@toc@subsectionindent
\newdimen\hllpdf@toc@subsubsectionindent
\newdimen\hllpdf@toc@paragraphindent
\newdimen\hllpdf@toc@subparagraphindent
\def\hllpdf@tocdefaults{%
  \hllpdf@toc@chapter=2em
  \hllpdf@toc@section=3.2em
  \hllpdf@toc@subsection=3.9em
  \@ifundefined{c@chapter}{%
    \hllpdf@toc@subsubsection=3.2em
    \hllpdf@toc@paragraph=4.1em
    \hllpdf@toc@subparagraph=5em
  }{%
    \hllpdf@toc@subsubsection=4.1em
    \hllpdf@toc@paragraph=5em
    \hllpdf@toc@subparagraph=6em
  }%
  \hllpdf@toc@figure=2.3em
  \hllpdf@toc@table=2.3em
  \hllpdf@tocpage=\@pnumwidth\relax
}
\def\hllpdf@tocindents{%
  % Preserve the established shallow layout when its slots already fit.
  \hllpdf@toc@sectionindent=1.5em
  \@tempdima=\hllpdf@toc@chapter \advance\@tempdima by -2em
  \advance\hllpdf@toc@sectionindent by \@tempdima
  \hllpdf@toc@subsectionindent=\dimexpr
    \hllpdf@toc@sectionindent+\hllpdf@toc@section\relax
  \hllpdf@toc@subsubsectionindent=\dimexpr
    \hllpdf@toc@subsectionindent+\hllpdf@toc@subsection\relax
  \hllpdf@toc@paragraphindent=\dimexpr
    \hllpdf@toc@subsubsectionindent+\hllpdf@toc@subsubsection\relax
  \hllpdf@toc@subparagraphindent=\dimexpr
    \hllpdf@toc@paragraphindent+\hllpdf@toc@paragraph\relax
}
\def\hllpdf@tocfont@chapter{\bfseries\boldmath}
\def\hllpdf@tocfont@part{\large\bfseries\boldmath}
\def\hllpdf@tocentryfont{%
  \ifcsname hllpdf@tocfont@\hllpdf@toclevel\endcsname
    \csname hllpdf@tocfont@\hllpdf@toclevel\endcsname
  \fi}
\long\def\hllpdf@tocdiscard#1\hllpdf@tocstop{}
\def\hllpdf@tocscanrecord{%
  \@ifnextchar\contentsline{\hllpdf@tocscanline}{%
    \@ifnextchar\babel@toc{\hllpdf@tocscanlanguage}{\hllpdf@tocdiscard}}}
\long\def\hllpdf@tocscanlanguage\babel@toc#1#2#3\hllpdf@tocstop{%
  % Babel records are declarations, not arbitrary code. Track the language
  % without executing the original directive or changing the surrounding list.
  \def\hllpdf@toclanguage{#1}}
\long\def\hllpdf@tocscanline\contentsline#1#2#3#4#5\hllpdf@tocstop{%
  \def\hllpdf@toclevel{#1}%
  \ifcsname hllpdf@toc@#1\endcsname
    \hllpdf@tocscantitle#2\hllpdf@tocstop
  \fi
  % The page-number column needs its own measurement (including bold chapter
  % rows). Its extra clearance separates the page label from dot leaders.
  \setbox\@tempboxa=\hbox{\foreignlanguage{\hllpdf@toclanguage}{\hllpdf@tocentryfont #3\kern.3em}}%
  \ifdim\wd\@tempboxa>\hllpdf@tocpage
    \hllpdf@tocpage=\wd\@tempboxa
  \fi}
\def\hllpdf@tocscantitle{%
  \@ifnextchar\numberline{\hllpdf@tocscanlabel}{\hllpdf@tocdiscard}}
\long\def\hllpdf@tocscanlabel\numberline#1#2\hllpdf@tocstop{%
  \setbox\@tempboxa=\hbox{\foreignlanguage{\hllpdf@toclanguage}{\hllpdf@tocentryfont #1\kern.5em}}%
  \ifdim\wd\@tempboxa>\csname hllpdf@toc@\hllpdf@toclevel\endcsname
    \csname hllpdf@toc@\hllpdf@toclevel\endcsname=\wd\@tempboxa
  \fi}
\def\hllpdf@tocread{%
  \ifeof\hllpdf@tocstream\else
    \read\hllpdf@tocstream to\hllpdf@tocrecord
    \expandafter\hllpdf@tocscanrecord\hllpdf@tocrecord\hllpdf@tocstop
    \expandafter\hllpdf@tocread
  \fi}
\def\hllpdf@tocmeasure#1{%
  \makeatletter
  \edef\hllpdf@toclanguage{\languagename}%
  \hllpdf@tocdefaults
  \openin\hllpdf@tocstream=\jobname.#1\relax
  \hllpdf@tocread
  \closein\hllpdf@tocstream
  \hllpdf@tocindents
  \edef\@pnumwidth{\the\hllpdf@tocpage}%
  \@tempdima=\dimexpr\hllpdf@tocpage+1em\relax
  \ifdim\@tempdima>\@tocrmarg\relax
    \edef\@tocrmarg{\the\@tempdima}%
  \fi}
\let\hllpdf@originalstarttoc\@starttoc
\def\@starttoc#1{%
  \begingroup
    \def\hllpdf@tocext{#1}%
    \def\hllpdf@tocwanted{toc}%
    \ifx\hllpdf@tocext\hllpdf@tocwanted \hllpdf@tocmeasure{#1}\else
      \def\hllpdf@tocwanted{lof}%
      \ifx\hllpdf@tocext\hllpdf@tocwanted \hllpdf@tocmeasure{#1}\else
        \def\hllpdf@tocwanted{lot}%
        \ifx\hllpdf@tocext\hllpdf@tocwanted \hllpdf@tocmeasure{#1}\fi
      \fi
    \fi
    \hllpdf@originalstarttoc{#1}%
  \endgroup}
\hllpdf@tocdefaults
\hllpdf@tocindents
\@ifundefined{l@chapter}{}{%
  \patchcmd{\l@chapter}{\setlength\@tempdima{1.5em}}%
    {\setlength\@tempdima{\hllpdf@toc@chapter}}{}%
    {\PackageError{hebrew-lualatex-pdf}{Cannot size chapter TOC numbers}%
       {Unsupported contents layout: the chapter number slot was not found}}}
\renewcommand*\l@section{\@dottedtocline{1}%
  {\hllpdf@toc@sectionindent}{\hllpdf@toc@section}}
\renewcommand*\l@subsection{\@dottedtocline{2}%
  {\hllpdf@toc@subsectionindent}{\hllpdf@toc@subsection}}
\renewcommand*\l@subsubsection{\@dottedtocline{3}%
  {\hllpdf@toc@subsubsectionindent}{\hllpdf@toc@subsubsection}}
\renewcommand*\l@paragraph{\@dottedtocline{4}%
  {\hllpdf@toc@paragraphindent}{\hllpdf@toc@paragraph}}
\renewcommand*\l@subparagraph{\@dottedtocline{5}%
  {\hllpdf@toc@subparagraphindent}{\hllpdf@toc@subparagraph}}
\renewcommand*\l@figure{\@dottedtocline{1}{1.5em}{\hllpdf@toc@figure}}
\renewcommand*\l@table{\@dottedtocline{1}{1.5em}{\hllpdf@toc@table}}
\makeatother
% </hllpdf:measured-toc>
\setcounter{secnumdepth}{\lakitsecnumdepth}
% Public preamble-only caption-numbering policy. The historical within-major
% default is untouched unless this command is called. "continuous" removes the
% reset edge as well as the printed prefix, so Figure 1/2 and Table 1/2 stay
% unique across sections/chapters. Returning to "within-major" restores the
% original class-aware displays and native destination conventions. The policy
% never replaces \figcap/\tabcap or their measured link-placement mechanism.
\makeatletter
\let\lakin@majorfig\thefig
\let\lakin@majortbl\thetbl
\let\lakin@majorHfig\theHfig
\let\lakin@majorHtbl\theHtbl
\newcommand{\lakitcaptionnumbering}[1]{%
  \edef\lakin@numbering{#1}%
  \def\lakin@continuous{continuous}\def\lakin@withinmajor{within-major}%
  \ifx\lakin@numbering\lakin@continuous
    \@ifundefined{c@chapter}{%
      \@removefromreset{fig}{section}\@removefromreset{tbl}{section}%
    }{\@removefromreset{fig}{chapter}\@removefromreset{tbl}{chapter}}%
    \renewcommand{\thefig}{\arabic{fig}}\renewcommand{\thetbl}{\arabic{tbl}}%
    \def\theHfig{\arabic{fig}}\def\theHtbl{\arabic{tbl}}%
  \else\ifx\lakin@numbering\lakin@withinmajor
    \@ifundefined{c@chapter}{%
      \@removefromreset{fig}{section}\@removefromreset{tbl}{section}%
      \@addtoreset{fig}{section}\@addtoreset{tbl}{section}%
    }{%
      \@removefromreset{fig}{chapter}\@removefromreset{tbl}{chapter}%
      \@addtoreset{fig}{chapter}\@addtoreset{tbl}{chapter}}%
    \let\thefig\lakin@majorfig\let\thetbl\lakin@majortbl
    \let\theHfig\lakin@majorHfig\let\theHtbl\lakin@majorHtbl
  \else
    \PackageError{hebrew-english-latex-studio}{Unknown caption-numbering policy '#1'}%
      {Use continuous or within-major before begin{document}.}%
  \fi\fi}
\@onlypreamble\lakitcaptionnumbering
\makeatother
% Architecture readiness sentinel. Optional runtime layers require the complete
% core, so the marker is defined only after every mechanism/baseline definition.
% <hllpdf:architecture-ready>
\def\hllpdfcoreloaded{1}
% </hllpdf:architecture-ready>

% v4.22: registered file resources, not an assertion about embedded subset bytes.
% The log framing avoids long-path wrapping and preserves UTF-8 as hex.
% <hllpdf:font-inventory>
\directlua{
  local function inventory()
    local names = {}
    local unnamed = 0
    for id, item in font.each() do
      local path = item.filename
      if path and not (path == "") then names[path] = true else unnamed = unnamed + 1 end
    end
    local sorted = {}
    for path in pairs(names) do table.insert(sorted,path) end
    table.sort(sorted)
    texio.write_nl("log","HLLFONT|BEGIN|1")
    for index,path in ipairs(sorted) do
      local hex = path:gsub(".",function(c) return string.format(string.char(37) .. "02x",string.byte(c)) end)
      local count = math.ceil(string.len(hex)/48)
      texio.write_nl("log","HLLFONT|FILE|" .. index .. "|" .. count)
      for offset=1,string.len(hex),48 do texio.write_nl("log","HLLFONT|HEX|" .. hex:sub(offset,offset+47)) end
    end
    texio.write_nl("log","HLLFONT|END|" .. rawlen(sorted) .. "|" .. unnamed)
    texio.write_nl("log","")
  end
  luatexbase.add_to_callback("finish_pdffile",inventory,"hllpdf.runtime.fonts")
}
% </hllpdf:font-inventory>
     % (or paste these lines at the top)
%   \begin{document}  ...  \end{document}
%
% Compile with:  lualatex -interaction=nonstopmode doc.tex
%   Repeat until two consecutive PDFs are byte-identical (scripts/build.sh
%   does this) — a TOC document converges only on pass 3+; after two passes
%   the TOC page numbers are STALE (measured; see SKILL.md step 3).
% Font binaries are external dependencies, pinned by assets/fonts/SHA256SUMS.
% scripts/setup_fonts.sh imports a verified local installation or prior skill.
%
% PUBLIC DOCUMENT-SHELL / LAYOUT CONTRACT (introduced v4.3.1; architecture-owned v4.7.1). These hooks may be
% defined BEFORE % ============================================================================
% preamble.tex — tested LuaLaTeX preamble for Hebrew + math + figures PDFs (v4.32.0)
% Battle-tested: compiled an 80-page Hebrew quantum-computing summary AND a
% 197-page Hebrew computer-architecture anthology (TikZ diagrams, code listings,
% colored boxes) — both with 0 missing characters.
%
% Usage in your document:
%   \input{preamble.tex}     % (or paste these lines at the top)
%   \begin{document}  ...  \end{document}
%
% Compile with:  lualatex -interaction=nonstopmode doc.tex
%   Repeat until two consecutive PDFs are byte-identical (scripts/build.sh
%   does this) — a TOC document converges only on pass 3+; after two passes
%   the TOC page numbers are STALE (measured; see SKILL.md step 3).
% Font binaries are external dependencies, pinned by assets/fonts/SHA256SUMS.
% scripts/setup_fonts.sh imports a verified local installation or prior skill.
%
% PUBLIC DOCUMENT-SHELL / LAYOUT CONTRACT (introduced v4.3.1; architecture-owned v4.7.1). These hooks may be
% defined BEFORE \input{preamble.tex}; if omitted, the historical canonical
% defaults below are exact.  They deliberately control only the document shell
% and page layout.  RTL/math semantics stay in this preamble, while visual
% presentation themes (for example clean_theme.tex) load AFTER it.
%
%   \def\lakitdocumentclass{article}
%   \def\lakitclassoptions{11pt,a4paper}
%   \def\lakitgeometryoptions{margin=1.8cm,headheight=14pt}
%   \def\lakitlinespread{1.15}
%   \def\lakitparagraphstyle{block} % letters/briefs; indented for continuous prose
%   \def\lakittocdepth{1}
%   \def\lakitsecnumdepth{2}
%   \input{preamble.tex}
%
% Supported class contract: report or article.  Alternate paper sizes must also
% be declared to build.sh through EXPECTED_PAPER; the PDF gate never infers them.
% Do NOT use these hooks as a theme mechanism: theme overlays must not own page
% geometry, class selection, or direction semantics. A declared presentation-policy
% overlay may intentionally change numbering visibility after the preamble.
% ============================================================================
% Runtime architecture guard: generated from scripts/architecture_contract.py.
% The readiness sentinel remains delayed until the END of this file.
% <hllpdf:architecture-guard>
\ifdefined\hllpdfcoreloaded
  \errmessage{hllpdf: preamble.tex loaded more than once}
  \expandafter\endinput
\fi
% </hllpdf:architecture-guard>

\ifdefined\lakitdocumentclass\else\def\lakitdocumentclass{report}\fi
\ifdefined\lakitclassoptions\else\def\lakitclassoptions{12pt,a4paper}\fi
\ifdefined\lakitgeometryoptions\else\def\lakitgeometryoptions{margin=2.3cm,headheight=14pt}\fi
\ifdefined\lakitlinespread\else\def\lakitlinespread{1.3}\fi
\ifdefined\lakittocdepth\else\def\lakittocdepth{1}\fi
\ifdefined\lakitsecnumdepth\else\def\lakitsecnumdepth{2}\fi
% Paragraph organization belongs to the document layout, independently of its
% language and theme. Preserve historical class dimensions unless explicitly
% selecting block paragraphs; never suppress indentation with an everypar hack.
\ifdefined\lakitparagraphstyle\else\def\lakitparagraphstyle{indented}\fi
\edef\lakitparagraphstyle{\lakitparagraphstyle}
\def\hllpdfparagraphblock{block}
\def\hllpdfparagraphindented{indented}
\ifx\lakitparagraphstyle\hllpdfparagraphblock\else
  \ifx\lakitparagraphstyle\hllpdfparagraphindented\else
    \errmessage{hllpdf: paragraph style must be block or indented}
  \fi
\fi
% v4.3.1: the public class hook has a deliberately narrow support contract.
% Before this guard, \def\lakitdocumentclass{book} silently compiled even though
% every heading/TOC regression test targets report/article; an unsupported shell
% could therefore look "supported" until a subtle class-specific defect shipped.
% Fail before \documentclass rather than letting an undocumented combination enter
% the mechanism core. Plain-TeX comparison is used because LaTeX packages are not
% available yet at this point.
\edef\lakitvalidatedclass{\lakitdocumentclass}
\def\lakitreportclass{report}
\def\lakitarticleclass{article}
\ifx\lakitvalidatedclass\lakitreportclass
\else\ifx\lakitvalidatedclass\lakitarticleclass
\else
  \errmessage{hllpdf: unsupported class \lakitdocumentclass; use report or article}
\fi\fi
\let\lakitvalidatedclass\relax
\let\lakitreportclass\relax
\let\lakitarticleclass\relax
% Language intent is independent of document class, theme and content domain.
% The two dominant-language profiles differ in authoring intent, not direction.
% A balanced bilingual document must explicitly choose a navigation/base locale.
% <hllpdf:language-profile>
\providecommand{\lakitlanguageprofile}{hebrew}
\edef\lakitlanguageprofile{\lakitlanguageprofile}
\def\hllpdfhebrew{hebrew}\def\hllpdfenglish{english}
\def\hllpdfhebrewmixed{hebrew-with-english}
\def\hllpdfenglishmixed{english-with-hebrew}
\def\hllpdfbilingual{bilingual}
\let\hllpdfexpectedlanguage\relax
\ifx\lakitlanguageprofile\hllpdfhebrew
  \let\hllpdfexpectedlanguage\hllpdfhebrew
\else\ifx\lakitlanguageprofile\hllpdfhebrewmixed
  \let\hllpdfexpectedlanguage\hllpdfhebrew
\else\ifx\lakitlanguageprofile\hllpdfenglish
  \let\hllpdfexpectedlanguage\hllpdfenglish
\else\ifx\lakitlanguageprofile\hllpdfenglishmixed
  \let\hllpdfexpectedlanguage\hllpdfenglish
\else\ifx\lakitlanguageprofile\hllpdfbilingual
  \ifdefined\lakitmainlanguage\else
    \errmessage{hllpdf: bilingual profile requires explicit lakitmainlanguage}
  \fi
\else
  \errmessage{hllpdf: unsupported language profile \lakitlanguageprofile}
\fi\fi\fi\fi\fi
\ifdefined\lakitmainlanguage
  \edef\lakitmainlanguage{\lakitmainlanguage}
\else
  \let\lakitmainlanguage\hllpdfexpectedlanguage
\fi
\ifx\lakitmainlanguage\hllpdfhebrew\else
  \ifx\lakitmainlanguage\hllpdfenglish\else
    \errmessage{hllpdf: main language must be hebrew or english}
  \fi
\fi
\ifx\hllpdfexpectedlanguage\relax\else
  \ifx\lakitmainlanguage\hllpdfexpectedlanguage\else
    \errmessage{hllpdf: main language conflicts with language profile}
  \fi
\fi
\typeout{HLLLANG|\lakitlanguageprofile|\lakitmainlanguage}
% </hllpdf:language-profile>
\documentclass[\lakitclassoptions]{\lakitdocumentclass}
\usepackage[no-math]{fontspec}        % (1) ★ no-math — THE critical line (see recipe.md)
% layout=lists asks Babel to adapt list parshape when paragraph and page
% directions differ. Without it, minority-language list bodies lose their
% indentation and labels hang outside the text block, despite a clean log.
% The upstream mechanism preserves ordinary same-direction list geometry.
\usepackage[bidi=basic,layout=lists]{babel}
\ifx\lakitmainlanguage\hllpdfenglish
  \babelprovide[import]{hebrew}
  \babelprovide[import,main]{english}
\else
  \babelprovide[import,main]{hebrew}
  \babelprovide[import]{english}
\fi
% v4.22: bind canonical families and fallbacks to the selected installation.
% Lookup by basename or family alone is not proof of the file actually used.
% <hllpdf:font-binding>
\directlua{
  local dir = os.getenv("HEBREW_LATEX_FONT_DIR")
  if not dir or dir == "" then
    local data = os.getenv("XDG_DATA_HOME")
    if not data or data == "" or not (data:sub(1,1) == "/") then
      data = (os.getenv("HOME") or "") .. "/.local/share"
    end
    dir = data .. "/fonts/hebrew-lualatex-pdf"
  end
  if not (dir:sub(1,1) == "/") then dir = lfs.currentdir() .. "/" .. dir end
  if dir:find("[",1,true) or dir:find("]",1,true) or dir:find(string.char(0),1,true)
     or dir:find(string.char(10),1,true) or dir:find(string.char(13),1,true) then
    tex.error("hllpdf: unsupported character in font directory")
  end
  if not (lfs.attributes(dir,"mode") == "directory") then
    tex.error("hllpdf: font directory is missing")
  end
  hllpdf_font_dir = dir:gsub("/+$", "") .. "/"
  tex.sprint(string.char(92) .. "def" .. string.char(92) .. "hllpdffontdir{")
  tex.sprint(-2,hllpdf_font_dir)
  tex.sprint("}")
  luaotfload.add_fallback("hebfb",{"[" .. hllpdf_font_dir .. "DejaVuSerif.ttf]:mode=node;"})
  luaotfload.add_fallback("hebfbbf",{"[" .. hllpdf_font_dir .. "DejaVuSerif-Bold.ttf]:mode=node;"})
}
% </hllpdf:font-binding>
% v3.22: Frank Ruhl Libre ships NO italic face. Instead of silently-upright
% \emph (v<=3.20) or bold-as-emphasis (v3.21), synthesize a slanted face with
% FakeSlant — real, VISIBLE italic emphasis with correct \emph nesting
% semantics (inner \emph toggles back upright). Verified empirically at
% 300dpi: slant clearly visible, 0 errors, 0 missing glyphs.
\babelfont{rm}[Path={\hllpdffontdir},RawFeature={fallback=hebfb},
  BoldFont={FrankRuhlLibre-Bold.ttf},
  BoldFeatures={RawFeature={fallback=hebfbbf}},   % v3.56 — see hebfbbf above
  ItalicFont={FrankRuhlLibre-Regular.ttf}, ItalicFeatures={FakeSlant=0.22},
  % Bold-italic uses the weight-matched fallback and a synthetic slant.
  BoldItalicFont={FrankRuhlLibre-Bold.ttf}, BoldItalicFeatures={RawFeature={fallback=hebfbbf},FakeSlant=0.22}
  ]{FrankRuhlLibre-Regular.ttf}
% Sans fallbacks preserve family and weight for the few missing codepoints.
\directlua{luaotfload.add_fallback("hebfbsans",{"[" .. hllpdf_font_dir .. "DejaVuSans.ttf]:mode=node;"})}
\directlua{luaotfload.add_fallback("hebfbsansbf",{"[" .. hllpdf_font_dir .. "DejaVuSans-Bold.ttf]:mode=node;"})}
\babelfont{sf}[Path={\hllpdffontdir},RawFeature={fallback=hebfbsans},
  BoldFont={Heebo-Bold.ttf},
  BoldFeatures={RawFeature={fallback=hebfbsansbf}}]{Heebo-Regular.ttf}
% v3.24: Scale=MatchLowercase — DejaVu Sans Mono has an outsized x-height, so
% unscaled inline code visually SHOUTS next to the delicate Frank Ruhl body
% (measured on a 197-page build: scaled, a code atom sits in the line at the
% body's weight, and the document tightened by ~5%). Applies to listings too.
\babelfont{tt}[Path={\hllpdffontdir},RawFeature={fallback=hebfb},Scale=MatchLowercase,
  BoldFont={DejaVuSansMono.ttf},
  BoldFeatures={RawFeature={fallback=hebfb},FakeBold=1.5},
  ItalicFont={DejaVuSansMono.ttf},
  ItalicFeatures={RawFeature={fallback=hebfb},FakeSlant=0.2},
  BoldItalicFont={DejaVuSansMono.ttf},
  BoldItalicFeatures={RawFeature={fallback=hebfb},FakeBold=1.5,FakeSlant=0.2}
  ]{DejaVuSansMono.ttf}

% Explicit multilingual/English profiles use real Latin faces and italics.
% The legacy/default Hebrew profile retains its established measured typography.
% Latin Modern is provided by the TeX distribution, not the pinned font bundle;
% recorder/runtime evidence reports it as an additional toolchain resource.
\ifx\lakitlanguageprofile\hllpdfhebrew\else
  \babelfont[english]{rm}[Ligatures=TeX,RawFeature={fallback=hebfb},
    BoldFeatures={RawFeature={fallback=hebfbbf}},
    BoldItalicFeatures={RawFeature={fallback=hebfbbf}}]{Latin Modern Roman}
  \babelfont[english]{sf}[Ligatures=TeX,RawFeature={fallback=hebfbsans},
    BoldFeatures={RawFeature={fallback=hebfbsansbf}},
    BoldItalicFeatures={RawFeature={fallback=hebfbsansbf}}]{Latin Modern Sans}
\fi

\usepackage{amsmath,amssymb,mathtools}
% amsthm supplies theorem/proof machinery. Environment names remain
% document-specific; the generic proof label is localized below.
\usepackage{amsthm}
% Proof names follow the surrounding prose, including return from an island.
\addto\extrashebrew{\renewcommand{\proofname}{הוכחה}}
\addto\extrasenglish{\renewcommand{\proofname}{Proof}}
\ifx\lakitmainlanguage\hllpdfenglish
  \renewcommand{\proofname}{Proof}
\else
  \renewcommand{\proofname}{הוכחה}
\fi
% v3.24.1: equation tags under babel-bidi. A tag re-emitted by mathtools'
% showonlyrefs inside align inherits the RTL base and its parens MIRROR —
% ")2.1(" instead of "(2.1)" (plain align/equation tags are unaffected;
% verified on a minimal case + a 198-page build). LTR-isolate at the tag-BOX
% level: mathtools' show/hide + numbering logic is untouched. Do NOT override
% \tagform@ instead — that clobbers the showonlyrefs mechanism (measured:
% every tag shows, numbering breaks).
% v3.28.1: the v3.24.1 wrap must apply ONLY to tags of display environments.
% Inline \eqref in Hebrew PROSE also goes through \maketag@@@, and there the
% english-wrap MIRRORS the parens — ")1(" — while the unwrapped original is
% correct (proven on a 2x2 isolation matrix: patch x showonlyrefs; regressed
% every prose \eqref since v3.24.1, caught by the render gate's R4). The two
% call sites are indistinguishable by \ifmmode (tag boxes are text mode), so
% the display environments themselves raise a flag; \eqref in prose sees it
% false and gets the original untouched \maketag@@@.
\usepackage{etoolbox}

% babel bidi=basic (including upstream 26.9) assumes every automatic
% discretionary has a one-glyph replacement. LuaTeX can legally produce one
% with replace=nil after a hyphenated Latin value; Babel then indexes nil and the
% whole build dies. Do the narrow upstream-compatible correction: while Babel
% classifies the list, present only that malformed disc as an ordinary
% discretionary, then restore its subtype. Babel still reorders the ENTIRE box;
% there is no pcall and no successful fallback to an unreordered node list.
% If Babel throws for any other reason, the build remains fatal. The synthetic
% nil-disc probe and the real table gallery in test_all.sh pin both layers.
\AtBeginDocument{\directlua{
  if Babel and Babel.bidi and not Babel.heb_nil_disc_compat then
    Babel.heb_nil_disc_compat = true
    local original = Babel.bidi
    local disc_id = node.id('disc')
    Babel.bidi = function(head, ispar, hdir)
      local hidden = {}
      for item in node.traverse_id(disc_id, head) do
        if item.subtype == 2 and (not item.replace or not item.replace.char) then
          table.insert(hidden, item)
          item.subtype = 0
        end
      end
      local result = original(head, ispar, hdir)
      for _, item in ipairs(hidden) do item.subtype = 2 end
      return result
    end
  end
}}
\makeatletter
\newif\ifHEBdisptag
\let\ORIGmaketag@@@\maketag@@@
\def\maketag@@@#1{\ifHEBdisptag\ORIGmaketag@@@{\foreignlanguage{english}{#1}}\else\ORIGmaketag@@@{#1}\fi}
\@for\HEB@env:=equation,equation*,align,align*,gather,gather*,multline,multline*,flalign,flalign*,alignat,alignat*,eqnarray,eqnarray*\do{%
  \expandafter\AtBeginEnvironment\expandafter{\HEB@env}{\HEBdisptagtrue}}
\makeatother
% English page + Hebrew prose: Babel's equation-only tag flip uses a primitive
% leqno path that shifts the formula to the left edge and the tag into the body
% of an English page. Align's full-width alignment is unaffected. Give only this
% display a physical LTR container and native amsmath eqno/leqno. The engine then
% owns real tag width, collision handling and labels; no proxy tag box hides ink
% from its fit calculation. Babel's saved \bbl@eqnodir still restores Hebrew text
% inside math, and the prose locale/math direction are never switched.
% Close the prior RTL paragraph before normalization. After the environment's
% group restores its directions, close TeX's empty resumed display paragraph and
% use native LaTeX continuation handling: same-paragraph prose has no indent,
% while an authored blank line starts an ordinary indented RTL paragraph.
% Register after Babel's setup; remain local to equation(*) in a main-English
% document's actual RTL prose. Native Hebrew, ordinary English and explicit
% fleqn retain their original paths. No page/display dimensions are assigned.
% <hllpdf:minority-rtl-equation-geometry>
\makeatletter
\newcommand{\HEB@minorityprosecontinuation}{\par\@doendpe}
\newcommand{\HEB@minorityequationgeometry}{%
  \ifx\lakitmainlanguage\hllpdfenglish
    \ifnum\pardirection=\@ne
      \if@fleqn\else
        \par\aftergroup\HEB@minorityprosecontinuation
        \iftagsleft@\let\veqno\@@leqno
        \else\let\veqno\@@eqno\fi
        \bodydir TLT\pardir TLT\textdir TLT\relax
      \fi
    \fi
  \fi}
\AtBeginDocument{%
  \AddToHook{env/equation/begin}{\HEB@minorityequationgeometry}%
  \AddToHook{env/equation*/begin}{\HEB@minorityequationgeometry}}
\makeatother
% </hllpdf:minority-rtl-equation-geometry>
% v3.63: \shortintertext rendered Hebrew MIRROR-REVERSED and flush LEFT — an entire
% commentary line inside an align printed backwards. babel's bidi=basic patches
% \intertext@ only (luababel.def), and mathtools' \shortintertext uses the identical
% \vbox{\normalbaselines … \noindent\ignorespaces#1\par} skeleton but is never
% patched, so its paragraph keeps pardir/textdir LTR: no bidi reordering (logical
% order on the page) and flush-left. Measured against \intertext with the same
% Applying babel's own replacement to \shortintertext@ makes the two identical.
% \h{}/\mbox fix the letter order but leave the line flush LEFT, so they are not
% substitutes. Guarded + \AtBeginDocument because mathtools \let-copies at load time.
% NB \ExplSyntaxOn is REQUIRED: mathtools' internals carry expl3 names, and without
% the expl3 catcodes the `_` is read as a subscript and the patch detonates far away
% with "! Illegal unit of measure (pt inserted)" at the \end{align} (measured).
\makeatletter
\AtBeginDocument{%
  \@ifundefined{shortintertext@}{}{%
    \ExplSyntaxOn
    \bbl@sreplace\shortintertext@{\normalbaselines}%
      {\normalbaselines\ifx\bbl@eqnodir\relax\else\bbl@pardir\@ne\bbl@eqnodir\fi}%
    \ExplSyntaxOff}}
\makeatother
% geometry inherits the canonical A4 sheet from \documentclass; build.sh then
% verifies every rendered MediaBox and CropBox, so a silent Letter fallback is
% impossible to publish as a clean default build.
\usepackage[\lakitgeometryoptions]{geometry}
\linespread{\lakitlinespread}       % layout hook; canonical default remains 1.3
\ifx\lakitparagraphstyle\hllpdfparagraphblock
  \setlength{\parindent}{0pt}
  \setlength{\parskip}{6pt plus 1pt minus 1pt}
\fi
% These dimensions remain ordinary scoped TeX lengths: authors can customize
% them after core input. Language switches inherit them; lists keep their own
% indentation and narrow reference layouts establish their local paragraph gaps.
\typeout{HLLPAR|\lakitparagraphstyle}
\usepackage{xcolor}
\usepackage{mdframed}
\usepackage{listings}                 % code listings (assembly / C / pseudo-code)
\usepackage{tikz}
\usetikzlibrary{arrows.meta,positioning,calc,shapes.geometric,fit,backgrounds}
\usepackage{pgfplots}\pgfplotsset{compat=1.18}
\usepackage{quantikz}                 % quantum circuits (remove if unused)
\usepackage{booktabs}
\usepackage{array}
\usepackage{float}                    % [H] = place table exactly here (reference doc: no drifting)
\usepackage{enumitem}
\usepackage{needspace}

% ============================================================================
% ★★ v3.4 — THE DIRECTION-ISOLATION TRIO (the heart of mixing Hebrew + Latin)
% Mental model:
%   \h{...}    Hebrew (RTL) island      — Hebrew text inside an LTR context
%                                          (TikZ nodes, the english `tikzpic`).
%   \en{...}   English (LTR) isolate    — Latin text / parentheses inside a
%                                          Hebrew (RTL) context. KEEPS PARENS
%                                          AND ORDER CORRECT.
%   \code{...} \en + monospace          — code inside Hebrew / inside a frametitle.
% CRITICAL: never wrap Latin in \h{} (it is for HEBREW). Inside a TikZ node,
% Latin in \h{} renders REVERSED (Datapath->htapataD, CPU->UPC) and parens
% mirror. Isolate every Latin/paren chunk with \en{} or \code{}. See
% references/figures-boxes-listings.md.
% ---- TEXT INSIDE MATH (v3.8) ----------------------------------------------
% Put ANY text that sits inside math in a \text{}. The ONLY thing that is
% invalid in math is a bare \h/\en (= \foreignlanguage) used DIRECTLY, outside
% \text -> "\rmfamily invalid in math mode". Inside \text:
%   Hebrew label           -> \text{...}            (bare Hebrew renders fine;
%                              \h is optional — BUT NOT ALWAYS HARMLESS (v3.67):
%                              if the \h group OPENS with a nested island
%                              ($..$, \en{}, \code{}, \texttt{}, \mbox{}) the
%                              WHOLE group loses bidi and the Hebrew prints in
%                              LOGICAL order = gibberish. Measured:
%                                \text{\h{$f$ קבועה}}  -> reads "העובק"  BROKEN
%                                \text{\h{קבועה $f$}}  -> correct
%                                \text{$f$ קבועה}      -> correct (no \h)
%                              Hebrew first inside \h, or drop the \h. Flagged
%                              H-OPENS-WITH-LTR-GROUP.)
%   plain Latin word        -> \text{...}            (no marker needed)
%   Latin WITH ( ) or [ ]   -> \text{\en{...}}       (\en forces LTR; a plain
%                              \text mirrors the parens in RTL -> "Spec (TNR(" )
%   f(Hebrew/text arg)      -> \text{f}(\text{...})  -- keep the parens at MATH
%                              level, NOT \text{f(arg)} (parens inside the box mirror)
%   code in math            -> \mathtt{...}  (or \texttt{}, math-guarded below)
% NEVER write \en{...} or \h{...} straight into $...$ / \[..\] (outside \text);
% always via \text{}. (\mbox{...} also works but does not scale in scripts.)
% ============================================================================
\newcommand{\h}[1]{\foreignlanguage{hebrew}{#1}}
% v3.42.1 — \en islands do NOT hyphenate (auto OR at explicit hyphens). English hyphenation patterns were
% tearing identifiers and terms ACROSS THE RTL SEAM ("IO Regis-/ters",
% "In-/put", "Band-/width", "Per-/formance", "Multi-/ple" — all shipped on
% real pages; probe-proven: "ALUS-/rcB", "Mem-/ToReg"). An island is a label
% or a term, not English prose — a hyphen tear there is never right. A/B on
% a full 197pp lakit: identical overfull sets, same page count, ~11
% paragraphs healed, one mild underfull (badness 3525) that replaced a
% mutilated "Regis-/ters". Long English PROSE belongs in the englishprose
% environment below, which restores explicit-hyphen breaks while its paragraph
% is still active. A raw otherlanguage block inherits the global prohibition.
% (\l@nohyphenation needs \csname:
% '@' is not a letter here.)
% v3.43.2 — \exhyphenpenalty moved OUT of \en to a GLOBAL assignment below.
% Render-proven: a group-local \exhyphenpenalty is DEAD CODE — TeX reads the
% line-breaking parameters at \par time, after the \en group has closed, so
% \en{write-through} / plain not-taken still tore at the "-" ("write-/
% through" shipped on a real exam page). \language is different: it is
% recorded in the node list as text is added, so the in-group nohyphenation
% assignment above DOES work (and stays). The global penalty was A/B-measured
% on a 197pp lakit: identical overfull set (25), identical page count, the
% not-/taken tear healed. Hebrew is unaffected in practice — Hebrew compounds
% use maqaf (־), not ASCII "-".
\newcommand{\en}[1]{\foreignlanguage{english}{\language=\csname l@nohyphenation\endcsname #1}}
% Babel runs extras locally for inline islands and paragraph language changes.
% An English main document must not inherit the Hebrew no-hyphen paragraph rule.
\addto\extrashebrew{\exhyphenpenalty=10000\relax}
\addto\extrasenglish{\exhyphenpenalty=50\relax}
\ifx\lakitmainlanguage\hllpdfenglish
  \exhyphenpenalty=50\relax
\else
  \exhyphenpenalty=10000\relax
\fi
% Full English paragraphs need ordinary explicit-hyphen breaks. Use this
% environment instead of a raw otherlanguage block; the paragraph ends while
% the local penalty is still active.
\newenvironment{englishprose}
  {\par\begin{otherlanguage}{english}\exhyphenpenalty=50\relax}
  {\par\end{otherlanguage}}
\newenvironment{hebrewprose}
  {\par\begin{otherlanguage}{hebrew}\exhyphenpenalty=10000\relax}
  {\par\end{otherlanguage}}
% Use for shell text whose locale must not depend on the last body paragraph.
\DeclareRobustCommand{\doclang}[1]{\foreignlanguage{\lakitmainlanguage}{#1}}

% ============================================================================
% ★★ — fixes for defects that are INVISIBLE in the log and in the checks
% and only surface on a VISUAL RENDER of the PDF. (Each verified empirically.)
% ----------------------------------------------------------------------------
% v3.22: \emph is now REAL slanted emphasis (FakeSlant italic face defined on
% the \babelfont{rm} line above), replacing the v3.21 \emph->\textbf hack: it
% preserves \emph's toggle semantics and keeps bold free for its own role.
% Prefer bold-as-emphasis for your doc? one line: \renewcommand{\emph}[1]{\textbf{#1}}
% page number: force the digits LTR so they never reverse (18 -> 81) in the RTL
%   footer. That reversal is a NON-DETERMINISTIC bidi glitch (observed once in
%   ~80 pp); this LuaLaTeX build is not byte-reproducible, so the defensive
%   LTR wrap is the only reliable guard. Latin digits are always LTR anyway.
\renewcommand{\thepage}{\foreignlanguage{english}{\arabic{page}}}
% v3.70: \pagenumbering \gdef's \thepage and silently DESTROYS the wrap above —
% after \pagenumbering{arabic} the folio guard is gone for the whole rest of the
% document. Two collections discovered this independently and hand-restore the
% wrap in main.tex with the same comment; that knowledge now lives here instead.
% The wrapper re-applies the LTR wrap after the kernel call, for any style.
\makeatletter
\let\lakin@orig@pagenumbering\pagenumbering
\renewcommand{\pagenumbering}[1]{%
  \lakin@orig@pagenumbering{#1}%
  \renewcommand{\thepage}{\foreignlanguage{english}{\csname @#1\endcsname\c@page}}}
\makeatother
% emergencystretch: Hebrew does not hyphenate, so an unbreakable LTR island
%   (\en{...} / an inline math run) landing at the margin can leave a small
%   (~2-13pt) overfull with no legal breakpoint. This rescues ONLY those
%   inline-prose tie-breaks; it is a ceiling used only on offending lines and
%   does NOT touch display-math / listing / table overflow (those still surface
%   in the log for the content-level overflow pass). Verified: baseline overfull
%   -> 0 with this set, with no change to pagination or underfull.
\setlength{\emergencystretch}{3em}
% v3.40: inline math (almost) never breaks across lines. TeX's default
%   (binop 700 / rel 500) lets an inline run split at an operator —
%   $(1.f)\cdot 2^{e}$ shipped as "(1.f)·" ending one line and "2^e" opening
%   the next (reader-reported). In LTR that split is merely ugly; under an
%   RTL base the carried operator lands on the visually WRONG side of the
%   break and the formula reads backwards.
%   WHY 10000 (v3.45.0 — the display-first resolution). Inline math is
%   FORBIDDEN from breaking, so a torn island under an RTL base is
%   structurally impossible. The v3.40 cycle used 9999 ("break only as a last
%   resort") because a bare 10000 left ~30-39 WIDE inline equations
%   overflowing (worst 90.7pt) — but 9999's "last resort" fires OFTEN in
%   tight, non-hyphenating Hebrew: a reader found torn islands across the
%   fleet (~150 candidates; e.g. "P(y |" ending one line and "x_0)" opening
%   the next, read backwards). The step the 9999 cycle missed: every one of
%   those ~30 overflows is a wide equation that BELONGS in display. Convert
%   them to display FIRST, then 10000 gives 0 tears AND 0 overflows (measured,
%   12/12 fleet). Doctrine: inline math is ATOMIC; a wide equation goes in a
%   display, never inline. See content-style ("inline math is atomic").
\binoppenalty=10000 \relpenalty=10000
% v3.38: uniform sentence spacing. TeX's default (\nonfrenchspacing, sfcode
% 3000) is the OLD Anglo-American wide-space-after-period convention — Hebrew
% typography never had it, and modern Latin practice (Chicago et al.) dropped
% it too. Measured on real pages: post-period gaps up to 1.9-2.1x the line's
% mean word gap — the dominant "rivers" source in justified Hebrew prose.
% Survives babel language switches (probed: sfcode stays french inside \en{}
% islands and otherlanguage blocks).
\frenchspacing
% v3.38: widow/orphan protection. LaTeX's default penalties (150/150/50)
% allow a lone line at a page top/bottom almost freely. 10000 = forbid;
% this is SAFE here precisely because the document is \raggedbottom (the
% demoted line just moves on, no vertical stretch fight) — under flushbottom
% these values would cause underfull pages. raggedbottom is the report-class
% default; made explicit so a future class/option change cannot silently
% flip it and turn the 10000s harmful.
\raggedbottom
\widowpenalty=10000 \clubpenalty=10000 \displaywidowpenalty=10000
% em-dash: --- does NOT become — under Frank Ruhl (the font lacks tlig; verified
%   the width is identical with/without Ligatures=TeX, while the mechanism works
%   on Latin Modern). TYPE THE CHARACTER — (U+2014) DIRECTLY; likewise – (U+2013).
%   check_bidi_figures.py flags a --/--- that is adjacent to Hebrew.
% ============================================================================

% ---- STRUCTURAL DIRECTION-SAFETY for \texttt (v3.7) -------------------------
% A parenthesis, bracket, or $ INSIDE \texttt{} mirrors/migrates in RTL prose
% (MIPS 0($s0)->"(s0$)0", $rs->"rs$", arr[i]->"[arr[i") because \texttt sets the
% font but NOT the direction. Redefining \texttt to LTR-isolate its content fixes
% the WHOLE class at once (parens, brackets, $) so plain \texttt{lw $t0,0($s0)}
% is safe — no need to remember \code{}. Guarded by \ifmmode because
% \foreignlanguage is invalid in math (there, plain monospace is already LTR).
% NB (corrected v3.56 — the old note here said "stays tofu", which is WRONG):
% \texttt{Hebrew} is a FATAL build failure. The tt family is DejaVu Sans Mono
% (zero Hebrew) and the tt fallback hebfb is DejaVuSerif (also zero Hebrew), so
% the subsetter gets an empty glyph set and LuaTeX dies with
%   "! error: (file .../DejaVuSansMono.ttf) (type 2): there are no glyphs in the
%    subset"  -> no output PDF.
% The message names a FONT FILE, never Hebrew/\texttt/the source line. And it is
% non-local: add an unrelated Latin \code{} elsewhere and the fatal degrades into
% a SILENT DELETION of the Hebrew word. Put Hebrew in \h{}/normal text, never in
% \texttt{}/\code{}; flagged HEBREW-IN-MONO. To disable, delete this block.
\let\origtexttt\texttt
\protected\long\def\texttt#1{\ifmmode\origtexttt{#1}\else\foreignlanguage{english}{\origtexttt{#1}}\fi}
% v4.8.0: defer the direction test until TeX has inserted an alignment cell's
% math template. Unprotected expansion here tested \ifmmode too early at the
% FIRST token of aligned/array/cases/matrix cells. Engine-level \protected also
% preserves the original \texttt implementation without the self-recursion
% observed when retrofitting \DeclareRobustCommand to its saved alias.
\protected\long\def\code#1{\ifmmode\origtexttt{#1}\else\foreignlanguage{english}{\origtexttt{#1}}\fi}

% ---- \codeline — display-code line (v3.24): the code analogue of display math.
% Reader-friction rule (content-style.md §5b): inline \code{} is for a SINGLE
% ATOM only (one register/variable/mnemonic, a tight no-space expression).
% A multi-token expression — an instruction with operands, a spaced assignment,
% a register-binding list — zigzags the eye RTL<->LTR mid-sentence and, being
% an unbreakable island, line-breaks badly. Give it its own centered line with
% display-math vertical rhythm. \centerline = one box => the expression never
% wraps; \nopagebreak keeps it glued to its introducing sentence.
%   ...the Hebrew sentence ends with a colon:
%   \codeline{lw \$t1, -4(\$t0)}
%   and the prose continues (escape dollars as \$).
% v3.66 (first line of the body): flush a PENDING \item label before the \par.
% LaTeX defers the label box (\@labels) until the next \everypar, and an \hbox in
% vertical mode never fires \everypar — so `\item \codeline{...}` printed the code
% at y=69.1 and the item's own "1." BELOW it at y=103.6 (the number appears after
% its content, attached to nothing). \if@inlabel\leavevmode\fi is CONDITIONAL on
% purpose: an unconditional \leavevmode would open an empty paragraph line when
% \codeline follows ordinary prose. Verified both ways — in a list the label now
% sits on the code's line at its normal x, and after prose the geometry is
% byte-identical to before.
\makeatletter
\newcommand{\codeline}[1]{%
  \if@inlabel\leavevmode\fi
  \par\nopagebreak
  \vspace{\abovedisplayskip}\nopagebreak
  % v3.62: was \centerline{...}, i.e. \hbox to\hsize{\hss #1 \hss}. \hss has
  % INFINITE SHRINK, so TeX could always "fit" the box and NEVER reported an
  % Overfull — the same \hss mechanism as BLD-27's \item label. Measured: a wide
  % \codeline in a defbox painted its first token ON the box's left border bar and
  % sat 9.5pt outside the frame interior on both sides, with 0 Overfull, 0 Underfull
  % and all four gates PASS; in bare prose it reached 21.0pt into BOTH margins.
  % \hfil has no shrink, so the overflow is reported and build.sh's ovf_vis/ovf_big
  % counters see it; \linewidth (not \hsize) makes it box-aware, so a \codeline
  % inside a colored box centres on the box interior. Verified: a FITTING \codeline
  % keeps byte-identical glyph positions, and the box-crossing case now emits
  % "Overfull \hbox (19.03pt too wide)". A wide \codeline should be split in two.
  \hbox to\linewidth{\hfil\code{#1}\hfil}%
  \nopagebreak\vspace{\belowdisplayskip}}
\makeatother

% ---- a thin decorative rule helper ----
\newcommand{\sep}{\par\vspace{2pt}\noindent\textcolor{cRule}{\rule{\linewidth}{0.6pt}}\par\vspace{4pt}}

% ---- colors ----
\definecolor{cBlue}{HTML}{1F4E79}   \definecolor{cBlueBg}{HTML}{EAF1F8}
\definecolor{cGreen}{HTML}{2E6B4F}  \definecolor{cGreenBg}{HTML}{EAF4EE}
\definecolor{cAmber}{HTML}{9C6A00}  \definecolor{cAmberBg}{HTML}{FBF3E2}
\definecolor{cPurple}{HTML}{5B2C6F} \definecolor{cPurpleBg}{HTML}{F3ECF7}
\definecolor{cTeal}{HTML}{0E6E6E}   \definecolor{cTealBg}{HTML}{E6F4F4}
\definecolor{cRed}{HTML}{8E2B2B}    \definecolor{cRedBg}{HTML}{FBEAEA}
\definecolor{cGray}{HTML}{555555}   \definecolor{cRule}{HTML}{1F4E79}

% ---- math operators ----
\DeclareMathOperator{\Tr}{Tr}
\DeclareMathOperator{\rank}{rank}
\DeclareMathOperator{\diag}{diag}
\DeclareMathOperator{\sgn}{sgn}

% ---- bra-ket (quantikz defines \ket/\bra/\braket; add the rest safely) ----
\providecommand{\ket}[1]{\left|#1\right\rangle}
\providecommand{\bra}[1]{\left\langle#1\right|}
\providecommand{\braket}[2]{\left\langle#1\middle|#2\right\rangle}
\newcommand{\dyad}[1]{\left|#1\right\rangle\!\left\langle#1\right|}
% v3.21: general outer product |a><b| and expectation value <A> — agents reach
% for these by reflex (esp. \expval); quantikz defines neither, so their absence
% threw "Undefined control sequence" in a real build. \providecommand = safe.
\providecommand{\ketbra}[2]{\left|#1\right\rangle\!\left\langle#2\right|}
\providecommand{\expval}[1]{\left\langle#1\right\rangle}
\newcommand{\abs}[1]{\left|#1\right|}
\newcommand{\norm}[1]{\left\lVert#1\right\rVert}
\newcommand{\half}{\tfrac12}
\newcommand{\R}{\mathbb{R}}\newcommand{\C}{\mathbb{C}}\newcommand{\Z}{\mathbb{Z}}
% identity operator. NB: \mathbb{1} silently renders a WRONG glyph (msbm has no
% blackboard digits). Use \mathbb{I} = 𝕀 (works with amssymb, no extra package).
% Do NOT "fix" this back to \mathbb{1}.
\newcommand{\id}{\mathbb{I}}

% ---- physical units (v3.35): \punit — upright unit with thin space ---------
% Usage: $300\punit{K}$, $80\punit{mV/decade}$, $10^{7}\punit{cm/s}$ — the
% NUMBER AND UNIT LIVE IN ONE MATH ISLAND. Promoted verbatim from a 199-page
% production lakit (advanced-electronic-devices, dozens of uses).
% Two measured traps this design avoids (both render-proven):
%   1. siunitx SCRAMBLES under RTL — \SI{3.5}{eV} renders "eV 3.5", exponents
%      detach ("Hz 10 x 10^9"): its quantity assembly is not bidi-isolated.
%      Do NOT load siunitx in this pipeline.
%   2. A text-mode unit next to a bare-prose number is TWO adjacent LTR runs
%      and reorders ("5 mm" -> "mm5" — the BIDI-2/3 class). \punit is therefore
%      math-only ON PURPOSE: prose misuse fails LOUDLY at compile (\mathrm
%      outside math) instead of scrambling silently.
% Domain shorthands (\MeV, \fm...) stay per-document: \newcommand{\MeV}{\,\text{\en{MeV}}}
\providecommand{\punit}[1]{\,\mathrm{#1}}

% ============================================================================
% Colored boxes — mdframed (NOT tcolorbox, which breaks under bidi).
% ★★ v3.4: innerleftmargin/innerrightmargin give code listings room so their
% frame/line-numbers do not spill out the box border (see listing styles below).
% ============================================================================
% v3.61: everyline/split skips were applied to `tablebox` only (BLD-14, v3.37) —
% the six colored boxes below never got them, so they still split OPEN across a page
% break: the page-N segment gets top+sides and NO bottom rule, the page-N+1 segment
% sides+bottom and NO top rule. Measured on the shipped fleet this is a COVERAGE GAP,
% not a re-report of BLD-14: 0 split-open `tablebox` segments (that fix holds) vs 62
% split-open COLORED boxes across 7 documents. Verified after the fix: both segments
% close, and the frametitle is NOT duplicated on the continuation. BLD-14's caveat
% still applies — a box taller than a full page grinds mdframed's split-trial loop, so
% split such content manually.
% ============================================================================
% v3.72: frametitle={} leads EVERY \newmdenv option list below. mdframed options
% are GROUP-scoped: a box nested inside \begin{exbox}[frametitle={...}] used to
% INHERIT the outer title — measured, the identical title string rendered twice,
% the second riding a title-bar fill inside the INNER box's frame, for all seven
% environments in both directions. Both natural forms were broken (= preamble
% bug, not an authoring rule). The explicit reset makes every box titleless
% unless ITS OWN [frametitle={...}] says otherwise; verified both directions:
% no duplicate on the nested box, an explicit title still renders.
% Geometry is shared with the optional semantic-width table container. Themes
% select their own existing semantic style below; neither path fixes columns.
\mdfdefinestyle{hllpdfsemanticgeometry}{linewidth=1.1pt,
  innertopmargin=6pt,innerbottommargin=7pt,skipabove=9pt,skipbelow=9pt,
  innerleftmargin=9pt,innerrightmargin=9pt,everyline=true,
  splitbottomskip=6pt,splittopskip=12pt,roundcorner=3pt}
\newcommand{\hllpdfsemanticframestyle}{hllpdfsemanticgeometry}
\newmdenv[frametitle={},style=hllpdfsemanticgeometry,linecolor=cBlue,backgroundcolor=cBlueBg,
  frametitlebackgroundcolor=cBlue,frametitlefontcolor=white,
  ]{defbox}   % הגדרה
\newmdenv[frametitle={},style=hllpdfsemanticgeometry,linecolor=cGreen,backgroundcolor=cGreenBg,
  frametitlebackgroundcolor=cGreen,frametitlefontcolor=white,
  ]{thmbox}   % משפט/טענה
\newmdenv[frametitle={},style=hllpdfsemanticgeometry,linecolor=cAmber,backgroundcolor=cAmberBg,
  frametitlebackgroundcolor=cAmber,frametitlefontcolor=white,
  ]{notebox}  % הערה
\newmdenv[frametitle={},style=hllpdfsemanticgeometry,linecolor=cPurple,backgroundcolor=cPurpleBg,
  frametitlebackgroundcolor=cPurple,frametitlefontcolor=white,
  ]{exbox}    % דוגמה
\newmdenv[frametitle={},style=hllpdfsemanticgeometry,linecolor=cRed,backgroundcolor=cRedBg,
  frametitlebackgroundcolor=cRed,frametitlefontcolor=white,
  ]{warnbox}  % טעות נפוצה
\newmdenv[frametitle={},style=hllpdfsemanticgeometry,linecolor=cTeal,backgroundcolor=cTealBg,
  frametitlebackgroundcolor=cTeal,frametitlefontcolor=white,
  ]{keybox}   % רעיון מרכזי

% ---- codebox (v4.5.0): the dedicated container for code -------------------
% An OPTION, not a policy: asmblock unboxed stays canonical for a full listing,
% and a snippet inside any of the six colored boxes still works as before.
% What makes codebox worth its own environment is that it is the ONE framed box
% that keeps its margin line numbers.
%
% v4.4.0 shipped it WITHOUT them, on the strength of v3.37's "no margin knob can
% save them" plus four fresh probes that agreed. That was wrong: every probe had
% moved the code, the box, or both — never the key that carries the NUMBER
% COLUMN. Measured on the page:
%   resetmargins=false   the listing follows THIS box instead of the page,
%                        moving the code from x=75.26 to x=97.54
%   xrightmargin=26pt    under RTL this is the PHYSICAL LEFT indent, and the
%                        only key that carries the \llap'ed digits with it:
%                        x=64.28 -> 90.18, clearing the rule at 65.64 by 23.6pt
%   xleftmargin=26pt     moves the code and LEAVES THE DIGITS on the rule —
%                        the asymmetry that hid the fix for eleven versions
% resetmargins is hoisted out of the four styles (see the \lstset below them) so
% this override is possible at all; every other listing still resolves to true.
% Pinned by tests/codebox_numbers.tex + probe_lstnum.py --inside-rule, whose
% negative control strips xrightmargin and watches the digits fall back.
\definecolor{cSlate}{HTML}{3F4A56}  \definecolor{cSlateBg}{HTML}{F2F4F7}
\newmdenv[frametitle={},leftline=true,rightline=true,topline=true,bottomline=true,
  linecolor=cSlate,backgroundcolor=cSlateBg,linewidth=0.9pt,
  frametitlebackgroundcolor=cSlate,frametitlefontcolor=white,
  innertopmargin=6pt,innerbottommargin=7pt,skipabove=9pt,skipbelow=9pt,
  innerleftmargin=9pt,innerrightmargin=9pt,everyline=true,splitbottomskip=6pt,splittopskip=12pt,
  roundcorner=3pt]{codebox}   % קוד
% The two keys that make in-box numbering work, both measured on the page:
%   resetmargins=false  the listing follows THIS box instead of the page
%   xrightmargin        under RTL this is the PHYSICAL LEFT indent, and it is
%                       the only one that carries the \llap'ed number column
%                       with it. xleftmargin moves the code and leaves the
%                       digits behind on the rule (measured: code 97.54,
%                       digits 64.28 — the v3.37 defect, reproduced); with
%                       xrightmargin=26pt the digits land at x=90.18, a
%                       clearance of 23.6pt from the rule at 65.64.
% This asymmetry is why v3.4's note had the semantics backwards and v3.37
% concluded the digits could not be saved.
% <hllpdf:listing-container-state>
% v4.29: all state is local to the current environment group. A semantic
% callout masks number text, while a nested codebox restores the formatter
% saved before that mask. Save it once across successive semantic wrappers;
% an explicit numbers=none (or an explicitly empty formatter) remains valid.
\makeatletter
\newif\ifhllpdf@incodebox
\newif\ifhllpdf@englishcodebox
\newif\ifhllpdf@lstnumbermasked
\let\hllpdf@savedlstnumber\thelstnumber
\makeatother
% </hllpdf:listing-container-state>
% <hllpdf:codebox-gutter>
\makeatletter
\AtBeginEnvironment{codebox}{%
  \hllpdf@incodeboxtrue
  % Snapshot the OUTER language before asmblock selects English.
  \iflanguage{english}{\hllpdf@englishcodeboxtrue}{\hllpdf@englishcodeboxfalse}%
  \ifhllpdf@lstnumbermasked
    \let\thelstnumber\hllpdf@savedlstnumber
    \hllpdf@lstnumbermaskedfalse
  \fi
}
% listings' AddToHook is GLOBAL. Register once, and gate on local box state;
% adding an unconditional hook at every codebox entry leaks into later listings.
\lst@AddToHook{PreInit}{%
  \ifhllpdf@incodebox
    % Apply every listing override here, inside the LISTING's own group.
    % Setting margins at box entry leaks them into nested semantic callouts.
    \lstset{resetmargins=false,xrightmargin=26pt}%
    % Styles resolve before PreInit. In LTR the physical number gutter follows
    % xleftmargin; in Hebrew retain the measured inherited xrightmargin path.
    \ifhllpdf@englishcodebox
      \lstset{xleftmargin=26pt,xrightmargin=0pt}%
    \fi
    % ONE box, not two: suppress only the frame/background/skips of a listing
    % inside this codebox. The next unboxed listing keeps its own style.
    \lstset{frame=none,backgroundcolor=\color{cSlateBg},aboveskip=0pt,belowskip=0pt}%
  \fi
}
\makeatother
% </hllpdf:codebox-gutter>

% v3.37: line numbers inside a COLORED BOX are sliced by the box border —
% measured: digit span x=63.5-67.3 vs border line at x=65.2 (the border runs
% THROUGH the digits), and the inter-border gap (3.5pt) is narrower than a
% digit (3.8pt), so no margin knob can save them (under resetmargins,
% xleftmargin is inert in-box — measured: code at x=75.26 with 2.2em AND
% 3.0em). A 3-6 line boxed snippet does not need line numbers; blank the
% number TEXT inside every colored box (\thelstnumber is what gets printed;
% the style's numbers=left survives, so nothing else shifts — code position
% byte-frozen at x=75.26, verified). Unboxed listings keep their numbers.
% v3.78: the blanking is coupled to the FRAME, not to the box name —
% clean_theme.tex registers LATER hooks (registration order wins) that
% restore the number text inside the boxes it renders frameless.
% v4.5.0 RETRACTS the "no margin knob can save them" clause above. It holds only
% under resetmargins=true (the shipped default): there xleftmargin is inert
% in-box, as measured. With resetmargins=false, xrightmargin DOES carry the
% \llap'ed digit column clear of the rule — measured; the codebox \lstset that
% does it is the one above, in the codebox block. The loop below therefore
% covers SIX boxes, not "every colored box": codebox opts out and keeps its
% numbers.
% <hllpdf:semantic-listing-mask>
\makeatletter
\@for\HEB@lstbox:=defbox,thmbox,notebox,exbox,warnbox,keybox\do{%
  \expandafter\AtBeginEnvironment\expandafter{\HEB@lstbox}{%
    \ifhllpdf@lstnumbermasked\else
      \let\hllpdf@savedlstnumber\thelstnumber
    \fi
    \hllpdf@lstnumbermaskedtrue
    \def\thelstnumber{}%
    % The nearest semantic container owns its listing presentation. On exit,
    % TeX's group restores the enclosing codebox state and formatter.
    \hllpdf@incodeboxfalse}}
\makeatother
% </hllpdf:semantic-listing-mask>

% ============================================================================
% Unified TABLE presentation: subtle uniform slate frame + auto-numbered
% caption rendered INSIDE the frame. Use:
%   \begin{tablebox}\centering  <tabular>  \tabcap{<description>}\end{tablebox}
% Replaces the buggy \caption* (which emitted a stray "טבלה N : *" line) and
% provides an explicit framed option + numbered caption "טבלה N.M: תיאור".
% Optional intro text may be placed at the TOP, before the tabular, inside the box.
% ============================================================================
\definecolor{cTblLine}{HTML}{B9C4D2}  \definecolor{cTblBg}{HTML}{FAFBFC}
% v3.23: class-aware — under `article` there is no chapter counter (the promised
% report->article switch used to die with "No counter 'chapter' defined").
\makeatletter
\@ifundefined{c@chapter}{%
  % v3.31: in article the counter resets per SECTION, so an unprefixed number
  % repeated ("טבלה 1." in every section — measured on a two-section render).
  % Same conditional idea as the report branch: prefix the section number once
  % a real \section exists; a pre-section table stays "טבלה 1".
  \newcounter{tbl}[section]
  \renewcommand{\thetbl}{\ifnum\value{section}>0 \thesection.\fi\arabic{tbl}}%
}{%
  \newcounter{tbl}[chapter]
  % v3.30: a report-class doc that uses SECTIONS ONLY (chapter stays 0 —
  % the common short-doc pattern) used to print "טבלה 0.1"; the chapter
  % prefix now appears only once a real \chapter exists (render-verified
  % both ways: "טבלה 1." chapterless, "טבלה 1.1" with chapters).
  \renewcommand{\thetbl}{\ifnum\value{chapter}>0 \thechapter.\fi\arabic{tbl}}%
  % v3.64: the SAME leak in \thesection itself, left behind by v3.30. That fix
  % guarded the CAPTION counter and stopped there, so a chapterless report still
  % printed the heading "0.1 סעיף ראשון" and \ref to it printed "0.2" — measured:
  %   p1 y=65.2  '…0.1'   p1 y=144.9 '…0.2'   p1 y=183.4 '…0.2.1'
  % report.cls hardcodes \thesection as \thechapter.\@arabic\c@section, and the
  % leading 0 is the chapter counter that no \chapter ever incremented. Fixing
  % \thesection fixes the whole chain — \thesubsection and \thesubsubsection are
  % defined in terms of it. \theequation and \thetbl/\thefig were already guarded,
  % which is exactly why this survived: every NUMBER around the heading looked
  % right, so the heading's own 0 read as a style choice rather than a bug.
  % \appendix is safe: report.cls resets chapter to 0 and redefines \thechapter to
  % \@Alph\c@chapter, so an appendix \chapter makes the counter 1 and the guard
  % emits "A.1" unchanged. Render-verified in all four modes (chapterless,
  % chaptered, appendix+chapter, appendix+section-only).
  \renewcommand{\thesection}{\ifnum\value{chapter}>0 \thechapter.\fi\arabic{section}}%
}
\makeatother
% v3.37: everyline=true — if the box DOES split across pages, every segment
% gets a CLOSED frame (measured before: the split ran open-ended into the page
% edge with no bottom border, and the caption could strand alone at the next
% page top — "טבלה A.2" shipped that way). The split skips give the segments
% breathing room. In practice, with the caption GLUE below, every previously-
% split table in the corpus now moves WHOLE instead (the white-gap tradeoff,
% documented benign); everyline stays as the safety net for a genuinely
% page-tall table. NB: a table taller than ONE FULL PAGE can grind mdframed's
% split-trial loop (reproduced: minimal doc, minutes of 1pt-step Overfull
% probes) — break such a monster into two tables.
\mdfdefinestyle{hllpdftablegeometry}{leftline=true,rightline=true,topline=true,bottomline=true,
  linewidth=0.9pt,
  innertopmargin=9pt,innerbottommargin=8pt,skipabove=11pt,skipbelow=11pt,
  innerleftmargin=11pt,innerrightmargin=11pt,roundcorner=4pt}
\newcommand{\hllpdftableframestyle}{hllpdftablegeometry}
\newmdenv[frametitle={},style=hllpdftablegeometry,
  linecolor=cTblLine,backgroundcolor=cTblBg,
  % v4.2.0 (adopted): tablebox is ATOMIC. Short tableboxes near page bottoms
  % repeatedly entered mdframed's split trial, emitting benign Overfull-vbox
  % bursts and occasionally separating the intro prose from the tabular.
  % nobreak=true moves the complete short table to the next page instead. The
  % output-active Overfull-vbox gate still catches a box taller than a page;
  % everyline/split skips remain as defensive styling should that path ever
  % re-arm, but the canonical path never publishes a split tablebox.
  nobreak=true,
  everyline=true,splitbottomskip=6pt,splittopskip=12pt,
  % v3.79 (BLD-37 table side, correction): startinnercode/endinnercode run
  % INSIDE mdframed's body vbox (mdframed.sty:848 / :856+863) -- the second one
  % after \par\unskip, i.e. after line breaking, which is the only moment the
  % \tabcap anchor's true offset from the top of the box EXISTS. See
  % \lakintabboxclose below. These two keys are the ONLY tablebox change.
  startinnercode={\lakintabboxopen},endinnercode={\lakintabboxclose}]{tablebox}
% v3.37: \nopagebreak GLUE — \vsplit honors penalties, so the caption can no
% longer strand alone at the top of the next page (it takes its rows along, or
% the whole box moves). Render-verified on the real shipped orphan.
% v3.59: \NR@gettitle is what nameref injects into the stdlib title-bearing
% commands (\@sect, \@chapter, \caption) to fill \@currentlabelname. \figcap and
% \tabcap are INDEPENDENT macros — they never called it, so \@currentlabelname kept
% whatever the last heading left and \nameref{fig:x} silently printed the enclosing
% SECTION's title instead of the caption (proven in the .aux: field 3 of \newlabel
% was the heading text while field 1 and the anchor `fig.1.1` were correct — so this
% is a THIRD, distinct head next to BLD-13's counter fix and BIDI-38's autoref-name
% fix, neither of which writes that field). \@ifundefined keeps it inert if
% hyperref/nameref is ever dropped. Guarded helper because \figcap is defined
% after \makeatother.
% NB both helpers below need \makeatletter — this part of the preamble is OUTSIDE
% any letter-catcode region, so an unwrapped \@currentlabel parses as LaTeX's \@
% and dies with "You can't use \spacefactor in vertical mode" (measured).
\makeatletter
\newcommand{\lakinrtitle}[1]{\@ifundefined{NR@gettitle}{}{\NR@gettitle{#1}}}
% v3.59: \hitem{<letter>}{<key>} — a REFERENCEABLE Hebrew-lettered list item.
%   \begin{enumerate} \hitem{א}{it:a} הסעיף הראשון. \end{enumerate}  ->  \ref{it:a} = (א)
% Why this exists: automatic Hebrew list labels fail silently, so SKILL.md mandates
% an EXPLICIT \item[(\textbf{א})] — but LaTeX puts \refstepcounter INSIDE
% \if@noitemarg (latex.ltx:12964), and \item[...] never sets that flag. So an
% explicit item steps no counter and never touches \@currentlabel: a \label there
% binds to the AMBIENT one and \ref prints the CHAPTER/SECTION number where the
% item letter belongs (proven in the .aux: {chapter.1}). Same doctrine as \figcap —
% own the label state before \label. \phantomsection matters: without it the
% hyperlink still lands on the chapter head (measured); with it the anchor is
% section*.N and the bookmark outline stays clean.
\newcommand{\hitem}[2]{\item[(\textbf{#1})]%
  \def\@currentlabel{(#1)}\def\@currentlabelname{#1}\phantomsection\label{#2}}
% NB (v3.66): \hitem is \ref-only. \autoref{key} calls it a SECTION ("סעיף (א)")
% because \phantomsection brands the anchor section*.N, and \nameref{key} prints
% the bare letter. For \autoref-able items prefer \hebenumlabels below.
%
% v3.66: \hebenumlabels — OPT-IN automatic Hebrew enumerate letters, done right.
%   {\hebenumlabels\begin{enumerate}\item\label{q:a} ...\end{enumerate}}
%   -> labels (א) (ב) (ג); \ref{q:a} = א; anchor Item.N (real, per-item).
% This retires the belief that made \item[(\textbf{א})] + \hitem necessary.
% "Automatic Hebrew labels fail silently" was TRUE but unexplained, and the
% mechanism is now bisected (plain article vs +enumitem, measured): ENUMITEM
% PRE-EXPANDS the label expression in an \edef BEFORE \refstepcounter, with the
% counter still at 0. Tokens it shields by name (\arabic,\alph,\roman) survive;
% a plain \newcommand \ifcase wrapper evaluates at 0 -> the EMPTY zero-branch ->
% the label prints as "()" with the letter GONE (and \the\c@enumi prints a
% literal 0, \number\value dies fatally pointing at \item). The cure is one
% property, not a doctrine: the wrapper must be ROBUST (survives the \edef
% unexpanded), while \theenumi stays EXPANDABLE (so \refstepcounter writes the
% real letter to the .aux). Verified: labels (א)(ב)(ג), \ref -> א/ב/ג each item
% its own value, .aux = {א}{...}{Item.1} — versus the explicit-\item doctrine
% where BIDI-50/\hitem exist purely to patch the counter that never stepped.
\def\@hebalph#1{\ifcase#1\or א\or ב\or ג\or ד\or ה\or ו\or ז\or ח\or ט\or י\or
  יא\or יב\or יג\or יד\or טו\or טז\or יז\or יח\or יט\or כ\else\@ctrerr\fi}
\DeclareRobustCommand*\hebalph[1]{\expandafter\@hebalph\csname c@#1\endcsname}
\newcommand{\hebenumlabels}{%
  \renewcommand{\theenumi}{\@hebalph\c@enumi}%
  \renewcommand{\labelenumi}{(\textbf{\hebalph{enumi}})}}
\makeatother
% v3.63: \mbox the LABEL. \thetbl ("7.1") and a caption description that STARTS
% WITH A DIGIT ("32 הרגיסטרים…") are two European-Number runs separated only by
% neutrals (the label's own "." and \enspace). Under the RTL paragraph base the UBA
% resolves those neutrals to the RTL level, so both numbers sit in ONE reversed span
% and EXCHANGE VISUAL POSITIONS. Measured on the shipped PDF: a table whose real
% number is 7.1 printed "טבלה 32" — a different, EXISTING table number — while \ref
% correctly printed 7.1, so the caption silently contradicted its own cross-reference,
% and the description lost its leading count. The source contains no adjacency at all;
% it is manufactured by this macro, which is why every source linter was blind.
% \mbox makes the label one atomic box so nothing can reorder across it. Verified
% zero-regression: an ordinary Hebrew caption keeps identical x-geometry.
% One paragraph style covers label, number, separator and description.
% Explicit project/template overrides remain possible; the label hook is kept
% for compatibility, with an identity default. Counters/anchors stay unchanged.
\newcommand{\lakitcaptionstyle}{\normalfont\small\color{black}\hypersetup{allcolors=black}}
\newcommand{\lakitcaptionlabelstyle}[1]{#1}
\newcommand{\lakitcaptionpunctuation}{:}
% Native floats, captionof and longtable use the same presentation policy.
% subcaption is optional and inherits the same font rather than a second size.
\usepackage{caption}
\DeclareCaptionFont{lakit}{\lakitcaptionstyle}
\DeclareCaptionLabelFormat{lakit}{\lakitcaptionlabelstyle{#1~\en{#2}}}
\DeclareCaptionLabelSeparator{lakit}{\lakitcaptionpunctuation\enspace}
\captionsetup{font=lakit,labelfont={},textfont={},labelformat=lakit,
  labelsep=lakit,justification=centering}
\makeatletter
\AtBeginDocument{\@ifpackageloaded{subcaption}{%
  \captionsetup[sub]{font=lakit,labelfont={},textfont={},labelsep=lakit}}{}}
\makeatother
% v3.64: OPT-IN list-of-figures/tables entry via an optional short caption.
% \listoffigures and \listoftables are inherited from report.cls and LOOK available,
% but \figcap/\tabcap step the custom \c@fig/\c@tbl counters and never call
% \addcontentsline — so both printed their heading and NOTHING ELSE. Measured: a doc
% with one \figcap and one \tabcap rendered page 1 as 'רשימת האיורים\n1\n', and the
% .lof held only \@chapter's \addvspace, zero \contentsline entries. An empty list
% under a correct Hebrew heading reads as "this document has no figures".
%
% WHY OPTIONAL, AND NOT AUTOMATIC. Writing #1 unconditionally was implemented first
% and REJECTED: a caption containing \codeline compiles CLEANLY today (verified on
% the pre-v3.64 preamble: 0 errors, PDF produced) but becomes FATAL once it is also
% written to the .lof — \codeline is \hbox to\linewidth{...}, and a full-width box in
% a list entry dies with "\pdfextension endlink cannot be used in vertical mode", no
% output PDF. Gaining a feature nobody uses (0 of 13 collections call \listoffigures)
% at the price of breaking a construct that works is a bad trade. Stock \caption has
% exactly this limitation and exactly this escape hatch, so the convention is familiar.
% \en{}, \code{}, \h{} and inline math in the short text are all safe — verified, they
% survive the write and typeset correctly in the list (RTL geometry confirmed: number
% rightmost at x=509.9, title left of it, dots and folio at the far left).
% With no optional argument the macro is byte-identical to v3.63 — zero risk to every
% existing document, and \listoffigures lists exactly the captions that opted in.
% BLD-37, TABLE side (the \figcap v3.74 idea, adapted): clicking a \ref to a
% table used to land MID-table — hyperref's automatic anchor inside
% \refstepcounter was emitted on the still-open horizontal line whose TAIL is
% the tabular (an hbox whose baseline is the table's vertical CENTER), so the
% jump scrolled the top half of the table above the viewport. The v3.74 entry
% deferred this side claiming "tablebox splits — no single box to measure";
% that deferral dissolves: the measurement target is the TABULAR itself at
% \tabcap ENTRY, while it is still the tail of the open list — the same
% situation \lakin@grabfig handles at an island's end (and the fleet has 0
% split tableboxes). So \lakin@grabtab runs as the FIRST token, before any
% typesetting can push the box away, and the anchor is re-emitted raised above
% the measured tabular. \lakintabraise is a SEPARATE register from
% \lakinfigraise (defined with it, below the LTR-island block): a captionless
% island intentionally leaves the fig register set, and a table must never
% inherit that stale figure height.
%
% WHY the anchor stays AT THE ENTRY SPOT and is NOT deferred into the caption
% paragraph the way \figcap defers it: at \par TeX DROPS a final glue item,
% and the automatic anchor's smash box sitting BEHIND the interword glue kept
% that glue alive — so every centered tabular has always sat half a space off
% true center. Deferring the anchor removed that box and was measured as 982
% differing glyph records: EVERY tabular line shifted 1.30pt (2.61pt for
% pattern 7's two minipages) while captions stayed put. Emitting the raised
% anchor at the old spot — through hyperref's own \Hy@raisedlink, exactly the
% macro the automatic path uses — keeps the line's node structure, and so the
% A/B glyph diff at 0, BY CONSTRUCTION. \Hy@raisedlink reads the raise from
% \HyperRaiseLinkDefault (default \baselineskip = the old mid-table dest); the
% group-local \edef freezes measured-height+42pt into it, and \smash inside
% \Hy@raisedlink hides any raise from the line's extent. \HyCnt@ProvideTheH-
% Counter keeps the dest NAME byte-identical to the automatic one (it defines
% \theHtbl, which \hyper@makecurrent prefers — without it the first measured
% caption was named tbl.1 where the automatic path said tbl.0.1, measured).
% vmode entry, or anything but a box at the tail (no tabular to measure),
% falls back to the plain \refstepcounter — byte-identical to the old path.
\makeatletter
\newcommand{\tabcap}[2][]{\lakin@grabtab
  % v3.79: arm the box-close correction. Set unconditionally and in BOTH
  % branches below's scope, so the plain-\refstepcounter fallback is corrected
  % too; \lakintabboxopen clears it at every \begin{tablebox}, so a \tabcap
  % outside a tablebox can never leave it armed for a later box.
  \global\lakintabcapseentrue
  \ifdim\lakintabraise>0pt
    \H@refstepcounter{tbl}\HyCnt@ProvideTheHCounter{tbl}\hyper@makecurrent{tbl}%
    \begingroup
      % A frameless table targets its measured top. The frame's historical
      % cushion belongs only to tablebox and is corrected at box close.
      \edef\HyperRaiseLinkDefault{\the\lakintabraise}%
      \iflakintabframed
        \edef\HyperRaiseLinkDefault{\the\dimexpr\lakintabraise+42pt\relax}%
      \fi
      \Hy@raisedlink{\hyper@anchorstart{\@currentHref}\hyper@anchorend}%
    \endgroup
    \global\lakintabraise 0pt\relax
  \else
    \refstepcounter{tbl}%
  \fi
  \lakinrtitle{#2}\par\nopagebreak\addvspace{6pt}%
  \def\lakin@short{#1}\ifx\lakin@short\@empty\else
    \addcontentsline{lot}{table}{\protect\numberline{\thetbl}{\protect\ignorespaces #1}}\fi
  {\centering\lakitcaptionstyle\mbox{\lakitcaptionlabelstyle{\tablename~\en{\thetbl}\lakitcaptionpunctuation}}\enspace #2\par}}
\makeatother

% ============================================================================
% Unified FIGURE caption — the \tabcap twin for TikZ/plot figures (v3.34.0).
% A figure (tikzpic) sits inline in the prose; put \figcap BELOW it for a
% numbered, \ref-able caption "איור N.M: תיאור" in the same style as \tabcap:
%   \begin{tikzpic} ... \end{tikzpic}
%   \figcap{\label{fig:foo}description}     % \label FIRST, inside the argument
% WHY a macro, not a hand-written gray line: a manual
%   {\centering\small\textcolor{cGray}{איור: ...}} with a \label inside has NO
% \refstepcounter, so the \label binds to the AMBIENT \@currentlabel (the last
% equation) — \ref then prints an equation number AS the figure number
% (measured on a real page: "איור 48" where 48 is equation 48). \figcap steps a
% real \c@fig counter first, so \label is correct, numbering is consistent, and
% refs are stable. Class-aware like \thetbl (chapter- or section-prefixed).
% ============================================================================
\makeatletter
\@ifundefined{c@chapter}{%
  \newcounter{fig}[section]
  \renewcommand{\thefig}{\ifnum\value{section}>0 \thesection.\fi\arabic{fig}}%
}{%
  \newcounter{fig}[chapter]
  \renewcommand{\thefig}{\ifnum\value{chapter}>0 \thechapter.\fi\arabic{fig}}%
}
\makeatother
\makeatletter
\newcommand{\figcap}[2][]{%
  \ifdim\lakinfigraise>0pt
    % v3.74: step WITHOUT hyperref's automatic anchor, then emit the anchor
    % raised above the measured figure. \H@refstepcounter is hyperref's saved
    % pre-hook \refstepcounter (sets the counter AND \@currentlabel — \label
    % in the caption arg binds exactly as before); \hyper@makecurrent supplies
    % \@currentHref, so the dest NAME is identical to the old automatic one —
    % \ref/\autoref/bookmarks are unchanged, only the dest'S POSITION moves.
    \H@refstepcounter{fig}\hyper@makecurrent{fig}%
    % The anchor is emitted INSIDE the caption paragraph below, as a zero-size
    % raisebox — NOT here in vertical mode, where a \raisebox starts an empty
    % paragraph and adds a blank line (measured: every glyph after the caption
    % shifted +18.8pt). Inside an existing line a [0pt][0pt] box is inert.
    \edef\lakin@anch{\@currentHref}%
    % \edef FREEZES the measured height into the pending macro: the deferred
    % \raisebox is typeset inside the caption paragraph AFTER the register is
    % reset to the 0pt sentinel below — a plain \def read 0 at use time
    % (measured: LAKINRAISE=0.0, dest moved 14pt instead of above the figure).
    \edef\lakin@pend{\noexpand\raisebox{\the\dimexpr\lakinfigraise+42pt\relax}[0pt][0pt]{%
      \noexpand\hyper@anchorstart{\lakin@anch}\noexpand\hyper@anchorend}}%
    \global\lakinfigraise 0pt\relax
  \else
    \refstepcounter{fig}\let\lakin@pend\relax
  \fi
  \lakinrtitle{#2}\par\nopagebreak\addvspace{4pt}%
  \def\lakin@short{#1}\ifx\lakin@short\@empty\else
    \addcontentsline{lof}{figure}{\protect\numberline{\thefig}{\protect\ignorespaces #1}}\fi
  {\centering\lakin@pend\lakitcaptionstyle\mbox{\lakitcaptionlabelstyle{\figurename~\en{\thefig}\lakitcaptionpunctuation}}\enspace #2\par}}
\makeatother

% ============================================================================
% LTR islands.  EVERY tikzpicture / quantikz MUST live inside one, or the whole
% figure's coordinates and Latin labels flip. (Hebrew node text still uses \h{}.)
% ============================================================================
% v3.74 (BLD-37): the \figcap hyperlink DESTINATION used to land at the caption
% line — BELOW the figure — so clicking a figure reference scrolled the figure
% off-screen (measured fleet-wide: 97/97 fig dests in one collection sat 19-54pt
% under their figure's ink; a 347pt figure was fully hidden after the jump).
% hypcap-style fix in two parts: (1) at the island's END, \lakin@grabfig grabs
% the just-typeset centered line with \lastbox, measures it, and puts it back
% with \nointerlineskip — the interline glue above the line was already
% contributed and is NOT re-added, so the page geometry is unchanged to the
% glyph (verified); (2) \figcap then steps the counter WITHOUT hyperref's
% automatic anchor (\H@refstepcounter = the pre-hyperref \refstepcounter) and
% emits the anchor itself, raised by the measured height — the dest lands ABOVE
% the figure top. A captionless island leaves numbering untouched (nothing is
% stepped at the island; the grab flag is cleared at every island begin, so a
% stale height from a captionless figure cannot leak into a later caption).
\newsavebox{\lakinfigbox}
\newdimen\lakinfigraise  % real \dimen: tex.setdimen cannot target a \newlength SKIP register (measured lua error)
% BLD-37 table side: a SEPARATE register for \tabcap. The two must not share —
% a captionless island INTENTIONALLY leaves \lakinfigraise set (the grab runs
% at every island END; only the NEXT island's begin clears it), so a shared
% register would hand a \tabcap after a captionless figure that stale FIGURE
% height. \lakintabraise is written only by \lakin@grabtab at \tabcap entry
% and consumed/reset inside \tabcap itself, so it can never go stale.
\newdimen\lakintabraise
\makeatletter
% v3.74b: measure WITHOUT touching the list. The first draft grabbed the box
% with \lastbox and put it back — and \lastbox DISCARDS the node's baseline
% shift, so a vertically-anchored quantikz circuit rose 2.57pt and everything
% after it sank 5.32pt (measured on the A/B glyph diff). LuaTeX lets us READ
% the tail node's dimensions without removing anything: geometry is identical
% by construction. \lakinfigraise=0pt is the "no figure measured" sentinel.
\def\lakin@grabfig{%
  \ifhmode
    \directlua{
      local n = tex.nest[tex.nest.ptr].tail
      local GLUE, PEN = node.id('glue'), node.id('penalty')
      while n and (n.id == GLUE or n.id == PEN) do n = n.prev end
      local HL, VL = node.id('hlist'), node.id('vlist')
      if n and (n.id == HL or n.id == VL) then
        tex.setdimen('global', 'lakinfigraise',
                     (n.height or 0) + (n.depth or 0))
      else
        tex.setdimen('global', 'lakinfigraise', 0)
      end}%
  \else
    \global\lakinfigraise 0pt\relax
  \fi}
% BLD-37 table side: \lakin@grabfig's twin, NOT a shared helper — the Lua body
% differs in exactly one place, the skip set. Between \end{tabular} and \tabcap
% an otherlanguage wrapper (patterns 3/5 of example_tables) leaves a \write
% whatsit and a dir node ON TOP of the interword glue (measured node dump:
% "whatsit/write <- dir <- glue(2.62) <- hlist"), so the table reader must
% also step over whatsit/dir markers — all zero-size — to reach the tabular.
% \lakin@grabfig deliberately does NOT get the wider skip set: its glue/penalty
% body is the one the v3.74 A/B proof froze, and widening it would newly
% measure (and so newly raise) anchors in island shapes that today take the
% plain path — an unproven behavior change on the FIG side. Anything that is
% not a box behind those markers still hits the 0pt sentinel and \tabcap falls
% back to the plain \refstepcounter.
% SECOND deliberate difference: this stores the box's HEIGHT, not height+depth.
% \figcap's anchor is emitted in the caption paragraph BELOW the figure, so it
% must climb the full extent; \tabcap's anchor sits ON the tabular's own line,
% whose baseline is the vcentered tabular's MIDPOINT — height alone reaches
% the top edge, and adding the depth was measured to overshoot by half the
% table (dests 80-135pt above their frame, INSIDE the preceding table's box).
\def\lakin@grabtab{%
  \ifhmode
    \directlua{
      local n = tex.nest[tex.nest.ptr].tail
      local GLUE, PEN = node.id('glue'), node.id('penalty')
      local WHAT, DIR = node.id('whatsit'), node.id('dir')
      while n and (n.id == GLUE or n.id == PEN
                   or n.id == WHAT or n.id == DIR) do n = n.prev end
      local HL, VL = node.id('hlist'), node.id('vlist')
      if n and (n.id == HL or n.id == VL) then
        tex.setdimen('global', 'lakintabraise', n.height or 0)
      else
        tex.setdimen('global', 'lakintabraise', 0)
      end}%
  \else
    \global\lakintabraise 0pt\relax
  \fi}
\makeatother
% ============================================================================
% v3.79 (BLD-37 table side, CORRECTION): put the anchor at the TOP OF THE BOX,
% instead of guessing the distance to it.
% ============================================================================
% v3.77 raised the \tabcap anchor by (height of the box at the tail of the open
% horizontal list) + 42pt. That is a MEASUREMENT taken BEFORE line breaking, and
% it is short by every term it cannot see. Fleet-measured on the 12 built
% collections (299 tbl.* dests, ownership established through each table's own
% CAPTION, never through "which frame encloses the dest"): 3 dests landed INSIDE
% their own tablebox, below the frame top —
%   computer-architecture       tbl.14.1  -46.4pt   three stacked tabulars
%   computer-architecture       tbl.30.5  -13.5pt   tabular 2.85pt over \linewidth
%   advanced-electronic-devices tbl.B.1   -59.2pt   tabular 2.56pt over \linewidth
% Two missing terms, both reproduced in a minimal document (probe P2/P3):
%   (1) VERTICAL LIST ALREADY CONTRIBUTED inside the box — a second and third
%       tabular, or a heading above the table. The tail read sees only the LAST
%       box, so 77.7pt of stacked tabulars is invisible to the raise.
%   (2) WHICH LINE THE ANCHOR LANDS ON. When the tabular alone overfills
%       \linewidth (459.6121pt in a tablebox), the interword glue after
%       \end{tabular} becomes the chosen breakpoint and \Hy@raisedlink is
%       emitted on the NEXT line — one (depth + lineskip) lower. That term does
%       not EXIST at \tabcap time: line breaking has not happened yet. No
%       measurement taken there can ever recover it. (Proof: probe P3, the same
%       table 2.84912pt too wide, fails at -13.5pt; probe P4, the identical
%       table with \tabcolsep 2pt so it fits, is correct at +32.0pt.)
% So the raise is not computed at \tabcap at all. It is CORRECTED at
% \end{tablebox}, from inside mdframed's own body vbox (endinnercode runs after
% \par\unskip, i.e. after line breaking, while tex.nest[tex.nest.ptr] is still
% that vbox's vertical list). There the anchor's true offset from the top of the
% box content is not estimated but WALKED, and the \Hy@raisedlink shift is set
% so the dest lands 42pt above it — which is 42pt - innertopmargin(9pt) -
% linewidth(0.9pt) = 32.0pt above the FRAME top, the exact position the 295
% already-correct dests have today. That is why this is a correction and not a
% re-placement: the good sites do not move.
%
% WHY THE GLYPH DIFF IS 0 BY CONSTRUCTION — the only thing written is the
% `shift` field of the hlist that \Hy@raisedlink already built, and the guard
% below refuses to write it unless that box is 0x0x0. A zero-size box carries no
% dimension into its parent's packing, the packing already happened, and shift
% is not an input to it. Nothing else in the list is touched: no node is added,
% removed, resized or reordered. Measured, not assumed: assets/example_tables.tex
% and the whole computer-architecture collection A/B at 0 differing pymupdf
% rawdict glyph records.
%
% WHY NOT the alternatives (all prototyped and measured before this one):
%   * accumulate the vertical list at \tabcap time — fixes term (1), CANNOT fix
%     term (2); 2 of the 3 real failures survive. A fix that leaves the majority
%     of the defect standing is the false all-clear this repo exists to avoid.
%   * \zref-savepos two-pass — correct on all shapes, but it makes the dest
%     POSITION aux-derived (DEFECTS BLD-1/2), a single-pass compile silently
%     ships an unraised anchor, and the bare y-difference has no page guard.
%   * emit the anchor at \begin{tablebox} and let \tabcap adopt it — needs the
%     tbl counter stepped at box OPEN to keep the dest NAME byte-identical; a
%     captionless tablebox then either shifts every later table number or leaves
%     an orphan dest that the next table duplicates (a HARD build gate). This
%     correction touches neither the counter nor the name: \tabcap's naming path
%     is v3.77's, unchanged, so resolve_names() is byte-identical (verified).
%   * suppressing the line break (a penalty before the glue) — that MOVES the
%     tabular: the break currently discards the interword glue, so forbidding it
%     re-centers the two overfull tables. Reflow, rejected.
%
% SCOPE. The correction is armed only by \tabcap ITSELF (\lakintabcapseen),
% cleared at every box open, so a captionless tablebox is a no-op and a \tabcap
% hosted OUTSIDE a tablebox (the fleet has exactly one, computer-architecture
% tbl.14.2 in an exbox) keeps v3.77's measured raise. Both fallbacks are
% therefore never worse than today; and because hyperref's AUTOMATIC anchor uses
% the same \Hy@raisedlink shape, the plain-\refstepcounter path (vmode entry, or
% prose between \end{tabular} and \tabcap — probe P6, -51.5pt) is corrected too.
% KNOWN LIMIT — a tablebox that SPLITS across pages. The offset is measured from
% the top of the UNSPLIT body, so a dest in segment 2 gets raised past its own
% segment. MEASURED on a forced split (probe: 40 filler lines + a table): the
% dest went from y=585.0 on page 2 (inside segment 2, at the caption — the old
% BLD-37 symptom) to y=-113.4, i.e. ABOVE the page top. Viewers clamp that to the
% page top, which happens to show segment 2 whole, but an out-of-page coordinate
% is not something to ship on purpose. Hence the \textheight guard in
% lakintab.fix: a body TALLER THAN THE TEXT BLOCK cannot fit on any page, so it
% must split — the correction is skipped and v3.77's behaviour stands unchanged
% (verified: the probe returns to y=585.0). The guard is deliberately EXACT, not
% a page-fitting heuristic: it is sound (a taller-than-page body always splits)
% but PARTIAL (a shorter body that starts low on a page could still split). That
% residue is unreachable in practice — v3.37's everyline + the caption's
% \nopagebreak glue move such boxes WHOLE instead (probe N7: a 22-row table
% pushed down by 430pt moved whole, 0 segments), and the fleet has 0 split
% tableboxes, re-counted after this change. Reimplementing mdframed's own
% page-fitting arithmetic to close the residue would be a heuristic with its own
% failure mode, which is the trade this repo has repeatedly refused.
\newif\iflakintabcapseen
\newif\iflakintabframed
% The Lua lives in ONE \directlua at load time, and is written with NO
% backslash, %, # or ~ : \directlua expands its argument like \edef, so an
% active ~ would expand to \nobreakspace and a % would eat the rest of the line
% (the second is how the first draft of this block died). "not (a == b)" stands
% in for Lua's inequality operator for the same reason.
\directlua{
  lakintab = { }
  lakintab.HL = node.id("hlist")
  lakintab.VL = node.id("vlist")
  lakintab.RU = node.id("rule")
  lakintab.GL = node.id("glue")
  lakintab.KE = node.id("kern")
  lakintab.WH = node.id("whatsit")
  lakintab.TM = node.id("temp")
  lakintab.WD = node.subtype("pdf_dest")
  function lakintab.d(v)
    if v == nil then return 0 end
    if v < -1000000000 then return 0 end
    return v
  end
  function lakintab.scan(list, y, horiz, parent, target)
    local x = list
    while x do
      local i = x.id
      if i == lakintab.HL or i == lakintab.VL then
        if horiz then
          local b = y + x.shift
          if i == lakintab.HL then
            lakintab.scan(x.head, b, true, x, target)
          else
            lakintab.scan(x.head, b - lakintab.d(x.height), false, x, target)
          end
        else
          if i == lakintab.HL then
            lakintab.scan(x.head, y + lakintab.d(x.height), true, x, target)
          else
            lakintab.scan(x.head, y, false, x, target)
          end
          y = y + lakintab.d(x.height) + lakintab.d(x.depth)
        end
      elseif i == lakintab.WH then
        if x.subtype == lakintab.WD then
          local id = x.dest_id
          if type(id) == "string" and string.sub(id, 1, 4) == "tbl." then
            if horiz and parent and parent.width == 0
               and parent.height == 0 and parent.depth == 0 then
              parent.shift = parent.shift + (target - y)
            end
          end
        end
      elseif not horiz then
        if i == lakintab.RU then
          y = y + lakintab.d(x.height) + lakintab.d(x.depth)
        elseif i == lakintab.GL or i == lakintab.KE then
          y = y + lakintab.d(x.width)
        end
      end
      x = x.next
    end
  end
  function lakintab.fix(target, maxh)
    local nst = tex.nest[tex.nest.ptr]
    if nst == nil then return end
    local h = nst.head
    if h == nil then return end
    if h.id == lakintab.TM then h = h.next end
    local tot = 0
    local x = h
    while x do
      local i = x.id
      if i == lakintab.HL or i == lakintab.VL or i == lakintab.RU then
        tot = tot + lakintab.d(x.height) + lakintab.d(x.depth)
      elseif i == lakintab.GL or i == lakintab.KE then
        tot = tot + lakintab.d(x.width)
      end
      x = x.next
    end
    if tot > maxh then return end
    lakintab.scan(h, 0, false, nil, target)
  end
}
% @-FREE NAMES ON PURPOSE: these two are named in the `tablebox` key list at
% \newmdenv above, which is tokenised OUTSIDE any \makeatletter region — a
% \lakin@... there would tokenise as \lakin followed by the letters @tabopen.
\newcommand{\lakintabboxopen}{\lakintabframedtrue\global\lakintabcapseenfalse}
\newcommand{\lakintabboxclose}{%
  \iflakintabcapseen
    \directlua{lakintab.fix(-\number\dimexpr\lakintabcushion\relax,
                            \number\dimexpr\textheight\relax)}%
    \global\lakintabcapseenfalse
  \fi\lakintabframedfalse}
% The cushion above the box's CONTENT top. 42pt puts the dest 32.0pt above the
% frame top (42pt - innertopmargin 9pt - linewidth 0.9pt), i.e. exactly where
% the 295 correct dests already are — the correction moves only the broken ones.
\newcommand{\lakintabcushion}{42pt}
\newenvironment{qcirc}{\par\medskip\global\lakinfigraise 0pt\relax\begin{otherlanguage}{english}\begin{center}}
  {\csname lakin@grabfig\endcsname\end{center}\end{otherlanguage}\par\medskip}
\newenvironment{tikzpic}{\par\medskip\global\lakinfigraise 0pt\relax\begin{otherlanguage}{english}\begin{center}}
  {\csname lakin@grabfig\endcsname\end{center}\end{otherlanguage}\par\medskip}

% ============================================================================
% Code listings.  ★★ v3.4 THE BOX-OVERFLOW FIX:
%   resetmargins=true   → the listing IGNORES the enclosing environment's
%                         indentation and measures from the PAGE. v3.4 stated
%                         this backwards ("measured from the ENCLOSING box"),
%                         and v3.37 then blanked in-box line numbers because
%                         the page-anchored listing put its hanging digits on
%                         top of the box rule. Measured v4.5.0: with the key
%                         FALSE the listing follows the box and the code moves
%                         from x=75.26 to x=97.74, clearing the rule at 72.07.
%                         The key now lives in one \lstset below instead of in
%                         each style, so a box can override it — codebox does.
%   framexleftmargin=0  → the frame (left bar) does not extend left of the code.
%   xleftmargin keeps a small indent so the line numbers have room INSIDE.
% asmblock wraps the listing in an LTR island (code is LTR).
% ============================================================================
\lstdefinestyle{asm}{
  basicstyle=\ttfamily\small,
  backgroundcolor=\color{cBlueBg!55},
  frame=leftline,framerule=2.5pt,rulecolor=\color{cBlue},
  numbers=left,numberstyle=\tiny\color{cGray},numbersep=8pt,
  xleftmargin=2.2em,framexleftmargin=0pt,framexrightmargin=0pt,
  breaklines=true,columns=fullflexible,keepspaces=true,
  showstringspaces=false,aboveskip=8pt,belowskip=8pt,
  % v3.36: without a comment DELIMITER the commentstyle below never fires —
  % listings has no `language` here, so it does not know '#' opens a comment,
  % and every asm listing rendered flat black (measured: a whole-page pixel
  % sweep found ZERO cGreen glyphs across 135 shipped listings). morecomment
  % gives it the delimiter; color-only change, glyph widths untouched.
  morecomment={[l]{\#}},
  commentstyle=\color{cGreen},
}
\lstdefinestyle{cstyle}{
  language=C,basicstyle=\ttfamily\small,
  backgroundcolor=\color{cGreenBg!55},
  frame=leftline,framerule=2.5pt,rulecolor=\color{cGreen},
  numbers=left,numberstyle=\tiny\color{cGray},numbersep=8pt,
  xleftmargin=2.2em,framexleftmargin=0pt,framexrightmargin=0pt,
  breaklines=true,columns=fullflexible,keepspaces=true,
  showstringspaces=false,keywordstyle=\color{cBlue}\bfseries,
  commentstyle=\color{cGray}\itshape,aboveskip=8pt,belowskip=8pt,
}
% v3.36: pystyle/htmlstyle — colored twins of cstyle for the Python/HTML
% listings that used to repeat a ~200-char options blob with NO color
% (keywords rendered bold-black, comments italic-black — measured; C alone
% was colored). Geometry is copied VERBATIM from the shipped blob (frame,
% margins, skips) so swapping a call site to style=pystyle moves NOTHING —
% only keywordstyle/commentstyle/stringstyle are added (bold stays bold,
% italic stays italic, mono metrics identical => zero reflow, proven by
% per-page word-bbox diff). breaklines/breakatwhitespace stay per-site.
% v3.78: numbers=left joins both styles, with the exact numberstyle/numbersep
% keys asm/cstyle carry — fleet-wide consistency (author decision: all four
% styles numbered; "asm/cstyle numbered, py/html not" was an accident of the
% v3.36 geometry freeze, not a design). The numbers hang LEFT of the code box
% (listings sets them with \llap, so no overfull hbox can arise and the code
% x is untouched — measured: code frozen, digits at code_x-11.8pt); inside a
% colored box they are blanked like everyone else's (v3.37 hook), and the
% clean variant restores them in its frameless boxes (v3.78, above).
\lstdefinestyle{pystyle}{
  language=Python,basicstyle=\ttfamily\small,
  columns=fullflexible,keepspaces=true,
  frame=single,rulecolor=\color{cTblLine},
  numbers=left,numberstyle=\tiny\color{cGray},numbersep=8pt,
  xleftmargin=0.5em,framexleftmargin=0pt,framexrightmargin=0pt,
  showstringspaces=false,aboveskip=8pt,belowskip=8pt,
  keywordstyle=\color{cBlue}\bfseries,
  commentstyle=\color{cGray}\itshape,
  stringstyle=\color{cGreen},
}
% v4.5.0: hoisted out of the four styles so an enclosing box can override it.
% Every listing still resolves to true unless something later says otherwise,
% so shipped output is unchanged (verified byte-for-byte on all eight PDFs).
\lstset{resetmargins=true}
\lstdefinestyle{htmlstyle}{
  language=HTML,basicstyle=\ttfamily\small,
  columns=fullflexible,keepspaces=true,
  frame=single,rulecolor=\color{cTblLine},
  numbers=left,numberstyle=\tiny\color{cGray},numbersep=8pt,
  xleftmargin=0.5em,framexleftmargin=0pt,framexrightmargin=0pt,
  showstringspaces=false,aboveskip=8pt,belowskip=8pt,
  keywordstyle=\color{cBlue}\bfseries,
  commentstyle=\color{cGray}\itshape,
  stringstyle=\color{cGreen},
}
\newenvironment{asmblock}{\par\begin{otherlanguage}{english}}{\end{otherlanguage}\par}

% <hllpdf:listing-smart-quotes>
% v4.21.0: four Unicode punctuation characters must enter listings' token
% buffer. Unregistered characters bypass it: a closing quote was printed
% BEFORE the pending word, and an opening quote BEFORE a pending brace.
% Extend the standard character table, then activate only these codepoints
% in the listing's local character-table scope. Do this BEFORE the package's
% delimiter/alsoother/literate hooks so explicit user mappings still win.
% This is not universal Unicode support, nor a conversion to ASCII quotes.
% See listings developer guide 15.3.1 and references/listing-smart-quotes.md.
\makeatletter
\lst@CCPut \lst@ProcessOther {"2018}{"2019}{"201C}{"201D}\z@
\lst@AddToHookAtTop{SelectCharTable}{%
  \lst@ifactivechars
    \catcode"2018=\active \catcode"2019=\active
    \catcode"201C=\active \catcode"201D=\active
  \fi}
\makeatother
% </hllpdf:listing-smart-quotes>

% <hllpdf:listing-technical-symbols>
% v4.30: bounded extension for common technical punctuation in LTR listings.
% Without character-table registration these glyphs bypass the token buffer
% and can move before the preceding word even inside quoted strings.
% Keep the original four-quote block separate for its standing mutation test.
% This is source-faithful printing, not syntax validation or general Unicode.
\makeatletter
\lst@CCPut \lst@ProcessOther
  {"2013}{"2014}{"2026}{"2190}{"2192}{"2194}{"21D2}
  {"2212}{"221E}{"2260}{"2264}{"2265}\z@
\lst@AddToHookAtTop{SelectCharTable}{%
  \lst@ifactivechars
    \catcode"2013=\active \catcode"2014=\active \catcode"2026=\active
    \catcode"2190=\active \catcode"2192=\active \catcode"2194=\active
    \catcode"21D2=\active \catcode"2212=\active \catcode"221E=\active
    \catcode"2260=\active \catcode"2264=\active \catcode"2265=\active
  \fi}
\makeatother
% </hllpdf:listing-technical-symbols>


\usepackage[hidelinks,unicode]{hyperref}   % LAST, after colors. unicode = Hebrew bookmarks.
% v3.62: THE ISOLATION MACROS CORRUPTED EVERY PDF BOOKMARK THEY APPEARED IN.
% \en is defined (line ~155) as
%     \foreignlanguage{english}{\language=\csname l@nohyphenation\endcsname #1}
% and in a heading that body is written verbatim to the .toc and handed to
% hyperref's \pdfstringdef. hyperref strips the two CONTROL SEQUENCES (\language,
% \l@nohyphenation) but the plain `=` CHARACTER token between them is not a control
% sequence, so it survives into the PDF string: \chapter{שער \en{CNOT} והצמדה}
% produced the bookmark "שער =CNOT והצמדה". The printed TOC and the heading on the
% page are CORRECT — only the navigation pane and the viewer window title are wrong,
% which is why it went unnoticed. Measured on the shipped fleet: 70 corrupted
% bookmarks across 5 PDFs, from 177 unprotected \en-in-heading sites. Same corruption
% reached /Title and /Subject via the \hypersetup{pdftitle=...} the docs recommend.
% This hook tells hyperref to use the plain argument for all three isolation macros
% when building a PDF string. It does NOT touch the .toc (the printed TOC is
% byte-unchanged) and it does NOT mask the math signal — a heading with $...$ still
% logs "removing math shift", so BIDI-15 detection is intact.
% BONUS, measured: it also makes \code{} in an UNSTARRED heading COMPILE. BLD-10
% recorded that as needing \texorpdfstring or a starred heading; the real root cause
% is hyperref's \pdfstringdef scan of \code's \ifmmode...\else...\fi body — which is
% exactly why the identically-bodied \texttt survived (hyperref pre-neutralizes
% \texttt in PDF strings) and why \DeclareRobustCommand could never help.
% v3.64: \def\\{ } added. A manual line break in a heading WELDS the two words
% together in the bookmark: \chapter{שורה ראשונה\\שורה שנייה} produced the outline
% entry 'שורה ראשונהשורה שנייה' — "ראשונהשורה" is not a word. hyperref drops the
% \\ and puts nothing in its place, so the break that looks like a space on the
% page is a zero-width join in the navigation pane. The spaced form
% 'מילה אחת \\ מילה שנייה' escapes with a harmless DOUBLE space ('אחת  מילה'),
% which is why this never got noticed. Both are measured on a real outline.
% 0 live instances in the corpus — but SKILL.md and content-style.md both RECOMMEND
% a manual heading break, so this is a primed trap rather than a dormant one, and a
% one-token preamble fix retires it for good instead of asking authors to remember.
% Mapping \\ to a space fixes the WELD, which is the actual defect (a nonexistent
% word in the navigation pane). It does not normalise whitespace: measured after the
% fix, the welded form becomes 'שורה ראשונה שורה שנייה' (correct) and the already-
% spaced form goes from two spaces to three. That is cosmetic — viewers render it as
% a wider gap, and no word is corrupted. A per-heading \texorpdfstring is still the
% way to get an exactly-tuned bookmark string.
\pdfstringdefDisableCommands{\def\en#1{#1}\def\h#1{#1}\def\doclang#1{#1}\def\code#1{#1}\def\codeline#1{#1}\def\\{ }}
% v3.70: RESTORE \url LINE-BREAKING. url.sty freezes its break penalties at LOAD
% time (\mathchardef\UrlBreakPenalty=\binoppenalty), and v3.45's global
% \binoppenalty=10000 (inline-math atomicity) therefore priced EVERY breakpoint
% inside \Url at 10000 = never break. Measured: a long \url in Hebrew prose is
% one unbroken line, ~100pt overfull, invading the margin band the ink gates
% deliberately tolerate. With the two lines below: 0 overfull, the URL breaks
% after '/', the RTL split is CORRECT (fragment order forward), and each line
% gets its own link annotation exactly covering its fragment (measured).
% Confined to \Url's group — inline-math atomicity (TYP-7) is untouched.
\mathchardef\UrlBreakPenalty=200
\mathchardef\UrlBigBreakPenalty=100
\ifx\lakitmainlanguage\hllpdfenglish
  \hypersetup{pdflang=en}
\else
  \hypersetup{pdflang=he}
\fi  % /Lang is the document base locale (not per-span accessibility tagging).
% Set a per-document \hypersetup{pdftitle={...}} too — see SKILL.md.

% v3.52 (round-2 audit): two RTL bibliography/cross-ref defects, both invisible to
% every gate and both affecting the FIRST author to use the feature.
% (a) \@biblabel's default [#1] brackets are bidi-neutral and mirror to ]1[ under
% the RTL bibliography base — every \thebibliography reference list shows ]1[.
% LTR-isolating the label restores [1]. (The render gate's mirror rule fingerprints
% round parens )1(, so square brackets slipped through.)
\makeatletter
\renewcommand*{\@biblabel}[1]{\foreignlanguage{english}{[#1]}}
% v3.54: the SAME neutral-bracket trap on the INLINE side — a multi-key \cite{a,b}
% under the RTL base lays its European-number list right-to-left, so [1, 2] renders
% "[2,1]" and [1, 2, 3] renders "[3,2,1]" (verified by glyph geometry — the reader
% sees the sources in REVERSE order; a note form \cite[p.5]{a} also reorders to
% "[p. 5,1]"). LTR-isolate the whole bracket group (number list AND optional note),
% exactly like \@biblabel above: [1, 2] / [1, 2, 3] / [1, p. 5] all render correct.
% Placed AFTER hyperref so it wins; hyperref hyperlinks each number inside #1 (via
% \@citex), which this wrapper leaves intact. Mirrors the LaTeX kernel \@cite shape.
% v4.9.1: directional isolation must NOT switch the language here. Babel's
% safe-active citation wrapper passes its key list through \bbl@tempa;
% \foreignlanguage reuses that scratch macro before the deferred citation
% loop expands it. The old line therefore cited the literal key "english".
% If such a key existed, even the undefined-reference gate stayed green.
% \babelsublr isolates the bracket group without selecting a locale or
% changing that key list. Preserve native links, key order and optional notes.
\def\@cite#1#2{\babelsublr{[{#1\if@tempswa , #2\fi}]}}
\makeatother

% v4.11.0: a reference inside align's \text inherited HEBdisptag=true,
% so its enclosing hbox mirrored the brackets despite a clean compile.
% Isolate the complete printed reference without selecting a Babel language.
% Also neutralize the display-only normalfont direction reset locally. A
% direction-only patch inverted 12 to 21 inside equation in the measured probe.
% Plain maketag is local to the reference: mathtools' display-side wrapper
% otherwise globally consumes manual_tag while formatting a nested reference.
% Preserve each original command (including the selective auxiliary write),
% current tag forms, native links, and showonlyrefs on/off transitions.
\makeatletter
\NewCommandCopy{\HEB@nativeeqref}{\eqref}
\newcommand{\HEB@eqrefisolate}[1]{%
  \babelsublr{\HEBdisptagfalse\let\bbl@eqnodir\relax\let\maketag@@@\ORIGmaketag@@@#1}}
\protected\def\eqref#1{\HEB@eqrefisolate{\HEB@nativeeqref{#1}}}
\ExplSyntaxOn
\NewCommandCopy{\HEB@nativeMTeqref}{\MT_extended_eqref:n}
\cs_set_protected:Npn \MT_extended_eqref:n #1
  {\HEB@eqrefisolate{\HEB@nativeMTeqref{#1}}}
\ExplSyntaxOff
\makeatother

% v4.11.0: custom caption counters are not part of hyperref's English
% locale, and Hebrew has no upstream hyperref locale in this toolchain.
% Restore the appropriate names at every Babel language switch, including
% return from a nested inline island. Do not localize the destination keys.
\makeatletter
\newcommand{\HEB@refnames}{%
\def\chapterautorefname{פרק}%
\def\sectionautorefname{סעיף}%
\def\subsectionautorefname{תת־סעיף}%
\def\subsubsectionautorefname{תת־תת־סעיף}%
\def\paragraphautorefname{פסקה}%
% v4.31: hyperref uses a distinct counter name at the deepest standard level.
\def\subparagraphautorefname{תת־פסקה}%
\def\equationautorefname{משוואה}%
\def\figureautorefname{איור}%
\def\tableautorefname{טבלה}%
\def\partautorefname{חלק}%
\def\appendixautorefname{נספח}%
\def\footnoteautorefname{הערת שוליים}%
\def\itemautorefname{סעיף}%
\def\Itemautorefname{סעיף}%
\def\figautorefname{איור}%
\def\tblautorefname{טבלה}%
\def\pageautorefname{עמוד}% 
}
\addto\extrashebrew{\HEB@refnames}
\newcommand{\ENG@refnames}{%
  \def\subparagraphautorefname{subparagraph}%
  \def\chapterautorefname{chapter}%
  \def\sectionautorefname{section}%
  \def\subsectionautorefname{subsection}%
  \def\subsubsectionautorefname{subsubsection}%
  \def\paragraphautorefname{paragraph}%
  \def\equationautorefname{Equation}%
  \def\figureautorefname{Figure}%
  \def\tableautorefname{Table}%
  \def\partautorefname{Part}%
  \def\appendixautorefname{Appendix}%
  \def\footnoteautorefname{footnote}%
  \def\itemautorefname{item}%
  \def\figautorefname{Figure}%
  \def\tblautorefname{Table}%
  \def\Itemautorefname{item}%
  \def\pageautorefname{page}%
}
\addto\extrasenglish{\ENG@refnames}
\ifx\lakitmainlanguage\hllpdfenglish\ENG@refnames\else\HEB@refnames\fi
\makeatother

% ============================================================================
% Optional running head (v3.39): \lakithead{<book title>} — OFF by default.
% A document opts in with ONE line after \input{preamble.tex}:
%   \lakithead{מבנה המחשב}
% Every content page then carries the series' orientation line: chapter name
% at the reading edge, book title opposite, thin rule. (fancyhdr's [L]/[R]
% are PHYSICAL; under the RTL page direction [L] lands at the RIGHT/reading
% edge — measured, consistent across the corpus.) Packaged like the topic
% rail: always loaded, inert without the call. This centralizes the 8-line
% block that 8 lakits used to copy by hand — the same duplication vector
% that once let a lakit preamble go ~100 lines stale. Class-aware: report
% fills the mark from \chapter, article from \section. NB: a titlepage env
% self-empties; front matter that sets \pagestyle{empty} must restore with
% \pagestyle{fancy} afterwards (NOT plain — that discards the head).
% ============================================================================
\usepackage{fancyhdr}
% v3.61: \chapter* left a STALE running head — pages of an UNNUMBERED chapter
% carried the PREVIOUS chapter's title. report.cls routes \chapter* through
% \@schapter, which never calls \chaptermark, so \leftmark keeps whatever the last
% numbered chapter left. Measured: with \chapter{פרק ממוספר} then
% \chapter*{מבוא לא ממוספר}, page 4 (inside the starred chapter) printed
% "פרק ממוספר". Invisible on the chapter's FIRST page (\thispagestyle{plain}), so it
% only shows from the second page on — the pages an author is least likely to check.
% An unnumbered front-matter chapter (\chapter*{מבוא}, \chapter*{הקדמה}) is everyday
% Hebrew-academic structure, and the corpus already carries HAND-WRITTEN \markboth
% lines right after \chapter* — the exact per-document duplication \lakithead exists
% to retire. Patched INSIDE \lakithead so the macro stays inert when unused; a
% document's own later \markboth still wins, so those hand-patches do not regress.
\makeatletter
% <hllpdf:starred-heading-scope>
\newif\iflakin@starmarkpatched
\newcommand{\lakistarmark}{%
  \iflakin@starmarkpatched\else
  % Dispatch on the actual class counter, not hyperref's shared internals.
  \ifdefined\c@chapter
    \let\HEBorigschapter\@schapter
    % Record only after the native heading has reached its page. A mark
    % before the native routine can advertise the next heading too early.
    \def\@schapter##1{\HEBorigschapter{##1}\chaptermark{##1}}%
  \else
    % @ssect receives layout arguments, NOT the section name or level.
    % Capture the name in @startsection without introducing a TeX group:
    % grouping sectioning would break the native after-heading state.
    \def\HEB@sectionkind{}%
    \def\HEB@mainsection{section}%
    \let\HEBorigstartsection\@startsection
    \def\@startsection##1##2##3##4##5##6{%
      \def\HEB@sectionkind{##1}%
      \HEBorigstartsection{##1}{##2}{##3}{##4}{##5}{##6}}%
    \let\HEBorigssect\@ssect
    \def\@ssect##1##2##3##4##5{%
      \HEBorigssect{##1}{##2}{##3}{##4}{##5}%
      % Lower starred headings must not replace the article's section mark.
      \ifx\HEB@sectionkind\HEB@mainsection
        \sectionmark{##5}%
      \fi
      \ignorespaces}%
  \fi
  \global\lakin@starmarkpatchedtrue
  \fi}
% </hllpdf:starred-heading-scope>
% <hllpdf:running-language>
% Choose the wrapper when the mark is written, not when a later page ships.
% English punctuation must remain LTR even after returning to Hebrew prose.
\newcommand{\HEB@runningmark}[1]{%
  \iflanguage{english}{\markboth{\en{#1}}{}}{\markboth{\h{#1}}{}}}
% </hllpdf:running-language>
\makeatother
% v4.2.1 (NAV): every collection already names itself in \lakithead — feed the
% PDF /Title from the same declaration, so the navigation gate's
% PDF_REQUIRE_TITLE has a universal source and no fleet member ships
% title-less by missing a second declaration site. Measured before this: 14
% of 16 collections shipped /Title-empty while every one carried a perfectly
% good \lakithead string (plain Hebrew in all 16 — no math, no formatting, so
% the PDF-string path is clean). An explicit title always wins: set it with
% \lakitpdftitle{...} anywhere before \lakithead, or with a plain
% \hypersetup{pdftitle=...} — the fallback detects that via hyperref's
% \@pdftitle, probe-verified \@empty-equal until the first explicit set.
\makeatletter
\newif\iflakit@pdftitleset
\newcommand{\lakitpdftitle}[1]{%
  \hypersetup{pdftitle={#1}}\global\lakit@pdftitlesettrue}
\newcommand{\lakithead}[1]{%
  \iflakit@pdftitleset\else
    \ifx\@pdftitle\@empty
      \hypersetup{pdftitle={#1}}\global\lakit@pdftitlesettrue
    \fi
  \fi
  \pagestyle{fancy}\fancyhf{}%
  \lakistarmark

  \fancyhead[L]{\small\textcolor{cGray}{\leftmark}}%
  % <hllpdf:fixed-head-language>
  % The fixed document title has its OWN language snapshot, independent of
  % the current section/chapter mark. Choose the branch NOW, not at shipout:
  % the page may be Hebrew by the time an English title is finally printed.
  \iflanguage{english}{%
    \fancyhead[R]{\small\textcolor{cGray}{\en{#1}}}%
  }{%
    \fancyhead[R]{\small\textcolor{cGray}{\h{#1}}}%
  }%
  % </hllpdf:fixed-head-language>
  \fancyfoot[C]{\thepage}%
  \renewcommand{\headrulewidth}{0.4pt}%
  \ifdefined\chaptermark
    \renewcommand{\chaptermark}[1]{\HEB@runningmark{##1}}%
  \else
    \renewcommand{\sectionmark}[1]{\HEB@runningmark{##1}}%
  \fi}
\makeatother

% v3.39.1: path-independent byte-reproducibility. LuaTeX derives the trailer
% /ID from the OUTPUT PATH, so two clones of the same source at different
% directories produced different sha256 (measured: the ONLY delta was the
% 62-byte /ID line). Bit 512 suppresses exactly that line — measured: Creator/
% Producer/CreationDate/ModDate and all content untouched. /ID is optional
% per ISO 32000 (required only for ENCRYPTED PDFs and PDF/A — neither is a
% target here). If a future workflow needs an ID (PDF/A, Acrobat revision
% matching), delete this line and it returns.
\pdfvariable suppressoptionalinfo \numexpr 512\relax

% v3.75: BOLD MATH IN BOLD HEADINGS. Heading text is \bfseries but math inside
% it stayed regular-weight CM — a hairline $V_T$ inside a heavy 24.8pt chapter
% title (measured: CMMI12 fl=6 next to FrankRuhlLibre-Bold fl=20; 155 live
% heading-with-math sites fleet-wide). \boldmath rides along wherever the class
% sets \bfseries for a heading; A/B measured: $V_T$ -> CMMIB10 fl=22, 0 missing
% characters (the bold CM faces ship with this TeX Live). Guarded per-macro so
% the article class (no \@makechapterhead) and future class changes stay safe.
% v4.0.1 (adoption fix): the SEARCH patterns above carry no parameter token.
% v4.0.0 made them specific by including the class's own #1/#2, and measured on
% TeX Live 2023 that never matches: a \patchcmd search pattern containing a
% parameter token cannot match while the call sits inside \@ifundefined{...}{}{...}
% (probed both #1 and ##1 -> FAIL; the same pattern at top level succeeds). Every
% one of these patches therefore failed, and because v4.0.0 also replaced the
% empty failure branch with \PackageError, the failure became FATAL: all four
% bundled example assets stopped building ("Cannot patch chapter-title math
% weight -> Fatal error, no output PDF"). The \PackageError branches are KEPT —
% a silent no-op here is what let v3.75's fix rot unnoticed — only the parameter
% tokens are dropped, so the patches match and a real class change still errors.
\makeatletter
\@ifundefined{@makechapterhead}{}{%
  \patchcmd{\@makechapterhead}{\huge\bfseries \@chapapp}
    {\huge\bfseries\boldmath \@chapapp}{}
    {\PackageError{hebrew-lualatex-pdf}{Cannot patch chapter-number math weight}{Unsupported document class}}%
  \patchcmd{\@makechapterhead}{\Huge \bfseries}
    {\Huge \bfseries\boldmath}{}
    {\PackageError{hebrew-lualatex-pdf}{Cannot patch chapter-title math weight}{Unsupported document class}}%
  \patchcmd{\@makeschapterhead}{\Huge \bfseries}
    {\Huge \bfseries\boldmath}{}
    {\PackageError{hebrew-lualatex-pdf}{Cannot patch starred chapter math weight}{Unsupported document class}}}
% v4.0.1: \part sits under a THREE-DEEP wrapper chain, so the class-based anchor
% is not in \@part at all. Measured with \show: hyperref rewrites \@part to
% `...\phantomsection\fi \H@old@part`, and nameref rewrites THAT to
% `\NR@gettitle{#1}\NR@part[{#1}]{#2}` — the class body ends up in \NR@part.
% v4.0.0 patched \@part and made failure fatal, which is why all seven bundled
% examples stopped building. \part is NOT a dead heading level here: 188 live uses
% across the fleet, so the bold-math fix is worth getting right rather than
% dropping. We walk the chain innermost-first and stop at the first hit. Failure is
% a WARNING, not an error: the chain is a property of hyperref/nameref versions, so
% a future package release must degrade to regular-weight part math — a cosmetic
% loss — instead of refusing to build. Chapter and section keep \PackageError,
% because those targets are in the CLASS and were verified to patch here.
\newif\iflakin@partnum \newif\iflakin@parttitle \newif\iflakin@sparttitle
\def\lakin@partnum#1#2{\iflakin@partnum\else
  \patchcmd{#1}{#2\bfseries \partname}{#2\bfseries\boldmath \partname}
    {\lakin@partnumtrue}{}\fi}
\def\lakin@parttitle#1#2{\iflakin@parttitle\else
  \patchcmd{#1}{#2 \bfseries}{#2 \bfseries\boldmath}{\lakin@parttitletrue}{}\fi}
\def\lakin@sparttitle#1#2{\iflakin@sparttitle\else
  \patchcmd{#1}{#2 \bfseries}{#2 \bfseries\boldmath}{\lakin@sparttitletrue}{}\fi}
\@ifundefined{c@chapter}{\def\lakin@pnum{\Large}\def\lakin@ptit{\huge}}%
                        {\def\lakin@pnum{\huge}\def\lakin@ptit{\Huge}}
\@ifundefined{NR@part}{}{\expandafter\lakin@partnum\expandafter\NR@part\expandafter{\lakin@pnum}%
                         \expandafter\lakin@parttitle\expandafter\NR@part\expandafter{\lakin@ptit}}
\@ifundefined{H@old@part}{}{\expandafter\lakin@partnum\expandafter\H@old@part\expandafter{\lakin@pnum}%
                            \expandafter\lakin@parttitle\expandafter\H@old@part\expandafter{\lakin@ptit}}
\@ifundefined{@part}{}{\expandafter\lakin@partnum\expandafter\@part\expandafter{\lakin@pnum}%
                       \expandafter\lakin@parttitle\expandafter\@part\expandafter{\lakin@ptit}}
\@ifundefined{NR@spart}{}{\expandafter\lakin@sparttitle\expandafter\NR@spart\expandafter{\lakin@ptit}}
\@ifundefined{H@old@spart}{}{\expandafter\lakin@sparttitle\expandafter\H@old@spart\expandafter{\lakin@ptit}}
\@ifundefined{@spart}{}{\expandafter\lakin@sparttitle\expandafter\@spart\expandafter{\lakin@ptit}}
\iflakin@partnum\else\PackageWarning{hebrew-lualatex-pdf}{Part-number math stays
  regular weight: no known \noexpand\part wrapper matched (hyperref/nameref
  version change). Cosmetic only}\fi
\iflakin@parttitle\else\PackageWarning{hebrew-lualatex-pdf}{Part-title math stays
  regular weight: no known \noexpand\part wrapper matched. Cosmetic only}\fi
% v4.12.0: THE PART NUMBER IN THE TOC WAS NOT DIRECTIONALLY SEPARATED FROM THE
% TITLE. report.cls/article.cls compose the part entry as ONE text stream --
% \addcontentsline{toc}{part}{\thepart\hspace{1em}#1} -- while EVERY other level
% goes through \numberline, which puts the number in an \hbox. A box is a bidi
% OBJECT and cannot join a neighbouring run; a bare "IV" can. So in the RTL TOC
% paragraph a Latin number and a Latin-INITIAL title merge into a SINGLE LTR run
% and are laid out left-to-right as a unit: the number lands on the wrong side of
% the row, and a mixed title is scrambled outright. Measured on \part{MCP ...}:
% the printed order was "<tail> V MCP" -- title fragment, number, title fragment.
% The chapter/section rows one line below, with the SAME Latin-initial titles,
% were correct on the same page. That contrast is the whole diagnosis.
% \en{} IS NOT A WORKAROUND, measured both ways on one fixture: it fixes a
% PURE-Latin title (\part{\en{Hooks}} renders correctly) and does NOT fix a MIXED
% one (\part{\en{MCP} <hebrew>} stays scrambled). babel's language switch changes
% font and hyphenation, not the bidi run structure -- so the defect is not
% reachable from the source side at all, which is why it belongs here.
% Boxing the number is the entire fix and is a STRICT NO-OP where the bug never
% bit: a Hebrew-titled document compiles BYTE-IDENTICAL before and after (md5,
% 3 clean runs), because there the neutral glue already sat between an L run and
% an R run. The PDF outline is byte-identical too -- hyperref's \pdfstringdef
% neutralises \mbox, so bookmark strings keep their part numbers.
% \numberline{\thepart} would also fix the direction, but it re-geometries every
% part row in the fleet and DROPS the number from the part bookmark (hyperref
% strips \numberline). \mbox changes nothing that was already right.
% The search pattern deliberately carries no # token: etoolbox cannot match a
% parameter token in a patch pattern (measured -- all three chain members
% reported "not patchable" while the pattern included #1). "\thepart \hspace
% {1em}" is brace-balanced and occurs exactly once in the wrapper body.
\newif\iflakin@parttocnum
\def\lakin@parttocnum#1{\iflakin@parttocnum\else
  \patchcmd{#1}{\thepart \hspace {1em}}{\mbox {\thepart }\hspace {1em}}
    {\lakin@parttocnumtrue}{}\fi}
\@ifundefined{NR@part}{}{\lakin@parttocnum\NR@part}
\@ifundefined{H@old@part}{}{\lakin@parttocnum\H@old@part}
\@ifundefined{@part}{}{\lakin@parttocnum\@part}
\iflakin@parttocnum\else\PackageWarning{hebrew-lualatex-pdf}{Part TOC number is
  not boxed: no known \noexpand\part wrapper matched. A part title that BEGINS
  with Latin will print its number on the wrong side of its TOC row}\fi
\patchcmd{\section}{\bfseries}{\bfseries\boldmath}{}
  {\PackageError{hebrew-lualatex-pdf}{Cannot patch section math weight}{Unsupported document class}}
\patchcmd{\subsection}{\bfseries}{\bfseries\boldmath}{}
  {\PackageError{hebrew-lualatex-pdf}{Cannot patch subsection math weight}{Unsupported document class}}
\patchcmd{\subsubsection}{\bfseries}{\bfseries\boldmath}{}
  {\PackageError{hebrew-lualatex-pdf}{Cannot patch subsubsection math weight}{Unsupported document class}}
\patchcmd{\paragraph}{\bfseries}{\bfseries\boldmath}{}{}
\patchcmd{\subparagraph}{\bfseries}{\bfseries\boldmath}{}{}
\@ifundefined{l@chapter}{}{%
  \patchcmd{\l@chapter}{\bfseries}{\bfseries\boldmath}{}{}}
\makeatother
\setcounter{tocdepth}{\lakittocdepth}
% <hllpdf:measured-toc>
% v4.16.0: Measure actual number labels, not a guessed digit count. Read only
% canonical contentsline records; never execute titles or addtocontents code
% during the measuring pass. Widths are local to this individual list and are
% available before its first row: no per-row drift and no extra aux pass.
\makeatletter
\newread\hllpdf@tocstream
\newdimen\hllpdf@toc@chapter
\newdimen\hllpdf@toc@section
\newdimen\hllpdf@toc@subsection
\newdimen\hllpdf@toc@subsubsection
\newdimen\hllpdf@toc@paragraph
\newdimen\hllpdf@toc@subparagraph
\newdimen\hllpdf@toc@figure
\newdimen\hllpdf@toc@table
\newdimen\hllpdf@tocpage
\newdimen\hllpdf@toc@sectionindent
\newdimen\hllpdf@toc@subsectionindent
\newdimen\hllpdf@toc@subsubsectionindent
\newdimen\hllpdf@toc@paragraphindent
\newdimen\hllpdf@toc@subparagraphindent
\def\hllpdf@tocdefaults{%
  \hllpdf@toc@chapter=2em
  \hllpdf@toc@section=3.2em
  \hllpdf@toc@subsection=3.9em
  \@ifundefined{c@chapter}{%
    \hllpdf@toc@subsubsection=3.2em
    \hllpdf@toc@paragraph=4.1em
    \hllpdf@toc@subparagraph=5em
  }{%
    \hllpdf@toc@subsubsection=4.1em
    \hllpdf@toc@paragraph=5em
    \hllpdf@toc@subparagraph=6em
  }%
  \hllpdf@toc@figure=2.3em
  \hllpdf@toc@table=2.3em
  \hllpdf@tocpage=\@pnumwidth\relax
}
\def\hllpdf@tocindents{%
  % Preserve the established shallow layout when its slots already fit.
  \hllpdf@toc@sectionindent=1.5em
  \@tempdima=\hllpdf@toc@chapter \advance\@tempdima by -2em
  \advance\hllpdf@toc@sectionindent by \@tempdima
  \hllpdf@toc@subsectionindent=\dimexpr
    \hllpdf@toc@sectionindent+\hllpdf@toc@section\relax
  \hllpdf@toc@subsubsectionindent=\dimexpr
    \hllpdf@toc@subsectionindent+\hllpdf@toc@subsection\relax
  \hllpdf@toc@paragraphindent=\dimexpr
    \hllpdf@toc@subsubsectionindent+\hllpdf@toc@subsubsection\relax
  \hllpdf@toc@subparagraphindent=\dimexpr
    \hllpdf@toc@paragraphindent+\hllpdf@toc@paragraph\relax
}
\def\hllpdf@tocfont@chapter{\bfseries\boldmath}
\def\hllpdf@tocfont@part{\large\bfseries\boldmath}
\def\hllpdf@tocentryfont{%
  \ifcsname hllpdf@tocfont@\hllpdf@toclevel\endcsname
    \csname hllpdf@tocfont@\hllpdf@toclevel\endcsname
  \fi}
\long\def\hllpdf@tocdiscard#1\hllpdf@tocstop{}
\def\hllpdf@tocscanrecord{%
  \@ifnextchar\contentsline{\hllpdf@tocscanline}{%
    \@ifnextchar\babel@toc{\hllpdf@tocscanlanguage}{\hllpdf@tocdiscard}}}
\long\def\hllpdf@tocscanlanguage\babel@toc#1#2#3\hllpdf@tocstop{%
  % Babel records are declarations, not arbitrary code. Track the language
  % without executing the original directive or changing the surrounding list.
  \def\hllpdf@toclanguage{#1}}
\long\def\hllpdf@tocscanline\contentsline#1#2#3#4#5\hllpdf@tocstop{%
  \def\hllpdf@toclevel{#1}%
  \ifcsname hllpdf@toc@#1\endcsname
    \hllpdf@tocscantitle#2\hllpdf@tocstop
  \fi
  % The page-number column needs its own measurement (including bold chapter
  % rows). Its extra clearance separates the page label from dot leaders.
  \setbox\@tempboxa=\hbox{\foreignlanguage{\hllpdf@toclanguage}{\hllpdf@tocentryfont #3\kern.3em}}%
  \ifdim\wd\@tempboxa>\hllpdf@tocpage
    \hllpdf@tocpage=\wd\@tempboxa
  \fi}
\def\hllpdf@tocscantitle{%
  \@ifnextchar\numberline{\hllpdf@tocscanlabel}{\hllpdf@tocdiscard}}
\long\def\hllpdf@tocscanlabel\numberline#1#2\hllpdf@tocstop{%
  \setbox\@tempboxa=\hbox{\foreignlanguage{\hllpdf@toclanguage}{\hllpdf@tocentryfont #1\kern.5em}}%
  \ifdim\wd\@tempboxa>\csname hllpdf@toc@\hllpdf@toclevel\endcsname
    \csname hllpdf@toc@\hllpdf@toclevel\endcsname=\wd\@tempboxa
  \fi}
\def\hllpdf@tocread{%
  \ifeof\hllpdf@tocstream\else
    \read\hllpdf@tocstream to\hllpdf@tocrecord
    \expandafter\hllpdf@tocscanrecord\hllpdf@tocrecord\hllpdf@tocstop
    \expandafter\hllpdf@tocread
  \fi}
\def\hllpdf@tocmeasure#1{%
  \makeatletter
  \edef\hllpdf@toclanguage{\languagename}%
  \hllpdf@tocdefaults
  \openin\hllpdf@tocstream=\jobname.#1\relax
  \hllpdf@tocread
  \closein\hllpdf@tocstream
  \hllpdf@tocindents
  \edef\@pnumwidth{\the\hllpdf@tocpage}%
  \@tempdima=\dimexpr\hllpdf@tocpage+1em\relax
  \ifdim\@tempdima>\@tocrmarg\relax
    \edef\@tocrmarg{\the\@tempdima}%
  \fi}
\let\hllpdf@originalstarttoc\@starttoc
\def\@starttoc#1{%
  \begingroup
    \def\hllpdf@tocext{#1}%
    \def\hllpdf@tocwanted{toc}%
    \ifx\hllpdf@tocext\hllpdf@tocwanted \hllpdf@tocmeasure{#1}\else
      \def\hllpdf@tocwanted{lof}%
      \ifx\hllpdf@tocext\hllpdf@tocwanted \hllpdf@tocmeasure{#1}\else
        \def\hllpdf@tocwanted{lot}%
        \ifx\hllpdf@tocext\hllpdf@tocwanted \hllpdf@tocmeasure{#1}\fi
      \fi
    \fi
    \hllpdf@originalstarttoc{#1}%
  \endgroup}
\hllpdf@tocdefaults
\hllpdf@tocindents
\@ifundefined{l@chapter}{}{%
  \patchcmd{\l@chapter}{\setlength\@tempdima{1.5em}}%
    {\setlength\@tempdima{\hllpdf@toc@chapter}}{}%
    {\PackageError{hebrew-lualatex-pdf}{Cannot size chapter TOC numbers}%
       {Unsupported contents layout: the chapter number slot was not found}}}
\renewcommand*\l@section{\@dottedtocline{1}%
  {\hllpdf@toc@sectionindent}{\hllpdf@toc@section}}
\renewcommand*\l@subsection{\@dottedtocline{2}%
  {\hllpdf@toc@subsectionindent}{\hllpdf@toc@subsection}}
\renewcommand*\l@subsubsection{\@dottedtocline{3}%
  {\hllpdf@toc@subsubsectionindent}{\hllpdf@toc@subsubsection}}
\renewcommand*\l@paragraph{\@dottedtocline{4}%
  {\hllpdf@toc@paragraphindent}{\hllpdf@toc@paragraph}}
\renewcommand*\l@subparagraph{\@dottedtocline{5}%
  {\hllpdf@toc@subparagraphindent}{\hllpdf@toc@subparagraph}}
\renewcommand*\l@figure{\@dottedtocline{1}{1.5em}{\hllpdf@toc@figure}}
\renewcommand*\l@table{\@dottedtocline{1}{1.5em}{\hllpdf@toc@table}}
\makeatother
% </hllpdf:measured-toc>
\setcounter{secnumdepth}{\lakitsecnumdepth}
% Public preamble-only caption-numbering policy. The historical within-major
% default is untouched unless this command is called. "continuous" removes the
% reset edge as well as the printed prefix, so Figure 1/2 and Table 1/2 stay
% unique across sections/chapters. Returning to "within-major" restores the
% original class-aware displays and native destination conventions. The policy
% never replaces \figcap/\tabcap or their measured link-placement mechanism.
\makeatletter
\let\lakin@majorfig\thefig
\let\lakin@majortbl\thetbl
\let\lakin@majorHfig\theHfig
\let\lakin@majorHtbl\theHtbl
\newcommand{\lakitcaptionnumbering}[1]{%
  \edef\lakin@numbering{#1}%
  \def\lakin@continuous{continuous}\def\lakin@withinmajor{within-major}%
  \ifx\lakin@numbering\lakin@continuous
    \@ifundefined{c@chapter}{%
      \@removefromreset{fig}{section}\@removefromreset{tbl}{section}%
    }{\@removefromreset{fig}{chapter}\@removefromreset{tbl}{chapter}}%
    \renewcommand{\thefig}{\arabic{fig}}\renewcommand{\thetbl}{\arabic{tbl}}%
    \def\theHfig{\arabic{fig}}\def\theHtbl{\arabic{tbl}}%
  \else\ifx\lakin@numbering\lakin@withinmajor
    \@ifundefined{c@chapter}{%
      \@removefromreset{fig}{section}\@removefromreset{tbl}{section}%
      \@addtoreset{fig}{section}\@addtoreset{tbl}{section}%
    }{%
      \@removefromreset{fig}{chapter}\@removefromreset{tbl}{chapter}%
      \@addtoreset{fig}{chapter}\@addtoreset{tbl}{chapter}}%
    \let\thefig\lakin@majorfig\let\thetbl\lakin@majortbl
    \let\theHfig\lakin@majorHfig\let\theHtbl\lakin@majorHtbl
  \else
    \PackageError{hebrew-english-latex-studio}{Unknown caption-numbering policy '#1'}%
      {Use continuous or within-major before begin{document}.}%
  \fi\fi}
\@onlypreamble\lakitcaptionnumbering
\makeatother
% Architecture readiness sentinel. Optional runtime layers require the complete
% core, so the marker is defined only after every mechanism/baseline definition.
% <hllpdf:architecture-ready>
\def\hllpdfcoreloaded{1}
% </hllpdf:architecture-ready>

% v4.22: registered file resources, not an assertion about embedded subset bytes.
% The log framing avoids long-path wrapping and preserves UTF-8 as hex.
% <hllpdf:font-inventory>
\directlua{
  local function inventory()
    local names = {}
    local unnamed = 0
    for id, item in font.each() do
      local path = item.filename
      if path and not (path == "") then names[path] = true else unnamed = unnamed + 1 end
    end
    local sorted = {}
    for path in pairs(names) do table.insert(sorted,path) end
    table.sort(sorted)
    texio.write_nl("log","HLLFONT|BEGIN|1")
    for index,path in ipairs(sorted) do
      local hex = path:gsub(".",function(c) return string.format(string.char(37) .. "02x",string.byte(c)) end)
      local count = math.ceil(string.len(hex)/48)
      texio.write_nl("log","HLLFONT|FILE|" .. index .. "|" .. count)
      for offset=1,string.len(hex),48 do texio.write_nl("log","HLLFONT|HEX|" .. hex:sub(offset,offset+47)) end
    end
    texio.write_nl("log","HLLFONT|END|" .. rawlen(sorted) .. "|" .. unnamed)
    texio.write_nl("log","")
  end
  luatexbase.add_to_callback("finish_pdffile",inventory,"hllpdf.runtime.fonts")
}
% </hllpdf:font-inventory>
; if omitted, the historical canonical
% defaults below are exact.  They deliberately control only the document shell
% and page layout.  RTL/math semantics stay in this preamble, while visual
% presentation themes (for example clean_theme.tex) load AFTER it.
%
%   \def\lakitdocumentclass{article}
%   \def\lakitclassoptions{11pt,a4paper}
%   \def\lakitgeometryoptions{margin=1.8cm,headheight=14pt}
%   \def\lakitlinespread{1.15}
%   \def\lakitparagraphstyle{block} % letters/briefs; indented for continuous prose
%   \def\lakittocdepth{1}
%   \def\lakitsecnumdepth{2}
%   % ============================================================================
% preamble.tex — tested LuaLaTeX preamble for Hebrew + math + figures PDFs (v4.32.0)
% Battle-tested: compiled an 80-page Hebrew quantum-computing summary AND a
% 197-page Hebrew computer-architecture anthology (TikZ diagrams, code listings,
% colored boxes) — both with 0 missing characters.
%
% Usage in your document:
%   \input{preamble.tex}     % (or paste these lines at the top)
%   \begin{document}  ...  \end{document}
%
% Compile with:  lualatex -interaction=nonstopmode doc.tex
%   Repeat until two consecutive PDFs are byte-identical (scripts/build.sh
%   does this) — a TOC document converges only on pass 3+; after two passes
%   the TOC page numbers are STALE (measured; see SKILL.md step 3).
% Font binaries are external dependencies, pinned by assets/fonts/SHA256SUMS.
% scripts/setup_fonts.sh imports a verified local installation or prior skill.
%
% PUBLIC DOCUMENT-SHELL / LAYOUT CONTRACT (introduced v4.3.1; architecture-owned v4.7.1). These hooks may be
% defined BEFORE \input{preamble.tex}; if omitted, the historical canonical
% defaults below are exact.  They deliberately control only the document shell
% and page layout.  RTL/math semantics stay in this preamble, while visual
% presentation themes (for example clean_theme.tex) load AFTER it.
%
%   \def\lakitdocumentclass{article}
%   \def\lakitclassoptions{11pt,a4paper}
%   \def\lakitgeometryoptions{margin=1.8cm,headheight=14pt}
%   \def\lakitlinespread{1.15}
%   \def\lakitparagraphstyle{block} % letters/briefs; indented for continuous prose
%   \def\lakittocdepth{1}
%   \def\lakitsecnumdepth{2}
%   \input{preamble.tex}
%
% Supported class contract: report or article.  Alternate paper sizes must also
% be declared to build.sh through EXPECTED_PAPER; the PDF gate never infers them.
% Do NOT use these hooks as a theme mechanism: theme overlays must not own page
% geometry, class selection, or direction semantics. A declared presentation-policy
% overlay may intentionally change numbering visibility after the preamble.
% ============================================================================
% Runtime architecture guard: generated from scripts/architecture_contract.py.
% The readiness sentinel remains delayed until the END of this file.
% <hllpdf:architecture-guard>
\ifdefined\hllpdfcoreloaded
  \errmessage{hllpdf: preamble.tex loaded more than once}
  \expandafter\endinput
\fi
% </hllpdf:architecture-guard>

\ifdefined\lakitdocumentclass\else\def\lakitdocumentclass{report}\fi
\ifdefined\lakitclassoptions\else\def\lakitclassoptions{12pt,a4paper}\fi
\ifdefined\lakitgeometryoptions\else\def\lakitgeometryoptions{margin=2.3cm,headheight=14pt}\fi
\ifdefined\lakitlinespread\else\def\lakitlinespread{1.3}\fi
\ifdefined\lakittocdepth\else\def\lakittocdepth{1}\fi
\ifdefined\lakitsecnumdepth\else\def\lakitsecnumdepth{2}\fi
% Paragraph organization belongs to the document layout, independently of its
% language and theme. Preserve historical class dimensions unless explicitly
% selecting block paragraphs; never suppress indentation with an everypar hack.
\ifdefined\lakitparagraphstyle\else\def\lakitparagraphstyle{indented}\fi
\edef\lakitparagraphstyle{\lakitparagraphstyle}
\def\hllpdfparagraphblock{block}
\def\hllpdfparagraphindented{indented}
\ifx\lakitparagraphstyle\hllpdfparagraphblock\else
  \ifx\lakitparagraphstyle\hllpdfparagraphindented\else
    \errmessage{hllpdf: paragraph style must be block or indented}
  \fi
\fi
% v4.3.1: the public class hook has a deliberately narrow support contract.
% Before this guard, \def\lakitdocumentclass{book} silently compiled even though
% every heading/TOC regression test targets report/article; an unsupported shell
% could therefore look "supported" until a subtle class-specific defect shipped.
% Fail before \documentclass rather than letting an undocumented combination enter
% the mechanism core. Plain-TeX comparison is used because LaTeX packages are not
% available yet at this point.
\edef\lakitvalidatedclass{\lakitdocumentclass}
\def\lakitreportclass{report}
\def\lakitarticleclass{article}
\ifx\lakitvalidatedclass\lakitreportclass
\else\ifx\lakitvalidatedclass\lakitarticleclass
\else
  \errmessage{hllpdf: unsupported class \lakitdocumentclass; use report or article}
\fi\fi
\let\lakitvalidatedclass\relax
\let\lakitreportclass\relax
\let\lakitarticleclass\relax
% Language intent is independent of document class, theme and content domain.
% The two dominant-language profiles differ in authoring intent, not direction.
% A balanced bilingual document must explicitly choose a navigation/base locale.
% <hllpdf:language-profile>
\providecommand{\lakitlanguageprofile}{hebrew}
\edef\lakitlanguageprofile{\lakitlanguageprofile}
\def\hllpdfhebrew{hebrew}\def\hllpdfenglish{english}
\def\hllpdfhebrewmixed{hebrew-with-english}
\def\hllpdfenglishmixed{english-with-hebrew}
\def\hllpdfbilingual{bilingual}
\let\hllpdfexpectedlanguage\relax
\ifx\lakitlanguageprofile\hllpdfhebrew
  \let\hllpdfexpectedlanguage\hllpdfhebrew
\else\ifx\lakitlanguageprofile\hllpdfhebrewmixed
  \let\hllpdfexpectedlanguage\hllpdfhebrew
\else\ifx\lakitlanguageprofile\hllpdfenglish
  \let\hllpdfexpectedlanguage\hllpdfenglish
\else\ifx\lakitlanguageprofile\hllpdfenglishmixed
  \let\hllpdfexpectedlanguage\hllpdfenglish
\else\ifx\lakitlanguageprofile\hllpdfbilingual
  \ifdefined\lakitmainlanguage\else
    \errmessage{hllpdf: bilingual profile requires explicit lakitmainlanguage}
  \fi
\else
  \errmessage{hllpdf: unsupported language profile \lakitlanguageprofile}
\fi\fi\fi\fi\fi
\ifdefined\lakitmainlanguage
  \edef\lakitmainlanguage{\lakitmainlanguage}
\else
  \let\lakitmainlanguage\hllpdfexpectedlanguage
\fi
\ifx\lakitmainlanguage\hllpdfhebrew\else
  \ifx\lakitmainlanguage\hllpdfenglish\else
    \errmessage{hllpdf: main language must be hebrew or english}
  \fi
\fi
\ifx\hllpdfexpectedlanguage\relax\else
  \ifx\lakitmainlanguage\hllpdfexpectedlanguage\else
    \errmessage{hllpdf: main language conflicts with language profile}
  \fi
\fi
\typeout{HLLLANG|\lakitlanguageprofile|\lakitmainlanguage}
% </hllpdf:language-profile>
\documentclass[\lakitclassoptions]{\lakitdocumentclass}
\usepackage[no-math]{fontspec}        % (1) ★ no-math — THE critical line (see recipe.md)
% layout=lists asks Babel to adapt list parshape when paragraph and page
% directions differ. Without it, minority-language list bodies lose their
% indentation and labels hang outside the text block, despite a clean log.
% The upstream mechanism preserves ordinary same-direction list geometry.
\usepackage[bidi=basic,layout=lists]{babel}
\ifx\lakitmainlanguage\hllpdfenglish
  \babelprovide[import]{hebrew}
  \babelprovide[import,main]{english}
\else
  \babelprovide[import,main]{hebrew}
  \babelprovide[import]{english}
\fi
% v4.22: bind canonical families and fallbacks to the selected installation.
% Lookup by basename or family alone is not proof of the file actually used.
% <hllpdf:font-binding>
\directlua{
  local dir = os.getenv("HEBREW_LATEX_FONT_DIR")
  if not dir or dir == "" then
    local data = os.getenv("XDG_DATA_HOME")
    if not data or data == "" or not (data:sub(1,1) == "/") then
      data = (os.getenv("HOME") or "") .. "/.local/share"
    end
    dir = data .. "/fonts/hebrew-lualatex-pdf"
  end
  if not (dir:sub(1,1) == "/") then dir = lfs.currentdir() .. "/" .. dir end
  if dir:find("[",1,true) or dir:find("]",1,true) or dir:find(string.char(0),1,true)
     or dir:find(string.char(10),1,true) or dir:find(string.char(13),1,true) then
    tex.error("hllpdf: unsupported character in font directory")
  end
  if not (lfs.attributes(dir,"mode") == "directory") then
    tex.error("hllpdf: font directory is missing")
  end
  hllpdf_font_dir = dir:gsub("/+$", "") .. "/"
  tex.sprint(string.char(92) .. "def" .. string.char(92) .. "hllpdffontdir{")
  tex.sprint(-2,hllpdf_font_dir)
  tex.sprint("}")
  luaotfload.add_fallback("hebfb",{"[" .. hllpdf_font_dir .. "DejaVuSerif.ttf]:mode=node;"})
  luaotfload.add_fallback("hebfbbf",{"[" .. hllpdf_font_dir .. "DejaVuSerif-Bold.ttf]:mode=node;"})
}
% </hllpdf:font-binding>
% v3.22: Frank Ruhl Libre ships NO italic face. Instead of silently-upright
% \emph (v<=3.20) or bold-as-emphasis (v3.21), synthesize a slanted face with
% FakeSlant — real, VISIBLE italic emphasis with correct \emph nesting
% semantics (inner \emph toggles back upright). Verified empirically at
% 300dpi: slant clearly visible, 0 errors, 0 missing glyphs.
\babelfont{rm}[Path={\hllpdffontdir},RawFeature={fallback=hebfb},
  BoldFont={FrankRuhlLibre-Bold.ttf},
  BoldFeatures={RawFeature={fallback=hebfbbf}},   % v3.56 — see hebfbbf above
  ItalicFont={FrankRuhlLibre-Regular.ttf}, ItalicFeatures={FakeSlant=0.22},
  % Bold-italic uses the weight-matched fallback and a synthetic slant.
  BoldItalicFont={FrankRuhlLibre-Bold.ttf}, BoldItalicFeatures={RawFeature={fallback=hebfbbf},FakeSlant=0.22}
  ]{FrankRuhlLibre-Regular.ttf}
% Sans fallbacks preserve family and weight for the few missing codepoints.
\directlua{luaotfload.add_fallback("hebfbsans",{"[" .. hllpdf_font_dir .. "DejaVuSans.ttf]:mode=node;"})}
\directlua{luaotfload.add_fallback("hebfbsansbf",{"[" .. hllpdf_font_dir .. "DejaVuSans-Bold.ttf]:mode=node;"})}
\babelfont{sf}[Path={\hllpdffontdir},RawFeature={fallback=hebfbsans},
  BoldFont={Heebo-Bold.ttf},
  BoldFeatures={RawFeature={fallback=hebfbsansbf}}]{Heebo-Regular.ttf}
% v3.24: Scale=MatchLowercase — DejaVu Sans Mono has an outsized x-height, so
% unscaled inline code visually SHOUTS next to the delicate Frank Ruhl body
% (measured on a 197-page build: scaled, a code atom sits in the line at the
% body's weight, and the document tightened by ~5%). Applies to listings too.
\babelfont{tt}[Path={\hllpdffontdir},RawFeature={fallback=hebfb},Scale=MatchLowercase,
  BoldFont={DejaVuSansMono.ttf},
  BoldFeatures={RawFeature={fallback=hebfb},FakeBold=1.5},
  ItalicFont={DejaVuSansMono.ttf},
  ItalicFeatures={RawFeature={fallback=hebfb},FakeSlant=0.2},
  BoldItalicFont={DejaVuSansMono.ttf},
  BoldItalicFeatures={RawFeature={fallback=hebfb},FakeBold=1.5,FakeSlant=0.2}
  ]{DejaVuSansMono.ttf}

% Explicit multilingual/English profiles use real Latin faces and italics.
% The legacy/default Hebrew profile retains its established measured typography.
% Latin Modern is provided by the TeX distribution, not the pinned font bundle;
% recorder/runtime evidence reports it as an additional toolchain resource.
\ifx\lakitlanguageprofile\hllpdfhebrew\else
  \babelfont[english]{rm}[Ligatures=TeX,RawFeature={fallback=hebfb},
    BoldFeatures={RawFeature={fallback=hebfbbf}},
    BoldItalicFeatures={RawFeature={fallback=hebfbbf}}]{Latin Modern Roman}
  \babelfont[english]{sf}[Ligatures=TeX,RawFeature={fallback=hebfbsans},
    BoldFeatures={RawFeature={fallback=hebfbsansbf}},
    BoldItalicFeatures={RawFeature={fallback=hebfbsansbf}}]{Latin Modern Sans}
\fi

\usepackage{amsmath,amssymb,mathtools}
% amsthm supplies theorem/proof machinery. Environment names remain
% document-specific; the generic proof label is localized below.
\usepackage{amsthm}
% Proof names follow the surrounding prose, including return from an island.
\addto\extrashebrew{\renewcommand{\proofname}{הוכחה}}
\addto\extrasenglish{\renewcommand{\proofname}{Proof}}
\ifx\lakitmainlanguage\hllpdfenglish
  \renewcommand{\proofname}{Proof}
\else
  \renewcommand{\proofname}{הוכחה}
\fi
% v3.24.1: equation tags under babel-bidi. A tag re-emitted by mathtools'
% showonlyrefs inside align inherits the RTL base and its parens MIRROR —
% ")2.1(" instead of "(2.1)" (plain align/equation tags are unaffected;
% verified on a minimal case + a 198-page build). LTR-isolate at the tag-BOX
% level: mathtools' show/hide + numbering logic is untouched. Do NOT override
% \tagform@ instead — that clobbers the showonlyrefs mechanism (measured:
% every tag shows, numbering breaks).
% v3.28.1: the v3.24.1 wrap must apply ONLY to tags of display environments.
% Inline \eqref in Hebrew PROSE also goes through \maketag@@@, and there the
% english-wrap MIRRORS the parens — ")1(" — while the unwrapped original is
% correct (proven on a 2x2 isolation matrix: patch x showonlyrefs; regressed
% every prose \eqref since v3.24.1, caught by the render gate's R4). The two
% call sites are indistinguishable by \ifmmode (tag boxes are text mode), so
% the display environments themselves raise a flag; \eqref in prose sees it
% false and gets the original untouched \maketag@@@.
\usepackage{etoolbox}

% babel bidi=basic (including upstream 26.9) assumes every automatic
% discretionary has a one-glyph replacement. LuaTeX can legally produce one
% with replace=nil after a hyphenated Latin value; Babel then indexes nil and the
% whole build dies. Do the narrow upstream-compatible correction: while Babel
% classifies the list, present only that malformed disc as an ordinary
% discretionary, then restore its subtype. Babel still reorders the ENTIRE box;
% there is no pcall and no successful fallback to an unreordered node list.
% If Babel throws for any other reason, the build remains fatal. The synthetic
% nil-disc probe and the real table gallery in test_all.sh pin both layers.
\AtBeginDocument{\directlua{
  if Babel and Babel.bidi and not Babel.heb_nil_disc_compat then
    Babel.heb_nil_disc_compat = true
    local original = Babel.bidi
    local disc_id = node.id('disc')
    Babel.bidi = function(head, ispar, hdir)
      local hidden = {}
      for item in node.traverse_id(disc_id, head) do
        if item.subtype == 2 and (not item.replace or not item.replace.char) then
          table.insert(hidden, item)
          item.subtype = 0
        end
      end
      local result = original(head, ispar, hdir)
      for _, item in ipairs(hidden) do item.subtype = 2 end
      return result
    end
  end
}}
\makeatletter
\newif\ifHEBdisptag
\let\ORIGmaketag@@@\maketag@@@
\def\maketag@@@#1{\ifHEBdisptag\ORIGmaketag@@@{\foreignlanguage{english}{#1}}\else\ORIGmaketag@@@{#1}\fi}
\@for\HEB@env:=equation,equation*,align,align*,gather,gather*,multline,multline*,flalign,flalign*,alignat,alignat*,eqnarray,eqnarray*\do{%
  \expandafter\AtBeginEnvironment\expandafter{\HEB@env}{\HEBdisptagtrue}}
\makeatother
% English page + Hebrew prose: Babel's equation-only tag flip uses a primitive
% leqno path that shifts the formula to the left edge and the tag into the body
% of an English page. Align's full-width alignment is unaffected. Give only this
% display a physical LTR container and native amsmath eqno/leqno. The engine then
% owns real tag width, collision handling and labels; no proxy tag box hides ink
% from its fit calculation. Babel's saved \bbl@eqnodir still restores Hebrew text
% inside math, and the prose locale/math direction are never switched.
% Close the prior RTL paragraph before normalization. After the environment's
% group restores its directions, close TeX's empty resumed display paragraph and
% use native LaTeX continuation handling: same-paragraph prose has no indent,
% while an authored blank line starts an ordinary indented RTL paragraph.
% Register after Babel's setup; remain local to equation(*) in a main-English
% document's actual RTL prose. Native Hebrew, ordinary English and explicit
% fleqn retain their original paths. No page/display dimensions are assigned.
% <hllpdf:minority-rtl-equation-geometry>
\makeatletter
\newcommand{\HEB@minorityprosecontinuation}{\par\@doendpe}
\newcommand{\HEB@minorityequationgeometry}{%
  \ifx\lakitmainlanguage\hllpdfenglish
    \ifnum\pardirection=\@ne
      \if@fleqn\else
        \par\aftergroup\HEB@minorityprosecontinuation
        \iftagsleft@\let\veqno\@@leqno
        \else\let\veqno\@@eqno\fi
        \bodydir TLT\pardir TLT\textdir TLT\relax
      \fi
    \fi
  \fi}
\AtBeginDocument{%
  \AddToHook{env/equation/begin}{\HEB@minorityequationgeometry}%
  \AddToHook{env/equation*/begin}{\HEB@minorityequationgeometry}}
\makeatother
% </hllpdf:minority-rtl-equation-geometry>
% v3.63: \shortintertext rendered Hebrew MIRROR-REVERSED and flush LEFT — an entire
% commentary line inside an align printed backwards. babel's bidi=basic patches
% \intertext@ only (luababel.def), and mathtools' \shortintertext uses the identical
% \vbox{\normalbaselines … \noindent\ignorespaces#1\par} skeleton but is never
% patched, so its paragraph keeps pardir/textdir LTR: no bidi reordering (logical
% order on the page) and flush-left. Measured against \intertext with the same
% Applying babel's own replacement to \shortintertext@ makes the two identical.
% \h{}/\mbox fix the letter order but leave the line flush LEFT, so they are not
% substitutes. Guarded + \AtBeginDocument because mathtools \let-copies at load time.
% NB \ExplSyntaxOn is REQUIRED: mathtools' internals carry expl3 names, and without
% the expl3 catcodes the `_` is read as a subscript and the patch detonates far away
% with "! Illegal unit of measure (pt inserted)" at the \end{align} (measured).
\makeatletter
\AtBeginDocument{%
  \@ifundefined{shortintertext@}{}{%
    \ExplSyntaxOn
    \bbl@sreplace\shortintertext@{\normalbaselines}%
      {\normalbaselines\ifx\bbl@eqnodir\relax\else\bbl@pardir\@ne\bbl@eqnodir\fi}%
    \ExplSyntaxOff}}
\makeatother
% geometry inherits the canonical A4 sheet from \documentclass; build.sh then
% verifies every rendered MediaBox and CropBox, so a silent Letter fallback is
% impossible to publish as a clean default build.
\usepackage[\lakitgeometryoptions]{geometry}
\linespread{\lakitlinespread}       % layout hook; canonical default remains 1.3
\ifx\lakitparagraphstyle\hllpdfparagraphblock
  \setlength{\parindent}{0pt}
  \setlength{\parskip}{6pt plus 1pt minus 1pt}
\fi
% These dimensions remain ordinary scoped TeX lengths: authors can customize
% them after core input. Language switches inherit them; lists keep their own
% indentation and narrow reference layouts establish their local paragraph gaps.
\typeout{HLLPAR|\lakitparagraphstyle}
\usepackage{xcolor}
\usepackage{mdframed}
\usepackage{listings}                 % code listings (assembly / C / pseudo-code)
\usepackage{tikz}
\usetikzlibrary{arrows.meta,positioning,calc,shapes.geometric,fit,backgrounds}
\usepackage{pgfplots}\pgfplotsset{compat=1.18}
\usepackage{quantikz}                 % quantum circuits (remove if unused)
\usepackage{booktabs}
\usepackage{array}
\usepackage{float}                    % [H] = place table exactly here (reference doc: no drifting)
\usepackage{enumitem}
\usepackage{needspace}

% ============================================================================
% ★★ v3.4 — THE DIRECTION-ISOLATION TRIO (the heart of mixing Hebrew + Latin)
% Mental model:
%   \h{...}    Hebrew (RTL) island      — Hebrew text inside an LTR context
%                                          (TikZ nodes, the english `tikzpic`).
%   \en{...}   English (LTR) isolate    — Latin text / parentheses inside a
%                                          Hebrew (RTL) context. KEEPS PARENS
%                                          AND ORDER CORRECT.
%   \code{...} \en + monospace          — code inside Hebrew / inside a frametitle.
% CRITICAL: never wrap Latin in \h{} (it is for HEBREW). Inside a TikZ node,
% Latin in \h{} renders REVERSED (Datapath->htapataD, CPU->UPC) and parens
% mirror. Isolate every Latin/paren chunk with \en{} or \code{}. See
% references/figures-boxes-listings.md.
% ---- TEXT INSIDE MATH (v3.8) ----------------------------------------------
% Put ANY text that sits inside math in a \text{}. The ONLY thing that is
% invalid in math is a bare \h/\en (= \foreignlanguage) used DIRECTLY, outside
% \text -> "\rmfamily invalid in math mode". Inside \text:
%   Hebrew label           -> \text{...}            (bare Hebrew renders fine;
%                              \h is optional — BUT NOT ALWAYS HARMLESS (v3.67):
%                              if the \h group OPENS with a nested island
%                              ($..$, \en{}, \code{}, \texttt{}, \mbox{}) the
%                              WHOLE group loses bidi and the Hebrew prints in
%                              LOGICAL order = gibberish. Measured:
%                                \text{\h{$f$ קבועה}}  -> reads "העובק"  BROKEN
%                                \text{\h{קבועה $f$}}  -> correct
%                                \text{$f$ קבועה}      -> correct (no \h)
%                              Hebrew first inside \h, or drop the \h. Flagged
%                              H-OPENS-WITH-LTR-GROUP.)
%   plain Latin word        -> \text{...}            (no marker needed)
%   Latin WITH ( ) or [ ]   -> \text{\en{...}}       (\en forces LTR; a plain
%                              \text mirrors the parens in RTL -> "Spec (TNR(" )
%   f(Hebrew/text arg)      -> \text{f}(\text{...})  -- keep the parens at MATH
%                              level, NOT \text{f(arg)} (parens inside the box mirror)
%   code in math            -> \mathtt{...}  (or \texttt{}, math-guarded below)
% NEVER write \en{...} or \h{...} straight into $...$ / \[..\] (outside \text);
% always via \text{}. (\mbox{...} also works but does not scale in scripts.)
% ============================================================================
\newcommand{\h}[1]{\foreignlanguage{hebrew}{#1}}
% v3.42.1 — \en islands do NOT hyphenate (auto OR at explicit hyphens). English hyphenation patterns were
% tearing identifiers and terms ACROSS THE RTL SEAM ("IO Regis-/ters",
% "In-/put", "Band-/width", "Per-/formance", "Multi-/ple" — all shipped on
% real pages; probe-proven: "ALUS-/rcB", "Mem-/ToReg"). An island is a label
% or a term, not English prose — a hyphen tear there is never right. A/B on
% a full 197pp lakit: identical overfull sets, same page count, ~11
% paragraphs healed, one mild underfull (badness 3525) that replaced a
% mutilated "Regis-/ters". Long English PROSE belongs in the englishprose
% environment below, which restores explicit-hyphen breaks while its paragraph
% is still active. A raw otherlanguage block inherits the global prohibition.
% (\l@nohyphenation needs \csname:
% '@' is not a letter here.)
% v3.43.2 — \exhyphenpenalty moved OUT of \en to a GLOBAL assignment below.
% Render-proven: a group-local \exhyphenpenalty is DEAD CODE — TeX reads the
% line-breaking parameters at \par time, after the \en group has closed, so
% \en{write-through} / plain not-taken still tore at the "-" ("write-/
% through" shipped on a real exam page). \language is different: it is
% recorded in the node list as text is added, so the in-group nohyphenation
% assignment above DOES work (and stays). The global penalty was A/B-measured
% on a 197pp lakit: identical overfull set (25), identical page count, the
% not-/taken tear healed. Hebrew is unaffected in practice — Hebrew compounds
% use maqaf (־), not ASCII "-".
\newcommand{\en}[1]{\foreignlanguage{english}{\language=\csname l@nohyphenation\endcsname #1}}
% Babel runs extras locally for inline islands and paragraph language changes.
% An English main document must not inherit the Hebrew no-hyphen paragraph rule.
\addto\extrashebrew{\exhyphenpenalty=10000\relax}
\addto\extrasenglish{\exhyphenpenalty=50\relax}
\ifx\lakitmainlanguage\hllpdfenglish
  \exhyphenpenalty=50\relax
\else
  \exhyphenpenalty=10000\relax
\fi
% Full English paragraphs need ordinary explicit-hyphen breaks. Use this
% environment instead of a raw otherlanguage block; the paragraph ends while
% the local penalty is still active.
\newenvironment{englishprose}
  {\par\begin{otherlanguage}{english}\exhyphenpenalty=50\relax}
  {\par\end{otherlanguage}}
\newenvironment{hebrewprose}
  {\par\begin{otherlanguage}{hebrew}\exhyphenpenalty=10000\relax}
  {\par\end{otherlanguage}}
% Use for shell text whose locale must not depend on the last body paragraph.
\DeclareRobustCommand{\doclang}[1]{\foreignlanguage{\lakitmainlanguage}{#1}}

% ============================================================================
% ★★ — fixes for defects that are INVISIBLE in the log and in the checks
% and only surface on a VISUAL RENDER of the PDF. (Each verified empirically.)
% ----------------------------------------------------------------------------
% v3.22: \emph is now REAL slanted emphasis (FakeSlant italic face defined on
% the \babelfont{rm} line above), replacing the v3.21 \emph->\textbf hack: it
% preserves \emph's toggle semantics and keeps bold free for its own role.
% Prefer bold-as-emphasis for your doc? one line: \renewcommand{\emph}[1]{\textbf{#1}}
% page number: force the digits LTR so they never reverse (18 -> 81) in the RTL
%   footer. That reversal is a NON-DETERMINISTIC bidi glitch (observed once in
%   ~80 pp); this LuaLaTeX build is not byte-reproducible, so the defensive
%   LTR wrap is the only reliable guard. Latin digits are always LTR anyway.
\renewcommand{\thepage}{\foreignlanguage{english}{\arabic{page}}}
% v3.70: \pagenumbering \gdef's \thepage and silently DESTROYS the wrap above —
% after \pagenumbering{arabic} the folio guard is gone for the whole rest of the
% document. Two collections discovered this independently and hand-restore the
% wrap in main.tex with the same comment; that knowledge now lives here instead.
% The wrapper re-applies the LTR wrap after the kernel call, for any style.
\makeatletter
\let\lakin@orig@pagenumbering\pagenumbering
\renewcommand{\pagenumbering}[1]{%
  \lakin@orig@pagenumbering{#1}%
  \renewcommand{\thepage}{\foreignlanguage{english}{\csname @#1\endcsname\c@page}}}
\makeatother
% emergencystretch: Hebrew does not hyphenate, so an unbreakable LTR island
%   (\en{...} / an inline math run) landing at the margin can leave a small
%   (~2-13pt) overfull with no legal breakpoint. This rescues ONLY those
%   inline-prose tie-breaks; it is a ceiling used only on offending lines and
%   does NOT touch display-math / listing / table overflow (those still surface
%   in the log for the content-level overflow pass). Verified: baseline overfull
%   -> 0 with this set, with no change to pagination or underfull.
\setlength{\emergencystretch}{3em}
% v3.40: inline math (almost) never breaks across lines. TeX's default
%   (binop 700 / rel 500) lets an inline run split at an operator —
%   $(1.f)\cdot 2^{e}$ shipped as "(1.f)·" ending one line and "2^e" opening
%   the next (reader-reported). In LTR that split is merely ugly; under an
%   RTL base the carried operator lands on the visually WRONG side of the
%   break and the formula reads backwards.
%   WHY 10000 (v3.45.0 — the display-first resolution). Inline math is
%   FORBIDDEN from breaking, so a torn island under an RTL base is
%   structurally impossible. The v3.40 cycle used 9999 ("break only as a last
%   resort") because a bare 10000 left ~30-39 WIDE inline equations
%   overflowing (worst 90.7pt) — but 9999's "last resort" fires OFTEN in
%   tight, non-hyphenating Hebrew: a reader found torn islands across the
%   fleet (~150 candidates; e.g. "P(y |" ending one line and "x_0)" opening
%   the next, read backwards). The step the 9999 cycle missed: every one of
%   those ~30 overflows is a wide equation that BELONGS in display. Convert
%   them to display FIRST, then 10000 gives 0 tears AND 0 overflows (measured,
%   12/12 fleet). Doctrine: inline math is ATOMIC; a wide equation goes in a
%   display, never inline. See content-style ("inline math is atomic").
\binoppenalty=10000 \relpenalty=10000
% v3.38: uniform sentence spacing. TeX's default (\nonfrenchspacing, sfcode
% 3000) is the OLD Anglo-American wide-space-after-period convention — Hebrew
% typography never had it, and modern Latin practice (Chicago et al.) dropped
% it too. Measured on real pages: post-period gaps up to 1.9-2.1x the line's
% mean word gap — the dominant "rivers" source in justified Hebrew prose.
% Survives babel language switches (probed: sfcode stays french inside \en{}
% islands and otherlanguage blocks).
\frenchspacing
% v3.38: widow/orphan protection. LaTeX's default penalties (150/150/50)
% allow a lone line at a page top/bottom almost freely. 10000 = forbid;
% this is SAFE here precisely because the document is \raggedbottom (the
% demoted line just moves on, no vertical stretch fight) — under flushbottom
% these values would cause underfull pages. raggedbottom is the report-class
% default; made explicit so a future class/option change cannot silently
% flip it and turn the 10000s harmful.
\raggedbottom
\widowpenalty=10000 \clubpenalty=10000 \displaywidowpenalty=10000
% em-dash: --- does NOT become — under Frank Ruhl (the font lacks tlig; verified
%   the width is identical with/without Ligatures=TeX, while the mechanism works
%   on Latin Modern). TYPE THE CHARACTER — (U+2014) DIRECTLY; likewise – (U+2013).
%   check_bidi_figures.py flags a --/--- that is adjacent to Hebrew.
% ============================================================================

% ---- STRUCTURAL DIRECTION-SAFETY for \texttt (v3.7) -------------------------
% A parenthesis, bracket, or $ INSIDE \texttt{} mirrors/migrates in RTL prose
% (MIPS 0($s0)->"(s0$)0", $rs->"rs$", arr[i]->"[arr[i") because \texttt sets the
% font but NOT the direction. Redefining \texttt to LTR-isolate its content fixes
% the WHOLE class at once (parens, brackets, $) so plain \texttt{lw $t0,0($s0)}
% is safe — no need to remember \code{}. Guarded by \ifmmode because
% \foreignlanguage is invalid in math (there, plain monospace is already LTR).
% NB (corrected v3.56 — the old note here said "stays tofu", which is WRONG):
% \texttt{Hebrew} is a FATAL build failure. The tt family is DejaVu Sans Mono
% (zero Hebrew) and the tt fallback hebfb is DejaVuSerif (also zero Hebrew), so
% the subsetter gets an empty glyph set and LuaTeX dies with
%   "! error: (file .../DejaVuSansMono.ttf) (type 2): there are no glyphs in the
%    subset"  -> no output PDF.
% The message names a FONT FILE, never Hebrew/\texttt/the source line. And it is
% non-local: add an unrelated Latin \code{} elsewhere and the fatal degrades into
% a SILENT DELETION of the Hebrew word. Put Hebrew in \h{}/normal text, never in
% \texttt{}/\code{}; flagged HEBREW-IN-MONO. To disable, delete this block.
\let\origtexttt\texttt
\protected\long\def\texttt#1{\ifmmode\origtexttt{#1}\else\foreignlanguage{english}{\origtexttt{#1}}\fi}
% v4.8.0: defer the direction test until TeX has inserted an alignment cell's
% math template. Unprotected expansion here tested \ifmmode too early at the
% FIRST token of aligned/array/cases/matrix cells. Engine-level \protected also
% preserves the original \texttt implementation without the self-recursion
% observed when retrofitting \DeclareRobustCommand to its saved alias.
\protected\long\def\code#1{\ifmmode\origtexttt{#1}\else\foreignlanguage{english}{\origtexttt{#1}}\fi}

% ---- \codeline — display-code line (v3.24): the code analogue of display math.
% Reader-friction rule (content-style.md §5b): inline \code{} is for a SINGLE
% ATOM only (one register/variable/mnemonic, a tight no-space expression).
% A multi-token expression — an instruction with operands, a spaced assignment,
% a register-binding list — zigzags the eye RTL<->LTR mid-sentence and, being
% an unbreakable island, line-breaks badly. Give it its own centered line with
% display-math vertical rhythm. \centerline = one box => the expression never
% wraps; \nopagebreak keeps it glued to its introducing sentence.
%   ...the Hebrew sentence ends with a colon:
%   \codeline{lw \$t1, -4(\$t0)}
%   and the prose continues (escape dollars as \$).
% v3.66 (first line of the body): flush a PENDING \item label before the \par.
% LaTeX defers the label box (\@labels) until the next \everypar, and an \hbox in
% vertical mode never fires \everypar — so `\item \codeline{...}` printed the code
% at y=69.1 and the item's own "1." BELOW it at y=103.6 (the number appears after
% its content, attached to nothing). \if@inlabel\leavevmode\fi is CONDITIONAL on
% purpose: an unconditional \leavevmode would open an empty paragraph line when
% \codeline follows ordinary prose. Verified both ways — in a list the label now
% sits on the code's line at its normal x, and after prose the geometry is
% byte-identical to before.
\makeatletter
\newcommand{\codeline}[1]{%
  \if@inlabel\leavevmode\fi
  \par\nopagebreak
  \vspace{\abovedisplayskip}\nopagebreak
  % v3.62: was \centerline{...}, i.e. \hbox to\hsize{\hss #1 \hss}. \hss has
  % INFINITE SHRINK, so TeX could always "fit" the box and NEVER reported an
  % Overfull — the same \hss mechanism as BLD-27's \item label. Measured: a wide
  % \codeline in a defbox painted its first token ON the box's left border bar and
  % sat 9.5pt outside the frame interior on both sides, with 0 Overfull, 0 Underfull
  % and all four gates PASS; in bare prose it reached 21.0pt into BOTH margins.
  % \hfil has no shrink, so the overflow is reported and build.sh's ovf_vis/ovf_big
  % counters see it; \linewidth (not \hsize) makes it box-aware, so a \codeline
  % inside a colored box centres on the box interior. Verified: a FITTING \codeline
  % keeps byte-identical glyph positions, and the box-crossing case now emits
  % "Overfull \hbox (19.03pt too wide)". A wide \codeline should be split in two.
  \hbox to\linewidth{\hfil\code{#1}\hfil}%
  \nopagebreak\vspace{\belowdisplayskip}}
\makeatother

% ---- a thin decorative rule helper ----
\newcommand{\sep}{\par\vspace{2pt}\noindent\textcolor{cRule}{\rule{\linewidth}{0.6pt}}\par\vspace{4pt}}

% ---- colors ----
\definecolor{cBlue}{HTML}{1F4E79}   \definecolor{cBlueBg}{HTML}{EAF1F8}
\definecolor{cGreen}{HTML}{2E6B4F}  \definecolor{cGreenBg}{HTML}{EAF4EE}
\definecolor{cAmber}{HTML}{9C6A00}  \definecolor{cAmberBg}{HTML}{FBF3E2}
\definecolor{cPurple}{HTML}{5B2C6F} \definecolor{cPurpleBg}{HTML}{F3ECF7}
\definecolor{cTeal}{HTML}{0E6E6E}   \definecolor{cTealBg}{HTML}{E6F4F4}
\definecolor{cRed}{HTML}{8E2B2B}    \definecolor{cRedBg}{HTML}{FBEAEA}
\definecolor{cGray}{HTML}{555555}   \definecolor{cRule}{HTML}{1F4E79}

% ---- math operators ----
\DeclareMathOperator{\Tr}{Tr}
\DeclareMathOperator{\rank}{rank}
\DeclareMathOperator{\diag}{diag}
\DeclareMathOperator{\sgn}{sgn}

% ---- bra-ket (quantikz defines \ket/\bra/\braket; add the rest safely) ----
\providecommand{\ket}[1]{\left|#1\right\rangle}
\providecommand{\bra}[1]{\left\langle#1\right|}
\providecommand{\braket}[2]{\left\langle#1\middle|#2\right\rangle}
\newcommand{\dyad}[1]{\left|#1\right\rangle\!\left\langle#1\right|}
% v3.21: general outer product |a><b| and expectation value <A> — agents reach
% for these by reflex (esp. \expval); quantikz defines neither, so their absence
% threw "Undefined control sequence" in a real build. \providecommand = safe.
\providecommand{\ketbra}[2]{\left|#1\right\rangle\!\left\langle#2\right|}
\providecommand{\expval}[1]{\left\langle#1\right\rangle}
\newcommand{\abs}[1]{\left|#1\right|}
\newcommand{\norm}[1]{\left\lVert#1\right\rVert}
\newcommand{\half}{\tfrac12}
\newcommand{\R}{\mathbb{R}}\newcommand{\C}{\mathbb{C}}\newcommand{\Z}{\mathbb{Z}}
% identity operator. NB: \mathbb{1} silently renders a WRONG glyph (msbm has no
% blackboard digits). Use \mathbb{I} = 𝕀 (works with amssymb, no extra package).
% Do NOT "fix" this back to \mathbb{1}.
\newcommand{\id}{\mathbb{I}}

% ---- physical units (v3.35): \punit — upright unit with thin space ---------
% Usage: $300\punit{K}$, $80\punit{mV/decade}$, $10^{7}\punit{cm/s}$ — the
% NUMBER AND UNIT LIVE IN ONE MATH ISLAND. Promoted verbatim from a 199-page
% production lakit (advanced-electronic-devices, dozens of uses).
% Two measured traps this design avoids (both render-proven):
%   1. siunitx SCRAMBLES under RTL — \SI{3.5}{eV} renders "eV 3.5", exponents
%      detach ("Hz 10 x 10^9"): its quantity assembly is not bidi-isolated.
%      Do NOT load siunitx in this pipeline.
%   2. A text-mode unit next to a bare-prose number is TWO adjacent LTR runs
%      and reorders ("5 mm" -> "mm5" — the BIDI-2/3 class). \punit is therefore
%      math-only ON PURPOSE: prose misuse fails LOUDLY at compile (\mathrm
%      outside math) instead of scrambling silently.
% Domain shorthands (\MeV, \fm...) stay per-document: \newcommand{\MeV}{\,\text{\en{MeV}}}
\providecommand{\punit}[1]{\,\mathrm{#1}}

% ============================================================================
% Colored boxes — mdframed (NOT tcolorbox, which breaks under bidi).
% ★★ v3.4: innerleftmargin/innerrightmargin give code listings room so their
% frame/line-numbers do not spill out the box border (see listing styles below).
% ============================================================================
% v3.61: everyline/split skips were applied to `tablebox` only (BLD-14, v3.37) —
% the six colored boxes below never got them, so they still split OPEN across a page
% break: the page-N segment gets top+sides and NO bottom rule, the page-N+1 segment
% sides+bottom and NO top rule. Measured on the shipped fleet this is a COVERAGE GAP,
% not a re-report of BLD-14: 0 split-open `tablebox` segments (that fix holds) vs 62
% split-open COLORED boxes across 7 documents. Verified after the fix: both segments
% close, and the frametitle is NOT duplicated on the continuation. BLD-14's caveat
% still applies — a box taller than a full page grinds mdframed's split-trial loop, so
% split such content manually.
% ============================================================================
% v3.72: frametitle={} leads EVERY \newmdenv option list below. mdframed options
% are GROUP-scoped: a box nested inside \begin{exbox}[frametitle={...}] used to
% INHERIT the outer title — measured, the identical title string rendered twice,
% the second riding a title-bar fill inside the INNER box's frame, for all seven
% environments in both directions. Both natural forms were broken (= preamble
% bug, not an authoring rule). The explicit reset makes every box titleless
% unless ITS OWN [frametitle={...}] says otherwise; verified both directions:
% no duplicate on the nested box, an explicit title still renders.
% Geometry is shared with the optional semantic-width table container. Themes
% select their own existing semantic style below; neither path fixes columns.
\mdfdefinestyle{hllpdfsemanticgeometry}{linewidth=1.1pt,
  innertopmargin=6pt,innerbottommargin=7pt,skipabove=9pt,skipbelow=9pt,
  innerleftmargin=9pt,innerrightmargin=9pt,everyline=true,
  splitbottomskip=6pt,splittopskip=12pt,roundcorner=3pt}
\newcommand{\hllpdfsemanticframestyle}{hllpdfsemanticgeometry}
\newmdenv[frametitle={},style=hllpdfsemanticgeometry,linecolor=cBlue,backgroundcolor=cBlueBg,
  frametitlebackgroundcolor=cBlue,frametitlefontcolor=white,
  ]{defbox}   % הגדרה
\newmdenv[frametitle={},style=hllpdfsemanticgeometry,linecolor=cGreen,backgroundcolor=cGreenBg,
  frametitlebackgroundcolor=cGreen,frametitlefontcolor=white,
  ]{thmbox}   % משפט/טענה
\newmdenv[frametitle={},style=hllpdfsemanticgeometry,linecolor=cAmber,backgroundcolor=cAmberBg,
  frametitlebackgroundcolor=cAmber,frametitlefontcolor=white,
  ]{notebox}  % הערה
\newmdenv[frametitle={},style=hllpdfsemanticgeometry,linecolor=cPurple,backgroundcolor=cPurpleBg,
  frametitlebackgroundcolor=cPurple,frametitlefontcolor=white,
  ]{exbox}    % דוגמה
\newmdenv[frametitle={},style=hllpdfsemanticgeometry,linecolor=cRed,backgroundcolor=cRedBg,
  frametitlebackgroundcolor=cRed,frametitlefontcolor=white,
  ]{warnbox}  % טעות נפוצה
\newmdenv[frametitle={},style=hllpdfsemanticgeometry,linecolor=cTeal,backgroundcolor=cTealBg,
  frametitlebackgroundcolor=cTeal,frametitlefontcolor=white,
  ]{keybox}   % רעיון מרכזי

% ---- codebox (v4.5.0): the dedicated container for code -------------------
% An OPTION, not a policy: asmblock unboxed stays canonical for a full listing,
% and a snippet inside any of the six colored boxes still works as before.
% What makes codebox worth its own environment is that it is the ONE framed box
% that keeps its margin line numbers.
%
% v4.4.0 shipped it WITHOUT them, on the strength of v3.37's "no margin knob can
% save them" plus four fresh probes that agreed. That was wrong: every probe had
% moved the code, the box, or both — never the key that carries the NUMBER
% COLUMN. Measured on the page:
%   resetmargins=false   the listing follows THIS box instead of the page,
%                        moving the code from x=75.26 to x=97.54
%   xrightmargin=26pt    under RTL this is the PHYSICAL LEFT indent, and the
%                        only key that carries the \llap'ed digits with it:
%                        x=64.28 -> 90.18, clearing the rule at 65.64 by 23.6pt
%   xleftmargin=26pt     moves the code and LEAVES THE DIGITS on the rule —
%                        the asymmetry that hid the fix for eleven versions
% resetmargins is hoisted out of the four styles (see the \lstset below them) so
% this override is possible at all; every other listing still resolves to true.
% Pinned by tests/codebox_numbers.tex + probe_lstnum.py --inside-rule, whose
% negative control strips xrightmargin and watches the digits fall back.
\definecolor{cSlate}{HTML}{3F4A56}  \definecolor{cSlateBg}{HTML}{F2F4F7}
\newmdenv[frametitle={},leftline=true,rightline=true,topline=true,bottomline=true,
  linecolor=cSlate,backgroundcolor=cSlateBg,linewidth=0.9pt,
  frametitlebackgroundcolor=cSlate,frametitlefontcolor=white,
  innertopmargin=6pt,innerbottommargin=7pt,skipabove=9pt,skipbelow=9pt,
  innerleftmargin=9pt,innerrightmargin=9pt,everyline=true,splitbottomskip=6pt,splittopskip=12pt,
  roundcorner=3pt]{codebox}   % קוד
% The two keys that make in-box numbering work, both measured on the page:
%   resetmargins=false  the listing follows THIS box instead of the page
%   xrightmargin        under RTL this is the PHYSICAL LEFT indent, and it is
%                       the only one that carries the \llap'ed number column
%                       with it. xleftmargin moves the code and leaves the
%                       digits behind on the rule (measured: code 97.54,
%                       digits 64.28 — the v3.37 defect, reproduced); with
%                       xrightmargin=26pt the digits land at x=90.18, a
%                       clearance of 23.6pt from the rule at 65.64.
% This asymmetry is why v3.4's note had the semantics backwards and v3.37
% concluded the digits could not be saved.
% <hllpdf:listing-container-state>
% v4.29: all state is local to the current environment group. A semantic
% callout masks number text, while a nested codebox restores the formatter
% saved before that mask. Save it once across successive semantic wrappers;
% an explicit numbers=none (or an explicitly empty formatter) remains valid.
\makeatletter
\newif\ifhllpdf@incodebox
\newif\ifhllpdf@englishcodebox
\newif\ifhllpdf@lstnumbermasked
\let\hllpdf@savedlstnumber\thelstnumber
\makeatother
% </hllpdf:listing-container-state>
% <hllpdf:codebox-gutter>
\makeatletter
\AtBeginEnvironment{codebox}{%
  \hllpdf@incodeboxtrue
  % Snapshot the OUTER language before asmblock selects English.
  \iflanguage{english}{\hllpdf@englishcodeboxtrue}{\hllpdf@englishcodeboxfalse}%
  \ifhllpdf@lstnumbermasked
    \let\thelstnumber\hllpdf@savedlstnumber
    \hllpdf@lstnumbermaskedfalse
  \fi
}
% listings' AddToHook is GLOBAL. Register once, and gate on local box state;
% adding an unconditional hook at every codebox entry leaks into later listings.
\lst@AddToHook{PreInit}{%
  \ifhllpdf@incodebox
    % Apply every listing override here, inside the LISTING's own group.
    % Setting margins at box entry leaks them into nested semantic callouts.
    \lstset{resetmargins=false,xrightmargin=26pt}%
    % Styles resolve before PreInit. In LTR the physical number gutter follows
    % xleftmargin; in Hebrew retain the measured inherited xrightmargin path.
    \ifhllpdf@englishcodebox
      \lstset{xleftmargin=26pt,xrightmargin=0pt}%
    \fi
    % ONE box, not two: suppress only the frame/background/skips of a listing
    % inside this codebox. The next unboxed listing keeps its own style.
    \lstset{frame=none,backgroundcolor=\color{cSlateBg},aboveskip=0pt,belowskip=0pt}%
  \fi
}
\makeatother
% </hllpdf:codebox-gutter>

% v3.37: line numbers inside a COLORED BOX are sliced by the box border —
% measured: digit span x=63.5-67.3 vs border line at x=65.2 (the border runs
% THROUGH the digits), and the inter-border gap (3.5pt) is narrower than a
% digit (3.8pt), so no margin knob can save them (under resetmargins,
% xleftmargin is inert in-box — measured: code at x=75.26 with 2.2em AND
% 3.0em). A 3-6 line boxed snippet does not need line numbers; blank the
% number TEXT inside every colored box (\thelstnumber is what gets printed;
% the style's numbers=left survives, so nothing else shifts — code position
% byte-frozen at x=75.26, verified). Unboxed listings keep their numbers.
% v3.78: the blanking is coupled to the FRAME, not to the box name —
% clean_theme.tex registers LATER hooks (registration order wins) that
% restore the number text inside the boxes it renders frameless.
% v4.5.0 RETRACTS the "no margin knob can save them" clause above. It holds only
% under resetmargins=true (the shipped default): there xleftmargin is inert
% in-box, as measured. With resetmargins=false, xrightmargin DOES carry the
% \llap'ed digit column clear of the rule — measured; the codebox \lstset that
% does it is the one above, in the codebox block. The loop below therefore
% covers SIX boxes, not "every colored box": codebox opts out and keeps its
% numbers.
% <hllpdf:semantic-listing-mask>
\makeatletter
\@for\HEB@lstbox:=defbox,thmbox,notebox,exbox,warnbox,keybox\do{%
  \expandafter\AtBeginEnvironment\expandafter{\HEB@lstbox}{%
    \ifhllpdf@lstnumbermasked\else
      \let\hllpdf@savedlstnumber\thelstnumber
    \fi
    \hllpdf@lstnumbermaskedtrue
    \def\thelstnumber{}%
    % The nearest semantic container owns its listing presentation. On exit,
    % TeX's group restores the enclosing codebox state and formatter.
    \hllpdf@incodeboxfalse}}
\makeatother
% </hllpdf:semantic-listing-mask>

% ============================================================================
% Unified TABLE presentation: subtle uniform slate frame + auto-numbered
% caption rendered INSIDE the frame. Use:
%   \begin{tablebox}\centering  <tabular>  \tabcap{<description>}\end{tablebox}
% Replaces the buggy \caption* (which emitted a stray "טבלה N : *" line) and
% provides an explicit framed option + numbered caption "טבלה N.M: תיאור".
% Optional intro text may be placed at the TOP, before the tabular, inside the box.
% ============================================================================
\definecolor{cTblLine}{HTML}{B9C4D2}  \definecolor{cTblBg}{HTML}{FAFBFC}
% v3.23: class-aware — under `article` there is no chapter counter (the promised
% report->article switch used to die with "No counter 'chapter' defined").
\makeatletter
\@ifundefined{c@chapter}{%
  % v3.31: in article the counter resets per SECTION, so an unprefixed number
  % repeated ("טבלה 1." in every section — measured on a two-section render).
  % Same conditional idea as the report branch: prefix the section number once
  % a real \section exists; a pre-section table stays "טבלה 1".
  \newcounter{tbl}[section]
  \renewcommand{\thetbl}{\ifnum\value{section}>0 \thesection.\fi\arabic{tbl}}%
}{%
  \newcounter{tbl}[chapter]
  % v3.30: a report-class doc that uses SECTIONS ONLY (chapter stays 0 —
  % the common short-doc pattern) used to print "טבלה 0.1"; the chapter
  % prefix now appears only once a real \chapter exists (render-verified
  % both ways: "טבלה 1." chapterless, "טבלה 1.1" with chapters).
  \renewcommand{\thetbl}{\ifnum\value{chapter}>0 \thechapter.\fi\arabic{tbl}}%
  % v3.64: the SAME leak in \thesection itself, left behind by v3.30. That fix
  % guarded the CAPTION counter and stopped there, so a chapterless report still
  % printed the heading "0.1 סעיף ראשון" and \ref to it printed "0.2" — measured:
  %   p1 y=65.2  '…0.1'   p1 y=144.9 '…0.2'   p1 y=183.4 '…0.2.1'
  % report.cls hardcodes \thesection as \thechapter.\@arabic\c@section, and the
  % leading 0 is the chapter counter that no \chapter ever incremented. Fixing
  % \thesection fixes the whole chain — \thesubsection and \thesubsubsection are
  % defined in terms of it. \theequation and \thetbl/\thefig were already guarded,
  % which is exactly why this survived: every NUMBER around the heading looked
  % right, so the heading's own 0 read as a style choice rather than a bug.
  % \appendix is safe: report.cls resets chapter to 0 and redefines \thechapter to
  % \@Alph\c@chapter, so an appendix \chapter makes the counter 1 and the guard
  % emits "A.1" unchanged. Render-verified in all four modes (chapterless,
  % chaptered, appendix+chapter, appendix+section-only).
  \renewcommand{\thesection}{\ifnum\value{chapter}>0 \thechapter.\fi\arabic{section}}%
}
\makeatother
% v3.37: everyline=true — if the box DOES split across pages, every segment
% gets a CLOSED frame (measured before: the split ran open-ended into the page
% edge with no bottom border, and the caption could strand alone at the next
% page top — "טבלה A.2" shipped that way). The split skips give the segments
% breathing room. In practice, with the caption GLUE below, every previously-
% split table in the corpus now moves WHOLE instead (the white-gap tradeoff,
% documented benign); everyline stays as the safety net for a genuinely
% page-tall table. NB: a table taller than ONE FULL PAGE can grind mdframed's
% split-trial loop (reproduced: minimal doc, minutes of 1pt-step Overfull
% probes) — break such a monster into two tables.
\mdfdefinestyle{hllpdftablegeometry}{leftline=true,rightline=true,topline=true,bottomline=true,
  linewidth=0.9pt,
  innertopmargin=9pt,innerbottommargin=8pt,skipabove=11pt,skipbelow=11pt,
  innerleftmargin=11pt,innerrightmargin=11pt,roundcorner=4pt}
\newcommand{\hllpdftableframestyle}{hllpdftablegeometry}
\newmdenv[frametitle={},style=hllpdftablegeometry,
  linecolor=cTblLine,backgroundcolor=cTblBg,
  % v4.2.0 (adopted): tablebox is ATOMIC. Short tableboxes near page bottoms
  % repeatedly entered mdframed's split trial, emitting benign Overfull-vbox
  % bursts and occasionally separating the intro prose from the tabular.
  % nobreak=true moves the complete short table to the next page instead. The
  % output-active Overfull-vbox gate still catches a box taller than a page;
  % everyline/split skips remain as defensive styling should that path ever
  % re-arm, but the canonical path never publishes a split tablebox.
  nobreak=true,
  everyline=true,splitbottomskip=6pt,splittopskip=12pt,
  % v3.79 (BLD-37 table side, correction): startinnercode/endinnercode run
  % INSIDE mdframed's body vbox (mdframed.sty:848 / :856+863) -- the second one
  % after \par\unskip, i.e. after line breaking, which is the only moment the
  % \tabcap anchor's true offset from the top of the box EXISTS. See
  % \lakintabboxclose below. These two keys are the ONLY tablebox change.
  startinnercode={\lakintabboxopen},endinnercode={\lakintabboxclose}]{tablebox}
% v3.37: \nopagebreak GLUE — \vsplit honors penalties, so the caption can no
% longer strand alone at the top of the next page (it takes its rows along, or
% the whole box moves). Render-verified on the real shipped orphan.
% v3.59: \NR@gettitle is what nameref injects into the stdlib title-bearing
% commands (\@sect, \@chapter, \caption) to fill \@currentlabelname. \figcap and
% \tabcap are INDEPENDENT macros — they never called it, so \@currentlabelname kept
% whatever the last heading left and \nameref{fig:x} silently printed the enclosing
% SECTION's title instead of the caption (proven in the .aux: field 3 of \newlabel
% was the heading text while field 1 and the anchor `fig.1.1` were correct — so this
% is a THIRD, distinct head next to BLD-13's counter fix and BIDI-38's autoref-name
% fix, neither of which writes that field). \@ifundefined keeps it inert if
% hyperref/nameref is ever dropped. Guarded helper because \figcap is defined
% after \makeatother.
% NB both helpers below need \makeatletter — this part of the preamble is OUTSIDE
% any letter-catcode region, so an unwrapped \@currentlabel parses as LaTeX's \@
% and dies with "You can't use \spacefactor in vertical mode" (measured).
\makeatletter
\newcommand{\lakinrtitle}[1]{\@ifundefined{NR@gettitle}{}{\NR@gettitle{#1}}}
% v3.59: \hitem{<letter>}{<key>} — a REFERENCEABLE Hebrew-lettered list item.
%   \begin{enumerate} \hitem{א}{it:a} הסעיף הראשון. \end{enumerate}  ->  \ref{it:a} = (א)
% Why this exists: automatic Hebrew list labels fail silently, so SKILL.md mandates
% an EXPLICIT \item[(\textbf{א})] — but LaTeX puts \refstepcounter INSIDE
% \if@noitemarg (latex.ltx:12964), and \item[...] never sets that flag. So an
% explicit item steps no counter and never touches \@currentlabel: a \label there
% binds to the AMBIENT one and \ref prints the CHAPTER/SECTION number where the
% item letter belongs (proven in the .aux: {chapter.1}). Same doctrine as \figcap —
% own the label state before \label. \phantomsection matters: without it the
% hyperlink still lands on the chapter head (measured); with it the anchor is
% section*.N and the bookmark outline stays clean.
\newcommand{\hitem}[2]{\item[(\textbf{#1})]%
  \def\@currentlabel{(#1)}\def\@currentlabelname{#1}\phantomsection\label{#2}}
% NB (v3.66): \hitem is \ref-only. \autoref{key} calls it a SECTION ("סעיף (א)")
% because \phantomsection brands the anchor section*.N, and \nameref{key} prints
% the bare letter. For \autoref-able items prefer \hebenumlabels below.
%
% v3.66: \hebenumlabels — OPT-IN automatic Hebrew enumerate letters, done right.
%   {\hebenumlabels\begin{enumerate}\item\label{q:a} ...\end{enumerate}}
%   -> labels (א) (ב) (ג); \ref{q:a} = א; anchor Item.N (real, per-item).
% This retires the belief that made \item[(\textbf{א})] + \hitem necessary.
% "Automatic Hebrew labels fail silently" was TRUE but unexplained, and the
% mechanism is now bisected (plain article vs +enumitem, measured): ENUMITEM
% PRE-EXPANDS the label expression in an \edef BEFORE \refstepcounter, with the
% counter still at 0. Tokens it shields by name (\arabic,\alph,\roman) survive;
% a plain \newcommand \ifcase wrapper evaluates at 0 -> the EMPTY zero-branch ->
% the label prints as "()" with the letter GONE (and \the\c@enumi prints a
% literal 0, \number\value dies fatally pointing at \item). The cure is one
% property, not a doctrine: the wrapper must be ROBUST (survives the \edef
% unexpanded), while \theenumi stays EXPANDABLE (so \refstepcounter writes the
% real letter to the .aux). Verified: labels (א)(ב)(ג), \ref -> א/ב/ג each item
% its own value, .aux = {א}{...}{Item.1} — versus the explicit-\item doctrine
% where BIDI-50/\hitem exist purely to patch the counter that never stepped.
\def\@hebalph#1{\ifcase#1\or א\or ב\or ג\or ד\or ה\or ו\or ז\or ח\or ט\or י\or
  יא\or יב\or יג\or יד\or טו\or טז\or יז\or יח\or יט\or כ\else\@ctrerr\fi}
\DeclareRobustCommand*\hebalph[1]{\expandafter\@hebalph\csname c@#1\endcsname}
\newcommand{\hebenumlabels}{%
  \renewcommand{\theenumi}{\@hebalph\c@enumi}%
  \renewcommand{\labelenumi}{(\textbf{\hebalph{enumi}})}}
\makeatother
% v3.63: \mbox the LABEL. \thetbl ("7.1") and a caption description that STARTS
% WITH A DIGIT ("32 הרגיסטרים…") are two European-Number runs separated only by
% neutrals (the label's own "." and \enspace). Under the RTL paragraph base the UBA
% resolves those neutrals to the RTL level, so both numbers sit in ONE reversed span
% and EXCHANGE VISUAL POSITIONS. Measured on the shipped PDF: a table whose real
% number is 7.1 printed "טבלה 32" — a different, EXISTING table number — while \ref
% correctly printed 7.1, so the caption silently contradicted its own cross-reference,
% and the description lost its leading count. The source contains no adjacency at all;
% it is manufactured by this macro, which is why every source linter was blind.
% \mbox makes the label one atomic box so nothing can reorder across it. Verified
% zero-regression: an ordinary Hebrew caption keeps identical x-geometry.
% One paragraph style covers label, number, separator and description.
% Explicit project/template overrides remain possible; the label hook is kept
% for compatibility, with an identity default. Counters/anchors stay unchanged.
\newcommand{\lakitcaptionstyle}{\normalfont\small\color{black}\hypersetup{allcolors=black}}
\newcommand{\lakitcaptionlabelstyle}[1]{#1}
\newcommand{\lakitcaptionpunctuation}{:}
% Native floats, captionof and longtable use the same presentation policy.
% subcaption is optional and inherits the same font rather than a second size.
\usepackage{caption}
\DeclareCaptionFont{lakit}{\lakitcaptionstyle}
\DeclareCaptionLabelFormat{lakit}{\lakitcaptionlabelstyle{#1~\en{#2}}}
\DeclareCaptionLabelSeparator{lakit}{\lakitcaptionpunctuation\enspace}
\captionsetup{font=lakit,labelfont={},textfont={},labelformat=lakit,
  labelsep=lakit,justification=centering}
\makeatletter
\AtBeginDocument{\@ifpackageloaded{subcaption}{%
  \captionsetup[sub]{font=lakit,labelfont={},textfont={},labelsep=lakit}}{}}
\makeatother
% v3.64: OPT-IN list-of-figures/tables entry via an optional short caption.
% \listoffigures and \listoftables are inherited from report.cls and LOOK available,
% but \figcap/\tabcap step the custom \c@fig/\c@tbl counters and never call
% \addcontentsline — so both printed their heading and NOTHING ELSE. Measured: a doc
% with one \figcap and one \tabcap rendered page 1 as 'רשימת האיורים\n1\n', and the
% .lof held only \@chapter's \addvspace, zero \contentsline entries. An empty list
% under a correct Hebrew heading reads as "this document has no figures".
%
% WHY OPTIONAL, AND NOT AUTOMATIC. Writing #1 unconditionally was implemented first
% and REJECTED: a caption containing \codeline compiles CLEANLY today (verified on
% the pre-v3.64 preamble: 0 errors, PDF produced) but becomes FATAL once it is also
% written to the .lof — \codeline is \hbox to\linewidth{...}, and a full-width box in
% a list entry dies with "\pdfextension endlink cannot be used in vertical mode", no
% output PDF. Gaining a feature nobody uses (0 of 13 collections call \listoffigures)
% at the price of breaking a construct that works is a bad trade. Stock \caption has
% exactly this limitation and exactly this escape hatch, so the convention is familiar.
% \en{}, \code{}, \h{} and inline math in the short text are all safe — verified, they
% survive the write and typeset correctly in the list (RTL geometry confirmed: number
% rightmost at x=509.9, title left of it, dots and folio at the far left).
% With no optional argument the macro is byte-identical to v3.63 — zero risk to every
% existing document, and \listoffigures lists exactly the captions that opted in.
% BLD-37, TABLE side (the \figcap v3.74 idea, adapted): clicking a \ref to a
% table used to land MID-table — hyperref's automatic anchor inside
% \refstepcounter was emitted on the still-open horizontal line whose TAIL is
% the tabular (an hbox whose baseline is the table's vertical CENTER), so the
% jump scrolled the top half of the table above the viewport. The v3.74 entry
% deferred this side claiming "tablebox splits — no single box to measure";
% that deferral dissolves: the measurement target is the TABULAR itself at
% \tabcap ENTRY, while it is still the tail of the open list — the same
% situation \lakin@grabfig handles at an island's end (and the fleet has 0
% split tableboxes). So \lakin@grabtab runs as the FIRST token, before any
% typesetting can push the box away, and the anchor is re-emitted raised above
% the measured tabular. \lakintabraise is a SEPARATE register from
% \lakinfigraise (defined with it, below the LTR-island block): a captionless
% island intentionally leaves the fig register set, and a table must never
% inherit that stale figure height.
%
% WHY the anchor stays AT THE ENTRY SPOT and is NOT deferred into the caption
% paragraph the way \figcap defers it: at \par TeX DROPS a final glue item,
% and the automatic anchor's smash box sitting BEHIND the interword glue kept
% that glue alive — so every centered tabular has always sat half a space off
% true center. Deferring the anchor removed that box and was measured as 982
% differing glyph records: EVERY tabular line shifted 1.30pt (2.61pt for
% pattern 7's two minipages) while captions stayed put. Emitting the raised
% anchor at the old spot — through hyperref's own \Hy@raisedlink, exactly the
% macro the automatic path uses — keeps the line's node structure, and so the
% A/B glyph diff at 0, BY CONSTRUCTION. \Hy@raisedlink reads the raise from
% \HyperRaiseLinkDefault (default \baselineskip = the old mid-table dest); the
% group-local \edef freezes measured-height+42pt into it, and \smash inside
% \Hy@raisedlink hides any raise from the line's extent. \HyCnt@ProvideTheH-
% Counter keeps the dest NAME byte-identical to the automatic one (it defines
% \theHtbl, which \hyper@makecurrent prefers — without it the first measured
% caption was named tbl.1 where the automatic path said tbl.0.1, measured).
% vmode entry, or anything but a box at the tail (no tabular to measure),
% falls back to the plain \refstepcounter — byte-identical to the old path.
\makeatletter
\newcommand{\tabcap}[2][]{\lakin@grabtab
  % v3.79: arm the box-close correction. Set unconditionally and in BOTH
  % branches below's scope, so the plain-\refstepcounter fallback is corrected
  % too; \lakintabboxopen clears it at every \begin{tablebox}, so a \tabcap
  % outside a tablebox can never leave it armed for a later box.
  \global\lakintabcapseentrue
  \ifdim\lakintabraise>0pt
    \H@refstepcounter{tbl}\HyCnt@ProvideTheHCounter{tbl}\hyper@makecurrent{tbl}%
    \begingroup
      % A frameless table targets its measured top. The frame's historical
      % cushion belongs only to tablebox and is corrected at box close.
      \edef\HyperRaiseLinkDefault{\the\lakintabraise}%
      \iflakintabframed
        \edef\HyperRaiseLinkDefault{\the\dimexpr\lakintabraise+42pt\relax}%
      \fi
      \Hy@raisedlink{\hyper@anchorstart{\@currentHref}\hyper@anchorend}%
    \endgroup
    \global\lakintabraise 0pt\relax
  \else
    \refstepcounter{tbl}%
  \fi
  \lakinrtitle{#2}\par\nopagebreak\addvspace{6pt}%
  \def\lakin@short{#1}\ifx\lakin@short\@empty\else
    \addcontentsline{lot}{table}{\protect\numberline{\thetbl}{\protect\ignorespaces #1}}\fi
  {\centering\lakitcaptionstyle\mbox{\lakitcaptionlabelstyle{\tablename~\en{\thetbl}\lakitcaptionpunctuation}}\enspace #2\par}}
\makeatother

% ============================================================================
% Unified FIGURE caption — the \tabcap twin for TikZ/plot figures (v3.34.0).
% A figure (tikzpic) sits inline in the prose; put \figcap BELOW it for a
% numbered, \ref-able caption "איור N.M: תיאור" in the same style as \tabcap:
%   \begin{tikzpic} ... \end{tikzpic}
%   \figcap{\label{fig:foo}description}     % \label FIRST, inside the argument
% WHY a macro, not a hand-written gray line: a manual
%   {\centering\small\textcolor{cGray}{איור: ...}} with a \label inside has NO
% \refstepcounter, so the \label binds to the AMBIENT \@currentlabel (the last
% equation) — \ref then prints an equation number AS the figure number
% (measured on a real page: "איור 48" where 48 is equation 48). \figcap steps a
% real \c@fig counter first, so \label is correct, numbering is consistent, and
% refs are stable. Class-aware like \thetbl (chapter- or section-prefixed).
% ============================================================================
\makeatletter
\@ifundefined{c@chapter}{%
  \newcounter{fig}[section]
  \renewcommand{\thefig}{\ifnum\value{section}>0 \thesection.\fi\arabic{fig}}%
}{%
  \newcounter{fig}[chapter]
  \renewcommand{\thefig}{\ifnum\value{chapter}>0 \thechapter.\fi\arabic{fig}}%
}
\makeatother
\makeatletter
\newcommand{\figcap}[2][]{%
  \ifdim\lakinfigraise>0pt
    % v3.74: step WITHOUT hyperref's automatic anchor, then emit the anchor
    % raised above the measured figure. \H@refstepcounter is hyperref's saved
    % pre-hook \refstepcounter (sets the counter AND \@currentlabel — \label
    % in the caption arg binds exactly as before); \hyper@makecurrent supplies
    % \@currentHref, so the dest NAME is identical to the old automatic one —
    % \ref/\autoref/bookmarks are unchanged, only the dest'S POSITION moves.
    \H@refstepcounter{fig}\hyper@makecurrent{fig}%
    % The anchor is emitted INSIDE the caption paragraph below, as a zero-size
    % raisebox — NOT here in vertical mode, where a \raisebox starts an empty
    % paragraph and adds a blank line (measured: every glyph after the caption
    % shifted +18.8pt). Inside an existing line a [0pt][0pt] box is inert.
    \edef\lakin@anch{\@currentHref}%
    % \edef FREEZES the measured height into the pending macro: the deferred
    % \raisebox is typeset inside the caption paragraph AFTER the register is
    % reset to the 0pt sentinel below — a plain \def read 0 at use time
    % (measured: LAKINRAISE=0.0, dest moved 14pt instead of above the figure).
    \edef\lakin@pend{\noexpand\raisebox{\the\dimexpr\lakinfigraise+42pt\relax}[0pt][0pt]{%
      \noexpand\hyper@anchorstart{\lakin@anch}\noexpand\hyper@anchorend}}%
    \global\lakinfigraise 0pt\relax
  \else
    \refstepcounter{fig}\let\lakin@pend\relax
  \fi
  \lakinrtitle{#2}\par\nopagebreak\addvspace{4pt}%
  \def\lakin@short{#1}\ifx\lakin@short\@empty\else
    \addcontentsline{lof}{figure}{\protect\numberline{\thefig}{\protect\ignorespaces #1}}\fi
  {\centering\lakin@pend\lakitcaptionstyle\mbox{\lakitcaptionlabelstyle{\figurename~\en{\thefig}\lakitcaptionpunctuation}}\enspace #2\par}}
\makeatother

% ============================================================================
% LTR islands.  EVERY tikzpicture / quantikz MUST live inside one, or the whole
% figure's coordinates and Latin labels flip. (Hebrew node text still uses \h{}.)
% ============================================================================
% v3.74 (BLD-37): the \figcap hyperlink DESTINATION used to land at the caption
% line — BELOW the figure — so clicking a figure reference scrolled the figure
% off-screen (measured fleet-wide: 97/97 fig dests in one collection sat 19-54pt
% under their figure's ink; a 347pt figure was fully hidden after the jump).
% hypcap-style fix in two parts: (1) at the island's END, \lakin@grabfig grabs
% the just-typeset centered line with \lastbox, measures it, and puts it back
% with \nointerlineskip — the interline glue above the line was already
% contributed and is NOT re-added, so the page geometry is unchanged to the
% glyph (verified); (2) \figcap then steps the counter WITHOUT hyperref's
% automatic anchor (\H@refstepcounter = the pre-hyperref \refstepcounter) and
% emits the anchor itself, raised by the measured height — the dest lands ABOVE
% the figure top. A captionless island leaves numbering untouched (nothing is
% stepped at the island; the grab flag is cleared at every island begin, so a
% stale height from a captionless figure cannot leak into a later caption).
\newsavebox{\lakinfigbox}
\newdimen\lakinfigraise  % real \dimen: tex.setdimen cannot target a \newlength SKIP register (measured lua error)
% BLD-37 table side: a SEPARATE register for \tabcap. The two must not share —
% a captionless island INTENTIONALLY leaves \lakinfigraise set (the grab runs
% at every island END; only the NEXT island's begin clears it), so a shared
% register would hand a \tabcap after a captionless figure that stale FIGURE
% height. \lakintabraise is written only by \lakin@grabtab at \tabcap entry
% and consumed/reset inside \tabcap itself, so it can never go stale.
\newdimen\lakintabraise
\makeatletter
% v3.74b: measure WITHOUT touching the list. The first draft grabbed the box
% with \lastbox and put it back — and \lastbox DISCARDS the node's baseline
% shift, so a vertically-anchored quantikz circuit rose 2.57pt and everything
% after it sank 5.32pt (measured on the A/B glyph diff). LuaTeX lets us READ
% the tail node's dimensions without removing anything: geometry is identical
% by construction. \lakinfigraise=0pt is the "no figure measured" sentinel.
\def\lakin@grabfig{%
  \ifhmode
    \directlua{
      local n = tex.nest[tex.nest.ptr].tail
      local GLUE, PEN = node.id('glue'), node.id('penalty')
      while n and (n.id == GLUE or n.id == PEN) do n = n.prev end
      local HL, VL = node.id('hlist'), node.id('vlist')
      if n and (n.id == HL or n.id == VL) then
        tex.setdimen('global', 'lakinfigraise',
                     (n.height or 0) + (n.depth or 0))
      else
        tex.setdimen('global', 'lakinfigraise', 0)
      end}%
  \else
    \global\lakinfigraise 0pt\relax
  \fi}
% BLD-37 table side: \lakin@grabfig's twin, NOT a shared helper — the Lua body
% differs in exactly one place, the skip set. Between \end{tabular} and \tabcap
% an otherlanguage wrapper (patterns 3/5 of example_tables) leaves a \write
% whatsit and a dir node ON TOP of the interword glue (measured node dump:
% "whatsit/write <- dir <- glue(2.62) <- hlist"), so the table reader must
% also step over whatsit/dir markers — all zero-size — to reach the tabular.
% \lakin@grabfig deliberately does NOT get the wider skip set: its glue/penalty
% body is the one the v3.74 A/B proof froze, and widening it would newly
% measure (and so newly raise) anchors in island shapes that today take the
% plain path — an unproven behavior change on the FIG side. Anything that is
% not a box behind those markers still hits the 0pt sentinel and \tabcap falls
% back to the plain \refstepcounter.
% SECOND deliberate difference: this stores the box's HEIGHT, not height+depth.
% \figcap's anchor is emitted in the caption paragraph BELOW the figure, so it
% must climb the full extent; \tabcap's anchor sits ON the tabular's own line,
% whose baseline is the vcentered tabular's MIDPOINT — height alone reaches
% the top edge, and adding the depth was measured to overshoot by half the
% table (dests 80-135pt above their frame, INSIDE the preceding table's box).
\def\lakin@grabtab{%
  \ifhmode
    \directlua{
      local n = tex.nest[tex.nest.ptr].tail
      local GLUE, PEN = node.id('glue'), node.id('penalty')
      local WHAT, DIR = node.id('whatsit'), node.id('dir')
      while n and (n.id == GLUE or n.id == PEN
                   or n.id == WHAT or n.id == DIR) do n = n.prev end
      local HL, VL = node.id('hlist'), node.id('vlist')
      if n and (n.id == HL or n.id == VL) then
        tex.setdimen('global', 'lakintabraise', n.height or 0)
      else
        tex.setdimen('global', 'lakintabraise', 0)
      end}%
  \else
    \global\lakintabraise 0pt\relax
  \fi}
\makeatother
% ============================================================================
% v3.79 (BLD-37 table side, CORRECTION): put the anchor at the TOP OF THE BOX,
% instead of guessing the distance to it.
% ============================================================================
% v3.77 raised the \tabcap anchor by (height of the box at the tail of the open
% horizontal list) + 42pt. That is a MEASUREMENT taken BEFORE line breaking, and
% it is short by every term it cannot see. Fleet-measured on the 12 built
% collections (299 tbl.* dests, ownership established through each table's own
% CAPTION, never through "which frame encloses the dest"): 3 dests landed INSIDE
% their own tablebox, below the frame top —
%   computer-architecture       tbl.14.1  -46.4pt   three stacked tabulars
%   computer-architecture       tbl.30.5  -13.5pt   tabular 2.85pt over \linewidth
%   advanced-electronic-devices tbl.B.1   -59.2pt   tabular 2.56pt over \linewidth
% Two missing terms, both reproduced in a minimal document (probe P2/P3):
%   (1) VERTICAL LIST ALREADY CONTRIBUTED inside the box — a second and third
%       tabular, or a heading above the table. The tail read sees only the LAST
%       box, so 77.7pt of stacked tabulars is invisible to the raise.
%   (2) WHICH LINE THE ANCHOR LANDS ON. When the tabular alone overfills
%       \linewidth (459.6121pt in a tablebox), the interword glue after
%       \end{tabular} becomes the chosen breakpoint and \Hy@raisedlink is
%       emitted on the NEXT line — one (depth + lineskip) lower. That term does
%       not EXIST at \tabcap time: line breaking has not happened yet. No
%       measurement taken there can ever recover it. (Proof: probe P3, the same
%       table 2.84912pt too wide, fails at -13.5pt; probe P4, the identical
%       table with \tabcolsep 2pt so it fits, is correct at +32.0pt.)
% So the raise is not computed at \tabcap at all. It is CORRECTED at
% \end{tablebox}, from inside mdframed's own body vbox (endinnercode runs after
% \par\unskip, i.e. after line breaking, while tex.nest[tex.nest.ptr] is still
% that vbox's vertical list). There the anchor's true offset from the top of the
% box content is not estimated but WALKED, and the \Hy@raisedlink shift is set
% so the dest lands 42pt above it — which is 42pt - innertopmargin(9pt) -
% linewidth(0.9pt) = 32.0pt above the FRAME top, the exact position the 295
% already-correct dests have today. That is why this is a correction and not a
% re-placement: the good sites do not move.
%
% WHY THE GLYPH DIFF IS 0 BY CONSTRUCTION — the only thing written is the
% `shift` field of the hlist that \Hy@raisedlink already built, and the guard
% below refuses to write it unless that box is 0x0x0. A zero-size box carries no
% dimension into its parent's packing, the packing already happened, and shift
% is not an input to it. Nothing else in the list is touched: no node is added,
% removed, resized or reordered. Measured, not assumed: assets/example_tables.tex
% and the whole computer-architecture collection A/B at 0 differing pymupdf
% rawdict glyph records.
%
% WHY NOT the alternatives (all prototyped and measured before this one):
%   * accumulate the vertical list at \tabcap time — fixes term (1), CANNOT fix
%     term (2); 2 of the 3 real failures survive. A fix that leaves the majority
%     of the defect standing is the false all-clear this repo exists to avoid.
%   * \zref-savepos two-pass — correct on all shapes, but it makes the dest
%     POSITION aux-derived (DEFECTS BLD-1/2), a single-pass compile silently
%     ships an unraised anchor, and the bare y-difference has no page guard.
%   * emit the anchor at \begin{tablebox} and let \tabcap adopt it — needs the
%     tbl counter stepped at box OPEN to keep the dest NAME byte-identical; a
%     captionless tablebox then either shifts every later table number or leaves
%     an orphan dest that the next table duplicates (a HARD build gate). This
%     correction touches neither the counter nor the name: \tabcap's naming path
%     is v3.77's, unchanged, so resolve_names() is byte-identical (verified).
%   * suppressing the line break (a penalty before the glue) — that MOVES the
%     tabular: the break currently discards the interword glue, so forbidding it
%     re-centers the two overfull tables. Reflow, rejected.
%
% SCOPE. The correction is armed only by \tabcap ITSELF (\lakintabcapseen),
% cleared at every box open, so a captionless tablebox is a no-op and a \tabcap
% hosted OUTSIDE a tablebox (the fleet has exactly one, computer-architecture
% tbl.14.2 in an exbox) keeps v3.77's measured raise. Both fallbacks are
% therefore never worse than today; and because hyperref's AUTOMATIC anchor uses
% the same \Hy@raisedlink shape, the plain-\refstepcounter path (vmode entry, or
% prose between \end{tabular} and \tabcap — probe P6, -51.5pt) is corrected too.
% KNOWN LIMIT — a tablebox that SPLITS across pages. The offset is measured from
% the top of the UNSPLIT body, so a dest in segment 2 gets raised past its own
% segment. MEASURED on a forced split (probe: 40 filler lines + a table): the
% dest went from y=585.0 on page 2 (inside segment 2, at the caption — the old
% BLD-37 symptom) to y=-113.4, i.e. ABOVE the page top. Viewers clamp that to the
% page top, which happens to show segment 2 whole, but an out-of-page coordinate
% is not something to ship on purpose. Hence the \textheight guard in
% lakintab.fix: a body TALLER THAN THE TEXT BLOCK cannot fit on any page, so it
% must split — the correction is skipped and v3.77's behaviour stands unchanged
% (verified: the probe returns to y=585.0). The guard is deliberately EXACT, not
% a page-fitting heuristic: it is sound (a taller-than-page body always splits)
% but PARTIAL (a shorter body that starts low on a page could still split). That
% residue is unreachable in practice — v3.37's everyline + the caption's
% \nopagebreak glue move such boxes WHOLE instead (probe N7: a 22-row table
% pushed down by 430pt moved whole, 0 segments), and the fleet has 0 split
% tableboxes, re-counted after this change. Reimplementing mdframed's own
% page-fitting arithmetic to close the residue would be a heuristic with its own
% failure mode, which is the trade this repo has repeatedly refused.
\newif\iflakintabcapseen
\newif\iflakintabframed
% The Lua lives in ONE \directlua at load time, and is written with NO
% backslash, %, # or ~ : \directlua expands its argument like \edef, so an
% active ~ would expand to \nobreakspace and a % would eat the rest of the line
% (the second is how the first draft of this block died). "not (a == b)" stands
% in for Lua's inequality operator for the same reason.
\directlua{
  lakintab = { }
  lakintab.HL = node.id("hlist")
  lakintab.VL = node.id("vlist")
  lakintab.RU = node.id("rule")
  lakintab.GL = node.id("glue")
  lakintab.KE = node.id("kern")
  lakintab.WH = node.id("whatsit")
  lakintab.TM = node.id("temp")
  lakintab.WD = node.subtype("pdf_dest")
  function lakintab.d(v)
    if v == nil then return 0 end
    if v < -1000000000 then return 0 end
    return v
  end
  function lakintab.scan(list, y, horiz, parent, target)
    local x = list
    while x do
      local i = x.id
      if i == lakintab.HL or i == lakintab.VL then
        if horiz then
          local b = y + x.shift
          if i == lakintab.HL then
            lakintab.scan(x.head, b, true, x, target)
          else
            lakintab.scan(x.head, b - lakintab.d(x.height), false, x, target)
          end
        else
          if i == lakintab.HL then
            lakintab.scan(x.head, y + lakintab.d(x.height), true, x, target)
          else
            lakintab.scan(x.head, y, false, x, target)
          end
          y = y + lakintab.d(x.height) + lakintab.d(x.depth)
        end
      elseif i == lakintab.WH then
        if x.subtype == lakintab.WD then
          local id = x.dest_id
          if type(id) == "string" and string.sub(id, 1, 4) == "tbl." then
            if horiz and parent and parent.width == 0
               and parent.height == 0 and parent.depth == 0 then
              parent.shift = parent.shift + (target - y)
            end
          end
        end
      elseif not horiz then
        if i == lakintab.RU then
          y = y + lakintab.d(x.height) + lakintab.d(x.depth)
        elseif i == lakintab.GL or i == lakintab.KE then
          y = y + lakintab.d(x.width)
        end
      end
      x = x.next
    end
  end
  function lakintab.fix(target, maxh)
    local nst = tex.nest[tex.nest.ptr]
    if nst == nil then return end
    local h = nst.head
    if h == nil then return end
    if h.id == lakintab.TM then h = h.next end
    local tot = 0
    local x = h
    while x do
      local i = x.id
      if i == lakintab.HL or i == lakintab.VL or i == lakintab.RU then
        tot = tot + lakintab.d(x.height) + lakintab.d(x.depth)
      elseif i == lakintab.GL or i == lakintab.KE then
        tot = tot + lakintab.d(x.width)
      end
      x = x.next
    end
    if tot > maxh then return end
    lakintab.scan(h, 0, false, nil, target)
  end
}
% @-FREE NAMES ON PURPOSE: these two are named in the `tablebox` key list at
% \newmdenv above, which is tokenised OUTSIDE any \makeatletter region — a
% \lakin@... there would tokenise as \lakin followed by the letters @tabopen.
\newcommand{\lakintabboxopen}{\lakintabframedtrue\global\lakintabcapseenfalse}
\newcommand{\lakintabboxclose}{%
  \iflakintabcapseen
    \directlua{lakintab.fix(-\number\dimexpr\lakintabcushion\relax,
                            \number\dimexpr\textheight\relax)}%
    \global\lakintabcapseenfalse
  \fi\lakintabframedfalse}
% The cushion above the box's CONTENT top. 42pt puts the dest 32.0pt above the
% frame top (42pt - innertopmargin 9pt - linewidth 0.9pt), i.e. exactly where
% the 295 correct dests already are — the correction moves only the broken ones.
\newcommand{\lakintabcushion}{42pt}
\newenvironment{qcirc}{\par\medskip\global\lakinfigraise 0pt\relax\begin{otherlanguage}{english}\begin{center}}
  {\csname lakin@grabfig\endcsname\end{center}\end{otherlanguage}\par\medskip}
\newenvironment{tikzpic}{\par\medskip\global\lakinfigraise 0pt\relax\begin{otherlanguage}{english}\begin{center}}
  {\csname lakin@grabfig\endcsname\end{center}\end{otherlanguage}\par\medskip}

% ============================================================================
% Code listings.  ★★ v3.4 THE BOX-OVERFLOW FIX:
%   resetmargins=true   → the listing IGNORES the enclosing environment's
%                         indentation and measures from the PAGE. v3.4 stated
%                         this backwards ("measured from the ENCLOSING box"),
%                         and v3.37 then blanked in-box line numbers because
%                         the page-anchored listing put its hanging digits on
%                         top of the box rule. Measured v4.5.0: with the key
%                         FALSE the listing follows the box and the code moves
%                         from x=75.26 to x=97.74, clearing the rule at 72.07.
%                         The key now lives in one \lstset below instead of in
%                         each style, so a box can override it — codebox does.
%   framexleftmargin=0  → the frame (left bar) does not extend left of the code.
%   xleftmargin keeps a small indent so the line numbers have room INSIDE.
% asmblock wraps the listing in an LTR island (code is LTR).
% ============================================================================
\lstdefinestyle{asm}{
  basicstyle=\ttfamily\small,
  backgroundcolor=\color{cBlueBg!55},
  frame=leftline,framerule=2.5pt,rulecolor=\color{cBlue},
  numbers=left,numberstyle=\tiny\color{cGray},numbersep=8pt,
  xleftmargin=2.2em,framexleftmargin=0pt,framexrightmargin=0pt,
  breaklines=true,columns=fullflexible,keepspaces=true,
  showstringspaces=false,aboveskip=8pt,belowskip=8pt,
  % v3.36: without a comment DELIMITER the commentstyle below never fires —
  % listings has no `language` here, so it does not know '#' opens a comment,
  % and every asm listing rendered flat black (measured: a whole-page pixel
  % sweep found ZERO cGreen glyphs across 135 shipped listings). morecomment
  % gives it the delimiter; color-only change, glyph widths untouched.
  morecomment={[l]{\#}},
  commentstyle=\color{cGreen},
}
\lstdefinestyle{cstyle}{
  language=C,basicstyle=\ttfamily\small,
  backgroundcolor=\color{cGreenBg!55},
  frame=leftline,framerule=2.5pt,rulecolor=\color{cGreen},
  numbers=left,numberstyle=\tiny\color{cGray},numbersep=8pt,
  xleftmargin=2.2em,framexleftmargin=0pt,framexrightmargin=0pt,
  breaklines=true,columns=fullflexible,keepspaces=true,
  showstringspaces=false,keywordstyle=\color{cBlue}\bfseries,
  commentstyle=\color{cGray}\itshape,aboveskip=8pt,belowskip=8pt,
}
% v3.36: pystyle/htmlstyle — colored twins of cstyle for the Python/HTML
% listings that used to repeat a ~200-char options blob with NO color
% (keywords rendered bold-black, comments italic-black — measured; C alone
% was colored). Geometry is copied VERBATIM from the shipped blob (frame,
% margins, skips) so swapping a call site to style=pystyle moves NOTHING —
% only keywordstyle/commentstyle/stringstyle are added (bold stays bold,
% italic stays italic, mono metrics identical => zero reflow, proven by
% per-page word-bbox diff). breaklines/breakatwhitespace stay per-site.
% v3.78: numbers=left joins both styles, with the exact numberstyle/numbersep
% keys asm/cstyle carry — fleet-wide consistency (author decision: all four
% styles numbered; "asm/cstyle numbered, py/html not" was an accident of the
% v3.36 geometry freeze, not a design). The numbers hang LEFT of the code box
% (listings sets them with \llap, so no overfull hbox can arise and the code
% x is untouched — measured: code frozen, digits at code_x-11.8pt); inside a
% colored box they are blanked like everyone else's (v3.37 hook), and the
% clean variant restores them in its frameless boxes (v3.78, above).
\lstdefinestyle{pystyle}{
  language=Python,basicstyle=\ttfamily\small,
  columns=fullflexible,keepspaces=true,
  frame=single,rulecolor=\color{cTblLine},
  numbers=left,numberstyle=\tiny\color{cGray},numbersep=8pt,
  xleftmargin=0.5em,framexleftmargin=0pt,framexrightmargin=0pt,
  showstringspaces=false,aboveskip=8pt,belowskip=8pt,
  keywordstyle=\color{cBlue}\bfseries,
  commentstyle=\color{cGray}\itshape,
  stringstyle=\color{cGreen},
}
% v4.5.0: hoisted out of the four styles so an enclosing box can override it.
% Every listing still resolves to true unless something later says otherwise,
% so shipped output is unchanged (verified byte-for-byte on all eight PDFs).
\lstset{resetmargins=true}
\lstdefinestyle{htmlstyle}{
  language=HTML,basicstyle=\ttfamily\small,
  columns=fullflexible,keepspaces=true,
  frame=single,rulecolor=\color{cTblLine},
  numbers=left,numberstyle=\tiny\color{cGray},numbersep=8pt,
  xleftmargin=0.5em,framexleftmargin=0pt,framexrightmargin=0pt,
  showstringspaces=false,aboveskip=8pt,belowskip=8pt,
  keywordstyle=\color{cBlue}\bfseries,
  commentstyle=\color{cGray}\itshape,
  stringstyle=\color{cGreen},
}
\newenvironment{asmblock}{\par\begin{otherlanguage}{english}}{\end{otherlanguage}\par}

% <hllpdf:listing-smart-quotes>
% v4.21.0: four Unicode punctuation characters must enter listings' token
% buffer. Unregistered characters bypass it: a closing quote was printed
% BEFORE the pending word, and an opening quote BEFORE a pending brace.
% Extend the standard character table, then activate only these codepoints
% in the listing's local character-table scope. Do this BEFORE the package's
% delimiter/alsoother/literate hooks so explicit user mappings still win.
% This is not universal Unicode support, nor a conversion to ASCII quotes.
% See listings developer guide 15.3.1 and references/listing-smart-quotes.md.
\makeatletter
\lst@CCPut \lst@ProcessOther {"2018}{"2019}{"201C}{"201D}\z@
\lst@AddToHookAtTop{SelectCharTable}{%
  \lst@ifactivechars
    \catcode"2018=\active \catcode"2019=\active
    \catcode"201C=\active \catcode"201D=\active
  \fi}
\makeatother
% </hllpdf:listing-smart-quotes>

% <hllpdf:listing-technical-symbols>
% v4.30: bounded extension for common technical punctuation in LTR listings.
% Without character-table registration these glyphs bypass the token buffer
% and can move before the preceding word even inside quoted strings.
% Keep the original four-quote block separate for its standing mutation test.
% This is source-faithful printing, not syntax validation or general Unicode.
\makeatletter
\lst@CCPut \lst@ProcessOther
  {"2013}{"2014}{"2026}{"2190}{"2192}{"2194}{"21D2}
  {"2212}{"221E}{"2260}{"2264}{"2265}\z@
\lst@AddToHookAtTop{SelectCharTable}{%
  \lst@ifactivechars
    \catcode"2013=\active \catcode"2014=\active \catcode"2026=\active
    \catcode"2190=\active \catcode"2192=\active \catcode"2194=\active
    \catcode"21D2=\active \catcode"2212=\active \catcode"221E=\active
    \catcode"2260=\active \catcode"2264=\active \catcode"2265=\active
  \fi}
\makeatother
% </hllpdf:listing-technical-symbols>


\usepackage[hidelinks,unicode]{hyperref}   % LAST, after colors. unicode = Hebrew bookmarks.
% v3.62: THE ISOLATION MACROS CORRUPTED EVERY PDF BOOKMARK THEY APPEARED IN.
% \en is defined (line ~155) as
%     \foreignlanguage{english}{\language=\csname l@nohyphenation\endcsname #1}
% and in a heading that body is written verbatim to the .toc and handed to
% hyperref's \pdfstringdef. hyperref strips the two CONTROL SEQUENCES (\language,
% \l@nohyphenation) but the plain `=` CHARACTER token between them is not a control
% sequence, so it survives into the PDF string: \chapter{שער \en{CNOT} והצמדה}
% produced the bookmark "שער =CNOT והצמדה". The printed TOC and the heading on the
% page are CORRECT — only the navigation pane and the viewer window title are wrong,
% which is why it went unnoticed. Measured on the shipped fleet: 70 corrupted
% bookmarks across 5 PDFs, from 177 unprotected \en-in-heading sites. Same corruption
% reached /Title and /Subject via the \hypersetup{pdftitle=...} the docs recommend.
% This hook tells hyperref to use the plain argument for all three isolation macros
% when building a PDF string. It does NOT touch the .toc (the printed TOC is
% byte-unchanged) and it does NOT mask the math signal — a heading with $...$ still
% logs "removing math shift", so BIDI-15 detection is intact.
% BONUS, measured: it also makes \code{} in an UNSTARRED heading COMPILE. BLD-10
% recorded that as needing \texorpdfstring or a starred heading; the real root cause
% is hyperref's \pdfstringdef scan of \code's \ifmmode...\else...\fi body — which is
% exactly why the identically-bodied \texttt survived (hyperref pre-neutralizes
% \texttt in PDF strings) and why \DeclareRobustCommand could never help.
% v3.64: \def\\{ } added. A manual line break in a heading WELDS the two words
% together in the bookmark: \chapter{שורה ראשונה\\שורה שנייה} produced the outline
% entry 'שורה ראשונהשורה שנייה' — "ראשונהשורה" is not a word. hyperref drops the
% \\ and puts nothing in its place, so the break that looks like a space on the
% page is a zero-width join in the navigation pane. The spaced form
% 'מילה אחת \\ מילה שנייה' escapes with a harmless DOUBLE space ('אחת  מילה'),
% which is why this never got noticed. Both are measured on a real outline.
% 0 live instances in the corpus — but SKILL.md and content-style.md both RECOMMEND
% a manual heading break, so this is a primed trap rather than a dormant one, and a
% one-token preamble fix retires it for good instead of asking authors to remember.
% Mapping \\ to a space fixes the WELD, which is the actual defect (a nonexistent
% word in the navigation pane). It does not normalise whitespace: measured after the
% fix, the welded form becomes 'שורה ראשונה שורה שנייה' (correct) and the already-
% spaced form goes from two spaces to three. That is cosmetic — viewers render it as
% a wider gap, and no word is corrupted. A per-heading \texorpdfstring is still the
% way to get an exactly-tuned bookmark string.
\pdfstringdefDisableCommands{\def\en#1{#1}\def\h#1{#1}\def\doclang#1{#1}\def\code#1{#1}\def\codeline#1{#1}\def\\{ }}
% v3.70: RESTORE \url LINE-BREAKING. url.sty freezes its break penalties at LOAD
% time (\mathchardef\UrlBreakPenalty=\binoppenalty), and v3.45's global
% \binoppenalty=10000 (inline-math atomicity) therefore priced EVERY breakpoint
% inside \Url at 10000 = never break. Measured: a long \url in Hebrew prose is
% one unbroken line, ~100pt overfull, invading the margin band the ink gates
% deliberately tolerate. With the two lines below: 0 overfull, the URL breaks
% after '/', the RTL split is CORRECT (fragment order forward), and each line
% gets its own link annotation exactly covering its fragment (measured).
% Confined to \Url's group — inline-math atomicity (TYP-7) is untouched.
\mathchardef\UrlBreakPenalty=200
\mathchardef\UrlBigBreakPenalty=100
\ifx\lakitmainlanguage\hllpdfenglish
  \hypersetup{pdflang=en}
\else
  \hypersetup{pdflang=he}
\fi  % /Lang is the document base locale (not per-span accessibility tagging).
% Set a per-document \hypersetup{pdftitle={...}} too — see SKILL.md.

% v3.52 (round-2 audit): two RTL bibliography/cross-ref defects, both invisible to
% every gate and both affecting the FIRST author to use the feature.
% (a) \@biblabel's default [#1] brackets are bidi-neutral and mirror to ]1[ under
% the RTL bibliography base — every \thebibliography reference list shows ]1[.
% LTR-isolating the label restores [1]. (The render gate's mirror rule fingerprints
% round parens )1(, so square brackets slipped through.)
\makeatletter
\renewcommand*{\@biblabel}[1]{\foreignlanguage{english}{[#1]}}
% v3.54: the SAME neutral-bracket trap on the INLINE side — a multi-key \cite{a,b}
% under the RTL base lays its European-number list right-to-left, so [1, 2] renders
% "[2,1]" and [1, 2, 3] renders "[3,2,1]" (verified by glyph geometry — the reader
% sees the sources in REVERSE order; a note form \cite[p.5]{a} also reorders to
% "[p. 5,1]"). LTR-isolate the whole bracket group (number list AND optional note),
% exactly like \@biblabel above: [1, 2] / [1, 2, 3] / [1, p. 5] all render correct.
% Placed AFTER hyperref so it wins; hyperref hyperlinks each number inside #1 (via
% \@citex), which this wrapper leaves intact. Mirrors the LaTeX kernel \@cite shape.
% v4.9.1: directional isolation must NOT switch the language here. Babel's
% safe-active citation wrapper passes its key list through \bbl@tempa;
% \foreignlanguage reuses that scratch macro before the deferred citation
% loop expands it. The old line therefore cited the literal key "english".
% If such a key existed, even the undefined-reference gate stayed green.
% \babelsublr isolates the bracket group without selecting a locale or
% changing that key list. Preserve native links, key order and optional notes.
\def\@cite#1#2{\babelsublr{[{#1\if@tempswa , #2\fi}]}}
\makeatother

% v4.11.0: a reference inside align's \text inherited HEBdisptag=true,
% so its enclosing hbox mirrored the brackets despite a clean compile.
% Isolate the complete printed reference without selecting a Babel language.
% Also neutralize the display-only normalfont direction reset locally. A
% direction-only patch inverted 12 to 21 inside equation in the measured probe.
% Plain maketag is local to the reference: mathtools' display-side wrapper
% otherwise globally consumes manual_tag while formatting a nested reference.
% Preserve each original command (including the selective auxiliary write),
% current tag forms, native links, and showonlyrefs on/off transitions.
\makeatletter
\NewCommandCopy{\HEB@nativeeqref}{\eqref}
\newcommand{\HEB@eqrefisolate}[1]{%
  \babelsublr{\HEBdisptagfalse\let\bbl@eqnodir\relax\let\maketag@@@\ORIGmaketag@@@#1}}
\protected\def\eqref#1{\HEB@eqrefisolate{\HEB@nativeeqref{#1}}}
\ExplSyntaxOn
\NewCommandCopy{\HEB@nativeMTeqref}{\MT_extended_eqref:n}
\cs_set_protected:Npn \MT_extended_eqref:n #1
  {\HEB@eqrefisolate{\HEB@nativeMTeqref{#1}}}
\ExplSyntaxOff
\makeatother

% v4.11.0: custom caption counters are not part of hyperref's English
% locale, and Hebrew has no upstream hyperref locale in this toolchain.
% Restore the appropriate names at every Babel language switch, including
% return from a nested inline island. Do not localize the destination keys.
\makeatletter
\newcommand{\HEB@refnames}{%
\def\chapterautorefname{פרק}%
\def\sectionautorefname{סעיף}%
\def\subsectionautorefname{תת־סעיף}%
\def\subsubsectionautorefname{תת־תת־סעיף}%
\def\paragraphautorefname{פסקה}%
% v4.31: hyperref uses a distinct counter name at the deepest standard level.
\def\subparagraphautorefname{תת־פסקה}%
\def\equationautorefname{משוואה}%
\def\figureautorefname{איור}%
\def\tableautorefname{טבלה}%
\def\partautorefname{חלק}%
\def\appendixautorefname{נספח}%
\def\footnoteautorefname{הערת שוליים}%
\def\itemautorefname{סעיף}%
\def\Itemautorefname{סעיף}%
\def\figautorefname{איור}%
\def\tblautorefname{טבלה}%
\def\pageautorefname{עמוד}% 
}
\addto\extrashebrew{\HEB@refnames}
\newcommand{\ENG@refnames}{%
  \def\subparagraphautorefname{subparagraph}%
  \def\chapterautorefname{chapter}%
  \def\sectionautorefname{section}%
  \def\subsectionautorefname{subsection}%
  \def\subsubsectionautorefname{subsubsection}%
  \def\paragraphautorefname{paragraph}%
  \def\equationautorefname{Equation}%
  \def\figureautorefname{Figure}%
  \def\tableautorefname{Table}%
  \def\partautorefname{Part}%
  \def\appendixautorefname{Appendix}%
  \def\footnoteautorefname{footnote}%
  \def\itemautorefname{item}%
  \def\figautorefname{Figure}%
  \def\tblautorefname{Table}%
  \def\Itemautorefname{item}%
  \def\pageautorefname{page}%
}
\addto\extrasenglish{\ENG@refnames}
\ifx\lakitmainlanguage\hllpdfenglish\ENG@refnames\else\HEB@refnames\fi
\makeatother

% ============================================================================
% Optional running head (v3.39): \lakithead{<book title>} — OFF by default.
% A document opts in with ONE line after \input{preamble.tex}:
%   \lakithead{מבנה המחשב}
% Every content page then carries the series' orientation line: chapter name
% at the reading edge, book title opposite, thin rule. (fancyhdr's [L]/[R]
% are PHYSICAL; under the RTL page direction [L] lands at the RIGHT/reading
% edge — measured, consistent across the corpus.) Packaged like the topic
% rail: always loaded, inert without the call. This centralizes the 8-line
% block that 8 lakits used to copy by hand — the same duplication vector
% that once let a lakit preamble go ~100 lines stale. Class-aware: report
% fills the mark from \chapter, article from \section. NB: a titlepage env
% self-empties; front matter that sets \pagestyle{empty} must restore with
% \pagestyle{fancy} afterwards (NOT plain — that discards the head).
% ============================================================================
\usepackage{fancyhdr}
% v3.61: \chapter* left a STALE running head — pages of an UNNUMBERED chapter
% carried the PREVIOUS chapter's title. report.cls routes \chapter* through
% \@schapter, which never calls \chaptermark, so \leftmark keeps whatever the last
% numbered chapter left. Measured: with \chapter{פרק ממוספר} then
% \chapter*{מבוא לא ממוספר}, page 4 (inside the starred chapter) printed
% "פרק ממוספר". Invisible on the chapter's FIRST page (\thispagestyle{plain}), so it
% only shows from the second page on — the pages an author is least likely to check.
% An unnumbered front-matter chapter (\chapter*{מבוא}, \chapter*{הקדמה}) is everyday
% Hebrew-academic structure, and the corpus already carries HAND-WRITTEN \markboth
% lines right after \chapter* — the exact per-document duplication \lakithead exists
% to retire. Patched INSIDE \lakithead so the macro stays inert when unused; a
% document's own later \markboth still wins, so those hand-patches do not regress.
\makeatletter
% <hllpdf:starred-heading-scope>
\newif\iflakin@starmarkpatched
\newcommand{\lakistarmark}{%
  \iflakin@starmarkpatched\else
  % Dispatch on the actual class counter, not hyperref's shared internals.
  \ifdefined\c@chapter
    \let\HEBorigschapter\@schapter
    % Record only after the native heading has reached its page. A mark
    % before the native routine can advertise the next heading too early.
    \def\@schapter##1{\HEBorigschapter{##1}\chaptermark{##1}}%
  \else
    % @ssect receives layout arguments, NOT the section name or level.
    % Capture the name in @startsection without introducing a TeX group:
    % grouping sectioning would break the native after-heading state.
    \def\HEB@sectionkind{}%
    \def\HEB@mainsection{section}%
    \let\HEBorigstartsection\@startsection
    \def\@startsection##1##2##3##4##5##6{%
      \def\HEB@sectionkind{##1}%
      \HEBorigstartsection{##1}{##2}{##3}{##4}{##5}{##6}}%
    \let\HEBorigssect\@ssect
    \def\@ssect##1##2##3##4##5{%
      \HEBorigssect{##1}{##2}{##3}{##4}{##5}%
      % Lower starred headings must not replace the article's section mark.
      \ifx\HEB@sectionkind\HEB@mainsection
        \sectionmark{##5}%
      \fi
      \ignorespaces}%
  \fi
  \global\lakin@starmarkpatchedtrue
  \fi}
% </hllpdf:starred-heading-scope>
% <hllpdf:running-language>
% Choose the wrapper when the mark is written, not when a later page ships.
% English punctuation must remain LTR even after returning to Hebrew prose.
\newcommand{\HEB@runningmark}[1]{%
  \iflanguage{english}{\markboth{\en{#1}}{}}{\markboth{\h{#1}}{}}}
% </hllpdf:running-language>
\makeatother
% v4.2.1 (NAV): every collection already names itself in \lakithead — feed the
% PDF /Title from the same declaration, so the navigation gate's
% PDF_REQUIRE_TITLE has a universal source and no fleet member ships
% title-less by missing a second declaration site. Measured before this: 14
% of 16 collections shipped /Title-empty while every one carried a perfectly
% good \lakithead string (plain Hebrew in all 16 — no math, no formatting, so
% the PDF-string path is clean). An explicit title always wins: set it with
% \lakitpdftitle{...} anywhere before \lakithead, or with a plain
% \hypersetup{pdftitle=...} — the fallback detects that via hyperref's
% \@pdftitle, probe-verified \@empty-equal until the first explicit set.
\makeatletter
\newif\iflakit@pdftitleset
\newcommand{\lakitpdftitle}[1]{%
  \hypersetup{pdftitle={#1}}\global\lakit@pdftitlesettrue}
\newcommand{\lakithead}[1]{%
  \iflakit@pdftitleset\else
    \ifx\@pdftitle\@empty
      \hypersetup{pdftitle={#1}}\global\lakit@pdftitlesettrue
    \fi
  \fi
  \pagestyle{fancy}\fancyhf{}%
  \lakistarmark

  \fancyhead[L]{\small\textcolor{cGray}{\leftmark}}%
  % <hllpdf:fixed-head-language>
  % The fixed document title has its OWN language snapshot, independent of
  % the current section/chapter mark. Choose the branch NOW, not at shipout:
  % the page may be Hebrew by the time an English title is finally printed.
  \iflanguage{english}{%
    \fancyhead[R]{\small\textcolor{cGray}{\en{#1}}}%
  }{%
    \fancyhead[R]{\small\textcolor{cGray}{\h{#1}}}%
  }%
  % </hllpdf:fixed-head-language>
  \fancyfoot[C]{\thepage}%
  \renewcommand{\headrulewidth}{0.4pt}%
  \ifdefined\chaptermark
    \renewcommand{\chaptermark}[1]{\HEB@runningmark{##1}}%
  \else
    \renewcommand{\sectionmark}[1]{\HEB@runningmark{##1}}%
  \fi}
\makeatother

% v3.39.1: path-independent byte-reproducibility. LuaTeX derives the trailer
% /ID from the OUTPUT PATH, so two clones of the same source at different
% directories produced different sha256 (measured: the ONLY delta was the
% 62-byte /ID line). Bit 512 suppresses exactly that line — measured: Creator/
% Producer/CreationDate/ModDate and all content untouched. /ID is optional
% per ISO 32000 (required only for ENCRYPTED PDFs and PDF/A — neither is a
% target here). If a future workflow needs an ID (PDF/A, Acrobat revision
% matching), delete this line and it returns.
\pdfvariable suppressoptionalinfo \numexpr 512\relax

% v3.75: BOLD MATH IN BOLD HEADINGS. Heading text is \bfseries but math inside
% it stayed regular-weight CM — a hairline $V_T$ inside a heavy 24.8pt chapter
% title (measured: CMMI12 fl=6 next to FrankRuhlLibre-Bold fl=20; 155 live
% heading-with-math sites fleet-wide). \boldmath rides along wherever the class
% sets \bfseries for a heading; A/B measured: $V_T$ -> CMMIB10 fl=22, 0 missing
% characters (the bold CM faces ship with this TeX Live). Guarded per-macro so
% the article class (no \@makechapterhead) and future class changes stay safe.
% v4.0.1 (adoption fix): the SEARCH patterns above carry no parameter token.
% v4.0.0 made them specific by including the class's own #1/#2, and measured on
% TeX Live 2023 that never matches: a \patchcmd search pattern containing a
% parameter token cannot match while the call sits inside \@ifundefined{...}{}{...}
% (probed both #1 and ##1 -> FAIL; the same pattern at top level succeeds). Every
% one of these patches therefore failed, and because v4.0.0 also replaced the
% empty failure branch with \PackageError, the failure became FATAL: all four
% bundled example assets stopped building ("Cannot patch chapter-title math
% weight -> Fatal error, no output PDF"). The \PackageError branches are KEPT —
% a silent no-op here is what let v3.75's fix rot unnoticed — only the parameter
% tokens are dropped, so the patches match and a real class change still errors.
\makeatletter
\@ifundefined{@makechapterhead}{}{%
  \patchcmd{\@makechapterhead}{\huge\bfseries \@chapapp}
    {\huge\bfseries\boldmath \@chapapp}{}
    {\PackageError{hebrew-lualatex-pdf}{Cannot patch chapter-number math weight}{Unsupported document class}}%
  \patchcmd{\@makechapterhead}{\Huge \bfseries}
    {\Huge \bfseries\boldmath}{}
    {\PackageError{hebrew-lualatex-pdf}{Cannot patch chapter-title math weight}{Unsupported document class}}%
  \patchcmd{\@makeschapterhead}{\Huge \bfseries}
    {\Huge \bfseries\boldmath}{}
    {\PackageError{hebrew-lualatex-pdf}{Cannot patch starred chapter math weight}{Unsupported document class}}}
% v4.0.1: \part sits under a THREE-DEEP wrapper chain, so the class-based anchor
% is not in \@part at all. Measured with \show: hyperref rewrites \@part to
% `...\phantomsection\fi \H@old@part`, and nameref rewrites THAT to
% `\NR@gettitle{#1}\NR@part[{#1}]{#2}` — the class body ends up in \NR@part.
% v4.0.0 patched \@part and made failure fatal, which is why all seven bundled
% examples stopped building. \part is NOT a dead heading level here: 188 live uses
% across the fleet, so the bold-math fix is worth getting right rather than
% dropping. We walk the chain innermost-first and stop at the first hit. Failure is
% a WARNING, not an error: the chain is a property of hyperref/nameref versions, so
% a future package release must degrade to regular-weight part math — a cosmetic
% loss — instead of refusing to build. Chapter and section keep \PackageError,
% because those targets are in the CLASS and were verified to patch here.
\newif\iflakin@partnum \newif\iflakin@parttitle \newif\iflakin@sparttitle
\def\lakin@partnum#1#2{\iflakin@partnum\else
  \patchcmd{#1}{#2\bfseries \partname}{#2\bfseries\boldmath \partname}
    {\lakin@partnumtrue}{}\fi}
\def\lakin@parttitle#1#2{\iflakin@parttitle\else
  \patchcmd{#1}{#2 \bfseries}{#2 \bfseries\boldmath}{\lakin@parttitletrue}{}\fi}
\def\lakin@sparttitle#1#2{\iflakin@sparttitle\else
  \patchcmd{#1}{#2 \bfseries}{#2 \bfseries\boldmath}{\lakin@sparttitletrue}{}\fi}
\@ifundefined{c@chapter}{\def\lakin@pnum{\Large}\def\lakin@ptit{\huge}}%
                        {\def\lakin@pnum{\huge}\def\lakin@ptit{\Huge}}
\@ifundefined{NR@part}{}{\expandafter\lakin@partnum\expandafter\NR@part\expandafter{\lakin@pnum}%
                         \expandafter\lakin@parttitle\expandafter\NR@part\expandafter{\lakin@ptit}}
\@ifundefined{H@old@part}{}{\expandafter\lakin@partnum\expandafter\H@old@part\expandafter{\lakin@pnum}%
                            \expandafter\lakin@parttitle\expandafter\H@old@part\expandafter{\lakin@ptit}}
\@ifundefined{@part}{}{\expandafter\lakin@partnum\expandafter\@part\expandafter{\lakin@pnum}%
                       \expandafter\lakin@parttitle\expandafter\@part\expandafter{\lakin@ptit}}
\@ifundefined{NR@spart}{}{\expandafter\lakin@sparttitle\expandafter\NR@spart\expandafter{\lakin@ptit}}
\@ifundefined{H@old@spart}{}{\expandafter\lakin@sparttitle\expandafter\H@old@spart\expandafter{\lakin@ptit}}
\@ifundefined{@spart}{}{\expandafter\lakin@sparttitle\expandafter\@spart\expandafter{\lakin@ptit}}
\iflakin@partnum\else\PackageWarning{hebrew-lualatex-pdf}{Part-number math stays
  regular weight: no known \noexpand\part wrapper matched (hyperref/nameref
  version change). Cosmetic only}\fi
\iflakin@parttitle\else\PackageWarning{hebrew-lualatex-pdf}{Part-title math stays
  regular weight: no known \noexpand\part wrapper matched. Cosmetic only}\fi
% v4.12.0: THE PART NUMBER IN THE TOC WAS NOT DIRECTIONALLY SEPARATED FROM THE
% TITLE. report.cls/article.cls compose the part entry as ONE text stream --
% \addcontentsline{toc}{part}{\thepart\hspace{1em}#1} -- while EVERY other level
% goes through \numberline, which puts the number in an \hbox. A box is a bidi
% OBJECT and cannot join a neighbouring run; a bare "IV" can. So in the RTL TOC
% paragraph a Latin number and a Latin-INITIAL title merge into a SINGLE LTR run
% and are laid out left-to-right as a unit: the number lands on the wrong side of
% the row, and a mixed title is scrambled outright. Measured on \part{MCP ...}:
% the printed order was "<tail> V MCP" -- title fragment, number, title fragment.
% The chapter/section rows one line below, with the SAME Latin-initial titles,
% were correct on the same page. That contrast is the whole diagnosis.
% \en{} IS NOT A WORKAROUND, measured both ways on one fixture: it fixes a
% PURE-Latin title (\part{\en{Hooks}} renders correctly) and does NOT fix a MIXED
% one (\part{\en{MCP} <hebrew>} stays scrambled). babel's language switch changes
% font and hyphenation, not the bidi run structure -- so the defect is not
% reachable from the source side at all, which is why it belongs here.
% Boxing the number is the entire fix and is a STRICT NO-OP where the bug never
% bit: a Hebrew-titled document compiles BYTE-IDENTICAL before and after (md5,
% 3 clean runs), because there the neutral glue already sat between an L run and
% an R run. The PDF outline is byte-identical too -- hyperref's \pdfstringdef
% neutralises \mbox, so bookmark strings keep their part numbers.
% \numberline{\thepart} would also fix the direction, but it re-geometries every
% part row in the fleet and DROPS the number from the part bookmark (hyperref
% strips \numberline). \mbox changes nothing that was already right.
% The search pattern deliberately carries no # token: etoolbox cannot match a
% parameter token in a patch pattern (measured -- all three chain members
% reported "not patchable" while the pattern included #1). "\thepart \hspace
% {1em}" is brace-balanced and occurs exactly once in the wrapper body.
\newif\iflakin@parttocnum
\def\lakin@parttocnum#1{\iflakin@parttocnum\else
  \patchcmd{#1}{\thepart \hspace {1em}}{\mbox {\thepart }\hspace {1em}}
    {\lakin@parttocnumtrue}{}\fi}
\@ifundefined{NR@part}{}{\lakin@parttocnum\NR@part}
\@ifundefined{H@old@part}{}{\lakin@parttocnum\H@old@part}
\@ifundefined{@part}{}{\lakin@parttocnum\@part}
\iflakin@parttocnum\else\PackageWarning{hebrew-lualatex-pdf}{Part TOC number is
  not boxed: no known \noexpand\part wrapper matched. A part title that BEGINS
  with Latin will print its number on the wrong side of its TOC row}\fi
\patchcmd{\section}{\bfseries}{\bfseries\boldmath}{}
  {\PackageError{hebrew-lualatex-pdf}{Cannot patch section math weight}{Unsupported document class}}
\patchcmd{\subsection}{\bfseries}{\bfseries\boldmath}{}
  {\PackageError{hebrew-lualatex-pdf}{Cannot patch subsection math weight}{Unsupported document class}}
\patchcmd{\subsubsection}{\bfseries}{\bfseries\boldmath}{}
  {\PackageError{hebrew-lualatex-pdf}{Cannot patch subsubsection math weight}{Unsupported document class}}
\patchcmd{\paragraph}{\bfseries}{\bfseries\boldmath}{}{}
\patchcmd{\subparagraph}{\bfseries}{\bfseries\boldmath}{}{}
\@ifundefined{l@chapter}{}{%
  \patchcmd{\l@chapter}{\bfseries}{\bfseries\boldmath}{}{}}
\makeatother
\setcounter{tocdepth}{\lakittocdepth}
% <hllpdf:measured-toc>
% v4.16.0: Measure actual number labels, not a guessed digit count. Read only
% canonical contentsline records; never execute titles or addtocontents code
% during the measuring pass. Widths are local to this individual list and are
% available before its first row: no per-row drift and no extra aux pass.
\makeatletter
\newread\hllpdf@tocstream
\newdimen\hllpdf@toc@chapter
\newdimen\hllpdf@toc@section
\newdimen\hllpdf@toc@subsection
\newdimen\hllpdf@toc@subsubsection
\newdimen\hllpdf@toc@paragraph
\newdimen\hllpdf@toc@subparagraph
\newdimen\hllpdf@toc@figure
\newdimen\hllpdf@toc@table
\newdimen\hllpdf@tocpage
\newdimen\hllpdf@toc@sectionindent
\newdimen\hllpdf@toc@subsectionindent
\newdimen\hllpdf@toc@subsubsectionindent
\newdimen\hllpdf@toc@paragraphindent
\newdimen\hllpdf@toc@subparagraphindent
\def\hllpdf@tocdefaults{%
  \hllpdf@toc@chapter=2em
  \hllpdf@toc@section=3.2em
  \hllpdf@toc@subsection=3.9em
  \@ifundefined{c@chapter}{%
    \hllpdf@toc@subsubsection=3.2em
    \hllpdf@toc@paragraph=4.1em
    \hllpdf@toc@subparagraph=5em
  }{%
    \hllpdf@toc@subsubsection=4.1em
    \hllpdf@toc@paragraph=5em
    \hllpdf@toc@subparagraph=6em
  }%
  \hllpdf@toc@figure=2.3em
  \hllpdf@toc@table=2.3em
  \hllpdf@tocpage=\@pnumwidth\relax
}
\def\hllpdf@tocindents{%
  % Preserve the established shallow layout when its slots already fit.
  \hllpdf@toc@sectionindent=1.5em
  \@tempdima=\hllpdf@toc@chapter \advance\@tempdima by -2em
  \advance\hllpdf@toc@sectionindent by \@tempdima
  \hllpdf@toc@subsectionindent=\dimexpr
    \hllpdf@toc@sectionindent+\hllpdf@toc@section\relax
  \hllpdf@toc@subsubsectionindent=\dimexpr
    \hllpdf@toc@subsectionindent+\hllpdf@toc@subsection\relax
  \hllpdf@toc@paragraphindent=\dimexpr
    \hllpdf@toc@subsubsectionindent+\hllpdf@toc@subsubsection\relax
  \hllpdf@toc@subparagraphindent=\dimexpr
    \hllpdf@toc@paragraphindent+\hllpdf@toc@paragraph\relax
}
\def\hllpdf@tocfont@chapter{\bfseries\boldmath}
\def\hllpdf@tocfont@part{\large\bfseries\boldmath}
\def\hllpdf@tocentryfont{%
  \ifcsname hllpdf@tocfont@\hllpdf@toclevel\endcsname
    \csname hllpdf@tocfont@\hllpdf@toclevel\endcsname
  \fi}
\long\def\hllpdf@tocdiscard#1\hllpdf@tocstop{}
\def\hllpdf@tocscanrecord{%
  \@ifnextchar\contentsline{\hllpdf@tocscanline}{%
    \@ifnextchar\babel@toc{\hllpdf@tocscanlanguage}{\hllpdf@tocdiscard}}}
\long\def\hllpdf@tocscanlanguage\babel@toc#1#2#3\hllpdf@tocstop{%
  % Babel records are declarations, not arbitrary code. Track the language
  % without executing the original directive or changing the surrounding list.
  \def\hllpdf@toclanguage{#1}}
\long\def\hllpdf@tocscanline\contentsline#1#2#3#4#5\hllpdf@tocstop{%
  \def\hllpdf@toclevel{#1}%
  \ifcsname hllpdf@toc@#1\endcsname
    \hllpdf@tocscantitle#2\hllpdf@tocstop
  \fi
  % The page-number column needs its own measurement (including bold chapter
  % rows). Its extra clearance separates the page label from dot leaders.
  \setbox\@tempboxa=\hbox{\foreignlanguage{\hllpdf@toclanguage}{\hllpdf@tocentryfont #3\kern.3em}}%
  \ifdim\wd\@tempboxa>\hllpdf@tocpage
    \hllpdf@tocpage=\wd\@tempboxa
  \fi}
\def\hllpdf@tocscantitle{%
  \@ifnextchar\numberline{\hllpdf@tocscanlabel}{\hllpdf@tocdiscard}}
\long\def\hllpdf@tocscanlabel\numberline#1#2\hllpdf@tocstop{%
  \setbox\@tempboxa=\hbox{\foreignlanguage{\hllpdf@toclanguage}{\hllpdf@tocentryfont #1\kern.5em}}%
  \ifdim\wd\@tempboxa>\csname hllpdf@toc@\hllpdf@toclevel\endcsname
    \csname hllpdf@toc@\hllpdf@toclevel\endcsname=\wd\@tempboxa
  \fi}
\def\hllpdf@tocread{%
  \ifeof\hllpdf@tocstream\else
    \read\hllpdf@tocstream to\hllpdf@tocrecord
    \expandafter\hllpdf@tocscanrecord\hllpdf@tocrecord\hllpdf@tocstop
    \expandafter\hllpdf@tocread
  \fi}
\def\hllpdf@tocmeasure#1{%
  \makeatletter
  \edef\hllpdf@toclanguage{\languagename}%
  \hllpdf@tocdefaults
  \openin\hllpdf@tocstream=\jobname.#1\relax
  \hllpdf@tocread
  \closein\hllpdf@tocstream
  \hllpdf@tocindents
  \edef\@pnumwidth{\the\hllpdf@tocpage}%
  \@tempdima=\dimexpr\hllpdf@tocpage+1em\relax
  \ifdim\@tempdima>\@tocrmarg\relax
    \edef\@tocrmarg{\the\@tempdima}%
  \fi}
\let\hllpdf@originalstarttoc\@starttoc
\def\@starttoc#1{%
  \begingroup
    \def\hllpdf@tocext{#1}%
    \def\hllpdf@tocwanted{toc}%
    \ifx\hllpdf@tocext\hllpdf@tocwanted \hllpdf@tocmeasure{#1}\else
      \def\hllpdf@tocwanted{lof}%
      \ifx\hllpdf@tocext\hllpdf@tocwanted \hllpdf@tocmeasure{#1}\else
        \def\hllpdf@tocwanted{lot}%
        \ifx\hllpdf@tocext\hllpdf@tocwanted \hllpdf@tocmeasure{#1}\fi
      \fi
    \fi
    \hllpdf@originalstarttoc{#1}%
  \endgroup}
\hllpdf@tocdefaults
\hllpdf@tocindents
\@ifundefined{l@chapter}{}{%
  \patchcmd{\l@chapter}{\setlength\@tempdima{1.5em}}%
    {\setlength\@tempdima{\hllpdf@toc@chapter}}{}%
    {\PackageError{hebrew-lualatex-pdf}{Cannot size chapter TOC numbers}%
       {Unsupported contents layout: the chapter number slot was not found}}}
\renewcommand*\l@section{\@dottedtocline{1}%
  {\hllpdf@toc@sectionindent}{\hllpdf@toc@section}}
\renewcommand*\l@subsection{\@dottedtocline{2}%
  {\hllpdf@toc@subsectionindent}{\hllpdf@toc@subsection}}
\renewcommand*\l@subsubsection{\@dottedtocline{3}%
  {\hllpdf@toc@subsubsectionindent}{\hllpdf@toc@subsubsection}}
\renewcommand*\l@paragraph{\@dottedtocline{4}%
  {\hllpdf@toc@paragraphindent}{\hllpdf@toc@paragraph}}
\renewcommand*\l@subparagraph{\@dottedtocline{5}%
  {\hllpdf@toc@subparagraphindent}{\hllpdf@toc@subparagraph}}
\renewcommand*\l@figure{\@dottedtocline{1}{1.5em}{\hllpdf@toc@figure}}
\renewcommand*\l@table{\@dottedtocline{1}{1.5em}{\hllpdf@toc@table}}
\makeatother
% </hllpdf:measured-toc>
\setcounter{secnumdepth}{\lakitsecnumdepth}
% Public preamble-only caption-numbering policy. The historical within-major
% default is untouched unless this command is called. "continuous" removes the
% reset edge as well as the printed prefix, so Figure 1/2 and Table 1/2 stay
% unique across sections/chapters. Returning to "within-major" restores the
% original class-aware displays and native destination conventions. The policy
% never replaces \figcap/\tabcap or their measured link-placement mechanism.
\makeatletter
\let\lakin@majorfig\thefig
\let\lakin@majortbl\thetbl
\let\lakin@majorHfig\theHfig
\let\lakin@majorHtbl\theHtbl
\newcommand{\lakitcaptionnumbering}[1]{%
  \edef\lakin@numbering{#1}%
  \def\lakin@continuous{continuous}\def\lakin@withinmajor{within-major}%
  \ifx\lakin@numbering\lakin@continuous
    \@ifundefined{c@chapter}{%
      \@removefromreset{fig}{section}\@removefromreset{tbl}{section}%
    }{\@removefromreset{fig}{chapter}\@removefromreset{tbl}{chapter}}%
    \renewcommand{\thefig}{\arabic{fig}}\renewcommand{\thetbl}{\arabic{tbl}}%
    \def\theHfig{\arabic{fig}}\def\theHtbl{\arabic{tbl}}%
  \else\ifx\lakin@numbering\lakin@withinmajor
    \@ifundefined{c@chapter}{%
      \@removefromreset{fig}{section}\@removefromreset{tbl}{section}%
      \@addtoreset{fig}{section}\@addtoreset{tbl}{section}%
    }{%
      \@removefromreset{fig}{chapter}\@removefromreset{tbl}{chapter}%
      \@addtoreset{fig}{chapter}\@addtoreset{tbl}{chapter}}%
    \let\thefig\lakin@majorfig\let\thetbl\lakin@majortbl
    \let\theHfig\lakin@majorHfig\let\theHtbl\lakin@majorHtbl
  \else
    \PackageError{hebrew-english-latex-studio}{Unknown caption-numbering policy '#1'}%
      {Use continuous or within-major before begin{document}.}%
  \fi\fi}
\@onlypreamble\lakitcaptionnumbering
\makeatother
% Architecture readiness sentinel. Optional runtime layers require the complete
% core, so the marker is defined only after every mechanism/baseline definition.
% <hllpdf:architecture-ready>
\def\hllpdfcoreloaded{1}
% </hllpdf:architecture-ready>

% v4.22: registered file resources, not an assertion about embedded subset bytes.
% The log framing avoids long-path wrapping and preserves UTF-8 as hex.
% <hllpdf:font-inventory>
\directlua{
  local function inventory()
    local names = {}
    local unnamed = 0
    for id, item in font.each() do
      local path = item.filename
      if path and not (path == "") then names[path] = true else unnamed = unnamed + 1 end
    end
    local sorted = {}
    for path in pairs(names) do table.insert(sorted,path) end
    table.sort(sorted)
    texio.write_nl("log","HLLFONT|BEGIN|1")
    for index,path in ipairs(sorted) do
      local hex = path:gsub(".",function(c) return string.format(string.char(37) .. "02x",string.byte(c)) end)
      local count = math.ceil(string.len(hex)/48)
      texio.write_nl("log","HLLFONT|FILE|" .. index .. "|" .. count)
      for offset=1,string.len(hex),48 do texio.write_nl("log","HLLFONT|HEX|" .. hex:sub(offset,offset+47)) end
    end
    texio.write_nl("log","HLLFONT|END|" .. rawlen(sorted) .. "|" .. unnamed)
    texio.write_nl("log","")
  end
  luatexbase.add_to_callback("finish_pdffile",inventory,"hllpdf.runtime.fonts")
}
% </hllpdf:font-inventory>

%
% Supported class contract: report or article.  Alternate paper sizes must also
% be declared to build.sh through EXPECTED_PAPER; the PDF gate never infers them.
% Do NOT use these hooks as a theme mechanism: theme overlays must not own page
% geometry, class selection, or direction semantics. A declared presentation-policy
% overlay may intentionally change numbering visibility after the preamble.
% ============================================================================
% Runtime architecture guard: generated from scripts/architecture_contract.py.
% The readiness sentinel remains delayed until the END of this file.
% <hllpdf:architecture-guard>
\ifdefined\hllpdfcoreloaded
  \errmessage{hllpdf: preamble.tex loaded more than once}
  \expandafter\endinput
\fi
% </hllpdf:architecture-guard>

\ifdefined\lakitdocumentclass\else\def\lakitdocumentclass{report}\fi
\ifdefined\lakitclassoptions\else\def\lakitclassoptions{12pt,a4paper}\fi
\ifdefined\lakitgeometryoptions\else\def\lakitgeometryoptions{margin=2.3cm,headheight=14pt}\fi
\ifdefined\lakitlinespread\else\def\lakitlinespread{1.3}\fi
\ifdefined\lakittocdepth\else\def\lakittocdepth{1}\fi
\ifdefined\lakitsecnumdepth\else\def\lakitsecnumdepth{2}\fi
% Paragraph organization belongs to the document layout, independently of its
% language and theme. Preserve historical class dimensions unless explicitly
% selecting block paragraphs; never suppress indentation with an everypar hack.
\ifdefined\lakitparagraphstyle\else\def\lakitparagraphstyle{indented}\fi
\edef\lakitparagraphstyle{\lakitparagraphstyle}
\def\hllpdfparagraphblock{block}
\def\hllpdfparagraphindented{indented}
\ifx\lakitparagraphstyle\hllpdfparagraphblock\else
  \ifx\lakitparagraphstyle\hllpdfparagraphindented\else
    \errmessage{hllpdf: paragraph style must be block or indented}
  \fi
\fi
% v4.3.1: the public class hook has a deliberately narrow support contract.
% Before this guard, \def\lakitdocumentclass{book} silently compiled even though
% every heading/TOC regression test targets report/article; an unsupported shell
% could therefore look "supported" until a subtle class-specific defect shipped.
% Fail before \documentclass rather than letting an undocumented combination enter
% the mechanism core. Plain-TeX comparison is used because LaTeX packages are not
% available yet at this point.
\edef\lakitvalidatedclass{\lakitdocumentclass}
\def\lakitreportclass{report}
\def\lakitarticleclass{article}
\ifx\lakitvalidatedclass\lakitreportclass
\else\ifx\lakitvalidatedclass\lakitarticleclass
\else
  \errmessage{hllpdf: unsupported class \lakitdocumentclass; use report or article}
\fi\fi
\let\lakitvalidatedclass\relax
\let\lakitreportclass\relax
\let\lakitarticleclass\relax
% Language intent is independent of document class, theme and content domain.
% The two dominant-language profiles differ in authoring intent, not direction.
% A balanced bilingual document must explicitly choose a navigation/base locale.
% <hllpdf:language-profile>
\providecommand{\lakitlanguageprofile}{hebrew}
\edef\lakitlanguageprofile{\lakitlanguageprofile}
\def\hllpdfhebrew{hebrew}\def\hllpdfenglish{english}
\def\hllpdfhebrewmixed{hebrew-with-english}
\def\hllpdfenglishmixed{english-with-hebrew}
\def\hllpdfbilingual{bilingual}
\let\hllpdfexpectedlanguage\relax
\ifx\lakitlanguageprofile\hllpdfhebrew
  \let\hllpdfexpectedlanguage\hllpdfhebrew
\else\ifx\lakitlanguageprofile\hllpdfhebrewmixed
  \let\hllpdfexpectedlanguage\hllpdfhebrew
\else\ifx\lakitlanguageprofile\hllpdfenglish
  \let\hllpdfexpectedlanguage\hllpdfenglish
\else\ifx\lakitlanguageprofile\hllpdfenglishmixed
  \let\hllpdfexpectedlanguage\hllpdfenglish
\else\ifx\lakitlanguageprofile\hllpdfbilingual
  \ifdefined\lakitmainlanguage\else
    \errmessage{hllpdf: bilingual profile requires explicit lakitmainlanguage}
  \fi
\else
  \errmessage{hllpdf: unsupported language profile \lakitlanguageprofile}
\fi\fi\fi\fi\fi
\ifdefined\lakitmainlanguage
  \edef\lakitmainlanguage{\lakitmainlanguage}
\else
  \let\lakitmainlanguage\hllpdfexpectedlanguage
\fi
\ifx\lakitmainlanguage\hllpdfhebrew\else
  \ifx\lakitmainlanguage\hllpdfenglish\else
    \errmessage{hllpdf: main language must be hebrew or english}
  \fi
\fi
\ifx\hllpdfexpectedlanguage\relax\else
  \ifx\lakitmainlanguage\hllpdfexpectedlanguage\else
    \errmessage{hllpdf: main language conflicts with language profile}
  \fi
\fi
\typeout{HLLLANG|\lakitlanguageprofile|\lakitmainlanguage}
% </hllpdf:language-profile>
\documentclass[\lakitclassoptions]{\lakitdocumentclass}
\usepackage[no-math]{fontspec}        % (1) ★ no-math — THE critical line (see recipe.md)
% layout=lists asks Babel to adapt list parshape when paragraph and page
% directions differ. Without it, minority-language list bodies lose their
% indentation and labels hang outside the text block, despite a clean log.
% The upstream mechanism preserves ordinary same-direction list geometry.
\usepackage[bidi=basic,layout=lists]{babel}
\ifx\lakitmainlanguage\hllpdfenglish
  \babelprovide[import]{hebrew}
  \babelprovide[import,main]{english}
\else
  \babelprovide[import,main]{hebrew}
  \babelprovide[import]{english}
\fi
% v4.22: bind canonical families and fallbacks to the selected installation.
% Lookup by basename or family alone is not proof of the file actually used.
% <hllpdf:font-binding>
\directlua{
  local dir = os.getenv("HEBREW_LATEX_FONT_DIR")
  if not dir or dir == "" then
    local data = os.getenv("XDG_DATA_HOME")
    if not data or data == "" or not (data:sub(1,1) == "/") then
      data = (os.getenv("HOME") or "") .. "/.local/share"
    end
    dir = data .. "/fonts/hebrew-lualatex-pdf"
  end
  if not (dir:sub(1,1) == "/") then dir = lfs.currentdir() .. "/" .. dir end
  if dir:find("[",1,true) or dir:find("]",1,true) or dir:find(string.char(0),1,true)
     or dir:find(string.char(10),1,true) or dir:find(string.char(13),1,true) then
    tex.error("hllpdf: unsupported character in font directory")
  end
  if not (lfs.attributes(dir,"mode") == "directory") then
    tex.error("hllpdf: font directory is missing")
  end
  hllpdf_font_dir = dir:gsub("/+$", "") .. "/"
  tex.sprint(string.char(92) .. "def" .. string.char(92) .. "hllpdffontdir{")
  tex.sprint(-2,hllpdf_font_dir)
  tex.sprint("}")
  luaotfload.add_fallback("hebfb",{"[" .. hllpdf_font_dir .. "DejaVuSerif.ttf]:mode=node;"})
  luaotfload.add_fallback("hebfbbf",{"[" .. hllpdf_font_dir .. "DejaVuSerif-Bold.ttf]:mode=node;"})
}
% </hllpdf:font-binding>
% v3.22: Frank Ruhl Libre ships NO italic face. Instead of silently-upright
% \emph (v<=3.20) or bold-as-emphasis (v3.21), synthesize a slanted face with
% FakeSlant — real, VISIBLE italic emphasis with correct \emph nesting
% semantics (inner \emph toggles back upright). Verified empirically at
% 300dpi: slant clearly visible, 0 errors, 0 missing glyphs.
\babelfont{rm}[Path={\hllpdffontdir},RawFeature={fallback=hebfb},
  BoldFont={FrankRuhlLibre-Bold.ttf},
  BoldFeatures={RawFeature={fallback=hebfbbf}},   % v3.56 — see hebfbbf above
  ItalicFont={FrankRuhlLibre-Regular.ttf}, ItalicFeatures={FakeSlant=0.22},
  % Bold-italic uses the weight-matched fallback and a synthetic slant.
  BoldItalicFont={FrankRuhlLibre-Bold.ttf}, BoldItalicFeatures={RawFeature={fallback=hebfbbf},FakeSlant=0.22}
  ]{FrankRuhlLibre-Regular.ttf}
% Sans fallbacks preserve family and weight for the few missing codepoints.
\directlua{luaotfload.add_fallback("hebfbsans",{"[" .. hllpdf_font_dir .. "DejaVuSans.ttf]:mode=node;"})}
\directlua{luaotfload.add_fallback("hebfbsansbf",{"[" .. hllpdf_font_dir .. "DejaVuSans-Bold.ttf]:mode=node;"})}
\babelfont{sf}[Path={\hllpdffontdir},RawFeature={fallback=hebfbsans},
  BoldFont={Heebo-Bold.ttf},
  BoldFeatures={RawFeature={fallback=hebfbsansbf}}]{Heebo-Regular.ttf}
% v3.24: Scale=MatchLowercase — DejaVu Sans Mono has an outsized x-height, so
% unscaled inline code visually SHOUTS next to the delicate Frank Ruhl body
% (measured on a 197-page build: scaled, a code atom sits in the line at the
% body's weight, and the document tightened by ~5%). Applies to listings too.
\babelfont{tt}[Path={\hllpdffontdir},RawFeature={fallback=hebfb},Scale=MatchLowercase,
  BoldFont={DejaVuSansMono.ttf},
  BoldFeatures={RawFeature={fallback=hebfb},FakeBold=1.5},
  ItalicFont={DejaVuSansMono.ttf},
  ItalicFeatures={RawFeature={fallback=hebfb},FakeSlant=0.2},
  BoldItalicFont={DejaVuSansMono.ttf},
  BoldItalicFeatures={RawFeature={fallback=hebfb},FakeBold=1.5,FakeSlant=0.2}
  ]{DejaVuSansMono.ttf}

% Explicit multilingual/English profiles use real Latin faces and italics.
% The legacy/default Hebrew profile retains its established measured typography.
% Latin Modern is provided by the TeX distribution, not the pinned font bundle;
% recorder/runtime evidence reports it as an additional toolchain resource.
\ifx\lakitlanguageprofile\hllpdfhebrew\else
  \babelfont[english]{rm}[Ligatures=TeX,RawFeature={fallback=hebfb},
    BoldFeatures={RawFeature={fallback=hebfbbf}},
    BoldItalicFeatures={RawFeature={fallback=hebfbbf}}]{Latin Modern Roman}
  \babelfont[english]{sf}[Ligatures=TeX,RawFeature={fallback=hebfbsans},
    BoldFeatures={RawFeature={fallback=hebfbsansbf}},
    BoldItalicFeatures={RawFeature={fallback=hebfbsansbf}}]{Latin Modern Sans}
\fi

\usepackage{amsmath,amssymb,mathtools}
% amsthm supplies theorem/proof machinery. Environment names remain
% document-specific; the generic proof label is localized below.
\usepackage{amsthm}
% Proof names follow the surrounding prose, including return from an island.
\addto\extrashebrew{\renewcommand{\proofname}{הוכחה}}
\addto\extrasenglish{\renewcommand{\proofname}{Proof}}
\ifx\lakitmainlanguage\hllpdfenglish
  \renewcommand{\proofname}{Proof}
\else
  \renewcommand{\proofname}{הוכחה}
\fi
% v3.24.1: equation tags under babel-bidi. A tag re-emitted by mathtools'
% showonlyrefs inside align inherits the RTL base and its parens MIRROR —
% ")2.1(" instead of "(2.1)" (plain align/equation tags are unaffected;
% verified on a minimal case + a 198-page build). LTR-isolate at the tag-BOX
% level: mathtools' show/hide + numbering logic is untouched. Do NOT override
% \tagform@ instead — that clobbers the showonlyrefs mechanism (measured:
% every tag shows, numbering breaks).
% v3.28.1: the v3.24.1 wrap must apply ONLY to tags of display environments.
% Inline \eqref in Hebrew PROSE also goes through \maketag@@@, and there the
% english-wrap MIRRORS the parens — ")1(" — while the unwrapped original is
% correct (proven on a 2x2 isolation matrix: patch x showonlyrefs; regressed
% every prose \eqref since v3.24.1, caught by the render gate's R4). The two
% call sites are indistinguishable by \ifmmode (tag boxes are text mode), so
% the display environments themselves raise a flag; \eqref in prose sees it
% false and gets the original untouched \maketag@@@.
\usepackage{etoolbox}

% babel bidi=basic (including upstream 26.9) assumes every automatic
% discretionary has a one-glyph replacement. LuaTeX can legally produce one
% with replace=nil after a hyphenated Latin value; Babel then indexes nil and the
% whole build dies. Do the narrow upstream-compatible correction: while Babel
% classifies the list, present only that malformed disc as an ordinary
% discretionary, then restore its subtype. Babel still reorders the ENTIRE box;
% there is no pcall and no successful fallback to an unreordered node list.
% If Babel throws for any other reason, the build remains fatal. The synthetic
% nil-disc probe and the real table gallery in test_all.sh pin both layers.
\AtBeginDocument{\directlua{
  if Babel and Babel.bidi and not Babel.heb_nil_disc_compat then
    Babel.heb_nil_disc_compat = true
    local original = Babel.bidi
    local disc_id = node.id('disc')
    Babel.bidi = function(head, ispar, hdir)
      local hidden = {}
      for item in node.traverse_id(disc_id, head) do
        if item.subtype == 2 and (not item.replace or not item.replace.char) then
          table.insert(hidden, item)
          item.subtype = 0
        end
      end
      local result = original(head, ispar, hdir)
      for _, item in ipairs(hidden) do item.subtype = 2 end
      return result
    end
  end
}}
\makeatletter
\newif\ifHEBdisptag
\let\ORIGmaketag@@@\maketag@@@
\def\maketag@@@#1{\ifHEBdisptag\ORIGmaketag@@@{\foreignlanguage{english}{#1}}\else\ORIGmaketag@@@{#1}\fi}
\@for\HEB@env:=equation,equation*,align,align*,gather,gather*,multline,multline*,flalign,flalign*,alignat,alignat*,eqnarray,eqnarray*\do{%
  \expandafter\AtBeginEnvironment\expandafter{\HEB@env}{\HEBdisptagtrue}}
\makeatother
% English page + Hebrew prose: Babel's equation-only tag flip uses a primitive
% leqno path that shifts the formula to the left edge and the tag into the body
% of an English page. Align's full-width alignment is unaffected. Give only this
% display a physical LTR container and native amsmath eqno/leqno. The engine then
% owns real tag width, collision handling and labels; no proxy tag box hides ink
% from its fit calculation. Babel's saved \bbl@eqnodir still restores Hebrew text
% inside math, and the prose locale/math direction are never switched.
% Close the prior RTL paragraph before normalization. After the environment's
% group restores its directions, close TeX's empty resumed display paragraph and
% use native LaTeX continuation handling: same-paragraph prose has no indent,
% while an authored blank line starts an ordinary indented RTL paragraph.
% Register after Babel's setup; remain local to equation(*) in a main-English
% document's actual RTL prose. Native Hebrew, ordinary English and explicit
% fleqn retain their original paths. No page/display dimensions are assigned.
% <hllpdf:minority-rtl-equation-geometry>
\makeatletter
\newcommand{\HEB@minorityprosecontinuation}{\par\@doendpe}
\newcommand{\HEB@minorityequationgeometry}{%
  \ifx\lakitmainlanguage\hllpdfenglish
    \ifnum\pardirection=\@ne
      \if@fleqn\else
        \par\aftergroup\HEB@minorityprosecontinuation
        \iftagsleft@\let\veqno\@@leqno
        \else\let\veqno\@@eqno\fi
        \bodydir TLT\pardir TLT\textdir TLT\relax
      \fi
    \fi
  \fi}
\AtBeginDocument{%
  \AddToHook{env/equation/begin}{\HEB@minorityequationgeometry}%
  \AddToHook{env/equation*/begin}{\HEB@minorityequationgeometry}}
\makeatother
% </hllpdf:minority-rtl-equation-geometry>
% v3.63: \shortintertext rendered Hebrew MIRROR-REVERSED and flush LEFT — an entire
% commentary line inside an align printed backwards. babel's bidi=basic patches
% \intertext@ only (luababel.def), and mathtools' \shortintertext uses the identical
% \vbox{\normalbaselines … \noindent\ignorespaces#1\par} skeleton but is never
% patched, so its paragraph keeps pardir/textdir LTR: no bidi reordering (logical
% order on the page) and flush-left. Measured against \intertext with the same
% Applying babel's own replacement to \shortintertext@ makes the two identical.
% \h{}/\mbox fix the letter order but leave the line flush LEFT, so they are not
% substitutes. Guarded + \AtBeginDocument because mathtools \let-copies at load time.
% NB \ExplSyntaxOn is REQUIRED: mathtools' internals carry expl3 names, and without
% the expl3 catcodes the `_` is read as a subscript and the patch detonates far away
% with "! Illegal unit of measure (pt inserted)" at the \end{align} (measured).
\makeatletter
\AtBeginDocument{%
  \@ifundefined{shortintertext@}{}{%
    \ExplSyntaxOn
    \bbl@sreplace\shortintertext@{\normalbaselines}%
      {\normalbaselines\ifx\bbl@eqnodir\relax\else\bbl@pardir\@ne\bbl@eqnodir\fi}%
    \ExplSyntaxOff}}
\makeatother
% geometry inherits the canonical A4 sheet from \documentclass; build.sh then
% verifies every rendered MediaBox and CropBox, so a silent Letter fallback is
% impossible to publish as a clean default build.
\usepackage[\lakitgeometryoptions]{geometry}
\linespread{\lakitlinespread}       % layout hook; canonical default remains 1.3
\ifx\lakitparagraphstyle\hllpdfparagraphblock
  \setlength{\parindent}{0pt}
  \setlength{\parskip}{6pt plus 1pt minus 1pt}
\fi
% These dimensions remain ordinary scoped TeX lengths: authors can customize
% them after core input. Language switches inherit them; lists keep their own
% indentation and narrow reference layouts establish their local paragraph gaps.
\typeout{HLLPAR|\lakitparagraphstyle}
\usepackage{xcolor}
\usepackage{mdframed}
\usepackage{listings}                 % code listings (assembly / C / pseudo-code)
\usepackage{tikz}
\usetikzlibrary{arrows.meta,positioning,calc,shapes.geometric,fit,backgrounds}
\usepackage{pgfplots}\pgfplotsset{compat=1.18}
\usepackage{quantikz}                 % quantum circuits (remove if unused)
\usepackage{booktabs}
\usepackage{array}
\usepackage{float}                    % [H] = place table exactly here (reference doc: no drifting)
\usepackage{enumitem}
\usepackage{needspace}

% ============================================================================
% ★★ v3.4 — THE DIRECTION-ISOLATION TRIO (the heart of mixing Hebrew + Latin)
% Mental model:
%   \h{...}    Hebrew (RTL) island      — Hebrew text inside an LTR context
%                                          (TikZ nodes, the english `tikzpic`).
%   \en{...}   English (LTR) isolate    — Latin text / parentheses inside a
%                                          Hebrew (RTL) context. KEEPS PARENS
%                                          AND ORDER CORRECT.
%   \code{...} \en + monospace          — code inside Hebrew / inside a frametitle.
% CRITICAL: never wrap Latin in \h{} (it is for HEBREW). Inside a TikZ node,
% Latin in \h{} renders REVERSED (Datapath->htapataD, CPU->UPC) and parens
% mirror. Isolate every Latin/paren chunk with \en{} or \code{}. See
% references/figures-boxes-listings.md.
% ---- TEXT INSIDE MATH (v3.8) ----------------------------------------------
% Put ANY text that sits inside math in a \text{}. The ONLY thing that is
% invalid in math is a bare \h/\en (= \foreignlanguage) used DIRECTLY, outside
% \text -> "\rmfamily invalid in math mode". Inside \text:
%   Hebrew label           -> \text{...}            (bare Hebrew renders fine;
%                              \h is optional — BUT NOT ALWAYS HARMLESS (v3.67):
%                              if the \h group OPENS with a nested island
%                              ($..$, \en{}, \code{}, \texttt{}, \mbox{}) the
%                              WHOLE group loses bidi and the Hebrew prints in
%                              LOGICAL order = gibberish. Measured:
%                                \text{\h{$f$ קבועה}}  -> reads "העובק"  BROKEN
%                                \text{\h{קבועה $f$}}  -> correct
%                                \text{$f$ קבועה}      -> correct (no \h)
%                              Hebrew first inside \h, or drop the \h. Flagged
%                              H-OPENS-WITH-LTR-GROUP.)
%   plain Latin word        -> \text{...}            (no marker needed)
%   Latin WITH ( ) or [ ]   -> \text{\en{...}}       (\en forces LTR; a plain
%                              \text mirrors the parens in RTL -> "Spec (TNR(" )
%   f(Hebrew/text arg)      -> \text{f}(\text{...})  -- keep the parens at MATH
%                              level, NOT \text{f(arg)} (parens inside the box mirror)
%   code in math            -> \mathtt{...}  (or \texttt{}, math-guarded below)
% NEVER write \en{...} or \h{...} straight into $...$ / \[..\] (outside \text);
% always via \text{}. (\mbox{...} also works but does not scale in scripts.)
% ============================================================================
\newcommand{\h}[1]{\foreignlanguage{hebrew}{#1}}
% v3.42.1 — \en islands do NOT hyphenate (auto OR at explicit hyphens). English hyphenation patterns were
% tearing identifiers and terms ACROSS THE RTL SEAM ("IO Regis-/ters",
% "In-/put", "Band-/width", "Per-/formance", "Multi-/ple" — all shipped on
% real pages; probe-proven: "ALUS-/rcB", "Mem-/ToReg"). An island is a label
% or a term, not English prose — a hyphen tear there is never right. A/B on
% a full 197pp lakit: identical overfull sets, same page count, ~11
% paragraphs healed, one mild underfull (badness 3525) that replaced a
% mutilated "Regis-/ters". Long English PROSE belongs in the englishprose
% environment below, which restores explicit-hyphen breaks while its paragraph
% is still active. A raw otherlanguage block inherits the global prohibition.
% (\l@nohyphenation needs \csname:
% '@' is not a letter here.)
% v3.43.2 — \exhyphenpenalty moved OUT of \en to a GLOBAL assignment below.
% Render-proven: a group-local \exhyphenpenalty is DEAD CODE — TeX reads the
% line-breaking parameters at \par time, after the \en group has closed, so
% \en{write-through} / plain not-taken still tore at the "-" ("write-/
% through" shipped on a real exam page). \language is different: it is
% recorded in the node list as text is added, so the in-group nohyphenation
% assignment above DOES work (and stays). The global penalty was A/B-measured
% on a 197pp lakit: identical overfull set (25), identical page count, the
% not-/taken tear healed. Hebrew is unaffected in practice — Hebrew compounds
% use maqaf (־), not ASCII "-".
\newcommand{\en}[1]{\foreignlanguage{english}{\language=\csname l@nohyphenation\endcsname #1}}
% Babel runs extras locally for inline islands and paragraph language changes.
% An English main document must not inherit the Hebrew no-hyphen paragraph rule.
\addto\extrashebrew{\exhyphenpenalty=10000\relax}
\addto\extrasenglish{\exhyphenpenalty=50\relax}
\ifx\lakitmainlanguage\hllpdfenglish
  \exhyphenpenalty=50\relax
\else
  \exhyphenpenalty=10000\relax
\fi
% Full English paragraphs need ordinary explicit-hyphen breaks. Use this
% environment instead of a raw otherlanguage block; the paragraph ends while
% the local penalty is still active.
\newenvironment{englishprose}
  {\par\begin{otherlanguage}{english}\exhyphenpenalty=50\relax}
  {\par\end{otherlanguage}}
\newenvironment{hebrewprose}
  {\par\begin{otherlanguage}{hebrew}\exhyphenpenalty=10000\relax}
  {\par\end{otherlanguage}}
% Use for shell text whose locale must not depend on the last body paragraph.
\DeclareRobustCommand{\doclang}[1]{\foreignlanguage{\lakitmainlanguage}{#1}}

% ============================================================================
% ★★ — fixes for defects that are INVISIBLE in the log and in the checks
% and only surface on a VISUAL RENDER of the PDF. (Each verified empirically.)
% ----------------------------------------------------------------------------
% v3.22: \emph is now REAL slanted emphasis (FakeSlant italic face defined on
% the \babelfont{rm} line above), replacing the v3.21 \emph->\textbf hack: it
% preserves \emph's toggle semantics and keeps bold free for its own role.
% Prefer bold-as-emphasis for your doc? one line: \renewcommand{\emph}[1]{\textbf{#1}}
% page number: force the digits LTR so they never reverse (18 -> 81) in the RTL
%   footer. That reversal is a NON-DETERMINISTIC bidi glitch (observed once in
%   ~80 pp); this LuaLaTeX build is not byte-reproducible, so the defensive
%   LTR wrap is the only reliable guard. Latin digits are always LTR anyway.
\renewcommand{\thepage}{\foreignlanguage{english}{\arabic{page}}}
% v3.70: \pagenumbering \gdef's \thepage and silently DESTROYS the wrap above —
% after \pagenumbering{arabic} the folio guard is gone for the whole rest of the
% document. Two collections discovered this independently and hand-restore the
% wrap in main.tex with the same comment; that knowledge now lives here instead.
% The wrapper re-applies the LTR wrap after the kernel call, for any style.
\makeatletter
\let\lakin@orig@pagenumbering\pagenumbering
\renewcommand{\pagenumbering}[1]{%
  \lakin@orig@pagenumbering{#1}%
  \renewcommand{\thepage}{\foreignlanguage{english}{\csname @#1\endcsname\c@page}}}
\makeatother
% emergencystretch: Hebrew does not hyphenate, so an unbreakable LTR island
%   (\en{...} / an inline math run) landing at the margin can leave a small
%   (~2-13pt) overfull with no legal breakpoint. This rescues ONLY those
%   inline-prose tie-breaks; it is a ceiling used only on offending lines and
%   does NOT touch display-math / listing / table overflow (those still surface
%   in the log for the content-level overflow pass). Verified: baseline overfull
%   -> 0 with this set, with no change to pagination or underfull.
\setlength{\emergencystretch}{3em}
% v3.40: inline math (almost) never breaks across lines. TeX's default
%   (binop 700 / rel 500) lets an inline run split at an operator —
%   $(1.f)\cdot 2^{e}$ shipped as "(1.f)·" ending one line and "2^e" opening
%   the next (reader-reported). In LTR that split is merely ugly; under an
%   RTL base the carried operator lands on the visually WRONG side of the
%   break and the formula reads backwards.
%   WHY 10000 (v3.45.0 — the display-first resolution). Inline math is
%   FORBIDDEN from breaking, so a torn island under an RTL base is
%   structurally impossible. The v3.40 cycle used 9999 ("break only as a last
%   resort") because a bare 10000 left ~30-39 WIDE inline equations
%   overflowing (worst 90.7pt) — but 9999's "last resort" fires OFTEN in
%   tight, non-hyphenating Hebrew: a reader found torn islands across the
%   fleet (~150 candidates; e.g. "P(y |" ending one line and "x_0)" opening
%   the next, read backwards). The step the 9999 cycle missed: every one of
%   those ~30 overflows is a wide equation that BELONGS in display. Convert
%   them to display FIRST, then 10000 gives 0 tears AND 0 overflows (measured,
%   12/12 fleet). Doctrine: inline math is ATOMIC; a wide equation goes in a
%   display, never inline. See content-style ("inline math is atomic").
\binoppenalty=10000 \relpenalty=10000
% v3.38: uniform sentence spacing. TeX's default (\nonfrenchspacing, sfcode
% 3000) is the OLD Anglo-American wide-space-after-period convention — Hebrew
% typography never had it, and modern Latin practice (Chicago et al.) dropped
% it too. Measured on real pages: post-period gaps up to 1.9-2.1x the line's
% mean word gap — the dominant "rivers" source in justified Hebrew prose.
% Survives babel language switches (probed: sfcode stays french inside \en{}
% islands and otherlanguage blocks).
\frenchspacing
% v3.38: widow/orphan protection. LaTeX's default penalties (150/150/50)
% allow a lone line at a page top/bottom almost freely. 10000 = forbid;
% this is SAFE here precisely because the document is \raggedbottom (the
% demoted line just moves on, no vertical stretch fight) — under flushbottom
% these values would cause underfull pages. raggedbottom is the report-class
% default; made explicit so a future class/option change cannot silently
% flip it and turn the 10000s harmful.
\raggedbottom
\widowpenalty=10000 \clubpenalty=10000 \displaywidowpenalty=10000
% em-dash: --- does NOT become — under Frank Ruhl (the font lacks tlig; verified
%   the width is identical with/without Ligatures=TeX, while the mechanism works
%   on Latin Modern). TYPE THE CHARACTER — (U+2014) DIRECTLY; likewise – (U+2013).
%   check_bidi_figures.py flags a --/--- that is adjacent to Hebrew.
% ============================================================================

% ---- STRUCTURAL DIRECTION-SAFETY for \texttt (v3.7) -------------------------
% A parenthesis, bracket, or $ INSIDE \texttt{} mirrors/migrates in RTL prose
% (MIPS 0($s0)->"(s0$)0", $rs->"rs$", arr[i]->"[arr[i") because \texttt sets the
% font but NOT the direction. Redefining \texttt to LTR-isolate its content fixes
% the WHOLE class at once (parens, brackets, $) so plain \texttt{lw $t0,0($s0)}
% is safe — no need to remember \code{}. Guarded by \ifmmode because
% \foreignlanguage is invalid in math (there, plain monospace is already LTR).
% NB (corrected v3.56 — the old note here said "stays tofu", which is WRONG):
% \texttt{Hebrew} is a FATAL build failure. The tt family is DejaVu Sans Mono
% (zero Hebrew) and the tt fallback hebfb is DejaVuSerif (also zero Hebrew), so
% the subsetter gets an empty glyph set and LuaTeX dies with
%   "! error: (file .../DejaVuSansMono.ttf) (type 2): there are no glyphs in the
%    subset"  -> no output PDF.
% The message names a FONT FILE, never Hebrew/\texttt/the source line. And it is
% non-local: add an unrelated Latin \code{} elsewhere and the fatal degrades into
% a SILENT DELETION of the Hebrew word. Put Hebrew in \h{}/normal text, never in
% \texttt{}/\code{}; flagged HEBREW-IN-MONO. To disable, delete this block.
\let\origtexttt\texttt
\protected\long\def\texttt#1{\ifmmode\origtexttt{#1}\else\foreignlanguage{english}{\origtexttt{#1}}\fi}
% v4.8.0: defer the direction test until TeX has inserted an alignment cell's
% math template. Unprotected expansion here tested \ifmmode too early at the
% FIRST token of aligned/array/cases/matrix cells. Engine-level \protected also
% preserves the original \texttt implementation without the self-recursion
% observed when retrofitting \DeclareRobustCommand to its saved alias.
\protected\long\def\code#1{\ifmmode\origtexttt{#1}\else\foreignlanguage{english}{\origtexttt{#1}}\fi}

% ---- \codeline — display-code line (v3.24): the code analogue of display math.
% Reader-friction rule (content-style.md §5b): inline \code{} is for a SINGLE
% ATOM only (one register/variable/mnemonic, a tight no-space expression).
% A multi-token expression — an instruction with operands, a spaced assignment,
% a register-binding list — zigzags the eye RTL<->LTR mid-sentence and, being
% an unbreakable island, line-breaks badly. Give it its own centered line with
% display-math vertical rhythm. \centerline = one box => the expression never
% wraps; \nopagebreak keeps it glued to its introducing sentence.
%   ...the Hebrew sentence ends with a colon:
%   \codeline{lw \$t1, -4(\$t0)}
%   and the prose continues (escape dollars as \$).
% v3.66 (first line of the body): flush a PENDING \item label before the \par.
% LaTeX defers the label box (\@labels) until the next \everypar, and an \hbox in
% vertical mode never fires \everypar — so `\item \codeline{...}` printed the code
% at y=69.1 and the item's own "1." BELOW it at y=103.6 (the number appears after
% its content, attached to nothing). \if@inlabel\leavevmode\fi is CONDITIONAL on
% purpose: an unconditional \leavevmode would open an empty paragraph line when
% \codeline follows ordinary prose. Verified both ways — in a list the label now
% sits on the code's line at its normal x, and after prose the geometry is
% byte-identical to before.
\makeatletter
\newcommand{\codeline}[1]{%
  \if@inlabel\leavevmode\fi
  \par\nopagebreak
  \vspace{\abovedisplayskip}\nopagebreak
  % v3.62: was \centerline{...}, i.e. \hbox to\hsize{\hss #1 \hss}. \hss has
  % INFINITE SHRINK, so TeX could always "fit" the box and NEVER reported an
  % Overfull — the same \hss mechanism as BLD-27's \item label. Measured: a wide
  % \codeline in a defbox painted its first token ON the box's left border bar and
  % sat 9.5pt outside the frame interior on both sides, with 0 Overfull, 0 Underfull
  % and all four gates PASS; in bare prose it reached 21.0pt into BOTH margins.
  % \hfil has no shrink, so the overflow is reported and build.sh's ovf_vis/ovf_big
  % counters see it; \linewidth (not \hsize) makes it box-aware, so a \codeline
  % inside a colored box centres on the box interior. Verified: a FITTING \codeline
  % keeps byte-identical glyph positions, and the box-crossing case now emits
  % "Overfull \hbox (19.03pt too wide)". A wide \codeline should be split in two.
  \hbox to\linewidth{\hfil\code{#1}\hfil}%
  \nopagebreak\vspace{\belowdisplayskip}}
\makeatother

% ---- a thin decorative rule helper ----
\newcommand{\sep}{\par\vspace{2pt}\noindent\textcolor{cRule}{\rule{\linewidth}{0.6pt}}\par\vspace{4pt}}

% ---- colors ----
\definecolor{cBlue}{HTML}{1F4E79}   \definecolor{cBlueBg}{HTML}{EAF1F8}
\definecolor{cGreen}{HTML}{2E6B4F}  \definecolor{cGreenBg}{HTML}{EAF4EE}
\definecolor{cAmber}{HTML}{9C6A00}  \definecolor{cAmberBg}{HTML}{FBF3E2}
\definecolor{cPurple}{HTML}{5B2C6F} \definecolor{cPurpleBg}{HTML}{F3ECF7}
\definecolor{cTeal}{HTML}{0E6E6E}   \definecolor{cTealBg}{HTML}{E6F4F4}
\definecolor{cRed}{HTML}{8E2B2B}    \definecolor{cRedBg}{HTML}{FBEAEA}
\definecolor{cGray}{HTML}{555555}   \definecolor{cRule}{HTML}{1F4E79}

% ---- math operators ----
\DeclareMathOperator{\Tr}{Tr}
\DeclareMathOperator{\rank}{rank}
\DeclareMathOperator{\diag}{diag}
\DeclareMathOperator{\sgn}{sgn}

% ---- bra-ket (quantikz defines \ket/\bra/\braket; add the rest safely) ----
\providecommand{\ket}[1]{\left|#1\right\rangle}
\providecommand{\bra}[1]{\left\langle#1\right|}
\providecommand{\braket}[2]{\left\langle#1\middle|#2\right\rangle}
\newcommand{\dyad}[1]{\left|#1\right\rangle\!\left\langle#1\right|}
% v3.21: general outer product |a><b| and expectation value <A> — agents reach
% for these by reflex (esp. \expval); quantikz defines neither, so their absence
% threw "Undefined control sequence" in a real build. \providecommand = safe.
\providecommand{\ketbra}[2]{\left|#1\right\rangle\!\left\langle#2\right|}
\providecommand{\expval}[1]{\left\langle#1\right\rangle}
\newcommand{\abs}[1]{\left|#1\right|}
\newcommand{\norm}[1]{\left\lVert#1\right\rVert}
\newcommand{\half}{\tfrac12}
\newcommand{\R}{\mathbb{R}}\newcommand{\C}{\mathbb{C}}\newcommand{\Z}{\mathbb{Z}}
% identity operator. NB: \mathbb{1} silently renders a WRONG glyph (msbm has no
% blackboard digits). Use \mathbb{I} = 𝕀 (works with amssymb, no extra package).
% Do NOT "fix" this back to \mathbb{1}.
\newcommand{\id}{\mathbb{I}}

% ---- physical units (v3.35): \punit — upright unit with thin space ---------
% Usage: $300\punit{K}$, $80\punit{mV/decade}$, $10^{7}\punit{cm/s}$ — the
% NUMBER AND UNIT LIVE IN ONE MATH ISLAND. Promoted verbatim from a 199-page
% production lakit (advanced-electronic-devices, dozens of uses).
% Two measured traps this design avoids (both render-proven):
%   1. siunitx SCRAMBLES under RTL — \SI{3.5}{eV} renders "eV 3.5", exponents
%      detach ("Hz 10 x 10^9"): its quantity assembly is not bidi-isolated.
%      Do NOT load siunitx in this pipeline.
%   2. A text-mode unit next to a bare-prose number is TWO adjacent LTR runs
%      and reorders ("5 mm" -> "mm5" — the BIDI-2/3 class). \punit is therefore
%      math-only ON PURPOSE: prose misuse fails LOUDLY at compile (\mathrm
%      outside math) instead of scrambling silently.
% Domain shorthands (\MeV, \fm...) stay per-document: \newcommand{\MeV}{\,\text{\en{MeV}}}
\providecommand{\punit}[1]{\,\mathrm{#1}}

% ============================================================================
% Colored boxes — mdframed (NOT tcolorbox, which breaks under bidi).
% ★★ v3.4: innerleftmargin/innerrightmargin give code listings room so their
% frame/line-numbers do not spill out the box border (see listing styles below).
% ============================================================================
% v3.61: everyline/split skips were applied to `tablebox` only (BLD-14, v3.37) —
% the six colored boxes below never got them, so they still split OPEN across a page
% break: the page-N segment gets top+sides and NO bottom rule, the page-N+1 segment
% sides+bottom and NO top rule. Measured on the shipped fleet this is a COVERAGE GAP,
% not a re-report of BLD-14: 0 split-open `tablebox` segments (that fix holds) vs 62
% split-open COLORED boxes across 7 documents. Verified after the fix: both segments
% close, and the frametitle is NOT duplicated on the continuation. BLD-14's caveat
% still applies — a box taller than a full page grinds mdframed's split-trial loop, so
% split such content manually.
% ============================================================================
% v3.72: frametitle={} leads EVERY \newmdenv option list below. mdframed options
% are GROUP-scoped: a box nested inside \begin{exbox}[frametitle={...}] used to
% INHERIT the outer title — measured, the identical title string rendered twice,
% the second riding a title-bar fill inside the INNER box's frame, for all seven
% environments in both directions. Both natural forms were broken (= preamble
% bug, not an authoring rule). The explicit reset makes every box titleless
% unless ITS OWN [frametitle={...}] says otherwise; verified both directions:
% no duplicate on the nested box, an explicit title still renders.
% Geometry is shared with the optional semantic-width table container. Themes
% select their own existing semantic style below; neither path fixes columns.
\mdfdefinestyle{hllpdfsemanticgeometry}{linewidth=1.1pt,
  innertopmargin=6pt,innerbottommargin=7pt,skipabove=9pt,skipbelow=9pt,
  innerleftmargin=9pt,innerrightmargin=9pt,everyline=true,
  splitbottomskip=6pt,splittopskip=12pt,roundcorner=3pt}
\newcommand{\hllpdfsemanticframestyle}{hllpdfsemanticgeometry}
\newmdenv[frametitle={},style=hllpdfsemanticgeometry,linecolor=cBlue,backgroundcolor=cBlueBg,
  frametitlebackgroundcolor=cBlue,frametitlefontcolor=white,
  ]{defbox}   % הגדרה
\newmdenv[frametitle={},style=hllpdfsemanticgeometry,linecolor=cGreen,backgroundcolor=cGreenBg,
  frametitlebackgroundcolor=cGreen,frametitlefontcolor=white,
  ]{thmbox}   % משפט/טענה
\newmdenv[frametitle={},style=hllpdfsemanticgeometry,linecolor=cAmber,backgroundcolor=cAmberBg,
  frametitlebackgroundcolor=cAmber,frametitlefontcolor=white,
  ]{notebox}  % הערה
\newmdenv[frametitle={},style=hllpdfsemanticgeometry,linecolor=cPurple,backgroundcolor=cPurpleBg,
  frametitlebackgroundcolor=cPurple,frametitlefontcolor=white,
  ]{exbox}    % דוגמה
\newmdenv[frametitle={},style=hllpdfsemanticgeometry,linecolor=cRed,backgroundcolor=cRedBg,
  frametitlebackgroundcolor=cRed,frametitlefontcolor=white,
  ]{warnbox}  % טעות נפוצה
\newmdenv[frametitle={},style=hllpdfsemanticgeometry,linecolor=cTeal,backgroundcolor=cTealBg,
  frametitlebackgroundcolor=cTeal,frametitlefontcolor=white,
  ]{keybox}   % רעיון מרכזי

% ---- codebox (v4.5.0): the dedicated container for code -------------------
% An OPTION, not a policy: asmblock unboxed stays canonical for a full listing,
% and a snippet inside any of the six colored boxes still works as before.
% What makes codebox worth its own environment is that it is the ONE framed box
% that keeps its margin line numbers.
%
% v4.4.0 shipped it WITHOUT them, on the strength of v3.37's "no margin knob can
% save them" plus four fresh probes that agreed. That was wrong: every probe had
% moved the code, the box, or both — never the key that carries the NUMBER
% COLUMN. Measured on the page:
%   resetmargins=false   the listing follows THIS box instead of the page,
%                        moving the code from x=75.26 to x=97.54
%   xrightmargin=26pt    under RTL this is the PHYSICAL LEFT indent, and the
%                        only key that carries the \llap'ed digits with it:
%                        x=64.28 -> 90.18, clearing the rule at 65.64 by 23.6pt
%   xleftmargin=26pt     moves the code and LEAVES THE DIGITS on the rule —
%                        the asymmetry that hid the fix for eleven versions
% resetmargins is hoisted out of the four styles (see the \lstset below them) so
% this override is possible at all; every other listing still resolves to true.
% Pinned by tests/codebox_numbers.tex + probe_lstnum.py --inside-rule, whose
% negative control strips xrightmargin and watches the digits fall back.
\definecolor{cSlate}{HTML}{3F4A56}  \definecolor{cSlateBg}{HTML}{F2F4F7}
\newmdenv[frametitle={},leftline=true,rightline=true,topline=true,bottomline=true,
  linecolor=cSlate,backgroundcolor=cSlateBg,linewidth=0.9pt,
  frametitlebackgroundcolor=cSlate,frametitlefontcolor=white,
  innertopmargin=6pt,innerbottommargin=7pt,skipabove=9pt,skipbelow=9pt,
  innerleftmargin=9pt,innerrightmargin=9pt,everyline=true,splitbottomskip=6pt,splittopskip=12pt,
  roundcorner=3pt]{codebox}   % קוד
% The two keys that make in-box numbering work, both measured on the page:
%   resetmargins=false  the listing follows THIS box instead of the page
%   xrightmargin        under RTL this is the PHYSICAL LEFT indent, and it is
%                       the only one that carries the \llap'ed number column
%                       with it. xleftmargin moves the code and leaves the
%                       digits behind on the rule (measured: code 97.54,
%                       digits 64.28 — the v3.37 defect, reproduced); with
%                       xrightmargin=26pt the digits land at x=90.18, a
%                       clearance of 23.6pt from the rule at 65.64.
% This asymmetry is why v3.4's note had the semantics backwards and v3.37
% concluded the digits could not be saved.
% <hllpdf:listing-container-state>
% v4.29: all state is local to the current environment group. A semantic
% callout masks number text, while a nested codebox restores the formatter
% saved before that mask. Save it once across successive semantic wrappers;
% an explicit numbers=none (or an explicitly empty formatter) remains valid.
\makeatletter
\newif\ifhllpdf@incodebox
\newif\ifhllpdf@englishcodebox
\newif\ifhllpdf@lstnumbermasked
\let\hllpdf@savedlstnumber\thelstnumber
\makeatother
% </hllpdf:listing-container-state>
% <hllpdf:codebox-gutter>
\makeatletter
\AtBeginEnvironment{codebox}{%
  \hllpdf@incodeboxtrue
  % Snapshot the OUTER language before asmblock selects English.
  \iflanguage{english}{\hllpdf@englishcodeboxtrue}{\hllpdf@englishcodeboxfalse}%
  \ifhllpdf@lstnumbermasked
    \let\thelstnumber\hllpdf@savedlstnumber
    \hllpdf@lstnumbermaskedfalse
  \fi
}
% listings' AddToHook is GLOBAL. Register once, and gate on local box state;
% adding an unconditional hook at every codebox entry leaks into later listings.
\lst@AddToHook{PreInit}{%
  \ifhllpdf@incodebox
    % Apply every listing override here, inside the LISTING's own group.
    % Setting margins at box entry leaks them into nested semantic callouts.
    \lstset{resetmargins=false,xrightmargin=26pt}%
    % Styles resolve before PreInit. In LTR the physical number gutter follows
    % xleftmargin; in Hebrew retain the measured inherited xrightmargin path.
    \ifhllpdf@englishcodebox
      \lstset{xleftmargin=26pt,xrightmargin=0pt}%
    \fi
    % ONE box, not two: suppress only the frame/background/skips of a listing
    % inside this codebox. The next unboxed listing keeps its own style.
    \lstset{frame=none,backgroundcolor=\color{cSlateBg},aboveskip=0pt,belowskip=0pt}%
  \fi
}
\makeatother
% </hllpdf:codebox-gutter>

% v3.37: line numbers inside a COLORED BOX are sliced by the box border —
% measured: digit span x=63.5-67.3 vs border line at x=65.2 (the border runs
% THROUGH the digits), and the inter-border gap (3.5pt) is narrower than a
% digit (3.8pt), so no margin knob can save them (under resetmargins,
% xleftmargin is inert in-box — measured: code at x=75.26 with 2.2em AND
% 3.0em). A 3-6 line boxed snippet does not need line numbers; blank the
% number TEXT inside every colored box (\thelstnumber is what gets printed;
% the style's numbers=left survives, so nothing else shifts — code position
% byte-frozen at x=75.26, verified). Unboxed listings keep their numbers.
% v3.78: the blanking is coupled to the FRAME, not to the box name —
% clean_theme.tex registers LATER hooks (registration order wins) that
% restore the number text inside the boxes it renders frameless.
% v4.5.0 RETRACTS the "no margin knob can save them" clause above. It holds only
% under resetmargins=true (the shipped default): there xleftmargin is inert
% in-box, as measured. With resetmargins=false, xrightmargin DOES carry the
% \llap'ed digit column clear of the rule — measured; the codebox \lstset that
% does it is the one above, in the codebox block. The loop below therefore
% covers SIX boxes, not "every colored box": codebox opts out and keeps its
% numbers.
% <hllpdf:semantic-listing-mask>
\makeatletter
\@for\HEB@lstbox:=defbox,thmbox,notebox,exbox,warnbox,keybox\do{%
  \expandafter\AtBeginEnvironment\expandafter{\HEB@lstbox}{%
    \ifhllpdf@lstnumbermasked\else
      \let\hllpdf@savedlstnumber\thelstnumber
    \fi
    \hllpdf@lstnumbermaskedtrue
    \def\thelstnumber{}%
    % The nearest semantic container owns its listing presentation. On exit,
    % TeX's group restores the enclosing codebox state and formatter.
    \hllpdf@incodeboxfalse}}
\makeatother
% </hllpdf:semantic-listing-mask>

% ============================================================================
% Unified TABLE presentation: subtle uniform slate frame + auto-numbered
% caption rendered INSIDE the frame. Use:
%   \begin{tablebox}\centering  <tabular>  \tabcap{<description>}\end{tablebox}
% Replaces the buggy \caption* (which emitted a stray "טבלה N : *" line) and
% provides an explicit framed option + numbered caption "טבלה N.M: תיאור".
% Optional intro text may be placed at the TOP, before the tabular, inside the box.
% ============================================================================
\definecolor{cTblLine}{HTML}{B9C4D2}  \definecolor{cTblBg}{HTML}{FAFBFC}
% v3.23: class-aware — under `article` there is no chapter counter (the promised
% report->article switch used to die with "No counter 'chapter' defined").
\makeatletter
\@ifundefined{c@chapter}{%
  % v3.31: in article the counter resets per SECTION, so an unprefixed number
  % repeated ("טבלה 1." in every section — measured on a two-section render).
  % Same conditional idea as the report branch: prefix the section number once
  % a real \section exists; a pre-section table stays "טבלה 1".
  \newcounter{tbl}[section]
  \renewcommand{\thetbl}{\ifnum\value{section}>0 \thesection.\fi\arabic{tbl}}%
}{%
  \newcounter{tbl}[chapter]
  % v3.30: a report-class doc that uses SECTIONS ONLY (chapter stays 0 —
  % the common short-doc pattern) used to print "טבלה 0.1"; the chapter
  % prefix now appears only once a real \chapter exists (render-verified
  % both ways: "טבלה 1." chapterless, "טבלה 1.1" with chapters).
  \renewcommand{\thetbl}{\ifnum\value{chapter}>0 \thechapter.\fi\arabic{tbl}}%
  % v3.64: the SAME leak in \thesection itself, left behind by v3.30. That fix
  % guarded the CAPTION counter and stopped there, so a chapterless report still
  % printed the heading "0.1 סעיף ראשון" and \ref to it printed "0.2" — measured:
  %   p1 y=65.2  '…0.1'   p1 y=144.9 '…0.2'   p1 y=183.4 '…0.2.1'
  % report.cls hardcodes \thesection as \thechapter.\@arabic\c@section, and the
  % leading 0 is the chapter counter that no \chapter ever incremented. Fixing
  % \thesection fixes the whole chain — \thesubsection and \thesubsubsection are
  % defined in terms of it. \theequation and \thetbl/\thefig were already guarded,
  % which is exactly why this survived: every NUMBER around the heading looked
  % right, so the heading's own 0 read as a style choice rather than a bug.
  % \appendix is safe: report.cls resets chapter to 0 and redefines \thechapter to
  % \@Alph\c@chapter, so an appendix \chapter makes the counter 1 and the guard
  % emits "A.1" unchanged. Render-verified in all four modes (chapterless,
  % chaptered, appendix+chapter, appendix+section-only).
  \renewcommand{\thesection}{\ifnum\value{chapter}>0 \thechapter.\fi\arabic{section}}%
}
\makeatother
% v3.37: everyline=true — if the box DOES split across pages, every segment
% gets a CLOSED frame (measured before: the split ran open-ended into the page
% edge with no bottom border, and the caption could strand alone at the next
% page top — "טבלה A.2" shipped that way). The split skips give the segments
% breathing room. In practice, with the caption GLUE below, every previously-
% split table in the corpus now moves WHOLE instead (the white-gap tradeoff,
% documented benign); everyline stays as the safety net for a genuinely
% page-tall table. NB: a table taller than ONE FULL PAGE can grind mdframed's
% split-trial loop (reproduced: minimal doc, minutes of 1pt-step Overfull
% probes) — break such a monster into two tables.
\mdfdefinestyle{hllpdftablegeometry}{leftline=true,rightline=true,topline=true,bottomline=true,
  linewidth=0.9pt,
  innertopmargin=9pt,innerbottommargin=8pt,skipabove=11pt,skipbelow=11pt,
  innerleftmargin=11pt,innerrightmargin=11pt,roundcorner=4pt}
\newcommand{\hllpdftableframestyle}{hllpdftablegeometry}
\newmdenv[frametitle={},style=hllpdftablegeometry,
  linecolor=cTblLine,backgroundcolor=cTblBg,
  % v4.2.0 (adopted): tablebox is ATOMIC. Short tableboxes near page bottoms
  % repeatedly entered mdframed's split trial, emitting benign Overfull-vbox
  % bursts and occasionally separating the intro prose from the tabular.
  % nobreak=true moves the complete short table to the next page instead. The
  % output-active Overfull-vbox gate still catches a box taller than a page;
  % everyline/split skips remain as defensive styling should that path ever
  % re-arm, but the canonical path never publishes a split tablebox.
  nobreak=true,
  everyline=true,splitbottomskip=6pt,splittopskip=12pt,
  % v3.79 (BLD-37 table side, correction): startinnercode/endinnercode run
  % INSIDE mdframed's body vbox (mdframed.sty:848 / :856+863) -- the second one
  % after \par\unskip, i.e. after line breaking, which is the only moment the
  % \tabcap anchor's true offset from the top of the box EXISTS. See
  % \lakintabboxclose below. These two keys are the ONLY tablebox change.
  startinnercode={\lakintabboxopen},endinnercode={\lakintabboxclose}]{tablebox}
% v3.37: \nopagebreak GLUE — \vsplit honors penalties, so the caption can no
% longer strand alone at the top of the next page (it takes its rows along, or
% the whole box moves). Render-verified on the real shipped orphan.
% v3.59: \NR@gettitle is what nameref injects into the stdlib title-bearing
% commands (\@sect, \@chapter, \caption) to fill \@currentlabelname. \figcap and
% \tabcap are INDEPENDENT macros — they never called it, so \@currentlabelname kept
% whatever the last heading left and \nameref{fig:x} silently printed the enclosing
% SECTION's title instead of the caption (proven in the .aux: field 3 of \newlabel
% was the heading text while field 1 and the anchor `fig.1.1` were correct — so this
% is a THIRD, distinct head next to BLD-13's counter fix and BIDI-38's autoref-name
% fix, neither of which writes that field). \@ifundefined keeps it inert if
% hyperref/nameref is ever dropped. Guarded helper because \figcap is defined
% after \makeatother.
% NB both helpers below need \makeatletter — this part of the preamble is OUTSIDE
% any letter-catcode region, so an unwrapped \@currentlabel parses as LaTeX's \@
% and dies with "You can't use \spacefactor in vertical mode" (measured).
\makeatletter
\newcommand{\lakinrtitle}[1]{\@ifundefined{NR@gettitle}{}{\NR@gettitle{#1}}}
% v3.59: \hitem{<letter>}{<key>} — a REFERENCEABLE Hebrew-lettered list item.
%   \begin{enumerate} \hitem{א}{it:a} הסעיף הראשון. \end{enumerate}  ->  \ref{it:a} = (א)
% Why this exists: automatic Hebrew list labels fail silently, so SKILL.md mandates
% an EXPLICIT \item[(\textbf{א})] — but LaTeX puts \refstepcounter INSIDE
% \if@noitemarg (latex.ltx:12964), and \item[...] never sets that flag. So an
% explicit item steps no counter and never touches \@currentlabel: a \label there
% binds to the AMBIENT one and \ref prints the CHAPTER/SECTION number where the
% item letter belongs (proven in the .aux: {chapter.1}). Same doctrine as \figcap —
% own the label state before \label. \phantomsection matters: without it the
% hyperlink still lands on the chapter head (measured); with it the anchor is
% section*.N and the bookmark outline stays clean.
\newcommand{\hitem}[2]{\item[(\textbf{#1})]%
  \def\@currentlabel{(#1)}\def\@currentlabelname{#1}\phantomsection\label{#2}}
% NB (v3.66): \hitem is \ref-only. \autoref{key} calls it a SECTION ("סעיף (א)")
% because \phantomsection brands the anchor section*.N, and \nameref{key} prints
% the bare letter. For \autoref-able items prefer \hebenumlabels below.
%
% v3.66: \hebenumlabels — OPT-IN automatic Hebrew enumerate letters, done right.
%   {\hebenumlabels\begin{enumerate}\item\label{q:a} ...\end{enumerate}}
%   -> labels (א) (ב) (ג); \ref{q:a} = א; anchor Item.N (real, per-item).
% This retires the belief that made \item[(\textbf{א})] + \hitem necessary.
% "Automatic Hebrew labels fail silently" was TRUE but unexplained, and the
% mechanism is now bisected (plain article vs +enumitem, measured): ENUMITEM
% PRE-EXPANDS the label expression in an \edef BEFORE \refstepcounter, with the
% counter still at 0. Tokens it shields by name (\arabic,\alph,\roman) survive;
% a plain \newcommand \ifcase wrapper evaluates at 0 -> the EMPTY zero-branch ->
% the label prints as "()" with the letter GONE (and \the\c@enumi prints a
% literal 0, \number\value dies fatally pointing at \item). The cure is one
% property, not a doctrine: the wrapper must be ROBUST (survives the \edef
% unexpanded), while \theenumi stays EXPANDABLE (so \refstepcounter writes the
% real letter to the .aux). Verified: labels (א)(ב)(ג), \ref -> א/ב/ג each item
% its own value, .aux = {א}{...}{Item.1} — versus the explicit-\item doctrine
% where BIDI-50/\hitem exist purely to patch the counter that never stepped.
\def\@hebalph#1{\ifcase#1\or א\or ב\or ג\or ד\or ה\or ו\or ז\or ח\or ט\or י\or
  יא\or יב\or יג\or יד\or טו\or טז\or יז\or יח\or יט\or כ\else\@ctrerr\fi}
\DeclareRobustCommand*\hebalph[1]{\expandafter\@hebalph\csname c@#1\endcsname}
\newcommand{\hebenumlabels}{%
  \renewcommand{\theenumi}{\@hebalph\c@enumi}%
  \renewcommand{\labelenumi}{(\textbf{\hebalph{enumi}})}}
\makeatother
% v3.63: \mbox the LABEL. \thetbl ("7.1") and a caption description that STARTS
% WITH A DIGIT ("32 הרגיסטרים…") are two European-Number runs separated only by
% neutrals (the label's own "." and \enspace). Under the RTL paragraph base the UBA
% resolves those neutrals to the RTL level, so both numbers sit in ONE reversed span
% and EXCHANGE VISUAL POSITIONS. Measured on the shipped PDF: a table whose real
% number is 7.1 printed "טבלה 32" — a different, EXISTING table number — while \ref
% correctly printed 7.1, so the caption silently contradicted its own cross-reference,
% and the description lost its leading count. The source contains no adjacency at all;
% it is manufactured by this macro, which is why every source linter was blind.
% \mbox makes the label one atomic box so nothing can reorder across it. Verified
% zero-regression: an ordinary Hebrew caption keeps identical x-geometry.
% One paragraph style covers label, number, separator and description.
% Explicit project/template overrides remain possible; the label hook is kept
% for compatibility, with an identity default. Counters/anchors stay unchanged.
\newcommand{\lakitcaptionstyle}{\normalfont\small\color{black}\hypersetup{allcolors=black}}
\newcommand{\lakitcaptionlabelstyle}[1]{#1}
\newcommand{\lakitcaptionpunctuation}{:}
% Native floats, captionof and longtable use the same presentation policy.
% subcaption is optional and inherits the same font rather than a second size.
\usepackage{caption}
\DeclareCaptionFont{lakit}{\lakitcaptionstyle}
\DeclareCaptionLabelFormat{lakit}{\lakitcaptionlabelstyle{#1~\en{#2}}}
\DeclareCaptionLabelSeparator{lakit}{\lakitcaptionpunctuation\enspace}
\captionsetup{font=lakit,labelfont={},textfont={},labelformat=lakit,
  labelsep=lakit,justification=centering}
\makeatletter
\AtBeginDocument{\@ifpackageloaded{subcaption}{%
  \captionsetup[sub]{font=lakit,labelfont={},textfont={},labelsep=lakit}}{}}
\makeatother
% v3.64: OPT-IN list-of-figures/tables entry via an optional short caption.
% \listoffigures and \listoftables are inherited from report.cls and LOOK available,
% but \figcap/\tabcap step the custom \c@fig/\c@tbl counters and never call
% \addcontentsline — so both printed their heading and NOTHING ELSE. Measured: a doc
% with one \figcap and one \tabcap rendered page 1 as 'רשימת האיורים\n1\n', and the
% .lof held only \@chapter's \addvspace, zero \contentsline entries. An empty list
% under a correct Hebrew heading reads as "this document has no figures".
%
% WHY OPTIONAL, AND NOT AUTOMATIC. Writing #1 unconditionally was implemented first
% and REJECTED: a caption containing \codeline compiles CLEANLY today (verified on
% the pre-v3.64 preamble: 0 errors, PDF produced) but becomes FATAL once it is also
% written to the .lof — \codeline is \hbox to\linewidth{...}, and a full-width box in
% a list entry dies with "\pdfextension endlink cannot be used in vertical mode", no
% output PDF. Gaining a feature nobody uses (0 of 13 collections call \listoffigures)
% at the price of breaking a construct that works is a bad trade. Stock \caption has
% exactly this limitation and exactly this escape hatch, so the convention is familiar.
% \en{}, \code{}, \h{} and inline math in the short text are all safe — verified, they
% survive the write and typeset correctly in the list (RTL geometry confirmed: number
% rightmost at x=509.9, title left of it, dots and folio at the far left).
% With no optional argument the macro is byte-identical to v3.63 — zero risk to every
% existing document, and \listoffigures lists exactly the captions that opted in.
% BLD-37, TABLE side (the \figcap v3.74 idea, adapted): clicking a \ref to a
% table used to land MID-table — hyperref's automatic anchor inside
% \refstepcounter was emitted on the still-open horizontal line whose TAIL is
% the tabular (an hbox whose baseline is the table's vertical CENTER), so the
% jump scrolled the top half of the table above the viewport. The v3.74 entry
% deferred this side claiming "tablebox splits — no single box to measure";
% that deferral dissolves: the measurement target is the TABULAR itself at
% \tabcap ENTRY, while it is still the tail of the open list — the same
% situation \lakin@grabfig handles at an island's end (and the fleet has 0
% split tableboxes). So \lakin@grabtab runs as the FIRST token, before any
% typesetting can push the box away, and the anchor is re-emitted raised above
% the measured tabular. \lakintabraise is a SEPARATE register from
% \lakinfigraise (defined with it, below the LTR-island block): a captionless
% island intentionally leaves the fig register set, and a table must never
% inherit that stale figure height.
%
% WHY the anchor stays AT THE ENTRY SPOT and is NOT deferred into the caption
% paragraph the way \figcap defers it: at \par TeX DROPS a final glue item,
% and the automatic anchor's smash box sitting BEHIND the interword glue kept
% that glue alive — so every centered tabular has always sat half a space off
% true center. Deferring the anchor removed that box and was measured as 982
% differing glyph records: EVERY tabular line shifted 1.30pt (2.61pt for
% pattern 7's two minipages) while captions stayed put. Emitting the raised
% anchor at the old spot — through hyperref's own \Hy@raisedlink, exactly the
% macro the automatic path uses — keeps the line's node structure, and so the
% A/B glyph diff at 0, BY CONSTRUCTION. \Hy@raisedlink reads the raise from
% \HyperRaiseLinkDefault (default \baselineskip = the old mid-table dest); the
% group-local \edef freezes measured-height+42pt into it, and \smash inside
% \Hy@raisedlink hides any raise from the line's extent. \HyCnt@ProvideTheH-
% Counter keeps the dest NAME byte-identical to the automatic one (it defines
% \theHtbl, which \hyper@makecurrent prefers — without it the first measured
% caption was named tbl.1 where the automatic path said tbl.0.1, measured).
% vmode entry, or anything but a box at the tail (no tabular to measure),
% falls back to the plain \refstepcounter — byte-identical to the old path.
\makeatletter
\newcommand{\tabcap}[2][]{\lakin@grabtab
  % v3.79: arm the box-close correction. Set unconditionally and in BOTH
  % branches below's scope, so the plain-\refstepcounter fallback is corrected
  % too; \lakintabboxopen clears it at every \begin{tablebox}, so a \tabcap
  % outside a tablebox can never leave it armed for a later box.
  \global\lakintabcapseentrue
  \ifdim\lakintabraise>0pt
    \H@refstepcounter{tbl}\HyCnt@ProvideTheHCounter{tbl}\hyper@makecurrent{tbl}%
    \begingroup
      % A frameless table targets its measured top. The frame's historical
      % cushion belongs only to tablebox and is corrected at box close.
      \edef\HyperRaiseLinkDefault{\the\lakintabraise}%
      \iflakintabframed
        \edef\HyperRaiseLinkDefault{\the\dimexpr\lakintabraise+42pt\relax}%
      \fi
      \Hy@raisedlink{\hyper@anchorstart{\@currentHref}\hyper@anchorend}%
    \endgroup
    \global\lakintabraise 0pt\relax
  \else
    \refstepcounter{tbl}%
  \fi
  \lakinrtitle{#2}\par\nopagebreak\addvspace{6pt}%
  \def\lakin@short{#1}\ifx\lakin@short\@empty\else
    \addcontentsline{lot}{table}{\protect\numberline{\thetbl}{\protect\ignorespaces #1}}\fi
  {\centering\lakitcaptionstyle\mbox{\lakitcaptionlabelstyle{\tablename~\en{\thetbl}\lakitcaptionpunctuation}}\enspace #2\par}}
\makeatother

% ============================================================================
% Unified FIGURE caption — the \tabcap twin for TikZ/plot figures (v3.34.0).
% A figure (tikzpic) sits inline in the prose; put \figcap BELOW it for a
% numbered, \ref-able caption "איור N.M: תיאור" in the same style as \tabcap:
%   \begin{tikzpic} ... \end{tikzpic}
%   \figcap{\label{fig:foo}description}     % \label FIRST, inside the argument
% WHY a macro, not a hand-written gray line: a manual
%   {\centering\small\textcolor{cGray}{איור: ...}} with a \label inside has NO
% \refstepcounter, so the \label binds to the AMBIENT \@currentlabel (the last
% equation) — \ref then prints an equation number AS the figure number
% (measured on a real page: "איור 48" where 48 is equation 48). \figcap steps a
% real \c@fig counter first, so \label is correct, numbering is consistent, and
% refs are stable. Class-aware like \thetbl (chapter- or section-prefixed).
% ============================================================================
\makeatletter
\@ifundefined{c@chapter}{%
  \newcounter{fig}[section]
  \renewcommand{\thefig}{\ifnum\value{section}>0 \thesection.\fi\arabic{fig}}%
}{%
  \newcounter{fig}[chapter]
  \renewcommand{\thefig}{\ifnum\value{chapter}>0 \thechapter.\fi\arabic{fig}}%
}
\makeatother
\makeatletter
\newcommand{\figcap}[2][]{%
  \ifdim\lakinfigraise>0pt
    % v3.74: step WITHOUT hyperref's automatic anchor, then emit the anchor
    % raised above the measured figure. \H@refstepcounter is hyperref's saved
    % pre-hook \refstepcounter (sets the counter AND \@currentlabel — \label
    % in the caption arg binds exactly as before); \hyper@makecurrent supplies
    % \@currentHref, so the dest NAME is identical to the old automatic one —
    % \ref/\autoref/bookmarks are unchanged, only the dest'S POSITION moves.
    \H@refstepcounter{fig}\hyper@makecurrent{fig}%
    % The anchor is emitted INSIDE the caption paragraph below, as a zero-size
    % raisebox — NOT here in vertical mode, where a \raisebox starts an empty
    % paragraph and adds a blank line (measured: every glyph after the caption
    % shifted +18.8pt). Inside an existing line a [0pt][0pt] box is inert.
    \edef\lakin@anch{\@currentHref}%
    % \edef FREEZES the measured height into the pending macro: the deferred
    % \raisebox is typeset inside the caption paragraph AFTER the register is
    % reset to the 0pt sentinel below — a plain \def read 0 at use time
    % (measured: LAKINRAISE=0.0, dest moved 14pt instead of above the figure).
    \edef\lakin@pend{\noexpand\raisebox{\the\dimexpr\lakinfigraise+42pt\relax}[0pt][0pt]{%
      \noexpand\hyper@anchorstart{\lakin@anch}\noexpand\hyper@anchorend}}%
    \global\lakinfigraise 0pt\relax
  \else
    \refstepcounter{fig}\let\lakin@pend\relax
  \fi
  \lakinrtitle{#2}\par\nopagebreak\addvspace{4pt}%
  \def\lakin@short{#1}\ifx\lakin@short\@empty\else
    \addcontentsline{lof}{figure}{\protect\numberline{\thefig}{\protect\ignorespaces #1}}\fi
  {\centering\lakin@pend\lakitcaptionstyle\mbox{\lakitcaptionlabelstyle{\figurename~\en{\thefig}\lakitcaptionpunctuation}}\enspace #2\par}}
\makeatother

% ============================================================================
% LTR islands.  EVERY tikzpicture / quantikz MUST live inside one, or the whole
% figure's coordinates and Latin labels flip. (Hebrew node text still uses \h{}.)
% ============================================================================
% v3.74 (BLD-37): the \figcap hyperlink DESTINATION used to land at the caption
% line — BELOW the figure — so clicking a figure reference scrolled the figure
% off-screen (measured fleet-wide: 97/97 fig dests in one collection sat 19-54pt
% under their figure's ink; a 347pt figure was fully hidden after the jump).
% hypcap-style fix in two parts: (1) at the island's END, \lakin@grabfig grabs
% the just-typeset centered line with \lastbox, measures it, and puts it back
% with \nointerlineskip — the interline glue above the line was already
% contributed and is NOT re-added, so the page geometry is unchanged to the
% glyph (verified); (2) \figcap then steps the counter WITHOUT hyperref's
% automatic anchor (\H@refstepcounter = the pre-hyperref \refstepcounter) and
% emits the anchor itself, raised by the measured height — the dest lands ABOVE
% the figure top. A captionless island leaves numbering untouched (nothing is
% stepped at the island; the grab flag is cleared at every island begin, so a
% stale height from a captionless figure cannot leak into a later caption).
\newsavebox{\lakinfigbox}
\newdimen\lakinfigraise  % real \dimen: tex.setdimen cannot target a \newlength SKIP register (measured lua error)
% BLD-37 table side: a SEPARATE register for \tabcap. The two must not share —
% a captionless island INTENTIONALLY leaves \lakinfigraise set (the grab runs
% at every island END; only the NEXT island's begin clears it), so a shared
% register would hand a \tabcap after a captionless figure that stale FIGURE
% height. \lakintabraise is written only by \lakin@grabtab at \tabcap entry
% and consumed/reset inside \tabcap itself, so it can never go stale.
\newdimen\lakintabraise
\makeatletter
% v3.74b: measure WITHOUT touching the list. The first draft grabbed the box
% with \lastbox and put it back — and \lastbox DISCARDS the node's baseline
% shift, so a vertically-anchored quantikz circuit rose 2.57pt and everything
% after it sank 5.32pt (measured on the A/B glyph diff). LuaTeX lets us READ
% the tail node's dimensions without removing anything: geometry is identical
% by construction. \lakinfigraise=0pt is the "no figure measured" sentinel.
\def\lakin@grabfig{%
  \ifhmode
    \directlua{
      local n = tex.nest[tex.nest.ptr].tail
      local GLUE, PEN = node.id('glue'), node.id('penalty')
      while n and (n.id == GLUE or n.id == PEN) do n = n.prev end
      local HL, VL = node.id('hlist'), node.id('vlist')
      if n and (n.id == HL or n.id == VL) then
        tex.setdimen('global', 'lakinfigraise',
                     (n.height or 0) + (n.depth or 0))
      else
        tex.setdimen('global', 'lakinfigraise', 0)
      end}%
  \else
    \global\lakinfigraise 0pt\relax
  \fi}
% BLD-37 table side: \lakin@grabfig's twin, NOT a shared helper — the Lua body
% differs in exactly one place, the skip set. Between \end{tabular} and \tabcap
% an otherlanguage wrapper (patterns 3/5 of example_tables) leaves a \write
% whatsit and a dir node ON TOP of the interword glue (measured node dump:
% "whatsit/write <- dir <- glue(2.62) <- hlist"), so the table reader must
% also step over whatsit/dir markers — all zero-size — to reach the tabular.
% \lakin@grabfig deliberately does NOT get the wider skip set: its glue/penalty
% body is the one the v3.74 A/B proof froze, and widening it would newly
% measure (and so newly raise) anchors in island shapes that today take the
% plain path — an unproven behavior change on the FIG side. Anything that is
% not a box behind those markers still hits the 0pt sentinel and \tabcap falls
% back to the plain \refstepcounter.
% SECOND deliberate difference: this stores the box's HEIGHT, not height+depth.
% \figcap's anchor is emitted in the caption paragraph BELOW the figure, so it
% must climb the full extent; \tabcap's anchor sits ON the tabular's own line,
% whose baseline is the vcentered tabular's MIDPOINT — height alone reaches
% the top edge, and adding the depth was measured to overshoot by half the
% table (dests 80-135pt above their frame, INSIDE the preceding table's box).
\def\lakin@grabtab{%
  \ifhmode
    \directlua{
      local n = tex.nest[tex.nest.ptr].tail
      local GLUE, PEN = node.id('glue'), node.id('penalty')
      local WHAT, DIR = node.id('whatsit'), node.id('dir')
      while n and (n.id == GLUE or n.id == PEN
                   or n.id == WHAT or n.id == DIR) do n = n.prev end
      local HL, VL = node.id('hlist'), node.id('vlist')
      if n and (n.id == HL or n.id == VL) then
        tex.setdimen('global', 'lakintabraise', n.height or 0)
      else
        tex.setdimen('global', 'lakintabraise', 0)
      end}%
  \else
    \global\lakintabraise 0pt\relax
  \fi}
\makeatother
% ============================================================================
% v3.79 (BLD-37 table side, CORRECTION): put the anchor at the TOP OF THE BOX,
% instead of guessing the distance to it.
% ============================================================================
% v3.77 raised the \tabcap anchor by (height of the box at the tail of the open
% horizontal list) + 42pt. That is a MEASUREMENT taken BEFORE line breaking, and
% it is short by every term it cannot see. Fleet-measured on the 12 built
% collections (299 tbl.* dests, ownership established through each table's own
% CAPTION, never through "which frame encloses the dest"): 3 dests landed INSIDE
% their own tablebox, below the frame top —
%   computer-architecture       tbl.14.1  -46.4pt   three stacked tabulars
%   computer-architecture       tbl.30.5  -13.5pt   tabular 2.85pt over \linewidth
%   advanced-electronic-devices tbl.B.1   -59.2pt   tabular 2.56pt over \linewidth
% Two missing terms, both reproduced in a minimal document (probe P2/P3):
%   (1) VERTICAL LIST ALREADY CONTRIBUTED inside the box — a second and third
%       tabular, or a heading above the table. The tail read sees only the LAST
%       box, so 77.7pt of stacked tabulars is invisible to the raise.
%   (2) WHICH LINE THE ANCHOR LANDS ON. When the tabular alone overfills
%       \linewidth (459.6121pt in a tablebox), the interword glue after
%       \end{tabular} becomes the chosen breakpoint and \Hy@raisedlink is
%       emitted on the NEXT line — one (depth + lineskip) lower. That term does
%       not EXIST at \tabcap time: line breaking has not happened yet. No
%       measurement taken there can ever recover it. (Proof: probe P3, the same
%       table 2.84912pt too wide, fails at -13.5pt; probe P4, the identical
%       table with \tabcolsep 2pt so it fits, is correct at +32.0pt.)
% So the raise is not computed at \tabcap at all. It is CORRECTED at
% \end{tablebox}, from inside mdframed's own body vbox (endinnercode runs after
% \par\unskip, i.e. after line breaking, while tex.nest[tex.nest.ptr] is still
% that vbox's vertical list). There the anchor's true offset from the top of the
% box content is not estimated but WALKED, and the \Hy@raisedlink shift is set
% so the dest lands 42pt above it — which is 42pt - innertopmargin(9pt) -
% linewidth(0.9pt) = 32.0pt above the FRAME top, the exact position the 295
% already-correct dests have today. That is why this is a correction and not a
% re-placement: the good sites do not move.
%
% WHY THE GLYPH DIFF IS 0 BY CONSTRUCTION — the only thing written is the
% `shift` field of the hlist that \Hy@raisedlink already built, and the guard
% below refuses to write it unless that box is 0x0x0. A zero-size box carries no
% dimension into its parent's packing, the packing already happened, and shift
% is not an input to it. Nothing else in the list is touched: no node is added,
% removed, resized or reordered. Measured, not assumed: assets/example_tables.tex
% and the whole computer-architecture collection A/B at 0 differing pymupdf
% rawdict glyph records.
%
% WHY NOT the alternatives (all prototyped and measured before this one):
%   * accumulate the vertical list at \tabcap time — fixes term (1), CANNOT fix
%     term (2); 2 of the 3 real failures survive. A fix that leaves the majority
%     of the defect standing is the false all-clear this repo exists to avoid.
%   * \zref-savepos two-pass — correct on all shapes, but it makes the dest
%     POSITION aux-derived (DEFECTS BLD-1/2), a single-pass compile silently
%     ships an unraised anchor, and the bare y-difference has no page guard.
%   * emit the anchor at \begin{tablebox} and let \tabcap adopt it — needs the
%     tbl counter stepped at box OPEN to keep the dest NAME byte-identical; a
%     captionless tablebox then either shifts every later table number or leaves
%     an orphan dest that the next table duplicates (a HARD build gate). This
%     correction touches neither the counter nor the name: \tabcap's naming path
%     is v3.77's, unchanged, so resolve_names() is byte-identical (verified).
%   * suppressing the line break (a penalty before the glue) — that MOVES the
%     tabular: the break currently discards the interword glue, so forbidding it
%     re-centers the two overfull tables. Reflow, rejected.
%
% SCOPE. The correction is armed only by \tabcap ITSELF (\lakintabcapseen),
% cleared at every box open, so a captionless tablebox is a no-op and a \tabcap
% hosted OUTSIDE a tablebox (the fleet has exactly one, computer-architecture
% tbl.14.2 in an exbox) keeps v3.77's measured raise. Both fallbacks are
% therefore never worse than today; and because hyperref's AUTOMATIC anchor uses
% the same \Hy@raisedlink shape, the plain-\refstepcounter path (vmode entry, or
% prose between \end{tabular} and \tabcap — probe P6, -51.5pt) is corrected too.
% KNOWN LIMIT — a tablebox that SPLITS across pages. The offset is measured from
% the top of the UNSPLIT body, so a dest in segment 2 gets raised past its own
% segment. MEASURED on a forced split (probe: 40 filler lines + a table): the
% dest went from y=585.0 on page 2 (inside segment 2, at the caption — the old
% BLD-37 symptom) to y=-113.4, i.e. ABOVE the page top. Viewers clamp that to the
% page top, which happens to show segment 2 whole, but an out-of-page coordinate
% is not something to ship on purpose. Hence the \textheight guard in
% lakintab.fix: a body TALLER THAN THE TEXT BLOCK cannot fit on any page, so it
% must split — the correction is skipped and v3.77's behaviour stands unchanged
% (verified: the probe returns to y=585.0). The guard is deliberately EXACT, not
% a page-fitting heuristic: it is sound (a taller-than-page body always splits)
% but PARTIAL (a shorter body that starts low on a page could still split). That
% residue is unreachable in practice — v3.37's everyline + the caption's
% \nopagebreak glue move such boxes WHOLE instead (probe N7: a 22-row table
% pushed down by 430pt moved whole, 0 segments), and the fleet has 0 split
% tableboxes, re-counted after this change. Reimplementing mdframed's own
% page-fitting arithmetic to close the residue would be a heuristic with its own
% failure mode, which is the trade this repo has repeatedly refused.
\newif\iflakintabcapseen
\newif\iflakintabframed
% The Lua lives in ONE \directlua at load time, and is written with NO
% backslash, %, # or ~ : \directlua expands its argument like \edef, so an
% active ~ would expand to \nobreakspace and a % would eat the rest of the line
% (the second is how the first draft of this block died). "not (a == b)" stands
% in for Lua's inequality operator for the same reason.
\directlua{
  lakintab = { }
  lakintab.HL = node.id("hlist")
  lakintab.VL = node.id("vlist")
  lakintab.RU = node.id("rule")
  lakintab.GL = node.id("glue")
  lakintab.KE = node.id("kern")
  lakintab.WH = node.id("whatsit")
  lakintab.TM = node.id("temp")
  lakintab.WD = node.subtype("pdf_dest")
  function lakintab.d(v)
    if v == nil then return 0 end
    if v < -1000000000 then return 0 end
    return v
  end
  function lakintab.scan(list, y, horiz, parent, target)
    local x = list
    while x do
      local i = x.id
      if i == lakintab.HL or i == lakintab.VL then
        if horiz then
          local b = y + x.shift
          if i == lakintab.HL then
            lakintab.scan(x.head, b, true, x, target)
          else
            lakintab.scan(x.head, b - lakintab.d(x.height), false, x, target)
          end
        else
          if i == lakintab.HL then
            lakintab.scan(x.head, y + lakintab.d(x.height), true, x, target)
          else
            lakintab.scan(x.head, y, false, x, target)
          end
          y = y + lakintab.d(x.height) + lakintab.d(x.depth)
        end
      elseif i == lakintab.WH then
        if x.subtype == lakintab.WD then
          local id = x.dest_id
          if type(id) == "string" and string.sub(id, 1, 4) == "tbl." then
            if horiz and parent and parent.width == 0
               and parent.height == 0 and parent.depth == 0 then
              parent.shift = parent.shift + (target - y)
            end
          end
        end
      elseif not horiz then
        if i == lakintab.RU then
          y = y + lakintab.d(x.height) + lakintab.d(x.depth)
        elseif i == lakintab.GL or i == lakintab.KE then
          y = y + lakintab.d(x.width)
        end
      end
      x = x.next
    end
  end
  function lakintab.fix(target, maxh)
    local nst = tex.nest[tex.nest.ptr]
    if nst == nil then return end
    local h = nst.head
    if h == nil then return end
    if h.id == lakintab.TM then h = h.next end
    local tot = 0
    local x = h
    while x do
      local i = x.id
      if i == lakintab.HL or i == lakintab.VL or i == lakintab.RU then
        tot = tot + lakintab.d(x.height) + lakintab.d(x.depth)
      elseif i == lakintab.GL or i == lakintab.KE then
        tot = tot + lakintab.d(x.width)
      end
      x = x.next
    end
    if tot > maxh then return end
    lakintab.scan(h, 0, false, nil, target)
  end
}
% @-FREE NAMES ON PURPOSE: these two are named in the `tablebox` key list at
% \newmdenv above, which is tokenised OUTSIDE any \makeatletter region — a
% \lakin@... there would tokenise as \lakin followed by the letters @tabopen.
\newcommand{\lakintabboxopen}{\lakintabframedtrue\global\lakintabcapseenfalse}
\newcommand{\lakintabboxclose}{%
  \iflakintabcapseen
    \directlua{lakintab.fix(-\number\dimexpr\lakintabcushion\relax,
                            \number\dimexpr\textheight\relax)}%
    \global\lakintabcapseenfalse
  \fi\lakintabframedfalse}
% The cushion above the box's CONTENT top. 42pt puts the dest 32.0pt above the
% frame top (42pt - innertopmargin 9pt - linewidth 0.9pt), i.e. exactly where
% the 295 correct dests already are — the correction moves only the broken ones.
\newcommand{\lakintabcushion}{42pt}
\newenvironment{qcirc}{\par\medskip\global\lakinfigraise 0pt\relax\begin{otherlanguage}{english}\begin{center}}
  {\csname lakin@grabfig\endcsname\end{center}\end{otherlanguage}\par\medskip}
\newenvironment{tikzpic}{\par\medskip\global\lakinfigraise 0pt\relax\begin{otherlanguage}{english}\begin{center}}
  {\csname lakin@grabfig\endcsname\end{center}\end{otherlanguage}\par\medskip}

% ============================================================================
% Code listings.  ★★ v3.4 THE BOX-OVERFLOW FIX:
%   resetmargins=true   → the listing IGNORES the enclosing environment's
%                         indentation and measures from the PAGE. v3.4 stated
%                         this backwards ("measured from the ENCLOSING box"),
%                         and v3.37 then blanked in-box line numbers because
%                         the page-anchored listing put its hanging digits on
%                         top of the box rule. Measured v4.5.0: with the key
%                         FALSE the listing follows the box and the code moves
%                         from x=75.26 to x=97.74, clearing the rule at 72.07.
%                         The key now lives in one \lstset below instead of in
%                         each style, so a box can override it — codebox does.
%   framexleftmargin=0  → the frame (left bar) does not extend left of the code.
%   xleftmargin keeps a small indent so the line numbers have room INSIDE.
% asmblock wraps the listing in an LTR island (code is LTR).
% ============================================================================
\lstdefinestyle{asm}{
  basicstyle=\ttfamily\small,
  backgroundcolor=\color{cBlueBg!55},
  frame=leftline,framerule=2.5pt,rulecolor=\color{cBlue},
  numbers=left,numberstyle=\tiny\color{cGray},numbersep=8pt,
  xleftmargin=2.2em,framexleftmargin=0pt,framexrightmargin=0pt,
  breaklines=true,columns=fullflexible,keepspaces=true,
  showstringspaces=false,aboveskip=8pt,belowskip=8pt,
  % v3.36: without a comment DELIMITER the commentstyle below never fires —
  % listings has no `language` here, so it does not know '#' opens a comment,
  % and every asm listing rendered flat black (measured: a whole-page pixel
  % sweep found ZERO cGreen glyphs across 135 shipped listings). morecomment
  % gives it the delimiter; color-only change, glyph widths untouched.
  morecomment={[l]{\#}},
  commentstyle=\color{cGreen},
}
\lstdefinestyle{cstyle}{
  language=C,basicstyle=\ttfamily\small,
  backgroundcolor=\color{cGreenBg!55},
  frame=leftline,framerule=2.5pt,rulecolor=\color{cGreen},
  numbers=left,numberstyle=\tiny\color{cGray},numbersep=8pt,
  xleftmargin=2.2em,framexleftmargin=0pt,framexrightmargin=0pt,
  breaklines=true,columns=fullflexible,keepspaces=true,
  showstringspaces=false,keywordstyle=\color{cBlue}\bfseries,
  commentstyle=\color{cGray}\itshape,aboveskip=8pt,belowskip=8pt,
}
% v3.36: pystyle/htmlstyle — colored twins of cstyle for the Python/HTML
% listings that used to repeat a ~200-char options blob with NO color
% (keywords rendered bold-black, comments italic-black — measured; C alone
% was colored). Geometry is copied VERBATIM from the shipped blob (frame,
% margins, skips) so swapping a call site to style=pystyle moves NOTHING —
% only keywordstyle/commentstyle/stringstyle are added (bold stays bold,
% italic stays italic, mono metrics identical => zero reflow, proven by
% per-page word-bbox diff). breaklines/breakatwhitespace stay per-site.
% v3.78: numbers=left joins both styles, with the exact numberstyle/numbersep
% keys asm/cstyle carry — fleet-wide consistency (author decision: all four
% styles numbered; "asm/cstyle numbered, py/html not" was an accident of the
% v3.36 geometry freeze, not a design). The numbers hang LEFT of the code box
% (listings sets them with \llap, so no overfull hbox can arise and the code
% x is untouched — measured: code frozen, digits at code_x-11.8pt); inside a
% colored box they are blanked like everyone else's (v3.37 hook), and the
% clean variant restores them in its frameless boxes (v3.78, above).
\lstdefinestyle{pystyle}{
  language=Python,basicstyle=\ttfamily\small,
  columns=fullflexible,keepspaces=true,
  frame=single,rulecolor=\color{cTblLine},
  numbers=left,numberstyle=\tiny\color{cGray},numbersep=8pt,
  xleftmargin=0.5em,framexleftmargin=0pt,framexrightmargin=0pt,
  showstringspaces=false,aboveskip=8pt,belowskip=8pt,
  keywordstyle=\color{cBlue}\bfseries,
  commentstyle=\color{cGray}\itshape,
  stringstyle=\color{cGreen},
}
% v4.5.0: hoisted out of the four styles so an enclosing box can override it.
% Every listing still resolves to true unless something later says otherwise,
% so shipped output is unchanged (verified byte-for-byte on all eight PDFs).
\lstset{resetmargins=true}
\lstdefinestyle{htmlstyle}{
  language=HTML,basicstyle=\ttfamily\small,
  columns=fullflexible,keepspaces=true,
  frame=single,rulecolor=\color{cTblLine},
  numbers=left,numberstyle=\tiny\color{cGray},numbersep=8pt,
  xleftmargin=0.5em,framexleftmargin=0pt,framexrightmargin=0pt,
  showstringspaces=false,aboveskip=8pt,belowskip=8pt,
  keywordstyle=\color{cBlue}\bfseries,
  commentstyle=\color{cGray}\itshape,
  stringstyle=\color{cGreen},
}
\newenvironment{asmblock}{\par\begin{otherlanguage}{english}}{\end{otherlanguage}\par}

% <hllpdf:listing-smart-quotes>
% v4.21.0: four Unicode punctuation characters must enter listings' token
% buffer. Unregistered characters bypass it: a closing quote was printed
% BEFORE the pending word, and an opening quote BEFORE a pending brace.
% Extend the standard character table, then activate only these codepoints
% in the listing's local character-table scope. Do this BEFORE the package's
% delimiter/alsoother/literate hooks so explicit user mappings still win.
% This is not universal Unicode support, nor a conversion to ASCII quotes.
% See listings developer guide 15.3.1 and references/listing-smart-quotes.md.
\makeatletter
\lst@CCPut \lst@ProcessOther {"2018}{"2019}{"201C}{"201D}\z@
\lst@AddToHookAtTop{SelectCharTable}{%
  \lst@ifactivechars
    \catcode"2018=\active \catcode"2019=\active
    \catcode"201C=\active \catcode"201D=\active
  \fi}
\makeatother
% </hllpdf:listing-smart-quotes>

% <hllpdf:listing-technical-symbols>
% v4.30: bounded extension for common technical punctuation in LTR listings.
% Without character-table registration these glyphs bypass the token buffer
% and can move before the preceding word even inside quoted strings.
% Keep the original four-quote block separate for its standing mutation test.
% This is source-faithful printing, not syntax validation or general Unicode.
\makeatletter
\lst@CCPut \lst@ProcessOther
  {"2013}{"2014}{"2026}{"2190}{"2192}{"2194}{"21D2}
  {"2212}{"221E}{"2260}{"2264}{"2265}\z@
\lst@AddToHookAtTop{SelectCharTable}{%
  \lst@ifactivechars
    \catcode"2013=\active \catcode"2014=\active \catcode"2026=\active
    \catcode"2190=\active \catcode"2192=\active \catcode"2194=\active
    \catcode"21D2=\active \catcode"2212=\active \catcode"221E=\active
    \catcode"2260=\active \catcode"2264=\active \catcode"2265=\active
  \fi}
\makeatother
% </hllpdf:listing-technical-symbols>


\usepackage[hidelinks,unicode]{hyperref}   % LAST, after colors. unicode = Hebrew bookmarks.
% v3.62: THE ISOLATION MACROS CORRUPTED EVERY PDF BOOKMARK THEY APPEARED IN.
% \en is defined (line ~155) as
%     \foreignlanguage{english}{\language=\csname l@nohyphenation\endcsname #1}
% and in a heading that body is written verbatim to the .toc and handed to
% hyperref's \pdfstringdef. hyperref strips the two CONTROL SEQUENCES (\language,
% \l@nohyphenation) but the plain `=` CHARACTER token between them is not a control
% sequence, so it survives into the PDF string: \chapter{שער \en{CNOT} והצמדה}
% produced the bookmark "שער =CNOT והצמדה". The printed TOC and the heading on the
% page are CORRECT — only the navigation pane and the viewer window title are wrong,
% which is why it went unnoticed. Measured on the shipped fleet: 70 corrupted
% bookmarks across 5 PDFs, from 177 unprotected \en-in-heading sites. Same corruption
% reached /Title and /Subject via the \hypersetup{pdftitle=...} the docs recommend.
% This hook tells hyperref to use the plain argument for all three isolation macros
% when building a PDF string. It does NOT touch the .toc (the printed TOC is
% byte-unchanged) and it does NOT mask the math signal — a heading with $...$ still
% logs "removing math shift", so BIDI-15 detection is intact.
% BONUS, measured: it also makes \code{} in an UNSTARRED heading COMPILE. BLD-10
% recorded that as needing \texorpdfstring or a starred heading; the real root cause
% is hyperref's \pdfstringdef scan of \code's \ifmmode...\else...\fi body — which is
% exactly why the identically-bodied \texttt survived (hyperref pre-neutralizes
% \texttt in PDF strings) and why \DeclareRobustCommand could never help.
% v3.64: \def\\{ } added. A manual line break in a heading WELDS the two words
% together in the bookmark: \chapter{שורה ראשונה\\שורה שנייה} produced the outline
% entry 'שורה ראשונהשורה שנייה' — "ראשונהשורה" is not a word. hyperref drops the
% \\ and puts nothing in its place, so the break that looks like a space on the
% page is a zero-width join in the navigation pane. The spaced form
% 'מילה אחת \\ מילה שנייה' escapes with a harmless DOUBLE space ('אחת  מילה'),
% which is why this never got noticed. Both are measured on a real outline.
% 0 live instances in the corpus — but SKILL.md and content-style.md both RECOMMEND
% a manual heading break, so this is a primed trap rather than a dormant one, and a
% one-token preamble fix retires it for good instead of asking authors to remember.
% Mapping \\ to a space fixes the WELD, which is the actual defect (a nonexistent
% word in the navigation pane). It does not normalise whitespace: measured after the
% fix, the welded form becomes 'שורה ראשונה שורה שנייה' (correct) and the already-
% spaced form goes from two spaces to three. That is cosmetic — viewers render it as
% a wider gap, and no word is corrupted. A per-heading \texorpdfstring is still the
% way to get an exactly-tuned bookmark string.
\pdfstringdefDisableCommands{\def\en#1{#1}\def\h#1{#1}\def\doclang#1{#1}\def\code#1{#1}\def\codeline#1{#1}\def\\{ }}
% v3.70: RESTORE \url LINE-BREAKING. url.sty freezes its break penalties at LOAD
% time (\mathchardef\UrlBreakPenalty=\binoppenalty), and v3.45's global
% \binoppenalty=10000 (inline-math atomicity) therefore priced EVERY breakpoint
% inside \Url at 10000 = never break. Measured: a long \url in Hebrew prose is
% one unbroken line, ~100pt overfull, invading the margin band the ink gates
% deliberately tolerate. With the two lines below: 0 overfull, the URL breaks
% after '/', the RTL split is CORRECT (fragment order forward), and each line
% gets its own link annotation exactly covering its fragment (measured).
% Confined to \Url's group — inline-math atomicity (TYP-7) is untouched.
\mathchardef\UrlBreakPenalty=200
\mathchardef\UrlBigBreakPenalty=100
\ifx\lakitmainlanguage\hllpdfenglish
  \hypersetup{pdflang=en}
\else
  \hypersetup{pdflang=he}
\fi  % /Lang is the document base locale (not per-span accessibility tagging).
% Set a per-document \hypersetup{pdftitle={...}} too — see SKILL.md.

% v3.52 (round-2 audit): two RTL bibliography/cross-ref defects, both invisible to
% every gate and both affecting the FIRST author to use the feature.
% (a) \@biblabel's default [#1] brackets are bidi-neutral and mirror to ]1[ under
% the RTL bibliography base — every \thebibliography reference list shows ]1[.
% LTR-isolating the label restores [1]. (The render gate's mirror rule fingerprints
% round parens )1(, so square brackets slipped through.)
\makeatletter
\renewcommand*{\@biblabel}[1]{\foreignlanguage{english}{[#1]}}
% v3.54: the SAME neutral-bracket trap on the INLINE side — a multi-key \cite{a,b}
% under the RTL base lays its European-number list right-to-left, so [1, 2] renders
% "[2,1]" and [1, 2, 3] renders "[3,2,1]" (verified by glyph geometry — the reader
% sees the sources in REVERSE order; a note form \cite[p.5]{a} also reorders to
% "[p. 5,1]"). LTR-isolate the whole bracket group (number list AND optional note),
% exactly like \@biblabel above: [1, 2] / [1, 2, 3] / [1, p. 5] all render correct.
% Placed AFTER hyperref so it wins; hyperref hyperlinks each number inside #1 (via
% \@citex), which this wrapper leaves intact. Mirrors the LaTeX kernel \@cite shape.
% v4.9.1: directional isolation must NOT switch the language here. Babel's
% safe-active citation wrapper passes its key list through \bbl@tempa;
% \foreignlanguage reuses that scratch macro before the deferred citation
% loop expands it. The old line therefore cited the literal key "english".
% If such a key existed, even the undefined-reference gate stayed green.
% \babelsublr isolates the bracket group without selecting a locale or
% changing that key list. Preserve native links, key order and optional notes.
\def\@cite#1#2{\babelsublr{[{#1\if@tempswa , #2\fi}]}}
\makeatother

% v4.11.0: a reference inside align's \text inherited HEBdisptag=true,
% so its enclosing hbox mirrored the brackets despite a clean compile.
% Isolate the complete printed reference without selecting a Babel language.
% Also neutralize the display-only normalfont direction reset locally. A
% direction-only patch inverted 12 to 21 inside equation in the measured probe.
% Plain maketag is local to the reference: mathtools' display-side wrapper
% otherwise globally consumes manual_tag while formatting a nested reference.
% Preserve each original command (including the selective auxiliary write),
% current tag forms, native links, and showonlyrefs on/off transitions.
\makeatletter
\NewCommandCopy{\HEB@nativeeqref}{\eqref}
\newcommand{\HEB@eqrefisolate}[1]{%
  \babelsublr{\HEBdisptagfalse\let\bbl@eqnodir\relax\let\maketag@@@\ORIGmaketag@@@#1}}
\protected\def\eqref#1{\HEB@eqrefisolate{\HEB@nativeeqref{#1}}}
\ExplSyntaxOn
\NewCommandCopy{\HEB@nativeMTeqref}{\MT_extended_eqref:n}
\cs_set_protected:Npn \MT_extended_eqref:n #1
  {\HEB@eqrefisolate{\HEB@nativeMTeqref{#1}}}
\ExplSyntaxOff
\makeatother

% v4.11.0: custom caption counters are not part of hyperref's English
% locale, and Hebrew has no upstream hyperref locale in this toolchain.
% Restore the appropriate names at every Babel language switch, including
% return from a nested inline island. Do not localize the destination keys.
\makeatletter
\newcommand{\HEB@refnames}{%
\def\chapterautorefname{פרק}%
\def\sectionautorefname{סעיף}%
\def\subsectionautorefname{תת־סעיף}%
\def\subsubsectionautorefname{תת־תת־סעיף}%
\def\paragraphautorefname{פסקה}%
% v4.31: hyperref uses a distinct counter name at the deepest standard level.
\def\subparagraphautorefname{תת־פסקה}%
\def\equationautorefname{משוואה}%
\def\figureautorefname{איור}%
\def\tableautorefname{טבלה}%
\def\partautorefname{חלק}%
\def\appendixautorefname{נספח}%
\def\footnoteautorefname{הערת שוליים}%
\def\itemautorefname{סעיף}%
\def\Itemautorefname{סעיף}%
\def\figautorefname{איור}%
\def\tblautorefname{טבלה}%
\def\pageautorefname{עמוד}% 
}
\addto\extrashebrew{\HEB@refnames}
\newcommand{\ENG@refnames}{%
  \def\subparagraphautorefname{subparagraph}%
  \def\chapterautorefname{chapter}%
  \def\sectionautorefname{section}%
  \def\subsectionautorefname{subsection}%
  \def\subsubsectionautorefname{subsubsection}%
  \def\paragraphautorefname{paragraph}%
  \def\equationautorefname{Equation}%
  \def\figureautorefname{Figure}%
  \def\tableautorefname{Table}%
  \def\partautorefname{Part}%
  \def\appendixautorefname{Appendix}%
  \def\footnoteautorefname{footnote}%
  \def\itemautorefname{item}%
  \def\figautorefname{Figure}%
  \def\tblautorefname{Table}%
  \def\Itemautorefname{item}%
  \def\pageautorefname{page}%
}
\addto\extrasenglish{\ENG@refnames}
\ifx\lakitmainlanguage\hllpdfenglish\ENG@refnames\else\HEB@refnames\fi
\makeatother

% ============================================================================
% Optional running head (v3.39): \lakithead{<book title>} — OFF by default.
% A document opts in with ONE line after % ============================================================================
% preamble.tex — tested LuaLaTeX preamble for Hebrew + math + figures PDFs (v4.32.0)
% Battle-tested: compiled an 80-page Hebrew quantum-computing summary AND a
% 197-page Hebrew computer-architecture anthology (TikZ diagrams, code listings,
% colored boxes) — both with 0 missing characters.
%
% Usage in your document:
%   \input{preamble.tex}     % (or paste these lines at the top)
%   \begin{document}  ...  \end{document}
%
% Compile with:  lualatex -interaction=nonstopmode doc.tex
%   Repeat until two consecutive PDFs are byte-identical (scripts/build.sh
%   does this) — a TOC document converges only on pass 3+; after two passes
%   the TOC page numbers are STALE (measured; see SKILL.md step 3).
% Font binaries are external dependencies, pinned by assets/fonts/SHA256SUMS.
% scripts/setup_fonts.sh imports a verified local installation or prior skill.
%
% PUBLIC DOCUMENT-SHELL / LAYOUT CONTRACT (introduced v4.3.1; architecture-owned v4.7.1). These hooks may be
% defined BEFORE \input{preamble.tex}; if omitted, the historical canonical
% defaults below are exact.  They deliberately control only the document shell
% and page layout.  RTL/math semantics stay in this preamble, while visual
% presentation themes (for example clean_theme.tex) load AFTER it.
%
%   \def\lakitdocumentclass{article}
%   \def\lakitclassoptions{11pt,a4paper}
%   \def\lakitgeometryoptions{margin=1.8cm,headheight=14pt}
%   \def\lakitlinespread{1.15}
%   \def\lakitparagraphstyle{block} % letters/briefs; indented for continuous prose
%   \def\lakittocdepth{1}
%   \def\lakitsecnumdepth{2}
%   \input{preamble.tex}
%
% Supported class contract: report or article.  Alternate paper sizes must also
% be declared to build.sh through EXPECTED_PAPER; the PDF gate never infers them.
% Do NOT use these hooks as a theme mechanism: theme overlays must not own page
% geometry, class selection, or direction semantics. A declared presentation-policy
% overlay may intentionally change numbering visibility after the preamble.
% ============================================================================
% Runtime architecture guard: generated from scripts/architecture_contract.py.
% The readiness sentinel remains delayed until the END of this file.
% <hllpdf:architecture-guard>
\ifdefined\hllpdfcoreloaded
  \errmessage{hllpdf: preamble.tex loaded more than once}
  \expandafter\endinput
\fi
% </hllpdf:architecture-guard>

\ifdefined\lakitdocumentclass\else\def\lakitdocumentclass{report}\fi
\ifdefined\lakitclassoptions\else\def\lakitclassoptions{12pt,a4paper}\fi
\ifdefined\lakitgeometryoptions\else\def\lakitgeometryoptions{margin=2.3cm,headheight=14pt}\fi
\ifdefined\lakitlinespread\else\def\lakitlinespread{1.3}\fi
\ifdefined\lakittocdepth\else\def\lakittocdepth{1}\fi
\ifdefined\lakitsecnumdepth\else\def\lakitsecnumdepth{2}\fi
% Paragraph organization belongs to the document layout, independently of its
% language and theme. Preserve historical class dimensions unless explicitly
% selecting block paragraphs; never suppress indentation with an everypar hack.
\ifdefined\lakitparagraphstyle\else\def\lakitparagraphstyle{indented}\fi
\edef\lakitparagraphstyle{\lakitparagraphstyle}
\def\hllpdfparagraphblock{block}
\def\hllpdfparagraphindented{indented}
\ifx\lakitparagraphstyle\hllpdfparagraphblock\else
  \ifx\lakitparagraphstyle\hllpdfparagraphindented\else
    \errmessage{hllpdf: paragraph style must be block or indented}
  \fi
\fi
% v4.3.1: the public class hook has a deliberately narrow support contract.
% Before this guard, \def\lakitdocumentclass{book} silently compiled even though
% every heading/TOC regression test targets report/article; an unsupported shell
% could therefore look "supported" until a subtle class-specific defect shipped.
% Fail before \documentclass rather than letting an undocumented combination enter
% the mechanism core. Plain-TeX comparison is used because LaTeX packages are not
% available yet at this point.
\edef\lakitvalidatedclass{\lakitdocumentclass}
\def\lakitreportclass{report}
\def\lakitarticleclass{article}
\ifx\lakitvalidatedclass\lakitreportclass
\else\ifx\lakitvalidatedclass\lakitarticleclass
\else
  \errmessage{hllpdf: unsupported class \lakitdocumentclass; use report or article}
\fi\fi
\let\lakitvalidatedclass\relax
\let\lakitreportclass\relax
\let\lakitarticleclass\relax
% Language intent is independent of document class, theme and content domain.
% The two dominant-language profiles differ in authoring intent, not direction.
% A balanced bilingual document must explicitly choose a navigation/base locale.
% <hllpdf:language-profile>
\providecommand{\lakitlanguageprofile}{hebrew}
\edef\lakitlanguageprofile{\lakitlanguageprofile}
\def\hllpdfhebrew{hebrew}\def\hllpdfenglish{english}
\def\hllpdfhebrewmixed{hebrew-with-english}
\def\hllpdfenglishmixed{english-with-hebrew}
\def\hllpdfbilingual{bilingual}
\let\hllpdfexpectedlanguage\relax
\ifx\lakitlanguageprofile\hllpdfhebrew
  \let\hllpdfexpectedlanguage\hllpdfhebrew
\else\ifx\lakitlanguageprofile\hllpdfhebrewmixed
  \let\hllpdfexpectedlanguage\hllpdfhebrew
\else\ifx\lakitlanguageprofile\hllpdfenglish
  \let\hllpdfexpectedlanguage\hllpdfenglish
\else\ifx\lakitlanguageprofile\hllpdfenglishmixed
  \let\hllpdfexpectedlanguage\hllpdfenglish
\else\ifx\lakitlanguageprofile\hllpdfbilingual
  \ifdefined\lakitmainlanguage\else
    \errmessage{hllpdf: bilingual profile requires explicit lakitmainlanguage}
  \fi
\else
  \errmessage{hllpdf: unsupported language profile \lakitlanguageprofile}
\fi\fi\fi\fi\fi
\ifdefined\lakitmainlanguage
  \edef\lakitmainlanguage{\lakitmainlanguage}
\else
  \let\lakitmainlanguage\hllpdfexpectedlanguage
\fi
\ifx\lakitmainlanguage\hllpdfhebrew\else
  \ifx\lakitmainlanguage\hllpdfenglish\else
    \errmessage{hllpdf: main language must be hebrew or english}
  \fi
\fi
\ifx\hllpdfexpectedlanguage\relax\else
  \ifx\lakitmainlanguage\hllpdfexpectedlanguage\else
    \errmessage{hllpdf: main language conflicts with language profile}
  \fi
\fi
\typeout{HLLLANG|\lakitlanguageprofile|\lakitmainlanguage}
% </hllpdf:language-profile>
\documentclass[\lakitclassoptions]{\lakitdocumentclass}
\usepackage[no-math]{fontspec}        % (1) ★ no-math — THE critical line (see recipe.md)
% layout=lists asks Babel to adapt list parshape when paragraph and page
% directions differ. Without it, minority-language list bodies lose their
% indentation and labels hang outside the text block, despite a clean log.
% The upstream mechanism preserves ordinary same-direction list geometry.
\usepackage[bidi=basic,layout=lists]{babel}
\ifx\lakitmainlanguage\hllpdfenglish
  \babelprovide[import]{hebrew}
  \babelprovide[import,main]{english}
\else
  \babelprovide[import,main]{hebrew}
  \babelprovide[import]{english}
\fi
% v4.22: bind canonical families and fallbacks to the selected installation.
% Lookup by basename or family alone is not proof of the file actually used.
% <hllpdf:font-binding>
\directlua{
  local dir = os.getenv("HEBREW_LATEX_FONT_DIR")
  if not dir or dir == "" then
    local data = os.getenv("XDG_DATA_HOME")
    if not data or data == "" or not (data:sub(1,1) == "/") then
      data = (os.getenv("HOME") or "") .. "/.local/share"
    end
    dir = data .. "/fonts/hebrew-lualatex-pdf"
  end
  if not (dir:sub(1,1) == "/") then dir = lfs.currentdir() .. "/" .. dir end
  if dir:find("[",1,true) or dir:find("]",1,true) or dir:find(string.char(0),1,true)
     or dir:find(string.char(10),1,true) or dir:find(string.char(13),1,true) then
    tex.error("hllpdf: unsupported character in font directory")
  end
  if not (lfs.attributes(dir,"mode") == "directory") then
    tex.error("hllpdf: font directory is missing")
  end
  hllpdf_font_dir = dir:gsub("/+$", "") .. "/"
  tex.sprint(string.char(92) .. "def" .. string.char(92) .. "hllpdffontdir{")
  tex.sprint(-2,hllpdf_font_dir)
  tex.sprint("}")
  luaotfload.add_fallback("hebfb",{"[" .. hllpdf_font_dir .. "DejaVuSerif.ttf]:mode=node;"})
  luaotfload.add_fallback("hebfbbf",{"[" .. hllpdf_font_dir .. "DejaVuSerif-Bold.ttf]:mode=node;"})
}
% </hllpdf:font-binding>
% v3.22: Frank Ruhl Libre ships NO italic face. Instead of silently-upright
% \emph (v<=3.20) or bold-as-emphasis (v3.21), synthesize a slanted face with
% FakeSlant — real, VISIBLE italic emphasis with correct \emph nesting
% semantics (inner \emph toggles back upright). Verified empirically at
% 300dpi: slant clearly visible, 0 errors, 0 missing glyphs.
\babelfont{rm}[Path={\hllpdffontdir},RawFeature={fallback=hebfb},
  BoldFont={FrankRuhlLibre-Bold.ttf},
  BoldFeatures={RawFeature={fallback=hebfbbf}},   % v3.56 — see hebfbbf above
  ItalicFont={FrankRuhlLibre-Regular.ttf}, ItalicFeatures={FakeSlant=0.22},
  % Bold-italic uses the weight-matched fallback and a synthetic slant.
  BoldItalicFont={FrankRuhlLibre-Bold.ttf}, BoldItalicFeatures={RawFeature={fallback=hebfbbf},FakeSlant=0.22}
  ]{FrankRuhlLibre-Regular.ttf}
% Sans fallbacks preserve family and weight for the few missing codepoints.
\directlua{luaotfload.add_fallback("hebfbsans",{"[" .. hllpdf_font_dir .. "DejaVuSans.ttf]:mode=node;"})}
\directlua{luaotfload.add_fallback("hebfbsansbf",{"[" .. hllpdf_font_dir .. "DejaVuSans-Bold.ttf]:mode=node;"})}
\babelfont{sf}[Path={\hllpdffontdir},RawFeature={fallback=hebfbsans},
  BoldFont={Heebo-Bold.ttf},
  BoldFeatures={RawFeature={fallback=hebfbsansbf}}]{Heebo-Regular.ttf}
% v3.24: Scale=MatchLowercase — DejaVu Sans Mono has an outsized x-height, so
% unscaled inline code visually SHOUTS next to the delicate Frank Ruhl body
% (measured on a 197-page build: scaled, a code atom sits in the line at the
% body's weight, and the document tightened by ~5%). Applies to listings too.
\babelfont{tt}[Path={\hllpdffontdir},RawFeature={fallback=hebfb},Scale=MatchLowercase,
  BoldFont={DejaVuSansMono.ttf},
  BoldFeatures={RawFeature={fallback=hebfb},FakeBold=1.5},
  ItalicFont={DejaVuSansMono.ttf},
  ItalicFeatures={RawFeature={fallback=hebfb},FakeSlant=0.2},
  BoldItalicFont={DejaVuSansMono.ttf},
  BoldItalicFeatures={RawFeature={fallback=hebfb},FakeBold=1.5,FakeSlant=0.2}
  ]{DejaVuSansMono.ttf}

% Explicit multilingual/English profiles use real Latin faces and italics.
% The legacy/default Hebrew profile retains its established measured typography.
% Latin Modern is provided by the TeX distribution, not the pinned font bundle;
% recorder/runtime evidence reports it as an additional toolchain resource.
\ifx\lakitlanguageprofile\hllpdfhebrew\else
  \babelfont[english]{rm}[Ligatures=TeX,RawFeature={fallback=hebfb},
    BoldFeatures={RawFeature={fallback=hebfbbf}},
    BoldItalicFeatures={RawFeature={fallback=hebfbbf}}]{Latin Modern Roman}
  \babelfont[english]{sf}[Ligatures=TeX,RawFeature={fallback=hebfbsans},
    BoldFeatures={RawFeature={fallback=hebfbsansbf}},
    BoldItalicFeatures={RawFeature={fallback=hebfbsansbf}}]{Latin Modern Sans}
\fi

\usepackage{amsmath,amssymb,mathtools}
% amsthm supplies theorem/proof machinery. Environment names remain
% document-specific; the generic proof label is localized below.
\usepackage{amsthm}
% Proof names follow the surrounding prose, including return from an island.
\addto\extrashebrew{\renewcommand{\proofname}{הוכחה}}
\addto\extrasenglish{\renewcommand{\proofname}{Proof}}
\ifx\lakitmainlanguage\hllpdfenglish
  \renewcommand{\proofname}{Proof}
\else
  \renewcommand{\proofname}{הוכחה}
\fi
% v3.24.1: equation tags under babel-bidi. A tag re-emitted by mathtools'
% showonlyrefs inside align inherits the RTL base and its parens MIRROR —
% ")2.1(" instead of "(2.1)" (plain align/equation tags are unaffected;
% verified on a minimal case + a 198-page build). LTR-isolate at the tag-BOX
% level: mathtools' show/hide + numbering logic is untouched. Do NOT override
% \tagform@ instead — that clobbers the showonlyrefs mechanism (measured:
% every tag shows, numbering breaks).
% v3.28.1: the v3.24.1 wrap must apply ONLY to tags of display environments.
% Inline \eqref in Hebrew PROSE also goes through \maketag@@@, and there the
% english-wrap MIRRORS the parens — ")1(" — while the unwrapped original is
% correct (proven on a 2x2 isolation matrix: patch x showonlyrefs; regressed
% every prose \eqref since v3.24.1, caught by the render gate's R4). The two
% call sites are indistinguishable by \ifmmode (tag boxes are text mode), so
% the display environments themselves raise a flag; \eqref in prose sees it
% false and gets the original untouched \maketag@@@.
\usepackage{etoolbox}

% babel bidi=basic (including upstream 26.9) assumes every automatic
% discretionary has a one-glyph replacement. LuaTeX can legally produce one
% with replace=nil after a hyphenated Latin value; Babel then indexes nil and the
% whole build dies. Do the narrow upstream-compatible correction: while Babel
% classifies the list, present only that malformed disc as an ordinary
% discretionary, then restore its subtype. Babel still reorders the ENTIRE box;
% there is no pcall and no successful fallback to an unreordered node list.
% If Babel throws for any other reason, the build remains fatal. The synthetic
% nil-disc probe and the real table gallery in test_all.sh pin both layers.
\AtBeginDocument{\directlua{
  if Babel and Babel.bidi and not Babel.heb_nil_disc_compat then
    Babel.heb_nil_disc_compat = true
    local original = Babel.bidi
    local disc_id = node.id('disc')
    Babel.bidi = function(head, ispar, hdir)
      local hidden = {}
      for item in node.traverse_id(disc_id, head) do
        if item.subtype == 2 and (not item.replace or not item.replace.char) then
          table.insert(hidden, item)
          item.subtype = 0
        end
      end
      local result = original(head, ispar, hdir)
      for _, item in ipairs(hidden) do item.subtype = 2 end
      return result
    end
  end
}}
\makeatletter
\newif\ifHEBdisptag
\let\ORIGmaketag@@@\maketag@@@
\def\maketag@@@#1{\ifHEBdisptag\ORIGmaketag@@@{\foreignlanguage{english}{#1}}\else\ORIGmaketag@@@{#1}\fi}
\@for\HEB@env:=equation,equation*,align,align*,gather,gather*,multline,multline*,flalign,flalign*,alignat,alignat*,eqnarray,eqnarray*\do{%
  \expandafter\AtBeginEnvironment\expandafter{\HEB@env}{\HEBdisptagtrue}}
\makeatother
% English page + Hebrew prose: Babel's equation-only tag flip uses a primitive
% leqno path that shifts the formula to the left edge and the tag into the body
% of an English page. Align's full-width alignment is unaffected. Give only this
% display a physical LTR container and native amsmath eqno/leqno. The engine then
% owns real tag width, collision handling and labels; no proxy tag box hides ink
% from its fit calculation. Babel's saved \bbl@eqnodir still restores Hebrew text
% inside math, and the prose locale/math direction are never switched.
% Close the prior RTL paragraph before normalization. After the environment's
% group restores its directions, close TeX's empty resumed display paragraph and
% use native LaTeX continuation handling: same-paragraph prose has no indent,
% while an authored blank line starts an ordinary indented RTL paragraph.
% Register after Babel's setup; remain local to equation(*) in a main-English
% document's actual RTL prose. Native Hebrew, ordinary English and explicit
% fleqn retain their original paths. No page/display dimensions are assigned.
% <hllpdf:minority-rtl-equation-geometry>
\makeatletter
\newcommand{\HEB@minorityprosecontinuation}{\par\@doendpe}
\newcommand{\HEB@minorityequationgeometry}{%
  \ifx\lakitmainlanguage\hllpdfenglish
    \ifnum\pardirection=\@ne
      \if@fleqn\else
        \par\aftergroup\HEB@minorityprosecontinuation
        \iftagsleft@\let\veqno\@@leqno
        \else\let\veqno\@@eqno\fi
        \bodydir TLT\pardir TLT\textdir TLT\relax
      \fi
    \fi
  \fi}
\AtBeginDocument{%
  \AddToHook{env/equation/begin}{\HEB@minorityequationgeometry}%
  \AddToHook{env/equation*/begin}{\HEB@minorityequationgeometry}}
\makeatother
% </hllpdf:minority-rtl-equation-geometry>
% v3.63: \shortintertext rendered Hebrew MIRROR-REVERSED and flush LEFT — an entire
% commentary line inside an align printed backwards. babel's bidi=basic patches
% \intertext@ only (luababel.def), and mathtools' \shortintertext uses the identical
% \vbox{\normalbaselines … \noindent\ignorespaces#1\par} skeleton but is never
% patched, so its paragraph keeps pardir/textdir LTR: no bidi reordering (logical
% order on the page) and flush-left. Measured against \intertext with the same
% Applying babel's own replacement to \shortintertext@ makes the two identical.
% \h{}/\mbox fix the letter order but leave the line flush LEFT, so they are not
% substitutes. Guarded + \AtBeginDocument because mathtools \let-copies at load time.
% NB \ExplSyntaxOn is REQUIRED: mathtools' internals carry expl3 names, and without
% the expl3 catcodes the `_` is read as a subscript and the patch detonates far away
% with "! Illegal unit of measure (pt inserted)" at the \end{align} (measured).
\makeatletter
\AtBeginDocument{%
  \@ifundefined{shortintertext@}{}{%
    \ExplSyntaxOn
    \bbl@sreplace\shortintertext@{\normalbaselines}%
      {\normalbaselines\ifx\bbl@eqnodir\relax\else\bbl@pardir\@ne\bbl@eqnodir\fi}%
    \ExplSyntaxOff}}
\makeatother
% geometry inherits the canonical A4 sheet from \documentclass; build.sh then
% verifies every rendered MediaBox and CropBox, so a silent Letter fallback is
% impossible to publish as a clean default build.
\usepackage[\lakitgeometryoptions]{geometry}
\linespread{\lakitlinespread}       % layout hook; canonical default remains 1.3
\ifx\lakitparagraphstyle\hllpdfparagraphblock
  \setlength{\parindent}{0pt}
  \setlength{\parskip}{6pt plus 1pt minus 1pt}
\fi
% These dimensions remain ordinary scoped TeX lengths: authors can customize
% them after core input. Language switches inherit them; lists keep their own
% indentation and narrow reference layouts establish their local paragraph gaps.
\typeout{HLLPAR|\lakitparagraphstyle}
\usepackage{xcolor}
\usepackage{mdframed}
\usepackage{listings}                 % code listings (assembly / C / pseudo-code)
\usepackage{tikz}
\usetikzlibrary{arrows.meta,positioning,calc,shapes.geometric,fit,backgrounds}
\usepackage{pgfplots}\pgfplotsset{compat=1.18}
\usepackage{quantikz}                 % quantum circuits (remove if unused)
\usepackage{booktabs}
\usepackage{array}
\usepackage{float}                    % [H] = place table exactly here (reference doc: no drifting)
\usepackage{enumitem}
\usepackage{needspace}

% ============================================================================
% ★★ v3.4 — THE DIRECTION-ISOLATION TRIO (the heart of mixing Hebrew + Latin)
% Mental model:
%   \h{...}    Hebrew (RTL) island      — Hebrew text inside an LTR context
%                                          (TikZ nodes, the english `tikzpic`).
%   \en{...}   English (LTR) isolate    — Latin text / parentheses inside a
%                                          Hebrew (RTL) context. KEEPS PARENS
%                                          AND ORDER CORRECT.
%   \code{...} \en + monospace          — code inside Hebrew / inside a frametitle.
% CRITICAL: never wrap Latin in \h{} (it is for HEBREW). Inside a TikZ node,
% Latin in \h{} renders REVERSED (Datapath->htapataD, CPU->UPC) and parens
% mirror. Isolate every Latin/paren chunk with \en{} or \code{}. See
% references/figures-boxes-listings.md.
% ---- TEXT INSIDE MATH (v3.8) ----------------------------------------------
% Put ANY text that sits inside math in a \text{}. The ONLY thing that is
% invalid in math is a bare \h/\en (= \foreignlanguage) used DIRECTLY, outside
% \text -> "\rmfamily invalid in math mode". Inside \text:
%   Hebrew label           -> \text{...}            (bare Hebrew renders fine;
%                              \h is optional — BUT NOT ALWAYS HARMLESS (v3.67):
%                              if the \h group OPENS with a nested island
%                              ($..$, \en{}, \code{}, \texttt{}, \mbox{}) the
%                              WHOLE group loses bidi and the Hebrew prints in
%                              LOGICAL order = gibberish. Measured:
%                                \text{\h{$f$ קבועה}}  -> reads "העובק"  BROKEN
%                                \text{\h{קבועה $f$}}  -> correct
%                                \text{$f$ קבועה}      -> correct (no \h)
%                              Hebrew first inside \h, or drop the \h. Flagged
%                              H-OPENS-WITH-LTR-GROUP.)
%   plain Latin word        -> \text{...}            (no marker needed)
%   Latin WITH ( ) or [ ]   -> \text{\en{...}}       (\en forces LTR; a plain
%                              \text mirrors the parens in RTL -> "Spec (TNR(" )
%   f(Hebrew/text arg)      -> \text{f}(\text{...})  -- keep the parens at MATH
%                              level, NOT \text{f(arg)} (parens inside the box mirror)
%   code in math            -> \mathtt{...}  (or \texttt{}, math-guarded below)
% NEVER write \en{...} or \h{...} straight into $...$ / \[..\] (outside \text);
% always via \text{}. (\mbox{...} also works but does not scale in scripts.)
% ============================================================================
\newcommand{\h}[1]{\foreignlanguage{hebrew}{#1}}
% v3.42.1 — \en islands do NOT hyphenate (auto OR at explicit hyphens). English hyphenation patterns were
% tearing identifiers and terms ACROSS THE RTL SEAM ("IO Regis-/ters",
% "In-/put", "Band-/width", "Per-/formance", "Multi-/ple" — all shipped on
% real pages; probe-proven: "ALUS-/rcB", "Mem-/ToReg"). An island is a label
% or a term, not English prose — a hyphen tear there is never right. A/B on
% a full 197pp lakit: identical overfull sets, same page count, ~11
% paragraphs healed, one mild underfull (badness 3525) that replaced a
% mutilated "Regis-/ters". Long English PROSE belongs in the englishprose
% environment below, which restores explicit-hyphen breaks while its paragraph
% is still active. A raw otherlanguage block inherits the global prohibition.
% (\l@nohyphenation needs \csname:
% '@' is not a letter here.)
% v3.43.2 — \exhyphenpenalty moved OUT of \en to a GLOBAL assignment below.
% Render-proven: a group-local \exhyphenpenalty is DEAD CODE — TeX reads the
% line-breaking parameters at \par time, after the \en group has closed, so
% \en{write-through} / plain not-taken still tore at the "-" ("write-/
% through" shipped on a real exam page). \language is different: it is
% recorded in the node list as text is added, so the in-group nohyphenation
% assignment above DOES work (and stays). The global penalty was A/B-measured
% on a 197pp lakit: identical overfull set (25), identical page count, the
% not-/taken tear healed. Hebrew is unaffected in practice — Hebrew compounds
% use maqaf (־), not ASCII "-".
\newcommand{\en}[1]{\foreignlanguage{english}{\language=\csname l@nohyphenation\endcsname #1}}
% Babel runs extras locally for inline islands and paragraph language changes.
% An English main document must not inherit the Hebrew no-hyphen paragraph rule.
\addto\extrashebrew{\exhyphenpenalty=10000\relax}
\addto\extrasenglish{\exhyphenpenalty=50\relax}
\ifx\lakitmainlanguage\hllpdfenglish
  \exhyphenpenalty=50\relax
\else
  \exhyphenpenalty=10000\relax
\fi
% Full English paragraphs need ordinary explicit-hyphen breaks. Use this
% environment instead of a raw otherlanguage block; the paragraph ends while
% the local penalty is still active.
\newenvironment{englishprose}
  {\par\begin{otherlanguage}{english}\exhyphenpenalty=50\relax}
  {\par\end{otherlanguage}}
\newenvironment{hebrewprose}
  {\par\begin{otherlanguage}{hebrew}\exhyphenpenalty=10000\relax}
  {\par\end{otherlanguage}}
% Use for shell text whose locale must not depend on the last body paragraph.
\DeclareRobustCommand{\doclang}[1]{\foreignlanguage{\lakitmainlanguage}{#1}}

% ============================================================================
% ★★ — fixes for defects that are INVISIBLE in the log and in the checks
% and only surface on a VISUAL RENDER of the PDF. (Each verified empirically.)
% ----------------------------------------------------------------------------
% v3.22: \emph is now REAL slanted emphasis (FakeSlant italic face defined on
% the \babelfont{rm} line above), replacing the v3.21 \emph->\textbf hack: it
% preserves \emph's toggle semantics and keeps bold free for its own role.
% Prefer bold-as-emphasis for your doc? one line: \renewcommand{\emph}[1]{\textbf{#1}}
% page number: force the digits LTR so they never reverse (18 -> 81) in the RTL
%   footer. That reversal is a NON-DETERMINISTIC bidi glitch (observed once in
%   ~80 pp); this LuaLaTeX build is not byte-reproducible, so the defensive
%   LTR wrap is the only reliable guard. Latin digits are always LTR anyway.
\renewcommand{\thepage}{\foreignlanguage{english}{\arabic{page}}}
% v3.70: \pagenumbering \gdef's \thepage and silently DESTROYS the wrap above —
% after \pagenumbering{arabic} the folio guard is gone for the whole rest of the
% document. Two collections discovered this independently and hand-restore the
% wrap in main.tex with the same comment; that knowledge now lives here instead.
% The wrapper re-applies the LTR wrap after the kernel call, for any style.
\makeatletter
\let\lakin@orig@pagenumbering\pagenumbering
\renewcommand{\pagenumbering}[1]{%
  \lakin@orig@pagenumbering{#1}%
  \renewcommand{\thepage}{\foreignlanguage{english}{\csname @#1\endcsname\c@page}}}
\makeatother
% emergencystretch: Hebrew does not hyphenate, so an unbreakable LTR island
%   (\en{...} / an inline math run) landing at the margin can leave a small
%   (~2-13pt) overfull with no legal breakpoint. This rescues ONLY those
%   inline-prose tie-breaks; it is a ceiling used only on offending lines and
%   does NOT touch display-math / listing / table overflow (those still surface
%   in the log for the content-level overflow pass). Verified: baseline overfull
%   -> 0 with this set, with no change to pagination or underfull.
\setlength{\emergencystretch}{3em}
% v3.40: inline math (almost) never breaks across lines. TeX's default
%   (binop 700 / rel 500) lets an inline run split at an operator —
%   $(1.f)\cdot 2^{e}$ shipped as "(1.f)·" ending one line and "2^e" opening
%   the next (reader-reported). In LTR that split is merely ugly; under an
%   RTL base the carried operator lands on the visually WRONG side of the
%   break and the formula reads backwards.
%   WHY 10000 (v3.45.0 — the display-first resolution). Inline math is
%   FORBIDDEN from breaking, so a torn island under an RTL base is
%   structurally impossible. The v3.40 cycle used 9999 ("break only as a last
%   resort") because a bare 10000 left ~30-39 WIDE inline equations
%   overflowing (worst 90.7pt) — but 9999's "last resort" fires OFTEN in
%   tight, non-hyphenating Hebrew: a reader found torn islands across the
%   fleet (~150 candidates; e.g. "P(y |" ending one line and "x_0)" opening
%   the next, read backwards). The step the 9999 cycle missed: every one of
%   those ~30 overflows is a wide equation that BELONGS in display. Convert
%   them to display FIRST, then 10000 gives 0 tears AND 0 overflows (measured,
%   12/12 fleet). Doctrine: inline math is ATOMIC; a wide equation goes in a
%   display, never inline. See content-style ("inline math is atomic").
\binoppenalty=10000 \relpenalty=10000
% v3.38: uniform sentence spacing. TeX's default (\nonfrenchspacing, sfcode
% 3000) is the OLD Anglo-American wide-space-after-period convention — Hebrew
% typography never had it, and modern Latin practice (Chicago et al.) dropped
% it too. Measured on real pages: post-period gaps up to 1.9-2.1x the line's
% mean word gap — the dominant "rivers" source in justified Hebrew prose.
% Survives babel language switches (probed: sfcode stays french inside \en{}
% islands and otherlanguage blocks).
\frenchspacing
% v3.38: widow/orphan protection. LaTeX's default penalties (150/150/50)
% allow a lone line at a page top/bottom almost freely. 10000 = forbid;
% this is SAFE here precisely because the document is \raggedbottom (the
% demoted line just moves on, no vertical stretch fight) — under flushbottom
% these values would cause underfull pages. raggedbottom is the report-class
% default; made explicit so a future class/option change cannot silently
% flip it and turn the 10000s harmful.
\raggedbottom
\widowpenalty=10000 \clubpenalty=10000 \displaywidowpenalty=10000
% em-dash: --- does NOT become — under Frank Ruhl (the font lacks tlig; verified
%   the width is identical with/without Ligatures=TeX, while the mechanism works
%   on Latin Modern). TYPE THE CHARACTER — (U+2014) DIRECTLY; likewise – (U+2013).
%   check_bidi_figures.py flags a --/--- that is adjacent to Hebrew.
% ============================================================================

% ---- STRUCTURAL DIRECTION-SAFETY for \texttt (v3.7) -------------------------
% A parenthesis, bracket, or $ INSIDE \texttt{} mirrors/migrates in RTL prose
% (MIPS 0($s0)->"(s0$)0", $rs->"rs$", arr[i]->"[arr[i") because \texttt sets the
% font but NOT the direction. Redefining \texttt to LTR-isolate its content fixes
% the WHOLE class at once (parens, brackets, $) so plain \texttt{lw $t0,0($s0)}
% is safe — no need to remember \code{}. Guarded by \ifmmode because
% \foreignlanguage is invalid in math (there, plain monospace is already LTR).
% NB (corrected v3.56 — the old note here said "stays tofu", which is WRONG):
% \texttt{Hebrew} is a FATAL build failure. The tt family is DejaVu Sans Mono
% (zero Hebrew) and the tt fallback hebfb is DejaVuSerif (also zero Hebrew), so
% the subsetter gets an empty glyph set and LuaTeX dies with
%   "! error: (file .../DejaVuSansMono.ttf) (type 2): there are no glyphs in the
%    subset"  -> no output PDF.
% The message names a FONT FILE, never Hebrew/\texttt/the source line. And it is
% non-local: add an unrelated Latin \code{} elsewhere and the fatal degrades into
% a SILENT DELETION of the Hebrew word. Put Hebrew in \h{}/normal text, never in
% \texttt{}/\code{}; flagged HEBREW-IN-MONO. To disable, delete this block.
\let\origtexttt\texttt
\protected\long\def\texttt#1{\ifmmode\origtexttt{#1}\else\foreignlanguage{english}{\origtexttt{#1}}\fi}
% v4.8.0: defer the direction test until TeX has inserted an alignment cell's
% math template. Unprotected expansion here tested \ifmmode too early at the
% FIRST token of aligned/array/cases/matrix cells. Engine-level \protected also
% preserves the original \texttt implementation without the self-recursion
% observed when retrofitting \DeclareRobustCommand to its saved alias.
\protected\long\def\code#1{\ifmmode\origtexttt{#1}\else\foreignlanguage{english}{\origtexttt{#1}}\fi}

% ---- \codeline — display-code line (v3.24): the code analogue of display math.
% Reader-friction rule (content-style.md §5b): inline \code{} is for a SINGLE
% ATOM only (one register/variable/mnemonic, a tight no-space expression).
% A multi-token expression — an instruction with operands, a spaced assignment,
% a register-binding list — zigzags the eye RTL<->LTR mid-sentence and, being
% an unbreakable island, line-breaks badly. Give it its own centered line with
% display-math vertical rhythm. \centerline = one box => the expression never
% wraps; \nopagebreak keeps it glued to its introducing sentence.
%   ...the Hebrew sentence ends with a colon:
%   \codeline{lw \$t1, -4(\$t0)}
%   and the prose continues (escape dollars as \$).
% v3.66 (first line of the body): flush a PENDING \item label before the \par.
% LaTeX defers the label box (\@labels) until the next \everypar, and an \hbox in
% vertical mode never fires \everypar — so `\item \codeline{...}` printed the code
% at y=69.1 and the item's own "1." BELOW it at y=103.6 (the number appears after
% its content, attached to nothing). \if@inlabel\leavevmode\fi is CONDITIONAL on
% purpose: an unconditional \leavevmode would open an empty paragraph line when
% \codeline follows ordinary prose. Verified both ways — in a list the label now
% sits on the code's line at its normal x, and after prose the geometry is
% byte-identical to before.
\makeatletter
\newcommand{\codeline}[1]{%
  \if@inlabel\leavevmode\fi
  \par\nopagebreak
  \vspace{\abovedisplayskip}\nopagebreak
  % v3.62: was \centerline{...}, i.e. \hbox to\hsize{\hss #1 \hss}. \hss has
  % INFINITE SHRINK, so TeX could always "fit" the box and NEVER reported an
  % Overfull — the same \hss mechanism as BLD-27's \item label. Measured: a wide
  % \codeline in a defbox painted its first token ON the box's left border bar and
  % sat 9.5pt outside the frame interior on both sides, with 0 Overfull, 0 Underfull
  % and all four gates PASS; in bare prose it reached 21.0pt into BOTH margins.
  % \hfil has no shrink, so the overflow is reported and build.sh's ovf_vis/ovf_big
  % counters see it; \linewidth (not \hsize) makes it box-aware, so a \codeline
  % inside a colored box centres on the box interior. Verified: a FITTING \codeline
  % keeps byte-identical glyph positions, and the box-crossing case now emits
  % "Overfull \hbox (19.03pt too wide)". A wide \codeline should be split in two.
  \hbox to\linewidth{\hfil\code{#1}\hfil}%
  \nopagebreak\vspace{\belowdisplayskip}}
\makeatother

% ---- a thin decorative rule helper ----
\newcommand{\sep}{\par\vspace{2pt}\noindent\textcolor{cRule}{\rule{\linewidth}{0.6pt}}\par\vspace{4pt}}

% ---- colors ----
\definecolor{cBlue}{HTML}{1F4E79}   \definecolor{cBlueBg}{HTML}{EAF1F8}
\definecolor{cGreen}{HTML}{2E6B4F}  \definecolor{cGreenBg}{HTML}{EAF4EE}
\definecolor{cAmber}{HTML}{9C6A00}  \definecolor{cAmberBg}{HTML}{FBF3E2}
\definecolor{cPurple}{HTML}{5B2C6F} \definecolor{cPurpleBg}{HTML}{F3ECF7}
\definecolor{cTeal}{HTML}{0E6E6E}   \definecolor{cTealBg}{HTML}{E6F4F4}
\definecolor{cRed}{HTML}{8E2B2B}    \definecolor{cRedBg}{HTML}{FBEAEA}
\definecolor{cGray}{HTML}{555555}   \definecolor{cRule}{HTML}{1F4E79}

% ---- math operators ----
\DeclareMathOperator{\Tr}{Tr}
\DeclareMathOperator{\rank}{rank}
\DeclareMathOperator{\diag}{diag}
\DeclareMathOperator{\sgn}{sgn}

% ---- bra-ket (quantikz defines \ket/\bra/\braket; add the rest safely) ----
\providecommand{\ket}[1]{\left|#1\right\rangle}
\providecommand{\bra}[1]{\left\langle#1\right|}
\providecommand{\braket}[2]{\left\langle#1\middle|#2\right\rangle}
\newcommand{\dyad}[1]{\left|#1\right\rangle\!\left\langle#1\right|}
% v3.21: general outer product |a><b| and expectation value <A> — agents reach
% for these by reflex (esp. \expval); quantikz defines neither, so their absence
% threw "Undefined control sequence" in a real build. \providecommand = safe.
\providecommand{\ketbra}[2]{\left|#1\right\rangle\!\left\langle#2\right|}
\providecommand{\expval}[1]{\left\langle#1\right\rangle}
\newcommand{\abs}[1]{\left|#1\right|}
\newcommand{\norm}[1]{\left\lVert#1\right\rVert}
\newcommand{\half}{\tfrac12}
\newcommand{\R}{\mathbb{R}}\newcommand{\C}{\mathbb{C}}\newcommand{\Z}{\mathbb{Z}}
% identity operator. NB: \mathbb{1} silently renders a WRONG glyph (msbm has no
% blackboard digits). Use \mathbb{I} = 𝕀 (works with amssymb, no extra package).
% Do NOT "fix" this back to \mathbb{1}.
\newcommand{\id}{\mathbb{I}}

% ---- physical units (v3.35): \punit — upright unit with thin space ---------
% Usage: $300\punit{K}$, $80\punit{mV/decade}$, $10^{7}\punit{cm/s}$ — the
% NUMBER AND UNIT LIVE IN ONE MATH ISLAND. Promoted verbatim from a 199-page
% production lakit (advanced-electronic-devices, dozens of uses).
% Two measured traps this design avoids (both render-proven):
%   1. siunitx SCRAMBLES under RTL — \SI{3.5}{eV} renders "eV 3.5", exponents
%      detach ("Hz 10 x 10^9"): its quantity assembly is not bidi-isolated.
%      Do NOT load siunitx in this pipeline.
%   2. A text-mode unit next to a bare-prose number is TWO adjacent LTR runs
%      and reorders ("5 mm" -> "mm5" — the BIDI-2/3 class). \punit is therefore
%      math-only ON PURPOSE: prose misuse fails LOUDLY at compile (\mathrm
%      outside math) instead of scrambling silently.
% Domain shorthands (\MeV, \fm...) stay per-document: \newcommand{\MeV}{\,\text{\en{MeV}}}
\providecommand{\punit}[1]{\,\mathrm{#1}}

% ============================================================================
% Colored boxes — mdframed (NOT tcolorbox, which breaks under bidi).
% ★★ v3.4: innerleftmargin/innerrightmargin give code listings room so their
% frame/line-numbers do not spill out the box border (see listing styles below).
% ============================================================================
% v3.61: everyline/split skips were applied to `tablebox` only (BLD-14, v3.37) —
% the six colored boxes below never got them, so they still split OPEN across a page
% break: the page-N segment gets top+sides and NO bottom rule, the page-N+1 segment
% sides+bottom and NO top rule. Measured on the shipped fleet this is a COVERAGE GAP,
% not a re-report of BLD-14: 0 split-open `tablebox` segments (that fix holds) vs 62
% split-open COLORED boxes across 7 documents. Verified after the fix: both segments
% close, and the frametitle is NOT duplicated on the continuation. BLD-14's caveat
% still applies — a box taller than a full page grinds mdframed's split-trial loop, so
% split such content manually.
% ============================================================================
% v3.72: frametitle={} leads EVERY \newmdenv option list below. mdframed options
% are GROUP-scoped: a box nested inside \begin{exbox}[frametitle={...}] used to
% INHERIT the outer title — measured, the identical title string rendered twice,
% the second riding a title-bar fill inside the INNER box's frame, for all seven
% environments in both directions. Both natural forms were broken (= preamble
% bug, not an authoring rule). The explicit reset makes every box titleless
% unless ITS OWN [frametitle={...}] says otherwise; verified both directions:
% no duplicate on the nested box, an explicit title still renders.
% Geometry is shared with the optional semantic-width table container. Themes
% select their own existing semantic style below; neither path fixes columns.
\mdfdefinestyle{hllpdfsemanticgeometry}{linewidth=1.1pt,
  innertopmargin=6pt,innerbottommargin=7pt,skipabove=9pt,skipbelow=9pt,
  innerleftmargin=9pt,innerrightmargin=9pt,everyline=true,
  splitbottomskip=6pt,splittopskip=12pt,roundcorner=3pt}
\newcommand{\hllpdfsemanticframestyle}{hllpdfsemanticgeometry}
\newmdenv[frametitle={},style=hllpdfsemanticgeometry,linecolor=cBlue,backgroundcolor=cBlueBg,
  frametitlebackgroundcolor=cBlue,frametitlefontcolor=white,
  ]{defbox}   % הגדרה
\newmdenv[frametitle={},style=hllpdfsemanticgeometry,linecolor=cGreen,backgroundcolor=cGreenBg,
  frametitlebackgroundcolor=cGreen,frametitlefontcolor=white,
  ]{thmbox}   % משפט/טענה
\newmdenv[frametitle={},style=hllpdfsemanticgeometry,linecolor=cAmber,backgroundcolor=cAmberBg,
  frametitlebackgroundcolor=cAmber,frametitlefontcolor=white,
  ]{notebox}  % הערה
\newmdenv[frametitle={},style=hllpdfsemanticgeometry,linecolor=cPurple,backgroundcolor=cPurpleBg,
  frametitlebackgroundcolor=cPurple,frametitlefontcolor=white,
  ]{exbox}    % דוגמה
\newmdenv[frametitle={},style=hllpdfsemanticgeometry,linecolor=cRed,backgroundcolor=cRedBg,
  frametitlebackgroundcolor=cRed,frametitlefontcolor=white,
  ]{warnbox}  % טעות נפוצה
\newmdenv[frametitle={},style=hllpdfsemanticgeometry,linecolor=cTeal,backgroundcolor=cTealBg,
  frametitlebackgroundcolor=cTeal,frametitlefontcolor=white,
  ]{keybox}   % רעיון מרכזי

% ---- codebox (v4.5.0): the dedicated container for code -------------------
% An OPTION, not a policy: asmblock unboxed stays canonical for a full listing,
% and a snippet inside any of the six colored boxes still works as before.
% What makes codebox worth its own environment is that it is the ONE framed box
% that keeps its margin line numbers.
%
% v4.4.0 shipped it WITHOUT them, on the strength of v3.37's "no margin knob can
% save them" plus four fresh probes that agreed. That was wrong: every probe had
% moved the code, the box, or both — never the key that carries the NUMBER
% COLUMN. Measured on the page:
%   resetmargins=false   the listing follows THIS box instead of the page,
%                        moving the code from x=75.26 to x=97.54
%   xrightmargin=26pt    under RTL this is the PHYSICAL LEFT indent, and the
%                        only key that carries the \llap'ed digits with it:
%                        x=64.28 -> 90.18, clearing the rule at 65.64 by 23.6pt
%   xleftmargin=26pt     moves the code and LEAVES THE DIGITS on the rule —
%                        the asymmetry that hid the fix for eleven versions
% resetmargins is hoisted out of the four styles (see the \lstset below them) so
% this override is possible at all; every other listing still resolves to true.
% Pinned by tests/codebox_numbers.tex + probe_lstnum.py --inside-rule, whose
% negative control strips xrightmargin and watches the digits fall back.
\definecolor{cSlate}{HTML}{3F4A56}  \definecolor{cSlateBg}{HTML}{F2F4F7}
\newmdenv[frametitle={},leftline=true,rightline=true,topline=true,bottomline=true,
  linecolor=cSlate,backgroundcolor=cSlateBg,linewidth=0.9pt,
  frametitlebackgroundcolor=cSlate,frametitlefontcolor=white,
  innertopmargin=6pt,innerbottommargin=7pt,skipabove=9pt,skipbelow=9pt,
  innerleftmargin=9pt,innerrightmargin=9pt,everyline=true,splitbottomskip=6pt,splittopskip=12pt,
  roundcorner=3pt]{codebox}   % קוד
% The two keys that make in-box numbering work, both measured on the page:
%   resetmargins=false  the listing follows THIS box instead of the page
%   xrightmargin        under RTL this is the PHYSICAL LEFT indent, and it is
%                       the only one that carries the \llap'ed number column
%                       with it. xleftmargin moves the code and leaves the
%                       digits behind on the rule (measured: code 97.54,
%                       digits 64.28 — the v3.37 defect, reproduced); with
%                       xrightmargin=26pt the digits land at x=90.18, a
%                       clearance of 23.6pt from the rule at 65.64.
% This asymmetry is why v3.4's note had the semantics backwards and v3.37
% concluded the digits could not be saved.
% <hllpdf:listing-container-state>
% v4.29: all state is local to the current environment group. A semantic
% callout masks number text, while a nested codebox restores the formatter
% saved before that mask. Save it once across successive semantic wrappers;
% an explicit numbers=none (or an explicitly empty formatter) remains valid.
\makeatletter
\newif\ifhllpdf@incodebox
\newif\ifhllpdf@englishcodebox
\newif\ifhllpdf@lstnumbermasked
\let\hllpdf@savedlstnumber\thelstnumber
\makeatother
% </hllpdf:listing-container-state>
% <hllpdf:codebox-gutter>
\makeatletter
\AtBeginEnvironment{codebox}{%
  \hllpdf@incodeboxtrue
  % Snapshot the OUTER language before asmblock selects English.
  \iflanguage{english}{\hllpdf@englishcodeboxtrue}{\hllpdf@englishcodeboxfalse}%
  \ifhllpdf@lstnumbermasked
    \let\thelstnumber\hllpdf@savedlstnumber
    \hllpdf@lstnumbermaskedfalse
  \fi
}
% listings' AddToHook is GLOBAL. Register once, and gate on local box state;
% adding an unconditional hook at every codebox entry leaks into later listings.
\lst@AddToHook{PreInit}{%
  \ifhllpdf@incodebox
    % Apply every listing override here, inside the LISTING's own group.
    % Setting margins at box entry leaks them into nested semantic callouts.
    \lstset{resetmargins=false,xrightmargin=26pt}%
    % Styles resolve before PreInit. In LTR the physical number gutter follows
    % xleftmargin; in Hebrew retain the measured inherited xrightmargin path.
    \ifhllpdf@englishcodebox
      \lstset{xleftmargin=26pt,xrightmargin=0pt}%
    \fi
    % ONE box, not two: suppress only the frame/background/skips of a listing
    % inside this codebox. The next unboxed listing keeps its own style.
    \lstset{frame=none,backgroundcolor=\color{cSlateBg},aboveskip=0pt,belowskip=0pt}%
  \fi
}
\makeatother
% </hllpdf:codebox-gutter>

% v3.37: line numbers inside a COLORED BOX are sliced by the box border —
% measured: digit span x=63.5-67.3 vs border line at x=65.2 (the border runs
% THROUGH the digits), and the inter-border gap (3.5pt) is narrower than a
% digit (3.8pt), so no margin knob can save them (under resetmargins,
% xleftmargin is inert in-box — measured: code at x=75.26 with 2.2em AND
% 3.0em). A 3-6 line boxed snippet does not need line numbers; blank the
% number TEXT inside every colored box (\thelstnumber is what gets printed;
% the style's numbers=left survives, so nothing else shifts — code position
% byte-frozen at x=75.26, verified). Unboxed listings keep their numbers.
% v3.78: the blanking is coupled to the FRAME, not to the box name —
% clean_theme.tex registers LATER hooks (registration order wins) that
% restore the number text inside the boxes it renders frameless.
% v4.5.0 RETRACTS the "no margin knob can save them" clause above. It holds only
% under resetmargins=true (the shipped default): there xleftmargin is inert
% in-box, as measured. With resetmargins=false, xrightmargin DOES carry the
% \llap'ed digit column clear of the rule — measured; the codebox \lstset that
% does it is the one above, in the codebox block. The loop below therefore
% covers SIX boxes, not "every colored box": codebox opts out and keeps its
% numbers.
% <hllpdf:semantic-listing-mask>
\makeatletter
\@for\HEB@lstbox:=defbox,thmbox,notebox,exbox,warnbox,keybox\do{%
  \expandafter\AtBeginEnvironment\expandafter{\HEB@lstbox}{%
    \ifhllpdf@lstnumbermasked\else
      \let\hllpdf@savedlstnumber\thelstnumber
    \fi
    \hllpdf@lstnumbermaskedtrue
    \def\thelstnumber{}%
    % The nearest semantic container owns its listing presentation. On exit,
    % TeX's group restores the enclosing codebox state and formatter.
    \hllpdf@incodeboxfalse}}
\makeatother
% </hllpdf:semantic-listing-mask>

% ============================================================================
% Unified TABLE presentation: subtle uniform slate frame + auto-numbered
% caption rendered INSIDE the frame. Use:
%   \begin{tablebox}\centering  <tabular>  \tabcap{<description>}\end{tablebox}
% Replaces the buggy \caption* (which emitted a stray "טבלה N : *" line) and
% provides an explicit framed option + numbered caption "טבלה N.M: תיאור".
% Optional intro text may be placed at the TOP, before the tabular, inside the box.
% ============================================================================
\definecolor{cTblLine}{HTML}{B9C4D2}  \definecolor{cTblBg}{HTML}{FAFBFC}
% v3.23: class-aware — under `article` there is no chapter counter (the promised
% report->article switch used to die with "No counter 'chapter' defined").
\makeatletter
\@ifundefined{c@chapter}{%
  % v3.31: in article the counter resets per SECTION, so an unprefixed number
  % repeated ("טבלה 1." in every section — measured on a two-section render).
  % Same conditional idea as the report branch: prefix the section number once
  % a real \section exists; a pre-section table stays "טבלה 1".
  \newcounter{tbl}[section]
  \renewcommand{\thetbl}{\ifnum\value{section}>0 \thesection.\fi\arabic{tbl}}%
}{%
  \newcounter{tbl}[chapter]
  % v3.30: a report-class doc that uses SECTIONS ONLY (chapter stays 0 —
  % the common short-doc pattern) used to print "טבלה 0.1"; the chapter
  % prefix now appears only once a real \chapter exists (render-verified
  % both ways: "טבלה 1." chapterless, "טבלה 1.1" with chapters).
  \renewcommand{\thetbl}{\ifnum\value{chapter}>0 \thechapter.\fi\arabic{tbl}}%
  % v3.64: the SAME leak in \thesection itself, left behind by v3.30. That fix
  % guarded the CAPTION counter and stopped there, so a chapterless report still
  % printed the heading "0.1 סעיף ראשון" and \ref to it printed "0.2" — measured:
  %   p1 y=65.2  '…0.1'   p1 y=144.9 '…0.2'   p1 y=183.4 '…0.2.1'
  % report.cls hardcodes \thesection as \thechapter.\@arabic\c@section, and the
  % leading 0 is the chapter counter that no \chapter ever incremented. Fixing
  % \thesection fixes the whole chain — \thesubsection and \thesubsubsection are
  % defined in terms of it. \theequation and \thetbl/\thefig were already guarded,
  % which is exactly why this survived: every NUMBER around the heading looked
  % right, so the heading's own 0 read as a style choice rather than a bug.
  % \appendix is safe: report.cls resets chapter to 0 and redefines \thechapter to
  % \@Alph\c@chapter, so an appendix \chapter makes the counter 1 and the guard
  % emits "A.1" unchanged. Render-verified in all four modes (chapterless,
  % chaptered, appendix+chapter, appendix+section-only).
  \renewcommand{\thesection}{\ifnum\value{chapter}>0 \thechapter.\fi\arabic{section}}%
}
\makeatother
% v3.37: everyline=true — if the box DOES split across pages, every segment
% gets a CLOSED frame (measured before: the split ran open-ended into the page
% edge with no bottom border, and the caption could strand alone at the next
% page top — "טבלה A.2" shipped that way). The split skips give the segments
% breathing room. In practice, with the caption GLUE below, every previously-
% split table in the corpus now moves WHOLE instead (the white-gap tradeoff,
% documented benign); everyline stays as the safety net for a genuinely
% page-tall table. NB: a table taller than ONE FULL PAGE can grind mdframed's
% split-trial loop (reproduced: minimal doc, minutes of 1pt-step Overfull
% probes) — break such a monster into two tables.
\mdfdefinestyle{hllpdftablegeometry}{leftline=true,rightline=true,topline=true,bottomline=true,
  linewidth=0.9pt,
  innertopmargin=9pt,innerbottommargin=8pt,skipabove=11pt,skipbelow=11pt,
  innerleftmargin=11pt,innerrightmargin=11pt,roundcorner=4pt}
\newcommand{\hllpdftableframestyle}{hllpdftablegeometry}
\newmdenv[frametitle={},style=hllpdftablegeometry,
  linecolor=cTblLine,backgroundcolor=cTblBg,
  % v4.2.0 (adopted): tablebox is ATOMIC. Short tableboxes near page bottoms
  % repeatedly entered mdframed's split trial, emitting benign Overfull-vbox
  % bursts and occasionally separating the intro prose from the tabular.
  % nobreak=true moves the complete short table to the next page instead. The
  % output-active Overfull-vbox gate still catches a box taller than a page;
  % everyline/split skips remain as defensive styling should that path ever
  % re-arm, but the canonical path never publishes a split tablebox.
  nobreak=true,
  everyline=true,splitbottomskip=6pt,splittopskip=12pt,
  % v3.79 (BLD-37 table side, correction): startinnercode/endinnercode run
  % INSIDE mdframed's body vbox (mdframed.sty:848 / :856+863) -- the second one
  % after \par\unskip, i.e. after line breaking, which is the only moment the
  % \tabcap anchor's true offset from the top of the box EXISTS. See
  % \lakintabboxclose below. These two keys are the ONLY tablebox change.
  startinnercode={\lakintabboxopen},endinnercode={\lakintabboxclose}]{tablebox}
% v3.37: \nopagebreak GLUE — \vsplit honors penalties, so the caption can no
% longer strand alone at the top of the next page (it takes its rows along, or
% the whole box moves). Render-verified on the real shipped orphan.
% v3.59: \NR@gettitle is what nameref injects into the stdlib title-bearing
% commands (\@sect, \@chapter, \caption) to fill \@currentlabelname. \figcap and
% \tabcap are INDEPENDENT macros — they never called it, so \@currentlabelname kept
% whatever the last heading left and \nameref{fig:x} silently printed the enclosing
% SECTION's title instead of the caption (proven in the .aux: field 3 of \newlabel
% was the heading text while field 1 and the anchor `fig.1.1` were correct — so this
% is a THIRD, distinct head next to BLD-13's counter fix and BIDI-38's autoref-name
% fix, neither of which writes that field). \@ifundefined keeps it inert if
% hyperref/nameref is ever dropped. Guarded helper because \figcap is defined
% after \makeatother.
% NB both helpers below need \makeatletter — this part of the preamble is OUTSIDE
% any letter-catcode region, so an unwrapped \@currentlabel parses as LaTeX's \@
% and dies with "You can't use \spacefactor in vertical mode" (measured).
\makeatletter
\newcommand{\lakinrtitle}[1]{\@ifundefined{NR@gettitle}{}{\NR@gettitle{#1}}}
% v3.59: \hitem{<letter>}{<key>} — a REFERENCEABLE Hebrew-lettered list item.
%   \begin{enumerate} \hitem{א}{it:a} הסעיף הראשון. \end{enumerate}  ->  \ref{it:a} = (א)
% Why this exists: automatic Hebrew list labels fail silently, so SKILL.md mandates
% an EXPLICIT \item[(\textbf{א})] — but LaTeX puts \refstepcounter INSIDE
% \if@noitemarg (latex.ltx:12964), and \item[...] never sets that flag. So an
% explicit item steps no counter and never touches \@currentlabel: a \label there
% binds to the AMBIENT one and \ref prints the CHAPTER/SECTION number where the
% item letter belongs (proven in the .aux: {chapter.1}). Same doctrine as \figcap —
% own the label state before \label. \phantomsection matters: without it the
% hyperlink still lands on the chapter head (measured); with it the anchor is
% section*.N and the bookmark outline stays clean.
\newcommand{\hitem}[2]{\item[(\textbf{#1})]%
  \def\@currentlabel{(#1)}\def\@currentlabelname{#1}\phantomsection\label{#2}}
% NB (v3.66): \hitem is \ref-only. \autoref{key} calls it a SECTION ("סעיף (א)")
% because \phantomsection brands the anchor section*.N, and \nameref{key} prints
% the bare letter. For \autoref-able items prefer \hebenumlabels below.
%
% v3.66: \hebenumlabels — OPT-IN automatic Hebrew enumerate letters, done right.
%   {\hebenumlabels\begin{enumerate}\item\label{q:a} ...\end{enumerate}}
%   -> labels (א) (ב) (ג); \ref{q:a} = א; anchor Item.N (real, per-item).
% This retires the belief that made \item[(\textbf{א})] + \hitem necessary.
% "Automatic Hebrew labels fail silently" was TRUE but unexplained, and the
% mechanism is now bisected (plain article vs +enumitem, measured): ENUMITEM
% PRE-EXPANDS the label expression in an \edef BEFORE \refstepcounter, with the
% counter still at 0. Tokens it shields by name (\arabic,\alph,\roman) survive;
% a plain \newcommand \ifcase wrapper evaluates at 0 -> the EMPTY zero-branch ->
% the label prints as "()" with the letter GONE (and \the\c@enumi prints a
% literal 0, \number\value dies fatally pointing at \item). The cure is one
% property, not a doctrine: the wrapper must be ROBUST (survives the \edef
% unexpanded), while \theenumi stays EXPANDABLE (so \refstepcounter writes the
% real letter to the .aux). Verified: labels (א)(ב)(ג), \ref -> א/ב/ג each item
% its own value, .aux = {א}{...}{Item.1} — versus the explicit-\item doctrine
% where BIDI-50/\hitem exist purely to patch the counter that never stepped.
\def\@hebalph#1{\ifcase#1\or א\or ב\or ג\or ד\or ה\or ו\or ז\or ח\or ט\or י\or
  יא\or יב\or יג\or יד\or טו\or טז\or יז\or יח\or יט\or כ\else\@ctrerr\fi}
\DeclareRobustCommand*\hebalph[1]{\expandafter\@hebalph\csname c@#1\endcsname}
\newcommand{\hebenumlabels}{%
  \renewcommand{\theenumi}{\@hebalph\c@enumi}%
  \renewcommand{\labelenumi}{(\textbf{\hebalph{enumi}})}}
\makeatother
% v3.63: \mbox the LABEL. \thetbl ("7.1") and a caption description that STARTS
% WITH A DIGIT ("32 הרגיסטרים…") are two European-Number runs separated only by
% neutrals (the label's own "." and \enspace). Under the RTL paragraph base the UBA
% resolves those neutrals to the RTL level, so both numbers sit in ONE reversed span
% and EXCHANGE VISUAL POSITIONS. Measured on the shipped PDF: a table whose real
% number is 7.1 printed "טבלה 32" — a different, EXISTING table number — while \ref
% correctly printed 7.1, so the caption silently contradicted its own cross-reference,
% and the description lost its leading count. The source contains no adjacency at all;
% it is manufactured by this macro, which is why every source linter was blind.
% \mbox makes the label one atomic box so nothing can reorder across it. Verified
% zero-regression: an ordinary Hebrew caption keeps identical x-geometry.
% One paragraph style covers label, number, separator and description.
% Explicit project/template overrides remain possible; the label hook is kept
% for compatibility, with an identity default. Counters/anchors stay unchanged.
\newcommand{\lakitcaptionstyle}{\normalfont\small\color{black}\hypersetup{allcolors=black}}
\newcommand{\lakitcaptionlabelstyle}[1]{#1}
\newcommand{\lakitcaptionpunctuation}{:}
% Native floats, captionof and longtable use the same presentation policy.
% subcaption is optional and inherits the same font rather than a second size.
\usepackage{caption}
\DeclareCaptionFont{lakit}{\lakitcaptionstyle}
\DeclareCaptionLabelFormat{lakit}{\lakitcaptionlabelstyle{#1~\en{#2}}}
\DeclareCaptionLabelSeparator{lakit}{\lakitcaptionpunctuation\enspace}
\captionsetup{font=lakit,labelfont={},textfont={},labelformat=lakit,
  labelsep=lakit,justification=centering}
\makeatletter
\AtBeginDocument{\@ifpackageloaded{subcaption}{%
  \captionsetup[sub]{font=lakit,labelfont={},textfont={},labelsep=lakit}}{}}
\makeatother
% v3.64: OPT-IN list-of-figures/tables entry via an optional short caption.
% \listoffigures and \listoftables are inherited from report.cls and LOOK available,
% but \figcap/\tabcap step the custom \c@fig/\c@tbl counters and never call
% \addcontentsline — so both printed their heading and NOTHING ELSE. Measured: a doc
% with one \figcap and one \tabcap rendered page 1 as 'רשימת האיורים\n1\n', and the
% .lof held only \@chapter's \addvspace, zero \contentsline entries. An empty list
% under a correct Hebrew heading reads as "this document has no figures".
%
% WHY OPTIONAL, AND NOT AUTOMATIC. Writing #1 unconditionally was implemented first
% and REJECTED: a caption containing \codeline compiles CLEANLY today (verified on
% the pre-v3.64 preamble: 0 errors, PDF produced) but becomes FATAL once it is also
% written to the .lof — \codeline is \hbox to\linewidth{...}, and a full-width box in
% a list entry dies with "\pdfextension endlink cannot be used in vertical mode", no
% output PDF. Gaining a feature nobody uses (0 of 13 collections call \listoffigures)
% at the price of breaking a construct that works is a bad trade. Stock \caption has
% exactly this limitation and exactly this escape hatch, so the convention is familiar.
% \en{}, \code{}, \h{} and inline math in the short text are all safe — verified, they
% survive the write and typeset correctly in the list (RTL geometry confirmed: number
% rightmost at x=509.9, title left of it, dots and folio at the far left).
% With no optional argument the macro is byte-identical to v3.63 — zero risk to every
% existing document, and \listoffigures lists exactly the captions that opted in.
% BLD-37, TABLE side (the \figcap v3.74 idea, adapted): clicking a \ref to a
% table used to land MID-table — hyperref's automatic anchor inside
% \refstepcounter was emitted on the still-open horizontal line whose TAIL is
% the tabular (an hbox whose baseline is the table's vertical CENTER), so the
% jump scrolled the top half of the table above the viewport. The v3.74 entry
% deferred this side claiming "tablebox splits — no single box to measure";
% that deferral dissolves: the measurement target is the TABULAR itself at
% \tabcap ENTRY, while it is still the tail of the open list — the same
% situation \lakin@grabfig handles at an island's end (and the fleet has 0
% split tableboxes). So \lakin@grabtab runs as the FIRST token, before any
% typesetting can push the box away, and the anchor is re-emitted raised above
% the measured tabular. \lakintabraise is a SEPARATE register from
% \lakinfigraise (defined with it, below the LTR-island block): a captionless
% island intentionally leaves the fig register set, and a table must never
% inherit that stale figure height.
%
% WHY the anchor stays AT THE ENTRY SPOT and is NOT deferred into the caption
% paragraph the way \figcap defers it: at \par TeX DROPS a final glue item,
% and the automatic anchor's smash box sitting BEHIND the interword glue kept
% that glue alive — so every centered tabular has always sat half a space off
% true center. Deferring the anchor removed that box and was measured as 982
% differing glyph records: EVERY tabular line shifted 1.30pt (2.61pt for
% pattern 7's two minipages) while captions stayed put. Emitting the raised
% anchor at the old spot — through hyperref's own \Hy@raisedlink, exactly the
% macro the automatic path uses — keeps the line's node structure, and so the
% A/B glyph diff at 0, BY CONSTRUCTION. \Hy@raisedlink reads the raise from
% \HyperRaiseLinkDefault (default \baselineskip = the old mid-table dest); the
% group-local \edef freezes measured-height+42pt into it, and \smash inside
% \Hy@raisedlink hides any raise from the line's extent. \HyCnt@ProvideTheH-
% Counter keeps the dest NAME byte-identical to the automatic one (it defines
% \theHtbl, which \hyper@makecurrent prefers — without it the first measured
% caption was named tbl.1 where the automatic path said tbl.0.1, measured).
% vmode entry, or anything but a box at the tail (no tabular to measure),
% falls back to the plain \refstepcounter — byte-identical to the old path.
\makeatletter
\newcommand{\tabcap}[2][]{\lakin@grabtab
  % v3.79: arm the box-close correction. Set unconditionally and in BOTH
  % branches below's scope, so the plain-\refstepcounter fallback is corrected
  % too; \lakintabboxopen clears it at every \begin{tablebox}, so a \tabcap
  % outside a tablebox can never leave it armed for a later box.
  \global\lakintabcapseentrue
  \ifdim\lakintabraise>0pt
    \H@refstepcounter{tbl}\HyCnt@ProvideTheHCounter{tbl}\hyper@makecurrent{tbl}%
    \begingroup
      % A frameless table targets its measured top. The frame's historical
      % cushion belongs only to tablebox and is corrected at box close.
      \edef\HyperRaiseLinkDefault{\the\lakintabraise}%
      \iflakintabframed
        \edef\HyperRaiseLinkDefault{\the\dimexpr\lakintabraise+42pt\relax}%
      \fi
      \Hy@raisedlink{\hyper@anchorstart{\@currentHref}\hyper@anchorend}%
    \endgroup
    \global\lakintabraise 0pt\relax
  \else
    \refstepcounter{tbl}%
  \fi
  \lakinrtitle{#2}\par\nopagebreak\addvspace{6pt}%
  \def\lakin@short{#1}\ifx\lakin@short\@empty\else
    \addcontentsline{lot}{table}{\protect\numberline{\thetbl}{\protect\ignorespaces #1}}\fi
  {\centering\lakitcaptionstyle\mbox{\lakitcaptionlabelstyle{\tablename~\en{\thetbl}\lakitcaptionpunctuation}}\enspace #2\par}}
\makeatother

% ============================================================================
% Unified FIGURE caption — the \tabcap twin for TikZ/plot figures (v3.34.0).
% A figure (tikzpic) sits inline in the prose; put \figcap BELOW it for a
% numbered, \ref-able caption "איור N.M: תיאור" in the same style as \tabcap:
%   \begin{tikzpic} ... \end{tikzpic}
%   \figcap{\label{fig:foo}description}     % \label FIRST, inside the argument
% WHY a macro, not a hand-written gray line: a manual
%   {\centering\small\textcolor{cGray}{איור: ...}} with a \label inside has NO
% \refstepcounter, so the \label binds to the AMBIENT \@currentlabel (the last
% equation) — \ref then prints an equation number AS the figure number
% (measured on a real page: "איור 48" where 48 is equation 48). \figcap steps a
% real \c@fig counter first, so \label is correct, numbering is consistent, and
% refs are stable. Class-aware like \thetbl (chapter- or section-prefixed).
% ============================================================================
\makeatletter
\@ifundefined{c@chapter}{%
  \newcounter{fig}[section]
  \renewcommand{\thefig}{\ifnum\value{section}>0 \thesection.\fi\arabic{fig}}%
}{%
  \newcounter{fig}[chapter]
  \renewcommand{\thefig}{\ifnum\value{chapter}>0 \thechapter.\fi\arabic{fig}}%
}
\makeatother
\makeatletter
\newcommand{\figcap}[2][]{%
  \ifdim\lakinfigraise>0pt
    % v3.74: step WITHOUT hyperref's automatic anchor, then emit the anchor
    % raised above the measured figure. \H@refstepcounter is hyperref's saved
    % pre-hook \refstepcounter (sets the counter AND \@currentlabel — \label
    % in the caption arg binds exactly as before); \hyper@makecurrent supplies
    % \@currentHref, so the dest NAME is identical to the old automatic one —
    % \ref/\autoref/bookmarks are unchanged, only the dest'S POSITION moves.
    \H@refstepcounter{fig}\hyper@makecurrent{fig}%
    % The anchor is emitted INSIDE the caption paragraph below, as a zero-size
    % raisebox — NOT here in vertical mode, where a \raisebox starts an empty
    % paragraph and adds a blank line (measured: every glyph after the caption
    % shifted +18.8pt). Inside an existing line a [0pt][0pt] box is inert.
    \edef\lakin@anch{\@currentHref}%
    % \edef FREEZES the measured height into the pending macro: the deferred
    % \raisebox is typeset inside the caption paragraph AFTER the register is
    % reset to the 0pt sentinel below — a plain \def read 0 at use time
    % (measured: LAKINRAISE=0.0, dest moved 14pt instead of above the figure).
    \edef\lakin@pend{\noexpand\raisebox{\the\dimexpr\lakinfigraise+42pt\relax}[0pt][0pt]{%
      \noexpand\hyper@anchorstart{\lakin@anch}\noexpand\hyper@anchorend}}%
    \global\lakinfigraise 0pt\relax
  \else
    \refstepcounter{fig}\let\lakin@pend\relax
  \fi
  \lakinrtitle{#2}\par\nopagebreak\addvspace{4pt}%
  \def\lakin@short{#1}\ifx\lakin@short\@empty\else
    \addcontentsline{lof}{figure}{\protect\numberline{\thefig}{\protect\ignorespaces #1}}\fi
  {\centering\lakin@pend\lakitcaptionstyle\mbox{\lakitcaptionlabelstyle{\figurename~\en{\thefig}\lakitcaptionpunctuation}}\enspace #2\par}}
\makeatother

% ============================================================================
% LTR islands.  EVERY tikzpicture / quantikz MUST live inside one, or the whole
% figure's coordinates and Latin labels flip. (Hebrew node text still uses \h{}.)
% ============================================================================
% v3.74 (BLD-37): the \figcap hyperlink DESTINATION used to land at the caption
% line — BELOW the figure — so clicking a figure reference scrolled the figure
% off-screen (measured fleet-wide: 97/97 fig dests in one collection sat 19-54pt
% under their figure's ink; a 347pt figure was fully hidden after the jump).
% hypcap-style fix in two parts: (1) at the island's END, \lakin@grabfig grabs
% the just-typeset centered line with \lastbox, measures it, and puts it back
% with \nointerlineskip — the interline glue above the line was already
% contributed and is NOT re-added, so the page geometry is unchanged to the
% glyph (verified); (2) \figcap then steps the counter WITHOUT hyperref's
% automatic anchor (\H@refstepcounter = the pre-hyperref \refstepcounter) and
% emits the anchor itself, raised by the measured height — the dest lands ABOVE
% the figure top. A captionless island leaves numbering untouched (nothing is
% stepped at the island; the grab flag is cleared at every island begin, so a
% stale height from a captionless figure cannot leak into a later caption).
\newsavebox{\lakinfigbox}
\newdimen\lakinfigraise  % real \dimen: tex.setdimen cannot target a \newlength SKIP register (measured lua error)
% BLD-37 table side: a SEPARATE register for \tabcap. The two must not share —
% a captionless island INTENTIONALLY leaves \lakinfigraise set (the grab runs
% at every island END; only the NEXT island's begin clears it), so a shared
% register would hand a \tabcap after a captionless figure that stale FIGURE
% height. \lakintabraise is written only by \lakin@grabtab at \tabcap entry
% and consumed/reset inside \tabcap itself, so it can never go stale.
\newdimen\lakintabraise
\makeatletter
% v3.74b: measure WITHOUT touching the list. The first draft grabbed the box
% with \lastbox and put it back — and \lastbox DISCARDS the node's baseline
% shift, so a vertically-anchored quantikz circuit rose 2.57pt and everything
% after it sank 5.32pt (measured on the A/B glyph diff). LuaTeX lets us READ
% the tail node's dimensions without removing anything: geometry is identical
% by construction. \lakinfigraise=0pt is the "no figure measured" sentinel.
\def\lakin@grabfig{%
  \ifhmode
    \directlua{
      local n = tex.nest[tex.nest.ptr].tail
      local GLUE, PEN = node.id('glue'), node.id('penalty')
      while n and (n.id == GLUE or n.id == PEN) do n = n.prev end
      local HL, VL = node.id('hlist'), node.id('vlist')
      if n and (n.id == HL or n.id == VL) then
        tex.setdimen('global', 'lakinfigraise',
                     (n.height or 0) + (n.depth or 0))
      else
        tex.setdimen('global', 'lakinfigraise', 0)
      end}%
  \else
    \global\lakinfigraise 0pt\relax
  \fi}
% BLD-37 table side: \lakin@grabfig's twin, NOT a shared helper — the Lua body
% differs in exactly one place, the skip set. Between \end{tabular} and \tabcap
% an otherlanguage wrapper (patterns 3/5 of example_tables) leaves a \write
% whatsit and a dir node ON TOP of the interword glue (measured node dump:
% "whatsit/write <- dir <- glue(2.62) <- hlist"), so the table reader must
% also step over whatsit/dir markers — all zero-size — to reach the tabular.
% \lakin@grabfig deliberately does NOT get the wider skip set: its glue/penalty
% body is the one the v3.74 A/B proof froze, and widening it would newly
% measure (and so newly raise) anchors in island shapes that today take the
% plain path — an unproven behavior change on the FIG side. Anything that is
% not a box behind those markers still hits the 0pt sentinel and \tabcap falls
% back to the plain \refstepcounter.
% SECOND deliberate difference: this stores the box's HEIGHT, not height+depth.
% \figcap's anchor is emitted in the caption paragraph BELOW the figure, so it
% must climb the full extent; \tabcap's anchor sits ON the tabular's own line,
% whose baseline is the vcentered tabular's MIDPOINT — height alone reaches
% the top edge, and adding the depth was measured to overshoot by half the
% table (dests 80-135pt above their frame, INSIDE the preceding table's box).
\def\lakin@grabtab{%
  \ifhmode
    \directlua{
      local n = tex.nest[tex.nest.ptr].tail
      local GLUE, PEN = node.id('glue'), node.id('penalty')
      local WHAT, DIR = node.id('whatsit'), node.id('dir')
      while n and (n.id == GLUE or n.id == PEN
                   or n.id == WHAT or n.id == DIR) do n = n.prev end
      local HL, VL = node.id('hlist'), node.id('vlist')
      if n and (n.id == HL or n.id == VL) then
        tex.setdimen('global', 'lakintabraise', n.height or 0)
      else
        tex.setdimen('global', 'lakintabraise', 0)
      end}%
  \else
    \global\lakintabraise 0pt\relax
  \fi}
\makeatother
% ============================================================================
% v3.79 (BLD-37 table side, CORRECTION): put the anchor at the TOP OF THE BOX,
% instead of guessing the distance to it.
% ============================================================================
% v3.77 raised the \tabcap anchor by (height of the box at the tail of the open
% horizontal list) + 42pt. That is a MEASUREMENT taken BEFORE line breaking, and
% it is short by every term it cannot see. Fleet-measured on the 12 built
% collections (299 tbl.* dests, ownership established through each table's own
% CAPTION, never through "which frame encloses the dest"): 3 dests landed INSIDE
% their own tablebox, below the frame top —
%   computer-architecture       tbl.14.1  -46.4pt   three stacked tabulars
%   computer-architecture       tbl.30.5  -13.5pt   tabular 2.85pt over \linewidth
%   advanced-electronic-devices tbl.B.1   -59.2pt   tabular 2.56pt over \linewidth
% Two missing terms, both reproduced in a minimal document (probe P2/P3):
%   (1) VERTICAL LIST ALREADY CONTRIBUTED inside the box — a second and third
%       tabular, or a heading above the table. The tail read sees only the LAST
%       box, so 77.7pt of stacked tabulars is invisible to the raise.
%   (2) WHICH LINE THE ANCHOR LANDS ON. When the tabular alone overfills
%       \linewidth (459.6121pt in a tablebox), the interword glue after
%       \end{tabular} becomes the chosen breakpoint and \Hy@raisedlink is
%       emitted on the NEXT line — one (depth + lineskip) lower. That term does
%       not EXIST at \tabcap time: line breaking has not happened yet. No
%       measurement taken there can ever recover it. (Proof: probe P3, the same
%       table 2.84912pt too wide, fails at -13.5pt; probe P4, the identical
%       table with \tabcolsep 2pt so it fits, is correct at +32.0pt.)
% So the raise is not computed at \tabcap at all. It is CORRECTED at
% \end{tablebox}, from inside mdframed's own body vbox (endinnercode runs after
% \par\unskip, i.e. after line breaking, while tex.nest[tex.nest.ptr] is still
% that vbox's vertical list). There the anchor's true offset from the top of the
% box content is not estimated but WALKED, and the \Hy@raisedlink shift is set
% so the dest lands 42pt above it — which is 42pt - innertopmargin(9pt) -
% linewidth(0.9pt) = 32.0pt above the FRAME top, the exact position the 295
% already-correct dests have today. That is why this is a correction and not a
% re-placement: the good sites do not move.
%
% WHY THE GLYPH DIFF IS 0 BY CONSTRUCTION — the only thing written is the
% `shift` field of the hlist that \Hy@raisedlink already built, and the guard
% below refuses to write it unless that box is 0x0x0. A zero-size box carries no
% dimension into its parent's packing, the packing already happened, and shift
% is not an input to it. Nothing else in the list is touched: no node is added,
% removed, resized or reordered. Measured, not assumed: assets/example_tables.tex
% and the whole computer-architecture collection A/B at 0 differing pymupdf
% rawdict glyph records.
%
% WHY NOT the alternatives (all prototyped and measured before this one):
%   * accumulate the vertical list at \tabcap time — fixes term (1), CANNOT fix
%     term (2); 2 of the 3 real failures survive. A fix that leaves the majority
%     of the defect standing is the false all-clear this repo exists to avoid.
%   * \zref-savepos two-pass — correct on all shapes, but it makes the dest
%     POSITION aux-derived (DEFECTS BLD-1/2), a single-pass compile silently
%     ships an unraised anchor, and the bare y-difference has no page guard.
%   * emit the anchor at \begin{tablebox} and let \tabcap adopt it — needs the
%     tbl counter stepped at box OPEN to keep the dest NAME byte-identical; a
%     captionless tablebox then either shifts every later table number or leaves
%     an orphan dest that the next table duplicates (a HARD build gate). This
%     correction touches neither the counter nor the name: \tabcap's naming path
%     is v3.77's, unchanged, so resolve_names() is byte-identical (verified).
%   * suppressing the line break (a penalty before the glue) — that MOVES the
%     tabular: the break currently discards the interword glue, so forbidding it
%     re-centers the two overfull tables. Reflow, rejected.
%
% SCOPE. The correction is armed only by \tabcap ITSELF (\lakintabcapseen),
% cleared at every box open, so a captionless tablebox is a no-op and a \tabcap
% hosted OUTSIDE a tablebox (the fleet has exactly one, computer-architecture
% tbl.14.2 in an exbox) keeps v3.77's measured raise. Both fallbacks are
% therefore never worse than today; and because hyperref's AUTOMATIC anchor uses
% the same \Hy@raisedlink shape, the plain-\refstepcounter path (vmode entry, or
% prose between \end{tabular} and \tabcap — probe P6, -51.5pt) is corrected too.
% KNOWN LIMIT — a tablebox that SPLITS across pages. The offset is measured from
% the top of the UNSPLIT body, so a dest in segment 2 gets raised past its own
% segment. MEASURED on a forced split (probe: 40 filler lines + a table): the
% dest went from y=585.0 on page 2 (inside segment 2, at the caption — the old
% BLD-37 symptom) to y=-113.4, i.e. ABOVE the page top. Viewers clamp that to the
% page top, which happens to show segment 2 whole, but an out-of-page coordinate
% is not something to ship on purpose. Hence the \textheight guard in
% lakintab.fix: a body TALLER THAN THE TEXT BLOCK cannot fit on any page, so it
% must split — the correction is skipped and v3.77's behaviour stands unchanged
% (verified: the probe returns to y=585.0). The guard is deliberately EXACT, not
% a page-fitting heuristic: it is sound (a taller-than-page body always splits)
% but PARTIAL (a shorter body that starts low on a page could still split). That
% residue is unreachable in practice — v3.37's everyline + the caption's
% \nopagebreak glue move such boxes WHOLE instead (probe N7: a 22-row table
% pushed down by 430pt moved whole, 0 segments), and the fleet has 0 split
% tableboxes, re-counted after this change. Reimplementing mdframed's own
% page-fitting arithmetic to close the residue would be a heuristic with its own
% failure mode, which is the trade this repo has repeatedly refused.
\newif\iflakintabcapseen
\newif\iflakintabframed
% The Lua lives in ONE \directlua at load time, and is written with NO
% backslash, %, # or ~ : \directlua expands its argument like \edef, so an
% active ~ would expand to \nobreakspace and a % would eat the rest of the line
% (the second is how the first draft of this block died). "not (a == b)" stands
% in for Lua's inequality operator for the same reason.
\directlua{
  lakintab = { }
  lakintab.HL = node.id("hlist")
  lakintab.VL = node.id("vlist")
  lakintab.RU = node.id("rule")
  lakintab.GL = node.id("glue")
  lakintab.KE = node.id("kern")
  lakintab.WH = node.id("whatsit")
  lakintab.TM = node.id("temp")
  lakintab.WD = node.subtype("pdf_dest")
  function lakintab.d(v)
    if v == nil then return 0 end
    if v < -1000000000 then return 0 end
    return v
  end
  function lakintab.scan(list, y, horiz, parent, target)
    local x = list
    while x do
      local i = x.id
      if i == lakintab.HL or i == lakintab.VL then
        if horiz then
          local b = y + x.shift
          if i == lakintab.HL then
            lakintab.scan(x.head, b, true, x, target)
          else
            lakintab.scan(x.head, b - lakintab.d(x.height), false, x, target)
          end
        else
          if i == lakintab.HL then
            lakintab.scan(x.head, y + lakintab.d(x.height), true, x, target)
          else
            lakintab.scan(x.head, y, false, x, target)
          end
          y = y + lakintab.d(x.height) + lakintab.d(x.depth)
        end
      elseif i == lakintab.WH then
        if x.subtype == lakintab.WD then
          local id = x.dest_id
          if type(id) == "string" and string.sub(id, 1, 4) == "tbl." then
            if horiz and parent and parent.width == 0
               and parent.height == 0 and parent.depth == 0 then
              parent.shift = parent.shift + (target - y)
            end
          end
        end
      elseif not horiz then
        if i == lakintab.RU then
          y = y + lakintab.d(x.height) + lakintab.d(x.depth)
        elseif i == lakintab.GL or i == lakintab.KE then
          y = y + lakintab.d(x.width)
        end
      end
      x = x.next
    end
  end
  function lakintab.fix(target, maxh)
    local nst = tex.nest[tex.nest.ptr]
    if nst == nil then return end
    local h = nst.head
    if h == nil then return end
    if h.id == lakintab.TM then h = h.next end
    local tot = 0
    local x = h
    while x do
      local i = x.id
      if i == lakintab.HL or i == lakintab.VL or i == lakintab.RU then
        tot = tot + lakintab.d(x.height) + lakintab.d(x.depth)
      elseif i == lakintab.GL or i == lakintab.KE then
        tot = tot + lakintab.d(x.width)
      end
      x = x.next
    end
    if tot > maxh then return end
    lakintab.scan(h, 0, false, nil, target)
  end
}
% @-FREE NAMES ON PURPOSE: these two are named in the `tablebox` key list at
% \newmdenv above, which is tokenised OUTSIDE any \makeatletter region — a
% \lakin@... there would tokenise as \lakin followed by the letters @tabopen.
\newcommand{\lakintabboxopen}{\lakintabframedtrue\global\lakintabcapseenfalse}
\newcommand{\lakintabboxclose}{%
  \iflakintabcapseen
    \directlua{lakintab.fix(-\number\dimexpr\lakintabcushion\relax,
                            \number\dimexpr\textheight\relax)}%
    \global\lakintabcapseenfalse
  \fi\lakintabframedfalse}
% The cushion above the box's CONTENT top. 42pt puts the dest 32.0pt above the
% frame top (42pt - innertopmargin 9pt - linewidth 0.9pt), i.e. exactly where
% the 295 correct dests already are — the correction moves only the broken ones.
\newcommand{\lakintabcushion}{42pt}
\newenvironment{qcirc}{\par\medskip\global\lakinfigraise 0pt\relax\begin{otherlanguage}{english}\begin{center}}
  {\csname lakin@grabfig\endcsname\end{center}\end{otherlanguage}\par\medskip}
\newenvironment{tikzpic}{\par\medskip\global\lakinfigraise 0pt\relax\begin{otherlanguage}{english}\begin{center}}
  {\csname lakin@grabfig\endcsname\end{center}\end{otherlanguage}\par\medskip}

% ============================================================================
% Code listings.  ★★ v3.4 THE BOX-OVERFLOW FIX:
%   resetmargins=true   → the listing IGNORES the enclosing environment's
%                         indentation and measures from the PAGE. v3.4 stated
%                         this backwards ("measured from the ENCLOSING box"),
%                         and v3.37 then blanked in-box line numbers because
%                         the page-anchored listing put its hanging digits on
%                         top of the box rule. Measured v4.5.0: with the key
%                         FALSE the listing follows the box and the code moves
%                         from x=75.26 to x=97.74, clearing the rule at 72.07.
%                         The key now lives in one \lstset below instead of in
%                         each style, so a box can override it — codebox does.
%   framexleftmargin=0  → the frame (left bar) does not extend left of the code.
%   xleftmargin keeps a small indent so the line numbers have room INSIDE.
% asmblock wraps the listing in an LTR island (code is LTR).
% ============================================================================
\lstdefinestyle{asm}{
  basicstyle=\ttfamily\small,
  backgroundcolor=\color{cBlueBg!55},
  frame=leftline,framerule=2.5pt,rulecolor=\color{cBlue},
  numbers=left,numberstyle=\tiny\color{cGray},numbersep=8pt,
  xleftmargin=2.2em,framexleftmargin=0pt,framexrightmargin=0pt,
  breaklines=true,columns=fullflexible,keepspaces=true,
  showstringspaces=false,aboveskip=8pt,belowskip=8pt,
  % v3.36: without a comment DELIMITER the commentstyle below never fires —
  % listings has no `language` here, so it does not know '#' opens a comment,
  % and every asm listing rendered flat black (measured: a whole-page pixel
  % sweep found ZERO cGreen glyphs across 135 shipped listings). morecomment
  % gives it the delimiter; color-only change, glyph widths untouched.
  morecomment={[l]{\#}},
  commentstyle=\color{cGreen},
}
\lstdefinestyle{cstyle}{
  language=C,basicstyle=\ttfamily\small,
  backgroundcolor=\color{cGreenBg!55},
  frame=leftline,framerule=2.5pt,rulecolor=\color{cGreen},
  numbers=left,numberstyle=\tiny\color{cGray},numbersep=8pt,
  xleftmargin=2.2em,framexleftmargin=0pt,framexrightmargin=0pt,
  breaklines=true,columns=fullflexible,keepspaces=true,
  showstringspaces=false,keywordstyle=\color{cBlue}\bfseries,
  commentstyle=\color{cGray}\itshape,aboveskip=8pt,belowskip=8pt,
}
% v3.36: pystyle/htmlstyle — colored twins of cstyle for the Python/HTML
% listings that used to repeat a ~200-char options blob with NO color
% (keywords rendered bold-black, comments italic-black — measured; C alone
% was colored). Geometry is copied VERBATIM from the shipped blob (frame,
% margins, skips) so swapping a call site to style=pystyle moves NOTHING —
% only keywordstyle/commentstyle/stringstyle are added (bold stays bold,
% italic stays italic, mono metrics identical => zero reflow, proven by
% per-page word-bbox diff). breaklines/breakatwhitespace stay per-site.
% v3.78: numbers=left joins both styles, with the exact numberstyle/numbersep
% keys asm/cstyle carry — fleet-wide consistency (author decision: all four
% styles numbered; "asm/cstyle numbered, py/html not" was an accident of the
% v3.36 geometry freeze, not a design). The numbers hang LEFT of the code box
% (listings sets them with \llap, so no overfull hbox can arise and the code
% x is untouched — measured: code frozen, digits at code_x-11.8pt); inside a
% colored box they are blanked like everyone else's (v3.37 hook), and the
% clean variant restores them in its frameless boxes (v3.78, above).
\lstdefinestyle{pystyle}{
  language=Python,basicstyle=\ttfamily\small,
  columns=fullflexible,keepspaces=true,
  frame=single,rulecolor=\color{cTblLine},
  numbers=left,numberstyle=\tiny\color{cGray},numbersep=8pt,
  xleftmargin=0.5em,framexleftmargin=0pt,framexrightmargin=0pt,
  showstringspaces=false,aboveskip=8pt,belowskip=8pt,
  keywordstyle=\color{cBlue}\bfseries,
  commentstyle=\color{cGray}\itshape,
  stringstyle=\color{cGreen},
}
% v4.5.0: hoisted out of the four styles so an enclosing box can override it.
% Every listing still resolves to true unless something later says otherwise,
% so shipped output is unchanged (verified byte-for-byte on all eight PDFs).
\lstset{resetmargins=true}
\lstdefinestyle{htmlstyle}{
  language=HTML,basicstyle=\ttfamily\small,
  columns=fullflexible,keepspaces=true,
  frame=single,rulecolor=\color{cTblLine},
  numbers=left,numberstyle=\tiny\color{cGray},numbersep=8pt,
  xleftmargin=0.5em,framexleftmargin=0pt,framexrightmargin=0pt,
  showstringspaces=false,aboveskip=8pt,belowskip=8pt,
  keywordstyle=\color{cBlue}\bfseries,
  commentstyle=\color{cGray}\itshape,
  stringstyle=\color{cGreen},
}
\newenvironment{asmblock}{\par\begin{otherlanguage}{english}}{\end{otherlanguage}\par}

% <hllpdf:listing-smart-quotes>
% v4.21.0: four Unicode punctuation characters must enter listings' token
% buffer. Unregistered characters bypass it: a closing quote was printed
% BEFORE the pending word, and an opening quote BEFORE a pending brace.
% Extend the standard character table, then activate only these codepoints
% in the listing's local character-table scope. Do this BEFORE the package's
% delimiter/alsoother/literate hooks so explicit user mappings still win.
% This is not universal Unicode support, nor a conversion to ASCII quotes.
% See listings developer guide 15.3.1 and references/listing-smart-quotes.md.
\makeatletter
\lst@CCPut \lst@ProcessOther {"2018}{"2019}{"201C}{"201D}\z@
\lst@AddToHookAtTop{SelectCharTable}{%
  \lst@ifactivechars
    \catcode"2018=\active \catcode"2019=\active
    \catcode"201C=\active \catcode"201D=\active
  \fi}
\makeatother
% </hllpdf:listing-smart-quotes>

% <hllpdf:listing-technical-symbols>
% v4.30: bounded extension for common technical punctuation in LTR listings.
% Without character-table registration these glyphs bypass the token buffer
% and can move before the preceding word even inside quoted strings.
% Keep the original four-quote block separate for its standing mutation test.
% This is source-faithful printing, not syntax validation or general Unicode.
\makeatletter
\lst@CCPut \lst@ProcessOther
  {"2013}{"2014}{"2026}{"2190}{"2192}{"2194}{"21D2}
  {"2212}{"221E}{"2260}{"2264}{"2265}\z@
\lst@AddToHookAtTop{SelectCharTable}{%
  \lst@ifactivechars
    \catcode"2013=\active \catcode"2014=\active \catcode"2026=\active
    \catcode"2190=\active \catcode"2192=\active \catcode"2194=\active
    \catcode"21D2=\active \catcode"2212=\active \catcode"221E=\active
    \catcode"2260=\active \catcode"2264=\active \catcode"2265=\active
  \fi}
\makeatother
% </hllpdf:listing-technical-symbols>


\usepackage[hidelinks,unicode]{hyperref}   % LAST, after colors. unicode = Hebrew bookmarks.
% v3.62: THE ISOLATION MACROS CORRUPTED EVERY PDF BOOKMARK THEY APPEARED IN.
% \en is defined (line ~155) as
%     \foreignlanguage{english}{\language=\csname l@nohyphenation\endcsname #1}
% and in a heading that body is written verbatim to the .toc and handed to
% hyperref's \pdfstringdef. hyperref strips the two CONTROL SEQUENCES (\language,
% \l@nohyphenation) but the plain `=` CHARACTER token between them is not a control
% sequence, so it survives into the PDF string: \chapter{שער \en{CNOT} והצמדה}
% produced the bookmark "שער =CNOT והצמדה". The printed TOC and the heading on the
% page are CORRECT — only the navigation pane and the viewer window title are wrong,
% which is why it went unnoticed. Measured on the shipped fleet: 70 corrupted
% bookmarks across 5 PDFs, from 177 unprotected \en-in-heading sites. Same corruption
% reached /Title and /Subject via the \hypersetup{pdftitle=...} the docs recommend.
% This hook tells hyperref to use the plain argument for all three isolation macros
% when building a PDF string. It does NOT touch the .toc (the printed TOC is
% byte-unchanged) and it does NOT mask the math signal — a heading with $...$ still
% logs "removing math shift", so BIDI-15 detection is intact.
% BONUS, measured: it also makes \code{} in an UNSTARRED heading COMPILE. BLD-10
% recorded that as needing \texorpdfstring or a starred heading; the real root cause
% is hyperref's \pdfstringdef scan of \code's \ifmmode...\else...\fi body — which is
% exactly why the identically-bodied \texttt survived (hyperref pre-neutralizes
% \texttt in PDF strings) and why \DeclareRobustCommand could never help.
% v3.64: \def\\{ } added. A manual line break in a heading WELDS the two words
% together in the bookmark: \chapter{שורה ראשונה\\שורה שנייה} produced the outline
% entry 'שורה ראשונהשורה שנייה' — "ראשונהשורה" is not a word. hyperref drops the
% \\ and puts nothing in its place, so the break that looks like a space on the
% page is a zero-width join in the navigation pane. The spaced form
% 'מילה אחת \\ מילה שנייה' escapes with a harmless DOUBLE space ('אחת  מילה'),
% which is why this never got noticed. Both are measured on a real outline.
% 0 live instances in the corpus — but SKILL.md and content-style.md both RECOMMEND
% a manual heading break, so this is a primed trap rather than a dormant one, and a
% one-token preamble fix retires it for good instead of asking authors to remember.
% Mapping \\ to a space fixes the WELD, which is the actual defect (a nonexistent
% word in the navigation pane). It does not normalise whitespace: measured after the
% fix, the welded form becomes 'שורה ראשונה שורה שנייה' (correct) and the already-
% spaced form goes from two spaces to three. That is cosmetic — viewers render it as
% a wider gap, and no word is corrupted. A per-heading \texorpdfstring is still the
% way to get an exactly-tuned bookmark string.
\pdfstringdefDisableCommands{\def\en#1{#1}\def\h#1{#1}\def\doclang#1{#1}\def\code#1{#1}\def\codeline#1{#1}\def\\{ }}
% v3.70: RESTORE \url LINE-BREAKING. url.sty freezes its break penalties at LOAD
% time (\mathchardef\UrlBreakPenalty=\binoppenalty), and v3.45's global
% \binoppenalty=10000 (inline-math atomicity) therefore priced EVERY breakpoint
% inside \Url at 10000 = never break. Measured: a long \url in Hebrew prose is
% one unbroken line, ~100pt overfull, invading the margin band the ink gates
% deliberately tolerate. With the two lines below: 0 overfull, the URL breaks
% after '/', the RTL split is CORRECT (fragment order forward), and each line
% gets its own link annotation exactly covering its fragment (measured).
% Confined to \Url's group — inline-math atomicity (TYP-7) is untouched.
\mathchardef\UrlBreakPenalty=200
\mathchardef\UrlBigBreakPenalty=100
\ifx\lakitmainlanguage\hllpdfenglish
  \hypersetup{pdflang=en}
\else
  \hypersetup{pdflang=he}
\fi  % /Lang is the document base locale (not per-span accessibility tagging).
% Set a per-document \hypersetup{pdftitle={...}} too — see SKILL.md.

% v3.52 (round-2 audit): two RTL bibliography/cross-ref defects, both invisible to
% every gate and both affecting the FIRST author to use the feature.
% (a) \@biblabel's default [#1] brackets are bidi-neutral and mirror to ]1[ under
% the RTL bibliography base — every \thebibliography reference list shows ]1[.
% LTR-isolating the label restores [1]. (The render gate's mirror rule fingerprints
% round parens )1(, so square brackets slipped through.)
\makeatletter
\renewcommand*{\@biblabel}[1]{\foreignlanguage{english}{[#1]}}
% v3.54: the SAME neutral-bracket trap on the INLINE side — a multi-key \cite{a,b}
% under the RTL base lays its European-number list right-to-left, so [1, 2] renders
% "[2,1]" and [1, 2, 3] renders "[3,2,1]" (verified by glyph geometry — the reader
% sees the sources in REVERSE order; a note form \cite[p.5]{a} also reorders to
% "[p. 5,1]"). LTR-isolate the whole bracket group (number list AND optional note),
% exactly like \@biblabel above: [1, 2] / [1, 2, 3] / [1, p. 5] all render correct.
% Placed AFTER hyperref so it wins; hyperref hyperlinks each number inside #1 (via
% \@citex), which this wrapper leaves intact. Mirrors the LaTeX kernel \@cite shape.
% v4.9.1: directional isolation must NOT switch the language here. Babel's
% safe-active citation wrapper passes its key list through \bbl@tempa;
% \foreignlanguage reuses that scratch macro before the deferred citation
% loop expands it. The old line therefore cited the literal key "english".
% If such a key existed, even the undefined-reference gate stayed green.
% \babelsublr isolates the bracket group without selecting a locale or
% changing that key list. Preserve native links, key order and optional notes.
\def\@cite#1#2{\babelsublr{[{#1\if@tempswa , #2\fi}]}}
\makeatother

% v4.11.0: a reference inside align's \text inherited HEBdisptag=true,
% so its enclosing hbox mirrored the brackets despite a clean compile.
% Isolate the complete printed reference without selecting a Babel language.
% Also neutralize the display-only normalfont direction reset locally. A
% direction-only patch inverted 12 to 21 inside equation in the measured probe.
% Plain maketag is local to the reference: mathtools' display-side wrapper
% otherwise globally consumes manual_tag while formatting a nested reference.
% Preserve each original command (including the selective auxiliary write),
% current tag forms, native links, and showonlyrefs on/off transitions.
\makeatletter
\NewCommandCopy{\HEB@nativeeqref}{\eqref}
\newcommand{\HEB@eqrefisolate}[1]{%
  \babelsublr{\HEBdisptagfalse\let\bbl@eqnodir\relax\let\maketag@@@\ORIGmaketag@@@#1}}
\protected\def\eqref#1{\HEB@eqrefisolate{\HEB@nativeeqref{#1}}}
\ExplSyntaxOn
\NewCommandCopy{\HEB@nativeMTeqref}{\MT_extended_eqref:n}
\cs_set_protected:Npn \MT_extended_eqref:n #1
  {\HEB@eqrefisolate{\HEB@nativeMTeqref{#1}}}
\ExplSyntaxOff
\makeatother

% v4.11.0: custom caption counters are not part of hyperref's English
% locale, and Hebrew has no upstream hyperref locale in this toolchain.
% Restore the appropriate names at every Babel language switch, including
% return from a nested inline island. Do not localize the destination keys.
\makeatletter
\newcommand{\HEB@refnames}{%
\def\chapterautorefname{פרק}%
\def\sectionautorefname{סעיף}%
\def\subsectionautorefname{תת־סעיף}%
\def\subsubsectionautorefname{תת־תת־סעיף}%
\def\paragraphautorefname{פסקה}%
% v4.31: hyperref uses a distinct counter name at the deepest standard level.
\def\subparagraphautorefname{תת־פסקה}%
\def\equationautorefname{משוואה}%
\def\figureautorefname{איור}%
\def\tableautorefname{טבלה}%
\def\partautorefname{חלק}%
\def\appendixautorefname{נספח}%
\def\footnoteautorefname{הערת שוליים}%
\def\itemautorefname{סעיף}%
\def\Itemautorefname{סעיף}%
\def\figautorefname{איור}%
\def\tblautorefname{טבלה}%
\def\pageautorefname{עמוד}% 
}
\addto\extrashebrew{\HEB@refnames}
\newcommand{\ENG@refnames}{%
  \def\subparagraphautorefname{subparagraph}%
  \def\chapterautorefname{chapter}%
  \def\sectionautorefname{section}%
  \def\subsectionautorefname{subsection}%
  \def\subsubsectionautorefname{subsubsection}%
  \def\paragraphautorefname{paragraph}%
  \def\equationautorefname{Equation}%
  \def\figureautorefname{Figure}%
  \def\tableautorefname{Table}%
  \def\partautorefname{Part}%
  \def\appendixautorefname{Appendix}%
  \def\footnoteautorefname{footnote}%
  \def\itemautorefname{item}%
  \def\figautorefname{Figure}%
  \def\tblautorefname{Table}%
  \def\Itemautorefname{item}%
  \def\pageautorefname{page}%
}
\addto\extrasenglish{\ENG@refnames}
\ifx\lakitmainlanguage\hllpdfenglish\ENG@refnames\else\HEB@refnames\fi
\makeatother

% ============================================================================
% Optional running head (v3.39): \lakithead{<book title>} — OFF by default.
% A document opts in with ONE line after \input{preamble.tex}:
%   \lakithead{מבנה המחשב}
% Every content page then carries the series' orientation line: chapter name
% at the reading edge, book title opposite, thin rule. (fancyhdr's [L]/[R]
% are PHYSICAL; under the RTL page direction [L] lands at the RIGHT/reading
% edge — measured, consistent across the corpus.) Packaged like the topic
% rail: always loaded, inert without the call. This centralizes the 8-line
% block that 8 lakits used to copy by hand — the same duplication vector
% that once let a lakit preamble go ~100 lines stale. Class-aware: report
% fills the mark from \chapter, article from \section. NB: a titlepage env
% self-empties; front matter that sets \pagestyle{empty} must restore with
% \pagestyle{fancy} afterwards (NOT plain — that discards the head).
% ============================================================================
\usepackage{fancyhdr}
% v3.61: \chapter* left a STALE running head — pages of an UNNUMBERED chapter
% carried the PREVIOUS chapter's title. report.cls routes \chapter* through
% \@schapter, which never calls \chaptermark, so \leftmark keeps whatever the last
% numbered chapter left. Measured: with \chapter{פרק ממוספר} then
% \chapter*{מבוא לא ממוספר}, page 4 (inside the starred chapter) printed
% "פרק ממוספר". Invisible on the chapter's FIRST page (\thispagestyle{plain}), so it
% only shows from the second page on — the pages an author is least likely to check.
% An unnumbered front-matter chapter (\chapter*{מבוא}, \chapter*{הקדמה}) is everyday
% Hebrew-academic structure, and the corpus already carries HAND-WRITTEN \markboth
% lines right after \chapter* — the exact per-document duplication \lakithead exists
% to retire. Patched INSIDE \lakithead so the macro stays inert when unused; a
% document's own later \markboth still wins, so those hand-patches do not regress.
\makeatletter
% <hllpdf:starred-heading-scope>
\newif\iflakin@starmarkpatched
\newcommand{\lakistarmark}{%
  \iflakin@starmarkpatched\else
  % Dispatch on the actual class counter, not hyperref's shared internals.
  \ifdefined\c@chapter
    \let\HEBorigschapter\@schapter
    % Record only after the native heading has reached its page. A mark
    % before the native routine can advertise the next heading too early.
    \def\@schapter##1{\HEBorigschapter{##1}\chaptermark{##1}}%
  \else
    % @ssect receives layout arguments, NOT the section name or level.
    % Capture the name in @startsection without introducing a TeX group:
    % grouping sectioning would break the native after-heading state.
    \def\HEB@sectionkind{}%
    \def\HEB@mainsection{section}%
    \let\HEBorigstartsection\@startsection
    \def\@startsection##1##2##3##4##5##6{%
      \def\HEB@sectionkind{##1}%
      \HEBorigstartsection{##1}{##2}{##3}{##4}{##5}{##6}}%
    \let\HEBorigssect\@ssect
    \def\@ssect##1##2##3##4##5{%
      \HEBorigssect{##1}{##2}{##3}{##4}{##5}%
      % Lower starred headings must not replace the article's section mark.
      \ifx\HEB@sectionkind\HEB@mainsection
        \sectionmark{##5}%
      \fi
      \ignorespaces}%
  \fi
  \global\lakin@starmarkpatchedtrue
  \fi}
% </hllpdf:starred-heading-scope>
% <hllpdf:running-language>
% Choose the wrapper when the mark is written, not when a later page ships.
% English punctuation must remain LTR even after returning to Hebrew prose.
\newcommand{\HEB@runningmark}[1]{%
  \iflanguage{english}{\markboth{\en{#1}}{}}{\markboth{\h{#1}}{}}}
% </hllpdf:running-language>
\makeatother
% v4.2.1 (NAV): every collection already names itself in \lakithead — feed the
% PDF /Title from the same declaration, so the navigation gate's
% PDF_REQUIRE_TITLE has a universal source and no fleet member ships
% title-less by missing a second declaration site. Measured before this: 14
% of 16 collections shipped /Title-empty while every one carried a perfectly
% good \lakithead string (plain Hebrew in all 16 — no math, no formatting, so
% the PDF-string path is clean). An explicit title always wins: set it with
% \lakitpdftitle{...} anywhere before \lakithead, or with a plain
% \hypersetup{pdftitle=...} — the fallback detects that via hyperref's
% \@pdftitle, probe-verified \@empty-equal until the first explicit set.
\makeatletter
\newif\iflakit@pdftitleset
\newcommand{\lakitpdftitle}[1]{%
  \hypersetup{pdftitle={#1}}\global\lakit@pdftitlesettrue}
\newcommand{\lakithead}[1]{%
  \iflakit@pdftitleset\else
    \ifx\@pdftitle\@empty
      \hypersetup{pdftitle={#1}}\global\lakit@pdftitlesettrue
    \fi
  \fi
  \pagestyle{fancy}\fancyhf{}%
  \lakistarmark

  \fancyhead[L]{\small\textcolor{cGray}{\leftmark}}%
  % <hllpdf:fixed-head-language>
  % The fixed document title has its OWN language snapshot, independent of
  % the current section/chapter mark. Choose the branch NOW, not at shipout:
  % the page may be Hebrew by the time an English title is finally printed.
  \iflanguage{english}{%
    \fancyhead[R]{\small\textcolor{cGray}{\en{#1}}}%
  }{%
    \fancyhead[R]{\small\textcolor{cGray}{\h{#1}}}%
  }%
  % </hllpdf:fixed-head-language>
  \fancyfoot[C]{\thepage}%
  \renewcommand{\headrulewidth}{0.4pt}%
  \ifdefined\chaptermark
    \renewcommand{\chaptermark}[1]{\HEB@runningmark{##1}}%
  \else
    \renewcommand{\sectionmark}[1]{\HEB@runningmark{##1}}%
  \fi}
\makeatother

% v3.39.1: path-independent byte-reproducibility. LuaTeX derives the trailer
% /ID from the OUTPUT PATH, so two clones of the same source at different
% directories produced different sha256 (measured: the ONLY delta was the
% 62-byte /ID line). Bit 512 suppresses exactly that line — measured: Creator/
% Producer/CreationDate/ModDate and all content untouched. /ID is optional
% per ISO 32000 (required only for ENCRYPTED PDFs and PDF/A — neither is a
% target here). If a future workflow needs an ID (PDF/A, Acrobat revision
% matching), delete this line and it returns.
\pdfvariable suppressoptionalinfo \numexpr 512\relax

% v3.75: BOLD MATH IN BOLD HEADINGS. Heading text is \bfseries but math inside
% it stayed regular-weight CM — a hairline $V_T$ inside a heavy 24.8pt chapter
% title (measured: CMMI12 fl=6 next to FrankRuhlLibre-Bold fl=20; 155 live
% heading-with-math sites fleet-wide). \boldmath rides along wherever the class
% sets \bfseries for a heading; A/B measured: $V_T$ -> CMMIB10 fl=22, 0 missing
% characters (the bold CM faces ship with this TeX Live). Guarded per-macro so
% the article class (no \@makechapterhead) and future class changes stay safe.
% v4.0.1 (adoption fix): the SEARCH patterns above carry no parameter token.
% v4.0.0 made them specific by including the class's own #1/#2, and measured on
% TeX Live 2023 that never matches: a \patchcmd search pattern containing a
% parameter token cannot match while the call sits inside \@ifundefined{...}{}{...}
% (probed both #1 and ##1 -> FAIL; the same pattern at top level succeeds). Every
% one of these patches therefore failed, and because v4.0.0 also replaced the
% empty failure branch with \PackageError, the failure became FATAL: all four
% bundled example assets stopped building ("Cannot patch chapter-title math
% weight -> Fatal error, no output PDF"). The \PackageError branches are KEPT —
% a silent no-op here is what let v3.75's fix rot unnoticed — only the parameter
% tokens are dropped, so the patches match and a real class change still errors.
\makeatletter
\@ifundefined{@makechapterhead}{}{%
  \patchcmd{\@makechapterhead}{\huge\bfseries \@chapapp}
    {\huge\bfseries\boldmath \@chapapp}{}
    {\PackageError{hebrew-lualatex-pdf}{Cannot patch chapter-number math weight}{Unsupported document class}}%
  \patchcmd{\@makechapterhead}{\Huge \bfseries}
    {\Huge \bfseries\boldmath}{}
    {\PackageError{hebrew-lualatex-pdf}{Cannot patch chapter-title math weight}{Unsupported document class}}%
  \patchcmd{\@makeschapterhead}{\Huge \bfseries}
    {\Huge \bfseries\boldmath}{}
    {\PackageError{hebrew-lualatex-pdf}{Cannot patch starred chapter math weight}{Unsupported document class}}}
% v4.0.1: \part sits under a THREE-DEEP wrapper chain, so the class-based anchor
% is not in \@part at all. Measured with \show: hyperref rewrites \@part to
% `...\phantomsection\fi \H@old@part`, and nameref rewrites THAT to
% `\NR@gettitle{#1}\NR@part[{#1}]{#2}` — the class body ends up in \NR@part.
% v4.0.0 patched \@part and made failure fatal, which is why all seven bundled
% examples stopped building. \part is NOT a dead heading level here: 188 live uses
% across the fleet, so the bold-math fix is worth getting right rather than
% dropping. We walk the chain innermost-first and stop at the first hit. Failure is
% a WARNING, not an error: the chain is a property of hyperref/nameref versions, so
% a future package release must degrade to regular-weight part math — a cosmetic
% loss — instead of refusing to build. Chapter and section keep \PackageError,
% because those targets are in the CLASS and were verified to patch here.
\newif\iflakin@partnum \newif\iflakin@parttitle \newif\iflakin@sparttitle
\def\lakin@partnum#1#2{\iflakin@partnum\else
  \patchcmd{#1}{#2\bfseries \partname}{#2\bfseries\boldmath \partname}
    {\lakin@partnumtrue}{}\fi}
\def\lakin@parttitle#1#2{\iflakin@parttitle\else
  \patchcmd{#1}{#2 \bfseries}{#2 \bfseries\boldmath}{\lakin@parttitletrue}{}\fi}
\def\lakin@sparttitle#1#2{\iflakin@sparttitle\else
  \patchcmd{#1}{#2 \bfseries}{#2 \bfseries\boldmath}{\lakin@sparttitletrue}{}\fi}
\@ifundefined{c@chapter}{\def\lakin@pnum{\Large}\def\lakin@ptit{\huge}}%
                        {\def\lakin@pnum{\huge}\def\lakin@ptit{\Huge}}
\@ifundefined{NR@part}{}{\expandafter\lakin@partnum\expandafter\NR@part\expandafter{\lakin@pnum}%
                         \expandafter\lakin@parttitle\expandafter\NR@part\expandafter{\lakin@ptit}}
\@ifundefined{H@old@part}{}{\expandafter\lakin@partnum\expandafter\H@old@part\expandafter{\lakin@pnum}%
                            \expandafter\lakin@parttitle\expandafter\H@old@part\expandafter{\lakin@ptit}}
\@ifundefined{@part}{}{\expandafter\lakin@partnum\expandafter\@part\expandafter{\lakin@pnum}%
                       \expandafter\lakin@parttitle\expandafter\@part\expandafter{\lakin@ptit}}
\@ifundefined{NR@spart}{}{\expandafter\lakin@sparttitle\expandafter\NR@spart\expandafter{\lakin@ptit}}
\@ifundefined{H@old@spart}{}{\expandafter\lakin@sparttitle\expandafter\H@old@spart\expandafter{\lakin@ptit}}
\@ifundefined{@spart}{}{\expandafter\lakin@sparttitle\expandafter\@spart\expandafter{\lakin@ptit}}
\iflakin@partnum\else\PackageWarning{hebrew-lualatex-pdf}{Part-number math stays
  regular weight: no known \noexpand\part wrapper matched (hyperref/nameref
  version change). Cosmetic only}\fi
\iflakin@parttitle\else\PackageWarning{hebrew-lualatex-pdf}{Part-title math stays
  regular weight: no known \noexpand\part wrapper matched. Cosmetic only}\fi
% v4.12.0: THE PART NUMBER IN THE TOC WAS NOT DIRECTIONALLY SEPARATED FROM THE
% TITLE. report.cls/article.cls compose the part entry as ONE text stream --
% \addcontentsline{toc}{part}{\thepart\hspace{1em}#1} -- while EVERY other level
% goes through \numberline, which puts the number in an \hbox. A box is a bidi
% OBJECT and cannot join a neighbouring run; a bare "IV" can. So in the RTL TOC
% paragraph a Latin number and a Latin-INITIAL title merge into a SINGLE LTR run
% and are laid out left-to-right as a unit: the number lands on the wrong side of
% the row, and a mixed title is scrambled outright. Measured on \part{MCP ...}:
% the printed order was "<tail> V MCP" -- title fragment, number, title fragment.
% The chapter/section rows one line below, with the SAME Latin-initial titles,
% were correct on the same page. That contrast is the whole diagnosis.
% \en{} IS NOT A WORKAROUND, measured both ways on one fixture: it fixes a
% PURE-Latin title (\part{\en{Hooks}} renders correctly) and does NOT fix a MIXED
% one (\part{\en{MCP} <hebrew>} stays scrambled). babel's language switch changes
% font and hyphenation, not the bidi run structure -- so the defect is not
% reachable from the source side at all, which is why it belongs here.
% Boxing the number is the entire fix and is a STRICT NO-OP where the bug never
% bit: a Hebrew-titled document compiles BYTE-IDENTICAL before and after (md5,
% 3 clean runs), because there the neutral glue already sat between an L run and
% an R run. The PDF outline is byte-identical too -- hyperref's \pdfstringdef
% neutralises \mbox, so bookmark strings keep their part numbers.
% \numberline{\thepart} would also fix the direction, but it re-geometries every
% part row in the fleet and DROPS the number from the part bookmark (hyperref
% strips \numberline). \mbox changes nothing that was already right.
% The search pattern deliberately carries no # token: etoolbox cannot match a
% parameter token in a patch pattern (measured -- all three chain members
% reported "not patchable" while the pattern included #1). "\thepart \hspace
% {1em}" is brace-balanced and occurs exactly once in the wrapper body.
\newif\iflakin@parttocnum
\def\lakin@parttocnum#1{\iflakin@parttocnum\else
  \patchcmd{#1}{\thepart \hspace {1em}}{\mbox {\thepart }\hspace {1em}}
    {\lakin@parttocnumtrue}{}\fi}
\@ifundefined{NR@part}{}{\lakin@parttocnum\NR@part}
\@ifundefined{H@old@part}{}{\lakin@parttocnum\H@old@part}
\@ifundefined{@part}{}{\lakin@parttocnum\@part}
\iflakin@parttocnum\else\PackageWarning{hebrew-lualatex-pdf}{Part TOC number is
  not boxed: no known \noexpand\part wrapper matched. A part title that BEGINS
  with Latin will print its number on the wrong side of its TOC row}\fi
\patchcmd{\section}{\bfseries}{\bfseries\boldmath}{}
  {\PackageError{hebrew-lualatex-pdf}{Cannot patch section math weight}{Unsupported document class}}
\patchcmd{\subsection}{\bfseries}{\bfseries\boldmath}{}
  {\PackageError{hebrew-lualatex-pdf}{Cannot patch subsection math weight}{Unsupported document class}}
\patchcmd{\subsubsection}{\bfseries}{\bfseries\boldmath}{}
  {\PackageError{hebrew-lualatex-pdf}{Cannot patch subsubsection math weight}{Unsupported document class}}
\patchcmd{\paragraph}{\bfseries}{\bfseries\boldmath}{}{}
\patchcmd{\subparagraph}{\bfseries}{\bfseries\boldmath}{}{}
\@ifundefined{l@chapter}{}{%
  \patchcmd{\l@chapter}{\bfseries}{\bfseries\boldmath}{}{}}
\makeatother
\setcounter{tocdepth}{\lakittocdepth}
% <hllpdf:measured-toc>
% v4.16.0: Measure actual number labels, not a guessed digit count. Read only
% canonical contentsline records; never execute titles or addtocontents code
% during the measuring pass. Widths are local to this individual list and are
% available before its first row: no per-row drift and no extra aux pass.
\makeatletter
\newread\hllpdf@tocstream
\newdimen\hllpdf@toc@chapter
\newdimen\hllpdf@toc@section
\newdimen\hllpdf@toc@subsection
\newdimen\hllpdf@toc@subsubsection
\newdimen\hllpdf@toc@paragraph
\newdimen\hllpdf@toc@subparagraph
\newdimen\hllpdf@toc@figure
\newdimen\hllpdf@toc@table
\newdimen\hllpdf@tocpage
\newdimen\hllpdf@toc@sectionindent
\newdimen\hllpdf@toc@subsectionindent
\newdimen\hllpdf@toc@subsubsectionindent
\newdimen\hllpdf@toc@paragraphindent
\newdimen\hllpdf@toc@subparagraphindent
\def\hllpdf@tocdefaults{%
  \hllpdf@toc@chapter=2em
  \hllpdf@toc@section=3.2em
  \hllpdf@toc@subsection=3.9em
  \@ifundefined{c@chapter}{%
    \hllpdf@toc@subsubsection=3.2em
    \hllpdf@toc@paragraph=4.1em
    \hllpdf@toc@subparagraph=5em
  }{%
    \hllpdf@toc@subsubsection=4.1em
    \hllpdf@toc@paragraph=5em
    \hllpdf@toc@subparagraph=6em
  }%
  \hllpdf@toc@figure=2.3em
  \hllpdf@toc@table=2.3em
  \hllpdf@tocpage=\@pnumwidth\relax
}
\def\hllpdf@tocindents{%
  % Preserve the established shallow layout when its slots already fit.
  \hllpdf@toc@sectionindent=1.5em
  \@tempdima=\hllpdf@toc@chapter \advance\@tempdima by -2em
  \advance\hllpdf@toc@sectionindent by \@tempdima
  \hllpdf@toc@subsectionindent=\dimexpr
    \hllpdf@toc@sectionindent+\hllpdf@toc@section\relax
  \hllpdf@toc@subsubsectionindent=\dimexpr
    \hllpdf@toc@subsectionindent+\hllpdf@toc@subsection\relax
  \hllpdf@toc@paragraphindent=\dimexpr
    \hllpdf@toc@subsubsectionindent+\hllpdf@toc@subsubsection\relax
  \hllpdf@toc@subparagraphindent=\dimexpr
    \hllpdf@toc@paragraphindent+\hllpdf@toc@paragraph\relax
}
\def\hllpdf@tocfont@chapter{\bfseries\boldmath}
\def\hllpdf@tocfont@part{\large\bfseries\boldmath}
\def\hllpdf@tocentryfont{%
  \ifcsname hllpdf@tocfont@\hllpdf@toclevel\endcsname
    \csname hllpdf@tocfont@\hllpdf@toclevel\endcsname
  \fi}
\long\def\hllpdf@tocdiscard#1\hllpdf@tocstop{}
\def\hllpdf@tocscanrecord{%
  \@ifnextchar\contentsline{\hllpdf@tocscanline}{%
    \@ifnextchar\babel@toc{\hllpdf@tocscanlanguage}{\hllpdf@tocdiscard}}}
\long\def\hllpdf@tocscanlanguage\babel@toc#1#2#3\hllpdf@tocstop{%
  % Babel records are declarations, not arbitrary code. Track the language
  % without executing the original directive or changing the surrounding list.
  \def\hllpdf@toclanguage{#1}}
\long\def\hllpdf@tocscanline\contentsline#1#2#3#4#5\hllpdf@tocstop{%
  \def\hllpdf@toclevel{#1}%
  \ifcsname hllpdf@toc@#1\endcsname
    \hllpdf@tocscantitle#2\hllpdf@tocstop
  \fi
  % The page-number column needs its own measurement (including bold chapter
  % rows). Its extra clearance separates the page label from dot leaders.
  \setbox\@tempboxa=\hbox{\foreignlanguage{\hllpdf@toclanguage}{\hllpdf@tocentryfont #3\kern.3em}}%
  \ifdim\wd\@tempboxa>\hllpdf@tocpage
    \hllpdf@tocpage=\wd\@tempboxa
  \fi}
\def\hllpdf@tocscantitle{%
  \@ifnextchar\numberline{\hllpdf@tocscanlabel}{\hllpdf@tocdiscard}}
\long\def\hllpdf@tocscanlabel\numberline#1#2\hllpdf@tocstop{%
  \setbox\@tempboxa=\hbox{\foreignlanguage{\hllpdf@toclanguage}{\hllpdf@tocentryfont #1\kern.5em}}%
  \ifdim\wd\@tempboxa>\csname hllpdf@toc@\hllpdf@toclevel\endcsname
    \csname hllpdf@toc@\hllpdf@toclevel\endcsname=\wd\@tempboxa
  \fi}
\def\hllpdf@tocread{%
  \ifeof\hllpdf@tocstream\else
    \read\hllpdf@tocstream to\hllpdf@tocrecord
    \expandafter\hllpdf@tocscanrecord\hllpdf@tocrecord\hllpdf@tocstop
    \expandafter\hllpdf@tocread
  \fi}
\def\hllpdf@tocmeasure#1{%
  \makeatletter
  \edef\hllpdf@toclanguage{\languagename}%
  \hllpdf@tocdefaults
  \openin\hllpdf@tocstream=\jobname.#1\relax
  \hllpdf@tocread
  \closein\hllpdf@tocstream
  \hllpdf@tocindents
  \edef\@pnumwidth{\the\hllpdf@tocpage}%
  \@tempdima=\dimexpr\hllpdf@tocpage+1em\relax
  \ifdim\@tempdima>\@tocrmarg\relax
    \edef\@tocrmarg{\the\@tempdima}%
  \fi}
\let\hllpdf@originalstarttoc\@starttoc
\def\@starttoc#1{%
  \begingroup
    \def\hllpdf@tocext{#1}%
    \def\hllpdf@tocwanted{toc}%
    \ifx\hllpdf@tocext\hllpdf@tocwanted \hllpdf@tocmeasure{#1}\else
      \def\hllpdf@tocwanted{lof}%
      \ifx\hllpdf@tocext\hllpdf@tocwanted \hllpdf@tocmeasure{#1}\else
        \def\hllpdf@tocwanted{lot}%
        \ifx\hllpdf@tocext\hllpdf@tocwanted \hllpdf@tocmeasure{#1}\fi
      \fi
    \fi
    \hllpdf@originalstarttoc{#1}%
  \endgroup}
\hllpdf@tocdefaults
\hllpdf@tocindents
\@ifundefined{l@chapter}{}{%
  \patchcmd{\l@chapter}{\setlength\@tempdima{1.5em}}%
    {\setlength\@tempdima{\hllpdf@toc@chapter}}{}%
    {\PackageError{hebrew-lualatex-pdf}{Cannot size chapter TOC numbers}%
       {Unsupported contents layout: the chapter number slot was not found}}}
\renewcommand*\l@section{\@dottedtocline{1}%
  {\hllpdf@toc@sectionindent}{\hllpdf@toc@section}}
\renewcommand*\l@subsection{\@dottedtocline{2}%
  {\hllpdf@toc@subsectionindent}{\hllpdf@toc@subsection}}
\renewcommand*\l@subsubsection{\@dottedtocline{3}%
  {\hllpdf@toc@subsubsectionindent}{\hllpdf@toc@subsubsection}}
\renewcommand*\l@paragraph{\@dottedtocline{4}%
  {\hllpdf@toc@paragraphindent}{\hllpdf@toc@paragraph}}
\renewcommand*\l@subparagraph{\@dottedtocline{5}%
  {\hllpdf@toc@subparagraphindent}{\hllpdf@toc@subparagraph}}
\renewcommand*\l@figure{\@dottedtocline{1}{1.5em}{\hllpdf@toc@figure}}
\renewcommand*\l@table{\@dottedtocline{1}{1.5em}{\hllpdf@toc@table}}
\makeatother
% </hllpdf:measured-toc>
\setcounter{secnumdepth}{\lakitsecnumdepth}
% Public preamble-only caption-numbering policy. The historical within-major
% default is untouched unless this command is called. "continuous" removes the
% reset edge as well as the printed prefix, so Figure 1/2 and Table 1/2 stay
% unique across sections/chapters. Returning to "within-major" restores the
% original class-aware displays and native destination conventions. The policy
% never replaces \figcap/\tabcap or their measured link-placement mechanism.
\makeatletter
\let\lakin@majorfig\thefig
\let\lakin@majortbl\thetbl
\let\lakin@majorHfig\theHfig
\let\lakin@majorHtbl\theHtbl
\newcommand{\lakitcaptionnumbering}[1]{%
  \edef\lakin@numbering{#1}%
  \def\lakin@continuous{continuous}\def\lakin@withinmajor{within-major}%
  \ifx\lakin@numbering\lakin@continuous
    \@ifundefined{c@chapter}{%
      \@removefromreset{fig}{section}\@removefromreset{tbl}{section}%
    }{\@removefromreset{fig}{chapter}\@removefromreset{tbl}{chapter}}%
    \renewcommand{\thefig}{\arabic{fig}}\renewcommand{\thetbl}{\arabic{tbl}}%
    \def\theHfig{\arabic{fig}}\def\theHtbl{\arabic{tbl}}%
  \else\ifx\lakin@numbering\lakin@withinmajor
    \@ifundefined{c@chapter}{%
      \@removefromreset{fig}{section}\@removefromreset{tbl}{section}%
      \@addtoreset{fig}{section}\@addtoreset{tbl}{section}%
    }{%
      \@removefromreset{fig}{chapter}\@removefromreset{tbl}{chapter}%
      \@addtoreset{fig}{chapter}\@addtoreset{tbl}{chapter}}%
    \let\thefig\lakin@majorfig\let\thetbl\lakin@majortbl
    \let\theHfig\lakin@majorHfig\let\theHtbl\lakin@majorHtbl
  \else
    \PackageError{hebrew-english-latex-studio}{Unknown caption-numbering policy '#1'}%
      {Use continuous or within-major before begin{document}.}%
  \fi\fi}
\@onlypreamble\lakitcaptionnumbering
\makeatother
% Architecture readiness sentinel. Optional runtime layers require the complete
% core, so the marker is defined only after every mechanism/baseline definition.
% <hllpdf:architecture-ready>
\def\hllpdfcoreloaded{1}
% </hllpdf:architecture-ready>

% v4.22: registered file resources, not an assertion about embedded subset bytes.
% The log framing avoids long-path wrapping and preserves UTF-8 as hex.
% <hllpdf:font-inventory>
\directlua{
  local function inventory()
    local names = {}
    local unnamed = 0
    for id, item in font.each() do
      local path = item.filename
      if path and not (path == "") then names[path] = true else unnamed = unnamed + 1 end
    end
    local sorted = {}
    for path in pairs(names) do table.insert(sorted,path) end
    table.sort(sorted)
    texio.write_nl("log","HLLFONT|BEGIN|1")
    for index,path in ipairs(sorted) do
      local hex = path:gsub(".",function(c) return string.format(string.char(37) .. "02x",string.byte(c)) end)
      local count = math.ceil(string.len(hex)/48)
      texio.write_nl("log","HLLFONT|FILE|" .. index .. "|" .. count)
      for offset=1,string.len(hex),48 do texio.write_nl("log","HLLFONT|HEX|" .. hex:sub(offset,offset+47)) end
    end
    texio.write_nl("log","HLLFONT|END|" .. rawlen(sorted) .. "|" .. unnamed)
    texio.write_nl("log","")
  end
  luatexbase.add_to_callback("finish_pdffile",inventory,"hllpdf.runtime.fonts")
}
% </hllpdf:font-inventory>
:
%   \lakithead{מבנה המחשב}
% Every content page then carries the series' orientation line: chapter name
% at the reading edge, book title opposite, thin rule. (fancyhdr's [L]/[R]
% are PHYSICAL; under the RTL page direction [L] lands at the RIGHT/reading
% edge — measured, consistent across the corpus.) Packaged like the topic
% rail: always loaded, inert without the call. This centralizes the 8-line
% block that 8 lakits used to copy by hand — the same duplication vector
% that once let a lakit preamble go ~100 lines stale. Class-aware: report
% fills the mark from \chapter, article from \section. NB: a titlepage env
% self-empties; front matter that sets \pagestyle{empty} must restore with
% \pagestyle{fancy} afterwards (NOT plain — that discards the head).
% ============================================================================
\usepackage{fancyhdr}
% v3.61: \chapter* left a STALE running head — pages of an UNNUMBERED chapter
% carried the PREVIOUS chapter's title. report.cls routes \chapter* through
% \@schapter, which never calls \chaptermark, so \leftmark keeps whatever the last
% numbered chapter left. Measured: with \chapter{פרק ממוספר} then
% \chapter*{מבוא לא ממוספר}, page 4 (inside the starred chapter) printed
% "פרק ממוספר". Invisible on the chapter's FIRST page (\thispagestyle{plain}), so it
% only shows from the second page on — the pages an author is least likely to check.
% An unnumbered front-matter chapter (\chapter*{מבוא}, \chapter*{הקדמה}) is everyday
% Hebrew-academic structure, and the corpus already carries HAND-WRITTEN \markboth
% lines right after \chapter* — the exact per-document duplication \lakithead exists
% to retire. Patched INSIDE \lakithead so the macro stays inert when unused; a
% document's own later \markboth still wins, so those hand-patches do not regress.
\makeatletter
% <hllpdf:starred-heading-scope>
\newif\iflakin@starmarkpatched
\newcommand{\lakistarmark}{%
  \iflakin@starmarkpatched\else
  % Dispatch on the actual class counter, not hyperref's shared internals.
  \ifdefined\c@chapter
    \let\HEBorigschapter\@schapter
    % Record only after the native heading has reached its page. A mark
    % before the native routine can advertise the next heading too early.
    \def\@schapter##1{\HEBorigschapter{##1}\chaptermark{##1}}%
  \else
    % @ssect receives layout arguments, NOT the section name or level.
    % Capture the name in @startsection without introducing a TeX group:
    % grouping sectioning would break the native after-heading state.
    \def\HEB@sectionkind{}%
    \def\HEB@mainsection{section}%
    \let\HEBorigstartsection\@startsection
    \def\@startsection##1##2##3##4##5##6{%
      \def\HEB@sectionkind{##1}%
      \HEBorigstartsection{##1}{##2}{##3}{##4}{##5}{##6}}%
    \let\HEBorigssect\@ssect
    \def\@ssect##1##2##3##4##5{%
      \HEBorigssect{##1}{##2}{##3}{##4}{##5}%
      % Lower starred headings must not replace the article's section mark.
      \ifx\HEB@sectionkind\HEB@mainsection
        \sectionmark{##5}%
      \fi
      \ignorespaces}%
  \fi
  \global\lakin@starmarkpatchedtrue
  \fi}
% </hllpdf:starred-heading-scope>
% <hllpdf:running-language>
% Choose the wrapper when the mark is written, not when a later page ships.
% English punctuation must remain LTR even after returning to Hebrew prose.
\newcommand{\HEB@runningmark}[1]{%
  \iflanguage{english}{\markboth{\en{#1}}{}}{\markboth{\h{#1}}{}}}
% </hllpdf:running-language>
\makeatother
% v4.2.1 (NAV): every collection already names itself in \lakithead — feed the
% PDF /Title from the same declaration, so the navigation gate's
% PDF_REQUIRE_TITLE has a universal source and no fleet member ships
% title-less by missing a second declaration site. Measured before this: 14
% of 16 collections shipped /Title-empty while every one carried a perfectly
% good \lakithead string (plain Hebrew in all 16 — no math, no formatting, so
% the PDF-string path is clean). An explicit title always wins: set it with
% \lakitpdftitle{...} anywhere before \lakithead, or with a plain
% \hypersetup{pdftitle=...} — the fallback detects that via hyperref's
% \@pdftitle, probe-verified \@empty-equal until the first explicit set.
\makeatletter
\newif\iflakit@pdftitleset
\newcommand{\lakitpdftitle}[1]{%
  \hypersetup{pdftitle={#1}}\global\lakit@pdftitlesettrue}
\newcommand{\lakithead}[1]{%
  \iflakit@pdftitleset\else
    \ifx\@pdftitle\@empty
      \hypersetup{pdftitle={#1}}\global\lakit@pdftitlesettrue
    \fi
  \fi
  \pagestyle{fancy}\fancyhf{}%
  \lakistarmark

  \fancyhead[L]{\small\textcolor{cGray}{\leftmark}}%
  % <hllpdf:fixed-head-language>
  % The fixed document title has its OWN language snapshot, independent of
  % the current section/chapter mark. Choose the branch NOW, not at shipout:
  % the page may be Hebrew by the time an English title is finally printed.
  \iflanguage{english}{%
    \fancyhead[R]{\small\textcolor{cGray}{\en{#1}}}%
  }{%
    \fancyhead[R]{\small\textcolor{cGray}{\h{#1}}}%
  }%
  % </hllpdf:fixed-head-language>
  \fancyfoot[C]{\thepage}%
  \renewcommand{\headrulewidth}{0.4pt}%
  \ifdefined\chaptermark
    \renewcommand{\chaptermark}[1]{\HEB@runningmark{##1}}%
  \else
    \renewcommand{\sectionmark}[1]{\HEB@runningmark{##1}}%
  \fi}
\makeatother

% v3.39.1: path-independent byte-reproducibility. LuaTeX derives the trailer
% /ID from the OUTPUT PATH, so two clones of the same source at different
% directories produced different sha256 (measured: the ONLY delta was the
% 62-byte /ID line). Bit 512 suppresses exactly that line — measured: Creator/
% Producer/CreationDate/ModDate and all content untouched. /ID is optional
% per ISO 32000 (required only for ENCRYPTED PDFs and PDF/A — neither is a
% target here). If a future workflow needs an ID (PDF/A, Acrobat revision
% matching), delete this line and it returns.
\pdfvariable suppressoptionalinfo \numexpr 512\relax

% v3.75: BOLD MATH IN BOLD HEADINGS. Heading text is \bfseries but math inside
% it stayed regular-weight CM — a hairline $V_T$ inside a heavy 24.8pt chapter
% title (measured: CMMI12 fl=6 next to FrankRuhlLibre-Bold fl=20; 155 live
% heading-with-math sites fleet-wide). \boldmath rides along wherever the class
% sets \bfseries for a heading; A/B measured: $V_T$ -> CMMIB10 fl=22, 0 missing
% characters (the bold CM faces ship with this TeX Live). Guarded per-macro so
% the article class (no \@makechapterhead) and future class changes stay safe.
% v4.0.1 (adoption fix): the SEARCH patterns above carry no parameter token.
% v4.0.0 made them specific by including the class's own #1/#2, and measured on
% TeX Live 2023 that never matches: a \patchcmd search pattern containing a
% parameter token cannot match while the call sits inside \@ifundefined{...}{}{...}
% (probed both #1 and ##1 -> FAIL; the same pattern at top level succeeds). Every
% one of these patches therefore failed, and because v4.0.0 also replaced the
% empty failure branch with \PackageError, the failure became FATAL: all four
% bundled example assets stopped building ("Cannot patch chapter-title math
% weight -> Fatal error, no output PDF"). The \PackageError branches are KEPT —
% a silent no-op here is what let v3.75's fix rot unnoticed — only the parameter
% tokens are dropped, so the patches match and a real class change still errors.
\makeatletter
\@ifundefined{@makechapterhead}{}{%
  \patchcmd{\@makechapterhead}{\huge\bfseries \@chapapp}
    {\huge\bfseries\boldmath \@chapapp}{}
    {\PackageError{hebrew-lualatex-pdf}{Cannot patch chapter-number math weight}{Unsupported document class}}%
  \patchcmd{\@makechapterhead}{\Huge \bfseries}
    {\Huge \bfseries\boldmath}{}
    {\PackageError{hebrew-lualatex-pdf}{Cannot patch chapter-title math weight}{Unsupported document class}}%
  \patchcmd{\@makeschapterhead}{\Huge \bfseries}
    {\Huge \bfseries\boldmath}{}
    {\PackageError{hebrew-lualatex-pdf}{Cannot patch starred chapter math weight}{Unsupported document class}}}
% v4.0.1: \part sits under a THREE-DEEP wrapper chain, so the class-based anchor
% is not in \@part at all. Measured with \show: hyperref rewrites \@part to
% `...\phantomsection\fi \H@old@part`, and nameref rewrites THAT to
% `\NR@gettitle{#1}\NR@part[{#1}]{#2}` — the class body ends up in \NR@part.
% v4.0.0 patched \@part and made failure fatal, which is why all seven bundled
% examples stopped building. \part is NOT a dead heading level here: 188 live uses
% across the fleet, so the bold-math fix is worth getting right rather than
% dropping. We walk the chain innermost-first and stop at the first hit. Failure is
% a WARNING, not an error: the chain is a property of hyperref/nameref versions, so
% a future package release must degrade to regular-weight part math — a cosmetic
% loss — instead of refusing to build. Chapter and section keep \PackageError,
% because those targets are in the CLASS and were verified to patch here.
\newif\iflakin@partnum \newif\iflakin@parttitle \newif\iflakin@sparttitle
\def\lakin@partnum#1#2{\iflakin@partnum\else
  \patchcmd{#1}{#2\bfseries \partname}{#2\bfseries\boldmath \partname}
    {\lakin@partnumtrue}{}\fi}
\def\lakin@parttitle#1#2{\iflakin@parttitle\else
  \patchcmd{#1}{#2 \bfseries}{#2 \bfseries\boldmath}{\lakin@parttitletrue}{}\fi}
\def\lakin@sparttitle#1#2{\iflakin@sparttitle\else
  \patchcmd{#1}{#2 \bfseries}{#2 \bfseries\boldmath}{\lakin@sparttitletrue}{}\fi}
\@ifundefined{c@chapter}{\def\lakin@pnum{\Large}\def\lakin@ptit{\huge}}%
                        {\def\lakin@pnum{\huge}\def\lakin@ptit{\Huge}}
\@ifundefined{NR@part}{}{\expandafter\lakin@partnum\expandafter\NR@part\expandafter{\lakin@pnum}%
                         \expandafter\lakin@parttitle\expandafter\NR@part\expandafter{\lakin@ptit}}
\@ifundefined{H@old@part}{}{\expandafter\lakin@partnum\expandafter\H@old@part\expandafter{\lakin@pnum}%
                            \expandafter\lakin@parttitle\expandafter\H@old@part\expandafter{\lakin@ptit}}
\@ifundefined{@part}{}{\expandafter\lakin@partnum\expandafter\@part\expandafter{\lakin@pnum}%
                       \expandafter\lakin@parttitle\expandafter\@part\expandafter{\lakin@ptit}}
\@ifundefined{NR@spart}{}{\expandafter\lakin@sparttitle\expandafter\NR@spart\expandafter{\lakin@ptit}}
\@ifundefined{H@old@spart}{}{\expandafter\lakin@sparttitle\expandafter\H@old@spart\expandafter{\lakin@ptit}}
\@ifundefined{@spart}{}{\expandafter\lakin@sparttitle\expandafter\@spart\expandafter{\lakin@ptit}}
\iflakin@partnum\else\PackageWarning{hebrew-lualatex-pdf}{Part-number math stays
  regular weight: no known \noexpand\part wrapper matched (hyperref/nameref
  version change). Cosmetic only}\fi
\iflakin@parttitle\else\PackageWarning{hebrew-lualatex-pdf}{Part-title math stays
  regular weight: no known \noexpand\part wrapper matched. Cosmetic only}\fi
% v4.12.0: THE PART NUMBER IN THE TOC WAS NOT DIRECTIONALLY SEPARATED FROM THE
% TITLE. report.cls/article.cls compose the part entry as ONE text stream --
% \addcontentsline{toc}{part}{\thepart\hspace{1em}#1} -- while EVERY other level
% goes through \numberline, which puts the number in an \hbox. A box is a bidi
% OBJECT and cannot join a neighbouring run; a bare "IV" can. So in the RTL TOC
% paragraph a Latin number and a Latin-INITIAL title merge into a SINGLE LTR run
% and are laid out left-to-right as a unit: the number lands on the wrong side of
% the row, and a mixed title is scrambled outright. Measured on \part{MCP ...}:
% the printed order was "<tail> V MCP" -- title fragment, number, title fragment.
% The chapter/section rows one line below, with the SAME Latin-initial titles,
% were correct on the same page. That contrast is the whole diagnosis.
% \en{} IS NOT A WORKAROUND, measured both ways on one fixture: it fixes a
% PURE-Latin title (\part{\en{Hooks}} renders correctly) and does NOT fix a MIXED
% one (\part{\en{MCP} <hebrew>} stays scrambled). babel's language switch changes
% font and hyphenation, not the bidi run structure -- so the defect is not
% reachable from the source side at all, which is why it belongs here.
% Boxing the number is the entire fix and is a STRICT NO-OP where the bug never
% bit: a Hebrew-titled document compiles BYTE-IDENTICAL before and after (md5,
% 3 clean runs), because there the neutral glue already sat between an L run and
% an R run. The PDF outline is byte-identical too -- hyperref's \pdfstringdef
% neutralises \mbox, so bookmark strings keep their part numbers.
% \numberline{\thepart} would also fix the direction, but it re-geometries every
% part row in the fleet and DROPS the number from the part bookmark (hyperref
% strips \numberline). \mbox changes nothing that was already right.
% The search pattern deliberately carries no # token: etoolbox cannot match a
% parameter token in a patch pattern (measured -- all three chain members
% reported "not patchable" while the pattern included #1). "\thepart \hspace
% {1em}" is brace-balanced and occurs exactly once in the wrapper body.
\newif\iflakin@parttocnum
\def\lakin@parttocnum#1{\iflakin@parttocnum\else
  \patchcmd{#1}{\thepart \hspace {1em}}{\mbox {\thepart }\hspace {1em}}
    {\lakin@parttocnumtrue}{}\fi}
\@ifundefined{NR@part}{}{\lakin@parttocnum\NR@part}
\@ifundefined{H@old@part}{}{\lakin@parttocnum\H@old@part}
\@ifundefined{@part}{}{\lakin@parttocnum\@part}
\iflakin@parttocnum\else\PackageWarning{hebrew-lualatex-pdf}{Part TOC number is
  not boxed: no known \noexpand\part wrapper matched. A part title that BEGINS
  with Latin will print its number on the wrong side of its TOC row}\fi
\patchcmd{\section}{\bfseries}{\bfseries\boldmath}{}
  {\PackageError{hebrew-lualatex-pdf}{Cannot patch section math weight}{Unsupported document class}}
\patchcmd{\subsection}{\bfseries}{\bfseries\boldmath}{}
  {\PackageError{hebrew-lualatex-pdf}{Cannot patch subsection math weight}{Unsupported document class}}
\patchcmd{\subsubsection}{\bfseries}{\bfseries\boldmath}{}
  {\PackageError{hebrew-lualatex-pdf}{Cannot patch subsubsection math weight}{Unsupported document class}}
\patchcmd{\paragraph}{\bfseries}{\bfseries\boldmath}{}{}
\patchcmd{\subparagraph}{\bfseries}{\bfseries\boldmath}{}{}
\@ifundefined{l@chapter}{}{%
  \patchcmd{\l@chapter}{\bfseries}{\bfseries\boldmath}{}{}}
\makeatother
\setcounter{tocdepth}{\lakittocdepth}
% <hllpdf:measured-toc>
% v4.16.0: Measure actual number labels, not a guessed digit count. Read only
% canonical contentsline records; never execute titles or addtocontents code
% during the measuring pass. Widths are local to this individual list and are
% available before its first row: no per-row drift and no extra aux pass.
\makeatletter
\newread\hllpdf@tocstream
\newdimen\hllpdf@toc@chapter
\newdimen\hllpdf@toc@section
\newdimen\hllpdf@toc@subsection
\newdimen\hllpdf@toc@subsubsection
\newdimen\hllpdf@toc@paragraph
\newdimen\hllpdf@toc@subparagraph
\newdimen\hllpdf@toc@figure
\newdimen\hllpdf@toc@table
\newdimen\hllpdf@tocpage
\newdimen\hllpdf@toc@sectionindent
\newdimen\hllpdf@toc@subsectionindent
\newdimen\hllpdf@toc@subsubsectionindent
\newdimen\hllpdf@toc@paragraphindent
\newdimen\hllpdf@toc@subparagraphindent
\def\hllpdf@tocdefaults{%
  \hllpdf@toc@chapter=2em
  \hllpdf@toc@section=3.2em
  \hllpdf@toc@subsection=3.9em
  \@ifundefined{c@chapter}{%
    \hllpdf@toc@subsubsection=3.2em
    \hllpdf@toc@paragraph=4.1em
    \hllpdf@toc@subparagraph=5em
  }{%
    \hllpdf@toc@subsubsection=4.1em
    \hllpdf@toc@paragraph=5em
    \hllpdf@toc@subparagraph=6em
  }%
  \hllpdf@toc@figure=2.3em
  \hllpdf@toc@table=2.3em
  \hllpdf@tocpage=\@pnumwidth\relax
}
\def\hllpdf@tocindents{%
  % Preserve the established shallow layout when its slots already fit.
  \hllpdf@toc@sectionindent=1.5em
  \@tempdima=\hllpdf@toc@chapter \advance\@tempdima by -2em
  \advance\hllpdf@toc@sectionindent by \@tempdima
  \hllpdf@toc@subsectionindent=\dimexpr
    \hllpdf@toc@sectionindent+\hllpdf@toc@section\relax
  \hllpdf@toc@subsubsectionindent=\dimexpr
    \hllpdf@toc@subsectionindent+\hllpdf@toc@subsection\relax
  \hllpdf@toc@paragraphindent=\dimexpr
    \hllpdf@toc@subsubsectionindent+\hllpdf@toc@subsubsection\relax
  \hllpdf@toc@subparagraphindent=\dimexpr
    \hllpdf@toc@paragraphindent+\hllpdf@toc@paragraph\relax
}
\def\hllpdf@tocfont@chapter{\bfseries\boldmath}
\def\hllpdf@tocfont@part{\large\bfseries\boldmath}
\def\hllpdf@tocentryfont{%
  \ifcsname hllpdf@tocfont@\hllpdf@toclevel\endcsname
    \csname hllpdf@tocfont@\hllpdf@toclevel\endcsname
  \fi}
\long\def\hllpdf@tocdiscard#1\hllpdf@tocstop{}
\def\hllpdf@tocscanrecord{%
  \@ifnextchar\contentsline{\hllpdf@tocscanline}{%
    \@ifnextchar\babel@toc{\hllpdf@tocscanlanguage}{\hllpdf@tocdiscard}}}
\long\def\hllpdf@tocscanlanguage\babel@toc#1#2#3\hllpdf@tocstop{%
  % Babel records are declarations, not arbitrary code. Track the language
  % without executing the original directive or changing the surrounding list.
  \def\hllpdf@toclanguage{#1}}
\long\def\hllpdf@tocscanline\contentsline#1#2#3#4#5\hllpdf@tocstop{%
  \def\hllpdf@toclevel{#1}%
  \ifcsname hllpdf@toc@#1\endcsname
    \hllpdf@tocscantitle#2\hllpdf@tocstop
  \fi
  % The page-number column needs its own measurement (including bold chapter
  % rows). Its extra clearance separates the page label from dot leaders.
  \setbox\@tempboxa=\hbox{\foreignlanguage{\hllpdf@toclanguage}{\hllpdf@tocentryfont #3\kern.3em}}%
  \ifdim\wd\@tempboxa>\hllpdf@tocpage
    \hllpdf@tocpage=\wd\@tempboxa
  \fi}
\def\hllpdf@tocscantitle{%
  \@ifnextchar\numberline{\hllpdf@tocscanlabel}{\hllpdf@tocdiscard}}
\long\def\hllpdf@tocscanlabel\numberline#1#2\hllpdf@tocstop{%
  \setbox\@tempboxa=\hbox{\foreignlanguage{\hllpdf@toclanguage}{\hllpdf@tocentryfont #1\kern.5em}}%
  \ifdim\wd\@tempboxa>\csname hllpdf@toc@\hllpdf@toclevel\endcsname
    \csname hllpdf@toc@\hllpdf@toclevel\endcsname=\wd\@tempboxa
  \fi}
\def\hllpdf@tocread{%
  \ifeof\hllpdf@tocstream\else
    \read\hllpdf@tocstream to\hllpdf@tocrecord
    \expandafter\hllpdf@tocscanrecord\hllpdf@tocrecord\hllpdf@tocstop
    \expandafter\hllpdf@tocread
  \fi}
\def\hllpdf@tocmeasure#1{%
  \makeatletter
  \edef\hllpdf@toclanguage{\languagename}%
  \hllpdf@tocdefaults
  \openin\hllpdf@tocstream=\jobname.#1\relax
  \hllpdf@tocread
  \closein\hllpdf@tocstream
  \hllpdf@tocindents
  \edef\@pnumwidth{\the\hllpdf@tocpage}%
  \@tempdima=\dimexpr\hllpdf@tocpage+1em\relax
  \ifdim\@tempdima>\@tocrmarg\relax
    \edef\@tocrmarg{\the\@tempdima}%
  \fi}
\let\hllpdf@originalstarttoc\@starttoc
\def\@starttoc#1{%
  \begingroup
    \def\hllpdf@tocext{#1}%
    \def\hllpdf@tocwanted{toc}%
    \ifx\hllpdf@tocext\hllpdf@tocwanted \hllpdf@tocmeasure{#1}\else
      \def\hllpdf@tocwanted{lof}%
      \ifx\hllpdf@tocext\hllpdf@tocwanted \hllpdf@tocmeasure{#1}\else
        \def\hllpdf@tocwanted{lot}%
        \ifx\hllpdf@tocext\hllpdf@tocwanted \hllpdf@tocmeasure{#1}\fi
      \fi
    \fi
    \hllpdf@originalstarttoc{#1}%
  \endgroup}
\hllpdf@tocdefaults
\hllpdf@tocindents
\@ifundefined{l@chapter}{}{%
  \patchcmd{\l@chapter}{\setlength\@tempdima{1.5em}}%
    {\setlength\@tempdima{\hllpdf@toc@chapter}}{}%
    {\PackageError{hebrew-lualatex-pdf}{Cannot size chapter TOC numbers}%
       {Unsupported contents layout: the chapter number slot was not found}}}
\renewcommand*\l@section{\@dottedtocline{1}%
  {\hllpdf@toc@sectionindent}{\hllpdf@toc@section}}
\renewcommand*\l@subsection{\@dottedtocline{2}%
  {\hllpdf@toc@subsectionindent}{\hllpdf@toc@subsection}}
\renewcommand*\l@subsubsection{\@dottedtocline{3}%
  {\hllpdf@toc@subsubsectionindent}{\hllpdf@toc@subsubsection}}
\renewcommand*\l@paragraph{\@dottedtocline{4}%
  {\hllpdf@toc@paragraphindent}{\hllpdf@toc@paragraph}}
\renewcommand*\l@subparagraph{\@dottedtocline{5}%
  {\hllpdf@toc@subparagraphindent}{\hllpdf@toc@subparagraph}}
\renewcommand*\l@figure{\@dottedtocline{1}{1.5em}{\hllpdf@toc@figure}}
\renewcommand*\l@table{\@dottedtocline{1}{1.5em}{\hllpdf@toc@table}}
\makeatother
% </hllpdf:measured-toc>
\setcounter{secnumdepth}{\lakitsecnumdepth}
% Public preamble-only caption-numbering policy. The historical within-major
% default is untouched unless this command is called. "continuous" removes the
% reset edge as well as the printed prefix, so Figure 1/2 and Table 1/2 stay
% unique across sections/chapters. Returning to "within-major" restores the
% original class-aware displays and native destination conventions. The policy
% never replaces \figcap/\tabcap or their measured link-placement mechanism.
\makeatletter
\let\lakin@majorfig\thefig
\let\lakin@majortbl\thetbl
\let\lakin@majorHfig\theHfig
\let\lakin@majorHtbl\theHtbl
\newcommand{\lakitcaptionnumbering}[1]{%
  \edef\lakin@numbering{#1}%
  \def\lakin@continuous{continuous}\def\lakin@withinmajor{within-major}%
  \ifx\lakin@numbering\lakin@continuous
    \@ifundefined{c@chapter}{%
      \@removefromreset{fig}{section}\@removefromreset{tbl}{section}%
    }{\@removefromreset{fig}{chapter}\@removefromreset{tbl}{chapter}}%
    \renewcommand{\thefig}{\arabic{fig}}\renewcommand{\thetbl}{\arabic{tbl}}%
    \def\theHfig{\arabic{fig}}\def\theHtbl{\arabic{tbl}}%
  \else\ifx\lakin@numbering\lakin@withinmajor
    \@ifundefined{c@chapter}{%
      \@removefromreset{fig}{section}\@removefromreset{tbl}{section}%
      \@addtoreset{fig}{section}\@addtoreset{tbl}{section}%
    }{%
      \@removefromreset{fig}{chapter}\@removefromreset{tbl}{chapter}%
      \@addtoreset{fig}{chapter}\@addtoreset{tbl}{chapter}}%
    \let\thefig\lakin@majorfig\let\thetbl\lakin@majortbl
    \let\theHfig\lakin@majorHfig\let\theHtbl\lakin@majorHtbl
  \else
    \PackageError{hebrew-english-latex-studio}{Unknown caption-numbering policy '#1'}%
      {Use continuous or within-major before begin{document}.}%
  \fi\fi}
\@onlypreamble\lakitcaptionnumbering
\makeatother
% Architecture readiness sentinel. Optional runtime layers require the complete
% core, so the marker is defined only after every mechanism/baseline definition.
% <hllpdf:architecture-ready>
\def\hllpdfcoreloaded{1}
% </hllpdf:architecture-ready>

% v4.22: registered file resources, not an assertion about embedded subset bytes.
% The log framing avoids long-path wrapping and preserves UTF-8 as hex.
% <hllpdf:font-inventory>
\directlua{
  local function inventory()
    local names = {}
    local unnamed = 0
    for id, item in font.each() do
      local path = item.filename
      if path and not (path == "") then names[path] = true else unnamed = unnamed + 1 end
    end
    local sorted = {}
    for path in pairs(names) do table.insert(sorted,path) end
    table.sort(sorted)
    texio.write_nl("log","HLLFONT|BEGIN|1")
    for index,path in ipairs(sorted) do
      local hex = path:gsub(".",function(c) return string.format(string.char(37) .. "02x",string.byte(c)) end)
      local count = math.ceil(string.len(hex)/48)
      texio.write_nl("log","HLLFONT|FILE|" .. index .. "|" .. count)
      for offset=1,string.len(hex),48 do texio.write_nl("log","HLLFONT|HEX|" .. hex:sub(offset,offset+47)) end
    end
    texio.write_nl("log","HLLFONT|END|" .. rawlen(sorted) .. "|" .. unnamed)
    texio.write_nl("log","")
  end
  luatexbase.add_to_callback("finish_pdffile",inventory,"hllpdf.runtime.fonts")
}
% </hllpdf:font-inventory>

%
% Supported class contract: report or article.  Alternate paper sizes must also
% be declared to build.sh through EXPECTED_PAPER; the PDF gate never infers them.
% Do NOT use these hooks as a theme mechanism: theme overlays must not own page
% geometry, class selection, or direction semantics. A declared presentation-policy
% overlay may intentionally change numbering visibility after the preamble.
% ============================================================================
% Runtime architecture guard: generated from scripts/architecture_contract.py.
% The readiness sentinel remains delayed until the END of this file.
% <hllpdf:architecture-guard>
\ifdefined\hllpdfcoreloaded
  \errmessage{hllpdf: preamble.tex loaded more than once}
  \expandafter\endinput
\fi
% </hllpdf:architecture-guard>

\ifdefined\lakitdocumentclass\else\def\lakitdocumentclass{report}\fi
\ifdefined\lakitclassoptions\else\def\lakitclassoptions{12pt,a4paper}\fi
\ifdefined\lakitgeometryoptions\else\def\lakitgeometryoptions{margin=2.3cm,headheight=14pt}\fi
\ifdefined\lakitlinespread\else\def\lakitlinespread{1.3}\fi
\ifdefined\lakittocdepth\else\def\lakittocdepth{1}\fi
\ifdefined\lakitsecnumdepth\else\def\lakitsecnumdepth{2}\fi
% Paragraph organization belongs to the document layout, independently of its
% language and theme. Preserve historical class dimensions unless explicitly
% selecting block paragraphs; never suppress indentation with an everypar hack.
\ifdefined\lakitparagraphstyle\else\def\lakitparagraphstyle{indented}\fi
\edef\lakitparagraphstyle{\lakitparagraphstyle}
\def\hllpdfparagraphblock{block}
\def\hllpdfparagraphindented{indented}
\ifx\lakitparagraphstyle\hllpdfparagraphblock\else
  \ifx\lakitparagraphstyle\hllpdfparagraphindented\else
    \errmessage{hllpdf: paragraph style must be block or indented}
  \fi
\fi
% v4.3.1: the public class hook has a deliberately narrow support contract.
% Before this guard, \def\lakitdocumentclass{book} silently compiled even though
% every heading/TOC regression test targets report/article; an unsupported shell
% could therefore look "supported" until a subtle class-specific defect shipped.
% Fail before \documentclass rather than letting an undocumented combination enter
% the mechanism core. Plain-TeX comparison is used because LaTeX packages are not
% available yet at this point.
\edef\lakitvalidatedclass{\lakitdocumentclass}
\def\lakitreportclass{report}
\def\lakitarticleclass{article}
\ifx\lakitvalidatedclass\lakitreportclass
\else\ifx\lakitvalidatedclass\lakitarticleclass
\else
  \errmessage{hllpdf: unsupported class \lakitdocumentclass; use report or article}
\fi\fi
\let\lakitvalidatedclass\relax
\let\lakitreportclass\relax
\let\lakitarticleclass\relax
% Language intent is independent of document class, theme and content domain.
% The two dominant-language profiles differ in authoring intent, not direction.
% A balanced bilingual document must explicitly choose a navigation/base locale.
% <hllpdf:language-profile>
\providecommand{\lakitlanguageprofile}{hebrew}
\edef\lakitlanguageprofile{\lakitlanguageprofile}
\def\hllpdfhebrew{hebrew}\def\hllpdfenglish{english}
\def\hllpdfhebrewmixed{hebrew-with-english}
\def\hllpdfenglishmixed{english-with-hebrew}
\def\hllpdfbilingual{bilingual}
\let\hllpdfexpectedlanguage\relax
\ifx\lakitlanguageprofile\hllpdfhebrew
  \let\hllpdfexpectedlanguage\hllpdfhebrew
\else\ifx\lakitlanguageprofile\hllpdfhebrewmixed
  \let\hllpdfexpectedlanguage\hllpdfhebrew
\else\ifx\lakitlanguageprofile\hllpdfenglish
  \let\hllpdfexpectedlanguage\hllpdfenglish
\else\ifx\lakitlanguageprofile\hllpdfenglishmixed
  \let\hllpdfexpectedlanguage\hllpdfenglish
\else\ifx\lakitlanguageprofile\hllpdfbilingual
  \ifdefined\lakitmainlanguage\else
    \errmessage{hllpdf: bilingual profile requires explicit lakitmainlanguage}
  \fi
\else
  \errmessage{hllpdf: unsupported language profile \lakitlanguageprofile}
\fi\fi\fi\fi\fi
\ifdefined\lakitmainlanguage
  \edef\lakitmainlanguage{\lakitmainlanguage}
\else
  \let\lakitmainlanguage\hllpdfexpectedlanguage
\fi
\ifx\lakitmainlanguage\hllpdfhebrew\else
  \ifx\lakitmainlanguage\hllpdfenglish\else
    \errmessage{hllpdf: main language must be hebrew or english}
  \fi
\fi
\ifx\hllpdfexpectedlanguage\relax\else
  \ifx\lakitmainlanguage\hllpdfexpectedlanguage\else
    \errmessage{hllpdf: main language conflicts with language profile}
  \fi
\fi
\typeout{HLLLANG|\lakitlanguageprofile|\lakitmainlanguage}
% </hllpdf:language-profile>
\documentclass[\lakitclassoptions]{\lakitdocumentclass}
\usepackage[no-math]{fontspec}        % (1) ★ no-math — THE critical line (see recipe.md)
% layout=lists asks Babel to adapt list parshape when paragraph and page
% directions differ. Without it, minority-language list bodies lose their
% indentation and labels hang outside the text block, despite a clean log.
% The upstream mechanism preserves ordinary same-direction list geometry.
\usepackage[bidi=basic,layout=lists]{babel}
\ifx\lakitmainlanguage\hllpdfenglish
  \babelprovide[import]{hebrew}
  \babelprovide[import,main]{english}
\else
  \babelprovide[import,main]{hebrew}
  \babelprovide[import]{english}
\fi
% v4.22: bind canonical families and fallbacks to the selected installation.
% Lookup by basename or family alone is not proof of the file actually used.
% <hllpdf:font-binding>
\directlua{
  local dir = os.getenv("HEBREW_LATEX_FONT_DIR")
  if not dir or dir == "" then
    local data = os.getenv("XDG_DATA_HOME")
    if not data or data == "" or not (data:sub(1,1) == "/") then
      data = (os.getenv("HOME") or "") .. "/.local/share"
    end
    dir = data .. "/fonts/hebrew-lualatex-pdf"
  end
  if not (dir:sub(1,1) == "/") then dir = lfs.currentdir() .. "/" .. dir end
  if dir:find("[",1,true) or dir:find("]",1,true) or dir:find(string.char(0),1,true)
     or dir:find(string.char(10),1,true) or dir:find(string.char(13),1,true) then
    tex.error("hllpdf: unsupported character in font directory")
  end
  if not (lfs.attributes(dir,"mode") == "directory") then
    tex.error("hllpdf: font directory is missing")
  end
  hllpdf_font_dir = dir:gsub("/+$", "") .. "/"
  tex.sprint(string.char(92) .. "def" .. string.char(92) .. "hllpdffontdir{")
  tex.sprint(-2,hllpdf_font_dir)
  tex.sprint("}")
  luaotfload.add_fallback("hebfb",{"[" .. hllpdf_font_dir .. "DejaVuSerif.ttf]:mode=node;"})
  luaotfload.add_fallback("hebfbbf",{"[" .. hllpdf_font_dir .. "DejaVuSerif-Bold.ttf]:mode=node;"})
}
% </hllpdf:font-binding>
% v3.22: Frank Ruhl Libre ships NO italic face. Instead of silently-upright
% \emph (v<=3.20) or bold-as-emphasis (v3.21), synthesize a slanted face with
% FakeSlant — real, VISIBLE italic emphasis with correct \emph nesting
% semantics (inner \emph toggles back upright). Verified empirically at
% 300dpi: slant clearly visible, 0 errors, 0 missing glyphs.
\babelfont{rm}[Path={\hllpdffontdir},RawFeature={fallback=hebfb},
  BoldFont={FrankRuhlLibre-Bold.ttf},
  BoldFeatures={RawFeature={fallback=hebfbbf}},   % v3.56 — see hebfbbf above
  ItalicFont={FrankRuhlLibre-Regular.ttf}, ItalicFeatures={FakeSlant=0.22},
  % Bold-italic uses the weight-matched fallback and a synthetic slant.
  BoldItalicFont={FrankRuhlLibre-Bold.ttf}, BoldItalicFeatures={RawFeature={fallback=hebfbbf},FakeSlant=0.22}
  ]{FrankRuhlLibre-Regular.ttf}
% Sans fallbacks preserve family and weight for the few missing codepoints.
\directlua{luaotfload.add_fallback("hebfbsans",{"[" .. hllpdf_font_dir .. "DejaVuSans.ttf]:mode=node;"})}
\directlua{luaotfload.add_fallback("hebfbsansbf",{"[" .. hllpdf_font_dir .. "DejaVuSans-Bold.ttf]:mode=node;"})}
\babelfont{sf}[Path={\hllpdffontdir},RawFeature={fallback=hebfbsans},
  BoldFont={Heebo-Bold.ttf},
  BoldFeatures={RawFeature={fallback=hebfbsansbf}}]{Heebo-Regular.ttf}
% v3.24: Scale=MatchLowercase — DejaVu Sans Mono has an outsized x-height, so
% unscaled inline code visually SHOUTS next to the delicate Frank Ruhl body
% (measured on a 197-page build: scaled, a code atom sits in the line at the
% body's weight, and the document tightened by ~5%). Applies to listings too.
\babelfont{tt}[Path={\hllpdffontdir},RawFeature={fallback=hebfb},Scale=MatchLowercase,
  BoldFont={DejaVuSansMono.ttf},
  BoldFeatures={RawFeature={fallback=hebfb},FakeBold=1.5},
  ItalicFont={DejaVuSansMono.ttf},
  ItalicFeatures={RawFeature={fallback=hebfb},FakeSlant=0.2},
  BoldItalicFont={DejaVuSansMono.ttf},
  BoldItalicFeatures={RawFeature={fallback=hebfb},FakeBold=1.5,FakeSlant=0.2}
  ]{DejaVuSansMono.ttf}

% Explicit multilingual/English profiles use real Latin faces and italics.
% The legacy/default Hebrew profile retains its established measured typography.
% Latin Modern is provided by the TeX distribution, not the pinned font bundle;
% recorder/runtime evidence reports it as an additional toolchain resource.
\ifx\lakitlanguageprofile\hllpdfhebrew\else
  \babelfont[english]{rm}[Ligatures=TeX,RawFeature={fallback=hebfb},
    BoldFeatures={RawFeature={fallback=hebfbbf}},
    BoldItalicFeatures={RawFeature={fallback=hebfbbf}}]{Latin Modern Roman}
  \babelfont[english]{sf}[Ligatures=TeX,RawFeature={fallback=hebfbsans},
    BoldFeatures={RawFeature={fallback=hebfbsansbf}},
    BoldItalicFeatures={RawFeature={fallback=hebfbsansbf}}]{Latin Modern Sans}
\fi

\usepackage{amsmath,amssymb,mathtools}
% amsthm supplies theorem/proof machinery. Environment names remain
% document-specific; the generic proof label is localized below.
\usepackage{amsthm}
% Proof names follow the surrounding prose, including return from an island.
\addto\extrashebrew{\renewcommand{\proofname}{הוכחה}}
\addto\extrasenglish{\renewcommand{\proofname}{Proof}}
\ifx\lakitmainlanguage\hllpdfenglish
  \renewcommand{\proofname}{Proof}
\else
  \renewcommand{\proofname}{הוכחה}
\fi
% v3.24.1: equation tags under babel-bidi. A tag re-emitted by mathtools'
% showonlyrefs inside align inherits the RTL base and its parens MIRROR —
% ")2.1(" instead of "(2.1)" (plain align/equation tags are unaffected;
% verified on a minimal case + a 198-page build). LTR-isolate at the tag-BOX
% level: mathtools' show/hide + numbering logic is untouched. Do NOT override
% \tagform@ instead — that clobbers the showonlyrefs mechanism (measured:
% every tag shows, numbering breaks).
% v3.28.1: the v3.24.1 wrap must apply ONLY to tags of display environments.
% Inline \eqref in Hebrew PROSE also goes through \maketag@@@, and there the
% english-wrap MIRRORS the parens — ")1(" — while the unwrapped original is
% correct (proven on a 2x2 isolation matrix: patch x showonlyrefs; regressed
% every prose \eqref since v3.24.1, caught by the render gate's R4). The two
% call sites are indistinguishable by \ifmmode (tag boxes are text mode), so
% the display environments themselves raise a flag; \eqref in prose sees it
% false and gets the original untouched \maketag@@@.
\usepackage{etoolbox}

% babel bidi=basic (including upstream 26.9) assumes every automatic
% discretionary has a one-glyph replacement. LuaTeX can legally produce one
% with replace=nil after a hyphenated Latin value; Babel then indexes nil and the
% whole build dies. Do the narrow upstream-compatible correction: while Babel
% classifies the list, present only that malformed disc as an ordinary
% discretionary, then restore its subtype. Babel still reorders the ENTIRE box;
% there is no pcall and no successful fallback to an unreordered node list.
% If Babel throws for any other reason, the build remains fatal. The synthetic
% nil-disc probe and the real table gallery in test_all.sh pin both layers.
\AtBeginDocument{\directlua{
  if Babel and Babel.bidi and not Babel.heb_nil_disc_compat then
    Babel.heb_nil_disc_compat = true
    local original = Babel.bidi
    local disc_id = node.id('disc')
    Babel.bidi = function(head, ispar, hdir)
      local hidden = {}
      for item in node.traverse_id(disc_id, head) do
        if item.subtype == 2 and (not item.replace or not item.replace.char) then
          table.insert(hidden, item)
          item.subtype = 0
        end
      end
      local result = original(head, ispar, hdir)
      for _, item in ipairs(hidden) do item.subtype = 2 end
      return result
    end
  end
}}
\makeatletter
\newif\ifHEBdisptag
\let\ORIGmaketag@@@\maketag@@@
\def\maketag@@@#1{\ifHEBdisptag\ORIGmaketag@@@{\foreignlanguage{english}{#1}}\else\ORIGmaketag@@@{#1}\fi}
\@for\HEB@env:=equation,equation*,align,align*,gather,gather*,multline,multline*,flalign,flalign*,alignat,alignat*,eqnarray,eqnarray*\do{%
  \expandafter\AtBeginEnvironment\expandafter{\HEB@env}{\HEBdisptagtrue}}
\makeatother
% English page + Hebrew prose: Babel's equation-only tag flip uses a primitive
% leqno path that shifts the formula to the left edge and the tag into the body
% of an English page. Align's full-width alignment is unaffected. Give only this
% display a physical LTR container and native amsmath eqno/leqno. The engine then
% owns real tag width, collision handling and labels; no proxy tag box hides ink
% from its fit calculation. Babel's saved \bbl@eqnodir still restores Hebrew text
% inside math, and the prose locale/math direction are never switched.
% Close the prior RTL paragraph before normalization. After the environment's
% group restores its directions, close TeX's empty resumed display paragraph and
% use native LaTeX continuation handling: same-paragraph prose has no indent,
% while an authored blank line starts an ordinary indented RTL paragraph.
% Register after Babel's setup; remain local to equation(*) in a main-English
% document's actual RTL prose. Native Hebrew, ordinary English and explicit
% fleqn retain their original paths. No page/display dimensions are assigned.
% <hllpdf:minority-rtl-equation-geometry>
\makeatletter
\newcommand{\HEB@minorityprosecontinuation}{\par\@doendpe}
\newcommand{\HEB@minorityequationgeometry}{%
  \ifx\lakitmainlanguage\hllpdfenglish
    \ifnum\pardirection=\@ne
      \if@fleqn\else
        \par\aftergroup\HEB@minorityprosecontinuation
        \iftagsleft@\let\veqno\@@leqno
        \else\let\veqno\@@eqno\fi
        \bodydir TLT\pardir TLT\textdir TLT\relax
      \fi
    \fi
  \fi}
\AtBeginDocument{%
  \AddToHook{env/equation/begin}{\HEB@minorityequationgeometry}%
  \AddToHook{env/equation*/begin}{\HEB@minorityequationgeometry}}
\makeatother
% </hllpdf:minority-rtl-equation-geometry>
% v3.63: \shortintertext rendered Hebrew MIRROR-REVERSED and flush LEFT — an entire
% commentary line inside an align printed backwards. babel's bidi=basic patches
% \intertext@ only (luababel.def), and mathtools' \shortintertext uses the identical
% \vbox{\normalbaselines … \noindent\ignorespaces#1\par} skeleton but is never
% patched, so its paragraph keeps pardir/textdir LTR: no bidi reordering (logical
% order on the page) and flush-left. Measured against \intertext with the same
% Applying babel's own replacement to \shortintertext@ makes the two identical.
% \h{}/\mbox fix the letter order but leave the line flush LEFT, so they are not
% substitutes. Guarded + \AtBeginDocument because mathtools \let-copies at load time.
% NB \ExplSyntaxOn is REQUIRED: mathtools' internals carry expl3 names, and without
% the expl3 catcodes the `_` is read as a subscript and the patch detonates far away
% with "! Illegal unit of measure (pt inserted)" at the \end{align} (measured).
\makeatletter
\AtBeginDocument{%
  \@ifundefined{shortintertext@}{}{%
    \ExplSyntaxOn
    \bbl@sreplace\shortintertext@{\normalbaselines}%
      {\normalbaselines\ifx\bbl@eqnodir\relax\else\bbl@pardir\@ne\bbl@eqnodir\fi}%
    \ExplSyntaxOff}}
\makeatother
% geometry inherits the canonical A4 sheet from \documentclass; build.sh then
% verifies every rendered MediaBox and CropBox, so a silent Letter fallback is
% impossible to publish as a clean default build.
\usepackage[\lakitgeometryoptions]{geometry}
\linespread{\lakitlinespread}       % layout hook; canonical default remains 1.3
\ifx\lakitparagraphstyle\hllpdfparagraphblock
  \setlength{\parindent}{0pt}
  \setlength{\parskip}{6pt plus 1pt minus 1pt}
\fi
% These dimensions remain ordinary scoped TeX lengths: authors can customize
% them after core input. Language switches inherit them; lists keep their own
% indentation and narrow reference layouts establish their local paragraph gaps.
\typeout{HLLPAR|\lakitparagraphstyle}
\usepackage{xcolor}
\usepackage{mdframed}
\usepackage{listings}                 % code listings (assembly / C / pseudo-code)
\usepackage{tikz}
\usetikzlibrary{arrows.meta,positioning,calc,shapes.geometric,fit,backgrounds}
\usepackage{pgfplots}\pgfplotsset{compat=1.18}
\usepackage{quantikz}                 % quantum circuits (remove if unused)
\usepackage{booktabs}
\usepackage{array}
\usepackage{float}                    % [H] = place table exactly here (reference doc: no drifting)
\usepackage{enumitem}
\usepackage{needspace}

% ============================================================================
% ★★ v3.4 — THE DIRECTION-ISOLATION TRIO (the heart of mixing Hebrew + Latin)
% Mental model:
%   \h{...}    Hebrew (RTL) island      — Hebrew text inside an LTR context
%                                          (TikZ nodes, the english `tikzpic`).
%   \en{...}   English (LTR) isolate    — Latin text / parentheses inside a
%                                          Hebrew (RTL) context. KEEPS PARENS
%                                          AND ORDER CORRECT.
%   \code{...} \en + monospace          — code inside Hebrew / inside a frametitle.
% CRITICAL: never wrap Latin in \h{} (it is for HEBREW). Inside a TikZ node,
% Latin in \h{} renders REVERSED (Datapath->htapataD, CPU->UPC) and parens
% mirror. Isolate every Latin/paren chunk with \en{} or \code{}. See
% references/figures-boxes-listings.md.
% ---- TEXT INSIDE MATH (v3.8) ----------------------------------------------
% Put ANY text that sits inside math in a \text{}. The ONLY thing that is
% invalid in math is a bare \h/\en (= \foreignlanguage) used DIRECTLY, outside
% \text -> "\rmfamily invalid in math mode". Inside \text:
%   Hebrew label           -> \text{...}            (bare Hebrew renders fine;
%                              \h is optional — BUT NOT ALWAYS HARMLESS (v3.67):
%                              if the \h group OPENS with a nested island
%                              ($..$, \en{}, \code{}, \texttt{}, \mbox{}) the
%                              WHOLE group loses bidi and the Hebrew prints in
%                              LOGICAL order = gibberish. Measured:
%                                \text{\h{$f$ קבועה}}  -> reads "העובק"  BROKEN
%                                \text{\h{קבועה $f$}}  -> correct
%                                \text{$f$ קבועה}      -> correct (no \h)
%                              Hebrew first inside \h, or drop the \h. Flagged
%                              H-OPENS-WITH-LTR-GROUP.)
%   plain Latin word        -> \text{...}            (no marker needed)
%   Latin WITH ( ) or [ ]   -> \text{\en{...}}       (\en forces LTR; a plain
%                              \text mirrors the parens in RTL -> "Spec (TNR(" )
%   f(Hebrew/text arg)      -> \text{f}(\text{...})  -- keep the parens at MATH
%                              level, NOT \text{f(arg)} (parens inside the box mirror)
%   code in math            -> \mathtt{...}  (or \texttt{}, math-guarded below)
% NEVER write \en{...} or \h{...} straight into $...$ / \[..\] (outside \text);
% always via \text{}. (\mbox{...} also works but does not scale in scripts.)
% ============================================================================
\newcommand{\h}[1]{\foreignlanguage{hebrew}{#1}}
% v3.42.1 — \en islands do NOT hyphenate (auto OR at explicit hyphens). English hyphenation patterns were
% tearing identifiers and terms ACROSS THE RTL SEAM ("IO Regis-/ters",
% "In-/put", "Band-/width", "Per-/formance", "Multi-/ple" — all shipped on
% real pages; probe-proven: "ALUS-/rcB", "Mem-/ToReg"). An island is a label
% or a term, not English prose — a hyphen tear there is never right. A/B on
% a full 197pp lakit: identical overfull sets, same page count, ~11
% paragraphs healed, one mild underfull (badness 3525) that replaced a
% mutilated "Regis-/ters". Long English PROSE belongs in the englishprose
% environment below, which restores explicit-hyphen breaks while its paragraph
% is still active. A raw otherlanguage block inherits the global prohibition.
% (\l@nohyphenation needs \csname:
% '@' is not a letter here.)
% v3.43.2 — \exhyphenpenalty moved OUT of \en to a GLOBAL assignment below.
% Render-proven: a group-local \exhyphenpenalty is DEAD CODE — TeX reads the
% line-breaking parameters at \par time, after the \en group has closed, so
% \en{write-through} / plain not-taken still tore at the "-" ("write-/
% through" shipped on a real exam page). \language is different: it is
% recorded in the node list as text is added, so the in-group nohyphenation
% assignment above DOES work (and stays). The global penalty was A/B-measured
% on a 197pp lakit: identical overfull set (25), identical page count, the
% not-/taken tear healed. Hebrew is unaffected in practice — Hebrew compounds
% use maqaf (־), not ASCII "-".
\newcommand{\en}[1]{\foreignlanguage{english}{\language=\csname l@nohyphenation\endcsname #1}}
% Babel runs extras locally for inline islands and paragraph language changes.
% An English main document must not inherit the Hebrew no-hyphen paragraph rule.
\addto\extrashebrew{\exhyphenpenalty=10000\relax}
\addto\extrasenglish{\exhyphenpenalty=50\relax}
\ifx\lakitmainlanguage\hllpdfenglish
  \exhyphenpenalty=50\relax
\else
  \exhyphenpenalty=10000\relax
\fi
% Full English paragraphs need ordinary explicit-hyphen breaks. Use this
% environment instead of a raw otherlanguage block; the paragraph ends while
% the local penalty is still active.
\newenvironment{englishprose}
  {\par\begin{otherlanguage}{english}\exhyphenpenalty=50\relax}
  {\par\end{otherlanguage}}
\newenvironment{hebrewprose}
  {\par\begin{otherlanguage}{hebrew}\exhyphenpenalty=10000\relax}
  {\par\end{otherlanguage}}
% Use for shell text whose locale must not depend on the last body paragraph.
\DeclareRobustCommand{\doclang}[1]{\foreignlanguage{\lakitmainlanguage}{#1}}

% ============================================================================
% ★★ — fixes for defects that are INVISIBLE in the log and in the checks
% and only surface on a VISUAL RENDER of the PDF. (Each verified empirically.)
% ----------------------------------------------------------------------------
% v3.22: \emph is now REAL slanted emphasis (FakeSlant italic face defined on
% the \babelfont{rm} line above), replacing the v3.21 \emph->\textbf hack: it
% preserves \emph's toggle semantics and keeps bold free for its own role.
% Prefer bold-as-emphasis for your doc? one line: \renewcommand{\emph}[1]{\textbf{#1}}
% page number: force the digits LTR so they never reverse (18 -> 81) in the RTL
%   footer. That reversal is a NON-DETERMINISTIC bidi glitch (observed once in
%   ~80 pp); this LuaLaTeX build is not byte-reproducible, so the defensive
%   LTR wrap is the only reliable guard. Latin digits are always LTR anyway.
\renewcommand{\thepage}{\foreignlanguage{english}{\arabic{page}}}
% v3.70: \pagenumbering \gdef's \thepage and silently DESTROYS the wrap above —
% after \pagenumbering{arabic} the folio guard is gone for the whole rest of the
% document. Two collections discovered this independently and hand-restore the
% wrap in main.tex with the same comment; that knowledge now lives here instead.
% The wrapper re-applies the LTR wrap after the kernel call, for any style.
\makeatletter
\let\lakin@orig@pagenumbering\pagenumbering
\renewcommand{\pagenumbering}[1]{%
  \lakin@orig@pagenumbering{#1}%
  \renewcommand{\thepage}{\foreignlanguage{english}{\csname @#1\endcsname\c@page}}}
\makeatother
% emergencystretch: Hebrew does not hyphenate, so an unbreakable LTR island
%   (\en{...} / an inline math run) landing at the margin can leave a small
%   (~2-13pt) overfull with no legal breakpoint. This rescues ONLY those
%   inline-prose tie-breaks; it is a ceiling used only on offending lines and
%   does NOT touch display-math / listing / table overflow (those still surface
%   in the log for the content-level overflow pass). Verified: baseline overfull
%   -> 0 with this set, with no change to pagination or underfull.
\setlength{\emergencystretch}{3em}
% v3.40: inline math (almost) never breaks across lines. TeX's default
%   (binop 700 / rel 500) lets an inline run split at an operator —
%   $(1.f)\cdot 2^{e}$ shipped as "(1.f)·" ending one line and "2^e" opening
%   the next (reader-reported). In LTR that split is merely ugly; under an
%   RTL base the carried operator lands on the visually WRONG side of the
%   break and the formula reads backwards.
%   WHY 10000 (v3.45.0 — the display-first resolution). Inline math is
%   FORBIDDEN from breaking, so a torn island under an RTL base is
%   structurally impossible. The v3.40 cycle used 9999 ("break only as a last
%   resort") because a bare 10000 left ~30-39 WIDE inline equations
%   overflowing (worst 90.7pt) — but 9999's "last resort" fires OFTEN in
%   tight, non-hyphenating Hebrew: a reader found torn islands across the
%   fleet (~150 candidates; e.g. "P(y |" ending one line and "x_0)" opening
%   the next, read backwards). The step the 9999 cycle missed: every one of
%   those ~30 overflows is a wide equation that BELONGS in display. Convert
%   them to display FIRST, then 10000 gives 0 tears AND 0 overflows (measured,
%   12/12 fleet). Doctrine: inline math is ATOMIC; a wide equation goes in a
%   display, never inline. See content-style ("inline math is atomic").
\binoppenalty=10000 \relpenalty=10000
% v3.38: uniform sentence spacing. TeX's default (\nonfrenchspacing, sfcode
% 3000) is the OLD Anglo-American wide-space-after-period convention — Hebrew
% typography never had it, and modern Latin practice (Chicago et al.) dropped
% it too. Measured on real pages: post-period gaps up to 1.9-2.1x the line's
% mean word gap — the dominant "rivers" source in justified Hebrew prose.
% Survives babel language switches (probed: sfcode stays french inside \en{}
% islands and otherlanguage blocks).
\frenchspacing
% v3.38: widow/orphan protection. LaTeX's default penalties (150/150/50)
% allow a lone line at a page top/bottom almost freely. 10000 = forbid;
% this is SAFE here precisely because the document is \raggedbottom (the
% demoted line just moves on, no vertical stretch fight) — under flushbottom
% these values would cause underfull pages. raggedbottom is the report-class
% default; made explicit so a future class/option change cannot silently
% flip it and turn the 10000s harmful.
\raggedbottom
\widowpenalty=10000 \clubpenalty=10000 \displaywidowpenalty=10000
% em-dash: --- does NOT become — under Frank Ruhl (the font lacks tlig; verified
%   the width is identical with/without Ligatures=TeX, while the mechanism works
%   on Latin Modern). TYPE THE CHARACTER — (U+2014) DIRECTLY; likewise – (U+2013).
%   check_bidi_figures.py flags a --/--- that is adjacent to Hebrew.
% ============================================================================

% ---- STRUCTURAL DIRECTION-SAFETY for \texttt (v3.7) -------------------------
% A parenthesis, bracket, or $ INSIDE \texttt{} mirrors/migrates in RTL prose
% (MIPS 0($s0)->"(s0$)0", $rs->"rs$", arr[i]->"[arr[i") because \texttt sets the
% font but NOT the direction. Redefining \texttt to LTR-isolate its content fixes
% the WHOLE class at once (parens, brackets, $) so plain \texttt{lw $t0,0($s0)}
% is safe — no need to remember \code{}. Guarded by \ifmmode because
% \foreignlanguage is invalid in math (there, plain monospace is already LTR).
% NB (corrected v3.56 — the old note here said "stays tofu", which is WRONG):
% \texttt{Hebrew} is a FATAL build failure. The tt family is DejaVu Sans Mono
% (zero Hebrew) and the tt fallback hebfb is DejaVuSerif (also zero Hebrew), so
% the subsetter gets an empty glyph set and LuaTeX dies with
%   "! error: (file .../DejaVuSansMono.ttf) (type 2): there are no glyphs in the
%    subset"  -> no output PDF.
% The message names a FONT FILE, never Hebrew/\texttt/the source line. And it is
% non-local: add an unrelated Latin \code{} elsewhere and the fatal degrades into
% a SILENT DELETION of the Hebrew word. Put Hebrew in \h{}/normal text, never in
% \texttt{}/\code{}; flagged HEBREW-IN-MONO. To disable, delete this block.
\let\origtexttt\texttt
\protected\long\def\texttt#1{\ifmmode\origtexttt{#1}\else\foreignlanguage{english}{\origtexttt{#1}}\fi}
% v4.8.0: defer the direction test until TeX has inserted an alignment cell's
% math template. Unprotected expansion here tested \ifmmode too early at the
% FIRST token of aligned/array/cases/matrix cells. Engine-level \protected also
% preserves the original \texttt implementation without the self-recursion
% observed when retrofitting \DeclareRobustCommand to its saved alias.
\protected\long\def\code#1{\ifmmode\origtexttt{#1}\else\foreignlanguage{english}{\origtexttt{#1}}\fi}

% ---- \codeline — display-code line (v3.24): the code analogue of display math.
% Reader-friction rule (content-style.md §5b): inline \code{} is for a SINGLE
% ATOM only (one register/variable/mnemonic, a tight no-space expression).
% A multi-token expression — an instruction with operands, a spaced assignment,
% a register-binding list — zigzags the eye RTL<->LTR mid-sentence and, being
% an unbreakable island, line-breaks badly. Give it its own centered line with
% display-math vertical rhythm. \centerline = one box => the expression never
% wraps; \nopagebreak keeps it glued to its introducing sentence.
%   ...the Hebrew sentence ends with a colon:
%   \codeline{lw \$t1, -4(\$t0)}
%   and the prose continues (escape dollars as \$).
% v3.66 (first line of the body): flush a PENDING \item label before the \par.
% LaTeX defers the label box (\@labels) until the next \everypar, and an \hbox in
% vertical mode never fires \everypar — so `\item \codeline{...}` printed the code
% at y=69.1 and the item's own "1." BELOW it at y=103.6 (the number appears after
% its content, attached to nothing). \if@inlabel\leavevmode\fi is CONDITIONAL on
% purpose: an unconditional \leavevmode would open an empty paragraph line when
% \codeline follows ordinary prose. Verified both ways — in a list the label now
% sits on the code's line at its normal x, and after prose the geometry is
% byte-identical to before.
\makeatletter
\newcommand{\codeline}[1]{%
  \if@inlabel\leavevmode\fi
  \par\nopagebreak
  \vspace{\abovedisplayskip}\nopagebreak
  % v3.62: was \centerline{...}, i.e. \hbox to\hsize{\hss #1 \hss}. \hss has
  % INFINITE SHRINK, so TeX could always "fit" the box and NEVER reported an
  % Overfull — the same \hss mechanism as BLD-27's \item label. Measured: a wide
  % \codeline in a defbox painted its first token ON the box's left border bar and
  % sat 9.5pt outside the frame interior on both sides, with 0 Overfull, 0 Underfull
  % and all four gates PASS; in bare prose it reached 21.0pt into BOTH margins.
  % \hfil has no shrink, so the overflow is reported and build.sh's ovf_vis/ovf_big
  % counters see it; \linewidth (not \hsize) makes it box-aware, so a \codeline
  % inside a colored box centres on the box interior. Verified: a FITTING \codeline
  % keeps byte-identical glyph positions, and the box-crossing case now emits
  % "Overfull \hbox (19.03pt too wide)". A wide \codeline should be split in two.
  \hbox to\linewidth{\hfil\code{#1}\hfil}%
  \nopagebreak\vspace{\belowdisplayskip}}
\makeatother

% ---- a thin decorative rule helper ----
\newcommand{\sep}{\par\vspace{2pt}\noindent\textcolor{cRule}{\rule{\linewidth}{0.6pt}}\par\vspace{4pt}}

% ---- colors ----
\definecolor{cBlue}{HTML}{1F4E79}   \definecolor{cBlueBg}{HTML}{EAF1F8}
\definecolor{cGreen}{HTML}{2E6B4F}  \definecolor{cGreenBg}{HTML}{EAF4EE}
\definecolor{cAmber}{HTML}{9C6A00}  \definecolor{cAmberBg}{HTML}{FBF3E2}
\definecolor{cPurple}{HTML}{5B2C6F} \definecolor{cPurpleBg}{HTML}{F3ECF7}
\definecolor{cTeal}{HTML}{0E6E6E}   \definecolor{cTealBg}{HTML}{E6F4F4}
\definecolor{cRed}{HTML}{8E2B2B}    \definecolor{cRedBg}{HTML}{FBEAEA}
\definecolor{cGray}{HTML}{555555}   \definecolor{cRule}{HTML}{1F4E79}

% ---- math operators ----
\DeclareMathOperator{\Tr}{Tr}
\DeclareMathOperator{\rank}{rank}
\DeclareMathOperator{\diag}{diag}
\DeclareMathOperator{\sgn}{sgn}

% ---- bra-ket (quantikz defines \ket/\bra/\braket; add the rest safely) ----
\providecommand{\ket}[1]{\left|#1\right\rangle}
\providecommand{\bra}[1]{\left\langle#1\right|}
\providecommand{\braket}[2]{\left\langle#1\middle|#2\right\rangle}
\newcommand{\dyad}[1]{\left|#1\right\rangle\!\left\langle#1\right|}
% v3.21: general outer product |a><b| and expectation value <A> — agents reach
% for these by reflex (esp. \expval); quantikz defines neither, so their absence
% threw "Undefined control sequence" in a real build. \providecommand = safe.
\providecommand{\ketbra}[2]{\left|#1\right\rangle\!\left\langle#2\right|}
\providecommand{\expval}[1]{\left\langle#1\right\rangle}
\newcommand{\abs}[1]{\left|#1\right|}
\newcommand{\norm}[1]{\left\lVert#1\right\rVert}
\newcommand{\half}{\tfrac12}
\newcommand{\R}{\mathbb{R}}\newcommand{\C}{\mathbb{C}}\newcommand{\Z}{\mathbb{Z}}
% identity operator. NB: \mathbb{1} silently renders a WRONG glyph (msbm has no
% blackboard digits). Use \mathbb{I} = 𝕀 (works with amssymb, no extra package).
% Do NOT "fix" this back to \mathbb{1}.
\newcommand{\id}{\mathbb{I}}

% ---- physical units (v3.35): \punit — upright unit with thin space ---------
% Usage: $300\punit{K}$, $80\punit{mV/decade}$, $10^{7}\punit{cm/s}$ — the
% NUMBER AND UNIT LIVE IN ONE MATH ISLAND. Promoted verbatim from a 199-page
% production lakit (advanced-electronic-devices, dozens of uses).
% Two measured traps this design avoids (both render-proven):
%   1. siunitx SCRAMBLES under RTL — \SI{3.5}{eV} renders "eV 3.5", exponents
%      detach ("Hz 10 x 10^9"): its quantity assembly is not bidi-isolated.
%      Do NOT load siunitx in this pipeline.
%   2. A text-mode unit next to a bare-prose number is TWO adjacent LTR runs
%      and reorders ("5 mm" -> "mm5" — the BIDI-2/3 class). \punit is therefore
%      math-only ON PURPOSE: prose misuse fails LOUDLY at compile (\mathrm
%      outside math) instead of scrambling silently.
% Domain shorthands (\MeV, \fm...) stay per-document: \newcommand{\MeV}{\,\text{\en{MeV}}}
\providecommand{\punit}[1]{\,\mathrm{#1}}

% ============================================================================
% Colored boxes — mdframed (NOT tcolorbox, which breaks under bidi).
% ★★ v3.4: innerleftmargin/innerrightmargin give code listings room so their
% frame/line-numbers do not spill out the box border (see listing styles below).
% ============================================================================
% v3.61: everyline/split skips were applied to `tablebox` only (BLD-14, v3.37) —
% the six colored boxes below never got them, so they still split OPEN across a page
% break: the page-N segment gets top+sides and NO bottom rule, the page-N+1 segment
% sides+bottom and NO top rule. Measured on the shipped fleet this is a COVERAGE GAP,
% not a re-report of BLD-14: 0 split-open `tablebox` segments (that fix holds) vs 62
% split-open COLORED boxes across 7 documents. Verified after the fix: both segments
% close, and the frametitle is NOT duplicated on the continuation. BLD-14's caveat
% still applies — a box taller than a full page grinds mdframed's split-trial loop, so
% split such content manually.
% ============================================================================
% v3.72: frametitle={} leads EVERY \newmdenv option list below. mdframed options
% are GROUP-scoped: a box nested inside \begin{exbox}[frametitle={...}] used to
% INHERIT the outer title — measured, the identical title string rendered twice,
% the second riding a title-bar fill inside the INNER box's frame, for all seven
% environments in both directions. Both natural forms were broken (= preamble
% bug, not an authoring rule). The explicit reset makes every box titleless
% unless ITS OWN [frametitle={...}] says otherwise; verified both directions:
% no duplicate on the nested box, an explicit title still renders.
% Geometry is shared with the optional semantic-width table container. Themes
% select their own existing semantic style below; neither path fixes columns.
\mdfdefinestyle{hllpdfsemanticgeometry}{linewidth=1.1pt,
  innertopmargin=6pt,innerbottommargin=7pt,skipabove=9pt,skipbelow=9pt,
  innerleftmargin=9pt,innerrightmargin=9pt,everyline=true,
  splitbottomskip=6pt,splittopskip=12pt,roundcorner=3pt}
\newcommand{\hllpdfsemanticframestyle}{hllpdfsemanticgeometry}
\newmdenv[frametitle={},style=hllpdfsemanticgeometry,linecolor=cBlue,backgroundcolor=cBlueBg,
  frametitlebackgroundcolor=cBlue,frametitlefontcolor=white,
  ]{defbox}   % הגדרה
\newmdenv[frametitle={},style=hllpdfsemanticgeometry,linecolor=cGreen,backgroundcolor=cGreenBg,
  frametitlebackgroundcolor=cGreen,frametitlefontcolor=white,
  ]{thmbox}   % משפט/טענה
\newmdenv[frametitle={},style=hllpdfsemanticgeometry,linecolor=cAmber,backgroundcolor=cAmberBg,
  frametitlebackgroundcolor=cAmber,frametitlefontcolor=white,
  ]{notebox}  % הערה
\newmdenv[frametitle={},style=hllpdfsemanticgeometry,linecolor=cPurple,backgroundcolor=cPurpleBg,
  frametitlebackgroundcolor=cPurple,frametitlefontcolor=white,
  ]{exbox}    % דוגמה
\newmdenv[frametitle={},style=hllpdfsemanticgeometry,linecolor=cRed,backgroundcolor=cRedBg,
  frametitlebackgroundcolor=cRed,frametitlefontcolor=white,
  ]{warnbox}  % טעות נפוצה
\newmdenv[frametitle={},style=hllpdfsemanticgeometry,linecolor=cTeal,backgroundcolor=cTealBg,
  frametitlebackgroundcolor=cTeal,frametitlefontcolor=white,
  ]{keybox}   % רעיון מרכזי

% ---- codebox (v4.5.0): the dedicated container for code -------------------
% An OPTION, not a policy: asmblock unboxed stays canonical for a full listing,
% and a snippet inside any of the six colored boxes still works as before.
% What makes codebox worth its own environment is that it is the ONE framed box
% that keeps its margin line numbers.
%
% v4.4.0 shipped it WITHOUT them, on the strength of v3.37's "no margin knob can
% save them" plus four fresh probes that agreed. That was wrong: every probe had
% moved the code, the box, or both — never the key that carries the NUMBER
% COLUMN. Measured on the page:
%   resetmargins=false   the listing follows THIS box instead of the page,
%                        moving the code from x=75.26 to x=97.54
%   xrightmargin=26pt    under RTL this is the PHYSICAL LEFT indent, and the
%                        only key that carries the \llap'ed digits with it:
%                        x=64.28 -> 90.18, clearing the rule at 65.64 by 23.6pt
%   xleftmargin=26pt     moves the code and LEAVES THE DIGITS on the rule —
%                        the asymmetry that hid the fix for eleven versions
% resetmargins is hoisted out of the four styles (see the \lstset below them) so
% this override is possible at all; every other listing still resolves to true.
% Pinned by tests/codebox_numbers.tex + probe_lstnum.py --inside-rule, whose
% negative control strips xrightmargin and watches the digits fall back.
\definecolor{cSlate}{HTML}{3F4A56}  \definecolor{cSlateBg}{HTML}{F2F4F7}
\newmdenv[frametitle={},leftline=true,rightline=true,topline=true,bottomline=true,
  linecolor=cSlate,backgroundcolor=cSlateBg,linewidth=0.9pt,
  frametitlebackgroundcolor=cSlate,frametitlefontcolor=white,
  innertopmargin=6pt,innerbottommargin=7pt,skipabove=9pt,skipbelow=9pt,
  innerleftmargin=9pt,innerrightmargin=9pt,everyline=true,splitbottomskip=6pt,splittopskip=12pt,
  roundcorner=3pt]{codebox}   % קוד
% The two keys that make in-box numbering work, both measured on the page:
%   resetmargins=false  the listing follows THIS box instead of the page
%   xrightmargin        under RTL this is the PHYSICAL LEFT indent, and it is
%                       the only one that carries the \llap'ed number column
%                       with it. xleftmargin moves the code and leaves the
%                       digits behind on the rule (measured: code 97.54,
%                       digits 64.28 — the v3.37 defect, reproduced); with
%                       xrightmargin=26pt the digits land at x=90.18, a
%                       clearance of 23.6pt from the rule at 65.64.
% This asymmetry is why v3.4's note had the semantics backwards and v3.37
% concluded the digits could not be saved.
% <hllpdf:listing-container-state>
% v4.29: all state is local to the current environment group. A semantic
% callout masks number text, while a nested codebox restores the formatter
% saved before that mask. Save it once across successive semantic wrappers;
% an explicit numbers=none (or an explicitly empty formatter) remains valid.
\makeatletter
\newif\ifhllpdf@incodebox
\newif\ifhllpdf@englishcodebox
\newif\ifhllpdf@lstnumbermasked
\let\hllpdf@savedlstnumber\thelstnumber
\makeatother
% </hllpdf:listing-container-state>
% <hllpdf:codebox-gutter>
\makeatletter
\AtBeginEnvironment{codebox}{%
  \hllpdf@incodeboxtrue
  % Snapshot the OUTER language before asmblock selects English.
  \iflanguage{english}{\hllpdf@englishcodeboxtrue}{\hllpdf@englishcodeboxfalse}%
  \ifhllpdf@lstnumbermasked
    \let\thelstnumber\hllpdf@savedlstnumber
    \hllpdf@lstnumbermaskedfalse
  \fi
}
% listings' AddToHook is GLOBAL. Register once, and gate on local box state;
% adding an unconditional hook at every codebox entry leaks into later listings.
\lst@AddToHook{PreInit}{%
  \ifhllpdf@incodebox
    % Apply every listing override here, inside the LISTING's own group.
    % Setting margins at box entry leaks them into nested semantic callouts.
    \lstset{resetmargins=false,xrightmargin=26pt}%
    % Styles resolve before PreInit. In LTR the physical number gutter follows
    % xleftmargin; in Hebrew retain the measured inherited xrightmargin path.
    \ifhllpdf@englishcodebox
      \lstset{xleftmargin=26pt,xrightmargin=0pt}%
    \fi
    % ONE box, not two: suppress only the frame/background/skips of a listing
    % inside this codebox. The next unboxed listing keeps its own style.
    \lstset{frame=none,backgroundcolor=\color{cSlateBg},aboveskip=0pt,belowskip=0pt}%
  \fi
}
\makeatother
% </hllpdf:codebox-gutter>

% v3.37: line numbers inside a COLORED BOX are sliced by the box border —
% measured: digit span x=63.5-67.3 vs border line at x=65.2 (the border runs
% THROUGH the digits), and the inter-border gap (3.5pt) is narrower than a
% digit (3.8pt), so no margin knob can save them (under resetmargins,
% xleftmargin is inert in-box — measured: code at x=75.26 with 2.2em AND
% 3.0em). A 3-6 line boxed snippet does not need line numbers; blank the
% number TEXT inside every colored box (\thelstnumber is what gets printed;
% the style's numbers=left survives, so nothing else shifts — code position
% byte-frozen at x=75.26, verified). Unboxed listings keep their numbers.
% v3.78: the blanking is coupled to the FRAME, not to the box name —
% clean_theme.tex registers LATER hooks (registration order wins) that
% restore the number text inside the boxes it renders frameless.
% v4.5.0 RETRACTS the "no margin knob can save them" clause above. It holds only
% under resetmargins=true (the shipped default): there xleftmargin is inert
% in-box, as measured. With resetmargins=false, xrightmargin DOES carry the
% \llap'ed digit column clear of the rule — measured; the codebox \lstset that
% does it is the one above, in the codebox block. The loop below therefore
% covers SIX boxes, not "every colored box": codebox opts out and keeps its
% numbers.
% <hllpdf:semantic-listing-mask>
\makeatletter
\@for\HEB@lstbox:=defbox,thmbox,notebox,exbox,warnbox,keybox\do{%
  \expandafter\AtBeginEnvironment\expandafter{\HEB@lstbox}{%
    \ifhllpdf@lstnumbermasked\else
      \let\hllpdf@savedlstnumber\thelstnumber
    \fi
    \hllpdf@lstnumbermaskedtrue
    \def\thelstnumber{}%
    % The nearest semantic container owns its listing presentation. On exit,
    % TeX's group restores the enclosing codebox state and formatter.
    \hllpdf@incodeboxfalse}}
\makeatother
% </hllpdf:semantic-listing-mask>

% ============================================================================
% Unified TABLE presentation: subtle uniform slate frame + auto-numbered
% caption rendered INSIDE the frame. Use:
%   \begin{tablebox}\centering  <tabular>  \tabcap{<description>}\end{tablebox}
% Replaces the buggy \caption* (which emitted a stray "טבלה N : *" line) and
% provides an explicit framed option + numbered caption "טבלה N.M: תיאור".
% Optional intro text may be placed at the TOP, before the tabular, inside the box.
% ============================================================================
\definecolor{cTblLine}{HTML}{B9C4D2}  \definecolor{cTblBg}{HTML}{FAFBFC}
% v3.23: class-aware — under `article` there is no chapter counter (the promised
% report->article switch used to die with "No counter 'chapter' defined").
\makeatletter
\@ifundefined{c@chapter}{%
  % v3.31: in article the counter resets per SECTION, so an unprefixed number
  % repeated ("טבלה 1." in every section — measured on a two-section render).
  % Same conditional idea as the report branch: prefix the section number once
  % a real \section exists; a pre-section table stays "טבלה 1".
  \newcounter{tbl}[section]
  \renewcommand{\thetbl}{\ifnum\value{section}>0 \thesection.\fi\arabic{tbl}}%
}{%
  \newcounter{tbl}[chapter]
  % v3.30: a report-class doc that uses SECTIONS ONLY (chapter stays 0 —
  % the common short-doc pattern) used to print "טבלה 0.1"; the chapter
  % prefix now appears only once a real \chapter exists (render-verified
  % both ways: "טבלה 1." chapterless, "טבלה 1.1" with chapters).
  \renewcommand{\thetbl}{\ifnum\value{chapter}>0 \thechapter.\fi\arabic{tbl}}%
  % v3.64: the SAME leak in \thesection itself, left behind by v3.30. That fix
  % guarded the CAPTION counter and stopped there, so a chapterless report still
  % printed the heading "0.1 סעיף ראשון" and \ref to it printed "0.2" — measured:
  %   p1 y=65.2  '…0.1'   p1 y=144.9 '…0.2'   p1 y=183.4 '…0.2.1'
  % report.cls hardcodes \thesection as \thechapter.\@arabic\c@section, and the
  % leading 0 is the chapter counter that no \chapter ever incremented. Fixing
  % \thesection fixes the whole chain — \thesubsection and \thesubsubsection are
  % defined in terms of it. \theequation and \thetbl/\thefig were already guarded,
  % which is exactly why this survived: every NUMBER around the heading looked
  % right, so the heading's own 0 read as a style choice rather than a bug.
  % \appendix is safe: report.cls resets chapter to 0 and redefines \thechapter to
  % \@Alph\c@chapter, so an appendix \chapter makes the counter 1 and the guard
  % emits "A.1" unchanged. Render-verified in all four modes (chapterless,
  % chaptered, appendix+chapter, appendix+section-only).
  \renewcommand{\thesection}{\ifnum\value{chapter}>0 \thechapter.\fi\arabic{section}}%
}
\makeatother
% v3.37: everyline=true — if the box DOES split across pages, every segment
% gets a CLOSED frame (measured before: the split ran open-ended into the page
% edge with no bottom border, and the caption could strand alone at the next
% page top — "טבלה A.2" shipped that way). The split skips give the segments
% breathing room. In practice, with the caption GLUE below, every previously-
% split table in the corpus now moves WHOLE instead (the white-gap tradeoff,
% documented benign); everyline stays as the safety net for a genuinely
% page-tall table. NB: a table taller than ONE FULL PAGE can grind mdframed's
% split-trial loop (reproduced: minimal doc, minutes of 1pt-step Overfull
% probes) — break such a monster into two tables.
\mdfdefinestyle{hllpdftablegeometry}{leftline=true,rightline=true,topline=true,bottomline=true,
  linewidth=0.9pt,
  innertopmargin=9pt,innerbottommargin=8pt,skipabove=11pt,skipbelow=11pt,
  innerleftmargin=11pt,innerrightmargin=11pt,roundcorner=4pt}
\newcommand{\hllpdftableframestyle}{hllpdftablegeometry}
\newmdenv[frametitle={},style=hllpdftablegeometry,
  linecolor=cTblLine,backgroundcolor=cTblBg,
  % v4.2.0 (adopted): tablebox is ATOMIC. Short tableboxes near page bottoms
  % repeatedly entered mdframed's split trial, emitting benign Overfull-vbox
  % bursts and occasionally separating the intro prose from the tabular.
  % nobreak=true moves the complete short table to the next page instead. The
  % output-active Overfull-vbox gate still catches a box taller than a page;
  % everyline/split skips remain as defensive styling should that path ever
  % re-arm, but the canonical path never publishes a split tablebox.
  nobreak=true,
  everyline=true,splitbottomskip=6pt,splittopskip=12pt,
  % v3.79 (BLD-37 table side, correction): startinnercode/endinnercode run
  % INSIDE mdframed's body vbox (mdframed.sty:848 / :856+863) -- the second one
  % after \par\unskip, i.e. after line breaking, which is the only moment the
  % \tabcap anchor's true offset from the top of the box EXISTS. See
  % \lakintabboxclose below. These two keys are the ONLY tablebox change.
  startinnercode={\lakintabboxopen},endinnercode={\lakintabboxclose}]{tablebox}
% v3.37: \nopagebreak GLUE — \vsplit honors penalties, so the caption can no
% longer strand alone at the top of the next page (it takes its rows along, or
% the whole box moves). Render-verified on the real shipped orphan.
% v3.59: \NR@gettitle is what nameref injects into the stdlib title-bearing
% commands (\@sect, \@chapter, \caption) to fill \@currentlabelname. \figcap and
% \tabcap are INDEPENDENT macros — they never called it, so \@currentlabelname kept
% whatever the last heading left and \nameref{fig:x} silently printed the enclosing
% SECTION's title instead of the caption (proven in the .aux: field 3 of \newlabel
% was the heading text while field 1 and the anchor `fig.1.1` were correct — so this
% is a THIRD, distinct head next to BLD-13's counter fix and BIDI-38's autoref-name
% fix, neither of which writes that field). \@ifundefined keeps it inert if
% hyperref/nameref is ever dropped. Guarded helper because \figcap is defined
% after \makeatother.
% NB both helpers below need \makeatletter — this part of the preamble is OUTSIDE
% any letter-catcode region, so an unwrapped \@currentlabel parses as LaTeX's \@
% and dies with "You can't use \spacefactor in vertical mode" (measured).
\makeatletter
\newcommand{\lakinrtitle}[1]{\@ifundefined{NR@gettitle}{}{\NR@gettitle{#1}}}
% v3.59: \hitem{<letter>}{<key>} — a REFERENCEABLE Hebrew-lettered list item.
%   \begin{enumerate} \hitem{א}{it:a} הסעיף הראשון. \end{enumerate}  ->  \ref{it:a} = (א)
% Why this exists: automatic Hebrew list labels fail silently, so SKILL.md mandates
% an EXPLICIT \item[(\textbf{א})] — but LaTeX puts \refstepcounter INSIDE
% \if@noitemarg (latex.ltx:12964), and \item[...] never sets that flag. So an
% explicit item steps no counter and never touches \@currentlabel: a \label there
% binds to the AMBIENT one and \ref prints the CHAPTER/SECTION number where the
% item letter belongs (proven in the .aux: {chapter.1}). Same doctrine as \figcap —
% own the label state before \label. \phantomsection matters: without it the
% hyperlink still lands on the chapter head (measured); with it the anchor is
% section*.N and the bookmark outline stays clean.
\newcommand{\hitem}[2]{\item[(\textbf{#1})]%
  \def\@currentlabel{(#1)}\def\@currentlabelname{#1}\phantomsection\label{#2}}
% NB (v3.66): \hitem is \ref-only. \autoref{key} calls it a SECTION ("סעיף (א)")
% because \phantomsection brands the anchor section*.N, and \nameref{key} prints
% the bare letter. For \autoref-able items prefer \hebenumlabels below.
%
% v3.66: \hebenumlabels — OPT-IN automatic Hebrew enumerate letters, done right.
%   {\hebenumlabels\begin{enumerate}\item\label{q:a} ...\end{enumerate}}
%   -> labels (א) (ב) (ג); \ref{q:a} = א; anchor Item.N (real, per-item).
% This retires the belief that made \item[(\textbf{א})] + \hitem necessary.
% "Automatic Hebrew labels fail silently" was TRUE but unexplained, and the
% mechanism is now bisected (plain article vs +enumitem, measured): ENUMITEM
% PRE-EXPANDS the label expression in an \edef BEFORE \refstepcounter, with the
% counter still at 0. Tokens it shields by name (\arabic,\alph,\roman) survive;
% a plain \newcommand \ifcase wrapper evaluates at 0 -> the EMPTY zero-branch ->
% the label prints as "()" with the letter GONE (and \the\c@enumi prints a
% literal 0, \number\value dies fatally pointing at \item). The cure is one
% property, not a doctrine: the wrapper must be ROBUST (survives the \edef
% unexpanded), while \theenumi stays EXPANDABLE (so \refstepcounter writes the
% real letter to the .aux). Verified: labels (א)(ב)(ג), \ref -> א/ב/ג each item
% its own value, .aux = {א}{...}{Item.1} — versus the explicit-\item doctrine
% where BIDI-50/\hitem exist purely to patch the counter that never stepped.
\def\@hebalph#1{\ifcase#1\or א\or ב\or ג\or ד\or ה\or ו\or ז\or ח\or ט\or י\or
  יא\or יב\or יג\or יד\or טו\or טז\or יז\or יח\or יט\or כ\else\@ctrerr\fi}
\DeclareRobustCommand*\hebalph[1]{\expandafter\@hebalph\csname c@#1\endcsname}
\newcommand{\hebenumlabels}{%
  \renewcommand{\theenumi}{\@hebalph\c@enumi}%
  \renewcommand{\labelenumi}{(\textbf{\hebalph{enumi}})}}
\makeatother
% v3.63: \mbox the LABEL. \thetbl ("7.1") and a caption description that STARTS
% WITH A DIGIT ("32 הרגיסטרים…") are two European-Number runs separated only by
% neutrals (the label's own "." and \enspace). Under the RTL paragraph base the UBA
% resolves those neutrals to the RTL level, so both numbers sit in ONE reversed span
% and EXCHANGE VISUAL POSITIONS. Measured on the shipped PDF: a table whose real
% number is 7.1 printed "טבלה 32" — a different, EXISTING table number — while \ref
% correctly printed 7.1, so the caption silently contradicted its own cross-reference,
% and the description lost its leading count. The source contains no adjacency at all;
% it is manufactured by this macro, which is why every source linter was blind.
% \mbox makes the label one atomic box so nothing can reorder across it. Verified
% zero-regression: an ordinary Hebrew caption keeps identical x-geometry.
% One paragraph style covers label, number, separator and description.
% Explicit project/template overrides remain possible; the label hook is kept
% for compatibility, with an identity default. Counters/anchors stay unchanged.
\newcommand{\lakitcaptionstyle}{\normalfont\small\color{black}\hypersetup{allcolors=black}}
\newcommand{\lakitcaptionlabelstyle}[1]{#1}
\newcommand{\lakitcaptionpunctuation}{:}
% Native floats, captionof and longtable use the same presentation policy.
% subcaption is optional and inherits the same font rather than a second size.
\usepackage{caption}
\DeclareCaptionFont{lakit}{\lakitcaptionstyle}
\DeclareCaptionLabelFormat{lakit}{\lakitcaptionlabelstyle{#1~\en{#2}}}
\DeclareCaptionLabelSeparator{lakit}{\lakitcaptionpunctuation\enspace}
\captionsetup{font=lakit,labelfont={},textfont={},labelformat=lakit,
  labelsep=lakit,justification=centering}
\makeatletter
\AtBeginDocument{\@ifpackageloaded{subcaption}{%
  \captionsetup[sub]{font=lakit,labelfont={},textfont={},labelsep=lakit}}{}}
\makeatother
% v3.64: OPT-IN list-of-figures/tables entry via an optional short caption.
% \listoffigures and \listoftables are inherited from report.cls and LOOK available,
% but \figcap/\tabcap step the custom \c@fig/\c@tbl counters and never call
% \addcontentsline — so both printed their heading and NOTHING ELSE. Measured: a doc
% with one \figcap and one \tabcap rendered page 1 as 'רשימת האיורים\n1\n', and the
% .lof held only \@chapter's \addvspace, zero \contentsline entries. An empty list
% under a correct Hebrew heading reads as "this document has no figures".
%
% WHY OPTIONAL, AND NOT AUTOMATIC. Writing #1 unconditionally was implemented first
% and REJECTED: a caption containing \codeline compiles CLEANLY today (verified on
% the pre-v3.64 preamble: 0 errors, PDF produced) but becomes FATAL once it is also
% written to the .lof — \codeline is \hbox to\linewidth{...}, and a full-width box in
% a list entry dies with "\pdfextension endlink cannot be used in vertical mode", no
% output PDF. Gaining a feature nobody uses (0 of 13 collections call \listoffigures)
% at the price of breaking a construct that works is a bad trade. Stock \caption has
% exactly this limitation and exactly this escape hatch, so the convention is familiar.
% \en{}, \code{}, \h{} and inline math in the short text are all safe — verified, they
% survive the write and typeset correctly in the list (RTL geometry confirmed: number
% rightmost at x=509.9, title left of it, dots and folio at the far left).
% With no optional argument the macro is byte-identical to v3.63 — zero risk to every
% existing document, and \listoffigures lists exactly the captions that opted in.
% BLD-37, TABLE side (the \figcap v3.74 idea, adapted): clicking a \ref to a
% table used to land MID-table — hyperref's automatic anchor inside
% \refstepcounter was emitted on the still-open horizontal line whose TAIL is
% the tabular (an hbox whose baseline is the table's vertical CENTER), so the
% jump scrolled the top half of the table above the viewport. The v3.74 entry
% deferred this side claiming "tablebox splits — no single box to measure";
% that deferral dissolves: the measurement target is the TABULAR itself at
% \tabcap ENTRY, while it is still the tail of the open list — the same
% situation \lakin@grabfig handles at an island's end (and the fleet has 0
% split tableboxes). So \lakin@grabtab runs as the FIRST token, before any
% typesetting can push the box away, and the anchor is re-emitted raised above
% the measured tabular. \lakintabraise is a SEPARATE register from
% \lakinfigraise (defined with it, below the LTR-island block): a captionless
% island intentionally leaves the fig register set, and a table must never
% inherit that stale figure height.
%
% WHY the anchor stays AT THE ENTRY SPOT and is NOT deferred into the caption
% paragraph the way \figcap defers it: at \par TeX DROPS a final glue item,
% and the automatic anchor's smash box sitting BEHIND the interword glue kept
% that glue alive — so every centered tabular has always sat half a space off
% true center. Deferring the anchor removed that box and was measured as 982
% differing glyph records: EVERY tabular line shifted 1.30pt (2.61pt for
% pattern 7's two minipages) while captions stayed put. Emitting the raised
% anchor at the old spot — through hyperref's own \Hy@raisedlink, exactly the
% macro the automatic path uses — keeps the line's node structure, and so the
% A/B glyph diff at 0, BY CONSTRUCTION. \Hy@raisedlink reads the raise from
% \HyperRaiseLinkDefault (default \baselineskip = the old mid-table dest); the
% group-local \edef freezes measured-height+42pt into it, and \smash inside
% \Hy@raisedlink hides any raise from the line's extent. \HyCnt@ProvideTheH-
% Counter keeps the dest NAME byte-identical to the automatic one (it defines
% \theHtbl, which \hyper@makecurrent prefers — without it the first measured
% caption was named tbl.1 where the automatic path said tbl.0.1, measured).
% vmode entry, or anything but a box at the tail (no tabular to measure),
% falls back to the plain \refstepcounter — byte-identical to the old path.
\makeatletter
\newcommand{\tabcap}[2][]{\lakin@grabtab
  % v3.79: arm the box-close correction. Set unconditionally and in BOTH
  % branches below's scope, so the plain-\refstepcounter fallback is corrected
  % too; \lakintabboxopen clears it at every \begin{tablebox}, so a \tabcap
  % outside a tablebox can never leave it armed for a later box.
  \global\lakintabcapseentrue
  \ifdim\lakintabraise>0pt
    \H@refstepcounter{tbl}\HyCnt@ProvideTheHCounter{tbl}\hyper@makecurrent{tbl}%
    \begingroup
      % A frameless table targets its measured top. The frame's historical
      % cushion belongs only to tablebox and is corrected at box close.
      \edef\HyperRaiseLinkDefault{\the\lakintabraise}%
      \iflakintabframed
        \edef\HyperRaiseLinkDefault{\the\dimexpr\lakintabraise+42pt\relax}%
      \fi
      \Hy@raisedlink{\hyper@anchorstart{\@currentHref}\hyper@anchorend}%
    \endgroup
    \global\lakintabraise 0pt\relax
  \else
    \refstepcounter{tbl}%
  \fi
  \lakinrtitle{#2}\par\nopagebreak\addvspace{6pt}%
  \def\lakin@short{#1}\ifx\lakin@short\@empty\else
    \addcontentsline{lot}{table}{\protect\numberline{\thetbl}{\protect\ignorespaces #1}}\fi
  {\centering\lakitcaptionstyle\mbox{\lakitcaptionlabelstyle{\tablename~\en{\thetbl}\lakitcaptionpunctuation}}\enspace #2\par}}
\makeatother

% ============================================================================
% Unified FIGURE caption — the \tabcap twin for TikZ/plot figures (v3.34.0).
% A figure (tikzpic) sits inline in the prose; put \figcap BELOW it for a
% numbered, \ref-able caption "איור N.M: תיאור" in the same style as \tabcap:
%   \begin{tikzpic} ... \end{tikzpic}
%   \figcap{\label{fig:foo}description}     % \label FIRST, inside the argument
% WHY a macro, not a hand-written gray line: a manual
%   {\centering\small\textcolor{cGray}{איור: ...}} with a \label inside has NO
% \refstepcounter, so the \label binds to the AMBIENT \@currentlabel (the last
% equation) — \ref then prints an equation number AS the figure number
% (measured on a real page: "איור 48" where 48 is equation 48). \figcap steps a
% real \c@fig counter first, so \label is correct, numbering is consistent, and
% refs are stable. Class-aware like \thetbl (chapter- or section-prefixed).
% ============================================================================
\makeatletter
\@ifundefined{c@chapter}{%
  \newcounter{fig}[section]
  \renewcommand{\thefig}{\ifnum\value{section}>0 \thesection.\fi\arabic{fig}}%
}{%
  \newcounter{fig}[chapter]
  \renewcommand{\thefig}{\ifnum\value{chapter}>0 \thechapter.\fi\arabic{fig}}%
}
\makeatother
\makeatletter
\newcommand{\figcap}[2][]{%
  \ifdim\lakinfigraise>0pt
    % v3.74: step WITHOUT hyperref's automatic anchor, then emit the anchor
    % raised above the measured figure. \H@refstepcounter is hyperref's saved
    % pre-hook \refstepcounter (sets the counter AND \@currentlabel — \label
    % in the caption arg binds exactly as before); \hyper@makecurrent supplies
    % \@currentHref, so the dest NAME is identical to the old automatic one —
    % \ref/\autoref/bookmarks are unchanged, only the dest'S POSITION moves.
    \H@refstepcounter{fig}\hyper@makecurrent{fig}%
    % The anchor is emitted INSIDE the caption paragraph below, as a zero-size
    % raisebox — NOT here in vertical mode, where a \raisebox starts an empty
    % paragraph and adds a blank line (measured: every glyph after the caption
    % shifted +18.8pt). Inside an existing line a [0pt][0pt] box is inert.
    \edef\lakin@anch{\@currentHref}%
    % \edef FREEZES the measured height into the pending macro: the deferred
    % \raisebox is typeset inside the caption paragraph AFTER the register is
    % reset to the 0pt sentinel below — a plain \def read 0 at use time
    % (measured: LAKINRAISE=0.0, dest moved 14pt instead of above the figure).
    \edef\lakin@pend{\noexpand\raisebox{\the\dimexpr\lakinfigraise+42pt\relax}[0pt][0pt]{%
      \noexpand\hyper@anchorstart{\lakin@anch}\noexpand\hyper@anchorend}}%
    \global\lakinfigraise 0pt\relax
  \else
    \refstepcounter{fig}\let\lakin@pend\relax
  \fi
  \lakinrtitle{#2}\par\nopagebreak\addvspace{4pt}%
  \def\lakin@short{#1}\ifx\lakin@short\@empty\else
    \addcontentsline{lof}{figure}{\protect\numberline{\thefig}{\protect\ignorespaces #1}}\fi
  {\centering\lakin@pend\lakitcaptionstyle\mbox{\lakitcaptionlabelstyle{\figurename~\en{\thefig}\lakitcaptionpunctuation}}\enspace #2\par}}
\makeatother

% ============================================================================
% LTR islands.  EVERY tikzpicture / quantikz MUST live inside one, or the whole
% figure's coordinates and Latin labels flip. (Hebrew node text still uses \h{}.)
% ============================================================================
% v3.74 (BLD-37): the \figcap hyperlink DESTINATION used to land at the caption
% line — BELOW the figure — so clicking a figure reference scrolled the figure
% off-screen (measured fleet-wide: 97/97 fig dests in one collection sat 19-54pt
% under their figure's ink; a 347pt figure was fully hidden after the jump).
% hypcap-style fix in two parts: (1) at the island's END, \lakin@grabfig grabs
% the just-typeset centered line with \lastbox, measures it, and puts it back
% with \nointerlineskip — the interline glue above the line was already
% contributed and is NOT re-added, so the page geometry is unchanged to the
% glyph (verified); (2) \figcap then steps the counter WITHOUT hyperref's
% automatic anchor (\H@refstepcounter = the pre-hyperref \refstepcounter) and
% emits the anchor itself, raised by the measured height — the dest lands ABOVE
% the figure top. A captionless island leaves numbering untouched (nothing is
% stepped at the island; the grab flag is cleared at every island begin, so a
% stale height from a captionless figure cannot leak into a later caption).
\newsavebox{\lakinfigbox}
\newdimen\lakinfigraise  % real \dimen: tex.setdimen cannot target a \newlength SKIP register (measured lua error)
% BLD-37 table side: a SEPARATE register for \tabcap. The two must not share —
% a captionless island INTENTIONALLY leaves \lakinfigraise set (the grab runs
% at every island END; only the NEXT island's begin clears it), so a shared
% register would hand a \tabcap after a captionless figure that stale FIGURE
% height. \lakintabraise is written only by \lakin@grabtab at \tabcap entry
% and consumed/reset inside \tabcap itself, so it can never go stale.
\newdimen\lakintabraise
\makeatletter
% v3.74b: measure WITHOUT touching the list. The first draft grabbed the box
% with \lastbox and put it back — and \lastbox DISCARDS the node's baseline
% shift, so a vertically-anchored quantikz circuit rose 2.57pt and everything
% after it sank 5.32pt (measured on the A/B glyph diff). LuaTeX lets us READ
% the tail node's dimensions without removing anything: geometry is identical
% by construction. \lakinfigraise=0pt is the "no figure measured" sentinel.
\def\lakin@grabfig{%
  \ifhmode
    \directlua{
      local n = tex.nest[tex.nest.ptr].tail
      local GLUE, PEN = node.id('glue'), node.id('penalty')
      while n and (n.id == GLUE or n.id == PEN) do n = n.prev end
      local HL, VL = node.id('hlist'), node.id('vlist')
      if n and (n.id == HL or n.id == VL) then
        tex.setdimen('global', 'lakinfigraise',
                     (n.height or 0) + (n.depth or 0))
      else
        tex.setdimen('global', 'lakinfigraise', 0)
      end}%
  \else
    \global\lakinfigraise 0pt\relax
  \fi}
% BLD-37 table side: \lakin@grabfig's twin, NOT a shared helper — the Lua body
% differs in exactly one place, the skip set. Between \end{tabular} and \tabcap
% an otherlanguage wrapper (patterns 3/5 of example_tables) leaves a \write
% whatsit and a dir node ON TOP of the interword glue (measured node dump:
% "whatsit/write <- dir <- glue(2.62) <- hlist"), so the table reader must
% also step over whatsit/dir markers — all zero-size — to reach the tabular.
% \lakin@grabfig deliberately does NOT get the wider skip set: its glue/penalty
% body is the one the v3.74 A/B proof froze, and widening it would newly
% measure (and so newly raise) anchors in island shapes that today take the
% plain path — an unproven behavior change on the FIG side. Anything that is
% not a box behind those markers still hits the 0pt sentinel and \tabcap falls
% back to the plain \refstepcounter.
% SECOND deliberate difference: this stores the box's HEIGHT, not height+depth.
% \figcap's anchor is emitted in the caption paragraph BELOW the figure, so it
% must climb the full extent; \tabcap's anchor sits ON the tabular's own line,
% whose baseline is the vcentered tabular's MIDPOINT — height alone reaches
% the top edge, and adding the depth was measured to overshoot by half the
% table (dests 80-135pt above their frame, INSIDE the preceding table's box).
\def\lakin@grabtab{%
  \ifhmode
    \directlua{
      local n = tex.nest[tex.nest.ptr].tail
      local GLUE, PEN = node.id('glue'), node.id('penalty')
      local WHAT, DIR = node.id('whatsit'), node.id('dir')
      while n and (n.id == GLUE or n.id == PEN
                   or n.id == WHAT or n.id == DIR) do n = n.prev end
      local HL, VL = node.id('hlist'), node.id('vlist')
      if n and (n.id == HL or n.id == VL) then
        tex.setdimen('global', 'lakintabraise', n.height or 0)
      else
        tex.setdimen('global', 'lakintabraise', 0)
      end}%
  \else
    \global\lakintabraise 0pt\relax
  \fi}
\makeatother
% ============================================================================
% v3.79 (BLD-37 table side, CORRECTION): put the anchor at the TOP OF THE BOX,
% instead of guessing the distance to it.
% ============================================================================
% v3.77 raised the \tabcap anchor by (height of the box at the tail of the open
% horizontal list) + 42pt. That is a MEASUREMENT taken BEFORE line breaking, and
% it is short by every term it cannot see. Fleet-measured on the 12 built
% collections (299 tbl.* dests, ownership established through each table's own
% CAPTION, never through "which frame encloses the dest"): 3 dests landed INSIDE
% their own tablebox, below the frame top —
%   computer-architecture       tbl.14.1  -46.4pt   three stacked tabulars
%   computer-architecture       tbl.30.5  -13.5pt   tabular 2.85pt over \linewidth
%   advanced-electronic-devices tbl.B.1   -59.2pt   tabular 2.56pt over \linewidth
% Two missing terms, both reproduced in a minimal document (probe P2/P3):
%   (1) VERTICAL LIST ALREADY CONTRIBUTED inside the box — a second and third
%       tabular, or a heading above the table. The tail read sees only the LAST
%       box, so 77.7pt of stacked tabulars is invisible to the raise.
%   (2) WHICH LINE THE ANCHOR LANDS ON. When the tabular alone overfills
%       \linewidth (459.6121pt in a tablebox), the interword glue after
%       \end{tabular} becomes the chosen breakpoint and \Hy@raisedlink is
%       emitted on the NEXT line — one (depth + lineskip) lower. That term does
%       not EXIST at \tabcap time: line breaking has not happened yet. No
%       measurement taken there can ever recover it. (Proof: probe P3, the same
%       table 2.84912pt too wide, fails at -13.5pt; probe P4, the identical
%       table with \tabcolsep 2pt so it fits, is correct at +32.0pt.)
% So the raise is not computed at \tabcap at all. It is CORRECTED at
% \end{tablebox}, from inside mdframed's own body vbox (endinnercode runs after
% \par\unskip, i.e. after line breaking, while tex.nest[tex.nest.ptr] is still
% that vbox's vertical list). There the anchor's true offset from the top of the
% box content is not estimated but WALKED, and the \Hy@raisedlink shift is set
% so the dest lands 42pt above it — which is 42pt - innertopmargin(9pt) -
% linewidth(0.9pt) = 32.0pt above the FRAME top, the exact position the 295
% already-correct dests have today. That is why this is a correction and not a
% re-placement: the good sites do not move.
%
% WHY THE GLYPH DIFF IS 0 BY CONSTRUCTION — the only thing written is the
% `shift` field of the hlist that \Hy@raisedlink already built, and the guard
% below refuses to write it unless that box is 0x0x0. A zero-size box carries no
% dimension into its parent's packing, the packing already happened, and shift
% is not an input to it. Nothing else in the list is touched: no node is added,
% removed, resized or reordered. Measured, not assumed: assets/example_tables.tex
% and the whole computer-architecture collection A/B at 0 differing pymupdf
% rawdict glyph records.
%
% WHY NOT the alternatives (all prototyped and measured before this one):
%   * accumulate the vertical list at \tabcap time — fixes term (1), CANNOT fix
%     term (2); 2 of the 3 real failures survive. A fix that leaves the majority
%     of the defect standing is the false all-clear this repo exists to avoid.
%   * \zref-savepos two-pass — correct on all shapes, but it makes the dest
%     POSITION aux-derived (DEFECTS BLD-1/2), a single-pass compile silently
%     ships an unraised anchor, and the bare y-difference has no page guard.
%   * emit the anchor at \begin{tablebox} and let \tabcap adopt it — needs the
%     tbl counter stepped at box OPEN to keep the dest NAME byte-identical; a
%     captionless tablebox then either shifts every later table number or leaves
%     an orphan dest that the next table duplicates (a HARD build gate). This
%     correction touches neither the counter nor the name: \tabcap's naming path
%     is v3.77's, unchanged, so resolve_names() is byte-identical (verified).
%   * suppressing the line break (a penalty before the glue) — that MOVES the
%     tabular: the break currently discards the interword glue, so forbidding it
%     re-centers the two overfull tables. Reflow, rejected.
%
% SCOPE. The correction is armed only by \tabcap ITSELF (\lakintabcapseen),
% cleared at every box open, so a captionless tablebox is a no-op and a \tabcap
% hosted OUTSIDE a tablebox (the fleet has exactly one, computer-architecture
% tbl.14.2 in an exbox) keeps v3.77's measured raise. Both fallbacks are
% therefore never worse than today; and because hyperref's AUTOMATIC anchor uses
% the same \Hy@raisedlink shape, the plain-\refstepcounter path (vmode entry, or
% prose between \end{tabular} and \tabcap — probe P6, -51.5pt) is corrected too.
% KNOWN LIMIT — a tablebox that SPLITS across pages. The offset is measured from
% the top of the UNSPLIT body, so a dest in segment 2 gets raised past its own
% segment. MEASURED on a forced split (probe: 40 filler lines + a table): the
% dest went from y=585.0 on page 2 (inside segment 2, at the caption — the old
% BLD-37 symptom) to y=-113.4, i.e. ABOVE the page top. Viewers clamp that to the
% page top, which happens to show segment 2 whole, but an out-of-page coordinate
% is not something to ship on purpose. Hence the \textheight guard in
% lakintab.fix: a body TALLER THAN THE TEXT BLOCK cannot fit on any page, so it
% must split — the correction is skipped and v3.77's behaviour stands unchanged
% (verified: the probe returns to y=585.0). The guard is deliberately EXACT, not
% a page-fitting heuristic: it is sound (a taller-than-page body always splits)
% but PARTIAL (a shorter body that starts low on a page could still split). That
% residue is unreachable in practice — v3.37's everyline + the caption's
% \nopagebreak glue move such boxes WHOLE instead (probe N7: a 22-row table
% pushed down by 430pt moved whole, 0 segments), and the fleet has 0 split
% tableboxes, re-counted after this change. Reimplementing mdframed's own
% page-fitting arithmetic to close the residue would be a heuristic with its own
% failure mode, which is the trade this repo has repeatedly refused.
\newif\iflakintabcapseen
\newif\iflakintabframed
% The Lua lives in ONE \directlua at load time, and is written with NO
% backslash, %, # or ~ : \directlua expands its argument like \edef, so an
% active ~ would expand to \nobreakspace and a % would eat the rest of the line
% (the second is how the first draft of this block died). "not (a == b)" stands
% in for Lua's inequality operator for the same reason.
\directlua{
  lakintab = { }
  lakintab.HL = node.id("hlist")
  lakintab.VL = node.id("vlist")
  lakintab.RU = node.id("rule")
  lakintab.GL = node.id("glue")
  lakintab.KE = node.id("kern")
  lakintab.WH = node.id("whatsit")
  lakintab.TM = node.id("temp")
  lakintab.WD = node.subtype("pdf_dest")
  function lakintab.d(v)
    if v == nil then return 0 end
    if v < -1000000000 then return 0 end
    return v
  end
  function lakintab.scan(list, y, horiz, parent, target)
    local x = list
    while x do
      local i = x.id
      if i == lakintab.HL or i == lakintab.VL then
        if horiz then
          local b = y + x.shift
          if i == lakintab.HL then
            lakintab.scan(x.head, b, true, x, target)
          else
            lakintab.scan(x.head, b - lakintab.d(x.height), false, x, target)
          end
        else
          if i == lakintab.HL then
            lakintab.scan(x.head, y + lakintab.d(x.height), true, x, target)
          else
            lakintab.scan(x.head, y, false, x, target)
          end
          y = y + lakintab.d(x.height) + lakintab.d(x.depth)
        end
      elseif i == lakintab.WH then
        if x.subtype == lakintab.WD then
          local id = x.dest_id
          if type(id) == "string" and string.sub(id, 1, 4) == "tbl." then
            if horiz and parent and parent.width == 0
               and parent.height == 0 and parent.depth == 0 then
              parent.shift = parent.shift + (target - y)
            end
          end
        end
      elseif not horiz then
        if i == lakintab.RU then
          y = y + lakintab.d(x.height) + lakintab.d(x.depth)
        elseif i == lakintab.GL or i == lakintab.KE then
          y = y + lakintab.d(x.width)
        end
      end
      x = x.next
    end
  end
  function lakintab.fix(target, maxh)
    local nst = tex.nest[tex.nest.ptr]
    if nst == nil then return end
    local h = nst.head
    if h == nil then return end
    if h.id == lakintab.TM then h = h.next end
    local tot = 0
    local x = h
    while x do
      local i = x.id
      if i == lakintab.HL or i == lakintab.VL or i == lakintab.RU then
        tot = tot + lakintab.d(x.height) + lakintab.d(x.depth)
      elseif i == lakintab.GL or i == lakintab.KE then
        tot = tot + lakintab.d(x.width)
      end
      x = x.next
    end
    if tot > maxh then return end
    lakintab.scan(h, 0, false, nil, target)
  end
}
% @-FREE NAMES ON PURPOSE: these two are named in the `tablebox` key list at
% \newmdenv above, which is tokenised OUTSIDE any \makeatletter region — a
% \lakin@... there would tokenise as \lakin followed by the letters @tabopen.
\newcommand{\lakintabboxopen}{\lakintabframedtrue\global\lakintabcapseenfalse}
\newcommand{\lakintabboxclose}{%
  \iflakintabcapseen
    \directlua{lakintab.fix(-\number\dimexpr\lakintabcushion\relax,
                            \number\dimexpr\textheight\relax)}%
    \global\lakintabcapseenfalse
  \fi\lakintabframedfalse}
% The cushion above the box's CONTENT top. 42pt puts the dest 32.0pt above the
% frame top (42pt - innertopmargin 9pt - linewidth 0.9pt), i.e. exactly where
% the 295 correct dests already are — the correction moves only the broken ones.
\newcommand{\lakintabcushion}{42pt}
\newenvironment{qcirc}{\par\medskip\global\lakinfigraise 0pt\relax\begin{otherlanguage}{english}\begin{center}}
  {\csname lakin@grabfig\endcsname\end{center}\end{otherlanguage}\par\medskip}
\newenvironment{tikzpic}{\par\medskip\global\lakinfigraise 0pt\relax\begin{otherlanguage}{english}\begin{center}}
  {\csname lakin@grabfig\endcsname\end{center}\end{otherlanguage}\par\medskip}

% ============================================================================
% Code listings.  ★★ v3.4 THE BOX-OVERFLOW FIX:
%   resetmargins=true   → the listing IGNORES the enclosing environment's
%                         indentation and measures from the PAGE. v3.4 stated
%                         this backwards ("measured from the ENCLOSING box"),
%                         and v3.37 then blanked in-box line numbers because
%                         the page-anchored listing put its hanging digits on
%                         top of the box rule. Measured v4.5.0: with the key
%                         FALSE the listing follows the box and the code moves
%                         from x=75.26 to x=97.74, clearing the rule at 72.07.
%                         The key now lives in one \lstset below instead of in
%                         each style, so a box can override it — codebox does.
%   framexleftmargin=0  → the frame (left bar) does not extend left of the code.
%   xleftmargin keeps a small indent so the line numbers have room INSIDE.
% asmblock wraps the listing in an LTR island (code is LTR).
% ============================================================================
\lstdefinestyle{asm}{
  basicstyle=\ttfamily\small,
  backgroundcolor=\color{cBlueBg!55},
  frame=leftline,framerule=2.5pt,rulecolor=\color{cBlue},
  numbers=left,numberstyle=\tiny\color{cGray},numbersep=8pt,
  xleftmargin=2.2em,framexleftmargin=0pt,framexrightmargin=0pt,
  breaklines=true,columns=fullflexible,keepspaces=true,
  showstringspaces=false,aboveskip=8pt,belowskip=8pt,
  % v3.36: without a comment DELIMITER the commentstyle below never fires —
  % listings has no `language` here, so it does not know '#' opens a comment,
  % and every asm listing rendered flat black (measured: a whole-page pixel
  % sweep found ZERO cGreen glyphs across 135 shipped listings). morecomment
  % gives it the delimiter; color-only change, glyph widths untouched.
  morecomment={[l]{\#}},
  commentstyle=\color{cGreen},
}
\lstdefinestyle{cstyle}{
  language=C,basicstyle=\ttfamily\small,
  backgroundcolor=\color{cGreenBg!55},
  frame=leftline,framerule=2.5pt,rulecolor=\color{cGreen},
  numbers=left,numberstyle=\tiny\color{cGray},numbersep=8pt,
  xleftmargin=2.2em,framexleftmargin=0pt,framexrightmargin=0pt,
  breaklines=true,columns=fullflexible,keepspaces=true,
  showstringspaces=false,keywordstyle=\color{cBlue}\bfseries,
  commentstyle=\color{cGray}\itshape,aboveskip=8pt,belowskip=8pt,
}
% v3.36: pystyle/htmlstyle — colored twins of cstyle for the Python/HTML
% listings that used to repeat a ~200-char options blob with NO color
% (keywords rendered bold-black, comments italic-black — measured; C alone
% was colored). Geometry is copied VERBATIM from the shipped blob (frame,
% margins, skips) so swapping a call site to style=pystyle moves NOTHING —
% only keywordstyle/commentstyle/stringstyle are added (bold stays bold,
% italic stays italic, mono metrics identical => zero reflow, proven by
% per-page word-bbox diff). breaklines/breakatwhitespace stay per-site.
% v3.78: numbers=left joins both styles, with the exact numberstyle/numbersep
% keys asm/cstyle carry — fleet-wide consistency (author decision: all four
% styles numbered; "asm/cstyle numbered, py/html not" was an accident of the
% v3.36 geometry freeze, not a design). The numbers hang LEFT of the code box
% (listings sets them with \llap, so no overfull hbox can arise and the code
% x is untouched — measured: code frozen, digits at code_x-11.8pt); inside a
% colored box they are blanked like everyone else's (v3.37 hook), and the
% clean variant restores them in its frameless boxes (v3.78, above).
\lstdefinestyle{pystyle}{
  language=Python,basicstyle=\ttfamily\small,
  columns=fullflexible,keepspaces=true,
  frame=single,rulecolor=\color{cTblLine},
  numbers=left,numberstyle=\tiny\color{cGray},numbersep=8pt,
  xleftmargin=0.5em,framexleftmargin=0pt,framexrightmargin=0pt,
  showstringspaces=false,aboveskip=8pt,belowskip=8pt,
  keywordstyle=\color{cBlue}\bfseries,
  commentstyle=\color{cGray}\itshape,
  stringstyle=\color{cGreen},
}
% v4.5.0: hoisted out of the four styles so an enclosing box can override it.
% Every listing still resolves to true unless something later says otherwise,
% so shipped output is unchanged (verified byte-for-byte on all eight PDFs).
\lstset{resetmargins=true}
\lstdefinestyle{htmlstyle}{
  language=HTML,basicstyle=\ttfamily\small,
  columns=fullflexible,keepspaces=true,
  frame=single,rulecolor=\color{cTblLine},
  numbers=left,numberstyle=\tiny\color{cGray},numbersep=8pt,
  xleftmargin=0.5em,framexleftmargin=0pt,framexrightmargin=0pt,
  showstringspaces=false,aboveskip=8pt,belowskip=8pt,
  keywordstyle=\color{cBlue}\bfseries,
  commentstyle=\color{cGray}\itshape,
  stringstyle=\color{cGreen},
}
\newenvironment{asmblock}{\par\begin{otherlanguage}{english}}{\end{otherlanguage}\par}

% <hllpdf:listing-smart-quotes>
% v4.21.0: four Unicode punctuation characters must enter listings' token
% buffer. Unregistered characters bypass it: a closing quote was printed
% BEFORE the pending word, and an opening quote BEFORE a pending brace.
% Extend the standard character table, then activate only these codepoints
% in the listing's local character-table scope. Do this BEFORE the package's
% delimiter/alsoother/literate hooks so explicit user mappings still win.
% This is not universal Unicode support, nor a conversion to ASCII quotes.
% See listings developer guide 15.3.1 and references/listing-smart-quotes.md.
\makeatletter
\lst@CCPut \lst@ProcessOther {"2018}{"2019}{"201C}{"201D}\z@
\lst@AddToHookAtTop{SelectCharTable}{%
  \lst@ifactivechars
    \catcode"2018=\active \catcode"2019=\active
    \catcode"201C=\active \catcode"201D=\active
  \fi}
\makeatother
% </hllpdf:listing-smart-quotes>

% <hllpdf:listing-technical-symbols>
% v4.30: bounded extension for common technical punctuation in LTR listings.
% Without character-table registration these glyphs bypass the token buffer
% and can move before the preceding word even inside quoted strings.
% Keep the original four-quote block separate for its standing mutation test.
% This is source-faithful printing, not syntax validation or general Unicode.
\makeatletter
\lst@CCPut \lst@ProcessOther
  {"2013}{"2014}{"2026}{"2190}{"2192}{"2194}{"21D2}
  {"2212}{"221E}{"2260}{"2264}{"2265}\z@
\lst@AddToHookAtTop{SelectCharTable}{%
  \lst@ifactivechars
    \catcode"2013=\active \catcode"2014=\active \catcode"2026=\active
    \catcode"2190=\active \catcode"2192=\active \catcode"2194=\active
    \catcode"21D2=\active \catcode"2212=\active \catcode"221E=\active
    \catcode"2260=\active \catcode"2264=\active \catcode"2265=\active
  \fi}
\makeatother
% </hllpdf:listing-technical-symbols>


\usepackage[hidelinks,unicode]{hyperref}   % LAST, after colors. unicode = Hebrew bookmarks.
% v3.62: THE ISOLATION MACROS CORRUPTED EVERY PDF BOOKMARK THEY APPEARED IN.
% \en is defined (line ~155) as
%     \foreignlanguage{english}{\language=\csname l@nohyphenation\endcsname #1}
% and in a heading that body is written verbatim to the .toc and handed to
% hyperref's \pdfstringdef. hyperref strips the two CONTROL SEQUENCES (\language,
% \l@nohyphenation) but the plain `=` CHARACTER token between them is not a control
% sequence, so it survives into the PDF string: \chapter{שער \en{CNOT} והצמדה}
% produced the bookmark "שער =CNOT והצמדה". The printed TOC and the heading on the
% page are CORRECT — only the navigation pane and the viewer window title are wrong,
% which is why it went unnoticed. Measured on the shipped fleet: 70 corrupted
% bookmarks across 5 PDFs, from 177 unprotected \en-in-heading sites. Same corruption
% reached /Title and /Subject via the \hypersetup{pdftitle=...} the docs recommend.
% This hook tells hyperref to use the plain argument for all three isolation macros
% when building a PDF string. It does NOT touch the .toc (the printed TOC is
% byte-unchanged) and it does NOT mask the math signal — a heading with $...$ still
% logs "removing math shift", so BIDI-15 detection is intact.
% BONUS, measured: it also makes \code{} in an UNSTARRED heading COMPILE. BLD-10
% recorded that as needing \texorpdfstring or a starred heading; the real root cause
% is hyperref's \pdfstringdef scan of \code's \ifmmode...\else...\fi body — which is
% exactly why the identically-bodied \texttt survived (hyperref pre-neutralizes
% \texttt in PDF strings) and why \DeclareRobustCommand could never help.
% v3.64: \def\\{ } added. A manual line break in a heading WELDS the two words
% together in the bookmark: \chapter{שורה ראשונה\\שורה שנייה} produced the outline
% entry 'שורה ראשונהשורה שנייה' — "ראשונהשורה" is not a word. hyperref drops the
% \\ and puts nothing in its place, so the break that looks like a space on the
% page is a zero-width join in the navigation pane. The spaced form
% 'מילה אחת \\ מילה שנייה' escapes with a harmless DOUBLE space ('אחת  מילה'),
% which is why this never got noticed. Both are measured on a real outline.
% 0 live instances in the corpus — but SKILL.md and content-style.md both RECOMMEND
% a manual heading break, so this is a primed trap rather than a dormant one, and a
% one-token preamble fix retires it for good instead of asking authors to remember.
% Mapping \\ to a space fixes the WELD, which is the actual defect (a nonexistent
% word in the navigation pane). It does not normalise whitespace: measured after the
% fix, the welded form becomes 'שורה ראשונה שורה שנייה' (correct) and the already-
% spaced form goes from two spaces to three. That is cosmetic — viewers render it as
% a wider gap, and no word is corrupted. A per-heading \texorpdfstring is still the
% way to get an exactly-tuned bookmark string.
\pdfstringdefDisableCommands{\def\en#1{#1}\def\h#1{#1}\def\doclang#1{#1}\def\code#1{#1}\def\codeline#1{#1}\def\\{ }}
% v3.70: RESTORE \url LINE-BREAKING. url.sty freezes its break penalties at LOAD
% time (\mathchardef\UrlBreakPenalty=\binoppenalty), and v3.45's global
% \binoppenalty=10000 (inline-math atomicity) therefore priced EVERY breakpoint
% inside \Url at 10000 = never break. Measured: a long \url in Hebrew prose is
% one unbroken line, ~100pt overfull, invading the margin band the ink gates
% deliberately tolerate. With the two lines below: 0 overfull, the URL breaks
% after '/', the RTL split is CORRECT (fragment order forward), and each line
% gets its own link annotation exactly covering its fragment (measured).
% Confined to \Url's group — inline-math atomicity (TYP-7) is untouched.
\mathchardef\UrlBreakPenalty=200
\mathchardef\UrlBigBreakPenalty=100
\ifx\lakitmainlanguage\hllpdfenglish
  \hypersetup{pdflang=en}
\else
  \hypersetup{pdflang=he}
\fi  % /Lang is the document base locale (not per-span accessibility tagging).
% Set a per-document \hypersetup{pdftitle={...}} too — see SKILL.md.

% v3.52 (round-2 audit): two RTL bibliography/cross-ref defects, both invisible to
% every gate and both affecting the FIRST author to use the feature.
% (a) \@biblabel's default [#1] brackets are bidi-neutral and mirror to ]1[ under
% the RTL bibliography base — every \thebibliography reference list shows ]1[.
% LTR-isolating the label restores [1]. (The render gate's mirror rule fingerprints
% round parens )1(, so square brackets slipped through.)
\makeatletter
\renewcommand*{\@biblabel}[1]{\foreignlanguage{english}{[#1]}}
% v3.54: the SAME neutral-bracket trap on the INLINE side — a multi-key \cite{a,b}
% under the RTL base lays its European-number list right-to-left, so [1, 2] renders
% "[2,1]" and [1, 2, 3] renders "[3,2,1]" (verified by glyph geometry — the reader
% sees the sources in REVERSE order; a note form \cite[p.5]{a} also reorders to
% "[p. 5,1]"). LTR-isolate the whole bracket group (number list AND optional note),
% exactly like \@biblabel above: [1, 2] / [1, 2, 3] / [1, p. 5] all render correct.
% Placed AFTER hyperref so it wins; hyperref hyperlinks each number inside #1 (via
% \@citex), which this wrapper leaves intact. Mirrors the LaTeX kernel \@cite shape.
% v4.9.1: directional isolation must NOT switch the language here. Babel's
% safe-active citation wrapper passes its key list through \bbl@tempa;
% \foreignlanguage reuses that scratch macro before the deferred citation
% loop expands it. The old line therefore cited the literal key "english".
% If such a key existed, even the undefined-reference gate stayed green.
% \babelsublr isolates the bracket group without selecting a locale or
% changing that key list. Preserve native links, key order and optional notes.
\def\@cite#1#2{\babelsublr{[{#1\if@tempswa , #2\fi}]}}
\makeatother

% v4.11.0: a reference inside align's \text inherited HEBdisptag=true,
% so its enclosing hbox mirrored the brackets despite a clean compile.
% Isolate the complete printed reference without selecting a Babel language.
% Also neutralize the display-only normalfont direction reset locally. A
% direction-only patch inverted 12 to 21 inside equation in the measured probe.
% Plain maketag is local to the reference: mathtools' display-side wrapper
% otherwise globally consumes manual_tag while formatting a nested reference.
% Preserve each original command (including the selective auxiliary write),
% current tag forms, native links, and showonlyrefs on/off transitions.
\makeatletter
\NewCommandCopy{\HEB@nativeeqref}{\eqref}
\newcommand{\HEB@eqrefisolate}[1]{%
  \babelsublr{\HEBdisptagfalse\let\bbl@eqnodir\relax\let\maketag@@@\ORIGmaketag@@@#1}}
\protected\def\eqref#1{\HEB@eqrefisolate{\HEB@nativeeqref{#1}}}
\ExplSyntaxOn
\NewCommandCopy{\HEB@nativeMTeqref}{\MT_extended_eqref:n}
\cs_set_protected:Npn \MT_extended_eqref:n #1
  {\HEB@eqrefisolate{\HEB@nativeMTeqref{#1}}}
\ExplSyntaxOff
\makeatother

% v4.11.0: custom caption counters are not part of hyperref's English
% locale, and Hebrew has no upstream hyperref locale in this toolchain.
% Restore the appropriate names at every Babel language switch, including
% return from a nested inline island. Do not localize the destination keys.
\makeatletter
\newcommand{\HEB@refnames}{%
\def\chapterautorefname{פרק}%
\def\sectionautorefname{סעיף}%
\def\subsectionautorefname{תת־סעיף}%
\def\subsubsectionautorefname{תת־תת־סעיף}%
\def\paragraphautorefname{פסקה}%
% v4.31: hyperref uses a distinct counter name at the deepest standard level.
\def\subparagraphautorefname{תת־פסקה}%
\def\equationautorefname{משוואה}%
\def\figureautorefname{איור}%
\def\tableautorefname{טבלה}%
\def\partautorefname{חלק}%
\def\appendixautorefname{נספח}%
\def\footnoteautorefname{הערת שוליים}%
\def\itemautorefname{סעיף}%
\def\Itemautorefname{סעיף}%
\def\figautorefname{איור}%
\def\tblautorefname{טבלה}%
\def\pageautorefname{עמוד}% 
}
\addto\extrashebrew{\HEB@refnames}
\newcommand{\ENG@refnames}{%
  \def\subparagraphautorefname{subparagraph}%
  \def\chapterautorefname{chapter}%
  \def\sectionautorefname{section}%
  \def\subsectionautorefname{subsection}%
  \def\subsubsectionautorefname{subsubsection}%
  \def\paragraphautorefname{paragraph}%
  \def\equationautorefname{Equation}%
  \def\figureautorefname{Figure}%
  \def\tableautorefname{Table}%
  \def\partautorefname{Part}%
  \def\appendixautorefname{Appendix}%
  \def\footnoteautorefname{footnote}%
  \def\itemautorefname{item}%
  \def\figautorefname{Figure}%
  \def\tblautorefname{Table}%
  \def\Itemautorefname{item}%
  \def\pageautorefname{page}%
}
\addto\extrasenglish{\ENG@refnames}
\ifx\lakitmainlanguage\hllpdfenglish\ENG@refnames\else\HEB@refnames\fi
\makeatother

% ============================================================================
% Optional running head (v3.39): \lakithead{<book title>} — OFF by default.
% A document opts in with ONE line after % ============================================================================
% preamble.tex — tested LuaLaTeX preamble for Hebrew + math + figures PDFs (v4.32.0)
% Battle-tested: compiled an 80-page Hebrew quantum-computing summary AND a
% 197-page Hebrew computer-architecture anthology (TikZ diagrams, code listings,
% colored boxes) — both with 0 missing characters.
%
% Usage in your document:
%   % ============================================================================
% preamble.tex — tested LuaLaTeX preamble for Hebrew + math + figures PDFs (v4.32.0)
% Battle-tested: compiled an 80-page Hebrew quantum-computing summary AND a
% 197-page Hebrew computer-architecture anthology (TikZ diagrams, code listings,
% colored boxes) — both with 0 missing characters.
%
% Usage in your document:
%   \input{preamble.tex}     % (or paste these lines at the top)
%   \begin{document}  ...  \end{document}
%
% Compile with:  lualatex -interaction=nonstopmode doc.tex
%   Repeat until two consecutive PDFs are byte-identical (scripts/build.sh
%   does this) — a TOC document converges only on pass 3+; after two passes
%   the TOC page numbers are STALE (measured; see SKILL.md step 3).
% Font binaries are external dependencies, pinned by assets/fonts/SHA256SUMS.
% scripts/setup_fonts.sh imports a verified local installation or prior skill.
%
% PUBLIC DOCUMENT-SHELL / LAYOUT CONTRACT (introduced v4.3.1; architecture-owned v4.7.1). These hooks may be
% defined BEFORE \input{preamble.tex}; if omitted, the historical canonical
% defaults below are exact.  They deliberately control only the document shell
% and page layout.  RTL/math semantics stay in this preamble, while visual
% presentation themes (for example clean_theme.tex) load AFTER it.
%
%   \def\lakitdocumentclass{article}
%   \def\lakitclassoptions{11pt,a4paper}
%   \def\lakitgeometryoptions{margin=1.8cm,headheight=14pt}
%   \def\lakitlinespread{1.15}
%   \def\lakitparagraphstyle{block} % letters/briefs; indented for continuous prose
%   \def\lakittocdepth{1}
%   \def\lakitsecnumdepth{2}
%   \input{preamble.tex}
%
% Supported class contract: report or article.  Alternate paper sizes must also
% be declared to build.sh through EXPECTED_PAPER; the PDF gate never infers them.
% Do NOT use these hooks as a theme mechanism: theme overlays must not own page
% geometry, class selection, or direction semantics. A declared presentation-policy
% overlay may intentionally change numbering visibility after the preamble.
% ============================================================================
% Runtime architecture guard: generated from scripts/architecture_contract.py.
% The readiness sentinel remains delayed until the END of this file.
% <hllpdf:architecture-guard>
\ifdefined\hllpdfcoreloaded
  \errmessage{hllpdf: preamble.tex loaded more than once}
  \expandafter\endinput
\fi
% </hllpdf:architecture-guard>

\ifdefined\lakitdocumentclass\else\def\lakitdocumentclass{report}\fi
\ifdefined\lakitclassoptions\else\def\lakitclassoptions{12pt,a4paper}\fi
\ifdefined\lakitgeometryoptions\else\def\lakitgeometryoptions{margin=2.3cm,headheight=14pt}\fi
\ifdefined\lakitlinespread\else\def\lakitlinespread{1.3}\fi
\ifdefined\lakittocdepth\else\def\lakittocdepth{1}\fi
\ifdefined\lakitsecnumdepth\else\def\lakitsecnumdepth{2}\fi
% Paragraph organization belongs to the document layout, independently of its
% language and theme. Preserve historical class dimensions unless explicitly
% selecting block paragraphs; never suppress indentation with an everypar hack.
\ifdefined\lakitparagraphstyle\else\def\lakitparagraphstyle{indented}\fi
\edef\lakitparagraphstyle{\lakitparagraphstyle}
\def\hllpdfparagraphblock{block}
\def\hllpdfparagraphindented{indented}
\ifx\lakitparagraphstyle\hllpdfparagraphblock\else
  \ifx\lakitparagraphstyle\hllpdfparagraphindented\else
    \errmessage{hllpdf: paragraph style must be block or indented}
  \fi
\fi
% v4.3.1: the public class hook has a deliberately narrow support contract.
% Before this guard, \def\lakitdocumentclass{book} silently compiled even though
% every heading/TOC regression test targets report/article; an unsupported shell
% could therefore look "supported" until a subtle class-specific defect shipped.
% Fail before \documentclass rather than letting an undocumented combination enter
% the mechanism core. Plain-TeX comparison is used because LaTeX packages are not
% available yet at this point.
\edef\lakitvalidatedclass{\lakitdocumentclass}
\def\lakitreportclass{report}
\def\lakitarticleclass{article}
\ifx\lakitvalidatedclass\lakitreportclass
\else\ifx\lakitvalidatedclass\lakitarticleclass
\else
  \errmessage{hllpdf: unsupported class \lakitdocumentclass; use report or article}
\fi\fi
\let\lakitvalidatedclass\relax
\let\lakitreportclass\relax
\let\lakitarticleclass\relax
% Language intent is independent of document class, theme and content domain.
% The two dominant-language profiles differ in authoring intent, not direction.
% A balanced bilingual document must explicitly choose a navigation/base locale.
% <hllpdf:language-profile>
\providecommand{\lakitlanguageprofile}{hebrew}
\edef\lakitlanguageprofile{\lakitlanguageprofile}
\def\hllpdfhebrew{hebrew}\def\hllpdfenglish{english}
\def\hllpdfhebrewmixed{hebrew-with-english}
\def\hllpdfenglishmixed{english-with-hebrew}
\def\hllpdfbilingual{bilingual}
\let\hllpdfexpectedlanguage\relax
\ifx\lakitlanguageprofile\hllpdfhebrew
  \let\hllpdfexpectedlanguage\hllpdfhebrew
\else\ifx\lakitlanguageprofile\hllpdfhebrewmixed
  \let\hllpdfexpectedlanguage\hllpdfhebrew
\else\ifx\lakitlanguageprofile\hllpdfenglish
  \let\hllpdfexpectedlanguage\hllpdfenglish
\else\ifx\lakitlanguageprofile\hllpdfenglishmixed
  \let\hllpdfexpectedlanguage\hllpdfenglish
\else\ifx\lakitlanguageprofile\hllpdfbilingual
  \ifdefined\lakitmainlanguage\else
    \errmessage{hllpdf: bilingual profile requires explicit lakitmainlanguage}
  \fi
\else
  \errmessage{hllpdf: unsupported language profile \lakitlanguageprofile}
\fi\fi\fi\fi\fi
\ifdefined\lakitmainlanguage
  \edef\lakitmainlanguage{\lakitmainlanguage}
\else
  \let\lakitmainlanguage\hllpdfexpectedlanguage
\fi
\ifx\lakitmainlanguage\hllpdfhebrew\else
  \ifx\lakitmainlanguage\hllpdfenglish\else
    \errmessage{hllpdf: main language must be hebrew or english}
  \fi
\fi
\ifx\hllpdfexpectedlanguage\relax\else
  \ifx\lakitmainlanguage\hllpdfexpectedlanguage\else
    \errmessage{hllpdf: main language conflicts with language profile}
  \fi
\fi
\typeout{HLLLANG|\lakitlanguageprofile|\lakitmainlanguage}
% </hllpdf:language-profile>
\documentclass[\lakitclassoptions]{\lakitdocumentclass}
\usepackage[no-math]{fontspec}        % (1) ★ no-math — THE critical line (see recipe.md)
% layout=lists asks Babel to adapt list parshape when paragraph and page
% directions differ. Without it, minority-language list bodies lose their
% indentation and labels hang outside the text block, despite a clean log.
% The upstream mechanism preserves ordinary same-direction list geometry.
\usepackage[bidi=basic,layout=lists]{babel}
\ifx\lakitmainlanguage\hllpdfenglish
  \babelprovide[import]{hebrew}
  \babelprovide[import,main]{english}
\else
  \babelprovide[import,main]{hebrew}
  \babelprovide[import]{english}
\fi
% v4.22: bind canonical families and fallbacks to the selected installation.
% Lookup by basename or family alone is not proof of the file actually used.
% <hllpdf:font-binding>
\directlua{
  local dir = os.getenv("HEBREW_LATEX_FONT_DIR")
  if not dir or dir == "" then
    local data = os.getenv("XDG_DATA_HOME")
    if not data or data == "" or not (data:sub(1,1) == "/") then
      data = (os.getenv("HOME") or "") .. "/.local/share"
    end
    dir = data .. "/fonts/hebrew-lualatex-pdf"
  end
  if not (dir:sub(1,1) == "/") then dir = lfs.currentdir() .. "/" .. dir end
  if dir:find("[",1,true) or dir:find("]",1,true) or dir:find(string.char(0),1,true)
     or dir:find(string.char(10),1,true) or dir:find(string.char(13),1,true) then
    tex.error("hllpdf: unsupported character in font directory")
  end
  if not (lfs.attributes(dir,"mode") == "directory") then
    tex.error("hllpdf: font directory is missing")
  end
  hllpdf_font_dir = dir:gsub("/+$", "") .. "/"
  tex.sprint(string.char(92) .. "def" .. string.char(92) .. "hllpdffontdir{")
  tex.sprint(-2,hllpdf_font_dir)
  tex.sprint("}")
  luaotfload.add_fallback("hebfb",{"[" .. hllpdf_font_dir .. "DejaVuSerif.ttf]:mode=node;"})
  luaotfload.add_fallback("hebfbbf",{"[" .. hllpdf_font_dir .. "DejaVuSerif-Bold.ttf]:mode=node;"})
}
% </hllpdf:font-binding>
% v3.22: Frank Ruhl Libre ships NO italic face. Instead of silently-upright
% \emph (v<=3.20) or bold-as-emphasis (v3.21), synthesize a slanted face with
% FakeSlant — real, VISIBLE italic emphasis with correct \emph nesting
% semantics (inner \emph toggles back upright). Verified empirically at
% 300dpi: slant clearly visible, 0 errors, 0 missing glyphs.
\babelfont{rm}[Path={\hllpdffontdir},RawFeature={fallback=hebfb},
  BoldFont={FrankRuhlLibre-Bold.ttf},
  BoldFeatures={RawFeature={fallback=hebfbbf}},   % v3.56 — see hebfbbf above
  ItalicFont={FrankRuhlLibre-Regular.ttf}, ItalicFeatures={FakeSlant=0.22},
  % Bold-italic uses the weight-matched fallback and a synthetic slant.
  BoldItalicFont={FrankRuhlLibre-Bold.ttf}, BoldItalicFeatures={RawFeature={fallback=hebfbbf},FakeSlant=0.22}
  ]{FrankRuhlLibre-Regular.ttf}
% Sans fallbacks preserve family and weight for the few missing codepoints.
\directlua{luaotfload.add_fallback("hebfbsans",{"[" .. hllpdf_font_dir .. "DejaVuSans.ttf]:mode=node;"})}
\directlua{luaotfload.add_fallback("hebfbsansbf",{"[" .. hllpdf_font_dir .. "DejaVuSans-Bold.ttf]:mode=node;"})}
\babelfont{sf}[Path={\hllpdffontdir},RawFeature={fallback=hebfbsans},
  BoldFont={Heebo-Bold.ttf},
  BoldFeatures={RawFeature={fallback=hebfbsansbf}}]{Heebo-Regular.ttf}
% v3.24: Scale=MatchLowercase — DejaVu Sans Mono has an outsized x-height, so
% unscaled inline code visually SHOUTS next to the delicate Frank Ruhl body
% (measured on a 197-page build: scaled, a code atom sits in the line at the
% body's weight, and the document tightened by ~5%). Applies to listings too.
\babelfont{tt}[Path={\hllpdffontdir},RawFeature={fallback=hebfb},Scale=MatchLowercase,
  BoldFont={DejaVuSansMono.ttf},
  BoldFeatures={RawFeature={fallback=hebfb},FakeBold=1.5},
  ItalicFont={DejaVuSansMono.ttf},
  ItalicFeatures={RawFeature={fallback=hebfb},FakeSlant=0.2},
  BoldItalicFont={DejaVuSansMono.ttf},
  BoldItalicFeatures={RawFeature={fallback=hebfb},FakeBold=1.5,FakeSlant=0.2}
  ]{DejaVuSansMono.ttf}

% Explicit multilingual/English profiles use real Latin faces and italics.
% The legacy/default Hebrew profile retains its established measured typography.
% Latin Modern is provided by the TeX distribution, not the pinned font bundle;
% recorder/runtime evidence reports it as an additional toolchain resource.
\ifx\lakitlanguageprofile\hllpdfhebrew\else
  \babelfont[english]{rm}[Ligatures=TeX,RawFeature={fallback=hebfb},
    BoldFeatures={RawFeature={fallback=hebfbbf}},
    BoldItalicFeatures={RawFeature={fallback=hebfbbf}}]{Latin Modern Roman}
  \babelfont[english]{sf}[Ligatures=TeX,RawFeature={fallback=hebfbsans},
    BoldFeatures={RawFeature={fallback=hebfbsansbf}},
    BoldItalicFeatures={RawFeature={fallback=hebfbsansbf}}]{Latin Modern Sans}
\fi

\usepackage{amsmath,amssymb,mathtools}
% amsthm supplies theorem/proof machinery. Environment names remain
% document-specific; the generic proof label is localized below.
\usepackage{amsthm}
% Proof names follow the surrounding prose, including return from an island.
\addto\extrashebrew{\renewcommand{\proofname}{הוכחה}}
\addto\extrasenglish{\renewcommand{\proofname}{Proof}}
\ifx\lakitmainlanguage\hllpdfenglish
  \renewcommand{\proofname}{Proof}
\else
  \renewcommand{\proofname}{הוכחה}
\fi
% v3.24.1: equation tags under babel-bidi. A tag re-emitted by mathtools'
% showonlyrefs inside align inherits the RTL base and its parens MIRROR —
% ")2.1(" instead of "(2.1)" (plain align/equation tags are unaffected;
% verified on a minimal case + a 198-page build). LTR-isolate at the tag-BOX
% level: mathtools' show/hide + numbering logic is untouched. Do NOT override
% \tagform@ instead — that clobbers the showonlyrefs mechanism (measured:
% every tag shows, numbering breaks).
% v3.28.1: the v3.24.1 wrap must apply ONLY to tags of display environments.
% Inline \eqref in Hebrew PROSE also goes through \maketag@@@, and there the
% english-wrap MIRRORS the parens — ")1(" — while the unwrapped original is
% correct (proven on a 2x2 isolation matrix: patch x showonlyrefs; regressed
% every prose \eqref since v3.24.1, caught by the render gate's R4). The two
% call sites are indistinguishable by \ifmmode (tag boxes are text mode), so
% the display environments themselves raise a flag; \eqref in prose sees it
% false and gets the original untouched \maketag@@@.
\usepackage{etoolbox}

% babel bidi=basic (including upstream 26.9) assumes every automatic
% discretionary has a one-glyph replacement. LuaTeX can legally produce one
% with replace=nil after a hyphenated Latin value; Babel then indexes nil and the
% whole build dies. Do the narrow upstream-compatible correction: while Babel
% classifies the list, present only that malformed disc as an ordinary
% discretionary, then restore its subtype. Babel still reorders the ENTIRE box;
% there is no pcall and no successful fallback to an unreordered node list.
% If Babel throws for any other reason, the build remains fatal. The synthetic
% nil-disc probe and the real table gallery in test_all.sh pin both layers.
\AtBeginDocument{\directlua{
  if Babel and Babel.bidi and not Babel.heb_nil_disc_compat then
    Babel.heb_nil_disc_compat = true
    local original = Babel.bidi
    local disc_id = node.id('disc')
    Babel.bidi = function(head, ispar, hdir)
      local hidden = {}
      for item in node.traverse_id(disc_id, head) do
        if item.subtype == 2 and (not item.replace or not item.replace.char) then
          table.insert(hidden, item)
          item.subtype = 0
        end
      end
      local result = original(head, ispar, hdir)
      for _, item in ipairs(hidden) do item.subtype = 2 end
      return result
    end
  end
}}
\makeatletter
\newif\ifHEBdisptag
\let\ORIGmaketag@@@\maketag@@@
\def\maketag@@@#1{\ifHEBdisptag\ORIGmaketag@@@{\foreignlanguage{english}{#1}}\else\ORIGmaketag@@@{#1}\fi}
\@for\HEB@env:=equation,equation*,align,align*,gather,gather*,multline,multline*,flalign,flalign*,alignat,alignat*,eqnarray,eqnarray*\do{%
  \expandafter\AtBeginEnvironment\expandafter{\HEB@env}{\HEBdisptagtrue}}
\makeatother
% English page + Hebrew prose: Babel's equation-only tag flip uses a primitive
% leqno path that shifts the formula to the left edge and the tag into the body
% of an English page. Align's full-width alignment is unaffected. Give only this
% display a physical LTR container and native amsmath eqno/leqno. The engine then
% owns real tag width, collision handling and labels; no proxy tag box hides ink
% from its fit calculation. Babel's saved \bbl@eqnodir still restores Hebrew text
% inside math, and the prose locale/math direction are never switched.
% Close the prior RTL paragraph before normalization. After the environment's
% group restores its directions, close TeX's empty resumed display paragraph and
% use native LaTeX continuation handling: same-paragraph prose has no indent,
% while an authored blank line starts an ordinary indented RTL paragraph.
% Register after Babel's setup; remain local to equation(*) in a main-English
% document's actual RTL prose. Native Hebrew, ordinary English and explicit
% fleqn retain their original paths. No page/display dimensions are assigned.
% <hllpdf:minority-rtl-equation-geometry>
\makeatletter
\newcommand{\HEB@minorityprosecontinuation}{\par\@doendpe}
\newcommand{\HEB@minorityequationgeometry}{%
  \ifx\lakitmainlanguage\hllpdfenglish
    \ifnum\pardirection=\@ne
      \if@fleqn\else
        \par\aftergroup\HEB@minorityprosecontinuation
        \iftagsleft@\let\veqno\@@leqno
        \else\let\veqno\@@eqno\fi
        \bodydir TLT\pardir TLT\textdir TLT\relax
      \fi
    \fi
  \fi}
\AtBeginDocument{%
  \AddToHook{env/equation/begin}{\HEB@minorityequationgeometry}%
  \AddToHook{env/equation*/begin}{\HEB@minorityequationgeometry}}
\makeatother
% </hllpdf:minority-rtl-equation-geometry>
% v3.63: \shortintertext rendered Hebrew MIRROR-REVERSED and flush LEFT — an entire
% commentary line inside an align printed backwards. babel's bidi=basic patches
% \intertext@ only (luababel.def), and mathtools' \shortintertext uses the identical
% \vbox{\normalbaselines … \noindent\ignorespaces#1\par} skeleton but is never
% patched, so its paragraph keeps pardir/textdir LTR: no bidi reordering (logical
% order on the page) and flush-left. Measured against \intertext with the same
% Applying babel's own replacement to \shortintertext@ makes the two identical.
% \h{}/\mbox fix the letter order but leave the line flush LEFT, so they are not
% substitutes. Guarded + \AtBeginDocument because mathtools \let-copies at load time.
% NB \ExplSyntaxOn is REQUIRED: mathtools' internals carry expl3 names, and without
% the expl3 catcodes the `_` is read as a subscript and the patch detonates far away
% with "! Illegal unit of measure (pt inserted)" at the \end{align} (measured).
\makeatletter
\AtBeginDocument{%
  \@ifundefined{shortintertext@}{}{%
    \ExplSyntaxOn
    \bbl@sreplace\shortintertext@{\normalbaselines}%
      {\normalbaselines\ifx\bbl@eqnodir\relax\else\bbl@pardir\@ne\bbl@eqnodir\fi}%
    \ExplSyntaxOff}}
\makeatother
% geometry inherits the canonical A4 sheet from \documentclass; build.sh then
% verifies every rendered MediaBox and CropBox, so a silent Letter fallback is
% impossible to publish as a clean default build.
\usepackage[\lakitgeometryoptions]{geometry}
\linespread{\lakitlinespread}       % layout hook; canonical default remains 1.3
\ifx\lakitparagraphstyle\hllpdfparagraphblock
  \setlength{\parindent}{0pt}
  \setlength{\parskip}{6pt plus 1pt minus 1pt}
\fi
% These dimensions remain ordinary scoped TeX lengths: authors can customize
% them after core input. Language switches inherit them; lists keep their own
% indentation and narrow reference layouts establish their local paragraph gaps.
\typeout{HLLPAR|\lakitparagraphstyle}
\usepackage{xcolor}
\usepackage{mdframed}
\usepackage{listings}                 % code listings (assembly / C / pseudo-code)
\usepackage{tikz}
\usetikzlibrary{arrows.meta,positioning,calc,shapes.geometric,fit,backgrounds}
\usepackage{pgfplots}\pgfplotsset{compat=1.18}
\usepackage{quantikz}                 % quantum circuits (remove if unused)
\usepackage{booktabs}
\usepackage{array}
\usepackage{float}                    % [H] = place table exactly here (reference doc: no drifting)
\usepackage{enumitem}
\usepackage{needspace}

% ============================================================================
% ★★ v3.4 — THE DIRECTION-ISOLATION TRIO (the heart of mixing Hebrew + Latin)
% Mental model:
%   \h{...}    Hebrew (RTL) island      — Hebrew text inside an LTR context
%                                          (TikZ nodes, the english `tikzpic`).
%   \en{...}   English (LTR) isolate    — Latin text / parentheses inside a
%                                          Hebrew (RTL) context. KEEPS PARENS
%                                          AND ORDER CORRECT.
%   \code{...} \en + monospace          — code inside Hebrew / inside a frametitle.
% CRITICAL: never wrap Latin in \h{} (it is for HEBREW). Inside a TikZ node,
% Latin in \h{} renders REVERSED (Datapath->htapataD, CPU->UPC) and parens
% mirror. Isolate every Latin/paren chunk with \en{} or \code{}. See
% references/figures-boxes-listings.md.
% ---- TEXT INSIDE MATH (v3.8) ----------------------------------------------
% Put ANY text that sits inside math in a \text{}. The ONLY thing that is
% invalid in math is a bare \h/\en (= \foreignlanguage) used DIRECTLY, outside
% \text -> "\rmfamily invalid in math mode". Inside \text:
%   Hebrew label           -> \text{...}            (bare Hebrew renders fine;
%                              \h is optional — BUT NOT ALWAYS HARMLESS (v3.67):
%                              if the \h group OPENS with a nested island
%                              ($..$, \en{}, \code{}, \texttt{}, \mbox{}) the
%                              WHOLE group loses bidi and the Hebrew prints in
%                              LOGICAL order = gibberish. Measured:
%                                \text{\h{$f$ קבועה}}  -> reads "העובק"  BROKEN
%                                \text{\h{קבועה $f$}}  -> correct
%                                \text{$f$ קבועה}      -> correct (no \h)
%                              Hebrew first inside \h, or drop the \h. Flagged
%                              H-OPENS-WITH-LTR-GROUP.)
%   plain Latin word        -> \text{...}            (no marker needed)
%   Latin WITH ( ) or [ ]   -> \text{\en{...}}       (\en forces LTR; a plain
%                              \text mirrors the parens in RTL -> "Spec (TNR(" )
%   f(Hebrew/text arg)      -> \text{f}(\text{...})  -- keep the parens at MATH
%                              level, NOT \text{f(arg)} (parens inside the box mirror)
%   code in math            -> \mathtt{...}  (or \texttt{}, math-guarded below)
% NEVER write \en{...} or \h{...} straight into $...$ / \[..\] (outside \text);
% always via \text{}. (\mbox{...} also works but does not scale in scripts.)
% ============================================================================
\newcommand{\h}[1]{\foreignlanguage{hebrew}{#1}}
% v3.42.1 — \en islands do NOT hyphenate (auto OR at explicit hyphens). English hyphenation patterns were
% tearing identifiers and terms ACROSS THE RTL SEAM ("IO Regis-/ters",
% "In-/put", "Band-/width", "Per-/formance", "Multi-/ple" — all shipped on
% real pages; probe-proven: "ALUS-/rcB", "Mem-/ToReg"). An island is a label
% or a term, not English prose — a hyphen tear there is never right. A/B on
% a full 197pp lakit: identical overfull sets, same page count, ~11
% paragraphs healed, one mild underfull (badness 3525) that replaced a
% mutilated "Regis-/ters". Long English PROSE belongs in the englishprose
% environment below, which restores explicit-hyphen breaks while its paragraph
% is still active. A raw otherlanguage block inherits the global prohibition.
% (\l@nohyphenation needs \csname:
% '@' is not a letter here.)
% v3.43.2 — \exhyphenpenalty moved OUT of \en to a GLOBAL assignment below.
% Render-proven: a group-local \exhyphenpenalty is DEAD CODE — TeX reads the
% line-breaking parameters at \par time, after the \en group has closed, so
% \en{write-through} / plain not-taken still tore at the "-" ("write-/
% through" shipped on a real exam page). \language is different: it is
% recorded in the node list as text is added, so the in-group nohyphenation
% assignment above DOES work (and stays). The global penalty was A/B-measured
% on a 197pp lakit: identical overfull set (25), identical page count, the
% not-/taken tear healed. Hebrew is unaffected in practice — Hebrew compounds
% use maqaf (־), not ASCII "-".
\newcommand{\en}[1]{\foreignlanguage{english}{\language=\csname l@nohyphenation\endcsname #1}}
% Babel runs extras locally for inline islands and paragraph language changes.
% An English main document must not inherit the Hebrew no-hyphen paragraph rule.
\addto\extrashebrew{\exhyphenpenalty=10000\relax}
\addto\extrasenglish{\exhyphenpenalty=50\relax}
\ifx\lakitmainlanguage\hllpdfenglish
  \exhyphenpenalty=50\relax
\else
  \exhyphenpenalty=10000\relax
\fi
% Full English paragraphs need ordinary explicit-hyphen breaks. Use this
% environment instead of a raw otherlanguage block; the paragraph ends while
% the local penalty is still active.
\newenvironment{englishprose}
  {\par\begin{otherlanguage}{english}\exhyphenpenalty=50\relax}
  {\par\end{otherlanguage}}
\newenvironment{hebrewprose}
  {\par\begin{otherlanguage}{hebrew}\exhyphenpenalty=10000\relax}
  {\par\end{otherlanguage}}
% Use for shell text whose locale must not depend on the last body paragraph.
\DeclareRobustCommand{\doclang}[1]{\foreignlanguage{\lakitmainlanguage}{#1}}

% ============================================================================
% ★★ — fixes for defects that are INVISIBLE in the log and in the checks
% and only surface on a VISUAL RENDER of the PDF. (Each verified empirically.)
% ----------------------------------------------------------------------------
% v3.22: \emph is now REAL slanted emphasis (FakeSlant italic face defined on
% the \babelfont{rm} line above), replacing the v3.21 \emph->\textbf hack: it
% preserves \emph's toggle semantics and keeps bold free for its own role.
% Prefer bold-as-emphasis for your doc? one line: \renewcommand{\emph}[1]{\textbf{#1}}
% page number: force the digits LTR so they never reverse (18 -> 81) in the RTL
%   footer. That reversal is a NON-DETERMINISTIC bidi glitch (observed once in
%   ~80 pp); this LuaLaTeX build is not byte-reproducible, so the defensive
%   LTR wrap is the only reliable guard. Latin digits are always LTR anyway.
\renewcommand{\thepage}{\foreignlanguage{english}{\arabic{page}}}
% v3.70: \pagenumbering \gdef's \thepage and silently DESTROYS the wrap above —
% after \pagenumbering{arabic} the folio guard is gone for the whole rest of the
% document. Two collections discovered this independently and hand-restore the
% wrap in main.tex with the same comment; that knowledge now lives here instead.
% The wrapper re-applies the LTR wrap after the kernel call, for any style.
\makeatletter
\let\lakin@orig@pagenumbering\pagenumbering
\renewcommand{\pagenumbering}[1]{%
  \lakin@orig@pagenumbering{#1}%
  \renewcommand{\thepage}{\foreignlanguage{english}{\csname @#1\endcsname\c@page}}}
\makeatother
% emergencystretch: Hebrew does not hyphenate, so an unbreakable LTR island
%   (\en{...} / an inline math run) landing at the margin can leave a small
%   (~2-13pt) overfull with no legal breakpoint. This rescues ONLY those
%   inline-prose tie-breaks; it is a ceiling used only on offending lines and
%   does NOT touch display-math / listing / table overflow (those still surface
%   in the log for the content-level overflow pass). Verified: baseline overfull
%   -> 0 with this set, with no change to pagination or underfull.
\setlength{\emergencystretch}{3em}
% v3.40: inline math (almost) never breaks across lines. TeX's default
%   (binop 700 / rel 500) lets an inline run split at an operator —
%   $(1.f)\cdot 2^{e}$ shipped as "(1.f)·" ending one line and "2^e" opening
%   the next (reader-reported). In LTR that split is merely ugly; under an
%   RTL base the carried operator lands on the visually WRONG side of the
%   break and the formula reads backwards.
%   WHY 10000 (v3.45.0 — the display-first resolution). Inline math is
%   FORBIDDEN from breaking, so a torn island under an RTL base is
%   structurally impossible. The v3.40 cycle used 9999 ("break only as a last
%   resort") because a bare 10000 left ~30-39 WIDE inline equations
%   overflowing (worst 90.7pt) — but 9999's "last resort" fires OFTEN in
%   tight, non-hyphenating Hebrew: a reader found torn islands across the
%   fleet (~150 candidates; e.g. "P(y |" ending one line and "x_0)" opening
%   the next, read backwards). The step the 9999 cycle missed: every one of
%   those ~30 overflows is a wide equation that BELONGS in display. Convert
%   them to display FIRST, then 10000 gives 0 tears AND 0 overflows (measured,
%   12/12 fleet). Doctrine: inline math is ATOMIC; a wide equation goes in a
%   display, never inline. See content-style ("inline math is atomic").
\binoppenalty=10000 \relpenalty=10000
% v3.38: uniform sentence spacing. TeX's default (\nonfrenchspacing, sfcode
% 3000) is the OLD Anglo-American wide-space-after-period convention — Hebrew
% typography never had it, and modern Latin practice (Chicago et al.) dropped
% it too. Measured on real pages: post-period gaps up to 1.9-2.1x the line's
% mean word gap — the dominant "rivers" source in justified Hebrew prose.
% Survives babel language switches (probed: sfcode stays french inside \en{}
% islands and otherlanguage blocks).
\frenchspacing
% v3.38: widow/orphan protection. LaTeX's default penalties (150/150/50)
% allow a lone line at a page top/bottom almost freely. 10000 = forbid;
% this is SAFE here precisely because the document is \raggedbottom (the
% demoted line just moves on, no vertical stretch fight) — under flushbottom
% these values would cause underfull pages. raggedbottom is the report-class
% default; made explicit so a future class/option change cannot silently
% flip it and turn the 10000s harmful.
\raggedbottom
\widowpenalty=10000 \clubpenalty=10000 \displaywidowpenalty=10000
% em-dash: --- does NOT become — under Frank Ruhl (the font lacks tlig; verified
%   the width is identical with/without Ligatures=TeX, while the mechanism works
%   on Latin Modern). TYPE THE CHARACTER — (U+2014) DIRECTLY; likewise – (U+2013).
%   check_bidi_figures.py flags a --/--- that is adjacent to Hebrew.
% ============================================================================

% ---- STRUCTURAL DIRECTION-SAFETY for \texttt (v3.7) -------------------------
% A parenthesis, bracket, or $ INSIDE \texttt{} mirrors/migrates in RTL prose
% (MIPS 0($s0)->"(s0$)0", $rs->"rs$", arr[i]->"[arr[i") because \texttt sets the
% font but NOT the direction. Redefining \texttt to LTR-isolate its content fixes
% the WHOLE class at once (parens, brackets, $) so plain \texttt{lw $t0,0($s0)}
% is safe — no need to remember \code{}. Guarded by \ifmmode because
% \foreignlanguage is invalid in math (there, plain monospace is already LTR).
% NB (corrected v3.56 — the old note here said "stays tofu", which is WRONG):
% \texttt{Hebrew} is a FATAL build failure. The tt family is DejaVu Sans Mono
% (zero Hebrew) and the tt fallback hebfb is DejaVuSerif (also zero Hebrew), so
% the subsetter gets an empty glyph set and LuaTeX dies with
%   "! error: (file .../DejaVuSansMono.ttf) (type 2): there are no glyphs in the
%    subset"  -> no output PDF.
% The message names a FONT FILE, never Hebrew/\texttt/the source line. And it is
% non-local: add an unrelated Latin \code{} elsewhere and the fatal degrades into
% a SILENT DELETION of the Hebrew word. Put Hebrew in \h{}/normal text, never in
% \texttt{}/\code{}; flagged HEBREW-IN-MONO. To disable, delete this block.
\let\origtexttt\texttt
\protected\long\def\texttt#1{\ifmmode\origtexttt{#1}\else\foreignlanguage{english}{\origtexttt{#1}}\fi}
% v4.8.0: defer the direction test until TeX has inserted an alignment cell's
% math template. Unprotected expansion here tested \ifmmode too early at the
% FIRST token of aligned/array/cases/matrix cells. Engine-level \protected also
% preserves the original \texttt implementation without the self-recursion
% observed when retrofitting \DeclareRobustCommand to its saved alias.
\protected\long\def\code#1{\ifmmode\origtexttt{#1}\else\foreignlanguage{english}{\origtexttt{#1}}\fi}

% ---- \codeline — display-code line (v3.24): the code analogue of display math.
% Reader-friction rule (content-style.md §5b): inline \code{} is for a SINGLE
% ATOM only (one register/variable/mnemonic, a tight no-space expression).
% A multi-token expression — an instruction with operands, a spaced assignment,
% a register-binding list — zigzags the eye RTL<->LTR mid-sentence and, being
% an unbreakable island, line-breaks badly. Give it its own centered line with
% display-math vertical rhythm. \centerline = one box => the expression never
% wraps; \nopagebreak keeps it glued to its introducing sentence.
%   ...the Hebrew sentence ends with a colon:
%   \codeline{lw \$t1, -4(\$t0)}
%   and the prose continues (escape dollars as \$).
% v3.66 (first line of the body): flush a PENDING \item label before the \par.
% LaTeX defers the label box (\@labels) until the next \everypar, and an \hbox in
% vertical mode never fires \everypar — so `\item \codeline{...}` printed the code
% at y=69.1 and the item's own "1." BELOW it at y=103.6 (the number appears after
% its content, attached to nothing). \if@inlabel\leavevmode\fi is CONDITIONAL on
% purpose: an unconditional \leavevmode would open an empty paragraph line when
% \codeline follows ordinary prose. Verified both ways — in a list the label now
% sits on the code's line at its normal x, and after prose the geometry is
% byte-identical to before.
\makeatletter
\newcommand{\codeline}[1]{%
  \if@inlabel\leavevmode\fi
  \par\nopagebreak
  \vspace{\abovedisplayskip}\nopagebreak
  % v3.62: was \centerline{...}, i.e. \hbox to\hsize{\hss #1 \hss}. \hss has
  % INFINITE SHRINK, so TeX could always "fit" the box and NEVER reported an
  % Overfull — the same \hss mechanism as BLD-27's \item label. Measured: a wide
  % \codeline in a defbox painted its first token ON the box's left border bar and
  % sat 9.5pt outside the frame interior on both sides, with 0 Overfull, 0 Underfull
  % and all four gates PASS; in bare prose it reached 21.0pt into BOTH margins.
  % \hfil has no shrink, so the overflow is reported and build.sh's ovf_vis/ovf_big
  % counters see it; \linewidth (not \hsize) makes it box-aware, so a \codeline
  % inside a colored box centres on the box interior. Verified: a FITTING \codeline
  % keeps byte-identical glyph positions, and the box-crossing case now emits
  % "Overfull \hbox (19.03pt too wide)". A wide \codeline should be split in two.
  \hbox to\linewidth{\hfil\code{#1}\hfil}%
  \nopagebreak\vspace{\belowdisplayskip}}
\makeatother

% ---- a thin decorative rule helper ----
\newcommand{\sep}{\par\vspace{2pt}\noindent\textcolor{cRule}{\rule{\linewidth}{0.6pt}}\par\vspace{4pt}}

% ---- colors ----
\definecolor{cBlue}{HTML}{1F4E79}   \definecolor{cBlueBg}{HTML}{EAF1F8}
\definecolor{cGreen}{HTML}{2E6B4F}  \definecolor{cGreenBg}{HTML}{EAF4EE}
\definecolor{cAmber}{HTML}{9C6A00}  \definecolor{cAmberBg}{HTML}{FBF3E2}
\definecolor{cPurple}{HTML}{5B2C6F} \definecolor{cPurpleBg}{HTML}{F3ECF7}
\definecolor{cTeal}{HTML}{0E6E6E}   \definecolor{cTealBg}{HTML}{E6F4F4}
\definecolor{cRed}{HTML}{8E2B2B}    \definecolor{cRedBg}{HTML}{FBEAEA}
\definecolor{cGray}{HTML}{555555}   \definecolor{cRule}{HTML}{1F4E79}

% ---- math operators ----
\DeclareMathOperator{\Tr}{Tr}
\DeclareMathOperator{\rank}{rank}
\DeclareMathOperator{\diag}{diag}
\DeclareMathOperator{\sgn}{sgn}

% ---- bra-ket (quantikz defines \ket/\bra/\braket; add the rest safely) ----
\providecommand{\ket}[1]{\left|#1\right\rangle}
\providecommand{\bra}[1]{\left\langle#1\right|}
\providecommand{\braket}[2]{\left\langle#1\middle|#2\right\rangle}
\newcommand{\dyad}[1]{\left|#1\right\rangle\!\left\langle#1\right|}
% v3.21: general outer product |a><b| and expectation value <A> — agents reach
% for these by reflex (esp. \expval); quantikz defines neither, so their absence
% threw "Undefined control sequence" in a real build. \providecommand = safe.
\providecommand{\ketbra}[2]{\left|#1\right\rangle\!\left\langle#2\right|}
\providecommand{\expval}[1]{\left\langle#1\right\rangle}
\newcommand{\abs}[1]{\left|#1\right|}
\newcommand{\norm}[1]{\left\lVert#1\right\rVert}
\newcommand{\half}{\tfrac12}
\newcommand{\R}{\mathbb{R}}\newcommand{\C}{\mathbb{C}}\newcommand{\Z}{\mathbb{Z}}
% identity operator. NB: \mathbb{1} silently renders a WRONG glyph (msbm has no
% blackboard digits). Use \mathbb{I} = 𝕀 (works with amssymb, no extra package).
% Do NOT "fix" this back to \mathbb{1}.
\newcommand{\id}{\mathbb{I}}

% ---- physical units (v3.35): \punit — upright unit with thin space ---------
% Usage: $300\punit{K}$, $80\punit{mV/decade}$, $10^{7}\punit{cm/s}$ — the
% NUMBER AND UNIT LIVE IN ONE MATH ISLAND. Promoted verbatim from a 199-page
% production lakit (advanced-electronic-devices, dozens of uses).
% Two measured traps this design avoids (both render-proven):
%   1. siunitx SCRAMBLES under RTL — \SI{3.5}{eV} renders "eV 3.5", exponents
%      detach ("Hz 10 x 10^9"): its quantity assembly is not bidi-isolated.
%      Do NOT load siunitx in this pipeline.
%   2. A text-mode unit next to a bare-prose number is TWO adjacent LTR runs
%      and reorders ("5 mm" -> "mm5" — the BIDI-2/3 class). \punit is therefore
%      math-only ON PURPOSE: prose misuse fails LOUDLY at compile (\mathrm
%      outside math) instead of scrambling silently.
% Domain shorthands (\MeV, \fm...) stay per-document: \newcommand{\MeV}{\,\text{\en{MeV}}}
\providecommand{\punit}[1]{\,\mathrm{#1}}

% ============================================================================
% Colored boxes — mdframed (NOT tcolorbox, which breaks under bidi).
% ★★ v3.4: innerleftmargin/innerrightmargin give code listings room so their
% frame/line-numbers do not spill out the box border (see listing styles below).
% ============================================================================
% v3.61: everyline/split skips were applied to `tablebox` only (BLD-14, v3.37) —
% the six colored boxes below never got them, so they still split OPEN across a page
% break: the page-N segment gets top+sides and NO bottom rule, the page-N+1 segment
% sides+bottom and NO top rule. Measured on the shipped fleet this is a COVERAGE GAP,
% not a re-report of BLD-14: 0 split-open `tablebox` segments (that fix holds) vs 62
% split-open COLORED boxes across 7 documents. Verified after the fix: both segments
% close, and the frametitle is NOT duplicated on the continuation. BLD-14's caveat
% still applies — a box taller than a full page grinds mdframed's split-trial loop, so
% split such content manually.
% ============================================================================
% v3.72: frametitle={} leads EVERY \newmdenv option list below. mdframed options
% are GROUP-scoped: a box nested inside \begin{exbox}[frametitle={...}] used to
% INHERIT the outer title — measured, the identical title string rendered twice,
% the second riding a title-bar fill inside the INNER box's frame, for all seven
% environments in both directions. Both natural forms were broken (= preamble
% bug, not an authoring rule). The explicit reset makes every box titleless
% unless ITS OWN [frametitle={...}] says otherwise; verified both directions:
% no duplicate on the nested box, an explicit title still renders.
% Geometry is shared with the optional semantic-width table container. Themes
% select their own existing semantic style below; neither path fixes columns.
\mdfdefinestyle{hllpdfsemanticgeometry}{linewidth=1.1pt,
  innertopmargin=6pt,innerbottommargin=7pt,skipabove=9pt,skipbelow=9pt,
  innerleftmargin=9pt,innerrightmargin=9pt,everyline=true,
  splitbottomskip=6pt,splittopskip=12pt,roundcorner=3pt}
\newcommand{\hllpdfsemanticframestyle}{hllpdfsemanticgeometry}
\newmdenv[frametitle={},style=hllpdfsemanticgeometry,linecolor=cBlue,backgroundcolor=cBlueBg,
  frametitlebackgroundcolor=cBlue,frametitlefontcolor=white,
  ]{defbox}   % הגדרה
\newmdenv[frametitle={},style=hllpdfsemanticgeometry,linecolor=cGreen,backgroundcolor=cGreenBg,
  frametitlebackgroundcolor=cGreen,frametitlefontcolor=white,
  ]{thmbox}   % משפט/טענה
\newmdenv[frametitle={},style=hllpdfsemanticgeometry,linecolor=cAmber,backgroundcolor=cAmberBg,
  frametitlebackgroundcolor=cAmber,frametitlefontcolor=white,
  ]{notebox}  % הערה
\newmdenv[frametitle={},style=hllpdfsemanticgeometry,linecolor=cPurple,backgroundcolor=cPurpleBg,
  frametitlebackgroundcolor=cPurple,frametitlefontcolor=white,
  ]{exbox}    % דוגמה
\newmdenv[frametitle={},style=hllpdfsemanticgeometry,linecolor=cRed,backgroundcolor=cRedBg,
  frametitlebackgroundcolor=cRed,frametitlefontcolor=white,
  ]{warnbox}  % טעות נפוצה
\newmdenv[frametitle={},style=hllpdfsemanticgeometry,linecolor=cTeal,backgroundcolor=cTealBg,
  frametitlebackgroundcolor=cTeal,frametitlefontcolor=white,
  ]{keybox}   % רעיון מרכזי

% ---- codebox (v4.5.0): the dedicated container for code -------------------
% An OPTION, not a policy: asmblock unboxed stays canonical for a full listing,
% and a snippet inside any of the six colored boxes still works as before.
% What makes codebox worth its own environment is that it is the ONE framed box
% that keeps its margin line numbers.
%
% v4.4.0 shipped it WITHOUT them, on the strength of v3.37's "no margin knob can
% save them" plus four fresh probes that agreed. That was wrong: every probe had
% moved the code, the box, or both — never the key that carries the NUMBER
% COLUMN. Measured on the page:
%   resetmargins=false   the listing follows THIS box instead of the page,
%                        moving the code from x=75.26 to x=97.54
%   xrightmargin=26pt    under RTL this is the PHYSICAL LEFT indent, and the
%                        only key that carries the \llap'ed digits with it:
%                        x=64.28 -> 90.18, clearing the rule at 65.64 by 23.6pt
%   xleftmargin=26pt     moves the code and LEAVES THE DIGITS on the rule —
%                        the asymmetry that hid the fix for eleven versions
% resetmargins is hoisted out of the four styles (see the \lstset below them) so
% this override is possible at all; every other listing still resolves to true.
% Pinned by tests/codebox_numbers.tex + probe_lstnum.py --inside-rule, whose
% negative control strips xrightmargin and watches the digits fall back.
\definecolor{cSlate}{HTML}{3F4A56}  \definecolor{cSlateBg}{HTML}{F2F4F7}
\newmdenv[frametitle={},leftline=true,rightline=true,topline=true,bottomline=true,
  linecolor=cSlate,backgroundcolor=cSlateBg,linewidth=0.9pt,
  frametitlebackgroundcolor=cSlate,frametitlefontcolor=white,
  innertopmargin=6pt,innerbottommargin=7pt,skipabove=9pt,skipbelow=9pt,
  innerleftmargin=9pt,innerrightmargin=9pt,everyline=true,splitbottomskip=6pt,splittopskip=12pt,
  roundcorner=3pt]{codebox}   % קוד
% The two keys that make in-box numbering work, both measured on the page:
%   resetmargins=false  the listing follows THIS box instead of the page
%   xrightmargin        under RTL this is the PHYSICAL LEFT indent, and it is
%                       the only one that carries the \llap'ed number column
%                       with it. xleftmargin moves the code and leaves the
%                       digits behind on the rule (measured: code 97.54,
%                       digits 64.28 — the v3.37 defect, reproduced); with
%                       xrightmargin=26pt the digits land at x=90.18, a
%                       clearance of 23.6pt from the rule at 65.64.
% This asymmetry is why v3.4's note had the semantics backwards and v3.37
% concluded the digits could not be saved.
% <hllpdf:listing-container-state>
% v4.29: all state is local to the current environment group. A semantic
% callout masks number text, while a nested codebox restores the formatter
% saved before that mask. Save it once across successive semantic wrappers;
% an explicit numbers=none (or an explicitly empty formatter) remains valid.
\makeatletter
\newif\ifhllpdf@incodebox
\newif\ifhllpdf@englishcodebox
\newif\ifhllpdf@lstnumbermasked
\let\hllpdf@savedlstnumber\thelstnumber
\makeatother
% </hllpdf:listing-container-state>
% <hllpdf:codebox-gutter>
\makeatletter
\AtBeginEnvironment{codebox}{%
  \hllpdf@incodeboxtrue
  % Snapshot the OUTER language before asmblock selects English.
  \iflanguage{english}{\hllpdf@englishcodeboxtrue}{\hllpdf@englishcodeboxfalse}%
  \ifhllpdf@lstnumbermasked
    \let\thelstnumber\hllpdf@savedlstnumber
    \hllpdf@lstnumbermaskedfalse
  \fi
}
% listings' AddToHook is GLOBAL. Register once, and gate on local box state;
% adding an unconditional hook at every codebox entry leaks into later listings.
\lst@AddToHook{PreInit}{%
  \ifhllpdf@incodebox
    % Apply every listing override here, inside the LISTING's own group.
    % Setting margins at box entry leaks them into nested semantic callouts.
    \lstset{resetmargins=false,xrightmargin=26pt}%
    % Styles resolve before PreInit. In LTR the physical number gutter follows
    % xleftmargin; in Hebrew retain the measured inherited xrightmargin path.
    \ifhllpdf@englishcodebox
      \lstset{xleftmargin=26pt,xrightmargin=0pt}%
    \fi
    % ONE box, not two: suppress only the frame/background/skips of a listing
    % inside this codebox. The next unboxed listing keeps its own style.
    \lstset{frame=none,backgroundcolor=\color{cSlateBg},aboveskip=0pt,belowskip=0pt}%
  \fi
}
\makeatother
% </hllpdf:codebox-gutter>

% v3.37: line numbers inside a COLORED BOX are sliced by the box border —
% measured: digit span x=63.5-67.3 vs border line at x=65.2 (the border runs
% THROUGH the digits), and the inter-border gap (3.5pt) is narrower than a
% digit (3.8pt), so no margin knob can save them (under resetmargins,
% xleftmargin is inert in-box — measured: code at x=75.26 with 2.2em AND
% 3.0em). A 3-6 line boxed snippet does not need line numbers; blank the
% number TEXT inside every colored box (\thelstnumber is what gets printed;
% the style's numbers=left survives, so nothing else shifts — code position
% byte-frozen at x=75.26, verified). Unboxed listings keep their numbers.
% v3.78: the blanking is coupled to the FRAME, not to the box name —
% clean_theme.tex registers LATER hooks (registration order wins) that
% restore the number text inside the boxes it renders frameless.
% v4.5.0 RETRACTS the "no margin knob can save them" clause above. It holds only
% under resetmargins=true (the shipped default): there xleftmargin is inert
% in-box, as measured. With resetmargins=false, xrightmargin DOES carry the
% \llap'ed digit column clear of the rule — measured; the codebox \lstset that
% does it is the one above, in the codebox block. The loop below therefore
% covers SIX boxes, not "every colored box": codebox opts out and keeps its
% numbers.
% <hllpdf:semantic-listing-mask>
\makeatletter
\@for\HEB@lstbox:=defbox,thmbox,notebox,exbox,warnbox,keybox\do{%
  \expandafter\AtBeginEnvironment\expandafter{\HEB@lstbox}{%
    \ifhllpdf@lstnumbermasked\else
      \let\hllpdf@savedlstnumber\thelstnumber
    \fi
    \hllpdf@lstnumbermaskedtrue
    \def\thelstnumber{}%
    % The nearest semantic container owns its listing presentation. On exit,
    % TeX's group restores the enclosing codebox state and formatter.
    \hllpdf@incodeboxfalse}}
\makeatother
% </hllpdf:semantic-listing-mask>

% ============================================================================
% Unified TABLE presentation: subtle uniform slate frame + auto-numbered
% caption rendered INSIDE the frame. Use:
%   \begin{tablebox}\centering  <tabular>  \tabcap{<description>}\end{tablebox}
% Replaces the buggy \caption* (which emitted a stray "טבלה N : *" line) and
% provides an explicit framed option + numbered caption "טבלה N.M: תיאור".
% Optional intro text may be placed at the TOP, before the tabular, inside the box.
% ============================================================================
\definecolor{cTblLine}{HTML}{B9C4D2}  \definecolor{cTblBg}{HTML}{FAFBFC}
% v3.23: class-aware — under `article` there is no chapter counter (the promised
% report->article switch used to die with "No counter 'chapter' defined").
\makeatletter
\@ifundefined{c@chapter}{%
  % v3.31: in article the counter resets per SECTION, so an unprefixed number
  % repeated ("טבלה 1." in every section — measured on a two-section render).
  % Same conditional idea as the report branch: prefix the section number once
  % a real \section exists; a pre-section table stays "טבלה 1".
  \newcounter{tbl}[section]
  \renewcommand{\thetbl}{\ifnum\value{section}>0 \thesection.\fi\arabic{tbl}}%
}{%
  \newcounter{tbl}[chapter]
  % v3.30: a report-class doc that uses SECTIONS ONLY (chapter stays 0 —
  % the common short-doc pattern) used to print "טבלה 0.1"; the chapter
  % prefix now appears only once a real \chapter exists (render-verified
  % both ways: "טבלה 1." chapterless, "טבלה 1.1" with chapters).
  \renewcommand{\thetbl}{\ifnum\value{chapter}>0 \thechapter.\fi\arabic{tbl}}%
  % v3.64: the SAME leak in \thesection itself, left behind by v3.30. That fix
  % guarded the CAPTION counter and stopped there, so a chapterless report still
  % printed the heading "0.1 סעיף ראשון" and \ref to it printed "0.2" — measured:
  %   p1 y=65.2  '…0.1'   p1 y=144.9 '…0.2'   p1 y=183.4 '…0.2.1'
  % report.cls hardcodes \thesection as \thechapter.\@arabic\c@section, and the
  % leading 0 is the chapter counter that no \chapter ever incremented. Fixing
  % \thesection fixes the whole chain — \thesubsection and \thesubsubsection are
  % defined in terms of it. \theequation and \thetbl/\thefig were already guarded,
  % which is exactly why this survived: every NUMBER around the heading looked
  % right, so the heading's own 0 read as a style choice rather than a bug.
  % \appendix is safe: report.cls resets chapter to 0 and redefines \thechapter to
  % \@Alph\c@chapter, so an appendix \chapter makes the counter 1 and the guard
  % emits "A.1" unchanged. Render-verified in all four modes (chapterless,
  % chaptered, appendix+chapter, appendix+section-only).
  \renewcommand{\thesection}{\ifnum\value{chapter}>0 \thechapter.\fi\arabic{section}}%
}
\makeatother
% v3.37: everyline=true — if the box DOES split across pages, every segment
% gets a CLOSED frame (measured before: the split ran open-ended into the page
% edge with no bottom border, and the caption could strand alone at the next
% page top — "טבלה A.2" shipped that way). The split skips give the segments
% breathing room. In practice, with the caption GLUE below, every previously-
% split table in the corpus now moves WHOLE instead (the white-gap tradeoff,
% documented benign); everyline stays as the safety net for a genuinely
% page-tall table. NB: a table taller than ONE FULL PAGE can grind mdframed's
% split-trial loop (reproduced: minimal doc, minutes of 1pt-step Overfull
% probes) — break such a monster into two tables.
\mdfdefinestyle{hllpdftablegeometry}{leftline=true,rightline=true,topline=true,bottomline=true,
  linewidth=0.9pt,
  innertopmargin=9pt,innerbottommargin=8pt,skipabove=11pt,skipbelow=11pt,
  innerleftmargin=11pt,innerrightmargin=11pt,roundcorner=4pt}
\newcommand{\hllpdftableframestyle}{hllpdftablegeometry}
\newmdenv[frametitle={},style=hllpdftablegeometry,
  linecolor=cTblLine,backgroundcolor=cTblBg,
  % v4.2.0 (adopted): tablebox is ATOMIC. Short tableboxes near page bottoms
  % repeatedly entered mdframed's split trial, emitting benign Overfull-vbox
  % bursts and occasionally separating the intro prose from the tabular.
  % nobreak=true moves the complete short table to the next page instead. The
  % output-active Overfull-vbox gate still catches a box taller than a page;
  % everyline/split skips remain as defensive styling should that path ever
  % re-arm, but the canonical path never publishes a split tablebox.
  nobreak=true,
  everyline=true,splitbottomskip=6pt,splittopskip=12pt,
  % v3.79 (BLD-37 table side, correction): startinnercode/endinnercode run
  % INSIDE mdframed's body vbox (mdframed.sty:848 / :856+863) -- the second one
  % after \par\unskip, i.e. after line breaking, which is the only moment the
  % \tabcap anchor's true offset from the top of the box EXISTS. See
  % \lakintabboxclose below. These two keys are the ONLY tablebox change.
  startinnercode={\lakintabboxopen},endinnercode={\lakintabboxclose}]{tablebox}
% v3.37: \nopagebreak GLUE — \vsplit honors penalties, so the caption can no
% longer strand alone at the top of the next page (it takes its rows along, or
% the whole box moves). Render-verified on the real shipped orphan.
% v3.59: \NR@gettitle is what nameref injects into the stdlib title-bearing
% commands (\@sect, \@chapter, \caption) to fill \@currentlabelname. \figcap and
% \tabcap are INDEPENDENT macros — they never called it, so \@currentlabelname kept
% whatever the last heading left and \nameref{fig:x} silently printed the enclosing
% SECTION's title instead of the caption (proven in the .aux: field 3 of \newlabel
% was the heading text while field 1 and the anchor `fig.1.1` were correct — so this
% is a THIRD, distinct head next to BLD-13's counter fix and BIDI-38's autoref-name
% fix, neither of which writes that field). \@ifundefined keeps it inert if
% hyperref/nameref is ever dropped. Guarded helper because \figcap is defined
% after \makeatother.
% NB both helpers below need \makeatletter — this part of the preamble is OUTSIDE
% any letter-catcode region, so an unwrapped \@currentlabel parses as LaTeX's \@
% and dies with "You can't use \spacefactor in vertical mode" (measured).
\makeatletter
\newcommand{\lakinrtitle}[1]{\@ifundefined{NR@gettitle}{}{\NR@gettitle{#1}}}
% v3.59: \hitem{<letter>}{<key>} — a REFERENCEABLE Hebrew-lettered list item.
%   \begin{enumerate} \hitem{א}{it:a} הסעיף הראשון. \end{enumerate}  ->  \ref{it:a} = (א)
% Why this exists: automatic Hebrew list labels fail silently, so SKILL.md mandates
% an EXPLICIT \item[(\textbf{א})] — but LaTeX puts \refstepcounter INSIDE
% \if@noitemarg (latex.ltx:12964), and \item[...] never sets that flag. So an
% explicit item steps no counter and never touches \@currentlabel: a \label there
% binds to the AMBIENT one and \ref prints the CHAPTER/SECTION number where the
% item letter belongs (proven in the .aux: {chapter.1}). Same doctrine as \figcap —
% own the label state before \label. \phantomsection matters: without it the
% hyperlink still lands on the chapter head (measured); with it the anchor is
% section*.N and the bookmark outline stays clean.
\newcommand{\hitem}[2]{\item[(\textbf{#1})]%
  \def\@currentlabel{(#1)}\def\@currentlabelname{#1}\phantomsection\label{#2}}
% NB (v3.66): \hitem is \ref-only. \autoref{key} calls it a SECTION ("סעיף (א)")
% because \phantomsection brands the anchor section*.N, and \nameref{key} prints
% the bare letter. For \autoref-able items prefer \hebenumlabels below.
%
% v3.66: \hebenumlabels — OPT-IN automatic Hebrew enumerate letters, done right.
%   {\hebenumlabels\begin{enumerate}\item\label{q:a} ...\end{enumerate}}
%   -> labels (א) (ב) (ג); \ref{q:a} = א; anchor Item.N (real, per-item).
% This retires the belief that made \item[(\textbf{א})] + \hitem necessary.
% "Automatic Hebrew labels fail silently" was TRUE but unexplained, and the
% mechanism is now bisected (plain article vs +enumitem, measured): ENUMITEM
% PRE-EXPANDS the label expression in an \edef BEFORE \refstepcounter, with the
% counter still at 0. Tokens it shields by name (\arabic,\alph,\roman) survive;
% a plain \newcommand \ifcase wrapper evaluates at 0 -> the EMPTY zero-branch ->
% the label prints as "()" with the letter GONE (and \the\c@enumi prints a
% literal 0, \number\value dies fatally pointing at \item). The cure is one
% property, not a doctrine: the wrapper must be ROBUST (survives the \edef
% unexpanded), while \theenumi stays EXPANDABLE (so \refstepcounter writes the
% real letter to the .aux). Verified: labels (א)(ב)(ג), \ref -> א/ב/ג each item
% its own value, .aux = {א}{...}{Item.1} — versus the explicit-\item doctrine
% where BIDI-50/\hitem exist purely to patch the counter that never stepped.
\def\@hebalph#1{\ifcase#1\or א\or ב\or ג\or ד\or ה\or ו\or ז\or ח\or ט\or י\or
  יא\or יב\or יג\or יד\or טו\or טז\or יז\or יח\or יט\or כ\else\@ctrerr\fi}
\DeclareRobustCommand*\hebalph[1]{\expandafter\@hebalph\csname c@#1\endcsname}
\newcommand{\hebenumlabels}{%
  \renewcommand{\theenumi}{\@hebalph\c@enumi}%
  \renewcommand{\labelenumi}{(\textbf{\hebalph{enumi}})}}
\makeatother
% v3.63: \mbox the LABEL. \thetbl ("7.1") and a caption description that STARTS
% WITH A DIGIT ("32 הרגיסטרים…") are two European-Number runs separated only by
% neutrals (the label's own "." and \enspace). Under the RTL paragraph base the UBA
% resolves those neutrals to the RTL level, so both numbers sit in ONE reversed span
% and EXCHANGE VISUAL POSITIONS. Measured on the shipped PDF: a table whose real
% number is 7.1 printed "טבלה 32" — a different, EXISTING table number — while \ref
% correctly printed 7.1, so the caption silently contradicted its own cross-reference,
% and the description lost its leading count. The source contains no adjacency at all;
% it is manufactured by this macro, which is why every source linter was blind.
% \mbox makes the label one atomic box so nothing can reorder across it. Verified
% zero-regression: an ordinary Hebrew caption keeps identical x-geometry.
% One paragraph style covers label, number, separator and description.
% Explicit project/template overrides remain possible; the label hook is kept
% for compatibility, with an identity default. Counters/anchors stay unchanged.
\newcommand{\lakitcaptionstyle}{\normalfont\small\color{black}\hypersetup{allcolors=black}}
\newcommand{\lakitcaptionlabelstyle}[1]{#1}
\newcommand{\lakitcaptionpunctuation}{:}
% Native floats, captionof and longtable use the same presentation policy.
% subcaption is optional and inherits the same font rather than a second size.
\usepackage{caption}
\DeclareCaptionFont{lakit}{\lakitcaptionstyle}
\DeclareCaptionLabelFormat{lakit}{\lakitcaptionlabelstyle{#1~\en{#2}}}
\DeclareCaptionLabelSeparator{lakit}{\lakitcaptionpunctuation\enspace}
\captionsetup{font=lakit,labelfont={},textfont={},labelformat=lakit,
  labelsep=lakit,justification=centering}
\makeatletter
\AtBeginDocument{\@ifpackageloaded{subcaption}{%
  \captionsetup[sub]{font=lakit,labelfont={},textfont={},labelsep=lakit}}{}}
\makeatother
% v3.64: OPT-IN list-of-figures/tables entry via an optional short caption.
% \listoffigures and \listoftables are inherited from report.cls and LOOK available,
% but \figcap/\tabcap step the custom \c@fig/\c@tbl counters and never call
% \addcontentsline — so both printed their heading and NOTHING ELSE. Measured: a doc
% with one \figcap and one \tabcap rendered page 1 as 'רשימת האיורים\n1\n', and the
% .lof held only \@chapter's \addvspace, zero \contentsline entries. An empty list
% under a correct Hebrew heading reads as "this document has no figures".
%
% WHY OPTIONAL, AND NOT AUTOMATIC. Writing #1 unconditionally was implemented first
% and REJECTED: a caption containing \codeline compiles CLEANLY today (verified on
% the pre-v3.64 preamble: 0 errors, PDF produced) but becomes FATAL once it is also
% written to the .lof — \codeline is \hbox to\linewidth{...}, and a full-width box in
% a list entry dies with "\pdfextension endlink cannot be used in vertical mode", no
% output PDF. Gaining a feature nobody uses (0 of 13 collections call \listoffigures)
% at the price of breaking a construct that works is a bad trade. Stock \caption has
% exactly this limitation and exactly this escape hatch, so the convention is familiar.
% \en{}, \code{}, \h{} and inline math in the short text are all safe — verified, they
% survive the write and typeset correctly in the list (RTL geometry confirmed: number
% rightmost at x=509.9, title left of it, dots and folio at the far left).
% With no optional argument the macro is byte-identical to v3.63 — zero risk to every
% existing document, and \listoffigures lists exactly the captions that opted in.
% BLD-37, TABLE side (the \figcap v3.74 idea, adapted): clicking a \ref to a
% table used to land MID-table — hyperref's automatic anchor inside
% \refstepcounter was emitted on the still-open horizontal line whose TAIL is
% the tabular (an hbox whose baseline is the table's vertical CENTER), so the
% jump scrolled the top half of the table above the viewport. The v3.74 entry
% deferred this side claiming "tablebox splits — no single box to measure";
% that deferral dissolves: the measurement target is the TABULAR itself at
% \tabcap ENTRY, while it is still the tail of the open list — the same
% situation \lakin@grabfig handles at an island's end (and the fleet has 0
% split tableboxes). So \lakin@grabtab runs as the FIRST token, before any
% typesetting can push the box away, and the anchor is re-emitted raised above
% the measured tabular. \lakintabraise is a SEPARATE register from
% \lakinfigraise (defined with it, below the LTR-island block): a captionless
% island intentionally leaves the fig register set, and a table must never
% inherit that stale figure height.
%
% WHY the anchor stays AT THE ENTRY SPOT and is NOT deferred into the caption
% paragraph the way \figcap defers it: at \par TeX DROPS a final glue item,
% and the automatic anchor's smash box sitting BEHIND the interword glue kept
% that glue alive — so every centered tabular has always sat half a space off
% true center. Deferring the anchor removed that box and was measured as 982
% differing glyph records: EVERY tabular line shifted 1.30pt (2.61pt for
% pattern 7's two minipages) while captions stayed put. Emitting the raised
% anchor at the old spot — through hyperref's own \Hy@raisedlink, exactly the
% macro the automatic path uses — keeps the line's node structure, and so the
% A/B glyph diff at 0, BY CONSTRUCTION. \Hy@raisedlink reads the raise from
% \HyperRaiseLinkDefault (default \baselineskip = the old mid-table dest); the
% group-local \edef freezes measured-height+42pt into it, and \smash inside
% \Hy@raisedlink hides any raise from the line's extent. \HyCnt@ProvideTheH-
% Counter keeps the dest NAME byte-identical to the automatic one (it defines
% \theHtbl, which \hyper@makecurrent prefers — without it the first measured
% caption was named tbl.1 where the automatic path said tbl.0.1, measured).
% vmode entry, or anything but a box at the tail (no tabular to measure),
% falls back to the plain \refstepcounter — byte-identical to the old path.
\makeatletter
\newcommand{\tabcap}[2][]{\lakin@grabtab
  % v3.79: arm the box-close correction. Set unconditionally and in BOTH
  % branches below's scope, so the plain-\refstepcounter fallback is corrected
  % too; \lakintabboxopen clears it at every \begin{tablebox}, so a \tabcap
  % outside a tablebox can never leave it armed for a later box.
  \global\lakintabcapseentrue
  \ifdim\lakintabraise>0pt
    \H@refstepcounter{tbl}\HyCnt@ProvideTheHCounter{tbl}\hyper@makecurrent{tbl}%
    \begingroup
      % A frameless table targets its measured top. The frame's historical
      % cushion belongs only to tablebox and is corrected at box close.
      \edef\HyperRaiseLinkDefault{\the\lakintabraise}%
      \iflakintabframed
        \edef\HyperRaiseLinkDefault{\the\dimexpr\lakintabraise+42pt\relax}%
      \fi
      \Hy@raisedlink{\hyper@anchorstart{\@currentHref}\hyper@anchorend}%
    \endgroup
    \global\lakintabraise 0pt\relax
  \else
    \refstepcounter{tbl}%
  \fi
  \lakinrtitle{#2}\par\nopagebreak\addvspace{6pt}%
  \def\lakin@short{#1}\ifx\lakin@short\@empty\else
    \addcontentsline{lot}{table}{\protect\numberline{\thetbl}{\protect\ignorespaces #1}}\fi
  {\centering\lakitcaptionstyle\mbox{\lakitcaptionlabelstyle{\tablename~\en{\thetbl}\lakitcaptionpunctuation}}\enspace #2\par}}
\makeatother

% ============================================================================
% Unified FIGURE caption — the \tabcap twin for TikZ/plot figures (v3.34.0).
% A figure (tikzpic) sits inline in the prose; put \figcap BELOW it for a
% numbered, \ref-able caption "איור N.M: תיאור" in the same style as \tabcap:
%   \begin{tikzpic} ... \end{tikzpic}
%   \figcap{\label{fig:foo}description}     % \label FIRST, inside the argument
% WHY a macro, not a hand-written gray line: a manual
%   {\centering\small\textcolor{cGray}{איור: ...}} with a \label inside has NO
% \refstepcounter, so the \label binds to the AMBIENT \@currentlabel (the last
% equation) — \ref then prints an equation number AS the figure number
% (measured on a real page: "איור 48" where 48 is equation 48). \figcap steps a
% real \c@fig counter first, so \label is correct, numbering is consistent, and
% refs are stable. Class-aware like \thetbl (chapter- or section-prefixed).
% ============================================================================
\makeatletter
\@ifundefined{c@chapter}{%
  \newcounter{fig}[section]
  \renewcommand{\thefig}{\ifnum\value{section}>0 \thesection.\fi\arabic{fig}}%
}{%
  \newcounter{fig}[chapter]
  \renewcommand{\thefig}{\ifnum\value{chapter}>0 \thechapter.\fi\arabic{fig}}%
}
\makeatother
\makeatletter
\newcommand{\figcap}[2][]{%
  \ifdim\lakinfigraise>0pt
    % v3.74: step WITHOUT hyperref's automatic anchor, then emit the anchor
    % raised above the measured figure. \H@refstepcounter is hyperref's saved
    % pre-hook \refstepcounter (sets the counter AND \@currentlabel — \label
    % in the caption arg binds exactly as before); \hyper@makecurrent supplies
    % \@currentHref, so the dest NAME is identical to the old automatic one —
    % \ref/\autoref/bookmarks are unchanged, only the dest'S POSITION moves.
    \H@refstepcounter{fig}\hyper@makecurrent{fig}%
    % The anchor is emitted INSIDE the caption paragraph below, as a zero-size
    % raisebox — NOT here in vertical mode, where a \raisebox starts an empty
    % paragraph and adds a blank line (measured: every glyph after the caption
    % shifted +18.8pt). Inside an existing line a [0pt][0pt] box is inert.
    \edef\lakin@anch{\@currentHref}%
    % \edef FREEZES the measured height into the pending macro: the deferred
    % \raisebox is typeset inside the caption paragraph AFTER the register is
    % reset to the 0pt sentinel below — a plain \def read 0 at use time
    % (measured: LAKINRAISE=0.0, dest moved 14pt instead of above the figure).
    \edef\lakin@pend{\noexpand\raisebox{\the\dimexpr\lakinfigraise+42pt\relax}[0pt][0pt]{%
      \noexpand\hyper@anchorstart{\lakin@anch}\noexpand\hyper@anchorend}}%
    \global\lakinfigraise 0pt\relax
  \else
    \refstepcounter{fig}\let\lakin@pend\relax
  \fi
  \lakinrtitle{#2}\par\nopagebreak\addvspace{4pt}%
  \def\lakin@short{#1}\ifx\lakin@short\@empty\else
    \addcontentsline{lof}{figure}{\protect\numberline{\thefig}{\protect\ignorespaces #1}}\fi
  {\centering\lakin@pend\lakitcaptionstyle\mbox{\lakitcaptionlabelstyle{\figurename~\en{\thefig}\lakitcaptionpunctuation}}\enspace #2\par}}
\makeatother

% ============================================================================
% LTR islands.  EVERY tikzpicture / quantikz MUST live inside one, or the whole
% figure's coordinates and Latin labels flip. (Hebrew node text still uses \h{}.)
% ============================================================================
% v3.74 (BLD-37): the \figcap hyperlink DESTINATION used to land at the caption
% line — BELOW the figure — so clicking a figure reference scrolled the figure
% off-screen (measured fleet-wide: 97/97 fig dests in one collection sat 19-54pt
% under their figure's ink; a 347pt figure was fully hidden after the jump).
% hypcap-style fix in two parts: (1) at the island's END, \lakin@grabfig grabs
% the just-typeset centered line with \lastbox, measures it, and puts it back
% with \nointerlineskip — the interline glue above the line was already
% contributed and is NOT re-added, so the page geometry is unchanged to the
% glyph (verified); (2) \figcap then steps the counter WITHOUT hyperref's
% automatic anchor (\H@refstepcounter = the pre-hyperref \refstepcounter) and
% emits the anchor itself, raised by the measured height — the dest lands ABOVE
% the figure top. A captionless island leaves numbering untouched (nothing is
% stepped at the island; the grab flag is cleared at every island begin, so a
% stale height from a captionless figure cannot leak into a later caption).
\newsavebox{\lakinfigbox}
\newdimen\lakinfigraise  % real \dimen: tex.setdimen cannot target a \newlength SKIP register (measured lua error)
% BLD-37 table side: a SEPARATE register for \tabcap. The two must not share —
% a captionless island INTENTIONALLY leaves \lakinfigraise set (the grab runs
% at every island END; only the NEXT island's begin clears it), so a shared
% register would hand a \tabcap after a captionless figure that stale FIGURE
% height. \lakintabraise is written only by \lakin@grabtab at \tabcap entry
% and consumed/reset inside \tabcap itself, so it can never go stale.
\newdimen\lakintabraise
\makeatletter
% v3.74b: measure WITHOUT touching the list. The first draft grabbed the box
% with \lastbox and put it back — and \lastbox DISCARDS the node's baseline
% shift, so a vertically-anchored quantikz circuit rose 2.57pt and everything
% after it sank 5.32pt (measured on the A/B glyph diff). LuaTeX lets us READ
% the tail node's dimensions without removing anything: geometry is identical
% by construction. \lakinfigraise=0pt is the "no figure measured" sentinel.
\def\lakin@grabfig{%
  \ifhmode
    \directlua{
      local n = tex.nest[tex.nest.ptr].tail
      local GLUE, PEN = node.id('glue'), node.id('penalty')
      while n and (n.id == GLUE or n.id == PEN) do n = n.prev end
      local HL, VL = node.id('hlist'), node.id('vlist')
      if n and (n.id == HL or n.id == VL) then
        tex.setdimen('global', 'lakinfigraise',
                     (n.height or 0) + (n.depth or 0))
      else
        tex.setdimen('global', 'lakinfigraise', 0)
      end}%
  \else
    \global\lakinfigraise 0pt\relax
  \fi}
% BLD-37 table side: \lakin@grabfig's twin, NOT a shared helper — the Lua body
% differs in exactly one place, the skip set. Between \end{tabular} and \tabcap
% an otherlanguage wrapper (patterns 3/5 of example_tables) leaves a \write
% whatsit and a dir node ON TOP of the interword glue (measured node dump:
% "whatsit/write <- dir <- glue(2.62) <- hlist"), so the table reader must
% also step over whatsit/dir markers — all zero-size — to reach the tabular.
% \lakin@grabfig deliberately does NOT get the wider skip set: its glue/penalty
% body is the one the v3.74 A/B proof froze, and widening it would newly
% measure (and so newly raise) anchors in island shapes that today take the
% plain path — an unproven behavior change on the FIG side. Anything that is
% not a box behind those markers still hits the 0pt sentinel and \tabcap falls
% back to the plain \refstepcounter.
% SECOND deliberate difference: this stores the box's HEIGHT, not height+depth.
% \figcap's anchor is emitted in the caption paragraph BELOW the figure, so it
% must climb the full extent; \tabcap's anchor sits ON the tabular's own line,
% whose baseline is the vcentered tabular's MIDPOINT — height alone reaches
% the top edge, and adding the depth was measured to overshoot by half the
% table (dests 80-135pt above their frame, INSIDE the preceding table's box).
\def\lakin@grabtab{%
  \ifhmode
    \directlua{
      local n = tex.nest[tex.nest.ptr].tail
      local GLUE, PEN = node.id('glue'), node.id('penalty')
      local WHAT, DIR = node.id('whatsit'), node.id('dir')
      while n and (n.id == GLUE or n.id == PEN
                   or n.id == WHAT or n.id == DIR) do n = n.prev end
      local HL, VL = node.id('hlist'), node.id('vlist')
      if n and (n.id == HL or n.id == VL) then
        tex.setdimen('global', 'lakintabraise', n.height or 0)
      else
        tex.setdimen('global', 'lakintabraise', 0)
      end}%
  \else
    \global\lakintabraise 0pt\relax
  \fi}
\makeatother
% ============================================================================
% v3.79 (BLD-37 table side, CORRECTION): put the anchor at the TOP OF THE BOX,
% instead of guessing the distance to it.
% ============================================================================
% v3.77 raised the \tabcap anchor by (height of the box at the tail of the open
% horizontal list) + 42pt. That is a MEASUREMENT taken BEFORE line breaking, and
% it is short by every term it cannot see. Fleet-measured on the 12 built
% collections (299 tbl.* dests, ownership established through each table's own
% CAPTION, never through "which frame encloses the dest"): 3 dests landed INSIDE
% their own tablebox, below the frame top —
%   computer-architecture       tbl.14.1  -46.4pt   three stacked tabulars
%   computer-architecture       tbl.30.5  -13.5pt   tabular 2.85pt over \linewidth
%   advanced-electronic-devices tbl.B.1   -59.2pt   tabular 2.56pt over \linewidth
% Two missing terms, both reproduced in a minimal document (probe P2/P3):
%   (1) VERTICAL LIST ALREADY CONTRIBUTED inside the box — a second and third
%       tabular, or a heading above the table. The tail read sees only the LAST
%       box, so 77.7pt of stacked tabulars is invisible to the raise.
%   (2) WHICH LINE THE ANCHOR LANDS ON. When the tabular alone overfills
%       \linewidth (459.6121pt in a tablebox), the interword glue after
%       \end{tabular} becomes the chosen breakpoint and \Hy@raisedlink is
%       emitted on the NEXT line — one (depth + lineskip) lower. That term does
%       not EXIST at \tabcap time: line breaking has not happened yet. No
%       measurement taken there can ever recover it. (Proof: probe P3, the same
%       table 2.84912pt too wide, fails at -13.5pt; probe P4, the identical
%       table with \tabcolsep 2pt so it fits, is correct at +32.0pt.)
% So the raise is not computed at \tabcap at all. It is CORRECTED at
% \end{tablebox}, from inside mdframed's own body vbox (endinnercode runs after
% \par\unskip, i.e. after line breaking, while tex.nest[tex.nest.ptr] is still
% that vbox's vertical list). There the anchor's true offset from the top of the
% box content is not estimated but WALKED, and the \Hy@raisedlink shift is set
% so the dest lands 42pt above it — which is 42pt - innertopmargin(9pt) -
% linewidth(0.9pt) = 32.0pt above the FRAME top, the exact position the 295
% already-correct dests have today. That is why this is a correction and not a
% re-placement: the good sites do not move.
%
% WHY THE GLYPH DIFF IS 0 BY CONSTRUCTION — the only thing written is the
% `shift` field of the hlist that \Hy@raisedlink already built, and the guard
% below refuses to write it unless that box is 0x0x0. A zero-size box carries no
% dimension into its parent's packing, the packing already happened, and shift
% is not an input to it. Nothing else in the list is touched: no node is added,
% removed, resized or reordered. Measured, not assumed: assets/example_tables.tex
% and the whole computer-architecture collection A/B at 0 differing pymupdf
% rawdict glyph records.
%
% WHY NOT the alternatives (all prototyped and measured before this one):
%   * accumulate the vertical list at \tabcap time — fixes term (1), CANNOT fix
%     term (2); 2 of the 3 real failures survive. A fix that leaves the majority
%     of the defect standing is the false all-clear this repo exists to avoid.
%   * \zref-savepos two-pass — correct on all shapes, but it makes the dest
%     POSITION aux-derived (DEFECTS BLD-1/2), a single-pass compile silently
%     ships an unraised anchor, and the bare y-difference has no page guard.
%   * emit the anchor at \begin{tablebox} and let \tabcap adopt it — needs the
%     tbl counter stepped at box OPEN to keep the dest NAME byte-identical; a
%     captionless tablebox then either shifts every later table number or leaves
%     an orphan dest that the next table duplicates (a HARD build gate). This
%     correction touches neither the counter nor the name: \tabcap's naming path
%     is v3.77's, unchanged, so resolve_names() is byte-identical (verified).
%   * suppressing the line break (a penalty before the glue) — that MOVES the
%     tabular: the break currently discards the interword glue, so forbidding it
%     re-centers the two overfull tables. Reflow, rejected.
%
% SCOPE. The correction is armed only by \tabcap ITSELF (\lakintabcapseen),
% cleared at every box open, so a captionless tablebox is a no-op and a \tabcap
% hosted OUTSIDE a tablebox (the fleet has exactly one, computer-architecture
% tbl.14.2 in an exbox) keeps v3.77's measured raise. Both fallbacks are
% therefore never worse than today; and because hyperref's AUTOMATIC anchor uses
% the same \Hy@raisedlink shape, the plain-\refstepcounter path (vmode entry, or
% prose between \end{tabular} and \tabcap — probe P6, -51.5pt) is corrected too.
% KNOWN LIMIT — a tablebox that SPLITS across pages. The offset is measured from
% the top of the UNSPLIT body, so a dest in segment 2 gets raised past its own
% segment. MEASURED on a forced split (probe: 40 filler lines + a table): the
% dest went from y=585.0 on page 2 (inside segment 2, at the caption — the old
% BLD-37 symptom) to y=-113.4, i.e. ABOVE the page top. Viewers clamp that to the
% page top, which happens to show segment 2 whole, but an out-of-page coordinate
% is not something to ship on purpose. Hence the \textheight guard in
% lakintab.fix: a body TALLER THAN THE TEXT BLOCK cannot fit on any page, so it
% must split — the correction is skipped and v3.77's behaviour stands unchanged
% (verified: the probe returns to y=585.0). The guard is deliberately EXACT, not
% a page-fitting heuristic: it is sound (a taller-than-page body always splits)
% but PARTIAL (a shorter body that starts low on a page could still split). That
% residue is unreachable in practice — v3.37's everyline + the caption's
% \nopagebreak glue move such boxes WHOLE instead (probe N7: a 22-row table
% pushed down by 430pt moved whole, 0 segments), and the fleet has 0 split
% tableboxes, re-counted after this change. Reimplementing mdframed's own
% page-fitting arithmetic to close the residue would be a heuristic with its own
% failure mode, which is the trade this repo has repeatedly refused.
\newif\iflakintabcapseen
\newif\iflakintabframed
% The Lua lives in ONE \directlua at load time, and is written with NO
% backslash, %, # or ~ : \directlua expands its argument like \edef, so an
% active ~ would expand to \nobreakspace and a % would eat the rest of the line
% (the second is how the first draft of this block died). "not (a == b)" stands
% in for Lua's inequality operator for the same reason.
\directlua{
  lakintab = { }
  lakintab.HL = node.id("hlist")
  lakintab.VL = node.id("vlist")
  lakintab.RU = node.id("rule")
  lakintab.GL = node.id("glue")
  lakintab.KE = node.id("kern")
  lakintab.WH = node.id("whatsit")
  lakintab.TM = node.id("temp")
  lakintab.WD = node.subtype("pdf_dest")
  function lakintab.d(v)
    if v == nil then return 0 end
    if v < -1000000000 then return 0 end
    return v
  end
  function lakintab.scan(list, y, horiz, parent, target)
    local x = list
    while x do
      local i = x.id
      if i == lakintab.HL or i == lakintab.VL then
        if horiz then
          local b = y + x.shift
          if i == lakintab.HL then
            lakintab.scan(x.head, b, true, x, target)
          else
            lakintab.scan(x.head, b - lakintab.d(x.height), false, x, target)
          end
        else
          if i == lakintab.HL then
            lakintab.scan(x.head, y + lakintab.d(x.height), true, x, target)
          else
            lakintab.scan(x.head, y, false, x, target)
          end
          y = y + lakintab.d(x.height) + lakintab.d(x.depth)
        end
      elseif i == lakintab.WH then
        if x.subtype == lakintab.WD then
          local id = x.dest_id
          if type(id) == "string" and string.sub(id, 1, 4) == "tbl." then
            if horiz and parent and parent.width == 0
               and parent.height == 0 and parent.depth == 0 then
              parent.shift = parent.shift + (target - y)
            end
          end
        end
      elseif not horiz then
        if i == lakintab.RU then
          y = y + lakintab.d(x.height) + lakintab.d(x.depth)
        elseif i == lakintab.GL or i == lakintab.KE then
          y = y + lakintab.d(x.width)
        end
      end
      x = x.next
    end
  end
  function lakintab.fix(target, maxh)
    local nst = tex.nest[tex.nest.ptr]
    if nst == nil then return end
    local h = nst.head
    if h == nil then return end
    if h.id == lakintab.TM then h = h.next end
    local tot = 0
    local x = h
    while x do
      local i = x.id
      if i == lakintab.HL or i == lakintab.VL or i == lakintab.RU then
        tot = tot + lakintab.d(x.height) + lakintab.d(x.depth)
      elseif i == lakintab.GL or i == lakintab.KE then
        tot = tot + lakintab.d(x.width)
      end
      x = x.next
    end
    if tot > maxh then return end
    lakintab.scan(h, 0, false, nil, target)
  end
}
% @-FREE NAMES ON PURPOSE: these two are named in the `tablebox` key list at
% \newmdenv above, which is tokenised OUTSIDE any \makeatletter region — a
% \lakin@... there would tokenise as \lakin followed by the letters @tabopen.
\newcommand{\lakintabboxopen}{\lakintabframedtrue\global\lakintabcapseenfalse}
\newcommand{\lakintabboxclose}{%
  \iflakintabcapseen
    \directlua{lakintab.fix(-\number\dimexpr\lakintabcushion\relax,
                            \number\dimexpr\textheight\relax)}%
    \global\lakintabcapseenfalse
  \fi\lakintabframedfalse}
% The cushion above the box's CONTENT top. 42pt puts the dest 32.0pt above the
% frame top (42pt - innertopmargin 9pt - linewidth 0.9pt), i.e. exactly where
% the 295 correct dests already are — the correction moves only the broken ones.
\newcommand{\lakintabcushion}{42pt}
\newenvironment{qcirc}{\par\medskip\global\lakinfigraise 0pt\relax\begin{otherlanguage}{english}\begin{center}}
  {\csname lakin@grabfig\endcsname\end{center}\end{otherlanguage}\par\medskip}
\newenvironment{tikzpic}{\par\medskip\global\lakinfigraise 0pt\relax\begin{otherlanguage}{english}\begin{center}}
  {\csname lakin@grabfig\endcsname\end{center}\end{otherlanguage}\par\medskip}

% ============================================================================
% Code listings.  ★★ v3.4 THE BOX-OVERFLOW FIX:
%   resetmargins=true   → the listing IGNORES the enclosing environment's
%                         indentation and measures from the PAGE. v3.4 stated
%                         this backwards ("measured from the ENCLOSING box"),
%                         and v3.37 then blanked in-box line numbers because
%                         the page-anchored listing put its hanging digits on
%                         top of the box rule. Measured v4.5.0: with the key
%                         FALSE the listing follows the box and the code moves
%                         from x=75.26 to x=97.74, clearing the rule at 72.07.
%                         The key now lives in one \lstset below instead of in
%                         each style, so a box can override it — codebox does.
%   framexleftmargin=0  → the frame (left bar) does not extend left of the code.
%   xleftmargin keeps a small indent so the line numbers have room INSIDE.
% asmblock wraps the listing in an LTR island (code is LTR).
% ============================================================================
\lstdefinestyle{asm}{
  basicstyle=\ttfamily\small,
  backgroundcolor=\color{cBlueBg!55},
  frame=leftline,framerule=2.5pt,rulecolor=\color{cBlue},
  numbers=left,numberstyle=\tiny\color{cGray},numbersep=8pt,
  xleftmargin=2.2em,framexleftmargin=0pt,framexrightmargin=0pt,
  breaklines=true,columns=fullflexible,keepspaces=true,
  showstringspaces=false,aboveskip=8pt,belowskip=8pt,
  % v3.36: without a comment DELIMITER the commentstyle below never fires —
  % listings has no `language` here, so it does not know '#' opens a comment,
  % and every asm listing rendered flat black (measured: a whole-page pixel
  % sweep found ZERO cGreen glyphs across 135 shipped listings). morecomment
  % gives it the delimiter; color-only change, glyph widths untouched.
  morecomment={[l]{\#}},
  commentstyle=\color{cGreen},
}
\lstdefinestyle{cstyle}{
  language=C,basicstyle=\ttfamily\small,
  backgroundcolor=\color{cGreenBg!55},
  frame=leftline,framerule=2.5pt,rulecolor=\color{cGreen},
  numbers=left,numberstyle=\tiny\color{cGray},numbersep=8pt,
  xleftmargin=2.2em,framexleftmargin=0pt,framexrightmargin=0pt,
  breaklines=true,columns=fullflexible,keepspaces=true,
  showstringspaces=false,keywordstyle=\color{cBlue}\bfseries,
  commentstyle=\color{cGray}\itshape,aboveskip=8pt,belowskip=8pt,
}
% v3.36: pystyle/htmlstyle — colored twins of cstyle for the Python/HTML
% listings that used to repeat a ~200-char options blob with NO color
% (keywords rendered bold-black, comments italic-black — measured; C alone
% was colored). Geometry is copied VERBATIM from the shipped blob (frame,
% margins, skips) so swapping a call site to style=pystyle moves NOTHING —
% only keywordstyle/commentstyle/stringstyle are added (bold stays bold,
% italic stays italic, mono metrics identical => zero reflow, proven by
% per-page word-bbox diff). breaklines/breakatwhitespace stay per-site.
% v3.78: numbers=left joins both styles, with the exact numberstyle/numbersep
% keys asm/cstyle carry — fleet-wide consistency (author decision: all four
% styles numbered; "asm/cstyle numbered, py/html not" was an accident of the
% v3.36 geometry freeze, not a design). The numbers hang LEFT of the code box
% (listings sets them with \llap, so no overfull hbox can arise and the code
% x is untouched — measured: code frozen, digits at code_x-11.8pt); inside a
% colored box they are blanked like everyone else's (v3.37 hook), and the
% clean variant restores them in its frameless boxes (v3.78, above).
\lstdefinestyle{pystyle}{
  language=Python,basicstyle=\ttfamily\small,
  columns=fullflexible,keepspaces=true,
  frame=single,rulecolor=\color{cTblLine},
  numbers=left,numberstyle=\tiny\color{cGray},numbersep=8pt,
  xleftmargin=0.5em,framexleftmargin=0pt,framexrightmargin=0pt,
  showstringspaces=false,aboveskip=8pt,belowskip=8pt,
  keywordstyle=\color{cBlue}\bfseries,
  commentstyle=\color{cGray}\itshape,
  stringstyle=\color{cGreen},
}
% v4.5.0: hoisted out of the four styles so an enclosing box can override it.
% Every listing still resolves to true unless something later says otherwise,
% so shipped output is unchanged (verified byte-for-byte on all eight PDFs).
\lstset{resetmargins=true}
\lstdefinestyle{htmlstyle}{
  language=HTML,basicstyle=\ttfamily\small,
  columns=fullflexible,keepspaces=true,
  frame=single,rulecolor=\color{cTblLine},
  numbers=left,numberstyle=\tiny\color{cGray},numbersep=8pt,
  xleftmargin=0.5em,framexleftmargin=0pt,framexrightmargin=0pt,
  showstringspaces=false,aboveskip=8pt,belowskip=8pt,
  keywordstyle=\color{cBlue}\bfseries,
  commentstyle=\color{cGray}\itshape,
  stringstyle=\color{cGreen},
}
\newenvironment{asmblock}{\par\begin{otherlanguage}{english}}{\end{otherlanguage}\par}

% <hllpdf:listing-smart-quotes>
% v4.21.0: four Unicode punctuation characters must enter listings' token
% buffer. Unregistered characters bypass it: a closing quote was printed
% BEFORE the pending word, and an opening quote BEFORE a pending brace.
% Extend the standard character table, then activate only these codepoints
% in the listing's local character-table scope. Do this BEFORE the package's
% delimiter/alsoother/literate hooks so explicit user mappings still win.
% This is not universal Unicode support, nor a conversion to ASCII quotes.
% See listings developer guide 15.3.1 and references/listing-smart-quotes.md.
\makeatletter
\lst@CCPut \lst@ProcessOther {"2018}{"2019}{"201C}{"201D}\z@
\lst@AddToHookAtTop{SelectCharTable}{%
  \lst@ifactivechars
    \catcode"2018=\active \catcode"2019=\active
    \catcode"201C=\active \catcode"201D=\active
  \fi}
\makeatother
% </hllpdf:listing-smart-quotes>

% <hllpdf:listing-technical-symbols>
% v4.30: bounded extension for common technical punctuation in LTR listings.
% Without character-table registration these glyphs bypass the token buffer
% and can move before the preceding word even inside quoted strings.
% Keep the original four-quote block separate for its standing mutation test.
% This is source-faithful printing, not syntax validation or general Unicode.
\makeatletter
\lst@CCPut \lst@ProcessOther
  {"2013}{"2014}{"2026}{"2190}{"2192}{"2194}{"21D2}
  {"2212}{"221E}{"2260}{"2264}{"2265}\z@
\lst@AddToHookAtTop{SelectCharTable}{%
  \lst@ifactivechars
    \catcode"2013=\active \catcode"2014=\active \catcode"2026=\active
    \catcode"2190=\active \catcode"2192=\active \catcode"2194=\active
    \catcode"21D2=\active \catcode"2212=\active \catcode"221E=\active
    \catcode"2260=\active \catcode"2264=\active \catcode"2265=\active
  \fi}
\makeatother
% </hllpdf:listing-technical-symbols>


\usepackage[hidelinks,unicode]{hyperref}   % LAST, after colors. unicode = Hebrew bookmarks.
% v3.62: THE ISOLATION MACROS CORRUPTED EVERY PDF BOOKMARK THEY APPEARED IN.
% \en is defined (line ~155) as
%     \foreignlanguage{english}{\language=\csname l@nohyphenation\endcsname #1}
% and in a heading that body is written verbatim to the .toc and handed to
% hyperref's \pdfstringdef. hyperref strips the two CONTROL SEQUENCES (\language,
% \l@nohyphenation) but the plain `=` CHARACTER token between them is not a control
% sequence, so it survives into the PDF string: \chapter{שער \en{CNOT} והצמדה}
% produced the bookmark "שער =CNOT והצמדה". The printed TOC and the heading on the
% page are CORRECT — only the navigation pane and the viewer window title are wrong,
% which is why it went unnoticed. Measured on the shipped fleet: 70 corrupted
% bookmarks across 5 PDFs, from 177 unprotected \en-in-heading sites. Same corruption
% reached /Title and /Subject via the \hypersetup{pdftitle=...} the docs recommend.
% This hook tells hyperref to use the plain argument for all three isolation macros
% when building a PDF string. It does NOT touch the .toc (the printed TOC is
% byte-unchanged) and it does NOT mask the math signal — a heading with $...$ still
% logs "removing math shift", so BIDI-15 detection is intact.
% BONUS, measured: it also makes \code{} in an UNSTARRED heading COMPILE. BLD-10
% recorded that as needing \texorpdfstring or a starred heading; the real root cause
% is hyperref's \pdfstringdef scan of \code's \ifmmode...\else...\fi body — which is
% exactly why the identically-bodied \texttt survived (hyperref pre-neutralizes
% \texttt in PDF strings) and why \DeclareRobustCommand could never help.
% v3.64: \def\\{ } added. A manual line break in a heading WELDS the two words
% together in the bookmark: \chapter{שורה ראשונה\\שורה שנייה} produced the outline
% entry 'שורה ראשונהשורה שנייה' — "ראשונהשורה" is not a word. hyperref drops the
% \\ and puts nothing in its place, so the break that looks like a space on the
% page is a zero-width join in the navigation pane. The spaced form
% 'מילה אחת \\ מילה שנייה' escapes with a harmless DOUBLE space ('אחת  מילה'),
% which is why this never got noticed. Both are measured on a real outline.
% 0 live instances in the corpus — but SKILL.md and content-style.md both RECOMMEND
% a manual heading break, so this is a primed trap rather than a dormant one, and a
% one-token preamble fix retires it for good instead of asking authors to remember.
% Mapping \\ to a space fixes the WELD, which is the actual defect (a nonexistent
% word in the navigation pane). It does not normalise whitespace: measured after the
% fix, the welded form becomes 'שורה ראשונה שורה שנייה' (correct) and the already-
% spaced form goes from two spaces to three. That is cosmetic — viewers render it as
% a wider gap, and no word is corrupted. A per-heading \texorpdfstring is still the
% way to get an exactly-tuned bookmark string.
\pdfstringdefDisableCommands{\def\en#1{#1}\def\h#1{#1}\def\doclang#1{#1}\def\code#1{#1}\def\codeline#1{#1}\def\\{ }}
% v3.70: RESTORE \url LINE-BREAKING. url.sty freezes its break penalties at LOAD
% time (\mathchardef\UrlBreakPenalty=\binoppenalty), and v3.45's global
% \binoppenalty=10000 (inline-math atomicity) therefore priced EVERY breakpoint
% inside \Url at 10000 = never break. Measured: a long \url in Hebrew prose is
% one unbroken line, ~100pt overfull, invading the margin band the ink gates
% deliberately tolerate. With the two lines below: 0 overfull, the URL breaks
% after '/', the RTL split is CORRECT (fragment order forward), and each line
% gets its own link annotation exactly covering its fragment (measured).
% Confined to \Url's group — inline-math atomicity (TYP-7) is untouched.
\mathchardef\UrlBreakPenalty=200
\mathchardef\UrlBigBreakPenalty=100
\ifx\lakitmainlanguage\hllpdfenglish
  \hypersetup{pdflang=en}
\else
  \hypersetup{pdflang=he}
\fi  % /Lang is the document base locale (not per-span accessibility tagging).
% Set a per-document \hypersetup{pdftitle={...}} too — see SKILL.md.

% v3.52 (round-2 audit): two RTL bibliography/cross-ref defects, both invisible to
% every gate and both affecting the FIRST author to use the feature.
% (a) \@biblabel's default [#1] brackets are bidi-neutral and mirror to ]1[ under
% the RTL bibliography base — every \thebibliography reference list shows ]1[.
% LTR-isolating the label restores [1]. (The render gate's mirror rule fingerprints
% round parens )1(, so square brackets slipped through.)
\makeatletter
\renewcommand*{\@biblabel}[1]{\foreignlanguage{english}{[#1]}}
% v3.54: the SAME neutral-bracket trap on the INLINE side — a multi-key \cite{a,b}
% under the RTL base lays its European-number list right-to-left, so [1, 2] renders
% "[2,1]" and [1, 2, 3] renders "[3,2,1]" (verified by glyph geometry — the reader
% sees the sources in REVERSE order; a note form \cite[p.5]{a} also reorders to
% "[p. 5,1]"). LTR-isolate the whole bracket group (number list AND optional note),
% exactly like \@biblabel above: [1, 2] / [1, 2, 3] / [1, p. 5] all render correct.
% Placed AFTER hyperref so it wins; hyperref hyperlinks each number inside #1 (via
% \@citex), which this wrapper leaves intact. Mirrors the LaTeX kernel \@cite shape.
% v4.9.1: directional isolation must NOT switch the language here. Babel's
% safe-active citation wrapper passes its key list through \bbl@tempa;
% \foreignlanguage reuses that scratch macro before the deferred citation
% loop expands it. The old line therefore cited the literal key "english".
% If such a key existed, even the undefined-reference gate stayed green.
% \babelsublr isolates the bracket group without selecting a locale or
% changing that key list. Preserve native links, key order and optional notes.
\def\@cite#1#2{\babelsublr{[{#1\if@tempswa , #2\fi}]}}
\makeatother

% v4.11.0: a reference inside align's \text inherited HEBdisptag=true,
% so its enclosing hbox mirrored the brackets despite a clean compile.
% Isolate the complete printed reference without selecting a Babel language.
% Also neutralize the display-only normalfont direction reset locally. A
% direction-only patch inverted 12 to 21 inside equation in the measured probe.
% Plain maketag is local to the reference: mathtools' display-side wrapper
% otherwise globally consumes manual_tag while formatting a nested reference.
% Preserve each original command (including the selective auxiliary write),
% current tag forms, native links, and showonlyrefs on/off transitions.
\makeatletter
\NewCommandCopy{\HEB@nativeeqref}{\eqref}
\newcommand{\HEB@eqrefisolate}[1]{%
  \babelsublr{\HEBdisptagfalse\let\bbl@eqnodir\relax\let\maketag@@@\ORIGmaketag@@@#1}}
\protected\def\eqref#1{\HEB@eqrefisolate{\HEB@nativeeqref{#1}}}
\ExplSyntaxOn
\NewCommandCopy{\HEB@nativeMTeqref}{\MT_extended_eqref:n}
\cs_set_protected:Npn \MT_extended_eqref:n #1
  {\HEB@eqrefisolate{\HEB@nativeMTeqref{#1}}}
\ExplSyntaxOff
\makeatother

% v4.11.0: custom caption counters are not part of hyperref's English
% locale, and Hebrew has no upstream hyperref locale in this toolchain.
% Restore the appropriate names at every Babel language switch, including
% return from a nested inline island. Do not localize the destination keys.
\makeatletter
\newcommand{\HEB@refnames}{%
\def\chapterautorefname{פרק}%
\def\sectionautorefname{סעיף}%
\def\subsectionautorefname{תת־סעיף}%
\def\subsubsectionautorefname{תת־תת־סעיף}%
\def\paragraphautorefname{פסקה}%
% v4.31: hyperref uses a distinct counter name at the deepest standard level.
\def\subparagraphautorefname{תת־פסקה}%
\def\equationautorefname{משוואה}%
\def\figureautorefname{איור}%
\def\tableautorefname{טבלה}%
\def\partautorefname{חלק}%
\def\appendixautorefname{נספח}%
\def\footnoteautorefname{הערת שוליים}%
\def\itemautorefname{סעיף}%
\def\Itemautorefname{סעיף}%
\def\figautorefname{איור}%
\def\tblautorefname{טבלה}%
\def\pageautorefname{עמוד}% 
}
\addto\extrashebrew{\HEB@refnames}
\newcommand{\ENG@refnames}{%
  \def\subparagraphautorefname{subparagraph}%
  \def\chapterautorefname{chapter}%
  \def\sectionautorefname{section}%
  \def\subsectionautorefname{subsection}%
  \def\subsubsectionautorefname{subsubsection}%
  \def\paragraphautorefname{paragraph}%
  \def\equationautorefname{Equation}%
  \def\figureautorefname{Figure}%
  \def\tableautorefname{Table}%
  \def\partautorefname{Part}%
  \def\appendixautorefname{Appendix}%
  \def\footnoteautorefname{footnote}%
  \def\itemautorefname{item}%
  \def\figautorefname{Figure}%
  \def\tblautorefname{Table}%
  \def\Itemautorefname{item}%
  \def\pageautorefname{page}%
}
\addto\extrasenglish{\ENG@refnames}
\ifx\lakitmainlanguage\hllpdfenglish\ENG@refnames\else\HEB@refnames\fi
\makeatother

% ============================================================================
% Optional running head (v3.39): \lakithead{<book title>} — OFF by default.
% A document opts in with ONE line after \input{preamble.tex}:
%   \lakithead{מבנה המחשב}
% Every content page then carries the series' orientation line: chapter name
% at the reading edge, book title opposite, thin rule. (fancyhdr's [L]/[R]
% are PHYSICAL; under the RTL page direction [L] lands at the RIGHT/reading
% edge — measured, consistent across the corpus.) Packaged like the topic
% rail: always loaded, inert without the call. This centralizes the 8-line
% block that 8 lakits used to copy by hand — the same duplication vector
% that once let a lakit preamble go ~100 lines stale. Class-aware: report
% fills the mark from \chapter, article from \section. NB: a titlepage env
% self-empties; front matter that sets \pagestyle{empty} must restore with
% \pagestyle{fancy} afterwards (NOT plain — that discards the head).
% ============================================================================
\usepackage{fancyhdr}
% v3.61: \chapter* left a STALE running head — pages of an UNNUMBERED chapter
% carried the PREVIOUS chapter's title. report.cls routes \chapter* through
% \@schapter, which never calls \chaptermark, so \leftmark keeps whatever the last
% numbered chapter left. Measured: with \chapter{פרק ממוספר} then
% \chapter*{מבוא לא ממוספר}, page 4 (inside the starred chapter) printed
% "פרק ממוספר". Invisible on the chapter's FIRST page (\thispagestyle{plain}), so it
% only shows from the second page on — the pages an author is least likely to check.
% An unnumbered front-matter chapter (\chapter*{מבוא}, \chapter*{הקדמה}) is everyday
% Hebrew-academic structure, and the corpus already carries HAND-WRITTEN \markboth
% lines right after \chapter* — the exact per-document duplication \lakithead exists
% to retire. Patched INSIDE \lakithead so the macro stays inert when unused; a
% document's own later \markboth still wins, so those hand-patches do not regress.
\makeatletter
% <hllpdf:starred-heading-scope>
\newif\iflakin@starmarkpatched
\newcommand{\lakistarmark}{%
  \iflakin@starmarkpatched\else
  % Dispatch on the actual class counter, not hyperref's shared internals.
  \ifdefined\c@chapter
    \let\HEBorigschapter\@schapter
    % Record only after the native heading has reached its page. A mark
    % before the native routine can advertise the next heading too early.
    \def\@schapter##1{\HEBorigschapter{##1}\chaptermark{##1}}%
  \else
    % @ssect receives layout arguments, NOT the section name or level.
    % Capture the name in @startsection without introducing a TeX group:
    % grouping sectioning would break the native after-heading state.
    \def\HEB@sectionkind{}%
    \def\HEB@mainsection{section}%
    \let\HEBorigstartsection\@startsection
    \def\@startsection##1##2##3##4##5##6{%
      \def\HEB@sectionkind{##1}%
      \HEBorigstartsection{##1}{##2}{##3}{##4}{##5}{##6}}%
    \let\HEBorigssect\@ssect
    \def\@ssect##1##2##3##4##5{%
      \HEBorigssect{##1}{##2}{##3}{##4}{##5}%
      % Lower starred headings must not replace the article's section mark.
      \ifx\HEB@sectionkind\HEB@mainsection
        \sectionmark{##5}%
      \fi
      \ignorespaces}%
  \fi
  \global\lakin@starmarkpatchedtrue
  \fi}
% </hllpdf:starred-heading-scope>
% <hllpdf:running-language>
% Choose the wrapper when the mark is written, not when a later page ships.
% English punctuation must remain LTR even after returning to Hebrew prose.
\newcommand{\HEB@runningmark}[1]{%
  \iflanguage{english}{\markboth{\en{#1}}{}}{\markboth{\h{#1}}{}}}
% </hllpdf:running-language>
\makeatother
% v4.2.1 (NAV): every collection already names itself in \lakithead — feed the
% PDF /Title from the same declaration, so the navigation gate's
% PDF_REQUIRE_TITLE has a universal source and no fleet member ships
% title-less by missing a second declaration site. Measured before this: 14
% of 16 collections shipped /Title-empty while every one carried a perfectly
% good \lakithead string (plain Hebrew in all 16 — no math, no formatting, so
% the PDF-string path is clean). An explicit title always wins: set it with
% \lakitpdftitle{...} anywhere before \lakithead, or with a plain
% \hypersetup{pdftitle=...} — the fallback detects that via hyperref's
% \@pdftitle, probe-verified \@empty-equal until the first explicit set.
\makeatletter
\newif\iflakit@pdftitleset
\newcommand{\lakitpdftitle}[1]{%
  \hypersetup{pdftitle={#1}}\global\lakit@pdftitlesettrue}
\newcommand{\lakithead}[1]{%
  \iflakit@pdftitleset\else
    \ifx\@pdftitle\@empty
      \hypersetup{pdftitle={#1}}\global\lakit@pdftitlesettrue
    \fi
  \fi
  \pagestyle{fancy}\fancyhf{}%
  \lakistarmark

  \fancyhead[L]{\small\textcolor{cGray}{\leftmark}}%
  % <hllpdf:fixed-head-language>
  % The fixed document title has its OWN language snapshot, independent of
  % the current section/chapter mark. Choose the branch NOW, not at shipout:
  % the page may be Hebrew by the time an English title is finally printed.
  \iflanguage{english}{%
    \fancyhead[R]{\small\textcolor{cGray}{\en{#1}}}%
  }{%
    \fancyhead[R]{\small\textcolor{cGray}{\h{#1}}}%
  }%
  % </hllpdf:fixed-head-language>
  \fancyfoot[C]{\thepage}%
  \renewcommand{\headrulewidth}{0.4pt}%
  \ifdefined\chaptermark
    \renewcommand{\chaptermark}[1]{\HEB@runningmark{##1}}%
  \else
    \renewcommand{\sectionmark}[1]{\HEB@runningmark{##1}}%
  \fi}
\makeatother

% v3.39.1: path-independent byte-reproducibility. LuaTeX derives the trailer
% /ID from the OUTPUT PATH, so two clones of the same source at different
% directories produced different sha256 (measured: the ONLY delta was the
% 62-byte /ID line). Bit 512 suppresses exactly that line — measured: Creator/
% Producer/CreationDate/ModDate and all content untouched. /ID is optional
% per ISO 32000 (required only for ENCRYPTED PDFs and PDF/A — neither is a
% target here). If a future workflow needs an ID (PDF/A, Acrobat revision
% matching), delete this line and it returns.
\pdfvariable suppressoptionalinfo \numexpr 512\relax

% v3.75: BOLD MATH IN BOLD HEADINGS. Heading text is \bfseries but math inside
% it stayed regular-weight CM — a hairline $V_T$ inside a heavy 24.8pt chapter
% title (measured: CMMI12 fl=6 next to FrankRuhlLibre-Bold fl=20; 155 live
% heading-with-math sites fleet-wide). \boldmath rides along wherever the class
% sets \bfseries for a heading; A/B measured: $V_T$ -> CMMIB10 fl=22, 0 missing
% characters (the bold CM faces ship with this TeX Live). Guarded per-macro so
% the article class (no \@makechapterhead) and future class changes stay safe.
% v4.0.1 (adoption fix): the SEARCH patterns above carry no parameter token.
% v4.0.0 made them specific by including the class's own #1/#2, and measured on
% TeX Live 2023 that never matches: a \patchcmd search pattern containing a
% parameter token cannot match while the call sits inside \@ifundefined{...}{}{...}
% (probed both #1 and ##1 -> FAIL; the same pattern at top level succeeds). Every
% one of these patches therefore failed, and because v4.0.0 also replaced the
% empty failure branch with \PackageError, the failure became FATAL: all four
% bundled example assets stopped building ("Cannot patch chapter-title math
% weight -> Fatal error, no output PDF"). The \PackageError branches are KEPT —
% a silent no-op here is what let v3.75's fix rot unnoticed — only the parameter
% tokens are dropped, so the patches match and a real class change still errors.
\makeatletter
\@ifundefined{@makechapterhead}{}{%
  \patchcmd{\@makechapterhead}{\huge\bfseries \@chapapp}
    {\huge\bfseries\boldmath \@chapapp}{}
    {\PackageError{hebrew-lualatex-pdf}{Cannot patch chapter-number math weight}{Unsupported document class}}%
  \patchcmd{\@makechapterhead}{\Huge \bfseries}
    {\Huge \bfseries\boldmath}{}
    {\PackageError{hebrew-lualatex-pdf}{Cannot patch chapter-title math weight}{Unsupported document class}}%
  \patchcmd{\@makeschapterhead}{\Huge \bfseries}
    {\Huge \bfseries\boldmath}{}
    {\PackageError{hebrew-lualatex-pdf}{Cannot patch starred chapter math weight}{Unsupported document class}}}
% v4.0.1: \part sits under a THREE-DEEP wrapper chain, so the class-based anchor
% is not in \@part at all. Measured with \show: hyperref rewrites \@part to
% `...\phantomsection\fi \H@old@part`, and nameref rewrites THAT to
% `\NR@gettitle{#1}\NR@part[{#1}]{#2}` — the class body ends up in \NR@part.
% v4.0.0 patched \@part and made failure fatal, which is why all seven bundled
% examples stopped building. \part is NOT a dead heading level here: 188 live uses
% across the fleet, so the bold-math fix is worth getting right rather than
% dropping. We walk the chain innermost-first and stop at the first hit. Failure is
% a WARNING, not an error: the chain is a property of hyperref/nameref versions, so
% a future package release must degrade to regular-weight part math — a cosmetic
% loss — instead of refusing to build. Chapter and section keep \PackageError,
% because those targets are in the CLASS and were verified to patch here.
\newif\iflakin@partnum \newif\iflakin@parttitle \newif\iflakin@sparttitle
\def\lakin@partnum#1#2{\iflakin@partnum\else
  \patchcmd{#1}{#2\bfseries \partname}{#2\bfseries\boldmath \partname}
    {\lakin@partnumtrue}{}\fi}
\def\lakin@parttitle#1#2{\iflakin@parttitle\else
  \patchcmd{#1}{#2 \bfseries}{#2 \bfseries\boldmath}{\lakin@parttitletrue}{}\fi}
\def\lakin@sparttitle#1#2{\iflakin@sparttitle\else
  \patchcmd{#1}{#2 \bfseries}{#2 \bfseries\boldmath}{\lakin@sparttitletrue}{}\fi}
\@ifundefined{c@chapter}{\def\lakin@pnum{\Large}\def\lakin@ptit{\huge}}%
                        {\def\lakin@pnum{\huge}\def\lakin@ptit{\Huge}}
\@ifundefined{NR@part}{}{\expandafter\lakin@partnum\expandafter\NR@part\expandafter{\lakin@pnum}%
                         \expandafter\lakin@parttitle\expandafter\NR@part\expandafter{\lakin@ptit}}
\@ifundefined{H@old@part}{}{\expandafter\lakin@partnum\expandafter\H@old@part\expandafter{\lakin@pnum}%
                            \expandafter\lakin@parttitle\expandafter\H@old@part\expandafter{\lakin@ptit}}
\@ifundefined{@part}{}{\expandafter\lakin@partnum\expandafter\@part\expandafter{\lakin@pnum}%
                       \expandafter\lakin@parttitle\expandafter\@part\expandafter{\lakin@ptit}}
\@ifundefined{NR@spart}{}{\expandafter\lakin@sparttitle\expandafter\NR@spart\expandafter{\lakin@ptit}}
\@ifundefined{H@old@spart}{}{\expandafter\lakin@sparttitle\expandafter\H@old@spart\expandafter{\lakin@ptit}}
\@ifundefined{@spart}{}{\expandafter\lakin@sparttitle\expandafter\@spart\expandafter{\lakin@ptit}}
\iflakin@partnum\else\PackageWarning{hebrew-lualatex-pdf}{Part-number math stays
  regular weight: no known \noexpand\part wrapper matched (hyperref/nameref
  version change). Cosmetic only}\fi
\iflakin@parttitle\else\PackageWarning{hebrew-lualatex-pdf}{Part-title math stays
  regular weight: no known \noexpand\part wrapper matched. Cosmetic only}\fi
% v4.12.0: THE PART NUMBER IN THE TOC WAS NOT DIRECTIONALLY SEPARATED FROM THE
% TITLE. report.cls/article.cls compose the part entry as ONE text stream --
% \addcontentsline{toc}{part}{\thepart\hspace{1em}#1} -- while EVERY other level
% goes through \numberline, which puts the number in an \hbox. A box is a bidi
% OBJECT and cannot join a neighbouring run; a bare "IV" can. So in the RTL TOC
% paragraph a Latin number and a Latin-INITIAL title merge into a SINGLE LTR run
% and are laid out left-to-right as a unit: the number lands on the wrong side of
% the row, and a mixed title is scrambled outright. Measured on \part{MCP ...}:
% the printed order was "<tail> V MCP" -- title fragment, number, title fragment.
% The chapter/section rows one line below, with the SAME Latin-initial titles,
% were correct on the same page. That contrast is the whole diagnosis.
% \en{} IS NOT A WORKAROUND, measured both ways on one fixture: it fixes a
% PURE-Latin title (\part{\en{Hooks}} renders correctly) and does NOT fix a MIXED
% one (\part{\en{MCP} <hebrew>} stays scrambled). babel's language switch changes
% font and hyphenation, not the bidi run structure -- so the defect is not
% reachable from the source side at all, which is why it belongs here.
% Boxing the number is the entire fix and is a STRICT NO-OP where the bug never
% bit: a Hebrew-titled document compiles BYTE-IDENTICAL before and after (md5,
% 3 clean runs), because there the neutral glue already sat between an L run and
% an R run. The PDF outline is byte-identical too -- hyperref's \pdfstringdef
% neutralises \mbox, so bookmark strings keep their part numbers.
% \numberline{\thepart} would also fix the direction, but it re-geometries every
% part row in the fleet and DROPS the number from the part bookmark (hyperref
% strips \numberline). \mbox changes nothing that was already right.
% The search pattern deliberately carries no # token: etoolbox cannot match a
% parameter token in a patch pattern (measured -- all three chain members
% reported "not patchable" while the pattern included #1). "\thepart \hspace
% {1em}" is brace-balanced and occurs exactly once in the wrapper body.
\newif\iflakin@parttocnum
\def\lakin@parttocnum#1{\iflakin@parttocnum\else
  \patchcmd{#1}{\thepart \hspace {1em}}{\mbox {\thepart }\hspace {1em}}
    {\lakin@parttocnumtrue}{}\fi}
\@ifundefined{NR@part}{}{\lakin@parttocnum\NR@part}
\@ifundefined{H@old@part}{}{\lakin@parttocnum\H@old@part}
\@ifundefined{@part}{}{\lakin@parttocnum\@part}
\iflakin@parttocnum\else\PackageWarning{hebrew-lualatex-pdf}{Part TOC number is
  not boxed: no known \noexpand\part wrapper matched. A part title that BEGINS
  with Latin will print its number on the wrong side of its TOC row}\fi
\patchcmd{\section}{\bfseries}{\bfseries\boldmath}{}
  {\PackageError{hebrew-lualatex-pdf}{Cannot patch section math weight}{Unsupported document class}}
\patchcmd{\subsection}{\bfseries}{\bfseries\boldmath}{}
  {\PackageError{hebrew-lualatex-pdf}{Cannot patch subsection math weight}{Unsupported document class}}
\patchcmd{\subsubsection}{\bfseries}{\bfseries\boldmath}{}
  {\PackageError{hebrew-lualatex-pdf}{Cannot patch subsubsection math weight}{Unsupported document class}}
\patchcmd{\paragraph}{\bfseries}{\bfseries\boldmath}{}{}
\patchcmd{\subparagraph}{\bfseries}{\bfseries\boldmath}{}{}
\@ifundefined{l@chapter}{}{%
  \patchcmd{\l@chapter}{\bfseries}{\bfseries\boldmath}{}{}}
\makeatother
\setcounter{tocdepth}{\lakittocdepth}
% <hllpdf:measured-toc>
% v4.16.0: Measure actual number labels, not a guessed digit count. Read only
% canonical contentsline records; never execute titles or addtocontents code
% during the measuring pass. Widths are local to this individual list and are
% available before its first row: no per-row drift and no extra aux pass.
\makeatletter
\newread\hllpdf@tocstream
\newdimen\hllpdf@toc@chapter
\newdimen\hllpdf@toc@section
\newdimen\hllpdf@toc@subsection
\newdimen\hllpdf@toc@subsubsection
\newdimen\hllpdf@toc@paragraph
\newdimen\hllpdf@toc@subparagraph
\newdimen\hllpdf@toc@figure
\newdimen\hllpdf@toc@table
\newdimen\hllpdf@tocpage
\newdimen\hllpdf@toc@sectionindent
\newdimen\hllpdf@toc@subsectionindent
\newdimen\hllpdf@toc@subsubsectionindent
\newdimen\hllpdf@toc@paragraphindent
\newdimen\hllpdf@toc@subparagraphindent
\def\hllpdf@tocdefaults{%
  \hllpdf@toc@chapter=2em
  \hllpdf@toc@section=3.2em
  \hllpdf@toc@subsection=3.9em
  \@ifundefined{c@chapter}{%
    \hllpdf@toc@subsubsection=3.2em
    \hllpdf@toc@paragraph=4.1em
    \hllpdf@toc@subparagraph=5em
  }{%
    \hllpdf@toc@subsubsection=4.1em
    \hllpdf@toc@paragraph=5em
    \hllpdf@toc@subparagraph=6em
  }%
  \hllpdf@toc@figure=2.3em
  \hllpdf@toc@table=2.3em
  \hllpdf@tocpage=\@pnumwidth\relax
}
\def\hllpdf@tocindents{%
  % Preserve the established shallow layout when its slots already fit.
  \hllpdf@toc@sectionindent=1.5em
  \@tempdima=\hllpdf@toc@chapter \advance\@tempdima by -2em
  \advance\hllpdf@toc@sectionindent by \@tempdima
  \hllpdf@toc@subsectionindent=\dimexpr
    \hllpdf@toc@sectionindent+\hllpdf@toc@section\relax
  \hllpdf@toc@subsubsectionindent=\dimexpr
    \hllpdf@toc@subsectionindent+\hllpdf@toc@subsection\relax
  \hllpdf@toc@paragraphindent=\dimexpr
    \hllpdf@toc@subsubsectionindent+\hllpdf@toc@subsubsection\relax
  \hllpdf@toc@subparagraphindent=\dimexpr
    \hllpdf@toc@paragraphindent+\hllpdf@toc@paragraph\relax
}
\def\hllpdf@tocfont@chapter{\bfseries\boldmath}
\def\hllpdf@tocfont@part{\large\bfseries\boldmath}
\def\hllpdf@tocentryfont{%
  \ifcsname hllpdf@tocfont@\hllpdf@toclevel\endcsname
    \csname hllpdf@tocfont@\hllpdf@toclevel\endcsname
  \fi}
\long\def\hllpdf@tocdiscard#1\hllpdf@tocstop{}
\def\hllpdf@tocscanrecord{%
  \@ifnextchar\contentsline{\hllpdf@tocscanline}{%
    \@ifnextchar\babel@toc{\hllpdf@tocscanlanguage}{\hllpdf@tocdiscard}}}
\long\def\hllpdf@tocscanlanguage\babel@toc#1#2#3\hllpdf@tocstop{%
  % Babel records are declarations, not arbitrary code. Track the language
  % without executing the original directive or changing the surrounding list.
  \def\hllpdf@toclanguage{#1}}
\long\def\hllpdf@tocscanline\contentsline#1#2#3#4#5\hllpdf@tocstop{%
  \def\hllpdf@toclevel{#1}%
  \ifcsname hllpdf@toc@#1\endcsname
    \hllpdf@tocscantitle#2\hllpdf@tocstop
  \fi
  % The page-number column needs its own measurement (including bold chapter
  % rows). Its extra clearance separates the page label from dot leaders.
  \setbox\@tempboxa=\hbox{\foreignlanguage{\hllpdf@toclanguage}{\hllpdf@tocentryfont #3\kern.3em}}%
  \ifdim\wd\@tempboxa>\hllpdf@tocpage
    \hllpdf@tocpage=\wd\@tempboxa
  \fi}
\def\hllpdf@tocscantitle{%
  \@ifnextchar\numberline{\hllpdf@tocscanlabel}{\hllpdf@tocdiscard}}
\long\def\hllpdf@tocscanlabel\numberline#1#2\hllpdf@tocstop{%
  \setbox\@tempboxa=\hbox{\foreignlanguage{\hllpdf@toclanguage}{\hllpdf@tocentryfont #1\kern.5em}}%
  \ifdim\wd\@tempboxa>\csname hllpdf@toc@\hllpdf@toclevel\endcsname
    \csname hllpdf@toc@\hllpdf@toclevel\endcsname=\wd\@tempboxa
  \fi}
\def\hllpdf@tocread{%
  \ifeof\hllpdf@tocstream\else
    \read\hllpdf@tocstream to\hllpdf@tocrecord
    \expandafter\hllpdf@tocscanrecord\hllpdf@tocrecord\hllpdf@tocstop
    \expandafter\hllpdf@tocread
  \fi}
\def\hllpdf@tocmeasure#1{%
  \makeatletter
  \edef\hllpdf@toclanguage{\languagename}%
  \hllpdf@tocdefaults
  \openin\hllpdf@tocstream=\jobname.#1\relax
  \hllpdf@tocread
  \closein\hllpdf@tocstream
  \hllpdf@tocindents
  \edef\@pnumwidth{\the\hllpdf@tocpage}%
  \@tempdima=\dimexpr\hllpdf@tocpage+1em\relax
  \ifdim\@tempdima>\@tocrmarg\relax
    \edef\@tocrmarg{\the\@tempdima}%
  \fi}
\let\hllpdf@originalstarttoc\@starttoc
\def\@starttoc#1{%
  \begingroup
    \def\hllpdf@tocext{#1}%
    \def\hllpdf@tocwanted{toc}%
    \ifx\hllpdf@tocext\hllpdf@tocwanted \hllpdf@tocmeasure{#1}\else
      \def\hllpdf@tocwanted{lof}%
      \ifx\hllpdf@tocext\hllpdf@tocwanted \hllpdf@tocmeasure{#1}\else
        \def\hllpdf@tocwanted{lot}%
        \ifx\hllpdf@tocext\hllpdf@tocwanted \hllpdf@tocmeasure{#1}\fi
      \fi
    \fi
    \hllpdf@originalstarttoc{#1}%
  \endgroup}
\hllpdf@tocdefaults
\hllpdf@tocindents
\@ifundefined{l@chapter}{}{%
  \patchcmd{\l@chapter}{\setlength\@tempdima{1.5em}}%
    {\setlength\@tempdima{\hllpdf@toc@chapter}}{}%
    {\PackageError{hebrew-lualatex-pdf}{Cannot size chapter TOC numbers}%
       {Unsupported contents layout: the chapter number slot was not found}}}
\renewcommand*\l@section{\@dottedtocline{1}%
  {\hllpdf@toc@sectionindent}{\hllpdf@toc@section}}
\renewcommand*\l@subsection{\@dottedtocline{2}%
  {\hllpdf@toc@subsectionindent}{\hllpdf@toc@subsection}}
\renewcommand*\l@subsubsection{\@dottedtocline{3}%
  {\hllpdf@toc@subsubsectionindent}{\hllpdf@toc@subsubsection}}
\renewcommand*\l@paragraph{\@dottedtocline{4}%
  {\hllpdf@toc@paragraphindent}{\hllpdf@toc@paragraph}}
\renewcommand*\l@subparagraph{\@dottedtocline{5}%
  {\hllpdf@toc@subparagraphindent}{\hllpdf@toc@subparagraph}}
\renewcommand*\l@figure{\@dottedtocline{1}{1.5em}{\hllpdf@toc@figure}}
\renewcommand*\l@table{\@dottedtocline{1}{1.5em}{\hllpdf@toc@table}}
\makeatother
% </hllpdf:measured-toc>
\setcounter{secnumdepth}{\lakitsecnumdepth}
% Public preamble-only caption-numbering policy. The historical within-major
% default is untouched unless this command is called. "continuous" removes the
% reset edge as well as the printed prefix, so Figure 1/2 and Table 1/2 stay
% unique across sections/chapters. Returning to "within-major" restores the
% original class-aware displays and native destination conventions. The policy
% never replaces \figcap/\tabcap or their measured link-placement mechanism.
\makeatletter
\let\lakin@majorfig\thefig
\let\lakin@majortbl\thetbl
\let\lakin@majorHfig\theHfig
\let\lakin@majorHtbl\theHtbl
\newcommand{\lakitcaptionnumbering}[1]{%
  \edef\lakin@numbering{#1}%
  \def\lakin@continuous{continuous}\def\lakin@withinmajor{within-major}%
  \ifx\lakin@numbering\lakin@continuous
    \@ifundefined{c@chapter}{%
      \@removefromreset{fig}{section}\@removefromreset{tbl}{section}%
    }{\@removefromreset{fig}{chapter}\@removefromreset{tbl}{chapter}}%
    \renewcommand{\thefig}{\arabic{fig}}\renewcommand{\thetbl}{\arabic{tbl}}%
    \def\theHfig{\arabic{fig}}\def\theHtbl{\arabic{tbl}}%
  \else\ifx\lakin@numbering\lakin@withinmajor
    \@ifundefined{c@chapter}{%
      \@removefromreset{fig}{section}\@removefromreset{tbl}{section}%
      \@addtoreset{fig}{section}\@addtoreset{tbl}{section}%
    }{%
      \@removefromreset{fig}{chapter}\@removefromreset{tbl}{chapter}%
      \@addtoreset{fig}{chapter}\@addtoreset{tbl}{chapter}}%
    \let\thefig\lakin@majorfig\let\thetbl\lakin@majortbl
    \let\theHfig\lakin@majorHfig\let\theHtbl\lakin@majorHtbl
  \else
    \PackageError{hebrew-english-latex-studio}{Unknown caption-numbering policy '#1'}%
      {Use continuous or within-major before begin{document}.}%
  \fi\fi}
\@onlypreamble\lakitcaptionnumbering
\makeatother
% Architecture readiness sentinel. Optional runtime layers require the complete
% core, so the marker is defined only after every mechanism/baseline definition.
% <hllpdf:architecture-ready>
\def\hllpdfcoreloaded{1}
% </hllpdf:architecture-ready>

% v4.22: registered file resources, not an assertion about embedded subset bytes.
% The log framing avoids long-path wrapping and preserves UTF-8 as hex.
% <hllpdf:font-inventory>
\directlua{
  local function inventory()
    local names = {}
    local unnamed = 0
    for id, item in font.each() do
      local path = item.filename
      if path and not (path == "") then names[path] = true else unnamed = unnamed + 1 end
    end
    local sorted = {}
    for path in pairs(names) do table.insert(sorted,path) end
    table.sort(sorted)
    texio.write_nl("log","HLLFONT|BEGIN|1")
    for index,path in ipairs(sorted) do
      local hex = path:gsub(".",function(c) return string.format(string.char(37) .. "02x",string.byte(c)) end)
      local count = math.ceil(string.len(hex)/48)
      texio.write_nl("log","HLLFONT|FILE|" .. index .. "|" .. count)
      for offset=1,string.len(hex),48 do texio.write_nl("log","HLLFONT|HEX|" .. hex:sub(offset,offset+47)) end
    end
    texio.write_nl("log","HLLFONT|END|" .. rawlen(sorted) .. "|" .. unnamed)
    texio.write_nl("log","")
  end
  luatexbase.add_to_callback("finish_pdffile",inventory,"hllpdf.runtime.fonts")
}
% </hllpdf:font-inventory>
     % (or paste these lines at the top)
%   \begin{document}  ...  \end{document}
%
% Compile with:  lualatex -interaction=nonstopmode doc.tex
%   Repeat until two consecutive PDFs are byte-identical (scripts/build.sh
%   does this) — a TOC document converges only on pass 3+; after two passes
%   the TOC page numbers are STALE (measured; see SKILL.md step 3).
% Font binaries are external dependencies, pinned by assets/fonts/SHA256SUMS.
% scripts/setup_fonts.sh imports a verified local installation or prior skill.
%
% PUBLIC DOCUMENT-SHELL / LAYOUT CONTRACT (introduced v4.3.1; architecture-owned v4.7.1). These hooks may be
% defined BEFORE % ============================================================================
% preamble.tex — tested LuaLaTeX preamble for Hebrew + math + figures PDFs (v4.32.0)
% Battle-tested: compiled an 80-page Hebrew quantum-computing summary AND a
% 197-page Hebrew computer-architecture anthology (TikZ diagrams, code listings,
% colored boxes) — both with 0 missing characters.
%
% Usage in your document:
%   \input{preamble.tex}     % (or paste these lines at the top)
%   \begin{document}  ...  \end{document}
%
% Compile with:  lualatex -interaction=nonstopmode doc.tex
%   Repeat until two consecutive PDFs are byte-identical (scripts/build.sh
%   does this) — a TOC document converges only on pass 3+; after two passes
%   the TOC page numbers are STALE (measured; see SKILL.md step 3).
% Font binaries are external dependencies, pinned by assets/fonts/SHA256SUMS.
% scripts/setup_fonts.sh imports a verified local installation or prior skill.
%
% PUBLIC DOCUMENT-SHELL / LAYOUT CONTRACT (introduced v4.3.1; architecture-owned v4.7.1). These hooks may be
% defined BEFORE \input{preamble.tex}; if omitted, the historical canonical
% defaults below are exact.  They deliberately control only the document shell
% and page layout.  RTL/math semantics stay in this preamble, while visual
% presentation themes (for example clean_theme.tex) load AFTER it.
%
%   \def\lakitdocumentclass{article}
%   \def\lakitclassoptions{11pt,a4paper}
%   \def\lakitgeometryoptions{margin=1.8cm,headheight=14pt}
%   \def\lakitlinespread{1.15}
%   \def\lakitparagraphstyle{block} % letters/briefs; indented for continuous prose
%   \def\lakittocdepth{1}
%   \def\lakitsecnumdepth{2}
%   \input{preamble.tex}
%
% Supported class contract: report or article.  Alternate paper sizes must also
% be declared to build.sh through EXPECTED_PAPER; the PDF gate never infers them.
% Do NOT use these hooks as a theme mechanism: theme overlays must not own page
% geometry, class selection, or direction semantics. A declared presentation-policy
% overlay may intentionally change numbering visibility after the preamble.
% ============================================================================
% Runtime architecture guard: generated from scripts/architecture_contract.py.
% The readiness sentinel remains delayed until the END of this file.
% <hllpdf:architecture-guard>
\ifdefined\hllpdfcoreloaded
  \errmessage{hllpdf: preamble.tex loaded more than once}
  \expandafter\endinput
\fi
% </hllpdf:architecture-guard>

\ifdefined\lakitdocumentclass\else\def\lakitdocumentclass{report}\fi
\ifdefined\lakitclassoptions\else\def\lakitclassoptions{12pt,a4paper}\fi
\ifdefined\lakitgeometryoptions\else\def\lakitgeometryoptions{margin=2.3cm,headheight=14pt}\fi
\ifdefined\lakitlinespread\else\def\lakitlinespread{1.3}\fi
\ifdefined\lakittocdepth\else\def\lakittocdepth{1}\fi
\ifdefined\lakitsecnumdepth\else\def\lakitsecnumdepth{2}\fi
% Paragraph organization belongs to the document layout, independently of its
% language and theme. Preserve historical class dimensions unless explicitly
% selecting block paragraphs; never suppress indentation with an everypar hack.
\ifdefined\lakitparagraphstyle\else\def\lakitparagraphstyle{indented}\fi
\edef\lakitparagraphstyle{\lakitparagraphstyle}
\def\hllpdfparagraphblock{block}
\def\hllpdfparagraphindented{indented}
\ifx\lakitparagraphstyle\hllpdfparagraphblock\else
  \ifx\lakitparagraphstyle\hllpdfparagraphindented\else
    \errmessage{hllpdf: paragraph style must be block or indented}
  \fi
\fi
% v4.3.1: the public class hook has a deliberately narrow support contract.
% Before this guard, \def\lakitdocumentclass{book} silently compiled even though
% every heading/TOC regression test targets report/article; an unsupported shell
% could therefore look "supported" until a subtle class-specific defect shipped.
% Fail before \documentclass rather than letting an undocumented combination enter
% the mechanism core. Plain-TeX comparison is used because LaTeX packages are not
% available yet at this point.
\edef\lakitvalidatedclass{\lakitdocumentclass}
\def\lakitreportclass{report}
\def\lakitarticleclass{article}
\ifx\lakitvalidatedclass\lakitreportclass
\else\ifx\lakitvalidatedclass\lakitarticleclass
\else
  \errmessage{hllpdf: unsupported class \lakitdocumentclass; use report or article}
\fi\fi
\let\lakitvalidatedclass\relax
\let\lakitreportclass\relax
\let\lakitarticleclass\relax
% Language intent is independent of document class, theme and content domain.
% The two dominant-language profiles differ in authoring intent, not direction.
% A balanced bilingual document must explicitly choose a navigation/base locale.
% <hllpdf:language-profile>
\providecommand{\lakitlanguageprofile}{hebrew}
\edef\lakitlanguageprofile{\lakitlanguageprofile}
\def\hllpdfhebrew{hebrew}\def\hllpdfenglish{english}
\def\hllpdfhebrewmixed{hebrew-with-english}
\def\hllpdfenglishmixed{english-with-hebrew}
\def\hllpdfbilingual{bilingual}
\let\hllpdfexpectedlanguage\relax
\ifx\lakitlanguageprofile\hllpdfhebrew
  \let\hllpdfexpectedlanguage\hllpdfhebrew
\else\ifx\lakitlanguageprofile\hllpdfhebrewmixed
  \let\hllpdfexpectedlanguage\hllpdfhebrew
\else\ifx\lakitlanguageprofile\hllpdfenglish
  \let\hllpdfexpectedlanguage\hllpdfenglish
\else\ifx\lakitlanguageprofile\hllpdfenglishmixed
  \let\hllpdfexpectedlanguage\hllpdfenglish
\else\ifx\lakitlanguageprofile\hllpdfbilingual
  \ifdefined\lakitmainlanguage\else
    \errmessage{hllpdf: bilingual profile requires explicit lakitmainlanguage}
  \fi
\else
  \errmessage{hllpdf: unsupported language profile \lakitlanguageprofile}
\fi\fi\fi\fi\fi
\ifdefined\lakitmainlanguage
  \edef\lakitmainlanguage{\lakitmainlanguage}
\else
  \let\lakitmainlanguage\hllpdfexpectedlanguage
\fi
\ifx\lakitmainlanguage\hllpdfhebrew\else
  \ifx\lakitmainlanguage\hllpdfenglish\else
    \errmessage{hllpdf: main language must be hebrew or english}
  \fi
\fi
\ifx\hllpdfexpectedlanguage\relax\else
  \ifx\lakitmainlanguage\hllpdfexpectedlanguage\else
    \errmessage{hllpdf: main language conflicts with language profile}
  \fi
\fi
\typeout{HLLLANG|\lakitlanguageprofile|\lakitmainlanguage}
% </hllpdf:language-profile>
\documentclass[\lakitclassoptions]{\lakitdocumentclass}
\usepackage[no-math]{fontspec}        % (1) ★ no-math — THE critical line (see recipe.md)
% layout=lists asks Babel to adapt list parshape when paragraph and page
% directions differ. Without it, minority-language list bodies lose their
% indentation and labels hang outside the text block, despite a clean log.
% The upstream mechanism preserves ordinary same-direction list geometry.
\usepackage[bidi=basic,layout=lists]{babel}
\ifx\lakitmainlanguage\hllpdfenglish
  \babelprovide[import]{hebrew}
  \babelprovide[import,main]{english}
\else
  \babelprovide[import,main]{hebrew}
  \babelprovide[import]{english}
\fi
% v4.22: bind canonical families and fallbacks to the selected installation.
% Lookup by basename or family alone is not proof of the file actually used.
% <hllpdf:font-binding>
\directlua{
  local dir = os.getenv("HEBREW_LATEX_FONT_DIR")
  if not dir or dir == "" then
    local data = os.getenv("XDG_DATA_HOME")
    if not data or data == "" or not (data:sub(1,1) == "/") then
      data = (os.getenv("HOME") or "") .. "/.local/share"
    end
    dir = data .. "/fonts/hebrew-lualatex-pdf"
  end
  if not (dir:sub(1,1) == "/") then dir = lfs.currentdir() .. "/" .. dir end
  if dir:find("[",1,true) or dir:find("]",1,true) or dir:find(string.char(0),1,true)
     or dir:find(string.char(10),1,true) or dir:find(string.char(13),1,true) then
    tex.error("hllpdf: unsupported character in font directory")
  end
  if not (lfs.attributes(dir,"mode") == "directory") then
    tex.error("hllpdf: font directory is missing")
  end
  hllpdf_font_dir = dir:gsub("/+$", "") .. "/"
  tex.sprint(string.char(92) .. "def" .. string.char(92) .. "hllpdffontdir{")
  tex.sprint(-2,hllpdf_font_dir)
  tex.sprint("}")
  luaotfload.add_fallback("hebfb",{"[" .. hllpdf_font_dir .. "DejaVuSerif.ttf]:mode=node;"})
  luaotfload.add_fallback("hebfbbf",{"[" .. hllpdf_font_dir .. "DejaVuSerif-Bold.ttf]:mode=node;"})
}
% </hllpdf:font-binding>
% v3.22: Frank Ruhl Libre ships NO italic face. Instead of silently-upright
% \emph (v<=3.20) or bold-as-emphasis (v3.21), synthesize a slanted face with
% FakeSlant — real, VISIBLE italic emphasis with correct \emph nesting
% semantics (inner \emph toggles back upright). Verified empirically at
% 300dpi: slant clearly visible, 0 errors, 0 missing glyphs.
\babelfont{rm}[Path={\hllpdffontdir},RawFeature={fallback=hebfb},
  BoldFont={FrankRuhlLibre-Bold.ttf},
  BoldFeatures={RawFeature={fallback=hebfbbf}},   % v3.56 — see hebfbbf above
  ItalicFont={FrankRuhlLibre-Regular.ttf}, ItalicFeatures={FakeSlant=0.22},
  % Bold-italic uses the weight-matched fallback and a synthetic slant.
  BoldItalicFont={FrankRuhlLibre-Bold.ttf}, BoldItalicFeatures={RawFeature={fallback=hebfbbf},FakeSlant=0.22}
  ]{FrankRuhlLibre-Regular.ttf}
% Sans fallbacks preserve family and weight for the few missing codepoints.
\directlua{luaotfload.add_fallback("hebfbsans",{"[" .. hllpdf_font_dir .. "DejaVuSans.ttf]:mode=node;"})}
\directlua{luaotfload.add_fallback("hebfbsansbf",{"[" .. hllpdf_font_dir .. "DejaVuSans-Bold.ttf]:mode=node;"})}
\babelfont{sf}[Path={\hllpdffontdir},RawFeature={fallback=hebfbsans},
  BoldFont={Heebo-Bold.ttf},
  BoldFeatures={RawFeature={fallback=hebfbsansbf}}]{Heebo-Regular.ttf}
% v3.24: Scale=MatchLowercase — DejaVu Sans Mono has an outsized x-height, so
% unscaled inline code visually SHOUTS next to the delicate Frank Ruhl body
% (measured on a 197-page build: scaled, a code atom sits in the line at the
% body's weight, and the document tightened by ~5%). Applies to listings too.
\babelfont{tt}[Path={\hllpdffontdir},RawFeature={fallback=hebfb},Scale=MatchLowercase,
  BoldFont={DejaVuSansMono.ttf},
  BoldFeatures={RawFeature={fallback=hebfb},FakeBold=1.5},
  ItalicFont={DejaVuSansMono.ttf},
  ItalicFeatures={RawFeature={fallback=hebfb},FakeSlant=0.2},
  BoldItalicFont={DejaVuSansMono.ttf},
  BoldItalicFeatures={RawFeature={fallback=hebfb},FakeBold=1.5,FakeSlant=0.2}
  ]{DejaVuSansMono.ttf}

% Explicit multilingual/English profiles use real Latin faces and italics.
% The legacy/default Hebrew profile retains its established measured typography.
% Latin Modern is provided by the TeX distribution, not the pinned font bundle;
% recorder/runtime evidence reports it as an additional toolchain resource.
\ifx\lakitlanguageprofile\hllpdfhebrew\else
  \babelfont[english]{rm}[Ligatures=TeX,RawFeature={fallback=hebfb},
    BoldFeatures={RawFeature={fallback=hebfbbf}},
    BoldItalicFeatures={RawFeature={fallback=hebfbbf}}]{Latin Modern Roman}
  \babelfont[english]{sf}[Ligatures=TeX,RawFeature={fallback=hebfbsans},
    BoldFeatures={RawFeature={fallback=hebfbsansbf}},
    BoldItalicFeatures={RawFeature={fallback=hebfbsansbf}}]{Latin Modern Sans}
\fi

\usepackage{amsmath,amssymb,mathtools}
% amsthm supplies theorem/proof machinery. Environment names remain
% document-specific; the generic proof label is localized below.
\usepackage{amsthm}
% Proof names follow the surrounding prose, including return from an island.
\addto\extrashebrew{\renewcommand{\proofname}{הוכחה}}
\addto\extrasenglish{\renewcommand{\proofname}{Proof}}
\ifx\lakitmainlanguage\hllpdfenglish
  \renewcommand{\proofname}{Proof}
\else
  \renewcommand{\proofname}{הוכחה}
\fi
% v3.24.1: equation tags under babel-bidi. A tag re-emitted by mathtools'
% showonlyrefs inside align inherits the RTL base and its parens MIRROR —
% ")2.1(" instead of "(2.1)" (plain align/equation tags are unaffected;
% verified on a minimal case + a 198-page build). LTR-isolate at the tag-BOX
% level: mathtools' show/hide + numbering logic is untouched. Do NOT override
% \tagform@ instead — that clobbers the showonlyrefs mechanism (measured:
% every tag shows, numbering breaks).
% v3.28.1: the v3.24.1 wrap must apply ONLY to tags of display environments.
% Inline \eqref in Hebrew PROSE also goes through \maketag@@@, and there the
% english-wrap MIRRORS the parens — ")1(" — while the unwrapped original is
% correct (proven on a 2x2 isolation matrix: patch x showonlyrefs; regressed
% every prose \eqref since v3.24.1, caught by the render gate's R4). The two
% call sites are indistinguishable by \ifmmode (tag boxes are text mode), so
% the display environments themselves raise a flag; \eqref in prose sees it
% false and gets the original untouched \maketag@@@.
\usepackage{etoolbox}

% babel bidi=basic (including upstream 26.9) assumes every automatic
% discretionary has a one-glyph replacement. LuaTeX can legally produce one
% with replace=nil after a hyphenated Latin value; Babel then indexes nil and the
% whole build dies. Do the narrow upstream-compatible correction: while Babel
% classifies the list, present only that malformed disc as an ordinary
% discretionary, then restore its subtype. Babel still reorders the ENTIRE box;
% there is no pcall and no successful fallback to an unreordered node list.
% If Babel throws for any other reason, the build remains fatal. The synthetic
% nil-disc probe and the real table gallery in test_all.sh pin both layers.
\AtBeginDocument{\directlua{
  if Babel and Babel.bidi and not Babel.heb_nil_disc_compat then
    Babel.heb_nil_disc_compat = true
    local original = Babel.bidi
    local disc_id = node.id('disc')
    Babel.bidi = function(head, ispar, hdir)
      local hidden = {}
      for item in node.traverse_id(disc_id, head) do
        if item.subtype == 2 and (not item.replace or not item.replace.char) then
          table.insert(hidden, item)
          item.subtype = 0
        end
      end
      local result = original(head, ispar, hdir)
      for _, item in ipairs(hidden) do item.subtype = 2 end
      return result
    end
  end
}}
\makeatletter
\newif\ifHEBdisptag
\let\ORIGmaketag@@@\maketag@@@
\def\maketag@@@#1{\ifHEBdisptag\ORIGmaketag@@@{\foreignlanguage{english}{#1}}\else\ORIGmaketag@@@{#1}\fi}
\@for\HEB@env:=equation,equation*,align,align*,gather,gather*,multline,multline*,flalign,flalign*,alignat,alignat*,eqnarray,eqnarray*\do{%
  \expandafter\AtBeginEnvironment\expandafter{\HEB@env}{\HEBdisptagtrue}}
\makeatother
% English page + Hebrew prose: Babel's equation-only tag flip uses a primitive
% leqno path that shifts the formula to the left edge and the tag into the body
% of an English page. Align's full-width alignment is unaffected. Give only this
% display a physical LTR container and native amsmath eqno/leqno. The engine then
% owns real tag width, collision handling and labels; no proxy tag box hides ink
% from its fit calculation. Babel's saved \bbl@eqnodir still restores Hebrew text
% inside math, and the prose locale/math direction are never switched.
% Close the prior RTL paragraph before normalization. After the environment's
% group restores its directions, close TeX's empty resumed display paragraph and
% use native LaTeX continuation handling: same-paragraph prose has no indent,
% while an authored blank line starts an ordinary indented RTL paragraph.
% Register after Babel's setup; remain local to equation(*) in a main-English
% document's actual RTL prose. Native Hebrew, ordinary English and explicit
% fleqn retain their original paths. No page/display dimensions are assigned.
% <hllpdf:minority-rtl-equation-geometry>
\makeatletter
\newcommand{\HEB@minorityprosecontinuation}{\par\@doendpe}
\newcommand{\HEB@minorityequationgeometry}{%
  \ifx\lakitmainlanguage\hllpdfenglish
    \ifnum\pardirection=\@ne
      \if@fleqn\else
        \par\aftergroup\HEB@minorityprosecontinuation
        \iftagsleft@\let\veqno\@@leqno
        \else\let\veqno\@@eqno\fi
        \bodydir TLT\pardir TLT\textdir TLT\relax
      \fi
    \fi
  \fi}
\AtBeginDocument{%
  \AddToHook{env/equation/begin}{\HEB@minorityequationgeometry}%
  \AddToHook{env/equation*/begin}{\HEB@minorityequationgeometry}}
\makeatother
% </hllpdf:minority-rtl-equation-geometry>
% v3.63: \shortintertext rendered Hebrew MIRROR-REVERSED and flush LEFT — an entire
% commentary line inside an align printed backwards. babel's bidi=basic patches
% \intertext@ only (luababel.def), and mathtools' \shortintertext uses the identical
% \vbox{\normalbaselines … \noindent\ignorespaces#1\par} skeleton but is never
% patched, so its paragraph keeps pardir/textdir LTR: no bidi reordering (logical
% order on the page) and flush-left. Measured against \intertext with the same
% Applying babel's own replacement to \shortintertext@ makes the two identical.
% \h{}/\mbox fix the letter order but leave the line flush LEFT, so they are not
% substitutes. Guarded + \AtBeginDocument because mathtools \let-copies at load time.
% NB \ExplSyntaxOn is REQUIRED: mathtools' internals carry expl3 names, and without
% the expl3 catcodes the `_` is read as a subscript and the patch detonates far away
% with "! Illegal unit of measure (pt inserted)" at the \end{align} (measured).
\makeatletter
\AtBeginDocument{%
  \@ifundefined{shortintertext@}{}{%
    \ExplSyntaxOn
    \bbl@sreplace\shortintertext@{\normalbaselines}%
      {\normalbaselines\ifx\bbl@eqnodir\relax\else\bbl@pardir\@ne\bbl@eqnodir\fi}%
    \ExplSyntaxOff}}
\makeatother
% geometry inherits the canonical A4 sheet from \documentclass; build.sh then
% verifies every rendered MediaBox and CropBox, so a silent Letter fallback is
% impossible to publish as a clean default build.
\usepackage[\lakitgeometryoptions]{geometry}
\linespread{\lakitlinespread}       % layout hook; canonical default remains 1.3
\ifx\lakitparagraphstyle\hllpdfparagraphblock
  \setlength{\parindent}{0pt}
  \setlength{\parskip}{6pt plus 1pt minus 1pt}
\fi
% These dimensions remain ordinary scoped TeX lengths: authors can customize
% them after core input. Language switches inherit them; lists keep their own
% indentation and narrow reference layouts establish their local paragraph gaps.
\typeout{HLLPAR|\lakitparagraphstyle}
\usepackage{xcolor}
\usepackage{mdframed}
\usepackage{listings}                 % code listings (assembly / C / pseudo-code)
\usepackage{tikz}
\usetikzlibrary{arrows.meta,positioning,calc,shapes.geometric,fit,backgrounds}
\usepackage{pgfplots}\pgfplotsset{compat=1.18}
\usepackage{quantikz}                 % quantum circuits (remove if unused)
\usepackage{booktabs}
\usepackage{array}
\usepackage{float}                    % [H] = place table exactly here (reference doc: no drifting)
\usepackage{enumitem}
\usepackage{needspace}

% ============================================================================
% ★★ v3.4 — THE DIRECTION-ISOLATION TRIO (the heart of mixing Hebrew + Latin)
% Mental model:
%   \h{...}    Hebrew (RTL) island      — Hebrew text inside an LTR context
%                                          (TikZ nodes, the english `tikzpic`).
%   \en{...}   English (LTR) isolate    — Latin text / parentheses inside a
%                                          Hebrew (RTL) context. KEEPS PARENS
%                                          AND ORDER CORRECT.
%   \code{...} \en + monospace          — code inside Hebrew / inside a frametitle.
% CRITICAL: never wrap Latin in \h{} (it is for HEBREW). Inside a TikZ node,
% Latin in \h{} renders REVERSED (Datapath->htapataD, CPU->UPC) and parens
% mirror. Isolate every Latin/paren chunk with \en{} or \code{}. See
% references/figures-boxes-listings.md.
% ---- TEXT INSIDE MATH (v3.8) ----------------------------------------------
% Put ANY text that sits inside math in a \text{}. The ONLY thing that is
% invalid in math is a bare \h/\en (= \foreignlanguage) used DIRECTLY, outside
% \text -> "\rmfamily invalid in math mode". Inside \text:
%   Hebrew label           -> \text{...}            (bare Hebrew renders fine;
%                              \h is optional — BUT NOT ALWAYS HARMLESS (v3.67):
%                              if the \h group OPENS with a nested island
%                              ($..$, \en{}, \code{}, \texttt{}, \mbox{}) the
%                              WHOLE group loses bidi and the Hebrew prints in
%                              LOGICAL order = gibberish. Measured:
%                                \text{\h{$f$ קבועה}}  -> reads "העובק"  BROKEN
%                                \text{\h{קבועה $f$}}  -> correct
%                                \text{$f$ קבועה}      -> correct (no \h)
%                              Hebrew first inside \h, or drop the \h. Flagged
%                              H-OPENS-WITH-LTR-GROUP.)
%   plain Latin word        -> \text{...}            (no marker needed)
%   Latin WITH ( ) or [ ]   -> \text{\en{...}}       (\en forces LTR; a plain
%                              \text mirrors the parens in RTL -> "Spec (TNR(" )
%   f(Hebrew/text arg)      -> \text{f}(\text{...})  -- keep the parens at MATH
%                              level, NOT \text{f(arg)} (parens inside the box mirror)
%   code in math            -> \mathtt{...}  (or \texttt{}, math-guarded below)
% NEVER write \en{...} or \h{...} straight into $...$ / \[..\] (outside \text);
% always via \text{}. (\mbox{...} also works but does not scale in scripts.)
% ============================================================================
\newcommand{\h}[1]{\foreignlanguage{hebrew}{#1}}
% v3.42.1 — \en islands do NOT hyphenate (auto OR at explicit hyphens). English hyphenation patterns were
% tearing identifiers and terms ACROSS THE RTL SEAM ("IO Regis-/ters",
% "In-/put", "Band-/width", "Per-/formance", "Multi-/ple" — all shipped on
% real pages; probe-proven: "ALUS-/rcB", "Mem-/ToReg"). An island is a label
% or a term, not English prose — a hyphen tear there is never right. A/B on
% a full 197pp lakit: identical overfull sets, same page count, ~11
% paragraphs healed, one mild underfull (badness 3525) that replaced a
% mutilated "Regis-/ters". Long English PROSE belongs in the englishprose
% environment below, which restores explicit-hyphen breaks while its paragraph
% is still active. A raw otherlanguage block inherits the global prohibition.
% (\l@nohyphenation needs \csname:
% '@' is not a letter here.)
% v3.43.2 — \exhyphenpenalty moved OUT of \en to a GLOBAL assignment below.
% Render-proven: a group-local \exhyphenpenalty is DEAD CODE — TeX reads the
% line-breaking parameters at \par time, after the \en group has closed, so
% \en{write-through} / plain not-taken still tore at the "-" ("write-/
% through" shipped on a real exam page). \language is different: it is
% recorded in the node list as text is added, so the in-group nohyphenation
% assignment above DOES work (and stays). The global penalty was A/B-measured
% on a 197pp lakit: identical overfull set (25), identical page count, the
% not-/taken tear healed. Hebrew is unaffected in practice — Hebrew compounds
% use maqaf (־), not ASCII "-".
\newcommand{\en}[1]{\foreignlanguage{english}{\language=\csname l@nohyphenation\endcsname #1}}
% Babel runs extras locally for inline islands and paragraph language changes.
% An English main document must not inherit the Hebrew no-hyphen paragraph rule.
\addto\extrashebrew{\exhyphenpenalty=10000\relax}
\addto\extrasenglish{\exhyphenpenalty=50\relax}
\ifx\lakitmainlanguage\hllpdfenglish
  \exhyphenpenalty=50\relax
\else
  \exhyphenpenalty=10000\relax
\fi
% Full English paragraphs need ordinary explicit-hyphen breaks. Use this
% environment instead of a raw otherlanguage block; the paragraph ends while
% the local penalty is still active.
\newenvironment{englishprose}
  {\par\begin{otherlanguage}{english}\exhyphenpenalty=50\relax}
  {\par\end{otherlanguage}}
\newenvironment{hebrewprose}
  {\par\begin{otherlanguage}{hebrew}\exhyphenpenalty=10000\relax}
  {\par\end{otherlanguage}}
% Use for shell text whose locale must not depend on the last body paragraph.
\DeclareRobustCommand{\doclang}[1]{\foreignlanguage{\lakitmainlanguage}{#1}}

% ============================================================================
% ★★ — fixes for defects that are INVISIBLE in the log and in the checks
% and only surface on a VISUAL RENDER of the PDF. (Each verified empirically.)
% ----------------------------------------------------------------------------
% v3.22: \emph is now REAL slanted emphasis (FakeSlant italic face defined on
% the \babelfont{rm} line above), replacing the v3.21 \emph->\textbf hack: it
% preserves \emph's toggle semantics and keeps bold free for its own role.
% Prefer bold-as-emphasis for your doc? one line: \renewcommand{\emph}[1]{\textbf{#1}}
% page number: force the digits LTR so they never reverse (18 -> 81) in the RTL
%   footer. That reversal is a NON-DETERMINISTIC bidi glitch (observed once in
%   ~80 pp); this LuaLaTeX build is not byte-reproducible, so the defensive
%   LTR wrap is the only reliable guard. Latin digits are always LTR anyway.
\renewcommand{\thepage}{\foreignlanguage{english}{\arabic{page}}}
% v3.70: \pagenumbering \gdef's \thepage and silently DESTROYS the wrap above —
% after \pagenumbering{arabic} the folio guard is gone for the whole rest of the
% document. Two collections discovered this independently and hand-restore the
% wrap in main.tex with the same comment; that knowledge now lives here instead.
% The wrapper re-applies the LTR wrap after the kernel call, for any style.
\makeatletter
\let\lakin@orig@pagenumbering\pagenumbering
\renewcommand{\pagenumbering}[1]{%
  \lakin@orig@pagenumbering{#1}%
  \renewcommand{\thepage}{\foreignlanguage{english}{\csname @#1\endcsname\c@page}}}
\makeatother
% emergencystretch: Hebrew does not hyphenate, so an unbreakable LTR island
%   (\en{...} / an inline math run) landing at the margin can leave a small
%   (~2-13pt) overfull with no legal breakpoint. This rescues ONLY those
%   inline-prose tie-breaks; it is a ceiling used only on offending lines and
%   does NOT touch display-math / listing / table overflow (those still surface
%   in the log for the content-level overflow pass). Verified: baseline overfull
%   -> 0 with this set, with no change to pagination or underfull.
\setlength{\emergencystretch}{3em}
% v3.40: inline math (almost) never breaks across lines. TeX's default
%   (binop 700 / rel 500) lets an inline run split at an operator —
%   $(1.f)\cdot 2^{e}$ shipped as "(1.f)·" ending one line and "2^e" opening
%   the next (reader-reported). In LTR that split is merely ugly; under an
%   RTL base the carried operator lands on the visually WRONG side of the
%   break and the formula reads backwards.
%   WHY 10000 (v3.45.0 — the display-first resolution). Inline math is
%   FORBIDDEN from breaking, so a torn island under an RTL base is
%   structurally impossible. The v3.40 cycle used 9999 ("break only as a last
%   resort") because a bare 10000 left ~30-39 WIDE inline equations
%   overflowing (worst 90.7pt) — but 9999's "last resort" fires OFTEN in
%   tight, non-hyphenating Hebrew: a reader found torn islands across the
%   fleet (~150 candidates; e.g. "P(y |" ending one line and "x_0)" opening
%   the next, read backwards). The step the 9999 cycle missed: every one of
%   those ~30 overflows is a wide equation that BELONGS in display. Convert
%   them to display FIRST, then 10000 gives 0 tears AND 0 overflows (measured,
%   12/12 fleet). Doctrine: inline math is ATOMIC; a wide equation goes in a
%   display, never inline. See content-style ("inline math is atomic").
\binoppenalty=10000 \relpenalty=10000
% v3.38: uniform sentence spacing. TeX's default (\nonfrenchspacing, sfcode
% 3000) is the OLD Anglo-American wide-space-after-period convention — Hebrew
% typography never had it, and modern Latin practice (Chicago et al.) dropped
% it too. Measured on real pages: post-period gaps up to 1.9-2.1x the line's
% mean word gap — the dominant "rivers" source in justified Hebrew prose.
% Survives babel language switches (probed: sfcode stays french inside \en{}
% islands and otherlanguage blocks).
\frenchspacing
% v3.38: widow/orphan protection. LaTeX's default penalties (150/150/50)
% allow a lone line at a page top/bottom almost freely. 10000 = forbid;
% this is SAFE here precisely because the document is \raggedbottom (the
% demoted line just moves on, no vertical stretch fight) — under flushbottom
% these values would cause underfull pages. raggedbottom is the report-class
% default; made explicit so a future class/option change cannot silently
% flip it and turn the 10000s harmful.
\raggedbottom
\widowpenalty=10000 \clubpenalty=10000 \displaywidowpenalty=10000
% em-dash: --- does NOT become — under Frank Ruhl (the font lacks tlig; verified
%   the width is identical with/without Ligatures=TeX, while the mechanism works
%   on Latin Modern). TYPE THE CHARACTER — (U+2014) DIRECTLY; likewise – (U+2013).
%   check_bidi_figures.py flags a --/--- that is adjacent to Hebrew.
% ============================================================================

% ---- STRUCTURAL DIRECTION-SAFETY for \texttt (v3.7) -------------------------
% A parenthesis, bracket, or $ INSIDE \texttt{} mirrors/migrates in RTL prose
% (MIPS 0($s0)->"(s0$)0", $rs->"rs$", arr[i]->"[arr[i") because \texttt sets the
% font but NOT the direction. Redefining \texttt to LTR-isolate its content fixes
% the WHOLE class at once (parens, brackets, $) so plain \texttt{lw $t0,0($s0)}
% is safe — no need to remember \code{}. Guarded by \ifmmode because
% \foreignlanguage is invalid in math (there, plain monospace is already LTR).
% NB (corrected v3.56 — the old note here said "stays tofu", which is WRONG):
% \texttt{Hebrew} is a FATAL build failure. The tt family is DejaVu Sans Mono
% (zero Hebrew) and the tt fallback hebfb is DejaVuSerif (also zero Hebrew), so
% the subsetter gets an empty glyph set and LuaTeX dies with
%   "! error: (file .../DejaVuSansMono.ttf) (type 2): there are no glyphs in the
%    subset"  -> no output PDF.
% The message names a FONT FILE, never Hebrew/\texttt/the source line. And it is
% non-local: add an unrelated Latin \code{} elsewhere and the fatal degrades into
% a SILENT DELETION of the Hebrew word. Put Hebrew in \h{}/normal text, never in
% \texttt{}/\code{}; flagged HEBREW-IN-MONO. To disable, delete this block.
\let\origtexttt\texttt
\protected\long\def\texttt#1{\ifmmode\origtexttt{#1}\else\foreignlanguage{english}{\origtexttt{#1}}\fi}
% v4.8.0: defer the direction test until TeX has inserted an alignment cell's
% math template. Unprotected expansion here tested \ifmmode too early at the
% FIRST token of aligned/array/cases/matrix cells. Engine-level \protected also
% preserves the original \texttt implementation without the self-recursion
% observed when retrofitting \DeclareRobustCommand to its saved alias.
\protected\long\def\code#1{\ifmmode\origtexttt{#1}\else\foreignlanguage{english}{\origtexttt{#1}}\fi}

% ---- \codeline — display-code line (v3.24): the code analogue of display math.
% Reader-friction rule (content-style.md §5b): inline \code{} is for a SINGLE
% ATOM only (one register/variable/mnemonic, a tight no-space expression).
% A multi-token expression — an instruction with operands, a spaced assignment,
% a register-binding list — zigzags the eye RTL<->LTR mid-sentence and, being
% an unbreakable island, line-breaks badly. Give it its own centered line with
% display-math vertical rhythm. \centerline = one box => the expression never
% wraps; \nopagebreak keeps it glued to its introducing sentence.
%   ...the Hebrew sentence ends with a colon:
%   \codeline{lw \$t1, -4(\$t0)}
%   and the prose continues (escape dollars as \$).
% v3.66 (first line of the body): flush a PENDING \item label before the \par.
% LaTeX defers the label box (\@labels) until the next \everypar, and an \hbox in
% vertical mode never fires \everypar — so `\item \codeline{...}` printed the code
% at y=69.1 and the item's own "1." BELOW it at y=103.6 (the number appears after
% its content, attached to nothing). \if@inlabel\leavevmode\fi is CONDITIONAL on
% purpose: an unconditional \leavevmode would open an empty paragraph line when
% \codeline follows ordinary prose. Verified both ways — in a list the label now
% sits on the code's line at its normal x, and after prose the geometry is
% byte-identical to before.
\makeatletter
\newcommand{\codeline}[1]{%
  \if@inlabel\leavevmode\fi
  \par\nopagebreak
  \vspace{\abovedisplayskip}\nopagebreak
  % v3.62: was \centerline{...}, i.e. \hbox to\hsize{\hss #1 \hss}. \hss has
  % INFINITE SHRINK, so TeX could always "fit" the box and NEVER reported an
  % Overfull — the same \hss mechanism as BLD-27's \item label. Measured: a wide
  % \codeline in a defbox painted its first token ON the box's left border bar and
  % sat 9.5pt outside the frame interior on both sides, with 0 Overfull, 0 Underfull
  % and all four gates PASS; in bare prose it reached 21.0pt into BOTH margins.
  % \hfil has no shrink, so the overflow is reported and build.sh's ovf_vis/ovf_big
  % counters see it; \linewidth (not \hsize) makes it box-aware, so a \codeline
  % inside a colored box centres on the box interior. Verified: a FITTING \codeline
  % keeps byte-identical glyph positions, and the box-crossing case now emits
  % "Overfull \hbox (19.03pt too wide)". A wide \codeline should be split in two.
  \hbox to\linewidth{\hfil\code{#1}\hfil}%
  \nopagebreak\vspace{\belowdisplayskip}}
\makeatother

% ---- a thin decorative rule helper ----
\newcommand{\sep}{\par\vspace{2pt}\noindent\textcolor{cRule}{\rule{\linewidth}{0.6pt}}\par\vspace{4pt}}

% ---- colors ----
\definecolor{cBlue}{HTML}{1F4E79}   \definecolor{cBlueBg}{HTML}{EAF1F8}
\definecolor{cGreen}{HTML}{2E6B4F}  \definecolor{cGreenBg}{HTML}{EAF4EE}
\definecolor{cAmber}{HTML}{9C6A00}  \definecolor{cAmberBg}{HTML}{FBF3E2}
\definecolor{cPurple}{HTML}{5B2C6F} \definecolor{cPurpleBg}{HTML}{F3ECF7}
\definecolor{cTeal}{HTML}{0E6E6E}   \definecolor{cTealBg}{HTML}{E6F4F4}
\definecolor{cRed}{HTML}{8E2B2B}    \definecolor{cRedBg}{HTML}{FBEAEA}
\definecolor{cGray}{HTML}{555555}   \definecolor{cRule}{HTML}{1F4E79}

% ---- math operators ----
\DeclareMathOperator{\Tr}{Tr}
\DeclareMathOperator{\rank}{rank}
\DeclareMathOperator{\diag}{diag}
\DeclareMathOperator{\sgn}{sgn}

% ---- bra-ket (quantikz defines \ket/\bra/\braket; add the rest safely) ----
\providecommand{\ket}[1]{\left|#1\right\rangle}
\providecommand{\bra}[1]{\left\langle#1\right|}
\providecommand{\braket}[2]{\left\langle#1\middle|#2\right\rangle}
\newcommand{\dyad}[1]{\left|#1\right\rangle\!\left\langle#1\right|}
% v3.21: general outer product |a><b| and expectation value <A> — agents reach
% for these by reflex (esp. \expval); quantikz defines neither, so their absence
% threw "Undefined control sequence" in a real build. \providecommand = safe.
\providecommand{\ketbra}[2]{\left|#1\right\rangle\!\left\langle#2\right|}
\providecommand{\expval}[1]{\left\langle#1\right\rangle}
\newcommand{\abs}[1]{\left|#1\right|}
\newcommand{\norm}[1]{\left\lVert#1\right\rVert}
\newcommand{\half}{\tfrac12}
\newcommand{\R}{\mathbb{R}}\newcommand{\C}{\mathbb{C}}\newcommand{\Z}{\mathbb{Z}}
% identity operator. NB: \mathbb{1} silently renders a WRONG glyph (msbm has no
% blackboard digits). Use \mathbb{I} = 𝕀 (works with amssymb, no extra package).
% Do NOT "fix" this back to \mathbb{1}.
\newcommand{\id}{\mathbb{I}}

% ---- physical units (v3.35): \punit — upright unit with thin space ---------
% Usage: $300\punit{K}$, $80\punit{mV/decade}$, $10^{7}\punit{cm/s}$ — the
% NUMBER AND UNIT LIVE IN ONE MATH ISLAND. Promoted verbatim from a 199-page
% production lakit (advanced-electronic-devices, dozens of uses).
% Two measured traps this design avoids (both render-proven):
%   1. siunitx SCRAMBLES under RTL — \SI{3.5}{eV} renders "eV 3.5", exponents
%      detach ("Hz 10 x 10^9"): its quantity assembly is not bidi-isolated.
%      Do NOT load siunitx in this pipeline.
%   2. A text-mode unit next to a bare-prose number is TWO adjacent LTR runs
%      and reorders ("5 mm" -> "mm5" — the BIDI-2/3 class). \punit is therefore
%      math-only ON PURPOSE: prose misuse fails LOUDLY at compile (\mathrm
%      outside math) instead of scrambling silently.
% Domain shorthands (\MeV, \fm...) stay per-document: \newcommand{\MeV}{\,\text{\en{MeV}}}
\providecommand{\punit}[1]{\,\mathrm{#1}}

% ============================================================================
% Colored boxes — mdframed (NOT tcolorbox, which breaks under bidi).
% ★★ v3.4: innerleftmargin/innerrightmargin give code listings room so their
% frame/line-numbers do not spill out the box border (see listing styles below).
% ============================================================================
% v3.61: everyline/split skips were applied to `tablebox` only (BLD-14, v3.37) —
% the six colored boxes below never got them, so they still split OPEN across a page
% break: the page-N segment gets top+sides and NO bottom rule, the page-N+1 segment
% sides+bottom and NO top rule. Measured on the shipped fleet this is a COVERAGE GAP,
% not a re-report of BLD-14: 0 split-open `tablebox` segments (that fix holds) vs 62
% split-open COLORED boxes across 7 documents. Verified after the fix: both segments
% close, and the frametitle is NOT duplicated on the continuation. BLD-14's caveat
% still applies — a box taller than a full page grinds mdframed's split-trial loop, so
% split such content manually.
% ============================================================================
% v3.72: frametitle={} leads EVERY \newmdenv option list below. mdframed options
% are GROUP-scoped: a box nested inside \begin{exbox}[frametitle={...}] used to
% INHERIT the outer title — measured, the identical title string rendered twice,
% the second riding a title-bar fill inside the INNER box's frame, for all seven
% environments in both directions. Both natural forms were broken (= preamble
% bug, not an authoring rule). The explicit reset makes every box titleless
% unless ITS OWN [frametitle={...}] says otherwise; verified both directions:
% no duplicate on the nested box, an explicit title still renders.
% Geometry is shared with the optional semantic-width table container. Themes
% select their own existing semantic style below; neither path fixes columns.
\mdfdefinestyle{hllpdfsemanticgeometry}{linewidth=1.1pt,
  innertopmargin=6pt,innerbottommargin=7pt,skipabove=9pt,skipbelow=9pt,
  innerleftmargin=9pt,innerrightmargin=9pt,everyline=true,
  splitbottomskip=6pt,splittopskip=12pt,roundcorner=3pt}
\newcommand{\hllpdfsemanticframestyle}{hllpdfsemanticgeometry}
\newmdenv[frametitle={},style=hllpdfsemanticgeometry,linecolor=cBlue,backgroundcolor=cBlueBg,
  frametitlebackgroundcolor=cBlue,frametitlefontcolor=white,
  ]{defbox}   % הגדרה
\newmdenv[frametitle={},style=hllpdfsemanticgeometry,linecolor=cGreen,backgroundcolor=cGreenBg,
  frametitlebackgroundcolor=cGreen,frametitlefontcolor=white,
  ]{thmbox}   % משפט/טענה
\newmdenv[frametitle={},style=hllpdfsemanticgeometry,linecolor=cAmber,backgroundcolor=cAmberBg,
  frametitlebackgroundcolor=cAmber,frametitlefontcolor=white,
  ]{notebox}  % הערה
\newmdenv[frametitle={},style=hllpdfsemanticgeometry,linecolor=cPurple,backgroundcolor=cPurpleBg,
  frametitlebackgroundcolor=cPurple,frametitlefontcolor=white,
  ]{exbox}    % דוגמה
\newmdenv[frametitle={},style=hllpdfsemanticgeometry,linecolor=cRed,backgroundcolor=cRedBg,
  frametitlebackgroundcolor=cRed,frametitlefontcolor=white,
  ]{warnbox}  % טעות נפוצה
\newmdenv[frametitle={},style=hllpdfsemanticgeometry,linecolor=cTeal,backgroundcolor=cTealBg,
  frametitlebackgroundcolor=cTeal,frametitlefontcolor=white,
  ]{keybox}   % רעיון מרכזי

% ---- codebox (v4.5.0): the dedicated container for code -------------------
% An OPTION, not a policy: asmblock unboxed stays canonical for a full listing,
% and a snippet inside any of the six colored boxes still works as before.
% What makes codebox worth its own environment is that it is the ONE framed box
% that keeps its margin line numbers.
%
% v4.4.0 shipped it WITHOUT them, on the strength of v3.37's "no margin knob can
% save them" plus four fresh probes that agreed. That was wrong: every probe had
% moved the code, the box, or both — never the key that carries the NUMBER
% COLUMN. Measured on the page:
%   resetmargins=false   the listing follows THIS box instead of the page,
%                        moving the code from x=75.26 to x=97.54
%   xrightmargin=26pt    under RTL this is the PHYSICAL LEFT indent, and the
%                        only key that carries the \llap'ed digits with it:
%                        x=64.28 -> 90.18, clearing the rule at 65.64 by 23.6pt
%   xleftmargin=26pt     moves the code and LEAVES THE DIGITS on the rule —
%                        the asymmetry that hid the fix for eleven versions
% resetmargins is hoisted out of the four styles (see the \lstset below them) so
% this override is possible at all; every other listing still resolves to true.
% Pinned by tests/codebox_numbers.tex + probe_lstnum.py --inside-rule, whose
% negative control strips xrightmargin and watches the digits fall back.
\definecolor{cSlate}{HTML}{3F4A56}  \definecolor{cSlateBg}{HTML}{F2F4F7}
\newmdenv[frametitle={},leftline=true,rightline=true,topline=true,bottomline=true,
  linecolor=cSlate,backgroundcolor=cSlateBg,linewidth=0.9pt,
  frametitlebackgroundcolor=cSlate,frametitlefontcolor=white,
  innertopmargin=6pt,innerbottommargin=7pt,skipabove=9pt,skipbelow=9pt,
  innerleftmargin=9pt,innerrightmargin=9pt,everyline=true,splitbottomskip=6pt,splittopskip=12pt,
  roundcorner=3pt]{codebox}   % קוד
% The two keys that make in-box numbering work, both measured on the page:
%   resetmargins=false  the listing follows THIS box instead of the page
%   xrightmargin        under RTL this is the PHYSICAL LEFT indent, and it is
%                       the only one that carries the \llap'ed number column
%                       with it. xleftmargin moves the code and leaves the
%                       digits behind on the rule (measured: code 97.54,
%                       digits 64.28 — the v3.37 defect, reproduced); with
%                       xrightmargin=26pt the digits land at x=90.18, a
%                       clearance of 23.6pt from the rule at 65.64.
% This asymmetry is why v3.4's note had the semantics backwards and v3.37
% concluded the digits could not be saved.
% <hllpdf:listing-container-state>
% v4.29: all state is local to the current environment group. A semantic
% callout masks number text, while a nested codebox restores the formatter
% saved before that mask. Save it once across successive semantic wrappers;
% an explicit numbers=none (or an explicitly empty formatter) remains valid.
\makeatletter
\newif\ifhllpdf@incodebox
\newif\ifhllpdf@englishcodebox
\newif\ifhllpdf@lstnumbermasked
\let\hllpdf@savedlstnumber\thelstnumber
\makeatother
% </hllpdf:listing-container-state>
% <hllpdf:codebox-gutter>
\makeatletter
\AtBeginEnvironment{codebox}{%
  \hllpdf@incodeboxtrue
  % Snapshot the OUTER language before asmblock selects English.
  \iflanguage{english}{\hllpdf@englishcodeboxtrue}{\hllpdf@englishcodeboxfalse}%
  \ifhllpdf@lstnumbermasked
    \let\thelstnumber\hllpdf@savedlstnumber
    \hllpdf@lstnumbermaskedfalse
  \fi
}
% listings' AddToHook is GLOBAL. Register once, and gate on local box state;
% adding an unconditional hook at every codebox entry leaks into later listings.
\lst@AddToHook{PreInit}{%
  \ifhllpdf@incodebox
    % Apply every listing override here, inside the LISTING's own group.
    % Setting margins at box entry leaks them into nested semantic callouts.
    \lstset{resetmargins=false,xrightmargin=26pt}%
    % Styles resolve before PreInit. In LTR the physical number gutter follows
    % xleftmargin; in Hebrew retain the measured inherited xrightmargin path.
    \ifhllpdf@englishcodebox
      \lstset{xleftmargin=26pt,xrightmargin=0pt}%
    \fi
    % ONE box, not two: suppress only the frame/background/skips of a listing
    % inside this codebox. The next unboxed listing keeps its own style.
    \lstset{frame=none,backgroundcolor=\color{cSlateBg},aboveskip=0pt,belowskip=0pt}%
  \fi
}
\makeatother
% </hllpdf:codebox-gutter>

% v3.37: line numbers inside a COLORED BOX are sliced by the box border —
% measured: digit span x=63.5-67.3 vs border line at x=65.2 (the border runs
% THROUGH the digits), and the inter-border gap (3.5pt) is narrower than a
% digit (3.8pt), so no margin knob can save them (under resetmargins,
% xleftmargin is inert in-box — measured: code at x=75.26 with 2.2em AND
% 3.0em). A 3-6 line boxed snippet does not need line numbers; blank the
% number TEXT inside every colored box (\thelstnumber is what gets printed;
% the style's numbers=left survives, so nothing else shifts — code position
% byte-frozen at x=75.26, verified). Unboxed listings keep their numbers.
% v3.78: the blanking is coupled to the FRAME, not to the box name —
% clean_theme.tex registers LATER hooks (registration order wins) that
% restore the number text inside the boxes it renders frameless.
% v4.5.0 RETRACTS the "no margin knob can save them" clause above. It holds only
% under resetmargins=true (the shipped default): there xleftmargin is inert
% in-box, as measured. With resetmargins=false, xrightmargin DOES carry the
% \llap'ed digit column clear of the rule — measured; the codebox \lstset that
% does it is the one above, in the codebox block. The loop below therefore
% covers SIX boxes, not "every colored box": codebox opts out and keeps its
% numbers.
% <hllpdf:semantic-listing-mask>
\makeatletter
\@for\HEB@lstbox:=defbox,thmbox,notebox,exbox,warnbox,keybox\do{%
  \expandafter\AtBeginEnvironment\expandafter{\HEB@lstbox}{%
    \ifhllpdf@lstnumbermasked\else
      \let\hllpdf@savedlstnumber\thelstnumber
    \fi
    \hllpdf@lstnumbermaskedtrue
    \def\thelstnumber{}%
    % The nearest semantic container owns its listing presentation. On exit,
    % TeX's group restores the enclosing codebox state and formatter.
    \hllpdf@incodeboxfalse}}
\makeatother
% </hllpdf:semantic-listing-mask>

% ============================================================================
% Unified TABLE presentation: subtle uniform slate frame + auto-numbered
% caption rendered INSIDE the frame. Use:
%   \begin{tablebox}\centering  <tabular>  \tabcap{<description>}\end{tablebox}
% Replaces the buggy \caption* (which emitted a stray "טבלה N : *" line) and
% provides an explicit framed option + numbered caption "טבלה N.M: תיאור".
% Optional intro text may be placed at the TOP, before the tabular, inside the box.
% ============================================================================
\definecolor{cTblLine}{HTML}{B9C4D2}  \definecolor{cTblBg}{HTML}{FAFBFC}
% v3.23: class-aware — under `article` there is no chapter counter (the promised
% report->article switch used to die with "No counter 'chapter' defined").
\makeatletter
\@ifundefined{c@chapter}{%
  % v3.31: in article the counter resets per SECTION, so an unprefixed number
  % repeated ("טבלה 1." in every section — measured on a two-section render).
  % Same conditional idea as the report branch: prefix the section number once
  % a real \section exists; a pre-section table stays "טבלה 1".
  \newcounter{tbl}[section]
  \renewcommand{\thetbl}{\ifnum\value{section}>0 \thesection.\fi\arabic{tbl}}%
}{%
  \newcounter{tbl}[chapter]
  % v3.30: a report-class doc that uses SECTIONS ONLY (chapter stays 0 —
  % the common short-doc pattern) used to print "טבלה 0.1"; the chapter
  % prefix now appears only once a real \chapter exists (render-verified
  % both ways: "טבלה 1." chapterless, "טבלה 1.1" with chapters).
  \renewcommand{\thetbl}{\ifnum\value{chapter}>0 \thechapter.\fi\arabic{tbl}}%
  % v3.64: the SAME leak in \thesection itself, left behind by v3.30. That fix
  % guarded the CAPTION counter and stopped there, so a chapterless report still
  % printed the heading "0.1 סעיף ראשון" and \ref to it printed "0.2" — measured:
  %   p1 y=65.2  '…0.1'   p1 y=144.9 '…0.2'   p1 y=183.4 '…0.2.1'
  % report.cls hardcodes \thesection as \thechapter.\@arabic\c@section, and the
  % leading 0 is the chapter counter that no \chapter ever incremented. Fixing
  % \thesection fixes the whole chain — \thesubsection and \thesubsubsection are
  % defined in terms of it. \theequation and \thetbl/\thefig were already guarded,
  % which is exactly why this survived: every NUMBER around the heading looked
  % right, so the heading's own 0 read as a style choice rather than a bug.
  % \appendix is safe: report.cls resets chapter to 0 and redefines \thechapter to
  % \@Alph\c@chapter, so an appendix \chapter makes the counter 1 and the guard
  % emits "A.1" unchanged. Render-verified in all four modes (chapterless,
  % chaptered, appendix+chapter, appendix+section-only).
  \renewcommand{\thesection}{\ifnum\value{chapter}>0 \thechapter.\fi\arabic{section}}%
}
\makeatother
% v3.37: everyline=true — if the box DOES split across pages, every segment
% gets a CLOSED frame (measured before: the split ran open-ended into the page
% edge with no bottom border, and the caption could strand alone at the next
% page top — "טבלה A.2" shipped that way). The split skips give the segments
% breathing room. In practice, with the caption GLUE below, every previously-
% split table in the corpus now moves WHOLE instead (the white-gap tradeoff,
% documented benign); everyline stays as the safety net for a genuinely
% page-tall table. NB: a table taller than ONE FULL PAGE can grind mdframed's
% split-trial loop (reproduced: minimal doc, minutes of 1pt-step Overfull
% probes) — break such a monster into two tables.
\mdfdefinestyle{hllpdftablegeometry}{leftline=true,rightline=true,topline=true,bottomline=true,
  linewidth=0.9pt,
  innertopmargin=9pt,innerbottommargin=8pt,skipabove=11pt,skipbelow=11pt,
  innerleftmargin=11pt,innerrightmargin=11pt,roundcorner=4pt}
\newcommand{\hllpdftableframestyle}{hllpdftablegeometry}
\newmdenv[frametitle={},style=hllpdftablegeometry,
  linecolor=cTblLine,backgroundcolor=cTblBg,
  % v4.2.0 (adopted): tablebox is ATOMIC. Short tableboxes near page bottoms
  % repeatedly entered mdframed's split trial, emitting benign Overfull-vbox
  % bursts and occasionally separating the intro prose from the tabular.
  % nobreak=true moves the complete short table to the next page instead. The
  % output-active Overfull-vbox gate still catches a box taller than a page;
  % everyline/split skips remain as defensive styling should that path ever
  % re-arm, but the canonical path never publishes a split tablebox.
  nobreak=true,
  everyline=true,splitbottomskip=6pt,splittopskip=12pt,
  % v3.79 (BLD-37 table side, correction): startinnercode/endinnercode run
  % INSIDE mdframed's body vbox (mdframed.sty:848 / :856+863) -- the second one
  % after \par\unskip, i.e. after line breaking, which is the only moment the
  % \tabcap anchor's true offset from the top of the box EXISTS. See
  % \lakintabboxclose below. These two keys are the ONLY tablebox change.
  startinnercode={\lakintabboxopen},endinnercode={\lakintabboxclose}]{tablebox}
% v3.37: \nopagebreak GLUE — \vsplit honors penalties, so the caption can no
% longer strand alone at the top of the next page (it takes its rows along, or
% the whole box moves). Render-verified on the real shipped orphan.
% v3.59: \NR@gettitle is what nameref injects into the stdlib title-bearing
% commands (\@sect, \@chapter, \caption) to fill \@currentlabelname. \figcap and
% \tabcap are INDEPENDENT macros — they never called it, so \@currentlabelname kept
% whatever the last heading left and \nameref{fig:x} silently printed the enclosing
% SECTION's title instead of the caption (proven in the .aux: field 3 of \newlabel
% was the heading text while field 1 and the anchor `fig.1.1` were correct — so this
% is a THIRD, distinct head next to BLD-13's counter fix and BIDI-38's autoref-name
% fix, neither of which writes that field). \@ifundefined keeps it inert if
% hyperref/nameref is ever dropped. Guarded helper because \figcap is defined
% after \makeatother.
% NB both helpers below need \makeatletter — this part of the preamble is OUTSIDE
% any letter-catcode region, so an unwrapped \@currentlabel parses as LaTeX's \@
% and dies with "You can't use \spacefactor in vertical mode" (measured).
\makeatletter
\newcommand{\lakinrtitle}[1]{\@ifundefined{NR@gettitle}{}{\NR@gettitle{#1}}}
% v3.59: \hitem{<letter>}{<key>} — a REFERENCEABLE Hebrew-lettered list item.
%   \begin{enumerate} \hitem{א}{it:a} הסעיף הראשון. \end{enumerate}  ->  \ref{it:a} = (א)
% Why this exists: automatic Hebrew list labels fail silently, so SKILL.md mandates
% an EXPLICIT \item[(\textbf{א})] — but LaTeX puts \refstepcounter INSIDE
% \if@noitemarg (latex.ltx:12964), and \item[...] never sets that flag. So an
% explicit item steps no counter and never touches \@currentlabel: a \label there
% binds to the AMBIENT one and \ref prints the CHAPTER/SECTION number where the
% item letter belongs (proven in the .aux: {chapter.1}). Same doctrine as \figcap —
% own the label state before \label. \phantomsection matters: without it the
% hyperlink still lands on the chapter head (measured); with it the anchor is
% section*.N and the bookmark outline stays clean.
\newcommand{\hitem}[2]{\item[(\textbf{#1})]%
  \def\@currentlabel{(#1)}\def\@currentlabelname{#1}\phantomsection\label{#2}}
% NB (v3.66): \hitem is \ref-only. \autoref{key} calls it a SECTION ("סעיף (א)")
% because \phantomsection brands the anchor section*.N, and \nameref{key} prints
% the bare letter. For \autoref-able items prefer \hebenumlabels below.
%
% v3.66: \hebenumlabels — OPT-IN automatic Hebrew enumerate letters, done right.
%   {\hebenumlabels\begin{enumerate}\item\label{q:a} ...\end{enumerate}}
%   -> labels (א) (ב) (ג); \ref{q:a} = א; anchor Item.N (real, per-item).
% This retires the belief that made \item[(\textbf{א})] + \hitem necessary.
% "Automatic Hebrew labels fail silently" was TRUE but unexplained, and the
% mechanism is now bisected (plain article vs +enumitem, measured): ENUMITEM
% PRE-EXPANDS the label expression in an \edef BEFORE \refstepcounter, with the
% counter still at 0. Tokens it shields by name (\arabic,\alph,\roman) survive;
% a plain \newcommand \ifcase wrapper evaluates at 0 -> the EMPTY zero-branch ->
% the label prints as "()" with the letter GONE (and \the\c@enumi prints a
% literal 0, \number\value dies fatally pointing at \item). The cure is one
% property, not a doctrine: the wrapper must be ROBUST (survives the \edef
% unexpanded), while \theenumi stays EXPANDABLE (so \refstepcounter writes the
% real letter to the .aux). Verified: labels (א)(ב)(ג), \ref -> א/ב/ג each item
% its own value, .aux = {א}{...}{Item.1} — versus the explicit-\item doctrine
% where BIDI-50/\hitem exist purely to patch the counter that never stepped.
\def\@hebalph#1{\ifcase#1\or א\or ב\or ג\or ד\or ה\or ו\or ז\or ח\or ט\or י\or
  יא\or יב\or יג\or יד\or טו\or טז\or יז\or יח\or יט\or כ\else\@ctrerr\fi}
\DeclareRobustCommand*\hebalph[1]{\expandafter\@hebalph\csname c@#1\endcsname}
\newcommand{\hebenumlabels}{%
  \renewcommand{\theenumi}{\@hebalph\c@enumi}%
  \renewcommand{\labelenumi}{(\textbf{\hebalph{enumi}})}}
\makeatother
% v3.63: \mbox the LABEL. \thetbl ("7.1") and a caption description that STARTS
% WITH A DIGIT ("32 הרגיסטרים…") are two European-Number runs separated only by
% neutrals (the label's own "." and \enspace). Under the RTL paragraph base the UBA
% resolves those neutrals to the RTL level, so both numbers sit in ONE reversed span
% and EXCHANGE VISUAL POSITIONS. Measured on the shipped PDF: a table whose real
% number is 7.1 printed "טבלה 32" — a different, EXISTING table number — while \ref
% correctly printed 7.1, so the caption silently contradicted its own cross-reference,
% and the description lost its leading count. The source contains no adjacency at all;
% it is manufactured by this macro, which is why every source linter was blind.
% \mbox makes the label one atomic box so nothing can reorder across it. Verified
% zero-regression: an ordinary Hebrew caption keeps identical x-geometry.
% One paragraph style covers label, number, separator and description.
% Explicit project/template overrides remain possible; the label hook is kept
% for compatibility, with an identity default. Counters/anchors stay unchanged.
\newcommand{\lakitcaptionstyle}{\normalfont\small\color{black}\hypersetup{allcolors=black}}
\newcommand{\lakitcaptionlabelstyle}[1]{#1}
\newcommand{\lakitcaptionpunctuation}{:}
% Native floats, captionof and longtable use the same presentation policy.
% subcaption is optional and inherits the same font rather than a second size.
\usepackage{caption}
\DeclareCaptionFont{lakit}{\lakitcaptionstyle}
\DeclareCaptionLabelFormat{lakit}{\lakitcaptionlabelstyle{#1~\en{#2}}}
\DeclareCaptionLabelSeparator{lakit}{\lakitcaptionpunctuation\enspace}
\captionsetup{font=lakit,labelfont={},textfont={},labelformat=lakit,
  labelsep=lakit,justification=centering}
\makeatletter
\AtBeginDocument{\@ifpackageloaded{subcaption}{%
  \captionsetup[sub]{font=lakit,labelfont={},textfont={},labelsep=lakit}}{}}
\makeatother
% v3.64: OPT-IN list-of-figures/tables entry via an optional short caption.
% \listoffigures and \listoftables are inherited from report.cls and LOOK available,
% but \figcap/\tabcap step the custom \c@fig/\c@tbl counters and never call
% \addcontentsline — so both printed their heading and NOTHING ELSE. Measured: a doc
% with one \figcap and one \tabcap rendered page 1 as 'רשימת האיורים\n1\n', and the
% .lof held only \@chapter's \addvspace, zero \contentsline entries. An empty list
% under a correct Hebrew heading reads as "this document has no figures".
%
% WHY OPTIONAL, AND NOT AUTOMATIC. Writing #1 unconditionally was implemented first
% and REJECTED: a caption containing \codeline compiles CLEANLY today (verified on
% the pre-v3.64 preamble: 0 errors, PDF produced) but becomes FATAL once it is also
% written to the .lof — \codeline is \hbox to\linewidth{...}, and a full-width box in
% a list entry dies with "\pdfextension endlink cannot be used in vertical mode", no
% output PDF. Gaining a feature nobody uses (0 of 13 collections call \listoffigures)
% at the price of breaking a construct that works is a bad trade. Stock \caption has
% exactly this limitation and exactly this escape hatch, so the convention is familiar.
% \en{}, \code{}, \h{} and inline math in the short text are all safe — verified, they
% survive the write and typeset correctly in the list (RTL geometry confirmed: number
% rightmost at x=509.9, title left of it, dots and folio at the far left).
% With no optional argument the macro is byte-identical to v3.63 — zero risk to every
% existing document, and \listoffigures lists exactly the captions that opted in.
% BLD-37, TABLE side (the \figcap v3.74 idea, adapted): clicking a \ref to a
% table used to land MID-table — hyperref's automatic anchor inside
% \refstepcounter was emitted on the still-open horizontal line whose TAIL is
% the tabular (an hbox whose baseline is the table's vertical CENTER), so the
% jump scrolled the top half of the table above the viewport. The v3.74 entry
% deferred this side claiming "tablebox splits — no single box to measure";
% that deferral dissolves: the measurement target is the TABULAR itself at
% \tabcap ENTRY, while it is still the tail of the open list — the same
% situation \lakin@grabfig handles at an island's end (and the fleet has 0
% split tableboxes). So \lakin@grabtab runs as the FIRST token, before any
% typesetting can push the box away, and the anchor is re-emitted raised above
% the measured tabular. \lakintabraise is a SEPARATE register from
% \lakinfigraise (defined with it, below the LTR-island block): a captionless
% island intentionally leaves the fig register set, and a table must never
% inherit that stale figure height.
%
% WHY the anchor stays AT THE ENTRY SPOT and is NOT deferred into the caption
% paragraph the way \figcap defers it: at \par TeX DROPS a final glue item,
% and the automatic anchor's smash box sitting BEHIND the interword glue kept
% that glue alive — so every centered tabular has always sat half a space off
% true center. Deferring the anchor removed that box and was measured as 982
% differing glyph records: EVERY tabular line shifted 1.30pt (2.61pt for
% pattern 7's two minipages) while captions stayed put. Emitting the raised
% anchor at the old spot — through hyperref's own \Hy@raisedlink, exactly the
% macro the automatic path uses — keeps the line's node structure, and so the
% A/B glyph diff at 0, BY CONSTRUCTION. \Hy@raisedlink reads the raise from
% \HyperRaiseLinkDefault (default \baselineskip = the old mid-table dest); the
% group-local \edef freezes measured-height+42pt into it, and \smash inside
% \Hy@raisedlink hides any raise from the line's extent. \HyCnt@ProvideTheH-
% Counter keeps the dest NAME byte-identical to the automatic one (it defines
% \theHtbl, which \hyper@makecurrent prefers — without it the first measured
% caption was named tbl.1 where the automatic path said tbl.0.1, measured).
% vmode entry, or anything but a box at the tail (no tabular to measure),
% falls back to the plain \refstepcounter — byte-identical to the old path.
\makeatletter
\newcommand{\tabcap}[2][]{\lakin@grabtab
  % v3.79: arm the box-close correction. Set unconditionally and in BOTH
  % branches below's scope, so the plain-\refstepcounter fallback is corrected
  % too; \lakintabboxopen clears it at every \begin{tablebox}, so a \tabcap
  % outside a tablebox can never leave it armed for a later box.
  \global\lakintabcapseentrue
  \ifdim\lakintabraise>0pt
    \H@refstepcounter{tbl}\HyCnt@ProvideTheHCounter{tbl}\hyper@makecurrent{tbl}%
    \begingroup
      % A frameless table targets its measured top. The frame's historical
      % cushion belongs only to tablebox and is corrected at box close.
      \edef\HyperRaiseLinkDefault{\the\lakintabraise}%
      \iflakintabframed
        \edef\HyperRaiseLinkDefault{\the\dimexpr\lakintabraise+42pt\relax}%
      \fi
      \Hy@raisedlink{\hyper@anchorstart{\@currentHref}\hyper@anchorend}%
    \endgroup
    \global\lakintabraise 0pt\relax
  \else
    \refstepcounter{tbl}%
  \fi
  \lakinrtitle{#2}\par\nopagebreak\addvspace{6pt}%
  \def\lakin@short{#1}\ifx\lakin@short\@empty\else
    \addcontentsline{lot}{table}{\protect\numberline{\thetbl}{\protect\ignorespaces #1}}\fi
  {\centering\lakitcaptionstyle\mbox{\lakitcaptionlabelstyle{\tablename~\en{\thetbl}\lakitcaptionpunctuation}}\enspace #2\par}}
\makeatother

% ============================================================================
% Unified FIGURE caption — the \tabcap twin for TikZ/plot figures (v3.34.0).
% A figure (tikzpic) sits inline in the prose; put \figcap BELOW it for a
% numbered, \ref-able caption "איור N.M: תיאור" in the same style as \tabcap:
%   \begin{tikzpic} ... \end{tikzpic}
%   \figcap{\label{fig:foo}description}     % \label FIRST, inside the argument
% WHY a macro, not a hand-written gray line: a manual
%   {\centering\small\textcolor{cGray}{איור: ...}} with a \label inside has NO
% \refstepcounter, so the \label binds to the AMBIENT \@currentlabel (the last
% equation) — \ref then prints an equation number AS the figure number
% (measured on a real page: "איור 48" where 48 is equation 48). \figcap steps a
% real \c@fig counter first, so \label is correct, numbering is consistent, and
% refs are stable. Class-aware like \thetbl (chapter- or section-prefixed).
% ============================================================================
\makeatletter
\@ifundefined{c@chapter}{%
  \newcounter{fig}[section]
  \renewcommand{\thefig}{\ifnum\value{section}>0 \thesection.\fi\arabic{fig}}%
}{%
  \newcounter{fig}[chapter]
  \renewcommand{\thefig}{\ifnum\value{chapter}>0 \thechapter.\fi\arabic{fig}}%
}
\makeatother
\makeatletter
\newcommand{\figcap}[2][]{%
  \ifdim\lakinfigraise>0pt
    % v3.74: step WITHOUT hyperref's automatic anchor, then emit the anchor
    % raised above the measured figure. \H@refstepcounter is hyperref's saved
    % pre-hook \refstepcounter (sets the counter AND \@currentlabel — \label
    % in the caption arg binds exactly as before); \hyper@makecurrent supplies
    % \@currentHref, so the dest NAME is identical to the old automatic one —
    % \ref/\autoref/bookmarks are unchanged, only the dest'S POSITION moves.
    \H@refstepcounter{fig}\hyper@makecurrent{fig}%
    % The anchor is emitted INSIDE the caption paragraph below, as a zero-size
    % raisebox — NOT here in vertical mode, where a \raisebox starts an empty
    % paragraph and adds a blank line (measured: every glyph after the caption
    % shifted +18.8pt). Inside an existing line a [0pt][0pt] box is inert.
    \edef\lakin@anch{\@currentHref}%
    % \edef FREEZES the measured height into the pending macro: the deferred
    % \raisebox is typeset inside the caption paragraph AFTER the register is
    % reset to the 0pt sentinel below — a plain \def read 0 at use time
    % (measured: LAKINRAISE=0.0, dest moved 14pt instead of above the figure).
    \edef\lakin@pend{\noexpand\raisebox{\the\dimexpr\lakinfigraise+42pt\relax}[0pt][0pt]{%
      \noexpand\hyper@anchorstart{\lakin@anch}\noexpand\hyper@anchorend}}%
    \global\lakinfigraise 0pt\relax
  \else
    \refstepcounter{fig}\let\lakin@pend\relax
  \fi
  \lakinrtitle{#2}\par\nopagebreak\addvspace{4pt}%
  \def\lakin@short{#1}\ifx\lakin@short\@empty\else
    \addcontentsline{lof}{figure}{\protect\numberline{\thefig}{\protect\ignorespaces #1}}\fi
  {\centering\lakin@pend\lakitcaptionstyle\mbox{\lakitcaptionlabelstyle{\figurename~\en{\thefig}\lakitcaptionpunctuation}}\enspace #2\par}}
\makeatother

% ============================================================================
% LTR islands.  EVERY tikzpicture / quantikz MUST live inside one, or the whole
% figure's coordinates and Latin labels flip. (Hebrew node text still uses \h{}.)
% ============================================================================
% v3.74 (BLD-37): the \figcap hyperlink DESTINATION used to land at the caption
% line — BELOW the figure — so clicking a figure reference scrolled the figure
% off-screen (measured fleet-wide: 97/97 fig dests in one collection sat 19-54pt
% under their figure's ink; a 347pt figure was fully hidden after the jump).
% hypcap-style fix in two parts: (1) at the island's END, \lakin@grabfig grabs
% the just-typeset centered line with \lastbox, measures it, and puts it back
% with \nointerlineskip — the interline glue above the line was already
% contributed and is NOT re-added, so the page geometry is unchanged to the
% glyph (verified); (2) \figcap then steps the counter WITHOUT hyperref's
% automatic anchor (\H@refstepcounter = the pre-hyperref \refstepcounter) and
% emits the anchor itself, raised by the measured height — the dest lands ABOVE
% the figure top. A captionless island leaves numbering untouched (nothing is
% stepped at the island; the grab flag is cleared at every island begin, so a
% stale height from a captionless figure cannot leak into a later caption).
\newsavebox{\lakinfigbox}
\newdimen\lakinfigraise  % real \dimen: tex.setdimen cannot target a \newlength SKIP register (measured lua error)
% BLD-37 table side: a SEPARATE register for \tabcap. The two must not share —
% a captionless island INTENTIONALLY leaves \lakinfigraise set (the grab runs
% at every island END; only the NEXT island's begin clears it), so a shared
% register would hand a \tabcap after a captionless figure that stale FIGURE
% height. \lakintabraise is written only by \lakin@grabtab at \tabcap entry
% and consumed/reset inside \tabcap itself, so it can never go stale.
\newdimen\lakintabraise
\makeatletter
% v3.74b: measure WITHOUT touching the list. The first draft grabbed the box
% with \lastbox and put it back — and \lastbox DISCARDS the node's baseline
% shift, so a vertically-anchored quantikz circuit rose 2.57pt and everything
% after it sank 5.32pt (measured on the A/B glyph diff). LuaTeX lets us READ
% the tail node's dimensions without removing anything: geometry is identical
% by construction. \lakinfigraise=0pt is the "no figure measured" sentinel.
\def\lakin@grabfig{%
  \ifhmode
    \directlua{
      local n = tex.nest[tex.nest.ptr].tail
      local GLUE, PEN = node.id('glue'), node.id('penalty')
      while n and (n.id == GLUE or n.id == PEN) do n = n.prev end
      local HL, VL = node.id('hlist'), node.id('vlist')
      if n and (n.id == HL or n.id == VL) then
        tex.setdimen('global', 'lakinfigraise',
                     (n.height or 0) + (n.depth or 0))
      else
        tex.setdimen('global', 'lakinfigraise', 0)
      end}%
  \else
    \global\lakinfigraise 0pt\relax
  \fi}
% BLD-37 table side: \lakin@grabfig's twin, NOT a shared helper — the Lua body
% differs in exactly one place, the skip set. Between \end{tabular} and \tabcap
% an otherlanguage wrapper (patterns 3/5 of example_tables) leaves a \write
% whatsit and a dir node ON TOP of the interword glue (measured node dump:
% "whatsit/write <- dir <- glue(2.62) <- hlist"), so the table reader must
% also step over whatsit/dir markers — all zero-size — to reach the tabular.
% \lakin@grabfig deliberately does NOT get the wider skip set: its glue/penalty
% body is the one the v3.74 A/B proof froze, and widening it would newly
% measure (and so newly raise) anchors in island shapes that today take the
% plain path — an unproven behavior change on the FIG side. Anything that is
% not a box behind those markers still hits the 0pt sentinel and \tabcap falls
% back to the plain \refstepcounter.
% SECOND deliberate difference: this stores the box's HEIGHT, not height+depth.
% \figcap's anchor is emitted in the caption paragraph BELOW the figure, so it
% must climb the full extent; \tabcap's anchor sits ON the tabular's own line,
% whose baseline is the vcentered tabular's MIDPOINT — height alone reaches
% the top edge, and adding the depth was measured to overshoot by half the
% table (dests 80-135pt above their frame, INSIDE the preceding table's box).
\def\lakin@grabtab{%
  \ifhmode
    \directlua{
      local n = tex.nest[tex.nest.ptr].tail
      local GLUE, PEN = node.id('glue'), node.id('penalty')
      local WHAT, DIR = node.id('whatsit'), node.id('dir')
      while n and (n.id == GLUE or n.id == PEN
                   or n.id == WHAT or n.id == DIR) do n = n.prev end
      local HL, VL = node.id('hlist'), node.id('vlist')
      if n and (n.id == HL or n.id == VL) then
        tex.setdimen('global', 'lakintabraise', n.height or 0)
      else
        tex.setdimen('global', 'lakintabraise', 0)
      end}%
  \else
    \global\lakintabraise 0pt\relax
  \fi}
\makeatother
% ============================================================================
% v3.79 (BLD-37 table side, CORRECTION): put the anchor at the TOP OF THE BOX,
% instead of guessing the distance to it.
% ============================================================================
% v3.77 raised the \tabcap anchor by (height of the box at the tail of the open
% horizontal list) + 42pt. That is a MEASUREMENT taken BEFORE line breaking, and
% it is short by every term it cannot see. Fleet-measured on the 12 built
% collections (299 tbl.* dests, ownership established through each table's own
% CAPTION, never through "which frame encloses the dest"): 3 dests landed INSIDE
% their own tablebox, below the frame top —
%   computer-architecture       tbl.14.1  -46.4pt   three stacked tabulars
%   computer-architecture       tbl.30.5  -13.5pt   tabular 2.85pt over \linewidth
%   advanced-electronic-devices tbl.B.1   -59.2pt   tabular 2.56pt over \linewidth
% Two missing terms, both reproduced in a minimal document (probe P2/P3):
%   (1) VERTICAL LIST ALREADY CONTRIBUTED inside the box — a second and third
%       tabular, or a heading above the table. The tail read sees only the LAST
%       box, so 77.7pt of stacked tabulars is invisible to the raise.
%   (2) WHICH LINE THE ANCHOR LANDS ON. When the tabular alone overfills
%       \linewidth (459.6121pt in a tablebox), the interword glue after
%       \end{tabular} becomes the chosen breakpoint and \Hy@raisedlink is
%       emitted on the NEXT line — one (depth + lineskip) lower. That term does
%       not EXIST at \tabcap time: line breaking has not happened yet. No
%       measurement taken there can ever recover it. (Proof: probe P3, the same
%       table 2.84912pt too wide, fails at -13.5pt; probe P4, the identical
%       table with \tabcolsep 2pt so it fits, is correct at +32.0pt.)
% So the raise is not computed at \tabcap at all. It is CORRECTED at
% \end{tablebox}, from inside mdframed's own body vbox (endinnercode runs after
% \par\unskip, i.e. after line breaking, while tex.nest[tex.nest.ptr] is still
% that vbox's vertical list). There the anchor's true offset from the top of the
% box content is not estimated but WALKED, and the \Hy@raisedlink shift is set
% so the dest lands 42pt above it — which is 42pt - innertopmargin(9pt) -
% linewidth(0.9pt) = 32.0pt above the FRAME top, the exact position the 295
% already-correct dests have today. That is why this is a correction and not a
% re-placement: the good sites do not move.
%
% WHY THE GLYPH DIFF IS 0 BY CONSTRUCTION — the only thing written is the
% `shift` field of the hlist that \Hy@raisedlink already built, and the guard
% below refuses to write it unless that box is 0x0x0. A zero-size box carries no
% dimension into its parent's packing, the packing already happened, and shift
% is not an input to it. Nothing else in the list is touched: no node is added,
% removed, resized or reordered. Measured, not assumed: assets/example_tables.tex
% and the whole computer-architecture collection A/B at 0 differing pymupdf
% rawdict glyph records.
%
% WHY NOT the alternatives (all prototyped and measured before this one):
%   * accumulate the vertical list at \tabcap time — fixes term (1), CANNOT fix
%     term (2); 2 of the 3 real failures survive. A fix that leaves the majority
%     of the defect standing is the false all-clear this repo exists to avoid.
%   * \zref-savepos two-pass — correct on all shapes, but it makes the dest
%     POSITION aux-derived (DEFECTS BLD-1/2), a single-pass compile silently
%     ships an unraised anchor, and the bare y-difference has no page guard.
%   * emit the anchor at \begin{tablebox} and let \tabcap adopt it — needs the
%     tbl counter stepped at box OPEN to keep the dest NAME byte-identical; a
%     captionless tablebox then either shifts every later table number or leaves
%     an orphan dest that the next table duplicates (a HARD build gate). This
%     correction touches neither the counter nor the name: \tabcap's naming path
%     is v3.77's, unchanged, so resolve_names() is byte-identical (verified).
%   * suppressing the line break (a penalty before the glue) — that MOVES the
%     tabular: the break currently discards the interword glue, so forbidding it
%     re-centers the two overfull tables. Reflow, rejected.
%
% SCOPE. The correction is armed only by \tabcap ITSELF (\lakintabcapseen),
% cleared at every box open, so a captionless tablebox is a no-op and a \tabcap
% hosted OUTSIDE a tablebox (the fleet has exactly one, computer-architecture
% tbl.14.2 in an exbox) keeps v3.77's measured raise. Both fallbacks are
% therefore never worse than today; and because hyperref's AUTOMATIC anchor uses
% the same \Hy@raisedlink shape, the plain-\refstepcounter path (vmode entry, or
% prose between \end{tabular} and \tabcap — probe P6, -51.5pt) is corrected too.
% KNOWN LIMIT — a tablebox that SPLITS across pages. The offset is measured from
% the top of the UNSPLIT body, so a dest in segment 2 gets raised past its own
% segment. MEASURED on a forced split (probe: 40 filler lines + a table): the
% dest went from y=585.0 on page 2 (inside segment 2, at the caption — the old
% BLD-37 symptom) to y=-113.4, i.e. ABOVE the page top. Viewers clamp that to the
% page top, which happens to show segment 2 whole, but an out-of-page coordinate
% is not something to ship on purpose. Hence the \textheight guard in
% lakintab.fix: a body TALLER THAN THE TEXT BLOCK cannot fit on any page, so it
% must split — the correction is skipped and v3.77's behaviour stands unchanged
% (verified: the probe returns to y=585.0). The guard is deliberately EXACT, not
% a page-fitting heuristic: it is sound (a taller-than-page body always splits)
% but PARTIAL (a shorter body that starts low on a page could still split). That
% residue is unreachable in practice — v3.37's everyline + the caption's
% \nopagebreak glue move such boxes WHOLE instead (probe N7: a 22-row table
% pushed down by 430pt moved whole, 0 segments), and the fleet has 0 split
% tableboxes, re-counted after this change. Reimplementing mdframed's own
% page-fitting arithmetic to close the residue would be a heuristic with its own
% failure mode, which is the trade this repo has repeatedly refused.
\newif\iflakintabcapseen
\newif\iflakintabframed
% The Lua lives in ONE \directlua at load time, and is written with NO
% backslash, %, # or ~ : \directlua expands its argument like \edef, so an
% active ~ would expand to \nobreakspace and a % would eat the rest of the line
% (the second is how the first draft of this block died). "not (a == b)" stands
% in for Lua's inequality operator for the same reason.
\directlua{
  lakintab = { }
  lakintab.HL = node.id("hlist")
  lakintab.VL = node.id("vlist")
  lakintab.RU = node.id("rule")
  lakintab.GL = node.id("glue")
  lakintab.KE = node.id("kern")
  lakintab.WH = node.id("whatsit")
  lakintab.TM = node.id("temp")
  lakintab.WD = node.subtype("pdf_dest")
  function lakintab.d(v)
    if v == nil then return 0 end
    if v < -1000000000 then return 0 end
    return v
  end
  function lakintab.scan(list, y, horiz, parent, target)
    local x = list
    while x do
      local i = x.id
      if i == lakintab.HL or i == lakintab.VL then
        if horiz then
          local b = y + x.shift
          if i == lakintab.HL then
            lakintab.scan(x.head, b, true, x, target)
          else
            lakintab.scan(x.head, b - lakintab.d(x.height), false, x, target)
          end
        else
          if i == lakintab.HL then
            lakintab.scan(x.head, y + lakintab.d(x.height), true, x, target)
          else
            lakintab.scan(x.head, y, false, x, target)
          end
          y = y + lakintab.d(x.height) + lakintab.d(x.depth)
        end
      elseif i == lakintab.WH then
        if x.subtype == lakintab.WD then
          local id = x.dest_id
          if type(id) == "string" and string.sub(id, 1, 4) == "tbl." then
            if horiz and parent and parent.width == 0
               and parent.height == 0 and parent.depth == 0 then
              parent.shift = parent.shift + (target - y)
            end
          end
        end
      elseif not horiz then
        if i == lakintab.RU then
          y = y + lakintab.d(x.height) + lakintab.d(x.depth)
        elseif i == lakintab.GL or i == lakintab.KE then
          y = y + lakintab.d(x.width)
        end
      end
      x = x.next
    end
  end
  function lakintab.fix(target, maxh)
    local nst = tex.nest[tex.nest.ptr]
    if nst == nil then return end
    local h = nst.head
    if h == nil then return end
    if h.id == lakintab.TM then h = h.next end
    local tot = 0
    local x = h
    while x do
      local i = x.id
      if i == lakintab.HL or i == lakintab.VL or i == lakintab.RU then
        tot = tot + lakintab.d(x.height) + lakintab.d(x.depth)
      elseif i == lakintab.GL or i == lakintab.KE then
        tot = tot + lakintab.d(x.width)
      end
      x = x.next
    end
    if tot > maxh then return end
    lakintab.scan(h, 0, false, nil, target)
  end
}
% @-FREE NAMES ON PURPOSE: these two are named in the `tablebox` key list at
% \newmdenv above, which is tokenised OUTSIDE any \makeatletter region — a
% \lakin@... there would tokenise as \lakin followed by the letters @tabopen.
\newcommand{\lakintabboxopen}{\lakintabframedtrue\global\lakintabcapseenfalse}
\newcommand{\lakintabboxclose}{%
  \iflakintabcapseen
    \directlua{lakintab.fix(-\number\dimexpr\lakintabcushion\relax,
                            \number\dimexpr\textheight\relax)}%
    \global\lakintabcapseenfalse
  \fi\lakintabframedfalse}
% The cushion above the box's CONTENT top. 42pt puts the dest 32.0pt above the
% frame top (42pt - innertopmargin 9pt - linewidth 0.9pt), i.e. exactly where
% the 295 correct dests already are — the correction moves only the broken ones.
\newcommand{\lakintabcushion}{42pt}
\newenvironment{qcirc}{\par\medskip\global\lakinfigraise 0pt\relax\begin{otherlanguage}{english}\begin{center}}
  {\csname lakin@grabfig\endcsname\end{center}\end{otherlanguage}\par\medskip}
\newenvironment{tikzpic}{\par\medskip\global\lakinfigraise 0pt\relax\begin{otherlanguage}{english}\begin{center}}
  {\csname lakin@grabfig\endcsname\end{center}\end{otherlanguage}\par\medskip}

% ============================================================================
% Code listings.  ★★ v3.4 THE BOX-OVERFLOW FIX:
%   resetmargins=true   → the listing IGNORES the enclosing environment's
%                         indentation and measures from the PAGE. v3.4 stated
%                         this backwards ("measured from the ENCLOSING box"),
%                         and v3.37 then blanked in-box line numbers because
%                         the page-anchored listing put its hanging digits on
%                         top of the box rule. Measured v4.5.0: with the key
%                         FALSE the listing follows the box and the code moves
%                         from x=75.26 to x=97.74, clearing the rule at 72.07.
%                         The key now lives in one \lstset below instead of in
%                         each style, so a box can override it — codebox does.
%   framexleftmargin=0  → the frame (left bar) does not extend left of the code.
%   xleftmargin keeps a small indent so the line numbers have room INSIDE.
% asmblock wraps the listing in an LTR island (code is LTR).
% ============================================================================
\lstdefinestyle{asm}{
  basicstyle=\ttfamily\small,
  backgroundcolor=\color{cBlueBg!55},
  frame=leftline,framerule=2.5pt,rulecolor=\color{cBlue},
  numbers=left,numberstyle=\tiny\color{cGray},numbersep=8pt,
  xleftmargin=2.2em,framexleftmargin=0pt,framexrightmargin=0pt,
  breaklines=true,columns=fullflexible,keepspaces=true,
  showstringspaces=false,aboveskip=8pt,belowskip=8pt,
  % v3.36: without a comment DELIMITER the commentstyle below never fires —
  % listings has no `language` here, so it does not know '#' opens a comment,
  % and every asm listing rendered flat black (measured: a whole-page pixel
  % sweep found ZERO cGreen glyphs across 135 shipped listings). morecomment
  % gives it the delimiter; color-only change, glyph widths untouched.
  morecomment={[l]{\#}},
  commentstyle=\color{cGreen},
}
\lstdefinestyle{cstyle}{
  language=C,basicstyle=\ttfamily\small,
  backgroundcolor=\color{cGreenBg!55},
  frame=leftline,framerule=2.5pt,rulecolor=\color{cGreen},
  numbers=left,numberstyle=\tiny\color{cGray},numbersep=8pt,
  xleftmargin=2.2em,framexleftmargin=0pt,framexrightmargin=0pt,
  breaklines=true,columns=fullflexible,keepspaces=true,
  showstringspaces=false,keywordstyle=\color{cBlue}\bfseries,
  commentstyle=\color{cGray}\itshape,aboveskip=8pt,belowskip=8pt,
}
% v3.36: pystyle/htmlstyle — colored twins of cstyle for the Python/HTML
% listings that used to repeat a ~200-char options blob with NO color
% (keywords rendered bold-black, comments italic-black — measured; C alone
% was colored). Geometry is copied VERBATIM from the shipped blob (frame,
% margins, skips) so swapping a call site to style=pystyle moves NOTHING —
% only keywordstyle/commentstyle/stringstyle are added (bold stays bold,
% italic stays italic, mono metrics identical => zero reflow, proven by
% per-page word-bbox diff). breaklines/breakatwhitespace stay per-site.
% v3.78: numbers=left joins both styles, with the exact numberstyle/numbersep
% keys asm/cstyle carry — fleet-wide consistency (author decision: all four
% styles numbered; "asm/cstyle numbered, py/html not" was an accident of the
% v3.36 geometry freeze, not a design). The numbers hang LEFT of the code box
% (listings sets them with \llap, so no overfull hbox can arise and the code
% x is untouched — measured: code frozen, digits at code_x-11.8pt); inside a
% colored box they are blanked like everyone else's (v3.37 hook), and the
% clean variant restores them in its frameless boxes (v3.78, above).
\lstdefinestyle{pystyle}{
  language=Python,basicstyle=\ttfamily\small,
  columns=fullflexible,keepspaces=true,
  frame=single,rulecolor=\color{cTblLine},
  numbers=left,numberstyle=\tiny\color{cGray},numbersep=8pt,
  xleftmargin=0.5em,framexleftmargin=0pt,framexrightmargin=0pt,
  showstringspaces=false,aboveskip=8pt,belowskip=8pt,
  keywordstyle=\color{cBlue}\bfseries,
  commentstyle=\color{cGray}\itshape,
  stringstyle=\color{cGreen},
}
% v4.5.0: hoisted out of the four styles so an enclosing box can override it.
% Every listing still resolves to true unless something later says otherwise,
% so shipped output is unchanged (verified byte-for-byte on all eight PDFs).
\lstset{resetmargins=true}
\lstdefinestyle{htmlstyle}{
  language=HTML,basicstyle=\ttfamily\small,
  columns=fullflexible,keepspaces=true,
  frame=single,rulecolor=\color{cTblLine},
  numbers=left,numberstyle=\tiny\color{cGray},numbersep=8pt,
  xleftmargin=0.5em,framexleftmargin=0pt,framexrightmargin=0pt,
  showstringspaces=false,aboveskip=8pt,belowskip=8pt,
  keywordstyle=\color{cBlue}\bfseries,
  commentstyle=\color{cGray}\itshape,
  stringstyle=\color{cGreen},
}
\newenvironment{asmblock}{\par\begin{otherlanguage}{english}}{\end{otherlanguage}\par}

% <hllpdf:listing-smart-quotes>
% v4.21.0: four Unicode punctuation characters must enter listings' token
% buffer. Unregistered characters bypass it: a closing quote was printed
% BEFORE the pending word, and an opening quote BEFORE a pending brace.
% Extend the standard character table, then activate only these codepoints
% in the listing's local character-table scope. Do this BEFORE the package's
% delimiter/alsoother/literate hooks so explicit user mappings still win.
% This is not universal Unicode support, nor a conversion to ASCII quotes.
% See listings developer guide 15.3.1 and references/listing-smart-quotes.md.
\makeatletter
\lst@CCPut \lst@ProcessOther {"2018}{"2019}{"201C}{"201D}\z@
\lst@AddToHookAtTop{SelectCharTable}{%
  \lst@ifactivechars
    \catcode"2018=\active \catcode"2019=\active
    \catcode"201C=\active \catcode"201D=\active
  \fi}
\makeatother
% </hllpdf:listing-smart-quotes>

% <hllpdf:listing-technical-symbols>
% v4.30: bounded extension for common technical punctuation in LTR listings.
% Without character-table registration these glyphs bypass the token buffer
% and can move before the preceding word even inside quoted strings.
% Keep the original four-quote block separate for its standing mutation test.
% This is source-faithful printing, not syntax validation or general Unicode.
\makeatletter
\lst@CCPut \lst@ProcessOther
  {"2013}{"2014}{"2026}{"2190}{"2192}{"2194}{"21D2}
  {"2212}{"221E}{"2260}{"2264}{"2265}\z@
\lst@AddToHookAtTop{SelectCharTable}{%
  \lst@ifactivechars
    \catcode"2013=\active \catcode"2014=\active \catcode"2026=\active
    \catcode"2190=\active \catcode"2192=\active \catcode"2194=\active
    \catcode"21D2=\active \catcode"2212=\active \catcode"221E=\active
    \catcode"2260=\active \catcode"2264=\active \catcode"2265=\active
  \fi}
\makeatother
% </hllpdf:listing-technical-symbols>


\usepackage[hidelinks,unicode]{hyperref}   % LAST, after colors. unicode = Hebrew bookmarks.
% v3.62: THE ISOLATION MACROS CORRUPTED EVERY PDF BOOKMARK THEY APPEARED IN.
% \en is defined (line ~155) as
%     \foreignlanguage{english}{\language=\csname l@nohyphenation\endcsname #1}
% and in a heading that body is written verbatim to the .toc and handed to
% hyperref's \pdfstringdef. hyperref strips the two CONTROL SEQUENCES (\language,
% \l@nohyphenation) but the plain `=` CHARACTER token between them is not a control
% sequence, so it survives into the PDF string: \chapter{שער \en{CNOT} והצמדה}
% produced the bookmark "שער =CNOT והצמדה". The printed TOC and the heading on the
% page are CORRECT — only the navigation pane and the viewer window title are wrong,
% which is why it went unnoticed. Measured on the shipped fleet: 70 corrupted
% bookmarks across 5 PDFs, from 177 unprotected \en-in-heading sites. Same corruption
% reached /Title and /Subject via the \hypersetup{pdftitle=...} the docs recommend.
% This hook tells hyperref to use the plain argument for all three isolation macros
% when building a PDF string. It does NOT touch the .toc (the printed TOC is
% byte-unchanged) and it does NOT mask the math signal — a heading with $...$ still
% logs "removing math shift", so BIDI-15 detection is intact.
% BONUS, measured: it also makes \code{} in an UNSTARRED heading COMPILE. BLD-10
% recorded that as needing \texorpdfstring or a starred heading; the real root cause
% is hyperref's \pdfstringdef scan of \code's \ifmmode...\else...\fi body — which is
% exactly why the identically-bodied \texttt survived (hyperref pre-neutralizes
% \texttt in PDF strings) and why \DeclareRobustCommand could never help.
% v3.64: \def\\{ } added. A manual line break in a heading WELDS the two words
% together in the bookmark: \chapter{שורה ראשונה\\שורה שנייה} produced the outline
% entry 'שורה ראשונהשורה שנייה' — "ראשונהשורה" is not a word. hyperref drops the
% \\ and puts nothing in its place, so the break that looks like a space on the
% page is a zero-width join in the navigation pane. The spaced form
% 'מילה אחת \\ מילה שנייה' escapes with a harmless DOUBLE space ('אחת  מילה'),
% which is why this never got noticed. Both are measured on a real outline.
% 0 live instances in the corpus — but SKILL.md and content-style.md both RECOMMEND
% a manual heading break, so this is a primed trap rather than a dormant one, and a
% one-token preamble fix retires it for good instead of asking authors to remember.
% Mapping \\ to a space fixes the WELD, which is the actual defect (a nonexistent
% word in the navigation pane). It does not normalise whitespace: measured after the
% fix, the welded form becomes 'שורה ראשונה שורה שנייה' (correct) and the already-
% spaced form goes from two spaces to three. That is cosmetic — viewers render it as
% a wider gap, and no word is corrupted. A per-heading \texorpdfstring is still the
% way to get an exactly-tuned bookmark string.
\pdfstringdefDisableCommands{\def\en#1{#1}\def\h#1{#1}\def\doclang#1{#1}\def\code#1{#1}\def\codeline#1{#1}\def\\{ }}
% v3.70: RESTORE \url LINE-BREAKING. url.sty freezes its break penalties at LOAD
% time (\mathchardef\UrlBreakPenalty=\binoppenalty), and v3.45's global
% \binoppenalty=10000 (inline-math atomicity) therefore priced EVERY breakpoint
% inside \Url at 10000 = never break. Measured: a long \url in Hebrew prose is
% one unbroken line, ~100pt overfull, invading the margin band the ink gates
% deliberately tolerate. With the two lines below: 0 overfull, the URL breaks
% after '/', the RTL split is CORRECT (fragment order forward), and each line
% gets its own link annotation exactly covering its fragment (measured).
% Confined to \Url's group — inline-math atomicity (TYP-7) is untouched.
\mathchardef\UrlBreakPenalty=200
\mathchardef\UrlBigBreakPenalty=100
\ifx\lakitmainlanguage\hllpdfenglish
  \hypersetup{pdflang=en}
\else
  \hypersetup{pdflang=he}
\fi  % /Lang is the document base locale (not per-span accessibility tagging).
% Set a per-document \hypersetup{pdftitle={...}} too — see SKILL.md.

% v3.52 (round-2 audit): two RTL bibliography/cross-ref defects, both invisible to
% every gate and both affecting the FIRST author to use the feature.
% (a) \@biblabel's default [#1] brackets are bidi-neutral and mirror to ]1[ under
% the RTL bibliography base — every \thebibliography reference list shows ]1[.
% LTR-isolating the label restores [1]. (The render gate's mirror rule fingerprints
% round parens )1(, so square brackets slipped through.)
\makeatletter
\renewcommand*{\@biblabel}[1]{\foreignlanguage{english}{[#1]}}
% v3.54: the SAME neutral-bracket trap on the INLINE side — a multi-key \cite{a,b}
% under the RTL base lays its European-number list right-to-left, so [1, 2] renders
% "[2,1]" and [1, 2, 3] renders "[3,2,1]" (verified by glyph geometry — the reader
% sees the sources in REVERSE order; a note form \cite[p.5]{a} also reorders to
% "[p. 5,1]"). LTR-isolate the whole bracket group (number list AND optional note),
% exactly like \@biblabel above: [1, 2] / [1, 2, 3] / [1, p. 5] all render correct.
% Placed AFTER hyperref so it wins; hyperref hyperlinks each number inside #1 (via
% \@citex), which this wrapper leaves intact. Mirrors the LaTeX kernel \@cite shape.
% v4.9.1: directional isolation must NOT switch the language here. Babel's
% safe-active citation wrapper passes its key list through \bbl@tempa;
% \foreignlanguage reuses that scratch macro before the deferred citation
% loop expands it. The old line therefore cited the literal key "english".
% If such a key existed, even the undefined-reference gate stayed green.
% \babelsublr isolates the bracket group without selecting a locale or
% changing that key list. Preserve native links, key order and optional notes.
\def\@cite#1#2{\babelsublr{[{#1\if@tempswa , #2\fi}]}}
\makeatother

% v4.11.0: a reference inside align's \text inherited HEBdisptag=true,
% so its enclosing hbox mirrored the brackets despite a clean compile.
% Isolate the complete printed reference without selecting a Babel language.
% Also neutralize the display-only normalfont direction reset locally. A
% direction-only patch inverted 12 to 21 inside equation in the measured probe.
% Plain maketag is local to the reference: mathtools' display-side wrapper
% otherwise globally consumes manual_tag while formatting a nested reference.
% Preserve each original command (including the selective auxiliary write),
% current tag forms, native links, and showonlyrefs on/off transitions.
\makeatletter
\NewCommandCopy{\HEB@nativeeqref}{\eqref}
\newcommand{\HEB@eqrefisolate}[1]{%
  \babelsublr{\HEBdisptagfalse\let\bbl@eqnodir\relax\let\maketag@@@\ORIGmaketag@@@#1}}
\protected\def\eqref#1{\HEB@eqrefisolate{\HEB@nativeeqref{#1}}}
\ExplSyntaxOn
\NewCommandCopy{\HEB@nativeMTeqref}{\MT_extended_eqref:n}
\cs_set_protected:Npn \MT_extended_eqref:n #1
  {\HEB@eqrefisolate{\HEB@nativeMTeqref{#1}}}
\ExplSyntaxOff
\makeatother

% v4.11.0: custom caption counters are not part of hyperref's English
% locale, and Hebrew has no upstream hyperref locale in this toolchain.
% Restore the appropriate names at every Babel language switch, including
% return from a nested inline island. Do not localize the destination keys.
\makeatletter
\newcommand{\HEB@refnames}{%
\def\chapterautorefname{פרק}%
\def\sectionautorefname{סעיף}%
\def\subsectionautorefname{תת־סעיף}%
\def\subsubsectionautorefname{תת־תת־סעיף}%
\def\paragraphautorefname{פסקה}%
% v4.31: hyperref uses a distinct counter name at the deepest standard level.
\def\subparagraphautorefname{תת־פסקה}%
\def\equationautorefname{משוואה}%
\def\figureautorefname{איור}%
\def\tableautorefname{טבלה}%
\def\partautorefname{חלק}%
\def\appendixautorefname{נספח}%
\def\footnoteautorefname{הערת שוליים}%
\def\itemautorefname{סעיף}%
\def\Itemautorefname{סעיף}%
\def\figautorefname{איור}%
\def\tblautorefname{טבלה}%
\def\pageautorefname{עמוד}% 
}
\addto\extrashebrew{\HEB@refnames}
\newcommand{\ENG@refnames}{%
  \def\subparagraphautorefname{subparagraph}%
  \def\chapterautorefname{chapter}%
  \def\sectionautorefname{section}%
  \def\subsectionautorefname{subsection}%
  \def\subsubsectionautorefname{subsubsection}%
  \def\paragraphautorefname{paragraph}%
  \def\equationautorefname{Equation}%
  \def\figureautorefname{Figure}%
  \def\tableautorefname{Table}%
  \def\partautorefname{Part}%
  \def\appendixautorefname{Appendix}%
  \def\footnoteautorefname{footnote}%
  \def\itemautorefname{item}%
  \def\figautorefname{Figure}%
  \def\tblautorefname{Table}%
  \def\Itemautorefname{item}%
  \def\pageautorefname{page}%
}
\addto\extrasenglish{\ENG@refnames}
\ifx\lakitmainlanguage\hllpdfenglish\ENG@refnames\else\HEB@refnames\fi
\makeatother

% ============================================================================
% Optional running head (v3.39): \lakithead{<book title>} — OFF by default.
% A document opts in with ONE line after \input{preamble.tex}:
%   \lakithead{מבנה המחשב}
% Every content page then carries the series' orientation line: chapter name
% at the reading edge, book title opposite, thin rule. (fancyhdr's [L]/[R]
% are PHYSICAL; under the RTL page direction [L] lands at the RIGHT/reading
% edge — measured, consistent across the corpus.) Packaged like the topic
% rail: always loaded, inert without the call. This centralizes the 8-line
% block that 8 lakits used to copy by hand — the same duplication vector
% that once let a lakit preamble go ~100 lines stale. Class-aware: report
% fills the mark from \chapter, article from \section. NB: a titlepage env
% self-empties; front matter that sets \pagestyle{empty} must restore with
% \pagestyle{fancy} afterwards (NOT plain — that discards the head).
% ============================================================================
\usepackage{fancyhdr}
% v3.61: \chapter* left a STALE running head — pages of an UNNUMBERED chapter
% carried the PREVIOUS chapter's title. report.cls routes \chapter* through
% \@schapter, which never calls \chaptermark, so \leftmark keeps whatever the last
% numbered chapter left. Measured: with \chapter{פרק ממוספר} then
% \chapter*{מבוא לא ממוספר}, page 4 (inside the starred chapter) printed
% "פרק ממוספר". Invisible on the chapter's FIRST page (\thispagestyle{plain}), so it
% only shows from the second page on — the pages an author is least likely to check.
% An unnumbered front-matter chapter (\chapter*{מבוא}, \chapter*{הקדמה}) is everyday
% Hebrew-academic structure, and the corpus already carries HAND-WRITTEN \markboth
% lines right after \chapter* — the exact per-document duplication \lakithead exists
% to retire. Patched INSIDE \lakithead so the macro stays inert when unused; a
% document's own later \markboth still wins, so those hand-patches do not regress.
\makeatletter
% <hllpdf:starred-heading-scope>
\newif\iflakin@starmarkpatched
\newcommand{\lakistarmark}{%
  \iflakin@starmarkpatched\else
  % Dispatch on the actual class counter, not hyperref's shared internals.
  \ifdefined\c@chapter
    \let\HEBorigschapter\@schapter
    % Record only after the native heading has reached its page. A mark
    % before the native routine can advertise the next heading too early.
    \def\@schapter##1{\HEBorigschapter{##1}\chaptermark{##1}}%
  \else
    % @ssect receives layout arguments, NOT the section name or level.
    % Capture the name in @startsection without introducing a TeX group:
    % grouping sectioning would break the native after-heading state.
    \def\HEB@sectionkind{}%
    \def\HEB@mainsection{section}%
    \let\HEBorigstartsection\@startsection
    \def\@startsection##1##2##3##4##5##6{%
      \def\HEB@sectionkind{##1}%
      \HEBorigstartsection{##1}{##2}{##3}{##4}{##5}{##6}}%
    \let\HEBorigssect\@ssect
    \def\@ssect##1##2##3##4##5{%
      \HEBorigssect{##1}{##2}{##3}{##4}{##5}%
      % Lower starred headings must not replace the article's section mark.
      \ifx\HEB@sectionkind\HEB@mainsection
        \sectionmark{##5}%
      \fi
      \ignorespaces}%
  \fi
  \global\lakin@starmarkpatchedtrue
  \fi}
% </hllpdf:starred-heading-scope>
% <hllpdf:running-language>
% Choose the wrapper when the mark is written, not when a later page ships.
% English punctuation must remain LTR even after returning to Hebrew prose.
\newcommand{\HEB@runningmark}[1]{%
  \iflanguage{english}{\markboth{\en{#1}}{}}{\markboth{\h{#1}}{}}}
% </hllpdf:running-language>
\makeatother
% v4.2.1 (NAV): every collection already names itself in \lakithead — feed the
% PDF /Title from the same declaration, so the navigation gate's
% PDF_REQUIRE_TITLE has a universal source and no fleet member ships
% title-less by missing a second declaration site. Measured before this: 14
% of 16 collections shipped /Title-empty while every one carried a perfectly
% good \lakithead string (plain Hebrew in all 16 — no math, no formatting, so
% the PDF-string path is clean). An explicit title always wins: set it with
% \lakitpdftitle{...} anywhere before \lakithead, or with a plain
% \hypersetup{pdftitle=...} — the fallback detects that via hyperref's
% \@pdftitle, probe-verified \@empty-equal until the first explicit set.
\makeatletter
\newif\iflakit@pdftitleset
\newcommand{\lakitpdftitle}[1]{%
  \hypersetup{pdftitle={#1}}\global\lakit@pdftitlesettrue}
\newcommand{\lakithead}[1]{%
  \iflakit@pdftitleset\else
    \ifx\@pdftitle\@empty
      \hypersetup{pdftitle={#1}}\global\lakit@pdftitlesettrue
    \fi
  \fi
  \pagestyle{fancy}\fancyhf{}%
  \lakistarmark

  \fancyhead[L]{\small\textcolor{cGray}{\leftmark}}%
  % <hllpdf:fixed-head-language>
  % The fixed document title has its OWN language snapshot, independent of
  % the current section/chapter mark. Choose the branch NOW, not at shipout:
  % the page may be Hebrew by the time an English title is finally printed.
  \iflanguage{english}{%
    \fancyhead[R]{\small\textcolor{cGray}{\en{#1}}}%
  }{%
    \fancyhead[R]{\small\textcolor{cGray}{\h{#1}}}%
  }%
  % </hllpdf:fixed-head-language>
  \fancyfoot[C]{\thepage}%
  \renewcommand{\headrulewidth}{0.4pt}%
  \ifdefined\chaptermark
    \renewcommand{\chaptermark}[1]{\HEB@runningmark{##1}}%
  \else
    \renewcommand{\sectionmark}[1]{\HEB@runningmark{##1}}%
  \fi}
\makeatother

% v3.39.1: path-independent byte-reproducibility. LuaTeX derives the trailer
% /ID from the OUTPUT PATH, so two clones of the same source at different
% directories produced different sha256 (measured: the ONLY delta was the
% 62-byte /ID line). Bit 512 suppresses exactly that line — measured: Creator/
% Producer/CreationDate/ModDate and all content untouched. /ID is optional
% per ISO 32000 (required only for ENCRYPTED PDFs and PDF/A — neither is a
% target here). If a future workflow needs an ID (PDF/A, Acrobat revision
% matching), delete this line and it returns.
\pdfvariable suppressoptionalinfo \numexpr 512\relax

% v3.75: BOLD MATH IN BOLD HEADINGS. Heading text is \bfseries but math inside
% it stayed regular-weight CM — a hairline $V_T$ inside a heavy 24.8pt chapter
% title (measured: CMMI12 fl=6 next to FrankRuhlLibre-Bold fl=20; 155 live
% heading-with-math sites fleet-wide). \boldmath rides along wherever the class
% sets \bfseries for a heading; A/B measured: $V_T$ -> CMMIB10 fl=22, 0 missing
% characters (the bold CM faces ship with this TeX Live). Guarded per-macro so
% the article class (no \@makechapterhead) and future class changes stay safe.
% v4.0.1 (adoption fix): the SEARCH patterns above carry no parameter token.
% v4.0.0 made them specific by including the class's own #1/#2, and measured on
% TeX Live 2023 that never matches: a \patchcmd search pattern containing a
% parameter token cannot match while the call sits inside \@ifundefined{...}{}{...}
% (probed both #1 and ##1 -> FAIL; the same pattern at top level succeeds). Every
% one of these patches therefore failed, and because v4.0.0 also replaced the
% empty failure branch with \PackageError, the failure became FATAL: all four
% bundled example assets stopped building ("Cannot patch chapter-title math
% weight -> Fatal error, no output PDF"). The \PackageError branches are KEPT —
% a silent no-op here is what let v3.75's fix rot unnoticed — only the parameter
% tokens are dropped, so the patches match and a real class change still errors.
\makeatletter
\@ifundefined{@makechapterhead}{}{%
  \patchcmd{\@makechapterhead}{\huge\bfseries \@chapapp}
    {\huge\bfseries\boldmath \@chapapp}{}
    {\PackageError{hebrew-lualatex-pdf}{Cannot patch chapter-number math weight}{Unsupported document class}}%
  \patchcmd{\@makechapterhead}{\Huge \bfseries}
    {\Huge \bfseries\boldmath}{}
    {\PackageError{hebrew-lualatex-pdf}{Cannot patch chapter-title math weight}{Unsupported document class}}%
  \patchcmd{\@makeschapterhead}{\Huge \bfseries}
    {\Huge \bfseries\boldmath}{}
    {\PackageError{hebrew-lualatex-pdf}{Cannot patch starred chapter math weight}{Unsupported document class}}}
% v4.0.1: \part sits under a THREE-DEEP wrapper chain, so the class-based anchor
% is not in \@part at all. Measured with \show: hyperref rewrites \@part to
% `...\phantomsection\fi \H@old@part`, and nameref rewrites THAT to
% `\NR@gettitle{#1}\NR@part[{#1}]{#2}` — the class body ends up in \NR@part.
% v4.0.0 patched \@part and made failure fatal, which is why all seven bundled
% examples stopped building. \part is NOT a dead heading level here: 188 live uses
% across the fleet, so the bold-math fix is worth getting right rather than
% dropping. We walk the chain innermost-first and stop at the first hit. Failure is
% a WARNING, not an error: the chain is a property of hyperref/nameref versions, so
% a future package release must degrade to regular-weight part math — a cosmetic
% loss — instead of refusing to build. Chapter and section keep \PackageError,
% because those targets are in the CLASS and were verified to patch here.
\newif\iflakin@partnum \newif\iflakin@parttitle \newif\iflakin@sparttitle
\def\lakin@partnum#1#2{\iflakin@partnum\else
  \patchcmd{#1}{#2\bfseries \partname}{#2\bfseries\boldmath \partname}
    {\lakin@partnumtrue}{}\fi}
\def\lakin@parttitle#1#2{\iflakin@parttitle\else
  \patchcmd{#1}{#2 \bfseries}{#2 \bfseries\boldmath}{\lakin@parttitletrue}{}\fi}
\def\lakin@sparttitle#1#2{\iflakin@sparttitle\else
  \patchcmd{#1}{#2 \bfseries}{#2 \bfseries\boldmath}{\lakin@sparttitletrue}{}\fi}
\@ifundefined{c@chapter}{\def\lakin@pnum{\Large}\def\lakin@ptit{\huge}}%
                        {\def\lakin@pnum{\huge}\def\lakin@ptit{\Huge}}
\@ifundefined{NR@part}{}{\expandafter\lakin@partnum\expandafter\NR@part\expandafter{\lakin@pnum}%
                         \expandafter\lakin@parttitle\expandafter\NR@part\expandafter{\lakin@ptit}}
\@ifundefined{H@old@part}{}{\expandafter\lakin@partnum\expandafter\H@old@part\expandafter{\lakin@pnum}%
                            \expandafter\lakin@parttitle\expandafter\H@old@part\expandafter{\lakin@ptit}}
\@ifundefined{@part}{}{\expandafter\lakin@partnum\expandafter\@part\expandafter{\lakin@pnum}%
                       \expandafter\lakin@parttitle\expandafter\@part\expandafter{\lakin@ptit}}
\@ifundefined{NR@spart}{}{\expandafter\lakin@sparttitle\expandafter\NR@spart\expandafter{\lakin@ptit}}
\@ifundefined{H@old@spart}{}{\expandafter\lakin@sparttitle\expandafter\H@old@spart\expandafter{\lakin@ptit}}
\@ifundefined{@spart}{}{\expandafter\lakin@sparttitle\expandafter\@spart\expandafter{\lakin@ptit}}
\iflakin@partnum\else\PackageWarning{hebrew-lualatex-pdf}{Part-number math stays
  regular weight: no known \noexpand\part wrapper matched (hyperref/nameref
  version change). Cosmetic only}\fi
\iflakin@parttitle\else\PackageWarning{hebrew-lualatex-pdf}{Part-title math stays
  regular weight: no known \noexpand\part wrapper matched. Cosmetic only}\fi
% v4.12.0: THE PART NUMBER IN THE TOC WAS NOT DIRECTIONALLY SEPARATED FROM THE
% TITLE. report.cls/article.cls compose the part entry as ONE text stream --
% \addcontentsline{toc}{part}{\thepart\hspace{1em}#1} -- while EVERY other level
% goes through \numberline, which puts the number in an \hbox. A box is a bidi
% OBJECT and cannot join a neighbouring run; a bare "IV" can. So in the RTL TOC
% paragraph a Latin number and a Latin-INITIAL title merge into a SINGLE LTR run
% and are laid out left-to-right as a unit: the number lands on the wrong side of
% the row, and a mixed title is scrambled outright. Measured on \part{MCP ...}:
% the printed order was "<tail> V MCP" -- title fragment, number, title fragment.
% The chapter/section rows one line below, with the SAME Latin-initial titles,
% were correct on the same page. That contrast is the whole diagnosis.
% \en{} IS NOT A WORKAROUND, measured both ways on one fixture: it fixes a
% PURE-Latin title (\part{\en{Hooks}} renders correctly) and does NOT fix a MIXED
% one (\part{\en{MCP} <hebrew>} stays scrambled). babel's language switch changes
% font and hyphenation, not the bidi run structure -- so the defect is not
% reachable from the source side at all, which is why it belongs here.
% Boxing the number is the entire fix and is a STRICT NO-OP where the bug never
% bit: a Hebrew-titled document compiles BYTE-IDENTICAL before and after (md5,
% 3 clean runs), because there the neutral glue already sat between an L run and
% an R run. The PDF outline is byte-identical too -- hyperref's \pdfstringdef
% neutralises \mbox, so bookmark strings keep their part numbers.
% \numberline{\thepart} would also fix the direction, but it re-geometries every
% part row in the fleet and DROPS the number from the part bookmark (hyperref
% strips \numberline). \mbox changes nothing that was already right.
% The search pattern deliberately carries no # token: etoolbox cannot match a
% parameter token in a patch pattern (measured -- all three chain members
% reported "not patchable" while the pattern included #1). "\thepart \hspace
% {1em}" is brace-balanced and occurs exactly once in the wrapper body.
\newif\iflakin@parttocnum
\def\lakin@parttocnum#1{\iflakin@parttocnum\else
  \patchcmd{#1}{\thepart \hspace {1em}}{\mbox {\thepart }\hspace {1em}}
    {\lakin@parttocnumtrue}{}\fi}
\@ifundefined{NR@part}{}{\lakin@parttocnum\NR@part}
\@ifundefined{H@old@part}{}{\lakin@parttocnum\H@old@part}
\@ifundefined{@part}{}{\lakin@parttocnum\@part}
\iflakin@parttocnum\else\PackageWarning{hebrew-lualatex-pdf}{Part TOC number is
  not boxed: no known \noexpand\part wrapper matched. A part title that BEGINS
  with Latin will print its number on the wrong side of its TOC row}\fi
\patchcmd{\section}{\bfseries}{\bfseries\boldmath}{}
  {\PackageError{hebrew-lualatex-pdf}{Cannot patch section math weight}{Unsupported document class}}
\patchcmd{\subsection}{\bfseries}{\bfseries\boldmath}{}
  {\PackageError{hebrew-lualatex-pdf}{Cannot patch subsection math weight}{Unsupported document class}}
\patchcmd{\subsubsection}{\bfseries}{\bfseries\boldmath}{}
  {\PackageError{hebrew-lualatex-pdf}{Cannot patch subsubsection math weight}{Unsupported document class}}
\patchcmd{\paragraph}{\bfseries}{\bfseries\boldmath}{}{}
\patchcmd{\subparagraph}{\bfseries}{\bfseries\boldmath}{}{}
\@ifundefined{l@chapter}{}{%
  \patchcmd{\l@chapter}{\bfseries}{\bfseries\boldmath}{}{}}
\makeatother
\setcounter{tocdepth}{\lakittocdepth}
% <hllpdf:measured-toc>
% v4.16.0: Measure actual number labels, not a guessed digit count. Read only
% canonical contentsline records; never execute titles or addtocontents code
% during the measuring pass. Widths are local to this individual list and are
% available before its first row: no per-row drift and no extra aux pass.
\makeatletter
\newread\hllpdf@tocstream
\newdimen\hllpdf@toc@chapter
\newdimen\hllpdf@toc@section
\newdimen\hllpdf@toc@subsection
\newdimen\hllpdf@toc@subsubsection
\newdimen\hllpdf@toc@paragraph
\newdimen\hllpdf@toc@subparagraph
\newdimen\hllpdf@toc@figure
\newdimen\hllpdf@toc@table
\newdimen\hllpdf@tocpage
\newdimen\hllpdf@toc@sectionindent
\newdimen\hllpdf@toc@subsectionindent
\newdimen\hllpdf@toc@subsubsectionindent
\newdimen\hllpdf@toc@paragraphindent
\newdimen\hllpdf@toc@subparagraphindent
\def\hllpdf@tocdefaults{%
  \hllpdf@toc@chapter=2em
  \hllpdf@toc@section=3.2em
  \hllpdf@toc@subsection=3.9em
  \@ifundefined{c@chapter}{%
    \hllpdf@toc@subsubsection=3.2em
    \hllpdf@toc@paragraph=4.1em
    \hllpdf@toc@subparagraph=5em
  }{%
    \hllpdf@toc@subsubsection=4.1em
    \hllpdf@toc@paragraph=5em
    \hllpdf@toc@subparagraph=6em
  }%
  \hllpdf@toc@figure=2.3em
  \hllpdf@toc@table=2.3em
  \hllpdf@tocpage=\@pnumwidth\relax
}
\def\hllpdf@tocindents{%
  % Preserve the established shallow layout when its slots already fit.
  \hllpdf@toc@sectionindent=1.5em
  \@tempdima=\hllpdf@toc@chapter \advance\@tempdima by -2em
  \advance\hllpdf@toc@sectionindent by \@tempdima
  \hllpdf@toc@subsectionindent=\dimexpr
    \hllpdf@toc@sectionindent+\hllpdf@toc@section\relax
  \hllpdf@toc@subsubsectionindent=\dimexpr
    \hllpdf@toc@subsectionindent+\hllpdf@toc@subsection\relax
  \hllpdf@toc@paragraphindent=\dimexpr
    \hllpdf@toc@subsubsectionindent+\hllpdf@toc@subsubsection\relax
  \hllpdf@toc@subparagraphindent=\dimexpr
    \hllpdf@toc@paragraphindent+\hllpdf@toc@paragraph\relax
}
\def\hllpdf@tocfont@chapter{\bfseries\boldmath}
\def\hllpdf@tocfont@part{\large\bfseries\boldmath}
\def\hllpdf@tocentryfont{%
  \ifcsname hllpdf@tocfont@\hllpdf@toclevel\endcsname
    \csname hllpdf@tocfont@\hllpdf@toclevel\endcsname
  \fi}
\long\def\hllpdf@tocdiscard#1\hllpdf@tocstop{}
\def\hllpdf@tocscanrecord{%
  \@ifnextchar\contentsline{\hllpdf@tocscanline}{%
    \@ifnextchar\babel@toc{\hllpdf@tocscanlanguage}{\hllpdf@tocdiscard}}}
\long\def\hllpdf@tocscanlanguage\babel@toc#1#2#3\hllpdf@tocstop{%
  % Babel records are declarations, not arbitrary code. Track the language
  % without executing the original directive or changing the surrounding list.
  \def\hllpdf@toclanguage{#1}}
\long\def\hllpdf@tocscanline\contentsline#1#2#3#4#5\hllpdf@tocstop{%
  \def\hllpdf@toclevel{#1}%
  \ifcsname hllpdf@toc@#1\endcsname
    \hllpdf@tocscantitle#2\hllpdf@tocstop
  \fi
  % The page-number column needs its own measurement (including bold chapter
  % rows). Its extra clearance separates the page label from dot leaders.
  \setbox\@tempboxa=\hbox{\foreignlanguage{\hllpdf@toclanguage}{\hllpdf@tocentryfont #3\kern.3em}}%
  \ifdim\wd\@tempboxa>\hllpdf@tocpage
    \hllpdf@tocpage=\wd\@tempboxa
  \fi}
\def\hllpdf@tocscantitle{%
  \@ifnextchar\numberline{\hllpdf@tocscanlabel}{\hllpdf@tocdiscard}}
\long\def\hllpdf@tocscanlabel\numberline#1#2\hllpdf@tocstop{%
  \setbox\@tempboxa=\hbox{\foreignlanguage{\hllpdf@toclanguage}{\hllpdf@tocentryfont #1\kern.5em}}%
  \ifdim\wd\@tempboxa>\csname hllpdf@toc@\hllpdf@toclevel\endcsname
    \csname hllpdf@toc@\hllpdf@toclevel\endcsname=\wd\@tempboxa
  \fi}
\def\hllpdf@tocread{%
  \ifeof\hllpdf@tocstream\else
    \read\hllpdf@tocstream to\hllpdf@tocrecord
    \expandafter\hllpdf@tocscanrecord\hllpdf@tocrecord\hllpdf@tocstop
    \expandafter\hllpdf@tocread
  \fi}
\def\hllpdf@tocmeasure#1{%
  \makeatletter
  \edef\hllpdf@toclanguage{\languagename}%
  \hllpdf@tocdefaults
  \openin\hllpdf@tocstream=\jobname.#1\relax
  \hllpdf@tocread
  \closein\hllpdf@tocstream
  \hllpdf@tocindents
  \edef\@pnumwidth{\the\hllpdf@tocpage}%
  \@tempdima=\dimexpr\hllpdf@tocpage+1em\relax
  \ifdim\@tempdima>\@tocrmarg\relax
    \edef\@tocrmarg{\the\@tempdima}%
  \fi}
\let\hllpdf@originalstarttoc\@starttoc
\def\@starttoc#1{%
  \begingroup
    \def\hllpdf@tocext{#1}%
    \def\hllpdf@tocwanted{toc}%
    \ifx\hllpdf@tocext\hllpdf@tocwanted \hllpdf@tocmeasure{#1}\else
      \def\hllpdf@tocwanted{lof}%
      \ifx\hllpdf@tocext\hllpdf@tocwanted \hllpdf@tocmeasure{#1}\else
        \def\hllpdf@tocwanted{lot}%
        \ifx\hllpdf@tocext\hllpdf@tocwanted \hllpdf@tocmeasure{#1}\fi
      \fi
    \fi
    \hllpdf@originalstarttoc{#1}%
  \endgroup}
\hllpdf@tocdefaults
\hllpdf@tocindents
\@ifundefined{l@chapter}{}{%
  \patchcmd{\l@chapter}{\setlength\@tempdima{1.5em}}%
    {\setlength\@tempdima{\hllpdf@toc@chapter}}{}%
    {\PackageError{hebrew-lualatex-pdf}{Cannot size chapter TOC numbers}%
       {Unsupported contents layout: the chapter number slot was not found}}}
\renewcommand*\l@section{\@dottedtocline{1}%
  {\hllpdf@toc@sectionindent}{\hllpdf@toc@section}}
\renewcommand*\l@subsection{\@dottedtocline{2}%
  {\hllpdf@toc@subsectionindent}{\hllpdf@toc@subsection}}
\renewcommand*\l@subsubsection{\@dottedtocline{3}%
  {\hllpdf@toc@subsubsectionindent}{\hllpdf@toc@subsubsection}}
\renewcommand*\l@paragraph{\@dottedtocline{4}%
  {\hllpdf@toc@paragraphindent}{\hllpdf@toc@paragraph}}
\renewcommand*\l@subparagraph{\@dottedtocline{5}%
  {\hllpdf@toc@subparagraphindent}{\hllpdf@toc@subparagraph}}
\renewcommand*\l@figure{\@dottedtocline{1}{1.5em}{\hllpdf@toc@figure}}
\renewcommand*\l@table{\@dottedtocline{1}{1.5em}{\hllpdf@toc@table}}
\makeatother
% </hllpdf:measured-toc>
\setcounter{secnumdepth}{\lakitsecnumdepth}
% Public preamble-only caption-numbering policy. The historical within-major
% default is untouched unless this command is called. "continuous" removes the
% reset edge as well as the printed prefix, so Figure 1/2 and Table 1/2 stay
% unique across sections/chapters. Returning to "within-major" restores the
% original class-aware displays and native destination conventions. The policy
% never replaces \figcap/\tabcap or their measured link-placement mechanism.
\makeatletter
\let\lakin@majorfig\thefig
\let\lakin@majortbl\thetbl
\let\lakin@majorHfig\theHfig
\let\lakin@majorHtbl\theHtbl
\newcommand{\lakitcaptionnumbering}[1]{%
  \edef\lakin@numbering{#1}%
  \def\lakin@continuous{continuous}\def\lakin@withinmajor{within-major}%
  \ifx\lakin@numbering\lakin@continuous
    \@ifundefined{c@chapter}{%
      \@removefromreset{fig}{section}\@removefromreset{tbl}{section}%
    }{\@removefromreset{fig}{chapter}\@removefromreset{tbl}{chapter}}%
    \renewcommand{\thefig}{\arabic{fig}}\renewcommand{\thetbl}{\arabic{tbl}}%
    \def\theHfig{\arabic{fig}}\def\theHtbl{\arabic{tbl}}%
  \else\ifx\lakin@numbering\lakin@withinmajor
    \@ifundefined{c@chapter}{%
      \@removefromreset{fig}{section}\@removefromreset{tbl}{section}%
      \@addtoreset{fig}{section}\@addtoreset{tbl}{section}%
    }{%
      \@removefromreset{fig}{chapter}\@removefromreset{tbl}{chapter}%
      \@addtoreset{fig}{chapter}\@addtoreset{tbl}{chapter}}%
    \let\thefig\lakin@majorfig\let\thetbl\lakin@majortbl
    \let\theHfig\lakin@majorHfig\let\theHtbl\lakin@majorHtbl
  \else
    \PackageError{hebrew-english-latex-studio}{Unknown caption-numbering policy '#1'}%
      {Use continuous or within-major before begin{document}.}%
  \fi\fi}
\@onlypreamble\lakitcaptionnumbering
\makeatother
% Architecture readiness sentinel. Optional runtime layers require the complete
% core, so the marker is defined only after every mechanism/baseline definition.
% <hllpdf:architecture-ready>
\def\hllpdfcoreloaded{1}
% </hllpdf:architecture-ready>

% v4.22: registered file resources, not an assertion about embedded subset bytes.
% The log framing avoids long-path wrapping and preserves UTF-8 as hex.
% <hllpdf:font-inventory>
\directlua{
  local function inventory()
    local names = {}
    local unnamed = 0
    for id, item in font.each() do
      local path = item.filename
      if path and not (path == "") then names[path] = true else unnamed = unnamed + 1 end
    end
    local sorted = {}
    for path in pairs(names) do table.insert(sorted,path) end
    table.sort(sorted)
    texio.write_nl("log","HLLFONT|BEGIN|1")
    for index,path in ipairs(sorted) do
      local hex = path:gsub(".",function(c) return string.format(string.char(37) .. "02x",string.byte(c)) end)
      local count = math.ceil(string.len(hex)/48)
      texio.write_nl("log","HLLFONT|FILE|" .. index .. "|" .. count)
      for offset=1,string.len(hex),48 do texio.write_nl("log","HLLFONT|HEX|" .. hex:sub(offset,offset+47)) end
    end
    texio.write_nl("log","HLLFONT|END|" .. rawlen(sorted) .. "|" .. unnamed)
    texio.write_nl("log","")
  end
  luatexbase.add_to_callback("finish_pdffile",inventory,"hllpdf.runtime.fonts")
}
% </hllpdf:font-inventory>
; if omitted, the historical canonical
% defaults below are exact.  They deliberately control only the document shell
% and page layout.  RTL/math semantics stay in this preamble, while visual
% presentation themes (for example clean_theme.tex) load AFTER it.
%
%   \def\lakitdocumentclass{article}
%   \def\lakitclassoptions{11pt,a4paper}
%   \def\lakitgeometryoptions{margin=1.8cm,headheight=14pt}
%   \def\lakitlinespread{1.15}
%   \def\lakitparagraphstyle{block} % letters/briefs; indented for continuous prose
%   \def\lakittocdepth{1}
%   \def\lakitsecnumdepth{2}
%   % ============================================================================
% preamble.tex — tested LuaLaTeX preamble for Hebrew + math + figures PDFs (v4.32.0)
% Battle-tested: compiled an 80-page Hebrew quantum-computing summary AND a
% 197-page Hebrew computer-architecture anthology (TikZ diagrams, code listings,
% colored boxes) — both with 0 missing characters.
%
% Usage in your document:
%   \input{preamble.tex}     % (or paste these lines at the top)
%   \begin{document}  ...  \end{document}
%
% Compile with:  lualatex -interaction=nonstopmode doc.tex
%   Repeat until two consecutive PDFs are byte-identical (scripts/build.sh
%   does this) — a TOC document converges only on pass 3+; after two passes
%   the TOC page numbers are STALE (measured; see SKILL.md step 3).
% Font binaries are external dependencies, pinned by assets/fonts/SHA256SUMS.
% scripts/setup_fonts.sh imports a verified local installation or prior skill.
%
% PUBLIC DOCUMENT-SHELL / LAYOUT CONTRACT (introduced v4.3.1; architecture-owned v4.7.1). These hooks may be
% defined BEFORE \input{preamble.tex}; if omitted, the historical canonical
% defaults below are exact.  They deliberately control only the document shell
% and page layout.  RTL/math semantics stay in this preamble, while visual
% presentation themes (for example clean_theme.tex) load AFTER it.
%
%   \def\lakitdocumentclass{article}
%   \def\lakitclassoptions{11pt,a4paper}
%   \def\lakitgeometryoptions{margin=1.8cm,headheight=14pt}
%   \def\lakitlinespread{1.15}
%   \def\lakitparagraphstyle{block} % letters/briefs; indented for continuous prose
%   \def\lakittocdepth{1}
%   \def\lakitsecnumdepth{2}
%   \input{preamble.tex}
%
% Supported class contract: report or article.  Alternate paper sizes must also
% be declared to build.sh through EXPECTED_PAPER; the PDF gate never infers them.
% Do NOT use these hooks as a theme mechanism: theme overlays must not own page
% geometry, class selection, or direction semantics. A declared presentation-policy
% overlay may intentionally change numbering visibility after the preamble.
% ============================================================================
% Runtime architecture guard: generated from scripts/architecture_contract.py.
% The readiness sentinel remains delayed until the END of this file.
% <hllpdf:architecture-guard>
\ifdefined\hllpdfcoreloaded
  \errmessage{hllpdf: preamble.tex loaded more than once}
  \expandafter\endinput
\fi
% </hllpdf:architecture-guard>

\ifdefined\lakitdocumentclass\else\def\lakitdocumentclass{report}\fi
\ifdefined\lakitclassoptions\else\def\lakitclassoptions{12pt,a4paper}\fi
\ifdefined\lakitgeometryoptions\else\def\lakitgeometryoptions{margin=2.3cm,headheight=14pt}\fi
\ifdefined\lakitlinespread\else\def\lakitlinespread{1.3}\fi
\ifdefined\lakittocdepth\else\def\lakittocdepth{1}\fi
\ifdefined\lakitsecnumdepth\else\def\lakitsecnumdepth{2}\fi
% Paragraph organization belongs to the document layout, independently of its
% language and theme. Preserve historical class dimensions unless explicitly
% selecting block paragraphs; never suppress indentation with an everypar hack.
\ifdefined\lakitparagraphstyle\else\def\lakitparagraphstyle{indented}\fi
\edef\lakitparagraphstyle{\lakitparagraphstyle}
\def\hllpdfparagraphblock{block}
\def\hllpdfparagraphindented{indented}
\ifx\lakitparagraphstyle\hllpdfparagraphblock\else
  \ifx\lakitparagraphstyle\hllpdfparagraphindented\else
    \errmessage{hllpdf: paragraph style must be block or indented}
  \fi
\fi
% v4.3.1: the public class hook has a deliberately narrow support contract.
% Before this guard, \def\lakitdocumentclass{book} silently compiled even though
% every heading/TOC regression test targets report/article; an unsupported shell
% could therefore look "supported" until a subtle class-specific defect shipped.
% Fail before \documentclass rather than letting an undocumented combination enter
% the mechanism core. Plain-TeX comparison is used because LaTeX packages are not
% available yet at this point.
\edef\lakitvalidatedclass{\lakitdocumentclass}
\def\lakitreportclass{report}
\def\lakitarticleclass{article}
\ifx\lakitvalidatedclass\lakitreportclass
\else\ifx\lakitvalidatedclass\lakitarticleclass
\else
  \errmessage{hllpdf: unsupported class \lakitdocumentclass; use report or article}
\fi\fi
\let\lakitvalidatedclass\relax
\let\lakitreportclass\relax
\let\lakitarticleclass\relax
% Language intent is independent of document class, theme and content domain.
% The two dominant-language profiles differ in authoring intent, not direction.
% A balanced bilingual document must explicitly choose a navigation/base locale.
% <hllpdf:language-profile>
\providecommand{\lakitlanguageprofile}{hebrew}
\edef\lakitlanguageprofile{\lakitlanguageprofile}
\def\hllpdfhebrew{hebrew}\def\hllpdfenglish{english}
\def\hllpdfhebrewmixed{hebrew-with-english}
\def\hllpdfenglishmixed{english-with-hebrew}
\def\hllpdfbilingual{bilingual}
\let\hllpdfexpectedlanguage\relax
\ifx\lakitlanguageprofile\hllpdfhebrew
  \let\hllpdfexpectedlanguage\hllpdfhebrew
\else\ifx\lakitlanguageprofile\hllpdfhebrewmixed
  \let\hllpdfexpectedlanguage\hllpdfhebrew
\else\ifx\lakitlanguageprofile\hllpdfenglish
  \let\hllpdfexpectedlanguage\hllpdfenglish
\else\ifx\lakitlanguageprofile\hllpdfenglishmixed
  \let\hllpdfexpectedlanguage\hllpdfenglish
\else\ifx\lakitlanguageprofile\hllpdfbilingual
  \ifdefined\lakitmainlanguage\else
    \errmessage{hllpdf: bilingual profile requires explicit lakitmainlanguage}
  \fi
\else
  \errmessage{hllpdf: unsupported language profile \lakitlanguageprofile}
\fi\fi\fi\fi\fi
\ifdefined\lakitmainlanguage
  \edef\lakitmainlanguage{\lakitmainlanguage}
\else
  \let\lakitmainlanguage\hllpdfexpectedlanguage
\fi
\ifx\lakitmainlanguage\hllpdfhebrew\else
  \ifx\lakitmainlanguage\hllpdfenglish\else
    \errmessage{hllpdf: main language must be hebrew or english}
  \fi
\fi
\ifx\hllpdfexpectedlanguage\relax\else
  \ifx\lakitmainlanguage\hllpdfexpectedlanguage\else
    \errmessage{hllpdf: main language conflicts with language profile}
  \fi
\fi
\typeout{HLLLANG|\lakitlanguageprofile|\lakitmainlanguage}
% </hllpdf:language-profile>
\documentclass[\lakitclassoptions]{\lakitdocumentclass}
\usepackage[no-math]{fontspec}        % (1) ★ no-math — THE critical line (see recipe.md)
% layout=lists asks Babel to adapt list parshape when paragraph and page
% directions differ. Without it, minority-language list bodies lose their
% indentation and labels hang outside the text block, despite a clean log.
% The upstream mechanism preserves ordinary same-direction list geometry.
\usepackage[bidi=basic,layout=lists]{babel}
\ifx\lakitmainlanguage\hllpdfenglish
  \babelprovide[import]{hebrew}
  \babelprovide[import,main]{english}
\else
  \babelprovide[import,main]{hebrew}
  \babelprovide[import]{english}
\fi
% v4.22: bind canonical families and fallbacks to the selected installation.
% Lookup by basename or family alone is not proof of the file actually used.
% <hllpdf:font-binding>
\directlua{
  local dir = os.getenv("HEBREW_LATEX_FONT_DIR")
  if not dir or dir == "" then
    local data = os.getenv("XDG_DATA_HOME")
    if not data or data == "" or not (data:sub(1,1) == "/") then
      data = (os.getenv("HOME") or "") .. "/.local/share"
    end
    dir = data .. "/fonts/hebrew-lualatex-pdf"
  end
  if not (dir:sub(1,1) == "/") then dir = lfs.currentdir() .. "/" .. dir end
  if dir:find("[",1,true) or dir:find("]",1,true) or dir:find(string.char(0),1,true)
     or dir:find(string.char(10),1,true) or dir:find(string.char(13),1,true) then
    tex.error("hllpdf: unsupported character in font directory")
  end
  if not (lfs.attributes(dir,"mode") == "directory") then
    tex.error("hllpdf: font directory is missing")
  end
  hllpdf_font_dir = dir:gsub("/+$", "") .. "/"
  tex.sprint(string.char(92) .. "def" .. string.char(92) .. "hllpdffontdir{")
  tex.sprint(-2,hllpdf_font_dir)
  tex.sprint("}")
  luaotfload.add_fallback("hebfb",{"[" .. hllpdf_font_dir .. "DejaVuSerif.ttf]:mode=node;"})
  luaotfload.add_fallback("hebfbbf",{"[" .. hllpdf_font_dir .. "DejaVuSerif-Bold.ttf]:mode=node;"})
}
% </hllpdf:font-binding>
% v3.22: Frank Ruhl Libre ships NO italic face. Instead of silently-upright
% \emph (v<=3.20) or bold-as-emphasis (v3.21), synthesize a slanted face with
% FakeSlant — real, VISIBLE italic emphasis with correct \emph nesting
% semantics (inner \emph toggles back upright). Verified empirically at
% 300dpi: slant clearly visible, 0 errors, 0 missing glyphs.
\babelfont{rm}[Path={\hllpdffontdir},RawFeature={fallback=hebfb},
  BoldFont={FrankRuhlLibre-Bold.ttf},
  BoldFeatures={RawFeature={fallback=hebfbbf}},   % v3.56 — see hebfbbf above
  ItalicFont={FrankRuhlLibre-Regular.ttf}, ItalicFeatures={FakeSlant=0.22},
  % Bold-italic uses the weight-matched fallback and a synthetic slant.
  BoldItalicFont={FrankRuhlLibre-Bold.ttf}, BoldItalicFeatures={RawFeature={fallback=hebfbbf},FakeSlant=0.22}
  ]{FrankRuhlLibre-Regular.ttf}
% Sans fallbacks preserve family and weight for the few missing codepoints.
\directlua{luaotfload.add_fallback("hebfbsans",{"[" .. hllpdf_font_dir .. "DejaVuSans.ttf]:mode=node;"})}
\directlua{luaotfload.add_fallback("hebfbsansbf",{"[" .. hllpdf_font_dir .. "DejaVuSans-Bold.ttf]:mode=node;"})}
\babelfont{sf}[Path={\hllpdffontdir},RawFeature={fallback=hebfbsans},
  BoldFont={Heebo-Bold.ttf},
  BoldFeatures={RawFeature={fallback=hebfbsansbf}}]{Heebo-Regular.ttf}
% v3.24: Scale=MatchLowercase — DejaVu Sans Mono has an outsized x-height, so
% unscaled inline code visually SHOUTS next to the delicate Frank Ruhl body
% (measured on a 197-page build: scaled, a code atom sits in the line at the
% body's weight, and the document tightened by ~5%). Applies to listings too.
\babelfont{tt}[Path={\hllpdffontdir},RawFeature={fallback=hebfb},Scale=MatchLowercase,
  BoldFont={DejaVuSansMono.ttf},
  BoldFeatures={RawFeature={fallback=hebfb},FakeBold=1.5},
  ItalicFont={DejaVuSansMono.ttf},
  ItalicFeatures={RawFeature={fallback=hebfb},FakeSlant=0.2},
  BoldItalicFont={DejaVuSansMono.ttf},
  BoldItalicFeatures={RawFeature={fallback=hebfb},FakeBold=1.5,FakeSlant=0.2}
  ]{DejaVuSansMono.ttf}

% Explicit multilingual/English profiles use real Latin faces and italics.
% The legacy/default Hebrew profile retains its established measured typography.
% Latin Modern is provided by the TeX distribution, not the pinned font bundle;
% recorder/runtime evidence reports it as an additional toolchain resource.
\ifx\lakitlanguageprofile\hllpdfhebrew\else
  \babelfont[english]{rm}[Ligatures=TeX,RawFeature={fallback=hebfb},
    BoldFeatures={RawFeature={fallback=hebfbbf}},
    BoldItalicFeatures={RawFeature={fallback=hebfbbf}}]{Latin Modern Roman}
  \babelfont[english]{sf}[Ligatures=TeX,RawFeature={fallback=hebfbsans},
    BoldFeatures={RawFeature={fallback=hebfbsansbf}},
    BoldItalicFeatures={RawFeature={fallback=hebfbsansbf}}]{Latin Modern Sans}
\fi

\usepackage{amsmath,amssymb,mathtools}
% amsthm supplies theorem/proof machinery. Environment names remain
% document-specific; the generic proof label is localized below.
\usepackage{amsthm}
% Proof names follow the surrounding prose, including return from an island.
\addto\extrashebrew{\renewcommand{\proofname}{הוכחה}}
\addto\extrasenglish{\renewcommand{\proofname}{Proof}}
\ifx\lakitmainlanguage\hllpdfenglish
  \renewcommand{\proofname}{Proof}
\else
  \renewcommand{\proofname}{הוכחה}
\fi
% v3.24.1: equation tags under babel-bidi. A tag re-emitted by mathtools'
% showonlyrefs inside align inherits the RTL base and its parens MIRROR —
% ")2.1(" instead of "(2.1)" (plain align/equation tags are unaffected;
% verified on a minimal case + a 198-page build). LTR-isolate at the tag-BOX
% level: mathtools' show/hide + numbering logic is untouched. Do NOT override
% \tagform@ instead — that clobbers the showonlyrefs mechanism (measured:
% every tag shows, numbering breaks).
% v3.28.1: the v3.24.1 wrap must apply ONLY to tags of display environments.
% Inline \eqref in Hebrew PROSE also goes through \maketag@@@, and there the
% english-wrap MIRRORS the parens — ")1(" — while the unwrapped original is
% correct (proven on a 2x2 isolation matrix: patch x showonlyrefs; regressed
% every prose \eqref since v3.24.1, caught by the render gate's R4). The two
% call sites are indistinguishable by \ifmmode (tag boxes are text mode), so
% the display environments themselves raise a flag; \eqref in prose sees it
% false and gets the original untouched \maketag@@@.
\usepackage{etoolbox}

% babel bidi=basic (including upstream 26.9) assumes every automatic
% discretionary has a one-glyph replacement. LuaTeX can legally produce one
% with replace=nil after a hyphenated Latin value; Babel then indexes nil and the
% whole build dies. Do the narrow upstream-compatible correction: while Babel
% classifies the list, present only that malformed disc as an ordinary
% discretionary, then restore its subtype. Babel still reorders the ENTIRE box;
% there is no pcall and no successful fallback to an unreordered node list.
% If Babel throws for any other reason, the build remains fatal. The synthetic
% nil-disc probe and the real table gallery in test_all.sh pin both layers.
\AtBeginDocument{\directlua{
  if Babel and Babel.bidi and not Babel.heb_nil_disc_compat then
    Babel.heb_nil_disc_compat = true
    local original = Babel.bidi
    local disc_id = node.id('disc')
    Babel.bidi = function(head, ispar, hdir)
      local hidden = {}
      for item in node.traverse_id(disc_id, head) do
        if item.subtype == 2 and (not item.replace or not item.replace.char) then
          table.insert(hidden, item)
          item.subtype = 0
        end
      end
      local result = original(head, ispar, hdir)
      for _, item in ipairs(hidden) do item.subtype = 2 end
      return result
    end
  end
}}
\makeatletter
\newif\ifHEBdisptag
\let\ORIGmaketag@@@\maketag@@@
\def\maketag@@@#1{\ifHEBdisptag\ORIGmaketag@@@{\foreignlanguage{english}{#1}}\else\ORIGmaketag@@@{#1}\fi}
\@for\HEB@env:=equation,equation*,align,align*,gather,gather*,multline,multline*,flalign,flalign*,alignat,alignat*,eqnarray,eqnarray*\do{%
  \expandafter\AtBeginEnvironment\expandafter{\HEB@env}{\HEBdisptagtrue}}
\makeatother
% English page + Hebrew prose: Babel's equation-only tag flip uses a primitive
% leqno path that shifts the formula to the left edge and the tag into the body
% of an English page. Align's full-width alignment is unaffected. Give only this
% display a physical LTR container and native amsmath eqno/leqno. The engine then
% owns real tag width, collision handling and labels; no proxy tag box hides ink
% from its fit calculation. Babel's saved \bbl@eqnodir still restores Hebrew text
% inside math, and the prose locale/math direction are never switched.
% Close the prior RTL paragraph before normalization. After the environment's
% group restores its directions, close TeX's empty resumed display paragraph and
% use native LaTeX continuation handling: same-paragraph prose has no indent,
% while an authored blank line starts an ordinary indented RTL paragraph.
% Register after Babel's setup; remain local to equation(*) in a main-English
% document's actual RTL prose. Native Hebrew, ordinary English and explicit
% fleqn retain their original paths. No page/display dimensions are assigned.
% <hllpdf:minority-rtl-equation-geometry>
\makeatletter
\newcommand{\HEB@minorityprosecontinuation}{\par\@doendpe}
\newcommand{\HEB@minorityequationgeometry}{%
  \ifx\lakitmainlanguage\hllpdfenglish
    \ifnum\pardirection=\@ne
      \if@fleqn\else
        \par\aftergroup\HEB@minorityprosecontinuation
        \iftagsleft@\let\veqno\@@leqno
        \else\let\veqno\@@eqno\fi
        \bodydir TLT\pardir TLT\textdir TLT\relax
      \fi
    \fi
  \fi}
\AtBeginDocument{%
  \AddToHook{env/equation/begin}{\HEB@minorityequationgeometry}%
  \AddToHook{env/equation*/begin}{\HEB@minorityequationgeometry}}
\makeatother
% </hllpdf:minority-rtl-equation-geometry>
% v3.63: \shortintertext rendered Hebrew MIRROR-REVERSED and flush LEFT — an entire
% commentary line inside an align printed backwards. babel's bidi=basic patches
% \intertext@ only (luababel.def), and mathtools' \shortintertext uses the identical
% \vbox{\normalbaselines … \noindent\ignorespaces#1\par} skeleton but is never
% patched, so its paragraph keeps pardir/textdir LTR: no bidi reordering (logical
% order on the page) and flush-left. Measured against \intertext with the same
% Applying babel's own replacement to \shortintertext@ makes the two identical.
% \h{}/\mbox fix the letter order but leave the line flush LEFT, so they are not
% substitutes. Guarded + \AtBeginDocument because mathtools \let-copies at load time.
% NB \ExplSyntaxOn is REQUIRED: mathtools' internals carry expl3 names, and without
% the expl3 catcodes the `_` is read as a subscript and the patch detonates far away
% with "! Illegal unit of measure (pt inserted)" at the \end{align} (measured).
\makeatletter
\AtBeginDocument{%
  \@ifundefined{shortintertext@}{}{%
    \ExplSyntaxOn
    \bbl@sreplace\shortintertext@{\normalbaselines}%
      {\normalbaselines\ifx\bbl@eqnodir\relax\else\bbl@pardir\@ne\bbl@eqnodir\fi}%
    \ExplSyntaxOff}}
\makeatother
% geometry inherits the canonical A4 sheet from \documentclass; build.sh then
% verifies every rendered MediaBox and CropBox, so a silent Letter fallback is
% impossible to publish as a clean default build.
\usepackage[\lakitgeometryoptions]{geometry}
\linespread{\lakitlinespread}       % layout hook; canonical default remains 1.3
\ifx\lakitparagraphstyle\hllpdfparagraphblock
  \setlength{\parindent}{0pt}
  \setlength{\parskip}{6pt plus 1pt minus 1pt}
\fi
% These dimensions remain ordinary scoped TeX lengths: authors can customize
% them after core input. Language switches inherit them; lists keep their own
% indentation and narrow reference layouts establish their local paragraph gaps.
\typeout{HLLPAR|\lakitparagraphstyle}
\usepackage{xcolor}
\usepackage{mdframed}
\usepackage{listings}                 % code listings (assembly / C / pseudo-code)
\usepackage{tikz}
\usetikzlibrary{arrows.meta,positioning,calc,shapes.geometric,fit,backgrounds}
\usepackage{pgfplots}\pgfplotsset{compat=1.18}
\usepackage{quantikz}                 % quantum circuits (remove if unused)
\usepackage{booktabs}
\usepackage{array}
\usepackage{float}                    % [H] = place table exactly here (reference doc: no drifting)
\usepackage{enumitem}
\usepackage{needspace}

% ============================================================================
% ★★ v3.4 — THE DIRECTION-ISOLATION TRIO (the heart of mixing Hebrew + Latin)
% Mental model:
%   \h{...}    Hebrew (RTL) island      — Hebrew text inside an LTR context
%                                          (TikZ nodes, the english `tikzpic`).
%   \en{...}   English (LTR) isolate    — Latin text / parentheses inside a
%                                          Hebrew (RTL) context. KEEPS PARENS
%                                          AND ORDER CORRECT.
%   \code{...} \en + monospace          — code inside Hebrew / inside a frametitle.
% CRITICAL: never wrap Latin in \h{} (it is for HEBREW). Inside a TikZ node,
% Latin in \h{} renders REVERSED (Datapath->htapataD, CPU->UPC) and parens
% mirror. Isolate every Latin/paren chunk with \en{} or \code{}. See
% references/figures-boxes-listings.md.
% ---- TEXT INSIDE MATH (v3.8) ----------------------------------------------
% Put ANY text that sits inside math in a \text{}. The ONLY thing that is
% invalid in math is a bare \h/\en (= \foreignlanguage) used DIRECTLY, outside
% \text -> "\rmfamily invalid in math mode". Inside \text:
%   Hebrew label           -> \text{...}            (bare Hebrew renders fine;
%                              \h is optional — BUT NOT ALWAYS HARMLESS (v3.67):
%                              if the \h group OPENS with a nested island
%                              ($..$, \en{}, \code{}, \texttt{}, \mbox{}) the
%                              WHOLE group loses bidi and the Hebrew prints in
%                              LOGICAL order = gibberish. Measured:
%                                \text{\h{$f$ קבועה}}  -> reads "העובק"  BROKEN
%                                \text{\h{קבועה $f$}}  -> correct
%                                \text{$f$ קבועה}      -> correct (no \h)
%                              Hebrew first inside \h, or drop the \h. Flagged
%                              H-OPENS-WITH-LTR-GROUP.)
%   plain Latin word        -> \text{...}            (no marker needed)
%   Latin WITH ( ) or [ ]   -> \text{\en{...}}       (\en forces LTR; a plain
%                              \text mirrors the parens in RTL -> "Spec (TNR(" )
%   f(Hebrew/text arg)      -> \text{f}(\text{...})  -- keep the parens at MATH
%                              level, NOT \text{f(arg)} (parens inside the box mirror)
%   code in math            -> \mathtt{...}  (or \texttt{}, math-guarded below)
% NEVER write \en{...} or \h{...} straight into $...$ / \[..\] (outside \text);
% always via \text{}. (\mbox{...} also works but does not scale in scripts.)
% ============================================================================
\newcommand{\h}[1]{\foreignlanguage{hebrew}{#1}}
% v3.42.1 — \en islands do NOT hyphenate (auto OR at explicit hyphens). English hyphenation patterns were
% tearing identifiers and terms ACROSS THE RTL SEAM ("IO Regis-/ters",
% "In-/put", "Band-/width", "Per-/formance", "Multi-/ple" — all shipped on
% real pages; probe-proven: "ALUS-/rcB", "Mem-/ToReg"). An island is a label
% or a term, not English prose — a hyphen tear there is never right. A/B on
% a full 197pp lakit: identical overfull sets, same page count, ~11
% paragraphs healed, one mild underfull (badness 3525) that replaced a
% mutilated "Regis-/ters". Long English PROSE belongs in the englishprose
% environment below, which restores explicit-hyphen breaks while its paragraph
% is still active. A raw otherlanguage block inherits the global prohibition.
% (\l@nohyphenation needs \csname:
% '@' is not a letter here.)
% v3.43.2 — \exhyphenpenalty moved OUT of \en to a GLOBAL assignment below.
% Render-proven: a group-local \exhyphenpenalty is DEAD CODE — TeX reads the
% line-breaking parameters at \par time, after the \en group has closed, so
% \en{write-through} / plain not-taken still tore at the "-" ("write-/
% through" shipped on a real exam page). \language is different: it is
% recorded in the node list as text is added, so the in-group nohyphenation
% assignment above DOES work (and stays). The global penalty was A/B-measured
% on a 197pp lakit: identical overfull set (25), identical page count, the
% not-/taken tear healed. Hebrew is unaffected in practice — Hebrew compounds
% use maqaf (־), not ASCII "-".
\newcommand{\en}[1]{\foreignlanguage{english}{\language=\csname l@nohyphenation\endcsname #1}}
% Babel runs extras locally for inline islands and paragraph language changes.
% An English main document must not inherit the Hebrew no-hyphen paragraph rule.
\addto\extrashebrew{\exhyphenpenalty=10000\relax}
\addto\extrasenglish{\exhyphenpenalty=50\relax}
\ifx\lakitmainlanguage\hllpdfenglish
  \exhyphenpenalty=50\relax
\else
  \exhyphenpenalty=10000\relax
\fi
% Full English paragraphs need ordinary explicit-hyphen breaks. Use this
% environment instead of a raw otherlanguage block; the paragraph ends while
% the local penalty is still active.
\newenvironment{englishprose}
  {\par\begin{otherlanguage}{english}\exhyphenpenalty=50\relax}
  {\par\end{otherlanguage}}
\newenvironment{hebrewprose}
  {\par\begin{otherlanguage}{hebrew}\exhyphenpenalty=10000\relax}
  {\par\end{otherlanguage}}
% Use for shell text whose locale must not depend on the last body paragraph.
\DeclareRobustCommand{\doclang}[1]{\foreignlanguage{\lakitmainlanguage}{#1}}

% ============================================================================
% ★★ — fixes for defects that are INVISIBLE in the log and in the checks
% and only surface on a VISUAL RENDER of the PDF. (Each verified empirically.)
% ----------------------------------------------------------------------------
% v3.22: \emph is now REAL slanted emphasis (FakeSlant italic face defined on
% the \babelfont{rm} line above), replacing the v3.21 \emph->\textbf hack: it
% preserves \emph's toggle semantics and keeps bold free for its own role.
% Prefer bold-as-emphasis for your doc? one line: \renewcommand{\emph}[1]{\textbf{#1}}
% page number: force the digits LTR so they never reverse (18 -> 81) in the RTL
%   footer. That reversal is a NON-DETERMINISTIC bidi glitch (observed once in
%   ~80 pp); this LuaLaTeX build is not byte-reproducible, so the defensive
%   LTR wrap is the only reliable guard. Latin digits are always LTR anyway.
\renewcommand{\thepage}{\foreignlanguage{english}{\arabic{page}}}
% v3.70: \pagenumbering \gdef's \thepage and silently DESTROYS the wrap above —
% after \pagenumbering{arabic} the folio guard is gone for the whole rest of the
% document. Two collections discovered this independently and hand-restore the
% wrap in main.tex with the same comment; that knowledge now lives here instead.
% The wrapper re-applies the LTR wrap after the kernel call, for any style.
\makeatletter
\let\lakin@orig@pagenumbering\pagenumbering
\renewcommand{\pagenumbering}[1]{%
  \lakin@orig@pagenumbering{#1}%
  \renewcommand{\thepage}{\foreignlanguage{english}{\csname @#1\endcsname\c@page}}}
\makeatother
% emergencystretch: Hebrew does not hyphenate, so an unbreakable LTR island
%   (\en{...} / an inline math run) landing at the margin can leave a small
%   (~2-13pt) overfull with no legal breakpoint. This rescues ONLY those
%   inline-prose tie-breaks; it is a ceiling used only on offending lines and
%   does NOT touch display-math / listing / table overflow (those still surface
%   in the log for the content-level overflow pass). Verified: baseline overfull
%   -> 0 with this set, with no change to pagination or underfull.
\setlength{\emergencystretch}{3em}
% v3.40: inline math (almost) never breaks across lines. TeX's default
%   (binop 700 / rel 500) lets an inline run split at an operator —
%   $(1.f)\cdot 2^{e}$ shipped as "(1.f)·" ending one line and "2^e" opening
%   the next (reader-reported). In LTR that split is merely ugly; under an
%   RTL base the carried operator lands on the visually WRONG side of the
%   break and the formula reads backwards.
%   WHY 10000 (v3.45.0 — the display-first resolution). Inline math is
%   FORBIDDEN from breaking, so a torn island under an RTL base is
%   structurally impossible. The v3.40 cycle used 9999 ("break only as a last
%   resort") because a bare 10000 left ~30-39 WIDE inline equations
%   overflowing (worst 90.7pt) — but 9999's "last resort" fires OFTEN in
%   tight, non-hyphenating Hebrew: a reader found torn islands across the
%   fleet (~150 candidates; e.g. "P(y |" ending one line and "x_0)" opening
%   the next, read backwards). The step the 9999 cycle missed: every one of
%   those ~30 overflows is a wide equation that BELONGS in display. Convert
%   them to display FIRST, then 10000 gives 0 tears AND 0 overflows (measured,
%   12/12 fleet). Doctrine: inline math is ATOMIC; a wide equation goes in a
%   display, never inline. See content-style ("inline math is atomic").
\binoppenalty=10000 \relpenalty=10000
% v3.38: uniform sentence spacing. TeX's default (\nonfrenchspacing, sfcode
% 3000) is the OLD Anglo-American wide-space-after-period convention — Hebrew
% typography never had it, and modern Latin practice (Chicago et al.) dropped
% it too. Measured on real pages: post-period gaps up to 1.9-2.1x the line's
% mean word gap — the dominant "rivers" source in justified Hebrew prose.
% Survives babel language switches (probed: sfcode stays french inside \en{}
% islands and otherlanguage blocks).
\frenchspacing
% v3.38: widow/orphan protection. LaTeX's default penalties (150/150/50)
% allow a lone line at a page top/bottom almost freely. 10000 = forbid;
% this is SAFE here precisely because the document is \raggedbottom (the
% demoted line just moves on, no vertical stretch fight) — under flushbottom
% these values would cause underfull pages. raggedbottom is the report-class
% default; made explicit so a future class/option change cannot silently
% flip it and turn the 10000s harmful.
\raggedbottom
\widowpenalty=10000 \clubpenalty=10000 \displaywidowpenalty=10000
% em-dash: --- does NOT become — under Frank Ruhl (the font lacks tlig; verified
%   the width is identical with/without Ligatures=TeX, while the mechanism works
%   on Latin Modern). TYPE THE CHARACTER — (U+2014) DIRECTLY; likewise – (U+2013).
%   check_bidi_figures.py flags a --/--- that is adjacent to Hebrew.
% ============================================================================

% ---- STRUCTURAL DIRECTION-SAFETY for \texttt (v3.7) -------------------------
% A parenthesis, bracket, or $ INSIDE \texttt{} mirrors/migrates in RTL prose
% (MIPS 0($s0)->"(s0$)0", $rs->"rs$", arr[i]->"[arr[i") because \texttt sets the
% font but NOT the direction. Redefining \texttt to LTR-isolate its content fixes
% the WHOLE class at once (parens, brackets, $) so plain \texttt{lw $t0,0($s0)}
% is safe — no need to remember \code{}. Guarded by \ifmmode because
% \foreignlanguage is invalid in math (there, plain monospace is already LTR).
% NB (corrected v3.56 — the old note here said "stays tofu", which is WRONG):
% \texttt{Hebrew} is a FATAL build failure. The tt family is DejaVu Sans Mono
% (zero Hebrew) and the tt fallback hebfb is DejaVuSerif (also zero Hebrew), so
% the subsetter gets an empty glyph set and LuaTeX dies with
%   "! error: (file .../DejaVuSansMono.ttf) (type 2): there are no glyphs in the
%    subset"  -> no output PDF.
% The message names a FONT FILE, never Hebrew/\texttt/the source line. And it is
% non-local: add an unrelated Latin \code{} elsewhere and the fatal degrades into
% a SILENT DELETION of the Hebrew word. Put Hebrew in \h{}/normal text, never in
% \texttt{}/\code{}; flagged HEBREW-IN-MONO. To disable, delete this block.
\let\origtexttt\texttt
\protected\long\def\texttt#1{\ifmmode\origtexttt{#1}\else\foreignlanguage{english}{\origtexttt{#1}}\fi}
% v4.8.0: defer the direction test until TeX has inserted an alignment cell's
% math template. Unprotected expansion here tested \ifmmode too early at the
% FIRST token of aligned/array/cases/matrix cells. Engine-level \protected also
% preserves the original \texttt implementation without the self-recursion
% observed when retrofitting \DeclareRobustCommand to its saved alias.
\protected\long\def\code#1{\ifmmode\origtexttt{#1}\else\foreignlanguage{english}{\origtexttt{#1}}\fi}

% ---- \codeline — display-code line (v3.24): the code analogue of display math.
% Reader-friction rule (content-style.md §5b): inline \code{} is for a SINGLE
% ATOM only (one register/variable/mnemonic, a tight no-space expression).
% A multi-token expression — an instruction with operands, a spaced assignment,
% a register-binding list — zigzags the eye RTL<->LTR mid-sentence and, being
% an unbreakable island, line-breaks badly. Give it its own centered line with
% display-math vertical rhythm. \centerline = one box => the expression never
% wraps; \nopagebreak keeps it glued to its introducing sentence.
%   ...the Hebrew sentence ends with a colon:
%   \codeline{lw \$t1, -4(\$t0)}
%   and the prose continues (escape dollars as \$).
% v3.66 (first line of the body): flush a PENDING \item label before the \par.
% LaTeX defers the label box (\@labels) until the next \everypar, and an \hbox in
% vertical mode never fires \everypar — so `\item \codeline{...}` printed the code
% at y=69.1 and the item's own "1." BELOW it at y=103.6 (the number appears after
% its content, attached to nothing). \if@inlabel\leavevmode\fi is CONDITIONAL on
% purpose: an unconditional \leavevmode would open an empty paragraph line when
% \codeline follows ordinary prose. Verified both ways — in a list the label now
% sits on the code's line at its normal x, and after prose the geometry is
% byte-identical to before.
\makeatletter
\newcommand{\codeline}[1]{%
  \if@inlabel\leavevmode\fi
  \par\nopagebreak
  \vspace{\abovedisplayskip}\nopagebreak
  % v3.62: was \centerline{...}, i.e. \hbox to\hsize{\hss #1 \hss}. \hss has
  % INFINITE SHRINK, so TeX could always "fit" the box and NEVER reported an
  % Overfull — the same \hss mechanism as BLD-27's \item label. Measured: a wide
  % \codeline in a defbox painted its first token ON the box's left border bar and
  % sat 9.5pt outside the frame interior on both sides, with 0 Overfull, 0 Underfull
  % and all four gates PASS; in bare prose it reached 21.0pt into BOTH margins.
  % \hfil has no shrink, so the overflow is reported and build.sh's ovf_vis/ovf_big
  % counters see it; \linewidth (not \hsize) makes it box-aware, so a \codeline
  % inside a colored box centres on the box interior. Verified: a FITTING \codeline
  % keeps byte-identical glyph positions, and the box-crossing case now emits
  % "Overfull \hbox (19.03pt too wide)". A wide \codeline should be split in two.
  \hbox to\linewidth{\hfil\code{#1}\hfil}%
  \nopagebreak\vspace{\belowdisplayskip}}
\makeatother

% ---- a thin decorative rule helper ----
\newcommand{\sep}{\par\vspace{2pt}\noindent\textcolor{cRule}{\rule{\linewidth}{0.6pt}}\par\vspace{4pt}}

% ---- colors ----
\definecolor{cBlue}{HTML}{1F4E79}   \definecolor{cBlueBg}{HTML}{EAF1F8}
\definecolor{cGreen}{HTML}{2E6B4F}  \definecolor{cGreenBg}{HTML}{EAF4EE}
\definecolor{cAmber}{HTML}{9C6A00}  \definecolor{cAmberBg}{HTML}{FBF3E2}
\definecolor{cPurple}{HTML}{5B2C6F} \definecolor{cPurpleBg}{HTML}{F3ECF7}
\definecolor{cTeal}{HTML}{0E6E6E}   \definecolor{cTealBg}{HTML}{E6F4F4}
\definecolor{cRed}{HTML}{8E2B2B}    \definecolor{cRedBg}{HTML}{FBEAEA}
\definecolor{cGray}{HTML}{555555}   \definecolor{cRule}{HTML}{1F4E79}

% ---- math operators ----
\DeclareMathOperator{\Tr}{Tr}
\DeclareMathOperator{\rank}{rank}
\DeclareMathOperator{\diag}{diag}
\DeclareMathOperator{\sgn}{sgn}

% ---- bra-ket (quantikz defines \ket/\bra/\braket; add the rest safely) ----
\providecommand{\ket}[1]{\left|#1\right\rangle}
\providecommand{\bra}[1]{\left\langle#1\right|}
\providecommand{\braket}[2]{\left\langle#1\middle|#2\right\rangle}
\newcommand{\dyad}[1]{\left|#1\right\rangle\!\left\langle#1\right|}
% v3.21: general outer product |a><b| and expectation value <A> — agents reach
% for these by reflex (esp. \expval); quantikz defines neither, so their absence
% threw "Undefined control sequence" in a real build. \providecommand = safe.
\providecommand{\ketbra}[2]{\left|#1\right\rangle\!\left\langle#2\right|}
\providecommand{\expval}[1]{\left\langle#1\right\rangle}
\newcommand{\abs}[1]{\left|#1\right|}
\newcommand{\norm}[1]{\left\lVert#1\right\rVert}
\newcommand{\half}{\tfrac12}
\newcommand{\R}{\mathbb{R}}\newcommand{\C}{\mathbb{C}}\newcommand{\Z}{\mathbb{Z}}
% identity operator. NB: \mathbb{1} silently renders a WRONG glyph (msbm has no
% blackboard digits). Use \mathbb{I} = 𝕀 (works with amssymb, no extra package).
% Do NOT "fix" this back to \mathbb{1}.
\newcommand{\id}{\mathbb{I}}

% ---- physical units (v3.35): \punit — upright unit with thin space ---------
% Usage: $300\punit{K}$, $80\punit{mV/decade}$, $10^{7}\punit{cm/s}$ — the
% NUMBER AND UNIT LIVE IN ONE MATH ISLAND. Promoted verbatim from a 199-page
% production lakit (advanced-electronic-devices, dozens of uses).
% Two measured traps this design avoids (both render-proven):
%   1. siunitx SCRAMBLES under RTL — \SI{3.5}{eV} renders "eV 3.5", exponents
%      detach ("Hz 10 x 10^9"): its quantity assembly is not bidi-isolated.
%      Do NOT load siunitx in this pipeline.
%   2. A text-mode unit next to a bare-prose number is TWO adjacent LTR runs
%      and reorders ("5 mm" -> "mm5" — the BIDI-2/3 class). \punit is therefore
%      math-only ON PURPOSE: prose misuse fails LOUDLY at compile (\mathrm
%      outside math) instead of scrambling silently.
% Domain shorthands (\MeV, \fm...) stay per-document: \newcommand{\MeV}{\,\text{\en{MeV}}}
\providecommand{\punit}[1]{\,\mathrm{#1}}

% ============================================================================
% Colored boxes — mdframed (NOT tcolorbox, which breaks under bidi).
% ★★ v3.4: innerleftmargin/innerrightmargin give code listings room so their
% frame/line-numbers do not spill out the box border (see listing styles below).
% ============================================================================
% v3.61: everyline/split skips were applied to `tablebox` only (BLD-14, v3.37) —
% the six colored boxes below never got them, so they still split OPEN across a page
% break: the page-N segment gets top+sides and NO bottom rule, the page-N+1 segment
% sides+bottom and NO top rule. Measured on the shipped fleet this is a COVERAGE GAP,
% not a re-report of BLD-14: 0 split-open `tablebox` segments (that fix holds) vs 62
% split-open COLORED boxes across 7 documents. Verified after the fix: both segments
% close, and the frametitle is NOT duplicated on the continuation. BLD-14's caveat
% still applies — a box taller than a full page grinds mdframed's split-trial loop, so
% split such content manually.
% ============================================================================
% v3.72: frametitle={} leads EVERY \newmdenv option list below. mdframed options
% are GROUP-scoped: a box nested inside \begin{exbox}[frametitle={...}] used to
% INHERIT the outer title — measured, the identical title string rendered twice,
% the second riding a title-bar fill inside the INNER box's frame, for all seven
% environments in both directions. Both natural forms were broken (= preamble
% bug, not an authoring rule). The explicit reset makes every box titleless
% unless ITS OWN [frametitle={...}] says otherwise; verified both directions:
% no duplicate on the nested box, an explicit title still renders.
% Geometry is shared with the optional semantic-width table container. Themes
% select their own existing semantic style below; neither path fixes columns.
\mdfdefinestyle{hllpdfsemanticgeometry}{linewidth=1.1pt,
  innertopmargin=6pt,innerbottommargin=7pt,skipabove=9pt,skipbelow=9pt,
  innerleftmargin=9pt,innerrightmargin=9pt,everyline=true,
  splitbottomskip=6pt,splittopskip=12pt,roundcorner=3pt}
\newcommand{\hllpdfsemanticframestyle}{hllpdfsemanticgeometry}
\newmdenv[frametitle={},style=hllpdfsemanticgeometry,linecolor=cBlue,backgroundcolor=cBlueBg,
  frametitlebackgroundcolor=cBlue,frametitlefontcolor=white,
  ]{defbox}   % הגדרה
\newmdenv[frametitle={},style=hllpdfsemanticgeometry,linecolor=cGreen,backgroundcolor=cGreenBg,
  frametitlebackgroundcolor=cGreen,frametitlefontcolor=white,
  ]{thmbox}   % משפט/טענה
\newmdenv[frametitle={},style=hllpdfsemanticgeometry,linecolor=cAmber,backgroundcolor=cAmberBg,
  frametitlebackgroundcolor=cAmber,frametitlefontcolor=white,
  ]{notebox}  % הערה
\newmdenv[frametitle={},style=hllpdfsemanticgeometry,linecolor=cPurple,backgroundcolor=cPurpleBg,
  frametitlebackgroundcolor=cPurple,frametitlefontcolor=white,
  ]{exbox}    % דוגמה
\newmdenv[frametitle={},style=hllpdfsemanticgeometry,linecolor=cRed,backgroundcolor=cRedBg,
  frametitlebackgroundcolor=cRed,frametitlefontcolor=white,
  ]{warnbox}  % טעות נפוצה
\newmdenv[frametitle={},style=hllpdfsemanticgeometry,linecolor=cTeal,backgroundcolor=cTealBg,
  frametitlebackgroundcolor=cTeal,frametitlefontcolor=white,
  ]{keybox}   % רעיון מרכזי

% ---- codebox (v4.5.0): the dedicated container for code -------------------
% An OPTION, not a policy: asmblock unboxed stays canonical for a full listing,
% and a snippet inside any of the six colored boxes still works as before.
% What makes codebox worth its own environment is that it is the ONE framed box
% that keeps its margin line numbers.
%
% v4.4.0 shipped it WITHOUT them, on the strength of v3.37's "no margin knob can
% save them" plus four fresh probes that agreed. That was wrong: every probe had
% moved the code, the box, or both — never the key that carries the NUMBER
% COLUMN. Measured on the page:
%   resetmargins=false   the listing follows THIS box instead of the page,
%                        moving the code from x=75.26 to x=97.54
%   xrightmargin=26pt    under RTL this is the PHYSICAL LEFT indent, and the
%                        only key that carries the \llap'ed digits with it:
%                        x=64.28 -> 90.18, clearing the rule at 65.64 by 23.6pt
%   xleftmargin=26pt     moves the code and LEAVES THE DIGITS on the rule —
%                        the asymmetry that hid the fix for eleven versions
% resetmargins is hoisted out of the four styles (see the \lstset below them) so
% this override is possible at all; every other listing still resolves to true.
% Pinned by tests/codebox_numbers.tex + probe_lstnum.py --inside-rule, whose
% negative control strips xrightmargin and watches the digits fall back.
\definecolor{cSlate}{HTML}{3F4A56}  \definecolor{cSlateBg}{HTML}{F2F4F7}
\newmdenv[frametitle={},leftline=true,rightline=true,topline=true,bottomline=true,
  linecolor=cSlate,backgroundcolor=cSlateBg,linewidth=0.9pt,
  frametitlebackgroundcolor=cSlate,frametitlefontcolor=white,
  innertopmargin=6pt,innerbottommargin=7pt,skipabove=9pt,skipbelow=9pt,
  innerleftmargin=9pt,innerrightmargin=9pt,everyline=true,splitbottomskip=6pt,splittopskip=12pt,
  roundcorner=3pt]{codebox}   % קוד
% The two keys that make in-box numbering work, both measured on the page:
%   resetmargins=false  the listing follows THIS box instead of the page
%   xrightmargin        under RTL this is the PHYSICAL LEFT indent, and it is
%                       the only one that carries the \llap'ed number column
%                       with it. xleftmargin moves the code and leaves the
%                       digits behind on the rule (measured: code 97.54,
%                       digits 64.28 — the v3.37 defect, reproduced); with
%                       xrightmargin=26pt the digits land at x=90.18, a
%                       clearance of 23.6pt from the rule at 65.64.
% This asymmetry is why v3.4's note had the semantics backwards and v3.37
% concluded the digits could not be saved.
% <hllpdf:listing-container-state>
% v4.29: all state is local to the current environment group. A semantic
% callout masks number text, while a nested codebox restores the formatter
% saved before that mask. Save it once across successive semantic wrappers;
% an explicit numbers=none (or an explicitly empty formatter) remains valid.
\makeatletter
\newif\ifhllpdf@incodebox
\newif\ifhllpdf@englishcodebox
\newif\ifhllpdf@lstnumbermasked
\let\hllpdf@savedlstnumber\thelstnumber
\makeatother
% </hllpdf:listing-container-state>
% <hllpdf:codebox-gutter>
\makeatletter
\AtBeginEnvironment{codebox}{%
  \hllpdf@incodeboxtrue
  % Snapshot the OUTER language before asmblock selects English.
  \iflanguage{english}{\hllpdf@englishcodeboxtrue}{\hllpdf@englishcodeboxfalse}%
  \ifhllpdf@lstnumbermasked
    \let\thelstnumber\hllpdf@savedlstnumber
    \hllpdf@lstnumbermaskedfalse
  \fi
}
% listings' AddToHook is GLOBAL. Register once, and gate on local box state;
% adding an unconditional hook at every codebox entry leaks into later listings.
\lst@AddToHook{PreInit}{%
  \ifhllpdf@incodebox
    % Apply every listing override here, inside the LISTING's own group.
    % Setting margins at box entry leaks them into nested semantic callouts.
    \lstset{resetmargins=false,xrightmargin=26pt}%
    % Styles resolve before PreInit. In LTR the physical number gutter follows
    % xleftmargin; in Hebrew retain the measured inherited xrightmargin path.
    \ifhllpdf@englishcodebox
      \lstset{xleftmargin=26pt,xrightmargin=0pt}%
    \fi
    % ONE box, not two: suppress only the frame/background/skips of a listing
    % inside this codebox. The next unboxed listing keeps its own style.
    \lstset{frame=none,backgroundcolor=\color{cSlateBg},aboveskip=0pt,belowskip=0pt}%
  \fi
}
\makeatother
% </hllpdf:codebox-gutter>

% v3.37: line numbers inside a COLORED BOX are sliced by the box border —
% measured: digit span x=63.5-67.3 vs border line at x=65.2 (the border runs
% THROUGH the digits), and the inter-border gap (3.5pt) is narrower than a
% digit (3.8pt), so no margin knob can save them (under resetmargins,
% xleftmargin is inert in-box — measured: code at x=75.26 with 2.2em AND
% 3.0em). A 3-6 line boxed snippet does not need line numbers; blank the
% number TEXT inside every colored box (\thelstnumber is what gets printed;
% the style's numbers=left survives, so nothing else shifts — code position
% byte-frozen at x=75.26, verified). Unboxed listings keep their numbers.
% v3.78: the blanking is coupled to the FRAME, not to the box name —
% clean_theme.tex registers LATER hooks (registration order wins) that
% restore the number text inside the boxes it renders frameless.
% v4.5.0 RETRACTS the "no margin knob can save them" clause above. It holds only
% under resetmargins=true (the shipped default): there xleftmargin is inert
% in-box, as measured. With resetmargins=false, xrightmargin DOES carry the
% \llap'ed digit column clear of the rule — measured; the codebox \lstset that
% does it is the one above, in the codebox block. The loop below therefore
% covers SIX boxes, not "every colored box": codebox opts out and keeps its
% numbers.
% <hllpdf:semantic-listing-mask>
\makeatletter
\@for\HEB@lstbox:=defbox,thmbox,notebox,exbox,warnbox,keybox\do{%
  \expandafter\AtBeginEnvironment\expandafter{\HEB@lstbox}{%
    \ifhllpdf@lstnumbermasked\else
      \let\hllpdf@savedlstnumber\thelstnumber
    \fi
    \hllpdf@lstnumbermaskedtrue
    \def\thelstnumber{}%
    % The nearest semantic container owns its listing presentation. On exit,
    % TeX's group restores the enclosing codebox state and formatter.
    \hllpdf@incodeboxfalse}}
\makeatother
% </hllpdf:semantic-listing-mask>

% ============================================================================
% Unified TABLE presentation: subtle uniform slate frame + auto-numbered
% caption rendered INSIDE the frame. Use:
%   \begin{tablebox}\centering  <tabular>  \tabcap{<description>}\end{tablebox}
% Replaces the buggy \caption* (which emitted a stray "טבלה N : *" line) and
% provides an explicit framed option + numbered caption "טבלה N.M: תיאור".
% Optional intro text may be placed at the TOP, before the tabular, inside the box.
% ============================================================================
\definecolor{cTblLine}{HTML}{B9C4D2}  \definecolor{cTblBg}{HTML}{FAFBFC}
% v3.23: class-aware — under `article` there is no chapter counter (the promised
% report->article switch used to die with "No counter 'chapter' defined").
\makeatletter
\@ifundefined{c@chapter}{%
  % v3.31: in article the counter resets per SECTION, so an unprefixed number
  % repeated ("טבלה 1." in every section — measured on a two-section render).
  % Same conditional idea as the report branch: prefix the section number once
  % a real \section exists; a pre-section table stays "טבלה 1".
  \newcounter{tbl}[section]
  \renewcommand{\thetbl}{\ifnum\value{section}>0 \thesection.\fi\arabic{tbl}}%
}{%
  \newcounter{tbl}[chapter]
  % v3.30: a report-class doc that uses SECTIONS ONLY (chapter stays 0 —
  % the common short-doc pattern) used to print "טבלה 0.1"; the chapter
  % prefix now appears only once a real \chapter exists (render-verified
  % both ways: "טבלה 1." chapterless, "טבלה 1.1" with chapters).
  \renewcommand{\thetbl}{\ifnum\value{chapter}>0 \thechapter.\fi\arabic{tbl}}%
  % v3.64: the SAME leak in \thesection itself, left behind by v3.30. That fix
  % guarded the CAPTION counter and stopped there, so a chapterless report still
  % printed the heading "0.1 סעיף ראשון" and \ref to it printed "0.2" — measured:
  %   p1 y=65.2  '…0.1'   p1 y=144.9 '…0.2'   p1 y=183.4 '…0.2.1'
  % report.cls hardcodes \thesection as \thechapter.\@arabic\c@section, and the
  % leading 0 is the chapter counter that no \chapter ever incremented. Fixing
  % \thesection fixes the whole chain — \thesubsection and \thesubsubsection are
  % defined in terms of it. \theequation and \thetbl/\thefig were already guarded,
  % which is exactly why this survived: every NUMBER around the heading looked
  % right, so the heading's own 0 read as a style choice rather than a bug.
  % \appendix is safe: report.cls resets chapter to 0 and redefines \thechapter to
  % \@Alph\c@chapter, so an appendix \chapter makes the counter 1 and the guard
  % emits "A.1" unchanged. Render-verified in all four modes (chapterless,
  % chaptered, appendix+chapter, appendix+section-only).
  \renewcommand{\thesection}{\ifnum\value{chapter}>0 \thechapter.\fi\arabic{section}}%
}
\makeatother
% v3.37: everyline=true — if the box DOES split across pages, every segment
% gets a CLOSED frame (measured before: the split ran open-ended into the page
% edge with no bottom border, and the caption could strand alone at the next
% page top — "טבלה A.2" shipped that way). The split skips give the segments
% breathing room. In practice, with the caption GLUE below, every previously-
% split table in the corpus now moves WHOLE instead (the white-gap tradeoff,
% documented benign); everyline stays as the safety net for a genuinely
% page-tall table. NB: a table taller than ONE FULL PAGE can grind mdframed's
% split-trial loop (reproduced: minimal doc, minutes of 1pt-step Overfull
% probes) — break such a monster into two tables.
\mdfdefinestyle{hllpdftablegeometry}{leftline=true,rightline=true,topline=true,bottomline=true,
  linewidth=0.9pt,
  innertopmargin=9pt,innerbottommargin=8pt,skipabove=11pt,skipbelow=11pt,
  innerleftmargin=11pt,innerrightmargin=11pt,roundcorner=4pt}
\newcommand{\hllpdftableframestyle}{hllpdftablegeometry}
\newmdenv[frametitle={},style=hllpdftablegeometry,
  linecolor=cTblLine,backgroundcolor=cTblBg,
  % v4.2.0 (adopted): tablebox is ATOMIC. Short tableboxes near page bottoms
  % repeatedly entered mdframed's split trial, emitting benign Overfull-vbox
  % bursts and occasionally separating the intro prose from the tabular.
  % nobreak=true moves the complete short table to the next page instead. The
  % output-active Overfull-vbox gate still catches a box taller than a page;
  % everyline/split skips remain as defensive styling should that path ever
  % re-arm, but the canonical path never publishes a split tablebox.
  nobreak=true,
  everyline=true,splitbottomskip=6pt,splittopskip=12pt,
  % v3.79 (BLD-37 table side, correction): startinnercode/endinnercode run
  % INSIDE mdframed's body vbox (mdframed.sty:848 / :856+863) -- the second one
  % after \par\unskip, i.e. after line breaking, which is the only moment the
  % \tabcap anchor's true offset from the top of the box EXISTS. See
  % \lakintabboxclose below. These two keys are the ONLY tablebox change.
  startinnercode={\lakintabboxopen},endinnercode={\lakintabboxclose}]{tablebox}
% v3.37: \nopagebreak GLUE — \vsplit honors penalties, so the caption can no
% longer strand alone at the top of the next page (it takes its rows along, or
% the whole box moves). Render-verified on the real shipped orphan.
% v3.59: \NR@gettitle is what nameref injects into the stdlib title-bearing
% commands (\@sect, \@chapter, \caption) to fill \@currentlabelname. \figcap and
% \tabcap are INDEPENDENT macros — they never called it, so \@currentlabelname kept
% whatever the last heading left and \nameref{fig:x} silently printed the enclosing
% SECTION's title instead of the caption (proven in the .aux: field 3 of \newlabel
% was the heading text while field 1 and the anchor `fig.1.1` were correct — so this
% is a THIRD, distinct head next to BLD-13's counter fix and BIDI-38's autoref-name
% fix, neither of which writes that field). \@ifundefined keeps it inert if
% hyperref/nameref is ever dropped. Guarded helper because \figcap is defined
% after \makeatother.
% NB both helpers below need \makeatletter — this part of the preamble is OUTSIDE
% any letter-catcode region, so an unwrapped \@currentlabel parses as LaTeX's \@
% and dies with "You can't use \spacefactor in vertical mode" (measured).
\makeatletter
\newcommand{\lakinrtitle}[1]{\@ifundefined{NR@gettitle}{}{\NR@gettitle{#1}}}
% v3.59: \hitem{<letter>}{<key>} — a REFERENCEABLE Hebrew-lettered list item.
%   \begin{enumerate} \hitem{א}{it:a} הסעיף הראשון. \end{enumerate}  ->  \ref{it:a} = (א)
% Why this exists: automatic Hebrew list labels fail silently, so SKILL.md mandates
% an EXPLICIT \item[(\textbf{א})] — but LaTeX puts \refstepcounter INSIDE
% \if@noitemarg (latex.ltx:12964), and \item[...] never sets that flag. So an
% explicit item steps no counter and never touches \@currentlabel: a \label there
% binds to the AMBIENT one and \ref prints the CHAPTER/SECTION number where the
% item letter belongs (proven in the .aux: {chapter.1}). Same doctrine as \figcap —
% own the label state before \label. \phantomsection matters: without it the
% hyperlink still lands on the chapter head (measured); with it the anchor is
% section*.N and the bookmark outline stays clean.
\newcommand{\hitem}[2]{\item[(\textbf{#1})]%
  \def\@currentlabel{(#1)}\def\@currentlabelname{#1}\phantomsection\label{#2}}
% NB (v3.66): \hitem is \ref-only. \autoref{key} calls it a SECTION ("סעיף (א)")
% because \phantomsection brands the anchor section*.N, and \nameref{key} prints
% the bare letter. For \autoref-able items prefer \hebenumlabels below.
%
% v3.66: \hebenumlabels — OPT-IN automatic Hebrew enumerate letters, done right.
%   {\hebenumlabels\begin{enumerate}\item\label{q:a} ...\end{enumerate}}
%   -> labels (א) (ב) (ג); \ref{q:a} = א; anchor Item.N (real, per-item).
% This retires the belief that made \item[(\textbf{א})] + \hitem necessary.
% "Automatic Hebrew labels fail silently" was TRUE but unexplained, and the
% mechanism is now bisected (plain article vs +enumitem, measured): ENUMITEM
% PRE-EXPANDS the label expression in an \edef BEFORE \refstepcounter, with the
% counter still at 0. Tokens it shields by name (\arabic,\alph,\roman) survive;
% a plain \newcommand \ifcase wrapper evaluates at 0 -> the EMPTY zero-branch ->
% the label prints as "()" with the letter GONE (and \the\c@enumi prints a
% literal 0, \number\value dies fatally pointing at \item). The cure is one
% property, not a doctrine: the wrapper must be ROBUST (survives the \edef
% unexpanded), while \theenumi stays EXPANDABLE (so \refstepcounter writes the
% real letter to the .aux). Verified: labels (א)(ב)(ג), \ref -> א/ב/ג each item
% its own value, .aux = {א}{...}{Item.1} — versus the explicit-\item doctrine
% where BIDI-50/\hitem exist purely to patch the counter that never stepped.
\def\@hebalph#1{\ifcase#1\or א\or ב\or ג\or ד\or ה\or ו\or ז\or ח\or ט\or י\or
  יא\or יב\or יג\or יד\or טו\or טז\or יז\or יח\or יט\or כ\else\@ctrerr\fi}
\DeclareRobustCommand*\hebalph[1]{\expandafter\@hebalph\csname c@#1\endcsname}
\newcommand{\hebenumlabels}{%
  \renewcommand{\theenumi}{\@hebalph\c@enumi}%
  \renewcommand{\labelenumi}{(\textbf{\hebalph{enumi}})}}
\makeatother
% v3.63: \mbox the LABEL. \thetbl ("7.1") and a caption description that STARTS
% WITH A DIGIT ("32 הרגיסטרים…") are two European-Number runs separated only by
% neutrals (the label's own "." and \enspace). Under the RTL paragraph base the UBA
% resolves those neutrals to the RTL level, so both numbers sit in ONE reversed span
% and EXCHANGE VISUAL POSITIONS. Measured on the shipped PDF: a table whose real
% number is 7.1 printed "טבלה 32" — a different, EXISTING table number — while \ref
% correctly printed 7.1, so the caption silently contradicted its own cross-reference,
% and the description lost its leading count. The source contains no adjacency at all;
% it is manufactured by this macro, which is why every source linter was blind.
% \mbox makes the label one atomic box so nothing can reorder across it. Verified
% zero-regression: an ordinary Hebrew caption keeps identical x-geometry.
% One paragraph style covers label, number, separator and description.
% Explicit project/template overrides remain possible; the label hook is kept
% for compatibility, with an identity default. Counters/anchors stay unchanged.
\newcommand{\lakitcaptionstyle}{\normalfont\small\color{black}\hypersetup{allcolors=black}}
\newcommand{\lakitcaptionlabelstyle}[1]{#1}
\newcommand{\lakitcaptionpunctuation}{:}
% Native floats, captionof and longtable use the same presentation policy.
% subcaption is optional and inherits the same font rather than a second size.
\usepackage{caption}
\DeclareCaptionFont{lakit}{\lakitcaptionstyle}
\DeclareCaptionLabelFormat{lakit}{\lakitcaptionlabelstyle{#1~\en{#2}}}
\DeclareCaptionLabelSeparator{lakit}{\lakitcaptionpunctuation\enspace}
\captionsetup{font=lakit,labelfont={},textfont={},labelformat=lakit,
  labelsep=lakit,justification=centering}
\makeatletter
\AtBeginDocument{\@ifpackageloaded{subcaption}{%
  \captionsetup[sub]{font=lakit,labelfont={},textfont={},labelsep=lakit}}{}}
\makeatother
% v3.64: OPT-IN list-of-figures/tables entry via an optional short caption.
% \listoffigures and \listoftables are inherited from report.cls and LOOK available,
% but \figcap/\tabcap step the custom \c@fig/\c@tbl counters and never call
% \addcontentsline — so both printed their heading and NOTHING ELSE. Measured: a doc
% with one \figcap and one \tabcap rendered page 1 as 'רשימת האיורים\n1\n', and the
% .lof held only \@chapter's \addvspace, zero \contentsline entries. An empty list
% under a correct Hebrew heading reads as "this document has no figures".
%
% WHY OPTIONAL, AND NOT AUTOMATIC. Writing #1 unconditionally was implemented first
% and REJECTED: a caption containing \codeline compiles CLEANLY today (verified on
% the pre-v3.64 preamble: 0 errors, PDF produced) but becomes FATAL once it is also
% written to the .lof — \codeline is \hbox to\linewidth{...}, and a full-width box in
% a list entry dies with "\pdfextension endlink cannot be used in vertical mode", no
% output PDF. Gaining a feature nobody uses (0 of 13 collections call \listoffigures)
% at the price of breaking a construct that works is a bad trade. Stock \caption has
% exactly this limitation and exactly this escape hatch, so the convention is familiar.
% \en{}, \code{}, \h{} and inline math in the short text are all safe — verified, they
% survive the write and typeset correctly in the list (RTL geometry confirmed: number
% rightmost at x=509.9, title left of it, dots and folio at the far left).
% With no optional argument the macro is byte-identical to v3.63 — zero risk to every
% existing document, and \listoffigures lists exactly the captions that opted in.
% BLD-37, TABLE side (the \figcap v3.74 idea, adapted): clicking a \ref to a
% table used to land MID-table — hyperref's automatic anchor inside
% \refstepcounter was emitted on the still-open horizontal line whose TAIL is
% the tabular (an hbox whose baseline is the table's vertical CENTER), so the
% jump scrolled the top half of the table above the viewport. The v3.74 entry
% deferred this side claiming "tablebox splits — no single box to measure";
% that deferral dissolves: the measurement target is the TABULAR itself at
% \tabcap ENTRY, while it is still the tail of the open list — the same
% situation \lakin@grabfig handles at an island's end (and the fleet has 0
% split tableboxes). So \lakin@grabtab runs as the FIRST token, before any
% typesetting can push the box away, and the anchor is re-emitted raised above
% the measured tabular. \lakintabraise is a SEPARATE register from
% \lakinfigraise (defined with it, below the LTR-island block): a captionless
% island intentionally leaves the fig register set, and a table must never
% inherit that stale figure height.
%
% WHY the anchor stays AT THE ENTRY SPOT and is NOT deferred into the caption
% paragraph the way \figcap defers it: at \par TeX DROPS a final glue item,
% and the automatic anchor's smash box sitting BEHIND the interword glue kept
% that glue alive — so every centered tabular has always sat half a space off
% true center. Deferring the anchor removed that box and was measured as 982
% differing glyph records: EVERY tabular line shifted 1.30pt (2.61pt for
% pattern 7's two minipages) while captions stayed put. Emitting the raised
% anchor at the old spot — through hyperref's own \Hy@raisedlink, exactly the
% macro the automatic path uses — keeps the line's node structure, and so the
% A/B glyph diff at 0, BY CONSTRUCTION. \Hy@raisedlink reads the raise from
% \HyperRaiseLinkDefault (default \baselineskip = the old mid-table dest); the
% group-local \edef freezes measured-height+42pt into it, and \smash inside
% \Hy@raisedlink hides any raise from the line's extent. \HyCnt@ProvideTheH-
% Counter keeps the dest NAME byte-identical to the automatic one (it defines
% \theHtbl, which \hyper@makecurrent prefers — without it the first measured
% caption was named tbl.1 where the automatic path said tbl.0.1, measured).
% vmode entry, or anything but a box at the tail (no tabular to measure),
% falls back to the plain \refstepcounter — byte-identical to the old path.
\makeatletter
\newcommand{\tabcap}[2][]{\lakin@grabtab
  % v3.79: arm the box-close correction. Set unconditionally and in BOTH
  % branches below's scope, so the plain-\refstepcounter fallback is corrected
  % too; \lakintabboxopen clears it at every \begin{tablebox}, so a \tabcap
  % outside a tablebox can never leave it armed for a later box.
  \global\lakintabcapseentrue
  \ifdim\lakintabraise>0pt
    \H@refstepcounter{tbl}\HyCnt@ProvideTheHCounter{tbl}\hyper@makecurrent{tbl}%
    \begingroup
      % A frameless table targets its measured top. The frame's historical
      % cushion belongs only to tablebox and is corrected at box close.
      \edef\HyperRaiseLinkDefault{\the\lakintabraise}%
      \iflakintabframed
        \edef\HyperRaiseLinkDefault{\the\dimexpr\lakintabraise+42pt\relax}%
      \fi
      \Hy@raisedlink{\hyper@anchorstart{\@currentHref}\hyper@anchorend}%
    \endgroup
    \global\lakintabraise 0pt\relax
  \else
    \refstepcounter{tbl}%
  \fi
  \lakinrtitle{#2}\par\nopagebreak\addvspace{6pt}%
  \def\lakin@short{#1}\ifx\lakin@short\@empty\else
    \addcontentsline{lot}{table}{\protect\numberline{\thetbl}{\protect\ignorespaces #1}}\fi
  {\centering\lakitcaptionstyle\mbox{\lakitcaptionlabelstyle{\tablename~\en{\thetbl}\lakitcaptionpunctuation}}\enspace #2\par}}
\makeatother

% ============================================================================
% Unified FIGURE caption — the \tabcap twin for TikZ/plot figures (v3.34.0).
% A figure (tikzpic) sits inline in the prose; put \figcap BELOW it for a
% numbered, \ref-able caption "איור N.M: תיאור" in the same style as \tabcap:
%   \begin{tikzpic} ... \end{tikzpic}
%   \figcap{\label{fig:foo}description}     % \label FIRST, inside the argument
% WHY a macro, not a hand-written gray line: a manual
%   {\centering\small\textcolor{cGray}{איור: ...}} with a \label inside has NO
% \refstepcounter, so the \label binds to the AMBIENT \@currentlabel (the last
% equation) — \ref then prints an equation number AS the figure number
% (measured on a real page: "איור 48" where 48 is equation 48). \figcap steps a
% real \c@fig counter first, so \label is correct, numbering is consistent, and
% refs are stable. Class-aware like \thetbl (chapter- or section-prefixed).
% ============================================================================
\makeatletter
\@ifundefined{c@chapter}{%
  \newcounter{fig}[section]
  \renewcommand{\thefig}{\ifnum\value{section}>0 \thesection.\fi\arabic{fig}}%
}{%
  \newcounter{fig}[chapter]
  \renewcommand{\thefig}{\ifnum\value{chapter}>0 \thechapter.\fi\arabic{fig}}%
}
\makeatother
\makeatletter
\newcommand{\figcap}[2][]{%
  \ifdim\lakinfigraise>0pt
    % v3.74: step WITHOUT hyperref's automatic anchor, then emit the anchor
    % raised above the measured figure. \H@refstepcounter is hyperref's saved
    % pre-hook \refstepcounter (sets the counter AND \@currentlabel — \label
    % in the caption arg binds exactly as before); \hyper@makecurrent supplies
    % \@currentHref, so the dest NAME is identical to the old automatic one —
    % \ref/\autoref/bookmarks are unchanged, only the dest'S POSITION moves.
    \H@refstepcounter{fig}\hyper@makecurrent{fig}%
    % The anchor is emitted INSIDE the caption paragraph below, as a zero-size
    % raisebox — NOT here in vertical mode, where a \raisebox starts an empty
    % paragraph and adds a blank line (measured: every glyph after the caption
    % shifted +18.8pt). Inside an existing line a [0pt][0pt] box is inert.
    \edef\lakin@anch{\@currentHref}%
    % \edef FREEZES the measured height into the pending macro: the deferred
    % \raisebox is typeset inside the caption paragraph AFTER the register is
    % reset to the 0pt sentinel below — a plain \def read 0 at use time
    % (measured: LAKINRAISE=0.0, dest moved 14pt instead of above the figure).
    \edef\lakin@pend{\noexpand\raisebox{\the\dimexpr\lakinfigraise+42pt\relax}[0pt][0pt]{%
      \noexpand\hyper@anchorstart{\lakin@anch}\noexpand\hyper@anchorend}}%
    \global\lakinfigraise 0pt\relax
  \else
    \refstepcounter{fig}\let\lakin@pend\relax
  \fi
  \lakinrtitle{#2}\par\nopagebreak\addvspace{4pt}%
  \def\lakin@short{#1}\ifx\lakin@short\@empty\else
    \addcontentsline{lof}{figure}{\protect\numberline{\thefig}{\protect\ignorespaces #1}}\fi
  {\centering\lakin@pend\lakitcaptionstyle\mbox{\lakitcaptionlabelstyle{\figurename~\en{\thefig}\lakitcaptionpunctuation}}\enspace #2\par}}
\makeatother

% ============================================================================
% LTR islands.  EVERY tikzpicture / quantikz MUST live inside one, or the whole
% figure's coordinates and Latin labels flip. (Hebrew node text still uses \h{}.)
% ============================================================================
% v3.74 (BLD-37): the \figcap hyperlink DESTINATION used to land at the caption
% line — BELOW the figure — so clicking a figure reference scrolled the figure
% off-screen (measured fleet-wide: 97/97 fig dests in one collection sat 19-54pt
% under their figure's ink; a 347pt figure was fully hidden after the jump).
% hypcap-style fix in two parts: (1) at the island's END, \lakin@grabfig grabs
% the just-typeset centered line with \lastbox, measures it, and puts it back
% with \nointerlineskip — the interline glue above the line was already
% contributed and is NOT re-added, so the page geometry is unchanged to the
% glyph (verified); (2) \figcap then steps the counter WITHOUT hyperref's
% automatic anchor (\H@refstepcounter = the pre-hyperref \refstepcounter) and
% emits the anchor itself, raised by the measured height — the dest lands ABOVE
% the figure top. A captionless island leaves numbering untouched (nothing is
% stepped at the island; the grab flag is cleared at every island begin, so a
% stale height from a captionless figure cannot leak into a later caption).
\newsavebox{\lakinfigbox}
\newdimen\lakinfigraise  % real \dimen: tex.setdimen cannot target a \newlength SKIP register (measured lua error)
% BLD-37 table side: a SEPARATE register for \tabcap. The two must not share —
% a captionless island INTENTIONALLY leaves \lakinfigraise set (the grab runs
% at every island END; only the NEXT island's begin clears it), so a shared
% register would hand a \tabcap after a captionless figure that stale FIGURE
% height. \lakintabraise is written only by \lakin@grabtab at \tabcap entry
% and consumed/reset inside \tabcap itself, so it can never go stale.
\newdimen\lakintabraise
\makeatletter
% v3.74b: measure WITHOUT touching the list. The first draft grabbed the box
% with \lastbox and put it back — and \lastbox DISCARDS the node's baseline
% shift, so a vertically-anchored quantikz circuit rose 2.57pt and everything
% after it sank 5.32pt (measured on the A/B glyph diff). LuaTeX lets us READ
% the tail node's dimensions without removing anything: geometry is identical
% by construction. \lakinfigraise=0pt is the "no figure measured" sentinel.
\def\lakin@grabfig{%
  \ifhmode
    \directlua{
      local n = tex.nest[tex.nest.ptr].tail
      local GLUE, PEN = node.id('glue'), node.id('penalty')
      while n and (n.id == GLUE or n.id == PEN) do n = n.prev end
      local HL, VL = node.id('hlist'), node.id('vlist')
      if n and (n.id == HL or n.id == VL) then
        tex.setdimen('global', 'lakinfigraise',
                     (n.height or 0) + (n.depth or 0))
      else
        tex.setdimen('global', 'lakinfigraise', 0)
      end}%
  \else
    \global\lakinfigraise 0pt\relax
  \fi}
% BLD-37 table side: \lakin@grabfig's twin, NOT a shared helper — the Lua body
% differs in exactly one place, the skip set. Between \end{tabular} and \tabcap
% an otherlanguage wrapper (patterns 3/5 of example_tables) leaves a \write
% whatsit and a dir node ON TOP of the interword glue (measured node dump:
% "whatsit/write <- dir <- glue(2.62) <- hlist"), so the table reader must
% also step over whatsit/dir markers — all zero-size — to reach the tabular.
% \lakin@grabfig deliberately does NOT get the wider skip set: its glue/penalty
% body is the one the v3.74 A/B proof froze, and widening it would newly
% measure (and so newly raise) anchors in island shapes that today take the
% plain path — an unproven behavior change on the FIG side. Anything that is
% not a box behind those markers still hits the 0pt sentinel and \tabcap falls
% back to the plain \refstepcounter.
% SECOND deliberate difference: this stores the box's HEIGHT, not height+depth.
% \figcap's anchor is emitted in the caption paragraph BELOW the figure, so it
% must climb the full extent; \tabcap's anchor sits ON the tabular's own line,
% whose baseline is the vcentered tabular's MIDPOINT — height alone reaches
% the top edge, and adding the depth was measured to overshoot by half the
% table (dests 80-135pt above their frame, INSIDE the preceding table's box).
\def\lakin@grabtab{%
  \ifhmode
    \directlua{
      local n = tex.nest[tex.nest.ptr].tail
      local GLUE, PEN = node.id('glue'), node.id('penalty')
      local WHAT, DIR = node.id('whatsit'), node.id('dir')
      while n and (n.id == GLUE or n.id == PEN
                   or n.id == WHAT or n.id == DIR) do n = n.prev end
      local HL, VL = node.id('hlist'), node.id('vlist')
      if n and (n.id == HL or n.id == VL) then
        tex.setdimen('global', 'lakintabraise', n.height or 0)
      else
        tex.setdimen('global', 'lakintabraise', 0)
      end}%
  \else
    \global\lakintabraise 0pt\relax
  \fi}
\makeatother
% ============================================================================
% v3.79 (BLD-37 table side, CORRECTION): put the anchor at the TOP OF THE BOX,
% instead of guessing the distance to it.
% ============================================================================
% v3.77 raised the \tabcap anchor by (height of the box at the tail of the open
% horizontal list) + 42pt. That is a MEASUREMENT taken BEFORE line breaking, and
% it is short by every term it cannot see. Fleet-measured on the 12 built
% collections (299 tbl.* dests, ownership established through each table's own
% CAPTION, never through "which frame encloses the dest"): 3 dests landed INSIDE
% their own tablebox, below the frame top —
%   computer-architecture       tbl.14.1  -46.4pt   three stacked tabulars
%   computer-architecture       tbl.30.5  -13.5pt   tabular 2.85pt over \linewidth
%   advanced-electronic-devices tbl.B.1   -59.2pt   tabular 2.56pt over \linewidth
% Two missing terms, both reproduced in a minimal document (probe P2/P3):
%   (1) VERTICAL LIST ALREADY CONTRIBUTED inside the box — a second and third
%       tabular, or a heading above the table. The tail read sees only the LAST
%       box, so 77.7pt of stacked tabulars is invisible to the raise.
%   (2) WHICH LINE THE ANCHOR LANDS ON. When the tabular alone overfills
%       \linewidth (459.6121pt in a tablebox), the interword glue after
%       \end{tabular} becomes the chosen breakpoint and \Hy@raisedlink is
%       emitted on the NEXT line — one (depth + lineskip) lower. That term does
%       not EXIST at \tabcap time: line breaking has not happened yet. No
%       measurement taken there can ever recover it. (Proof: probe P3, the same
%       table 2.84912pt too wide, fails at -13.5pt; probe P4, the identical
%       table with \tabcolsep 2pt so it fits, is correct at +32.0pt.)
% So the raise is not computed at \tabcap at all. It is CORRECTED at
% \end{tablebox}, from inside mdframed's own body vbox (endinnercode runs after
% \par\unskip, i.e. after line breaking, while tex.nest[tex.nest.ptr] is still
% that vbox's vertical list). There the anchor's true offset from the top of the
% box content is not estimated but WALKED, and the \Hy@raisedlink shift is set
% so the dest lands 42pt above it — which is 42pt - innertopmargin(9pt) -
% linewidth(0.9pt) = 32.0pt above the FRAME top, the exact position the 295
% already-correct dests have today. That is why this is a correction and not a
% re-placement: the good sites do not move.
%
% WHY THE GLYPH DIFF IS 0 BY CONSTRUCTION — the only thing written is the
% `shift` field of the hlist that \Hy@raisedlink already built, and the guard
% below refuses to write it unless that box is 0x0x0. A zero-size box carries no
% dimension into its parent's packing, the packing already happened, and shift
% is not an input to it. Nothing else in the list is touched: no node is added,
% removed, resized or reordered. Measured, not assumed: assets/example_tables.tex
% and the whole computer-architecture collection A/B at 0 differing pymupdf
% rawdict glyph records.
%
% WHY NOT the alternatives (all prototyped and measured before this one):
%   * accumulate the vertical list at \tabcap time — fixes term (1), CANNOT fix
%     term (2); 2 of the 3 real failures survive. A fix that leaves the majority
%     of the defect standing is the false all-clear this repo exists to avoid.
%   * \zref-savepos two-pass — correct on all shapes, but it makes the dest
%     POSITION aux-derived (DEFECTS BLD-1/2), a single-pass compile silently
%     ships an unraised anchor, and the bare y-difference has no page guard.
%   * emit the anchor at \begin{tablebox} and let \tabcap adopt it — needs the
%     tbl counter stepped at box OPEN to keep the dest NAME byte-identical; a
%     captionless tablebox then either shifts every later table number or leaves
%     an orphan dest that the next table duplicates (a HARD build gate). This
%     correction touches neither the counter nor the name: \tabcap's naming path
%     is v3.77's, unchanged, so resolve_names() is byte-identical (verified).
%   * suppressing the line break (a penalty before the glue) — that MOVES the
%     tabular: the break currently discards the interword glue, so forbidding it
%     re-centers the two overfull tables. Reflow, rejected.
%
% SCOPE. The correction is armed only by \tabcap ITSELF (\lakintabcapseen),
% cleared at every box open, so a captionless tablebox is a no-op and a \tabcap
% hosted OUTSIDE a tablebox (the fleet has exactly one, computer-architecture
% tbl.14.2 in an exbox) keeps v3.77's measured raise. Both fallbacks are
% therefore never worse than today; and because hyperref's AUTOMATIC anchor uses
% the same \Hy@raisedlink shape, the plain-\refstepcounter path (vmode entry, or
% prose between \end{tabular} and \tabcap — probe P6, -51.5pt) is corrected too.
% KNOWN LIMIT — a tablebox that SPLITS across pages. The offset is measured from
% the top of the UNSPLIT body, so a dest in segment 2 gets raised past its own
% segment. MEASURED on a forced split (probe: 40 filler lines + a table): the
% dest went from y=585.0 on page 2 (inside segment 2, at the caption — the old
% BLD-37 symptom) to y=-113.4, i.e. ABOVE the page top. Viewers clamp that to the
% page top, which happens to show segment 2 whole, but an out-of-page coordinate
% is not something to ship on purpose. Hence the \textheight guard in
% lakintab.fix: a body TALLER THAN THE TEXT BLOCK cannot fit on any page, so it
% must split — the correction is skipped and v3.77's behaviour stands unchanged
% (verified: the probe returns to y=585.0). The guard is deliberately EXACT, not
% a page-fitting heuristic: it is sound (a taller-than-page body always splits)
% but PARTIAL (a shorter body that starts low on a page could still split). That
% residue is unreachable in practice — v3.37's everyline + the caption's
% \nopagebreak glue move such boxes WHOLE instead (probe N7: a 22-row table
% pushed down by 430pt moved whole, 0 segments), and the fleet has 0 split
% tableboxes, re-counted after this change. Reimplementing mdframed's own
% page-fitting arithmetic to close the residue would be a heuristic with its own
% failure mode, which is the trade this repo has repeatedly refused.
\newif\iflakintabcapseen
\newif\iflakintabframed
% The Lua lives in ONE \directlua at load time, and is written with NO
% backslash, %, # or ~ : \directlua expands its argument like \edef, so an
% active ~ would expand to \nobreakspace and a % would eat the rest of the line
% (the second is how the first draft of this block died). "not (a == b)" stands
% in for Lua's inequality operator for the same reason.
\directlua{
  lakintab = { }
  lakintab.HL = node.id("hlist")
  lakintab.VL = node.id("vlist")
  lakintab.RU = node.id("rule")
  lakintab.GL = node.id("glue")
  lakintab.KE = node.id("kern")
  lakintab.WH = node.id("whatsit")
  lakintab.TM = node.id("temp")
  lakintab.WD = node.subtype("pdf_dest")
  function lakintab.d(v)
    if v == nil then return 0 end
    if v < -1000000000 then return 0 end
    return v
  end
  function lakintab.scan(list, y, horiz, parent, target)
    local x = list
    while x do
      local i = x.id
      if i == lakintab.HL or i == lakintab.VL then
        if horiz then
          local b = y + x.shift
          if i == lakintab.HL then
            lakintab.scan(x.head, b, true, x, target)
          else
            lakintab.scan(x.head, b - lakintab.d(x.height), false, x, target)
          end
        else
          if i == lakintab.HL then
            lakintab.scan(x.head, y + lakintab.d(x.height), true, x, target)
          else
            lakintab.scan(x.head, y, false, x, target)
          end
          y = y + lakintab.d(x.height) + lakintab.d(x.depth)
        end
      elseif i == lakintab.WH then
        if x.subtype == lakintab.WD then
          local id = x.dest_id
          if type(id) == "string" and string.sub(id, 1, 4) == "tbl." then
            if horiz and parent and parent.width == 0
               and parent.height == 0 and parent.depth == 0 then
              parent.shift = parent.shift + (target - y)
            end
          end
        end
      elseif not horiz then
        if i == lakintab.RU then
          y = y + lakintab.d(x.height) + lakintab.d(x.depth)
        elseif i == lakintab.GL or i == lakintab.KE then
          y = y + lakintab.d(x.width)
        end
      end
      x = x.next
    end
  end
  function lakintab.fix(target, maxh)
    local nst = tex.nest[tex.nest.ptr]
    if nst == nil then return end
    local h = nst.head
    if h == nil then return end
    if h.id == lakintab.TM then h = h.next end
    local tot = 0
    local x = h
    while x do
      local i = x.id
      if i == lakintab.HL or i == lakintab.VL or i == lakintab.RU then
        tot = tot + lakintab.d(x.height) + lakintab.d(x.depth)
      elseif i == lakintab.GL or i == lakintab.KE then
        tot = tot + lakintab.d(x.width)
      end
      x = x.next
    end
    if tot > maxh then return end
    lakintab.scan(h, 0, false, nil, target)
  end
}
% @-FREE NAMES ON PURPOSE: these two are named in the `tablebox` key list at
% \newmdenv above, which is tokenised OUTSIDE any \makeatletter region — a
% \lakin@... there would tokenise as \lakin followed by the letters @tabopen.
\newcommand{\lakintabboxopen}{\lakintabframedtrue\global\lakintabcapseenfalse}
\newcommand{\lakintabboxclose}{%
  \iflakintabcapseen
    \directlua{lakintab.fix(-\number\dimexpr\lakintabcushion\relax,
                            \number\dimexpr\textheight\relax)}%
    \global\lakintabcapseenfalse
  \fi\lakintabframedfalse}
% The cushion above the box's CONTENT top. 42pt puts the dest 32.0pt above the
% frame top (42pt - innertopmargin 9pt - linewidth 0.9pt), i.e. exactly where
% the 295 correct dests already are — the correction moves only the broken ones.
\newcommand{\lakintabcushion}{42pt}
\newenvironment{qcirc}{\par\medskip\global\lakinfigraise 0pt\relax\begin{otherlanguage}{english}\begin{center}}
  {\csname lakin@grabfig\endcsname\end{center}\end{otherlanguage}\par\medskip}
\newenvironment{tikzpic}{\par\medskip\global\lakinfigraise 0pt\relax\begin{otherlanguage}{english}\begin{center}}
  {\csname lakin@grabfig\endcsname\end{center}\end{otherlanguage}\par\medskip}

% ============================================================================
% Code listings.  ★★ v3.4 THE BOX-OVERFLOW FIX:
%   resetmargins=true   → the listing IGNORES the enclosing environment's
%                         indentation and measures from the PAGE. v3.4 stated
%                         this backwards ("measured from the ENCLOSING box"),
%                         and v3.37 then blanked in-box line numbers because
%                         the page-anchored listing put its hanging digits on
%                         top of the box rule. Measured v4.5.0: with the key
%                         FALSE the listing follows the box and the code moves
%                         from x=75.26 to x=97.74, clearing the rule at 72.07.
%                         The key now lives in one \lstset below instead of in
%                         each style, so a box can override it — codebox does.
%   framexleftmargin=0  → the frame (left bar) does not extend left of the code.
%   xleftmargin keeps a small indent so the line numbers have room INSIDE.
% asmblock wraps the listing in an LTR island (code is LTR).
% ============================================================================
\lstdefinestyle{asm}{
  basicstyle=\ttfamily\small,
  backgroundcolor=\color{cBlueBg!55},
  frame=leftline,framerule=2.5pt,rulecolor=\color{cBlue},
  numbers=left,numberstyle=\tiny\color{cGray},numbersep=8pt,
  xleftmargin=2.2em,framexleftmargin=0pt,framexrightmargin=0pt,
  breaklines=true,columns=fullflexible,keepspaces=true,
  showstringspaces=false,aboveskip=8pt,belowskip=8pt,
  % v3.36: without a comment DELIMITER the commentstyle below never fires —
  % listings has no `language` here, so it does not know '#' opens a comment,
  % and every asm listing rendered flat black (measured: a whole-page pixel
  % sweep found ZERO cGreen glyphs across 135 shipped listings). morecomment
  % gives it the delimiter; color-only change, glyph widths untouched.
  morecomment={[l]{\#}},
  commentstyle=\color{cGreen},
}
\lstdefinestyle{cstyle}{
  language=C,basicstyle=\ttfamily\small,
  backgroundcolor=\color{cGreenBg!55},
  frame=leftline,framerule=2.5pt,rulecolor=\color{cGreen},
  numbers=left,numberstyle=\tiny\color{cGray},numbersep=8pt,
  xleftmargin=2.2em,framexleftmargin=0pt,framexrightmargin=0pt,
  breaklines=true,columns=fullflexible,keepspaces=true,
  showstringspaces=false,keywordstyle=\color{cBlue}\bfseries,
  commentstyle=\color{cGray}\itshape,aboveskip=8pt,belowskip=8pt,
}
% v3.36: pystyle/htmlstyle — colored twins of cstyle for the Python/HTML
% listings that used to repeat a ~200-char options blob with NO color
% (keywords rendered bold-black, comments italic-black — measured; C alone
% was colored). Geometry is copied VERBATIM from the shipped blob (frame,
% margins, skips) so swapping a call site to style=pystyle moves NOTHING —
% only keywordstyle/commentstyle/stringstyle are added (bold stays bold,
% italic stays italic, mono metrics identical => zero reflow, proven by
% per-page word-bbox diff). breaklines/breakatwhitespace stay per-site.
% v3.78: numbers=left joins both styles, with the exact numberstyle/numbersep
% keys asm/cstyle carry — fleet-wide consistency (author decision: all four
% styles numbered; "asm/cstyle numbered, py/html not" was an accident of the
% v3.36 geometry freeze, not a design). The numbers hang LEFT of the code box
% (listings sets them with \llap, so no overfull hbox can arise and the code
% x is untouched — measured: code frozen, digits at code_x-11.8pt); inside a
% colored box they are blanked like everyone else's (v3.37 hook), and the
% clean variant restores them in its frameless boxes (v3.78, above).
\lstdefinestyle{pystyle}{
  language=Python,basicstyle=\ttfamily\small,
  columns=fullflexible,keepspaces=true,
  frame=single,rulecolor=\color{cTblLine},
  numbers=left,numberstyle=\tiny\color{cGray},numbersep=8pt,
  xleftmargin=0.5em,framexleftmargin=0pt,framexrightmargin=0pt,
  showstringspaces=false,aboveskip=8pt,belowskip=8pt,
  keywordstyle=\color{cBlue}\bfseries,
  commentstyle=\color{cGray}\itshape,
  stringstyle=\color{cGreen},
}
% v4.5.0: hoisted out of the four styles so an enclosing box can override it.
% Every listing still resolves to true unless something later says otherwise,
% so shipped output is unchanged (verified byte-for-byte on all eight PDFs).
\lstset{resetmargins=true}
\lstdefinestyle{htmlstyle}{
  language=HTML,basicstyle=\ttfamily\small,
  columns=fullflexible,keepspaces=true,
  frame=single,rulecolor=\color{cTblLine},
  numbers=left,numberstyle=\tiny\color{cGray},numbersep=8pt,
  xleftmargin=0.5em,framexleftmargin=0pt,framexrightmargin=0pt,
  showstringspaces=false,aboveskip=8pt,belowskip=8pt,
  keywordstyle=\color{cBlue}\bfseries,
  commentstyle=\color{cGray}\itshape,
  stringstyle=\color{cGreen},
}
\newenvironment{asmblock}{\par\begin{otherlanguage}{english}}{\end{otherlanguage}\par}

% <hllpdf:listing-smart-quotes>
% v4.21.0: four Unicode punctuation characters must enter listings' token
% buffer. Unregistered characters bypass it: a closing quote was printed
% BEFORE the pending word, and an opening quote BEFORE a pending brace.
% Extend the standard character table, then activate only these codepoints
% in the listing's local character-table scope. Do this BEFORE the package's
% delimiter/alsoother/literate hooks so explicit user mappings still win.
% This is not universal Unicode support, nor a conversion to ASCII quotes.
% See listings developer guide 15.3.1 and references/listing-smart-quotes.md.
\makeatletter
\lst@CCPut \lst@ProcessOther {"2018}{"2019}{"201C}{"201D}\z@
\lst@AddToHookAtTop{SelectCharTable}{%
  \lst@ifactivechars
    \catcode"2018=\active \catcode"2019=\active
    \catcode"201C=\active \catcode"201D=\active
  \fi}
\makeatother
% </hllpdf:listing-smart-quotes>

% <hllpdf:listing-technical-symbols>
% v4.30: bounded extension for common technical punctuation in LTR listings.
% Without character-table registration these glyphs bypass the token buffer
% and can move before the preceding word even inside quoted strings.
% Keep the original four-quote block separate for its standing mutation test.
% This is source-faithful printing, not syntax validation or general Unicode.
\makeatletter
\lst@CCPut \lst@ProcessOther
  {"2013}{"2014}{"2026}{"2190}{"2192}{"2194}{"21D2}
  {"2212}{"221E}{"2260}{"2264}{"2265}\z@
\lst@AddToHookAtTop{SelectCharTable}{%
  \lst@ifactivechars
    \catcode"2013=\active \catcode"2014=\active \catcode"2026=\active
    \catcode"2190=\active \catcode"2192=\active \catcode"2194=\active
    \catcode"21D2=\active \catcode"2212=\active \catcode"221E=\active
    \catcode"2260=\active \catcode"2264=\active \catcode"2265=\active
  \fi}
\makeatother
% </hllpdf:listing-technical-symbols>


\usepackage[hidelinks,unicode]{hyperref}   % LAST, after colors. unicode = Hebrew bookmarks.
% v3.62: THE ISOLATION MACROS CORRUPTED EVERY PDF BOOKMARK THEY APPEARED IN.
% \en is defined (line ~155) as
%     \foreignlanguage{english}{\language=\csname l@nohyphenation\endcsname #1}
% and in a heading that body is written verbatim to the .toc and handed to
% hyperref's \pdfstringdef. hyperref strips the two CONTROL SEQUENCES (\language,
% \l@nohyphenation) but the plain `=` CHARACTER token between them is not a control
% sequence, so it survives into the PDF string: \chapter{שער \en{CNOT} והצמדה}
% produced the bookmark "שער =CNOT והצמדה". The printed TOC and the heading on the
% page are CORRECT — only the navigation pane and the viewer window title are wrong,
% which is why it went unnoticed. Measured on the shipped fleet: 70 corrupted
% bookmarks across 5 PDFs, from 177 unprotected \en-in-heading sites. Same corruption
% reached /Title and /Subject via the \hypersetup{pdftitle=...} the docs recommend.
% This hook tells hyperref to use the plain argument for all three isolation macros
% when building a PDF string. It does NOT touch the .toc (the printed TOC is
% byte-unchanged) and it does NOT mask the math signal — a heading with $...$ still
% logs "removing math shift", so BIDI-15 detection is intact.
% BONUS, measured: it also makes \code{} in an UNSTARRED heading COMPILE. BLD-10
% recorded that as needing \texorpdfstring or a starred heading; the real root cause
% is hyperref's \pdfstringdef scan of \code's \ifmmode...\else...\fi body — which is
% exactly why the identically-bodied \texttt survived (hyperref pre-neutralizes
% \texttt in PDF strings) and why \DeclareRobustCommand could never help.
% v3.64: \def\\{ } added. A manual line break in a heading WELDS the two words
% together in the bookmark: \chapter{שורה ראשונה\\שורה שנייה} produced the outline
% entry 'שורה ראשונהשורה שנייה' — "ראשונהשורה" is not a word. hyperref drops the
% \\ and puts nothing in its place, so the break that looks like a space on the
% page is a zero-width join in the navigation pane. The spaced form
% 'מילה אחת \\ מילה שנייה' escapes with a harmless DOUBLE space ('אחת  מילה'),
% which is why this never got noticed. Both are measured on a real outline.
% 0 live instances in the corpus — but SKILL.md and content-style.md both RECOMMEND
% a manual heading break, so this is a primed trap rather than a dormant one, and a
% one-token preamble fix retires it for good instead of asking authors to remember.
% Mapping \\ to a space fixes the WELD, which is the actual defect (a nonexistent
% word in the navigation pane). It does not normalise whitespace: measured after the
% fix, the welded form becomes 'שורה ראשונה שורה שנייה' (correct) and the already-
% spaced form goes from two spaces to three. That is cosmetic — viewers render it as
% a wider gap, and no word is corrupted. A per-heading \texorpdfstring is still the
% way to get an exactly-tuned bookmark string.
\pdfstringdefDisableCommands{\def\en#1{#1}\def\h#1{#1}\def\doclang#1{#1}\def\code#1{#1}\def\codeline#1{#1}\def\\{ }}
% v3.70: RESTORE \url LINE-BREAKING. url.sty freezes its break penalties at LOAD
% time (\mathchardef\UrlBreakPenalty=\binoppenalty), and v3.45's global
% \binoppenalty=10000 (inline-math atomicity) therefore priced EVERY breakpoint
% inside \Url at 10000 = never break. Measured: a long \url in Hebrew prose is
% one unbroken line, ~100pt overfull, invading the margin band the ink gates
% deliberately tolerate. With the two lines below: 0 overfull, the URL breaks
% after '/', the RTL split is CORRECT (fragment order forward), and each line
% gets its own link annotation exactly covering its fragment (measured).
% Confined to \Url's group — inline-math atomicity (TYP-7) is untouched.
\mathchardef\UrlBreakPenalty=200
\mathchardef\UrlBigBreakPenalty=100
\ifx\lakitmainlanguage\hllpdfenglish
  \hypersetup{pdflang=en}
\else
  \hypersetup{pdflang=he}
\fi  % /Lang is the document base locale (not per-span accessibility tagging).
% Set a per-document \hypersetup{pdftitle={...}} too — see SKILL.md.

% v3.52 (round-2 audit): two RTL bibliography/cross-ref defects, both invisible to
% every gate and both affecting the FIRST author to use the feature.
% (a) \@biblabel's default [#1] brackets are bidi-neutral and mirror to ]1[ under
% the RTL bibliography base — every \thebibliography reference list shows ]1[.
% LTR-isolating the label restores [1]. (The render gate's mirror rule fingerprints
% round parens )1(, so square brackets slipped through.)
\makeatletter
\renewcommand*{\@biblabel}[1]{\foreignlanguage{english}{[#1]}}
% v3.54: the SAME neutral-bracket trap on the INLINE side — a multi-key \cite{a,b}
% under the RTL base lays its European-number list right-to-left, so [1, 2] renders
% "[2,1]" and [1, 2, 3] renders "[3,2,1]" (verified by glyph geometry — the reader
% sees the sources in REVERSE order; a note form \cite[p.5]{a} also reorders to
% "[p. 5,1]"). LTR-isolate the whole bracket group (number list AND optional note),
% exactly like \@biblabel above: [1, 2] / [1, 2, 3] / [1, p. 5] all render correct.
% Placed AFTER hyperref so it wins; hyperref hyperlinks each number inside #1 (via
% \@citex), which this wrapper leaves intact. Mirrors the LaTeX kernel \@cite shape.
% v4.9.1: directional isolation must NOT switch the language here. Babel's
% safe-active citation wrapper passes its key list through \bbl@tempa;
% \foreignlanguage reuses that scratch macro before the deferred citation
% loop expands it. The old line therefore cited the literal key "english".
% If such a key existed, even the undefined-reference gate stayed green.
% \babelsublr isolates the bracket group without selecting a locale or
% changing that key list. Preserve native links, key order and optional notes.
\def\@cite#1#2{\babelsublr{[{#1\if@tempswa , #2\fi}]}}
\makeatother

% v4.11.0: a reference inside align's \text inherited HEBdisptag=true,
% so its enclosing hbox mirrored the brackets despite a clean compile.
% Isolate the complete printed reference without selecting a Babel language.
% Also neutralize the display-only normalfont direction reset locally. A
% direction-only patch inverted 12 to 21 inside equation in the measured probe.
% Plain maketag is local to the reference: mathtools' display-side wrapper
% otherwise globally consumes manual_tag while formatting a nested reference.
% Preserve each original command (including the selective auxiliary write),
% current tag forms, native links, and showonlyrefs on/off transitions.
\makeatletter
\NewCommandCopy{\HEB@nativeeqref}{\eqref}
\newcommand{\HEB@eqrefisolate}[1]{%
  \babelsublr{\HEBdisptagfalse\let\bbl@eqnodir\relax\let\maketag@@@\ORIGmaketag@@@#1}}
\protected\def\eqref#1{\HEB@eqrefisolate{\HEB@nativeeqref{#1}}}
\ExplSyntaxOn
\NewCommandCopy{\HEB@nativeMTeqref}{\MT_extended_eqref:n}
\cs_set_protected:Npn \MT_extended_eqref:n #1
  {\HEB@eqrefisolate{\HEB@nativeMTeqref{#1}}}
\ExplSyntaxOff
\makeatother

% v4.11.0: custom caption counters are not part of hyperref's English
% locale, and Hebrew has no upstream hyperref locale in this toolchain.
% Restore the appropriate names at every Babel language switch, including
% return from a nested inline island. Do not localize the destination keys.
\makeatletter
\newcommand{\HEB@refnames}{%
\def\chapterautorefname{פרק}%
\def\sectionautorefname{סעיף}%
\def\subsectionautorefname{תת־סעיף}%
\def\subsubsectionautorefname{תת־תת־סעיף}%
\def\paragraphautorefname{פסקה}%
% v4.31: hyperref uses a distinct counter name at the deepest standard level.
\def\subparagraphautorefname{תת־פסקה}%
\def\equationautorefname{משוואה}%
\def\figureautorefname{איור}%
\def\tableautorefname{טבלה}%
\def\partautorefname{חלק}%
\def\appendixautorefname{נספח}%
\def\footnoteautorefname{הערת שוליים}%
\def\itemautorefname{סעיף}%
\def\Itemautorefname{סעיף}%
\def\figautorefname{איור}%
\def\tblautorefname{טבלה}%
\def\pageautorefname{עמוד}% 
}
\addto\extrashebrew{\HEB@refnames}
\newcommand{\ENG@refnames}{%
  \def\subparagraphautorefname{subparagraph}%
  \def\chapterautorefname{chapter}%
  \def\sectionautorefname{section}%
  \def\subsectionautorefname{subsection}%
  \def\subsubsectionautorefname{subsubsection}%
  \def\paragraphautorefname{paragraph}%
  \def\equationautorefname{Equation}%
  \def\figureautorefname{Figure}%
  \def\tableautorefname{Table}%
  \def\partautorefname{Part}%
  \def\appendixautorefname{Appendix}%
  \def\footnoteautorefname{footnote}%
  \def\itemautorefname{item}%
  \def\figautorefname{Figure}%
  \def\tblautorefname{Table}%
  \def\Itemautorefname{item}%
  \def\pageautorefname{page}%
}
\addto\extrasenglish{\ENG@refnames}
\ifx\lakitmainlanguage\hllpdfenglish\ENG@refnames\else\HEB@refnames\fi
\makeatother

% ============================================================================
% Optional running head (v3.39): \lakithead{<book title>} — OFF by default.
% A document opts in with ONE line after \input{preamble.tex}:
%   \lakithead{מבנה המחשב}
% Every content page then carries the series' orientation line: chapter name
% at the reading edge, book title opposite, thin rule. (fancyhdr's [L]/[R]
% are PHYSICAL; under the RTL page direction [L] lands at the RIGHT/reading
% edge — measured, consistent across the corpus.) Packaged like the topic
% rail: always loaded, inert without the call. This centralizes the 8-line
% block that 8 lakits used to copy by hand — the same duplication vector
% that once let a lakit preamble go ~100 lines stale. Class-aware: report
% fills the mark from \chapter, article from \section. NB: a titlepage env
% self-empties; front matter that sets \pagestyle{empty} must restore with
% \pagestyle{fancy} afterwards (NOT plain — that discards the head).
% ============================================================================
\usepackage{fancyhdr}
% v3.61: \chapter* left a STALE running head — pages of an UNNUMBERED chapter
% carried the PREVIOUS chapter's title. report.cls routes \chapter* through
% \@schapter, which never calls \chaptermark, so \leftmark keeps whatever the last
% numbered chapter left. Measured: with \chapter{פרק ממוספר} then
% \chapter*{מבוא לא ממוספר}, page 4 (inside the starred chapter) printed
% "פרק ממוספר". Invisible on the chapter's FIRST page (\thispagestyle{plain}), so it
% only shows from the second page on — the pages an author is least likely to check.
% An unnumbered front-matter chapter (\chapter*{מבוא}, \chapter*{הקדמה}) is everyday
% Hebrew-academic structure, and the corpus already carries HAND-WRITTEN \markboth
% lines right after \chapter* — the exact per-document duplication \lakithead exists
% to retire. Patched INSIDE \lakithead so the macro stays inert when unused; a
% document's own later \markboth still wins, so those hand-patches do not regress.
\makeatletter
% <hllpdf:starred-heading-scope>
\newif\iflakin@starmarkpatched
\newcommand{\lakistarmark}{%
  \iflakin@starmarkpatched\else
  % Dispatch on the actual class counter, not hyperref's shared internals.
  \ifdefined\c@chapter
    \let\HEBorigschapter\@schapter
    % Record only after the native heading has reached its page. A mark
    % before the native routine can advertise the next heading too early.
    \def\@schapter##1{\HEBorigschapter{##1}\chaptermark{##1}}%
  \else
    % @ssect receives layout arguments, NOT the section name or level.
    % Capture the name in @startsection without introducing a TeX group:
    % grouping sectioning would break the native after-heading state.
    \def\HEB@sectionkind{}%
    \def\HEB@mainsection{section}%
    \let\HEBorigstartsection\@startsection
    \def\@startsection##1##2##3##4##5##6{%
      \def\HEB@sectionkind{##1}%
      \HEBorigstartsection{##1}{##2}{##3}{##4}{##5}{##6}}%
    \let\HEBorigssect\@ssect
    \def\@ssect##1##2##3##4##5{%
      \HEBorigssect{##1}{##2}{##3}{##4}{##5}%
      % Lower starred headings must not replace the article's section mark.
      \ifx\HEB@sectionkind\HEB@mainsection
        \sectionmark{##5}%
      \fi
      \ignorespaces}%
  \fi
  \global\lakin@starmarkpatchedtrue
  \fi}
% </hllpdf:starred-heading-scope>
% <hllpdf:running-language>
% Choose the wrapper when the mark is written, not when a later page ships.
% English punctuation must remain LTR even after returning to Hebrew prose.
\newcommand{\HEB@runningmark}[1]{%
  \iflanguage{english}{\markboth{\en{#1}}{}}{\markboth{\h{#1}}{}}}
% </hllpdf:running-language>
\makeatother
% v4.2.1 (NAV): every collection already names itself in \lakithead — feed the
% PDF /Title from the same declaration, so the navigation gate's
% PDF_REQUIRE_TITLE has a universal source and no fleet member ships
% title-less by missing a second declaration site. Measured before this: 14
% of 16 collections shipped /Title-empty while every one carried a perfectly
% good \lakithead string (plain Hebrew in all 16 — no math, no formatting, so
% the PDF-string path is clean). An explicit title always wins: set it with
% \lakitpdftitle{...} anywhere before \lakithead, or with a plain
% \hypersetup{pdftitle=...} — the fallback detects that via hyperref's
% \@pdftitle, probe-verified \@empty-equal until the first explicit set.
\makeatletter
\newif\iflakit@pdftitleset
\newcommand{\lakitpdftitle}[1]{%
  \hypersetup{pdftitle={#1}}\global\lakit@pdftitlesettrue}
\newcommand{\lakithead}[1]{%
  \iflakit@pdftitleset\else
    \ifx\@pdftitle\@empty
      \hypersetup{pdftitle={#1}}\global\lakit@pdftitlesettrue
    \fi
  \fi
  \pagestyle{fancy}\fancyhf{}%
  \lakistarmark

  \fancyhead[L]{\small\textcolor{cGray}{\leftmark}}%
  % <hllpdf:fixed-head-language>
  % The fixed document title has its OWN language snapshot, independent of
  % the current section/chapter mark. Choose the branch NOW, not at shipout:
  % the page may be Hebrew by the time an English title is finally printed.
  \iflanguage{english}{%
    \fancyhead[R]{\small\textcolor{cGray}{\en{#1}}}%
  }{%
    \fancyhead[R]{\small\textcolor{cGray}{\h{#1}}}%
  }%
  % </hllpdf:fixed-head-language>
  \fancyfoot[C]{\thepage}%
  \renewcommand{\headrulewidth}{0.4pt}%
  \ifdefined\chaptermark
    \renewcommand{\chaptermark}[1]{\HEB@runningmark{##1}}%
  \else
    \renewcommand{\sectionmark}[1]{\HEB@runningmark{##1}}%
  \fi}
\makeatother

% v3.39.1: path-independent byte-reproducibility. LuaTeX derives the trailer
% /ID from the OUTPUT PATH, so two clones of the same source at different
% directories produced different sha256 (measured: the ONLY delta was the
% 62-byte /ID line). Bit 512 suppresses exactly that line — measured: Creator/
% Producer/CreationDate/ModDate and all content untouched. /ID is optional
% per ISO 32000 (required only for ENCRYPTED PDFs and PDF/A — neither is a
% target here). If a future workflow needs an ID (PDF/A, Acrobat revision
% matching), delete this line and it returns.
\pdfvariable suppressoptionalinfo \numexpr 512\relax

% v3.75: BOLD MATH IN BOLD HEADINGS. Heading text is \bfseries but math inside
% it stayed regular-weight CM — a hairline $V_T$ inside a heavy 24.8pt chapter
% title (measured: CMMI12 fl=6 next to FrankRuhlLibre-Bold fl=20; 155 live
% heading-with-math sites fleet-wide). \boldmath rides along wherever the class
% sets \bfseries for a heading; A/B measured: $V_T$ -> CMMIB10 fl=22, 0 missing
% characters (the bold CM faces ship with this TeX Live). Guarded per-macro so
% the article class (no \@makechapterhead) and future class changes stay safe.
% v4.0.1 (adoption fix): the SEARCH patterns above carry no parameter token.
% v4.0.0 made them specific by including the class's own #1/#2, and measured on
% TeX Live 2023 that never matches: a \patchcmd search pattern containing a
% parameter token cannot match while the call sits inside \@ifundefined{...}{}{...}
% (probed both #1 and ##1 -> FAIL; the same pattern at top level succeeds). Every
% one of these patches therefore failed, and because v4.0.0 also replaced the
% empty failure branch with \PackageError, the failure became FATAL: all four
% bundled example assets stopped building ("Cannot patch chapter-title math
% weight -> Fatal error, no output PDF"). The \PackageError branches are KEPT —
% a silent no-op here is what let v3.75's fix rot unnoticed — only the parameter
% tokens are dropped, so the patches match and a real class change still errors.
\makeatletter
\@ifundefined{@makechapterhead}{}{%
  \patchcmd{\@makechapterhead}{\huge\bfseries \@chapapp}
    {\huge\bfseries\boldmath \@chapapp}{}
    {\PackageError{hebrew-lualatex-pdf}{Cannot patch chapter-number math weight}{Unsupported document class}}%
  \patchcmd{\@makechapterhead}{\Huge \bfseries}
    {\Huge \bfseries\boldmath}{}
    {\PackageError{hebrew-lualatex-pdf}{Cannot patch chapter-title math weight}{Unsupported document class}}%
  \patchcmd{\@makeschapterhead}{\Huge \bfseries}
    {\Huge \bfseries\boldmath}{}
    {\PackageError{hebrew-lualatex-pdf}{Cannot patch starred chapter math weight}{Unsupported document class}}}
% v4.0.1: \part sits under a THREE-DEEP wrapper chain, so the class-based anchor
% is not in \@part at all. Measured with \show: hyperref rewrites \@part to
% `...\phantomsection\fi \H@old@part`, and nameref rewrites THAT to
% `\NR@gettitle{#1}\NR@part[{#1}]{#2}` — the class body ends up in \NR@part.
% v4.0.0 patched \@part and made failure fatal, which is why all seven bundled
% examples stopped building. \part is NOT a dead heading level here: 188 live uses
% across the fleet, so the bold-math fix is worth getting right rather than
% dropping. We walk the chain innermost-first and stop at the first hit. Failure is
% a WARNING, not an error: the chain is a property of hyperref/nameref versions, so
% a future package release must degrade to regular-weight part math — a cosmetic
% loss — instead of refusing to build. Chapter and section keep \PackageError,
% because those targets are in the CLASS and were verified to patch here.
\newif\iflakin@partnum \newif\iflakin@parttitle \newif\iflakin@sparttitle
\def\lakin@partnum#1#2{\iflakin@partnum\else
  \patchcmd{#1}{#2\bfseries \partname}{#2\bfseries\boldmath \partname}
    {\lakin@partnumtrue}{}\fi}
\def\lakin@parttitle#1#2{\iflakin@parttitle\else
  \patchcmd{#1}{#2 \bfseries}{#2 \bfseries\boldmath}{\lakin@parttitletrue}{}\fi}
\def\lakin@sparttitle#1#2{\iflakin@sparttitle\else
  \patchcmd{#1}{#2 \bfseries}{#2 \bfseries\boldmath}{\lakin@sparttitletrue}{}\fi}
\@ifundefined{c@chapter}{\def\lakin@pnum{\Large}\def\lakin@ptit{\huge}}%
                        {\def\lakin@pnum{\huge}\def\lakin@ptit{\Huge}}
\@ifundefined{NR@part}{}{\expandafter\lakin@partnum\expandafter\NR@part\expandafter{\lakin@pnum}%
                         \expandafter\lakin@parttitle\expandafter\NR@part\expandafter{\lakin@ptit}}
\@ifundefined{H@old@part}{}{\expandafter\lakin@partnum\expandafter\H@old@part\expandafter{\lakin@pnum}%
                            \expandafter\lakin@parttitle\expandafter\H@old@part\expandafter{\lakin@ptit}}
\@ifundefined{@part}{}{\expandafter\lakin@partnum\expandafter\@part\expandafter{\lakin@pnum}%
                       \expandafter\lakin@parttitle\expandafter\@part\expandafter{\lakin@ptit}}
\@ifundefined{NR@spart}{}{\expandafter\lakin@sparttitle\expandafter\NR@spart\expandafter{\lakin@ptit}}
\@ifundefined{H@old@spart}{}{\expandafter\lakin@sparttitle\expandafter\H@old@spart\expandafter{\lakin@ptit}}
\@ifundefined{@spart}{}{\expandafter\lakin@sparttitle\expandafter\@spart\expandafter{\lakin@ptit}}
\iflakin@partnum\else\PackageWarning{hebrew-lualatex-pdf}{Part-number math stays
  regular weight: no known \noexpand\part wrapper matched (hyperref/nameref
  version change). Cosmetic only}\fi
\iflakin@parttitle\else\PackageWarning{hebrew-lualatex-pdf}{Part-title math stays
  regular weight: no known \noexpand\part wrapper matched. Cosmetic only}\fi
% v4.12.0: THE PART NUMBER IN THE TOC WAS NOT DIRECTIONALLY SEPARATED FROM THE
% TITLE. report.cls/article.cls compose the part entry as ONE text stream --
% \addcontentsline{toc}{part}{\thepart\hspace{1em}#1} -- while EVERY other level
% goes through \numberline, which puts the number in an \hbox. A box is a bidi
% OBJECT and cannot join a neighbouring run; a bare "IV" can. So in the RTL TOC
% paragraph a Latin number and a Latin-INITIAL title merge into a SINGLE LTR run
% and are laid out left-to-right as a unit: the number lands on the wrong side of
% the row, and a mixed title is scrambled outright. Measured on \part{MCP ...}:
% the printed order was "<tail> V MCP" -- title fragment, number, title fragment.
% The chapter/section rows one line below, with the SAME Latin-initial titles,
% were correct on the same page. That contrast is the whole diagnosis.
% \en{} IS NOT A WORKAROUND, measured both ways on one fixture: it fixes a
% PURE-Latin title (\part{\en{Hooks}} renders correctly) and does NOT fix a MIXED
% one (\part{\en{MCP} <hebrew>} stays scrambled). babel's language switch changes
% font and hyphenation, not the bidi run structure -- so the defect is not
% reachable from the source side at all, which is why it belongs here.
% Boxing the number is the entire fix and is a STRICT NO-OP where the bug never
% bit: a Hebrew-titled document compiles BYTE-IDENTICAL before and after (md5,
% 3 clean runs), because there the neutral glue already sat between an L run and
% an R run. The PDF outline is byte-identical too -- hyperref's \pdfstringdef
% neutralises \mbox, so bookmark strings keep their part numbers.
% \numberline{\thepart} would also fix the direction, but it re-geometries every
% part row in the fleet and DROPS the number from the part bookmark (hyperref
% strips \numberline). \mbox changes nothing that was already right.
% The search pattern deliberately carries no # token: etoolbox cannot match a
% parameter token in a patch pattern (measured -- all three chain members
% reported "not patchable" while the pattern included #1). "\thepart \hspace
% {1em}" is brace-balanced and occurs exactly once in the wrapper body.
\newif\iflakin@parttocnum
\def\lakin@parttocnum#1{\iflakin@parttocnum\else
  \patchcmd{#1}{\thepart \hspace {1em}}{\mbox {\thepart }\hspace {1em}}
    {\lakin@parttocnumtrue}{}\fi}
\@ifundefined{NR@part}{}{\lakin@parttocnum\NR@part}
\@ifundefined{H@old@part}{}{\lakin@parttocnum\H@old@part}
\@ifundefined{@part}{}{\lakin@parttocnum\@part}
\iflakin@parttocnum\else\PackageWarning{hebrew-lualatex-pdf}{Part TOC number is
  not boxed: no known \noexpand\part wrapper matched. A part title that BEGINS
  with Latin will print its number on the wrong side of its TOC row}\fi
\patchcmd{\section}{\bfseries}{\bfseries\boldmath}{}
  {\PackageError{hebrew-lualatex-pdf}{Cannot patch section math weight}{Unsupported document class}}
\patchcmd{\subsection}{\bfseries}{\bfseries\boldmath}{}
  {\PackageError{hebrew-lualatex-pdf}{Cannot patch subsection math weight}{Unsupported document class}}
\patchcmd{\subsubsection}{\bfseries}{\bfseries\boldmath}{}
  {\PackageError{hebrew-lualatex-pdf}{Cannot patch subsubsection math weight}{Unsupported document class}}
\patchcmd{\paragraph}{\bfseries}{\bfseries\boldmath}{}{}
\patchcmd{\subparagraph}{\bfseries}{\bfseries\boldmath}{}{}
\@ifundefined{l@chapter}{}{%
  \patchcmd{\l@chapter}{\bfseries}{\bfseries\boldmath}{}{}}
\makeatother
\setcounter{tocdepth}{\lakittocdepth}
% <hllpdf:measured-toc>
% v4.16.0: Measure actual number labels, not a guessed digit count. Read only
% canonical contentsline records; never execute titles or addtocontents code
% during the measuring pass. Widths are local to this individual list and are
% available before its first row: no per-row drift and no extra aux pass.
\makeatletter
\newread\hllpdf@tocstream
\newdimen\hllpdf@toc@chapter
\newdimen\hllpdf@toc@section
\newdimen\hllpdf@toc@subsection
\newdimen\hllpdf@toc@subsubsection
\newdimen\hllpdf@toc@paragraph
\newdimen\hllpdf@toc@subparagraph
\newdimen\hllpdf@toc@figure
\newdimen\hllpdf@toc@table
\newdimen\hllpdf@tocpage
\newdimen\hllpdf@toc@sectionindent
\newdimen\hllpdf@toc@subsectionindent
\newdimen\hllpdf@toc@subsubsectionindent
\newdimen\hllpdf@toc@paragraphindent
\newdimen\hllpdf@toc@subparagraphindent
\def\hllpdf@tocdefaults{%
  \hllpdf@toc@chapter=2em
  \hllpdf@toc@section=3.2em
  \hllpdf@toc@subsection=3.9em
  \@ifundefined{c@chapter}{%
    \hllpdf@toc@subsubsection=3.2em
    \hllpdf@toc@paragraph=4.1em
    \hllpdf@toc@subparagraph=5em
  }{%
    \hllpdf@toc@subsubsection=4.1em
    \hllpdf@toc@paragraph=5em
    \hllpdf@toc@subparagraph=6em
  }%
  \hllpdf@toc@figure=2.3em
  \hllpdf@toc@table=2.3em
  \hllpdf@tocpage=\@pnumwidth\relax
}
\def\hllpdf@tocindents{%
  % Preserve the established shallow layout when its slots already fit.
  \hllpdf@toc@sectionindent=1.5em
  \@tempdima=\hllpdf@toc@chapter \advance\@tempdima by -2em
  \advance\hllpdf@toc@sectionindent by \@tempdima
  \hllpdf@toc@subsectionindent=\dimexpr
    \hllpdf@toc@sectionindent+\hllpdf@toc@section\relax
  \hllpdf@toc@subsubsectionindent=\dimexpr
    \hllpdf@toc@subsectionindent+\hllpdf@toc@subsection\relax
  \hllpdf@toc@paragraphindent=\dimexpr
    \hllpdf@toc@subsubsectionindent+\hllpdf@toc@subsubsection\relax
  \hllpdf@toc@subparagraphindent=\dimexpr
    \hllpdf@toc@paragraphindent+\hllpdf@toc@paragraph\relax
}
\def\hllpdf@tocfont@chapter{\bfseries\boldmath}
\def\hllpdf@tocfont@part{\large\bfseries\boldmath}
\def\hllpdf@tocentryfont{%
  \ifcsname hllpdf@tocfont@\hllpdf@toclevel\endcsname
    \csname hllpdf@tocfont@\hllpdf@toclevel\endcsname
  \fi}
\long\def\hllpdf@tocdiscard#1\hllpdf@tocstop{}
\def\hllpdf@tocscanrecord{%
  \@ifnextchar\contentsline{\hllpdf@tocscanline}{%
    \@ifnextchar\babel@toc{\hllpdf@tocscanlanguage}{\hllpdf@tocdiscard}}}
\long\def\hllpdf@tocscanlanguage\babel@toc#1#2#3\hllpdf@tocstop{%
  % Babel records are declarations, not arbitrary code. Track the language
  % without executing the original directive or changing the surrounding list.
  \def\hllpdf@toclanguage{#1}}
\long\def\hllpdf@tocscanline\contentsline#1#2#3#4#5\hllpdf@tocstop{%
  \def\hllpdf@toclevel{#1}%
  \ifcsname hllpdf@toc@#1\endcsname
    \hllpdf@tocscantitle#2\hllpdf@tocstop
  \fi
  % The page-number column needs its own measurement (including bold chapter
  % rows). Its extra clearance separates the page label from dot leaders.
  \setbox\@tempboxa=\hbox{\foreignlanguage{\hllpdf@toclanguage}{\hllpdf@tocentryfont #3\kern.3em}}%
  \ifdim\wd\@tempboxa>\hllpdf@tocpage
    \hllpdf@tocpage=\wd\@tempboxa
  \fi}
\def\hllpdf@tocscantitle{%
  \@ifnextchar\numberline{\hllpdf@tocscanlabel}{\hllpdf@tocdiscard}}
\long\def\hllpdf@tocscanlabel\numberline#1#2\hllpdf@tocstop{%
  \setbox\@tempboxa=\hbox{\foreignlanguage{\hllpdf@toclanguage}{\hllpdf@tocentryfont #1\kern.5em}}%
  \ifdim\wd\@tempboxa>\csname hllpdf@toc@\hllpdf@toclevel\endcsname
    \csname hllpdf@toc@\hllpdf@toclevel\endcsname=\wd\@tempboxa
  \fi}
\def\hllpdf@tocread{%
  \ifeof\hllpdf@tocstream\else
    \read\hllpdf@tocstream to\hllpdf@tocrecord
    \expandafter\hllpdf@tocscanrecord\hllpdf@tocrecord\hllpdf@tocstop
    \expandafter\hllpdf@tocread
  \fi}
\def\hllpdf@tocmeasure#1{%
  \makeatletter
  \edef\hllpdf@toclanguage{\languagename}%
  \hllpdf@tocdefaults
  \openin\hllpdf@tocstream=\jobname.#1\relax
  \hllpdf@tocread
  \closein\hllpdf@tocstream
  \hllpdf@tocindents
  \edef\@pnumwidth{\the\hllpdf@tocpage}%
  \@tempdima=\dimexpr\hllpdf@tocpage+1em\relax
  \ifdim\@tempdima>\@tocrmarg\relax
    \edef\@tocrmarg{\the\@tempdima}%
  \fi}
\let\hllpdf@originalstarttoc\@starttoc
\def\@starttoc#1{%
  \begingroup
    \def\hllpdf@tocext{#1}%
    \def\hllpdf@tocwanted{toc}%
    \ifx\hllpdf@tocext\hllpdf@tocwanted \hllpdf@tocmeasure{#1}\else
      \def\hllpdf@tocwanted{lof}%
      \ifx\hllpdf@tocext\hllpdf@tocwanted \hllpdf@tocmeasure{#1}\else
        \def\hllpdf@tocwanted{lot}%
        \ifx\hllpdf@tocext\hllpdf@tocwanted \hllpdf@tocmeasure{#1}\fi
      \fi
    \fi
    \hllpdf@originalstarttoc{#1}%
  \endgroup}
\hllpdf@tocdefaults
\hllpdf@tocindents
\@ifundefined{l@chapter}{}{%
  \patchcmd{\l@chapter}{\setlength\@tempdima{1.5em}}%
    {\setlength\@tempdima{\hllpdf@toc@chapter}}{}%
    {\PackageError{hebrew-lualatex-pdf}{Cannot size chapter TOC numbers}%
       {Unsupported contents layout: the chapter number slot was not found}}}
\renewcommand*\l@section{\@dottedtocline{1}%
  {\hllpdf@toc@sectionindent}{\hllpdf@toc@section}}
\renewcommand*\l@subsection{\@dottedtocline{2}%
  {\hllpdf@toc@subsectionindent}{\hllpdf@toc@subsection}}
\renewcommand*\l@subsubsection{\@dottedtocline{3}%
  {\hllpdf@toc@subsubsectionindent}{\hllpdf@toc@subsubsection}}
\renewcommand*\l@paragraph{\@dottedtocline{4}%
  {\hllpdf@toc@paragraphindent}{\hllpdf@toc@paragraph}}
\renewcommand*\l@subparagraph{\@dottedtocline{5}%
  {\hllpdf@toc@subparagraphindent}{\hllpdf@toc@subparagraph}}
\renewcommand*\l@figure{\@dottedtocline{1}{1.5em}{\hllpdf@toc@figure}}
\renewcommand*\l@table{\@dottedtocline{1}{1.5em}{\hllpdf@toc@table}}
\makeatother
% </hllpdf:measured-toc>
\setcounter{secnumdepth}{\lakitsecnumdepth}
% Public preamble-only caption-numbering policy. The historical within-major
% default is untouched unless this command is called. "continuous" removes the
% reset edge as well as the printed prefix, so Figure 1/2 and Table 1/2 stay
% unique across sections/chapters. Returning to "within-major" restores the
% original class-aware displays and native destination conventions. The policy
% never replaces \figcap/\tabcap or their measured link-placement mechanism.
\makeatletter
\let\lakin@majorfig\thefig
\let\lakin@majortbl\thetbl
\let\lakin@majorHfig\theHfig
\let\lakin@majorHtbl\theHtbl
\newcommand{\lakitcaptionnumbering}[1]{%
  \edef\lakin@numbering{#1}%
  \def\lakin@continuous{continuous}\def\lakin@withinmajor{within-major}%
  \ifx\lakin@numbering\lakin@continuous
    \@ifundefined{c@chapter}{%
      \@removefromreset{fig}{section}\@removefromreset{tbl}{section}%
    }{\@removefromreset{fig}{chapter}\@removefromreset{tbl}{chapter}}%
    \renewcommand{\thefig}{\arabic{fig}}\renewcommand{\thetbl}{\arabic{tbl}}%
    \def\theHfig{\arabic{fig}}\def\theHtbl{\arabic{tbl}}%
  \else\ifx\lakin@numbering\lakin@withinmajor
    \@ifundefined{c@chapter}{%
      \@removefromreset{fig}{section}\@removefromreset{tbl}{section}%
      \@addtoreset{fig}{section}\@addtoreset{tbl}{section}%
    }{%
      \@removefromreset{fig}{chapter}\@removefromreset{tbl}{chapter}%
      \@addtoreset{fig}{chapter}\@addtoreset{tbl}{chapter}}%
    \let\thefig\lakin@majorfig\let\thetbl\lakin@majortbl
    \let\theHfig\lakin@majorHfig\let\theHtbl\lakin@majorHtbl
  \else
    \PackageError{hebrew-english-latex-studio}{Unknown caption-numbering policy '#1'}%
      {Use continuous or within-major before begin{document}.}%
  \fi\fi}
\@onlypreamble\lakitcaptionnumbering
\makeatother
% Architecture readiness sentinel. Optional runtime layers require the complete
% core, so the marker is defined only after every mechanism/baseline definition.
% <hllpdf:architecture-ready>
\def\hllpdfcoreloaded{1}
% </hllpdf:architecture-ready>

% v4.22: registered file resources, not an assertion about embedded subset bytes.
% The log framing avoids long-path wrapping and preserves UTF-8 as hex.
% <hllpdf:font-inventory>
\directlua{
  local function inventory()
    local names = {}
    local unnamed = 0
    for id, item in font.each() do
      local path = item.filename
      if path and not (path == "") then names[path] = true else unnamed = unnamed + 1 end
    end
    local sorted = {}
    for path in pairs(names) do table.insert(sorted,path) end
    table.sort(sorted)
    texio.write_nl("log","HLLFONT|BEGIN|1")
    for index,path in ipairs(sorted) do
      local hex = path:gsub(".",function(c) return string.format(string.char(37) .. "02x",string.byte(c)) end)
      local count = math.ceil(string.len(hex)/48)
      texio.write_nl("log","HLLFONT|FILE|" .. index .. "|" .. count)
      for offset=1,string.len(hex),48 do texio.write_nl("log","HLLFONT|HEX|" .. hex:sub(offset,offset+47)) end
    end
    texio.write_nl("log","HLLFONT|END|" .. rawlen(sorted) .. "|" .. unnamed)
    texio.write_nl("log","")
  end
  luatexbase.add_to_callback("finish_pdffile",inventory,"hllpdf.runtime.fonts")
}
% </hllpdf:font-inventory>

%
% Supported class contract: report or article.  Alternate paper sizes must also
% be declared to build.sh through EXPECTED_PAPER; the PDF gate never infers them.
% Do NOT use these hooks as a theme mechanism: theme overlays must not own page
% geometry, class selection, or direction semantics. A declared presentation-policy
% overlay may intentionally change numbering visibility after the preamble.
% ============================================================================
% Runtime architecture guard: generated from scripts/architecture_contract.py.
% The readiness sentinel remains delayed until the END of this file.
% <hllpdf:architecture-guard>
\ifdefined\hllpdfcoreloaded
  \errmessage{hllpdf: preamble.tex loaded more than once}
  \expandafter\endinput
\fi
% </hllpdf:architecture-guard>

\ifdefined\lakitdocumentclass\else\def\lakitdocumentclass{report}\fi
\ifdefined\lakitclassoptions\else\def\lakitclassoptions{12pt,a4paper}\fi
\ifdefined\lakitgeometryoptions\else\def\lakitgeometryoptions{margin=2.3cm,headheight=14pt}\fi
\ifdefined\lakitlinespread\else\def\lakitlinespread{1.3}\fi
\ifdefined\lakittocdepth\else\def\lakittocdepth{1}\fi
\ifdefined\lakitsecnumdepth\else\def\lakitsecnumdepth{2}\fi
% Paragraph organization belongs to the document layout, independently of its
% language and theme. Preserve historical class dimensions unless explicitly
% selecting block paragraphs; never suppress indentation with an everypar hack.
\ifdefined\lakitparagraphstyle\else\def\lakitparagraphstyle{indented}\fi
\edef\lakitparagraphstyle{\lakitparagraphstyle}
\def\hllpdfparagraphblock{block}
\def\hllpdfparagraphindented{indented}
\ifx\lakitparagraphstyle\hllpdfparagraphblock\else
  \ifx\lakitparagraphstyle\hllpdfparagraphindented\else
    \errmessage{hllpdf: paragraph style must be block or indented}
  \fi
\fi
% v4.3.1: the public class hook has a deliberately narrow support contract.
% Before this guard, \def\lakitdocumentclass{book} silently compiled even though
% every heading/TOC regression test targets report/article; an unsupported shell
% could therefore look "supported" until a subtle class-specific defect shipped.
% Fail before \documentclass rather than letting an undocumented combination enter
% the mechanism core. Plain-TeX comparison is used because LaTeX packages are not
% available yet at this point.
\edef\lakitvalidatedclass{\lakitdocumentclass}
\def\lakitreportclass{report}
\def\lakitarticleclass{article}
\ifx\lakitvalidatedclass\lakitreportclass
\else\ifx\lakitvalidatedclass\lakitarticleclass
\else
  \errmessage{hllpdf: unsupported class \lakitdocumentclass; use report or article}
\fi\fi
\let\lakitvalidatedclass\relax
\let\lakitreportclass\relax
\let\lakitarticleclass\relax
% Language intent is independent of document class, theme and content domain.
% The two dominant-language profiles differ in authoring intent, not direction.
% A balanced bilingual document must explicitly choose a navigation/base locale.
% <hllpdf:language-profile>
\providecommand{\lakitlanguageprofile}{hebrew}
\edef\lakitlanguageprofile{\lakitlanguageprofile}
\def\hllpdfhebrew{hebrew}\def\hllpdfenglish{english}
\def\hllpdfhebrewmixed{hebrew-with-english}
\def\hllpdfenglishmixed{english-with-hebrew}
\def\hllpdfbilingual{bilingual}
\let\hllpdfexpectedlanguage\relax
\ifx\lakitlanguageprofile\hllpdfhebrew
  \let\hllpdfexpectedlanguage\hllpdfhebrew
\else\ifx\lakitlanguageprofile\hllpdfhebrewmixed
  \let\hllpdfexpectedlanguage\hllpdfhebrew
\else\ifx\lakitlanguageprofile\hllpdfenglish
  \let\hllpdfexpectedlanguage\hllpdfenglish
\else\ifx\lakitlanguageprofile\hllpdfenglishmixed
  \let\hllpdfexpectedlanguage\hllpdfenglish
\else\ifx\lakitlanguageprofile\hllpdfbilingual
  \ifdefined\lakitmainlanguage\else
    \errmessage{hllpdf: bilingual profile requires explicit lakitmainlanguage}
  \fi
\else
  \errmessage{hllpdf: unsupported language profile \lakitlanguageprofile}
\fi\fi\fi\fi\fi
\ifdefined\lakitmainlanguage
  \edef\lakitmainlanguage{\lakitmainlanguage}
\else
  \let\lakitmainlanguage\hllpdfexpectedlanguage
\fi
\ifx\lakitmainlanguage\hllpdfhebrew\else
  \ifx\lakitmainlanguage\hllpdfenglish\else
    \errmessage{hllpdf: main language must be hebrew or english}
  \fi
\fi
\ifx\hllpdfexpectedlanguage\relax\else
  \ifx\lakitmainlanguage\hllpdfexpectedlanguage\else
    \errmessage{hllpdf: main language conflicts with language profile}
  \fi
\fi
\typeout{HLLLANG|\lakitlanguageprofile|\lakitmainlanguage}
% </hllpdf:language-profile>
\documentclass[\lakitclassoptions]{\lakitdocumentclass}
\usepackage[no-math]{fontspec}        % (1) ★ no-math — THE critical line (see recipe.md)
% layout=lists asks Babel to adapt list parshape when paragraph and page
% directions differ. Without it, minority-language list bodies lose their
% indentation and labels hang outside the text block, despite a clean log.
% The upstream mechanism preserves ordinary same-direction list geometry.
\usepackage[bidi=basic,layout=lists]{babel}
\ifx\lakitmainlanguage\hllpdfenglish
  \babelprovide[import]{hebrew}
  \babelprovide[import,main]{english}
\else
  \babelprovide[import,main]{hebrew}
  \babelprovide[import]{english}
\fi
% v4.22: bind canonical families and fallbacks to the selected installation.
% Lookup by basename or family alone is not proof of the file actually used.
% <hllpdf:font-binding>
\directlua{
  local dir = os.getenv("HEBREW_LATEX_FONT_DIR")
  if not dir or dir == "" then
    local data = os.getenv("XDG_DATA_HOME")
    if not data or data == "" or not (data:sub(1,1) == "/") then
      data = (os.getenv("HOME") or "") .. "/.local/share"
    end
    dir = data .. "/fonts/hebrew-lualatex-pdf"
  end
  if not (dir:sub(1,1) == "/") then dir = lfs.currentdir() .. "/" .. dir end
  if dir:find("[",1,true) or dir:find("]",1,true) or dir:find(string.char(0),1,true)
     or dir:find(string.char(10),1,true) or dir:find(string.char(13),1,true) then
    tex.error("hllpdf: unsupported character in font directory")
  end
  if not (lfs.attributes(dir,"mode") == "directory") then
    tex.error("hllpdf: font directory is missing")
  end
  hllpdf_font_dir = dir:gsub("/+$", "") .. "/"
  tex.sprint(string.char(92) .. "def" .. string.char(92) .. "hllpdffontdir{")
  tex.sprint(-2,hllpdf_font_dir)
  tex.sprint("}")
  luaotfload.add_fallback("hebfb",{"[" .. hllpdf_font_dir .. "DejaVuSerif.ttf]:mode=node;"})
  luaotfload.add_fallback("hebfbbf",{"[" .. hllpdf_font_dir .. "DejaVuSerif-Bold.ttf]:mode=node;"})
}
% </hllpdf:font-binding>
% v3.22: Frank Ruhl Libre ships NO italic face. Instead of silently-upright
% \emph (v<=3.20) or bold-as-emphasis (v3.21), synthesize a slanted face with
% FakeSlant — real, VISIBLE italic emphasis with correct \emph nesting
% semantics (inner \emph toggles back upright). Verified empirically at
% 300dpi: slant clearly visible, 0 errors, 0 missing glyphs.
\babelfont{rm}[Path={\hllpdffontdir},RawFeature={fallback=hebfb},
  BoldFont={FrankRuhlLibre-Bold.ttf},
  BoldFeatures={RawFeature={fallback=hebfbbf}},   % v3.56 — see hebfbbf above
  ItalicFont={FrankRuhlLibre-Regular.ttf}, ItalicFeatures={FakeSlant=0.22},
  % Bold-italic uses the weight-matched fallback and a synthetic slant.
  BoldItalicFont={FrankRuhlLibre-Bold.ttf}, BoldItalicFeatures={RawFeature={fallback=hebfbbf},FakeSlant=0.22}
  ]{FrankRuhlLibre-Regular.ttf}
% Sans fallbacks preserve family and weight for the few missing codepoints.
\directlua{luaotfload.add_fallback("hebfbsans",{"[" .. hllpdf_font_dir .. "DejaVuSans.ttf]:mode=node;"})}
\directlua{luaotfload.add_fallback("hebfbsansbf",{"[" .. hllpdf_font_dir .. "DejaVuSans-Bold.ttf]:mode=node;"})}
\babelfont{sf}[Path={\hllpdffontdir},RawFeature={fallback=hebfbsans},
  BoldFont={Heebo-Bold.ttf},
  BoldFeatures={RawFeature={fallback=hebfbsansbf}}]{Heebo-Regular.ttf}
% v3.24: Scale=MatchLowercase — DejaVu Sans Mono has an outsized x-height, so
% unscaled inline code visually SHOUTS next to the delicate Frank Ruhl body
% (measured on a 197-page build: scaled, a code atom sits in the line at the
% body's weight, and the document tightened by ~5%). Applies to listings too.
\babelfont{tt}[Path={\hllpdffontdir},RawFeature={fallback=hebfb},Scale=MatchLowercase,
  BoldFont={DejaVuSansMono.ttf},
  BoldFeatures={RawFeature={fallback=hebfb},FakeBold=1.5},
  ItalicFont={DejaVuSansMono.ttf},
  ItalicFeatures={RawFeature={fallback=hebfb},FakeSlant=0.2},
  BoldItalicFont={DejaVuSansMono.ttf},
  BoldItalicFeatures={RawFeature={fallback=hebfb},FakeBold=1.5,FakeSlant=0.2}
  ]{DejaVuSansMono.ttf}

% Explicit multilingual/English profiles use real Latin faces and italics.
% The legacy/default Hebrew profile retains its established measured typography.
% Latin Modern is provided by the TeX distribution, not the pinned font bundle;
% recorder/runtime evidence reports it as an additional toolchain resource.
\ifx\lakitlanguageprofile\hllpdfhebrew\else
  \babelfont[english]{rm}[Ligatures=TeX,RawFeature={fallback=hebfb},
    BoldFeatures={RawFeature={fallback=hebfbbf}},
    BoldItalicFeatures={RawFeature={fallback=hebfbbf}}]{Latin Modern Roman}
  \babelfont[english]{sf}[Ligatures=TeX,RawFeature={fallback=hebfbsans},
    BoldFeatures={RawFeature={fallback=hebfbsansbf}},
    BoldItalicFeatures={RawFeature={fallback=hebfbsansbf}}]{Latin Modern Sans}
\fi

\usepackage{amsmath,amssymb,mathtools}
% amsthm supplies theorem/proof machinery. Environment names remain
% document-specific; the generic proof label is localized below.
\usepackage{amsthm}
% Proof names follow the surrounding prose, including return from an island.
\addto\extrashebrew{\renewcommand{\proofname}{הוכחה}}
\addto\extrasenglish{\renewcommand{\proofname}{Proof}}
\ifx\lakitmainlanguage\hllpdfenglish
  \renewcommand{\proofname}{Proof}
\else
  \renewcommand{\proofname}{הוכחה}
\fi
% v3.24.1: equation tags under babel-bidi. A tag re-emitted by mathtools'
% showonlyrefs inside align inherits the RTL base and its parens MIRROR —
% ")2.1(" instead of "(2.1)" (plain align/equation tags are unaffected;
% verified on a minimal case + a 198-page build). LTR-isolate at the tag-BOX
% level: mathtools' show/hide + numbering logic is untouched. Do NOT override
% \tagform@ instead — that clobbers the showonlyrefs mechanism (measured:
% every tag shows, numbering breaks).
% v3.28.1: the v3.24.1 wrap must apply ONLY to tags of display environments.
% Inline \eqref in Hebrew PROSE also goes through \maketag@@@, and there the
% english-wrap MIRRORS the parens — ")1(" — while the unwrapped original is
% correct (proven on a 2x2 isolation matrix: patch x showonlyrefs; regressed
% every prose \eqref since v3.24.1, caught by the render gate's R4). The two
% call sites are indistinguishable by \ifmmode (tag boxes are text mode), so
% the display environments themselves raise a flag; \eqref in prose sees it
% false and gets the original untouched \maketag@@@.
\usepackage{etoolbox}

% babel bidi=basic (including upstream 26.9) assumes every automatic
% discretionary has a one-glyph replacement. LuaTeX can legally produce one
% with replace=nil after a hyphenated Latin value; Babel then indexes nil and the
% whole build dies. Do the narrow upstream-compatible correction: while Babel
% classifies the list, present only that malformed disc as an ordinary
% discretionary, then restore its subtype. Babel still reorders the ENTIRE box;
% there is no pcall and no successful fallback to an unreordered node list.
% If Babel throws for any other reason, the build remains fatal. The synthetic
% nil-disc probe and the real table gallery in test_all.sh pin both layers.
\AtBeginDocument{\directlua{
  if Babel and Babel.bidi and not Babel.heb_nil_disc_compat then
    Babel.heb_nil_disc_compat = true
    local original = Babel.bidi
    local disc_id = node.id('disc')
    Babel.bidi = function(head, ispar, hdir)
      local hidden = {}
      for item in node.traverse_id(disc_id, head) do
        if item.subtype == 2 and (not item.replace or not item.replace.char) then
          table.insert(hidden, item)
          item.subtype = 0
        end
      end
      local result = original(head, ispar, hdir)
      for _, item in ipairs(hidden) do item.subtype = 2 end
      return result
    end
  end
}}
\makeatletter
\newif\ifHEBdisptag
\let\ORIGmaketag@@@\maketag@@@
\def\maketag@@@#1{\ifHEBdisptag\ORIGmaketag@@@{\foreignlanguage{english}{#1}}\else\ORIGmaketag@@@{#1}\fi}
\@for\HEB@env:=equation,equation*,align,align*,gather,gather*,multline,multline*,flalign,flalign*,alignat,alignat*,eqnarray,eqnarray*\do{%
  \expandafter\AtBeginEnvironment\expandafter{\HEB@env}{\HEBdisptagtrue}}
\makeatother
% English page + Hebrew prose: Babel's equation-only tag flip uses a primitive
% leqno path that shifts the formula to the left edge and the tag into the body
% of an English page. Align's full-width alignment is unaffected. Give only this
% display a physical LTR container and native amsmath eqno/leqno. The engine then
% owns real tag width, collision handling and labels; no proxy tag box hides ink
% from its fit calculation. Babel's saved \bbl@eqnodir still restores Hebrew text
% inside math, and the prose locale/math direction are never switched.
% Close the prior RTL paragraph before normalization. After the environment's
% group restores its directions, close TeX's empty resumed display paragraph and
% use native LaTeX continuation handling: same-paragraph prose has no indent,
% while an authored blank line starts an ordinary indented RTL paragraph.
% Register after Babel's setup; remain local to equation(*) in a main-English
% document's actual RTL prose. Native Hebrew, ordinary English and explicit
% fleqn retain their original paths. No page/display dimensions are assigned.
% <hllpdf:minority-rtl-equation-geometry>
\makeatletter
\newcommand{\HEB@minorityprosecontinuation}{\par\@doendpe}
\newcommand{\HEB@minorityequationgeometry}{%
  \ifx\lakitmainlanguage\hllpdfenglish
    \ifnum\pardirection=\@ne
      \if@fleqn\else
        \par\aftergroup\HEB@minorityprosecontinuation
        \iftagsleft@\let\veqno\@@leqno
        \else\let\veqno\@@eqno\fi
        \bodydir TLT\pardir TLT\textdir TLT\relax
      \fi
    \fi
  \fi}
\AtBeginDocument{%
  \AddToHook{env/equation/begin}{\HEB@minorityequationgeometry}%
  \AddToHook{env/equation*/begin}{\HEB@minorityequationgeometry}}
\makeatother
% </hllpdf:minority-rtl-equation-geometry>
% v3.63: \shortintertext rendered Hebrew MIRROR-REVERSED and flush LEFT — an entire
% commentary line inside an align printed backwards. babel's bidi=basic patches
% \intertext@ only (luababel.def), and mathtools' \shortintertext uses the identical
% \vbox{\normalbaselines … \noindent\ignorespaces#1\par} skeleton but is never
% patched, so its paragraph keeps pardir/textdir LTR: no bidi reordering (logical
% order on the page) and flush-left. Measured against \intertext with the same
% Applying babel's own replacement to \shortintertext@ makes the two identical.
% \h{}/\mbox fix the letter order but leave the line flush LEFT, so they are not
% substitutes. Guarded + \AtBeginDocument because mathtools \let-copies at load time.
% NB \ExplSyntaxOn is REQUIRED: mathtools' internals carry expl3 names, and without
% the expl3 catcodes the `_` is read as a subscript and the patch detonates far away
% with "! Illegal unit of measure (pt inserted)" at the \end{align} (measured).
\makeatletter
\AtBeginDocument{%
  \@ifundefined{shortintertext@}{}{%
    \ExplSyntaxOn
    \bbl@sreplace\shortintertext@{\normalbaselines}%
      {\normalbaselines\ifx\bbl@eqnodir\relax\else\bbl@pardir\@ne\bbl@eqnodir\fi}%
    \ExplSyntaxOff}}
\makeatother
% geometry inherits the canonical A4 sheet from \documentclass; build.sh then
% verifies every rendered MediaBox and CropBox, so a silent Letter fallback is
% impossible to publish as a clean default build.
\usepackage[\lakitgeometryoptions]{geometry}
\linespread{\lakitlinespread}       % layout hook; canonical default remains 1.3
\ifx\lakitparagraphstyle\hllpdfparagraphblock
  \setlength{\parindent}{0pt}
  \setlength{\parskip}{6pt plus 1pt minus 1pt}
\fi
% These dimensions remain ordinary scoped TeX lengths: authors can customize
% them after core input. Language switches inherit them; lists keep their own
% indentation and narrow reference layouts establish their local paragraph gaps.
\typeout{HLLPAR|\lakitparagraphstyle}
\usepackage{xcolor}
\usepackage{mdframed}
\usepackage{listings}                 % code listings (assembly / C / pseudo-code)
\usepackage{tikz}
\usetikzlibrary{arrows.meta,positioning,calc,shapes.geometric,fit,backgrounds}
\usepackage{pgfplots}\pgfplotsset{compat=1.18}
\usepackage{quantikz}                 % quantum circuits (remove if unused)
\usepackage{booktabs}
\usepackage{array}
\usepackage{float}                    % [H] = place table exactly here (reference doc: no drifting)
\usepackage{enumitem}
\usepackage{needspace}

% ============================================================================
% ★★ v3.4 — THE DIRECTION-ISOLATION TRIO (the heart of mixing Hebrew + Latin)
% Mental model:
%   \h{...}    Hebrew (RTL) island      — Hebrew text inside an LTR context
%                                          (TikZ nodes, the english `tikzpic`).
%   \en{...}   English (LTR) isolate    — Latin text / parentheses inside a
%                                          Hebrew (RTL) context. KEEPS PARENS
%                                          AND ORDER CORRECT.
%   \code{...} \en + monospace          — code inside Hebrew / inside a frametitle.
% CRITICAL: never wrap Latin in \h{} (it is for HEBREW). Inside a TikZ node,
% Latin in \h{} renders REVERSED (Datapath->htapataD, CPU->UPC) and parens
% mirror. Isolate every Latin/paren chunk with \en{} or \code{}. See
% references/figures-boxes-listings.md.
% ---- TEXT INSIDE MATH (v3.8) ----------------------------------------------
% Put ANY text that sits inside math in a \text{}. The ONLY thing that is
% invalid in math is a bare \h/\en (= \foreignlanguage) used DIRECTLY, outside
% \text -> "\rmfamily invalid in math mode". Inside \text:
%   Hebrew label           -> \text{...}            (bare Hebrew renders fine;
%                              \h is optional — BUT NOT ALWAYS HARMLESS (v3.67):
%                              if the \h group OPENS with a nested island
%                              ($..$, \en{}, \code{}, \texttt{}, \mbox{}) the
%                              WHOLE group loses bidi and the Hebrew prints in
%                              LOGICAL order = gibberish. Measured:
%                                \text{\h{$f$ קבועה}}  -> reads "העובק"  BROKEN
%                                \text{\h{קבועה $f$}}  -> correct
%                                \text{$f$ קבועה}      -> correct (no \h)
%                              Hebrew first inside \h, or drop the \h. Flagged
%                              H-OPENS-WITH-LTR-GROUP.)
%   plain Latin word        -> \text{...}            (no marker needed)
%   Latin WITH ( ) or [ ]   -> \text{\en{...}}       (\en forces LTR; a plain
%                              \text mirrors the parens in RTL -> "Spec (TNR(" )
%   f(Hebrew/text arg)      -> \text{f}(\text{...})  -- keep the parens at MATH
%                              level, NOT \text{f(arg)} (parens inside the box mirror)
%   code in math            -> \mathtt{...}  (or \texttt{}, math-guarded below)
% NEVER write \en{...} or \h{...} straight into $...$ / \[..\] (outside \text);
% always via \text{}. (\mbox{...} also works but does not scale in scripts.)
% ============================================================================
\newcommand{\h}[1]{\foreignlanguage{hebrew}{#1}}
% v3.42.1 — \en islands do NOT hyphenate (auto OR at explicit hyphens). English hyphenation patterns were
% tearing identifiers and terms ACROSS THE RTL SEAM ("IO Regis-/ters",
% "In-/put", "Band-/width", "Per-/formance", "Multi-/ple" — all shipped on
% real pages; probe-proven: "ALUS-/rcB", "Mem-/ToReg"). An island is a label
% or a term, not English prose — a hyphen tear there is never right. A/B on
% a full 197pp lakit: identical overfull sets, same page count, ~11
% paragraphs healed, one mild underfull (badness 3525) that replaced a
% mutilated "Regis-/ters". Long English PROSE belongs in the englishprose
% environment below, which restores explicit-hyphen breaks while its paragraph
% is still active. A raw otherlanguage block inherits the global prohibition.
% (\l@nohyphenation needs \csname:
% '@' is not a letter here.)
% v3.43.2 — \exhyphenpenalty moved OUT of \en to a GLOBAL assignment below.
% Render-proven: a group-local \exhyphenpenalty is DEAD CODE — TeX reads the
% line-breaking parameters at \par time, after the \en group has closed, so
% \en{write-through} / plain not-taken still tore at the "-" ("write-/
% through" shipped on a real exam page). \language is different: it is
% recorded in the node list as text is added, so the in-group nohyphenation
% assignment above DOES work (and stays). The global penalty was A/B-measured
% on a 197pp lakit: identical overfull set (25), identical page count, the
% not-/taken tear healed. Hebrew is unaffected in practice — Hebrew compounds
% use maqaf (־), not ASCII "-".
\newcommand{\en}[1]{\foreignlanguage{english}{\language=\csname l@nohyphenation\endcsname #1}}
% Babel runs extras locally for inline islands and paragraph language changes.
% An English main document must not inherit the Hebrew no-hyphen paragraph rule.
\addto\extrashebrew{\exhyphenpenalty=10000\relax}
\addto\extrasenglish{\exhyphenpenalty=50\relax}
\ifx\lakitmainlanguage\hllpdfenglish
  \exhyphenpenalty=50\relax
\else
  \exhyphenpenalty=10000\relax
\fi
% Full English paragraphs need ordinary explicit-hyphen breaks. Use this
% environment instead of a raw otherlanguage block; the paragraph ends while
% the local penalty is still active.
\newenvironment{englishprose}
  {\par\begin{otherlanguage}{english}\exhyphenpenalty=50\relax}
  {\par\end{otherlanguage}}
\newenvironment{hebrewprose}
  {\par\begin{otherlanguage}{hebrew}\exhyphenpenalty=10000\relax}
  {\par\end{otherlanguage}}
% Use for shell text whose locale must not depend on the last body paragraph.
\DeclareRobustCommand{\doclang}[1]{\foreignlanguage{\lakitmainlanguage}{#1}}

% ============================================================================
% ★★ — fixes for defects that are INVISIBLE in the log and in the checks
% and only surface on a VISUAL RENDER of the PDF. (Each verified empirically.)
% ----------------------------------------------------------------------------
% v3.22: \emph is now REAL slanted emphasis (FakeSlant italic face defined on
% the \babelfont{rm} line above), replacing the v3.21 \emph->\textbf hack: it
% preserves \emph's toggle semantics and keeps bold free for its own role.
% Prefer bold-as-emphasis for your doc? one line: \renewcommand{\emph}[1]{\textbf{#1}}
% page number: force the digits LTR so they never reverse (18 -> 81) in the RTL
%   footer. That reversal is a NON-DETERMINISTIC bidi glitch (observed once in
%   ~80 pp); this LuaLaTeX build is not byte-reproducible, so the defensive
%   LTR wrap is the only reliable guard. Latin digits are always LTR anyway.
\renewcommand{\thepage}{\foreignlanguage{english}{\arabic{page}}}
% v3.70: \pagenumbering \gdef's \thepage and silently DESTROYS the wrap above —
% after \pagenumbering{arabic} the folio guard is gone for the whole rest of the
% document. Two collections discovered this independently and hand-restore the
% wrap in main.tex with the same comment; that knowledge now lives here instead.
% The wrapper re-applies the LTR wrap after the kernel call, for any style.
\makeatletter
\let\lakin@orig@pagenumbering\pagenumbering
\renewcommand{\pagenumbering}[1]{%
  \lakin@orig@pagenumbering{#1}%
  \renewcommand{\thepage}{\foreignlanguage{english}{\csname @#1\endcsname\c@page}}}
\makeatother
% emergencystretch: Hebrew does not hyphenate, so an unbreakable LTR island
%   (\en{...} / an inline math run) landing at the margin can leave a small
%   (~2-13pt) overfull with no legal breakpoint. This rescues ONLY those
%   inline-prose tie-breaks; it is a ceiling used only on offending lines and
%   does NOT touch display-math / listing / table overflow (those still surface
%   in the log for the content-level overflow pass). Verified: baseline overfull
%   -> 0 with this set, with no change to pagination or underfull.
\setlength{\emergencystretch}{3em}
% v3.40: inline math (almost) never breaks across lines. TeX's default
%   (binop 700 / rel 500) lets an inline run split at an operator —
%   $(1.f)\cdot 2^{e}$ shipped as "(1.f)·" ending one line and "2^e" opening
%   the next (reader-reported). In LTR that split is merely ugly; under an
%   RTL base the carried operator lands on the visually WRONG side of the
%   break and the formula reads backwards.
%   WHY 10000 (v3.45.0 — the display-first resolution). Inline math is
%   FORBIDDEN from breaking, so a torn island under an RTL base is
%   structurally impossible. The v3.40 cycle used 9999 ("break only as a last
%   resort") because a bare 10000 left ~30-39 WIDE inline equations
%   overflowing (worst 90.7pt) — but 9999's "last resort" fires OFTEN in
%   tight, non-hyphenating Hebrew: a reader found torn islands across the
%   fleet (~150 candidates; e.g. "P(y |" ending one line and "x_0)" opening
%   the next, read backwards). The step the 9999 cycle missed: every one of
%   those ~30 overflows is a wide equation that BELONGS in display. Convert
%   them to display FIRST, then 10000 gives 0 tears AND 0 overflows (measured,
%   12/12 fleet). Doctrine: inline math is ATOMIC; a wide equation goes in a
%   display, never inline. See content-style ("inline math is atomic").
\binoppenalty=10000 \relpenalty=10000
% v3.38: uniform sentence spacing. TeX's default (\nonfrenchspacing, sfcode
% 3000) is the OLD Anglo-American wide-space-after-period convention — Hebrew
% typography never had it, and modern Latin practice (Chicago et al.) dropped
% it too. Measured on real pages: post-period gaps up to 1.9-2.1x the line's
% mean word gap — the dominant "rivers" source in justified Hebrew prose.
% Survives babel language switches (probed: sfcode stays french inside \en{}
% islands and otherlanguage blocks).
\frenchspacing
% v3.38: widow/orphan protection. LaTeX's default penalties (150/150/50)
% allow a lone line at a page top/bottom almost freely. 10000 = forbid;
% this is SAFE here precisely because the document is \raggedbottom (the
% demoted line just moves on, no vertical stretch fight) — under flushbottom
% these values would cause underfull pages. raggedbottom is the report-class
% default; made explicit so a future class/option change cannot silently
% flip it and turn the 10000s harmful.
\raggedbottom
\widowpenalty=10000 \clubpenalty=10000 \displaywidowpenalty=10000
% em-dash: --- does NOT become — under Frank Ruhl (the font lacks tlig; verified
%   the width is identical with/without Ligatures=TeX, while the mechanism works
%   on Latin Modern). TYPE THE CHARACTER — (U+2014) DIRECTLY; likewise – (U+2013).
%   check_bidi_figures.py flags a --/--- that is adjacent to Hebrew.
% ============================================================================

% ---- STRUCTURAL DIRECTION-SAFETY for \texttt (v3.7) -------------------------
% A parenthesis, bracket, or $ INSIDE \texttt{} mirrors/migrates in RTL prose
% (MIPS 0($s0)->"(s0$)0", $rs->"rs$", arr[i]->"[arr[i") because \texttt sets the
% font but NOT the direction. Redefining \texttt to LTR-isolate its content fixes
% the WHOLE class at once (parens, brackets, $) so plain \texttt{lw $t0,0($s0)}
% is safe — no need to remember \code{}. Guarded by \ifmmode because
% \foreignlanguage is invalid in math (there, plain monospace is already LTR).
% NB (corrected v3.56 — the old note here said "stays tofu", which is WRONG):
% \texttt{Hebrew} is a FATAL build failure. The tt family is DejaVu Sans Mono
% (zero Hebrew) and the tt fallback hebfb is DejaVuSerif (also zero Hebrew), so
% the subsetter gets an empty glyph set and LuaTeX dies with
%   "! error: (file .../DejaVuSansMono.ttf) (type 2): there are no glyphs in the
%    subset"  -> no output PDF.
% The message names a FONT FILE, never Hebrew/\texttt/the source line. And it is
% non-local: add an unrelated Latin \code{} elsewhere and the fatal degrades into
% a SILENT DELETION of the Hebrew word. Put Hebrew in \h{}/normal text, never in
% \texttt{}/\code{}; flagged HEBREW-IN-MONO. To disable, delete this block.
\let\origtexttt\texttt
\protected\long\def\texttt#1{\ifmmode\origtexttt{#1}\else\foreignlanguage{english}{\origtexttt{#1}}\fi}
% v4.8.0: defer the direction test until TeX has inserted an alignment cell's
% math template. Unprotected expansion here tested \ifmmode too early at the
% FIRST token of aligned/array/cases/matrix cells. Engine-level \protected also
% preserves the original \texttt implementation without the self-recursion
% observed when retrofitting \DeclareRobustCommand to its saved alias.
\protected\long\def\code#1{\ifmmode\origtexttt{#1}\else\foreignlanguage{english}{\origtexttt{#1}}\fi}

% ---- \codeline — display-code line (v3.24): the code analogue of display math.
% Reader-friction rule (content-style.md §5b): inline \code{} is for a SINGLE
% ATOM only (one register/variable/mnemonic, a tight no-space expression).
% A multi-token expression — an instruction with operands, a spaced assignment,
% a register-binding list — zigzags the eye RTL<->LTR mid-sentence and, being
% an unbreakable island, line-breaks badly. Give it its own centered line with
% display-math vertical rhythm. \centerline = one box => the expression never
% wraps; \nopagebreak keeps it glued to its introducing sentence.
%   ...the Hebrew sentence ends with a colon:
%   \codeline{lw \$t1, -4(\$t0)}
%   and the prose continues (escape dollars as \$).
% v3.66 (first line of the body): flush a PENDING \item label before the \par.
% LaTeX defers the label box (\@labels) until the next \everypar, and an \hbox in
% vertical mode never fires \everypar — so `\item \codeline{...}` printed the code
% at y=69.1 and the item's own "1." BELOW it at y=103.6 (the number appears after
% its content, attached to nothing). \if@inlabel\leavevmode\fi is CONDITIONAL on
% purpose: an unconditional \leavevmode would open an empty paragraph line when
% \codeline follows ordinary prose. Verified both ways — in a list the label now
% sits on the code's line at its normal x, and after prose the geometry is
% byte-identical to before.
\makeatletter
\newcommand{\codeline}[1]{%
  \if@inlabel\leavevmode\fi
  \par\nopagebreak
  \vspace{\abovedisplayskip}\nopagebreak
  % v3.62: was \centerline{...}, i.e. \hbox to\hsize{\hss #1 \hss}. \hss has
  % INFINITE SHRINK, so TeX could always "fit" the box and NEVER reported an
  % Overfull — the same \hss mechanism as BLD-27's \item label. Measured: a wide
  % \codeline in a defbox painted its first token ON the box's left border bar and
  % sat 9.5pt outside the frame interior on both sides, with 0 Overfull, 0 Underfull
  % and all four gates PASS; in bare prose it reached 21.0pt into BOTH margins.
  % \hfil has no shrink, so the overflow is reported and build.sh's ovf_vis/ovf_big
  % counters see it; \linewidth (not \hsize) makes it box-aware, so a \codeline
  % inside a colored box centres on the box interior. Verified: a FITTING \codeline
  % keeps byte-identical glyph positions, and the box-crossing case now emits
  % "Overfull \hbox (19.03pt too wide)". A wide \codeline should be split in two.
  \hbox to\linewidth{\hfil\code{#1}\hfil}%
  \nopagebreak\vspace{\belowdisplayskip}}
\makeatother

% ---- a thin decorative rule helper ----
\newcommand{\sep}{\par\vspace{2pt}\noindent\textcolor{cRule}{\rule{\linewidth}{0.6pt}}\par\vspace{4pt}}

% ---- colors ----
\definecolor{cBlue}{HTML}{1F4E79}   \definecolor{cBlueBg}{HTML}{EAF1F8}
\definecolor{cGreen}{HTML}{2E6B4F}  \definecolor{cGreenBg}{HTML}{EAF4EE}
\definecolor{cAmber}{HTML}{9C6A00}  \definecolor{cAmberBg}{HTML}{FBF3E2}
\definecolor{cPurple}{HTML}{5B2C6F} \definecolor{cPurpleBg}{HTML}{F3ECF7}
\definecolor{cTeal}{HTML}{0E6E6E}   \definecolor{cTealBg}{HTML}{E6F4F4}
\definecolor{cRed}{HTML}{8E2B2B}    \definecolor{cRedBg}{HTML}{FBEAEA}
\definecolor{cGray}{HTML}{555555}   \definecolor{cRule}{HTML}{1F4E79}

% ---- math operators ----
\DeclareMathOperator{\Tr}{Tr}
\DeclareMathOperator{\rank}{rank}
\DeclareMathOperator{\diag}{diag}
\DeclareMathOperator{\sgn}{sgn}

% ---- bra-ket (quantikz defines \ket/\bra/\braket; add the rest safely) ----
\providecommand{\ket}[1]{\left|#1\right\rangle}
\providecommand{\bra}[1]{\left\langle#1\right|}
\providecommand{\braket}[2]{\left\langle#1\middle|#2\right\rangle}
\newcommand{\dyad}[1]{\left|#1\right\rangle\!\left\langle#1\right|}
% v3.21: general outer product |a><b| and expectation value <A> — agents reach
% for these by reflex (esp. \expval); quantikz defines neither, so their absence
% threw "Undefined control sequence" in a real build. \providecommand = safe.
\providecommand{\ketbra}[2]{\left|#1\right\rangle\!\left\langle#2\right|}
\providecommand{\expval}[1]{\left\langle#1\right\rangle}
\newcommand{\abs}[1]{\left|#1\right|}
\newcommand{\norm}[1]{\left\lVert#1\right\rVert}
\newcommand{\half}{\tfrac12}
\newcommand{\R}{\mathbb{R}}\newcommand{\C}{\mathbb{C}}\newcommand{\Z}{\mathbb{Z}}
% identity operator. NB: \mathbb{1} silently renders a WRONG glyph (msbm has no
% blackboard digits). Use \mathbb{I} = 𝕀 (works with amssymb, no extra package).
% Do NOT "fix" this back to \mathbb{1}.
\newcommand{\id}{\mathbb{I}}

% ---- physical units (v3.35): \punit — upright unit with thin space ---------
% Usage: $300\punit{K}$, $80\punit{mV/decade}$, $10^{7}\punit{cm/s}$ — the
% NUMBER AND UNIT LIVE IN ONE MATH ISLAND. Promoted verbatim from a 199-page
% production lakit (advanced-electronic-devices, dozens of uses).
% Two measured traps this design avoids (both render-proven):
%   1. siunitx SCRAMBLES under RTL — \SI{3.5}{eV} renders "eV 3.5", exponents
%      detach ("Hz 10 x 10^9"): its quantity assembly is not bidi-isolated.
%      Do NOT load siunitx in this pipeline.
%   2. A text-mode unit next to a bare-prose number is TWO adjacent LTR runs
%      and reorders ("5 mm" -> "mm5" — the BIDI-2/3 class). \punit is therefore
%      math-only ON PURPOSE: prose misuse fails LOUDLY at compile (\mathrm
%      outside math) instead of scrambling silently.
% Domain shorthands (\MeV, \fm...) stay per-document: \newcommand{\MeV}{\,\text{\en{MeV}}}
\providecommand{\punit}[1]{\,\mathrm{#1}}

% ============================================================================
% Colored boxes — mdframed (NOT tcolorbox, which breaks under bidi).
% ★★ v3.4: innerleftmargin/innerrightmargin give code listings room so their
% frame/line-numbers do not spill out the box border (see listing styles below).
% ============================================================================
% v3.61: everyline/split skips were applied to `tablebox` only (BLD-14, v3.37) —
% the six colored boxes below never got them, so they still split OPEN across a page
% break: the page-N segment gets top+sides and NO bottom rule, the page-N+1 segment
% sides+bottom and NO top rule. Measured on the shipped fleet this is a COVERAGE GAP,
% not a re-report of BLD-14: 0 split-open `tablebox` segments (that fix holds) vs 62
% split-open COLORED boxes across 7 documents. Verified after the fix: both segments
% close, and the frametitle is NOT duplicated on the continuation. BLD-14's caveat
% still applies — a box taller than a full page grinds mdframed's split-trial loop, so
% split such content manually.
% ============================================================================
% v3.72: frametitle={} leads EVERY \newmdenv option list below. mdframed options
% are GROUP-scoped: a box nested inside \begin{exbox}[frametitle={...}] used to
% INHERIT the outer title — measured, the identical title string rendered twice,
% the second riding a title-bar fill inside the INNER box's frame, for all seven
% environments in both directions. Both natural forms were broken (= preamble
% bug, not an authoring rule). The explicit reset makes every box titleless
% unless ITS OWN [frametitle={...}] says otherwise; verified both directions:
% no duplicate on the nested box, an explicit title still renders.
% Geometry is shared with the optional semantic-width table container. Themes
% select their own existing semantic style below; neither path fixes columns.
\mdfdefinestyle{hllpdfsemanticgeometry}{linewidth=1.1pt,
  innertopmargin=6pt,innerbottommargin=7pt,skipabove=9pt,skipbelow=9pt,
  innerleftmargin=9pt,innerrightmargin=9pt,everyline=true,
  splitbottomskip=6pt,splittopskip=12pt,roundcorner=3pt}
\newcommand{\hllpdfsemanticframestyle}{hllpdfsemanticgeometry}
\newmdenv[frametitle={},style=hllpdfsemanticgeometry,linecolor=cBlue,backgroundcolor=cBlueBg,
  frametitlebackgroundcolor=cBlue,frametitlefontcolor=white,
  ]{defbox}   % הגדרה
\newmdenv[frametitle={},style=hllpdfsemanticgeometry,linecolor=cGreen,backgroundcolor=cGreenBg,
  frametitlebackgroundcolor=cGreen,frametitlefontcolor=white,
  ]{thmbox}   % משפט/טענה
\newmdenv[frametitle={},style=hllpdfsemanticgeometry,linecolor=cAmber,backgroundcolor=cAmberBg,
  frametitlebackgroundcolor=cAmber,frametitlefontcolor=white,
  ]{notebox}  % הערה
\newmdenv[frametitle={},style=hllpdfsemanticgeometry,linecolor=cPurple,backgroundcolor=cPurpleBg,
  frametitlebackgroundcolor=cPurple,frametitlefontcolor=white,
  ]{exbox}    % דוגמה
\newmdenv[frametitle={},style=hllpdfsemanticgeometry,linecolor=cRed,backgroundcolor=cRedBg,
  frametitlebackgroundcolor=cRed,frametitlefontcolor=white,
  ]{warnbox}  % טעות נפוצה
\newmdenv[frametitle={},style=hllpdfsemanticgeometry,linecolor=cTeal,backgroundcolor=cTealBg,
  frametitlebackgroundcolor=cTeal,frametitlefontcolor=white,
  ]{keybox}   % רעיון מרכזי

% ---- codebox (v4.5.0): the dedicated container for code -------------------
% An OPTION, not a policy: asmblock unboxed stays canonical for a full listing,
% and a snippet inside any of the six colored boxes still works as before.
% What makes codebox worth its own environment is that it is the ONE framed box
% that keeps its margin line numbers.
%
% v4.4.0 shipped it WITHOUT them, on the strength of v3.37's "no margin knob can
% save them" plus four fresh probes that agreed. That was wrong: every probe had
% moved the code, the box, or both — never the key that carries the NUMBER
% COLUMN. Measured on the page:
%   resetmargins=false   the listing follows THIS box instead of the page,
%                        moving the code from x=75.26 to x=97.54
%   xrightmargin=26pt    under RTL this is the PHYSICAL LEFT indent, and the
%                        only key that carries the \llap'ed digits with it:
%                        x=64.28 -> 90.18, clearing the rule at 65.64 by 23.6pt
%   xleftmargin=26pt     moves the code and LEAVES THE DIGITS on the rule —
%                        the asymmetry that hid the fix for eleven versions
% resetmargins is hoisted out of the four styles (see the \lstset below them) so
% this override is possible at all; every other listing still resolves to true.
% Pinned by tests/codebox_numbers.tex + probe_lstnum.py --inside-rule, whose
% negative control strips xrightmargin and watches the digits fall back.
\definecolor{cSlate}{HTML}{3F4A56}  \definecolor{cSlateBg}{HTML}{F2F4F7}
\newmdenv[frametitle={},leftline=true,rightline=true,topline=true,bottomline=true,
  linecolor=cSlate,backgroundcolor=cSlateBg,linewidth=0.9pt,
  frametitlebackgroundcolor=cSlate,frametitlefontcolor=white,
  innertopmargin=6pt,innerbottommargin=7pt,skipabove=9pt,skipbelow=9pt,
  innerleftmargin=9pt,innerrightmargin=9pt,everyline=true,splitbottomskip=6pt,splittopskip=12pt,
  roundcorner=3pt]{codebox}   % קוד
% The two keys that make in-box numbering work, both measured on the page:
%   resetmargins=false  the listing follows THIS box instead of the page
%   xrightmargin        under RTL this is the PHYSICAL LEFT indent, and it is
%                       the only one that carries the \llap'ed number column
%                       with it. xleftmargin moves the code and leaves the
%                       digits behind on the rule (measured: code 97.54,
%                       digits 64.28 — the v3.37 defect, reproduced); with
%                       xrightmargin=26pt the digits land at x=90.18, a
%                       clearance of 23.6pt from the rule at 65.64.
% This asymmetry is why v3.4's note had the semantics backwards and v3.37
% concluded the digits could not be saved.
% <hllpdf:listing-container-state>
% v4.29: all state is local to the current environment group. A semantic
% callout masks number text, while a nested codebox restores the formatter
% saved before that mask. Save it once across successive semantic wrappers;
% an explicit numbers=none (or an explicitly empty formatter) remains valid.
\makeatletter
\newif\ifhllpdf@incodebox
\newif\ifhllpdf@englishcodebox
\newif\ifhllpdf@lstnumbermasked
\let\hllpdf@savedlstnumber\thelstnumber
\makeatother
% </hllpdf:listing-container-state>
% <hllpdf:codebox-gutter>
\makeatletter
\AtBeginEnvironment{codebox}{%
  \hllpdf@incodeboxtrue
  % Snapshot the OUTER language before asmblock selects English.
  \iflanguage{english}{\hllpdf@englishcodeboxtrue}{\hllpdf@englishcodeboxfalse}%
  \ifhllpdf@lstnumbermasked
    \let\thelstnumber\hllpdf@savedlstnumber
    \hllpdf@lstnumbermaskedfalse
  \fi
}
% listings' AddToHook is GLOBAL. Register once, and gate on local box state;
% adding an unconditional hook at every codebox entry leaks into later listings.
\lst@AddToHook{PreInit}{%
  \ifhllpdf@incodebox
    % Apply every listing override here, inside the LISTING's own group.
    % Setting margins at box entry leaks them into nested semantic callouts.
    \lstset{resetmargins=false,xrightmargin=26pt}%
    % Styles resolve before PreInit. In LTR the physical number gutter follows
    % xleftmargin; in Hebrew retain the measured inherited xrightmargin path.
    \ifhllpdf@englishcodebox
      \lstset{xleftmargin=26pt,xrightmargin=0pt}%
    \fi
    % ONE box, not two: suppress only the frame/background/skips of a listing
    % inside this codebox. The next unboxed listing keeps its own style.
    \lstset{frame=none,backgroundcolor=\color{cSlateBg},aboveskip=0pt,belowskip=0pt}%
  \fi
}
\makeatother
% </hllpdf:codebox-gutter>

% v3.37: line numbers inside a COLORED BOX are sliced by the box border —
% measured: digit span x=63.5-67.3 vs border line at x=65.2 (the border runs
% THROUGH the digits), and the inter-border gap (3.5pt) is narrower than a
% digit (3.8pt), so no margin knob can save them (under resetmargins,
% xleftmargin is inert in-box — measured: code at x=75.26 with 2.2em AND
% 3.0em). A 3-6 line boxed snippet does not need line numbers; blank the
% number TEXT inside every colored box (\thelstnumber is what gets printed;
% the style's numbers=left survives, so nothing else shifts — code position
% byte-frozen at x=75.26, verified). Unboxed listings keep their numbers.
% v3.78: the blanking is coupled to the FRAME, not to the box name —
% clean_theme.tex registers LATER hooks (registration order wins) that
% restore the number text inside the boxes it renders frameless.
% v4.5.0 RETRACTS the "no margin knob can save them" clause above. It holds only
% under resetmargins=true (the shipped default): there xleftmargin is inert
% in-box, as measured. With resetmargins=false, xrightmargin DOES carry the
% \llap'ed digit column clear of the rule — measured; the codebox \lstset that
% does it is the one above, in the codebox block. The loop below therefore
% covers SIX boxes, not "every colored box": codebox opts out and keeps its
% numbers.
% <hllpdf:semantic-listing-mask>
\makeatletter
\@for\HEB@lstbox:=defbox,thmbox,notebox,exbox,warnbox,keybox\do{%
  \expandafter\AtBeginEnvironment\expandafter{\HEB@lstbox}{%
    \ifhllpdf@lstnumbermasked\else
      \let\hllpdf@savedlstnumber\thelstnumber
    \fi
    \hllpdf@lstnumbermaskedtrue
    \def\thelstnumber{}%
    % The nearest semantic container owns its listing presentation. On exit,
    % TeX's group restores the enclosing codebox state and formatter.
    \hllpdf@incodeboxfalse}}
\makeatother
% </hllpdf:semantic-listing-mask>

% ============================================================================
% Unified TABLE presentation: subtle uniform slate frame + auto-numbered
% caption rendered INSIDE the frame. Use:
%   \begin{tablebox}\centering  <tabular>  \tabcap{<description>}\end{tablebox}
% Replaces the buggy \caption* (which emitted a stray "טבלה N : *" line) and
% provides an explicit framed option + numbered caption "טבלה N.M: תיאור".
% Optional intro text may be placed at the TOP, before the tabular, inside the box.
% ============================================================================
\definecolor{cTblLine}{HTML}{B9C4D2}  \definecolor{cTblBg}{HTML}{FAFBFC}
% v3.23: class-aware — under `article` there is no chapter counter (the promised
% report->article switch used to die with "No counter 'chapter' defined").
\makeatletter
\@ifundefined{c@chapter}{%
  % v3.31: in article the counter resets per SECTION, so an unprefixed number
  % repeated ("טבלה 1." in every section — measured on a two-section render).
  % Same conditional idea as the report branch: prefix the section number once
  % a real \section exists; a pre-section table stays "טבלה 1".
  \newcounter{tbl}[section]
  \renewcommand{\thetbl}{\ifnum\value{section}>0 \thesection.\fi\arabic{tbl}}%
}{%
  \newcounter{tbl}[chapter]
  % v3.30: a report-class doc that uses SECTIONS ONLY (chapter stays 0 —
  % the common short-doc pattern) used to print "טבלה 0.1"; the chapter
  % prefix now appears only once a real \chapter exists (render-verified
  % both ways: "טבלה 1." chapterless, "טבלה 1.1" with chapters).
  \renewcommand{\thetbl}{\ifnum\value{chapter}>0 \thechapter.\fi\arabic{tbl}}%
  % v3.64: the SAME leak in \thesection itself, left behind by v3.30. That fix
  % guarded the CAPTION counter and stopped there, so a chapterless report still
  % printed the heading "0.1 סעיף ראשון" and \ref to it printed "0.2" — measured:
  %   p1 y=65.2  '…0.1'   p1 y=144.9 '…0.2'   p1 y=183.4 '…0.2.1'
  % report.cls hardcodes \thesection as \thechapter.\@arabic\c@section, and the
  % leading 0 is the chapter counter that no \chapter ever incremented. Fixing
  % \thesection fixes the whole chain — \thesubsection and \thesubsubsection are
  % defined in terms of it. \theequation and \thetbl/\thefig were already guarded,
  % which is exactly why this survived: every NUMBER around the heading looked
  % right, so the heading's own 0 read as a style choice rather than a bug.
  % \appendix is safe: report.cls resets chapter to 0 and redefines \thechapter to
  % \@Alph\c@chapter, so an appendix \chapter makes the counter 1 and the guard
  % emits "A.1" unchanged. Render-verified in all four modes (chapterless,
  % chaptered, appendix+chapter, appendix+section-only).
  \renewcommand{\thesection}{\ifnum\value{chapter}>0 \thechapter.\fi\arabic{section}}%
}
\makeatother
% v3.37: everyline=true — if the box DOES split across pages, every segment
% gets a CLOSED frame (measured before: the split ran open-ended into the page
% edge with no bottom border, and the caption could strand alone at the next
% page top — "טבלה A.2" shipped that way). The split skips give the segments
% breathing room. In practice, with the caption GLUE below, every previously-
% split table in the corpus now moves WHOLE instead (the white-gap tradeoff,
% documented benign); everyline stays as the safety net for a genuinely
% page-tall table. NB: a table taller than ONE FULL PAGE can grind mdframed's
% split-trial loop (reproduced: minimal doc, minutes of 1pt-step Overfull
% probes) — break such a monster into two tables.
\mdfdefinestyle{hllpdftablegeometry}{leftline=true,rightline=true,topline=true,bottomline=true,
  linewidth=0.9pt,
  innertopmargin=9pt,innerbottommargin=8pt,skipabove=11pt,skipbelow=11pt,
  innerleftmargin=11pt,innerrightmargin=11pt,roundcorner=4pt}
\newcommand{\hllpdftableframestyle}{hllpdftablegeometry}
\newmdenv[frametitle={},style=hllpdftablegeometry,
  linecolor=cTblLine,backgroundcolor=cTblBg,
  % v4.2.0 (adopted): tablebox is ATOMIC. Short tableboxes near page bottoms
  % repeatedly entered mdframed's split trial, emitting benign Overfull-vbox
  % bursts and occasionally separating the intro prose from the tabular.
  % nobreak=true moves the complete short table to the next page instead. The
  % output-active Overfull-vbox gate still catches a box taller than a page;
  % everyline/split skips remain as defensive styling should that path ever
  % re-arm, but the canonical path never publishes a split tablebox.
  nobreak=true,
  everyline=true,splitbottomskip=6pt,splittopskip=12pt,
  % v3.79 (BLD-37 table side, correction): startinnercode/endinnercode run
  % INSIDE mdframed's body vbox (mdframed.sty:848 / :856+863) -- the second one
  % after \par\unskip, i.e. after line breaking, which is the only moment the
  % \tabcap anchor's true offset from the top of the box EXISTS. See
  % \lakintabboxclose below. These two keys are the ONLY tablebox change.
  startinnercode={\lakintabboxopen},endinnercode={\lakintabboxclose}]{tablebox}
% v3.37: \nopagebreak GLUE — \vsplit honors penalties, so the caption can no
% longer strand alone at the top of the next page (it takes its rows along, or
% the whole box moves). Render-verified on the real shipped orphan.
% v3.59: \NR@gettitle is what nameref injects into the stdlib title-bearing
% commands (\@sect, \@chapter, \caption) to fill \@currentlabelname. \figcap and
% \tabcap are INDEPENDENT macros — they never called it, so \@currentlabelname kept
% whatever the last heading left and \nameref{fig:x} silently printed the enclosing
% SECTION's title instead of the caption (proven in the .aux: field 3 of \newlabel
% was the heading text while field 1 and the anchor `fig.1.1` were correct — so this
% is a THIRD, distinct head next to BLD-13's counter fix and BIDI-38's autoref-name
% fix, neither of which writes that field). \@ifundefined keeps it inert if
% hyperref/nameref is ever dropped. Guarded helper because \figcap is defined
% after \makeatother.
% NB both helpers below need \makeatletter — this part of the preamble is OUTSIDE
% any letter-catcode region, so an unwrapped \@currentlabel parses as LaTeX's \@
% and dies with "You can't use \spacefactor in vertical mode" (measured).
\makeatletter
\newcommand{\lakinrtitle}[1]{\@ifundefined{NR@gettitle}{}{\NR@gettitle{#1}}}
% v3.59: \hitem{<letter>}{<key>} — a REFERENCEABLE Hebrew-lettered list item.
%   \begin{enumerate} \hitem{א}{it:a} הסעיף הראשון. \end{enumerate}  ->  \ref{it:a} = (א)
% Why this exists: automatic Hebrew list labels fail silently, so SKILL.md mandates
% an EXPLICIT \item[(\textbf{א})] — but LaTeX puts \refstepcounter INSIDE
% \if@noitemarg (latex.ltx:12964), and \item[...] never sets that flag. So an
% explicit item steps no counter and never touches \@currentlabel: a \label there
% binds to the AMBIENT one and \ref prints the CHAPTER/SECTION number where the
% item letter belongs (proven in the .aux: {chapter.1}). Same doctrine as \figcap —
% own the label state before \label. \phantomsection matters: without it the
% hyperlink still lands on the chapter head (measured); with it the anchor is
% section*.N and the bookmark outline stays clean.
\newcommand{\hitem}[2]{\item[(\textbf{#1})]%
  \def\@currentlabel{(#1)}\def\@currentlabelname{#1}\phantomsection\label{#2}}
% NB (v3.66): \hitem is \ref-only. \autoref{key} calls it a SECTION ("סעיף (א)")
% because \phantomsection brands the anchor section*.N, and \nameref{key} prints
% the bare letter. For \autoref-able items prefer \hebenumlabels below.
%
% v3.66: \hebenumlabels — OPT-IN automatic Hebrew enumerate letters, done right.
%   {\hebenumlabels\begin{enumerate}\item\label{q:a} ...\end{enumerate}}
%   -> labels (א) (ב) (ג); \ref{q:a} = א; anchor Item.N (real, per-item).
% This retires the belief that made \item[(\textbf{א})] + \hitem necessary.
% "Automatic Hebrew labels fail silently" was TRUE but unexplained, and the
% mechanism is now bisected (plain article vs +enumitem, measured): ENUMITEM
% PRE-EXPANDS the label expression in an \edef BEFORE \refstepcounter, with the
% counter still at 0. Tokens it shields by name (\arabic,\alph,\roman) survive;
% a plain \newcommand \ifcase wrapper evaluates at 0 -> the EMPTY zero-branch ->
% the label prints as "()" with the letter GONE (and \the\c@enumi prints a
% literal 0, \number\value dies fatally pointing at \item). The cure is one
% property, not a doctrine: the wrapper must be ROBUST (survives the \edef
% unexpanded), while \theenumi stays EXPANDABLE (so \refstepcounter writes the
% real letter to the .aux). Verified: labels (א)(ב)(ג), \ref -> א/ב/ג each item
% its own value, .aux = {א}{...}{Item.1} — versus the explicit-\item doctrine
% where BIDI-50/\hitem exist purely to patch the counter that never stepped.
\def\@hebalph#1{\ifcase#1\or א\or ב\or ג\or ד\or ה\or ו\or ז\or ח\or ט\or י\or
  יא\or יב\or יג\or יד\or טו\or טז\or יז\or יח\or יט\or כ\else\@ctrerr\fi}
\DeclareRobustCommand*\hebalph[1]{\expandafter\@hebalph\csname c@#1\endcsname}
\newcommand{\hebenumlabels}{%
  \renewcommand{\theenumi}{\@hebalph\c@enumi}%
  \renewcommand{\labelenumi}{(\textbf{\hebalph{enumi}})}}
\makeatother
% v3.63: \mbox the LABEL. \thetbl ("7.1") and a caption description that STARTS
% WITH A DIGIT ("32 הרגיסטרים…") are two European-Number runs separated only by
% neutrals (the label's own "." and \enspace). Under the RTL paragraph base the UBA
% resolves those neutrals to the RTL level, so both numbers sit in ONE reversed span
% and EXCHANGE VISUAL POSITIONS. Measured on the shipped PDF: a table whose real
% number is 7.1 printed "טבלה 32" — a different, EXISTING table number — while \ref
% correctly printed 7.1, so the caption silently contradicted its own cross-reference,
% and the description lost its leading count. The source contains no adjacency at all;
% it is manufactured by this macro, which is why every source linter was blind.
% \mbox makes the label one atomic box so nothing can reorder across it. Verified
% zero-regression: an ordinary Hebrew caption keeps identical x-geometry.
% One paragraph style covers label, number, separator and description.
% Explicit project/template overrides remain possible; the label hook is kept
% for compatibility, with an identity default. Counters/anchors stay unchanged.
\newcommand{\lakitcaptionstyle}{\normalfont\small\color{black}\hypersetup{allcolors=black}}
\newcommand{\lakitcaptionlabelstyle}[1]{#1}
\newcommand{\lakitcaptionpunctuation}{:}
% Native floats, captionof and longtable use the same presentation policy.
% subcaption is optional and inherits the same font rather than a second size.
\usepackage{caption}
\DeclareCaptionFont{lakit}{\lakitcaptionstyle}
\DeclareCaptionLabelFormat{lakit}{\lakitcaptionlabelstyle{#1~\en{#2}}}
\DeclareCaptionLabelSeparator{lakit}{\lakitcaptionpunctuation\enspace}
\captionsetup{font=lakit,labelfont={},textfont={},labelformat=lakit,
  labelsep=lakit,justification=centering}
\makeatletter
\AtBeginDocument{\@ifpackageloaded{subcaption}{%
  \captionsetup[sub]{font=lakit,labelfont={},textfont={},labelsep=lakit}}{}}
\makeatother
% v3.64: OPT-IN list-of-figures/tables entry via an optional short caption.
% \listoffigures and \listoftables are inherited from report.cls and LOOK available,
% but \figcap/\tabcap step the custom \c@fig/\c@tbl counters and never call
% \addcontentsline — so both printed their heading and NOTHING ELSE. Measured: a doc
% with one \figcap and one \tabcap rendered page 1 as 'רשימת האיורים\n1\n', and the
% .lof held only \@chapter's \addvspace, zero \contentsline entries. An empty list
% under a correct Hebrew heading reads as "this document has no figures".
%
% WHY OPTIONAL, AND NOT AUTOMATIC. Writing #1 unconditionally was implemented first
% and REJECTED: a caption containing \codeline compiles CLEANLY today (verified on
% the pre-v3.64 preamble: 0 errors, PDF produced) but becomes FATAL once it is also
% written to the .lof — \codeline is \hbox to\linewidth{...}, and a full-width box in
% a list entry dies with "\pdfextension endlink cannot be used in vertical mode", no
% output PDF. Gaining a feature nobody uses (0 of 13 collections call \listoffigures)
% at the price of breaking a construct that works is a bad trade. Stock \caption has
% exactly this limitation and exactly this escape hatch, so the convention is familiar.
% \en{}, \code{}, \h{} and inline math in the short text are all safe — verified, they
% survive the write and typeset correctly in the list (RTL geometry confirmed: number
% rightmost at x=509.9, title left of it, dots and folio at the far left).
% With no optional argument the macro is byte-identical to v3.63 — zero risk to every
% existing document, and \listoffigures lists exactly the captions that opted in.
% BLD-37, TABLE side (the \figcap v3.74 idea, adapted): clicking a \ref to a
% table used to land MID-table — hyperref's automatic anchor inside
% \refstepcounter was emitted on the still-open horizontal line whose TAIL is
% the tabular (an hbox whose baseline is the table's vertical CENTER), so the
% jump scrolled the top half of the table above the viewport. The v3.74 entry
% deferred this side claiming "tablebox splits — no single box to measure";
% that deferral dissolves: the measurement target is the TABULAR itself at
% \tabcap ENTRY, while it is still the tail of the open list — the same
% situation \lakin@grabfig handles at an island's end (and the fleet has 0
% split tableboxes). So \lakin@grabtab runs as the FIRST token, before any
% typesetting can push the box away, and the anchor is re-emitted raised above
% the measured tabular. \lakintabraise is a SEPARATE register from
% \lakinfigraise (defined with it, below the LTR-island block): a captionless
% island intentionally leaves the fig register set, and a table must never
% inherit that stale figure height.
%
% WHY the anchor stays AT THE ENTRY SPOT and is NOT deferred into the caption
% paragraph the way \figcap defers it: at \par TeX DROPS a final glue item,
% and the automatic anchor's smash box sitting BEHIND the interword glue kept
% that glue alive — so every centered tabular has always sat half a space off
% true center. Deferring the anchor removed that box and was measured as 982
% differing glyph records: EVERY tabular line shifted 1.30pt (2.61pt for
% pattern 7's two minipages) while captions stayed put. Emitting the raised
% anchor at the old spot — through hyperref's own \Hy@raisedlink, exactly the
% macro the automatic path uses — keeps the line's node structure, and so the
% A/B glyph diff at 0, BY CONSTRUCTION. \Hy@raisedlink reads the raise from
% \HyperRaiseLinkDefault (default \baselineskip = the old mid-table dest); the
% group-local \edef freezes measured-height+42pt into it, and \smash inside
% \Hy@raisedlink hides any raise from the line's extent. \HyCnt@ProvideTheH-
% Counter keeps the dest NAME byte-identical to the automatic one (it defines
% \theHtbl, which \hyper@makecurrent prefers — without it the first measured
% caption was named tbl.1 where the automatic path said tbl.0.1, measured).
% vmode entry, or anything but a box at the tail (no tabular to measure),
% falls back to the plain \refstepcounter — byte-identical to the old path.
\makeatletter
\newcommand{\tabcap}[2][]{\lakin@grabtab
  % v3.79: arm the box-close correction. Set unconditionally and in BOTH
  % branches below's scope, so the plain-\refstepcounter fallback is corrected
  % too; \lakintabboxopen clears it at every \begin{tablebox}, so a \tabcap
  % outside a tablebox can never leave it armed for a later box.
  \global\lakintabcapseentrue
  \ifdim\lakintabraise>0pt
    \H@refstepcounter{tbl}\HyCnt@ProvideTheHCounter{tbl}\hyper@makecurrent{tbl}%
    \begingroup
      % A frameless table targets its measured top. The frame's historical
      % cushion belongs only to tablebox and is corrected at box close.
      \edef\HyperRaiseLinkDefault{\the\lakintabraise}%
      \iflakintabframed
        \edef\HyperRaiseLinkDefault{\the\dimexpr\lakintabraise+42pt\relax}%
      \fi
      \Hy@raisedlink{\hyper@anchorstart{\@currentHref}\hyper@anchorend}%
    \endgroup
    \global\lakintabraise 0pt\relax
  \else
    \refstepcounter{tbl}%
  \fi
  \lakinrtitle{#2}\par\nopagebreak\addvspace{6pt}%
  \def\lakin@short{#1}\ifx\lakin@short\@empty\else
    \addcontentsline{lot}{table}{\protect\numberline{\thetbl}{\protect\ignorespaces #1}}\fi
  {\centering\lakitcaptionstyle\mbox{\lakitcaptionlabelstyle{\tablename~\en{\thetbl}\lakitcaptionpunctuation}}\enspace #2\par}}
\makeatother

% ============================================================================
% Unified FIGURE caption — the \tabcap twin for TikZ/plot figures (v3.34.0).
% A figure (tikzpic) sits inline in the prose; put \figcap BELOW it for a
% numbered, \ref-able caption "איור N.M: תיאור" in the same style as \tabcap:
%   \begin{tikzpic} ... \end{tikzpic}
%   \figcap{\label{fig:foo}description}     % \label FIRST, inside the argument
% WHY a macro, not a hand-written gray line: a manual
%   {\centering\small\textcolor{cGray}{איור: ...}} with a \label inside has NO
% \refstepcounter, so the \label binds to the AMBIENT \@currentlabel (the last
% equation) — \ref then prints an equation number AS the figure number
% (measured on a real page: "איור 48" where 48 is equation 48). \figcap steps a
% real \c@fig counter first, so \label is correct, numbering is consistent, and
% refs are stable. Class-aware like \thetbl (chapter- or section-prefixed).
% ============================================================================
\makeatletter
\@ifundefined{c@chapter}{%
  \newcounter{fig}[section]
  \renewcommand{\thefig}{\ifnum\value{section}>0 \thesection.\fi\arabic{fig}}%
}{%
  \newcounter{fig}[chapter]
  \renewcommand{\thefig}{\ifnum\value{chapter}>0 \thechapter.\fi\arabic{fig}}%
}
\makeatother
\makeatletter
\newcommand{\figcap}[2][]{%
  \ifdim\lakinfigraise>0pt
    % v3.74: step WITHOUT hyperref's automatic anchor, then emit the anchor
    % raised above the measured figure. \H@refstepcounter is hyperref's saved
    % pre-hook \refstepcounter (sets the counter AND \@currentlabel — \label
    % in the caption arg binds exactly as before); \hyper@makecurrent supplies
    % \@currentHref, so the dest NAME is identical to the old automatic one —
    % \ref/\autoref/bookmarks are unchanged, only the dest'S POSITION moves.
    \H@refstepcounter{fig}\hyper@makecurrent{fig}%
    % The anchor is emitted INSIDE the caption paragraph below, as a zero-size
    % raisebox — NOT here in vertical mode, where a \raisebox starts an empty
    % paragraph and adds a blank line (measured: every glyph after the caption
    % shifted +18.8pt). Inside an existing line a [0pt][0pt] box is inert.
    \edef\lakin@anch{\@currentHref}%
    % \edef FREEZES the measured height into the pending macro: the deferred
    % \raisebox is typeset inside the caption paragraph AFTER the register is
    % reset to the 0pt sentinel below — a plain \def read 0 at use time
    % (measured: LAKINRAISE=0.0, dest moved 14pt instead of above the figure).
    \edef\lakin@pend{\noexpand\raisebox{\the\dimexpr\lakinfigraise+42pt\relax}[0pt][0pt]{%
      \noexpand\hyper@anchorstart{\lakin@anch}\noexpand\hyper@anchorend}}%
    \global\lakinfigraise 0pt\relax
  \else
    \refstepcounter{fig}\let\lakin@pend\relax
  \fi
  \lakinrtitle{#2}\par\nopagebreak\addvspace{4pt}%
  \def\lakin@short{#1}\ifx\lakin@short\@empty\else
    \addcontentsline{lof}{figure}{\protect\numberline{\thefig}{\protect\ignorespaces #1}}\fi
  {\centering\lakin@pend\lakitcaptionstyle\mbox{\lakitcaptionlabelstyle{\figurename~\en{\thefig}\lakitcaptionpunctuation}}\enspace #2\par}}
\makeatother

% ============================================================================
% LTR islands.  EVERY tikzpicture / quantikz MUST live inside one, or the whole
% figure's coordinates and Latin labels flip. (Hebrew node text still uses \h{}.)
% ============================================================================
% v3.74 (BLD-37): the \figcap hyperlink DESTINATION used to land at the caption
% line — BELOW the figure — so clicking a figure reference scrolled the figure
% off-screen (measured fleet-wide: 97/97 fig dests in one collection sat 19-54pt
% under their figure's ink; a 347pt figure was fully hidden after the jump).
% hypcap-style fix in two parts: (1) at the island's END, \lakin@grabfig grabs
% the just-typeset centered line with \lastbox, measures it, and puts it back
% with \nointerlineskip — the interline glue above the line was already
% contributed and is NOT re-added, so the page geometry is unchanged to the
% glyph (verified); (2) \figcap then steps the counter WITHOUT hyperref's
% automatic anchor (\H@refstepcounter = the pre-hyperref \refstepcounter) and
% emits the anchor itself, raised by the measured height — the dest lands ABOVE
% the figure top. A captionless island leaves numbering untouched (nothing is
% stepped at the island; the grab flag is cleared at every island begin, so a
% stale height from a captionless figure cannot leak into a later caption).
\newsavebox{\lakinfigbox}
\newdimen\lakinfigraise  % real \dimen: tex.setdimen cannot target a \newlength SKIP register (measured lua error)
% BLD-37 table side: a SEPARATE register for \tabcap. The two must not share —
% a captionless island INTENTIONALLY leaves \lakinfigraise set (the grab runs
% at every island END; only the NEXT island's begin clears it), so a shared
% register would hand a \tabcap after a captionless figure that stale FIGURE
% height. \lakintabraise is written only by \lakin@grabtab at \tabcap entry
% and consumed/reset inside \tabcap itself, so it can never go stale.
\newdimen\lakintabraise
\makeatletter
% v3.74b: measure WITHOUT touching the list. The first draft grabbed the box
% with \lastbox and put it back — and \lastbox DISCARDS the node's baseline
% shift, so a vertically-anchored quantikz circuit rose 2.57pt and everything
% after it sank 5.32pt (measured on the A/B glyph diff). LuaTeX lets us READ
% the tail node's dimensions without removing anything: geometry is identical
% by construction. \lakinfigraise=0pt is the "no figure measured" sentinel.
\def\lakin@grabfig{%
  \ifhmode
    \directlua{
      local n = tex.nest[tex.nest.ptr].tail
      local GLUE, PEN = node.id('glue'), node.id('penalty')
      while n and (n.id == GLUE or n.id == PEN) do n = n.prev end
      local HL, VL = node.id('hlist'), node.id('vlist')
      if n and (n.id == HL or n.id == VL) then
        tex.setdimen('global', 'lakinfigraise',
                     (n.height or 0) + (n.depth or 0))
      else
        tex.setdimen('global', 'lakinfigraise', 0)
      end}%
  \else
    \global\lakinfigraise 0pt\relax
  \fi}
% BLD-37 table side: \lakin@grabfig's twin, NOT a shared helper — the Lua body
% differs in exactly one place, the skip set. Between \end{tabular} and \tabcap
% an otherlanguage wrapper (patterns 3/5 of example_tables) leaves a \write
% whatsit and a dir node ON TOP of the interword glue (measured node dump:
% "whatsit/write <- dir <- glue(2.62) <- hlist"), so the table reader must
% also step over whatsit/dir markers — all zero-size — to reach the tabular.
% \lakin@grabfig deliberately does NOT get the wider skip set: its glue/penalty
% body is the one the v3.74 A/B proof froze, and widening it would newly
% measure (and so newly raise) anchors in island shapes that today take the
% plain path — an unproven behavior change on the FIG side. Anything that is
% not a box behind those markers still hits the 0pt sentinel and \tabcap falls
% back to the plain \refstepcounter.
% SECOND deliberate difference: this stores the box's HEIGHT, not height+depth.
% \figcap's anchor is emitted in the caption paragraph BELOW the figure, so it
% must climb the full extent; \tabcap's anchor sits ON the tabular's own line,
% whose baseline is the vcentered tabular's MIDPOINT — height alone reaches
% the top edge, and adding the depth was measured to overshoot by half the
% table (dests 80-135pt above their frame, INSIDE the preceding table's box).
\def\lakin@grabtab{%
  \ifhmode
    \directlua{
      local n = tex.nest[tex.nest.ptr].tail
      local GLUE, PEN = node.id('glue'), node.id('penalty')
      local WHAT, DIR = node.id('whatsit'), node.id('dir')
      while n and (n.id == GLUE or n.id == PEN
                   or n.id == WHAT or n.id == DIR) do n = n.prev end
      local HL, VL = node.id('hlist'), node.id('vlist')
      if n and (n.id == HL or n.id == VL) then
        tex.setdimen('global', 'lakintabraise', n.height or 0)
      else
        tex.setdimen('global', 'lakintabraise', 0)
      end}%
  \else
    \global\lakintabraise 0pt\relax
  \fi}
\makeatother
% ============================================================================
% v3.79 (BLD-37 table side, CORRECTION): put the anchor at the TOP OF THE BOX,
% instead of guessing the distance to it.
% ============================================================================
% v3.77 raised the \tabcap anchor by (height of the box at the tail of the open
% horizontal list) + 42pt. That is a MEASUREMENT taken BEFORE line breaking, and
% it is short by every term it cannot see. Fleet-measured on the 12 built
% collections (299 tbl.* dests, ownership established through each table's own
% CAPTION, never through "which frame encloses the dest"): 3 dests landed INSIDE
% their own tablebox, below the frame top —
%   computer-architecture       tbl.14.1  -46.4pt   three stacked tabulars
%   computer-architecture       tbl.30.5  -13.5pt   tabular 2.85pt over \linewidth
%   advanced-electronic-devices tbl.B.1   -59.2pt   tabular 2.56pt over \linewidth
% Two missing terms, both reproduced in a minimal document (probe P2/P3):
%   (1) VERTICAL LIST ALREADY CONTRIBUTED inside the box — a second and third
%       tabular, or a heading above the table. The tail read sees only the LAST
%       box, so 77.7pt of stacked tabulars is invisible to the raise.
%   (2) WHICH LINE THE ANCHOR LANDS ON. When the tabular alone overfills
%       \linewidth (459.6121pt in a tablebox), the interword glue after
%       \end{tabular} becomes the chosen breakpoint and \Hy@raisedlink is
%       emitted on the NEXT line — one (depth + lineskip) lower. That term does
%       not EXIST at \tabcap time: line breaking has not happened yet. No
%       measurement taken there can ever recover it. (Proof: probe P3, the same
%       table 2.84912pt too wide, fails at -13.5pt; probe P4, the identical
%       table with \tabcolsep 2pt so it fits, is correct at +32.0pt.)
% So the raise is not computed at \tabcap at all. It is CORRECTED at
% \end{tablebox}, from inside mdframed's own body vbox (endinnercode runs after
% \par\unskip, i.e. after line breaking, while tex.nest[tex.nest.ptr] is still
% that vbox's vertical list). There the anchor's true offset from the top of the
% box content is not estimated but WALKED, and the \Hy@raisedlink shift is set
% so the dest lands 42pt above it — which is 42pt - innertopmargin(9pt) -
% linewidth(0.9pt) = 32.0pt above the FRAME top, the exact position the 295
% already-correct dests have today. That is why this is a correction and not a
% re-placement: the good sites do not move.
%
% WHY THE GLYPH DIFF IS 0 BY CONSTRUCTION — the only thing written is the
% `shift` field of the hlist that \Hy@raisedlink already built, and the guard
% below refuses to write it unless that box is 0x0x0. A zero-size box carries no
% dimension into its parent's packing, the packing already happened, and shift
% is not an input to it. Nothing else in the list is touched: no node is added,
% removed, resized or reordered. Measured, not assumed: assets/example_tables.tex
% and the whole computer-architecture collection A/B at 0 differing pymupdf
% rawdict glyph records.
%
% WHY NOT the alternatives (all prototyped and measured before this one):
%   * accumulate the vertical list at \tabcap time — fixes term (1), CANNOT fix
%     term (2); 2 of the 3 real failures survive. A fix that leaves the majority
%     of the defect standing is the false all-clear this repo exists to avoid.
%   * \zref-savepos two-pass — correct on all shapes, but it makes the dest
%     POSITION aux-derived (DEFECTS BLD-1/2), a single-pass compile silently
%     ships an unraised anchor, and the bare y-difference has no page guard.
%   * emit the anchor at \begin{tablebox} and let \tabcap adopt it — needs the
%     tbl counter stepped at box OPEN to keep the dest NAME byte-identical; a
%     captionless tablebox then either shifts every later table number or leaves
%     an orphan dest that the next table duplicates (a HARD build gate). This
%     correction touches neither the counter nor the name: \tabcap's naming path
%     is v3.77's, unchanged, so resolve_names() is byte-identical (verified).
%   * suppressing the line break (a penalty before the glue) — that MOVES the
%     tabular: the break currently discards the interword glue, so forbidding it
%     re-centers the two overfull tables. Reflow, rejected.
%
% SCOPE. The correction is armed only by \tabcap ITSELF (\lakintabcapseen),
% cleared at every box open, so a captionless tablebox is a no-op and a \tabcap
% hosted OUTSIDE a tablebox (the fleet has exactly one, computer-architecture
% tbl.14.2 in an exbox) keeps v3.77's measured raise. Both fallbacks are
% therefore never worse than today; and because hyperref's AUTOMATIC anchor uses
% the same \Hy@raisedlink shape, the plain-\refstepcounter path (vmode entry, or
% prose between \end{tabular} and \tabcap — probe P6, -51.5pt) is corrected too.
% KNOWN LIMIT — a tablebox that SPLITS across pages. The offset is measured from
% the top of the UNSPLIT body, so a dest in segment 2 gets raised past its own
% segment. MEASURED on a forced split (probe: 40 filler lines + a table): the
% dest went from y=585.0 on page 2 (inside segment 2, at the caption — the old
% BLD-37 symptom) to y=-113.4, i.e. ABOVE the page top. Viewers clamp that to the
% page top, which happens to show segment 2 whole, but an out-of-page coordinate
% is not something to ship on purpose. Hence the \textheight guard in
% lakintab.fix: a body TALLER THAN THE TEXT BLOCK cannot fit on any page, so it
% must split — the correction is skipped and v3.77's behaviour stands unchanged
% (verified: the probe returns to y=585.0). The guard is deliberately EXACT, not
% a page-fitting heuristic: it is sound (a taller-than-page body always splits)
% but PARTIAL (a shorter body that starts low on a page could still split). That
% residue is unreachable in practice — v3.37's everyline + the caption's
% \nopagebreak glue move such boxes WHOLE instead (probe N7: a 22-row table
% pushed down by 430pt moved whole, 0 segments), and the fleet has 0 split
% tableboxes, re-counted after this change. Reimplementing mdframed's own
% page-fitting arithmetic to close the residue would be a heuristic with its own
% failure mode, which is the trade this repo has repeatedly refused.
\newif\iflakintabcapseen
\newif\iflakintabframed
% The Lua lives in ONE \directlua at load time, and is written with NO
% backslash, %, # or ~ : \directlua expands its argument like \edef, so an
% active ~ would expand to \nobreakspace and a % would eat the rest of the line
% (the second is how the first draft of this block died). "not (a == b)" stands
% in for Lua's inequality operator for the same reason.
\directlua{
  lakintab = { }
  lakintab.HL = node.id("hlist")
  lakintab.VL = node.id("vlist")
  lakintab.RU = node.id("rule")
  lakintab.GL = node.id("glue")
  lakintab.KE = node.id("kern")
  lakintab.WH = node.id("whatsit")
  lakintab.TM = node.id("temp")
  lakintab.WD = node.subtype("pdf_dest")
  function lakintab.d(v)
    if v == nil then return 0 end
    if v < -1000000000 then return 0 end
    return v
  end
  function lakintab.scan(list, y, horiz, parent, target)
    local x = list
    while x do
      local i = x.id
      if i == lakintab.HL or i == lakintab.VL then
        if horiz then
          local b = y + x.shift
          if i == lakintab.HL then
            lakintab.scan(x.head, b, true, x, target)
          else
            lakintab.scan(x.head, b - lakintab.d(x.height), false, x, target)
          end
        else
          if i == lakintab.HL then
            lakintab.scan(x.head, y + lakintab.d(x.height), true, x, target)
          else
            lakintab.scan(x.head, y, false, x, target)
          end
          y = y + lakintab.d(x.height) + lakintab.d(x.depth)
        end
      elseif i == lakintab.WH then
        if x.subtype == lakintab.WD then
          local id = x.dest_id
          if type(id) == "string" and string.sub(id, 1, 4) == "tbl." then
            if horiz and parent and parent.width == 0
               and parent.height == 0 and parent.depth == 0 then
              parent.shift = parent.shift + (target - y)
            end
          end
        end
      elseif not horiz then
        if i == lakintab.RU then
          y = y + lakintab.d(x.height) + lakintab.d(x.depth)
        elseif i == lakintab.GL or i == lakintab.KE then
          y = y + lakintab.d(x.width)
        end
      end
      x = x.next
    end
  end
  function lakintab.fix(target, maxh)
    local nst = tex.nest[tex.nest.ptr]
    if nst == nil then return end
    local h = nst.head
    if h == nil then return end
    if h.id == lakintab.TM then h = h.next end
    local tot = 0
    local x = h
    while x do
      local i = x.id
      if i == lakintab.HL or i == lakintab.VL or i == lakintab.RU then
        tot = tot + lakintab.d(x.height) + lakintab.d(x.depth)
      elseif i == lakintab.GL or i == lakintab.KE then
        tot = tot + lakintab.d(x.width)
      end
      x = x.next
    end
    if tot > maxh then return end
    lakintab.scan(h, 0, false, nil, target)
  end
}
% @-FREE NAMES ON PURPOSE: these two are named in the `tablebox` key list at
% \newmdenv above, which is tokenised OUTSIDE any \makeatletter region — a
% \lakin@... there would tokenise as \lakin followed by the letters @tabopen.
\newcommand{\lakintabboxopen}{\lakintabframedtrue\global\lakintabcapseenfalse}
\newcommand{\lakintabboxclose}{%
  \iflakintabcapseen
    \directlua{lakintab.fix(-\number\dimexpr\lakintabcushion\relax,
                            \number\dimexpr\textheight\relax)}%
    \global\lakintabcapseenfalse
  \fi\lakintabframedfalse}
% The cushion above the box's CONTENT top. 42pt puts the dest 32.0pt above the
% frame top (42pt - innertopmargin 9pt - linewidth 0.9pt), i.e. exactly where
% the 295 correct dests already are — the correction moves only the broken ones.
\newcommand{\lakintabcushion}{42pt}
\newenvironment{qcirc}{\par\medskip\global\lakinfigraise 0pt\relax\begin{otherlanguage}{english}\begin{center}}
  {\csname lakin@grabfig\endcsname\end{center}\end{otherlanguage}\par\medskip}
\newenvironment{tikzpic}{\par\medskip\global\lakinfigraise 0pt\relax\begin{otherlanguage}{english}\begin{center}}
  {\csname lakin@grabfig\endcsname\end{center}\end{otherlanguage}\par\medskip}

% ============================================================================
% Code listings.  ★★ v3.4 THE BOX-OVERFLOW FIX:
%   resetmargins=true   → the listing IGNORES the enclosing environment's
%                         indentation and measures from the PAGE. v3.4 stated
%                         this backwards ("measured from the ENCLOSING box"),
%                         and v3.37 then blanked in-box line numbers because
%                         the page-anchored listing put its hanging digits on
%                         top of the box rule. Measured v4.5.0: with the key
%                         FALSE the listing follows the box and the code moves
%                         from x=75.26 to x=97.74, clearing the rule at 72.07.
%                         The key now lives in one \lstset below instead of in
%                         each style, so a box can override it — codebox does.
%   framexleftmargin=0  → the frame (left bar) does not extend left of the code.
%   xleftmargin keeps a small indent so the line numbers have room INSIDE.
% asmblock wraps the listing in an LTR island (code is LTR).
% ============================================================================
\lstdefinestyle{asm}{
  basicstyle=\ttfamily\small,
  backgroundcolor=\color{cBlueBg!55},
  frame=leftline,framerule=2.5pt,rulecolor=\color{cBlue},
  numbers=left,numberstyle=\tiny\color{cGray},numbersep=8pt,
  xleftmargin=2.2em,framexleftmargin=0pt,framexrightmargin=0pt,
  breaklines=true,columns=fullflexible,keepspaces=true,
  showstringspaces=false,aboveskip=8pt,belowskip=8pt,
  % v3.36: without a comment DELIMITER the commentstyle below never fires —
  % listings has no `language` here, so it does not know '#' opens a comment,
  % and every asm listing rendered flat black (measured: a whole-page pixel
  % sweep found ZERO cGreen glyphs across 135 shipped listings). morecomment
  % gives it the delimiter; color-only change, glyph widths untouched.
  morecomment={[l]{\#}},
  commentstyle=\color{cGreen},
}
\lstdefinestyle{cstyle}{
  language=C,basicstyle=\ttfamily\small,
  backgroundcolor=\color{cGreenBg!55},
  frame=leftline,framerule=2.5pt,rulecolor=\color{cGreen},
  numbers=left,numberstyle=\tiny\color{cGray},numbersep=8pt,
  xleftmargin=2.2em,framexleftmargin=0pt,framexrightmargin=0pt,
  breaklines=true,columns=fullflexible,keepspaces=true,
  showstringspaces=false,keywordstyle=\color{cBlue}\bfseries,
  commentstyle=\color{cGray}\itshape,aboveskip=8pt,belowskip=8pt,
}
% v3.36: pystyle/htmlstyle — colored twins of cstyle for the Python/HTML
% listings that used to repeat a ~200-char options blob with NO color
% (keywords rendered bold-black, comments italic-black — measured; C alone
% was colored). Geometry is copied VERBATIM from the shipped blob (frame,
% margins, skips) so swapping a call site to style=pystyle moves NOTHING —
% only keywordstyle/commentstyle/stringstyle are added (bold stays bold,
% italic stays italic, mono metrics identical => zero reflow, proven by
% per-page word-bbox diff). breaklines/breakatwhitespace stay per-site.
% v3.78: numbers=left joins both styles, with the exact numberstyle/numbersep
% keys asm/cstyle carry — fleet-wide consistency (author decision: all four
% styles numbered; "asm/cstyle numbered, py/html not" was an accident of the
% v3.36 geometry freeze, not a design). The numbers hang LEFT of the code box
% (listings sets them with \llap, so no overfull hbox can arise and the code
% x is untouched — measured: code frozen, digits at code_x-11.8pt); inside a
% colored box they are blanked like everyone else's (v3.37 hook), and the
% clean variant restores them in its frameless boxes (v3.78, above).
\lstdefinestyle{pystyle}{
  language=Python,basicstyle=\ttfamily\small,
  columns=fullflexible,keepspaces=true,
  frame=single,rulecolor=\color{cTblLine},
  numbers=left,numberstyle=\tiny\color{cGray},numbersep=8pt,
  xleftmargin=0.5em,framexleftmargin=0pt,framexrightmargin=0pt,
  showstringspaces=false,aboveskip=8pt,belowskip=8pt,
  keywordstyle=\color{cBlue}\bfseries,
  commentstyle=\color{cGray}\itshape,
  stringstyle=\color{cGreen},
}
% v4.5.0: hoisted out of the four styles so an enclosing box can override it.
% Every listing still resolves to true unless something later says otherwise,
% so shipped output is unchanged (verified byte-for-byte on all eight PDFs).
\lstset{resetmargins=true}
\lstdefinestyle{htmlstyle}{
  language=HTML,basicstyle=\ttfamily\small,
  columns=fullflexible,keepspaces=true,
  frame=single,rulecolor=\color{cTblLine},
  numbers=left,numberstyle=\tiny\color{cGray},numbersep=8pt,
  xleftmargin=0.5em,framexleftmargin=0pt,framexrightmargin=0pt,
  showstringspaces=false,aboveskip=8pt,belowskip=8pt,
  keywordstyle=\color{cBlue}\bfseries,
  commentstyle=\color{cGray}\itshape,
  stringstyle=\color{cGreen},
}
\newenvironment{asmblock}{\par\begin{otherlanguage}{english}}{\end{otherlanguage}\par}

% <hllpdf:listing-smart-quotes>
% v4.21.0: four Unicode punctuation characters must enter listings' token
% buffer. Unregistered characters bypass it: a closing quote was printed
% BEFORE the pending word, and an opening quote BEFORE a pending brace.
% Extend the standard character table, then activate only these codepoints
% in the listing's local character-table scope. Do this BEFORE the package's
% delimiter/alsoother/literate hooks so explicit user mappings still win.
% This is not universal Unicode support, nor a conversion to ASCII quotes.
% See listings developer guide 15.3.1 and references/listing-smart-quotes.md.
\makeatletter
\lst@CCPut \lst@ProcessOther {"2018}{"2019}{"201C}{"201D}\z@
\lst@AddToHookAtTop{SelectCharTable}{%
  \lst@ifactivechars
    \catcode"2018=\active \catcode"2019=\active
    \catcode"201C=\active \catcode"201D=\active
  \fi}
\makeatother
% </hllpdf:listing-smart-quotes>

% <hllpdf:listing-technical-symbols>
% v4.30: bounded extension for common technical punctuation in LTR listings.
% Without character-table registration these glyphs bypass the token buffer
% and can move before the preceding word even inside quoted strings.
% Keep the original four-quote block separate for its standing mutation test.
% This is source-faithful printing, not syntax validation or general Unicode.
\makeatletter
\lst@CCPut \lst@ProcessOther
  {"2013}{"2014}{"2026}{"2190}{"2192}{"2194}{"21D2}
  {"2212}{"221E}{"2260}{"2264}{"2265}\z@
\lst@AddToHookAtTop{SelectCharTable}{%
  \lst@ifactivechars
    \catcode"2013=\active \catcode"2014=\active \catcode"2026=\active
    \catcode"2190=\active \catcode"2192=\active \catcode"2194=\active
    \catcode"21D2=\active \catcode"2212=\active \catcode"221E=\active
    \catcode"2260=\active \catcode"2264=\active \catcode"2265=\active
  \fi}
\makeatother
% </hllpdf:listing-technical-symbols>


\usepackage[hidelinks,unicode]{hyperref}   % LAST, after colors. unicode = Hebrew bookmarks.
% v3.62: THE ISOLATION MACROS CORRUPTED EVERY PDF BOOKMARK THEY APPEARED IN.
% \en is defined (line ~155) as
%     \foreignlanguage{english}{\language=\csname l@nohyphenation\endcsname #1}
% and in a heading that body is written verbatim to the .toc and handed to
% hyperref's \pdfstringdef. hyperref strips the two CONTROL SEQUENCES (\language,
% \l@nohyphenation) but the plain `=` CHARACTER token between them is not a control
% sequence, so it survives into the PDF string: \chapter{שער \en{CNOT} והצמדה}
% produced the bookmark "שער =CNOT והצמדה". The printed TOC and the heading on the
% page are CORRECT — only the navigation pane and the viewer window title are wrong,
% which is why it went unnoticed. Measured on the shipped fleet: 70 corrupted
% bookmarks across 5 PDFs, from 177 unprotected \en-in-heading sites. Same corruption
% reached /Title and /Subject via the \hypersetup{pdftitle=...} the docs recommend.
% This hook tells hyperref to use the plain argument for all three isolation macros
% when building a PDF string. It does NOT touch the .toc (the printed TOC is
% byte-unchanged) and it does NOT mask the math signal — a heading with $...$ still
% logs "removing math shift", so BIDI-15 detection is intact.
% BONUS, measured: it also makes \code{} in an UNSTARRED heading COMPILE. BLD-10
% recorded that as needing \texorpdfstring or a starred heading; the real root cause
% is hyperref's \pdfstringdef scan of \code's \ifmmode...\else...\fi body — which is
% exactly why the identically-bodied \texttt survived (hyperref pre-neutralizes
% \texttt in PDF strings) and why \DeclareRobustCommand could never help.
% v3.64: \def\\{ } added. A manual line break in a heading WELDS the two words
% together in the bookmark: \chapter{שורה ראשונה\\שורה שנייה} produced the outline
% entry 'שורה ראשונהשורה שנייה' — "ראשונהשורה" is not a word. hyperref drops the
% \\ and puts nothing in its place, so the break that looks like a space on the
% page is a zero-width join in the navigation pane. The spaced form
% 'מילה אחת \\ מילה שנייה' escapes with a harmless DOUBLE space ('אחת  מילה'),
% which is why this never got noticed. Both are measured on a real outline.
% 0 live instances in the corpus — but SKILL.md and content-style.md both RECOMMEND
% a manual heading break, so this is a primed trap rather than a dormant one, and a
% one-token preamble fix retires it for good instead of asking authors to remember.
% Mapping \\ to a space fixes the WELD, which is the actual defect (a nonexistent
% word in the navigation pane). It does not normalise whitespace: measured after the
% fix, the welded form becomes 'שורה ראשונה שורה שנייה' (correct) and the already-
% spaced form goes from two spaces to three. That is cosmetic — viewers render it as
% a wider gap, and no word is corrupted. A per-heading \texorpdfstring is still the
% way to get an exactly-tuned bookmark string.
\pdfstringdefDisableCommands{\def\en#1{#1}\def\h#1{#1}\def\doclang#1{#1}\def\code#1{#1}\def\codeline#1{#1}\def\\{ }}
% v3.70: RESTORE \url LINE-BREAKING. url.sty freezes its break penalties at LOAD
% time (\mathchardef\UrlBreakPenalty=\binoppenalty), and v3.45's global
% \binoppenalty=10000 (inline-math atomicity) therefore priced EVERY breakpoint
% inside \Url at 10000 = never break. Measured: a long \url in Hebrew prose is
% one unbroken line, ~100pt overfull, invading the margin band the ink gates
% deliberately tolerate. With the two lines below: 0 overfull, the URL breaks
% after '/', the RTL split is CORRECT (fragment order forward), and each line
% gets its own link annotation exactly covering its fragment (measured).
% Confined to \Url's group — inline-math atomicity (TYP-7) is untouched.
\mathchardef\UrlBreakPenalty=200
\mathchardef\UrlBigBreakPenalty=100
\ifx\lakitmainlanguage\hllpdfenglish
  \hypersetup{pdflang=en}
\else
  \hypersetup{pdflang=he}
\fi  % /Lang is the document base locale (not per-span accessibility tagging).
% Set a per-document \hypersetup{pdftitle={...}} too — see SKILL.md.

% v3.52 (round-2 audit): two RTL bibliography/cross-ref defects, both invisible to
% every gate and both affecting the FIRST author to use the feature.
% (a) \@biblabel's default [#1] brackets are bidi-neutral and mirror to ]1[ under
% the RTL bibliography base — every \thebibliography reference list shows ]1[.
% LTR-isolating the label restores [1]. (The render gate's mirror rule fingerprints
% round parens )1(, so square brackets slipped through.)
\makeatletter
\renewcommand*{\@biblabel}[1]{\foreignlanguage{english}{[#1]}}
% v3.54: the SAME neutral-bracket trap on the INLINE side — a multi-key \cite{a,b}
% under the RTL base lays its European-number list right-to-left, so [1, 2] renders
% "[2,1]" and [1, 2, 3] renders "[3,2,1]" (verified by glyph geometry — the reader
% sees the sources in REVERSE order; a note form \cite[p.5]{a} also reorders to
% "[p. 5,1]"). LTR-isolate the whole bracket group (number list AND optional note),
% exactly like \@biblabel above: [1, 2] / [1, 2, 3] / [1, p. 5] all render correct.
% Placed AFTER hyperref so it wins; hyperref hyperlinks each number inside #1 (via
% \@citex), which this wrapper leaves intact. Mirrors the LaTeX kernel \@cite shape.
% v4.9.1: directional isolation must NOT switch the language here. Babel's
% safe-active citation wrapper passes its key list through \bbl@tempa;
% \foreignlanguage reuses that scratch macro before the deferred citation
% loop expands it. The old line therefore cited the literal key "english".
% If such a key existed, even the undefined-reference gate stayed green.
% \babelsublr isolates the bracket group without selecting a locale or
% changing that key list. Preserve native links, key order and optional notes.
\def\@cite#1#2{\babelsublr{[{#1\if@tempswa , #2\fi}]}}
\makeatother

% v4.11.0: a reference inside align's \text inherited HEBdisptag=true,
% so its enclosing hbox mirrored the brackets despite a clean compile.
% Isolate the complete printed reference without selecting a Babel language.
% Also neutralize the display-only normalfont direction reset locally. A
% direction-only patch inverted 12 to 21 inside equation in the measured probe.
% Plain maketag is local to the reference: mathtools' display-side wrapper
% otherwise globally consumes manual_tag while formatting a nested reference.
% Preserve each original command (including the selective auxiliary write),
% current tag forms, native links, and showonlyrefs on/off transitions.
\makeatletter
\NewCommandCopy{\HEB@nativeeqref}{\eqref}
\newcommand{\HEB@eqrefisolate}[1]{%
  \babelsublr{\HEBdisptagfalse\let\bbl@eqnodir\relax\let\maketag@@@\ORIGmaketag@@@#1}}
\protected\def\eqref#1{\HEB@eqrefisolate{\HEB@nativeeqref{#1}}}
\ExplSyntaxOn
\NewCommandCopy{\HEB@nativeMTeqref}{\MT_extended_eqref:n}
\cs_set_protected:Npn \MT_extended_eqref:n #1
  {\HEB@eqrefisolate{\HEB@nativeMTeqref{#1}}}
\ExplSyntaxOff
\makeatother

% v4.11.0: custom caption counters are not part of hyperref's English
% locale, and Hebrew has no upstream hyperref locale in this toolchain.
% Restore the appropriate names at every Babel language switch, including
% return from a nested inline island. Do not localize the destination keys.
\makeatletter
\newcommand{\HEB@refnames}{%
\def\chapterautorefname{פרק}%
\def\sectionautorefname{סעיף}%
\def\subsectionautorefname{תת־סעיף}%
\def\subsubsectionautorefname{תת־תת־סעיף}%
\def\paragraphautorefname{פסקה}%
% v4.31: hyperref uses a distinct counter name at the deepest standard level.
\def\subparagraphautorefname{תת־פסקה}%
\def\equationautorefname{משוואה}%
\def\figureautorefname{איור}%
\def\tableautorefname{טבלה}%
\def\partautorefname{חלק}%
\def\appendixautorefname{נספח}%
\def\footnoteautorefname{הערת שוליים}%
\def\itemautorefname{סעיף}%
\def\Itemautorefname{סעיף}%
\def\figautorefname{איור}%
\def\tblautorefname{טבלה}%
\def\pageautorefname{עמוד}% 
}
\addto\extrashebrew{\HEB@refnames}
\newcommand{\ENG@refnames}{%
  \def\subparagraphautorefname{subparagraph}%
  \def\chapterautorefname{chapter}%
  \def\sectionautorefname{section}%
  \def\subsectionautorefname{subsection}%
  \def\subsubsectionautorefname{subsubsection}%
  \def\paragraphautorefname{paragraph}%
  \def\equationautorefname{Equation}%
  \def\figureautorefname{Figure}%
  \def\tableautorefname{Table}%
  \def\partautorefname{Part}%
  \def\appendixautorefname{Appendix}%
  \def\footnoteautorefname{footnote}%
  \def\itemautorefname{item}%
  \def\figautorefname{Figure}%
  \def\tblautorefname{Table}%
  \def\Itemautorefname{item}%
  \def\pageautorefname{page}%
}
\addto\extrasenglish{\ENG@refnames}
\ifx\lakitmainlanguage\hllpdfenglish\ENG@refnames\else\HEB@refnames\fi
\makeatother

% ============================================================================
% Optional running head (v3.39): \lakithead{<book title>} — OFF by default.
% A document opts in with ONE line after % ============================================================================
% preamble.tex — tested LuaLaTeX preamble for Hebrew + math + figures PDFs (v4.32.0)
% Battle-tested: compiled an 80-page Hebrew quantum-computing summary AND a
% 197-page Hebrew computer-architecture anthology (TikZ diagrams, code listings,
% colored boxes) — both with 0 missing characters.
%
% Usage in your document:
%   \input{preamble.tex}     % (or paste these lines at the top)
%   \begin{document}  ...  \end{document}
%
% Compile with:  lualatex -interaction=nonstopmode doc.tex
%   Repeat until two consecutive PDFs are byte-identical (scripts/build.sh
%   does this) — a TOC document converges only on pass 3+; after two passes
%   the TOC page numbers are STALE (measured; see SKILL.md step 3).
% Font binaries are external dependencies, pinned by assets/fonts/SHA256SUMS.
% scripts/setup_fonts.sh imports a verified local installation or prior skill.
%
% PUBLIC DOCUMENT-SHELL / LAYOUT CONTRACT (introduced v4.3.1; architecture-owned v4.7.1). These hooks may be
% defined BEFORE \input{preamble.tex}; if omitted, the historical canonical
% defaults below are exact.  They deliberately control only the document shell
% and page layout.  RTL/math semantics stay in this preamble, while visual
% presentation themes (for example clean_theme.tex) load AFTER it.
%
%   \def\lakitdocumentclass{article}
%   \def\lakitclassoptions{11pt,a4paper}
%   \def\lakitgeometryoptions{margin=1.8cm,headheight=14pt}
%   \def\lakitlinespread{1.15}
%   \def\lakitparagraphstyle{block} % letters/briefs; indented for continuous prose
%   \def\lakittocdepth{1}
%   \def\lakitsecnumdepth{2}
%   \input{preamble.tex}
%
% Supported class contract: report or article.  Alternate paper sizes must also
% be declared to build.sh through EXPECTED_PAPER; the PDF gate never infers them.
% Do NOT use these hooks as a theme mechanism: theme overlays must not own page
% geometry, class selection, or direction semantics. A declared presentation-policy
% overlay may intentionally change numbering visibility after the preamble.
% ============================================================================
% Runtime architecture guard: generated from scripts/architecture_contract.py.
% The readiness sentinel remains delayed until the END of this file.
% <hllpdf:architecture-guard>
\ifdefined\hllpdfcoreloaded
  \errmessage{hllpdf: preamble.tex loaded more than once}
  \expandafter\endinput
\fi
% </hllpdf:architecture-guard>

\ifdefined\lakitdocumentclass\else\def\lakitdocumentclass{report}\fi
\ifdefined\lakitclassoptions\else\def\lakitclassoptions{12pt,a4paper}\fi
\ifdefined\lakitgeometryoptions\else\def\lakitgeometryoptions{margin=2.3cm,headheight=14pt}\fi
\ifdefined\lakitlinespread\else\def\lakitlinespread{1.3}\fi
\ifdefined\lakittocdepth\else\def\lakittocdepth{1}\fi
\ifdefined\lakitsecnumdepth\else\def\lakitsecnumdepth{2}\fi
% Paragraph organization belongs to the document layout, independently of its
% language and theme. Preserve historical class dimensions unless explicitly
% selecting block paragraphs; never suppress indentation with an everypar hack.
\ifdefined\lakitparagraphstyle\else\def\lakitparagraphstyle{indented}\fi
\edef\lakitparagraphstyle{\lakitparagraphstyle}
\def\hllpdfparagraphblock{block}
\def\hllpdfparagraphindented{indented}
\ifx\lakitparagraphstyle\hllpdfparagraphblock\else
  \ifx\lakitparagraphstyle\hllpdfparagraphindented\else
    \errmessage{hllpdf: paragraph style must be block or indented}
  \fi
\fi
% v4.3.1: the public class hook has a deliberately narrow support contract.
% Before this guard, \def\lakitdocumentclass{book} silently compiled even though
% every heading/TOC regression test targets report/article; an unsupported shell
% could therefore look "supported" until a subtle class-specific defect shipped.
% Fail before \documentclass rather than letting an undocumented combination enter
% the mechanism core. Plain-TeX comparison is used because LaTeX packages are not
% available yet at this point.
\edef\lakitvalidatedclass{\lakitdocumentclass}
\def\lakitreportclass{report}
\def\lakitarticleclass{article}
\ifx\lakitvalidatedclass\lakitreportclass
\else\ifx\lakitvalidatedclass\lakitarticleclass
\else
  \errmessage{hllpdf: unsupported class \lakitdocumentclass; use report or article}
\fi\fi
\let\lakitvalidatedclass\relax
\let\lakitreportclass\relax
\let\lakitarticleclass\relax
% Language intent is independent of document class, theme and content domain.
% The two dominant-language profiles differ in authoring intent, not direction.
% A balanced bilingual document must explicitly choose a navigation/base locale.
% <hllpdf:language-profile>
\providecommand{\lakitlanguageprofile}{hebrew}
\edef\lakitlanguageprofile{\lakitlanguageprofile}
\def\hllpdfhebrew{hebrew}\def\hllpdfenglish{english}
\def\hllpdfhebrewmixed{hebrew-with-english}
\def\hllpdfenglishmixed{english-with-hebrew}
\def\hllpdfbilingual{bilingual}
\let\hllpdfexpectedlanguage\relax
\ifx\lakitlanguageprofile\hllpdfhebrew
  \let\hllpdfexpectedlanguage\hllpdfhebrew
\else\ifx\lakitlanguageprofile\hllpdfhebrewmixed
  \let\hllpdfexpectedlanguage\hllpdfhebrew
\else\ifx\lakitlanguageprofile\hllpdfenglish
  \let\hllpdfexpectedlanguage\hllpdfenglish
\else\ifx\lakitlanguageprofile\hllpdfenglishmixed
  \let\hllpdfexpectedlanguage\hllpdfenglish
\else\ifx\lakitlanguageprofile\hllpdfbilingual
  \ifdefined\lakitmainlanguage\else
    \errmessage{hllpdf: bilingual profile requires explicit lakitmainlanguage}
  \fi
\else
  \errmessage{hllpdf: unsupported language profile \lakitlanguageprofile}
\fi\fi\fi\fi\fi
\ifdefined\lakitmainlanguage
  \edef\lakitmainlanguage{\lakitmainlanguage}
\else
  \let\lakitmainlanguage\hllpdfexpectedlanguage
\fi
\ifx\lakitmainlanguage\hllpdfhebrew\else
  \ifx\lakitmainlanguage\hllpdfenglish\else
    \errmessage{hllpdf: main language must be hebrew or english}
  \fi
\fi
\ifx\hllpdfexpectedlanguage\relax\else
  \ifx\lakitmainlanguage\hllpdfexpectedlanguage\else
    \errmessage{hllpdf: main language conflicts with language profile}
  \fi
\fi
\typeout{HLLLANG|\lakitlanguageprofile|\lakitmainlanguage}
% </hllpdf:language-profile>
\documentclass[\lakitclassoptions]{\lakitdocumentclass}
\usepackage[no-math]{fontspec}        % (1) ★ no-math — THE critical line (see recipe.md)
% layout=lists asks Babel to adapt list parshape when paragraph and page
% directions differ. Without it, minority-language list bodies lose their
% indentation and labels hang outside the text block, despite a clean log.
% The upstream mechanism preserves ordinary same-direction list geometry.
\usepackage[bidi=basic,layout=lists]{babel}
\ifx\lakitmainlanguage\hllpdfenglish
  \babelprovide[import]{hebrew}
  \babelprovide[import,main]{english}
\else
  \babelprovide[import,main]{hebrew}
  \babelprovide[import]{english}
\fi
% v4.22: bind canonical families and fallbacks to the selected installation.
% Lookup by basename or family alone is not proof of the file actually used.
% <hllpdf:font-binding>
\directlua{
  local dir = os.getenv("HEBREW_LATEX_FONT_DIR")
  if not dir or dir == "" then
    local data = os.getenv("XDG_DATA_HOME")
    if not data or data == "" or not (data:sub(1,1) == "/") then
      data = (os.getenv("HOME") or "") .. "/.local/share"
    end
    dir = data .. "/fonts/hebrew-lualatex-pdf"
  end
  if not (dir:sub(1,1) == "/") then dir = lfs.currentdir() .. "/" .. dir end
  if dir:find("[",1,true) or dir:find("]",1,true) or dir:find(string.char(0),1,true)
     or dir:find(string.char(10),1,true) or dir:find(string.char(13),1,true) then
    tex.error("hllpdf: unsupported character in font directory")
  end
  if not (lfs.attributes(dir,"mode") == "directory") then
    tex.error("hllpdf: font directory is missing")
  end
  hllpdf_font_dir = dir:gsub("/+$", "") .. "/"
  tex.sprint(string.char(92) .. "def" .. string.char(92) .. "hllpdffontdir{")
  tex.sprint(-2,hllpdf_font_dir)
  tex.sprint("}")
  luaotfload.add_fallback("hebfb",{"[" .. hllpdf_font_dir .. "DejaVuSerif.ttf]:mode=node;"})
  luaotfload.add_fallback("hebfbbf",{"[" .. hllpdf_font_dir .. "DejaVuSerif-Bold.ttf]:mode=node;"})
}
% </hllpdf:font-binding>
% v3.22: Frank Ruhl Libre ships NO italic face. Instead of silently-upright
% \emph (v<=3.20) or bold-as-emphasis (v3.21), synthesize a slanted face with
% FakeSlant — real, VISIBLE italic emphasis with correct \emph nesting
% semantics (inner \emph toggles back upright). Verified empirically at
% 300dpi: slant clearly visible, 0 errors, 0 missing glyphs.
\babelfont{rm}[Path={\hllpdffontdir},RawFeature={fallback=hebfb},
  BoldFont={FrankRuhlLibre-Bold.ttf},
  BoldFeatures={RawFeature={fallback=hebfbbf}},   % v3.56 — see hebfbbf above
  ItalicFont={FrankRuhlLibre-Regular.ttf}, ItalicFeatures={FakeSlant=0.22},
  % Bold-italic uses the weight-matched fallback and a synthetic slant.
  BoldItalicFont={FrankRuhlLibre-Bold.ttf}, BoldItalicFeatures={RawFeature={fallback=hebfbbf},FakeSlant=0.22}
  ]{FrankRuhlLibre-Regular.ttf}
% Sans fallbacks preserve family and weight for the few missing codepoints.
\directlua{luaotfload.add_fallback("hebfbsans",{"[" .. hllpdf_font_dir .. "DejaVuSans.ttf]:mode=node;"})}
\directlua{luaotfload.add_fallback("hebfbsansbf",{"[" .. hllpdf_font_dir .. "DejaVuSans-Bold.ttf]:mode=node;"})}
\babelfont{sf}[Path={\hllpdffontdir},RawFeature={fallback=hebfbsans},
  BoldFont={Heebo-Bold.ttf},
  BoldFeatures={RawFeature={fallback=hebfbsansbf}}]{Heebo-Regular.ttf}
% v3.24: Scale=MatchLowercase — DejaVu Sans Mono has an outsized x-height, so
% unscaled inline code visually SHOUTS next to the delicate Frank Ruhl body
% (measured on a 197-page build: scaled, a code atom sits in the line at the
% body's weight, and the document tightened by ~5%). Applies to listings too.
\babelfont{tt}[Path={\hllpdffontdir},RawFeature={fallback=hebfb},Scale=MatchLowercase,
  BoldFont={DejaVuSansMono.ttf},
  BoldFeatures={RawFeature={fallback=hebfb},FakeBold=1.5},
  ItalicFont={DejaVuSansMono.ttf},
  ItalicFeatures={RawFeature={fallback=hebfb},FakeSlant=0.2},
  BoldItalicFont={DejaVuSansMono.ttf},
  BoldItalicFeatures={RawFeature={fallback=hebfb},FakeBold=1.5,FakeSlant=0.2}
  ]{DejaVuSansMono.ttf}

% Explicit multilingual/English profiles use real Latin faces and italics.
% The legacy/default Hebrew profile retains its established measured typography.
% Latin Modern is provided by the TeX distribution, not the pinned font bundle;
% recorder/runtime evidence reports it as an additional toolchain resource.
\ifx\lakitlanguageprofile\hllpdfhebrew\else
  \babelfont[english]{rm}[Ligatures=TeX,RawFeature={fallback=hebfb},
    BoldFeatures={RawFeature={fallback=hebfbbf}},
    BoldItalicFeatures={RawFeature={fallback=hebfbbf}}]{Latin Modern Roman}
  \babelfont[english]{sf}[Ligatures=TeX,RawFeature={fallback=hebfbsans},
    BoldFeatures={RawFeature={fallback=hebfbsansbf}},
    BoldItalicFeatures={RawFeature={fallback=hebfbsansbf}}]{Latin Modern Sans}
\fi

\usepackage{amsmath,amssymb,mathtools}
% amsthm supplies theorem/proof machinery. Environment names remain
% document-specific; the generic proof label is localized below.
\usepackage{amsthm}
% Proof names follow the surrounding prose, including return from an island.
\addto\extrashebrew{\renewcommand{\proofname}{הוכחה}}
\addto\extrasenglish{\renewcommand{\proofname}{Proof}}
\ifx\lakitmainlanguage\hllpdfenglish
  \renewcommand{\proofname}{Proof}
\else
  \renewcommand{\proofname}{הוכחה}
\fi
% v3.24.1: equation tags under babel-bidi. A tag re-emitted by mathtools'
% showonlyrefs inside align inherits the RTL base and its parens MIRROR —
% ")2.1(" instead of "(2.1)" (plain align/equation tags are unaffected;
% verified on a minimal case + a 198-page build). LTR-isolate at the tag-BOX
% level: mathtools' show/hide + numbering logic is untouched. Do NOT override
% \tagform@ instead — that clobbers the showonlyrefs mechanism (measured:
% every tag shows, numbering breaks).
% v3.28.1: the v3.24.1 wrap must apply ONLY to tags of display environments.
% Inline \eqref in Hebrew PROSE also goes through \maketag@@@, and there the
% english-wrap MIRRORS the parens — ")1(" — while the unwrapped original is
% correct (proven on a 2x2 isolation matrix: patch x showonlyrefs; regressed
% every prose \eqref since v3.24.1, caught by the render gate's R4). The two
% call sites are indistinguishable by \ifmmode (tag boxes are text mode), so
% the display environments themselves raise a flag; \eqref in prose sees it
% false and gets the original untouched \maketag@@@.
\usepackage{etoolbox}

% babel bidi=basic (including upstream 26.9) assumes every automatic
% discretionary has a one-glyph replacement. LuaTeX can legally produce one
% with replace=nil after a hyphenated Latin value; Babel then indexes nil and the
% whole build dies. Do the narrow upstream-compatible correction: while Babel
% classifies the list, present only that malformed disc as an ordinary
% discretionary, then restore its subtype. Babel still reorders the ENTIRE box;
% there is no pcall and no successful fallback to an unreordered node list.
% If Babel throws for any other reason, the build remains fatal. The synthetic
% nil-disc probe and the real table gallery in test_all.sh pin both layers.
\AtBeginDocument{\directlua{
  if Babel and Babel.bidi and not Babel.heb_nil_disc_compat then
    Babel.heb_nil_disc_compat = true
    local original = Babel.bidi
    local disc_id = node.id('disc')
    Babel.bidi = function(head, ispar, hdir)
      local hidden = {}
      for item in node.traverse_id(disc_id, head) do
        if item.subtype == 2 and (not item.replace or not item.replace.char) then
          table.insert(hidden, item)
          item.subtype = 0
        end
      end
      local result = original(head, ispar, hdir)
      for _, item in ipairs(hidden) do item.subtype = 2 end
      return result
    end
  end
}}
\makeatletter
\newif\ifHEBdisptag
\let\ORIGmaketag@@@\maketag@@@
\def\maketag@@@#1{\ifHEBdisptag\ORIGmaketag@@@{\foreignlanguage{english}{#1}}\else\ORIGmaketag@@@{#1}\fi}
\@for\HEB@env:=equation,equation*,align,align*,gather,gather*,multline,multline*,flalign,flalign*,alignat,alignat*,eqnarray,eqnarray*\do{%
  \expandafter\AtBeginEnvironment\expandafter{\HEB@env}{\HEBdisptagtrue}}
\makeatother
% English page + Hebrew prose: Babel's equation-only tag flip uses a primitive
% leqno path that shifts the formula to the left edge and the tag into the body
% of an English page. Align's full-width alignment is unaffected. Give only this
% display a physical LTR container and native amsmath eqno/leqno. The engine then
% owns real tag width, collision handling and labels; no proxy tag box hides ink
% from its fit calculation. Babel's saved \bbl@eqnodir still restores Hebrew text
% inside math, and the prose locale/math direction are never switched.
% Close the prior RTL paragraph before normalization. After the environment's
% group restores its directions, close TeX's empty resumed display paragraph and
% use native LaTeX continuation handling: same-paragraph prose has no indent,
% while an authored blank line starts an ordinary indented RTL paragraph.
% Register after Babel's setup; remain local to equation(*) in a main-English
% document's actual RTL prose. Native Hebrew, ordinary English and explicit
% fleqn retain their original paths. No page/display dimensions are assigned.
% <hllpdf:minority-rtl-equation-geometry>
\makeatletter
\newcommand{\HEB@minorityprosecontinuation}{\par\@doendpe}
\newcommand{\HEB@minorityequationgeometry}{%
  \ifx\lakitmainlanguage\hllpdfenglish
    \ifnum\pardirection=\@ne
      \if@fleqn\else
        \par\aftergroup\HEB@minorityprosecontinuation
        \iftagsleft@\let\veqno\@@leqno
        \else\let\veqno\@@eqno\fi
        \bodydir TLT\pardir TLT\textdir TLT\relax
      \fi
    \fi
  \fi}
\AtBeginDocument{%
  \AddToHook{env/equation/begin}{\HEB@minorityequationgeometry}%
  \AddToHook{env/equation*/begin}{\HEB@minorityequationgeometry}}
\makeatother
% </hllpdf:minority-rtl-equation-geometry>
% v3.63: \shortintertext rendered Hebrew MIRROR-REVERSED and flush LEFT — an entire
% commentary line inside an align printed backwards. babel's bidi=basic patches
% \intertext@ only (luababel.def), and mathtools' \shortintertext uses the identical
% \vbox{\normalbaselines … \noindent\ignorespaces#1\par} skeleton but is never
% patched, so its paragraph keeps pardir/textdir LTR: no bidi reordering (logical
% order on the page) and flush-left. Measured against \intertext with the same
% Applying babel's own replacement to \shortintertext@ makes the two identical.
% \h{}/\mbox fix the letter order but leave the line flush LEFT, so they are not
% substitutes. Guarded + \AtBeginDocument because mathtools \let-copies at load time.
% NB \ExplSyntaxOn is REQUIRED: mathtools' internals carry expl3 names, and without
% the expl3 catcodes the `_` is read as a subscript and the patch detonates far away
% with "! Illegal unit of measure (pt inserted)" at the \end{align} (measured).
\makeatletter
\AtBeginDocument{%
  \@ifundefined{shortintertext@}{}{%
    \ExplSyntaxOn
    \bbl@sreplace\shortintertext@{\normalbaselines}%
      {\normalbaselines\ifx\bbl@eqnodir\relax\else\bbl@pardir\@ne\bbl@eqnodir\fi}%
    \ExplSyntaxOff}}
\makeatother
% geometry inherits the canonical A4 sheet from \documentclass; build.sh then
% verifies every rendered MediaBox and CropBox, so a silent Letter fallback is
% impossible to publish as a clean default build.
\usepackage[\lakitgeometryoptions]{geometry}
\linespread{\lakitlinespread}       % layout hook; canonical default remains 1.3
\ifx\lakitparagraphstyle\hllpdfparagraphblock
  \setlength{\parindent}{0pt}
  \setlength{\parskip}{6pt plus 1pt minus 1pt}
\fi
% These dimensions remain ordinary scoped TeX lengths: authors can customize
% them after core input. Language switches inherit them; lists keep their own
% indentation and narrow reference layouts establish their local paragraph gaps.
\typeout{HLLPAR|\lakitparagraphstyle}
\usepackage{xcolor}
\usepackage{mdframed}
\usepackage{listings}                 % code listings (assembly / C / pseudo-code)
\usepackage{tikz}
\usetikzlibrary{arrows.meta,positioning,calc,shapes.geometric,fit,backgrounds}
\usepackage{pgfplots}\pgfplotsset{compat=1.18}
\usepackage{quantikz}                 % quantum circuits (remove if unused)
\usepackage{booktabs}
\usepackage{array}
\usepackage{float}                    % [H] = place table exactly here (reference doc: no drifting)
\usepackage{enumitem}
\usepackage{needspace}

% ============================================================================
% ★★ v3.4 — THE DIRECTION-ISOLATION TRIO (the heart of mixing Hebrew + Latin)
% Mental model:
%   \h{...}    Hebrew (RTL) island      — Hebrew text inside an LTR context
%                                          (TikZ nodes, the english `tikzpic`).
%   \en{...}   English (LTR) isolate    — Latin text / parentheses inside a
%                                          Hebrew (RTL) context. KEEPS PARENS
%                                          AND ORDER CORRECT.
%   \code{...} \en + monospace          — code inside Hebrew / inside a frametitle.
% CRITICAL: never wrap Latin in \h{} (it is for HEBREW). Inside a TikZ node,
% Latin in \h{} renders REVERSED (Datapath->htapataD, CPU->UPC) and parens
% mirror. Isolate every Latin/paren chunk with \en{} or \code{}. See
% references/figures-boxes-listings.md.
% ---- TEXT INSIDE MATH (v3.8) ----------------------------------------------
% Put ANY text that sits inside math in a \text{}. The ONLY thing that is
% invalid in math is a bare \h/\en (= \foreignlanguage) used DIRECTLY, outside
% \text -> "\rmfamily invalid in math mode". Inside \text:
%   Hebrew label           -> \text{...}            (bare Hebrew renders fine;
%                              \h is optional — BUT NOT ALWAYS HARMLESS (v3.67):
%                              if the \h group OPENS with a nested island
%                              ($..$, \en{}, \code{}, \texttt{}, \mbox{}) the
%                              WHOLE group loses bidi and the Hebrew prints in
%                              LOGICAL order = gibberish. Measured:
%                                \text{\h{$f$ קבועה}}  -> reads "העובק"  BROKEN
%                                \text{\h{קבועה $f$}}  -> correct
%                                \text{$f$ קבועה}      -> correct (no \h)
%                              Hebrew first inside \h, or drop the \h. Flagged
%                              H-OPENS-WITH-LTR-GROUP.)
%   plain Latin word        -> \text{...}            (no marker needed)
%   Latin WITH ( ) or [ ]   -> \text{\en{...}}       (\en forces LTR; a plain
%                              \text mirrors the parens in RTL -> "Spec (TNR(" )
%   f(Hebrew/text arg)      -> \text{f}(\text{...})  -- keep the parens at MATH
%                              level, NOT \text{f(arg)} (parens inside the box mirror)
%   code in math            -> \mathtt{...}  (or \texttt{}, math-guarded below)
% NEVER write \en{...} or \h{...} straight into $...$ / \[..\] (outside \text);
% always via \text{}. (\mbox{...} also works but does not scale in scripts.)
% ============================================================================
\newcommand{\h}[1]{\foreignlanguage{hebrew}{#1}}
% v3.42.1 — \en islands do NOT hyphenate (auto OR at explicit hyphens). English hyphenation patterns were
% tearing identifiers and terms ACROSS THE RTL SEAM ("IO Regis-/ters",
% "In-/put", "Band-/width", "Per-/formance", "Multi-/ple" — all shipped on
% real pages; probe-proven: "ALUS-/rcB", "Mem-/ToReg"). An island is a label
% or a term, not English prose — a hyphen tear there is never right. A/B on
% a full 197pp lakit: identical overfull sets, same page count, ~11
% paragraphs healed, one mild underfull (badness 3525) that replaced a
% mutilated "Regis-/ters". Long English PROSE belongs in the englishprose
% environment below, which restores explicit-hyphen breaks while its paragraph
% is still active. A raw otherlanguage block inherits the global prohibition.
% (\l@nohyphenation needs \csname:
% '@' is not a letter here.)
% v3.43.2 — \exhyphenpenalty moved OUT of \en to a GLOBAL assignment below.
% Render-proven: a group-local \exhyphenpenalty is DEAD CODE — TeX reads the
% line-breaking parameters at \par time, after the \en group has closed, so
% \en{write-through} / plain not-taken still tore at the "-" ("write-/
% through" shipped on a real exam page). \language is different: it is
% recorded in the node list as text is added, so the in-group nohyphenation
% assignment above DOES work (and stays). The global penalty was A/B-measured
% on a 197pp lakit: identical overfull set (25), identical page count, the
% not-/taken tear healed. Hebrew is unaffected in practice — Hebrew compounds
% use maqaf (־), not ASCII "-".
\newcommand{\en}[1]{\foreignlanguage{english}{\language=\csname l@nohyphenation\endcsname #1}}
% Babel runs extras locally for inline islands and paragraph language changes.
% An English main document must not inherit the Hebrew no-hyphen paragraph rule.
\addto\extrashebrew{\exhyphenpenalty=10000\relax}
\addto\extrasenglish{\exhyphenpenalty=50\relax}
\ifx\lakitmainlanguage\hllpdfenglish
  \exhyphenpenalty=50\relax
\else
  \exhyphenpenalty=10000\relax
\fi
% Full English paragraphs need ordinary explicit-hyphen breaks. Use this
% environment instead of a raw otherlanguage block; the paragraph ends while
% the local penalty is still active.
\newenvironment{englishprose}
  {\par\begin{otherlanguage}{english}\exhyphenpenalty=50\relax}
  {\par\end{otherlanguage}}
\newenvironment{hebrewprose}
  {\par\begin{otherlanguage}{hebrew}\exhyphenpenalty=10000\relax}
  {\par\end{otherlanguage}}
% Use for shell text whose locale must not depend on the last body paragraph.
\DeclareRobustCommand{\doclang}[1]{\foreignlanguage{\lakitmainlanguage}{#1}}

% ============================================================================
% ★★ — fixes for defects that are INVISIBLE in the log and in the checks
% and only surface on a VISUAL RENDER of the PDF. (Each verified empirically.)
% ----------------------------------------------------------------------------
% v3.22: \emph is now REAL slanted emphasis (FakeSlant italic face defined on
% the \babelfont{rm} line above), replacing the v3.21 \emph->\textbf hack: it
% preserves \emph's toggle semantics and keeps bold free for its own role.
% Prefer bold-as-emphasis for your doc? one line: \renewcommand{\emph}[1]{\textbf{#1}}
% page number: force the digits LTR so they never reverse (18 -> 81) in the RTL
%   footer. That reversal is a NON-DETERMINISTIC bidi glitch (observed once in
%   ~80 pp); this LuaLaTeX build is not byte-reproducible, so the defensive
%   LTR wrap is the only reliable guard. Latin digits are always LTR anyway.
\renewcommand{\thepage}{\foreignlanguage{english}{\arabic{page}}}
% v3.70: \pagenumbering \gdef's \thepage and silently DESTROYS the wrap above —
% after \pagenumbering{arabic} the folio guard is gone for the whole rest of the
% document. Two collections discovered this independently and hand-restore the
% wrap in main.tex with the same comment; that knowledge now lives here instead.
% The wrapper re-applies the LTR wrap after the kernel call, for any style.
\makeatletter
\let\lakin@orig@pagenumbering\pagenumbering
\renewcommand{\pagenumbering}[1]{%
  \lakin@orig@pagenumbering{#1}%
  \renewcommand{\thepage}{\foreignlanguage{english}{\csname @#1\endcsname\c@page}}}
\makeatother
% emergencystretch: Hebrew does not hyphenate, so an unbreakable LTR island
%   (\en{...} / an inline math run) landing at the margin can leave a small
%   (~2-13pt) overfull with no legal breakpoint. This rescues ONLY those
%   inline-prose tie-breaks; it is a ceiling used only on offending lines and
%   does NOT touch display-math / listing / table overflow (those still surface
%   in the log for the content-level overflow pass). Verified: baseline overfull
%   -> 0 with this set, with no change to pagination or underfull.
\setlength{\emergencystretch}{3em}
% v3.40: inline math (almost) never breaks across lines. TeX's default
%   (binop 700 / rel 500) lets an inline run split at an operator —
%   $(1.f)\cdot 2^{e}$ shipped as "(1.f)·" ending one line and "2^e" opening
%   the next (reader-reported). In LTR that split is merely ugly; under an
%   RTL base the carried operator lands on the visually WRONG side of the
%   break and the formula reads backwards.
%   WHY 10000 (v3.45.0 — the display-first resolution). Inline math is
%   FORBIDDEN from breaking, so a torn island under an RTL base is
%   structurally impossible. The v3.40 cycle used 9999 ("break only as a last
%   resort") because a bare 10000 left ~30-39 WIDE inline equations
%   overflowing (worst 90.7pt) — but 9999's "last resort" fires OFTEN in
%   tight, non-hyphenating Hebrew: a reader found torn islands across the
%   fleet (~150 candidates; e.g. "P(y |" ending one line and "x_0)" opening
%   the next, read backwards). The step the 9999 cycle missed: every one of
%   those ~30 overflows is a wide equation that BELONGS in display. Convert
%   them to display FIRST, then 10000 gives 0 tears AND 0 overflows (measured,
%   12/12 fleet). Doctrine: inline math is ATOMIC; a wide equation goes in a
%   display, never inline. See content-style ("inline math is atomic").
\binoppenalty=10000 \relpenalty=10000
% v3.38: uniform sentence spacing. TeX's default (\nonfrenchspacing, sfcode
% 3000) is the OLD Anglo-American wide-space-after-period convention — Hebrew
% typography never had it, and modern Latin practice (Chicago et al.) dropped
% it too. Measured on real pages: post-period gaps up to 1.9-2.1x the line's
% mean word gap — the dominant "rivers" source in justified Hebrew prose.
% Survives babel language switches (probed: sfcode stays french inside \en{}
% islands and otherlanguage blocks).
\frenchspacing
% v3.38: widow/orphan protection. LaTeX's default penalties (150/150/50)
% allow a lone line at a page top/bottom almost freely. 10000 = forbid;
% this is SAFE here precisely because the document is \raggedbottom (the
% demoted line just moves on, no vertical stretch fight) — under flushbottom
% these values would cause underfull pages. raggedbottom is the report-class
% default; made explicit so a future class/option change cannot silently
% flip it and turn the 10000s harmful.
\raggedbottom
\widowpenalty=10000 \clubpenalty=10000 \displaywidowpenalty=10000
% em-dash: --- does NOT become — under Frank Ruhl (the font lacks tlig; verified
%   the width is identical with/without Ligatures=TeX, while the mechanism works
%   on Latin Modern). TYPE THE CHARACTER — (U+2014) DIRECTLY; likewise – (U+2013).
%   check_bidi_figures.py flags a --/--- that is adjacent to Hebrew.
% ============================================================================

% ---- STRUCTURAL DIRECTION-SAFETY for \texttt (v3.7) -------------------------
% A parenthesis, bracket, or $ INSIDE \texttt{} mirrors/migrates in RTL prose
% (MIPS 0($s0)->"(s0$)0", $rs->"rs$", arr[i]->"[arr[i") because \texttt sets the
% font but NOT the direction. Redefining \texttt to LTR-isolate its content fixes
% the WHOLE class at once (parens, brackets, $) so plain \texttt{lw $t0,0($s0)}
% is safe — no need to remember \code{}. Guarded by \ifmmode because
% \foreignlanguage is invalid in math (there, plain monospace is already LTR).
% NB (corrected v3.56 — the old note here said "stays tofu", which is WRONG):
% \texttt{Hebrew} is a FATAL build failure. The tt family is DejaVu Sans Mono
% (zero Hebrew) and the tt fallback hebfb is DejaVuSerif (also zero Hebrew), so
% the subsetter gets an empty glyph set and LuaTeX dies with
%   "! error: (file .../DejaVuSansMono.ttf) (type 2): there are no glyphs in the
%    subset"  -> no output PDF.
% The message names a FONT FILE, never Hebrew/\texttt/the source line. And it is
% non-local: add an unrelated Latin \code{} elsewhere and the fatal degrades into
% a SILENT DELETION of the Hebrew word. Put Hebrew in \h{}/normal text, never in
% \texttt{}/\code{}; flagged HEBREW-IN-MONO. To disable, delete this block.
\let\origtexttt\texttt
\protected\long\def\texttt#1{\ifmmode\origtexttt{#1}\else\foreignlanguage{english}{\origtexttt{#1}}\fi}
% v4.8.0: defer the direction test until TeX has inserted an alignment cell's
% math template. Unprotected expansion here tested \ifmmode too early at the
% FIRST token of aligned/array/cases/matrix cells. Engine-level \protected also
% preserves the original \texttt implementation without the self-recursion
% observed when retrofitting \DeclareRobustCommand to its saved alias.
\protected\long\def\code#1{\ifmmode\origtexttt{#1}\else\foreignlanguage{english}{\origtexttt{#1}}\fi}

% ---- \codeline — display-code line (v3.24): the code analogue of display math.
% Reader-friction rule (content-style.md §5b): inline \code{} is for a SINGLE
% ATOM only (one register/variable/mnemonic, a tight no-space expression).
% A multi-token expression — an instruction with operands, a spaced assignment,
% a register-binding list — zigzags the eye RTL<->LTR mid-sentence and, being
% an unbreakable island, line-breaks badly. Give it its own centered line with
% display-math vertical rhythm. \centerline = one box => the expression never
% wraps; \nopagebreak keeps it glued to its introducing sentence.
%   ...the Hebrew sentence ends with a colon:
%   \codeline{lw \$t1, -4(\$t0)}
%   and the prose continues (escape dollars as \$).
% v3.66 (first line of the body): flush a PENDING \item label before the \par.
% LaTeX defers the label box (\@labels) until the next \everypar, and an \hbox in
% vertical mode never fires \everypar — so `\item \codeline{...}` printed the code
% at y=69.1 and the item's own "1." BELOW it at y=103.6 (the number appears after
% its content, attached to nothing). \if@inlabel\leavevmode\fi is CONDITIONAL on
% purpose: an unconditional \leavevmode would open an empty paragraph line when
% \codeline follows ordinary prose. Verified both ways — in a list the label now
% sits on the code's line at its normal x, and after prose the geometry is
% byte-identical to before.
\makeatletter
\newcommand{\codeline}[1]{%
  \if@inlabel\leavevmode\fi
  \par\nopagebreak
  \vspace{\abovedisplayskip}\nopagebreak
  % v3.62: was \centerline{...}, i.e. \hbox to\hsize{\hss #1 \hss}. \hss has
  % INFINITE SHRINK, so TeX could always "fit" the box and NEVER reported an
  % Overfull — the same \hss mechanism as BLD-27's \item label. Measured: a wide
  % \codeline in a defbox painted its first token ON the box's left border bar and
  % sat 9.5pt outside the frame interior on both sides, with 0 Overfull, 0 Underfull
  % and all four gates PASS; in bare prose it reached 21.0pt into BOTH margins.
  % \hfil has no shrink, so the overflow is reported and build.sh's ovf_vis/ovf_big
  % counters see it; \linewidth (not \hsize) makes it box-aware, so a \codeline
  % inside a colored box centres on the box interior. Verified: a FITTING \codeline
  % keeps byte-identical glyph positions, and the box-crossing case now emits
  % "Overfull \hbox (19.03pt too wide)". A wide \codeline should be split in two.
  \hbox to\linewidth{\hfil\code{#1}\hfil}%
  \nopagebreak\vspace{\belowdisplayskip}}
\makeatother

% ---- a thin decorative rule helper ----
\newcommand{\sep}{\par\vspace{2pt}\noindent\textcolor{cRule}{\rule{\linewidth}{0.6pt}}\par\vspace{4pt}}

% ---- colors ----
\definecolor{cBlue}{HTML}{1F4E79}   \definecolor{cBlueBg}{HTML}{EAF1F8}
\definecolor{cGreen}{HTML}{2E6B4F}  \definecolor{cGreenBg}{HTML}{EAF4EE}
\definecolor{cAmber}{HTML}{9C6A00}  \definecolor{cAmberBg}{HTML}{FBF3E2}
\definecolor{cPurple}{HTML}{5B2C6F} \definecolor{cPurpleBg}{HTML}{F3ECF7}
\definecolor{cTeal}{HTML}{0E6E6E}   \definecolor{cTealBg}{HTML}{E6F4F4}
\definecolor{cRed}{HTML}{8E2B2B}    \definecolor{cRedBg}{HTML}{FBEAEA}
\definecolor{cGray}{HTML}{555555}   \definecolor{cRule}{HTML}{1F4E79}

% ---- math operators ----
\DeclareMathOperator{\Tr}{Tr}
\DeclareMathOperator{\rank}{rank}
\DeclareMathOperator{\diag}{diag}
\DeclareMathOperator{\sgn}{sgn}

% ---- bra-ket (quantikz defines \ket/\bra/\braket; add the rest safely) ----
\providecommand{\ket}[1]{\left|#1\right\rangle}
\providecommand{\bra}[1]{\left\langle#1\right|}
\providecommand{\braket}[2]{\left\langle#1\middle|#2\right\rangle}
\newcommand{\dyad}[1]{\left|#1\right\rangle\!\left\langle#1\right|}
% v3.21: general outer product |a><b| and expectation value <A> — agents reach
% for these by reflex (esp. \expval); quantikz defines neither, so their absence
% threw "Undefined control sequence" in a real build. \providecommand = safe.
\providecommand{\ketbra}[2]{\left|#1\right\rangle\!\left\langle#2\right|}
\providecommand{\expval}[1]{\left\langle#1\right\rangle}
\newcommand{\abs}[1]{\left|#1\right|}
\newcommand{\norm}[1]{\left\lVert#1\right\rVert}
\newcommand{\half}{\tfrac12}
\newcommand{\R}{\mathbb{R}}\newcommand{\C}{\mathbb{C}}\newcommand{\Z}{\mathbb{Z}}
% identity operator. NB: \mathbb{1} silently renders a WRONG glyph (msbm has no
% blackboard digits). Use \mathbb{I} = 𝕀 (works with amssymb, no extra package).
% Do NOT "fix" this back to \mathbb{1}.
\newcommand{\id}{\mathbb{I}}

% ---- physical units (v3.35): \punit — upright unit with thin space ---------
% Usage: $300\punit{K}$, $80\punit{mV/decade}$, $10^{7}\punit{cm/s}$ — the
% NUMBER AND UNIT LIVE IN ONE MATH ISLAND. Promoted verbatim from a 199-page
% production lakit (advanced-electronic-devices, dozens of uses).
% Two measured traps this design avoids (both render-proven):
%   1. siunitx SCRAMBLES under RTL — \SI{3.5}{eV} renders "eV 3.5", exponents
%      detach ("Hz 10 x 10^9"): its quantity assembly is not bidi-isolated.
%      Do NOT load siunitx in this pipeline.
%   2. A text-mode unit next to a bare-prose number is TWO adjacent LTR runs
%      and reorders ("5 mm" -> "mm5" — the BIDI-2/3 class). \punit is therefore
%      math-only ON PURPOSE: prose misuse fails LOUDLY at compile (\mathrm
%      outside math) instead of scrambling silently.
% Domain shorthands (\MeV, \fm...) stay per-document: \newcommand{\MeV}{\,\text{\en{MeV}}}
\providecommand{\punit}[1]{\,\mathrm{#1}}

% ============================================================================
% Colored boxes — mdframed (NOT tcolorbox, which breaks under bidi).
% ★★ v3.4: innerleftmargin/innerrightmargin give code listings room so their
% frame/line-numbers do not spill out the box border (see listing styles below).
% ============================================================================
% v3.61: everyline/split skips were applied to `tablebox` only (BLD-14, v3.37) —
% the six colored boxes below never got them, so they still split OPEN across a page
% break: the page-N segment gets top+sides and NO bottom rule, the page-N+1 segment
% sides+bottom and NO top rule. Measured on the shipped fleet this is a COVERAGE GAP,
% not a re-report of BLD-14: 0 split-open `tablebox` segments (that fix holds) vs 62
% split-open COLORED boxes across 7 documents. Verified after the fix: both segments
% close, and the frametitle is NOT duplicated on the continuation. BLD-14's caveat
% still applies — a box taller than a full page grinds mdframed's split-trial loop, so
% split such content manually.
% ============================================================================
% v3.72: frametitle={} leads EVERY \newmdenv option list below. mdframed options
% are GROUP-scoped: a box nested inside \begin{exbox}[frametitle={...}] used to
% INHERIT the outer title — measured, the identical title string rendered twice,
% the second riding a title-bar fill inside the INNER box's frame, for all seven
% environments in both directions. Both natural forms were broken (= preamble
% bug, not an authoring rule). The explicit reset makes every box titleless
% unless ITS OWN [frametitle={...}] says otherwise; verified both directions:
% no duplicate on the nested box, an explicit title still renders.
% Geometry is shared with the optional semantic-width table container. Themes
% select their own existing semantic style below; neither path fixes columns.
\mdfdefinestyle{hllpdfsemanticgeometry}{linewidth=1.1pt,
  innertopmargin=6pt,innerbottommargin=7pt,skipabove=9pt,skipbelow=9pt,
  innerleftmargin=9pt,innerrightmargin=9pt,everyline=true,
  splitbottomskip=6pt,splittopskip=12pt,roundcorner=3pt}
\newcommand{\hllpdfsemanticframestyle}{hllpdfsemanticgeometry}
\newmdenv[frametitle={},style=hllpdfsemanticgeometry,linecolor=cBlue,backgroundcolor=cBlueBg,
  frametitlebackgroundcolor=cBlue,frametitlefontcolor=white,
  ]{defbox}   % הגדרה
\newmdenv[frametitle={},style=hllpdfsemanticgeometry,linecolor=cGreen,backgroundcolor=cGreenBg,
  frametitlebackgroundcolor=cGreen,frametitlefontcolor=white,
  ]{thmbox}   % משפט/טענה
\newmdenv[frametitle={},style=hllpdfsemanticgeometry,linecolor=cAmber,backgroundcolor=cAmberBg,
  frametitlebackgroundcolor=cAmber,frametitlefontcolor=white,
  ]{notebox}  % הערה
\newmdenv[frametitle={},style=hllpdfsemanticgeometry,linecolor=cPurple,backgroundcolor=cPurpleBg,
  frametitlebackgroundcolor=cPurple,frametitlefontcolor=white,
  ]{exbox}    % דוגמה
\newmdenv[frametitle={},style=hllpdfsemanticgeometry,linecolor=cRed,backgroundcolor=cRedBg,
  frametitlebackgroundcolor=cRed,frametitlefontcolor=white,
  ]{warnbox}  % טעות נפוצה
\newmdenv[frametitle={},style=hllpdfsemanticgeometry,linecolor=cTeal,backgroundcolor=cTealBg,
  frametitlebackgroundcolor=cTeal,frametitlefontcolor=white,
  ]{keybox}   % רעיון מרכזי

% ---- codebox (v4.5.0): the dedicated container for code -------------------
% An OPTION, not a policy: asmblock unboxed stays canonical for a full listing,
% and a snippet inside any of the six colored boxes still works as before.
% What makes codebox worth its own environment is that it is the ONE framed box
% that keeps its margin line numbers.
%
% v4.4.0 shipped it WITHOUT them, on the strength of v3.37's "no margin knob can
% save them" plus four fresh probes that agreed. That was wrong: every probe had
% moved the code, the box, or both — never the key that carries the NUMBER
% COLUMN. Measured on the page:
%   resetmargins=false   the listing follows THIS box instead of the page,
%                        moving the code from x=75.26 to x=97.54
%   xrightmargin=26pt    under RTL this is the PHYSICAL LEFT indent, and the
%                        only key that carries the \llap'ed digits with it:
%                        x=64.28 -> 90.18, clearing the rule at 65.64 by 23.6pt
%   xleftmargin=26pt     moves the code and LEAVES THE DIGITS on the rule —
%                        the asymmetry that hid the fix for eleven versions
% resetmargins is hoisted out of the four styles (see the \lstset below them) so
% this override is possible at all; every other listing still resolves to true.
% Pinned by tests/codebox_numbers.tex + probe_lstnum.py --inside-rule, whose
% negative control strips xrightmargin and watches the digits fall back.
\definecolor{cSlate}{HTML}{3F4A56}  \definecolor{cSlateBg}{HTML}{F2F4F7}
\newmdenv[frametitle={},leftline=true,rightline=true,topline=true,bottomline=true,
  linecolor=cSlate,backgroundcolor=cSlateBg,linewidth=0.9pt,
  frametitlebackgroundcolor=cSlate,frametitlefontcolor=white,
  innertopmargin=6pt,innerbottommargin=7pt,skipabove=9pt,skipbelow=9pt,
  innerleftmargin=9pt,innerrightmargin=9pt,everyline=true,splitbottomskip=6pt,splittopskip=12pt,
  roundcorner=3pt]{codebox}   % קוד
% The two keys that make in-box numbering work, both measured on the page:
%   resetmargins=false  the listing follows THIS box instead of the page
%   xrightmargin        under RTL this is the PHYSICAL LEFT indent, and it is
%                       the only one that carries the \llap'ed number column
%                       with it. xleftmargin moves the code and leaves the
%                       digits behind on the rule (measured: code 97.54,
%                       digits 64.28 — the v3.37 defect, reproduced); with
%                       xrightmargin=26pt the digits land at x=90.18, a
%                       clearance of 23.6pt from the rule at 65.64.
% This asymmetry is why v3.4's note had the semantics backwards and v3.37
% concluded the digits could not be saved.
% <hllpdf:listing-container-state>
% v4.29: all state is local to the current environment group. A semantic
% callout masks number text, while a nested codebox restores the formatter
% saved before that mask. Save it once across successive semantic wrappers;
% an explicit numbers=none (or an explicitly empty formatter) remains valid.
\makeatletter
\newif\ifhllpdf@incodebox
\newif\ifhllpdf@englishcodebox
\newif\ifhllpdf@lstnumbermasked
\let\hllpdf@savedlstnumber\thelstnumber
\makeatother
% </hllpdf:listing-container-state>
% <hllpdf:codebox-gutter>
\makeatletter
\AtBeginEnvironment{codebox}{%
  \hllpdf@incodeboxtrue
  % Snapshot the OUTER language before asmblock selects English.
  \iflanguage{english}{\hllpdf@englishcodeboxtrue}{\hllpdf@englishcodeboxfalse}%
  \ifhllpdf@lstnumbermasked
    \let\thelstnumber\hllpdf@savedlstnumber
    \hllpdf@lstnumbermaskedfalse
  \fi
}
% listings' AddToHook is GLOBAL. Register once, and gate on local box state;
% adding an unconditional hook at every codebox entry leaks into later listings.
\lst@AddToHook{PreInit}{%
  \ifhllpdf@incodebox
    % Apply every listing override here, inside the LISTING's own group.
    % Setting margins at box entry leaks them into nested semantic callouts.
    \lstset{resetmargins=false,xrightmargin=26pt}%
    % Styles resolve before PreInit. In LTR the physical number gutter follows
    % xleftmargin; in Hebrew retain the measured inherited xrightmargin path.
    \ifhllpdf@englishcodebox
      \lstset{xleftmargin=26pt,xrightmargin=0pt}%
    \fi
    % ONE box, not two: suppress only the frame/background/skips of a listing
    % inside this codebox. The next unboxed listing keeps its own style.
    \lstset{frame=none,backgroundcolor=\color{cSlateBg},aboveskip=0pt,belowskip=0pt}%
  \fi
}
\makeatother
% </hllpdf:codebox-gutter>

% v3.37: line numbers inside a COLORED BOX are sliced by the box border —
% measured: digit span x=63.5-67.3 vs border line at x=65.2 (the border runs
% THROUGH the digits), and the inter-border gap (3.5pt) is narrower than a
% digit (3.8pt), so no margin knob can save them (under resetmargins,
% xleftmargin is inert in-box — measured: code at x=75.26 with 2.2em AND
% 3.0em). A 3-6 line boxed snippet does not need line numbers; blank the
% number TEXT inside every colored box (\thelstnumber is what gets printed;
% the style's numbers=left survives, so nothing else shifts — code position
% byte-frozen at x=75.26, verified). Unboxed listings keep their numbers.
% v3.78: the blanking is coupled to the FRAME, not to the box name —
% clean_theme.tex registers LATER hooks (registration order wins) that
% restore the number text inside the boxes it renders frameless.
% v4.5.0 RETRACTS the "no margin knob can save them" clause above. It holds only
% under resetmargins=true (the shipped default): there xleftmargin is inert
% in-box, as measured. With resetmargins=false, xrightmargin DOES carry the
% \llap'ed digit column clear of the rule — measured; the codebox \lstset that
% does it is the one above, in the codebox block. The loop below therefore
% covers SIX boxes, not "every colored box": codebox opts out and keeps its
% numbers.
% <hllpdf:semantic-listing-mask>
\makeatletter
\@for\HEB@lstbox:=defbox,thmbox,notebox,exbox,warnbox,keybox\do{%
  \expandafter\AtBeginEnvironment\expandafter{\HEB@lstbox}{%
    \ifhllpdf@lstnumbermasked\else
      \let\hllpdf@savedlstnumber\thelstnumber
    \fi
    \hllpdf@lstnumbermaskedtrue
    \def\thelstnumber{}%
    % The nearest semantic container owns its listing presentation. On exit,
    % TeX's group restores the enclosing codebox state and formatter.
    \hllpdf@incodeboxfalse}}
\makeatother
% </hllpdf:semantic-listing-mask>

% ============================================================================
% Unified TABLE presentation: subtle uniform slate frame + auto-numbered
% caption rendered INSIDE the frame. Use:
%   \begin{tablebox}\centering  <tabular>  \tabcap{<description>}\end{tablebox}
% Replaces the buggy \caption* (which emitted a stray "טבלה N : *" line) and
% provides an explicit framed option + numbered caption "טבלה N.M: תיאור".
% Optional intro text may be placed at the TOP, before the tabular, inside the box.
% ============================================================================
\definecolor{cTblLine}{HTML}{B9C4D2}  \definecolor{cTblBg}{HTML}{FAFBFC}
% v3.23: class-aware — under `article` there is no chapter counter (the promised
% report->article switch used to die with "No counter 'chapter' defined").
\makeatletter
\@ifundefined{c@chapter}{%
  % v3.31: in article the counter resets per SECTION, so an unprefixed number
  % repeated ("טבלה 1." in every section — measured on a two-section render).
  % Same conditional idea as the report branch: prefix the section number once
  % a real \section exists; a pre-section table stays "טבלה 1".
  \newcounter{tbl}[section]
  \renewcommand{\thetbl}{\ifnum\value{section}>0 \thesection.\fi\arabic{tbl}}%
}{%
  \newcounter{tbl}[chapter]
  % v3.30: a report-class doc that uses SECTIONS ONLY (chapter stays 0 —
  % the common short-doc pattern) used to print "טבלה 0.1"; the chapter
  % prefix now appears only once a real \chapter exists (render-verified
  % both ways: "טבלה 1." chapterless, "טבלה 1.1" with chapters).
  \renewcommand{\thetbl}{\ifnum\value{chapter}>0 \thechapter.\fi\arabic{tbl}}%
  % v3.64: the SAME leak in \thesection itself, left behind by v3.30. That fix
  % guarded the CAPTION counter and stopped there, so a chapterless report still
  % printed the heading "0.1 סעיף ראשון" and \ref to it printed "0.2" — measured:
  %   p1 y=65.2  '…0.1'   p1 y=144.9 '…0.2'   p1 y=183.4 '…0.2.1'
  % report.cls hardcodes \thesection as \thechapter.\@arabic\c@section, and the
  % leading 0 is the chapter counter that no \chapter ever incremented. Fixing
  % \thesection fixes the whole chain — \thesubsection and \thesubsubsection are
  % defined in terms of it. \theequation and \thetbl/\thefig were already guarded,
  % which is exactly why this survived: every NUMBER around the heading looked
  % right, so the heading's own 0 read as a style choice rather than a bug.
  % \appendix is safe: report.cls resets chapter to 0 and redefines \thechapter to
  % \@Alph\c@chapter, so an appendix \chapter makes the counter 1 and the guard
  % emits "A.1" unchanged. Render-verified in all four modes (chapterless,
  % chaptered, appendix+chapter, appendix+section-only).
  \renewcommand{\thesection}{\ifnum\value{chapter}>0 \thechapter.\fi\arabic{section}}%
}
\makeatother
% v3.37: everyline=true — if the box DOES split across pages, every segment
% gets a CLOSED frame (measured before: the split ran open-ended into the page
% edge with no bottom border, and the caption could strand alone at the next
% page top — "טבלה A.2" shipped that way). The split skips give the segments
% breathing room. In practice, with the caption GLUE below, every previously-
% split table in the corpus now moves WHOLE instead (the white-gap tradeoff,
% documented benign); everyline stays as the safety net for a genuinely
% page-tall table. NB: a table taller than ONE FULL PAGE can grind mdframed's
% split-trial loop (reproduced: minimal doc, minutes of 1pt-step Overfull
% probes) — break such a monster into two tables.
\mdfdefinestyle{hllpdftablegeometry}{leftline=true,rightline=true,topline=true,bottomline=true,
  linewidth=0.9pt,
  innertopmargin=9pt,innerbottommargin=8pt,skipabove=11pt,skipbelow=11pt,
  innerleftmargin=11pt,innerrightmargin=11pt,roundcorner=4pt}
\newcommand{\hllpdftableframestyle}{hllpdftablegeometry}
\newmdenv[frametitle={},style=hllpdftablegeometry,
  linecolor=cTblLine,backgroundcolor=cTblBg,
  % v4.2.0 (adopted): tablebox is ATOMIC. Short tableboxes near page bottoms
  % repeatedly entered mdframed's split trial, emitting benign Overfull-vbox
  % bursts and occasionally separating the intro prose from the tabular.
  % nobreak=true moves the complete short table to the next page instead. The
  % output-active Overfull-vbox gate still catches a box taller than a page;
  % everyline/split skips remain as defensive styling should that path ever
  % re-arm, but the canonical path never publishes a split tablebox.
  nobreak=true,
  everyline=true,splitbottomskip=6pt,splittopskip=12pt,
  % v3.79 (BLD-37 table side, correction): startinnercode/endinnercode run
  % INSIDE mdframed's body vbox (mdframed.sty:848 / :856+863) -- the second one
  % after \par\unskip, i.e. after line breaking, which is the only moment the
  % \tabcap anchor's true offset from the top of the box EXISTS. See
  % \lakintabboxclose below. These two keys are the ONLY tablebox change.
  startinnercode={\lakintabboxopen},endinnercode={\lakintabboxclose}]{tablebox}
% v3.37: \nopagebreak GLUE — \vsplit honors penalties, so the caption can no
% longer strand alone at the top of the next page (it takes its rows along, or
% the whole box moves). Render-verified on the real shipped orphan.
% v3.59: \NR@gettitle is what nameref injects into the stdlib title-bearing
% commands (\@sect, \@chapter, \caption) to fill \@currentlabelname. \figcap and
% \tabcap are INDEPENDENT macros — they never called it, so \@currentlabelname kept
% whatever the last heading left and \nameref{fig:x} silently printed the enclosing
% SECTION's title instead of the caption (proven in the .aux: field 3 of \newlabel
% was the heading text while field 1 and the anchor `fig.1.1` were correct — so this
% is a THIRD, distinct head next to BLD-13's counter fix and BIDI-38's autoref-name
% fix, neither of which writes that field). \@ifundefined keeps it inert if
% hyperref/nameref is ever dropped. Guarded helper because \figcap is defined
% after \makeatother.
% NB both helpers below need \makeatletter — this part of the preamble is OUTSIDE
% any letter-catcode region, so an unwrapped \@currentlabel parses as LaTeX's \@
% and dies with "You can't use \spacefactor in vertical mode" (measured).
\makeatletter
\newcommand{\lakinrtitle}[1]{\@ifundefined{NR@gettitle}{}{\NR@gettitle{#1}}}
% v3.59: \hitem{<letter>}{<key>} — a REFERENCEABLE Hebrew-lettered list item.
%   \begin{enumerate} \hitem{א}{it:a} הסעיף הראשון. \end{enumerate}  ->  \ref{it:a} = (א)
% Why this exists: automatic Hebrew list labels fail silently, so SKILL.md mandates
% an EXPLICIT \item[(\textbf{א})] — but LaTeX puts \refstepcounter INSIDE
% \if@noitemarg (latex.ltx:12964), and \item[...] never sets that flag. So an
% explicit item steps no counter and never touches \@currentlabel: a \label there
% binds to the AMBIENT one and \ref prints the CHAPTER/SECTION number where the
% item letter belongs (proven in the .aux: {chapter.1}). Same doctrine as \figcap —
% own the label state before \label. \phantomsection matters: without it the
% hyperlink still lands on the chapter head (measured); with it the anchor is
% section*.N and the bookmark outline stays clean.
\newcommand{\hitem}[2]{\item[(\textbf{#1})]%
  \def\@currentlabel{(#1)}\def\@currentlabelname{#1}\phantomsection\label{#2}}
% NB (v3.66): \hitem is \ref-only. \autoref{key} calls it a SECTION ("סעיף (א)")
% because \phantomsection brands the anchor section*.N, and \nameref{key} prints
% the bare letter. For \autoref-able items prefer \hebenumlabels below.
%
% v3.66: \hebenumlabels — OPT-IN automatic Hebrew enumerate letters, done right.
%   {\hebenumlabels\begin{enumerate}\item\label{q:a} ...\end{enumerate}}
%   -> labels (א) (ב) (ג); \ref{q:a} = א; anchor Item.N (real, per-item).
% This retires the belief that made \item[(\textbf{א})] + \hitem necessary.
% "Automatic Hebrew labels fail silently" was TRUE but unexplained, and the
% mechanism is now bisected (plain article vs +enumitem, measured): ENUMITEM
% PRE-EXPANDS the label expression in an \edef BEFORE \refstepcounter, with the
% counter still at 0. Tokens it shields by name (\arabic,\alph,\roman) survive;
% a plain \newcommand \ifcase wrapper evaluates at 0 -> the EMPTY zero-branch ->
% the label prints as "()" with the letter GONE (and \the\c@enumi prints a
% literal 0, \number\value dies fatally pointing at \item). The cure is one
% property, not a doctrine: the wrapper must be ROBUST (survives the \edef
% unexpanded), while \theenumi stays EXPANDABLE (so \refstepcounter writes the
% real letter to the .aux). Verified: labels (א)(ב)(ג), \ref -> א/ב/ג each item
% its own value, .aux = {א}{...}{Item.1} — versus the explicit-\item doctrine
% where BIDI-50/\hitem exist purely to patch the counter that never stepped.
\def\@hebalph#1{\ifcase#1\or א\or ב\or ג\or ד\or ה\or ו\or ז\or ח\or ט\or י\or
  יא\or יב\or יג\or יד\or טו\or טז\or יז\or יח\or יט\or כ\else\@ctrerr\fi}
\DeclareRobustCommand*\hebalph[1]{\expandafter\@hebalph\csname c@#1\endcsname}
\newcommand{\hebenumlabels}{%
  \renewcommand{\theenumi}{\@hebalph\c@enumi}%
  \renewcommand{\labelenumi}{(\textbf{\hebalph{enumi}})}}
\makeatother
% v3.63: \mbox the LABEL. \thetbl ("7.1") and a caption description that STARTS
% WITH A DIGIT ("32 הרגיסטרים…") are two European-Number runs separated only by
% neutrals (the label's own "." and \enspace). Under the RTL paragraph base the UBA
% resolves those neutrals to the RTL level, so both numbers sit in ONE reversed span
% and EXCHANGE VISUAL POSITIONS. Measured on the shipped PDF: a table whose real
% number is 7.1 printed "טבלה 32" — a different, EXISTING table number — while \ref
% correctly printed 7.1, so the caption silently contradicted its own cross-reference,
% and the description lost its leading count. The source contains no adjacency at all;
% it is manufactured by this macro, which is why every source linter was blind.
% \mbox makes the label one atomic box so nothing can reorder across it. Verified
% zero-regression: an ordinary Hebrew caption keeps identical x-geometry.
% One paragraph style covers label, number, separator and description.
% Explicit project/template overrides remain possible; the label hook is kept
% for compatibility, with an identity default. Counters/anchors stay unchanged.
\newcommand{\lakitcaptionstyle}{\normalfont\small\color{black}\hypersetup{allcolors=black}}
\newcommand{\lakitcaptionlabelstyle}[1]{#1}
\newcommand{\lakitcaptionpunctuation}{:}
% Native floats, captionof and longtable use the same presentation policy.
% subcaption is optional and inherits the same font rather than a second size.
\usepackage{caption}
\DeclareCaptionFont{lakit}{\lakitcaptionstyle}
\DeclareCaptionLabelFormat{lakit}{\lakitcaptionlabelstyle{#1~\en{#2}}}
\DeclareCaptionLabelSeparator{lakit}{\lakitcaptionpunctuation\enspace}
\captionsetup{font=lakit,labelfont={},textfont={},labelformat=lakit,
  labelsep=lakit,justification=centering}
\makeatletter
\AtBeginDocument{\@ifpackageloaded{subcaption}{%
  \captionsetup[sub]{font=lakit,labelfont={},textfont={},labelsep=lakit}}{}}
\makeatother
% v3.64: OPT-IN list-of-figures/tables entry via an optional short caption.
% \listoffigures and \listoftables are inherited from report.cls and LOOK available,
% but \figcap/\tabcap step the custom \c@fig/\c@tbl counters and never call
% \addcontentsline — so both printed their heading and NOTHING ELSE. Measured: a doc
% with one \figcap and one \tabcap rendered page 1 as 'רשימת האיורים\n1\n', and the
% .lof held only \@chapter's \addvspace, zero \contentsline entries. An empty list
% under a correct Hebrew heading reads as "this document has no figures".
%
% WHY OPTIONAL, AND NOT AUTOMATIC. Writing #1 unconditionally was implemented first
% and REJECTED: a caption containing \codeline compiles CLEANLY today (verified on
% the pre-v3.64 preamble: 0 errors, PDF produced) but becomes FATAL once it is also
% written to the .lof — \codeline is \hbox to\linewidth{...}, and a full-width box in
% a list entry dies with "\pdfextension endlink cannot be used in vertical mode", no
% output PDF. Gaining a feature nobody uses (0 of 13 collections call \listoffigures)
% at the price of breaking a construct that works is a bad trade. Stock \caption has
% exactly this limitation and exactly this escape hatch, so the convention is familiar.
% \en{}, \code{}, \h{} and inline math in the short text are all safe — verified, they
% survive the write and typeset correctly in the list (RTL geometry confirmed: number
% rightmost at x=509.9, title left of it, dots and folio at the far left).
% With no optional argument the macro is byte-identical to v3.63 — zero risk to every
% existing document, and \listoffigures lists exactly the captions that opted in.
% BLD-37, TABLE side (the \figcap v3.74 idea, adapted): clicking a \ref to a
% table used to land MID-table — hyperref's automatic anchor inside
% \refstepcounter was emitted on the still-open horizontal line whose TAIL is
% the tabular (an hbox whose baseline is the table's vertical CENTER), so the
% jump scrolled the top half of the table above the viewport. The v3.74 entry
% deferred this side claiming "tablebox splits — no single box to measure";
% that deferral dissolves: the measurement target is the TABULAR itself at
% \tabcap ENTRY, while it is still the tail of the open list — the same
% situation \lakin@grabfig handles at an island's end (and the fleet has 0
% split tableboxes). So \lakin@grabtab runs as the FIRST token, before any
% typesetting can push the box away, and the anchor is re-emitted raised above
% the measured tabular. \lakintabraise is a SEPARATE register from
% \lakinfigraise (defined with it, below the LTR-island block): a captionless
% island intentionally leaves the fig register set, and a table must never
% inherit that stale figure height.
%
% WHY the anchor stays AT THE ENTRY SPOT and is NOT deferred into the caption
% paragraph the way \figcap defers it: at \par TeX DROPS a final glue item,
% and the automatic anchor's smash box sitting BEHIND the interword glue kept
% that glue alive — so every centered tabular has always sat half a space off
% true center. Deferring the anchor removed that box and was measured as 982
% differing glyph records: EVERY tabular line shifted 1.30pt (2.61pt for
% pattern 7's two minipages) while captions stayed put. Emitting the raised
% anchor at the old spot — through hyperref's own \Hy@raisedlink, exactly the
% macro the automatic path uses — keeps the line's node structure, and so the
% A/B glyph diff at 0, BY CONSTRUCTION. \Hy@raisedlink reads the raise from
% \HyperRaiseLinkDefault (default \baselineskip = the old mid-table dest); the
% group-local \edef freezes measured-height+42pt into it, and \smash inside
% \Hy@raisedlink hides any raise from the line's extent. \HyCnt@ProvideTheH-
% Counter keeps the dest NAME byte-identical to the automatic one (it defines
% \theHtbl, which \hyper@makecurrent prefers — without it the first measured
% caption was named tbl.1 where the automatic path said tbl.0.1, measured).
% vmode entry, or anything but a box at the tail (no tabular to measure),
% falls back to the plain \refstepcounter — byte-identical to the old path.
\makeatletter
\newcommand{\tabcap}[2][]{\lakin@grabtab
  % v3.79: arm the box-close correction. Set unconditionally and in BOTH
  % branches below's scope, so the plain-\refstepcounter fallback is corrected
  % too; \lakintabboxopen clears it at every \begin{tablebox}, so a \tabcap
  % outside a tablebox can never leave it armed for a later box.
  \global\lakintabcapseentrue
  \ifdim\lakintabraise>0pt
    \H@refstepcounter{tbl}\HyCnt@ProvideTheHCounter{tbl}\hyper@makecurrent{tbl}%
    \begingroup
      % A frameless table targets its measured top. The frame's historical
      % cushion belongs only to tablebox and is corrected at box close.
      \edef\HyperRaiseLinkDefault{\the\lakintabraise}%
      \iflakintabframed
        \edef\HyperRaiseLinkDefault{\the\dimexpr\lakintabraise+42pt\relax}%
      \fi
      \Hy@raisedlink{\hyper@anchorstart{\@currentHref}\hyper@anchorend}%
    \endgroup
    \global\lakintabraise 0pt\relax
  \else
    \refstepcounter{tbl}%
  \fi
  \lakinrtitle{#2}\par\nopagebreak\addvspace{6pt}%
  \def\lakin@short{#1}\ifx\lakin@short\@empty\else
    \addcontentsline{lot}{table}{\protect\numberline{\thetbl}{\protect\ignorespaces #1}}\fi
  {\centering\lakitcaptionstyle\mbox{\lakitcaptionlabelstyle{\tablename~\en{\thetbl}\lakitcaptionpunctuation}}\enspace #2\par}}
\makeatother

% ============================================================================
% Unified FIGURE caption — the \tabcap twin for TikZ/plot figures (v3.34.0).
% A figure (tikzpic) sits inline in the prose; put \figcap BELOW it for a
% numbered, \ref-able caption "איור N.M: תיאור" in the same style as \tabcap:
%   \begin{tikzpic} ... \end{tikzpic}
%   \figcap{\label{fig:foo}description}     % \label FIRST, inside the argument
% WHY a macro, not a hand-written gray line: a manual
%   {\centering\small\textcolor{cGray}{איור: ...}} with a \label inside has NO
% \refstepcounter, so the \label binds to the AMBIENT \@currentlabel (the last
% equation) — \ref then prints an equation number AS the figure number
% (measured on a real page: "איור 48" where 48 is equation 48). \figcap steps a
% real \c@fig counter first, so \label is correct, numbering is consistent, and
% refs are stable. Class-aware like \thetbl (chapter- or section-prefixed).
% ============================================================================
\makeatletter
\@ifundefined{c@chapter}{%
  \newcounter{fig}[section]
  \renewcommand{\thefig}{\ifnum\value{section}>0 \thesection.\fi\arabic{fig}}%
}{%
  \newcounter{fig}[chapter]
  \renewcommand{\thefig}{\ifnum\value{chapter}>0 \thechapter.\fi\arabic{fig}}%
}
\makeatother
\makeatletter
\newcommand{\figcap}[2][]{%
  \ifdim\lakinfigraise>0pt
    % v3.74: step WITHOUT hyperref's automatic anchor, then emit the anchor
    % raised above the measured figure. \H@refstepcounter is hyperref's saved
    % pre-hook \refstepcounter (sets the counter AND \@currentlabel — \label
    % in the caption arg binds exactly as before); \hyper@makecurrent supplies
    % \@currentHref, so the dest NAME is identical to the old automatic one —
    % \ref/\autoref/bookmarks are unchanged, only the dest'S POSITION moves.
    \H@refstepcounter{fig}\hyper@makecurrent{fig}%
    % The anchor is emitted INSIDE the caption paragraph below, as a zero-size
    % raisebox — NOT here in vertical mode, where a \raisebox starts an empty
    % paragraph and adds a blank line (measured: every glyph after the caption
    % shifted +18.8pt). Inside an existing line a [0pt][0pt] box is inert.
    \edef\lakin@anch{\@currentHref}%
    % \edef FREEZES the measured height into the pending macro: the deferred
    % \raisebox is typeset inside the caption paragraph AFTER the register is
    % reset to the 0pt sentinel below — a plain \def read 0 at use time
    % (measured: LAKINRAISE=0.0, dest moved 14pt instead of above the figure).
    \edef\lakin@pend{\noexpand\raisebox{\the\dimexpr\lakinfigraise+42pt\relax}[0pt][0pt]{%
      \noexpand\hyper@anchorstart{\lakin@anch}\noexpand\hyper@anchorend}}%
    \global\lakinfigraise 0pt\relax
  \else
    \refstepcounter{fig}\let\lakin@pend\relax
  \fi
  \lakinrtitle{#2}\par\nopagebreak\addvspace{4pt}%
  \def\lakin@short{#1}\ifx\lakin@short\@empty\else
    \addcontentsline{lof}{figure}{\protect\numberline{\thefig}{\protect\ignorespaces #1}}\fi
  {\centering\lakin@pend\lakitcaptionstyle\mbox{\lakitcaptionlabelstyle{\figurename~\en{\thefig}\lakitcaptionpunctuation}}\enspace #2\par}}
\makeatother

% ============================================================================
% LTR islands.  EVERY tikzpicture / quantikz MUST live inside one, or the whole
% figure's coordinates and Latin labels flip. (Hebrew node text still uses \h{}.)
% ============================================================================
% v3.74 (BLD-37): the \figcap hyperlink DESTINATION used to land at the caption
% line — BELOW the figure — so clicking a figure reference scrolled the figure
% off-screen (measured fleet-wide: 97/97 fig dests in one collection sat 19-54pt
% under their figure's ink; a 347pt figure was fully hidden after the jump).
% hypcap-style fix in two parts: (1) at the island's END, \lakin@grabfig grabs
% the just-typeset centered line with \lastbox, measures it, and puts it back
% with \nointerlineskip — the interline glue above the line was already
% contributed and is NOT re-added, so the page geometry is unchanged to the
% glyph (verified); (2) \figcap then steps the counter WITHOUT hyperref's
% automatic anchor (\H@refstepcounter = the pre-hyperref \refstepcounter) and
% emits the anchor itself, raised by the measured height — the dest lands ABOVE
% the figure top. A captionless island leaves numbering untouched (nothing is
% stepped at the island; the grab flag is cleared at every island begin, so a
% stale height from a captionless figure cannot leak into a later caption).
\newsavebox{\lakinfigbox}
\newdimen\lakinfigraise  % real \dimen: tex.setdimen cannot target a \newlength SKIP register (measured lua error)
% BLD-37 table side: a SEPARATE register for \tabcap. The two must not share —
% a captionless island INTENTIONALLY leaves \lakinfigraise set (the grab runs
% at every island END; only the NEXT island's begin clears it), so a shared
% register would hand a \tabcap after a captionless figure that stale FIGURE
% height. \lakintabraise is written only by \lakin@grabtab at \tabcap entry
% and consumed/reset inside \tabcap itself, so it can never go stale.
\newdimen\lakintabraise
\makeatletter
% v3.74b: measure WITHOUT touching the list. The first draft grabbed the box
% with \lastbox and put it back — and \lastbox DISCARDS the node's baseline
% shift, so a vertically-anchored quantikz circuit rose 2.57pt and everything
% after it sank 5.32pt (measured on the A/B glyph diff). LuaTeX lets us READ
% the tail node's dimensions without removing anything: geometry is identical
% by construction. \lakinfigraise=0pt is the "no figure measured" sentinel.
\def\lakin@grabfig{%
  \ifhmode
    \directlua{
      local n = tex.nest[tex.nest.ptr].tail
      local GLUE, PEN = node.id('glue'), node.id('penalty')
      while n and (n.id == GLUE or n.id == PEN) do n = n.prev end
      local HL, VL = node.id('hlist'), node.id('vlist')
      if n and (n.id == HL or n.id == VL) then
        tex.setdimen('global', 'lakinfigraise',
                     (n.height or 0) + (n.depth or 0))
      else
        tex.setdimen('global', 'lakinfigraise', 0)
      end}%
  \else
    \global\lakinfigraise 0pt\relax
  \fi}
% BLD-37 table side: \lakin@grabfig's twin, NOT a shared helper — the Lua body
% differs in exactly one place, the skip set. Between \end{tabular} and \tabcap
% an otherlanguage wrapper (patterns 3/5 of example_tables) leaves a \write
% whatsit and a dir node ON TOP of the interword glue (measured node dump:
% "whatsit/write <- dir <- glue(2.62) <- hlist"), so the table reader must
% also step over whatsit/dir markers — all zero-size — to reach the tabular.
% \lakin@grabfig deliberately does NOT get the wider skip set: its glue/penalty
% body is the one the v3.74 A/B proof froze, and widening it would newly
% measure (and so newly raise) anchors in island shapes that today take the
% plain path — an unproven behavior change on the FIG side. Anything that is
% not a box behind those markers still hits the 0pt sentinel and \tabcap falls
% back to the plain \refstepcounter.
% SECOND deliberate difference: this stores the box's HEIGHT, not height+depth.
% \figcap's anchor is emitted in the caption paragraph BELOW the figure, so it
% must climb the full extent; \tabcap's anchor sits ON the tabular's own line,
% whose baseline is the vcentered tabular's MIDPOINT — height alone reaches
% the top edge, and adding the depth was measured to overshoot by half the
% table (dests 80-135pt above their frame, INSIDE the preceding table's box).
\def\lakin@grabtab{%
  \ifhmode
    \directlua{
      local n = tex.nest[tex.nest.ptr].tail
      local GLUE, PEN = node.id('glue'), node.id('penalty')
      local WHAT, DIR = node.id('whatsit'), node.id('dir')
      while n and (n.id == GLUE or n.id == PEN
                   or n.id == WHAT or n.id == DIR) do n = n.prev end
      local HL, VL = node.id('hlist'), node.id('vlist')
      if n and (n.id == HL or n.id == VL) then
        tex.setdimen('global', 'lakintabraise', n.height or 0)
      else
        tex.setdimen('global', 'lakintabraise', 0)
      end}%
  \else
    \global\lakintabraise 0pt\relax
  \fi}
\makeatother
% ============================================================================
% v3.79 (BLD-37 table side, CORRECTION): put the anchor at the TOP OF THE BOX,
% instead of guessing the distance to it.
% ============================================================================
% v3.77 raised the \tabcap anchor by (height of the box at the tail of the open
% horizontal list) + 42pt. That is a MEASUREMENT taken BEFORE line breaking, and
% it is short by every term it cannot see. Fleet-measured on the 12 built
% collections (299 tbl.* dests, ownership established through each table's own
% CAPTION, never through "which frame encloses the dest"): 3 dests landed INSIDE
% their own tablebox, below the frame top —
%   computer-architecture       tbl.14.1  -46.4pt   three stacked tabulars
%   computer-architecture       tbl.30.5  -13.5pt   tabular 2.85pt over \linewidth
%   advanced-electronic-devices tbl.B.1   -59.2pt   tabular 2.56pt over \linewidth
% Two missing terms, both reproduced in a minimal document (probe P2/P3):
%   (1) VERTICAL LIST ALREADY CONTRIBUTED inside the box — a second and third
%       tabular, or a heading above the table. The tail read sees only the LAST
%       box, so 77.7pt of stacked tabulars is invisible to the raise.
%   (2) WHICH LINE THE ANCHOR LANDS ON. When the tabular alone overfills
%       \linewidth (459.6121pt in a tablebox), the interword glue after
%       \end{tabular} becomes the chosen breakpoint and \Hy@raisedlink is
%       emitted on the NEXT line — one (depth + lineskip) lower. That term does
%       not EXIST at \tabcap time: line breaking has not happened yet. No
%       measurement taken there can ever recover it. (Proof: probe P3, the same
%       table 2.84912pt too wide, fails at -13.5pt; probe P4, the identical
%       table with \tabcolsep 2pt so it fits, is correct at +32.0pt.)
% So the raise is not computed at \tabcap at all. It is CORRECTED at
% \end{tablebox}, from inside mdframed's own body vbox (endinnercode runs after
% \par\unskip, i.e. after line breaking, while tex.nest[tex.nest.ptr] is still
% that vbox's vertical list). There the anchor's true offset from the top of the
% box content is not estimated but WALKED, and the \Hy@raisedlink shift is set
% so the dest lands 42pt above it — which is 42pt - innertopmargin(9pt) -
% linewidth(0.9pt) = 32.0pt above the FRAME top, the exact position the 295
% already-correct dests have today. That is why this is a correction and not a
% re-placement: the good sites do not move.
%
% WHY THE GLYPH DIFF IS 0 BY CONSTRUCTION — the only thing written is the
% `shift` field of the hlist that \Hy@raisedlink already built, and the guard
% below refuses to write it unless that box is 0x0x0. A zero-size box carries no
% dimension into its parent's packing, the packing already happened, and shift
% is not an input to it. Nothing else in the list is touched: no node is added,
% removed, resized or reordered. Measured, not assumed: assets/example_tables.tex
% and the whole computer-architecture collection A/B at 0 differing pymupdf
% rawdict glyph records.
%
% WHY NOT the alternatives (all prototyped and measured before this one):
%   * accumulate the vertical list at \tabcap time — fixes term (1), CANNOT fix
%     term (2); 2 of the 3 real failures survive. A fix that leaves the majority
%     of the defect standing is the false all-clear this repo exists to avoid.
%   * \zref-savepos two-pass — correct on all shapes, but it makes the dest
%     POSITION aux-derived (DEFECTS BLD-1/2), a single-pass compile silently
%     ships an unraised anchor, and the bare y-difference has no page guard.
%   * emit the anchor at \begin{tablebox} and let \tabcap adopt it — needs the
%     tbl counter stepped at box OPEN to keep the dest NAME byte-identical; a
%     captionless tablebox then either shifts every later table number or leaves
%     an orphan dest that the next table duplicates (a HARD build gate). This
%     correction touches neither the counter nor the name: \tabcap's naming path
%     is v3.77's, unchanged, so resolve_names() is byte-identical (verified).
%   * suppressing the line break (a penalty before the glue) — that MOVES the
%     tabular: the break currently discards the interword glue, so forbidding it
%     re-centers the two overfull tables. Reflow, rejected.
%
% SCOPE. The correction is armed only by \tabcap ITSELF (\lakintabcapseen),
% cleared at every box open, so a captionless tablebox is a no-op and a \tabcap
% hosted OUTSIDE a tablebox (the fleet has exactly one, computer-architecture
% tbl.14.2 in an exbox) keeps v3.77's measured raise. Both fallbacks are
% therefore never worse than today; and because hyperref's AUTOMATIC anchor uses
% the same \Hy@raisedlink shape, the plain-\refstepcounter path (vmode entry, or
% prose between \end{tabular} and \tabcap — probe P6, -51.5pt) is corrected too.
% KNOWN LIMIT — a tablebox that SPLITS across pages. The offset is measured from
% the top of the UNSPLIT body, so a dest in segment 2 gets raised past its own
% segment. MEASURED on a forced split (probe: 40 filler lines + a table): the
% dest went from y=585.0 on page 2 (inside segment 2, at the caption — the old
% BLD-37 symptom) to y=-113.4, i.e. ABOVE the page top. Viewers clamp that to the
% page top, which happens to show segment 2 whole, but an out-of-page coordinate
% is not something to ship on purpose. Hence the \textheight guard in
% lakintab.fix: a body TALLER THAN THE TEXT BLOCK cannot fit on any page, so it
% must split — the correction is skipped and v3.77's behaviour stands unchanged
% (verified: the probe returns to y=585.0). The guard is deliberately EXACT, not
% a page-fitting heuristic: it is sound (a taller-than-page body always splits)
% but PARTIAL (a shorter body that starts low on a page could still split). That
% residue is unreachable in practice — v3.37's everyline + the caption's
% \nopagebreak glue move such boxes WHOLE instead (probe N7: a 22-row table
% pushed down by 430pt moved whole, 0 segments), and the fleet has 0 split
% tableboxes, re-counted after this change. Reimplementing mdframed's own
% page-fitting arithmetic to close the residue would be a heuristic with its own
% failure mode, which is the trade this repo has repeatedly refused.
\newif\iflakintabcapseen
\newif\iflakintabframed
% The Lua lives in ONE \directlua at load time, and is written with NO
% backslash, %, # or ~ : \directlua expands its argument like \edef, so an
% active ~ would expand to \nobreakspace and a % would eat the rest of the line
% (the second is how the first draft of this block died). "not (a == b)" stands
% in for Lua's inequality operator for the same reason.
\directlua{
  lakintab = { }
  lakintab.HL = node.id("hlist")
  lakintab.VL = node.id("vlist")
  lakintab.RU = node.id("rule")
  lakintab.GL = node.id("glue")
  lakintab.KE = node.id("kern")
  lakintab.WH = node.id("whatsit")
  lakintab.TM = node.id("temp")
  lakintab.WD = node.subtype("pdf_dest")
  function lakintab.d(v)
    if v == nil then return 0 end
    if v < -1000000000 then return 0 end
    return v
  end
  function lakintab.scan(list, y, horiz, parent, target)
    local x = list
    while x do
      local i = x.id
      if i == lakintab.HL or i == lakintab.VL then
        if horiz then
          local b = y + x.shift
          if i == lakintab.HL then
            lakintab.scan(x.head, b, true, x, target)
          else
            lakintab.scan(x.head, b - lakintab.d(x.height), false, x, target)
          end
        else
          if i == lakintab.HL then
            lakintab.scan(x.head, y + lakintab.d(x.height), true, x, target)
          else
            lakintab.scan(x.head, y, false, x, target)
          end
          y = y + lakintab.d(x.height) + lakintab.d(x.depth)
        end
      elseif i == lakintab.WH then
        if x.subtype == lakintab.WD then
          local id = x.dest_id
          if type(id) == "string" and string.sub(id, 1, 4) == "tbl." then
            if horiz and parent and parent.width == 0
               and parent.height == 0 and parent.depth == 0 then
              parent.shift = parent.shift + (target - y)
            end
          end
        end
      elseif not horiz then
        if i == lakintab.RU then
          y = y + lakintab.d(x.height) + lakintab.d(x.depth)
        elseif i == lakintab.GL or i == lakintab.KE then
          y = y + lakintab.d(x.width)
        end
      end
      x = x.next
    end
  end
  function lakintab.fix(target, maxh)
    local nst = tex.nest[tex.nest.ptr]
    if nst == nil then return end
    local h = nst.head
    if h == nil then return end
    if h.id == lakintab.TM then h = h.next end
    local tot = 0
    local x = h
    while x do
      local i = x.id
      if i == lakintab.HL or i == lakintab.VL or i == lakintab.RU then
        tot = tot + lakintab.d(x.height) + lakintab.d(x.depth)
      elseif i == lakintab.GL or i == lakintab.KE then
        tot = tot + lakintab.d(x.width)
      end
      x = x.next
    end
    if tot > maxh then return end
    lakintab.scan(h, 0, false, nil, target)
  end
}
% @-FREE NAMES ON PURPOSE: these two are named in the `tablebox` key list at
% \newmdenv above, which is tokenised OUTSIDE any \makeatletter region — a
% \lakin@... there would tokenise as \lakin followed by the letters @tabopen.
\newcommand{\lakintabboxopen}{\lakintabframedtrue\global\lakintabcapseenfalse}
\newcommand{\lakintabboxclose}{%
  \iflakintabcapseen
    \directlua{lakintab.fix(-\number\dimexpr\lakintabcushion\relax,
                            \number\dimexpr\textheight\relax)}%
    \global\lakintabcapseenfalse
  \fi\lakintabframedfalse}
% The cushion above the box's CONTENT top. 42pt puts the dest 32.0pt above the
% frame top (42pt - innertopmargin 9pt - linewidth 0.9pt), i.e. exactly where
% the 295 correct dests already are — the correction moves only the broken ones.
\newcommand{\lakintabcushion}{42pt}
\newenvironment{qcirc}{\par\medskip\global\lakinfigraise 0pt\relax\begin{otherlanguage}{english}\begin{center}}
  {\csname lakin@grabfig\endcsname\end{center}\end{otherlanguage}\par\medskip}
\newenvironment{tikzpic}{\par\medskip\global\lakinfigraise 0pt\relax\begin{otherlanguage}{english}\begin{center}}
  {\csname lakin@grabfig\endcsname\end{center}\end{otherlanguage}\par\medskip}

% ============================================================================
% Code listings.  ★★ v3.4 THE BOX-OVERFLOW FIX:
%   resetmargins=true   → the listing IGNORES the enclosing environment's
%                         indentation and measures from the PAGE. v3.4 stated
%                         this backwards ("measured from the ENCLOSING box"),
%                         and v3.37 then blanked in-box line numbers because
%                         the page-anchored listing put its hanging digits on
%                         top of the box rule. Measured v4.5.0: with the key
%                         FALSE the listing follows the box and the code moves
%                         from x=75.26 to x=97.74, clearing the rule at 72.07.
%                         The key now lives in one \lstset below instead of in
%                         each style, so a box can override it — codebox does.
%   framexleftmargin=0  → the frame (left bar) does not extend left of the code.
%   xleftmargin keeps a small indent so the line numbers have room INSIDE.
% asmblock wraps the listing in an LTR island (code is LTR).
% ============================================================================
\lstdefinestyle{asm}{
  basicstyle=\ttfamily\small,
  backgroundcolor=\color{cBlueBg!55},
  frame=leftline,framerule=2.5pt,rulecolor=\color{cBlue},
  numbers=left,numberstyle=\tiny\color{cGray},numbersep=8pt,
  xleftmargin=2.2em,framexleftmargin=0pt,framexrightmargin=0pt,
  breaklines=true,columns=fullflexible,keepspaces=true,
  showstringspaces=false,aboveskip=8pt,belowskip=8pt,
  % v3.36: without a comment DELIMITER the commentstyle below never fires —
  % listings has no `language` here, so it does not know '#' opens a comment,
  % and every asm listing rendered flat black (measured: a whole-page pixel
  % sweep found ZERO cGreen glyphs across 135 shipped listings). morecomment
  % gives it the delimiter; color-only change, glyph widths untouched.
  morecomment={[l]{\#}},
  commentstyle=\color{cGreen},
}
\lstdefinestyle{cstyle}{
  language=C,basicstyle=\ttfamily\small,
  backgroundcolor=\color{cGreenBg!55},
  frame=leftline,framerule=2.5pt,rulecolor=\color{cGreen},
  numbers=left,numberstyle=\tiny\color{cGray},numbersep=8pt,
  xleftmargin=2.2em,framexleftmargin=0pt,framexrightmargin=0pt,
  breaklines=true,columns=fullflexible,keepspaces=true,
  showstringspaces=false,keywordstyle=\color{cBlue}\bfseries,
  commentstyle=\color{cGray}\itshape,aboveskip=8pt,belowskip=8pt,
}
% v3.36: pystyle/htmlstyle — colored twins of cstyle for the Python/HTML
% listings that used to repeat a ~200-char options blob with NO color
% (keywords rendered bold-black, comments italic-black — measured; C alone
% was colored). Geometry is copied VERBATIM from the shipped blob (frame,
% margins, skips) so swapping a call site to style=pystyle moves NOTHING —
% only keywordstyle/commentstyle/stringstyle are added (bold stays bold,
% italic stays italic, mono metrics identical => zero reflow, proven by
% per-page word-bbox diff). breaklines/breakatwhitespace stay per-site.
% v3.78: numbers=left joins both styles, with the exact numberstyle/numbersep
% keys asm/cstyle carry — fleet-wide consistency (author decision: all four
% styles numbered; "asm/cstyle numbered, py/html not" was an accident of the
% v3.36 geometry freeze, not a design). The numbers hang LEFT of the code box
% (listings sets them with \llap, so no overfull hbox can arise and the code
% x is untouched — measured: code frozen, digits at code_x-11.8pt); inside a
% colored box they are blanked like everyone else's (v3.37 hook), and the
% clean variant restores them in its frameless boxes (v3.78, above).
\lstdefinestyle{pystyle}{
  language=Python,basicstyle=\ttfamily\small,
  columns=fullflexible,keepspaces=true,
  frame=single,rulecolor=\color{cTblLine},
  numbers=left,numberstyle=\tiny\color{cGray},numbersep=8pt,
  xleftmargin=0.5em,framexleftmargin=0pt,framexrightmargin=0pt,
  showstringspaces=false,aboveskip=8pt,belowskip=8pt,
  keywordstyle=\color{cBlue}\bfseries,
  commentstyle=\color{cGray}\itshape,
  stringstyle=\color{cGreen},
}
% v4.5.0: hoisted out of the four styles so an enclosing box can override it.
% Every listing still resolves to true unless something later says otherwise,
% so shipped output is unchanged (verified byte-for-byte on all eight PDFs).
\lstset{resetmargins=true}
\lstdefinestyle{htmlstyle}{
  language=HTML,basicstyle=\ttfamily\small,
  columns=fullflexible,keepspaces=true,
  frame=single,rulecolor=\color{cTblLine},
  numbers=left,numberstyle=\tiny\color{cGray},numbersep=8pt,
  xleftmargin=0.5em,framexleftmargin=0pt,framexrightmargin=0pt,
  showstringspaces=false,aboveskip=8pt,belowskip=8pt,
  keywordstyle=\color{cBlue}\bfseries,
  commentstyle=\color{cGray}\itshape,
  stringstyle=\color{cGreen},
}
\newenvironment{asmblock}{\par\begin{otherlanguage}{english}}{\end{otherlanguage}\par}

% <hllpdf:listing-smart-quotes>
% v4.21.0: four Unicode punctuation characters must enter listings' token
% buffer. Unregistered characters bypass it: a closing quote was printed
% BEFORE the pending word, and an opening quote BEFORE a pending brace.
% Extend the standard character table, then activate only these codepoints
% in the listing's local character-table scope. Do this BEFORE the package's
% delimiter/alsoother/literate hooks so explicit user mappings still win.
% This is not universal Unicode support, nor a conversion to ASCII quotes.
% See listings developer guide 15.3.1 and references/listing-smart-quotes.md.
\makeatletter
\lst@CCPut \lst@ProcessOther {"2018}{"2019}{"201C}{"201D}\z@
\lst@AddToHookAtTop{SelectCharTable}{%
  \lst@ifactivechars
    \catcode"2018=\active \catcode"2019=\active
    \catcode"201C=\active \catcode"201D=\active
  \fi}
\makeatother
% </hllpdf:listing-smart-quotes>

% <hllpdf:listing-technical-symbols>
% v4.30: bounded extension for common technical punctuation in LTR listings.
% Without character-table registration these glyphs bypass the token buffer
% and can move before the preceding word even inside quoted strings.
% Keep the original four-quote block separate for its standing mutation test.
% This is source-faithful printing, not syntax validation or general Unicode.
\makeatletter
\lst@CCPut \lst@ProcessOther
  {"2013}{"2014}{"2026}{"2190}{"2192}{"2194}{"21D2}
  {"2212}{"221E}{"2260}{"2264}{"2265}\z@
\lst@AddToHookAtTop{SelectCharTable}{%
  \lst@ifactivechars
    \catcode"2013=\active \catcode"2014=\active \catcode"2026=\active
    \catcode"2190=\active \catcode"2192=\active \catcode"2194=\active
    \catcode"21D2=\active \catcode"2212=\active \catcode"221E=\active
    \catcode"2260=\active \catcode"2264=\active \catcode"2265=\active
  \fi}
\makeatother
% </hllpdf:listing-technical-symbols>


\usepackage[hidelinks,unicode]{hyperref}   % LAST, after colors. unicode = Hebrew bookmarks.
% v3.62: THE ISOLATION MACROS CORRUPTED EVERY PDF BOOKMARK THEY APPEARED IN.
% \en is defined (line ~155) as
%     \foreignlanguage{english}{\language=\csname l@nohyphenation\endcsname #1}
% and in a heading that body is written verbatim to the .toc and handed to
% hyperref's \pdfstringdef. hyperref strips the two CONTROL SEQUENCES (\language,
% \l@nohyphenation) but the plain `=` CHARACTER token between them is not a control
% sequence, so it survives into the PDF string: \chapter{שער \en{CNOT} והצמדה}
% produced the bookmark "שער =CNOT והצמדה". The printed TOC and the heading on the
% page are CORRECT — only the navigation pane and the viewer window title are wrong,
% which is why it went unnoticed. Measured on the shipped fleet: 70 corrupted
% bookmarks across 5 PDFs, from 177 unprotected \en-in-heading sites. Same corruption
% reached /Title and /Subject via the \hypersetup{pdftitle=...} the docs recommend.
% This hook tells hyperref to use the plain argument for all three isolation macros
% when building a PDF string. It does NOT touch the .toc (the printed TOC is
% byte-unchanged) and it does NOT mask the math signal — a heading with $...$ still
% logs "removing math shift", so BIDI-15 detection is intact.
% BONUS, measured: it also makes \code{} in an UNSTARRED heading COMPILE. BLD-10
% recorded that as needing \texorpdfstring or a starred heading; the real root cause
% is hyperref's \pdfstringdef scan of \code's \ifmmode...\else...\fi body — which is
% exactly why the identically-bodied \texttt survived (hyperref pre-neutralizes
% \texttt in PDF strings) and why \DeclareRobustCommand could never help.
% v3.64: \def\\{ } added. A manual line break in a heading WELDS the two words
% together in the bookmark: \chapter{שורה ראשונה\\שורה שנייה} produced the outline
% entry 'שורה ראשונהשורה שנייה' — "ראשונהשורה" is not a word. hyperref drops the
% \\ and puts nothing in its place, so the break that looks like a space on the
% page is a zero-width join in the navigation pane. The spaced form
% 'מילה אחת \\ מילה שנייה' escapes with a harmless DOUBLE space ('אחת  מילה'),
% which is why this never got noticed. Both are measured on a real outline.
% 0 live instances in the corpus — but SKILL.md and content-style.md both RECOMMEND
% a manual heading break, so this is a primed trap rather than a dormant one, and a
% one-token preamble fix retires it for good instead of asking authors to remember.
% Mapping \\ to a space fixes the WELD, which is the actual defect (a nonexistent
% word in the navigation pane). It does not normalise whitespace: measured after the
% fix, the welded form becomes 'שורה ראשונה שורה שנייה' (correct) and the already-
% spaced form goes from two spaces to three. That is cosmetic — viewers render it as
% a wider gap, and no word is corrupted. A per-heading \texorpdfstring is still the
% way to get an exactly-tuned bookmark string.
\pdfstringdefDisableCommands{\def\en#1{#1}\def\h#1{#1}\def\doclang#1{#1}\def\code#1{#1}\def\codeline#1{#1}\def\\{ }}
% v3.70: RESTORE \url LINE-BREAKING. url.sty freezes its break penalties at LOAD
% time (\mathchardef\UrlBreakPenalty=\binoppenalty), and v3.45's global
% \binoppenalty=10000 (inline-math atomicity) therefore priced EVERY breakpoint
% inside \Url at 10000 = never break. Measured: a long \url in Hebrew prose is
% one unbroken line, ~100pt overfull, invading the margin band the ink gates
% deliberately tolerate. With the two lines below: 0 overfull, the URL breaks
% after '/', the RTL split is CORRECT (fragment order forward), and each line
% gets its own link annotation exactly covering its fragment (measured).
% Confined to \Url's group — inline-math atomicity (TYP-7) is untouched.
\mathchardef\UrlBreakPenalty=200
\mathchardef\UrlBigBreakPenalty=100
\ifx\lakitmainlanguage\hllpdfenglish
  \hypersetup{pdflang=en}
\else
  \hypersetup{pdflang=he}
\fi  % /Lang is the document base locale (not per-span accessibility tagging).
% Set a per-document \hypersetup{pdftitle={...}} too — see SKILL.md.

% v3.52 (round-2 audit): two RTL bibliography/cross-ref defects, both invisible to
% every gate and both affecting the FIRST author to use the feature.
% (a) \@biblabel's default [#1] brackets are bidi-neutral and mirror to ]1[ under
% the RTL bibliography base — every \thebibliography reference list shows ]1[.
% LTR-isolating the label restores [1]. (The render gate's mirror rule fingerprints
% round parens )1(, so square brackets slipped through.)
\makeatletter
\renewcommand*{\@biblabel}[1]{\foreignlanguage{english}{[#1]}}
% v3.54: the SAME neutral-bracket trap on the INLINE side — a multi-key \cite{a,b}
% under the RTL base lays its European-number list right-to-left, so [1, 2] renders
% "[2,1]" and [1, 2, 3] renders "[3,2,1]" (verified by glyph geometry — the reader
% sees the sources in REVERSE order; a note form \cite[p.5]{a} also reorders to
% "[p. 5,1]"). LTR-isolate the whole bracket group (number list AND optional note),
% exactly like \@biblabel above: [1, 2] / [1, 2, 3] / [1, p. 5] all render correct.
% Placed AFTER hyperref so it wins; hyperref hyperlinks each number inside #1 (via
% \@citex), which this wrapper leaves intact. Mirrors the LaTeX kernel \@cite shape.
% v4.9.1: directional isolation must NOT switch the language here. Babel's
% safe-active citation wrapper passes its key list through \bbl@tempa;
% \foreignlanguage reuses that scratch macro before the deferred citation
% loop expands it. The old line therefore cited the literal key "english".
% If such a key existed, even the undefined-reference gate stayed green.
% \babelsublr isolates the bracket group without selecting a locale or
% changing that key list. Preserve native links, key order and optional notes.
\def\@cite#1#2{\babelsublr{[{#1\if@tempswa , #2\fi}]}}
\makeatother

% v4.11.0: a reference inside align's \text inherited HEBdisptag=true,
% so its enclosing hbox mirrored the brackets despite a clean compile.
% Isolate the complete printed reference without selecting a Babel language.
% Also neutralize the display-only normalfont direction reset locally. A
% direction-only patch inverted 12 to 21 inside equation in the measured probe.
% Plain maketag is local to the reference: mathtools' display-side wrapper
% otherwise globally consumes manual_tag while formatting a nested reference.
% Preserve each original command (including the selective auxiliary write),
% current tag forms, native links, and showonlyrefs on/off transitions.
\makeatletter
\NewCommandCopy{\HEB@nativeeqref}{\eqref}
\newcommand{\HEB@eqrefisolate}[1]{%
  \babelsublr{\HEBdisptagfalse\let\bbl@eqnodir\relax\let\maketag@@@\ORIGmaketag@@@#1}}
\protected\def\eqref#1{\HEB@eqrefisolate{\HEB@nativeeqref{#1}}}
\ExplSyntaxOn
\NewCommandCopy{\HEB@nativeMTeqref}{\MT_extended_eqref:n}
\cs_set_protected:Npn \MT_extended_eqref:n #1
  {\HEB@eqrefisolate{\HEB@nativeMTeqref{#1}}}
\ExplSyntaxOff
\makeatother

% v4.11.0: custom caption counters are not part of hyperref's English
% locale, and Hebrew has no upstream hyperref locale in this toolchain.
% Restore the appropriate names at every Babel language switch, including
% return from a nested inline island. Do not localize the destination keys.
\makeatletter
\newcommand{\HEB@refnames}{%
\def\chapterautorefname{פרק}%
\def\sectionautorefname{סעיף}%
\def\subsectionautorefname{תת־סעיף}%
\def\subsubsectionautorefname{תת־תת־סעיף}%
\def\paragraphautorefname{פסקה}%
% v4.31: hyperref uses a distinct counter name at the deepest standard level.
\def\subparagraphautorefname{תת־פסקה}%
\def\equationautorefname{משוואה}%
\def\figureautorefname{איור}%
\def\tableautorefname{טבלה}%
\def\partautorefname{חלק}%
\def\appendixautorefname{נספח}%
\def\footnoteautorefname{הערת שוליים}%
\def\itemautorefname{סעיף}%
\def\Itemautorefname{סעיף}%
\def\figautorefname{איור}%
\def\tblautorefname{טבלה}%
\def\pageautorefname{עמוד}% 
}
\addto\extrashebrew{\HEB@refnames}
\newcommand{\ENG@refnames}{%
  \def\subparagraphautorefname{subparagraph}%
  \def\chapterautorefname{chapter}%
  \def\sectionautorefname{section}%
  \def\subsectionautorefname{subsection}%
  \def\subsubsectionautorefname{subsubsection}%
  \def\paragraphautorefname{paragraph}%
  \def\equationautorefname{Equation}%
  \def\figureautorefname{Figure}%
  \def\tableautorefname{Table}%
  \def\partautorefname{Part}%
  \def\appendixautorefname{Appendix}%
  \def\footnoteautorefname{footnote}%
  \def\itemautorefname{item}%
  \def\figautorefname{Figure}%
  \def\tblautorefname{Table}%
  \def\Itemautorefname{item}%
  \def\pageautorefname{page}%
}
\addto\extrasenglish{\ENG@refnames}
\ifx\lakitmainlanguage\hllpdfenglish\ENG@refnames\else\HEB@refnames\fi
\makeatother

% ============================================================================
% Optional running head (v3.39): \lakithead{<book title>} — OFF by default.
% A document opts in with ONE line after \input{preamble.tex}:
%   \lakithead{מבנה המחשב}
% Every content page then carries the series' orientation line: chapter name
% at the reading edge, book title opposite, thin rule. (fancyhdr's [L]/[R]
% are PHYSICAL; under the RTL page direction [L] lands at the RIGHT/reading
% edge — measured, consistent across the corpus.) Packaged like the topic
% rail: always loaded, inert without the call. This centralizes the 8-line
% block that 8 lakits used to copy by hand — the same duplication vector
% that once let a lakit preamble go ~100 lines stale. Class-aware: report
% fills the mark from \chapter, article from \section. NB: a titlepage env
% self-empties; front matter that sets \pagestyle{empty} must restore with
% \pagestyle{fancy} afterwards (NOT plain — that discards the head).
% ============================================================================
\usepackage{fancyhdr}
% v3.61: \chapter* left a STALE running head — pages of an UNNUMBERED chapter
% carried the PREVIOUS chapter's title. report.cls routes \chapter* through
% \@schapter, which never calls \chaptermark, so \leftmark keeps whatever the last
% numbered chapter left. Measured: with \chapter{פרק ממוספר} then
% \chapter*{מבוא לא ממוספר}, page 4 (inside the starred chapter) printed
% "פרק ממוספר". Invisible on the chapter's FIRST page (\thispagestyle{plain}), so it
% only shows from the second page on — the pages an author is least likely to check.
% An unnumbered front-matter chapter (\chapter*{מבוא}, \chapter*{הקדמה}) is everyday
% Hebrew-academic structure, and the corpus already carries HAND-WRITTEN \markboth
% lines right after \chapter* — the exact per-document duplication \lakithead exists
% to retire. Patched INSIDE \lakithead so the macro stays inert when unused; a
% document's own later \markboth still wins, so those hand-patches do not regress.
\makeatletter
% <hllpdf:starred-heading-scope>
\newif\iflakin@starmarkpatched
\newcommand{\lakistarmark}{%
  \iflakin@starmarkpatched\else
  % Dispatch on the actual class counter, not hyperref's shared internals.
  \ifdefined\c@chapter
    \let\HEBorigschapter\@schapter
    % Record only after the native heading has reached its page. A mark
    % before the native routine can advertise the next heading too early.
    \def\@schapter##1{\HEBorigschapter{##1}\chaptermark{##1}}%
  \else
    % @ssect receives layout arguments, NOT the section name or level.
    % Capture the name in @startsection without introducing a TeX group:
    % grouping sectioning would break the native after-heading state.
    \def\HEB@sectionkind{}%
    \def\HEB@mainsection{section}%
    \let\HEBorigstartsection\@startsection
    \def\@startsection##1##2##3##4##5##6{%
      \def\HEB@sectionkind{##1}%
      \HEBorigstartsection{##1}{##2}{##3}{##4}{##5}{##6}}%
    \let\HEBorigssect\@ssect
    \def\@ssect##1##2##3##4##5{%
      \HEBorigssect{##1}{##2}{##3}{##4}{##5}%
      % Lower starred headings must not replace the article's section mark.
      \ifx\HEB@sectionkind\HEB@mainsection
        \sectionmark{##5}%
      \fi
      \ignorespaces}%
  \fi
  \global\lakin@starmarkpatchedtrue
  \fi}
% </hllpdf:starred-heading-scope>
% <hllpdf:running-language>
% Choose the wrapper when the mark is written, not when a later page ships.
% English punctuation must remain LTR even after returning to Hebrew prose.
\newcommand{\HEB@runningmark}[1]{%
  \iflanguage{english}{\markboth{\en{#1}}{}}{\markboth{\h{#1}}{}}}
% </hllpdf:running-language>
\makeatother
% v4.2.1 (NAV): every collection already names itself in \lakithead — feed the
% PDF /Title from the same declaration, so the navigation gate's
% PDF_REQUIRE_TITLE has a universal source and no fleet member ships
% title-less by missing a second declaration site. Measured before this: 14
% of 16 collections shipped /Title-empty while every one carried a perfectly
% good \lakithead string (plain Hebrew in all 16 — no math, no formatting, so
% the PDF-string path is clean). An explicit title always wins: set it with
% \lakitpdftitle{...} anywhere before \lakithead, or with a plain
% \hypersetup{pdftitle=...} — the fallback detects that via hyperref's
% \@pdftitle, probe-verified \@empty-equal until the first explicit set.
\makeatletter
\newif\iflakit@pdftitleset
\newcommand{\lakitpdftitle}[1]{%
  \hypersetup{pdftitle={#1}}\global\lakit@pdftitlesettrue}
\newcommand{\lakithead}[1]{%
  \iflakit@pdftitleset\else
    \ifx\@pdftitle\@empty
      \hypersetup{pdftitle={#1}}\global\lakit@pdftitlesettrue
    \fi
  \fi
  \pagestyle{fancy}\fancyhf{}%
  \lakistarmark

  \fancyhead[L]{\small\textcolor{cGray}{\leftmark}}%
  % <hllpdf:fixed-head-language>
  % The fixed document title has its OWN language snapshot, independent of
  % the current section/chapter mark. Choose the branch NOW, not at shipout:
  % the page may be Hebrew by the time an English title is finally printed.
  \iflanguage{english}{%
    \fancyhead[R]{\small\textcolor{cGray}{\en{#1}}}%
  }{%
    \fancyhead[R]{\small\textcolor{cGray}{\h{#1}}}%
  }%
  % </hllpdf:fixed-head-language>
  \fancyfoot[C]{\thepage}%
  \renewcommand{\headrulewidth}{0.4pt}%
  \ifdefined\chaptermark
    \renewcommand{\chaptermark}[1]{\HEB@runningmark{##1}}%
  \else
    \renewcommand{\sectionmark}[1]{\HEB@runningmark{##1}}%
  \fi}
\makeatother

% v3.39.1: path-independent byte-reproducibility. LuaTeX derives the trailer
% /ID from the OUTPUT PATH, so two clones of the same source at different
% directories produced different sha256 (measured: the ONLY delta was the
% 62-byte /ID line). Bit 512 suppresses exactly that line — measured: Creator/
% Producer/CreationDate/ModDate and all content untouched. /ID is optional
% per ISO 32000 (required only for ENCRYPTED PDFs and PDF/A — neither is a
% target here). If a future workflow needs an ID (PDF/A, Acrobat revision
% matching), delete this line and it returns.
\pdfvariable suppressoptionalinfo \numexpr 512\relax

% v3.75: BOLD MATH IN BOLD HEADINGS. Heading text is \bfseries but math inside
% it stayed regular-weight CM — a hairline $V_T$ inside a heavy 24.8pt chapter
% title (measured: CMMI12 fl=6 next to FrankRuhlLibre-Bold fl=20; 155 live
% heading-with-math sites fleet-wide). \boldmath rides along wherever the class
% sets \bfseries for a heading; A/B measured: $V_T$ -> CMMIB10 fl=22, 0 missing
% characters (the bold CM faces ship with this TeX Live). Guarded per-macro so
% the article class (no \@makechapterhead) and future class changes stay safe.
% v4.0.1 (adoption fix): the SEARCH patterns above carry no parameter token.
% v4.0.0 made them specific by including the class's own #1/#2, and measured on
% TeX Live 2023 that never matches: a \patchcmd search pattern containing a
% parameter token cannot match while the call sits inside \@ifundefined{...}{}{...}
% (probed both #1 and ##1 -> FAIL; the same pattern at top level succeeds). Every
% one of these patches therefore failed, and because v4.0.0 also replaced the
% empty failure branch with \PackageError, the failure became FATAL: all four
% bundled example assets stopped building ("Cannot patch chapter-title math
% weight -> Fatal error, no output PDF"). The \PackageError branches are KEPT —
% a silent no-op here is what let v3.75's fix rot unnoticed — only the parameter
% tokens are dropped, so the patches match and a real class change still errors.
\makeatletter
\@ifundefined{@makechapterhead}{}{%
  \patchcmd{\@makechapterhead}{\huge\bfseries \@chapapp}
    {\huge\bfseries\boldmath \@chapapp}{}
    {\PackageError{hebrew-lualatex-pdf}{Cannot patch chapter-number math weight}{Unsupported document class}}%
  \patchcmd{\@makechapterhead}{\Huge \bfseries}
    {\Huge \bfseries\boldmath}{}
    {\PackageError{hebrew-lualatex-pdf}{Cannot patch chapter-title math weight}{Unsupported document class}}%
  \patchcmd{\@makeschapterhead}{\Huge \bfseries}
    {\Huge \bfseries\boldmath}{}
    {\PackageError{hebrew-lualatex-pdf}{Cannot patch starred chapter math weight}{Unsupported document class}}}
% v4.0.1: \part sits under a THREE-DEEP wrapper chain, so the class-based anchor
% is not in \@part at all. Measured with \show: hyperref rewrites \@part to
% `...\phantomsection\fi \H@old@part`, and nameref rewrites THAT to
% `\NR@gettitle{#1}\NR@part[{#1}]{#2}` — the class body ends up in \NR@part.
% v4.0.0 patched \@part and made failure fatal, which is why all seven bundled
% examples stopped building. \part is NOT a dead heading level here: 188 live uses
% across the fleet, so the bold-math fix is worth getting right rather than
% dropping. We walk the chain innermost-first and stop at the first hit. Failure is
% a WARNING, not an error: the chain is a property of hyperref/nameref versions, so
% a future package release must degrade to regular-weight part math — a cosmetic
% loss — instead of refusing to build. Chapter and section keep \PackageError,
% because those targets are in the CLASS and were verified to patch here.
\newif\iflakin@partnum \newif\iflakin@parttitle \newif\iflakin@sparttitle
\def\lakin@partnum#1#2{\iflakin@partnum\else
  \patchcmd{#1}{#2\bfseries \partname}{#2\bfseries\boldmath \partname}
    {\lakin@partnumtrue}{}\fi}
\def\lakin@parttitle#1#2{\iflakin@parttitle\else
  \patchcmd{#1}{#2 \bfseries}{#2 \bfseries\boldmath}{\lakin@parttitletrue}{}\fi}
\def\lakin@sparttitle#1#2{\iflakin@sparttitle\else
  \patchcmd{#1}{#2 \bfseries}{#2 \bfseries\boldmath}{\lakin@sparttitletrue}{}\fi}
\@ifundefined{c@chapter}{\def\lakin@pnum{\Large}\def\lakin@ptit{\huge}}%
                        {\def\lakin@pnum{\huge}\def\lakin@ptit{\Huge}}
\@ifundefined{NR@part}{}{\expandafter\lakin@partnum\expandafter\NR@part\expandafter{\lakin@pnum}%
                         \expandafter\lakin@parttitle\expandafter\NR@part\expandafter{\lakin@ptit}}
\@ifundefined{H@old@part}{}{\expandafter\lakin@partnum\expandafter\H@old@part\expandafter{\lakin@pnum}%
                            \expandafter\lakin@parttitle\expandafter\H@old@part\expandafter{\lakin@ptit}}
\@ifundefined{@part}{}{\expandafter\lakin@partnum\expandafter\@part\expandafter{\lakin@pnum}%
                       \expandafter\lakin@parttitle\expandafter\@part\expandafter{\lakin@ptit}}
\@ifundefined{NR@spart}{}{\expandafter\lakin@sparttitle\expandafter\NR@spart\expandafter{\lakin@ptit}}
\@ifundefined{H@old@spart}{}{\expandafter\lakin@sparttitle\expandafter\H@old@spart\expandafter{\lakin@ptit}}
\@ifundefined{@spart}{}{\expandafter\lakin@sparttitle\expandafter\@spart\expandafter{\lakin@ptit}}
\iflakin@partnum\else\PackageWarning{hebrew-lualatex-pdf}{Part-number math stays
  regular weight: no known \noexpand\part wrapper matched (hyperref/nameref
  version change). Cosmetic only}\fi
\iflakin@parttitle\else\PackageWarning{hebrew-lualatex-pdf}{Part-title math stays
  regular weight: no known \noexpand\part wrapper matched. Cosmetic only}\fi
% v4.12.0: THE PART NUMBER IN THE TOC WAS NOT DIRECTIONALLY SEPARATED FROM THE
% TITLE. report.cls/article.cls compose the part entry as ONE text stream --
% \addcontentsline{toc}{part}{\thepart\hspace{1em}#1} -- while EVERY other level
% goes through \numberline, which puts the number in an \hbox. A box is a bidi
% OBJECT and cannot join a neighbouring run; a bare "IV" can. So in the RTL TOC
% paragraph a Latin number and a Latin-INITIAL title merge into a SINGLE LTR run
% and are laid out left-to-right as a unit: the number lands on the wrong side of
% the row, and a mixed title is scrambled outright. Measured on \part{MCP ...}:
% the printed order was "<tail> V MCP" -- title fragment, number, title fragment.
% The chapter/section rows one line below, with the SAME Latin-initial titles,
% were correct on the same page. That contrast is the whole diagnosis.
% \en{} IS NOT A WORKAROUND, measured both ways on one fixture: it fixes a
% PURE-Latin title (\part{\en{Hooks}} renders correctly) and does NOT fix a MIXED
% one (\part{\en{MCP} <hebrew>} stays scrambled). babel's language switch changes
% font and hyphenation, not the bidi run structure -- so the defect is not
% reachable from the source side at all, which is why it belongs here.
% Boxing the number is the entire fix and is a STRICT NO-OP where the bug never
% bit: a Hebrew-titled document compiles BYTE-IDENTICAL before and after (md5,
% 3 clean runs), because there the neutral glue already sat between an L run and
% an R run. The PDF outline is byte-identical too -- hyperref's \pdfstringdef
% neutralises \mbox, so bookmark strings keep their part numbers.
% \numberline{\thepart} would also fix the direction, but it re-geometries every
% part row in the fleet and DROPS the number from the part bookmark (hyperref
% strips \numberline). \mbox changes nothing that was already right.
% The search pattern deliberately carries no # token: etoolbox cannot match a
% parameter token in a patch pattern (measured -- all three chain members
% reported "not patchable" while the pattern included #1). "\thepart \hspace
% {1em}" is brace-balanced and occurs exactly once in the wrapper body.
\newif\iflakin@parttocnum
\def\lakin@parttocnum#1{\iflakin@parttocnum\else
  \patchcmd{#1}{\thepart \hspace {1em}}{\mbox {\thepart }\hspace {1em}}
    {\lakin@parttocnumtrue}{}\fi}
\@ifundefined{NR@part}{}{\lakin@parttocnum\NR@part}
\@ifundefined{H@old@part}{}{\lakin@parttocnum\H@old@part}
\@ifundefined{@part}{}{\lakin@parttocnum\@part}
\iflakin@parttocnum\else\PackageWarning{hebrew-lualatex-pdf}{Part TOC number is
  not boxed: no known \noexpand\part wrapper matched. A part title that BEGINS
  with Latin will print its number on the wrong side of its TOC row}\fi
\patchcmd{\section}{\bfseries}{\bfseries\boldmath}{}
  {\PackageError{hebrew-lualatex-pdf}{Cannot patch section math weight}{Unsupported document class}}
\patchcmd{\subsection}{\bfseries}{\bfseries\boldmath}{}
  {\PackageError{hebrew-lualatex-pdf}{Cannot patch subsection math weight}{Unsupported document class}}
\patchcmd{\subsubsection}{\bfseries}{\bfseries\boldmath}{}
  {\PackageError{hebrew-lualatex-pdf}{Cannot patch subsubsection math weight}{Unsupported document class}}
\patchcmd{\paragraph}{\bfseries}{\bfseries\boldmath}{}{}
\patchcmd{\subparagraph}{\bfseries}{\bfseries\boldmath}{}{}
\@ifundefined{l@chapter}{}{%
  \patchcmd{\l@chapter}{\bfseries}{\bfseries\boldmath}{}{}}
\makeatother
\setcounter{tocdepth}{\lakittocdepth}
% <hllpdf:measured-toc>
% v4.16.0: Measure actual number labels, not a guessed digit count. Read only
% canonical contentsline records; never execute titles or addtocontents code
% during the measuring pass. Widths are local to this individual list and are
% available before its first row: no per-row drift and no extra aux pass.
\makeatletter
\newread\hllpdf@tocstream
\newdimen\hllpdf@toc@chapter
\newdimen\hllpdf@toc@section
\newdimen\hllpdf@toc@subsection
\newdimen\hllpdf@toc@subsubsection
\newdimen\hllpdf@toc@paragraph
\newdimen\hllpdf@toc@subparagraph
\newdimen\hllpdf@toc@figure
\newdimen\hllpdf@toc@table
\newdimen\hllpdf@tocpage
\newdimen\hllpdf@toc@sectionindent
\newdimen\hllpdf@toc@subsectionindent
\newdimen\hllpdf@toc@subsubsectionindent
\newdimen\hllpdf@toc@paragraphindent
\newdimen\hllpdf@toc@subparagraphindent
\def\hllpdf@tocdefaults{%
  \hllpdf@toc@chapter=2em
  \hllpdf@toc@section=3.2em
  \hllpdf@toc@subsection=3.9em
  \@ifundefined{c@chapter}{%
    \hllpdf@toc@subsubsection=3.2em
    \hllpdf@toc@paragraph=4.1em
    \hllpdf@toc@subparagraph=5em
  }{%
    \hllpdf@toc@subsubsection=4.1em
    \hllpdf@toc@paragraph=5em
    \hllpdf@toc@subparagraph=6em
  }%
  \hllpdf@toc@figure=2.3em
  \hllpdf@toc@table=2.3em
  \hllpdf@tocpage=\@pnumwidth\relax
}
\def\hllpdf@tocindents{%
  % Preserve the established shallow layout when its slots already fit.
  \hllpdf@toc@sectionindent=1.5em
  \@tempdima=\hllpdf@toc@chapter \advance\@tempdima by -2em
  \advance\hllpdf@toc@sectionindent by \@tempdima
  \hllpdf@toc@subsectionindent=\dimexpr
    \hllpdf@toc@sectionindent+\hllpdf@toc@section\relax
  \hllpdf@toc@subsubsectionindent=\dimexpr
    \hllpdf@toc@subsectionindent+\hllpdf@toc@subsection\relax
  \hllpdf@toc@paragraphindent=\dimexpr
    \hllpdf@toc@subsubsectionindent+\hllpdf@toc@subsubsection\relax
  \hllpdf@toc@subparagraphindent=\dimexpr
    \hllpdf@toc@paragraphindent+\hllpdf@toc@paragraph\relax
}
\def\hllpdf@tocfont@chapter{\bfseries\boldmath}
\def\hllpdf@tocfont@part{\large\bfseries\boldmath}
\def\hllpdf@tocentryfont{%
  \ifcsname hllpdf@tocfont@\hllpdf@toclevel\endcsname
    \csname hllpdf@tocfont@\hllpdf@toclevel\endcsname
  \fi}
\long\def\hllpdf@tocdiscard#1\hllpdf@tocstop{}
\def\hllpdf@tocscanrecord{%
  \@ifnextchar\contentsline{\hllpdf@tocscanline}{%
    \@ifnextchar\babel@toc{\hllpdf@tocscanlanguage}{\hllpdf@tocdiscard}}}
\long\def\hllpdf@tocscanlanguage\babel@toc#1#2#3\hllpdf@tocstop{%
  % Babel records are declarations, not arbitrary code. Track the language
  % without executing the original directive or changing the surrounding list.
  \def\hllpdf@toclanguage{#1}}
\long\def\hllpdf@tocscanline\contentsline#1#2#3#4#5\hllpdf@tocstop{%
  \def\hllpdf@toclevel{#1}%
  \ifcsname hllpdf@toc@#1\endcsname
    \hllpdf@tocscantitle#2\hllpdf@tocstop
  \fi
  % The page-number column needs its own measurement (including bold chapter
  % rows). Its extra clearance separates the page label from dot leaders.
  \setbox\@tempboxa=\hbox{\foreignlanguage{\hllpdf@toclanguage}{\hllpdf@tocentryfont #3\kern.3em}}%
  \ifdim\wd\@tempboxa>\hllpdf@tocpage
    \hllpdf@tocpage=\wd\@tempboxa
  \fi}
\def\hllpdf@tocscantitle{%
  \@ifnextchar\numberline{\hllpdf@tocscanlabel}{\hllpdf@tocdiscard}}
\long\def\hllpdf@tocscanlabel\numberline#1#2\hllpdf@tocstop{%
  \setbox\@tempboxa=\hbox{\foreignlanguage{\hllpdf@toclanguage}{\hllpdf@tocentryfont #1\kern.5em}}%
  \ifdim\wd\@tempboxa>\csname hllpdf@toc@\hllpdf@toclevel\endcsname
    \csname hllpdf@toc@\hllpdf@toclevel\endcsname=\wd\@tempboxa
  \fi}
\def\hllpdf@tocread{%
  \ifeof\hllpdf@tocstream\else
    \read\hllpdf@tocstream to\hllpdf@tocrecord
    \expandafter\hllpdf@tocscanrecord\hllpdf@tocrecord\hllpdf@tocstop
    \expandafter\hllpdf@tocread
  \fi}
\def\hllpdf@tocmeasure#1{%
  \makeatletter
  \edef\hllpdf@toclanguage{\languagename}%
  \hllpdf@tocdefaults
  \openin\hllpdf@tocstream=\jobname.#1\relax
  \hllpdf@tocread
  \closein\hllpdf@tocstream
  \hllpdf@tocindents
  \edef\@pnumwidth{\the\hllpdf@tocpage}%
  \@tempdima=\dimexpr\hllpdf@tocpage+1em\relax
  \ifdim\@tempdima>\@tocrmarg\relax
    \edef\@tocrmarg{\the\@tempdima}%
  \fi}
\let\hllpdf@originalstarttoc\@starttoc
\def\@starttoc#1{%
  \begingroup
    \def\hllpdf@tocext{#1}%
    \def\hllpdf@tocwanted{toc}%
    \ifx\hllpdf@tocext\hllpdf@tocwanted \hllpdf@tocmeasure{#1}\else
      \def\hllpdf@tocwanted{lof}%
      \ifx\hllpdf@tocext\hllpdf@tocwanted \hllpdf@tocmeasure{#1}\else
        \def\hllpdf@tocwanted{lot}%
        \ifx\hllpdf@tocext\hllpdf@tocwanted \hllpdf@tocmeasure{#1}\fi
      \fi
    \fi
    \hllpdf@originalstarttoc{#1}%
  \endgroup}
\hllpdf@tocdefaults
\hllpdf@tocindents
\@ifundefined{l@chapter}{}{%
  \patchcmd{\l@chapter}{\setlength\@tempdima{1.5em}}%
    {\setlength\@tempdima{\hllpdf@toc@chapter}}{}%
    {\PackageError{hebrew-lualatex-pdf}{Cannot size chapter TOC numbers}%
       {Unsupported contents layout: the chapter number slot was not found}}}
\renewcommand*\l@section{\@dottedtocline{1}%
  {\hllpdf@toc@sectionindent}{\hllpdf@toc@section}}
\renewcommand*\l@subsection{\@dottedtocline{2}%
  {\hllpdf@toc@subsectionindent}{\hllpdf@toc@subsection}}
\renewcommand*\l@subsubsection{\@dottedtocline{3}%
  {\hllpdf@toc@subsubsectionindent}{\hllpdf@toc@subsubsection}}
\renewcommand*\l@paragraph{\@dottedtocline{4}%
  {\hllpdf@toc@paragraphindent}{\hllpdf@toc@paragraph}}
\renewcommand*\l@subparagraph{\@dottedtocline{5}%
  {\hllpdf@toc@subparagraphindent}{\hllpdf@toc@subparagraph}}
\renewcommand*\l@figure{\@dottedtocline{1}{1.5em}{\hllpdf@toc@figure}}
\renewcommand*\l@table{\@dottedtocline{1}{1.5em}{\hllpdf@toc@table}}
\makeatother
% </hllpdf:measured-toc>
\setcounter{secnumdepth}{\lakitsecnumdepth}
% Public preamble-only caption-numbering policy. The historical within-major
% default is untouched unless this command is called. "continuous" removes the
% reset edge as well as the printed prefix, so Figure 1/2 and Table 1/2 stay
% unique across sections/chapters. Returning to "within-major" restores the
% original class-aware displays and native destination conventions. The policy
% never replaces \figcap/\tabcap or their measured link-placement mechanism.
\makeatletter
\let\lakin@majorfig\thefig
\let\lakin@majortbl\thetbl
\let\lakin@majorHfig\theHfig
\let\lakin@majorHtbl\theHtbl
\newcommand{\lakitcaptionnumbering}[1]{%
  \edef\lakin@numbering{#1}%
  \def\lakin@continuous{continuous}\def\lakin@withinmajor{within-major}%
  \ifx\lakin@numbering\lakin@continuous
    \@ifundefined{c@chapter}{%
      \@removefromreset{fig}{section}\@removefromreset{tbl}{section}%
    }{\@removefromreset{fig}{chapter}\@removefromreset{tbl}{chapter}}%
    \renewcommand{\thefig}{\arabic{fig}}\renewcommand{\thetbl}{\arabic{tbl}}%
    \def\theHfig{\arabic{fig}}\def\theHtbl{\arabic{tbl}}%
  \else\ifx\lakin@numbering\lakin@withinmajor
    \@ifundefined{c@chapter}{%
      \@removefromreset{fig}{section}\@removefromreset{tbl}{section}%
      \@addtoreset{fig}{section}\@addtoreset{tbl}{section}%
    }{%
      \@removefromreset{fig}{chapter}\@removefromreset{tbl}{chapter}%
      \@addtoreset{fig}{chapter}\@addtoreset{tbl}{chapter}}%
    \let\thefig\lakin@majorfig\let\thetbl\lakin@majortbl
    \let\theHfig\lakin@majorHfig\let\theHtbl\lakin@majorHtbl
  \else
    \PackageError{hebrew-english-latex-studio}{Unknown caption-numbering policy '#1'}%
      {Use continuous or within-major before begin{document}.}%
  \fi\fi}
\@onlypreamble\lakitcaptionnumbering
\makeatother
% Architecture readiness sentinel. Optional runtime layers require the complete
% core, so the marker is defined only after every mechanism/baseline definition.
% <hllpdf:architecture-ready>
\def\hllpdfcoreloaded{1}
% </hllpdf:architecture-ready>

% v4.22: registered file resources, not an assertion about embedded subset bytes.
% The log framing avoids long-path wrapping and preserves UTF-8 as hex.
% <hllpdf:font-inventory>
\directlua{
  local function inventory()
    local names = {}
    local unnamed = 0
    for id, item in font.each() do
      local path = item.filename
      if path and not (path == "") then names[path] = true else unnamed = unnamed + 1 end
    end
    local sorted = {}
    for path in pairs(names) do table.insert(sorted,path) end
    table.sort(sorted)
    texio.write_nl("log","HLLFONT|BEGIN|1")
    for index,path in ipairs(sorted) do
      local hex = path:gsub(".",function(c) return string.format(string.char(37) .. "02x",string.byte(c)) end)
      local count = math.ceil(string.len(hex)/48)
      texio.write_nl("log","HLLFONT|FILE|" .. index .. "|" .. count)
      for offset=1,string.len(hex),48 do texio.write_nl("log","HLLFONT|HEX|" .. hex:sub(offset,offset+47)) end
    end
    texio.write_nl("log","HLLFONT|END|" .. rawlen(sorted) .. "|" .. unnamed)
    texio.write_nl("log","")
  end
  luatexbase.add_to_callback("finish_pdffile",inventory,"hllpdf.runtime.fonts")
}
% </hllpdf:font-inventory>
:
%   \lakithead{מבנה המחשב}
% Every content page then carries the series' orientation line: chapter name
% at the reading edge, book title opposite, thin rule. (fancyhdr's [L]/[R]
% are PHYSICAL; under the RTL page direction [L] lands at the RIGHT/reading
% edge — measured, consistent across the corpus.) Packaged like the topic
% rail: always loaded, inert without the call. This centralizes the 8-line
% block that 8 lakits used to copy by hand — the same duplication vector
% that once let a lakit preamble go ~100 lines stale. Class-aware: report
% fills the mark from \chapter, article from \section. NB: a titlepage env
% self-empties; front matter that sets \pagestyle{empty} must restore with
% \pagestyle{fancy} afterwards (NOT plain — that discards the head).
% ============================================================================
\usepackage{fancyhdr}
% v3.61: \chapter* left a STALE running head — pages of an UNNUMBERED chapter
% carried the PREVIOUS chapter's title. report.cls routes \chapter* through
% \@schapter, which never calls \chaptermark, so \leftmark keeps whatever the last
% numbered chapter left. Measured: with \chapter{פרק ממוספר} then
% \chapter*{מבוא לא ממוספר}, page 4 (inside the starred chapter) printed
% "פרק ממוספר". Invisible on the chapter's FIRST page (\thispagestyle{plain}), so it
% only shows from the second page on — the pages an author is least likely to check.
% An unnumbered front-matter chapter (\chapter*{מבוא}, \chapter*{הקדמה}) is everyday
% Hebrew-academic structure, and the corpus already carries HAND-WRITTEN \markboth
% lines right after \chapter* — the exact per-document duplication \lakithead exists
% to retire. Patched INSIDE \lakithead so the macro stays inert when unused; a
% document's own later \markboth still wins, so those hand-patches do not regress.
\makeatletter
% <hllpdf:starred-heading-scope>
\newif\iflakin@starmarkpatched
\newcommand{\lakistarmark}{%
  \iflakin@starmarkpatched\else
  % Dispatch on the actual class counter, not hyperref's shared internals.
  \ifdefined\c@chapter
    \let\HEBorigschapter\@schapter
    % Record only after the native heading has reached its page. A mark
    % before the native routine can advertise the next heading too early.
    \def\@schapter##1{\HEBorigschapter{##1}\chaptermark{##1}}%
  \else
    % @ssect receives layout arguments, NOT the section name or level.
    % Capture the name in @startsection without introducing a TeX group:
    % grouping sectioning would break the native after-heading state.
    \def\HEB@sectionkind{}%
    \def\HEB@mainsection{section}%
    \let\HEBorigstartsection\@startsection
    \def\@startsection##1##2##3##4##5##6{%
      \def\HEB@sectionkind{##1}%
      \HEBorigstartsection{##1}{##2}{##3}{##4}{##5}{##6}}%
    \let\HEBorigssect\@ssect
    \def\@ssect##1##2##3##4##5{%
      \HEBorigssect{##1}{##2}{##3}{##4}{##5}%
      % Lower starred headings must not replace the article's section mark.
      \ifx\HEB@sectionkind\HEB@mainsection
        \sectionmark{##5}%
      \fi
      \ignorespaces}%
  \fi
  \global\lakin@starmarkpatchedtrue
  \fi}
% </hllpdf:starred-heading-scope>
% <hllpdf:running-language>
% Choose the wrapper when the mark is written, not when a later page ships.
% English punctuation must remain LTR even after returning to Hebrew prose.
\newcommand{\HEB@runningmark}[1]{%
  \iflanguage{english}{\markboth{\en{#1}}{}}{\markboth{\h{#1}}{}}}
% </hllpdf:running-language>
\makeatother
% v4.2.1 (NAV): every collection already names itself in \lakithead — feed the
% PDF /Title from the same declaration, so the navigation gate's
% PDF_REQUIRE_TITLE has a universal source and no fleet member ships
% title-less by missing a second declaration site. Measured before this: 14
% of 16 collections shipped /Title-empty while every one carried a perfectly
% good \lakithead string (plain Hebrew in all 16 — no math, no formatting, so
% the PDF-string path is clean). An explicit title always wins: set it with
% \lakitpdftitle{...} anywhere before \lakithead, or with a plain
% \hypersetup{pdftitle=...} — the fallback detects that via hyperref's
% \@pdftitle, probe-verified \@empty-equal until the first explicit set.
\makeatletter
\newif\iflakit@pdftitleset
\newcommand{\lakitpdftitle}[1]{%
  \hypersetup{pdftitle={#1}}\global\lakit@pdftitlesettrue}
\newcommand{\lakithead}[1]{%
  \iflakit@pdftitleset\else
    \ifx\@pdftitle\@empty
      \hypersetup{pdftitle={#1}}\global\lakit@pdftitlesettrue
    \fi
  \fi
  \pagestyle{fancy}\fancyhf{}%
  \lakistarmark

  \fancyhead[L]{\small\textcolor{cGray}{\leftmark}}%
  % <hllpdf:fixed-head-language>
  % The fixed document title has its OWN language snapshot, independent of
  % the current section/chapter mark. Choose the branch NOW, not at shipout:
  % the page may be Hebrew by the time an English title is finally printed.
  \iflanguage{english}{%
    \fancyhead[R]{\small\textcolor{cGray}{\en{#1}}}%
  }{%
    \fancyhead[R]{\small\textcolor{cGray}{\h{#1}}}%
  }%
  % </hllpdf:fixed-head-language>
  \fancyfoot[C]{\thepage}%
  \renewcommand{\headrulewidth}{0.4pt}%
  \ifdefined\chaptermark
    \renewcommand{\chaptermark}[1]{\HEB@runningmark{##1}}%
  \else
    \renewcommand{\sectionmark}[1]{\HEB@runningmark{##1}}%
  \fi}
\makeatother

% v3.39.1: path-independent byte-reproducibility. LuaTeX derives the trailer
% /ID from the OUTPUT PATH, so two clones of the same source at different
% directories produced different sha256 (measured: the ONLY delta was the
% 62-byte /ID line). Bit 512 suppresses exactly that line — measured: Creator/
% Producer/CreationDate/ModDate and all content untouched. /ID is optional
% per ISO 32000 (required only for ENCRYPTED PDFs and PDF/A — neither is a
% target here). If a future workflow needs an ID (PDF/A, Acrobat revision
% matching), delete this line and it returns.
\pdfvariable suppressoptionalinfo \numexpr 512\relax

% v3.75: BOLD MATH IN BOLD HEADINGS. Heading text is \bfseries but math inside
% it stayed regular-weight CM — a hairline $V_T$ inside a heavy 24.8pt chapter
% title (measured: CMMI12 fl=6 next to FrankRuhlLibre-Bold fl=20; 155 live
% heading-with-math sites fleet-wide). \boldmath rides along wherever the class
% sets \bfseries for a heading; A/B measured: $V_T$ -> CMMIB10 fl=22, 0 missing
% characters (the bold CM faces ship with this TeX Live). Guarded per-macro so
% the article class (no \@makechapterhead) and future class changes stay safe.
% v4.0.1 (adoption fix): the SEARCH patterns above carry no parameter token.
% v4.0.0 made them specific by including the class's own #1/#2, and measured on
% TeX Live 2023 that never matches: a \patchcmd search pattern containing a
% parameter token cannot match while the call sits inside \@ifundefined{...}{}{...}
% (probed both #1 and ##1 -> FAIL; the same pattern at top level succeeds). Every
% one of these patches therefore failed, and because v4.0.0 also replaced the
% empty failure branch with \PackageError, the failure became FATAL: all four
% bundled example assets stopped building ("Cannot patch chapter-title math
% weight -> Fatal error, no output PDF"). The \PackageError branches are KEPT —
% a silent no-op here is what let v3.75's fix rot unnoticed — only the parameter
% tokens are dropped, so the patches match and a real class change still errors.
\makeatletter
\@ifundefined{@makechapterhead}{}{%
  \patchcmd{\@makechapterhead}{\huge\bfseries \@chapapp}
    {\huge\bfseries\boldmath \@chapapp}{}
    {\PackageError{hebrew-lualatex-pdf}{Cannot patch chapter-number math weight}{Unsupported document class}}%
  \patchcmd{\@makechapterhead}{\Huge \bfseries}
    {\Huge \bfseries\boldmath}{}
    {\PackageError{hebrew-lualatex-pdf}{Cannot patch chapter-title math weight}{Unsupported document class}}%
  \patchcmd{\@makeschapterhead}{\Huge \bfseries}
    {\Huge \bfseries\boldmath}{}
    {\PackageError{hebrew-lualatex-pdf}{Cannot patch starred chapter math weight}{Unsupported document class}}}
% v4.0.1: \part sits under a THREE-DEEP wrapper chain, so the class-based anchor
% is not in \@part at all. Measured with \show: hyperref rewrites \@part to
% `...\phantomsection\fi \H@old@part`, and nameref rewrites THAT to
% `\NR@gettitle{#1}\NR@part[{#1}]{#2}` — the class body ends up in \NR@part.
% v4.0.0 patched \@part and made failure fatal, which is why all seven bundled
% examples stopped building. \part is NOT a dead heading level here: 188 live uses
% across the fleet, so the bold-math fix is worth getting right rather than
% dropping. We walk the chain innermost-first and stop at the first hit. Failure is
% a WARNING, not an error: the chain is a property of hyperref/nameref versions, so
% a future package release must degrade to regular-weight part math — a cosmetic
% loss — instead of refusing to build. Chapter and section keep \PackageError,
% because those targets are in the CLASS and were verified to patch here.
\newif\iflakin@partnum \newif\iflakin@parttitle \newif\iflakin@sparttitle
\def\lakin@partnum#1#2{\iflakin@partnum\else
  \patchcmd{#1}{#2\bfseries \partname}{#2\bfseries\boldmath \partname}
    {\lakin@partnumtrue}{}\fi}
\def\lakin@parttitle#1#2{\iflakin@parttitle\else
  \patchcmd{#1}{#2 \bfseries}{#2 \bfseries\boldmath}{\lakin@parttitletrue}{}\fi}
\def\lakin@sparttitle#1#2{\iflakin@sparttitle\else
  \patchcmd{#1}{#2 \bfseries}{#2 \bfseries\boldmath}{\lakin@sparttitletrue}{}\fi}
\@ifundefined{c@chapter}{\def\lakin@pnum{\Large}\def\lakin@ptit{\huge}}%
                        {\def\lakin@pnum{\huge}\def\lakin@ptit{\Huge}}
\@ifundefined{NR@part}{}{\expandafter\lakin@partnum\expandafter\NR@part\expandafter{\lakin@pnum}%
                         \expandafter\lakin@parttitle\expandafter\NR@part\expandafter{\lakin@ptit}}
\@ifundefined{H@old@part}{}{\expandafter\lakin@partnum\expandafter\H@old@part\expandafter{\lakin@pnum}%
                            \expandafter\lakin@parttitle\expandafter\H@old@part\expandafter{\lakin@ptit}}
\@ifundefined{@part}{}{\expandafter\lakin@partnum\expandafter\@part\expandafter{\lakin@pnum}%
                       \expandafter\lakin@parttitle\expandafter\@part\expandafter{\lakin@ptit}}
\@ifundefined{NR@spart}{}{\expandafter\lakin@sparttitle\expandafter\NR@spart\expandafter{\lakin@ptit}}
\@ifundefined{H@old@spart}{}{\expandafter\lakin@sparttitle\expandafter\H@old@spart\expandafter{\lakin@ptit}}
\@ifundefined{@spart}{}{\expandafter\lakin@sparttitle\expandafter\@spart\expandafter{\lakin@ptit}}
\iflakin@partnum\else\PackageWarning{hebrew-lualatex-pdf}{Part-number math stays
  regular weight: no known \noexpand\part wrapper matched (hyperref/nameref
  version change). Cosmetic only}\fi
\iflakin@parttitle\else\PackageWarning{hebrew-lualatex-pdf}{Part-title math stays
  regular weight: no known \noexpand\part wrapper matched. Cosmetic only}\fi
% v4.12.0: THE PART NUMBER IN THE TOC WAS NOT DIRECTIONALLY SEPARATED FROM THE
% TITLE. report.cls/article.cls compose the part entry as ONE text stream --
% \addcontentsline{toc}{part}{\thepart\hspace{1em}#1} -- while EVERY other level
% goes through \numberline, which puts the number in an \hbox. A box is a bidi
% OBJECT and cannot join a neighbouring run; a bare "IV" can. So in the RTL TOC
% paragraph a Latin number and a Latin-INITIAL title merge into a SINGLE LTR run
% and are laid out left-to-right as a unit: the number lands on the wrong side of
% the row, and a mixed title is scrambled outright. Measured on \part{MCP ...}:
% the printed order was "<tail> V MCP" -- title fragment, number, title fragment.
% The chapter/section rows one line below, with the SAME Latin-initial titles,
% were correct on the same page. That contrast is the whole diagnosis.
% \en{} IS NOT A WORKAROUND, measured both ways on one fixture: it fixes a
% PURE-Latin title (\part{\en{Hooks}} renders correctly) and does NOT fix a MIXED
% one (\part{\en{MCP} <hebrew>} stays scrambled). babel's language switch changes
% font and hyphenation, not the bidi run structure -- so the defect is not
% reachable from the source side at all, which is why it belongs here.
% Boxing the number is the entire fix and is a STRICT NO-OP where the bug never
% bit: a Hebrew-titled document compiles BYTE-IDENTICAL before and after (md5,
% 3 clean runs), because there the neutral glue already sat between an L run and
% an R run. The PDF outline is byte-identical too -- hyperref's \pdfstringdef
% neutralises \mbox, so bookmark strings keep their part numbers.
% \numberline{\thepart} would also fix the direction, but it re-geometries every
% part row in the fleet and DROPS the number from the part bookmark (hyperref
% strips \numberline). \mbox changes nothing that was already right.
% The search pattern deliberately carries no # token: etoolbox cannot match a
% parameter token in a patch pattern (measured -- all three chain members
% reported "not patchable" while the pattern included #1). "\thepart \hspace
% {1em}" is brace-balanced and occurs exactly once in the wrapper body.
\newif\iflakin@parttocnum
\def\lakin@parttocnum#1{\iflakin@parttocnum\else
  \patchcmd{#1}{\thepart \hspace {1em}}{\mbox {\thepart }\hspace {1em}}
    {\lakin@parttocnumtrue}{}\fi}
\@ifundefined{NR@part}{}{\lakin@parttocnum\NR@part}
\@ifundefined{H@old@part}{}{\lakin@parttocnum\H@old@part}
\@ifundefined{@part}{}{\lakin@parttocnum\@part}
\iflakin@parttocnum\else\PackageWarning{hebrew-lualatex-pdf}{Part TOC number is
  not boxed: no known \noexpand\part wrapper matched. A part title that BEGINS
  with Latin will print its number on the wrong side of its TOC row}\fi
\patchcmd{\section}{\bfseries}{\bfseries\boldmath}{}
  {\PackageError{hebrew-lualatex-pdf}{Cannot patch section math weight}{Unsupported document class}}
\patchcmd{\subsection}{\bfseries}{\bfseries\boldmath}{}
  {\PackageError{hebrew-lualatex-pdf}{Cannot patch subsection math weight}{Unsupported document class}}
\patchcmd{\subsubsection}{\bfseries}{\bfseries\boldmath}{}
  {\PackageError{hebrew-lualatex-pdf}{Cannot patch subsubsection math weight}{Unsupported document class}}
\patchcmd{\paragraph}{\bfseries}{\bfseries\boldmath}{}{}
\patchcmd{\subparagraph}{\bfseries}{\bfseries\boldmath}{}{}
\@ifundefined{l@chapter}{}{%
  \patchcmd{\l@chapter}{\bfseries}{\bfseries\boldmath}{}{}}
\makeatother
\setcounter{tocdepth}{\lakittocdepth}
% <hllpdf:measured-toc>
% v4.16.0: Measure actual number labels, not a guessed digit count. Read only
% canonical contentsline records; never execute titles or addtocontents code
% during the measuring pass. Widths are local to this individual list and are
% available before its first row: no per-row drift and no extra aux pass.
\makeatletter
\newread\hllpdf@tocstream
\newdimen\hllpdf@toc@chapter
\newdimen\hllpdf@toc@section
\newdimen\hllpdf@toc@subsection
\newdimen\hllpdf@toc@subsubsection
\newdimen\hllpdf@toc@paragraph
\newdimen\hllpdf@toc@subparagraph
\newdimen\hllpdf@toc@figure
\newdimen\hllpdf@toc@table
\newdimen\hllpdf@tocpage
\newdimen\hllpdf@toc@sectionindent
\newdimen\hllpdf@toc@subsectionindent
\newdimen\hllpdf@toc@subsubsectionindent
\newdimen\hllpdf@toc@paragraphindent
\newdimen\hllpdf@toc@subparagraphindent
\def\hllpdf@tocdefaults{%
  \hllpdf@toc@chapter=2em
  \hllpdf@toc@section=3.2em
  \hllpdf@toc@subsection=3.9em
  \@ifundefined{c@chapter}{%
    \hllpdf@toc@subsubsection=3.2em
    \hllpdf@toc@paragraph=4.1em
    \hllpdf@toc@subparagraph=5em
  }{%
    \hllpdf@toc@subsubsection=4.1em
    \hllpdf@toc@paragraph=5em
    \hllpdf@toc@subparagraph=6em
  }%
  \hllpdf@toc@figure=2.3em
  \hllpdf@toc@table=2.3em
  \hllpdf@tocpage=\@pnumwidth\relax
}
\def\hllpdf@tocindents{%
  % Preserve the established shallow layout when its slots already fit.
  \hllpdf@toc@sectionindent=1.5em
  \@tempdima=\hllpdf@toc@chapter \advance\@tempdima by -2em
  \advance\hllpdf@toc@sectionindent by \@tempdima
  \hllpdf@toc@subsectionindent=\dimexpr
    \hllpdf@toc@sectionindent+\hllpdf@toc@section\relax
  \hllpdf@toc@subsubsectionindent=\dimexpr
    \hllpdf@toc@subsectionindent+\hllpdf@toc@subsection\relax
  \hllpdf@toc@paragraphindent=\dimexpr
    \hllpdf@toc@subsubsectionindent+\hllpdf@toc@subsubsection\relax
  \hllpdf@toc@subparagraphindent=\dimexpr
    \hllpdf@toc@paragraphindent+\hllpdf@toc@paragraph\relax
}
\def\hllpdf@tocfont@chapter{\bfseries\boldmath}
\def\hllpdf@tocfont@part{\large\bfseries\boldmath}
\def\hllpdf@tocentryfont{%
  \ifcsname hllpdf@tocfont@\hllpdf@toclevel\endcsname
    \csname hllpdf@tocfont@\hllpdf@toclevel\endcsname
  \fi}
\long\def\hllpdf@tocdiscard#1\hllpdf@tocstop{}
\def\hllpdf@tocscanrecord{%
  \@ifnextchar\contentsline{\hllpdf@tocscanline}{%
    \@ifnextchar\babel@toc{\hllpdf@tocscanlanguage}{\hllpdf@tocdiscard}}}
\long\def\hllpdf@tocscanlanguage\babel@toc#1#2#3\hllpdf@tocstop{%
  % Babel records are declarations, not arbitrary code. Track the language
  % without executing the original directive or changing the surrounding list.
  \def\hllpdf@toclanguage{#1}}
\long\def\hllpdf@tocscanline\contentsline#1#2#3#4#5\hllpdf@tocstop{%
  \def\hllpdf@toclevel{#1}%
  \ifcsname hllpdf@toc@#1\endcsname
    \hllpdf@tocscantitle#2\hllpdf@tocstop
  \fi
  % The page-number column needs its own measurement (including bold chapter
  % rows). Its extra clearance separates the page label from dot leaders.
  \setbox\@tempboxa=\hbox{\foreignlanguage{\hllpdf@toclanguage}{\hllpdf@tocentryfont #3\kern.3em}}%
  \ifdim\wd\@tempboxa>\hllpdf@tocpage
    \hllpdf@tocpage=\wd\@tempboxa
  \fi}
\def\hllpdf@tocscantitle{%
  \@ifnextchar\numberline{\hllpdf@tocscanlabel}{\hllpdf@tocdiscard}}
\long\def\hllpdf@tocscanlabel\numberline#1#2\hllpdf@tocstop{%
  \setbox\@tempboxa=\hbox{\foreignlanguage{\hllpdf@toclanguage}{\hllpdf@tocentryfont #1\kern.5em}}%
  \ifdim\wd\@tempboxa>\csname hllpdf@toc@\hllpdf@toclevel\endcsname
    \csname hllpdf@toc@\hllpdf@toclevel\endcsname=\wd\@tempboxa
  \fi}
\def\hllpdf@tocread{%
  \ifeof\hllpdf@tocstream\else
    \read\hllpdf@tocstream to\hllpdf@tocrecord
    \expandafter\hllpdf@tocscanrecord\hllpdf@tocrecord\hllpdf@tocstop
    \expandafter\hllpdf@tocread
  \fi}
\def\hllpdf@tocmeasure#1{%
  \makeatletter
  \edef\hllpdf@toclanguage{\languagename}%
  \hllpdf@tocdefaults
  \openin\hllpdf@tocstream=\jobname.#1\relax
  \hllpdf@tocread
  \closein\hllpdf@tocstream
  \hllpdf@tocindents
  \edef\@pnumwidth{\the\hllpdf@tocpage}%
  \@tempdima=\dimexpr\hllpdf@tocpage+1em\relax
  \ifdim\@tempdima>\@tocrmarg\relax
    \edef\@tocrmarg{\the\@tempdima}%
  \fi}
\let\hllpdf@originalstarttoc\@starttoc
\def\@starttoc#1{%
  \begingroup
    \def\hllpdf@tocext{#1}%
    \def\hllpdf@tocwanted{toc}%
    \ifx\hllpdf@tocext\hllpdf@tocwanted \hllpdf@tocmeasure{#1}\else
      \def\hllpdf@tocwanted{lof}%
      \ifx\hllpdf@tocext\hllpdf@tocwanted \hllpdf@tocmeasure{#1}\else
        \def\hllpdf@tocwanted{lot}%
        \ifx\hllpdf@tocext\hllpdf@tocwanted \hllpdf@tocmeasure{#1}\fi
      \fi
    \fi
    \hllpdf@originalstarttoc{#1}%
  \endgroup}
\hllpdf@tocdefaults
\hllpdf@tocindents
\@ifundefined{l@chapter}{}{%
  \patchcmd{\l@chapter}{\setlength\@tempdima{1.5em}}%
    {\setlength\@tempdima{\hllpdf@toc@chapter}}{}%
    {\PackageError{hebrew-lualatex-pdf}{Cannot size chapter TOC numbers}%
       {Unsupported contents layout: the chapter number slot was not found}}}
\renewcommand*\l@section{\@dottedtocline{1}%
  {\hllpdf@toc@sectionindent}{\hllpdf@toc@section}}
\renewcommand*\l@subsection{\@dottedtocline{2}%
  {\hllpdf@toc@subsectionindent}{\hllpdf@toc@subsection}}
\renewcommand*\l@subsubsection{\@dottedtocline{3}%
  {\hllpdf@toc@subsubsectionindent}{\hllpdf@toc@subsubsection}}
\renewcommand*\l@paragraph{\@dottedtocline{4}%
  {\hllpdf@toc@paragraphindent}{\hllpdf@toc@paragraph}}
\renewcommand*\l@subparagraph{\@dottedtocline{5}%
  {\hllpdf@toc@subparagraphindent}{\hllpdf@toc@subparagraph}}
\renewcommand*\l@figure{\@dottedtocline{1}{1.5em}{\hllpdf@toc@figure}}
\renewcommand*\l@table{\@dottedtocline{1}{1.5em}{\hllpdf@toc@table}}
\makeatother
% </hllpdf:measured-toc>
\setcounter{secnumdepth}{\lakitsecnumdepth}
% Public preamble-only caption-numbering policy. The historical within-major
% default is untouched unless this command is called. "continuous" removes the
% reset edge as well as the printed prefix, so Figure 1/2 and Table 1/2 stay
% unique across sections/chapters. Returning to "within-major" restores the
% original class-aware displays and native destination conventions. The policy
% never replaces \figcap/\tabcap or their measured link-placement mechanism.
\makeatletter
\let\lakin@majorfig\thefig
\let\lakin@majortbl\thetbl
\let\lakin@majorHfig\theHfig
\let\lakin@majorHtbl\theHtbl
\newcommand{\lakitcaptionnumbering}[1]{%
  \edef\lakin@numbering{#1}%
  \def\lakin@continuous{continuous}\def\lakin@withinmajor{within-major}%
  \ifx\lakin@numbering\lakin@continuous
    \@ifundefined{c@chapter}{%
      \@removefromreset{fig}{section}\@removefromreset{tbl}{section}%
    }{\@removefromreset{fig}{chapter}\@removefromreset{tbl}{chapter}}%
    \renewcommand{\thefig}{\arabic{fig}}\renewcommand{\thetbl}{\arabic{tbl}}%
    \def\theHfig{\arabic{fig}}\def\theHtbl{\arabic{tbl}}%
  \else\ifx\lakin@numbering\lakin@withinmajor
    \@ifundefined{c@chapter}{%
      \@removefromreset{fig}{section}\@removefromreset{tbl}{section}%
      \@addtoreset{fig}{section}\@addtoreset{tbl}{section}%
    }{%
      \@removefromreset{fig}{chapter}\@removefromreset{tbl}{chapter}%
      \@addtoreset{fig}{chapter}\@addtoreset{tbl}{chapter}}%
    \let\thefig\lakin@majorfig\let\thetbl\lakin@majortbl
    \let\theHfig\lakin@majorHfig\let\theHtbl\lakin@majorHtbl
  \else
    \PackageError{hebrew-english-latex-studio}{Unknown caption-numbering policy '#1'}%
      {Use continuous or within-major before begin{document}.}%
  \fi\fi}
\@onlypreamble\lakitcaptionnumbering
\makeatother
% Architecture readiness sentinel. Optional runtime layers require the complete
% core, so the marker is defined only after every mechanism/baseline definition.
% <hllpdf:architecture-ready>
\def\hllpdfcoreloaded{1}
% </hllpdf:architecture-ready>

% v4.22: registered file resources, not an assertion about embedded subset bytes.
% The log framing avoids long-path wrapping and preserves UTF-8 as hex.
% <hllpdf:font-inventory>
\directlua{
  local function inventory()
    local names = {}
    local unnamed = 0
    for id, item in font.each() do
      local path = item.filename
      if path and not (path == "") then names[path] = true else unnamed = unnamed + 1 end
    end
    local sorted = {}
    for path in pairs(names) do table.insert(sorted,path) end
    table.sort(sorted)
    texio.write_nl("log","HLLFONT|BEGIN|1")
    for index,path in ipairs(sorted) do
      local hex = path:gsub(".",function(c) return string.format(string.char(37) .. "02x",string.byte(c)) end)
      local count = math.ceil(string.len(hex)/48)
      texio.write_nl("log","HLLFONT|FILE|" .. index .. "|" .. count)
      for offset=1,string.len(hex),48 do texio.write_nl("log","HLLFONT|HEX|" .. hex:sub(offset,offset+47)) end
    end
    texio.write_nl("log","HLLFONT|END|" .. rawlen(sorted) .. "|" .. unnamed)
    texio.write_nl("log","")
  end
  luatexbase.add_to_callback("finish_pdffile",inventory,"hllpdf.runtime.fonts")
}
% </hllpdf:font-inventory>
:
%   \lakithead{מבנה המחשב}
% Every content page then carries the series' orientation line: chapter name
% at the reading edge, book title opposite, thin rule. (fancyhdr's [L]/[R]
% are PHYSICAL; under the RTL page direction [L] lands at the RIGHT/reading
% edge — measured, consistent across the corpus.) Packaged like the topic
% rail: always loaded, inert without the call. This centralizes the 8-line
% block that 8 lakits used to copy by hand — the same duplication vector
% that once let a lakit preamble go ~100 lines stale. Class-aware: report
% fills the mark from \chapter, article from \section. NB: a titlepage env
% self-empties; front matter that sets \pagestyle{empty} must restore with
% \pagestyle{fancy} afterwards (NOT plain — that discards the head).
% ============================================================================
\usepackage{fancyhdr}
% v3.61: \chapter* left a STALE running head — pages of an UNNUMBERED chapter
% carried the PREVIOUS chapter's title. report.cls routes \chapter* through
% \@schapter, which never calls \chaptermark, so \leftmark keeps whatever the last
% numbered chapter left. Measured: with \chapter{פרק ממוספר} then
% \chapter*{מבוא לא ממוספר}, page 4 (inside the starred chapter) printed
% "פרק ממוספר". Invisible on the chapter's FIRST page (\thispagestyle{plain}), so it
% only shows from the second page on — the pages an author is least likely to check.
% An unnumbered front-matter chapter (\chapter*{מבוא}, \chapter*{הקדמה}) is everyday
% Hebrew-academic structure, and the corpus already carries HAND-WRITTEN \markboth
% lines right after \chapter* — the exact per-document duplication \lakithead exists
% to retire. Patched INSIDE \lakithead so the macro stays inert when unused; a
% document's own later \markboth still wins, so those hand-patches do not regress.
\makeatletter
% <hllpdf:starred-heading-scope>
\newif\iflakin@starmarkpatched
\newcommand{\lakistarmark}{%
  \iflakin@starmarkpatched\else
  % Dispatch on the actual class counter, not hyperref's shared internals.
  \ifdefined\c@chapter
    \let\HEBorigschapter\@schapter
    % Record only after the native heading has reached its page. A mark
    % before the native routine can advertise the next heading too early.
    \def\@schapter##1{\HEBorigschapter{##1}\chaptermark{##1}}%
  \else
    % @ssect receives layout arguments, NOT the section name or level.
    % Capture the name in @startsection without introducing a TeX group:
    % grouping sectioning would break the native after-heading state.
    \def\HEB@sectionkind{}%
    \def\HEB@mainsection{section}%
    \let\HEBorigstartsection\@startsection
    \def\@startsection##1##2##3##4##5##6{%
      \def\HEB@sectionkind{##1}%
      \HEBorigstartsection{##1}{##2}{##3}{##4}{##5}{##6}}%
    \let\HEBorigssect\@ssect
    \def\@ssect##1##2##3##4##5{%
      \HEBorigssect{##1}{##2}{##3}{##4}{##5}%
      % Lower starred headings must not replace the article's section mark.
      \ifx\HEB@sectionkind\HEB@mainsection
        \sectionmark{##5}%
      \fi
      \ignorespaces}%
  \fi
  \global\lakin@starmarkpatchedtrue
  \fi}
% </hllpdf:starred-heading-scope>
% <hllpdf:running-language>
% Choose the wrapper when the mark is written, not when a later page ships.
% English punctuation must remain LTR even after returning to Hebrew prose.
\newcommand{\HEB@runningmark}[1]{%
  \iflanguage{english}{\markboth{\en{#1}}{}}{\markboth{\h{#1}}{}}}
% </hllpdf:running-language>
\makeatother
% v4.2.1 (NAV): every collection already names itself in \lakithead — feed the
% PDF /Title from the same declaration, so the navigation gate's
% PDF_REQUIRE_TITLE has a universal source and no fleet member ships
% title-less by missing a second declaration site. Measured before this: 14
% of 16 collections shipped /Title-empty while every one carried a perfectly
% good \lakithead string (plain Hebrew in all 16 — no math, no formatting, so
% the PDF-string path is clean). An explicit title always wins: set it with
% \lakitpdftitle{...} anywhere before \lakithead, or with a plain
% \hypersetup{pdftitle=...} — the fallback detects that via hyperref's
% \@pdftitle, probe-verified \@empty-equal until the first explicit set.
\makeatletter
\newif\iflakit@pdftitleset
\newcommand{\lakitpdftitle}[1]{%
  \hypersetup{pdftitle={#1}}\global\lakit@pdftitlesettrue}
\newcommand{\lakithead}[1]{%
  \iflakit@pdftitleset\else
    \ifx\@pdftitle\@empty
      \hypersetup{pdftitle={#1}}\global\lakit@pdftitlesettrue
    \fi
  \fi
  \pagestyle{fancy}\fancyhf{}%
  \lakistarmark

  \fancyhead[L]{\small\textcolor{cGray}{\leftmark}}%
  % <hllpdf:fixed-head-language>
  % The fixed document title has its OWN language snapshot, independent of
  % the current section/chapter mark. Choose the branch NOW, not at shipout:
  % the page may be Hebrew by the time an English title is finally printed.
  \iflanguage{english}{%
    \fancyhead[R]{\small\textcolor{cGray}{\en{#1}}}%
  }{%
    \fancyhead[R]{\small\textcolor{cGray}{\h{#1}}}%
  }%
  % </hllpdf:fixed-head-language>
  \fancyfoot[C]{\thepage}%
  \renewcommand{\headrulewidth}{0.4pt}%
  \ifdefined\chaptermark
    \renewcommand{\chaptermark}[1]{\HEB@runningmark{##1}}%
  \else
    \renewcommand{\sectionmark}[1]{\HEB@runningmark{##1}}%
  \fi}
\makeatother

% v3.39.1: path-independent byte-reproducibility. LuaTeX derives the trailer
% /ID from the OUTPUT PATH, so two clones of the same source at different
% directories produced different sha256 (measured: the ONLY delta was the
% 62-byte /ID line). Bit 512 suppresses exactly that line — measured: Creator/
% Producer/CreationDate/ModDate and all content untouched. /ID is optional
% per ISO 32000 (required only for ENCRYPTED PDFs and PDF/A — neither is a
% target here). If a future workflow needs an ID (PDF/A, Acrobat revision
% matching), delete this line and it returns.
\pdfvariable suppressoptionalinfo \numexpr 512\relax

% v3.75: BOLD MATH IN BOLD HEADINGS. Heading text is \bfseries but math inside
% it stayed regular-weight CM — a hairline $V_T$ inside a heavy 24.8pt chapter
% title (measured: CMMI12 fl=6 next to FrankRuhlLibre-Bold fl=20; 155 live
% heading-with-math sites fleet-wide). \boldmath rides along wherever the class
% sets \bfseries for a heading; A/B measured: $V_T$ -> CMMIB10 fl=22, 0 missing
% characters (the bold CM faces ship with this TeX Live). Guarded per-macro so
% the article class (no \@makechapterhead) and future class changes stay safe.
% v4.0.1 (adoption fix): the SEARCH patterns above carry no parameter token.
% v4.0.0 made them specific by including the class's own #1/#2, and measured on
% TeX Live 2023 that never matches: a \patchcmd search pattern containing a
% parameter token cannot match while the call sits inside \@ifundefined{...}{}{...}
% (probed both #1 and ##1 -> FAIL; the same pattern at top level succeeds). Every
% one of these patches therefore failed, and because v4.0.0 also replaced the
% empty failure branch with \PackageError, the failure became FATAL: all four
% bundled example assets stopped building ("Cannot patch chapter-title math
% weight -> Fatal error, no output PDF"). The \PackageError branches are KEPT —
% a silent no-op here is what let v3.75's fix rot unnoticed — only the parameter
% tokens are dropped, so the patches match and a real class change still errors.
\makeatletter
\@ifundefined{@makechapterhead}{}{%
  \patchcmd{\@makechapterhead}{\huge\bfseries \@chapapp}
    {\huge\bfseries\boldmath \@chapapp}{}
    {\PackageError{hebrew-lualatex-pdf}{Cannot patch chapter-number math weight}{Unsupported document class}}%
  \patchcmd{\@makechapterhead}{\Huge \bfseries}
    {\Huge \bfseries\boldmath}{}
    {\PackageError{hebrew-lualatex-pdf}{Cannot patch chapter-title math weight}{Unsupported document class}}%
  \patchcmd{\@makeschapterhead}{\Huge \bfseries}
    {\Huge \bfseries\boldmath}{}
    {\PackageError{hebrew-lualatex-pdf}{Cannot patch starred chapter math weight}{Unsupported document class}}}
% v4.0.1: \part sits under a THREE-DEEP wrapper chain, so the class-based anchor
% is not in \@part at all. Measured with \show: hyperref rewrites \@part to
% `...\phantomsection\fi \H@old@part`, and nameref rewrites THAT to
% `\NR@gettitle{#1}\NR@part[{#1}]{#2}` — the class body ends up in \NR@part.
% v4.0.0 patched \@part and made failure fatal, which is why all seven bundled
% examples stopped building. \part is NOT a dead heading level here: 188 live uses
% across the fleet, so the bold-math fix is worth getting right rather than
% dropping. We walk the chain innermost-first and stop at the first hit. Failure is
% a WARNING, not an error: the chain is a property of hyperref/nameref versions, so
% a future package release must degrade to regular-weight part math — a cosmetic
% loss — instead of refusing to build. Chapter and section keep \PackageError,
% because those targets are in the CLASS and were verified to patch here.
\newif\iflakin@partnum \newif\iflakin@parttitle \newif\iflakin@sparttitle
\def\lakin@partnum#1#2{\iflakin@partnum\else
  \patchcmd{#1}{#2\bfseries \partname}{#2\bfseries\boldmath \partname}
    {\lakin@partnumtrue}{}\fi}
\def\lakin@parttitle#1#2{\iflakin@parttitle\else
  \patchcmd{#1}{#2 \bfseries}{#2 \bfseries\boldmath}{\lakin@parttitletrue}{}\fi}
\def\lakin@sparttitle#1#2{\iflakin@sparttitle\else
  \patchcmd{#1}{#2 \bfseries}{#2 \bfseries\boldmath}{\lakin@sparttitletrue}{}\fi}
\@ifundefined{c@chapter}{\def\lakin@pnum{\Large}\def\lakin@ptit{\huge}}%
                        {\def\lakin@pnum{\huge}\def\lakin@ptit{\Huge}}
\@ifundefined{NR@part}{}{\expandafter\lakin@partnum\expandafter\NR@part\expandafter{\lakin@pnum}%
                         \expandafter\lakin@parttitle\expandafter\NR@part\expandafter{\lakin@ptit}}
\@ifundefined{H@old@part}{}{\expandafter\lakin@partnum\expandafter\H@old@part\expandafter{\lakin@pnum}%
                            \expandafter\lakin@parttitle\expandafter\H@old@part\expandafter{\lakin@ptit}}
\@ifundefined{@part}{}{\expandafter\lakin@partnum\expandafter\@part\expandafter{\lakin@pnum}%
                       \expandafter\lakin@parttitle\expandafter\@part\expandafter{\lakin@ptit}}
\@ifundefined{NR@spart}{}{\expandafter\lakin@sparttitle\expandafter\NR@spart\expandafter{\lakin@ptit}}
\@ifundefined{H@old@spart}{}{\expandafter\lakin@sparttitle\expandafter\H@old@spart\expandafter{\lakin@ptit}}
\@ifundefined{@spart}{}{\expandafter\lakin@sparttitle\expandafter\@spart\expandafter{\lakin@ptit}}
\iflakin@partnum\else\PackageWarning{hebrew-lualatex-pdf}{Part-number math stays
  regular weight: no known \noexpand\part wrapper matched (hyperref/nameref
  version change). Cosmetic only}\fi
\iflakin@parttitle\else\PackageWarning{hebrew-lualatex-pdf}{Part-title math stays
  regular weight: no known \noexpand\part wrapper matched. Cosmetic only}\fi
% v4.12.0: THE PART NUMBER IN THE TOC WAS NOT DIRECTIONALLY SEPARATED FROM THE
% TITLE. report.cls/article.cls compose the part entry as ONE text stream --
% \addcontentsline{toc}{part}{\thepart\hspace{1em}#1} -- while EVERY other level
% goes through \numberline, which puts the number in an \hbox. A box is a bidi
% OBJECT and cannot join a neighbouring run; a bare "IV" can. So in the RTL TOC
% paragraph a Latin number and a Latin-INITIAL title merge into a SINGLE LTR run
% and are laid out left-to-right as a unit: the number lands on the wrong side of
% the row, and a mixed title is scrambled outright. Measured on \part{MCP ...}:
% the printed order was "<tail> V MCP" -- title fragment, number, title fragment.
% The chapter/section rows one line below, with the SAME Latin-initial titles,
% were correct on the same page. That contrast is the whole diagnosis.
% \en{} IS NOT A WORKAROUND, measured both ways on one fixture: it fixes a
% PURE-Latin title (\part{\en{Hooks}} renders correctly) and does NOT fix a MIXED
% one (\part{\en{MCP} <hebrew>} stays scrambled). babel's language switch changes
% font and hyphenation, not the bidi run structure -- so the defect is not
% reachable from the source side at all, which is why it belongs here.
% Boxing the number is the entire fix and is a STRICT NO-OP where the bug never
% bit: a Hebrew-titled document compiles BYTE-IDENTICAL before and after (md5,
% 3 clean runs), because there the neutral glue already sat between an L run and
% an R run. The PDF outline is byte-identical too -- hyperref's \pdfstringdef
% neutralises \mbox, so bookmark strings keep their part numbers.
% \numberline{\thepart} would also fix the direction, but it re-geometries every
% part row in the fleet and DROPS the number from the part bookmark (hyperref
% strips \numberline). \mbox changes nothing that was already right.
% The search pattern deliberately carries no # token: etoolbox cannot match a
% parameter token in a patch pattern (measured -- all three chain members
% reported "not patchable" while the pattern included #1). "\thepart \hspace
% {1em}" is brace-balanced and occurs exactly once in the wrapper body.
\newif\iflakin@parttocnum
\def\lakin@parttocnum#1{\iflakin@parttocnum\else
  \patchcmd{#1}{\thepart \hspace {1em}}{\mbox {\thepart }\hspace {1em}}
    {\lakin@parttocnumtrue}{}\fi}
\@ifundefined{NR@part}{}{\lakin@parttocnum\NR@part}
\@ifundefined{H@old@part}{}{\lakin@parttocnum\H@old@part}
\@ifundefined{@part}{}{\lakin@parttocnum\@part}
\iflakin@parttocnum\else\PackageWarning{hebrew-lualatex-pdf}{Part TOC number is
  not boxed: no known \noexpand\part wrapper matched. A part title that BEGINS
  with Latin will print its number on the wrong side of its TOC row}\fi
\patchcmd{\section}{\bfseries}{\bfseries\boldmath}{}
  {\PackageError{hebrew-lualatex-pdf}{Cannot patch section math weight}{Unsupported document class}}
\patchcmd{\subsection}{\bfseries}{\bfseries\boldmath}{}
  {\PackageError{hebrew-lualatex-pdf}{Cannot patch subsection math weight}{Unsupported document class}}
\patchcmd{\subsubsection}{\bfseries}{\bfseries\boldmath}{}
  {\PackageError{hebrew-lualatex-pdf}{Cannot patch subsubsection math weight}{Unsupported document class}}
\patchcmd{\paragraph}{\bfseries}{\bfseries\boldmath}{}{}
\patchcmd{\subparagraph}{\bfseries}{\bfseries\boldmath}{}{}
\@ifundefined{l@chapter}{}{%
  \patchcmd{\l@chapter}{\bfseries}{\bfseries\boldmath}{}{}}
\makeatother
\setcounter{tocdepth}{\lakittocdepth}
% <hllpdf:measured-toc>
% v4.16.0: Measure actual number labels, not a guessed digit count. Read only
% canonical contentsline records; never execute titles or addtocontents code
% during the measuring pass. Widths are local to this individual list and are
% available before its first row: no per-row drift and no extra aux pass.
\makeatletter
\newread\hllpdf@tocstream
\newdimen\hllpdf@toc@chapter
\newdimen\hllpdf@toc@section
\newdimen\hllpdf@toc@subsection
\newdimen\hllpdf@toc@subsubsection
\newdimen\hllpdf@toc@paragraph
\newdimen\hllpdf@toc@subparagraph
\newdimen\hllpdf@toc@figure
\newdimen\hllpdf@toc@table
\newdimen\hllpdf@tocpage
\newdimen\hllpdf@toc@sectionindent
\newdimen\hllpdf@toc@subsectionindent
\newdimen\hllpdf@toc@subsubsectionindent
\newdimen\hllpdf@toc@paragraphindent
\newdimen\hllpdf@toc@subparagraphindent
\def\hllpdf@tocdefaults{%
  \hllpdf@toc@chapter=2em
  \hllpdf@toc@section=3.2em
  \hllpdf@toc@subsection=3.9em
  \@ifundefined{c@chapter}{%
    \hllpdf@toc@subsubsection=3.2em
    \hllpdf@toc@paragraph=4.1em
    \hllpdf@toc@subparagraph=5em
  }{%
    \hllpdf@toc@subsubsection=4.1em
    \hllpdf@toc@paragraph=5em
    \hllpdf@toc@subparagraph=6em
  }%
  \hllpdf@toc@figure=2.3em
  \hllpdf@toc@table=2.3em
  \hllpdf@tocpage=\@pnumwidth\relax
}
\def\hllpdf@tocindents{%
  % Preserve the established shallow layout when its slots already fit.
  \hllpdf@toc@sectionindent=1.5em
  \@tempdima=\hllpdf@toc@chapter \advance\@tempdima by -2em
  \advance\hllpdf@toc@sectionindent by \@tempdima
  \hllpdf@toc@subsectionindent=\dimexpr
    \hllpdf@toc@sectionindent+\hllpdf@toc@section\relax
  \hllpdf@toc@subsubsectionindent=\dimexpr
    \hllpdf@toc@subsectionindent+\hllpdf@toc@subsection\relax
  \hllpdf@toc@paragraphindent=\dimexpr
    \hllpdf@toc@subsubsectionindent+\hllpdf@toc@subsubsection\relax
  \hllpdf@toc@subparagraphindent=\dimexpr
    \hllpdf@toc@paragraphindent+\hllpdf@toc@paragraph\relax
}
\def\hllpdf@tocfont@chapter{\bfseries\boldmath}
\def\hllpdf@tocfont@part{\large\bfseries\boldmath}
\def\hllpdf@tocentryfont{%
  \ifcsname hllpdf@tocfont@\hllpdf@toclevel\endcsname
    \csname hllpdf@tocfont@\hllpdf@toclevel\endcsname
  \fi}
\long\def\hllpdf@tocdiscard#1\hllpdf@tocstop{}
\def\hllpdf@tocscanrecord{%
  \@ifnextchar\contentsline{\hllpdf@tocscanline}{%
    \@ifnextchar\babel@toc{\hllpdf@tocscanlanguage}{\hllpdf@tocdiscard}}}
\long\def\hllpdf@tocscanlanguage\babel@toc#1#2#3\hllpdf@tocstop{%
  % Babel records are declarations, not arbitrary code. Track the language
  % without executing the original directive or changing the surrounding list.
  \def\hllpdf@toclanguage{#1}}
\long\def\hllpdf@tocscanline\contentsline#1#2#3#4#5\hllpdf@tocstop{%
  \def\hllpdf@toclevel{#1}%
  \ifcsname hllpdf@toc@#1\endcsname
    \hllpdf@tocscantitle#2\hllpdf@tocstop
  \fi
  % The page-number column needs its own measurement (including bold chapter
  % rows). Its extra clearance separates the page label from dot leaders.
  \setbox\@tempboxa=\hbox{\foreignlanguage{\hllpdf@toclanguage}{\hllpdf@tocentryfont #3\kern.3em}}%
  \ifdim\wd\@tempboxa>\hllpdf@tocpage
    \hllpdf@tocpage=\wd\@tempboxa
  \fi}
\def\hllpdf@tocscantitle{%
  \@ifnextchar\numberline{\hllpdf@tocscanlabel}{\hllpdf@tocdiscard}}
\long\def\hllpdf@tocscanlabel\numberline#1#2\hllpdf@tocstop{%
  \setbox\@tempboxa=\hbox{\foreignlanguage{\hllpdf@toclanguage}{\hllpdf@tocentryfont #1\kern.5em}}%
  \ifdim\wd\@tempboxa>\csname hllpdf@toc@\hllpdf@toclevel\endcsname
    \csname hllpdf@toc@\hllpdf@toclevel\endcsname=\wd\@tempboxa
  \fi}
\def\hllpdf@tocread{%
  \ifeof\hllpdf@tocstream\else
    \read\hllpdf@tocstream to\hllpdf@tocrecord
    \expandafter\hllpdf@tocscanrecord\hllpdf@tocrecord\hllpdf@tocstop
    \expandafter\hllpdf@tocread
  \fi}
\def\hllpdf@tocmeasure#1{%
  \makeatletter
  \edef\hllpdf@toclanguage{\languagename}%
  \hllpdf@tocdefaults
  \openin\hllpdf@tocstream=\jobname.#1\relax
  \hllpdf@tocread
  \closein\hllpdf@tocstream
  \hllpdf@tocindents
  \edef\@pnumwidth{\the\hllpdf@tocpage}%
  \@tempdima=\dimexpr\hllpdf@tocpage+1em\relax
  \ifdim\@tempdima>\@tocrmarg\relax
    \edef\@tocrmarg{\the\@tempdima}%
  \fi}
\let\hllpdf@originalstarttoc\@starttoc
\def\@starttoc#1{%
  \begingroup
    \def\hllpdf@tocext{#1}%
    \def\hllpdf@tocwanted{toc}%
    \ifx\hllpdf@tocext\hllpdf@tocwanted \hllpdf@tocmeasure{#1}\else
      \def\hllpdf@tocwanted{lof}%
      \ifx\hllpdf@tocext\hllpdf@tocwanted \hllpdf@tocmeasure{#1}\else
        \def\hllpdf@tocwanted{lot}%
        \ifx\hllpdf@tocext\hllpdf@tocwanted \hllpdf@tocmeasure{#1}\fi
      \fi
    \fi
    \hllpdf@originalstarttoc{#1}%
  \endgroup}
\hllpdf@tocdefaults
\hllpdf@tocindents
\@ifundefined{l@chapter}{}{%
  \patchcmd{\l@chapter}{\setlength\@tempdima{1.5em}}%
    {\setlength\@tempdima{\hllpdf@toc@chapter}}{}%
    {\PackageError{hebrew-lualatex-pdf}{Cannot size chapter TOC numbers}%
       {Unsupported contents layout: the chapter number slot was not found}}}
\renewcommand*\l@section{\@dottedtocline{1}%
  {\hllpdf@toc@sectionindent}{\hllpdf@toc@section}}
\renewcommand*\l@subsection{\@dottedtocline{2}%
  {\hllpdf@toc@subsectionindent}{\hllpdf@toc@subsection}}
\renewcommand*\l@subsubsection{\@dottedtocline{3}%
  {\hllpdf@toc@subsubsectionindent}{\hllpdf@toc@subsubsection}}
\renewcommand*\l@paragraph{\@dottedtocline{4}%
  {\hllpdf@toc@paragraphindent}{\hllpdf@toc@paragraph}}
\renewcommand*\l@subparagraph{\@dottedtocline{5}%
  {\hllpdf@toc@subparagraphindent}{\hllpdf@toc@subparagraph}}
\renewcommand*\l@figure{\@dottedtocline{1}{1.5em}{\hllpdf@toc@figure}}
\renewcommand*\l@table{\@dottedtocline{1}{1.5em}{\hllpdf@toc@table}}
\makeatother
% </hllpdf:measured-toc>
\setcounter{secnumdepth}{\lakitsecnumdepth}
% Public preamble-only caption-numbering policy. The historical within-major
% default is untouched unless this command is called. "continuous" removes the
% reset edge as well as the printed prefix, so Figure 1/2 and Table 1/2 stay
% unique across sections/chapters. Returning to "within-major" restores the
% original class-aware displays and native destination conventions. The policy
% never replaces \figcap/\tabcap or their measured link-placement mechanism.
\makeatletter
\let\lakin@majorfig\thefig
\let\lakin@majortbl\thetbl
\let\lakin@majorHfig\theHfig
\let\lakin@majorHtbl\theHtbl
\newcommand{\lakitcaptionnumbering}[1]{%
  \edef\lakin@numbering{#1}%
  \def\lakin@continuous{continuous}\def\lakin@withinmajor{within-major}%
  \ifx\lakin@numbering\lakin@continuous
    \@ifundefined{c@chapter}{%
      \@removefromreset{fig}{section}\@removefromreset{tbl}{section}%
    }{\@removefromreset{fig}{chapter}\@removefromreset{tbl}{chapter}}%
    \renewcommand{\thefig}{\arabic{fig}}\renewcommand{\thetbl}{\arabic{tbl}}%
    \def\theHfig{\arabic{fig}}\def\theHtbl{\arabic{tbl}}%
  \else\ifx\lakin@numbering\lakin@withinmajor
    \@ifundefined{c@chapter}{%
      \@removefromreset{fig}{section}\@removefromreset{tbl}{section}%
      \@addtoreset{fig}{section}\@addtoreset{tbl}{section}%
    }{%
      \@removefromreset{fig}{chapter}\@removefromreset{tbl}{chapter}%
      \@addtoreset{fig}{chapter}\@addtoreset{tbl}{chapter}}%
    \let\thefig\lakin@majorfig\let\thetbl\lakin@majortbl
    \let\theHfig\lakin@majorHfig\let\theHtbl\lakin@majorHtbl
  \else
    \PackageError{hebrew-english-latex-studio}{Unknown caption-numbering policy '#1'}%
      {Use continuous or within-major before begin{document}.}%
  \fi\fi}
\@onlypreamble\lakitcaptionnumbering
\makeatother
% Architecture readiness sentinel. Optional runtime layers require the complete
% core, so the marker is defined only after every mechanism/baseline definition.
% <hllpdf:architecture-ready>
\def\hllpdfcoreloaded{1}
% </hllpdf:architecture-ready>

% v4.22: registered file resources, not an assertion about embedded subset bytes.
% The log framing avoids long-path wrapping and preserves UTF-8 as hex.
% <hllpdf:font-inventory>
\directlua{
  local function inventory()
    local names = {}
    local unnamed = 0
    for id, item in font.each() do
      local path = item.filename
      if path and not (path == "") then names[path] = true else unnamed = unnamed + 1 end
    end
    local sorted = {}
    for path in pairs(names) do table.insert(sorted,path) end
    table.sort(sorted)
    texio.write_nl("log","HLLFONT|BEGIN|1")
    for index,path in ipairs(sorted) do
      local hex = path:gsub(".",function(c) return string.format(string.char(37) .. "02x",string.byte(c)) end)
      local count = math.ceil(string.len(hex)/48)
      texio.write_nl("log","HLLFONT|FILE|" .. index .. "|" .. count)
      for offset=1,string.len(hex),48 do texio.write_nl("log","HLLFONT|HEX|" .. hex:sub(offset,offset+47)) end
    end
    texio.write_nl("log","HLLFONT|END|" .. rawlen(sorted) .. "|" .. unnamed)
    texio.write_nl("log","")
  end
  luatexbase.add_to_callback("finish_pdffile",inventory,"hllpdf.runtime.fonts")
}
% </hllpdf:font-inventory>

% document_components.tex — semantic document helpers, independent of any theme.
% <hllpdf:architecture-guard>
\ifdefined\hllpdfcoreloaded\else
  \errmessage{hllpdf: document_components.tex requires preamble.tex first}
  \expandafter\endinput
\fi
\ifdefined\hllpdfdocumentcomponentsloaded
  \errmessage{hllpdf: document_components.tex loaded more than once}
  \expandafter\endinput
\fi
\def\hllpdfdocumentcomponentsloaded{1}
% </hllpdf:architecture-guard>
% \sourcelink{URL}{descriptive label}{why this source is here}
% Put English labels in \en{...}; the destination is not printed as raw prose.
% The role follows the surrounding language, not the link title.
% The source role is textual (not conveyed by color alone). Long labels can wrap.
% This is a SOURCE ANNOTATION, not a substitute for claim-level scientific citations.
\newcommand{\sourcelink}[3]{%
  \par\addvspace{6pt}\begingroup
  \small\noindent\textbf{\iflanguage{english}{Source:}{מקור:}}\enspace
  \href{#1}{\textcolor{cBlue}{\textbf{#2}}}\par\nobreak
  \noindent #3\par
  \endgroup\addvspace{6pt}}

% reader_theme.tex — optional, restrained presentation for everyday reading.
% No exam wording, suppressed titles, equation policy or numbering changes.
% Measured with this Babel/mdframed backend: leftline is the LEADING edge,
% not an absolute physical-left promise. It paints right in Hebrew and left
% in englishprose. Keep this choice tied to a rendered geometry test.
% <hllpdf:architecture-guard>
\ifdefined\hllpdfcoreloaded\else
  \errmessage{hllpdf: reader_theme.tex requires preamble.tex first}
  \expandafter\endinput
\fi
\ifdefined\hllpdfcleanthemeloaded
  \errmessage{hllpdf: reader_theme.tex conflicts with clean_theme.tex}
  \expandafter\endinput
\fi
\ifdefined\hllpdfcleanvariantloaded
  \errmessage{hllpdf: reader_theme.tex conflicts with clean_variant.tex}
  \expandafter\endinput
\fi
\ifdefined\hllpdfacademicthemeloaded
  \errmessage{hllpdf: reader_theme.tex conflicts with academic_theme.tex}
  \expandafter\endinput
\fi
\ifdefined\hllpdfreferencethemeloaded
  \errmessage{hllpdf: reader_theme.tex conflicts with reference_theme.tex}
  \expandafter\endinput
\fi
\ifdefined\hllpdfreaderthemeloaded
  \errmessage{hllpdf: reader_theme.tex loaded more than once}
  \expandafter\endinput
\fi
\def\hllpdfreaderthemeloaded{1}
% </hllpdf:architecture-guard>
\definecolor{cBlue}{HTML}{243F50}   \definecolor{cBlueBg}{HTML}{F3F6F8}
\definecolor{cGreen}{HTML}{35544B}  \definecolor{cGreenBg}{HTML}{F3F7F5}
\definecolor{cAmber}{HTML}{765815}  \definecolor{cAmberBg}{HTML}{FAF7F0}
\definecolor{cPurple}{HTML}{574465} \definecolor{cPurpleBg}{HTML}{F6F4F8}
\definecolor{cTeal}{HTML}{285858}   \definecolor{cTealBg}{HTML}{F1F7F7}
\definecolor{cRed}{HTML}{853A36}    \definecolor{cRedBg}{HTML}{FCF5F4}
\definecolor{cRule}{HTML}{243F50}
\definecolor{cSlate}{HTML}{3F4A56}  \definecolor{cSlateBg}{HTML}{F4F6F8}
\mdfdefinestyle{hllpdfreader}{
  frametitle={},leftline=true,rightline=false,topline=false,bottomline=false,
  linewidth=2.2pt,roundcorner=0pt,
  innerleftmargin=10pt,innerrightmargin=11pt,
  innertopmargin=7pt,innerbottommargin=8pt,
  skipabove=10pt,skipbelow=10pt,
  everyline=true,splitbottomskip=7pt,splittopskip=10pt,
  frametitlefont=\normalfont\sffamily\bfseries,
  frametitleaboveskip=5pt,frametitlebelowskip=3pt
}
\renewmdenv[style=hllpdfreader,linecolor=cBlue,backgroundcolor=cBlueBg,
  frametitlebackgroundcolor=cBlueBg,frametitlefontcolor=cBlue]{defbox}
\renewmdenv[style=hllpdfreader,linecolor=cGreen,backgroundcolor=cGreenBg,
  frametitlebackgroundcolor=cGreenBg,frametitlefontcolor=cGreen]{thmbox}
\renewmdenv[style=hllpdfreader,linecolor=cAmber,backgroundcolor=cAmberBg,
  frametitlebackgroundcolor=cAmberBg,frametitlefontcolor=cAmber]{notebox}
\renewmdenv[style=hllpdfreader,linecolor=cPurple,backgroundcolor=cPurpleBg,
  frametitlebackgroundcolor=cPurpleBg,frametitlefontcolor=cPurple]{exbox}
\renewmdenv[style=hllpdfreader,linecolor=cRed,backgroundcolor=cRedBg,
  frametitlebackgroundcolor=cRedBg,frametitlefontcolor=cRed]{warnbox}
\renewmdenv[style=hllpdfreader,linecolor=cTeal,backgroundcolor=cTealBg,
  frametitlebackgroundcolor=cTealBg,frametitlefontcolor=cTeal]{keybox}
% Keep the core's tested codebox number geometry; only its palette changes.

\hypersetup{pdftitle={מקורות שקוראים נכון},pdfauthor={},pdfsubject={הדגמה דו־לשונית למקורות והפניות}}
\begin{document}
{\small\sffamily\noindent דוגמה למקורות והפניות\par}\medskip
{\LARGE\sffamily\bfseries\noindent מקורות שקוראים נכון\par}\smallskip
{\large\noindent הפניה קצרה, רשומה ברורה וקישור למקום הנכון\par}\bigskip
המספר שבסוגריים מפנה לרשומה ברשימת המקורות. שם המקור מסביר מהו המסמך,
והקישור מאפשר להגיע אליו. בדוגמה זו משתמשים בשלושה עמודי תיעוד רשמיים;
אין כאן מקורות מדעיים מומצאים או מספרי עמודים שלא נבדקו.
\section{הפניה צמודה לטענה}
חבילת \en{Babel} מספקת תשתית רב־לשונית למסמכי \en{LaTeX}, ובה גם תמיכה
בכתיבה מימין לשמאל \cite{babel}. חבילת \en{hyperref} מוסיפה קישורים
להפניות במסמך \cite{hyperref}.

שני המקורות ניתנים לציטוט יחד \cite{babel,hyperref}. אפשר לציין גם את
ההקשר במילים: \cite[\h{תיאור החבילה}]{babel}. ההפניה החוזרת נשארת קשורה
לאותו מקור, בלי להחליף את המפתח או להפוך את סדר המספרים.
\section{קישור שמסביר למה הוא כאן}
תיעוד הליבה של \en{LaTeX} מרכז מדריכים למחברים ולכותבי חבילות \cite{latex}.
ההערה שמתחת לפסקה מציגה את הקישור כמקור ולא כהמשך מקרי של ההסבר.
\sourcelink{https://www.latex-project.org/help/documentation/}{תיעוד הליבה של \en{LaTeX}}{נקודת פתיחה למדריכי מחברים, למתמטיקה ולתיעוד של המנגנון.}
\begin{englishprose}
A source annotation follows the surrounding language. Its role remains explicit,
including in monochrome printing.
\sourcelink{https://ctan.org/pkg/babel}{Babel package documentation}{Package overview and links to the user guide and source repository.}
\end{englishprose}
\begin{notebox}[frametitle={מה עדיין צריך לבדוק}]
הפניה תקינה מבחינה טכנית אינה מוכיחה שהמקור תומך בטענה. יש לקרוא את המקור,
לבדוק את פרטי הרשומה ולוודא שהקישור מוביל לעמוד שהתכוונו לצטט.
\end{notebox}
\begin{thebibliography}{9}
\bibitem{babel}\bibentryprose{english}{Babel. Multilingual support for LaTeX. CTAN package page.
\href{https://ctan.org/pkg/babel}{Open package documentation}.}
\bibitem{hyperref}\bibentryprose{english}{Hyperref. Extensive support for hypertext in LaTeX. CTAN package page.
\href{https://ctan.org/pkg/hyperref}{Open package documentation}.}
\bibitem{latex}\bibentryprose{english}{The LaTeX Project. Core Documentation.
\href{https://www.latex-project.org/help/documentation/}{Open author and developer guides}.}
\end{thebibliography}
\end{document}
